% Generated mechanically from the frozen more-algebra.tex target.
% Do not edit this generated file; edit the bound locale source instead.

% OK, start here.
%

\setcounter{chapter}{14}
\chapter{代数进阶}



\phantomsection
\label{more-algebra-section-phantom}




\section{引言}
\label{more-algebra-section-introduction}

\noindent
本章证明交换代数中的若干结果；它们比第一章交换代数中的结果更为深入，见
Algebra, Section \ref{algebra-section-introduction}。
可参见 \cite{MatCA}。












\section{给读者的建议}
\label{more-algebra-section-advice}

\noindent
与交换代数一章相比，本章各节的独立性更强。从
Section \ref{more-algebra-section-derived-modules} 起，我们将自由使用环上模的（无界）
导出范畴及其所带的全部工具。





\section{稳定自由模}
\label{more-algebra-section-stably-free}

\noindent
下面给出看来已获普遍接受的定义。

\begin{definition}
\label{more-algebra-definition-stably-free}
设 $R$ 为环。
\begin{enumerate}
\item 称两个 $R$-模 $M$、$N$ {\it 稳定同构}，如果存在
$n, m \geq 0$，使得作为 $R$-模有
$M \oplus R^{\oplus m} \cong N \oplus R^{\oplus n}$。
\item 若模 $M$ 稳定同构于某个自由模，则称其 {\it 稳定自由}。
\end{enumerate}
\end{definition}
\noindent
注意，稳定自由模是投射模。

\begin{lemma}
\label{more-algebra-lemma-exact-category-stably-free}
设 $R$ 为环。设
$0 \to P' \to P \to P'' \to 0$ 是有限投射 $R$-模的短正合序列。
若这 $3$ 个模中的 $2$ 个稳定自由，则第三个也稳定自由。
\end{lemma}

\begin{proof}
由于这些模都是投射模，该序列分裂。因此可选取同构
$P = P' \oplus P''$。若 $P' \oplus R^{\oplus n}$ 与
$P'' \oplus R^{\oplus m}$ 均为自由模，则
$P \oplus R^{\oplus n + m}$ 为自由模。再假设 $P'$ 与 $P$
稳定自由，例如 $P \oplus R^{\oplus n}$ 和
$P' \oplus R^{\oplus m}$ 均为自由模。于是
$$
P'' \oplus (P' \oplus R^{\oplus m}) \oplus R^{\oplus n} =
(P'' \oplus P') \oplus R^{\oplus m} \oplus R^{\oplus n} =
(P \oplus R^{\oplus n}) \oplus R^{\oplus m}
$$
是自由模，故 $P''$ 稳定自由。由对称性也得到剩余一种情形。
\end{proof}
 
\begin{lemma}
\label{more-algebra-lemma-lift-stably-free}
设 $R$ 为环，$I \subset R$ 为理想。假设 $1 + I$ 中的每个元素都是
单位元（换言之，$I$ 包含于 $R$ 的 Jacobson 根）。对每个有限稳定自由
$R/I$-模 $E$，存在有限稳定自由 $R$-模 $M$，使得 $M/IM \cong E$。
\end{lemma}

\begin{proof}
选取 $n$、$m$ 以及一个同构
$E \oplus (R/I)^{\oplus n} \cong (R/I)^{\oplus m}$。
选取 $R$-线性映射 $\varphi : R^{\oplus m} \to R^{\oplus n}$ 和
$\psi : R^{\oplus n} \to R^{\oplus m}$，分别提升投影
$(R/I)^{\oplus m} \to (R/I)^{\oplus n}$ 与嵌入
$(R/I)^{\oplus n} \to (R/I)^{\oplus m}$。
于是 $\varphi \circ \psi : R^{\oplus n} \to R^{\oplus n}$
模 $I$ 后化为恒等映射。由关于 $I$ 的假设，该映射的行列式可逆。因此
$P = \Ker(\varphi)$ 稳定自由，并且提升 $E$。
\end{proof}

\begin{lemma}
\label{more-algebra-lemma-lift-projective}
设 $R$ 为环，$I \subset R$ 为理想。
假设 $1 + I$ 中的每个元素都是单位元
（换言之，$I$ 包含于 $R$ 的 Jacobson 根）。
设 $M$ 为有限平坦 $R$-模，并且 $M/IM$ 是投射 $R/I$-模。
则 $M$ 是有限投射 $R$-模。
\end{lemma}

\begin{proof}
由 Algebra, Lemma \ref{algebra-lemma-finite-flat-local}，
对所有素理想 $\mathfrak p \subset R$，$M_\mathfrak p$ 都是有限自由模。
由 Algebra, Lemma \ref{algebra-lemma-finite-projective}，
只需证明函数
$\rho_M : \mathfrak p \mapsto
\dim_{\kappa(\mathfrak p)} M \otimes_R \kappa(\mathfrak p)$
在 $\Spec(R)$ 上局部常值。由于 $M/IM$ 是有限投射模，这在
$V(I) \subset \Spec(R)$ 上成立。又因 $\Spec(R)$ 的每个闭点都属于
$V(I)$，且每当 $\mathfrak p \subset \mathfrak q \subset R$
为素理想时都有 $\rho_M(\mathfrak p) = \rho_M(\mathfrak q)$，
由一个此处略去的初等拓扑空间论证即可得出结论。
\end{proof}

\noindent
引理 \ref{more-algebra-lemma-lift-stably-free} 中的提升由下一个引理可知在同构意义下唯一。

\begin{lemma}
\label{more-algebra-lemma-isomorphic-finite-projective-lifts}
设 $R$ 为环，$I \subset R$ 为理想。假设 $1 + I$ 中的每个元素都是
单位元（换言之，$I$ 包含于 $R$ 的 Jacobson 根）。若 $P$ 和 $P'$
是有限投射 $R$-模，则
\begin{enumerate}
\item 若 $\varphi : P \to P'$ 是一个 $R$-模映射，并诱导同构
$\overline{\varphi} : P/IP \to P'/IP'$，则 $\varphi$ 是同构；
\item 若 $P/IP \cong P'/IP'$，则 $P \cong P'$。
\end{enumerate}
\end{lemma}

\begin{proof}
证明 (1)。由于 $P'$ 是投射 $R$-模，可选取映射
$P' \to P'/IP' \xrightarrow{\overline{\varphi}^{-1}} P/IP$
的一个提升 $\psi : P' \to P$。由 Nakayama 引理
（Algebra, Lemma \ref{algebra-lemma-NAK}），
$\psi \circ \varphi$ 和 $\varphi \circ \psi$ 都是满射。
因此这些映射都是同构
（Algebra, Lemma \ref{algebra-lemma-fun}），从而 $\varphi$ 是同构。

\medskip\noindent
证明 (2)。选取同构 $P/IP \cong P'/IP'$。由于 $P$ 是投射模，
可选取映射 $P \to P/IP \to P'/IP'$ 的一个提升
$\varphi : P \to P'$。由 (1)，$\varphi$ 是同构。
\end{proof}





\section{关于 Artin--Rees 性质的一点说明}
\label{more-algebra-section-artin-rees}

\noindent
本节部分材料取自 \cite{conrad-dejong}。在
\cite[Section 1]{Eis} 中可以找到含有更多参考文献的一般性讨论。

\medskip\noindent
设 $A$ 为 Noether 环，$I \subset A$ 为理想。给定有限 $A$-模的同态
$f : M \to N$，存在 $c \geq 0$，使得对所有 $n \geq c$ 都有
$$
f(M) \cap I^nN \subset f(I^{n - c}M)
$$
见 Algebra, Lemma \ref{algebra-lemma-map-AR}。在此情形下，我们称
{\it $c$ 对 $f$ 的 Artin--Rees 引理有效}。

\begin{lemma}
\label{more-algebra-lemma-approximate-complex}
设 $A$ 为 Noether 环，$I \subset A$ 为包含于 $A$ 的 Jacobson 根的理想。设
$$
S : L \xrightarrow{f} M \xrightarrow{g} N
\quad\text{且}\quad
S' : L \xrightarrow{f'} M \xrightarrow{g'} N
$$
为如上所示的两个有限 $A$-模复形。假设
\begin{enumerate}
\item $c$ 对 $f$ 和 $g$ 的 Artin--Rees 引理有效；
\item 复形 $S$ 正合；
\item $f' = f \bmod I^{c + 1}M$ 且 $g' = g \bmod I^{c + 1}N$。
\end{enumerate}
则 $c$ 对 $g'$ 的 Artin--Rees 引理有效，并且复形 $S'$ 正合。
\end{lemma}

\begin{proof}
先证明当 $n \geq c$ 时
$g'(M) \cap I^nN \subset g'(I^{n - c}M)$。设 $a$ 是 $M$ 的元素，
满足 $g'(a) \in I^nN$。我们希望用 $f'(L)$ 中的元素调整 $a$，
即不改变 $g'(a)$，使得 $a \in I^{n-c}M$。假设
$a \in I^rM$，其中 $r < n - c$。则
$$
g(a) = g'(a) + (g - g')(a) \in
I^n N + I^{r + c + 1}N = I^{r + c + 1}N.
$$
由 $g$ 的 Artin--Rees 性质，$g(a) \in g(I^{r + 1}M)$。设
$g(a) = g(a_1)$，其中 $a_1 \in I^{r + 1}M$。由于序列 $S$ 正合，
$a - a_1 \in f(L)$。相应地，对某个 $b \in L$，可写成
$a = f(b) + a_1$。于是 $f(b) = a - a_1 \in I^rM$。
$f$ 的 Artin--Rees 性质表明，若 $r \geq c$，则可用
$I^{r - c}L$ 中的元素替换 $b$。于是无论哪种情形，都有
$a = f'(b) + a_2$，其中
$a_2 = (f - f')(b) + a_1 \in I^{r + 1}M$。（确切地说，或是
$c \geq r$，由假设有 $(f - f')(b) \in I^{r + 1}M$；或是
$c < r$ 且 $b \in I^{r - c}$，从而仍有
$(f - f')(b) \in I^{c + 1} I^{r - c} M = I^{r + 1}M$。）
因此可用元素 $f'(b) \in f'(L)$ 调整 $a$，使 $r$ 增加 $1$。

\medskip\noindent
事实上，上述论证表明，对所有 $n \geq c$，都有
$(g')^{-1}(I^nN) \subset f'(L) + I^{n - c}M$。因此 $S'$ 正合，因为
$$
(g')^{-1}(0) = (g')^{-1}(\bigcap I^nN) \subset
\bigcap f'(L) + I^{n - c}M = f'(L)
$$
这里使用了 $I$ 包含于 $A$ 的 Jacobson 根这一事实，见
Algebra, Lemma \ref{algebra-lemma-intersection-powers-ideal-module}。
\end{proof}

\noindent
给定环 $A$ 的理想 $I \subset A$ 及 $A$-模 $M$，令
$$
\text{Gr}_I(M) = \bigoplus I^nM/I^{n + 1}M.
$$
我们把它视为一个分次 $\text{Gr}_I(A)$-模。

\begin{lemma}
\label{more-algebra-lemma-approximate-complex-graded}
假设同引理 \ref{more-algebra-lemma-approximate-complex}。令
$Q = \Coker(g)$ 和 $Q' = \Coker(g')$。则作为分次
$\text{Gr}_I(A)$-模，有
$\text{Gr}_I(Q) \cong \text{Gr}_I(Q')$。
\end{lemma}

\begin{proof}
在次数 $n$，有
$\text{Gr}_I(Q)_n = I^nN/(I^{n + 1}N + g(M) \cap I^nN)$，
$Q'$ 的情形也类似。我们断言
$$
g(M) \cap I^nN \subset I^{n + 1}N + g'(M) \cap I^nN.
$$
由对称性（断言的证明只会用到 $c$ 对 $g$ 有效，而由引理它对 $g'$
也有效），这将推出
$$
I^{n + 1}N + g(M) \cap I^nN = I^{n + 1}N + g'(M) \cap I^nN
$$
从而 $\text{Gr}_I(Q)_n$ 与 $\text{Gr}_I(Q')_n$ 作为 $N$ 的子商相同，
这就证明了引理。注意，当 $n \leq c$ 时，由
$g = g' \bmod I^{c + 1}N$，断言显然成立。若 $n > c$，设
$b \in g(M) \cap I^nN$。将其写为 $b = g(a)$，其中
$a \in I^{n - c}M$。令 $b' = g'(a)$。于是
$b - b' = (g - g')(a) \in I^{n + 1}N$，正合所需。
\end{proof}

\begin{lemma}
\label{more-algebra-lemma-works-flat-extension}
设 $A \to B$ 是 Noether 环的平坦映射，$I \subset A$ 为理想，
$f : M \to N$ 为有限 $A$-模的同态。假设 $c$ 对 $f$ 的
Artin--Rees 引理有效。则对于理想 $IB$，$c$ 对
$f \otimes 1 : M \otimes_A B \to N \otimes_A B$
的 Artin--Rees 引理也有效。
\end{lemma}

\begin{proof}
注意
$$
(f \otimes 1)(M) \cap I^n N \otimes_A B
= (f \otimes 1)\left((f \otimes 1)^{-1}(I^n N \otimes_A B)\right)
$$
另一方面，
\begin{align*}
(f \otimes 1)^{-1}(I^n N \otimes_A B) &
= \Ker(M \otimes_A B \to N \otimes_A B/(I^n N \otimes_A B)) \\
& =
\Ker(M \otimes_A B \to (N/I^nN) \otimes_A B).
\end{align*}
由于 $A \to B$ 平坦，取核与余核同张量 $B$ 可交换，故上式等于
$f^{-1}(I^nN) \otimes_A B$。由假设，$f^{-1}(I^nN)$ 包含于
$\Ker(f) + I^{n - c}M$，引理得证。
\end{proof}








\section{环的纤维积 I}
\label{more-algebra-section-fibre-products-rings}

\noindent
环的纤维积与概形的推出有关。More on Morphisms,
Section \ref{more-morphisms-section-pushouts} 讨论了概形推出的若干情形。

\begin{lemma}
\label{more-algebra-lemma-fibre-product-finite-type}
设 $R$ 为环，$A \to B$ 和 $C \to B$ 为 $R$-代数映射。假设
\begin{enumerate}
\item $R$ 是 Noether 环；
\item $A$、$B$、$C$ 在 $R$ 上为有限型；
\item $A \to B$ 是满射；
\item $B$ 在 $C$ 上有限。
\end{enumerate}
则 $A \times_B C$ 在 $R$ 上为有限型。
\end{lemma}

\begin{proof}
令 $D = A \times_B C$。有行正合的交换图
$$
\xymatrix{
0 &
B \ar[l] &
A \ar[l] &
I \ar[l] &
0 \ar[l] \\
0 &
C \ar[l] \ar[u] &
D \ar[l] \ar[u] &
I \ar[l] \ar[u] &
0 \ar[l]
}
$$
。选取 $y_1, \ldots, y_n \in B$，使其作为 $C$-模生成 $B$；
再选取映到 $y_i$ 的 $x_i \in A$。于是
$1, x_1, \ldots, x_n$ 作为 $D$-模生成 $A$。映射
$D \to A \times C$ 是单射，而环 $A \times C$ 作为 $D$-模有限
（因为它是两个有限 $D$-模 $A$ 和 $C$ 的直和）。故引理由
Artin--Tate 引理
（Algebra, Lemma \ref{algebra-lemma-Artin-Tate}）推出。
\end{proof}

\begin{lemma}
\label{more-algebra-lemma-formal-consequence}
设 $R$ 为 Noether 环，$I$ 为有限集。假设给定笛卡尔图
$$
\xymatrix{
\prod B_i &
\prod A_i \ar[l]^{\prod \varphi_i} \\
Q \ar[u]^{\prod \psi_i} &
P \ar[u] \ar[l]
}
$$
，其中 $\psi_i$ 和 $\varphi_i$ 都是满射，且
$Q$、$A_i$、$B_i$ 在 $R$ 上为有限型。则 $P$ 在 $R$ 上为有限型。
\end{lemma}

\begin{proof}
由引理 \ref{more-algebra-lemma-fibre-product-finite-type} 及对 $I$ 的大小作归纳即得。
确切地说，令 $I = I' \amalg \{i_0\}$，并令 $P'$ 为以 $I'$
代入引理中的图所定义的环。由归纳假设，$P'$ 为有限型。
最后，$P$ 位于纤维积图
$$
\xymatrix{
B_{i_0} &
A_{i_0} \ar[l] \\
P' \ar[u] &
P \ar[u] \ar[l] &
}
$$
中，可对该图应用上述引理。
\end{proof}

\begin{lemma}
\label{more-algebra-lemma-diagram-localize}
假设给定环的笛卡尔图
$$
\xymatrix{
R &
R' \ar[l]^t \\
B \ar[u]_s &
B'\ar[u] \ar[l]
}
$$
，即 $B' = B \times_R R'$。若 $h \in B'$ 对应于
$g \in B$ 与 $f \in R'$，并且 $s(g) = t(f)$，则图
$$
\xymatrix{
R_{s(g)} = R_{t(f)} &
(R')_f \ar[l]^-t \\
B_g \ar[u]_s &
(B')_h \ar[u] \ar[l]
}
$$
也是笛卡尔图。
\end{lemma}

\begin{proof}
等式 $B' = B \times_R R'$ 表明
$$
0 \to B' \to B \oplus R' \xrightarrow{s, -t} R
$$
是 $B'$-模的正合序列。作为 $B'$-模，有
$B_g = B_h$、$R'_f = R'_h$ 和 $R_{s(g)} = R_{t(f)} = R_h$。
由局部化的正合性
（Algebra, Proposition \ref{algebra-proposition-localization-exact}），得到正合序列
$$
0 \to B'_h \to B_g \oplus R'_f \xrightarrow{s, -t} R_{s(g)} = R_{t(f)}
$$
这就证明了引理。
\end{proof}

\noindent
考虑环的交换图
$$
\xymatrix{
R &
R' \ar[l] \\
B \ar[u] &
B' \ar[u] \ar[l]
}
$$
。考虑函子（其中范畴的纤维积按
Categories, Example \ref{categories-example-2-fibre-product-categories}
中的方式构造）
\begin{equation}
\label{more-algebra-equation-modules}
\text{Mod}_{B'} \longrightarrow
\text{Mod}_B \times_{\text{Mod}_R} \text{Mod}_{R'},\quad
L' \longmapsto (L' \otimes_{B'} B, L' \otimes_{B'} R', can)
\end{equation}
，其中 $can$ 是规范等同
$L' \otimes_{B'} B \otimes_B R = L' \otimes_{B'} R' \otimes_{R'} R$。
下文以 $(N, M', \varphi)$ 表示右端的对象，即 $N$ 是 $B$-模，
$M'$ 是 $R'$-模，而
$\varphi : N \otimes_B R \to M' \otimes_{R'} R$ 是同构。

\begin{lemma}
\label{more-algebra-lemma-modules}
给定环的交换图
$$
\xymatrix{
R &
R' \ar[l] \\
B \ar[u] &
B' \ar[u] \ar[l]
}
$$
，函子 (\ref{more-algebra-equation-modules}) 有右伴随，即函子
$$
F : (N, M', \varphi) \longmapsto N \times_\varphi M'
$$
（具体说明见证明）。
\end{lemma}

\begin{proof}
给定范畴
$\text{Mod}_B \times_{\text{Mod}_R} \text{Mod}_{R'}$
的对象 $(N, M', \varphi)$，令
$$
N \times_\varphi M' = \{(n, m') \in N \times M' \mid
\varphi(n \otimes 1) = m' \otimes 1\text{ 于 }M' \otimes_{R'} R\}
$$
并将其视为 $B'$-模。伴随性断言是：对 $B'$-模 $L'$ 和三元组
$(N, M', \varphi)$，有
$$
\Hom_{B'}(L', N \times_\varphi M') =
\Hom_B(L' \otimes_{B'} B, N)
\times_{\Hom_R(L' \otimes_{B'} R, M' \otimes_{R'} R)}
\Hom_{R'}(L' \otimes_{B'} R', M')
$$
由 Algebra, Lemma \ref{algebra-lemma-adjoint-tensor-restrict}，
右端等于
$$
\Hom_{B'}(L', N) \times_{\Hom_{B'}(L', M' \otimes_{R'} R)} \Hom_{B'}(L', M')
$$
因此很清楚，对此纤维积中的一对元素 $(g, f')$，可得到 $B'$-线性映射
$L' \to N \times_\varphi M'$，$l' \mapsto (g(l'), f'(l'))$。
反之，给定 $B'$-线性映射 $g' : L' \to N \times_\varphi M'$，
可令 $g$ 为复合 $L' \to N \times_\varphi M' \to N$，令 $f'$
为复合 $L' \to N \times_\varphi M' \to M'$。
这些构造互为逆，并定义所需同构。
\end{proof}






\section{环的纤维积 II}
\label{more-algebra-section-fibre-products-rings-II}

\noindent
本节讨论如下情形中的纤维积。

\begin{situation}
\label{more-algebra-situation-module-over-fibre-product}
下文考虑环映射
$$
\xymatrix{
B \ar[r] & A & A' \ar[l]
}
$$
，其中假设 $A' \to A$ 是核为 $I$ 的满射。在此情形中，令
$B' = B \times_A A'$，从而得到笛卡尔方块
$$
\xymatrix{
A & A' \ar[l] \\
B \ar[u] & B' \ar[l] \ar[u]
}
$$
。
\end{situation}

\begin{lemma}
\label{more-algebra-lemma-points-of-fibre-product}
在情形 \ref{more-algebra-situation-module-over-fibre-product} 中，作为拓扑空间有
$$
\Spec(B') = \Spec(B) \amalg_{\Spec(A)} \Spec(A')
$$
\end{lemma}

\begin{proof}
由 $B' = B \times_A A'$，得到谱的一个交换方块。由于源是拓扑空间范畴中的
推出（它由 Topology, Section \ref{topology-section-colimits} 存在），
该方块诱导连续映射
$$
can : \Spec(B) \amalg_{\Spec(A)} \Spec(A') \longrightarrow \Spec(B')
$$

\medskip\noindent
为证明映射 $can$ 是满射，设 $\mathfrak q' \subset B'$ 为素理想。
若 $I \subset \mathfrak q'$（此处及下文，我们也把 $I$ 视为 $B'$ 的理想，
而不仅是 $A'$ 的理想），则 $\mathfrak q'$ 对应于 $B$ 的一个素理想，因而
属于像。否则，选取 $h \in I$ 且 $h \not \in \mathfrak q'$。此时
$B_h = A_h = 0$，并且环映射 $B'_h \to A'_h$ 是同构，见
引理 \ref{more-algebra-lemma-diagram-localize}。因此 $\mathfrak q'$ 对应于
不包含 $I$ 的唯一素理想 $\mathfrak p' \subset A'$。

\medskip\noindent
由于 $B' \to B$ 是满射，$can$ 在分支 $\Spec(B)$ 上是单射。
上面已经看到，$\Spec(A') \to \Spec(B')$ 在
$V(I) \subset \Spec(A')$ 的补集上是单射。又因为
$V(I) \subset \Spec(A')$ 恰为 $\Spec(A) \to \Spec(A')$ 的像，
一个平凡的集合论论证表明 $can$ 是单射。

\medskip\noindent
最后还需证明 $can$ 是开映射。为此，注意推出中的开集形如
$V \amalg U'$，其中 $V \subset \Spec(B)$ 与
$U' \subset \Spec(A')$ 为开集，且它们在 $\Spec(A)$ 中的逆像相同。
设 $v \in V$。可选取 $g \in B$，使得 $v \in D(g) \subset V$。
令 $f \in A$ 为其像，并选取映到 $f$ 的 $f' \in A'$。则
$D(f') \cap U' \cap V(I) = D(f') \cap V(I)$。
因此 $V(I) \cap D(f')$ 与 $D(f') \cap (U')^c$ 是
$D(f') = \Spec(A'_{f'})$ 中互不相交的闭子集。对某个理想
$J \subset A'$，写 $(U')^c = V(J)$。由刚证明的不相交性，
映射
$A'_{f'} \to A'_{f'}/IA'_{f'} \times A'_{f'}/JA'_{f'}$
是满射，故可选取 $a'' \in A'_{f'}$，使其在
$A'_{f'}/IA'_{f'}$ 中映到 $1$，在 $A'_{f'}/JA'_{f'}$
中映到零。清除分母后，得到 $a' \in J$，它在 $A$ 中映到 $f^n$。
于是 $D(a'f') \subset U'$。令
$h' = (g^{n + 1}, a'f') \in B'$。由先前引用的引理，
$B'_{h'} = B_{g^{n + 1}} \times_{A_{f^{n + 1}}} A'_{a'f'}$，
故 $D(h')$ 在推出中的逆像是 $v$ 的一个开邻域；换言之，
$V \amalg U'$ 的像包含 $v$ 的像的一个开邻域。对于
$u' \in U'$ 且 $u' \not \in V(I)$ 的情形，结论也成立，其较为容易的证明略去。
\end{proof}

\begin{lemma}
\label{more-algebra-lemma-fibre-product-integral}
在情形 \ref{more-algebra-situation-module-over-fibre-product} 中，若
$B \to A$ 是整映射，则 $B' \to A'$ 是整映射。
\end{lemma}

\begin{proof}
设 $a' \in A'$，其像为 $a \in A$。设
$x^d + b_1 x^{d - 1} + \ldots + b_d$ 是系数在 $B$ 中且以 $a$
为根的首一多项式。选取映到 $b_i \in B$ 的 $b'_i \in B'$
（这是可能的）。于是
$(a')^d + b'_1 (a')^{d - 1} + \ldots + b'_d$
属于 $A' \to A$ 的核。由于
$\Ker(B' \to B) = \Ker(A' \to A)$，可修改 $b'_d$ 的选择，使得
$(a')^d + b'_1 (a')^{d - 1} + \ldots + b'_d = 0$，正合所需。
\end{proof}

\noindent
在情形 \ref{more-algebra-situation-module-over-fibre-product} 中，我们希望借助
$A'$、$A$ 与 $B$ 上的模来理解 $B'$-模。为此考虑函子
（其中范畴的纤维积按
Categories, Example \ref{categories-example-2-fibre-product-categories}
中的方式构造）
\begin{equation}
\label{more-algebra-equation-functor}
\text{Mod}_{B'} \longrightarrow
\text{Mod}_B \times_{\text{Mod}_A} \text{Mod}_{A'},\quad
L' \longmapsto (L' \otimes_{B'} B, L' \otimes_{B'} A', can)
\end{equation}
，其中 $can$ 是规范等同
$L' \otimes_{B'} B \otimes_B A = L' \otimes_{B'} A' \otimes_{A'} A$。
下文以 $(N, M', \varphi)$ 表示右端的对象，即 $N$ 是 $B$-模，
$M'$ 是 $A'$-模，而
$\varphi : N \otimes_B A \to M' \otimes_{A'} A$ 是同构。
不过，把 $\varphi$ 看成 $B$-线性映射
$\varphi : N \to M'/IM'$ 往往更为方便；它诱导同构
$N \otimes_B A \to M' \otimes_{A'} A = M'/IM'$。

\begin{lemma}
\label{more-algebra-lemma-module-over-fibre-product}
在情形 \ref{more-algebra-situation-module-over-fibre-product} 中，
函子 (\ref{more-algebra-equation-functor}) 有右伴随，即函子
$$
F : (N, M', \varphi) \longmapsto N \times_{\varphi, M} M'
$$
其中 $M = M'/IM'$。此外，$F$ 与 (\ref{more-algebra-equation-functor}) 的复合
是 $\text{Mod}_B \times_{\text{Mod}_A} \text{Mod}_{A'}$
上的恒等函子。换言之，令
$N' = N \times_{\varphi, M} M'$，则
$N' \otimes_{B'} B = N$ 且 $N' \otimes_{B'} A' = M'$。
\end{lemma}

\begin{proof}
伴随性断言由更一般的引理 \ref{more-algebra-lemma-modules} 推出。
为证明最后一个断言，回忆
$B' = B \times_A A'$ 和 $N' = N \times_{\varphi, M} M'$，
并把这些等式扩充为
$$
\vcenter{
\xymatrix{
A & A' \ar[l] & I \ar[l] \\
B \ar[u] & B' \ar[l] \ar[u] & J \ar[l] \ar[u]
}
}
\quad\text{且}\quad
\vcenter{
\xymatrix{
M & M' \ar[l] & K \ar[l] \\
N \ar[u]_\varphi & N' \ar[l] \ar[u] & L \ar[l] \ar[u]
}
}
$$
，其中 $I, J, K, L$ 分别是原来两个图中各水平映射的核。
把证明列为一系列观察：
\begin{enumerate}
\item $K = IM'$（见引理的陈述）；
\item $B' \to B$ 是核为 $J$ 的满射，且 $J \to I$ 是双射；
\item $N' \to N$ 是核为 $L$ 的满射，且 $L \to K$ 是双射；
\item $JN' \subset L$；
\item $\Im(N \to M)$ 作为 $A$-模生成 $M$
（因为 $N \otimes_B A = M$）；
\item $\Im(N' \to M')$ 作为 $A'$-模生成 $M'$
（因为模 $K$ 后这一点成立，且 $L$ 同构地映到 $K$）；
\item $JN' = L$（因为由上一项，$L \cong K = I M'$
由形如 $x n'$ 的元素的像生成，其中 $x \in I$ 且 $n' \in N'$）；
\item $N' \otimes_{B'} B = N$（因为 $N = N'/L$、$B = B'/J$
及上一项）；
\item 存在映射 $\gamma : N' \otimes_{B'} A' \to M'$；
\item $\gamma$ 是满射（见上文）；
\item 复合映射 $N' \otimes_{B'} A' \to M' \to M$ 的核
由形如 $l \otimes 1$ 和 $n' \otimes x$ 的元素生成，其中
$l \in K$、$n' \in N'$、$x \in I$（因为由假设
$M = N \otimes_B A$，并且 $N' \to N$ 与 $A' \to A$
分别是核为 $L$ 和 $I$ 的满射）；
\item $N' \otimes_{B'} A'$ 中由形如
$l \otimes 1$ 和 $n' \otimes x$ 的元素生成的子模的任意元素，
其中 $l \in L$、$n' \in N'$、$x \in I$，都可写成某个
$l \in L$ 的 $l \otimes 1$（因为 $J$ 同构地映到 $I$，
在 $N' \otimes_{B'} A'$ 中有
$n' \otimes x = n'x \otimes 1$；类似地，当
$n' \in N'$、$x \in J$、$a' \in A'$ 时，有
$x n' \otimes a' = n' \otimes xa' = n'(xa') \otimes 1$；
该等式成立于 $N' \otimes_{B'} A'$ 中；
再由已经证明的 $JN' = L$，便得断言）；
\item $\gamma$ 的核为零（因为由 (10) 和 (11)，核的每个元素都形如
$l \otimes 1$，其中 $l \in L$，而 $\gamma$ 将它映到
$l \in K \subset M'$）。
\end{enumerate}
证明完成。
\end{proof}

\begin{lemma}
\label{more-algebra-lemma-module-over-fibre-product-bis}
在引理 \ref{more-algebra-lemma-module-over-fibre-product} 的情形中，对 $B'$-模
$L'$，伴随映射
$$
L' \longrightarrow
(L' \otimes_{B'} B) \times_{(L' \otimes_{B'} A)} (L' \otimes_{B'} A')
$$
是满射，但一般不是单射。
\end{lemma}

\begin{proof}
与引理 \ref{more-algebra-lemma-module-over-fibre-product} 的证明相同，令
$J \subset B'$ 为映射 $B' \to B$ 的核。则
$L' \otimes_{B'} B = L'/JL'$。因此，为证明满性，只需说明纤维积中
形如 $(0, z)$ 的元素属于该引理映射的像。
映射 $L' \otimes_{B'} A' \to L' \otimes_{B'} A$ 的核是
$L' \otimes_{B'} I \to L' \otimes_{B'} A'$ 的像。
由于 $B' \to A'$ 诱导的映射 $J \to I$ 是同构，复合
$$
L' \otimes_{B'} J \to L' \to
(L' \otimes_{B'} B) \times_{(L' \otimes_{B'} A)} (L' \otimes_{B'} A')
$$
将 $L' \otimes_{B'} J$ 满射到形如 $(0, z)$ 的元素集合。
为看出此映射一般不是单射，给出一个简单例子。取域 $k$，令
$B' = k[x, y]/(xy)$、$A' = B'/(x)$、$B = B'/(y)$、
$A = B'/(x, y)$ 以及 $L' = B'/(x - y)$。此时 $L'$ 中 $x$
的类非零，但在上面展示的箭头下映到零。
\end{proof}

\begin{lemma}
\label{more-algebra-lemma-surjection-module-over-fibre-product}
在情形 \ref{more-algebra-situation-module-over-fibre-product} 中，设
$(N_1, M'_1, \varphi_1) \to (N_2, M'_2, \varphi_2)$
是 $\text{Mod}_B \times_{\text{Mod}_A} \text{Mod}_{A'}$
中的态射，并且 $N_1 \to N_2$ 和 $M'_1 \to M'_2$ 都是满射。则
$$
N_1 \times_{\varphi_1, M_1} M'_1 \to N_2 \times_{\varphi_2, M_2} M'_2
$$
是满射，其中 $M_1 = M'_1/IM'_1$ 且 $M_2 = M'_2/IM'_2$。
\end{lemma}

\begin{proof}
选取 $(x_2, y_2) \in N_2 \times_{\varphi_2, M_2} M'_2$。
选取映到 $x_2$ 的 $x_1 \in N_1$。由于 $M'_1 \to M_1$ 是满射，
可选取映到 $\varphi_1(x_1)$ 的 $y_1 \in M'_1$。于是
$(x_1, y_1)$ 在 $N_2 \times_{\varphi_2, M_2} M'_2$ 中映到
$(x_2, y'_2)$。因此只需证明形如 $(0, y_2)$ 的元素属于该映射的像。
这里 $y_2 \in IM'_2$。写
$y_2 = \sum t_i y_{2, i}$，其中 $t_i \in I$。选取映到
$y_{2, i}$ 的 $y_{1, i} \in M'_1$。于是
$y_1 = \sum t_iy_{1, i} \in IM'_1$，而元素 $(0, y_1)$ 满足要求。
\end{proof}

\begin{lemma}
\label{more-algebra-lemma-finite-module-over-fibre-product}
设 $A, A', B, B', I, M, M', N, \varphi$ 如引理
\ref{more-algebra-lemma-module-over-fibre-product}。若 $N$ 在 $B$ 上有限，且
$M'$ 在 $A'$ 上有限，则
$N' = N \times_{\varphi, M} M'$ 在 $B'$ 上有限。
\end{lemma}

\begin{proof}
下文将不再说明地使用引理 \ref{more-algebra-lemma-module-over-fibre-product}
的结论。选取 $N$ 在 $B$ 上的一组生成元 $y_1, \ldots, y_r$，
以及 $M'$ 在 $A'$ 上的一组生成元 $x_1, \ldots, x_s$。
利用 $N = N' \otimes_{B'} B$ 以及 $B' \to B$ 为满射，
可选取 $u_1, \ldots, u_r \in N'$，它们在 $N$ 中分别映到
$y_1, \ldots, y_r$。利用 $M' = N' \otimes_{B'} A'$，
可选取 $v_1, \ldots, v_t \in N'$，使得对某些
$a'_{ij} \in A'$ 有 $x_i = \sum v_j \otimes a'_{ij}$。
特别地，$v_j$ 的像 $\overline{v}_j \in M'$ 作为 $A'$-模生成 $M'$。
我们断言 $u_1, \ldots, u_r, v_1, \ldots, v_t$
作为 $B'$-模生成 $N'$。事实上，取 $\xi \in N'$。
先选取 $b'_1, \ldots, b'_r \in B'$，使得 $\xi$ 与
$\sum b'_i u_i$ 在 $N$ 中映到同一个元素。这是可能的，因为
$B' \to B$ 是满射，且 $y_1, \ldots, y_r$ 在 $B$ 上生成 $N$。
差 $\xi - \sum b'_i u_i$ 形如 $(0, \theta)$，其中某个
$\theta$ 属于 $IM'$。设 $\theta$ 为 $\sum t_j\overline{v}_j$，其中
$t_j \in I$。由于 $J = \Ker(B' \to B)$ 同构地映到 $I$，
可选取映到 $t_j$ 的 $s_j \in J \subset B'$。
由 $N' = N \times_{\varphi, M} M'$，得到
$\xi = \sum b'_i u_i + \sum s_j v_j$，正合所需。
\end{proof}

\begin{lemma}
\label{more-algebra-lemma-flat-module-over-fibre-product}
设 $A, A', B, B', I$ 如情形
\ref{more-algebra-situation-module-over-fibre-product}。
\begin{enumerate}
\item 设 $(N, M', \varphi)$ 是
$\text{Mod}_B \times_{\text{Mod}_A} \text{Mod}_{A'}$ 的对象。
若 $M'$ 在 $A'$ 上平坦，且 $N$ 在 $B$ 上平坦，则
$N' = N \times_{\varphi, M} M'$ 在 $B'$ 上平坦。
\item 若 $L'$ 是平坦 $B'$-模，则
$L' = (L \otimes_{B'} B) \times_{(L \otimes_{B'} A)} (L \otimes_{B'} A')$。
\item 平坦 $B'$-模的范畴等价于
$\text{Mod}_B \times_{\text{Mod}_A} \text{Mod}_{A'}$
的满子范畴；该满子范畴由满足 $N$ 在 $B$ 上平坦且 $M'$
在 $A'$ 上平坦的三元组 $(N, M', \varphi)$ 组成。
\end{enumerate}
\end{lemma}

\begin{proof}
证明中将不再说明地使用引理
\ref{more-algebra-lemma-module-over-fibre-product}。

\medskip\noindent
证明 (1)。令 $J = \Ker(B' \to B)$。这是 $B'$ 的理想，并同构地映到
$I = \Ker(A' \to A)$。设 $\mathfrak b' \subset B'$ 为理想。
必须证明 $\mathfrak b' \otimes_{B'} N' \to N'$ 是单射，见
Algebra, Lemma \ref{algebra-lemma-flat}。我们知道
$$
\mathfrak b'/(\mathfrak b' \cap J) \otimes_{B'} N' =
\mathfrak b'/(\mathfrak b' \cap J) \otimes_B N \to N
$$
是单射，因为 $N$ 在 $B$ 上平坦。由于
$\mathfrak b' \cap J \to \mathfrak b' \to
\mathfrak b'/(\mathfrak b' \cap J) \to 0$ 正合，可知只需证明
$(\mathfrak b' \cap J) \otimes_{B'} N' \to N'$ 是单射。
因此可假设 $\mathfrak b' \subset J$。其次，由于
$J \to I$ 是同构，有
$$
J \otimes_{B'} N' =
I \otimes_{A'} A' \otimes_{B'} N' =
I \otimes_{A'} M'
$$
而它单射地映入 $M'$，因为 $M'$ 是平坦 $A'$-模。因此
$J \otimes_{B'} N' \to N'$ 是单射，从而
$\text{Tor}_1^{B'}(B'/J, N') = 0$，见
Algebra, Remark \ref{algebra-remark-Tor-ring-mod-ideal}。
于是可将 Algebra, Lemma
\ref{algebra-lemma-what-does-it-mean} 应用于 $B'$ 上的 $N'$
及理想 $J$。
回到理想 $\mathfrak b' \subset J$，令
$\mathfrak b' \subset \mathfrak b'' \subset J$ 为其像在 $I$ 中
是 $I$ 的一个 $A'$-子模的最小理想。换言之，若把 $J = I$ 视为
$A'$-模，则 $\mathfrak b'' = A' \mathfrak b'$。于是
$\mathfrak b''/\mathfrak b'$ 被 $J$ 零化，并且由应用上述引理所得的
$\text{Tor}_1^{B'}(\mathfrak b''/\mathfrak b', N')$ 的消失，得到短正合序列
$$
0 \to \mathfrak b' \otimes_{B'} N' \to
\mathfrak b'' \otimes_{B'} N' \to
\mathfrak b''/\mathfrak b' \otimes_{B'} N' \to 0
$$
因此可用 $\mathfrak b''$ 代替 $\mathfrak b'$。特别地，可假设
$\mathfrak b'$ 是 $A'$-模，并映到 $A'$ 的一个理想。此时
$$
\mathfrak b' \otimes_{B'} N' =
\mathfrak b' \otimes_{A'} A' \otimes_{B'} N' =
\mathfrak b' \otimes_{A'} M'
$$
由 $M'$ 在 $A'$ 上平坦这一假设，该张量积单射地映入 $M'$。
由此得到 $\mathfrak b' \otimes_{B'} N' \to N' \to M'$ 是单射，
故第一个映射如所需也是单射。

\medskip\noindent
证明 (2)。将短正合序列
$0 \to B' \to B \oplus A' \to A \to 0$
在 $B'$ 上与 $L'$ 张量即可。

\medskip\noindent
证明 (3)。这是 (1) 与 (2) 的直接推论。
\end{proof}

\begin{lemma}
\label{more-algebra-lemma-finitely-presented-module-over-fibre-product}
设 $A, A', B, B', I$ 如情形
\ref{more-algebra-situation-module-over-fibre-product}。有限投射 $B'$-模的范畴
等价于 $\text{Mod}_B \times_{\text{Mod}_A} \text{Mod}_{A'}$
的满子范畴；该满子范畴由满足 $N$ 在 $B$ 上有限投射且 $M'$
在 $A'$ 上有限投射的三元组 $(N, M', \varphi)$ 组成。
\end{lemma}

\begin{proof}
回忆模有限投射当且仅当它有限表示且平坦，见
Algebra, Lemma \ref{algebra-lemma-finite-projective}。
利用引理 \ref{more-algebra-lemma-flat-module-over-fibre-product} 和
\ref{more-algebra-lemma-finite-module-over-fibre-product}，问题归结为证明：
若 $N$ 在 $B$ 上有限投射，且 $M'$ 在 $A'$ 上有限投射，则
$N' = N \times_{\varphi, M} M'$ 是有限表示的 $B'$-模。

\medskip\noindent
由引理 \ref{more-algebra-lemma-finite-module-over-fibre-product}，模 $N'$
在 $B'$ 上有限。选取核为 $K'$ 的满射
$(B')^{\oplus n} \to N'$。基变换给出映射
$B^{\oplus n} \to N$、$(A')^{\oplus n} \to M'$ 和
$A^{\oplus n} \to M$，其核分别为 $K_B$、$K_{A'}$ 和 $K_A$。
存在规范映射
$$
K' \longrightarrow K_B \times_{K_A} K_{A'}
$$
另一方面，由 $N' = N \times_{\varphi, M} M'$ 及
$B' = B \times_A A'$，也存在规范映射
$K_B \times_{K_A} K_{A'} \to K'$，它是上述箭头的逆。
故所展示的映射是同构。由
Algebra, Lemma \ref{algebra-lemma-extension}，模 $K_B$ 与
$K_{A'}$ 都有限。由引理
\ref{more-algebra-lemma-finite-module-over-fibre-product} 可知，只要
$K_B \to K_A$ 与 $K_{A'} \to K_A$ 诱导同构
$K_B \otimes_B A = K_A = K_{A'} \otimes_{A'} A$，
$K'$ 就是有限 $B'$-模。这一点成立，因为平坦性假设蕴含序列
$$
0 \to K_B \to B^{\oplus n} \to N \to 0
\quad\text{且}\quad
0 \to K_{A'} \to (A')^{\oplus n} \to M' \to 0
$$
在张量后仍然正合，见
Algebra, Lemma \ref{algebra-lemma-flat-tor-zero}。
\end{proof}

\section{环的纤维积 III}
\label{more-algebra-section-fibre-products-rings-III}

\noindent
本节讨论如下情形中的纤维积。

\begin{situation}
\label{more-algebra-situation-relative-module-over-fibre-product}
设 $A, A', B, B', I$ 如情形
\ref{more-algebra-situation-module-over-fibre-product}。设 $B' \to D'$ 为环映射。
令 $D = D' \otimes_{B'} B$、
$C' = D' \otimes_{B'} A'$ 和 $C = D' \otimes_{B'} A$。
由此得到环的一个大交换图
$$
\xymatrix{
C & & & C' \ar[lll] \\
& A \ar[ul] & A' \ar[l] \ar[ru] \\
& B \ar[u] \ar[ld] & B' \ar[l] \ar[u] \ar[rd] \\
D \ar[uuu] & & & D' \ar[lll] \ar[uuu]
}
$$
。注意，我们{\bf 不}假设映射 $D' \to D \times_C C'$ 是同构
\footnote{但由引理 \ref{more-algebra-lemma-module-over-fibre-product-bis}，
$D' \to D \times_C C'$ 是满射。}。
在此情形中，有与 (\ref{more-algebra-equation-functor}) 类似的函子
\begin{equation}
\label{more-algebra-equation-relative-functor}
\text{Mod}_{D'} \longrightarrow
\text{Mod}_D \times_{\text{Mod}_C} \text{Mod}_{C'},\quad
L' \longmapsto (L' \otimes_{D'} D, L' \otimes_{D'} C', can).
\end{equation}
注意
$L' \otimes_{D'} D = L \otimes_{D'} (D' \otimes_{B'} B) = L \otimes_{B'} B$
以及类似地
$L' \otimes_{D'} C' = L \otimes_{D'} (D' \otimes_{B'} A') = L \otimes_{B'} A'$，
故图
$$
\xymatrix{
\text{Mod}_{D'} \ar[r] \ar[d] &
\text{Mod}_D \times_{\text{Mod}_C} \text{Mod}_{C'} \ar[d] \\
\text{Mod}_{B'} \ar[r] &
\text{Mod}_B \times_{\text{Mod}_A} \text{Mod}_{A'}
}
$$
交换。下文以 $(N, M', \varphi)$ 表示
$\text{Mod}_D \times_{\text{Mod}_C} \text{Mod}_{C'}$ 的一个对象，
即 $N$ 是 $D$-模，$M'$ 是 $C'$-模，而
$\varphi : N \otimes_B A \to M' \otimes_{A'} A$
是 $C$-模同构。
不过，把 $\varphi$ 看成 $D$-线性映射
$\varphi : N \to M'/IM'$ 往往更为方便；它诱导同构
$N \otimes_B A \to M' \otimes_{A'} A = M'/IM'$。
\end{situation}

\begin{lemma}
\label{more-algebra-lemma-relative-module-over-fibre-product}
在情形 \ref{more-algebra-situation-relative-module-over-fibre-product} 中，
函子 (\ref{more-algebra-equation-relative-functor}) 有右伴随，即函子
$$
F : (N, M', \varphi) \longmapsto N \times_{\varphi, M} M'
$$
，其中 $M = M'/IM'$。此外，$F$ 与
(\ref{more-algebra-equation-relative-functor}) 的复合是
$\text{Mod}_D \times_{\text{Mod}_C} \text{Mod}_{C'}$
上的恒等函子。换言之，令
$N' = N \times_{\varphi, M} M'$，则
$N' \otimes_{D'} D = N$ 且 $N' \otimes_{D'} C' = M'$。
\end{lemma}

\begin{proof}
伴随性断言由更一般的引理 \ref{more-algebra-lemma-modules} 推出。
最后一个断言由引理 \ref{more-algebra-lemma-module-over-fibre-product}
中的相应断言推出，因为
$N' \otimes_{D'} D = N' \otimes_{D'} D' \otimes_{B'} B = N' \otimes_{B'} B$
且
$N' \otimes_{D'} C' = N' \otimes_{D'} D' \otimes_{B'} A' = N' \otimes_{B'} A'$。
\end{proof}

\begin{lemma}
\label{more-algebra-lemma-relative-surjection-ideals}
在情形 \ref{more-algebra-situation-relative-module-over-fibre-product} 中，映射
$JD' \to IC'$ 是满射，其中 $J = \Ker(B' \to B)$。
\end{lemma}

\begin{proof}
由于 $C' = D' \otimes_{B'} A'$，$IC'$ 是映射
$D' \otimes_{B'} I = C' \otimes_{A'} I \to C'$ 的像。
环映射 $B' \to A'$ 诱导同构 $J \to I$，故结论成立。
\end{proof}

\begin{lemma}
\label{more-algebra-lemma-relative-finite-module-over-fibre-product}
设 $A, A', B, B', C, C', D, D', I, M', M, N, \varphi$
如引理 \ref{more-algebra-lemma-relative-module-over-fibre-product}。
若 $N$ 在 $D$ 上有限，且 $M'$ 在 $C'$ 上有限，则
$N' = N \times_{\varphi, M} M'$ 在 $D'$ 上有限。
\end{lemma}

\begin{proof}
回忆由引理 \ref{more-algebra-lemma-module-over-fibre-product-bis}，
$D' \to D \times_C C'$ 是满射。注意
$N' = N \times_{\varphi, M} M'$ 是 $D \times_C C'$ 上的模。
把引理 \ref{more-algebra-lemma-finite-module-over-fibre-product} 应用于数据
$C, C', D, D', IC', M', M, N, \varphi$，可见
$N' = N \times_{\varphi, M} M'$ 在 $D \times_C C'$ 上有限，
故在 $D'$ 上也有限。
\end{proof}

\begin{lemma}
\label{more-algebra-lemma-relative-flat-module-over-fibre-product}
设 $A, A', B, B', C, C', D, D', I$ 如情形
\ref{more-algebra-situation-relative-module-over-fibre-product}。
\begin{enumerate}
\item 设 $(N, M', \varphi)$ 是
$\text{Mod}_D \times_{\text{Mod}_C} \text{Mod}_{C'}$ 的对象。
若 $M'$ 在 $A'$ 上平坦，且 $N$ 在 $B$ 上平坦，则
$N' = N \times_{\varphi, M} M'$ 在 $B'$ 上平坦。
\item 若 $L'$ 是在 $B'$ 上平坦的 $D'$-模，则
$L' = (L \otimes_{D'} D) \times_{(L \otimes_{D'} C)} (L \otimes_{D'} C')$。
\item 在 $B'$ 上平坦的 $D'$-模的范畴等价于
$\text{Mod}_D \times_{\text{Mod}_C} \text{Mod}_{C'}$
中满足 $N$ 在 $B$ 上平坦且 $M'$ 在 $A'$ 上平坦的对象
$(N, M', \varphi)$ 所成的范畴。
\end{enumerate}
\end{lemma}

\begin{proof}
第 (1) 项由引理 \ref{more-algebra-lemma-flat-module-over-fibre-product}
的第 (1) 项推出。

\medskip\noindent
第 (2) 项由引理 \ref{more-algebra-lemma-flat-module-over-fibre-product}
的第 (2) 项推出；这里使用
$L' \otimes_{D'} D = L' \otimes_{B'} B$、
$L' \otimes_{D'} C' = L' \otimes_{B'} A'$ 以及
$L' \otimes_{D'} C = L' \otimes_{B'} A$，见情形
\ref{more-algebra-situation-relative-module-over-fibre-product} 中的讨论。

\medskip\noindent
第 (3) 项是 (1) 与 (2) 的直接推论。
\end{proof}

\noindent
下面的引理比其绝对情形中的对应结论
（引理 \ref{more-algebra-lemma-finitely-presented-module-over-fibre-product}）
有趣得多，尽管证明本质上相同。

\begin{lemma}
\label{more-algebra-lemma-relative-finitely-presented-module-over-fibre-product}
设 $A, A', B, B', C, C', D, D', I, M', M, N, \varphi$
如引理 \ref{more-algebra-lemma-relative-module-over-fibre-product}。若
\begin{enumerate}
\item $N$ 在 $D$ 上有限表示，并在 $B$ 上平坦；
\item $M'$ 在 $C'$ 上有限表示，并在 $A'$ 上平坦；
\item 环映射 $B' \to D'$ 可分解为 $B' \to D'' \to D'$，
其中 $B' \to D''$ 平坦，而 $D'' \to D'$ 为有限表示；
\end{enumerate}
则 $N' = N \times_M M'$ 在 $D'$ 上有限表示。
\end{lemma}

\begin{proof}
选取核为有限生成理想 $J$ 的满射
$D''' = D''[x_1, \ldots, x_n] \to D'$。
由 Algebra, Lemma
\ref{algebra-lemma-finite-finitely-presented-extension}，
只需证明 $N'$ 作为 $D'''$-模有限表示。此外，
$D''' \otimes_{B'} B \to D' \otimes_{B'} B = D$ 以及
$D''' \otimes_{B'} A' \to D' \otimes_{B'} A' = C'$ 都是满射，
其核由 $J$ 的像生成。因此，再次应用
Algebra, Lemma
\ref{algebra-lemma-finite-finitely-presented-extension}，
$N$ 是有限表示的 $D''' \otimes_{B'} B$-模，而
$M'$ 是有限表示的 $D''' \otimes_{B'} A'$-模。
故可用 $D'''$ 代替 $D'$，并用 $D''' \otimes_{B'} B$
代替 $D$，其余类似。由于 $D'''$ 在 $B'$ 上平坦，
可归结到 $B' \to D'$ 平坦的情形。

\medskip\noindent
假设 $B' \to D'$ 平坦。由引理
\ref{more-algebra-lemma-relative-finite-module-over-fibre-product}，
模 $N'$ 在 $D'$ 上有限。选取核为 $K'$ 的满射
$(D')^{\oplus n} \to N'$。基变换给出映射
$D^{\oplus n} \to N$、$(C')^{\oplus n} \to M'$ 和
$C^{\oplus n} \to M$，其核分别为 $K_D$、$K_{C'}$ 和 $K_C$。
存在规范映射
$$
K' \longrightarrow K_D \times_{K_C} K_{C'}
$$
。另一方面，由于 $N' = N \times_M M'$，且
$D' = D \times_C C'$（把引理
\ref{more-algebra-lemma-flat-module-over-fibre-product} 应用于平坦 $B'$-模 $D'$），
还存在规范映射 $K_D \times_{K_C} K_{C'} \to K'$，
它是所展示箭头的逆。故该映射是同构。
由 Algebra, Lemma \ref{algebra-lemma-extension}，
模 $K_D$ 与 $K_{C'}$ 都有限。由引理
\ref{more-algebra-lemma-relative-finite-module-over-fibre-product}，只要
$K_D \to K_C$ 与 $K_{C'} \to K_C$ 诱导同构
$K_D \otimes_B A = K_C = K_{C'} \otimes_{A'} A$，
$K'$ 就是有限 $D'$-模。这一点成立，因为平坦性假设蕴含序列
$$
0 \to K_D \to D^{\oplus n} \to N \to 0
\quad\text{且}\quad
0 \to K_{C'} \to (C')^{\oplus n} \to M' \to 0
$$
在张量后仍然正合，见
Algebra, Lemma \ref{algebra-lemma-flat-tor-zero}。
\end{proof}

\begin{lemma}
\label{more-algebra-lemma-properties-algebras-over-fibre-product}
设 $A, A', B, B', I$ 如情形
\ref{more-algebra-situation-module-over-fibre-product}。设 $(D, C', \varphi)$
是由一个 $B$-代数 $D$、一个 $A'$-代数 $C'$ 以及一个同构
$D \otimes_B A \to C'/IC' = C$ 所组成的系统。令
$D' = D \times_C C'$（如引理
\ref{more-algebra-lemma-module-over-fibre-product}）。则
\begin{enumerate}
\item $B' \to D'$ 为有限型，当且仅当 $B \to D$ 与
$A' \to C'$ 为有限型；
\item $B' \to D'$ 平坦，当且仅当 $B \to D$ 与 $A' \to C'$ 平坦；
\item $B' \to D'$ 平坦且为有限表示，当且仅当
$B \to D$ 与 $A' \to C'$ 平坦且为有限表示；
\item $B' \to D'$ 光滑，当且仅当 $B \to D$ 与 $A' \to C'$ 光滑；
\item $B' \to D'$ 是 \'etale 的，当且仅当 $B \to D$ 与
$A' \to C'$ 是 \'etale 的。
\end{enumerate}
此外，若 $D'$ 是平坦 $B'$-代数，则
$D' \to (D' \otimes_{B'} B) \times_{(D' \otimes_{B'} A)} (D' \otimes_{B'} A')$
是同构。如此，平坦 $B'$-代数的范畴等价于上述系统
$(D, C', \varphi)$ 所成的范畴，其中 $D$ 在 $B$ 上平坦，而
$C'$ 在 $A'$ 上平坦。
\end{lemma}

\begin{proof}
蕴含关系“$\Rightarrow$”由下列引理推出：
Algebra, Lemmas \ref{algebra-lemma-base-change-finiteness},
\ref{algebra-lemma-flat-base-change},
\ref{algebra-lemma-base-change-smooth}, and
\ref{algebra-lemma-etale}；这是因为由引理
\ref{more-algebra-lemma-module-over-fibre-product} 有
$D' \otimes_{B'} B = D$ 及 $D' \otimes_{B'} A' = C'$。
因此只需证明反方向的蕴含关系。

\medskip\noindent
关于 (1)。假设 $D$ 在 $B$ 上为有限型，且 $C'$ 在 $A'$ 上为有限型。
下文将不再说明地使用引理
\ref{more-algebra-lemma-module-over-fibre-product} 的结论。
选取 $D$ 在 $B$ 上的一组生成元 $x_1, \ldots, x_r$，
以及 $C'$ 在 $A'$ 上的一组生成元 $y_1, \ldots, y_s$。
利用 $D = D' \otimes_{B'} B$ 以及 $B' \to B$ 为满射，
可选取 $u_1, \ldots, u_r \in D'$，它们在 $D$ 中分别映到
$x_1, \ldots, x_r$。利用 $C' = D' \otimes_{B'} A'$，
可选取 $v_1, \ldots, v_t \in D'$，使得对某些
$a'_{ij} \in A'$ 有 $y_i = \sum v_j \otimes a'_{ij}$。
特别地，$v_j$ 在 $C'$ 中的像作为 $A'$-代数生成 $C'$。
令 $N = r + t$，并考虑环的立方图
$$
\xymatrix{
A[x_1, \ldots, x_N] & & A'[x_1, \ldots, x_N] \ar[ll] \\
& A \ar[lu] & & A' \ar[ll] \ar[lu] \\
B[x_1, \ldots, x_N] \ar[uu] & & B'[x_1, \ldots, x_N] \ar[uu] \ar[ll] \\
& B \ar[uu] \ar[lu] & & B' \ar[ll] \ar[uu] \ar[lu]
}
$$
。注意后面的方块也是笛卡尔方块。考虑环映射
$$
B'[x_1, \ldots, x_N] \to D',\quad
x_i \mapsto u_i \quad\text{且}\quad x_{r + j} \mapsto v_j.
$$
于是所诱导的映射 $B[x_1, \ldots, x_N] \to D$ 与
$A'[x_1, \ldots, x_N] \to C'$ 都是满射，特别地都是有限映射。
由引理 \ref{more-algebra-lemma-relative-finite-module-over-fibre-product}，
$B'[x_1, \ldots, x_N] \to D'$ 是有限映射。因而 $D'$
在 $B'$ 上为有限型，例如可由
Algebra, Lemma \ref{algebra-lemma-compose-finite-type} 得到。

\medskip\noindent
关于 (2)。蕴含关系“$\Leftarrow$”由引理
\ref{more-algebra-lemma-relative-flat-module-over-fibre-product} 推出。
此外，最后一个断言由引理
\ref{more-algebra-lemma-relative-flat-module-over-fibre-product} 的最后一个断言推出。

\medskip\noindent
关于 (3)。假设 $B \to D$ 和 $A' \to C'$ 平坦且为有限表示。
$B' \to D'$ 的平坦性已经在 (2) 中得到。由 (1)，
$B' \to D'$ 为有限型。选取满射
$B'[x_1, \ldots, x_N] \to D'$。
由 Algebra, Lemma
\ref{algebra-lemma-finite-presentation-independent}，环 $D$
作为 $B[x_1, \ldots, x_N]$-模有限表示，而环 $C'$
作为 $A'[x_1, \ldots, x_N]$-模有限表示。
由引理
\ref{more-algebra-lemma-relative-finitely-presented-module-over-fibre-product}，
$D'$ 作为 $B'[x_1, \ldots, x_N]$-模有限表示，即
$B' \to D'$ 为有限表示。

\medskip\noindent
关于 (4)。假设 $B \to D$ 与 $A' \to C'$ 光滑。由 (3)，
$B' \to D'$ 平坦且为有限表示。由 Algebra, Lemma
\ref{algebra-lemma-flat-fibre-smooth}，只需对 $B'$ 上任意域 $k$
验证 $D' \otimes_{B'} k$ 光滑。若复合 $J \to B' \to k$ 为零，
则 $B' \to k$ 分解为 $B' \to B \to k$，且
$$
D' \otimes_{B'} k = D' \otimes_{B'} B \otimes_B k
= D \otimes_B k
$$
光滑，因为 $B \to D$ 光滑。若复合 $J \to B' \to k$ 非零，
则存在 $h \in J$，其在 $k$ 中的像非零。于是 $B' \to k$
分解为 $B' \to B'_h \to k$。注意 $h$ 在 $B$ 中映到零，故
$B_h = 0$。因此由引理 \ref{more-algebra-lemma-diagram-localize}，
$B'_h = A'_h$，并且
$$
D' \otimes_{B'} k = D' \otimes_{B'} B'_h \otimes_{B'_h} k
= C'_h \otimes_{A'_h} k
$$
光滑，因为 $A' \to C'$ 光滑。

\medskip\noindent
关于 (5)。假设 $B \to D$ 与 $A' \to C'$ 是 \'etale 的。
由 (4)，$B' \to D'$ 光滑。光滑映射是否为 \'etale 可由纤维的维数读出，
故 (5) 成立（仿照 (4) 的证明识别纤维；略去若干细节）。
\end{proof}

\begin{remark}
\label{more-algebra-remark-relative-modules-over-fibre-product}
在情形 \ref{more-algebra-situation-relative-module-over-fibre-product} 中，
假设 $B' \to D'$ 为有限表示，并设给定 $D'$-模 $L'$。
我们断言下列两类对象之间存在双射：
\begin{enumerate}
\item $D'$-模的满射 $L' \to Q'$，其中 $Q'$ 在 $D'$ 上有限表示，
且在 $B'$ 上平坦；
\item 模的成对满射
$(L' \otimes_{D'} D \to Q_1, L' \otimes_{D'} C' \to Q_2)$，
满足
\begin{enumerate}
\item $Q_1$ 在 $D$ 上有限表示，并在 $B$ 上平坦；
\item $Q_2$ 在 $C'$ 上有限表示，并在 $A'$ 上平坦；
\item 作为 $L' \otimes_{D'} C$ 的商，有
$Q_1 \otimes_D C = Q_2 \otimes_{C'} C$。
\end{enumerate}
\end{enumerate}
从前者到后者的对应由 $Q \mapsto (Q_1, Q_2)$ 给出，其中
$Q_1 = Q \otimes_{D'} D$ 且 $Q_2 = Q \otimes_{D'} C'$。
反过来，使用 $Q = Q_1 \times_{Q_{12}} Q_2$，其中 $Q_{12}$ 是
$L' \otimes_{D'} C$ 的公共商
$Q_1 \otimes_D C = Q_2 \otimes_{C'} C$。作为商映射，使用
$$
L' \longrightarrow
(L' \otimes_{D'} D) \times_{(L' \otimes_{D'} C)} (L' \otimes_{D'} C')
\longrightarrow Q_1 \times_{Q_{12}} Q_2 = Q
$$
，其中第一个箭头由引理
\ref{more-algebra-lemma-module-over-fibre-product-bis} 为满射，第二个箭头由引理
\ref{more-algebra-lemma-surjection-module-over-fibre-product} 为满射。
断言由引理 \ref{more-algebra-lemma-relative-flat-module-over-fibre-product} 与
\ref{more-algebra-lemma-relative-finitely-presented-module-over-fibre-product} 推出。
\end{remark}








\section{Fitting 理想}
\label{more-algebra-section-fitting-ideals}

\noindent
有限模的 Fitting 理想是由引理 \ref{more-algebra-lemma-fitting-ideal}
中的构造所确定的理想。

\begin{lemma}
\label{more-algebra-lemma-ideals-generated-by-minors}
设 $R$ 为环，$A$ 为系数在 $R$ 中的 $n \times m$ 矩阵。
令 $I_r(A)$ 为由 $A$ 的 $r \times r$ 子式生成的理想，并约定
$I_0(A) = R$；若 $r > \min(n, m)$，则约定 $I_r(A) = 0$。则
\begin{enumerate}
\item $I_0(A) \supset I_1(A) \supset I_2(A) \supset \ldots$；
\item 若 $B$ 是 $(n + n') \times m$ 矩阵，而 $A$ 是 $B$ 的前
$n$ 行，则 $I_{r + n'}(B) \subset I_r(A)$；
\item 若 $C$ 是 $n \times n$ 矩阵，则 $I_r(CA) \subset I_r(A)$。
\item 若 $A$ 是分块矩阵
$$
\left(
\begin{matrix}
A_1 & 0 \\
0 & A_2 
\end{matrix}
\right)
$$
，则 $I_r(A) = \sum_{r_1 + r_2 = r} I_{r_1}(A_1) I_{r_2}(A_2)$。
\item 此处待补充。
\end{enumerate}
\end{lemma}

\begin{proof}
略。（提示：行列式可沿一列或一行展开计算。）
\end{proof}

\begin{lemma}
\label{more-algebra-lemma-fitting-ideal}
设 $R$ 为环，$M$ 为有限 $R$-模。选取 $M$ 的一个表示
$$
\bigoplus\nolimits_{j \in J} R \longrightarrow R^{\oplus n}
\longrightarrow M \longrightarrow 0.
$$
令 $A = (a_{ij})_{i = 1, \ldots, n, j \in J}$ 为映射
$\bigoplus_{j \in J} R \to R^{\oplus n}$ 的矩阵。由 $A$ 的
$(n - k) \times (n - k)$ 子式生成的理想
$\text{Fit}_k(M)$ 不依赖于表示的选择。
\end{lemma}

\begin{proof}
设 $K \subset R^{\oplus n}$ 为满射 $R^{\oplus n} \to M$ 的核。
选取 $z_1, \ldots, z_{n - k} \in K$，并写
$z_j = (z_{1j}, \ldots, z_{nj})$。理想 $\text{Fit}_k(M)$
的另一种描述是：它由按这种方式得到的所有矩阵 $(z_{ij})$ 的
$(n - k) \times (n - k)$ 子式生成。

\medskip\noindent
假设把该满射换成核为 $K'$ 的满射
$R^{\oplus n + n'} \to M$，其中在 $R^{\oplus n + n'}$ 的前
$n$ 个标准基向量上使用原映射，而在最后 $n'$ 个基向量上取 $0$。
则相应的理想相同。事实上，若
$z_1, \ldots, z_{n - k} \in K$ 如上，令
$z'_j = (z_{1j}, \ldots, z_{nj}, 0, \ldots, 0) \in K'$，
其中 $j = 1, \ldots, n - k$，并令
$z'_{n + j'} = (0, \ldots, 0, 1, 0, \ldots, 0) \in K'$。
于是矩阵 $(z_{ij})$ 的 $(n - k) \times (n - k)$ 子式所生成的理想，
与 $(z'_{ij})$ 的 $(n + n' - k) \times (n + n' - k)$
子式所生成的理想相同。这给出一个包含关系。
反过来，给定 $K'$ 中的
$z'_1, \ldots, z'_{n + n' - k}$，可把它们投影到
$R^{\oplus n}$，得到 $K$ 中的
$z_1, \ldots, z_{n + n' - k}$。由引理
\ref{more-algebra-lemma-ideals-generated-by-minors}，$(z'_{ij})$ 的
$(n + n' - k) \times (n + n' - k)$ 子式所生成的理想，
包含于 $(z_{ij})$ 的 $(n - k) \times (n - k)$ 子式所生成的理想。
这给出另一个包含关系。

\medskip\noindent
设 $R^{\oplus m} \to M$ 是另一个核为 $L$ 的满射。
由 Schanuel 引理（Algebra, Lemma \ref{algebra-lemma-Schanuel}）
及上一段的结论，可假设 $m = n$，并且存在同构
$R^{\oplus n} \to R^{\oplus m}$，它与到 $M$ 的满射交换。
令 $C = (c_{li})$ 为该映射的（可逆）矩阵
（由于 $n = m$，这是方阵）。于是给定如上的
$z'_1, \ldots, z'_{n - k} \in L$，可选取
$z_1, \ldots, z_{n - k} \in K$，使得
$z_1' = Cz_1, \ldots, z'_{n - k} = Cz_{n - k}$。
由引理 \ref{more-algebra-lemma-ideals-generated-by-minors} 得到一个包含关系，
再由对称性得到另一个。
\end{proof}

\begin{definition}
\label{more-algebra-definition-fitting-ideal}
设 $R$ 为环，$M$ 为有限 $R$-模，$k \geq 0$。
$M$ 的{\it 第 $k$ 个 Fitting 理想}是引理
\ref{more-algebra-lemma-fitting-ideal} 中构造的理想 $\text{Fit}_k(M)$。
令 $\text{Fit}_{-1}(M) = 0$。
\end{definition}

\noindent
由于 Fitting 理想是一个大矩阵的子式理想
（其编号顺序与引理 \ref{more-algebra-lemma-ideals-generated-by-minors}
中的顺序相反），可见对某个 $t \gg 0$ 有
$$
0 = \text{Fit}_{-1}(M) \subset \text{Fit}_0(M) \subset \text{Fit}_1(M)
\subset \ldots \subset \text{Fit}_t(M) = R
$$
。下面给出 Fitting 理想的一些基本性质。

\begin{lemma}
\label{more-algebra-lemma-fitting-ideal-basics}
设 $R$ 为环，$M$ 为有限 $R$-模。
\begin{enumerate}
\item 若 $M$ 可由 $n$ 个元素生成，则 $\text{Fit}_n(M) = R$。
\item 给定另一个有限 $R$-模 $M'$，有
$\text{Fit}_0(M \oplus M') = \text{Fit}_0(M)\text{Fit}_0(M')$，
且更一般地，
$$
\text{Fit}_l(M \oplus M') =
\sum\nolimits_{k + k' = l} \text{Fit}_k(M)\text{Fit}_{k'}(M')
$$
\item 若 $R \to R'$ 为环映射，则 $\text{Fit}_k(M \otimes_R R')$
是 $R'$ 中由 $\text{Fit}_k(M)$ 的像生成的理想。
\item 若 $M$ 有限表示，则 $\text{Fit}_k(M)$ 是有限生成理想。
\item 若 $M \to M'$ 是满射，则
$\text{Fit}_k(M) \subset \text{Fit}_k(M')$。
\item 有 $\text{Fit}_0(M) \subset \text{Ann}_R(M)$。
\item 有 $V(\text{Fit}_0(M)) = \text{Supp}(M)$。
\item 若 $I$ 是 $R$ 的理想，则 $\text{Fit}_0(R/I) = I$。
\item 此处待补充。
\end{enumerate}
\end{lemma}

\begin{proof}
第 (1) 项由引理 \ref{more-algebra-lemma-ideals-generated-by-minors} 中
$I_0(A) = R$ 这一事实推出。

\medskip\noindent
第 (2) 项由引理 \ref{more-algebra-lemma-ideals-generated-by-minors} 中的相应断言推出。

\medskip\noindent
第 (3) 项由 $\otimes_R R'$ 右正合这一事实推出，因此 $M$ 的一个表示
作基变换后是 $M \otimes_R R'$ 的一个表示。

\medskip\noindent
证明 (4)。设
$R^{\oplus m} \xrightarrow{A} R^{\oplus n} \to M \to 0$
是一个表示。则 $\text{Fit}_k(M)$ 是由矩阵 $A$ 的
$(n - k) \times (n - k)$ 子式生成的理想。

\medskip\noindent
第 (5) 项由定义立即得到。

\medskip\noindent
证明 (6)。选取由矩阵 $A$ 给出的 $M$ 的一个表示，如引理
\ref{more-algebra-lemma-fitting-ideal}。设 $J' \subset J$ 是基数为 $n$ 的子集。
只需证明
$f = \det(a_{ij})_{i = 1, \ldots, n, j \in J'}$
零化 $M$。这是清楚的，因为
$$
R^{\oplus n} \xrightarrow{A' = (a_{ij})_{i = 1, \ldots, n, j \in J'}}
R^{\oplus n} \to M \to 0
$$
的余核被 $f$ 零化：存在矩阵 $B$，满足
$A' B = f1_{n \times n}$。

\medskip\noindent
证明 (7)。选取由矩阵 $A$ 给出的 $M$ 的一个表示，如引理
\ref{more-algebra-lemma-fitting-ideal}。由 Nakayama 引理
（Algebra, Lemma \ref{algebra-lemma-NAK}），有
$$
M_\mathfrak p \not = 0
\Leftrightarrow
M \otimes_R \kappa(\mathfrak p) \not = 0
\Leftrightarrow
\text{rank}(\text{像 }A\text{ 于 }\kappa(\mathfrak p)) < n
$$
。显然，$\text{Fit}_0(M)$ 恰好切出具有该性质的素理想集合。
\end{proof}

\begin{example}
\label{more-algebra-example-fitting-free}
设 $R$ 为环。有限自由模 $M = R^{\oplus n}$ 的 Fitting 理想满足：
当 $k < n$ 时 $\text{Fit}_k(M) = 0$，当 $k \geq n$ 时
$\text{Fit}_k(M) = R$。
\end{example}






\begin{example}
\label{more-algebra-example-laurent}
设 $R$ 为环，$I, J\subset R$ 为理想。则
$\text{Fit}_0(R/I) = I$、
$\text{Fit}_0(R/I \oplus R/J) = IJ$、
$\text{Fit}_1(R/I\oplus R/J) = I+J$。
\end{example}

\begin{lemma}
\label{more-algebra-lemma-fitting-ideal-generate-locally}
设 $R$ 为环，$M$ 为有限 $R$-模，$k \geq 0$。
设 $\mathfrak p \subset R$ 为素理想。下列条件等价：
\begin{enumerate}
\item $\text{Fit}_k(M) \not \subset \mathfrak p$；
\item $\dim_{\kappa(\mathfrak p)} M \otimes_R \kappa(\mathfrak p) \leq k$；
\item $M_\mathfrak p$ 作为 $R_\mathfrak p$-模可由 $k$ 个元素生成；
\item 对某个 $f \in R$、$f \not \in \mathfrak p$，
$M_f$ 作为 $R_f$-模可由 $k$ 个元素生成。
\end{enumerate}
\end{lemma}

\begin{proof}
由 Nakayama 引理（Algebra, Lemma \ref{algebra-lemma-NAK}），
若 $M \otimes_R \kappa(\mathfrak p)$ 可由 $k$ 个元素生成，
则对某个 $f \in R$、$f \not \in \mathfrak p$，
$M_f$ 作为 $R_f$-模可由 $k$ 个元素生成。因此 (2)、(3) 与 (4) 等价。
利用引理 \ref{more-algebra-lemma-fitting-ideal-basics} 第 (3) 项，
问题归结为 $R$ 是域且 $\mathfrak p = (0)$ 的情形。
此时结论由例 \ref{more-algebra-example-fitting-free} 推出。
\end{proof}

\begin{lemma}
\label{more-algebra-lemma-fitting-ideal-finite-locally-free}
设 $R$ 为环，$M$ 为有限 $R$-模，$r \geq 0$。下列条件等价：
\begin{enumerate}
\item $M$ 是秩为 $r$ 的有限局部自由模
（Algebra, Definition \ref{algebra-definition-locally-free}）；
\item $\text{Fit}_{r - 1}(M) = 0$ 且 $\text{Fit}_r(M) = R$；
\item 当 $k < r$ 时 $\text{Fit}_k(M) = 0$，而当 $k \geq r$ 时
$\text{Fit}_k(M) = R$。
\end{enumerate}
\end{lemma}

\begin{proof}
由于 Fitting 理想组成递增理想序列，(2) 与 (3) 显然等价。
由于 $\text{Fit}_k(M)$ 的形成与基变换可交换
（引理 \ref{more-algebra-lemma-fitting-ideal-basics}），例
\ref{more-algebra-example-fitting-free} 及粘合结果
（Algebra, Section \ref{algebra-section-more-glueing}）
表明 (1) 蕴含 (2)。反过来，假设 (2)。由引理
\ref{more-algebra-lemma-fitting-ideal-generate-locally}，可假设 $M$ 由 $r$
个元素生成。因此有一个表示
$\bigoplus_{j \in J} R \to R^{\oplus r} \to M \to 0$。
现在，假设 $\text{Fit}_{r - 1}(M) = 0$ 蕴含映射
$\bigoplus_{j \in J} R \to R^{\oplus r}$ 的矩阵的所有元素均为零。
因此 $M$ 是自由模。
\end{proof}

\begin{lemma}
\label{more-algebra-lemma-principal-fitting-ideal}
设 $R$ 为局部环，$M$ 为有限 $R$-模，$k \geq 0$。
假设对某个 $f \in R$ 有 $\text{Fit}_k(M) = (f)$。
令 $M'$ 为 $M$ 关于 $\{x \in M \mid fx = 0\}$ 的商。
则 $M'$ 可由 $k$ 个元素生成。
\end{lemma}

\begin{proof}
选取对应于满射 $R^{\oplus n} \to M$ 的生成元
$x_1, \ldots, x_n \in M$。由于 $R$ 是局部环，若元素集合
$E \subset (f)$ 生成 $(f)$，则某个 $e \in E$ 生成 $(f)$，见
Algebra, Lemma \ref{algebra-lemma-NAK}。因此可在
$R^{\oplus n} \to M$ 的核中选取 $z_1, \ldots, z_{n - k}$，
使得 $n \times (n - k)$ 矩阵 $A = (z_{ij})$ 的某个
$(n - k) \times (n - k)$ 子式生成 $(f)$。重新编号 $x_i$ 后，
可假设第一个子式
$\det(z_{ij})_{1 \leq i, j \leq n - k}$ 生成 $(f)$，即对某个
单位元 $u \in R$ 有
$\det(z_{ij})_{1 \leq i, j \leq n - k} = uf$。
其他每个子式都是 $f$ 的倍数。由 Algebra, Lemma
\ref{algebra-lemma-matrix-right-inverse}，存在一个
$n - k \times n - k$ 矩阵 $B$，使得
$$
AB = f
\left(
\begin{matrix}
u 1_{n - k \times n - k} \\
C
\end{matrix}
\right)
$$
，其中 $C$ 是系数在 $R$ 中的某个矩阵。这表明对每个
$i \leq n - k$，元素 $y_i = ux_i + \sum_j c_{ji}x_j$ 被 $f$ 零化。
由于 $M/\sum Ry_i$ 由 $x_{n - k + 1}, \ldots, x_n$ 的像生成，
结论成立。
\end{proof}

\begin{lemma}
\label{more-algebra-lemma-fitting-ideals-and-pd1}
设 $R$ 为环，$M$ 为有限表示 $R$-模，$k \geq 0$。
假设对某个非零因子 $f \in R$ 有
$\text{Fit}_k(M) = (f)$，并且 $\text{Fit}_{k - 1}(M) = 0$。则
\begin{enumerate}
\item $M$ 的投射维数 $\leq 1$；
\item $M' = \Ker(f : M \to M)$ 是 $M$ 的 $f$-幂挠子模；
\item $M'$ 的投射维数 $\leq 1$；
\item $M/M'$ 是秩为 $k$ 的有限局部自由模；
\item $M \cong M/M' \oplus M'$。
\end{enumerate}
\end{lemma}

\begin{proof}
选取由系数在 $R$ 中的某个矩阵 $A$ 给出的表示
$$
R^{\oplus m} \xrightarrow{A} R^{\oplus n} \to M \to 0
$$
。

\medskip\noindent
先在 $R$ 为局部环时证明该引理。如陈述中一样，令
$M' = \{x \in M \mid fx = 0\}$。由引理
\ref{more-algebra-lemma-principal-fitting-ideal}，可选取
$x_1, \ldots, x_k \in M$ 生成 $M/M'$。于是
$x_1, \ldots, x_k$ 生成 $M_f = (M/M')_f$。因此，若 $M$ 中存在关系
$\sum a_ix_i = 0$，则 $a_1, \ldots, a_k$ 在 $R_f$ 中映到零；
否则 $\text{Fit}_{k - 1}(M) R_f = \text{Fit}_{k - 1}(M_f)$
将非零。由于 $f$ 是非零因子，得到
$a_1 = \ldots = a_k = 0$。因此
$M \cong R^{\oplus k} \oplus M'$。对上述表示作基变换后，
可假设 $R^{\oplus n}$ 的前 $n - k$ 个基向量映入 $M$ 的直和分支
$M'$，而 $R^{\oplus n}$ 的后 $k$ 个基向量映到 $M$ 的直和分支 $R^{\oplus k}$
的基元素。如此，矩阵 $A$ 的后 $k$ 行为零。于是用 $M'$ 替换 $M$、
用 $0$ 替换 $k$、用 $n - k$ 替换 $n$，并用删除后 $k$ 行所得的
子矩阵替换 $A$，便归结为下一段讨论的情形。

\medskip\noindent
假设 $R$ 为局部环、$k = 0$，且 $M$ 被 $f$ 零化。
此时 $M$ 的第 $0$ 个 Fitting 理想是 $(f)$，并由大小为
$n \times m$ 的矩阵 $A$ 的 $n \times n$ 子式生成
（这特别蕴含 $m \geq n$）。由于 $R$ 是局部环，$A$ 的某个
$n \times n$ 子式等于 $uf$，其中 $u \in R$ 为单位元。
重新编号后，可假设这是第一个子式。此外，已知 $A$ 的所有其他
$n \times n$ 子式都被 $f$ 整除。以分块形式写
$A = (A_1 A_2)$，其中 $A_1$ 是 $n \times n$ 矩阵，
$A_2$ 是 $n \times (m - n)$ 矩阵。把
Algebra, Lemma \ref{algebra-lemma-matrix-right-inverse}
应用于 $A$ 的转置（！），得到一个 $n \times n$ 矩阵 $B$，使得
$$
BA = B(A_1 A_2) = f
\left(
\begin{matrix}
u 1_{n \times n} & C
\end{matrix}
\right)
$$
，其中 $C$ 是系数在 $R$ 中的某个 $n \times (m - n)$ 矩阵。
先得到 $BA_1 = fu 1_{n \times n}$。于是
$$
BA_2 = fC = u^{-1}fuC = u^{-1}BA_1C
$$
。由于 $B$ 的行列式是非零因子，得到
$A_2 = u^{-1}A_1C$。因此 $A$ 的像等于 $A_1$ 的像；
由于 $A_1$ 的行列式是非零因子，该像同构于 $R^{\oplus n}$。
故 $M$ 的投射维数 $\leq 1$。

\medskip\noindent
回到一般环 $R$ 的情形。由局部情形可知 $M/M'$ 是秩为 $k$
的有限局部自由模，见
Algebra, Lemma \ref{algebra-lemma-finite-projective}。
因此扩张 $0 \to M' \to M \to M/M' \to 0$ 分裂。
从而 $M'$ 是有限表示模。选取短正合序列
$0 \to K \to R^{\oplus a} \to M' \to 0$。则 $K$ 是有限 $R$-模，
见 Algebra, Lemma \ref{algebra-lemma-extension}。
由局部情形，对所有素理想都有
$K_\mathfrak p \cong R_\mathfrak p^{\oplus a}$。因此再次由
Algebra, Lemma \ref{algebra-lemma-finite-projective}，
$K$ 是秩为 $a$ 的有限局部自由模。由此 $M'$ 的投射维数
$\leq 1$，引理得证。
\end{proof}










\section{提升}
\label{more-algebra-section-lifting}

\noindent
本节汇集若干关于如下类型的提升陈述的引理：若 $A$ 是环、
$I \subset A$ 是理想，而 $\overline{\xi}$ 是 $A/I$ 上的某种结构，
则可把 $\overline{\xi}$ 提升为 $A$ 上或 $A$ 的某个 \'etale
扩张上的同类结构 $\xi$。对于下列结构，我们已经证明了一些结果：
\begin{enumerate}
\item 幂等元，见
Algebra, Lemmas \ref{algebra-lemma-lift-idempotents} and
\ref{algebra-lemma-lift-idempotents-noncommutative}；
\item 投射模，见
Algebra, Lemmas \ref{algebra-lemma-lift-projective-module} and
\ref{algebra-lemma-lift-finite-projective-module}；
\item 有限稳定自由模，见引理 \ref{more-algebra-lemma-lift-stably-free}；
\item 基元素，见
Algebra, Lemmas \ref{algebra-lemma-local-artinian-basis-when-flat} and
\ref{algebra-lemma-lift-basis}；
\item 环映射，即证明某些代数形式光滑，见
Algebra, Lemma \ref{algebra-lemma-polynomial-ring-formally-smooth},
Proposition \ref{algebra-proposition-smooth-formally-smooth}, and
Lemma \ref{algebra-lemma-smooth-strong-lift}；
\item syntomic 环映射，见
Algebra, Lemma \ref{algebra-lemma-lift-syntomic}；
\item 光滑环映射，见
Algebra, Lemma \ref{algebra-lemma-lift-smooth}；
\item \'etale 环映射，见
Algebra, Lemma \ref{algebra-lemma-lift-etale}；
\item 多项式的分解，见
Algebra, Lemma \ref{algebra-lemma-factor-mod-lift-etale}；
\item Algebra, Section \ref{algebra-section-henselian}
讨论 Hensel 局部环。
\end{enumerate}
感兴趣的读者可在 Smoothing Ring Maps,
Section \ref{smoothing-section-presentations} 中找到更多同类结果，
特别是 Smoothing Ring Maps, Proposition
\ref{smoothing-proposition-lift-smooth}。

\medskip\noindent
设 $A$ 为环，$I \subset A$ 为理想，$\overline{\xi}$
为 $A/I$ 上的某种结构。在下面的引理中，我们寻找诱导同构
$A/I \to A'/IA'$ 的 \'etale 环映射 $A \to A'$，以及 $A'$
上的对象 $\xi'$，使其提升 $\overline{\xi}$。一个一般性说明是：
给定 \'etale 环映射 $A \to A' \to A''$，若
$A/I \cong A'/IA'$ 且 $A'/IA' \cong A''/IA''$，则复合
$A \to A''$ 仍为 \'etale
（Algebra, Lemma \ref{algebra-lemma-etale}），并且仍满足
$A/I \cong A''/IA''$。下面将经常使用这一点而不再说明。
先给出一个我们所寻求类型的平凡例子。

\begin{lemma}
\label{more-algebra-lemma-lift-invertible-element}
设 $A$ 为环，$I \subset A$ 为理想，
$\overline{u} \in A/I$ 为可逆元素。存在一个诱导同构
$A/I \to A'/IA'$ 的 \'etale 环映射 $A \to A'$，以及提升
$\overline{u}$ 的可逆元素 $u' \in A'$。
\end{lemma}

\begin{proof}
任取 $\overline{u}$ 的提升 $f \in A$，并令 $A' = A_f$，
令 $u$ 为 $f$ 在 $A'$ 中的像。
\end{proof}

\begin{lemma}
\label{more-algebra-lemma-lift-idempotent}
设 $A$ 为环，$I \subset A$ 为理想，
$\overline{e} \in A/I$ 为幂等元。存在一个诱导同构
$A/I \to A'/IA'$ 的 \'etale 环映射 $A \to A'$，以及提升
$\overline{e}$ 的幂等元 $e' \in A'$。
\end{lemma}

\begin{proof}
任取 $\overline{e}$ 的提升 $x \in A$。令
$$
A' = A[t]/(t^2 - t)\left[\frac{1}{t - 1 + x}\right].
$$
环映射 $A \to A'$ 是 \'etale 的，因为
$(2t - 1)\text{d}t = 0$，且
$(2t - 1)(2t - 1) = 1$ 可逆。又有
$A'/IA' = A/I[t]/(t^2 - t)[\frac{1}{t - 1 + \overline{e}}] \cong A/I$；
最后一个映射把 $t$ 送到 $\overline{e}$，这是可行的，因为
$\overline{e}$ 是 $t^2 - t$ 的根。这也表明，令 $e'$ 为 $t$
在 $A'$ 中的类即可。
\end{proof}

\begin{lemma}
\label{more-algebra-lemma-lift-open-covering}
设 $A$ 为环，$I \subset A$ 为理想。设
$\Spec(A/I) = \coprod_{j \in J} \overline{U}_j$
为有限不交开覆盖。则存在一个诱导同构
$A/I \to A'/IA'$ 的 \'etale 环映射 $A \to A'$，以及提升给定覆盖的
有限不交开覆盖 $\Spec(A') = \coprod_{j \in J} U'_j$。
\end{lemma}

\begin{proof}
这由引理 \ref{more-algebra-lemma-lift-idempotent} 以及谱的开闭子集与幂等元对应这一事实
推出，见 Algebra, Lemma
\ref{algebra-lemma-disjoint-decomposition}。
\end{proof}

\begin{lemma}
\label{more-algebra-lemma-localize-upstairs}
设 $A \to B$ 为环映射，$J \subset B$ 为理想。
若 $A \to B$ 在 $V(J)$ 的每个素理想处都是 \'etale 的，
则存在 $g \in B$，它在 $B/J$ 中映到可逆元素，并且
$A' = B_g$ 在 $A$ 上为 \'etale。
\end{lemma}

\begin{proof}
$\Spec(B)$ 中 $A \to B$ 不是 \'etale 的点所成集合是
$\Spec(B)$ 的闭子集，见
Algebra, Definition \ref{algebra-definition-etale}。
把它写成 $V(J')$，其中 $J' \subset B$ 为某个理想。
由假设 $V(J') \cap V(J) = \emptyset$，故由
Algebra, Lemma \ref{algebra-lemma-Zariski-topology} 有
$J + J' = B$。写 $1 = f + g$，其中 $f \in J$ 且 $g \in J'$。
则 $g$ 满足要求。
\end{proof}

\noindent
下面三个引理说明可以提升多项式的分解。

\begin{lemma}
\label{more-algebra-lemma-lift-factorization-monic}
设 $A$ 为环，$I \subset A$ 为理想，$f \in A[x]$ 为首一多项式。
设 $\overline{f} = \overline{g} \overline{h}$ 是 $f$ 在
$A/I[x]$ 中的一个分解，其中 $\overline{g}$ 和 $\overline{h}$
首一且在 $A/I[x]$ 中生成单位理想。则存在一个诱导同构
$A/I \to A'/IA'$ 的 \'etale 环映射 $A \to A'$，以及 $A'[x]$
中的分解 $f = g' h'$，其中 $g'$、$h'$ 首一，并提升 $A/I$
上的给定分解。
\end{lemma}

\begin{proof}
我们将从先前证明的关于泛分解的结果推出这一点；不过，仍鼓励读者寻找一个
不使用该技巧的证明。设
$\deg(\overline{g}) = n$ 且 $\deg(\overline{h}) = m$，
从而 $\deg(f) = n + m$。对某些
$\alpha_1, \ldots, \alpha_{n + m} \in A$，写
$f = x^{n + m} + \sum \alpha_i x^{n + m - i}$。考虑
Algebra, Example \ref{algebra-example-factor-polynomials-etale} 中的环映射
$$
R = \mathbf{Z}[a_1, \ldots, a_{n + m}]
\longrightarrow
S = \mathbf{Z}[b_1, \ldots, b_n, c_1, \ldots, c_m]
$$
。令 $R \to A$ 为把 $a_i$ 送到 $\alpha_i$ 的环映射，并令
$$
B = A \otimes_R S
$$
。由构造，$f$ 在 $B[x]$ 中的像 $f_B$ 有一个分解，记为
$f_B = g_B h_B$，其中
$g_B = x^n + \sum (1 \otimes b_i) x^{n - i}$，而 $h_B$ 类似。
写
$\overline{g} = x^n + \sum \overline{\beta}_i x^{n - i}$ 以及
$\overline{h} = x^m + \sum \overline{\gamma}_i x^{m - i}$。
$A$-代数映射
$$
B \longrightarrow A/I, \quad
1 \otimes b_i \mapsto \overline{\beta}_i, \quad
1 \otimes c_i \mapsto \overline{\gamma}_i
$$
把 $g_B$ 和 $h_B$ 分别映到 $A/I[x]$ 中的
$\overline{g}$ 和 $\overline{h}$。所展示的映射是满射；
以 $J \subset B$ 表示其核。由
Algebra, Example \ref{algebra-example-factor-polynomials-etale}
中的讨论，显然 $A \to B$ 在 $V(J) \subset \Spec(B)$
的所有点处都是 etale 的。按引理 \ref{more-algebra-lemma-localize-upstairs}
选取 $g \in B$，并考虑 $A$-代数 $B_g$。由于 $g$ 在
$B/J = A/I$ 中映到单位元，还得到一个 $A/I$-代数映射
$B_g/I B_g \to A/I$。由于 $A/I \to B_g/I B_g$ 是 \'etale 的，
$B_g/IB_g \to A/I$ 也是 \'etale 的
（Algebra, Lemma \ref{algebra-lemma-map-between-etale}）。
因此存在幂等元 $e \in B_g/I B_g$，使得
$A/I = (B_g/I B_g)_e$
（Algebra, Lemma
\ref{algebra-lemma-surjective-flat-finitely-presented}）。
选取 $e$ 的提升 $h \in B_g$。于是 $A \to A' = (B_g)_h$，
连同 $f_B = g_B h_B$ 的分解在 $A'$ 中的像，给出了该引理所提问题的解。
\end{proof}

\begin{lemma}
\label{more-algebra-lemma-lift-factorization-easy}
设 $A$ 为环，$I \subset A$ 为理想，$f \in A[x]$ 为首一多项式。
设 $\overline{f} = \overline{g} \overline{h}$ 是 $f$ 在
$A/I[x]$ 中的一个分解，并假设
\begin{enumerate}
\item $\overline{g}$ 的首项系数是 $A/I$ 中的可逆元素；
\item $\overline{g}$、$\overline{h}$ 在 $A/I[x]$ 中生成单位理想。
\end{enumerate}
则存在一个诱导同构 $A/I \to A'/IA'$ 的 \'etale 环映射
$A \to A'$，以及 $A'[x]$ 中的分解 $f = g' h'$，
它提升 $A/I$ 上的给定分解。
\end{lemma}

\begin{proof}
应用引理 \ref{more-algebra-lemma-lift-invertible-element}，可假设
$\overline{g}$ 的首项系数是某个可逆元素 $u \in A$ 的约化。
于是可以用 $\overline{u}^{-1}\overline{g}$ 替换
$\overline{g}$，并用 $\overline{u}\overline{h}$ 替换
$\overline{h}$。故可假设 $\overline{g}$ 首一。由于 $f$ 首一，
得到 $\overline{h}$ 也首一。此时结论由引理
\ref{more-algebra-lemma-lift-factorization-monic} 推出。
\end{proof}

\begin{example}
\label{more-algebra-example-cannot-lift-factorization}
本例说明，引理 \ref{more-algebra-lemma-lift-factorization-easy} 中关于首项系数的
条件不可去掉。确切地说，令 $A = \mathbf{Z}$、
$I = 4\mathbf{Z}$、$f = 1$，并在 $A/I[x]$ 中令
$\bar{g} = \bar{h} = 2x^2 + 2x + 1$。若 $A \to A'$ 是 \'etale 的，
且 $A'/I A' = A/I$，则 $A'$ 的 $2$-进完备化同构于
$\mathbf{Z}_2$。然而，$\mathbf{Z}_2[x]$ 中任意模 $4$ 同余于
$2x^2 + 2x + 1$ 的多项式 $g'$ 都没有 $\mathbf{Z}_2[x]$ 中的
多项式乘法逆元，故该引理的结论不成立。
\end{example}

\begin{lemma}
\label{more-algebra-lemma-separate-image-closed-from-closed}
设 $R \to S$ 为环映射，$I \subset R$ 为 $R$ 的理想，
$J \subset S$ 为 $S$ 的理想。若 $V(J)$ 在 $\Spec(R)$ 中的像的闭包
与 $V(I)$ 不交，则存在元素 $f \in R$，它在 $R/I$ 中映到 $1$，
而在 $S$ 中映到 $J$ 的一个元素。
\end{lemma}

\begin{proof}
设 $I' \subset R$ 为理想，使 $V(I')$ 是 $V(J)$ 的像的闭包。
由假设 $V(I) \cap V(I') = \emptyset$，故由
Algebra, Lemma \ref{algebra-lemma-Zariski-topology} 有
$I + I' = R$。写 $1 = g + f$，其中 $g \in I$ 且 $f \in I'$。
设 $f'$ 是 $f$ 在 $S$ 中的像，则
$V(f') \supset V(J)$。因此对某个 $n$ 有 $(f')^n \in J$，见
Algebra, Lemma \ref{algebra-lemma-Zariski-topology}。
用 $f^n$ 替换 $f$ 即得结论。
\end{proof}

\begin{lemma}
\label{more-algebra-lemma-helper-integral}
设 $I$ 为环 $A$ 的理想，$A \to B$ 为整环映射。
设 $b \in B$ 在 $B/IB$ 中映到幂等元。则存在首一
$f \in A[x]$，使得 $f(b) = 0$，且对某个 $d \geq 1$ 有
$f \bmod I = x^d(x - 1)^d$。
\end{lemma}

\begin{proof}
注意 $z = b^2 - b$ 是 $IB$ 的元素。由 Algebra, Lemma
\ref{algebra-lemma-integral-integral-over-ideal}，存在次数为 $d$
的首一多项式 $g(x) = x^d + \sum a_j x^j$，其中 $a_j \in I$，
使得在 $B$ 中 $g(z) = 0$。因此
$f(x) = g(x^2 - x) \in A[x]$ 是首一多项式，满足
$f(x) \equiv x^d(x - 1)^d \bmod I$，并且在 $B$ 中
$f(b) = 0$。
\end{proof}

\begin{lemma}
\label{more-algebra-lemma-lift-idempotent-upstairs}
设 $A$ 为环，$I \subset A$ 为理想，$A \to B$ 为整环映射，
$\overline{e} \in B/IB$ 为幂等元。则存在一个诱导同构
$A/I \to A'/IA'$ 的 \'etale 环映射 $A \to A'$，以及提升
$\overline{e}$ 的幂等元 $e' \in B \otimes_A A'$。
\end{lemma}

\begin{proof}
选取 $\overline{e}$ 的提升 $y \in B$，并按引理
\ref{more-algebra-lemma-helper-integral} 为 $y$ 选取 $f \in A[x]$。
由引理 \ref{more-algebra-lemma-lift-factorization-easy}，可找到一个诱导同构
$A/I \to A'/IA'$ 的 \'etale 环映射 $A \to A'$，并使
$f = gh$ 在 $A[x]$ 中成立，其中
$g(x) = x^d \bmod IA'$ 且
$h(x) = (x - 1)^d \bmod IA'$。用 $A'$ 替换 $A$ 后，
可假设该分解定义在 $A$ 上。此时
$b_1 = g(y) \in B$ 是 $\overline{e}^d = \overline{e}$ 的提升，而
$b_2 = h(y) \in B$ 是
$(\overline{e} - 1)^d = (-1)^d (1 - \overline{e})^d = (-1)^d(1 - \overline{e})$
的提升，并且 $b_1b_2 = 0$。于是
$(b_1, b_2)B/IB = B/IB$，而
$V(b_1, b_2) \subset \Spec(B)$ 与 $V(IB)$ 不交。
由于 $\Spec(B) \to \Spec(A)$ 是闭映射（见
Algebra, Lemmas \ref{algebra-lemma-integral-going-up} and
\ref{algebra-lemma-going-up-closed}），由引理
\ref{more-algebra-lemma-separate-image-closed-from-closed} 可选取 $a \in A$，
使其在 $A/I$ 中映到可逆元素，而其在 $B$ 中的像属于
$(b_1, b_2)$。用局部化 $A_a$ 替换 $A$ 后，有
$(b_1, b_2) = B$。于是
$\Spec(B) = D(b_1) \amalg D(b_2)$；这是不交并，因为
$b_1b_2 = 0$；又覆盖 $\Spec(B)$，因为 $(b_1, b_2) = B$。
令 $e \in B$ 为对应于开闭子集 $D(b_1)$ 的幂等元，见
Algebra, Lemma \ref{algebra-lemma-disjoint-decomposition}。
由于 $b_1$ 是 $\overline{e}$ 的提升，而 $b_2$ 是
$\pm (1 - \overline{e})$ 的提升，由
Algebra, Lemma \ref{algebra-lemma-disjoint-decomposition}
中的唯一性断言，$e$ 是 $\overline{e}$ 的提升。
\end{proof}

\begin{lemma}
\label{more-algebra-lemma-lift-projective-module}
设 $A$ 为环，$I \subset A$ 为理想，
$\overline{P}$ 为有限投射 $A/I$-模。则存在一个诱导同构
$A/I \to A'/IA'$ 的 \'etale 环映射 $A \to A'$，以及提升
$\overline{P}$ 的有限投射 $A'$-模 $P'$。
\end{lemma}

\begin{proof}
可选取整数 $n$ 以及直和分解
$(A/I)^{\oplus n} = \overline{P} \oplus \overline{K}$，
其中 $\overline{K}$ 是某个 $R/I$-模。选取与直和分支
$\overline{P}$ 对应的投影算子 $\overline{p}$ 的一个提升
$\varphi : A^{\oplus n} \to A^{\oplus n}$。
令 $f \in A[x]$ 为 $\varphi$ 的特征多项式，并令
$B = A[x]/(f)$。由 Cayley--Hamilton
（Algebra, Lemma \ref{algebra-lemma-charpoly}），存在映射
$B \to \text{End}_A(A^{\oplus n})$，它把 $x$ 送到 $\varphi$。
对每个 $\mathfrak p \supset I$，$f$ 在 $\kappa(\mathfrak p)$
中的像是 $(x - 1)^rx^{n - r}$，其中 $r$ 是
$\overline{P} \otimes_{A/I} \kappa(\mathfrak p)$ 的维数。
因此对所有 $\mathfrak p \supset I$，
$(x - 1)^nx^n$ 在 $B \otimes_A \kappa(\mathfrak p)$ 中映到零。
故 $x(1 - x)$ 包含于 $B/IB$ 的每个素理想，从而对某个
$N \geq 1$，$x^N(1 - x)^N$ 包含于 $IB$。
由此 $x^N + (1 - x)^N$ 是 $B/IB$ 中的单位元，并且
$$
\overline{e} = \text{image of }\frac{x^N}{x^N + (1 - x)^N}\text{ 于 }B/IB
$$
是幂等元，因为这两个断言在
$\mathbf{Z}[x]/(x^N(x - 1)^N)$ 中都成立。
$\overline{e}$ 在 $\text{End}_{A/I}((A/I)^{\oplus n})$ 中的像是
$$
\frac{\overline{p}^N}{\overline{p}^N + (1 - \overline{p})^N} = \overline{p}
$$
，因为 $\overline{p}$ 是幂等元。按引理中的方式用一个 \'etale
扩张 $A'$ 替换 $A$ 后，可假设存在幂等元 $e \in B$，其在
$B/IB$ 中映到 $\overline{e}$，见引理
\ref{more-algebra-lemma-lift-idempotent-upstairs}。于是 $e$ 在映射
$$
B = A[x]/(f) \longrightarrow \text{End}_A(A^{\oplus n}).
$$
下的像是提升 $\overline{p}$ 的幂等元 $p$。令
$P = \Im(p)$ 即得结论。
\end{proof}

\begin{lemma}
\label{more-algebra-lemma-cotangent-complex-symmetric-algebra}
设 $A$ 为环，并设
$0 \to K \to A^{\oplus m} \to M \to 0$
为 $A$-模序列。考虑 $A$-代数
$C = \text{Sym}^*_A(M)$，以及由满射
$A^{\oplus m} \to M$ 给出的表示
$\alpha : A[y_1, \ldots, y_m] \to C$。则
$$
\NL(\alpha) =
(K \otimes_A C \to \bigoplus\nolimits_{j = 1, \ldots, m} C \text{d}y_j)
$$
（见 Algebra, Section \ref{algebra-section-netherlander}），
特别地 $\Omega_{C/A} = M \otimes_A C$。
\end{lemma}

\begin{proof}
令 $J = \Ker(\alpha)$。该引理断言
$J/J^2 \cong K \otimes_A C$。注意 $\alpha$ 是分次代数同态。
我们将证明在次数 $d$ 有
$(J/J^2)_d = K \otimes_A C_{d - 1}$。注意
$$
J_d = \Ker(\text{Sym}^d_A(A^{\oplus m}) \to \text{Sym}^d_A(M))
= \Im(K \otimes_A \text{Sym}^{d - 1}_A(A^{\oplus m})
\to \text{Sym}^d_A(A^{\oplus m})),
$$
见 Algebra, Lemma \ref{algebra-lemma-presentation-sym-exterior}。
由此 $(J^2)_d = \sum_{a + b = d} J_a \cdot J_b$ 是映射
$$
K \otimes_A K \otimes_A \text{Sym}^{d - 2}_A(A^{\otimes m})
\to \text{Sym}^d_A(A^{\oplus m}).
$$
的像。由上面引用的引理，映射
$K \otimes_A \text{Sym}^{d - 2}_A(A^{\otimes m}) \to
\text{Sym}^{d - 1}_A(A^{\oplus m})$
的余核是 $\text{Sym}^{d - 1}_A(M)$。
因此显然 $(J/J^2)_d = J_d/(J^2)_d$ 等于
\begin{align*}
\Coker(
K \otimes_A K \otimes_A \text{Sym}^{d - 2}_A(A^{\otimes m})
\to K \otimes_A \text{Sym}^{d - 1}_A(A^{\otimes m}))
& = K \otimes_A \text{Sym}^{d - 1}_A(M) \\
& = K \otimes_A C_{d -1}
\end{align*}
，正合所需。
\end{proof}

\begin{lemma}
\label{more-algebra-lemma-symmetric-algebra-smooth}
设 $A$ 为环，$M$ 为 $A$-模。则
$C = \text{Sym}_A^*(M)$ 在 $A$ 上光滑，当且仅当 $M$
是有限投射 $A$-模。
\end{lemma}

\begin{proof}
令 $\sigma : C \to A$ 为到 $C$ 的次数 $0$ 部分的投影。
于是 $J = \Ker(\sigma)$ 是次数 $> 0$ 的部分，并且作为 $A$-模有
$J/J^2 = M$。因此，若 $A \to C$ 光滑，则由
Algebra, Lemma \ref{algebra-lemma-section-smooth}，
$M$ 是有限投射 $A$-模。

\medskip\noindent
反过来，假设 $M$ 有限投射，并选取核为 $K$ 的满射
$A^{\oplus n} \to M$。当然，序列
$0 \to K \to A^{\oplus n} \to M \to 0$ 分裂，因为 $M$ 是投射模。
特别地，$K$ 是有限 $A$-模。因此 $C$ 在 $A$ 上为有限表示，
因为 $C$ 是 $A[x_1, \ldots, x_n]$ 关于由
$K \subset \bigoplus Ax_i$ 生成的理想的商。
引理 \ref{more-algebra-lemma-cotangent-complex-symmetric-algebra} 的计算表明，
$\NL_{C/A}$ 同伦等价于 $(K \to M) \otimes_A C$。
因此 $\NL_{C/A}$ 拟同构于置于次数 $0$ 的
$C \otimes_A M$；由 Algebra, Definition
\ref{algebra-definition-smooth}，这意味着 $C$ 在 $A$ 上光滑。
\end{proof}

\begin{lemma}
\label{more-algebra-lemma-lift-section-smooth-morphism}
设 $A$ 为环，$I \subset A$ 为理想。考虑交换图
$$
\xymatrix{
B \ar[rd] \\
A \ar[u] \ar[r] & A/I
}
$$
，其中 $B$ 是光滑 $A$-代数。则存在一个诱导同构
$A/I \to A'/IA'$ 的 \'etale 环映射 $A \to A'$，以及提升环映射
$B \to A/I$ 的 $A$-代数映射 $B \to A'$。
\end{lemma}

\begin{proof}
令 $J \subset B$ 为 $B \to A/I$ 的核，从而 $B/J = A/I$。
由 Algebra, Lemma \ref{algebra-lemma-application-NL-smooth}，序列
$$
0 \to I/I^2 \to J/J^2 \to \Omega_{B/A} \otimes_B B/J \to 0
$$
分裂正合。因此
$\overline{P} = J/(J^2 + IB) = \Omega_{B/A} \otimes_B B/J$
是有限投射 $A/I$-模。选取整数 $n$ 及直和分解
$A/I^{\oplus n} = \overline{P} \oplus \overline{K}$。
由引理 \ref{more-algebra-lemma-lift-projective-module}，可找到一个诱导同构
$A/I \to A'/IA'$ 的 \'etale 环映射 $A \to A'$，以及提升
$\overline{K}$ 的有限投射 $A$-模 $K$。我们可以且确实用 $A'$
替换 $A$。令
$B' = B \otimes_A \text{Sym}_A^*(K)$。由引理
\ref{more-algebra-lemma-symmetric-algebra-smooth}，
$A \to \text{Sym}_A^*(K)$ 光滑，故 $B \to B'$ 光滑，
进而 $A \to B'$ 光滑（见
Algebra, Lemmas \ref{algebra-lemma-base-change-smooth} and
\ref{algebra-lemma-locally-smooth}）。
此外，截面 $\text{Sym}^*_A(K) \to A$ 确定截面 $B' \to B$，
并令 $B' \to A/I$ 为复合 $B' \to B \to A/I$。
令 $J' \subset B'$ 为 $B' \to A/I$ 的核。有
$JB' \subset J'$ 且 $B \otimes_A K \subset J'$。
这些映射合并成同构
$$
(A/I)^{\oplus n} \cong J/(J^2 + IB) \oplus \overline{K}
\longrightarrow
J'/((J')^2 + IB')
$$
。因此，用 $B'$ 替换 $B$ 后，可假设
$J/(J^2 + IB) = \Omega_{B/A} \otimes_B B/J$
是秩为 $n$ 的自由 $A/I$-模。

\medskip\noindent
在此情形中，选取 $f_1, \ldots, f_n \in J$，使其像构成
$J/(J^2 + IB)$ 的一组基。考虑有限表示 $A$-代数
$C = B/(f_1, \ldots, f_n)$。注意有正合序列
$$
0 \to H_1(L_{C/A}) \to (f_1, \ldots, f_n)/(f_1, \ldots, f_n)^2
\to \Omega_{B/A} \otimes_B C \to \Omega_{C/A} \to 0
$$
，见 Algebra, Lemma \ref{algebra-lemma-exact-sequence-NL}
（注意 $H_1(L_{B/A}) = 0$，而 $\Omega_{B/A}$ 有限投射，
特别地平坦，故 Tor 群消失）。对 $B$ 的任意素理想
$\mathfrak q \supset J$，模 $\Omega_{B/A, \mathfrak q}$ 是秩为
$n$ 的自由模，因为 $\Omega_{B/A}$ 有限投射，且
$\Omega_{B/A} \otimes_B B/J$ 是秩为 $n$ 的自由模
（见 Algebra, Lemma \ref{algebra-lemma-finite-projective}）。
由 $f_1, \ldots, f_n$ 的选择，映射
$$
\left((f_1, \ldots, f_n)/(f_1, \ldots, f_n)^2\right)_{\mathfrak q}
\to
\Omega_{B/A, \mathfrak q}
$$
模 $J$ 后为满射。因此，这个局部环 $C_{\mathfrak q}$ 上的模映射
必为同构（由 Algebra, Lemma \ref{algebra-lemma-NAK}，该映射为满射；
而例如由 Algebra, Lemma \ref{algebra-lemma-fun}，由于
$$
((f_1, \ldots, f_n)/(f_1, \ldots, f_n)^2)_{\mathfrak q}
$$
可由 $n$ 个元素生成，该映射为单射）。故
$H_1(L_{C/A})_{\mathfrak q} = 0$ 且
$\Omega_{C/A, \mathfrak q} = 0$。由 Algebra, Lemma
\ref{algebra-lemma-smooth-at-point}，$A \to C$ 在 $C$ 中对应于
$\mathfrak q$ 的素理想 $\overline{\mathfrak q}$ 处光滑。
由于 $\Omega_{C/A, \mathfrak q} = 0$，它实际上在
$\overline{\mathfrak q}$ 处为 \'etale。因此 $A \to C$
在 $C$ 的所有包含 $JC$ 的素理想处都是 \'etale 的。
由引理 \ref{more-algebra-lemma-localize-upstairs}，可选取 $f \in C$，
它在 $C/JC$ 中映到单位元，并使 $A \to C_f$ 为 \'etale。
由 $f$ 的选择，仍有 $C_f/JC_f = A/I$。映射
$C_f/IC_f \to A/I$ 是满且 \'etale 的，见
Algebra, Lemma \ref{algebra-lemma-map-between-etale}。
因此 $A/I$ 同构于 $C_f/IC_f$ 在某个元素 $g \in C$ 处的局部化，
见 Algebra, Lemma
\ref{algebra-lemma-surjective-flat-finitely-presented}。
令 $A' = C_{fg}$，证明完成。
\end{proof}








\section{Zariski 对}
\label{more-algebra-section-zariski-pairs}

\noindent
在本节及下一节中，一个{\it 对}是指一个对 $(A, I)$，其中 $A$
是环，$I \subset A$ 是理想。{\it 对的态射}
$(A, I) \to (B, J)$ 是满足 $\varphi(I) \subset J$ 的环映射
$\varphi : A \to B$。

\begin{definition}
\label{more-algebra-definition-zariski-pair}
一个{\it Zariski 对}是满足 $I$ 包含于 $A$ 的 Jacobson 根中的对
$(A, I)$。
\end{definition}

\begin{lemma}
\label{more-algebra-lemma-idempotents-determined-modulo-radical}
设 $(A, I)$ 为 Zariski 对。则从 $A$ 的幂等元到 $A/I$ 的幂等元的映射
是单射。
\end{lemma}

\begin{proof}
局部环的幂等元只能是 $0$ 或 $1$。因此，由
Algebra, Lemma \ref{algebra-lemma-characterize-zero-local}，
一个幂等元由它在哪些极大理想处为零所确定。
\end{proof}

\begin{lemma}
\label{more-algebra-lemma-check-isomorphism-zariski}
设 $(A, I)$ 为 Zariski 对。设 $A \to B$ 为平坦、整且有限表示的
环映射，并且 $A/I \to B/IB$ 是同构。则 $A \to B$ 是同构。
\end{lemma}

\begin{proof}
由 Algebra, Lemma
\ref{algebra-lemma-characterize-finite-in-terms-of-integral}，环映射
$A \to B$ 是有限的。因此，由 Algebra, Lemma
\ref{algebra-lemma-finite-finitely-presented-extension}，$B$ 作为
$A$-模是有限表示的。进而，由 Algebra, Lemma
\ref{algebra-lemma-finite-projective}，$B$ 是有限局部自由 $A$-模。
由于模 $B$ 沿 $V(I)$ 的秩为 $1$（见 Algebra, Lemma
\ref{algebra-lemma-finite-projective} 中描述的秩函数），并且
$(A, I)$ 是 Zariski 对，我们推出它处处秩为 $1$。
由此 $A \to B$ 是同构：证明如下事实是一个有趣的练习：若环映射
$R \to S$ 使得 $S$ 是秩为 $1$ 的局部自由 $R$-模，则该映射是同构
（提示：考察局部环）。
\end{proof}

\begin{lemma}
\label{more-algebra-lemma-helper-finite}
设 $(A, I)$ 为 Zariski 对。设 $A \to B$ 为有限环映射。假设
\begin{enumerate}
\item $B/IB = B_1 \times B_2$ 是 $A/I$-代数的乘积；
\item $A/I \to B_1/IB_1$ 是满射；
\item $b \in B$ 在该乘积中的像为 $(1, 0)$。
\end{enumerate}
则存在首一多项式 $f \in A[x]$，使得 $f(b) = 0$，且对某个
$d \geq 1$ 有 $f \bmod I = (x - 1)x^d$。
\end{lemma}

\begin{proof}
由引理 \ref{more-algebra-lemma-lift-idempotent-upstairs}，可找到诱导同构
$A/I \to A'/IA'$ 的 \'etale 环映射 $A \to A'$，使得
$B' = B \otimes_A A'$ 包含一个幂等元 $e'$，它提升 $b$ 在
$B'/IB'$ 中的像。考虑相应的 $A'$-代数分解
$$
B' = B'_1 \times B'_2
$$
它在模 $I$ 约化后与引理中给定的分解相容。
映射 $A' \to B'_1$ 模 $IA'$ 后为满射。由 Nakayama 引理
（Algebra, Lemma \ref{algebra-lemma-NAK}），可找到 $i \in IA'$，
使得以 $A'_{1 + i}$ 替换 $A'$ 后，映射 $A' \to B'_1$ 为满射。
注意，$b$ 在第一分量中的像 $b'_1 \in B'_1$ 满足
$b'_1 - 1 \in IB'_1$。因此可选取映到 $b'_1 - 1$ 的
$a' \in IA'$。另一方面，$b$ 在第二分量中的像
$b'_2 \in B'_2$ 属于 $IB'_2$。由 Algebra, Lemma
\ref{algebra-lemma-integral-integral-over-ideal}，存在一个次数为
$d$ 的首一多项式 $g(x) = x^d + \sum a'_j x^j$，其中
$a'_j \in IA'$，并且 $g(b'_2) = 0$ 在 $B'_2$ 中成立。
因此，$b$ 的像 $b' = (b'_1, b'_2) \in B'$ 是多项式
$(x - 1 - a')g(x)$ 的根。我们得到
$$
(b' - 1)(b')^d \in \sum\nolimits_{j = 0, \ldots, d} IA' \cdot (b')^j
$$
我们断言，这蕴含 $B$ 中有
$$
(b - 1)b^d \in \sum\nolimits_{j = 0, \ldots, d} I \cdot b^j
$$
为此，只需证明环映射 $A \to A'$ 忠实平坦，因为该条件等价于
$(b - 1)b^d$ 的像在
$B/\sum_{j = 0, \ldots, d} Ib^j$ 中为零（使用 Algebra, Lemma
\ref{algebra-lemma-faithfully-flat-universally-injective}）。
映射 $A \to A'$ 是平坦的，因为它是 \'etale 的
（Algebra, Lemma \ref{algebra-lemma-etale}）。另一方面，诱导的谱映射
是开映射（见 Algebra, Proposition \ref{algebra-proposition-fppf-open}，
并使用前面引用的引理），而且其像包含 $V(I)$。由于 $I$ 包含于
$A$ 的 Jacobson 根中，结论随即成立。
\end{proof}

\begin{lemma}
\label{more-algebra-lemma-noetherian-zariski-jacobson-complement}
设 $(A, I)$ 为 Zariski 对，且 $A$ 为 Noether 环。设 $f \in I$。
则 $A_f$ 是 Jacobson 环。
\end{lemma}

\begin{proof}
我们使用 Algebra, Lemma
\ref{algebra-lemma-noetherian-dim-1-Jacobson} 的判别准则。
设 $\mathfrak p \subset A$ 为素理想，使得
$\mathfrak p_f = \mathfrak p A_f$ 是非极大的素理想。我们必须证明
$A_f/\mathfrak p_f = (A/\mathfrak p)_f$ 有无穷多个素理想。
以 $A/\mathfrak p$ 替换 $A$ 后，可以假设 $A$ 是整环，
$\dim A_f > 0$，而我们的目标是证明 $\Spec(A_f)$ 是无限集。
由于 $\dim A_f > 0$，可找到一个不包含 $f$ 的非零素理想
$\mathfrak q \subset A$。选取包含 $\mathfrak q$ 的极大理想
$\mathfrak m \subset A$。由于 $(A, I)$ 是 Zariski 对，有
$I \subset \mathfrak m$。故 $\mathfrak m \not = \mathfrak q$，且
$\dim(A_\mathfrak m) > 1$。于是由 Algebra, Lemma
\ref{algebra-lemma-Noetherian-local-domain-dim-2-infinite-opens}，
$\Spec((A_\mathfrak m)_f) \subset \Spec(A_f)$ 是无限集，证毕。
\end{proof}











\section{Hensel 对}
\label{more-algebra-section-henselian-pairs}

\noindent
第 \ref{more-algebra-section-lifting} 节的一些结果可以看作关于 Hensel 对的结果。
本节中，一个{\it 对}是一个对 $(A, I)$，其中 $A$ 是环且
$I \subset A$ 是理想。{\it 对的态射} $(A, I) \to (B, J)$ 是满足
$\varphi(I) \subset J$ 的环映射 $\varphi : A \to B$。与第
\ref{more-algebra-section-lifting} 节一样，给定 $A$ 上的对象 $\xi$，我们以
$\overline{\xi}$ 表示 $\xi$ 到 $A/I$ 上对象的``基变换''
（如果这有意义）。

\begin{definition}
\label{more-algebra-definition-henselian-pair}
一个{\it Hensel 对}是满足下列条件的对 $(A, I)$：
\begin{enumerate}
\item $I$ 包含于 $A$ 的 Jacobson 根中；
\item 对任意首一多项式 $f \in A[T]$ 以及分解
$\overline{f} = g_0h_0$，其中 $g_0, h_0 \in A/I[T]$ 首一且在
$A/I[T]$ 中生成单位理想，都存在 $A[T]$ 中的分解 $f = gh$，
其中 $g, h$ 首一，并且 $g_0 = \overline{g}$、
$h_0 = \overline{h}$。
\end{enumerate}
\end{definition}

\noindent
注意，若 $A$ 是局部环且 $I = \mathfrak m$ 是其极大理想，则
$(A, I)$ 是 Hensel 对当且仅当 $A$ 是 Hensel 局部环，见
Algebra, Lemma \ref{algebra-lemma-characterize-henselian}。
在引理 \ref{more-algebra-lemma-characterize-henselian-pair} 中，我们给出 Hensel 对的
若干等价刻画（并且以后还会继续补充）。

\begin{lemma}
\label{more-algebra-lemma-locally-nilpotent-henselian}
设 $(A, I)$ 为一个对，且 $I$ 局部幂零。则函子
$B \mapsto B/IB$ 诱导 $A$ 上 \'etale 代数范畴与 $A/I$ 上
\'etale 代数范畴之间的等价。此外，该对是 Hensel 对。
\end{lemma}

\begin{proof}
本质满性由 Algebra, Lemma \ref{algebra-lemma-lift-etale} 给出。
若 $B$、$B'$ 在 $A$ 上 \'etale，且 $B/IB \to B'/IB'$ 是
$A/I$-代数态射，则由 Algebra, Lemma
\ref{algebra-lemma-smooth-strong-lift} 可以提升该态射。
最后，设 $f, g : B \to B'$ 是两个满足
$f \bmod I = g \bmod I$ 的 $A$-代数映射。选取生成乘法映射
$B \otimes_A B \to B$ 之核的幂等元 $e \in B \otimes_A B$，见
Algebra, Lemmas \ref{algebra-lemma-diagonal-unramified} and
\ref{algebra-lemma-unramified}（后者用于说明 \'etale 蕴含非分歧）。
于是 $(f \otimes g)(e) \in IB'$。由于 $IB'$ 局部幂零
（Algebra, Lemma \ref{algebra-lemma-locally-nilpotent}），
由 Algebra, Lemma \ref{algebra-lemma-lift-idempotents} 可知
$(f \otimes g)(e) = 0$。故 $f = g$。

\medskip\noindent
显然，$I$ 包含于 $A$ 的 Jacobson 根中。设 $f \in A[T]$ 为首一
多项式，并设 $\overline{f} = g_0h_0$ 是
$\overline{f} = f \bmod I$ 的一个分解，其中
$g_0, h_0 \in A/I[T]$ 首一且在 $A/I[T]$ 中生成单位理想。
由引理 \ref{more-algebra-lemma-lift-factorization-monic}，存在诱导同构
$A/I \to A'/IA'$ 的 \'etale 环映射 $A \to A'$，使得该分解提升为
$A'$ 上首一多项式的分解。由上文可知 $A = A'$，故该分解定义在
$A$ 上。
\end{proof}

\begin{lemma}
\label{more-algebra-lemma-limit-henselian}
设 $A = \lim A_n$，其中 $(A_n)$ 是由满的转移映射组成的环的逆向系统，
并且这些转移映射的核局部幂零。令 $I_n = \Ker(A \to A_n)$。
则 $(A, I_n)$ 是 Hensel 对。
\end{lemma}

\begin{proof}
固定 $n$。设 $a \in A$ 是在 $A_n$ 中映到 $1$ 的元素。
由 Algebra, Lemma \ref{algebra-lemma-locally-nilpotent-unit}，
$a$ 在所有 $m \geq n$ 的 $A_m$ 中都映到单位元。因此 $a$ 是
$A$ 中的单位元。故由 Algebra, Lemma
\ref{algebra-lemma-contained-in-radical}，理想 $I_n$ 包含于
$A$ 的 Jacobson 根中。设 $f \in A[T]$ 为首一多项式，并设
$\overline{f} = g_nh_n$ 是 $\overline{f} = f \bmod I_n$ 的一个分解，
其中 $g_n, h_n \in A_n[T]$ 首一且在 $A_n[T]$ 中生成单位理想。
由引理 \ref{more-algebra-lemma-locally-nilpotent-henselian}，可以逐次把该分解提升为
$f \bmod I_m = g_m h_m$，其中对所有 $m \geq n$，
$g_m, h_m$ 都是 $A_m[T]$ 中的首一多项式。在每一步中，我们必须验证
所取提升 $g_m, h_m$ 在 $A_n[T]$ 中生成单位理想；这由
$g_n, h_n$ 的相应事实以及
$\Spec(A_n[T]) = \Spec(A_m[T])$ 得出，因为 $A_m \to A_n$ 的核
局部幂零。由于 $A = \lim A_m$，证明完成。
\end{proof}

\begin{lemma}
\label{more-algebra-lemma-complete-henselian}
设 $(A, I)$ 为一个对。若 $A$ 是 $I$-进完备的，则该对是 Hensel 对。
\end{lemma}

\begin{proof}
由 Algebra, Lemma \ref{algebra-lemma-radical-completion}，理想 $I$
包含于 $A$ 的 Jacobson 根中。设 $f \in A[T]$ 为首一多项式，且
$\overline{f} = g_0h_0$ 是 $\overline{f} = f \bmod I$ 的分解，
其中 $g_0, h_0 \in A/I[T]$ 首一且在 $A/I[T]$ 中生成单位理想。
由引理 \ref{more-algebra-lemma-locally-nilpotent-henselian}，可逐次将其提升为
$f \bmod I^n = g_n h_n$，其中对所有 $n \geq 1$，
$g_n, h_n$ 都是 $A/I^n[T]$ 中的首一多项式。
由于 $A = \lim A/I^n$，证明完成。
\end{proof}

\begin{lemma}
\label{more-algebra-lemma-helper-finite-type}
设 $(A, I)$ 为一个对。设 $A \to B$ 为有限型环映射，并且
$B/IB = C_1 \times C_2$，其中 $A/I \to C_1$ 是有限的。
令 $B'$ 为 $A$ 在 $B$ 中的整闭包。则可写成
$B'/IB' = C_1 \times C'_2$，使得映射 $B'/IB' \to B/IB$ 保持乘积分解，
并且存在映到 $C_1 \times C'_2$ 中 $(1, 0)$ 的 $g \in B'$，使
$B'_g \to B_g$ 为同构。
\end{lemma}

\begin{proof}
注意，$A \to B$ 在闭子集
$T = \Spec(C_1) \subset \Spec(B)$ 的每个素理想处都是拟有限的
（这可通过考察纤维环看出，见 Algebra, Definition
\ref{algebra-definition-quasi-finite}）。考虑拓扑空间的交换图
$$
\xymatrix{
\Spec(B) \ar[rr]_\phi \ar[rd]_\psi & & \Spec(B') \ar[ld]^{\psi'} \\
& \Spec(A)
}
$$
由 Algebra, Theorem \ref{algebra-theorem-main-theorem}，对每个
$\mathfrak p \in T$，存在 $h_\mathfrak p \in B'$，满足
$h_\mathfrak p \not \in \mathfrak p$，且 $B'_h \to B_h$ 是同构。
并集 $U = \bigcup D(h_\mathfrak p)$ 给出开集
$U \subset \Spec(B')$，使得 $\phi^{-1}(U) \to U$ 为同胚且
$T \subset \phi^{-1}(U)$。由于 $T$ 在 $\psi^{-1}(V(I))$ 中为开集，
可知 $\phi(T)$ 在 $U \cap (\psi')^{-1}(V(I))$ 中为开集。
因此 $\phi(T)$ 在 $(\psi')^{-1}(V(I))$ 中为开集。
另一方面，由于 $C_1$ 在 $A/I$ 上有限，它在 $B'$ 上也是有限的。
故由 Algebra, Lemmas \ref{algebra-lemma-going-up-closed} and
\ref{algebra-lemma-integral-going-up}，$\phi(T)$ 是 $\Spec(B')$ 的闭子集。
我们推出 $\Spec(B'/IB') \supset \phi(T)$ 既开又闭。由
Algebra, Lemma \ref{algebra-lemma-disjoint-implies-product}，得到相应的
乘积分解 $B'/IB' = C'_1 \times C'_2$。通过考察谱上的情形可见，映射
$B'/IB' \to B/IB$ 把 $C'_1$ 映入 $C_1$，把 $C'_2$ 映入 $C_2$
（提示：$\phi(T)$ 的逆像恰为 $T$；略去一些细节）。
选取映到 $C'_1 \times C'_2$ 中 $(1, 0)$ 的 $g \in B'$，使得
$D(g) \subset U$；这是可能的，因为 $\Spec(C'_1)$ 与
$\Spec(C'_2)$ 在 $\Spec(B')$ 中是不交闭集，且 $\Spec(C'_1)$ 包含于
$U$。于是 $B'_g \to B_g$ 在谱上定义同胚，并在局部环上定义同构
（由上面对 $U$ 的选取）。因此它是同构，例如这可由
Algebra, Lemma \ref{algebra-lemma-characterize-zero-local} 推出。
最后，由此有 $C'_1 = C_1$，证明完成。
\end{proof}

\begin{lemma}
\label{more-algebra-lemma-characterize-henselian-pair}
\begin{reference}
\cite[Chapter XI]{Henselian} and \cite[Proposition 1]{Gabber-henselian}
\end{reference}
设 $(A, I)$ 为一个对。下列条件等价：
\begin{enumerate}
\item $(A, I)$ 是 Hensel 对；
\item 给定 \'etale 环映射 $A \to A'$ 以及 $A$-代数映射
$\sigma : A' \to A/I$，存在提升 $\sigma$ 的 $A$-代数映射
$A' \to A$；
\item 对任意有限 $A$-代数 $B$，映射 $B \to B/IB$ 在幂等元上
诱导双射；
\item 对任意整 $A$-代数 $B$，映射 $B \to B/IB$ 在幂等元上
诱导双射；
\item （Gabber）$I$ 包含于 $A$ 的 Jacobson 根中，并且每个形如
$$
f(T) = T^n(T - 1) + a_n T^n + \ldots + a_1 T + a_0
$$
的首一多项式 $f(T) \in A[T]$，其中
$a_n, \ldots, a_0 \in I$ 且 $n \ge 1$，都有一个根
$\alpha \in 1 + I$。
\end{enumerate}
此外，第 (5) 项中的根是唯一的。
\end{lemma}

\begin{proof}
假设 (2) 成立。则 $I$ 包含于 $A$ 的 Jacobson 根中；否则，存在一个
模 $I$ 同余于 $1$ 的非单位元 $f \in A$，而映射 $A \to A_f$
将与 (2) 矛盾。若 $B$ 在 $A$ 上整，则 $IB \subset B$ 包含于
$B$ 的 Jacobson 根中，因为由 Algebra, Lemmas
\ref{algebra-lemma-going-up-closed} and
\ref{algebra-lemma-integral-going-up}，$\Spec(B) \to \Spec(A)$ 是闭映射。
因此，由引理 \ref{more-algebra-lemma-idempotents-determined-modulo-radical}，
从 $B$ 的幂等元到 $B/IB$ 的幂等元的映射是单射。另一方面，
由于 (2) 成立，由引理 \ref{more-algebra-lemma-lift-idempotent-upstairs}，
$B/IB$ 的每个幂等元都提升为 $B$ 的幂等元。因此 (2) 蕴含 (4)。

\medskip\noindent
(4) $\Rightarrow$ (3) 是平凡的。

\medskip\noindent
假设 (3)。设 $\mathfrak m$ 为极大理想，并考虑有限映射
$A \to B = A/(I \cap \mathfrak m)$。$B \to B/IB$ 在幂等元上诱导
双射这一条件蕴含 $I \subset \mathfrak m$（否则
$B = A/I \times A/\mathfrak m$ 且 $B/IB = A/I$）。因此 $I$ 包含于
$A$ 的 Jacobson 根中。设 $f \in A[T]$ 首一，并假设给定分解
$\overline{f} = g_0h_0$，其中 $g_0, h_0 \in A/I[T]$ 首一且在
$A/I[T]$ 中生成单位理想。令 $B = A[T]/(f)$。令
$\overline{e}$ 为 $B/IB$ 中对应于 $A$-代数分解
$$
B/IB = A/I[T]/(g_0) \times A/I[T]/(h_0)
$$
的幂等元。由假设 (3)，存在提升 $\overline{e}$ 的幂等元
$e \in B$。这给出乘积分解
$$
B = eB \times (1 - e)B
$$
注意，$B$ 作为 $A$-模是秩为 $\deg(f)$ 的自由模。因此 $eB$ 和
$(1 - e)B$ 是有限局部自由 $A$-模。然而，由于 $eB$ 和
$(1 - e)B$ 在 $A/I$ 上分别有常秩 $\deg(g_0)$ 和 $\deg(h_0)$，
我们发现它们在 $\Spec(A)$ 上也如此。由此
\begin{align*}
f & = \text{CharPol}_A(T : B \to B) \\
& =
\text{CharPol}_A(T : eB \to eB)
\text{CharPol}_A(T : (1 - e)B \to (1 - e)B)
\end{align*}
是首一多项式的分解，模 $I$ 后约化为给定的分解。这里
$\text{CharPol}_A$ 表示有限局部自由 $A$-模之自同态的特征多项式。
若该模自由，则 $\text{CharPol}_A$ 定义为相应矩阵的特征多项式；
一般情形使用 Algebra, Lemma \ref{algebra-lemma-standard-covering}
进行粘合。略去细节。故 (3) 蕴含 (1)。

\medskip\noindent
假设 (1)。令 $f$ 如 (5) 所述。$f \bmod I$ 分解为 $T^n$ 与
$T - 1$ 的乘积；由定义 \ref{more-algebra-definition-henselian-pair}，它提升为
首一多项式的分解 $f = gh$，其中 $g$ 与 $h$ 首一。于是 $h$ 必须次数为 $1$，从而
$f$ 有一个模 $I$ 约化为 $1$ 的根。最后，由 Hensel 对的定义，
$I$ 包含于 Jacobson 根中。因此 (1) 蕴含 (5)。

\medskip\noindent
在证明最后一步前，先说明 (5) 中的根 $\alpha$ 若存在则唯一。
事实上，由 $f(T)$ 的明确形状，有 $f'(\alpha) \in 1 + I$，其中
$f'$ 是 $f$ 关于 $T$ 的导数。初等论证给出
$$
f(T) = f(\alpha + T - \alpha) =
f(\alpha) + f'(\alpha) \cdot (T - \alpha) \bmod
(T - \alpha)^2 A[T]
$$
这表明，$f(T)$ 的任意另一个根 $\alpha' \in 1 + I$ 都满足
$0 = f(\alpha') - f(\alpha) = (\alpha' - \alpha)(1 + i)$，其中某个
$i \in I$。由于 $1 + i$ 是 $A$ 中的单位元，故
$\alpha = \alpha'$。

\medskip\noindent
假设 (5)。我们证明 (2) 成立；换言之，对每个 \'etale 映射
$A \to A'$，模 $I$ 的每个截面 $\sigma : A' \to A/I$ 都提升为
截面 $A' \to A$。由于 $A \to A'$ 是 \'etale 的，截面 $\sigma$
确定 $A/I$-代数的分解
\begin{equation}
\label{more-algebra-equation-GCHP}
A'/IA' \cong A/I \times C
\end{equation}
事实上，由 Algebra, Lemma \ref{algebra-lemma-map-between-etale}，满环映射
$A'/IA' \to A/I$ 是 \'etale 的；继而由 Algebra, Lemma
\ref{algebra-lemma-surjective-flat-finitely-presented} 得到所需幂等元。
我们将证明，该分解提升为 $A$-代数的分解
\begin{equation}
\label{more-algebra-equation-GCHP-want}
A' \cong A'_1 \times A'_2
\end{equation}
并且 $A'_1$ 在 $A$ 上整。于是 $A \to A'_1$ 既整又 \'etale，且
$A/I \to A'_1/IA'_1$ 是同构；故由引理
\ref{more-algebra-lemma-check-isomorphism-zariski}，$A \to A'_1$ 是同构
（这里还用到 \'etale 环映射平坦且有限表示，见 Algebra, Lemma
\ref{algebra-lemma-etale}）。

\medskip\noindent
令 $B'$ 为 $A$ 在 $A'$ 中的整闭包。由引理
\ref{more-algebra-lemma-helper-finite-type}，可把它相容于
(\ref{more-algebra-equation-GCHP}) 地分解为 $A/I$-代数
\begin{equation}
\label{more-algebra-equation-dec-mod-I}
B'/IB' \cong A/I \times C'
\end{equation}
并可找到提升 $(1, 0)$ 的 $b \in B'$，使得
$B'_b \to A'_b$ 是同构。若分解
(\ref{more-algebra-equation-dec-mod-I}) 提升为 $A$-代数的分解
\begin{equation}
\label{more-algebra-equation-want-2}
B' \cong B'_1 \times B'_2
\end{equation}
则诱导的分解 $A' = A'_1 \times A'_2$ 给出所需的
(\ref{more-algebra-equation-GCHP-want})：事实上，由于 $b$ 是 $B'_1$ 中的单位元
（略去细节），有 $B'_1 \cong A'_1$，因而 $A'_1$ 在 $A$ 上整。

\medskip\noindent
选取包含 $b$ 的有限 $A$-子代数 $B'' \subset B'$（注意，$B'$ 的
任意有限生成 $A$-子代数都在 $A$ 上有限）。扩大 $B''$ 后，可假设
$b$ 在 $B''/IB''$ 中映到给出如下分解的幂等元：
\begin{equation}
\label{more-algebra-equation-again-dec-mod-I}
B''/IB'' \cong C''_1 \times C''_2
\end{equation}
由于 $B'_b \cong A'_b$，可知 $B'_b$ 在 $A$ 上是有限型的。
设 $B'_b$ 由 $b_1/b^n, \ldots, b_t/b^n$ 在 $A$ 上生成；扩大
$B''$，使得 $b_1, \ldots, b_t \in B''$。于是
$B''_b \to B'_b$ 既满又单，故是同构。特别地，
$C''_1 = A/I$！因此 $A/I \to C''_1$ 是同构，尤其是满射。
由引理 \ref{more-algebra-lemma-helper-finite}，可找到形如
$$
f(T) = T^n(T - 1) + a_n T^n + \ldots + a_1 T + a_0
$$
的 $f(T) \in A[T]$，其中 $a_n, \ldots, a_0 \in I$、
$n \ge 1$，并且 $f(b) = 0$。特别地，$B'$ 是一个
$A[T]/(f)$-代数。由 (5)，存在 $f$ 的根 $a \in 1 + I$。
这产生乘积分解 $A[T]/(f) = A[T]/(T - a) \times D$，它与
$B'/IB'$ 的分裂 (\ref{more-algebra-equation-dec-mod-I}) 相容。由此诱导的
$B'$ 的分裂即为所需的 (\ref{more-algebra-equation-want-2})。
\end{proof}

\begin{lemma}
\label{more-algebra-lemma-change-ideal-henselian-pair}
设 $A$ 为环。设 $I, J \subset A$ 为满足 $V(I) = V(J)$ 的理想。
则 $(A, I)$ 是 Hensel 对当且仅当 $(A, J)$ 是 Hensel 对。
\end{lemma}

\begin{proof}
对任意整环映射 $A \to B$，有 $V(IB) = V(JB)$。故 $B/IB$ 与
$B/JB$ 的幂等元一一对应
（Algebra, Lemma \ref{algebra-lemma-disjoint-decomposition}）。
因此，$B \to B/IB$ 在幂等元集合上诱导双射当且仅当
$B \to B/JB$ 在幂等元集合上诱导双射。结论由引理
\ref{more-algebra-lemma-characterize-henselian-pair} 得出。
\end{proof}

\begin{lemma}
\label{more-algebra-lemma-integral-over-henselian-pair}
设 $(A, I)$ 为 Hensel 对，并设 $A \to B$ 为整环映射。
则 $(B, IB)$ 是 Hensel 对。
\end{lemma}

\begin{proof}
这立即由引理 \ref{more-algebra-lemma-characterize-henselian-pair} 中 Hensel 对的
第四个刻画以及整环映射的复合仍为整映射这一事实得出。
\end{proof}

\begin{lemma}
\label{more-algebra-lemma-henselian-henselian-pair}
设 $I \subset J \subset A$ 为环 $A$ 的理想。下列条件等价：
\begin{enumerate}
\item $(A, I)$ 与 $(A/I, J/I)$ 都是 Hensel 对；
\item $(A, J)$ 是 Hensel 对。
\end{enumerate}
\end{lemma}

\begin{proof}
假设 (1)。令 $B$ 为整 $A$-代数。考虑环映射
$$
B \to B/IB \to B/JB
$$
由引理 \ref{more-algebra-lemma-characterize-henselian-pair}，两个箭头在幂等元上都
诱导双射，故其复合也如此。由引理
\ref{more-algebra-lemma-characterize-henselian-pair}，$(A, J)$ 是 Hensel 对。

\medskip\noindent
反之，假设 (2) 成立。由引理 \ref{more-algebra-lemma-integral-over-henselian-pair}，
$(A/I, J/I)$ 是 Hensel 对。令 $B$ 为整 $A$-代数。考虑环映射
$$
B \to B/IB \to B/JB
$$
由引理 \ref{more-algebra-lemma-characterize-henselian-pair}，复合与第二个箭头在
幂等元上诱导双射，故第一个箭头也是如此。因此，再由该引理，
$(A, I)$ 是 Hensel 对。
\end{proof}

\begin{lemma}
\label{more-algebra-lemma-sum-henselian}
设 $A$ 为环，且 $(A, I)$ 与 $(A, I')$ 都是 Hensel 对。
则 $(A, I + I')$ 是 Hensel 对。
\end{lemma}

\begin{proof}
由引理 \ref{more-algebra-lemma-integral-over-henselian-pair}，对
$(A/I, (I' + I)/I)$ 是 Hensel 对。因此结论由引理
\ref{more-algebra-lemma-henselian-henselian-pair} 得出。
\end{proof}

\begin{lemma}
\label{more-algebra-lemma-product-henselian-pairs}
设 $J$ 为集合，且 $\{ (A_j, I_j)\}_{j \in J}$ 为一族对。
则 $(\prod_{j \in J} A_j, \prod_{j\in J} I_j)$ 是 Hensel 对，
当且仅当每个 $(A_j, I_j)$ 都是 Hensel 对。
\end{lemma}
 
\begin{proof}
对每个 $j \in J$，投影 $\prod_{j \in J} A_j \rightarrow A_j$
是整环映射。因此，若
$(\prod_{j \in J} A_j, \prod_{j\in J} I_j)$ 是 Hensel 对，
则引理 \ref{more-algebra-lemma-integral-over-henselian-pair} 证明每个
$(A_j, I_j)$ 都是 Hensel 对。

\medskip\noindent
反之，假设每个 $(A_j, I_j)$ 都是 Hensel 对。对任意
$x \in \prod_{j \in J} I_j$，元素 $1 + x$ 在
$\prod_{j \in J} A_j$ 中是单位元，因为由 Algebra, Lemma
\ref{algebra-lemma-contained-in-radical} 和定义
\ref{more-algebra-definition-henselian-pair}，这在每个分量上都成立。
再次由 Algebra, Lemma \ref{algebra-lemma-contained-in-radical}，
$\prod_{j \in J} I_j$ 包含于 $\prod_{j \in J} A_j$ 的 Jacobson 根中。
继续逐分量考察可知，对每个首一多项式
$f \in (\prod_{j \in J} A_j)[T]$，以及每个分解
$\overline{f} = g_0h_0$，其中首一多项式
$g_0, h_0 \in (\prod_{j \in J} A_j / \prod_{j \in J} I_j)[T] =
(\prod_{j \in J} A_j/I_j)[T]$ 在
$(\prod_{j \in J} A_j / \prod_{j \in J} I_j)[T]$ 中生成单位理想，
都存在 $(\prod_{j \in J} A_j)[T]$ 中的分解 $f = gh$，其中
$g$、$h$ 首一且分别约化为 $g_0$、$h_0$。综上，由定义
\ref{more-algebra-definition-henselian-pair}，
$(\prod_{j \in J} A_j, \prod_{j\in J} I_j)$ 是 Hensel 对。
\end{proof}

\begin{lemma}
\label{more-algebra-lemma-limits-henselian}
Hensel 性在对的极限下保持。更确切地说，设 $J$ 为预序集，并设
$(A_j, I_j)$ 是 $J$ 上 Hensel 对的逆向系统。则
$A = \lim A_j$ 配以理想 $I = \lim I_j$ 所得的对 $(A, I)$ 是
Hensel 对。
\end{lemma}

\begin{proof}
由 Categories, Lemma
\ref{categories-lemma-limits-products-equalizers}，只需考虑乘积和等化子。
对乘积，断言由引理 \ref{more-algebra-lemma-product-henselian-pairs} 得出。
因此，考虑等化子图
$$
\xymatrix{
(A, I) \ar[r] & (A', I')
\ar@<1ex>[r]^{\varphi} \ar@<-1ex>[r]_{\psi}
&
(A'', I'')  
}
$$
其中对 $(A', I')$ 与 $(A'', I'')$ 都是 Hensel 对。为验证
$(A, I)$ 也是 Hensel 对，我们使用引理
\ref{more-algebra-lemma-characterize-henselian-pair} 中 Gabber 的判别准则。
$1 + I$ 的每个元素都是 $A$ 中的单位元，因为由单位元之逆的唯一性，
这一点可在 $(A', I')$ 中检验。因此 $I$ 包含于 $A$ 的 Jacobson 根中，
见 Algebra, Lemma \ref{algebra-lemma-contained-in-radical}。
于是，令
$$
f(T) = T^{N - 1}(T - 1) + a_{N - 1} T^{N - 1} + \dotsb + a_1 T + a_0
$$
为 $A[T]$ 中的多项式，其中
$a_{N - 1}, \dotsc, a_0 \in I$ 且 $N \ge 1$。$f(T)$ 在 $A'[T]$
中的像有唯一根 $\alpha' \in 1 + I'$，其在 $A''[T]$ 中的进一步像也
同样如此。因此，由唯一性，$\varphi(\alpha') = \psi(\alpha')$；
也就是说，$\alpha'$ 在 $1 + I$ 中定义了 $f(T)$ 的一个根，正合所需。
\end{proof}

\begin{lemma}
\label{more-algebra-lemma-filtered-colimits-henselian}
Hensel 性在对的滤过余极限下保持。更确切地说，设 $J$ 为有向集，
并设 $(A_j, I_j)$ 是 $J$ 上 Hensel 对的系统。则
$A = \colim A_j$ 配以理想 $I = \colim I_j$ 所得的对 $(A, I)$
是 Hensel 对。
\end{lemma}

\begin{proof}
若 $u \in 1 + I$，则对某个 $j \in J$，$u$ 是某个
$u_j \in 1 + I_j$ 的像。由 Algebra, Lemma
\ref{algebra-lemma-contained-in-radical} 以及 $I_j$ 包含于 $A_j$ 的
Jacobson 根中这一假设，$u_j$ 在 $A_j$ 中可逆。因此 $u$ 在 $A$ 中
可逆。故由该引理，$I$ 包含于 $A$ 的 Jacobson 根中。

\medskip\noindent
设 $f \in A[T]$ 为首一多项式，并设 $\overline{f} = g_0 h_0$ 是一个
分解，其中 $g_0, h_0 \in A/I[T]$ 首一且生成 $A/I[T]$ 中的单位理想。
对某些 $g'_0, h'_0 \in A/I[T]$，写成
$1 = g_0 g'_0 + h_0 h'_0$。由于 $A = \colim A_j$ 与
$A/I = \colim A_j/I_j$ 都是滤过余极限，可以找到 $j \in J$、
$f_j \in A_j$ 以及分解 $\overline{f}_j = g_{j, 0} h_{j, 0}$，其中
$g_{j, 0}, h_{j, 0} \in A_j/I_j[T]$ 首一，并且对某些
$g'_{j, 0}, h'_{j, 0} \in A_j/I_j[T]$ 有
$1 = g_{j, 0} g'_{j, 0} + h_{j, 0} h'_{j, 0}$；这里
$f_j, g_{j, 0}, h_{j, 0}, g'_{j, 0}, h'_{j, 0}$ 分别映到
$f, g_0, h_0, g'_0, h'_0$。由于 $(A_j, I_j)$ 是 Hensel 对，
可把 $\overline{f}_j = g_{j, 0} h_{j, 0}$ 提升为 $A_j$ 上的分解；
取其在 $A$ 中的像，得到 $A$ 中的相应分解。因此 $(A, I)$ 是 Hensel 对。
\end{proof}

\begin{example}[Moret-Bailly]
\label{more-algebra-example-moret-bailly}
若余极限不是滤过的，则引理 \ref{more-algebra-lemma-filtered-colimits-henselian} 不成立。
例如，取 Hensel 对 $(\mathbf{Z}_p, (p))$ 与
$(\mathbf{Z}_p, (p))$ 的余积，则得到 $(A, pA)$，其中
$A = \mathbf{Z}_p \otimes_\mathbf{Z} \mathbf{Z}_p$。这不是 Hensel 对：
$A/pA = \mathbf{F}_p$，故若 $(A, pA)$ 是 Hensel 对，则 $A$ 必须是
局部环。然而，$\Spec(A)$ 不连通；例如，对奇素数 $p$，有非平凡幂等元
$$
(1/2 \otimes 1)
\left(1 \otimes 1 - (1 + p)^{-1}u \otimes u\right)
$$
其中 $u \in \mathbf{Z}_p$ 是 $1 + p$ 的一个平方根。略去一些细节。
\end{example}

\begin{lemma}
\label{more-algebra-lemma-largest-ideal-henselian}
设 $A$ 为环。存在最大的理想 $I \subset A$，使得 $(A, I)$ 是 Hensel 对。
\end{lemma}

\begin{proof}
结合引理 \ref{more-algebra-lemma-henselian-henselian-pair}、
\ref{more-algebra-lemma-sum-henselian} 和 \ref{more-algebra-lemma-filtered-colimits-henselian}。
\end{proof}

\begin{lemma}
\label{more-algebra-lemma-irreducible-henselian-pair-connected}
设 $(A, I)$ 为 Hensel 对。设 $\mathfrak p \subset A$ 为素理想。
则 $V(\mathfrak p + I)$ 连通。
\end{lemma}

\begin{proof}
由引理 \ref{more-algebra-lemma-integral-over-henselian-pair}，
$(A/\mathfrak p, I + \mathfrak p/\mathfrak p)$ 是 Hensel 对。
故只需证明：若 $(A, I)$ 是 Hensel 对且 $A$ 是整环，则
$\Spec(A/I) = V(I)$ 连通。否则，由 Algebra, Lemma
\ref{algebra-lemma-characterize-spec-connected}，$A/I$ 有非平凡幂等元。
由引理 \ref{more-algebra-lemma-characterize-henselian-pair}，这将蕴含 $A$ 有非平凡
幂等元，矛盾。
\end{proof}








\section{对的 Hensel 化}
\label{more-algebra-section-henselization}

\noindent
我们继续第 \ref{more-algebra-section-henselian-pairs} 节开始的讨论。

\begin{lemma}
\label{more-algebra-lemma-henselization}
包含函子
$$
\text{category of henselian pairs}
\longrightarrow
\text{category of pairs}
$$
有左伴随 $(A, I) \mapsto (A^h, I^h)$。
\end{lemma}

\begin{proof}
设 $(A, I)$ 为一个对。考虑由 \'etale 环映射 $A \to B$ 组成的范畴
$\mathcal{C}$，其中 $A/I \to B/IB$ 是同构。我们将证明范畴
$\mathcal{C}$ 是有向的，并且令
$A^h = \colim_{B \in \mathcal{C}} B$、$I^h = IA^h$ 给出所需伴随。

\medskip\noindent
先证明 $\mathcal{C}$ 是有向的
（Categories, Definition \ref{categories-definition-directed}）。
它非空，因为 $\text{id} : A \to A$ 是一个对象。若 $B$ 与 $B'$ 是
$\mathcal{C}$ 的两个对象，则 $B'' = B \otimes_A B'$ 是
$\mathcal{C}$ 的对象（使用 Algebra, Lemma \ref{algebra-lemma-etale}），
并且存在态射 $B \to B''$ 与 $B' \to B''$。设
$f, g : B \to B'$ 是 $\mathcal{C}$ 中两个对象之间的映射。则一个
余等化子是
$$
(B' \otimes_{f, B, g} B') \otimes_{(B' \otimes_A B')} B'
$$
由 Algebra, Lemmas \ref{algebra-lemma-etale} and
\ref{algebra-lemma-map-between-etale}，它在 $A$ 上 \'etale。
因此范畴 $\mathcal{C}$ 是有向的。

\medskip\noindent
由于对 $\mathcal{C}$ 的所有对象 $B$ 都有 $B/IB = A/I$，故
$A^h/I^h = A^h/IA^h = \colim B/IB = \colim A/I = A/I$。

\medskip\noindent
下面证明 $A^h = \colim_{B \in \mathcal{C}} B$ 配以
$I^h = IA^h$ 是 Hensel 对。为此，验证引理
\ref{more-algebra-lemma-characterize-henselian-pair} 的条件 (2)。具体地，设给定
\'etale 环映射 $A^h \to A'$ 以及 $A^h$-代数映射
$\sigma : A' \to A^h/I^h$。则存在 $B \in \mathcal{C}$ 及
\'etale 环映射 $B \to B'$，使得 $A' = B' \otimes_B A^h$，见
Algebra, Lemma \ref{algebra-lemma-etale}。由于 $A^h/I^h = A/I$，
映射 $\sigma$ 诱导 $A$-代数映射 $s : B' \to A/I$。于是作为
$A/I$-代数有 $B'/IB' = A/I \times C$，其中 $C$ 是由 $s$ 诱导的
映射 $B'/IB' \to A/I$ 的核。令 $g \in B'$ 映到
$(1, 0) \in A/I \times C$。则 $B \to B'_g$ 是 \'etale 的，且
$A/I \to B'_g/IB'_g$ 是同构；即 $B'_g$ 是 $\mathcal{C}$ 的对象。
因此得到典范映射 $B'_g \to A^h$，使得
$$
\vcenter{
\xymatrix{
B'_g \ar[r] & A^h \\
B \ar[u] \ar[ur]
}
}
\quad\text{且}\quad
\vcenter{
\xymatrix{
B' \ar[r] \ar[rrd]_s & B'_g \ar[r] & A^h \ar[d] \\
& & A/I
}
}
$$
交换。这诱导与 $\sigma$ 相容的映射
$A' = B' \otimes_B A^h \to A^h$，正合所需。

\medskip\noindent
设 $(A, I) \to (A', I')$ 为对的态射，且 $(A', I')$ 是 Hensel 对。
我们证明存在唯一分解 $A \to A^h \to A'$，从而完成证明。
事实上，对 $\mathcal{C}$ 中每个 $A \to B$，环映射
$A' \to B' = A' \otimes_A B$ 是 \'etale 的，并诱导同构
$A'/I' \to B'/I'B'$。故由引理
\ref{more-algebra-lemma-characterize-henselian-pair}，存在截面
$\sigma_B : B' \to A'$。给定 $\mathcal{C}$ 中的态射
$B_1 \to B_2$，我们断言下图交换：
$$
\xymatrix{
B'_1 \ar[rr] \ar[rd]_{\sigma_{B_1}} & &
B'_2 \ar[ld]^{\sigma_{B_2}} \\
& A'
}
$$
只需证明，对 $\mathcal{C}$ 中每个 $B$，截面 $\sigma_B$ 是唯一的
$A'$-代数映射 $B' \to A'$。由 Algebra, Lemma
\ref{algebra-lemma-diagonal-unramified}，对某个环 $R$ 有
$B' \otimes_{A'} B' = B' \times R$。在本情形中，由
$B'/I'B' = A'/I'$ 有 $R/I'R = 0$。因此，给定两个 $A'$-代数映射
$\sigma_B, \sigma_B' : B' \to A'$，则
$e = (\sigma_B \otimes \sigma_B')(0, 1) \in A'$ 是包含于 $I'$ 的
幂等元。由引理 \ref{more-algebra-lemma-idempotents-determined-modulo-radical}，
$e = 0$。故 $\sigma_B = \sigma_B'$，正合所需。
利用交换性，得到
$$
A^h = \colim_{B \in \mathcal{C}} B \to
\colim_{B \in \mathcal{C}} A' \otimes_A B \xrightarrow{\colim \sigma_B} A'
$$
正合所需。映射 $\sigma_B$ 的唯一性还保证该映射唯一。
故 $(A, I) \mapsto (A^h, I^h)$ 是所需伴随。
\end{proof}

\begin{lemma}
\label{more-algebra-lemma-henselization-flat}
设 $(A, I)$ 为一个对。令 $(A^h, I^h)$ 如引理
\ref{more-algebra-lemma-henselization} 所述。则 $A \to A^h$ 平坦，
$I^h = IA^h$，并且对所有 $n$，$A/I^n \to A^h/I^nA^h$ 都是同构。
\end{lemma}

\begin{proof}
在引理 \ref{more-algebra-lemma-henselization} 的证明中，我们已经看到，$A^h$ 是
满足 $A/I \to B/IB$ 为同构的 \'etale $A$-代数 $B$ 的滤过余极限，
并且 $I^h = IA^h$。由于 \'etale 环映射平坦
（Algebra, Lemma \ref{algebra-lemma-etale}），由 Algebra, Lemma
\ref{algebra-lemma-colimit-flat} 可知 $A \to A^h$ 平坦。
由于每个 $A \to B$ 都平坦，映射 $A/I^n \to B/I^nB$ 也都是同构
（例如由 Algebra, Lemma \ref{algebra-lemma-lift-basis}）。
取余极限，得到 $A/I^n = A^h/I^nA^h$，正合所需。
\end{proof}

\begin{lemma}
\label{more-algebra-lemma-henselization-local-ring}
\begin{slogan}
对的 Hensel 化与局部环的 Hensel 化相容。
\end{slogan}
引理 \ref{more-algebra-lemma-henselization} 的函子把局部环
$(A, \mathfrak m)$ 映为其 Hensel 化。
\end{lemma}

\begin{proof}
令 $(A^h, \mathfrak m^h)$ 为引理 \ref{more-algebra-lemma-henselization} 中构造的
对 $(A, \mathfrak m)$ 的 Hensel 化。由引理
\ref{more-algebra-lemma-henselization-flat}，$\mathfrak m^h = \mathfrak m A^h$ 是极大
理想；又因它包含于 Jacobson 根中，我们推出 $A^h$ 是以
$\mathfrak m^h$ 为极大理想的局部环。说完这些，有两种方法完成证明。

\medskip\noindent
第一种证明：注意，在 Algebra, Lemma
\ref{algebra-lemma-henselization} 的证明中作为余极限的构造，与引理
\ref{more-algebra-lemma-henselization} 中用于构造 $A^h$ 的余极限相同。
第二种证明：Hensel 化 $A \to S$ 与引理
\ref{more-algebra-lemma-henselization} 的 $A \to A^h$ 都是局部环同态；
$S$ 与 $A^h$ 都是 \'etale $A$-代数的滤过余极限；$S$ 与 $A^h$
都是 Hensel 局部环；并且 $S$ 与 $A^h$ 的剩余域都等于
$\kappa(\mathfrak m)$（第二种情形使用引理
\ref{more-algebra-lemma-henselization-flat}）。因此，由 Algebra, Lemma
\ref{algebra-lemma-uniqueness-henselian}，它们典范同构。
\end{proof}

\begin{lemma}
\label{more-algebra-lemma-henselization-Noetherian-pair}
\begin{slogan}
Noether 对的 Hensel 化仍是 Noether 对，并且具有相同的完备化。
\end{slogan}
设 $(A, I)$ 为一个对，且 $A$ 为 Noether 环。令 $(A^h, I^h)$ 如引理
\ref{more-algebra-lemma-henselization} 所述。则 $I$-进完备化之间的映射
$$
A^\wedge \to (A^h)^\wedge
$$
是同构。此外，$A^h$ 是 Noether 环，映射
$A \to A^h \to A^\wedge$ 都平坦，且
$A^h \to A^\wedge$ 忠实平坦。
\end{lemma}

\begin{proof}
第一个断言立即由引理 \ref{more-algebra-lemma-henselization-flat} 得出，事实上无需
假设 $A$ 是 Noether 环。在引理 \ref{more-algebra-lemma-henselization} 的证明中，
我们已经看到，$A^h$ 是满足 $A/I \to B/IB$ 为同构的 \'etale
$A$-代数 $B$ 的滤过余极限。对每个这样的 $A \to B$，诱导映射
$A^\wedge \to B^\wedge$ 是同构
（见引理 \ref{more-algebra-lemma-henselization-flat} 的证明）。
由 Algebra, Lemma \ref{algebra-lemma-completion-flat}，对每个 $B$，
环映射 $B \to A^\wedge = B^\wedge = (A^h)^\wedge$ 都平坦。
因此，由 Algebra, Lemma \ref{algebra-lemma-colimit-rings-flat}，
$A^h \to A^\wedge = (A^h)^\wedge$ 平坦。由于
$I^h = IA^h$ 包含于 $A^h$ 的 Jacobson 根中，并且
$A^h \to A^\wedge$ 诱导同构 $A^h/I^h \to A/I$，由
Algebra, Lemma \ref{algebra-lemma-ff}，$A^h \to A^\wedge$ 忠实平坦。
由 Algebra, Lemma \ref{algebra-lemma-completion-Noetherian-Noetherian}，
环 $A^\wedge$ 是 Noether 环。故由 Algebra, Lemma
\ref{algebra-lemma-descent-Noetherian}，$A^h$ 是 Noether 环。
\end{proof}

\begin{lemma}
\label{more-algebra-lemma-henselization-colimit}
设 $(A, I) = \colim (A_i, I_i)$ 是对的滤过余极限。引理
\ref{more-algebra-lemma-henselization} 的函子给出 $A^h = \colim A_i^h$ 和
$I^h = \colim I_i^h$。
\end{lemma}

\noindent
对非滤过余极限，此引理不成立，见例 \ref{more-algebra-example-moret-bailly}。

\begin{proof}
由 Categories, Lemma \ref{categories-lemma-adjoint-exact}，
$(A^h, I^h)$ 是系统 $(A_i^h, I_i^h)$ 在 Hensel 对范畴中的余极限。
因此，对 Hensel 对 $(B, J)$，有
$$
\Mor((A^h, I^h), (B, J)) =
\lim \Mor((A_i^h, I_i^h), (B, J)) =
\Mor(\colim (A_i^h, I_i^h), (B, J))
$$
这里余极限取在对的范畴中。由于该余极限滤过，在对的范畴中有
$\colim (A_i^h, I_i^h) = (\colim A_i^h, \colim I_i^h)$；略去细节。
再次利用余极限滤过，可知这是 Hensel 对（引理
\ref{more-algebra-lemma-filtered-colimits-henselian}）。故由 Yoneda 引理，
$(A^h, I^h) = (\colim A_i^h, \colim I_i^h)$。
\end{proof}

\begin{lemma}
\label{more-algebra-lemma-henselization-change-ideal}
\begin{slogan}
一个对的 Hensel 化只依赖于理想的根。
\end{slogan}
设 $A$ 为环，$I$ 和 $J$ 是其理想。若 $V(I) = V(J)$，则引理
\ref{more-algebra-lemma-henselization} 的函子对 $(A, I)$ 与 $(A, J)$ 产生同一个环。
\end{lemma}

\begin{proof}
令 $(A', IA')$ 为从对 $(A, I)$ 出发由引理
\ref{more-algebra-lemma-henselization} 产生的对，见引理
\ref{more-algebra-lemma-henselization-flat}。令 $(A'', JA'')$ 为从对 $(A, J)$ 出发由
引理 \ref{more-algebra-lemma-henselization} 产生的对。由引理
\ref{more-algebra-lemma-change-ideal-henselian-pair}，$(A', JA')$ 与
$(A'', IA'')$ 都是 Hensel 对。由该构造的泛性质，得到唯一的
$A$-代数映射 $A'' \to A'$ 与 $A' \to A''$。唯一性表明它们互为逆映射。
\end{proof}

\begin{lemma}
\label{more-algebra-lemma-henselization-integral}
\begin{slogan}
Hensel 化与整基变换交换。
\end{slogan}
设 $(A, I) \to (B, J)$ 为满足 $V(J) = V(IB)$ 的对的映射。
令 $(A^h , I^h) \to (B^h, J^h)$ 为 Hensel 化上诱导的映射
（引理 \ref{more-algebra-lemma-henselization}）。若 $A \to B$ 是整的，则诱导映射
$A^h \otimes_A B \to B^h$ 是同构。
\end{lemma}

\begin{proof}
由引理 \ref{more-algebra-lemma-henselization-change-ideal}，可假设 $J = IB$。
由引理 \ref{more-algebra-lemma-integral-over-henselian-pair}，对
$(A^h \otimes_A B, I^h(A^h \otimes_A B))$ 是 Hensel 对。
由 $(B^h, IB^h)$ 的泛性质，得到映射
$B^h \to A^h \otimes_A B$。我们略去证明该映射是引理中映射的逆。
\end{proof}


\begin{lemma}
\label{more-algebra-lemma-CRT-henselization}
设 $I_1, I_2, \dots, I_n \subset A$ 为环 $A$ 的理想。假设对
$i \not = j$、$1 \leq i, j \leq n$ 有 $I_i + I_j = A$。
以 $A^h$ 表示对 $(A, I_1 \cap \ldots \cap I_n)$ 的 Hensel 化，
以 $A_i^h$ 表示对 $(A, I_i)$ 的 Hensel 化。则自然态射
$$
A^h \to \prod\nolimits_{i = 1, \ldots, n} A^h_i
$$
是同构。
\end{lemma}

\begin{proof}
该映射来自构造的函子性。若 $n = 1$，结论立即成立。若 $n > 1$，
注意，由理想所满足的条件，
$(I_1 \cap \ldots \cap I_{n - 1}) + I_n = A$。因此，只要证明
$n = 2$ 的情形，引理对所有 $n$ 即由归纳法成立。

\medskip\noindent
考虑 $n = 2$。选取 $f_i \in I_i$，使得 $f_1 + f_2 = 1$。
回顾 $A^h$ 构造为如下 $A$-代数 $B$ 的余极限：$A \to B$ 是
\'etale 的，且 $A/I_1 \cap I_2 \to B/(I_1 \cap I_2)B$ 是同构。
见引理 \ref{more-algebra-lemma-henselization} 的证明。以 $Hens$ 表示这样的
$A$-代数 $B$ 所成的范畴。类似地，以 $Hens_i$、$i = 1, 2$ 表示如下
\'etale $A$-代数 $B$ 的范畴：$A/I_i \to B/I_iB$ 是同构。

\medskip\noindent
考虑由满足 $f_2$ 可逆的那些 $B$ 组成的子范畴
$Hens_1(f_2) \subset Hens_1$。这是余终子范畴，因为若 $B$ 属于
$Hens_1$，则 $B_{f_2}$ 也属于 $Hens_1$
（提示：$B_{f_2}/I_1B_{f_2} = (B/I_1B)_{f_2} = B/I_1B$，因为
$f_2$ 模 $I_1$ 是单位元）。类似地，子范畴
$Hens_2(f_1) \subset Hens_2$ 是余终的。因此，为计算 $A_1^h$ 与
$A_2^h$，可以使用这些较小的范畴，见 Categories, Lemma
\ref{categories-lemma-cofinal}。

\medskip\noindent
考虑函子
$$
Hens_1(f_2) \times Hens_2(f_1) \to Hens,\quad (B_1, B_2) \mapsto B_1 \times B_2
$$
该函子有定义，因为
\begin{align*}
(B_1 \times B_2)/(I_1 \cap I_2)(B_1 \times B_2)
& =
B_1/(I_1 \cap I_2)B_1 \times B_2/(I_1 \cap I_2)B_2 \\
& =
B_1/I_1B_1 \times B_2/I_2B_2 \\
& = A/I_1 \times A/I_2 \\
& = A/(I_1 \cap I_2)
\end{align*}
第二个等号成立，是因为 $f_2 \in I_2$ 在 $B_1$ 中可逆，从而
$(I_1 \cap I_2)B_1 = I_1B_1 \cap I_2B_1 = I_1B_1$
（这里还使用了 Algebra, Lemma
\ref{algebra-lemma-flat-intersect-ideals}）。最后一个等号由
Algebra, Lemma \ref{algebra-lemma-chinese-remainder} 得出。
要完成证明，只需证明所显示的函子余终（见上面引用的引理）。
Categories, Definition \ref{categories-definition-cofinal} 的第一个条件成立，
因为 $Hens$ 的对象 $B$ 映到 $B_{f_2} \times B_{f_1}$，而后者在该函子的
像中。第二个条件总是成立，因为若有
$B \to B_1 \times B_2$ 和 $B \to B'_1 \times B'_2$，其中
$B_1, B'_1 \in Hens_1(f_2)$ 且 $B_2, B'_2 \in Hens_2(f_1)$，则可令
$B''_1 = B_1 \otimes_B B'_1$，它属于 $Hens_1(f_2)$；类似地，令
$B''_2 = B_2 \otimes_B B''_2$，它属于 $Hens_2(f_1)$，从而找到适配于
下列交换图的映射 $B \to B''_1 \times B''_2$：
$$
\xymatrix{
& B \ar[ld] \ar[d] \ar[rd] \\
B_1 \times B_2 \ar[r] & B_1'' \times B_2'' & B_1' \times B_2' \ar[l]
}
$$
这就完成了证明。
\end{proof}








\section{提升与 Hensel 对}
\label{more-algebra-section-lifting-henselian-pairs}

\noindent
本节主要结合第 \ref{more-algebra-section-lifting} 节和第
\ref{more-algebra-section-henselian-pairs} 节的结果。

\begin{lemma}
\label{more-algebra-lemma-lift-finite-projective-module}
设 $(R, I)$ 为 Hensel 对。映射
$$
P \longrightarrow P/IP
$$
诱导有限投射 $R$-模的同构类集合与有限投射 $R/I$-模的同构类集合
之间的双射。特别地，任意有限投射 $R/I$-模都同构于某个有限投射
$R$-模 $P$ 的 $P/IP$。
\end{lemma}

\begin{proof}
先证明最后一个断言。令 $\overline{P}$ 为有限投射 $R/I$-模。
由引理 \ref{more-algebra-lemma-lift-projective-module}，可以找到某个在 $R$ 上
\'etale 的 $R'$ 上的有限投射模 $P'$，其中 $R/I = R'/IR'$，且
$P'/IP'$ 同构于 $\overline{P}$。由于 $(R, I)$ 是 Hensel 对，
\'etale 环映射 $R \to R'$ 有截面 $\tau : R' \to R$
（引理 \ref{more-algebra-lemma-characterize-henselian-pair}）。令
$P = P' \otimes_{R', \tau} R$，便推出 $P/IP$ 同构于
$\overline{P}$。当然，这说明引理陈述中的映射是满的。

\medskip\noindent
单射性。假设 $P_1$ 与 $P_2$ 是有限投射 $R$-模，并且作为
$R/I$-模有 $P_1/IP_1 \cong P_2/IP_2$。由于 $P_1$ 投射，可找到提升
该同构的 $R$-模映射 $u : P_1 \to P_2$。由 Nakayama 引理
（Algebra, Lemma \ref{algebra-lemma-NAK}），$u$ 是满射。
类似地，可找到满射 $v : P_2 \to P_1$。由 Algebra, Lemma
\ref{algebra-lemma-fun}，映射 $v \circ u$ 是同构，因此 $u$ 是同构。
\end{proof}

\begin{lemma}
\label{more-algebra-lemma-finite-etale-equivalence}
设 $(A, I)$ 为 Hensel 对。函子 $B \to B/IB$ 确定有限 \'etale
$A$-代数与有限 \'etale $A/I$-代数之间的等价。
\end{lemma}

\begin{proof}
设 $B, B'$ 为两个在 $A$ 上有限 \'etale 的 $A$-代数。
则 $B' \to B'' = B \otimes_A B'$ 也有限 \'etale
（Algebra, Lemmas \ref{algebra-lemma-etale} and
\ref{algebra-lemma-base-change-integral}）。现在，下列对象之间存在
$1$ 对 $1$ 的对应：
\begin{enumerate}
\item $A$-代数映射 $B \to B'$；
\item $B' \to B''$ 的截面；
\item $B''$ 的幂等元 $e$，使得 $B' \to B'' \to eB''$ 是同构。
\end{enumerate}
(2) 与 (3) 之间的双射把 $\sigma : B'' \to B'$ 映到 $e$，其中
$(1 - e)$ 是生成 $\sigma$ 之核的幂等元；它的存在性由
Algebra, Lemmas \ref{algebra-lemma-map-between-etale} and
\ref{algebra-lemma-surjective-flat-finitely-presented} 给出。
类似地，$A/I$-代数映射 $B/IB \to B'/IB'$ 与 $B''/IB''$ 的幂等元
$\overline{e}$ 之间也有对应，其中
$B'/IB' \to B''/IB'' \to \overline{e}(B''/IB'')$ 是同构。
然而，$B''/IB''$ 的每个幂等元 $\overline{e}$ 都唯一提升为
$B''$ 的幂等元 $e$（引理 \ref{more-algebra-lemma-characterize-henselian-pair}）。
此外，若 $B'/IB' \to \overline{e}(B''/IB'')$ 是同构，则由 Nakayama
引理（Algebra, Lemma \ref{algebra-lemma-NAK}），
$B' \to eB''$ 也是同构。由此可见，该函子是全忠实的。

\medskip\noindent
本质满性。设 $A/I \to C$ 为有限 \'etale 映射。由 Algebra, Lemma
\ref{algebra-lemma-lift-etale}，存在 \'etale 映射 $A \to B$，使得
$B/IB \cong C$。令 $B'$ 为 $A$ 在 $B$ 中的整闭包。由引理
\ref{more-algebra-lemma-helper-finite-type}，对某个环 $C'$ 有
$B'/IB' = C \times C'$，并且对某个映到
$(1, 0) \in C \times C'$ 的 $g \in B'$，有 $B'_g \cong B_g$。
由于幂等元可以提升（引理 \ref{more-algebra-lemma-characterize-henselian-pair}），
得到 $B' = B'_1 \times B'_2$，其中 $C = B'_1/IB'_1$ 且
$C' = B'_2/IB'_2$。$g$ 在 $B'_1$ 中的像可逆。于是
$B_g = B'_g = B'_1 \times (B_2)_g$，这蕴含 $A \to B'_1$ 是
\'etale 的。我们推出 $B'_1$ 在 $A$ 上有限 \'etale
（例如，由 Algebra, Lemma
\ref{algebra-lemma-characterize-finite-in-terms-of-integral}，整且
\'etale 蕴含有限 \'etale），证明完成。
\end{proof}

\begin{lemma}
\label{more-algebra-lemma-lift-smooth-henselian}
设 $R$ 为环，$S$ 为光滑 $R$-代数。假设 $A$ 是 $R$-代数，且
$(A,I)$ 是 Hensel 对。则任意 $R$-代数映射 $S \to A/I$ 都可提升为
$R$-代数映射 $S \to A$。
\end{lemma}

\begin{proof}
设 $\tau : S \to A/I$ 为 $R$-代数映射。注意，由 Algebra, Lemma
\ref{algebra-lemma-base-change-smooth}，$S \otimes_R A$ 是光滑
$A$-代数。因此，由引理 \ref{more-algebra-lemma-lift-section-smooth-morphism}，
可将诱导映射 $S \otimes_R A \to A/I$ 提升为 $A$-代数同态
$S \otimes_R A \to A'$，其中 $A \to A'$ 是 \'etale 的，并诱导同构
$A/I \to A'/IA'$。由于 $(A, I)$ 是 Hensel 对，存在 $A$-代数映射
$A' \to A$，见引理 \ref{more-algebra-lemma-characterize-henselian-pair}。
复合 $S \to S \otimes_R A \to A' \to A$ 即为所需提升。
\end{proof}

\begin{lemma}
\label{more-algebra-lemma-lim-finite-projective-gives-finite-projective}
设 $A = \lim A_n$ 是环的逆向系统 $(A_n)$ 的极限。假设给定
$A_n$-模 $M_n$ 以及 $A_{n + 1}$-模映射 $M_{n + 1} \to M_n$。假设：
\begin{enumerate}
\item 转移映射 $A_{n + 1} \to A_n$ 是满射且核局部幂零；
\item $M_1$ 是有限投射 $A_1$-模；
\item $M_n$ 是有限平坦 $A_n$-模；
\item 这些映射诱导同构
$M_{n + 1} \otimes_{A_{n + 1}} A_n \to M_n$。
\end{enumerate}
则 $M = \lim M_n$ 是有限投射 $A$-模，并且对所有 $n$，
$M \otimes_A A_n \to M_n$ 都是同构。
\end{lemma}

\begin{proof}
由引理 \ref{more-algebra-lemma-limit-henselian}，对
$(A, \Ker(A \to A_1))$ 是 Hensel 对。由引理
\ref{more-algebra-lemma-lift-finite-projective-module}，可选取有限投射 $A$-模 $P$
以及同构 $P \otimes_A A_1 \to M_1$。由于 $P$ 投射，可逐次将
$A$-模映射 $P \to M_1$ 提升为 $A$-模映射
$P \to M_2$、$P \to M_3$，依此类推。由此得到映射
$$
P \longrightarrow M
$$
由于 $P$ 有限投射，对某个 $m \geq 0$ 和某个 $A$-模 $Q$，可写成
$A^{\oplus m} = P \oplus Q$。由于 $A = \lim A_n$，得到
$P = \lim P \otimes_A A_n$。因此，为证明所显示的 $A$-模映射是同构，
只需证明映射 $P \otimes_A A_n \to M_n$ 是同构。由引理
\ref{more-algebra-lemma-lift-projective} 可知 $M_n$ 是有限投射模。由引理
\ref{more-algebra-lemma-isomorphic-finite-projective-lifts}，映射
$P \otimes_A A_n \to M_n$ 是同构。
\end{proof}






\section{绝对整闭包}
\label{more-algebra-section-absolute-integral-closure}

\noindent
下面给出我们的定义。

\begin{definition}
\label{more-algebra-definition-absolutely-integrally-closed}
若每个首一多项式 $f \in A[T]$ 都是一次因子的乘积，则称环 $A$
{\it 绝对整闭}。
\end{definition}

\noindent
请注意：$f$ 可能能够以许多不同方式写成一次因子的乘积。

\begin{lemma}
\label{more-algebra-lemma-absolutely-integrally-closed}
设 $A$ 为环。下列条件等价：
\begin{enumerate}
\item $A$ 绝对整闭；
\item 任意首一多项式 $f \in A[T]$ 在 $A$ 中都有根。
\end{enumerate}
\end{lemma}

\begin{proof}
略。
\end{proof}


\begin{lemma}
\label{more-algebra-lemma-absolutely-integrally-closed-quotient-localization}
设 $A$ 绝对整闭。
\begin{enumerate}
\item $A$ 的任意商环 $A/I$ 都绝对整闭。
\item 任意局部化 $S^{-1}A$ 都绝对整闭。
\end{enumerate}
\end{lemma}

\begin{proof}
略。
\end{proof}

\begin{lemma}
\label{more-algebra-lemma-integrally-closed-in-absolutely-integrally-closed}
设 $A$ 为环。设 $S \subset A$ 为由非零因子组成的乘法子集。
若 $S^{-1}A$ 绝对整闭，且 $A \subset S^{-1}A$ 在 $S^{-1}A$ 中整闭，
则 $A$ 绝对整闭。
\end{lemma}

\begin{proof}
略。
\end{proof}

\begin{lemma}
\label{more-algebra-lemma-normal-domain-absolutely-integrally-closed}
设 $A$ 为正规整环。则 $A$ 绝对整闭当且仅当其分式域代数闭。
\end{lemma}

\begin{proof}
注意，一个域代数闭当且仅当它作为环绝对整闭。因此该引理由引理
\ref{more-algebra-lemma-absolutely-integrally-closed-quotient-localization} 和
\ref{more-algebra-lemma-integrally-closed-in-absolutely-integrally-closed} 得出。
\end{proof}

\begin{lemma}
\label{more-algebra-lemma-construct-absolute-integral-closure}
对任意环 $A$，存在扩张 $A \subset B$，使得：
\begin{enumerate}
\item $B$ 是有限自由 $A$-代数的滤过余极限；
\item $B$ 作为 $A$-模自由；
\item $B$ 绝对整闭。
\end{enumerate}
\end{lemma}

\begin{proof}
令 $I$ 为 $A$ 上所有首一多项式所成的集合。对 $i \in I$，以
$x_i$ 表示一个变量，以 $P_i$ 表示变量 $x_i$ 中相应的首一多项式。
然后令
$$
F(A) = A[x_i; i \in I]/(P_i; i \in I)
$$
正如记号所表明的，$F$ 是从环范畴到自身的函子。注意
$A \subset F(A)$，$F(A)$ 作为 $A$-模自由，并且 $F(A)$ 是有限自由
$A$-代数的滤过余极限。然后取
$$
B = \colim F^n(A)
$$
其中转移映射是包含映射
$F^n(A) \subset F(F^n(A)) = F^{n + 1}(A)$。系数属于 $B$ 的任意
首一多项式实际上对某个 $n$ 具有属于 $F^n(A)$ 的系数，因而由构造，
它在 $F^{n + 1}(A)$ 中有解。由引理
\ref{more-algebra-lemma-absolutely-integrally-closed}，这蕴含 $B$ 绝对整闭。
其他性质的证明从略。
\end{proof}

\begin{lemma}
\label{more-algebra-lemma-absolutely-integrally-closed-strictly-henselian}
设 $A$ 绝对整闭。设 $\mathfrak p \subset A$ 为素理想。
则局部环 $A_\mathfrak p$ 是严格 Hensel 的。
\end{lemma}

\begin{proof}
由引理 \ref{more-algebra-lemma-absolutely-integrally-closed-quotient-localization}，
可假设 $A$ 是局部环，且 $\mathfrak p$ 是其极大理想。再由引理
\ref{more-algebra-lemma-absolutely-integrally-closed-quotient-localization}，剩余域代数闭。
每个首一多项式都完全分解为一次因子，故可直接应用
Algebra, Definition \ref{algebra-definition-henselian}。
\end{proof}

\begin{lemma}
\label{more-algebra-lemma-absolutely-integrally-closed-henselian-pair}
设 $A$ 绝对整闭。设 $I \subset A$ 为理想。若（且仅若）下列条件成立，
则 $(A, I)$ 是 Hensel 对：
\begin{enumerate}
\item $I$ 包含于 $A$ 的 Jacobson 根中；
\item $A \to A/I$ 在幂等元上诱导双射。
\end{enumerate}
\end{lemma}

\begin{proof}
设 $f \in A[T]$ 为首一多项式，并设
$f \bmod I = g_0 h_0$ 是 $A/I$ 上的一个分解，其中 $g_0$、$h_0$ 首一，
且 $g_0$ 与 $h_0$ 在 $A/I[T]$ 中生成单位理想。这意味着
$$
A/I[T]/(f) = A/I[T]/(g_0) \times A/I[T]/(h_0)
$$
以 $e \in A/I[T]/(f)$ 表示对应于右侧环中幂等元 $(1, 0)$ 的元素。
写成 $f = (T - a_1) \ldots (T - a_d)$，其中 $a_i \in A$。
对每个 $i \in \{1, \ldots, d\}$，得到 $A$-代数映射
$\varphi_i : A[T]/(f) \to A$、$T \mapsto a_i$，它诱导类似的
$A/I$-代数映射 $\overline{\varphi}_i : A/I[T]/(f) \to A/I$。
令 $e_i = \overline{\varphi}_i(e) \in A/I$。这些都是幂等元。
由假设 (2)，可将 $e_i$ 提升为 $A$ 中的幂等元。这意味着可将
$A = \prod A_j$ 写成环的有限乘积，使得在 $A_j/IA_j$ 中，每个
$e_i$ 都是 $0$ 或 $1$。略去一些细节。注意，$A_j$ 作为 $A$ 的
因子环绝对整闭。只需把 $f$ 在 $A_j/IA_j$ 上的分解提升到 $A_j$。
这将我们归结到下一段讨论的情形。

\medskip\noindent
假设对 $i = 1, \ldots, r$ 有 $e_i = 1$，而对
$i = r + 1, \ldots, d$ 有 $e_i = 0$。由
$(g_0, h_0) = A/I[T]$，存在 $k_0, l_0 \in A/I[T]$，使得
$g_0 k_0 + h_0 l_0 = 1$。可见 $e = h_0 l_0$ 且
$e_i = h_0(a_i) l_0(a_i)$。于是对 $i = 1, \dots ,r$，
$h_0(a_i)$ 是单位元。由于 $f(a_i) = 0$，有
$0 = h_0(a_i)g_0(a_i)$，从而对 $i = 1, \ldots, r$ 有
$g_0(a_i) = 0$。因此 $(T - a_1)$ 在 $A/I[T]$ 中整除 $g_0$；设
$g_0 = (T - a_1) g_0'$。令
$f' = (T - a_2) \ldots (T - a_d)$ 且 $h'_0 = h_0$。
对 $d$ 归纳，可将分解 $f' \bmod I = g'_0 h'_0$ 提升为 $A$ 上的分解
$f' = g' h'$；它给出分解 $f = (T - a_1) g' h'$，提升给定分解
$f \bmod I = g_0 h_0$，正合所需。
\end{proof}








\section{自伴随环}
\label{more-algebra-section-auto-ass}

\noindent
部分材料见 \cite{Autour}。

\begin{definition}
\label{more-algebra-definition-auto-ass}
若 $R$ 是局部环，且其极大理想 $\mathfrak m$ 弱伴随于 $R$，则称环
$R$ 是{\it 自伴随的}。
\end{definition}

\begin{lemma}
\label{more-algebra-lemma-auto-ass-implies-P}
自伴随环 $R$ 具有下列性质：(P) 每个真有限生成理想
$I \subset R$ 都有非零零化子。
\end{lemma}

\begin{proof}
由假设，存在非零元素 $x \in R$，使得对每个
$f \in \mathfrak m$ 都有 $f^n x = 0$。设 $I = (f_1, \ldots, f_r)$。
则 $x$ 属于 $R \to \bigoplus R_{f_i}$ 的核。因此，由 Algebra, Lemma
\ref{algebra-lemma-when-injective-covering}，存在非零 $y \in R$，使得对
所有 $i$ 都有 $f_i y = 0$。由于 $y \in \text{Ann}_R(I)$，证毕。
\end{proof}

\begin{lemma}
\label{more-algebra-lemma-P-universally-injective}
设 $R$ 为具有引理 \ref{more-algebra-lemma-auto-ass-implies-P} 中性质 (P) 的环。
设 $u : N \to M$ 为投射 $R$-模的同态。则 $u$ 普遍单射当且仅当
$u$ 是单射。
\end{lemma}

\begin{proof}
假设 $u$ 是单射。我们的目标是证明 $u$ 普遍单射。先选取模 $Q$，
使得 $N \oplus Q$ 自由。考察映射
$N \oplus Q \to M \oplus Q$ 可知，只需在 $N$ 自由时证明该引理。
在此情形中，$N$ 是有限自由 $R$-模的有向余极限。因此归结到
$N$ 为有限自由 $R$-模的情形，设 $N = R^{\oplus n}$。对 $n$ 归纳
证明引理。$n = 0$ 的情形是平凡的。

\medskip\noindent
设 $u : R^{\oplus n} \to M$ 为单射模映射，其中 $M$ 投射。
选取 $R$-模 $Q$，使得 $M \oplus Q$ 自由。以复合
$R^{\oplus n} \to M \to M \oplus Q$ 替换 $u$ 后，可假设 $M$ 自由。
于是可找到直和因子 $R^{\oplus m} \subset M$，使得
$u(R^{\oplus n}) \subset R^{\oplus m}$。故可假设
$M = R^{\oplus m}$。此时 $u$ 由矩阵 $A = (a_{ij})$ 给出，满足
$u(x_1, \ldots, x_n) = (\sum x_i a_{i1}, \ldots, \sum x_i a_{im})$。
由于 $u$ 是单射，特别地，若 $x \not = 0$，则
$u(x, 0, \ldots, 0) = (xa_{11}, xa_{12}, \ldots, xa_{1m}) \not = 0$；
又因 $R$ 具有性质 (P)，可知
$a_{11}R + a_{12}R + \ldots + a_{1m}R = R$。由此可见，
$R(a_{11}, \ldots, a_{1m}) \subset R^{\oplus m}$ 是
$R^{\oplus m}$ 的直和因子；特别地，
$R^{\oplus m}/R(a_{11}, \ldots, a_{1m})$ 是投射 $R$-模。
得到交换图
$$
\xymatrix{
0 \ar[r] &
R \ar[rr] \ar[d]^1 & & R^{\oplus n} \ar[r] \ar[d]^u &
R^{\oplus n - 1} \ar[r] \ar[d] & 0 \\
0 \ar[r] & R \ar[rr]^{(a_{11}, \ldots, a_{1m})} & &
R^{\oplus m} \ar[r] & R^{\oplus m}/R(a_{11}, \ldots, a_{1m}) \ar[r] & 0
}
$$
其两行均分裂正合。因此右侧竖直箭头是单射；应用归纳假设，得到
右侧竖直箭头普遍单射。由此中间的竖直箭头普遍单射。
\end{proof}

\begin{lemma}
\label{more-algebra-lemma-P-fPD-zero}
设 $R$ 为环。下列条件等价：
\begin{enumerate}
\item $R$ 具有引理 \ref{more-algebra-lemma-auto-ass-implies-P} 中的性质 (P)；
\item 投射 $R$-模的任意单射都是普遍单射；
\item 若 $u : N \to M$ 是单射，且 $N$、$M$ 是有限投射 $R$-模，
则 $\Coker(u)$ 是有限投射 $R$-模；
\item 若 $N \subset M$，且 $N$、$M$ 作为 $R$-模有限投射，
则 $N$ 是 $M$ 的直和因子；
\item 任意单射 $R \to R^{\oplus n}$ 都是分裂单射。
\end{enumerate}
\end{lemma}

\begin{proof}
(1) $\Rightarrow$ (2) 是引理 \ref{more-algebra-lemma-P-universally-injective}。
显然，(3) 与 (4) 等价。由 Algebra, Lemma
\ref{algebra-lemma-universally-exact-split}，有
(2) $\Rightarrow$ (3)、(4)。第 (5) 项是 (4) 的特殊情形。
假设 (5)。设 $I = (a_1, \ldots, a_n)$ 为 $R$ 的真有限生成理想。
由于 $I \not = R$，映射 $R \to R^{\oplus n}$、
$x \mapsto (xa_1, \ldots, xa_n)$ 不是分裂单射。因此它有非零核，
从而 $\text{Ann}_R(I) \not = 0$。故 (1) 成立。
\end{proof}

\begin{example}
\label{more-algebra-example-auto-ass-weird-flat}
即使引理 \ref{more-algebra-lemma-P-fPD-zero} 中的等价条件成立，也不总是每个自由
$R$-模之间的单射都是分裂单射。例如，设
$R = k[x_1, x_2, x_3, \ldots]/(x_i^2)$。这是一个自伴随环。考虑自由
$R$-模之间的映射
$$
u :
\bigoplus\nolimits_{i \geq 1} Re_i
\longrightarrow
\bigoplus\nolimits_{i \geq 1} Rf_i, \quad
e_i \longmapsto f_i - x_if_{i + 1}.
$$
对任意整数 $n$，$u$ 在 $\bigoplus_{i = 1, \ldots, n} Re_i$ 上的限制
是单射，因为像 $u(e_1), \ldots, u(e_n)$ 在 $R$ 上线性无关。
因此 $u$ 是单射，从而由该引理它普遍单射。由于
$u \otimes \text{id}_k$ 是双射，可知若 $u$ 是分裂单射，则 $u$ 将为
满射。但 $u$ 不是满射，因为 $f_1$ 的逆像将是元素
$$
\sum\nolimits_{i \geq 0} x_1 \ldots x_ie_{i + 1} =
e_1 + x_1e_2 + x_1x_2e_3 + \ldots
$$
而它不属于该直和。顺便指出，$\Coker(u)$ 是可数生成的平坦
$R$-模（因为 $u$ 普遍单射），但它不是投射模（因为 $u$ 不分裂），
因而不是 Mittag-Leffler 模（见 Algebra, Lemma
\ref{algebra-lemma-countgen-projective}）。
\end{example}

\noindent
在局部环为 Noether 环的情形，下列引理是 Algebra, Proposition
\ref{algebra-proposition-what-exact} 的特殊情形。

\begin{lemma}
\label{more-algebra-lemma-exact-length-1}
设 $(R, \mathfrak m)$ 为局部环。假设
$\varphi : R^m \to R^n$ 是有限自由模之间的映射。下列条件等价：
\begin{enumerate}
\item $\varphi$ 是单射；
\item $\varphi$ 的秩为 $m$，且 $I(\varphi)$ 在 $R$ 中的零化子为零。
\end{enumerate}
若 $R$ 为 Noether 环，则它们还等价于：
\begin{enumerate}
\item[(3)] $\varphi$ 的秩为 $m$，且 $I(\varphi) = R$，或它包含一个非零因子。
\end{enumerate}
这里 $\varphi$ 的秩与 $I(\varphi)$ 如 Algebra, Definition
\ref{algebra-definition-rank} 中所定义。
\end{lemma}

\begin{proof}
若 $\varphi$ 的某个矩阵系数不属于 $\mathfrak m$，则应用
Algebra, Lemma \ref{algebra-lemma-add-trivial-complex}，把 $\varphi$
写成 $1 : R \to R$ 与映射
$\varphi' : R^{m-1} \to R^{n-1}$ 的和。容易看出，对 $\varphi'$
成立的引理蕴含对 $\varphi$ 成立的引理。因此从一开始即可假设
$\varphi$ 的所有矩阵系数都属于 $\mathfrak m$。

\medskip\noindent
假设 $\varphi$ 是单射。可假设 $m > 0$。设
$\mathfrak q \in \text{WeakAss}(R)$，从而 $R_\mathfrak q$ 是
自伴随环。于是 $\varphi$ 诱导单射
$R_\mathfrak q^m \to R_\mathfrak q^n$；由引理
\ref{more-algebra-lemma-auto-ass-implies-P} 和 \ref{more-algebra-lemma-P-universally-injective}，
该映射普遍单射。因此
$\varphi : \kappa(\mathfrak q)^m \to \kappa(\mathfrak q)^n$ 是单射。
故 $\varphi \bmod \mathfrak q$ 的秩为 $m$，且
$I(\varphi \otimes \kappa(\mathfrak q))$ 不是零理想。由于 $m$ 是
$\varphi$ 可能具有的最大秩，我们推出 $\varphi$ 的秩也为 $m$
（矩阵的秩在基变换后只能下降）。因此
$I(\varphi) \cdot \kappa(\mathfrak q) =
I(\varphi \otimes \kappa(\mathfrak q))$ 非零。
所以 $I(\varphi)$ 不包含于 $\mathfrak q$。于是 $R$ 的弱伴随素理想
没有一个是 $R$-模 $\text{Ann}_R I(\varphi)$ 的弱伴随素理想。
故 $\text{Ann}_R I(\varphi)$ 没有弱伴随素理想，见 Algebra, Lemma
\ref{algebra-lemma-weakly-ass}。由 Algebra, Lemma
\ref{algebra-lemma-weakly-ass-zero}，$\text{Ann}_R I(\varphi)$ 为零。

\medskip\noindent
反之，假设 (2)。秩为 $m$ 蕴含 $n \geq m$。写成
$I(\varphi) = (f_1, \ldots, f_r)$；这是可能的，因为
$I(\varphi)$ 有限生成。由 Algebra, Lemma
\ref{algebra-lemma-matrix-left-inverse}，可以找到映射
$\psi_i : R^n \to R^m$，使得
$\psi \circ \varphi = f_i \text{id}_{R^m}$。因此
$\varphi(x) = 0$ 蕴含对 $i = 1, \ldots, r$ 有 $f_i x = 0$。
这蕴含 $x = 0$，故 $\varphi$ 是单射。

\medskip\noindent
关于 Noether 局部情形中 (1) 与 (3) 的等价性，参见
Algebra, Proposition \ref{algebra-proposition-what-exact}。
若环 $R$ 是 Noether 环但非局部环，读者可由局部情形推出它；
略去细节。另一种方法是使用伴随素理想重做上述论证，并使用下列事实：
伴随素理想只有有限多个；素理想回避成立；Noether 环的非零因子恰为
不包含于任何伴随素理想的元素。
\end{proof}

\begin{lemma}
\label{more-algebra-lemma-coker-injective-free}
设 $R$ 为环。假设 $\varphi : R^n \to R^n$ 是相同秩的有限自由模之间的
单射。则 $\Hom_R(\Coker(\varphi), R) = 0$。
\end{lemma}

\begin{proof}
令 $\varphi^t : R^n \to R^n$ 为 $\varphi$ 的转置。引理断言
$\varphi^t$ 是单射。使用引理 \ref{more-algebra-lemma-exact-length-1} 的记号可知，
$\varphi^t$ 的秩为 $n$，且 $I(\varphi) = I(\varphi^t)$。
因此结论由该引理中 (1) 与 (2) 的等价性得出。
\end{proof}








\section{平坦化分层}
\label{more-algebra-section-flattening}

\noindent
设 $R \to S$ 为环映射，$M$ 为 $S$-模。对任意 $R$-代数 $R'$，
可以考虑基变换 $S' = S \otimes_R R'$ 与
$M' = M \otimes_R R'$。若模 $M'$ 在 $R'$ 上平坦，则称
$R \to R'$ {\it 使 $M$ 平坦}。我们希望理解使 $M$ 平坦的环映射
$R \to R'$ 所成集合的结构。特别地，我们希望知道是否存在
{\it $M$ 的普遍平坦化 $R \to R_{univ}$}；也就是说，是否存在
使 $M$ 平坦的环映射 $R \to R_{univ}$，并且任何使 $M$ 平坦的环映射
$R \to R'$ 都通过 $R \to R_{univ}$ 分解。事实表明，这样的普遍解
通常不存在。

\medskip\noindent
我们将在 More on Flatness, Section \ref{flat-section-flattening} 中，
在概形论设定 $\mathcal{F}/X/S$ 下讨论{\it 普遍平坦化}与
{\it 平坦化分层}。若普遍平坦化 $R \to R_{univ}$ 存在，则概形态射
$\Spec(R_{univ}) \to \Spec(R)$ 是 $\Spec(S)$ 上拟凝聚模
$\widetilde{M}$ 的普遍平坦化。

\medskip\noindent
在本节及接下来的几节中，我们证明一些与此有关的基本代数事实。
最基本的结果也许是下面这个。

\begin{lemma}
\label{more-algebra-lemma-intersection-flat}
设 $R$ 为环，$M$ 为 $R$-模，$I_1$、$I_2$ 为 $R$ 的理想。若
$M/I_1M$ 在 $R/I_1$ 上平坦，且 $M/I_2M$ 在 $R/I_2$ 上平坦，
则 $M/(I_1 \cap I_2)M$ 在 $R/(I_1 \cap I_2)$ 上平坦。
\end{lemma}

\begin{proof}
以 $R/(I_1 \cap I_2)$ 替换 $R$，并以
$M/(I_1 \cap I_2)M$ 替换 $M$，可假设 $I_1 \cap I_2 = 0$。
设 $J \subset R$ 为理想。为证明 $M$ 在 $R$ 上平坦，必须证明
$J \otimes_R M \to M$ 是单射，见 Algebra, Lemma
\ref{algebra-lemma-flat}。由 $M/I_1M$ 在 $R/I_1$ 上的平坦性，映射
$$
J/(J \cap I_1) \otimes_R M =
(J + I_1)/I_1 \otimes_{R/I_1} M/I_1M
\longrightarrow M/I_1M
$$
是单射。由于 $0 \to (J \cap I_1) \to J \to J/(J \cap I_1) \to 0$
正合，得到交换图
$$
\xymatrix{
(J \cap I_1) \otimes_R M \ar[r] \ar[d] &
J \otimes_R M \ar[r] \ar[d] &
J/(J \cap I_1) \otimes_R M \ar[r] \ar[d] & 0 \\
M \ar@{=}[r] &
M \ar[r] &
M/I_1M
}
$$
因此，只需证明 $(J \cap I_1) \otimes_R M \to M$ 是单射。
由于 $I_1 \cap I_2 = 0$，理想 $J \cap I_1$ 同构地映到
一个理想 $J' \subset R/I_2$，并且
$(J \cap I_1) \otimes_R M = J' \otimes_{R/I_2} M/I_2M$。
由 $M/I_2M$ 在 $R/I_2$ 上的平坦性，映射
$J' \otimes_{R/I_2} M/I_2M \to M/I_2M$ 是单射；这显然蕴含
$(J \cap I_1) \otimes_R M \to M$ 是单射。
\end{proof}





\section{Artin 环上的平坦化}
\label{more-algebra-section-flattening-artinian}

\noindent
当基环是 Artin 局部环时，普遍平坦化存在。对任意模，它都存在。
因此，正如稍后将看到的，当基概形是 Artin 局部环的谱时，平坦化分层
存在。

\begin{lemma}
\label{more-algebra-lemma-flattening-artinian}
设 $R$ 为 Artin 环，$M$ 为 $R$-模。则存在最小理想
$I \subset R$，使得 $M/IM$ 在 $R/I$ 上平坦。
\end{lemma}

\begin{proof}
这直接由引理 \ref{more-algebra-lemma-intersection-flat} 与 Artin 性得出。
\end{proof}

\noindent
该理想具有下列泛性质。

\begin{lemma}
\label{more-algebra-lemma-flattening-artinian-universal-property}
设 $R$ 为 Artin 环，$M$ 为 $R$-模。设 $I \subset R$ 是使得
$M/IM$ 在 $R/I$ 上平坦的最小理想 $I \subset R$。则 $I$ 具有下列
泛性质：对每个环映射 $\varphi : R \to R'$，有
$$
R' \otimes_R M\text{ is flat over }R'
\Leftrightarrow
\text{有 }\varphi(I) = 0.
$$
\end{lemma}

\begin{proof}
注意，由引理 \ref{more-algebra-lemma-flattening-artinian}，$I$ 存在。
蕴含关系 $\Rightarrow$ 由 Algebra, Lemma
\ref{algebra-lemma-flat-base-change} 得出。设 $\varphi : R \to R'$ 满足
$M \otimes_R R'$ 在 $R'$ 上平坦。令 $J = \Ker(\varphi)$。
由 Algebra, Lemma
\ref{algebra-lemma-descent-flatness-injective-map-artinian-rings}，并利用
$R' \otimes_R M = R' \otimes_{R/J} M/JM$ 在 $R'$ 上平坦，推出
$M/JM$ 在 $R/J$ 上平坦。故 $I \subset J$，正合所需。
\end{proof}











\section{基的闭子集上的平坦化}
\label{more-algebra-section-flattening-local-base}

\noindent
设 $R \to S$ 为环映射，$I \subset R$ 为理想，$M$ 为 $S$-模。
下文考虑条件
\begin{equation}
\label{more-algebra-equation-flat-at-primes-over}
\forall \mathfrak q \in V(IS) \subset \Spec(S) :
M_{\mathfrak q}\text{ is flat over }R.
\end{equation}
在几何上，这意味着 $M$ 沿 $V(I)$ 在 $\Spec(S)$ 中的逆像在 $R$
上平坦。若 $R$ 与 $S$ 是 Noether 环，且 $M$ 是有限 $S$-模，则
(\ref{more-algebra-equation-flat-at-primes-over}) 等价于：对所有 $n \geq 1$，
$M/I^nM$ 在 $R/I^n$ 上平坦；见 Algebra, Lemma
\ref{algebra-lemma-flat-module-powers}。

\begin{lemma}
\label{more-algebra-lemma-base-change-flat-at-primes-over}
设 $R \to S$ 为环映射，$I \subset R$ 为理想，$M$ 为 $S$-模。
设 $R \to R'$ 为环映射，$IR' \subset I' \subset R'$ 为理想。
若 (\ref{more-algebra-equation-flat-at-primes-over}) 对
$(R \to S, I, M)$ 成立，则 (\ref{more-algebra-equation-flat-at-primes-over}) 对
$(R' \to S \otimes_R R', I', M \otimes_R R')$ 成立。
\end{lemma}

\begin{proof}
假设 (\ref{more-algebra-equation-flat-at-primes-over}) 对
$(R \to S, I \subset R, M)$ 成立。设
$I'(S \otimes_R R') \subset \mathfrak q'$ 是 $S \otimes_R R'$ 的
素理想。令 $\mathfrak q \subset S$ 为 $S$ 中对应的素理想。
则 $IS \subset \mathfrak q$。注意，
$(M \otimes_R R')_{\mathfrak q'}$ 是基变换
$M_{\mathfrak q} \otimes_R R'$ 的局部化。因此，作为平坦模的局部化，
$(M \otimes_R R')_{\mathfrak q'}$ 在 $R'$ 上平坦；见
Algebra, Lemmas \ref{algebra-lemma-flat-base-change} and
\ref{algebra-lemma-flat-localization}。
\end{proof}

\begin{lemma}
\label{more-algebra-lemma-flat-descent-flat-at-primes-over}
设 $R \to S$ 为环映射，$I \subset R$ 为理想，$M$ 为 $S$-模。
设 $R \to R'$ 为环映射，$IR' \subset I' \subset R'$ 为理想，并且：
\begin{enumerate}
\item 由 $\Spec(R') \to \Spec(R)$ 诱导的映射
$V(I') \to V(I)$ 是满射；
\item 对所有素理想 $\mathfrak p' \in V(I')$，
$R'_{\mathfrak p'}$ 在 $R$ 上平坦。
\end{enumerate}
若 (\ref{more-algebra-equation-flat-at-primes-over}) 对
$(R' \to S \otimes_R R', I', M \otimes_R R')$ 成立，则
(\ref{more-algebra-equation-flat-at-primes-over}) 对 $(R \to S, I, M)$ 成立。
\end{lemma}

\begin{proof}
假设 (\ref{more-algebra-equation-flat-at-primes-over}) 对
$(R' \to S \otimes_R R', IR', M \otimes_R R')$ 成立。
选取素理想 $IS \subset \mathfrak q \subset S$。令
$I \subset \mathfrak p \subset R$ 为 $R$ 中对应的素理想。
由假设，存在位于 $\mathfrak p$ 之上的 $R'$ 的素理想
$\mathfrak p' \in V(I')$，并且
$R_{\mathfrak p} \to R'_{\mathfrak p'}$ 平坦。选取素理想
$\overline{\mathfrak q}' \subset
\kappa(\mathfrak q) \otimes_{\kappa(\mathfrak p)} \kappa(\mathfrak p')$；
它对应于素理想 $\mathfrak q' \subset S \otimes_R R'$，后者位于
$\mathfrak q$ 与 $\mathfrak p'$ 之上。注意，
$(S \otimes_R R')_{\mathfrak q'}$ 是
$S_{\mathfrak q} \otimes_{R_{\mathfrak p}} R'_{\mathfrak p'}$ 的局部化。
由假设，模 $(M \otimes_R R')_{\mathfrak q'}$ 在
$R'_{\mathfrak p'}$ 上平坦。因此，Algebra, Lemma
\ref{algebra-lemma-base-change-flat-up-down} 蕴含
$M_{\mathfrak q}$ 在 $R_{\mathfrak p}$ 上平坦，正是所需结论。
\end{proof}

\begin{lemma}
\label{more-algebra-lemma-limit-preserving-flat-at-primes-over}
设 $R \to S$ 为有限表示环映射，$M$ 为有限表示 $S$-模。
设 $R' = \colim_{\lambda \in \Lambda} R_\lambda$ 是 $R$-代数的有向余极限。
设 $I_\lambda \subset R_\lambda$ 为理想，使得对所有
$\mu \geq \lambda$ 有 $I_\lambda R_\mu \subset I_\mu$，并令
$I' = \colim_\lambda I_\lambda$。若
(\ref{more-algebra-equation-flat-at-primes-over}) 对
$(R' \to S \otimes_R R', I', M \otimes_R R')$ 成立，则存在
$\lambda \in \Lambda$，使得 (\ref{more-algebra-equation-flat-at-primes-over}) 对
$(R_\lambda \to S \otimes_R R_\lambda, I_\lambda, M \otimes_R R_\lambda)$
成立。
\end{lemma}

\begin{proof}
记 $S_\lambda = S \otimes_R R_\lambda$、$S' = S \otimes_R R'$、
$M_\lambda = M \otimes_R R_\lambda$ 和 $M' = M \otimes_R R'$。
基变换 $S'$ 在 $R'$ 上有限表示，$M'$ 在 $S'$ 上有限表示；带下标
$\lambda$ 的版本也同样如此，见 Algebra, Lemma
\ref{algebra-lemma-base-change-finiteness}。由 Algebra, Theorem
\ref{algebra-theorem-openness-flatness}，集合
$$
U' = \{\mathfrak q' \in \Spec(S') \mid
M'_{\mathfrak q'}\text{ is flat over }R'\}
$$
在 $\Spec(S')$ 中为开集。注意，$V(I'S')$ 是拟紧空间，且由假设它
包含于 $U'$。因此存在有限多个 $g'_j \in S'$、
$j = 1, \ldots, m$，使得 $D(g'_j) \subset U'$，并且
$V(I'S') \subset \bigcup D(g'_j)$。特别地，
$(M')_{g'_j}$ 在 $R'$ 上平坦。

\medskip\noindent
我们将依次选取越来越大的 $\lambda \in \Lambda$。首先选得足够大，
使得可找到映到 $g'_j$ 的 $g_{j, \lambda} \in S_{\lambda}$。
包含关系 $V(I'S') \subset \bigcup D(g'_j)$ 意味着
$I'S' + (g'_1, \ldots, g'_m) = S'$；它可表示为
$1 = \sum z_sh_s + \sum f_jg'_j$，其中
$z_s \in I'$、$h_s, f_j \in S'$。增大 $\lambda$ 后，可假设这样的等式
在 $S_\lambda$ 中成立。因此可假设
$V(I_\lambda S_\lambda) \subset \bigcup D(g_{j, \lambda})$。
由 Algebra, Lemma \ref{algebra-lemma-flat-finite-presentation-limit-flat}，
对某个足够大的 $\lambda$，模 $(M_\lambda)_{g_{j, \lambda}}$ 在
$R_\lambda$ 上平坦。特别地，$M_\lambda$ 在位于理想 $I_\lambda$
之上的所有素理想处都在 $R_\lambda$ 上平坦。
\end{proof}









\section{源与基的闭子集上的平坦化}
\label{more-algebra-section-flattening-local-source-base}

\noindent
本节略微推广第 \ref{more-algebra-section-flattening-local-base} 节的讨论。
强烈建议读者先阅读并理解该节。

\begin{situation}
\label{more-algebra-situation-flattening-general}
设 $R \to S$ 为环映射，$J \subset S$ 为理想，$M$ 为 $S$-模。
\end{situation}

\noindent
在此情形中，给定 $R$-代数 $R'$ 与理想 $I' \subset R'$，令
$S' = S \otimes_R R'$、$M' = M \otimes_R R'$。我们考虑条件
\begin{equation}
\label{more-algebra-equation-flat-at-primes}
\forall \mathfrak q' \in V(I'S' + JS') \subset \Spec(S') :
M'_{\mathfrak q'}\text{ is flat over }R'.
\end{equation}
在几何上，这意味着 $M'$ 在 $V(I')$ 的逆像与 $V(J)$ 的逆像之交上
相对于 $R'$ 平坦。由于 $(R \to S, J, M)$ 已固定，条件
(\ref{more-algebra-equation-flat-at-primes}) 只依赖于对 $(R', I')$，其中 $R'$ 被视为
$R$-代数。

\begin{lemma}
\label{more-algebra-lemma-base-change-flat-at-primes}
在情形 \ref{more-algebra-situation-flattening-general} 中，设
$R' \to R''$ 为 $R$-代数映射。设 $I' \subset R'$ 以及
$I'R'' \subset I'' \subset R''$ 为理想。若
(\ref{more-algebra-equation-flat-at-primes}) 对 $(R', I')$ 成立，则
(\ref{more-algebra-equation-flat-at-primes}) 对
$(R'', I'')$ 成立。
\end{lemma}

\begin{proof}
假设 (\ref{more-algebra-equation-flat-at-primes}) 对 $(R', I')$ 成立。设
$I''S'' + JS'' \subset \mathfrak q''$ 是 $S''$ 的素理想。
令 $\mathfrak q' \subset S'$ 为 $S'$ 中对应的素理想。于是
$I'S' \subset \mathfrak q'$ 与 $JS' \subset \mathfrak q'$ 都成立，
因为相应条件对 $\mathfrak q''$ 成立。注意，
$(M'')_{\mathfrak q''}$ 是基变换
$M'_{\mathfrak q'} \otimes_R R''$ 的局部化。因此，作为平坦模的局部化，
$(M'')_{\mathfrak q''}$ 在 $R''$ 上平坦；见 Algebra, Lemmas
\ref{algebra-lemma-flat-base-change} and
\ref{algebra-lemma-flat-localization}。
\end{proof}

\begin{lemma}
\label{more-algebra-lemma-flat-descent-flat-at-primes}
在情形 \ref{more-algebra-situation-flattening-general} 中，设
$R' \to R''$ 为 $R$-代数映射。设 $I' \subset R'$ 以及
$I'R'' \subset I'' \subset R''$ 为理想。假设：
\begin{enumerate}
\item 由 $\Spec(R'') \to \Spec(R')$ 诱导的映射
$V(I'') \to V(I')$ 是满射；
\item 对所有素理想 $\mathfrak p'' \in V(I'')$，
$R''_{\mathfrak p''}$ 在 $R'$ 上平坦。
\end{enumerate}
若 (\ref{more-algebra-equation-flat-at-primes}) 对 $(R'', I'')$ 成立，则
(\ref{more-algebra-equation-flat-at-primes}) 对
$(R', I')$ 成立。
\end{lemma}

\begin{proof}
假设 (\ref{more-algebra-equation-flat-at-primes}) 对 $(R'', I'')$ 成立。
选取素理想 $I'S' + JS' \subset \mathfrak q' \subset S'$。
令 $I' \subset \mathfrak p' \subset R'$ 为 $R'$ 中对应的素理想。
由假设，存在位于 $\mathfrak p'$ 之上的 $R''$ 的素理想
$\mathfrak p'' \in V(I'')$，并且
$R'_{\mathfrak p'} \to R''_{\mathfrak p''}$ 平坦。选取素理想
$\overline{\mathfrak q}'' \subset
\kappa(\mathfrak q') \otimes_{\kappa(\mathfrak p')} \kappa(\mathfrak p'')$。
这对应于素理想 $\mathfrak q'' \subset S'' = S' \otimes_{R'} R''$，
后者位于 $\mathfrak q'$ 与 $\mathfrak p''$ 之上。特别地，
$I''S'' \subset \mathfrak q''$ 且 $JS'' \subset \mathfrak q''$。
注意，$(S' \otimes_{R'} R'')_{\mathfrak q''}$ 是
$S'_{\mathfrak q'} \otimes_{R'_{\mathfrak p'}} R''_{\mathfrak p''}$ 的
局部化。由假设，模 $(M' \otimes_{R'} R'')_{\mathfrak q''}$ 在
$R''_{\mathfrak p''}$ 上平坦。因此，Algebra, Lemma
\ref{algebra-lemma-base-change-flat-up-down} 蕴含
$M'_{\mathfrak q'}$ 在 $R'_{\mathfrak p'}$ 上平坦，正是所需结论。
\end{proof}

\begin{lemma}
\label{more-algebra-lemma-limit-preserving-flat-at-primes}
在情形 \ref{more-algebra-situation-flattening-general} 中，假设 $R \to S$ 本质有限表示，
且 $M$ 是有限表示 $S$-模。设
$R' = \colim_{\lambda \in \Lambda} R_\lambda$ 是 $R$-代数的有向余极限。
设 $I_\lambda \subset R_\lambda$ 为理想，使得对所有
$\mu \geq \lambda$ 有 $I_\lambda R_\mu \subset I_\mu$，并令
$I' = \colim_\lambda I_\lambda$。若 (\ref{more-algebra-equation-flat-at-primes}) 对
$(R', I')$ 成立，则存在 $\lambda \in \Lambda$，使得
(\ref{more-algebra-equation-flat-at-primes}) 对 $(R_\lambda, I_\lambda)$ 成立。
\end{lemma}

\begin{proof}
先证明 $R \to S$ 有限表示的情形，再说明一般情形需要作何修改。
记 $S_\lambda = S \otimes_R R_\lambda$、$S' = S \otimes_R R'$、
$M_\lambda = M \otimes_R R_\lambda$ 及 $M' = M \otimes_R R'$。
基变换 $S'$ 在 $R'$ 上有限表示，$M'$ 在 $S'$ 上有限表示；带下标
$\lambda$ 的版本也同样如此，见 Algebra, Lemma
\ref{algebra-lemma-base-change-finiteness}。由 Algebra, Theorem
\ref{algebra-theorem-openness-flatness}，集合
$$
U' = \{\mathfrak q' \in \Spec(S') \mid
M'_{\mathfrak q'}\text{ is flat over }R'\}
$$
在 $\Spec(S')$ 中为开集。注意，$V(I'S' + JS')$ 是拟紧空间，且由
假设它包含于 $U'$。因此存在有限多个
$g'_j \in S'$、$j = 1, \ldots, m$，使得 $D(g'_j) \subset U'$，且
$V(I'S' + JS') \subset \bigcup D(g'_j)$。特别地，
$(M')_{g'_j}$ 在 $R'$ 上平坦。

\medskip\noindent
我们将依次选取越来越大的 $\lambda \in \Lambda$。先选得足够大，
使得可以找到映到 $g'_j$ 的 $g_{j, \lambda} \in S_{\lambda}$。
包含关系 $V(I'S' + JS') \subset \bigcup D(g'_j)$ 意味着
$I'S' + JS' + (g'_1, \ldots, g'_m) = S'$；它可表示为
$$
1 = \sum y_tk_t + \sum z_sh_s + \sum f_jg'_j
$$
其中 $z_s \in I'$、$y_t \in J$、$k_t, h_s, f_j \in S'$。
增大 $\lambda$ 后，可假设这样的等式在 $S_\lambda$ 中成立。
因此可假设
$V(I_\lambda S_\lambda + J S_\lambda) \subset \bigcup D(g_{j, \lambda})$。
由 Algebra, Lemma \ref{algebra-lemma-flat-finite-presentation-limit-flat}，
对某个足够大的 $\lambda$，模 $(M_\lambda)_{g_{j, \lambda}}$ 在
$R_\lambda$ 上平坦。特别地，$M_\lambda$ 在对应于
$V(I_\lambda S_\lambda + J S_\lambda)$ 中各点的所有素理想处，都在
$R_\lambda$ 上平坦。

\medskip\noindent
若 $S$ 本质有限表示，则可写成 $S = \Sigma^{-1}C$，其中
$R \to C$ 有限表示，$\Sigma \subset C$ 是乘法子集。还可写成
$M = \Sigma^{-1}N$，其中 $N$ 是有限表示 $C$-模，见 Algebra, Lemma
\ref{algebra-lemma-construct-fp-module}。此时引入
$C_\lambda$、$C'$、$N_\lambda$、$N'$。在上述讨论中，我们得到开集
$U' \subset \Spec(C')$，在其上 $N'$ 相对于 $R'$ 平坦。条件
(\ref{more-algebra-equation-flat-at-primes}) 成立意味着 $V(I'S' + JS')$ 映入
$U'$，因为对素理想 $\mathfrak q' \subset S'$ 及其对应的素理想
$\mathfrak r' \subset C'$，有 $M'_{\mathfrak q'} = N'_{\mathfrak r'}$。
因此可找到 $g'_j \in C'$，使得 $\bigcup D(g'_j)$ 包含
$V(I'S' + JS')$ 的像。证明的其余部分与前面完全相同。
\end{proof}

\begin{lemma}
\label{more-algebra-lemma-flat-module-powers-variant}
在情形 \ref{more-algebra-situation-flattening-general} 中，设 $I \subset R$ 为理想。
假设：
\begin{enumerate}
\item $R$ 是 Noether 环；
\item $S$ 是 Noether 环；
\item $M$ 是有限 $S$-模；
\item 对每个 $n \geq 1$ 以及任意素理想
$\mathfrak q \in V(J + IS)$，模 $(M/I^n M)_{\mathfrak q}$ 在
$R/I^n$ 上平坦。
\end{enumerate}
则 (\ref{more-algebra-equation-flat-at-primes}) 对 $(R, I)$ 成立；即对每个素理想
$\mathfrak q \in V(J + IS)$，局部化 $M_{\mathfrak q}$ 在 $R$ 上平坦。
\end{lemma}

\begin{proof}
设 $\mathfrak q \in V(J + IS)$。将 Algebra, Lemma
\ref{algebra-lemma-flat-module-powers} 应用于 $R \to S_{\mathfrak q}$ 与
$M_{\mathfrak q}$，即可推出 $M_{\mathfrak q}$ 在 $R$ 上平坦。
\end{proof}










\section{Noether 完备局部环上的平坦化}
\label{more-algebra-section-flattening-Noetherian-complete-local}

\noindent
当基为完备 Noether 局部环 $R$，且给定模是在有限型 $R$-代数 $S$ 上
有限的模时，下面三个引理给出``局部''平坦化分层存在性的纯代数证明。

\begin{lemma}
\label{more-algebra-lemma-flattening-complete-local-noetherian}
设 $R \to S$ 为环映射，$M$ 为 $S$-模。假设：
\begin{enumerate}
\item $(R, \mathfrak m)$ 是完备 Noether 局部环；
\item $S$ 是 Noether 环；
\item $M$ 在 $S$ 上有限。
\end{enumerate}
则存在理想 $I \subset \mathfrak m$，使得：
\begin{enumerate}
\item 对 $S/IS$ 中位于 $\mathfrak m$ 之上的所有素理想
$\mathfrak q$，$(M/IM)_{\mathfrak q}$ 在 $R/I$ 上平坦；
\item 若 $J \subset R$ 是理想，并且对位于 $\mathfrak m$ 之上的所有
素理想 $\mathfrak q$，$(M/JM)_{\mathfrak q}$ 在 $R/J$ 上平坦，
则 $I \subset J$。
\end{enumerate}
换言之，$I$ 是使 (\ref{more-algebra-equation-flat-at-primes-over}) 对
$(\overline{R} \to \overline{S}, \overline{\mathfrak m}, \overline{M})$
成立的 $R$ 的最小理想，其中 $\overline{R} = R/I$、
$\overline{S} = S/IS$、$\overline{\mathfrak m} = \mathfrak m/I$ 且
$\overline{M} = M/IM$。
\end{lemma}

\begin{proof}
设 $J \subset R$ 为理想。将 Algebra, Lemma
\ref{algebra-lemma-flat-module-powers} 应用于环 $R/J$ 上的模 $M/JM$。
于是可见，对 $S/JS$ 中所有素理想 $\mathfrak q$，
$(M/JM)_{\mathfrak q}$ 在 $R/J$ 上平坦，当且仅当对所有
$n \geq 1$，$M/(J + \mathfrak m^n)M$ 在
$R/(J + \mathfrak m^n)$ 上平坦。下文使用这一观察。

\medskip\noindent
对每个 $n \geq 1$，局部环 $R/\mathfrak m^n$ 是 Artin 环。
故由引理 \ref{more-algebra-lemma-flattening-artinian}，存在最小理想
$I_n \supset \mathfrak m^n$，使得 $M/I_nM$ 在 $R/I_n$ 上平坦。
显然，$I_{n + 1} + \mathfrak m^n$ 包含 $I_n$；应用引理
\ref{more-algebra-lemma-intersection-flat}，可知
$I_n = I_{n + 1} + \mathfrak m^n$。由于
$R = \lim_n\ R/\mathfrak m^n$，可知
$I = \lim_n\ I_n/\mathfrak m^n$ 是 $R$ 中的理想，并且对所有
$n \geq 1$ 有 $I_n = I + \mathfrak m^n$。由证明开头的观察，
$I$ 满足 (1) 与 (2)。略去一些细节。
\end{proof}

\begin{lemma}
\label{more-algebra-lemma-flattening-complete-local-noetherian-property-by-finite-type}
采用引理 \ref{more-algebra-lemma-flattening-complete-local-noetherian} 中的记号
$R \to S$、$M$、$I$ 及假设。考虑局部环的局部同态
$\varphi : (R, \mathfrak m) \to (R', \mathfrak m')$，并假设
$R'$ 是 Noether 环。则下列条件等价：
\begin{enumerate}
\item (\ref{more-algebra-equation-flat-at-primes-over}) 对
$(R' \to S \otimes_R R', \mathfrak m', M \otimes_R R')$ 成立；
\item $\varphi(I) = 0$。
\end{enumerate}
\end{lemma}

\begin{proof}
(2) $\Rightarrow$ (1) 由引理
\ref{more-algebra-lemma-base-change-flat-at-primes-over} 得出。设
$\varphi : R \to R'$ 如引理所述并满足 (1)。必须证明
$\varphi(I) = 0$。这等价于 $\varphi(I)R' = 0$。由 Noether 局部环
$R'$ 中的 Artin--Rees（见 Algebra, Lemma
\ref{algebra-lemma-intersect-powers-ideal-module-zero}），这等价于：对所有
$n > 0$，有
$\varphi(I)R' + (\mathfrak m')^n = (\mathfrak m')^n$。
继而，这等价于：对每个 $n$，复合
$\varphi_n : R \to R' \to R'/(\mathfrak m')^n$ 零化 $I$。
由引理 \ref{more-algebra-lemma-base-change-flat-at-primes-over}，$\varphi$ 的假设 (1)
蕴含 $\varphi_n$ 的假设 (1)。这将我们归结到 $R'$ 为 Artin 局部环的
情形。

\medskip\noindent
假设 $R'$ 为 Artin 环。令 $J = \Ker(\varphi)$。必须证明
$I \subset J$。由引理 \ref{more-algebra-lemma-flattening-complete-local-noetherian}
中 $I$ 的构造，只需证明：对 $S/JS$ 中位于 $\mathfrak m$ 之上的每个
素理想 $\mathfrak q$，$(M/JM)_{\mathfrak q}$ 在 $R/J$ 上平坦。
由于 $R'$ 是 Artin 环，条件 (1) 表示 $M \otimes_R R'$ 在 $R'$ 上
平坦。由于 $R'$ 是 Artin 环，且 $R/J \to R'$ 是单射局部环映射，
$R/J$ 也是 Artin 环。因此，由 Algebra, Lemma
\ref{algebra-lemma-descent-flatness-injective-map-artinian-rings}，
$M \otimes_R R' = M/JM \otimes_{R/J} R'$ 在 $R'$ 上的平坦性蕴含
$M/JM$ 在 $R/J$ 上平坦。这就完成了证明。
\end{proof}

\begin{lemma}
\label{more-algebra-lemma-flattening-complete-local-universal-property}
采用引理 \ref{more-algebra-lemma-flattening-complete-local-noetherian} 中的记号
$R \to S$、$M$、$I$ 及假设，并另外假设 $R \to S$ 是有限型的。
则对任意局部环的局部同态
$\varphi : (R, \mathfrak m) \to (R', \mathfrak m')$，下列条件等价：
\begin{enumerate}
\item (\ref{more-algebra-equation-flat-at-primes-over}) 对
$(R' \to S \otimes_R R', \mathfrak m', M \otimes_R R')$ 成立；
\item $\varphi(I) = 0$。
\end{enumerate}
\end{lemma}

\begin{proof}
(2) $\Rightarrow$ (1) 由引理
\ref{more-algebra-lemma-base-change-flat-at-primes-over} 得出。设
$\varphi : R \to R'$ 如引理所述并满足 (1)。由于 $R$ 是 Noether 环，
$R \to S$ 有限表示，且 $M$ 是有限表示 $S$-模。把
$R' = \colim_\lambda R_\lambda$ 写成局部 $R$-子代数
$R_\lambda \subset R'$ 的有向余极限，其中极大理想
$\mathfrak m_\lambda = R_\lambda \cap \mathfrak m'$，且每个
$R_\lambda$ 在 $R$ 上本质有限型。由引理
\ref{more-algebra-lemma-limit-preserving-flat-at-primes-over}，对某个 $\lambda$，条件
(\ref{more-algebra-equation-flat-at-primes-over}) 对
$(R_\lambda \to S \otimes_R R_\lambda, \mathfrak m_\lambda,
M \otimes_R R_\lambda)$ 成立。因此，引理
\ref{more-algebra-lemma-flattening-complete-local-noetherian-property-by-finite-type}
可应用于环映射 $R \to R_\lambda$；可知 $I$ 在 $R_\lambda$ 中映为零，
从而更在 $R'$ 中映为零。
\end{proof}





\section{沿整映射的平坦性下降}
\label{more-algebra-section-descent-flatness-integral}

\noindent
先给出几个简单引理。

\begin{lemma}
\label{more-algebra-lemma-have-one-root}
设 $R$ 为环，$P(T)$ 为系数属于 $R$ 的首一多项式。设
$\alpha \in R$ 满足 $P(\alpha) = 0$。则对某个首一多项式
$Q(T) \in R[T]$，有 $P(T) = (T - \alpha)Q(T)$。
\end{lemma}

\begin{proof}
对 $P$ 的次数归纳。若 $\deg(P) = 1$，则
$P(T) = T - \alpha$，结论成立。若 $\deg(P) > 1$，则由多项式长除法，
可对某个次数 $< \deg(P)$ 的多项式 $Q \in R[T]$ 与某个 $r \in R$
写成 $P(T) = (T - \alpha)Q(T) + r$。由假设，
$0 = P(\alpha) = (\alpha - \alpha)Q(\alpha) + r = r$，故
$r = 0$，正合所需。
\end{proof}

\begin{lemma}
\label{more-algebra-lemma-adjoin-one-root}
设 $R$ 为环，$P(T)$ 为系数属于 $R$ 的首一多项式。存在有限自由环映射
$R \to R'$，使得对某个 $\alpha \in R'$ 及某个首一多项式
$Q(T) \in R'[T]$，有 $P(T) = (T - \alpha)Q(T)$。
\end{lemma}

\begin{proof}
写成 $P(T) = T^d + a_1T^{d - 1} + \ldots + a_0$。令
$R' = R[x]/(x^d + a_1x^{d - 1} + \ldots + a_0)$。令 $\alpha$ 等于
$x$ 的同余类。显然 $P(\alpha) = 0$。因此结论由引理
\ref{more-algebra-lemma-have-one-root} 得出。
\end{proof}

\begin{lemma}
\label{more-algebra-lemma-finite-split}
设 $R \to S$ 为有限环映射。存在有限自由环扩张 $R \subset R'$，
使得 $S \otimes_R R'$ 是下列形状的环的商环：
$$
R'[T_1, \ldots, T_n]/(P_1(T_1), \ldots, P_n(T_n))
$$
其中对某些 $\alpha_{ij} \in R'$，有
$P_i(T) = \prod_{j = 1, \ldots, d_i} (T - \alpha_{ij})$。
\end{lemma}

\begin{proof}
设 $x_1, \ldots, x_n \in S$ 是 $S$ 在 $R$ 上的一组生成元。
对每个 $i$，可选取首一多项式 $P_i(T) \in R[T]$，使得
$P_i(x_i) = 0$ 在 $S$ 中成立，见 Algebra, Lemma
\ref{algebra-lemma-finite-is-integral}。设 $\deg(P_i) = d_i$。
将引理 \ref{more-algebra-lemma-adjoin-one-root} 应用 $\sum d_i$ 次，得到有限自由环
扩张 $R \subset R'$，使得每个 $P_i$ 完全分裂：
$$
P_i(T) = \prod\nolimits_{j = 1, \ldots, d_i} (T - \alpha_{ij})
$$
其中某些 $\alpha_{ik} \in R'$。令
$R'[T_1, \ldots, T_n] \to S \otimes_R R'$ 为把 $T_i$ 映到
$x_i \otimes 1$ 的 $R'$-代数映射。由于它把 $P_i(T_i)$ 映为零，
便诱导所需满射。
\end{proof}

\begin{lemma}
\label{more-algebra-lemma-split-image}
设 $R$ 为环。令 $S = R[T_1, \ldots, T_n]/J$。假设 $J$ 包含形如
$P_i(T_i)$ 的元素，其中对某些 $\alpha_{ij} \in R$，有
$P_i(T) = \prod_{j = 1, \ldots, d_i} (T - \alpha_{ij})$。
对满足 $1 \leq k_i \leq d_i$ 的
$\underline{k} = (k_1, \ldots, k_n)$，考虑环映射
$$
\Phi_{\underline{k}} : R[T_1, \ldots, T_n] \to R,
\quad
T_i \longmapsto \alpha_{ik_i}
$$
令 $J_{\underline{k}} = \Phi_{\underline{k}}(J)$。则
$\Spec(S) \to \Spec(R)$ 的像等于 $V(\bigcap J_{\underline{k}})$。
\end{lemma}

\begin{proof}
该引理不证自明。提示：
$V(\bigcap J_{\underline{k}}) = \bigcup V(J_{\underline{k}})$。
\end{proof}

\noindent
下列结果归功于 Ferrand，见 \cite{Ferrand}。

\begin{lemma}
\label{more-algebra-lemma-descent-flatness-injective-finite-Noetherian-rings}
设 $R \to S$ 为 Noether 环的有限单射同态。设 $M$ 为 $R$-模。
若 $M \otimes_R S$ 是平坦 $S$-模，则 $M$ 是平坦 $R$-模。
\end{lemma}

\begin{proof}
设 $M$ 为满足 $M \otimes_R S$ 在 $S$ 上平坦的 $R$-模。
由 Algebra, Lemma \ref{algebra-lemma-flatness-descends}，为证明 $M$
平坦，可用任意忠实平坦环扩张替换 $R$。由引理
\ref{more-algebra-lemma-finite-split}，可找到有限局部自由环扩张 $R \subset R'$，
使得对某个理想 $J \subset R'[T_1, \ldots, T_n]$，有
$S' = S \otimes_R R' = R'[T_1, \ldots, T_n]/J$；该理想包含形如
$P_i(T_i)$ 的元素，其中对某些 $\alpha_{ij} \in R'$，有
$P_i(T) = \prod_{j = 1, \ldots, d_i} (T - \alpha_{ij})$。
注意，$R'$ 是 Noether 环，且 $R' \subset S'$ 是有限环扩张。
因此，可用 $R'$ 替换 $R$，并假设 $S$ 具有引理
\ref{more-algebra-lemma-split-image} 中的表示。注意，
$\Spec(S) \to \Spec(R)$ 是满射，见 Algebra, Lemma
\ref{algebra-lemma-integral-overring-surjective}。于是利用引理
\ref{more-algebra-lemma-split-image}，推出 $I = \bigcap J_{\underline{k}}$ 是满足
$V(I) = \Spec(R)$ 的理想。这意味着 $I \subset \sqrt{(0)}$；由于
$R$ 是 Noether 环，$I$ 幂零。映射 $\Phi_{\underline{k}}$ 诱导交换图
$$
\xymatrix{
S \ar[rr] & & R/J_{\underline{k}} \\
& R \ar[lu] \ar[ru]
}
$$
由此推出 $M/J_{\underline{k}}M$ 在 $R/J_{\underline{k}}$ 上平坦。
由引理 \ref{more-algebra-lemma-intersection-flat}，$M/IM$ 在 $R/I$ 上平坦。
最后应用 Algebra, Lemma \ref{algebra-lemma-lift-flatness}，推出 $M$
在 $R$ 上平坦。
\end{proof}

\begin{lemma}
\label{more-algebra-lemma-descent-flatness-injective-integral}
设 $R \to S$ 为单射整环映射。设 $M$ 为
$R[x_1, \ldots, x_n]$ 上的有限表示模。若 $M \otimes_R S$ 在 $S$ 上
平坦，则 $M$ 在 $R$ 上平坦。
\end{lemma}

\begin{proof}
选取表示
$$
R[x_1, \ldots, x_n]^{\oplus t} \to R[x_1, \ldots, x_n]^{\oplus r} \to M \to 0.
$$
设第一个映射由 $r \times t$ 矩阵 $T = (f_{ij})$ 给出，其中
$f_{ij} \in R[x_1, \ldots, x_n]$。写成
$f_{ij} = \sum f_{ij, I} x^I$，其中 $f_{ij, I} \in R$
（多重指标记号）。考虑交换图
$$
\xymatrix{
R \ar[r] & S \\
R_\lambda \ar[u] \ar[r] & S_\lambda \ar[u]
}
$$
其中 $R_\lambda$ 是 $R$ 的有限生成 $\mathbf{Z}$-子代数，包含所有
$f_{ij, I}$；$S_\lambda$ 是 $S$ 的有限
$R_\lambda$-子代数。令 $M_\lambda$ 为有限
$R_\lambda[x_1, \ldots, x_n]$-模，它由上述同样的表示定义，使用同一
矩阵 $T$，但现在把它视为 $R_\lambda[x_1, \ldots, x_n]$ 上的矩阵。
注意，$S$ 是 $S_\lambda$ 的有向余极限（略去细节）。由
Algebra, Lemma \ref{algebra-lemma-flat-finite-presentation-limit-flat}，
对某个 $\lambda$，模
$M_\lambda \otimes_{R_\lambda} S_\lambda$ 在 $S_\lambda$ 上平坦。
由引理 \ref{more-algebra-lemma-descent-flatness-injective-finite-Noetherian-rings}，
$M_\lambda$ 在 $R_\lambda$ 上平坦。由于
$M = M_\lambda \otimes_{R_\lambda} R$，结论由 Algebra, Lemma
\ref{algebra-lemma-flat-base-change} 得出。
\end{proof}

\begin{lemma}
\label{more-algebra-lemma-descent-projective-injective-finite-Noetherian-rings}
设 $R \to S$ 为 Noether 环的有限单射同态。设 $P$ 为 $R$-模。
若 $P \otimes_R S$ 是投射 $S$-模，则 $P$ 是投射 $R$-模。
\end{lemma}

\begin{proof}
设 $P$ 为满足 $P \otimes_R S$ 在 $S$ 上投射的 $R$-模。
由 Algebra, Theorem \ref{algebra-theorem-ffdescent-projectivity}，
为证明 $P$ 投射，可用任意忠实平坦环扩张替换 $R$。由引理
\ref{more-algebra-lemma-finite-split}，可找到有限局部自由环扩张 $R \subset R'$，
使得对某个理想 $J \subset R'[T_1, \ldots, T_n]$，有
$S' = S \otimes_R R' = R'[T_1, \ldots, T_n]/J$；该理想包含形如
$P_i(T_i)$ 的元素，其中对某些 $\alpha_{ij} \in R'$，有
$P_i(T) = \prod_{j = 1, \ldots, d_i} (T - \alpha_{ij})$。
注意，$R'$ 是 Noether 环，且 $R' \subset S'$ 是有限环扩张。
因此，可用 $R'$ 替换 $R$，并假设 $S$ 具有引理
\ref{more-algebra-lemma-split-image} 中的表示。注意，
$\Spec(S) \to \Spec(R)$ 是满射，见 Algebra, Lemma
\ref{algebra-lemma-integral-overring-surjective}。于是利用引理
\ref{more-algebra-lemma-split-image}，推出 $I = \bigcap J_{\underline{k}}$ 是满足
$V(I) = \Spec(R)$ 的理想。这意味着 $I \subset \sqrt{(0)}$；由于
$R$ 是 Noether 环，$I$ 幂零。映射 $\Phi_{\underline{k}}$ 诱导交换图
$$
\xymatrix{
S \ar[rr] & & R/J_{\underline{k}} \\
& R \ar[lu] \ar[ru]
}
$$
由此推出 $P/J_{\underline{k}}P$ 在 $R/J_{\underline{k}}$ 上投射。
由 Algebra, Lemma \ref{algebra-lemma-projective-modulo-two-ideals}，
$P/IP$ 在 $R/I$ 上投射。由引理
\ref{more-algebra-lemma-descent-flatness-injective-finite-Noetherian-rings}，$P$ 在
$R$ 上平坦；故可应用 Algebra, Lemma
\ref{algebra-lemma-lift-projective}，推出 $P$ 在 $R$ 上投射。
\end{proof}








\section{无挠模}
\label{more-algebra-section-torsion-flat}

\noindent
本节讨论无挠模及其与平坦性的关系（尤其是在一维环上）。

\begin{definition}
\label{more-algebra-definition-torsion}
设 $R$ 为整环，$M$ 为 $R$-模。
\begin{enumerate}
\item 若存在非零 $f \in R$ 使得 $fx = 0$，则称元素 $x \in M$
为{\it 挠元}；
\item 若 $M$ 的唯一挠元是 $0$，则称 $M$ {\it 无挠}；
\item 若 $M$ 的每个元素都是挠元，则称 $M$ 为{\it 挠模}。
\end{enumerate}
\end{definition}

\noindent
设 $R$ 为整环，$S = R \setminus \{0\}$ 为 $R$ 中非零元素所成的
乘法集。则 $R$-模 $M$ 无挠，当且仅当 $M \to S^{-1}M$ 是单射。
换言之，当且仅当映射 $M \to M \otimes_R K$ 是单射，其中
$K = S^{-1}R$ 是 $R$ 的分式域。

\begin{lemma}
\label{more-algebra-lemma-torsion}
设 $R$ 为整环，$M$ 为 $R$-模。$M$ 的挠元集合 $M_{tors}$ 是映射
$M \to M \otimes_R K$ 的核。因此 $M_{tors}$ 是 $M$ 的
$R$-子模。商模 $M/M_{tors}$ 无挠。
\end{lemma}

\begin{proof}
见上面的讨论。
\end{proof}

\begin{lemma}
\label{more-algebra-lemma-localize-torsion}
设 $R$ 为整环，$M$ 为无挠 $R$-模。对任意乘法集
$S \subset R$，模 $S^{-1}M$ 是无挠 $S^{-1}R$-模。
\end{lemma}

\begin{proof}
略。
\end{proof}

\begin{lemma}
\label{more-algebra-lemma-flat-pullback-torsion}
设 $R \to R'$ 为整环之间的平坦同态。若 $M$ 是无挠 $R$-模，
则 $M \otimes_R R'$ 是无挠 $R'$-模。
\end{lemma}

\begin{proof}
若 $M$ 无挠，则 $M \subset M \otimes_R K$ 是单射，其中 $K$ 是
$R$ 的分式域。由于 $R'$ 在 $R$ 上平坦，映射
$M \otimes_R R' \to (M \otimes_R K) \otimes_R R'$ 是单射。
由于 $M \otimes_R K$ 同构于若干个 $K$ 的直和，只需说明
$K \otimes_R R'$ 无挠。这成立，因为它是 $R'$ 的局部化。
\end{proof}

\begin{lemma}
\label{more-algebra-lemma-extension-torsion-free}
设 $R$ 为整环。设 $0 \to M \to M' \to M'' \to 0$ 为
$R$-模的短正合列。若 $M$ 与 $M''$ 无挠，则 $M'$ 无挠。
\end{lemma}

\begin{proof}
略。
\end{proof}

\begin{lemma}
\label{more-algebra-lemma-check-torsion}
设 $R$ 为整环，$M$ 为 $R$-模。则 $M$ 无挠，当且仅当对 $R$ 的所有
极大理想 $\mathfrak m$，$M_\mathfrak m$ 都是无挠
$R_\mathfrak m$-模。
\end{lemma}

\begin{proof}
略。提示：使用引理 \ref{more-algebra-lemma-localize-torsion} 与 Algebra, Lemma
\ref{algebra-lemma-characterize-zero-local}。
\end{proof}

\begin{lemma}
\label{more-algebra-lemma-finite-torsion-free-submodule-free}
设 $R$ 为整环，$M$ 为有限 $R$-模。则 $M$ 无挠，当且仅当 $M$ 是某个
有限自由模的子模。
\end{lemma}

\begin{proof}
若 $M$ 是 $R^{\oplus n}$ 的子模，则 $M$ 无挠。反之，假设 $M$ 无挠。
令 $K$ 为 $R$ 的分式域。则 $M \otimes_R K$ 是有限维 $K$-向量空间。
选取该向量空间的一组基 $e_1, \ldots, e_r$。令
$x_1, \ldots, x_n$ 为 $M$ 的生成元。对某些
$a_{ij}, b_{ij} \in R$ 且 $b_{ij} \not = 0$，写成
$x_i = \sum (a_{ij}/b_{ij}) e_j$。令 $b = \prod_{i, j} b_{ij}$。
由于 $M$ 无挠，映射 $M \to M \otimes_R K$ 是单射，并且其像包含于
$R^{\oplus r} = R e_1/b \oplus \ldots \oplus Re_r/b$。
\end{proof}

\begin{lemma}
\label{more-algebra-lemma-torsion-free-finite-noetherian-domain}
设 $R$ 为 Noether 整环，$M$ 为非零有限 $R$-模。下列条件等价：
\begin{enumerate}
\item $M$ 无挠；
\item $M$ 是有限自由模的子模；
\item $(0)$ 是 $M$ 的唯一伴随素理想；
\item $(0)$ 属于 $M$ 的支撑，且 $M$ 具有性质 $(S_1)$；
\item $(0)$ 属于 $M$ 的支撑，且 $M$ 没有嵌入伴随素理想。
\end{enumerate}
\end{lemma}

\begin{proof}
我们已在引理 \ref{more-algebra-lemma-finite-torsion-free-submodule-free} 中看到 (1)
与 (2) 等价，并在 Algebra, Lemma
\ref{algebra-lemma-criterion-no-embedded-primes} 中看到 (4) 与 (5) 等价。
(3) 与 (5) 的等价性由定义立即得出。无挠模的局部化无挠
（引理 \ref{more-algebra-lemma-localize-torsion}），因此显然 $M$ 没有不同于
$(0)$ 的伴随素理想。故 (1) 蕴含 (5)。反之，假设 (5)。
若 $M$ 有挠，则对 $R$ 的某个非零理想 $I$，存在嵌入
$R/I \subset M$。因此 $M$ 有不同于 $(0)$ 的伴随素理想
（见 Algebra, Lemmas \ref{algebra-lemma-ass} and
\ref{algebra-lemma-ass-zero}）。这是嵌入伴随素理想，与假设矛盾。
\end{proof}

\begin{lemma}
\label{more-algebra-lemma-flat-torsion-free}
设 $R$ 为整环。任意平坦 $R$-模都无挠。
\end{lemma}

\begin{proof}
若 $x \in R$ 非零，则 $x : R \to R$ 是单射；因此，若 $M$ 在 $R$ 上
平坦，则 $x : M \to M$ 是单射。故若 $M$ 在 $R$ 上平坦，则 $M$ 无挠。
\end{proof}

\begin{lemma}
\label{more-algebra-lemma-valuation-ring-torsion-free-flat}
设 $A$ 为赋值环。$A$-模 $M$ 在 $A$ 上平坦当且仅当 $M$ 无挠。
\end{lemma}

\begin{proof}
蕴含关系``平坦 $\Rightarrow$ 无挠''是引理
\ref{more-algebra-lemma-flat-torsion-free}。反之，假设 $M$ 无挠。由平坦性的方程判别
准则（见 Algebra, Lemma \ref{algebra-lemma-flat-eq}），必须证明 $M$
中的每个关系都是平凡的。为此，假设
$\sum_{i = 1, \ldots, n} a_i x_i = 0$，其中
$x_i \in M$ 且 $a_i \in A$。重新编号后，可假设对所有 $i$ 都有
$v(a_1) \leq v(a_i)$。故对某个 $a'_i \in A$，可写成
$a_i = a'_i a_1$。注意 $a'_1 = 1$。由于 $M$ 无挠，有
$x_1 = - \sum_{i \geq 2} a'_i x_i$。因此，若选取
$y_i = x_i$、$i = 2, \ldots, n$，则
$$
x_1 = \sum\nolimits_{j \geq 2} -a'_j y_j, \quad
x_i = y_i, (i \geq 2)\quad
0 = a_1 \cdot (-a'_j) + a_j \cdot 1 (j \geq 2)
$$
表明该关系平凡（明确地说，元素 $a_{ij}$ 定义如下：令
$a_{11} = 0$；对 $j > 1$，令 $a_{1j} = -a'_j$；对
$i, j \geq 2$，令 $a_{ij} = \delta_{ij}$）。
\end{proof}

\begin{lemma}
\label{more-algebra-lemma-dedekind-torsion-free-flat}
设 $A$ 为 Dedekind 整环（例如离散赋值环，或更一般的 PID）。
\begin{enumerate}
\item $A$-模平坦当且仅当它无挠。
\item 有限无挠 $A$-模有限局部自由。
\item 若 $A$ 是 PID，则有限无挠 $A$-模有限自由。
\end{enumerate}
\end{lemma}

\begin{proof}
（关于引理陈述中的括号说明，见 Algebra, Lemma
\ref{algebra-lemma-PID-dedekind}。）证明 (1)。由引理
\ref{more-algebra-lemma-check-torsion} 与 Algebra, Lemma
\ref{algebra-lemma-flat-localization}，只需对极大理想
$\mathfrak m \subset A$ 在 $A_\mathfrak m$ 上检验该陈述。
由于 $A_\mathfrak m$ 是离散赋值环
（Algebra, Lemma \ref{algebra-lemma-characterize-Dedekind}），
结论由引理 \ref{more-algebra-lemma-valuation-ring-torsion-free-flat} 得出。

\medskip\noindent
证明 (2)。它由 Algebra, Lemma \ref{algebra-lemma-finite-projective}
与 (1) 得出。

\medskip\noindent
证明 (3)。设 $A$ 为 PID，$M$ 为有限无挠模。由引理
\ref{more-algebra-lemma-finite-torsion-free-submodule-free}，对某个 $n$ 有
$M \subset A^{\oplus n}$。对 $n$ 归纳证明 $M$ 自由。
$n = 1$ 的情形恰好表述 $A$ 是 PID 这一事实。若 $n > 1$，令
$M' \subset A^{\oplus n - 1}$ 为其在 $A^{\oplus n}$ 后
$n - 1$ 个直和分量的投影之像。于是得到短正合列
$0 \to I \to M \to M' \to 0$，其中 $I$ 是 $M$ 与
$A^{\oplus n}$ 的第一个直和分量 $A$ 的交。由归纳法，$M$ 是有限自由
$A$-模的扩张，因而有限自由。
\end{proof}

\begin{lemma}
\label{more-algebra-lemma-hom-into-torsion-free}
设 $R$ 为整环，$M$、$N$ 为 $R$-模。若 $N$ 无挠，则
$\Hom_R(M, N)$ 也无挠。
\end{lemma}

\begin{proof}
选取满射 $\bigoplus_{i \in I} R \to M$。则
$\Hom_R(M, N) \subset \prod_{i \in I} N$。
\end{proof}





\section{模的秩}
\label{more-algebra-section-rank}

\noindent
下面给出我们的定义。

\begin{definition}
\label{more-algebra-definition-rank}
设 $R$ 为分式域为 $K$ 的整环。$R$-模 $M$ 的{\it 秩}为
$\dim_K(M \otimes_R K)$。
\end{definition}

\noindent
当两个定义都适用时，该定义与 Algebra, Definition
\ref{algebra-definition-locally-free} 中秩为 $r$ 的局部自由模概念不冲突。

\begin{lemma}
\label{more-algebra-lemma-torsion-does-not-matter-rank}
设 $R$ 为整环。若 $R$-模映射 $M \to M'$ 的核与余核都是挠模，
则 $M$ 的秩等于 $M'$ 的秩。
\end{lemma}

\begin{proof}
略。提示：若 $K$ 是 $R$ 的分式域，则诱导映射
$M \otimes_R K \to M' \otimes_R K$ 是同构。
\end{proof}

\begin{lemma}
\label{more-algebra-lemma-rank-additive}
设 $R$ 为整环。设 $0 \to M \to M' \to M'' \to 0$ 为 $R$-模的
短正合列。则 $M'$ 的秩等于 $M$ 与 $M''$ 的秩之和。
\end{lemma}

\begin{proof}
略。
\end{proof}

\begin{lemma}
\label{more-algebra-lemma-rank-tensor}
设 $R$ 为整环，$M$ 和 $N$ 为 $R$-模。
\begin{enumerate}
\item $M \otimes_R N$ 的秩等于 $M$ 与 $N$ 的秩之积。
\item 若 $M$ 是有限表示 $R$-模，则 $\Hom_R(M, N)$ 的秩等于
$M$ 与 $N$ 的秩之积。
\end{enumerate}
\end{lemma}

\begin{proof}
第 (1) 项由向量空间的相应事实得出。假设 $M$ 作为 $R$-模有限表示。
由 Algebra, Lemma \ref{algebra-lemma-hom-from-finitely-presented}，
作为 $K$-向量空间，$\Hom_R(M, N) \otimes_R K$ 同构于
$\Hom_K(M \otimes_R K, N \otimes_R K)$。又由于
$M \otimes_R K$ 有限维，结论由向量空间的相应事实得出。
\end{proof}

\begin{lemma}
\label{more-algebra-lemma-dominant-pullback-rank}
设 $R \subset R'$ 为整环扩张。若 $M$ 是 $R$-模，则 $M$ 在 $R$ 上的秩
等于 $M \otimes_R R'$ 在 $R'$ 上的秩。
\end{lemma}

\begin{proof}
这是因为向量空间的维数在基域扩张下不变。
\end{proof}

\begin{lemma}
\label{more-algebra-lemma-finite-module-rank}
设 $R$ 为整环，$M$ 为有限 $R$-模。则：
\begin{enumerate}
\item $M$ 有有限秩 $r \geq 0$；
\item 存在映射 $M \to R^{\oplus r}$，其核与余核都是挠模；
\item 存在 $f \in R$、$f \not = 0$，使得 $M_f$ 是秩为 $r$ 的自由模；
\item 存在单射 $R^{\oplus r} \to M$，其余核是挠模。
\end{enumerate}
\end{lemma}

\begin{proof}
由于 $M$ 有限，$M \otimes_R K$ 是 $R$ 的分式域 $K$ 上的有限维向量
空间。因此秩 $r$ 有限。

\medskip\noindent
选取 $M$ 的生成元 $x_1, \ldots, x_n$，并选取
$M \otimes_R K$ 的一组基 $e_1, \ldots, e_r$。对某些
$a_{ij} , b_{ij} \in R$ 且 $b_{ij} \not = 0$，可写成
$x_i \otimes 1 = \sum (a_{ij}/b_{ij}) e_j$。令
$b = \prod b_{ij}$，并以 $a_{ij}b/b_{ij}$ 替换 $a_{ij}$，可假设
$x_i \otimes 1 = \sum (a_{ij}/b) e_j$。注意，若 $M$ 中有关系
$\sum r_i x_i = 0$，则对所有 $j = 1, \ldots, r$，
$\sum r_ia_{ij} = 0$ 在 $R$ 中成立。

\medskip\noindent
证明 (2)。上面的最后一个观察说明，把 $x_i$ 映到
$\sum a_{ij} e_j$ 的映射 $M \to \bigoplus Re_j$ 良定义。
$e_1, \ldots, e_r$ 构成 $M \otimes_R K$ 的一组基这一事实说明，
该映射与 $K$ 张量后成为同构。因此这是 (2) 中所述的映射。

\medskip\noindent
证明 (3)。在 (2) 中构造的映射 $M \to R^{\oplus r}$ 的余核有限且有挠。
故可找到零化该余核的 $g \in R$、$g \not = 0$。换言之，映射
$M_g \to R_g^{\oplus r}$ 是满射。于是可选取分裂
$M_g = N \oplus R_g^{\oplus r}$。此时 $N$ 是有限 $R_g$-模，且
$N \otimes_{R_g} K = 0$。换言之，$N$ 是 $R_g$ 上的有限挠模。
故可找到零化 $N$ 的 $g' \in R$。令 $f= gg'$，便得到 (3)。

\medskip\noindent
为证明 (4)，选取 $M_f$ 的一组基 $y_1, \ldots, y_r$，其中 $f$
如 (3) 所述。然后对某些 $m_j \in M$ 与 $t_j \geq 0$ 写成
$y_j = m_j/f^{t_j}$。把第 $j$ 个基向量映到 $m_j$ 的映射
$R^{\oplus r} \to M$ 即如 (4) 所述。略去细节。
\end{proof}





\section{自反模}
\label{more-algebra-section-reflexive}

\noindent
下面给出我们的定义。

\begin{definition}
\label{more-algebra-definition-reflexive}
设 $R$ 为整环。若自然映射
$$
j : M \longrightarrow \Hom_R(\Hom_R(M, R), R)
$$
是同构，则称 $R$-模 $M$ {\it 自反}；该映射把 $m \in M$ 映为如下
映射：它把 $\varphi \in \Hom_R(M, R)$ 映到 $\varphi(m) \in R$。
\end{definition}

\noindent
可以对更一般的环作此定义，但即使上面的定义也有缺点。明智的做法是
限于 Noether 整环与有限无挠模，并（或许）对 $R$ 加上某些正则性条件
（例如 $R$ 正规）。

\begin{lemma}
\label{more-algebra-lemma-reflexive-torsion-free}
设 $R$ 为整环，$M$ 为 $R$-模。
\begin{enumerate}
\item 若 $M$ 自反，则 $M$ 无挠。
\item 若 $M$ 有限，则
$j : M \to \Hom_R(\Hom_R(M, R), R)$ 的核与余核都是挠模。
\item 若 $M$ 有限，则 $j$ 是单射当且仅当 $M$ 无挠。
\end{enumerate}
\end{lemma}

\begin{proof}
为证明 (1)，由引理 \ref{more-algebra-lemma-hom-into-torsion-free}，$M$ 的双对偶无挠。
假设 $M$ 有限。选取 $M$ 在 $R$ 上的生成元 $x_1, \ldots, x_n$。
设 $K$ 为 $R$ 的分式域。重新编号后，可以假设对某个
$0 \leq r \leq n$，元 $x_1, \ldots, x_r$ 在 $M \otimes_R K$
中的像构成它在 $K$ 上的一组基。对 $r < i \leq n$，写成
$x_i \otimes 1 = \sum_{1 \leq j \leq r} x_j \otimes k_{ij}$，其中
$k_{ij} \in K$。选取 $c \in R$，使 $a_{ij} = ck_{ij} \in R$，
并使 $cx_i =  \sum_{1 \leq j \leq r} a_{ij} x_j$ 在 $M$ 中成立
（略去一个小细节）。于是得到映射
$$
R^{\oplus r} \xrightarrow{\alpha} M \xrightarrow{\beta} R^{\oplus r}
\xrightarrow{\alpha} M
$$
其中 $\alpha \circ \beta$ 与 $\beta \circ \alpha$ 都是乘以 $c$。
这里 $\alpha$ 由 $x_1, \ldots, x_r$ 给出；$\beta$ 对
$1 \leq i \leq r$ 把 $x_i$ 映到第 $i$ 个基向量的 $c$ 倍，
对 $r < i \leq n$ 则映到向量 $(a_{i1}, \ldots, a_{ir})$。
考虑交换图
$$
\xymatrix{
R^{\oplus r} \ar[r] \ar[d]^{\text{id}} &
M \ar[r] \ar[d]^j &
R^{\oplus r} \ar[d]^{\text{id}} \ar[r] &
M \ar[d]^j \\
R^{\oplus r} \ar[r] &
\Hom_R(\Hom_R(M, R), R) \ar[r] &
R^{\oplus r} \ar[r] &
\Hom_R(\Hom_R(M, R), R)
}
$$
当然，底行中相继箭头的复合都等于乘以 $c$。此图表明
$j$ 的核与余核都被 $c$ 零化，这证明了 (2)。
对 (3)，注意若 $M$ 无挠，则 $\beta$ 是单射，从而 $j$ 是单射。
逆命题是清楚的，因为我们已经看到 $M$ 的双对偶无挠。
\end{proof}

\begin{lemma}
\label{more-algebra-lemma-cokernel-map-double-dual-dvr}
设 $R$ 为离散赋值环，$M$ 为有限 $R$-模。则映射
$j : M \to \Hom_R(\Hom_R(M, R), R)$ 是满射。
\end{lemma}

\begin{proof}
设 $M_{tors} \subset M$ 为挠子模。则
$\Hom_R(M, R) = \Hom_R(M/M_{tors}, R)$（对任意整环都成立）。
因此可以假设 $M$ 无挠。由引理
\ref{more-algebra-lemma-dedekind-torsion-free-flat}，$M$ 自由，引理便是清楚的。
\end{proof}

\begin{lemma}
\label{more-algebra-lemma-check-reflexive}
设 $R$ 为 Noether 整环，$M$ 为有限 $R$-模。下列条件等价：
\begin{enumerate}
\item $M$ 自反；
\item 对每个素理想 $\mathfrak p \subset R$，$M_\mathfrak p$ 都是自反
$R_\mathfrak p$-模；
\item 对 $R$ 的每个极大理想 $\mathfrak m$，$M_\mathfrak m$ 都是自反
$R_\mathfrak m$-模。
\end{enumerate}
\end{lemma}

\begin{proof}
映射 $j : M \to \Hom_R(\Hom_R(M, R), R)$ 在素理想
$\mathfrak p$ 处的局部化，就是 Noether 局部整环 $R_\mathfrak p$ 上的模
$M_\mathfrak p$ 所对应的映射。见 Algebra, Lemma
\ref{algebra-lemma-hom-from-finitely-presented}。因此由 Algebra, Lemma
\ref{algebra-lemma-characterize-zero-local} 即得本引理。
\end{proof}

\begin{lemma}
\label{more-algebra-lemma-sequence-reflexive}
设 $R$ 为 Noether 整环。设 $0 \to M \to M' \to M''$ 是有限
$R$-模的正合列。若 $M'$ 自反且 $M''$ 无挠，则 $M$ 自反。
\end{lemma}

\begin{proof}
下面不再说明地使用这一事实：对任意有限 $R$-模 $N$ 与 $N'$，
$\Hom_R(N, N')$ 都是有限 $R$-模；见 Algebra, Lemma
\ref{algebra-lemma-ext-noetherian}。取对偶得到序列
$$
\Hom_R(M, R) \leftarrow \Hom_R(M', R) \leftarrow \Hom_R(M'', R)
$$
再次取对偶，得到交换图
{\small
$$
\xymatrix{
\Hom_R(\Hom_R(M, R), R) \ar[r]_j &
\Hom_R(\Hom_R(M', R), R) \ar[r] &
\Hom_R(\Hom_R(M'', R), R) \\
M \ar[u] \ar[r] & M' \ar[u] \ar[r] & M'' \ar[u]
}
$$
}
我们并不知道顶行是否正合。但由假设，中间的竖直箭头是同构，
右边的竖直箭头是单射（引理 \ref{more-algebra-lemma-reflexive-torsion-free}）。
我们断言 $j$ 是单射。若承认该断言，作图追迹即知左边的竖直箭头是同构，
也就是说 $M$ 自反。

\medskip\noindent
断言的证明。考虑定义 $Q$ 的正合列
$\Hom_R(M',R)\to \Hom_R(M,R)\to Q \to 0$。应用
Algebra, Lemma \ref{algebra-lemma-hom-from-finitely-presented} 得
$$
\Hom_K(M'\otimes_R K,K) \to
\Hom_K(M\otimes_R K,K) \to
Q\otimes_R K\to 0
$$
但 $M \otimes_R K \to M' \otimes_R K$ 是向量空间之间的单射，
因而是分裂单射，所以 $Q \otimes_R K = 0$，即 $Q$ 为挠模。
于是得到正合列
$$
0 \to \Hom_R(Q,R) \to
\Hom_R(\Hom_R(M,R),R) \to
\Hom_R(\Hom_R(M',R),R)
$$
而因 $Q$ 为挠模，$\Hom_R(Q,R)=0$。
\end{proof}

\begin{lemma}
\label{more-algebra-lemma-characterize-reflexive}
设 $R$ 为 Noether 整环，$M$ 为有限 $R$-模。下列条件等价：
\begin{enumerate}
\item $M$ 自反；
\item 存在短正合列 $0 \to M \to F \to N \to 0$，其中 $F$ 有限自由，
$N$ 无挠。
\end{enumerate}
\end{lemma}

\begin{proof}
注意有限自由模自反。由引理 \ref{more-algebra-lemma-sequence-reflexive}，
(2) 推出 (1)。假设 $M$ 自反。选取一个表示
$R^{\oplus m} \to R^{\oplus n} \to \Hom_R(M, R) \to 0$。
取对偶后得到正合列
$$
0 \to \Hom_R(\Hom_R(M, R), R) \to R^{\oplus n} \to N \to 0
$$
其中 $N = \Im(R^{\oplus n} \to R^{\oplus m})$ 为无挠模。因为
$M = \Hom_R(\Hom_R(M, R), R)$，所以得到 (2) 中所述的正合列。
\end{proof}

\begin{lemma}
\label{more-algebra-lemma-flat-pullback-reflexive}
设 $R \to R'$ 为 Noether 整环之间的平坦同态。若 $M$ 是有限自反
$R$-模，则 $M \otimes_R R'$ 是有限自反 $R'$-模。
\end{lemma}

\begin{proof}
选取短正合列 $0 \to M \to F \to N \to 0$，其中 $F$ 有限自由，
$N$ 无挠；见引理 \ref{more-algebra-lemma-characterize-reflexive}。由于 $R \to R'$ 平坦，
我们得到短正合列
$0 \to M \otimes_R R' \to F \otimes_R R' \to N \otimes_R R' \to 0$，
其中 $F \otimes_R R'$ 有限自由，$N \otimes_R R'$ 无挠
（引理 \ref{more-algebra-lemma-flat-pullback-torsion}）。因此由引理
\ref{more-algebra-lemma-characterize-reflexive}，$M \otimes_R R'$ 自反。
\end{proof}

\begin{lemma}
\label{more-algebra-lemma-dual-reflexive}
设 $R$ 为 Noether 整环，$M$ 为有限 $R$-模，$N$ 为有限自反
$R$-模。则 $\Hom_R(M, N)$ 自反。
\end{lemma}

\begin{proof}
选取一个表示 $R^{\oplus m} \to R^{\oplus n} \to M \to 0$。
于是得到
$$
0 \to \Hom_R(M, N) \to N^{\oplus n} \to N' \to 0
$$
其中 $N' = \Im(N^{\oplus n} \to N^{\oplus m})$ 无挠。由引理
\ref{more-algebra-lemma-sequence-reflexive} 得到结论。
\end{proof}

\begin{definition}
\label{more-algebra-definition-reflexive-hull}
设 $R$ 为 Noether 整环，$M$ 为有限 $R$-模。称模
$M^{**} = \Hom_R(\Hom_R(M, R), R)$ 为 $M$ 的{\it 自反包}。
\end{definition}

\noindent
这一定义是合理的，因为由引理 \ref{more-algebra-lemma-dual-reflexive}，自反包自反。
对应 $M \mapsto M^{**}$ 是一个函子。若 $\varphi : M \to N$ 是到自反
$R$-模 $N$ 的 $R$-模映射，则 $\varphi$ 经过 $M$ 的自反包分解为
$M \to M^{**} \to N$。换言之，在 Noether 整环 $R$ 上，取自反包是包含函子
$$
\text{finite reflexive modules} \subset
\text{finite modules}
$$
的左伴随。

\begin{lemma}
\label{more-algebra-lemma-hom-into-depth}
设 $R$ 为 Noether 局部环，$M$、$N$ 为有限 $R$-模。
\begin{enumerate}
\item 若 $N$ 的深度 $\geq 1$，则 $\Hom_R(M, N)$ 的深度 $\geq 1$。
\item 若 $N$ 的深度 $\geq 2$，则 $\Hom_R(M, N)$ 的深度 $\geq 2$。
\end{enumerate}
\end{lemma}

\begin{proof}
选取一个表示 $R^{\oplus m} \to R^{\oplus n} \to M \to 0$。
取对偶后得到正合列
$$
0 \to \Hom_R(M, N) \to N^{\oplus n} \to N' \to 0
$$
其中 $N' = \Im(N^{\oplus n} \to N^{\oplus m})$。深度 $\geq 1$ 的模的子模深度
$\geq 1$；这直接由定义得出。因此 (1) 是清楚的。对 (2)，注意这里的
假设与上一观察说明 $N'$ 的深度 $\geq 1$。模 $N^{\oplus n}$ 的深度
$\geq 2$。由 Algebra, Lemma \ref{algebra-lemma-depth-in-ses}，可知
$\Hom_R(M, N)$ 的深度 $\geq 2$。
\end{proof}

\begin{lemma}
\label{more-algebra-lemma-hom-into-S2}
设 $R$ 为 Noether 环，$M$、$N$ 为有限 $R$-模。
\begin{enumerate}
\item 若 $N$ 具有性质 $(S_1)$，则 $\Hom_R(M, N)$ 具有性质 $(S_1)$。
\item 若 $N$ 具有性质 $(S_2)$，则 $\Hom_R(M, N)$ 具有性质 $(S_2)$。
\item 若 $R$ 是整环，$N$ 无挠且具有 $(S_2)$，则
$\Hom_R(M, N)$ 无挠且具有性质 $(S_2)$。
\end{enumerate}
\end{lemma}

\begin{proof}
对有限 $R$-模，在素理想处局部化与取 $\Hom_R$ 可交换
（Algebra, Lemma \ref{algebra-lemma-ext-noetherian}），因此 (1) 与 (2)
直接由引理 \ref{more-algebra-lemma-hom-into-depth} 得出。(3) 由 (2) 与引理
\ref{more-algebra-lemma-hom-into-torsion-free} 得出。
\end{proof}

\begin{lemma}
\label{more-algebra-lemma-check-injective-on-ass}
设 $R$ 为 Noether 环，$\varphi : M \to N$ 为 $R$-模的映射。假设对 $R$
的每个素理想 $\mathfrak p$，下列情形至少有一个发生：
\begin{enumerate}
\item $M_\mathfrak p \to N_\mathfrak p$ 是单射；或
\item $\mathfrak p \not \in \text{Ass}(M)$。
\end{enumerate}
则 $\varphi$ 是单射。
\end{lemma}

\begin{proof}
设 $\mathfrak p$ 为 $\Ker(\varphi)$ 的一个伴随素理想。则存在元
$x \in M_\mathfrak p$，它属于 $M_\mathfrak p \to N_\mathfrak p$ 的核，且其零化子为
$\mathfrak pR_\mathfrak p$（Algebra, Lemma
\ref{algebra-lemma-associated-primes-localize}）。在上述两种情形中这都不可能。
因此 $\text{Ass}(\Ker(\varphi)) = \emptyset$，再由 Algebra, Lemma
\ref{algebra-lemma-ass-zero} 得 $\Ker(\varphi) = 0$。
\end{proof}

\begin{lemma}
\label{more-algebra-lemma-check-isomorphism-via-depth-and-ass}
设 $R$ 为 Noether 环，$\varphi : M \to N$ 为 $R$-模的映射。假设 $M$ 有限，
且对 $R$ 的每个素理想 $\mathfrak p$，下列情形之一发生：
\begin{enumerate}
\item $M_\mathfrak p \to N_\mathfrak p$ 是同构；或
\item $\text{depth}(M_\mathfrak p) \geq 2$ 且
$\mathfrak p \not \in \text{Ass}(N)$。
\end{enumerate}
则 $\varphi$ 是同构。
\end{lemma}

\begin{proof}
由引理 \ref{more-algebra-lemma-check-injective-on-ass}，$\varphi$ 是单射。
设 $N' \subset N$ 是包含 $M$ 的像的有限生成 $R$-模。那么
$\text{Ass}(N_\mathfrak p) = \emptyset$ 推出
$\text{Ass}(N'_\mathfrak p) = \emptyset$。因此本引理的假设对 $M \to N'$ 成立。
为了证明 $\varphi$ 是同构，只须对每个这样的 $N' \subset N$ 证明同一结论。
因此可以假设 $N$ 是有限 $R$-模。在此情形下，
$\mathfrak p \not \in \text{Ass}(N) \Rightarrow
\text{depth}(N_\mathfrak p) \geq 1$；见 Algebra, Lemma
\ref{algebra-lemma-ideal-nonzerodivisor}。考虑定义 $Q$ 的短正合列
$$
0 \to M \to N \to Q \to 0
$$
检视各条件可知：在情形 (1) 中 $Q_\mathfrak p = 0$，而在情形 (2) 中，
由 Algebra, Lemma \ref{algebra-lemma-depth-in-ses} 有
$\text{depth}(Q_\mathfrak p) \geq 1$。这说明 $Q$ 没有任何伴随素理想，
因此由 Algebra, Lemma \ref{algebra-lemma-ass-zero} 得 $Q = 0$。
\end{proof}

\begin{lemma}
\label{more-algebra-lemma-isom-depth-2-torsion-free}
设 $R$ 为 Noether 整环，$\varphi : M \to N$ 为 $R$-模的映射。假设
$M$ 有限、$N$ 无挠，且对 $R$ 的每个素理想 $\mathfrak p$，下列情形之一发生：
\begin{enumerate}
\item $M_\mathfrak p \to N_\mathfrak p$ 是同构；或
\item $\text{depth}(M_\mathfrak p) \geq 2$。
\end{enumerate}
则 $\varphi$ 是同构。
\end{lemma}

\begin{proof}
这是引理 \ref{more-algebra-lemma-check-isomorphism-via-depth-and-ass} 的特殊情形。
\end{proof}

\begin{lemma}
\label{more-algebra-lemma-reflexive-depth-2}
设 $R$ 为 Noether 整环，$M$ 为有限 $R$-模。下列条件等价：
\begin{enumerate}
\item $M$ 自反；
\item 对 $R$ 的每个素理想 $\mathfrak p$，下列情形之一发生：
\begin{enumerate}
\item $M_\mathfrak p$ 是自反 $R_\mathfrak p$-模；或
\item $\text{depth}(M_\mathfrak p) \geq 2$。
\end{enumerate}
\end{enumerate}
\end{lemma}

\begin{proof}
若 (1) 成立，则由引理 \ref{more-algebra-lemma-check-reflexive}，对每个素理想
$\mathfrak p$，$M_\mathfrak p$ 都是自反模。因此 (1) $\Rightarrow$ (2)。
假设 (2)。令 $N = \Hom_R(\Hom_R(M, R), R)$，则对 $R$ 的每个素理想
$\mathfrak p$ 有
$$
N_\mathfrak p =
\Hom_{R_\mathfrak p}(
\Hom_{R_\mathfrak p}(M_\mathfrak p, R_\mathfrak p), R_\mathfrak p)
$$
见 Algebra, Lemma \ref{algebra-lemma-hom-from-finitely-presented}。
对映射 $j : M \to N$ 应用引理
\ref{more-algebra-lemma-isom-depth-2-torsion-free}。这是允许的，因为 $M$ 有限，而由引理
\ref{more-algebra-lemma-hom-into-torsion-free}，$N$ 无挠。在情形 (2)(a) 中，映射
$M_\mathfrak p \to N_\mathfrak p$ 是同构；在情形 (2)(b) 中，有
$\text{depth}(M_\mathfrak p) \geq 2$。
\end{proof}

\begin{lemma}
\label{more-algebra-lemma-reflexive-S2}
设 $R$ 为 Noether 整环，$M$ 为有限自反 $R$-模。设
$\mathfrak p \subset R$ 为素理想。
\begin{enumerate}
\item 若 $\text{depth}(R_\mathfrak p) \geq 2$，则
$\text{depth}(M_\mathfrak p) \geq 2$。
\item 若 $R$ 是 $(S_2)$，则 $M$ 是 $(S_2)$。
\end{enumerate}
\end{lemma}

\begin{proof}
由于形成自反包 $\Hom_R(\Hom_R(M, R), R)$ 与局部化可交换
（Algebra, Lemma \ref{algebra-lemma-hom-from-finitely-presented}），
(1) 由引理 \ref{more-algebra-lemma-hom-into-depth} 得出。(2) 直接由引理
\ref{more-algebra-lemma-hom-into-S2} 得出。
\end{proof}

\begin{example}
\label{more-algebra-example-ring-not-S2}
上下的结果使人猜想自反性与 $(S_2)$ 条件有关；下面的例子用来阻止过于乐观的猜想。
设 $k$ 为域。设 $R$ 是 $k[x, y]$ 由 $1, y, x^2, xy, x^3$ 生成的
$k$-子代数。则 $R$ 不是 $(S_2)$。所以把 $R$ 看作 $R$-模，它是一个自反但不是
$(S_2)$ 的 $R$-模之例。设 $M = k[x, y]$，并将其视为 $R$-模。
则 $M$ 是自反 $R$-模，因为
$$
\Hom_R(M, R) = \mathfrak m = (y, x^2, xy, x^3)
\quad\text{且}\quad
\Hom_R(\mathfrak m, R) = M
$$
且 $M$ 作为 $R$-模是 $(S_2)$（略去计算）。因此 $R$ 是一个 Noether 整环，
它拥有自反 $(S_2)$ 模，但 $R$ 本身不是 $(S_2)$。
\end{example}

\begin{lemma}
\label{more-algebra-lemma-reflexive-over-normal}
设 $R$ 为 Noether 正规整环，其分式域为 $K$。设 $M$ 为有限 $R$-模。
下列条件等价：
\begin{enumerate}
\item $M$ 自反；
\item $M$ 无挠且具有性质 $(S_2)$；
\item $M$ 无挠，且
$M = \bigcap_{\text{高度}(\mathfrak p) = 1} M_{\mathfrak p}$，
其中交取在 $M_K = M \otimes_R K$ 中。
\end{enumerate}
\end{lemma}

\begin{proof}
由 Algebra, Lemma \ref{algebra-lemma-criterion-normal}，$R$ 满足 $(R_1)$ 与 $(S_2)$。

\medskip\noindent
假设 (1)。由引理 \ref{more-algebra-lemma-reflexive-torsion-free}，$M$ 无挠；
由引理 \ref{more-algebra-lemma-reflexive-S2}，它满足 $(S_2)$。因此 (2) 成立。

\medskip\noindent
假设 (2)。依定义，
$M' = \bigcap_{\text{高度}(\mathfrak p) = 1} M_{\mathfrak p}$
是映射
$$
M_K
\longrightarrow
\bigoplus\nolimits_{\text{高度}(\mathfrak p) = 1} M_K/M_\mathfrak p
\subset
\prod\nolimits_{\text{高度}(\mathfrak p) = 1} M_K/M_\mathfrak p
$$
的核。注意，我们的映射确实如所示经过直和分解，因为给定
$a/b \in K$，至多只有有限多个高为 $1$ 的素理想 $\mathfrak p$ 满足
$b \in \mathfrak p$。设 $\mathfrak p_0$ 为一个高为 $1$ 的素理想。则除非两者相等，
否则 $(M_K/M_\mathfrak p)_{\mathfrak p_0} = 0$；
而在 $\mathfrak p = \mathfrak p_0$ 时，我们得到
$(M_K/M_\mathfrak p)_{\mathfrak p_0} = M_K/M_{\mathfrak p_0}$。
因此，由局部化的正合性以及局部化与直和可交换，有
$M'_{\mathfrak p_0} = M_{\mathfrak p_0}$。因为 $M$ 在高 $> 1$ 的素理想处深度
$\geq 2$，由引理 \ref{more-algebra-lemma-isom-depth-2-torsion-free}，$M \to M'$ 是同构。
因此 (3) 成立。

\medskip\noindent
假设 (3)。设 $\mathfrak p$ 为高 $1$ 的素理想。由 $(R_1)$，
$R_\mathfrak p$ 是离散赋值环。由引理 \ref{more-algebra-lemma-dedekind-torsion-free-flat}，
$M_\mathfrak p$ 有限自由，特别地它自反。因此映射 $M \to M^{**}$ 在所有高为
$1$ 的素理想 $\mathfrak p$ 处诱导出同构。于是条件
$M = \bigcap_{\text{高度}(\mathfrak p) = 1} M_{\mathfrak p}$
推出 $M = M^{**}$，从而 (1) 成立。
\end{proof}

\begin{lemma}
\label{more-algebra-lemma-describe-reflexive-hull}
设 $R$ 为 Noether 正规整环，$M$ 为有限 $R$-模。则 $M$ 的自反包是交
$$
M^{**} =
\bigcap\nolimits_{\text{高度}(\mathfrak p) = 1}
M_{\mathfrak p}/(M_\mathfrak p)_{tors} =
\bigcap\nolimits_{\text{高度}(\mathfrak p) = 1}
(M/M_{tors})_\mathfrak p
$$
，其交取在 $M \otimes_R K$ 中。
\end{lemma}

\begin{proof}
设 $\mathfrak p$ 为高 $1$ 的素理想。映射
$M_\mathfrak p \to M \otimes_R K$ 的核是 $M_\mathfrak p$ 的挠子模
$(M_\mathfrak p)_{tors}$。此外，
$(M/M_{tors})_\mathfrak p = M_\mathfrak p/(M_\mathfrak p)_{tors}$，
并且它是离散赋值环 $R_\mathfrak p$ 上的有限自由模
（引理 \ref{more-algebra-lemma-dedekind-torsion-free-flat}）。于是
$M_\mathfrak p/(M_\mathfrak p)_{tors} \to
(M_\mathfrak p)^{**} = (M^{**})_\mathfrak p$
是同构，因而本引理是引理 \ref{more-algebra-lemma-reflexive-over-normal} 的结论。
\end{proof}

\begin{lemma}
\label{more-algebra-lemma-integral-closure-reflexive}
设 $A$ 为 Noether 正规整环，分式域为 $K$。设 $L$ 是 $K$ 的有限扩张。
若 $A$ 在 $L$ 中的整闭包 $B$ 在 $A$ 上有限，则 $B$ 作为 $A$-模自反。
\end{lemma}

\begin{proof}
只须证明 $B = \bigcap B_\mathfrak p$，其中交遍历高为 $1$ 的素理想
$\mathfrak p \subset A$；见引理 \ref{more-algebra-lemma-reflexive-over-normal}。
设 $b \in \bigcap B_\mathfrak p$。设
$x^d + a_1x^{d - 1} + \ldots + a_d$ 是 $b$ 在 $K$ 上的最小多项式。
我们要证明 $a_i \in A$。由 Algebra, Lemma
\ref{algebra-lemma-minimal-polynomial-normal-domain}，对所有 $i$ 及所有高一素理想
$\mathfrak p$，有 $a_i \in A_\mathfrak p$。因此由 Algebra, Lemma
\ref{algebra-lemma-normal-domain-intersection-localizations-height-1}
得到所需结论（也可使用已引用的引理，因为 $A$ 是其自身上的自反模）。
\end{proof}








\section{内容理想}
\label{more-algebra-section-content}

\noindent
这个定义可能与你预期的不同。

\begin{definition}
\label{more-algebra-definition-content-ideal}
设 $A$ 为环，$M$ 为平坦 $A$-模，$x \in M$。若 $A$ 中所有满足
$x \in IM$ 的理想 $I$ 所成的集合有最小元，则称该最小元为
{\it $x$ 的内容理想}。
\end{definition}

\noindent
注意，因为 $M$ 在 $A$ 上平坦，对 $A$ 的一对理想 $I, I'$，有
$IM \cap I'M = (I \cap I')M$。这可由把正合列
$0 \to I \cap I' \to I \oplus I' \to I + I' \to 0$ 与 $M$ 张量而得。

\begin{lemma}
\label{more-algebra-lemma-content-finitely-generated}
设 $A$ 为环，$M$ 为平坦 $A$-模，$x \in M$。若 $x$ 的内容理想存在，
则它有限生成。
\end{lemma}

\begin{proof}
设 $x \in IM$。则可写成 $x = \sum_{i = 1, \ldots, n} f_i x_i$，其中
$f_i \in I$ 且 $x_i \in M$。因此 $x \in I'M$，其中
$I' = (f_1, \ldots, f_n)$。
\end{proof}

\begin{lemma}
\label{more-algebra-lemma-equal-content}
设 $(A, \mathfrak m)$ 为局部环。设 $u : M \to N$ 为平坦 $A$-模之间的映射，
且 $\overline{u} : M/\mathfrak m M \to N/\mathfrak m N$ 是单射。若 $x \in M$
的内容理想为 $I$，则 $u(x)$ 的内容理想也为 $I$。
\end{lemma}

\begin{proof}
显然 $u(x) \in IN$。若 $u(x) \in I'N$，则
$u(x) \in (I' \cap I)N$；见定义 \ref{more-algebra-definition-content-ideal} 后的讨论。
因此只须证明：若 $x \in I'N$ 且 $I' \subset I$、$I' \not =  I$，则
$u(x) \not \in I'N$。因为 $I/I'$ 是非零有限 $A$-模
（引理 \ref{more-algebra-lemma-content-finitely-generated}），由 Nakayama 引理
（Algebra, Lemma \ref{algebra-lemma-NAK}），存在 $A$-模的非零映射
$\chi : I/I' \to A/\mathfrak m$。由于 $I$ 是 $x$ 的内容理想，有
$x \not \in I''M$，其中 $I'' = \Ker(\chi)$。因而 $x$ 不在映射
$$
IM = I \otimes_A M \xrightarrow{\chi \otimes 1}
A/\mathfrak m \otimes M \cong M/\mathfrak m M
$$
的核中。应用关于 $\overline{u}$ 的假设，可知 $u(x)$ 在映射
$$
IN = I \otimes_A N \xrightarrow{\chi \otimes 1}
A/\mathfrak m \otimes N \cong N/\mathfrak m N
$$
下的像不为零，于是得到结论。
\end{proof}

\begin{lemma}
\label{more-algebra-lemma-content-exists-flat-Mittag-Leffler}
设 $A$ 为环，$M$ 为平坦 Mittag-Leffler 模。则 $M$ 的每个元都有内容理想。
\end{lemma}

\begin{proof}
这是 Algebra, Lemma \ref{algebra-lemma-flat-ML-criterion} 的特殊情形。
\end{proof}








\section{平坦性与有限性条件}
\label{more-algebra-section-flat-finite}

\noindent
本节讨论“平坦 $+$ 有限型 $\Rightarrow$ 有限表示”这一类型的一些推论。
我们将在平坦性一章重新讨论这一结果；见 More on Flatness, Section
\ref{flat-section-introduction}。这种类型的第一个结果已在 Algebra, Lemma
\ref{algebra-lemma-finite-flat-module-finitely-presented} 中证明。

\begin{lemma}
\label{more-algebra-lemma-flat-finite-type-finite-presentation-local-module}
设 $R$ 为环。设 $S = R[x_1, \ldots, x_n]$ 为 $R$ 上的多项式环。
设 $M$ 为 $S$-模。假设：
\begin{enumerate}
\item $R$ 有有限多个素理想 $\mathfrak p_1, \ldots, \mathfrak p_m$，使映射
$R \to \prod R_{\mathfrak p_j}$ 为单射；
\item $M$ 是有限 $S$-模；
\item $M$ 在 $R$ 上平坦；以及
\item 对 $R$ 的每个素理想 $\mathfrak p$，模 $M_{\mathfrak p}$ 在
$S_{\mathfrak p}$ 上有限表示。
\end{enumerate}
则 $M$ 在 $S$ 上有限表示。
\end{lemma}

\begin{proof}
选取 $M$ 作为 $S$-模的表示
$$
0 \to K \to S^{\oplus r} \to M \to 0
$$
设 $\mathfrak q$ 为 $S$ 的一个素理想，它位于 $R$ 的素理想
$\mathfrak p$ 之上。由假设，存在有限多个元 $k_1, \ldots, k_t \in K$，
使得若令 $K' = \sum Sk_j \subset K$，则
$K'_{\mathfrak p} = K_{\mathfrak p}$ 且对 $j = 1, \ldots, m$ 有
$K'_{\mathfrak p_j} = K_{\mathfrak p_j}$。令 $M' = S^{\oplus r}/K'$，特别可得
$M'_{\mathfrak q} = M_{\mathfrak q}$。由平坦性的开性，见 Algebra, Theorem
\ref{algebra-theorem-openness-flatness}，存在 $g \in S$、$g \not \in \mathfrak q$，
使 $M'_g$ 在 $R$ 上平坦。因此 $M'_g \to M_g$ 是平坦 $R$-模之间的满射。
考虑交换图
$$
\xymatrix{
M'_g \ar[r] \ar[d] & M_g \ar[d] \\
\prod (M'_g)_{\mathfrak p_j} \ar[r] & \prod (M_g)_{\mathfrak p_j}
}
$$
由 $k_1, \ldots, k_t$ 的选取，底部箭头是同构。由于
$R \to \prod R_{\mathfrak p_j}$ 是单射，而 $M'_g$ 在 $R$ 上平坦，
左边竖直箭头是单射。因此顶部水平箭头是单射，从而是同构。
这证明 $M_g$ 在 $S_g$ 上有限表示。应用 Algebra, Lemma
\ref{algebra-lemma-cover} 得到结论。
\end{proof}

\begin{lemma}
\label{more-algebra-lemma-flat-finite-type-finite-presentation-local}
设 $R \to S$ 为环同态。假设：
\begin{enumerate}
\item $R$ 有有限多个素理想 $\mathfrak p_1, \ldots, \mathfrak p_m$，使映射
$R \to \prod R_{\mathfrak p_j}$ 为单射；
\item $R \to S$ 为有限型；
\item $S$ 在 $R$ 上平坦；以及
\item 对 $R$ 的每个素理想 $\mathfrak p$，环 $S_{\mathfrak p}$ 在
$R_{\mathfrak p}$ 上有限表示。
\end{enumerate}
则 $S$ 在 $R$ 上有限表示。
\end{lemma}

\begin{proof}
由假设，$S$ 是 $R$ 上某个多项式环的商。因此结果直接由引理
\ref{more-algebra-lemma-flat-finite-type-finite-presentation-local-module} 得出。
\end{proof}

\begin{lemma}
\label{more-algebra-lemma-flat-graded-finite-type-finite-presentation-module}
设 $R$ 为环。设 $S = R[x_1, \ldots, x_n]$ 为 $R$ 上的分次多项式代数，
也就是说 $\deg(x_i) > 0$ 但不一定等于 $1$。设 $M$ 为分次 $S$-模。
假设：
\begin{enumerate}
\item $R$ 是局部环；
\item $M$ 是有限 $S$-模；以及
\item $M$ 在 $R$ 上平坦。
\end{enumerate}
则 $M$ 作为 $S$-模有限表示。
\end{lemma}

\begin{proof}
设 $M = \bigoplus M_d$ 为 $M$ 上的分次。选取 $M$ 的齐次生成元
$m_1, \ldots, m_r \in M$。设 $\deg(m_i) = d_i \in \mathbf{Z}$。
这给出表示
$$
0 \to K \to \bigoplus\nolimits_{i = 1, \ldots, r} S(-d_i) \to M \to 0
$$
它在每个次数 $d$ 上给出短正合列
$$
0 \to K_d \to \bigoplus\nolimits_{i = 1, \ldots, r} S_{d - d_i} \to
M_d \to 0.
$$
由假设，每个 $M_d$ 都是有限平坦 $R$-模。由 Algebra, Lemma
\ref{algebra-lemma-finite-flat-local}，每个 $M_d$ 都是有限自由 $R$-模。
因此每个 $K_d$ 都是有限 $R$-模。另外，由 Algebra, Lemma
\ref{algebra-lemma-flat-ses}，每个 $K_d$ 都在 $R$ 上平坦。再次由 Algebra, Lemma
\ref{algebra-lemma-finite-flat-local}，得出每个 $K_d$ 有限自由。

\medskip\noindent
设 $\mathfrak m$ 为 $R$ 的极大理想。由于 $M$ 在 $R$ 上平坦，上述短正合列与
$\kappa = \kappa(\mathfrak m)$ 张量后仍是短正合列。由于环
$S \otimes_R \kappa$ 是 Noether 环，存在齐次元
$k_1, \ldots, k_t \in K$，使其像 $\overline{k}_j$ 在
$S \otimes_R \kappa$ 上生成 $K \otimes_R \kappa$。设 $\deg(k_j) = e_j$。
因此对任意 $d$，映射
$$
\bigoplus\nolimits_{j = 1, \ldots, t} S_{d - e_j}
\longrightarrow
K_d
$$
与 $\kappa$ 张量后成为满射。由 Nakayama 引理
（Algebra, Lemma \ref{algebra-lemma-NAK}），该映射在 $R$ 上是满射。
因此 $K$ 由 $k_1, \ldots, k_t$ 在 $S$ 上生成，于是得证。
\end{proof}

\begin{lemma}
\label{more-algebra-lemma-flat-graded-finite-type-finite-presentation}
设 $R$ 为环。设 $S = \bigoplus_{n \geq 0} S_n$ 为分次 $R$-代数，
$M = \bigoplus_{d \in \mathbf{Z}} M_d$ 为分次 $S$-模。假设 $S$ 作为
$R$-代数有限生成，$S_0$ 是有限 $R$-代数，且存在有限多个素理想
$\mathfrak p_j$、$i = 1, \ldots, m$，使
$R \to \prod R_{\mathfrak p_j}$ 是单射。
\begin{enumerate}
\item 若 $S$ 在 $R$ 上平坦，则 $S$ 是有限表示 $R$-代数。
\item 若 $M$ 作为 $R$-模平坦，作为 $S$-模有限，则 $M$ 作为
$S$-模有限表示。
\end{enumerate}
\end{lemma}

\begin{proof}
由于 $S$ 作为 $R$-代数有限生成，它作为 $S_0$-代数也有限生成；
设由次数为 $d_1, \ldots, d_n > 0$ 的齐次元
$t_1, \ldots, t_n \in S$ 生成。令 $P = R[x_1, \ldots, x_n]$，并令
$\deg(x_i) = d_i$。由于 $S_0$ 是有限 $R$-模，环映射
$P \to S$、$x_i \to t_i$ 是有限的。为证明 (1)，只须证明 $S$ 是有限表示
$P$-模。为证明 (2)，只须证明 $M$ 是有限表示 $P$-模。因此，
只须证明如下情形：$S = P$ 是分次多项式环，$M$ 是有限 $S$-模且在
$R$ 上平坦，则 $M$ 作为 $S$-模有限表示。由引理
\ref{more-algebra-lemma-flat-graded-finite-type-finite-presentation-module}，对 $R$ 的每个素理想
$\mathfrak p$，$M_{\mathfrak p}$ 都是有限表示 $S_{\mathfrak p}$-模。
所以结果由引理 \ref{more-algebra-lemma-flat-finite-type-finite-presentation-local-module} 得出。
\end{proof}

\begin{remark}
\label{more-algebra-remark-when-does-condition-hold}
设 $R$ 为环。$R$ 何时满足引理
\ref{more-algebra-lemma-flat-finite-type-finite-presentation-local-module}、
\ref{more-algebra-lemma-flat-finite-type-finite-presentation-local} 与
\ref{more-algebra-lemma-flat-graded-finite-type-finite-presentation} 中提及的条件？
下列情形中该条件成立：
\begin{enumerate}
\item $R$ 是局部环；
\item $R$ 是 Noether 环；
\item $R$ 是整环；
\item $R$ 是只有有限多个极小素理想的既约环；或
\item $R$ 只有有限多个弱伴随素理想；见 Algebra, Lemma
\ref{algebra-lemma-zero-at-weakly-ass-zero}。
\end{enumerate}
因此这些引理在上述所有情形中都成立。
\end{remark}

\noindent
下一引理将在 More on Flatness, Proposition
\ref{flat-proposition-flat-finite-type-finite-presentation-domain} 中得到改进。

\begin{lemma}
\label{more-algebra-lemma-flat-finite-type-valuation-ring-finite-presentation}
\begin{reference}
\cite[Theorem 3]{Nagata-Finitely}
\end{reference}
设 $A$ 为赋值环。设 $A \to B$ 为有限型环映射，$M$ 为有限 $B$-模。
\begin{enumerate}
\item 若 $B$ 在 $A$ 上平坦，则 $B$ 是有限表示 $A$-代数。
\item 若 $M$ 作为 $A$-模平坦，则 $M$ 作为 $B$-模有限表示。
\end{enumerate}
\end{lemma}

\begin{proof}
我们将使用这一事实：$A$-模平坦当且仅当它无挠；见引理
\ref{more-algebra-lemma-valuation-ring-torsion-free-flat}。由 Algebra, Lemma
\ref{algebra-lemma-homogenize}，可找到分次 $A$-代数 $S$，其中 $S_0 = A$，
且它由有限多个次数为 $1$ 的元生成；还可找到一个元
$f \in S_1$ 和一个有限分次 $S$-模 $N$，使
$B \cong S_{(f)}$ 且 $M \cong N_{(f)}$。若 $M$ 无挠，则用 $N/N_{tors}$
作此替换后，可以取 $N$ 无挠；见引理 \ref{more-algebra-lemma-torsion}。
同样，若 $B$ 无挠，则用 $S/S_{tors}$ 替换后，可以取 $S$ 无挠。
因此在情形 (1) 中，可以应用引理
\ref{more-algebra-lemma-flat-graded-finite-type-finite-presentation}，得到 $S$ 是有限表示
$A$-代数，这推出 $B = S_{(f)}$ 是有限表示 $A$-代数。
为证明 (2)，可以先用分次多项式环替换 $S$，然后应用引理
\ref{more-algebra-lemma-flat-graded-finite-type-finite-presentation-module} 得到结论。
\end{proof}

\begin{lemma}
\label{more-algebra-lemma-valuation-ring-flat-essentially-finite-type}
设 $A$ 为赋值环。设 $A \to B$ 为本质有限型的局部同态，$M$ 为有限
$B$-模。
\begin{enumerate}
\item 若 $B$ 在 $A$ 上平坦，则 $B$ 在 $A$ 上本质有限表示。
\item 若 $M$ 作为 $A$-模平坦，则 $M$ 作为 $B$-模有限表示。
\end{enumerate}
\end{lemma}

\begin{proof}
由假设，可以把 $B$ 写成多项式代数
$P = A[x_1, \ldots, x_n]$ 在素理想 $\mathfrak q$ 处局部化的一个商。
在情形 (1) 中，我们把 $M = B$ 视为 $P_\mathfrak q$ 上的有限模；
在情形 (2) 中，我们把 $M$ 视为 $P_\mathfrak q$ 上的有限模。
在两种情形中，都要证明这是有限表示 $P_\mathfrak q$-模；
对情形 (2)，见 Algebra, Lemma
\ref{algebra-lemma-finitely-presented-over-subring}。

\medskip\noindent
选取表示 $0 \to K \to P_\mathfrak q^{\oplus r} \to M \to 0$；
这是可能的，因为 $M$ 在 $P_\mathfrak q$ 上有限。
令 $L = P^{\oplus r} \cap K$。则 $K = L_\mathfrak q$；见 Algebra, Lemma
\ref{algebra-lemma-submodule-localization}。此时 $N = P^{\oplus r}/L$ 是 $M$ 的子模，
从而由引理 \ref{more-algebra-lemma-valuation-ring-torsion-free-flat} 它平坦。
又因为 $N$ 是有限 $P$-模，由引理
\ref{more-algebra-lemma-flat-finite-type-valuation-ring-finite-presentation}，$N$ 作为 $P$-模有限表示。
由局部化的正合性（Algebra, Proposition
\ref{algebra-proposition-localization-exact}），可知 $N_\mathfrak q = M$，于是得证。
\end{proof}











\section{爆破与平坦性}
\label{more-algebra-section-blowup-flat}

\noindent
本节开始讨论如下形式的结果：“经过爆破后，严格变换变为平坦”。
这种类型的更多结果见 Divisors, Section
\ref{divisors-section-blowup-flat} 与 More on Flatness, Section
\ref{flat-section-blowup-flat}。

\begin{definition}
\label{more-algebra-definition-strict-transform}
设 $R$ 为环，$I \subset R$ 为理想，$a \in I$。设
$R[\frac{I}{a}]$ 为仿射爆破代数；见 Algebra, Definition
\ref{algebra-definition-blow-up}。设 $M$ 为 $R$-模。称 $R[\frac{I}{a}]$-模
$$
M' = \left(M \otimes_R R[\textstyle{\frac{I}{a}}]\right)/a\text{-幂挠子}
$$
为 {\it $M$ 沿 $R \to R[\frac{I}{a}]$ 的严格变换}。
\end{definition}

\noindent
下一结果是通过爆破实现平坦化的一个非常弱的版本，但有时已经很有用。

\begin{lemma}
\label{more-algebra-lemma-flatten-on-affine-blowup}
设 $(R, \mathfrak m)$ 为分式域为 $K$ 的局部整环。设 $S$ 为有限型
$R$-代数，$M$ 为有限 $S$-模。对每个支配 $R$ 的赋值环
$A \subset K$，存在理想 $I \subset \mathfrak m$ 与非零元 $a \in I$，使得：
\begin{enumerate}
\item $I$ 有限生成；
\item $A$ 在 $R[\frac{I}{a}]$ 上有中心；
\item $R \to R[\frac{I}{a}]$ 在 $\mathfrak m$ 处的纤维环不为零；以及
\item $S$ 沿 $R \to R[\frac{I}{a}]$ 的严格变换 $S_{I, a}$ 在 $R$ 上平坦且有限表示，
并且 $M$ 沿 $R \to R[\frac{I}{a}]$ 的严格变换 $M_{I, a}$ 在 $R$ 上平坦，
并在 $S_{I, a}$ 上有限表示。
\end{enumerate}
\end{lemma}

\begin{proof}
写成 $S = R[x_1, \ldots, x_n]/J$，并把 $N = S \oplus M$ 视为
$P = R[x_1, \ldots, x_n]$ 上的模。若能在 $S$ 是 $R$ 上多项式代数的情形证明本引理，
则可找到满足 (1)、(2)、(3) 的 $I, a$，使 $N$ 沿
$R \to R[\frac{I}{a}]$ 的严格变换 $N_{I, a}$ 在 $R$ 上平坦，并且作为
$P$ 的严格变换 $P_{I, a}$ 上的模有限表示。因为
$P_{I, a} = R[\frac{I}{a}][x_1, \ldots, x_n]$（略去一个小细节），
可知直和项 $S_{I, a} \subset N_{I, a}$ 在 $R$ 上平坦，并且作为
$R[\frac{I}{a}][x_1, \ldots, x_n]$-模有限表示。因此 $S_{I, a}$ 是有限表示
$R[\frac{I}{a}]$-代数。此外，直和项 $M_{I, a} \subset N_{I, a}$ 在 $R$ 上平坦，
并且作为 $P_{I, a}$-模有限表示，因而作为 $S_{I, a}$-模也有限表示；见
Algebra, Lemma \ref{algebra-lemma-finitely-presented-over-subring}。
这将我们化约到下一段讨论的情形。

\medskip\noindent
假设 $S = R[x_1, \ldots, x_n]$。选取表示
$$
0 \to K \to S^{\oplus r} \to M \to 0.
$$
设 $M_A$ 是 $M \otimes_R A$ 模去其挠子模得到的商；见引理
\ref{more-algebra-lemma-torsion}。则 $M_A$ 是 $S_A = A[x_1, \ldots, x_n]$ 上的有限模。
由引理 \ref{more-algebra-lemma-valuation-ring-torsion-free-flat}，$M_A$ 在 $A$ 上平坦。
由引理 \ref{more-algebra-lemma-flat-finite-type-valuation-ring-finite-presentation}，
$M_A$ 有限表示。因此存在有限多个元
$k_1, \ldots, k_t \in S_A^{\oplus r}$，它们作为 $S_A$-模生成表示
$S_A^{\oplus r} \to M_A$ 的核。对任意满足 (1)、(2)、(3) 的选取
$a \in I \subset \mathfrak m$，用 $M_{I, a}$ 表示 $M$ 沿
$R \to R[\frac{I}{a}]$ 的严格变换。它是
$S_{I, a} = R[\frac{I}{a}][x_1, \ldots, x_n]$ 上的有限模。
由 Algebra, Lemma \ref{algebra-lemma-valuation-ring-colimit-affine-blowups}，有
$A = \colim_{I, a} R[\frac{I}{a}]$。这推出
$S_A = \colim S_{I, a}$ 且
$$
\colim M \otimes_R R[\textstyle{\frac{I}{a}}] = M \otimes_R A
$$
选取 $I, a$ 以及提升 $k_1, \ldots, k_t \in S_{I, a}^{\oplus r}$。
由于 $M_A$ 是 $M \otimes_R A$ 模去挠得到的商，可知
$k_1, \ldots, k_t$ 在 $M \otimes_R A$ 中的像被某个非零元
$\alpha \in A$ 零化。用另一对替换 $I, a$（回忆该余极限是滤过的）后，
可以假设对某个非零 $x \in I^n$ 有 $\alpha = x/a^n$。
于是 $x k_1, \ldots, x k_t$ 在 $M \otimes_R A$ 中映到零。
再次用另一对替换 $I, a$ 后，可以假设对某个非零 $x \in R$，
$x k_1, \ldots, x k_t$ 在 $M \otimes_R R[\frac{I}{a}]$ 中映到零。
最后用 $xI, xa$ 替换 $I, a$，可以假设 $k_1, \ldots, k_t$ 在
$M \otimes_R R[\frac{I}{a}]$ 中映到被 $a$ 的某个幂零化的挠元。
对任意这样的对 $(I, a)$，令
$$
M'_{I, a} = S_{I, a}^{\oplus r}/ \sum S_{I, a}k_j.
$$
因为 $M_A = S_A^{\oplus r}/ \sum S_Ak_j$，可知
$M_A = \colim_{I, a} M'_{I, a}$。此时终于可以应用 Algebra, Lemma
\ref{algebra-lemma-flat-finite-presentation-limit-flat} (3)，得出对某个如上的对
$(I, a)$，$M'_{I, a}$ 平坦。该引理先验上不能应用于严格变换组成的系统
$$
M_{I, a} = (M \otimes_R R[\textstyle{\frac{I}{a}}])/a\text{-幂挠子}
$$
因为转移映射可能不满足该引理的假设。但现在，由于平坦推出无挠
（引理 \ref{more-algebra-lemma-flat-torsion-free}），并且 $M_{I, a}$ 是 $M'_{I, a}$ 的商
（因为我们已安排使元 $k_1, \ldots, k_t$ 在 $M_{I, a}$ 中映到零），
所商去的是被 $a$ 的某个幂零化的挠子模，所以对这样的对还有
$M'_{I, a} = M_{I, a}$，从而得证。
\end{proof}

\begin{lemma}
\label{more-algebra-lemma-blowup-fitting-ideal}
设 $R$ 为环，$M$ 为有限 $R$-模。设 $k \geq 0$ 且
$I = \text{Fit}_k(M)$。对每个 $a \in I$，令 $R' = R[\frac{I}{a}]$，则严格变换
$$
M' = (M \otimes_R R')/a\text{-幂挠子}
$$
满足 $\text{Fit}_k(M') = R'$。
\end{lemma}

\begin{proof}
首先注意 $\text{Fit}_k(M \otimes_R R') = IR' = aR'$。第一个等式由引理
\ref{more-algebra-lemma-fitting-ideal-basics} 的 (3) 得出，第二个等式由 Algebra, Lemma
\ref{algebra-lemma-affine-blowup} 得出。由引理
\ref{more-algebra-lemma-principal-fitting-ideal} 与局部化的正合性，可知对 $R'$ 的每个素理想
$\mathfrak p'$，$M'_{\mathfrak p'}$ 可由 $\leq k$ 个元生成。
例如由引理 \ref{more-algebra-lemma-fitting-ideal-generate-locally}，可得
$\text{Fit}_k(M') = R'$。
\end{proof}

\begin{lemma}
\label{more-algebra-lemma-blowup-fitting-ideal-locally-free}
设 $R$ 为环，$M$ 为有限 $R$-模。设 $k \geq 0$ 且
$I = \text{Fit}_k(M)$。假设对每个 $\mathfrak p \not \in V(I)$，
$M_\mathfrak p$ 都是秩为 $k$ 的自由模。则对每个 $a \in I$，令
$R' = R[\frac{I}{a}]$，严格变换
$$
M' = (M \otimes_R R')/a\text{-幂挠子}
$$
都局部自由且秩为 $k$。
\end{lemma}

\begin{proof}
由引理 \ref{more-algebra-lemma-blowup-fitting-ideal}，有
$\text{Fit}_k(M') = R'$。由引理 \ref{more-algebra-lemma-fitting-ideal-finite-locally-free}，
只须证明 $\text{Fit}_{k - 1}(M') = 0$。回忆
$R' \subset R'_a = R_a$；见 Algebra, Lemma \ref{algebra-lemma-affine-blowup}。
因此只须证明 $\text{Fit}_{k - 1}(M')$ 在 $R'_a = R_a$ 中映到零。
由于显然 $(M')_a = M_a$，而由引理 \ref{more-algebra-lemma-fitting-ideal-basics} 的 (3)，
形成 Fitting 理想与基变换可交换，问题归结为证明
$\text{Fit}_{k - 1}(M_a) = 0$。由假设，$M_a$ 是秩为 $k$ 的有限局部自由模
（见 Algebra, Lemma \ref{algebra-lemma-finite-projective}），故此结论由已引用的引理
\ref{more-algebra-lemma-fitting-ideal-finite-locally-free} 得出。
\end{proof}

\begin{lemma}
\label{more-algebra-lemma-blowup-module}
设 $R$ 为环，$M$ 为有限 $R$-模。设 $f \in R$ 为元，使 $M_f$ 是秩为
$r$ 的有限局部自由模。则存在有限生成理想 $I \subset R$，满足
$V(f) = V(I)$，且对所有 $a \in I$，令 $R' = R[\frac{I}{a}]$，
严格变换
$$
M' = (M \otimes_R R')/a\text{-幂挠子}
$$
局部自由且秩为 $r$。
\end{lemma}

\begin{proof}
选取满射 $R^{\oplus n} \to M$。选取有限子模
$K \subset \Ker(R^{\oplus n} \to M)$，使 $R^{\oplus n}/K \to M$ 在逆转 $f$ 后成为同构。
这是可能的，因为例如由 Algebra, Lemma
\ref{algebra-lemma-finite-projective}，$M_f$ 有限表示。
令 $M_1 = R^{\oplus n}/K$，并假设能对 $M_1$ 证明本引理。
设 $I \subset R$ 为相应的理想。则对 $a \in I$，映射
$$
M_1' = (M_1 \otimes_R R')/a\text{-幂挠子}
\longrightarrow
M' = (M \otimes_R R')/a\text{-幂挠子}
$$
是满射。在 $R'$ 中逆转 $a$ 后，它也是同构，因为
$R'_a = R_f$；见 Algebra, Lemma \ref{algebra-lemma-blowup-in-principal}。
但 $a$ 在 $M'_1$ 上是非零因子，因此所显示的映射是同构。
所以只须在 $M$ 是有限表示 $R$-模的情形证明本引理。

\medskip\noindent
假设 $M$ 是有限表示 $R$-模。则
$J = \text{Fit}_r(M) \subset S$ 是有限生成理想。我们断言 $I = fJ$ 可行。

\medskip\noindent
首先验证 $V(f) = V(I)$。包含关系 $V(f) \subset V(I)$ 是显然的。
反之，若 $f \not \in \mathfrak p$，则由引理
\ref{more-algebra-lemma-fitting-ideal-generate-locally}，$\mathfrak p$ 不是 $V(J)$ 的元。
因此 $\mathfrak p \not \in V(fJ) = V(I)$。

\medskip\noindent
设 $a \in I$，并令 $R' = R[\frac{I}{a}]$。可以对某个 $b \in J$ 写成
$a = fb$。由 Algebra, Lemmas \ref{algebra-lemma-affine-blowup} 与
\ref{algebra-lemma-blowup-add-principal}，有 $J R' = b R'$，且 $b$ 是 $R'$ 中的非零因子。
设 $\mathfrak p' \subset R' = R[\frac{I}{a}]$ 为素理想。则
$JR'_{\mathfrak p'}$ 由 $b$ 生成。由引理 \ref{more-algebra-lemma-principal-fitting-ideal}，
$M'_{\mathfrak p'}$ 可由 $r$ 个元生成。由于 $M'$ 有限，存在
$m_1, \ldots, m_r \in M'$ 与 $g \in R'$、$g \not \in \mathfrak p'$，使相应的映射
$(R')^{\oplus r} \to M'$ 在逆转 $g$ 后成为满射。

\medskip\noindent
最后，考虑理想 $J' = \text{Fit}_{k - 1}(M')$。注意 $J' R'_g$ 由
$m_1, \ldots, m_r$ 之间各关系的系数生成（Fitting 理想与基变换的相容性）。
因此只须证明 $J' = 0$；见引理 \ref{more-algebra-lemma-fitting-ideal-finite-locally-free}。
由于 $R'_a = R_f$（Algebra, Lemma \ref{algebra-lemma-blowup-in-principal}），而
$M'_a = M_f$ 自由且秩为 $r$，可知 $J'_a = 0$。因为 $a$ 是 $R'$ 中的非零因子，
所以 $J' = 0$，于是得证。
\end{proof}













\section{完备化与平坦性}
\label{more-algebra-section-completion-flat}

\noindent
本节讨论一个“大”平坦模的完备化何时平坦。

\begin{lemma}
\label{more-algebra-lemma-ui-completion-direct-sum-into-product}
设 $R$ 为环，$I \subset R$ 为理想，$A$ 为集合。假设 $R$ 是 Noether 环，
且对 $I$ 完备。存在从直和的 $I$-进完备化到直积的规范映射
$$
\left(\bigoplus\nolimits_{\alpha \in A} R\right)^\wedge
\longrightarrow
\prod\nolimits_{\alpha \in A} R
$$
，且该映射是普遍单射。
\end{lemma}

\begin{proof}
依定义，左边的元 $x$ 为 $x = (x_n)$，其中
$x_n = (x_{n, \alpha}) \in \bigoplus\nolimits_{\alpha \in A} R/I^n$，并满足
$x_{n, \alpha} = x_{n + 1, \alpha} \bmod I^n$。由于 $R = R^\wedge$，
对任意 $\alpha$ 存在 $y_\alpha \in R$，使
$x_{n, \alpha} = y_\alpha \bmod I^n$。注意，对每个 $n$，只有有限多个
$\alpha$ 使 $x_{n, \alpha}$ 不为零。反之，给定
$(y_\alpha) \in \prod_\alpha R$，若对每个 $n$，只有有限多个 $\alpha$ 使
$y_{\alpha} \bmod I^n$ 不为零，则它定义了左边的一个元。
所以可以把左边的元看作无限“收敛和”
$\sum_\alpha y_\alpha$，其中 $y_\alpha \in R$，且对每个 $n$，
只有有限多个 $y_\alpha$ 模 $I^n$ 不为零。所显示的映射把此元映到直积中的元
$(y_\alpha)$。特别地，该映射是单射。

\medskip\noindent
设 $Q$ 为有限 $R$-模。我们必须证明映射
$$
Q \otimes_R \left(\bigoplus\nolimits_{\alpha \in A} R\right)^\wedge
\longrightarrow
Q \otimes_R \left(\prod\nolimits_{\alpha \in A} R\right)
$$
是单射；见 Algebra, Theorem
\ref{algebra-theorem-universally-exact-criteria}。选取表示
$R^{\oplus k} \to R^{\oplus m} \to Q \to 0$，并用
$q_1, \ldots, q_m \in Q$ 表示 $Q$ 的相应生成元。由 Artin-Rees
（Algebra, Lemma \ref{algebra-lemma-Artin-Rees}），存在常数 $c$，使
$\Im(R^{\oplus k} \to R^{\oplus m}) \cap (I^N)^{\oplus m}
\subset \Im((I^{N - c})^{\oplus k} \to R^{\oplus m})$。考虑图
$$
\xymatrix{
\bigoplus_{l = 1}^k \left(\bigoplus\nolimits_{\alpha \in A} R\right)^\wedge
\ar[r] \ar[d] &
\bigoplus_{j = 1}^m \left(\bigoplus\nolimits_{\alpha \in A} R\right)^\wedge
\ar[r] \ar[d] &
Q \otimes_R \left(\bigoplus\nolimits_{\alpha \in A} R\right)^\wedge
\ar[r] \ar[d] &
0 \\
\bigoplus_{l = 1}^k \left(\prod\nolimits_{\alpha \in A} R\right)
\ar[r] &
\bigoplus_{j = 1}^m \left(\prod\nolimits_{\alpha \in A} R\right)
\ar[r] &
Q \otimes_R \left(\prod\nolimits_{\alpha \in A} R\right)
\ar[r] &
0
}
$$
其两行正合。选取元
$\sum_j \sum_\alpha y_{j, \alpha}$，它属于
$\bigoplus_{j = 1, \ldots, m}
\left(\bigoplus\nolimits_{\alpha \in A} R\right)^\wedge$。
若此元在模
$Q \otimes_R \left(\prod\nolimits_{\alpha \in A} R\right)$ 中映到零，
则特别地，对每个 $\alpha$，有
$\sum_j q_j \otimes y_{j, \alpha} = 0$ 在 $Q$ 中成立。
因此可找到元
$(z_{1, \alpha}, \ldots, z_{k, \alpha}) \in \bigoplus_{l = 1, \ldots, k} R$，
它映到
$(y_{1, \alpha}, \ldots, y_{m, \alpha}) \in \bigoplus_{j = 1, \ldots, m} R$。
此外，若对 $j = 1, \ldots, m$ 有 $y_{j, \alpha} \in I^{N_\alpha}$，
则可以假设对 $l = 1, \ldots, k$ 有
$z_{l, \alpha} \in I^{N_\alpha - c}$。因此和
$\sum_l \sum_\alpha z_{l, \alpha}$ 是“收敛的”，并定义元
$\bigoplus_{l = 1, \ldots, k}
\left(\bigoplus\nolimits_{\alpha \in A} R\right)^\wedge$，
该元映到我们开始所取的元 $\sum_j \sum_\alpha y_{j, \alpha}$。
因此右边的竖直箭头是单射，于是得证。
\end{proof}

\noindent
下一引理也可以由下面的引理 \ref{more-algebra-lemma-limit-flat} 推出。

\begin{lemma}
\label{more-algebra-lemma-completed-direct-sum-flat}
设 $R$ 为环，$I \subset R$ 为理想，$A$ 为集合。假设 $R$ 是 Noether 环。
完备化 $(\bigoplus\nolimits_{\alpha \in A} R)^\wedge$ 是平坦 $R$-模。
\end{lemma}

\begin{proof}
用 $R^\wedge$ 表示 $R$ 关于 $I$ 的完备化。由 Algebra, Lemma
\ref{algebra-lemma-completion-flat}，$R \to R^\wedge$ 平坦，所以只须证明
$(\bigoplus\nolimits_{\alpha \in A} R)^\wedge$ 是平坦 $R^\wedge$-模
（使用 Algebra, Lemma \ref{algebra-lemma-composition-flat}）。因为
$$
(\bigoplus\nolimits_{\alpha \in A} R)^\wedge
=
(\bigoplus\nolimits_{\alpha \in A} R^\wedge)^\wedge
$$
可以用 $R^\wedge$ 替换 $R$，并假设 $R$ 对 $I$ 完备
（见 Algebra, Lemma \ref{algebra-lemma-completion-complete}）。
在这种情形下，引理 \ref{more-algebra-lemma-ui-completion-direct-sum-into-product} 告诉我们，映射
$(\bigoplus\nolimits_{\alpha \in A} R)^\wedge \to \prod_{\alpha \in A} R$
是普遍单射。因此由 Algebra, Lemma \ref{algebra-lemma-ui-flat-domain}，
只须证明 $\prod_{\alpha \in A} R$ 平坦。由 Algebra, Proposition
\ref{algebra-proposition-characterize-coherent}（以及 Algebra, Lemma
\ref{algebra-lemma-Noetherian-coherent}），$\prod_{\alpha \in A} R$ 平坦。
\end{proof}

\begin{lemma}
\label{more-algebra-lemma-tor-strictly-pro-zero}
\begin{reference}
这是 \cite[Lemma 9.9]{quillenhomology}；注意，作者在陈述中遗漏了“严格”一词，
尽管显然原意包含该词。
\end{reference}
设 $A$ 为 Noether 环，$I$ 为 $A$ 的理想，$M$ 为有限 $A$-模。
对每个 $p > 0$，存在 $c > 0$，使对所有 $n \geq c$，映射
$\text{Tor}_p^A(M, A/I^n) \to \text{Tor}_p^A(M, A/I^{n - c})$ 为零。
\end{lemma}

\begin{proof}
先证明 $p = 1$ 的情形。选取短正合列
$0 \to K \to A^{\oplus t} \to M \to 0$。则
$\text{Tor}_1^A(M, A/I^n) = K \cap (I^n)^{\oplus t}/I^nK$。
由 Artin-Rees（Algebra, Lemma \ref{algebra-lemma-Artin-Rees}），
存在常数 $c \geq 0$，使对 $n \geq c$ 有
$K \cap (I^n)^{\oplus t} \subset I^{n - c}K$。这证明了 $p = 1$ 的结论。
对 $p > 1$，有
$\text{Tor}_p^A(M, A/I^n) = \text{Tor}^A_{p - 1}(K, A/I^n)$。
因此本引理由归纳得证。
\end{proof}

\begin{lemma}
\label{more-algebra-lemma-limit-flat}
设 $A$ 为 Noether 环，$I$ 为 $A$ 的理想。设 $(M_n)$ 为 $A$-模的逆系统，且：
\begin{enumerate}
\item $M_n$ 是平坦 $A/I^n$-模；
\item $M_{n + 1} \to M_n$ 是满射。
\end{enumerate}
则 $M = \lim M_n$ 是平坦 $A$-模，并且对每个有限 $A$-模 $Q$，有
$Q \otimes_A M = \lim Q \otimes_A M_n$。
\end{lemma}

\begin{proof}
首先对每个有限 $A$-模 $Q$ 证明
$Q \otimes_A M = \lim Q \otimes_A M_n$。选取由有限自由 $A$-模 $F_i$ 组成的解消
$F_2 \to F_1 \to F_0 \to Q \to 0$。则
$$
F_2 \otimes_A M_n \to F_1 \otimes_A M_n \to F_0 \otimes_A M_n
$$
是一个链复形，其 $0$ 次同调为 $Q \otimes_A M_n$，而其 $1$ 次同调为
$$
\text{Tor}_1^A(Q, M_n) = \text{Tor}_1^A(Q, A/I^n) \otimes_{A/I^n} M_n
$$
，因为 $M_n$ 在 $A/I^n$ 上平坦。由引理
\ref{more-algebra-lemma-tor-strictly-pro-zero}，该系统本质上是值为 $0$ 的常系统。
由 Homology, Lemma \ref{homology-lemma-apply-Mittag-Leffler-again}，得
$\lim Q \otimes_A M_n =
\Coker(\lim F_1 \otimes_A M_n \to \lim F_0 \otimes_A M_n)$。
由于 $F_i$ 有限自由，这等于
$\Coker(F_1 \otimes_A M \to F_0 \otimes_A M) = Q \otimes_A M$。

\medskip\noindent
其次，设 $Q \to Q'$ 是有限 $A$-模之间的单射。我们必须证明
$Q \otimes_A M \to Q' \otimes_A M$ 是单射
（Algebra, Lemma \ref{algebra-lemma-flat}）。由上面的结果，有
$$
\Ker(Q \otimes_A M \to Q' \otimes_A M) =
\Ker(\lim Q \otimes_A M_n \to \lim Q' \otimes_A M_n).
$$
对每个 $n$，有正合列
$$
\text{Tor}_1^A(Q', M_n) \to \text{Tor}_1^A(Q'', M_n) \to
Q \otimes_A M_n \to Q' \otimes_A M_n
$$
其中 $Q'' = \Coker(Q \to Q')$。上面已经看到，这些 Tor 的逆系统本质上是值为
$0$ 的常系统。由 Homology, Lemma
\ref{homology-lemma-apply-Mittag-Leffler-again}，最右边映射的逆极限是单射。
\end{proof}

\begin{lemma}
\label{more-algebra-lemma-flat-after-completion}
设 $R$ 为环，$I \subset R$ 为理想，$M$ 为 $R$-模。假设：
\begin{enumerate}
\item $I$ 有限生成；
\item $R/I$ 是 Noether 环；
\item $M/IM$ 在 $R/I$ 上平坦；
\item $\text{Tor}_1^R(M, R/I) = 0$。
\end{enumerate}
则 $I$-进完备化 $R^\wedge$ 是 Noether 环，且 $M^\wedge$ 在
$R^\wedge$ 上平坦。
\end{lemma}

\begin{proof}
由 Algebra, Lemma \ref{algebra-lemma-what-does-it-mean}，对所有 $n$，
模 $M/I^nM$ 在 $R/I^n$ 上平坦。由 Algebra, Lemma
\ref{algebra-lemma-hathat-finitely-generated}，有：(a) $R^\wedge$ 与 $M^\wedge$
都是 $I$-进完备的；(b) 对所有 $n$，有
$R/I^n = R^\wedge/I^nR^\wedge$。由 Algebra, Lemma
\ref{algebra-lemma-completion-Noetherian}，环 $R^\wedge$ 是 Noether 环。
应用引理 \ref{more-algebra-lemma-limit-flat}，可得
$M^\wedge = \lim M/I^nM$ 作为 $R^\wedge$-模平坦。
\end{proof}









\section{Koszul 复形}
\label{more-algebra-section-koszul}

\noindent
我们如下定义 Koszul 复形。

\begin{definition}
\label{more-algebra-definition-koszul}
设 $R$ 为环，$\varphi : E \to R$ 为 $R$-模映射。与 $\varphi$ 相关联的
{\it Koszul 复形} $K_\bullet(\varphi)$ 是如下定义的交换微分分次代数：
\begin{enumerate}
\item 其底分次代数是外代数
$K_\bullet(\varphi) = \wedge(E)$；
\item 微分 $d : K_\bullet(\varphi) \to K_\bullet(\varphi)$ 是唯一的导子，满足对所有
$e \in E = K_1(\varphi)$ 有 $d(e) = \varphi(e)$。
\end{enumerate}
\end{definition}

\noindent
明确地，若 $e_1 \wedge \ldots \wedge e_n$ 是 $K_\bullet(\varphi)$ 中次数为 $n$
的生成元之一，则
$$
d(e_1 \wedge \ldots \wedge e_n) =
\sum\nolimits_{i = 1, \ldots, n} (-1)^{i + 1}
\varphi(e_i)e_1 \wedge \ldots \wedge \widehat{e_i} \wedge \ldots \wedge e_n.
$$
容易验证，这在张量代数上给出良定义的导子，该导子零化
$e \otimes e$，从而经过外代数分解。

\medskip\noindent
我们常假设 $E$ 是有限自由模，设 $E = R^{\oplus n}$。
在这种情形下，映射 $\varphi$ 由一列元 $f_1, \ldots, f_n \in R$ 给出。

\begin{definition}
\label{more-algebra-definition-koszul-complex}
设 $R$ 为环，$f_1, \ldots, f_r \in R$。{\it $f_1, \ldots, f_r$ 上的 Koszul 复形}
是与映射 $(f_1, \ldots, f_r) : R^{\oplus r} \to R$ 相关联的 Koszul 复形。
记为 $K_\bullet(f_\bullet)$、$K_\bullet(f_1, \ldots, f_r)$、
$K_\bullet(R, f_1, \ldots, f_r)$ 或 $K_\bullet(R, f_\bullet)$。
\end{definition}

\noindent
当然，若 $E$ 有限局部自由，则 $K_\bullet(\varphi)$ 在 $\Spec(R)$ 上局部同构于
Koszul 复形 $K_\bullet(f_1, \ldots, f_r)$。该复形具有许多有趣的形式性质。

\begin{lemma}
\label{more-algebra-lemma-functorial}
设 $\varphi : E \to R$ 与 $\varphi' : E' \to R$ 为 $R$-模映射。设
$\psi : E \to E'$ 为满足 $\varphi' \circ \psi = \varphi$ 的 $R$-模映射。
则 $\psi$ 诱导微分分次代数同态
$K_\bullet(\varphi) \to K_\bullet(\varphi')$。
\end{lemma}

\begin{proof}
这直接由定义得出。
\end{proof}

\begin{lemma}
\label{more-algebra-lemma-change-basis}
设 $f_1, \ldots, f_r \in R$ 为一列元，$(x_{ij})$ 为系数在 $R$ 中的可逆
$r \times r$-矩阵。则复形 $K_\bullet(f_\bullet)$ 与
$$
K_\bullet(\sum x_{1j}f_j, \sum x_{2j}f_j, \ldots, \sum x_{rj}f_j)
$$
同构。
\end{lemma}

\begin{proof}
令 $g_i = \sum x_{ij}f_j$。矩阵 $(x_{ji})$ 给出同构
$x : R^{\oplus r} \to R^{\oplus r}$，使
$(g_1, \ldots, g_r) = (f_1, \ldots, f_r) \circ x$。
因此结论由引理 \ref{more-algebra-lemma-functorial} 中所述 Koszul 复形的函子性得出。
\end{proof}

\begin{lemma}
\label{more-algebra-lemma-homotopy-koszul-abstract}
设 $R$ 为环，$\varphi : E \to R$ 为 $R$-模映射。设 $e \in E$ 在 $R$ 中的像为
$f = \varphi(e)$。则作为 $K_\bullet(\varphi)$ 的自同态，有
$$
f = de + ed
$$
。
\end{lemma}

\begin{proof}
这是因为 $d(ea) = d(e)a - ed(a) = fa - ed(a)$。
\end{proof}

\begin{lemma}
\label{more-algebra-lemma-homotopy-koszul}
设 $R$ 为环，$f_1, \ldots, f_r \in R$ 为一列元。在 $K_\bullet(f_\bullet)$ 上乘以
$f_i$ 的映射零同伦，特别地，上同调模 $H_i(K_\bullet(f_\bullet))$ 被理想
$(f_1, \ldots, f_r)$ 零化。
\end{lemma}

\begin{proof}
这是引理 \ref{more-algebra-lemma-homotopy-koszul-abstract} 的特殊情形。
\end{proof}

\noindent
在 Derived Categories, Section \ref{derived-section-cones} 中，我们定义了上链复形态射的锥。
链复形态射 $f : A_\bullet \to B_\bullet$ 的锥记为 $C(f)_\bullet$。复形
$C(f)_\bullet$ 由下式给出：
$C(f)_n = B_n \oplus A_{n - 1}$，微分为
\begin{equation}
\label{more-algebra-equation-differential-cone}
d_{C(f), n} =
\left(
\begin{matrix}
d_{B, n} & f_{n - 1} \\
0 & -d_{A, n - 1}
\end{matrix}
\right)
\end{equation}
它带有规范复形态射 $i : B_\bullet \to C(f)_\bullet$ 与
$p : C(f)_\bullet \to A_\bullet[-1]$，它们由显然的映射
$B_n \to C(f)_n \to A_{n - 1}$ 诱导。

\begin{lemma}
\label{more-algebra-lemma-cone-koszul-abstract}
设 $R$ 为环，$\varphi : E \to R$ 为 $R$-模映射，$f \in R$。令
$E' = E \oplus R$，并定义 $\varphi' : E' \to R$：在 $E$ 上等于 $\varphi$，
在 $R$ 上等于乘以 $f$。复形 $K_\bullet(\varphi')$ 同构于复形映射
$$
f :
K_\bullet(\varphi)
\longrightarrow
K_\bullet(\varphi).
$$
的锥。
\end{lemma}

\begin{proof}
用 $e_0 \in E'$ 表示元 $1 \in R \subset R \oplus E$。由上面对锥的定义，有
$$
C(f)_n = K_n(\varphi) \oplus K_{n - 1}(\varphi) =
\wedge^n(E) \oplus \wedge^{n - 1}(E) = \wedge^n(E')
$$
其中在最后一个 $=$ 中，把 $(0, e_1 \wedge \ldots \wedge e_{n - 1})$
映到 $\wedge^n(E')$ 中的 $e_0 \wedge e_1 \wedge \ldots \wedge e_{n - 1}$。
计算表明此同构与微分相容。具体地，对第一个直和项中的元，这是清楚的，
因为 $\varphi'|_E = \varphi$，而 $d_{C(f)}$ 限制在第一个直和项上就是
$d_{K_\bullet(\varphi)}$。另一方面，若 $e_1 \wedge \ldots \wedge e_{n - 1}$
在第二个直和项中，则
$$
d_{C(f)}(0, e_1 \wedge \ldots \wedge e_{n - 1}) =
fe_1 \wedge \ldots \wedge e_{n - 1}
- d_{K_\bullet(\varphi)}(e_1 \wedge \ldots \wedge e_{n - 1})
$$
而另一方面
\begin{align*}
& d_{K_\bullet(\varphi')}(0, e_0 \wedge e_1 \wedge \ldots \wedge e_{n - 1}) \\
& =
\sum\nolimits_{i = 0, \ldots, n - 1}
(-1)^i \varphi'(e_i)e_0 \wedge \ldots \wedge \widehat{e_i}
\wedge \ldots \wedge e_{n - 1} \\
& =
fe_1 \wedge \ldots \wedge e_{n - 1} +
\sum\nolimits_{i = 1, \ldots, n - 1}
(-1)^i \varphi(e_i)e_0 \wedge \ldots \wedge \widehat{e_i}
\wedge \ldots \wedge e_{n - 1} \\
& =
fe_1 \wedge \ldots \wedge e_{n - 1} -
e_0 \left(\sum\nolimits_{i = 1, \ldots, n - 1}
(-1)^{i + 1} \varphi(e_i)e_1 \wedge \ldots \wedge \widehat{e_i}
\wedge \ldots \wedge e_{n - 1}\right)
\end{align*}
这正是上一计算结果的像。
\end{proof}

\begin{lemma}
\label{more-algebra-lemma-cone-koszul}
设 $R$ 为环，$f_1, \ldots, f_r$ 为 $R$ 中的一列元。复形
$K_\bullet(f_1, \ldots, f_r)$ 同构于复形映射
$$
f_r :
K_\bullet(f_1, \ldots, f_{r - 1})
\longrightarrow
K_\bullet(f_1, \ldots, f_{r - 1}).
$$
的锥。
\end{lemma}

\begin{proof}
这是引理 \ref{more-algebra-lemma-cone-koszul-abstract} 的特殊情形。
\end{proof}

\begin{lemma}
\label{more-algebra-lemma-cone-squared}
设 $R$ 为环，$A_\bullet$ 为 $R$-模的复形，$f, g \in R$。设
$C(f)_\bullet$ 为 $f : A_\bullet \to A_\bullet$ 的锥。类似定义
$C(g)_\bullet$ 与 $C(fg)_\bullet$。则 $C(fg)_\bullet$ 同伦等价于某个映射
$$
C(f)_\bullet[1] \longrightarrow C(g)_\bullet
$$
的锥。
\end{lemma}

\begin{proof}
先在 $A_\bullet$ 是仅在次数 $0$ 放置 $R$ 所得的复形时证明此结论。
在这种情形下，复形 $C(f)_\bullet$ 是
$$
\ldots \to 0 \to R \xrightarrow{f} R \to 0 \to \ldots
$$
其中 $R$ 位于（同调）次数 $1$ 与 $0$。我们使用的复形映射是
$$
\xymatrix{
0 \ar[r] \ar[d] &
0 \ar[r] \ar[d] &
R \ar[r]^f \ar[d]^1 &
R \ar[r] \ar[d] & 0 \ar[d] \\
0 \ar[r] & R \ar[r]^g & R \ar[r] & 0 \ar[r] & 0
}
$$
它的锥是由位于次数 $1$ 与 $0$ 的 $R^{\oplus 2}$ 组成的链复形，
其微分（\ref{more-algebra-equation-differential-cone}）为
$$
\left(
\begin{matrix}
g & 1 \\
0 & -f
\end{matrix}
\right) :
R^{\oplus 2} \longrightarrow R^{\oplus 2}
$$
为证明此链复形同伦等价于 $C(fg)_\bullet$，也就是等价于
$R \xrightarrow{fg} R$，考虑复形映射
$$
\xymatrix{
R \ar[d]_{(1, -g)} \ar[r]_{fg} & R \ar[d]^{(0, 1)} \\
R^{\oplus 2} \ar[r] & R^{\oplus 2}
}
\quad\quad
\xymatrix{
R^{\oplus 2} \ar[d]_{(1, 0)} \ar[r] & R^{\oplus 2} \ar[d]^{(f, 1)} \\
R \ar[r]^{fg} & R
}
$$
其记号的含义是显然的。这两个映射在一个方向上的复合是
$C(fg)_\bullet$ 上的恒等映射，但在另一方向上的复合不是恒等映射。
我们略去写出所需同伦的过程。

\medskip\noindent
为说明一般情形下结果成立，我们使用复形范畴上的函子
$K_\bullet \mapsto \text{Tot}(A_\bullet \otimes_R K_\bullet)$，它与同伦及锥相容。
然后把 $C(f)_\bullet$ 与 $C(g)_\bullet$ 写成双复形
$$
(R \xrightarrow{f} R) \otimes_R A_\bullet
\quad\text{且}\quad
(R \xrightarrow{g} R) \otimes_R A_\bullet
$$
的总复形，并由此从上面讨论的特殊情形推出结论。略去一些细节。
\end{proof}

\begin{lemma}
\label{more-algebra-lemma-koszul-mult-abstract}
设 $R$ 为环，$\varphi : E \to R$ 为 $R$-模映射，$f, g \in R$。令
$E' = E \oplus R$，并定义 $\varphi'_f, \varphi'_g, \varphi'_{fg} : E' \to R$：
在 $E$ 上都等于 $\varphi$，在 $R$ 上分别等于乘以 $f, g, fg$。
复形 $K_\bullet(\varphi'_{fg})$ 同伦等价于某个复形映射
$$
K_\bullet(\varphi'_f)[1]
\longrightarrow
K_\bullet(\varphi'_g).
$$
的锥。
\end{lemma}

\begin{proof}
由引理 \ref{more-algebra-lemma-cone-koszul-abstract}，复形 $K_\bullet(\varphi'_f)$
同构于在 $K_\bullet(\varphi)$ 上乘以 $f$ 的映射的锥；另外两种情形同理。
所以本引理由引理 \ref{more-algebra-lemma-cone-squared} 得出。
\end{proof}

\begin{lemma}
\label{more-algebra-lemma-koszul-mult}
设 $R$ 为环，$f_1, \ldots, f_{r - 1}$ 为 $R$ 中的一列元，$f, g \in R$。
复形 $K_\bullet(f_1, \ldots, f_{r - 1}, fg)$ 同伦等价于某个复形映射
$$
K_\bullet(f_1, \ldots, f_{r - 1}, f)[1]
\longrightarrow
K_\bullet(f_1, \ldots, f_{r - 1}, g)
$$
的锥。
\end{lemma}

\begin{proof}
这是引理 \ref{more-algebra-lemma-koszul-mult-abstract} 的特殊情形。
\end{proof}

\begin{lemma}
\label{more-algebra-lemma-join-sequences-koszul-complex}
设 $R$ 为环，$f_1, \ldots, f_r$、$g_1, \ldots, g_s$ 为 $R$ 中的元。
则存在 Koszul 复形的同构
$$
K_\bullet(R, f_1, \ldots, f_r, g_1, \ldots, g_s) =
\text{Tot}(K_\bullet(R, f_1, \ldots, f_r) \otimes_R
K_\bullet(R, g_1, \ldots, g_s)).
$$
\end{lemma}

\begin{proof}
略。提示：若 $K_\bullet(R, f_1, \ldots, f_r)$ 作为微分分次代数由
$x_1, \ldots, x_r$ 生成，且 $\text{d}(x_i) = f_i$，而
$K_\bullet(R, g_1, \ldots, g_s)$ 作为微分分次代数由
$y_1, \ldots, y_s$ 生成，且 $\text{d}(y_j) = g_j$，则可以把
$K_\bullet(R, f_1, \ldots, f_r, g_1, \ldots, g_s)$ 看作由元列
$x_1, \ldots, x_r, y_1, \ldots, y_s$ 生成的微分分次代数，其中
$\text{d}(x_i) = f_i$ 且 $\text{d}(y_j) = g_j$。
\end{proof}






\section{扩张交错 {\v C}ech 复形}
\label{more-algebra-section-alternating-cech}

\noindent
设 $R$ 为环，$f_1, \ldots, f_r \in R$。$R$ 的扩张交错 {\v C}ech 复形是上链复形
$$
R \to \bigoplus\nolimits_{i_0} R_{f_{i_0}} \to
\bigoplus\nolimits_{i_0 < i_1} R_{f_{i_0}f_{i_1}} \to
\ldots \to R_{f_1\ldots f_r}
$$
其中 $R$ 位于次数 $0$，项 $\bigoplus_{i_0} R_{f_{i_0}}$ 位于次数 $1$，依此类推。
各映射定义如下：
\begin{enumerate}
\item 映射 $R \to \bigoplus\nolimits_{i_0} R_{f_{i_0}}$ 由规范映射
$R \to R_{f_{i_0}}$ 给出。
\item 给定 $1 \leq i_0 < \ldots < i_{p + 1} \leq r$ 与
$0 \leq j \leq p + 1$，有规范局部化映射
$$
R_{f_{i_0} \ldots \hat f_{i_j} \ldots f_{i_{p + 1}}} \to
R_{f_{i_0} \ldots f_{i_{p + 1}}}
$$
\item 微分使用 (2) 中的规范映射，其符号为 $(-1)^j$。
\end{enumerate}
若 $M$ 是任意 $R$-模，则 $M$ 的扩张交错 {\v C}ech 复形是按同样方式构造的上链复形
$$
M \to \bigoplus\nolimits_{i_0} M_{f_{i_0}} \to
\bigoplus\nolimits_{i_0 < i_1} M_{f_{i_0}f_{i_1}} \to
\ldots \to M_{f_1\ldots f_r}
$$
其中像上面一样，$M$ 位于次数 $0$。

\begin{lemma}
\label{more-algebra-lemma-extended-alternating-is-complex}
上面定义的扩张交错 {\v C}ech 复形是 $R$-模的复形。
\end{lemma}

\begin{proof}
略。
\end{proof}

\begin{lemma}
\label{more-algebra-lemma-extended-alternating-form-module}
设 $R$ 为环，$f_1, \ldots, f_r \in R$，$M$ 为 $R$-模。$M$ 的扩张交错
{\v C}ech 复形是 $M$ 与 $R$ 的扩张交错 {\v C}ech 复形在 $R$ 上的张量积。
\end{lemma}

\begin{proof}
略。
\end{proof}

\begin{lemma}
\label{more-algebra-lemma-extended-alternating-base-change}
设 $R$ 为环，$f_1, \ldots, f_r \in R$，$M$ 为 $R$-模。设 $R \to S$ 为环映射，
用 $g_1, \ldots, g_r \in S$ 表示 $f_1, \ldots, f_r$ 的像，并令
$N = M \otimes_R S$。使用 $S$、$g_1, \ldots, g_r$ 与 $N$ 构造的扩张交错
{\v C}ech 复形，是 $M$ 的扩张交错 {\v C}ech 复形与 $S$ 在 $R$ 上的张量积。
\end{lemma}

\begin{proof}
略。
\end{proof}

\begin{lemma}
\label{more-algebra-lemma-extended-alternating-homotopy-zero}
设 $R$ 为环，$f_1, \ldots, f_r \in R$，$M$ 为 $R$-模。若存在
$i \in \{1, \ldots, r\}$ 使 $f_i$ 为单位元，则 $M$ 的扩张交错 {\v C}ech 复形与
$0$ 同伦等价。
\end{lemma}

\begin{proof}
我们将使用如下记号：$M$ 的扩张交错 {\v C}ech 复形中次数为
$p + 1$ 的上链 $x$ 为 $x = (x_{i_0 \ldots i_p})$，其中
$x_{i_0 \ldots i_p}$ 属于 $M_{f_{i_0} \ldots f_{i_p}}$。用此记号，有
$$
d(x)_{i_0 \ldots i_{p + 1}} =
\sum\nolimits_j (-1)^j x_{i_0 \ldots \hat i_j \ldots i_{p + 1}}
$$
作为同伦，使用映射
$$
h : \text{cochains of degree }p + 2 \to \text{cochains of degree }p + 1
$$
它由下列规则给出：
$$
h(x)_{i_0 \ldots i_p} = 0 \text{ 若 } i \in \{i_0, \ldots, i_p\}
\text{ 且 }
h(x)_{i_0 \ldots i_p} = 
(-1)^j x_{i_0 \ldots i_j i i_{j + 1} \ldots i_p} \text{ if not}
$$
在第二种情形中，$j$ 是满足 $i_j < i < i_{j + 1}$ 的唯一指标；
另外，因为 $f_i$ 是单位元，有等式
$$
M_{f_{i_0} \ldots f_{i_p}} =
M_{f_{i_0} \ldots f_{i_j} f_i f_{i_{j + 1}} \ldots f_{i_p}}
$$
借助此等式，可以把
$(-1)^j x_{i_0 \ldots i_j i i_{j + 1} \ldots i_p}$
理解为 $M_{f_{i_0} \ldots f_{i_p}}$ 的元。我们将通过计算证明
$d h + h d$ 等于恒等映射的负，从而完成证明。
为此，固定次数为 $p + 1$ 的上链 $x$，并令
$1 \leq i_0 < \ldots < i_p \leq r$。

\medskip\noindent
情形 I：$i \in \{i_0, \ldots, i_p\}$。设 $i = i_t$。则
$h(d(x))_{i_0 \ldots i_p} = 0$。另一方面，有
$$
d(h(x))_{i_0 \ldots i_p} =
\sum (-1)^j h(x)_{i_0 \ldots \hat i_j \ldots i_p} =
(-1)^t h(x)_{i_0 \ldots \hat i \ldots i_p} =
(-1)^t (-1)^{t - 1} x_{i_0 \ldots i_p}
$$
因此如所需，$(dh + hd)(x)_{i_0 \ldots i_p} = -x_{i_0 \ldots i_p}$。

\medskip\noindent
情形 II：$i \not \in \{i_0, \ldots, i_p\}$。设 $j$ 满足 $i_j < i < i_{j + 1}$。则
\begin{align*}
h(d(x))_{i_0 \ldots i_p}
& =
(-1)^j d(x)_{i_0 \ldots i_j i i_{j + 1} \ldots i_p} \\
& =
\sum\nolimits_{j' \leq j} (-1)^{j + j'}
x_{i_0 \ldots \hat i_{j'} \ldots i_j i i_{j + 1} \ldots i_p} -
x_{i_0 \ldots i_p} \\
&
+ \sum\nolimits_{j' > j} (-1)^{j + j' + 1}
x_{i_0 \ldots i_j i i_{j + 1} \ldots \hat i_{j'} \ldots i_p}
\end{align*}
另一方面，
\begin{align*}
d(h(x))_{i_0 \ldots i_p}
& =
\sum\nolimits_{j'} (-1)^{j'} h(x)_{i_0 \ldots \hat i_{j'} \ldots i_p} \\
& =
\sum\nolimits_{j' \leq j} (-1)^{j' + j - 1}
x_{i_0 \ldots \hat i_{j'} \ldots i_j i i_{j + 1} \ldots i_p} \\
& +
\sum\nolimits_{j' > j} (-1)^{j' + j}
x_{i_0 \ldots i_j i i_{j + 1} \ldots \hat i_{j'} \ldots i_p}
\end{align*}
将它们相加，如所需，得
$(dh + hd)(x)_{i_0 \ldots i_p} = - x_{i_0 \ldots i_p}$。
\end{proof}

\begin{lemma}
\label{more-algebra-lemma-extended-alternating-torsion}
设 $R$ 为环，$f_1, \ldots, f_r \in R$，$M$ 为 $R$-模。设 $H^q$ 为 $M$ 的扩张交错
{\v C}ech 复形的第 $q$ 个上同调模。则
\begin{enumerate}
\item 若 $q \not \in [0, r]$，则 $H^q = 0$；
\item 对 $x \in H^i$，存在 $n \geq 1$，使对 $i = 1, \ldots, r$ 有 $f_i^n x = 0$；
\item $H^q$ 的支撑包含于 $V(f_1, \ldots, f_r)$；
\item 若存在 $f \in (f_1, \ldots, f_r)$，它在 $M$ 上可逆地作用，则 $H^q = 0$。
\end{enumerate}
\end{lemma}

\begin{proof}
(1) 由扩张交错 {\v C}ech 复形在次数 $< 0$ 与 $> r$ 上为零这一事实得出。
为证明 (2)，只须证明对每个 $i$，存在 $n \geq 1$ 使 $f_i^n x = 0$。
为此，只须证明 $(H^q)_{f_i} = 0$。因为局部化正合，
$(H^q)_{f_i}$ 是 $M$ 的扩张交错复形在 $f_i$ 处局部化的第 $q$ 个上同调模。
由引理 \ref{more-algebra-lemma-extended-alternating-base-change}，该局部化是 $M_{f_i}$ 在 $R_{f_i}$ 上关于
$f_1, \ldots, f_r$ 在 $R_{f_i}$ 中的像所得的扩张交错 {\v C}ech 复形。
因此问题化约为：若 $f_i$ 可逆，则 $H^q$ 为零；这由引理
\ref{more-algebra-lemma-extended-alternating-homotopy-zero} 得出。(3) 由刚才证明的对所有 $i$
都有 $(H^q)_{f_i} = 0$ 得出。为证明 (4)，注意此时 $f$ 在 $H^q$ 上可逆地作用，
而由 (3)，$H^q$ 支撑在 $V(f)$ 上。这强制 $H^q$ 为零（略去一个小细节）。
\end{proof}

\begin{lemma}
\label{more-algebra-lemma-extended-alternating-Cech-is-colimit-koszul}
设 $R$ 为环，$f_1, \ldots, f_r \in R$。扩张交错 {\v C}ech 复形
$$
R \to \bigoplus\nolimits_{i_0} R_{f_{i_0}} \to
\bigoplus\nolimits_{i_0 < i_1} R_{f_{i_0}f_{i_1}} \to
\ldots \to R_{f_1\ldots f_r}
$$
是 Koszul 复形 $K(R, f_1^n, \ldots, f_r^n)$ 的余极限；精确陈述见证明。
\end{lemma}

\begin{proof}
我们强烈建议读者自行证明此结论。用 $K(R, f_1^n, \ldots, f_r^n)$
表示定义 \ref{more-algebra-definition-koszul-complex} 的 Koszul 复形，并将其视为位于次数
$0, \ldots, r$ 的上链复形。因此有
$$
K(R, f_1^n, \ldots, f_r^n) :
0 \to \wedge^r(R^{\oplus r}) \to
\wedge^{r - 1}(R^{\oplus r}) \to \ldots \to
R^{\oplus r} \to R \to 0
$$
其中项 $\wedge^r(R^{\oplus r})$ 位于次数 $0$。设
$e^n_1, \ldots, e^n_r$ 为 $R^{\oplus r}$ 的标准基。则元
$e^n_{j_1} \wedge \ldots \wedge e^n_{j_{r - p}}$，其中
$1 \leq j_1 < \ldots < j_{r - p} \leq r$，构成 Koszul 复形次数 $p$ 项的一组基。
此外，由 Section \ref{more-algebra-section-koszul} 中对 Koszul 复形的构造，注意
$$
d(e^n_{j_1} \wedge \ldots \wedge e^n_{j_{r - p}}) =
\sum (-1)^{a + 1} f_{j_a}^n e^n_{j_1} \wedge \ldots \wedge \hat e^n_{j_a} \wedge
\ldots \wedge e^n_{j_{r - p}}
$$
我们系统的转移映射
$$
K(R, f_1^n, \ldots, f_r^n) \to K(R, f_1^{n + 1}, \ldots, f_r^{n + 1})
$$
由如下规则给出：
$$
e^n_{j_1} \wedge \ldots \wedge e^n_{j_{r - p}}
\longmapsto
f_{i_0} \ldots f_{i_{p - 1}}
e^{n + 1}_{j_1} \wedge \ldots \wedge e^{n + 1}_{j_{r - p}}
$$
其中指标 $1 \leq i_0 < \ldots < i_{p - 1} \leq r$ 满足
$\{1, \ldots r\} =
\{i_0, \ldots, i_{p - 1}\} \amalg \{j_1, \ldots, j_{r - p}\}$。
我们略去说明此映射与微分相容的简短计算。注意，转移映射在次数
$0$ 上总是 $1$，在次数 $r$ 上等于 $f_1 \ldots f_r$。

\medskip\noindent
用 $K^p(R, f_1^n, \ldots, f_r^n)$ 表示 Koszul 复形的次数 $p$ 项。
注意，对任意 $f \in R$，有
$$
R_f = \colim (R \xrightarrow{f} R \xrightarrow{f} R \to \ldots )
$$
因此在次数 $p$ 上得到
$$
\colim K^p(R, f_1^n, \ldots f_r^n) =
\bigoplus\nolimits_{1 \leq i_0 < \ldots < i_{p - 1} \leq r}
R_{f_{i_0} \ldots f_{i_{p - 1}}}
$$
这里，上述 Koszul 复形的元
$e^n_{j_1} \wedge \ldots \wedge e^n_{j_{r - p}}$ 在余极限中映到直和项
$R_{f_{i_0} \ldots f_{i_{p - 1}}}$ 中的元
$(f_{i_0} \ldots f_{i_{p - 1}})^{-n}$，其中指标选取为满足
$\{1, \ldots r\} = \{i_0, \ldots, i_{p - 1}\}
\amalg \{j_1, \ldots, j_{r - p}\}$。因此此复形上的微分由下式给出：
$$
d(1\text{ 于 }R_{f_{i_0} \ldots f_{i_{p - 1}}}) =
\sum\nolimits_{i \not \in \{i_0, \ldots, i_{p - 1}\}}
(-1)^{i - t}\text{ 于 }
R_{f_{i_0} \ldots f_{i_t} f_i f_{i_{t + 1}} \ldots f_{i_{p - 1}}}
$$
因此，若考虑在次数 $p$ 上由映射
$$
\bigoplus\nolimits_{1 \leq i_0 < \ldots < i_{p - 1} \leq r}
R_{f_{i_0} \ldots f_{i_{p - 1}}}
\longrightarrow
\bigoplus\nolimits_{1 \leq i_0 < \ldots < i_{p - 1} \leq r}
R_{f_{i_0} \ldots f_{i_{p - 1}}}
$$
给出的复形映射，且该映射由规则
$$
1\text{ 于 }R_{f_{i_0} \ldots f_{i_{p - 1}}}
\longmapsto
(-1)^{i_0 + \ldots + i_{p - 1} + p}\text{ 于 }R_{f_{i_0} \ldots f_{i_{p - 1}}}
$$
确定，则得到从 $\colim K(R, f_1^n, \ldots, f_r^n)$ 到本节定义的扩张交错
{\v C}ech 复形的同构。我们略去对符号相符的验证。
\end{proof}








\section{Koszul 正则序列}
\label{more-algebra-section-koszul-regular}

\noindent
在阅读本节前，请先参看 Algebra, Sections
\ref{algebra-section-regular-sequences}、
\ref{algebra-section-quasi-regular} 与
\ref{algebra-section-depth}。

\begin{definition}
\label{more-algebra-definition-koszul-regular-sequence}
设 $R$ 为环，$r \geq 0$，$f_1, \ldots, f_r \in R$ 为一列元。设 $M$ 为 $R$-模。
称序列 $f_1, \ldots, f_r$：
\begin{enumerate}
\item 若对所有 $i \not = 0$ 有
$H_i(K_\bullet(f_1, \ldots, f_r) \otimes_R M) = 0$，则它是
{\it $M$-Koszul 正则的}；
\item 若 $H_1(K_\bullet(f_1, \ldots, f_r) \otimes_R M) = 0$，则它是
{\it $M$-$H_1$ 正则的}；
\item 若对所有 $i \not = 0$ 有 $H_i(K_\bullet(f_1, \ldots, f_r)) = 0$，
则它是{\it Koszul 正则的}；以及
\item 若 $H_1(K_\bullet(f_1, \ldots, f_r)) = 0$，则它是{\it $H_1$ 正则的}。
\end{enumerate}
\end{definition}

\noindent
我们将在引理 \ref{more-algebra-lemma-regular-koszul-regular}、
\ref{more-algebra-lemma-koszul-regular-H1-regular} 与
\ref{more-algebra-lemma-H1-regular-quasi-regular} 中看到，对环 $R$ 的元
$f_1, \ldots, f_r$，有下列推含：
\begin{align*}
f_1, \ldots, f_r\text{ is a regular sequence}
& \Rightarrow f_1, \ldots, f_r\text{ is a Koszul-regular sequence} \\
& \Rightarrow f_1, \ldots, f_r\text{ is an }H_1\text{-regular sequence} \\
& \Rightarrow f_1, \ldots, f_r\text{ is a quasi-regular sequence.}
\end{align*}
一般而言，这些推含均不可逆；但若 $R$ 是 Noether 局部环，且
$f_1, \ldots, f_r \in \mathfrak m_R$，则四个条件全部等价
（引理 \ref{more-algebra-lemma-noetherian-finite-all-equivalent}）。
若 $f = f_1 \in R$ 是长度为 $1$ 的序列，且 $f$ 不是 $R$ 的单位元，
则下列条件显然全部等价：
\begin{enumerate}
\item $f$ 是长度为一的正则序列；
\item $f$ 是长度为一的 Koszul 正则序列；以及
\item $f$ 是长度为一的 $H_1$ 正则序列。
\end{enumerate}
这些条件也显然推出 $f$ 是长度为一的拟正则序列。但确实存在长度为 $1$
而非正则序列的拟正则序列。具体地，令
$$
R = k[x, y_0, y_1, \ldots]/(xy_0, xy_1 - y_0, xy_2 - y_1, \ldots)
$$
并令 $f$ 为 $x$ 在 $R$ 中的像。则 $f$ 是零因子，但
$\bigoplus_{n \geq 0} (f^n)/(f^{n + 1}) \cong k[x]$ 是多项式环。

\begin{lemma}
\label{more-algebra-lemma-regular-koszul-regular}
设 $R$ 为环，$M$ 为 $R$-模，$f_1, \ldots, f_r \in R$，且对
$i = 1, \ldots, r$，乘以 $f_i$ 在 $M/(f_1, \ldots, f_{i - 1})M$ 上是单射。
则 $f_1, \ldots, f_r$ 是 $M$-Koszul 正则的。特别地，$M$-正则序列是
$M$-Koszul 正则的，且任意正则序列都是 Koszul 正则的。
\end{lemma}

\begin{proof}
设 $R$、$M$、$f_1, \ldots, f_r$ 如引理首句所述。若 $r = 1$，则
$f_1$ 是 $M$-Koszul 正则的这一结论是立即的。假设 $r > 1$。
因为 $f_1$ 在 $M$ 上是非零因子，我们得到复形的短正合列：
$$
0 \to K_\bullet(f_2, \ldots, f_r) \otimes M
\xrightarrow{f_1}
K_\bullet(f_2, \ldots, f_r) \otimes M \to
K_\bullet(\overline{f}_2, \ldots, \overline{f}_r) \otimes M/f_1M \to 0
$$
这里 $\overline{f}_i$ 是 $f_i$ 在 $R/(f_1)$ 中的像。由引理
\ref{more-algebra-lemma-cone-koszul}，复形 $K_\bullet(f_1, \ldots, f_r)$ 同构于在
$K_\bullet(f_2, \ldots, f_r)$ 上乘以 $f_1$ 的映射的锥。因此
$K_\bullet(R, f_1, \ldots, f_r) \otimes M$ 同构于第一个映射的锥。
从而 $K_\bullet(\overline{f}_2, \ldots, \overline{f}_r) \otimes M/f_1M$
与 $K_\bullet(f_1, \ldots, f_r) \otimes M$ 拟同构。因为
$R/(f_1)$、$M/f_1M$、$\overline{f}_2, \ldots, \overline{f}_r$ 满足引理的条件，
由归纳法，该复形在正次数上无同调。这完成第一个陈述的证明。
第二个陈述立即由第一个陈述得出。
\end{proof}

\begin{lemma}
\label{more-algebra-lemma-koszul-regular-H1-regular}
$M$-Koszul 正则序列是 $M$-$H_1$ 正则的。
Koszul 正则序列是 $H_1$ 正则的。
\end{lemma}

\begin{proof}
这直接由定义得出。
\end{proof}

\begin{lemma}
\label{more-algebra-lemma-mult-koszul-regular}
设 $f_1, \ldots, f_{r - 1} \in R$ 为一列元，$f, g \in R$，$M$ 为 $R$-模。
\begin{enumerate}
\item 若 $f_1, \ldots, f_{r - 1}, f$ 与 $f_1, \ldots, f_{r - 1}, g$
都是 $M$-$H_1$ 正则的，则 $f_1, \ldots, f_{r - 1}, fg$ 也是
$M$-$H_1$ 正则的。
\item 若 $f_1, \ldots, f_{r - 1}, f$ 与 $f_1, \ldots, f_{r - 1}, g$
都是 $M$-Koszul 正则的，则 $f_1, \ldots, f_{r - 1}, fg$ 也是
$M$-Koszul 正则的。
\end{enumerate}
\end{lemma}

\begin{proof}
由引理 \ref{more-algebra-lemma-koszul-mult}，对所有 $i$，有正合列
{\small
$$
H_i(K_\bullet(f_1, \ldots, f_{r - 1}, f) \otimes M) \to
H_i(K_\bullet(f_1, \ldots, f_{r - 1}, fg) \otimes M) \to
H_i(K_\bullet(f_1, \ldots, f_{r - 1}, g) \otimes M)
$$
}
\end{proof}

\begin{lemma}
\label{more-algebra-lemma-koszul-regular-flat-base-change}
设 $\varphi : R \to S$ 为平坦环映射，$f_1, \ldots, f_r \in R$。设 $M$ 为
$R$-模，并令 $N = M \otimes_R S$。
\begin{enumerate}
\item 若 $R$ 中的 $f_1, \ldots, f_r$ 是 $M$-$H_1$ 正则序列，则
$\varphi(f_1), \ldots, \varphi(f_r)$ 是 $S$ 中的 $N$-$H_1$ 正则序列。
\item 若 $f_1, \ldots, f_r$ 是 $R$ 中的 $M$-Koszul 正则序列，则
$\varphi(f_1), \ldots, \varphi(f_r)$ 是 $S$ 中的 $N$-Koszul 正则序列。
\end{enumerate}
\end{lemma}

\begin{proof}
这是因为
$K_\bullet(f_1, \ldots, f_r) \otimes_R S =
K_\bullet(\varphi(f_1), \ldots, \varphi(f_r))$，从而
$(K_\bullet(f_1, \ldots, f_r) \otimes_R M) \otimes_R S =
K_\bullet(\varphi(f_1), \ldots, \varphi(f_r)) \otimes_S N$。
\end{proof}

\begin{lemma}
\label{more-algebra-lemma-H1-regular-quasi-regular}
$M$-$H_1$ 正则序列是 $M$-拟正则的。
\end{lemma}

\begin{proof}
设 $R$ 为环，$M$ 为 $R$-模。设 $f_1, \ldots, f_r$ 为 $M$-$H_1$ 正则序列。
记 $J = (f_1, \ldots, f_r)$。该假设意味着有正合列
$$
\wedge^2(R^r) \otimes M \to R^{\oplus r} \otimes M \to JM \to 0
$$
其中第一个箭头由
$e_i \wedge e_j \otimes m \mapsto (f_ie_j - f_je_i) \otimes m$ 给出。
将此序列与 $R/J$ 张量，可知
$$
JM/J^2M = (R/J)^{\oplus r} \otimes_R M = (M/JM)^{\oplus r}
$$
是有限自由模。为完成证明，我们必须对每个 $n \geq 2$ 证明：
若
$$
\xi = \sum\nolimits_{|I| = n, I = (i_1, \ldots, i_r)}
m_I f_1^{i_1} \ldots f_r^{i_r} \in J^{n + 1}M
$$
则对所有 $I$ 有 $m_I \in JM$。在下一段中，我们对
$I = (0, \ldots, 0, n)$ 证明 $m_I \in JM$；最后一段将一般情形化归到此特殊情形。

\medskip\noindent
设 $I = (0, \ldots, 0, n)$，$\xi$ 如上所述。可以写成
$\xi = m_1 f_1 + \ldots + m_{r - 1}f_{r - 1} + m_I f_r^n$。
由于已假设 $\xi \in J^{n + 1}M$，也可写成
$\xi = \sum_{1 \leq i \leq j \leq r - 1} m_{ij}f_if_j +
\sum_{1 \leq i \leq r - 1}m'_i f_if_r^n + m'' f_r^{n + 1}$。
于是
$$
\begin{matrix}
(m_1 - m_{11}f_1 - m'_1f_r^n)f_1 + \\
(m_2 - m_{12}f_1 - m_{22}f_2 - m'_2f_r^n)f_2 + \\
\ldots + \\
(m_{r - 1} - m_{1 r - 1}f_1 - \ldots - m_{r - 1 r - 1}f_{r - 1}
- m'_{r - 1}f_r^n)f_{r - 1} + \\
(m_I - m'' f_r)f_r^n = 0
\end{matrix}
$$
由引理 \ref{more-algebra-lemma-mult-koszul-regular}，$f_1, \ldots, f_{r - 1}, f_r^n$
是 $M$-$H_1$ 正则的，因此
$m_I - m'' f_r$ 属于子模 $f_1M + \ldots + f_{r - 1}M + f_r^nM$。
从而 $m_I \in f_1M + \ldots + f_rM$。

\medskip\noindent
设 $S = R[x_1, x_2, \ldots, x_r, 1/x_r]$。环映射 $R \to S$ 是忠实平坦的，
因此 $f_1, \ldots, f_r$ 是 $S$ 中的 $M$-$H_1$ 正则序列；见引理
\ref{more-algebra-lemma-koszul-regular-flat-base-change}。由引理 \ref{more-algebra-lemma-change-basis}，可知
$$
g_1 = f_1 - \frac{x_1}{x_r} f_r,
\ \ldots,
\ g_{r - 1} = f_{r - 1} - \frac{x_{r - 1}}{x_r} f_r,
\ g_r = \frac{1}{x_r}f_r
$$
是 $S$ 中的 $M$-$H_1$ 正则序列。最后，注意我们的元 $\xi$ 可重写为
$$
\xi = \sum\nolimits_{|I| = n, I = (i_1, \ldots, i_r)}
m_I (g_1 + x_1 g_r)^{i_1} \ldots (g_{r - 1} + x_{r - 1} g_r)^{i_{r - 1}}
(x_rg_r)^{i_r}
$$
该表达式中 $g_r^n$ 的系数是
$$
\sum m_I x_1^{i_1} \ldots x_r^{i_r}
$$
由上一段讨论的情形，该和属于 $J(M \otimes_R S)$。因为单项式
$x_1^{i_1} \ldots x_r^{i_r}$ 构成 $S$ 在 $R$ 上的一组 $R$-基的一部分，所以如所需，
对所有 $I$ 有 $m_I \in J$。
\end{proof}

\noindent
对 Noether 局部环上的非零有限模，目前引入的所有各类正则序列彼此等价。

\begin{lemma}
\label{more-algebra-lemma-noetherian-finite-all-equivalent}
设 $(R, \mathfrak m)$ 为 Noether 局部环，$M$ 为非零有限 $R$-模，
$f_1, \ldots, f_r \in \mathfrak m$。下列条件等价：
\begin{enumerate}
\item $f_1, \ldots, f_r$ 是 $M$-正则序列；
\item $f_1, \ldots, f_r$ 是 $M$-Koszul 正则序列；
\item $f_1, \ldots, f_r$ 是 $M$-$H_1$ 正则序列；
\item $f_1, \ldots, f_r$ 是 $M$-拟正则序列。
\end{enumerate}
特别地，序列 $f_1, \ldots, f_r$ 是 $R$ 中的正则序列，当且仅当它是
Koszul 正则序列，当且仅当它是 $H_1$ 正则序列，当且仅当它是拟正则序列。
\end{lemma}

\begin{proof}
推含 (1) $\Rightarrow$ (2) 是引理 \ref{more-algebra-lemma-regular-koszul-regular}。
推含 (2) $\Rightarrow$ (3) 是引理 \ref{more-algebra-lemma-koszul-regular-H1-regular}。
推含 (3) $\Rightarrow$ (4) 是引理 \ref{more-algebra-lemma-H1-regular-quasi-regular}。
推含 (4) $\Rightarrow$ (1) 是 Algebra, Lemma
\ref{algebra-lemma-quasi-regular-regular}。
\end{proof}

\begin{lemma}
\label{more-algebra-lemma-H1-regular-in-quotient}
设 $A$ 为环，$I \subset A$ 为理想。设 $g_1, \ldots, g_m$ 是 $A$ 中的序列，
其在 $A/I$ 中的像是 $H_1$ 正则的。则
$I \cap (g_1, \ldots, g_m) = I(g_1, \ldots, g_m)$。
\end{lemma}

\begin{proof}
考虑复形的正合列
$$
0 \to I \otimes_A K_\bullet(A, g_1, \ldots, g_m)
\to K_\bullet(A, g_1, \ldots, g_m) \to
K_\bullet(A/I, g_1, \ldots, g_m) \to 0
$$
由假设，右边复形的 $H_1 = 0$，所以映射
$$
\Coker(I^{\oplus m} \to I)
\longrightarrow
\Coker(A^{\oplus m} \to A)
$$
是单射。这等价于本引理的断言。
\end{proof}

\begin{lemma}
\label{more-algebra-lemma-conormal-sequence-H1-regular}
设 $A$ 为环，$I \subset J \subset A$ 为理想。假设
$J/I \subset A/I$ 由一个 $H_1$ 正则序列生成。则 $I \cap J^2 = IJ$。
\end{lemma}

\begin{proof}
为证明此结论，选取 $g_1, \ldots, g_m \in J$，使其在 $A/I$ 中的像构成
生成 $J/I$ 的 $H_1$ 正则序列。特别地，$J = I + (g_1, \ldots, g_m)$。
设 $x \in I \cap J^2$。由于 $x \in J^2$，可写成
$$
x =
\sum a_{ij} g_ig_j +
\sum a_j g_j +
a
$$
其中 $a_{ij} \in A$、$a_j \in I$ 且 $a \in I^2$。于是
$\sum a_{ij}g_ig_j \in I \cap (g_1, \ldots, g_m)$，故由引理
\ref{more-algebra-lemma-H1-regular-in-quotient} 可知
$\sum a_{ij}g_ig_j \in I(g_1, \ldots, g_m)$。从而如所需，$x \in IJ$。
\end{proof}

\begin{lemma}
\label{more-algebra-lemma-join-quasi-regular-H1-regular}
设 $A$ 为环，$I$ 是由 $A$ 中的拟正则序列 $f_1, \ldots, f_n$ 生成的理想。
设 $g_1, \ldots, g_m \in A$，且其像
$\overline{g}_1, \ldots, \overline{g}_m$ 在 $A/I$ 中构成 $H_1$ 正则序列。
则 $f_1, \ldots, f_n, g_1, \ldots, g_m$ 是 $A$ 中的拟正则序列。
\end{lemma}

\begin{proof}
我们断言，对每个 $d$，$g_1, \ldots, g_m$ 在 $A/I^d$ 中构成
$H_1$ 正则序列。归纳地假设该结论在 $A/I^{d - 1}$ 中成立。我们有复形的短正合列
$$
0 \to K_\bullet(A, g_\bullet) \otimes_A I^{d - 1}/I^d
\to K_\bullet(A/I^d, g_\bullet) \to
K_\bullet(A/I^{d - 1}, g_\bullet) \to 0
$$
由于 $f_1, \ldots, f_n$ 拟正则，第一个复形是
$K_\bullet(A/I, g_1, \ldots, g_m)$ 的若干副本的直和，因而在次数 $1$ 上无同调。
由归纳假设，最后一个复形在次数 $1$ 上无同调，故中间复形也是如此。
特别地，对每个 $d \geq 1$，序列 $g_1, \ldots, g_m$ 在 $A/I^d$ 中拟正则；
见引理 \ref{more-algebra-lemma-H1-regular-quasi-regular}。
现在可以证明 $f_1, \ldots, f_n, g_1, \ldots, g_m$ 是 $A$ 中的拟正则序列。
事实上，令 $J = (f_1, \ldots, f_n, g_1, \ldots, g_m)$，并设（采用多重指标记号）
$$
\sum\nolimits_{|N| + |M| = d} a_{N, M} f^N g^M \in J^{d + 1}
$$
其中某些 $a_{N, M} \in A$。我们必须证明对所有 $N, M$ 都有
$a_{N, M} \in J$。令 $e \in \{0, 1, \ldots, d\}$。则
$$
\sum\nolimits_{|N| = d - e, \ |M| = e} a_{N, M} f^N g^M \in
(g_1, \ldots, g_m)^{e + 1} + I^{d - e + 1}
$$
因为 $g_1, \ldots, g_m$ 在 $A/I^{d - e + 1}$ 中是拟正则序列，所以对每个满足
$|M| = e$ 的 $M$，可推出
$$
\sum\nolimits_{|N| = d - e} a_{N, M} f^N \in
(g_1, \ldots, g_m) + I^{d - e + 1}
$$
在环 $A/I^{d - e + 1}$ 中将引理 \ref{more-algebra-lemma-H1-regular-in-quotient}
应用于 $I^{d - e}/I^{d - e + 1}$，可知
$\sum_{|N| = d - e} a_{N, M} f^N \in I^{d - e}(g_1, \ldots, g_m)$。
由于 $f_1, \ldots, f_n$ 在 $A$ 中拟正则，这又推出：对每个满足
$|N| = d - e$ 且 $|M| = e$ 的 $N, M$，都有 $a_{N, M} \in J$。
引理得证。
\end{proof}

\begin{lemma}
\label{more-algebra-lemma-join-H1-regular-sequences}
设 $A$ 为环，$I$ 是由 $A$ 中的 $H_1$ 正则序列
$f_1, \ldots, f_n$ 生成的理想。设 $g_1, \ldots, g_m \in A$，且其像
$\overline{g}_1, \ldots, \overline{g}_m$ 在 $A/I$ 中构成 $H_1$ 正则序列。
则 $f_1, \ldots, f_n, g_1, \ldots, g_m$ 是 $A$ 中的 $H_1$ 正则序列。
\end{lemma}

\begin{proof}
我们必须证明 $H_1(A, f_1, \ldots, f_n, g_1, \ldots, g_m) = 0$。
为此考虑交换图
$$
\xymatrix{
\wedge^2(A^{\oplus n + m}) \ar[r] \ar[d] &
A^{\oplus n + m} \ar[r] \ar[d] &
A \ar[r] \ar[d] & 0 \\
\wedge^2(A/I^{\oplus m}) \ar[r] &
A/I^{\oplus m} \ar[r] &
A/I \ar[r] & 0
}
$$
考虑映到 $A$ 中零元的一个元
$(a_1, \ldots, a_{n + m}) \in A^{\oplus n + m}$。
因为 $\overline{g}_1, \ldots, \overline{g}_m$ 在 $A/I$ 中构成
$H_1$ 正则序列，故
$(\overline{a}_{n + 1}, \ldots, \overline{a}_{n + m})$ 是
$\wedge^2(A/I^{\oplus m})$ 中某个元 $\overline{\alpha}$ 的像。
可将 $\overline{\alpha}$ 提升为一个元
$\alpha \in \wedge^2(A^{\oplus n + m})$，并从
$(a_1, \ldots, a_{n + m})$ 中减去它在 $A^{\oplus n + m}$ 中的像。
因此可以假设 $a_{n + 1}, \ldots, a_{n + m} \in I$。
由于 $I = (f_1, \ldots, f_n)$，可用顶部左侧水平箭头的像中诸元
$$
(0, \ldots, g_j, 0, \ldots, 0, f_i, 0, \ldots, 0)
$$
的线性组合修改 $(a_1, \ldots, a_{n + m})$，从而化归到
$a_{n + 1}, \ldots, a_{n + m}$ 均为零的情形。在此情形下，
$(a_1, \ldots, a_n, 0, \ldots, 0)$ 定义了
$H_1(A, f_1, \ldots, f_n)$ 的一个元，而我们已假设该模为零。
\end{proof}

\begin{lemma}
\label{more-algebra-lemma-truncate-H1-regular}
设 $A$ 为环，$f_1, \ldots, f_n, g_1, \ldots, g_m \in A$ 是
$H_1$ 正则序列。则 $\overline{g}_1, \ldots, \overline{g}_m$ 在
$A/(f_1, \ldots, f_n)$ 中的像构成 $H_1$ 正则序列。
\end{lemma}

\begin{proof}
令 $I = (f_1, \ldots, f_n)$。我们必须证明，$A/I$ 中的任一关系
$\sum_{j = 1, \ldots, m} \overline{a}_j \overline{g}_j$
都是平凡关系的线性组合。由于 $I = (f_1, \ldots, f_n)$，可将此关系提升为
$A$ 中的关系
$$
\sum\nolimits_{j = 1, \ldots, m} a_j g_j +
\sum\nolimits_{i = 1, \ldots, n} b_if_i = 0
$$
由假设，$A$ 中的这一关系是平凡关系的线性组合。取其在 $A/I$ 中的像，
便得到所需结论。
\end{proof}

\begin{lemma}
\label{more-algebra-lemma-join-koszul-regular-sequences}
设 $A$ 为环，$I$ 是由 $A$ 中的 Koszul 正则序列
$f_1, \ldots, f_n$ 生成的理想。设 $g_1, \ldots, g_m \in A$，且其像
$\overline{g}_1, \ldots, \overline{g}_m$ 在 $A/I$ 中构成 Koszul 正则序列。
则 $f_1, \ldots, f_n, g_1, \ldots, g_m$ 是 $A$ 中的 Koszul 正则序列。
\end{lemma}

\begin{proof}
我们的假设说明 $K_\bullet(A, f_1, \ldots, f_n)$ 是 $A/I$ 的有限自由解消，且
$K_\bullet(A/I, \overline{g}_1, \ldots, \overline{g}_m)$ 是
$A/(f_i, g_j)$ 在 $A/I$ 上的有限自由解消。于是
\begin{align*}
K_\bullet(A, f_1, \ldots, f_n, g_1, \ldots, g_m)
& = \text{Tot}(K_\bullet(A, f_1, \ldots, f_n) \otimes_A
K_\bullet(A, g_1, \ldots, g_m)) \\
& \cong A/I \otimes_A K_\bullet(A, g_1, \ldots, g_m) \\
& = K_\bullet(A/I, \overline{g}_1, \ldots, \overline{g}_m) \\
& \cong A/(f_i, g_j)
\end{align*}
第一个等式由引理 \ref{more-algebra-lemma-join-sequences-koszul-complex} 得出。
由于双复形
$K_\bullet(A, f_1, \ldots, f_n) \otimes_A K_\bullet(A, g_1, \ldots, g_m)$
的第 $q$ 行是 $A/I \otimes_A K_q(A, g_1, \ldots, g_m)$ 的解消，
第一个拟同构 $\cong$ 由 Homology, Lemma
\ref{homology-lemma-double-complex-gives-resolution}（的对偶）得出。
第二个等式是清楚的。最后一个拟同构由假设得出。故结论成立。
\end{proof}

\noindent
为得到下一个引理的结论，有必要同时假设
$f_1, \ldots, f_n$ 与 $f_1, \ldots, f_n, g_1, \ldots, g_m$
都是 Koszul 正则的。若删去 $f_1, \ldots, f_n$ Koszul 正则这一假设，
反例见 Examples, Lemma \ref{examples-lemma-strange-regular-sequence}。

\begin{lemma}
\label{more-algebra-lemma-truncate-koszul-regular}
设 $A$ 为环，$f_1, \ldots, f_n, g_1, \ldots, g_m \in A$。
若 $f_1, \ldots, f_n$ 与 $f_1, \ldots, f_n, g_1, \ldots, g_m$
都是 $A$ 中的 Koszul 正则序列，则
$\overline{g}_1, \ldots, \overline{g}_m$ 在 $A/(f_1, \ldots, f_n)$ 中
构成 Koszul 正则序列。
\end{lemma}

\begin{proof}
令 $I = (f_1, \ldots, f_n)$。我们的假设说明
$K_\bullet(A, f_1, \ldots, f_n)$ 是 $A/I$ 的有限自由解消，且
$K_\bullet(A, f_1, \ldots, f_n, g_1, \ldots, g_m)$ 是
$A/(f_i, g_j)$ 在 $A$ 上的有限自由解消。于是
\begin{align*}
A/(f_i, g_j) & \cong K_\bullet(A, f_1, \ldots, f_n, g_1, \ldots, g_m) \\
& = \text{Tot}(K_\bullet(A, f_1, \ldots, f_n) \otimes_A
K_\bullet(A, g_1, \ldots, g_m)) \\
& \cong A/I \otimes_A K_\bullet(A, g_1, \ldots, g_m) \\
& = K_\bullet(A/I, \overline{g}_1, \ldots, \overline{g}_m)
\end{align*}
第一个拟同构 $\cong$ 由假设得出。第一个等式由引理
\ref{more-algebra-lemma-join-sequences-koszul-complex} 得出。由于双复形
$K_\bullet(A, f_1, \ldots, f_n) \otimes_A K_\bullet(A, g_1, \ldots, g_m)$
的第 $q$ 行是 $A/I \otimes_A K_q(A, g_1, \ldots, g_m)$ 的解消，
第二个拟同构由 Homology, Lemma
\ref{homology-lemma-double-complex-gives-resolution}（的对偶）得出。
第二个等式是清楚的。故结论成立。
\end{proof}

\begin{lemma}
\label{more-algebra-lemma-independence-of-generators}
设 $R$ 为环，$I$ 是由 $f_1, \ldots, f_r \in R$ 生成的理想。
\begin{enumerate}
\item 若 $I$ 可由长度为 $r$ 的拟正则序列生成，则
$f_1, \ldots, f_r$ 是拟正则序列。
\item 若 $I$ 可由长度为 $r$ 的 $H_1$ 正则序列生成，则
$f_1, \ldots, f_r$ 是 $H_1$ 正则序列。
\item 若 $I$ 可由长度为 $r$ 的 Koszul 正则序列生成，则
$f_1, \ldots, f_r$ 是 Koszul 正则序列。
\end{enumerate}
\end{lemma}

\begin{proof}
若 $I$ 可由长度为 $r$ 的拟正则序列生成，则 $I/I^2$ 是秩为 $r$ 的自由
$R/I$-模。由于按假设 $f_1, \ldots, f_r$ 是一组生成元，可知其像
$\overline{f}_i$ 构成 $I/I^2$ 在 $R/I$ 上的一组基。因此
$f_1, \ldots, f_r$ 是拟正则序列，因为除 $I/I^2$ 的自由性外，这一切恰好意味着映射
$\text{Sym}^n_{R/I}(I/I^2) \to I^n/I^{n + 1}$ 均为同构。

\medskip\noindent
继续假设 $I$ 可由某个拟正则序列生成，记为 $g_1, \ldots, g_r$。
写成 $g_j = \sum a_{ij}f_i$。由上一段，$f_1, \ldots, f_r$ 拟正则，
故 $\det(a_{ij})$ 模 $I$ 可逆。矩阵 $a_{ij}$ 给出映射
$R^{\oplus r} \to R^{\oplus r}$，并诱导 Koszul 复形之间的映射
$\alpha : K_\bullet(R, f_1, \ldots, f_r) \to K_\bullet(R, g_1, \ldots, g_r)$；
见引理 \ref{more-algebra-lemma-functorial}。将 $\det(a_{ij})$ 可逆化后，该映射成为同构。
因为 $K_\bullet(R, f_1, \ldots, f_r)$ 与
$K_\bullet(R, g_1, \ldots, g_r)$ 的上同调模均被 $I$ 零化；见引理
\ref{more-algebra-lemma-homotopy-koszul}，可知 $\alpha$ 是拟同构。

\medskip\noindent
现在假设 $g_1, \ldots, g_r$ 是生成 $I$ 的 $H_1$ 正则序列。
由引理 \ref{more-algebra-lemma-H1-regular-quasi-regular}，$g_1, \ldots, g_r$ 是拟正则序列。
由上一段可知 $f_1, \ldots, f_r$ 是 $H_1$ 正则序列。
对 Koszul 正则序列同理。
\end{proof}

\begin{lemma}
\label{more-algebra-lemma-make-nonzero-divisor}
\begin{reference}
这是 \cite[Corollary]{McCoy} 的一个特殊情形。
\end{reference}
设 $R$ 为环，$a_1, \ldots, a_n \in R$ 为诸元，并且
$R \to R^{\oplus n}$，$x \mapsto (xa_1, \ldots, xa_n)$ 是单射。
则多项式环 $R[t_1, \ldots, t_n]$ 中的元 $\sum a_i t_i$ 是非零因子。
\end{lemma}

\begin{proof}
若某个 $a_i$ 是单位元，这恰是说，在 $R$ 上的多项式环中，任意形如
$t_1 + a_2 t_2 + \ldots + a_n t_n$ 的元都是非零因子。

\medskip\noindent
情形 I：$R$ 是 Noether 环。设 $\mathfrak q_j$，$j = 1, \ldots, m$，
为 $R$ 的伴随素理想。我们必须证明各映射
$$
\sum a_i t_i :
\text{Sym}^d(R^{\oplus n})
\longrightarrow
\text{Sym}^{d + 1}(R^{\oplus n})
$$
都是单射。由于 $\text{Sym}^d(R^{\oplus n})$ 是自由 $R$-模，
其伴随素理想为 $\mathfrak q_j$，$j = 1, \ldots, m$。
对每个 $j$，存在 $i = i(j)$ 使 $a_i \not \in \mathfrak q_j$，
这是因为存在 $x \in R$ 满足 $\mathfrak q_jx = 0$，而由假设，
对某个 $i$ 有 $a_i x \not = 0$。因此 $a_i$ 在 $R_{\mathfrak q_j}$ 中是单位元，
且该映射在 $\mathfrak q_j$ 处局部化后是单射。故该映射本身是单射；见
Algebra, Lemma \ref{algebra-lemma-zero-at-ass-zero}。

\medskip\noindent
情形 II：一般的 $R$。可以将 $R$ 写为 Noether 环 $R_\lambda$ 的并，且
$a_1, \ldots, a_n \in R_\lambda$。结论对每个 $R_\lambda$ 成立，
因而也对 $R$ 成立。
\end{proof}

\begin{lemma}
\label{more-algebra-lemma-Koszul-regular-flat-locally-regular}
设 $R$ 为环，$f_1, \ldots, f_n$ 是 $R$ 中的 Koszul 正则序列，且
$(f_1, \ldots, f_n) \not = R$。考虑忠实平坦的光滑环映射
$$
R \longrightarrow
S = R[\{t_{ij}\}_{i \leq j}, t_{11}^{-1}, t_{22}^{-1}, \ldots, t_{nn}^{-1}]
$$
对 $1 \leq i \leq n$，令
$$
g_i = \sum\nolimits_{i \leq j} t_{ij} f_j \in S.
$$
则 $g_1, \ldots, g_n$ 是 $S$ 中的正则序列，且
$(f_1, \ldots, f_n)S = (g_1, \ldots, g_n)$。
\end{lemma}

\begin{proof}
理想相等是显然的，因为矩阵
$$
\left(
\begin{matrix}
t_{11} & t_{12} & t_{13} & \ldots \\
0 & t_{22} & t_{23} & \ldots \\
0 & 0 & t_{33} & \ldots \\
\ldots & \ldots & \ldots & \ldots
\end{matrix}
\right)
$$
在 $S$ 中可逆。因为 $f_1, \ldots, f_n$ 是 Koszul 正则序列，可知
$R \to R^{\oplus n}$，$x \mapsto (xf_1, \ldots, xf_n)$ 的核为零
（因为它计算 $R$ 关于 $f_1, \ldots, f_n$ 的第 $n$ 个 Koszul 同调）。
因此由引理 \ref{more-algebra-lemma-make-nonzero-divisor} 可知
$g_1 = f_1 t_{11} + \ldots + f_n t_{1n}$ 在
$S' = R[t_{11}, t_{12}, \ldots, t_{1n}, t_{11}^{-1}]$ 中是非零因子。
由引理 \ref{more-algebra-lemma-koszul-regular-flat-base-change} 与
\ref{more-algebra-lemma-independence-of-generators}，可知
$g_1, f_2, \ldots, f_n$ 是 $S'$ 中的 Koszul 序列。
由引理 \ref{more-algebra-lemma-truncate-koszul-regular}，可推出
$\overline{f}_2, \ldots, \overline{f}_n$ 是 $S'/(g_1)$ 中的 Koszul 正则序列。
因此，对 $n$ 归纳可知，$g_2, \ldots, g_n$ 在
$S'/(g_1)[\{t_{ij}\}_{2 \leq i \leq j}, t_{22}^{-1}, \ldots, t_{nn}^{-1}]$
中的像 $\overline{g}_2, \ldots, \overline{g}_n$ 构成正则序列。
这又意味着 $g_1, \ldots, g_n$ 在 $S$ 中构成正则序列。
\end{proof}







\section{Koszul 正则序列的进一步结果}
\label{more-algebra-section-more-koszul-regular}

\noindent
我们继续第 \ref{more-algebra-section-koszul-regular} 节中的讨论。

\begin{lemma}
\label{more-algebra-lemma-vanishing-extended-alternating-koszul}
设 $R$ 为环，$f_1, \ldots, f_r \in R$ 是 Koszul 正则序列。则第
\ref{more-algebra-section-alternating-cech} 节中的增广交错 {\v C}ech 复形
$R \to \bigoplus\nolimits_{i_0} R_{f_{i_0}} \to
\bigoplus\nolimits_{i_0 < i_1} R_{f_{i_0}f_{i_1}} \to
\ldots \to R_{f_1\ldots f_r}$ 仅在次数 $r$ 上有上同调。
\end{lemma}

\begin{proof}
由引理 \ref{more-algebra-lemma-mult-koszul-regular} 与归纳法，对所有 $n \geq 1$，
序列 $f_1, \ldots, f_{r - 1}, f_r^n$ 都是 Koszul 正则的。
由引理 \ref{more-algebra-lemma-change-basis}，Koszul 正则序列的任一置换仍是
Koszul 正则序列。因此可将任意一个（或全部）$f_i$ 替换为其第 $n$ 次幂，
所得序列仍然 Koszul 正则。从而 $K_\bullet(R, f_1^n, \ldots, f_r^n)$
仅在同调次数 $0$ 上有非零上同调。由引理
\ref{more-algebra-lemma-extended-alternating-Cech-is-colimit-koszul}，这推出所需结论。
\end{proof}

\begin{lemma}
\label{more-algebra-lemma-blowup-regular-sequence}
设 $a, a_2, \ldots, a_r$ 是环 $R$ 中的 $H_1$ 正则序列
（例如 Koszul 正则序列或正则序列；见引理
\ref{more-algebra-lemma-regular-koszul-regular} 与
\ref{more-algebra-lemma-koszul-regular-H1-regular}）。令 $I = (a, a_2, \ldots, a_r)$，
则爆破代数 $R' = R[\frac{I}{a}]$ 同构于
$R'' = R[y_2, \ldots, y_r]/(a y_i - a_i)$。
\end{lemma}

\begin{proof}
由 Algebra, Lemma \ref{algebra-lemma-affine-blowup-quotient-description}，
只需证明 $R''$ 无 $a$-挠。

\medskip\noindent
我们断言 $a, ay_2 - a_2, \ldots, ay_n - a_r$ 在
$R[y_2, \ldots, y_r]$ 中是 $H_1$ 正则序列。事实上，用来定义
$a, ay_2 - a_2, \ldots, ay_n - a_r$ 上 Koszul 复形的映射
$$
(a, ay_2 - a_2, \ldots, ay_n - a_r) :
R[y_2, \ldots, y_r]^{\oplus r}
\longrightarrow
R[y_2, \ldots, y_r]
$$
通过同构
$$
R[y_2, \ldots, y_r]^{\oplus r}
\longrightarrow
R[y_2, \ldots, y_r]^{\oplus r}
$$
同构于用来定义 $a, a_2, \ldots, a_r$ 上 Koszul 复形的映射
$$
(a, a_2, \ldots, a_r) :
R[y_2, \ldots, y_r]^{\oplus r} \longrightarrow
R[y_2, \ldots, y_r]
$$
其中该同构将 $(b_1, \ldots, b_r)$ 映为
$(b_1 - b_2y_2 \ldots - b_ry_r, -b_2, \ldots, - b_r)$。
由引理 \ref{more-algebra-lemma-functorial}，这些 Koszul 复形同构。
将引理 \ref{more-algebra-lemma-koszul-regular-flat-base-change} 应用于平坦环映射
$R \to R[y_2, \ldots, y_r]$，便证明了断言。由引理
\ref{more-algebra-lemma-cone-koszul}，$a, ay_2 - a_2, \ldots, ay_n - a_r$
上的 Koszul 复形 $K$ 是映射 $a : L \to L$ 的锥，其中 $L$ 是
$ay_2 - a_2, \ldots, ay_n - a_r$ 上的 Koszul 复形。
由断言，$H_1(K) = 0$，故 $a : H_0(L) \to H_0(L)$ 是单射；换言之，
$R'' = R[y_2, \ldots, y_r]/(a y_i - a_i)$ 没有非零的 $a$-挠元，正合所需。
\end{proof}

\begin{lemma}
\label{more-algebra-lemma-base-change-H1-regular}
设 $A \to B$ 为环映射。设 $f_1, \ldots, f_r$ 是 $B$ 中的序列，且
$B/(f_1, \ldots, f_r)$ 在 $A$ 上平坦。设 $A \to A'$ 为环映射。
则典范映射
$$
H_1(K_\bullet(B, f_1, \ldots, f_r)) \otimes_A A'
\longrightarrow
H_1(K_\bullet(B', f'_1, \ldots, f'_r))
$$
是满射。这里 $B' = B \otimes_A A'$，而 $f_i' \in B'$ 是 $f_i$ 的像。
\end{lemma}

\begin{proof}
序列
$$
\wedge^2(B^{\oplus r}) \to B^{\oplus r} \to B \to B/J \to 0
$$
是 $A$-模的复形，其中 $B/J$ 在 $A$ 上平坦，并且在位置
$B^{\oplus r}$ 处的上同调群为
$H_1 = H_1(K_\bullet(B, f_1, \ldots, f_r))$。将它与 $A'$ 张量，得到复形
$$
\wedge^2((B')^{\oplus r}) \to (B')^{\oplus r} \to B' \to B'/J' \to 0
$$
该复形在 $B'$ 与 $B'/J'$ 处正合。为了计算它在 $(B')^{\oplus r}$ 处的上同调群
$$
H'_1 = H_1(K_\bullet(B', f'_1, \ldots, f'_r))
$$
。将上面的第一个序列拆成如下三个正合列：
$$
0 \to J \to B \to B/J \to 0
$$
$$
0 \to K \to B^{\oplus r} \to J \to 0
$$
$$
\wedge^2(B^{\oplus r}) \to K \to H_1 \to 0
$$
在 $A$ 上与 $A'$ 张量，得到正合列
$$
\begin{matrix}
0 \to J \otimes_A A' \to B \otimes_A A' \to (B/J) \otimes_A A' \to 0 \\
K \otimes_A A' \to B^{\oplus r} \otimes_A A' \to J \otimes_A A' \to 0 \\
\wedge^2(B^{\oplus r}) \otimes_A A' \to K \otimes_A A' \to H_1 \otimes_A A'
\to 0
\end{matrix}
$$
其中第一个序列正合，因为 $B/J$ 在 $A$ 上平坦；见 Algebra, Lemma
\ref{algebra-lemma-flat-tor-zero}。由此可知 $J' = J \otimes_A A' \subset B'$，
且 $K \otimes_A A' \to \Ker((B')^{\oplus r} \to B')$ 是满射。因此
\begin{align*}
H_1 \otimes_A A'
& =
\Coker\left(\wedge^2(B^{\oplus r}) \otimes_A A' \to K \otimes_A A'\right) \\
& \to
\Coker\left(
\wedge^2((B')^{\oplus r})  \to \Ker((B')^{\oplus r} \to B')
\right) = H'_1
\end{align*}
也是满射。
\end{proof}

\begin{lemma}
\label{more-algebra-lemma-relative-regular-immersion-algebra}
设 $A \to B$ 与 $A \to A'$ 为环映射，令 $B' = B \otimes_A A'$。
设 $f_1, \ldots, f_r \in B$。假设 $B/(f_1, \ldots, f_r)B$ 在 $A$ 上平坦。
\begin{enumerate}
\item 若 $f_1, \ldots, f_r$ 是拟正则序列，则其在 $B'$ 中的像是拟正则序列。
\item 若 $f_1, \ldots, f_r$ 是 $H_1$ 正则序列，则其在 $B'$ 中的像是
$H_1$ 正则序列。
\end{enumerate}
\end{lemma}

\begin{proof}
假设 $f_1, \ldots, f_r$ 拟正则。令 $J = (f_1, \ldots, f_r)$。
由假设，$J^n/J^{n + 1}$ 同构于 $B/J$ 的若干副本的直和，因而在 $A$ 上平坦。
由归纳法及 Algebra, Lemma \ref{algebra-lemma-flat-ses}，可知
$B/J^n$ 在 $A$ 上平坦。理想 $(J')^n$ 等于 $J^n \otimes_A A'$；见
Algebra, Lemma \ref{algebra-lemma-flat-tor-zero}。因此
$(J')^n/(J')^{n + 1} = J^n/J^{n + 1} \otimes_A A'$，这显然推出
$f_1, \ldots, f_r$ 在 $B'$ 中是拟正则序列。

\medskip\noindent
假设 $f_1, \ldots, f_r$ 是 $H_1$ 正则的。由引理
\ref{more-algebra-lemma-base-change-H1-regular}，Koszul 同调群
$H_1(K_\bullet(B, f_1, \ldots, f_r))$ 的消失推出
$H_1(K_\bullet(B', f'_1, \ldots, f'_r))$ 的消失，故结论成立。
\end{proof}

\begin{lemma}
\label{more-algebra-lemma-cut-by-koszul-flat}
设 $A' \to B'$ 为环映射，$I \subset A'$ 为理想。令
$A = A'/I$ 且 $B = B'/IB'$。设 $f'_1, \ldots, f'_r \in B'$。假设：
\begin{enumerate}
\item $A' \to B'$ 平坦且为有限表现的；
\item $I$ 是局部幂零的；
\item 诸像 $f_1, \ldots, f_r \in B$ 构成拟正则序列；
\item $B/(f_1, \ldots, f_r)$ 在 $A$ 上平坦。
\end{enumerate}
则 $B'/(f'_1, \ldots, f'_r)$ 在 $A'$ 上平坦。
\end{lemma}

\begin{proof}
令 $C' = B'/(f'_1, \ldots, f'_r)$。我们必须证明 $A' \to C'$ 平坦。
设 $\mathfrak r' \subset C'$ 是位于 $\mathfrak p' \subset A'$ 之上的素理想，
并令 $\mathfrak q' \subset B'$ 为 $\mathfrak r'$ 的逆像。
由 Algebra, Lemma \ref{algebra-lemma-flat-localization}，只需证明
$A'_{\mathfrak p'} \to C'_{\mathfrak q'}$ 平坦。
Algebra, Lemma \ref{algebra-lemma-grothendieck-regular-sequence-general}
说明，只需证明 $f'_1, \ldots, f'_r$ 映到
$$
B'_{\mathfrak q'}/\mathfrak p'B'_{\mathfrak q'} =
B_\mathfrak q/\mathfrak p B_\mathfrak q =
(B \otimes_A \kappa(\mathfrak p))_\mathfrak q
$$
中的一个正则序列；此处记号含义显然。已知 $f_1, \ldots, f_r$
在 $B$ 中是拟正则序列，且 $B/(f_1, \ldots, f_r)$ 在 $A$ 上平坦。
由引理 \ref{more-algebra-lemma-relative-regular-immersion-algebra}，$f'_1, \ldots, f'_r$
在 $B \otimes_A \kappa(\mathfrak p)$ 中的像
$\overline{f}_1, \ldots, \overline{f}_r$ 构成拟正则序列。
由于 $(B \otimes_A \kappa(\mathfrak p))_\mathfrak q$ 是 Noether 局部环，
由引理 \ref{more-algebra-lemma-noetherian-finite-all-equivalent} 即得结论。
\end{proof}

\begin{lemma}
\label{more-algebra-lemma-cut-by-koszul}
设 $A' \to B'$ 为环映射，$I \subset A'$ 为理想。令
$A = A'/I$ 且 $B = B'/IB'$。设 $f'_1, \ldots, f'_r \in B'$。假设：
\begin{enumerate}
\item $A' \to B'$ 平坦且为有限表现的（例如光滑）；
\item $I$ 是局部幂零的；
\item 诸像 $f_1, \ldots, f_r \in B$ 构成拟正则序列；
\item $B/(f_1, \ldots, f_r)$ 在 $A$ 上光滑。
\end{enumerate}
则 $B'/(f'_1, \ldots, f'_r)$ 在 $A'$ 上光滑。
\end{lemma}

\begin{proof}
令 $C' = B'/(f'_1, \ldots, f'_r)$ 且 $C = B/(f_1, \ldots, f_r)$。
则 $A' \to C'$ 为有限表现的。由引理 \ref{more-algebra-lemma-cut-by-koszul-flat}，
$A' \to C'$ 平坦。映射 $A' \to C'$ 的纤维环等于 $A \to C$ 的纤维环，
故由假设 (4) 它们光滑。因此由 Algebra, Lemma
\ref{algebra-lemma-flat-fibre-smooth}，$A' \to C'$ 光滑。
\end{proof}








\section{正则理想}
\label{more-algebra-section-ideals}

\noindent
我们将在 Divisors, Section \ref{divisors-section-regular-ideal-sheaves}
中以很大的一般性讨论正则理想层的概念。这里定义仿射情形中的相应概念，
即环中理想这一情形。

\begin{definition}
\label{more-algebra-definition-regular-ideal}
设 $R$ 为环，$I \subset R$ 为理想。
\begin{enumerate}
\item 若对每个 $\mathfrak p \in V(I)$，存在 $g \in R$、
$g \not \in \mathfrak p$ 以及正则序列
$f_1, \ldots, f_r \in R_g$，使 $I_g$ 由 $f_1, \ldots, f_r$ 生成，
则称 $I$ 为{\it 正则理想}。
\item 若对每个 $\mathfrak p \in V(I)$，存在 $g \in R$、
$g \not \in \mathfrak p$ 以及 Koszul 正则序列
$f_1, \ldots, f_r \in R_g$，使 $I_g$ 由 $f_1, \ldots, f_r$ 生成，
则称 $I$ 为{\it Koszul 正则理想}。
\item 若对每个 $\mathfrak p \in V(I)$，存在 $g \in R$、
$g \not \in \mathfrak p$ 以及 $H_1$ 正则序列
$f_1, \ldots, f_r \in R_g$，使 $I_g$ 由 $f_1, \ldots, f_r$ 生成，
则称 $I$ 为{\it $H_1$ 正则理想}。
\item 若对每个 $\mathfrak p \in V(I)$，存在 $g \in R$、
$g \not \in \mathfrak p$ 以及拟正则序列
$f_1, \ldots, f_r \in R_g$，使 $I_g$ 由 $f_1, \ldots, f_r$ 生成，
则称 $I$ 为{\it 拟正则理想}。
\end{enumerate}
\end{definition}

\noindent
显然，给定 $I \subset R$，有推含关系
\begin{align*}
I\text{ is a regular ideal}
& \Rightarrow
I\text{ is a Koszul-regular ideal} \\
& \Rightarrow
I\text{ is a }H_1\text{-regular ideal} \\
& \Rightarrow
I\text{ is a quasi-regular ideal}
\end{align*}
见引理 \ref{more-algebra-lemma-regular-koszul-regular}、
\ref{more-algebra-lemma-koszul-regular-H1-regular} 与
\ref{more-algebra-lemma-H1-regular-quasi-regular}。这样的理想总是有限生成的。

\begin{lemma}
\label{more-algebra-lemma-quasi-regular-ideal-finite}
拟正则理想是有限生成的。
\end{lemma}

\begin{proof}
设 $I \subset R$ 为拟正则理想。由于 $V(I)$ 拟紧，存在
$g_1, \ldots, g_m \in R$，使
$V(I) \subset D(g_1) \cup \ldots \cup D(g_m)$，并且 $I_{g_j}$ 由拟正则序列
$g_{j1}, \ldots, g_{jr_j} \in R_{g_j}$ 生成。对某些 $g'_{ij} \in I$，写成
$g_{ji} = g'_{ji}/g_j^{e_{ij}}$。写 $1 + x = \sum g_j h_j$，其中
$x \in I$；这是可能的，因为 $V(I) \subset D(g_1) \cup \ldots \cup D(g_m)$。
注意
$\Spec(R) = D(g_1) \cup \ldots \cup D(g_m) \bigcup D(x)$。
于是 $I$ 由诸元 $g'_{ij}$ 与 $x$ 生成，因为这些元在该覆盖的每一片上生成；
见 Algebra, Lemma \ref{algebra-lemma-cover}。
\end{proof}

\begin{lemma}
\label{more-algebra-lemma-quasi-regular-ideal-finite-projective}
设 $I \subset R$ 是环中的拟正则理想。则 $I/I^2$ 是有限投射 $R/I$-模。
\end{lemma}

\begin{proof}
这由 Algebra, Lemma \ref{algebra-lemma-finite-projective} 与定义得出。
\end{proof}

\noindent
我们证明 Koszul 正则理想、$H_1$ 正则理想和拟正则理想的平坦下降。

\begin{lemma}
\label{more-algebra-lemma-flat-descent-regular-ideal}
设 $A \to B$ 为忠实平坦环映射，$I \subset A$ 为理想。
若 $IB$ 是 $B$ 中的 Koszul 正则（相应地，$H_1$ 正则；拟正则）理想，
则 $I$ 是 $A$ 中的 Koszul 正则（相应地，$H_1$ 正则；拟正则）理想。
\end{lemma}

\begin{proof}
在整个证明中固定素理想 $\mathfrak p \supset I$。假设 $IB$ 拟正则。
由引理 \ref{more-algebra-lemma-quasi-regular-ideal-finite}，$IB$ 是有限模，故由
Algebra, Lemma \ref{algebra-lemma-descend-properties-modules}，
$I$ 是有限 $A$-模。因为 $A \to B$ 平坦，可知
$$
I/I^2 \otimes_{A/I} B/IB = I/I^2 \otimes_A B = IB/(IB)^2.
$$
由于 $IB$ 拟正则，$B/IB$-模 $IB/(IB)^2$ 有限局部自由。因此由
Algebra, Proposition \ref{algebra-proposition-ffdescent-finite-projectivity}，
$I/I^2$ 有限投射。特别地，对某个 $f \in A$、$f \not \in \mathfrak p$，
以 $A_f$ 替换 $A$ 后，可假设 $I/I^2$ 是秩为 $r$ 的自由模。
选取 $f_1, \ldots, f_r \in I$，使其给出 $I/I^2$ 的一组基。
由 Nakayama 引理（见 Algebra, Lemma \ref{algebra-lemma-NAK}），
再次进行上述替换 $A \leadsto A_f$ 后，可知 $I$ 由 $f_1, \ldots, f_r$ 生成。

\medskip\noindent
“拟正则”情形的证明。上面已经看到，$I/I^2$ 以 $r$ 个生成元
$f_1, \ldots, f_r$ 为自由基。为完成此情形的证明，必须证明对每个
$d \geq 2$，映射 $\text{Sym}^d(I/I^2) \to I^d/I^{d + 1}$ 都是同构。
这是清楚的，因为其忠实平坦基变换
$\text{Sym}^d(IB/(IB)^2) \to (IB)^d/(IB)^{d + 1}$ 由假设在 $B$ 上局部为同构。
略去细节。

\medskip\noindent
“$H_1$ 正则”与“Koszul 正则”情形的证明。考虑上面构造的生成 $I$ 的序列
$f_1, \ldots, f_r$。由引理 \ref{more-algebra-lemma-independence-of-generators}，
对任意 $g \in B$，只要 $IB$ 由某个 $H_1$ 正则或 Koszul 正则序列生成，
$f_1, \ldots, f_r$ 在 $B_g$ 中的像便是 $H_1$ 正则或 Koszul 正则序列。
因此 $K_\bullet(A, f_1, \ldots, f_r) \otimes_A B_g$ 的
$H_1$ 或 $H_i$（$i > 0$）消失。由于
$K_\bullet(B, f_1, \ldots, f_r) = K_\bullet(A, f_1, \ldots, f_r) \otimes_A B$
的同调被 $IB$ 零化（见引理 \ref{more-algebra-lemma-homotopy-koszul}），且
$V(IB) \subset \bigcup_{g\text{ 如上}} D(g)$，可知
$K_\bullet(A, f_1, \ldots, f_r) \otimes_A B$ 在次数 $1$ 或所有正次数上的同调消失。
再利用 $A \to B$ 忠实平坦，可知
$K_\bullet(A, f_1, \ldots, f_r)$ 也具有同样的性质。
\end{proof}

\begin{lemma}
\label{more-algebra-lemma-conormal-sequence-H1-regular-ideal}
设 $A$ 为环，$I \subset J \subset A$ 为理想。假设
$J/I \subset A/I$ 是 $H_1$ 正则理想。则 $I \cap J^2 = IJ$。
\end{lemma}

\begin{proof}
局部化后，这立即由引理 \ref{more-algebra-lemma-conormal-sequence-H1-regular} 得出。
\end{proof}







\section{局部完全交映射}
\label{more-algebra-section-lci}

\noindent
可以利用上面的材料，通过（有限）多项式代数的表现来定义环之间的局部完全交映射。

\begin{lemma}
\label{more-algebra-lemma-koszul-independence-presentation}
设 $A \to B$ 是有限型环映射。若对某个表现
$\alpha : A[x_1, \ldots, x_n] \to B$，其核 $I$ 是 Koszul 正则理想，
则对任一表现 $\beta : A[y_1, \ldots, y_m] \to B$，其核 $J$ 都是
Koszul 正则理想。
\end{lemma}

\begin{proof}
选取 $f_j \in A[x_1, \ldots, x_n]$ 使 $\alpha(f_j) = \beta(y_j)$，并选取
$g_i \in A[y_1, \ldots, y_m]$ 使 $\beta(g_i) = \alpha(x_i)$。
于是得到交换图
$$
\xymatrix{
A[x_1, \ldots, x_n, y_1, \ldots, y_m]
\ar[d]^{x_i \mapsto g_i} \ar[rr]_-{y_j \mapsto f_j} & &
A[x_1, \ldots, x_n] \ar[d] \\
A[y_1, \ldots, y_m] \ar[rr] & & B
}
$$
注意，$A[x_i, y_j] \to B$ 的核 $K$ 等于
$K = (I, y_j - f_j) = (J, x_i - g_i)$。特别地，由引理
\ref{more-algebra-lemma-quasi-regular-ideal-finite}，$I$ 有限生成，故
$J = K/(x_i - g_i)$ 也有限生成。

\medskip\noindent
选取素理想 $\mathfrak q \subset B$。由于
$I/I^2 \oplus B^{\oplus m} = J/J^2 \oplus B^{\oplus n}$
（Algebra, Lemma \ref{algebra-lemma-conormal-module}），可知
$$
\dim J/J^2 \otimes_B \kappa(\mathfrak q) + n =
\dim I/I^2 \otimes_B \kappa(\mathfrak q) + m.
$$
选取 $p_1, \ldots, p_t \in I$，使其映到
$I/I^2 \otimes \kappa(\mathfrak q) = I \otimes_{A[x_i]} \kappa(\mathfrak q)$
的一组基。选取 $q_1, \ldots, q_s \in J$，使其映到
$J/J^2 \otimes \kappa(\mathfrak q) = J \otimes_{A[y_j]} \kappa(\mathfrak q)$
的一组基。于是 $s + n = t + m$。由 Nakayama 引理，存在
$h \in A[x_i]$ 与 $h' \in A[y_j]$，两者都映到
$\kappa(\mathfrak q)$ 中的非零元，并且 $A[x_i, 1/h]$ 中
$I_h = (p_1, \ldots, p_t)$，$A[y_j, 1/h']$ 中
$J_{h'} = (q_1, \ldots, q_s)$。由于 $I$ Koszul 正则，还可假设
$I_h$ 由某个 Koszul 正则序列生成。该序列的长度必为
$t = \dim I/I^2 \otimes_B \kappa(\mathfrak q)$，故由引理
\ref{more-algebra-lemma-independence-of-generators}，$p_1, \ldots, p_t$ 是 Koszul 正则序列。
又因为 $y_1 - f_1, \ldots, y_m - f_m$ 是正则序列，由引理
\ref{more-algebra-lemma-join-koszul-regular-sequences} 可知
$$
y_1 - f_1, \ldots, y_m - f_m, p_1, \ldots, p_t
$$
是 $A[x_i, y_j, 1/h]$ 中的 Koszul 正则序列。该序列生成理想 $K_h$。
因此，理想 $K_{hh'}$ 由长度为 $m + t = n + s$ 的 Koszul 正则序列生成。
另一方面，它也由同样长度的序列
$$
x_1 - g_1, \ldots, x_n - g_n, q_1, \ldots, q_s
$$
生成，故由引理 \ref{more-algebra-lemma-independence-of-generators}，该序列也是 Koszul 正则的。
最后，由引理 \ref{more-algebra-lemma-truncate-koszul-regular}，可知
$q_1, \ldots, q_s$ 在
$$
A[x_i, y_j, 1/hh']/(x_1 - g_1, \ldots, x_n - g_n)
\cong A[y_j, 1/h'']
$$
中的像构成生成 $J_{h''}$ 的 Koszul 正则序列。由于 $h''$ 是 $hh'$ 的像，
它在 $\kappa(\mathfrak q)$ 中的像非零，故结论成立。
\end{proof}

\noindent
该引理使我们可以作如下定义。

\begin{definition}
\label{more-algebra-definition-local-complete-intersection}
若环映射 $A \to B$ 为有限型的，并且对某个（等价地，对任一）表现
$B = A[x_1, \ldots, x_n]/I$，理想 $I$ Koszul 正则，
则称它为{\it 局部完全交}。
\end{definition}

\noindent
这一概念是局部的。

\begin{lemma}
\label{more-algebra-lemma-lci-local}
设 $R \to S$ 为环映射，$g_1, \ldots, g_m \in S$ 生成单位理想。
若每个 $R \to S_{g_j}$ 都是局部完全交，则 $R \to S$ 也是局部完全交。
\end{lemma}

\begin{proof}
设 $S = R[x_1, \ldots, x_n]/I$ 为一个表现。选取
$h_j \in R[x_1, \ldots, x_n]$，使其在 $S$ 中映到 $g_j$。
则 $R[x_1, \ldots, x_n, x_{n + 1}]/(I, x_{n + 1}h_j - 1)$
是 $S_{g_j}$ 的表现。因此
$I_j = (I, x_{n + 1}h_j - 1)$ 是
$R[x_1, \ldots, x_n, x_{n + 1}]$ 中的 Koszul 正则理想。
选取素理想 $I \subset \mathfrak q \subset R[x_1, \ldots, x_n]$。
则对某个 $j$ 有 $h_j \not \in \mathfrak q$，且
$\mathfrak q_j = (\mathfrak q, x_{n + 1}h_j - 1)$ 是 $V(I_j)$ 中位于
$\mathfrak q$ 之上的素理想。选取 $f_1, \ldots, f_r \in I$，使其映到
$I/I^2 \otimes \kappa(\mathfrak q)$ 的一组基。于是序列
$x_{n + 1}h_j - 1, f_1, \ldots, f_r$ 的各元属于 $I_j$，并映到
$I_j \otimes \kappa(\mathfrak q_j)$ 的一组基；见 Algebra, Lemma
\ref{algebra-lemma-principal-localization-NL}。由 Nakayama 引理，存在
$h \in R[x_1, \ldots, x_n, x_{n + 1}]$，使 $(I_j)_h$ 由
$x_{n + 1}h_j - 1, f_1, \ldots, f_r$ 生成。还可假设 $(I_j)_h$
由某个长度为 $e$ 的 Koszul 正则序列生成。考察
$I_j \otimes \kappa(\mathfrak q_j)$ 的维数，可知 $e = r + 1$。
因此由引理 \ref{more-algebra-lemma-independence-of-generators}，
$x_{n + 1}h_j - 1, f_1, \ldots, f_r$ 是生成 $(I_j)_h$ 的 Koszul 正则序列，
其中某个 $h \in R[x_1, \ldots, x_n, x_{n + 1}]$ 满足
$h \not \in \mathfrak q_j$。由引理 \ref{more-algebra-lemma-truncate-koszul-regular}，
可知对某个 $h' \in R[x_1, \ldots, x_n]$、$h' \not \in \mathfrak q$，
$I_{h'}$ 由 Koszul 正则序列生成，正合所需。
\end{proof}

\begin{lemma}
\label{more-algebra-lemma-relative-global-complete-intersection-koszul}
设 $R$ 为环。若 $R[x_1, \ldots, x_n]/(f_1, \ldots, f_c)$
是相对全局完全交，则 $f_1, \ldots, f_c$ 是 Koszul 正则序列。
\end{lemma}

\begin{proof}
回忆同调群 $H_i(K_\bullet(f_\bullet))$ 被理想 $(f_1, \ldots, f_c)$ 零化。
因此只需证明：对所有包含 $(f_1, \ldots, f_c)$ 的素理想
$\mathfrak q \subset R[x_1, \ldots, x_n]$，
$H_i(K_\bullet(f_\bullet))_\mathfrak q$ 为零。这由 Algebra, Lemma
\ref{algebra-lemma-relative-global-complete-intersection-conormal}
以及正则序列是 Koszul 正则这一事实得出
（引理 \ref{more-algebra-lemma-regular-koszul-regular}）。
\end{proof}

\begin{lemma}
\label{more-algebra-lemma-syntomic-lci}
设 $R \to S$ 为环映射。下列条件等价：
\begin{enumerate}
\item $R \to S$ 是 syntomic 的
（Algebra, Definition \ref{algebra-definition-lci}）；以及
\item $R \to S$ 平坦且为局部完全交。
\end{enumerate}
\end{lemma}

\begin{proof}
假设 (1)。则按定义，$R \to S$ 平坦。由 Algebra, Lemma
\ref{algebra-lemma-syntomic} 与引理 \ref{more-algebra-lemma-lci-local}，
只需证明相对全局完全交同态是局部完全交同态；这正是引理
\ref{more-algebra-lemma-relative-global-complete-intersection-koszul}。

\medskip\noindent
假设 (2)。局部完全交为有限表现的，因为 Koszul 正则理想有限生成。
设 $R \to k$ 为到一个域的映射。只需证明
$S' = S \otimes_R k$ 是 $k$ 上的局部完全交；见 Algebra, Definition
\ref{algebra-definition-lci-field}。选取素理想 $\mathfrak q' \subset S'$。
写 $S = R[x_1, \ldots, x_n]/I$。则
$S' = k[x_1, \ldots, x_n]/I'$，其中
$I' \subset k[x_1, \ldots, x_n]$ 是 $I$ 的像。设
$\mathfrak p' \subset k[x_1, \ldots, x_n]$、$\mathfrak q \subset S$ 与
$\mathfrak p \subset R[x_1, \ldots, x_n]$ 为相应的素理想。
由定义 \ref{more-algebra-definition-regular-ideal}，存在
$g \in R[x_1, \ldots, x_n]$、$g \not \in \mathfrak p$ 与
$f_1, \ldots, f_r \in R[x_1, \ldots, x_n]_g$，它们构成生成 $I_g$ 的
Koszul 正则序列。由于 $S$、因而 $S_g$ 在 $R$ 上平坦，
这些元在 $k[x_1, \ldots, x_n]_g$ 中的像
$f'_1, \ldots, f'_r$ 构成生成 $I'_g$ 的 $H_1$ 正则序列；见引理
\ref{more-algebra-lemma-relative-regular-immersion-algebra}。因此，由引理
\ref{more-algebra-lemma-noetherian-finite-all-equivalent}，$f'_1, \ldots, f'_r$
在 $k[x_1, \ldots, x_n]_{\mathfrak p'}$ 中映到生成
$I'_{\mathfrak p'}$ 的正则序列。应用 Algebra, Lemma
\ref{algebra-lemma-lci}，可知对某个 $g' \in S$、
$g' \not \in \mathfrak q'$，$S'_{gg'}$ 是 $k$ 上的全局完全交，正合所需。
\end{proof}

\noindent
对局部完全交 $R \to S$，当 $n \geq 2$ 时有 $H_n(L_{S/R}) = 0$。
由于我们（尚）未定义完整的余切复形，无法陈述并证明这一点，
但可以推出它的一个后果。

\begin{lemma}
\label{more-algebra-lemma-transitive-lci-at-end}
设 $A \to B \to C$ 为环映射。假设 $B \to C$ 是局部完全交同态。
选取核为 $I$ 的表现 $\alpha : A[x_s, s \in S] \to B$，
以及核为 $J$ 的表现 $\beta : B[y_1, \ldots, y_m] \to C$。
设 $\gamma : A[x_s, y_t] \to C$ 为所诱导的 $C$ 的表现，其核为 $K$。
则得到典范交换图
$$
\xymatrix{
0 \ar[r] &
\Omega_{A[x_s]/A} \otimes C \ar[r] &
\Omega_{A[x_s, y_t]/A} \otimes C \ar[r] &
\Omega_{B[y_t]/B} \otimes C \ar[r] &
0 \\
0 \ar[r] &
I/I^2 \otimes C \ar[r] \ar[u] &
K/K^2 \ar[r] \ar[u] &
J/J^2 \ar[r] \ar[u] &
0
}
$$
其两行均正合。特别地，Algebra, Lemma
\ref{algebra-lemma-exact-sequence-NL} 中的六项正合列可在左侧补上一个零；
即序列
$$
0 \to H_1(\NL_{B/A} \otimes_B C) \to
H_1(L_{C/A}) \to
H_1(L_{C/B}) \to
\Omega_{B/A} \otimes_B C \to
\Omega_{C/A} \to
\Omega_{C/B} \to 0
$$
正合。
\end{lemma}

\begin{proof}
唯一需要证明的是映射 $I/I^2 \otimes C \to K/K^2$ 的单射性。
由假设，理想 $J$ Koszul 正则。因此由引理
\ref{more-algebra-lemma-conormal-sequence-H1-regular-ideal}，有
$IA[x_s, y_j] \cap K^2 = IK$。这意味着 $K/K^2 \to J/J^2$ 的核同构于
$IA[x_s, y_j]/IK$。由于张量积右正合，
$I/I^2 \otimes_A C = IA[x_s, y_j]/IK$，这便给出了所需的
$I/I^2 \otimes_A C \to K/K^2$ 的单射性。
\end{proof}

\begin{lemma}
\label{more-algebra-lemma-transitive-colimit-lci-at-end}
设 $A \to B \to C$ 为环映射。若 $B \to C$ 是局部完全交同态的滤过余极限，
则引理 \ref{more-algebra-lemma-transitive-lci-at-end} 的结论仍成立。
\end{lemma}

\begin{proof}
这由引理 \ref{more-algebra-lemma-transitive-lci-at-end} 与 Algebra, Lemma
\ref{algebra-lemma-colimits-NL} 得出。
\end{proof}

\begin{lemma}
\label{more-algebra-lemma-henselization-NL}
设 $A \to B$ 为局部环的局部同态。设 $A^h \to B^h$（相应地，
$A^{sh} \to B^{sh}$）为所诱导的 Hensel 化（相应地，严格 Hensel 化）之间的映射
（Algebra, Lemma \ref{algebra-lemma-henselian-functorial}，相应地，
Lemma \ref{algebra-lemma-strictly-henselian-functorial}）。则
$\NL_{B/A} \otimes_B B^h \to \NL_{B^h/A^h}$ 与
$\NL_{B/A} \otimes_B B^{sh} \to \NL_{B^{sh}/A^{sh}}$
在上同调群上诱导同构。
\end{lemma}

\begin{proof}
因为 $A^h$ 是 $A$ 上 étale 代数的滤过余极限，由 Algebra, Lemma
\ref{algebra-lemma-colimits-NL} 与 Algebra, Definition
\ref{algebra-definition-etale}，可知 $\NL_{A^h/A}$ 是非循环复形。
$B^h/B$ 也有同样性质。对 $A \to A^h \to B^h$ 使用 Jacobi--Zariski 序列
（Algebra, Lemma \ref{algebra-lemma-exact-sequence-NL}），可知
$\NL_{B^h/A} \to \NL_{B^h/A^h}$ 在上同调群上诱导同构。
此外，étale 环映射是局部完全交，因为它甚至是全局完全交；见
Algebra, Lemma \ref{algebra-lemma-etale-standard-smooth}。
由引理 \ref{more-algebra-lemma-transitive-colimit-lci-at-end}，得到与
$A \to B \to B^h$ 相伴的六项正合 Jacobi--Zariski 序列，它证明
$\NL_{B/A} \otimes_B B^h \to \NL_{B^h/A}$ 在上同调群上诱导同构。
这完成 Hensel 化之间映射的情形。严格 Hensel 化的情形以完全相同的方式证明。
\end{proof}







\section{Cartier 等式与几何正则性}
\label{more-algebra-section-cartier-equality}

\noindent
本节及下一节的参考文献是 \cite[Section 39]{MatCA}。
为顺利阅读本节，读者应熟悉朴素余切复形及其性质；见 Algebra, Section
\ref{algebra-section-netherlander}。

\begin{lemma}[Cartier 等式]
\label{more-algebra-lemma-cartier-equality}
设 $K/k$ 为有限生成域扩张。则 $\Omega_{K/k}$ 与 $H_1(L_{K/k})$
都是有限维的，并且
$\text{trdeg}_k(K) = \dim_K \Omega_{K/k} - \dim_K H_1(L_{K/k})$。
\end{lemma}

\begin{proof}
可以找到 $k$ 上的全局完全交
$A = k[x_1, \ldots, x_n]/(f_1, \ldots, f_c)$，使 $K$ 同构于 $A$ 的分式域；
见 Algebra, Lemma \ref{algebra-lemma-colimit-syntomic} 及其证明。
在此情形下，由 Algebra, Lemmas \ref{algebra-lemma-NL-homotopy} 与
\ref{algebra-lemma-localize-NL}，可知 $\NL_{K/k}$ 同伦等价于复形
$$
\bigoplus\nolimits_{j = 1, \ldots, c} K \longrightarrow
\bigoplus\nolimits_{i = 1, \ldots, n} K\text{d}x_i
$$
$K$ 在 $k$ 上的超越次数是 $A$ 的维数
（由 Algebra, Lemma \ref{algebra-lemma-dimension-prime-polynomial-ring}），
即 $n - c$，故结论成立。
\end{proof}

\begin{lemma}
\label{more-algebra-lemma-transitivity-gamma}
设 $M/L/K$ 为域扩张。则 Jacobi--Zariski 序列
$$
0 \to H_1(L_{L/K}) \otimes_L M \to
H_1(L_{M/K}) \to
H_1(L_{M/L}) \to
\Omega_{L/K} \otimes_L M \to
\Omega_{M/K} \to
\Omega_{M/L} \to 0
$$
正合。
\end{lemma}

\begin{proof}
结合引理 \ref{more-algebra-lemma-transitive-colimit-lci-at-end} 与 Algebra, Lemma
\ref{algebra-lemma-colimit-syntomic} 即得。
\end{proof}

\begin{lemma}
\label{more-algebra-lemma-gamma-commutative-diagram}
给定域的交换图
$$
\xymatrix{
K \ar[r] & K' \\
k \ar[u] \ar[r] & k' \ar[u]
}
$$
其中 $k'/k$ 与 $K'/K$ 是有限生成域扩张，则映射
$$
\alpha : \Omega_{K/k} \otimes_K K' \to \Omega_{K'/k'}
\quad\text{且}\quad
\beta : H_1(L_{K/k}) \otimes_K K' \to H_1(L_{K'/k'})
$$
的核与余核均为有限维的，并且
$$
\dim \Ker(\alpha) - \dim \Coker(\alpha)
-\dim \Ker(\beta) + \dim \Coker(\beta)
=
\text{trdeg}_k(k') - \text{trdeg}_K(K')
$$
\end{lemma}

\begin{proof}
$k \subset k' \subset K'$ 与 $k \subset K \subset K'$ 的
Jacobi--Zariski 序列分别为
$$
0 \to H_1(L_{k'/k}) \otimes K'  \to
H_1(L_{K'/k}) \to
H_1(L_{K'/k'}) \to
\Omega_{k'/k} \otimes K' \to
\Omega_{K'/k} \to
\Omega_{K'/k'} \to 0
$$
和
$$
0 \to H_1(L_{K/k}) \otimes K' \to
H_1(L_{K'/k}) \to
H_1(L_{K'/K}) \to
\Omega_{K/k} \otimes K' \to
\Omega_{K'/k} \to
\Omega_{K'/K} \to 0
$$
由引理 \ref{more-algebra-lemma-cartier-equality}，向量空间
$\Omega_{k'/k}$、$\Omega_{K'/K}$、$H_1(L_{K'/K})$ 与
$H_1(L_{k'/k})$ 均为有限维的，且其维数的交错和为
$\text{trdeg}_k(k') - \text{trdeg}_K(K')$。引理得证。
\end{proof}





\section{几何正则性}
\label{more-algebra-section-geometrically-regular}

\noindent
设 $k$ 为域，$(A, \mathfrak m, K)$ 为 Noether 局部 $k$-代数。
Jacobi--Zariski 序列（Algebra, Lemma
\ref{algebra-lemma-exact-sequence-NL}）是典范正合列
$$
H_1(L_{K/k}) \to \mathfrak m/\mathfrak m^2 \to \Omega_{A/k} \otimes_A K \to
\Omega_{K/k} \to 0
$$
这是因为由 Algebra, Lemma \ref{algebra-lemma-NL-surjection}，
$H_1(L_{K/A}) = \mathfrak m/\mathfrak m^2$。
我们将证明，该序列左端的正合性刻画了正则局部环 $A$ 是否在 $k$ 上几何正则。
在第 \ref{more-algebra-section-regular-fs} 节中，我们会把它与形式光滑的概念联系起来。

\begin{proposition}
\label{more-algebra-proposition-characterization-geometrically-regular}
设 $k$ 为特征 $p > 0$ 的域，$(A, \mathfrak m, K)$ 为 Noether 局部
$k$-代数。下列条件等价：
\begin{enumerate}
\item $A$ 在 $k$ 上几何正则；
\item 对所有在 $k$ 上有限的 $k \subset k' \subset k^{1/p}$，
环 $A \otimes_k k'$ 正则；
\item $A$ 正则，且典范映射
$H_1(L_{K/k}) \to \mathfrak m/\mathfrak m^2$ 是单射；以及
\item $A$ 正则，且映射
$\Omega_{k/\mathbf{F}_p} \otimes_k K \to \Omega_{A/\mathbf{F}_p} \otimes_A K$
是单射。
\end{enumerate}
\end{proposition}

\begin{proof}
(3) $\Rightarrow$ (1) 的证明。假设 (3)。设 $k'/k$ 为有限纯不可分扩张。
令 $A' = A \otimes_k k'$。这是极大理想为 $\mathfrak m'$ 的局部环。
令 $K' = A'/\mathfrak m'$。得到交换图
{\small
$$
\xymatrix{
0 \ar[r] &
H_1(L_{K/k}) \otimes K' \ar[r] \ar[d]_\beta &
\mathfrak m/\mathfrak m^2 \otimes K' \ar[r] \ar[d] &
\Omega_{A/k} \otimes_A K' \ar[r] \ar[d]_{\cong} &
\Omega_{K/k} \otimes K' \ar[r] \ar[d]_\alpha &
0
\\
 &
H_1(L_{K'/k'}) \ar[r] &
\mathfrak m'/(\mathfrak m')^2 \ar[r] &
\Omega_{A'/k'} \otimes_{A'} K' \ar[r] &
\Omega_{K'/k'} \ar[r] &
0
}
$$
}
其两行均正合。由微分模的基变换
（Algebra, Lemma \ref{algebra-lemma-differentials-base-change}），
第三个竖直箭头是同构。因此 $\alpha$ 是满射。由引理
\ref{more-algebra-lemma-gamma-commutative-diagram}，有
$$
\dim \Ker(\alpha) - \dim \Ker(\beta) + \dim \Coker(\beta) = 0
$$
（且这些维数均有限）。追图可得
$\dim \mathfrak m'/(\mathfrak m')^2 \leq \dim \mathfrak m/\mathfrak m^2$。
另一方面，因为 $A \to A'$ 有限平坦，由 Algebra, Lemma
\ref{algebra-lemma-dimension-base-fibre-total}，可知
$\dim(A) = \dim(A')$。故按定义，$A'$ 正则。

\medskip\noindent
(3) 与 (4) 的等价性。考虑交换图各行的 Jacobi--Zariski 序列
$$
\xymatrix{
\mathbf{F}_p \ar[r] & A \ar[r] & K \\
\mathbf{F}_p \ar[r] \ar[u] & k \ar[r] \ar[u] & K \ar[u]
}
$$
得到交换图
$$
\xymatrix{
0 \ar[r] &
\mathfrak m/\mathfrak m^2 \ar[r] &
\Omega_{A/\mathbf{F}_p} \otimes_A K \ar[r] &
\Omega_{K/\mathbf{F}_p} \ar[r] & 0 & \\
0 \ar[r] &
H_1(L_{K/k}) \ar[r] \ar[u] &
\Omega_{k/\mathbf{F}_p} \otimes_k K \ar[r] \ar[u] &
\Omega_{K/\mathbf{F}_p} \ar[r] \ar[u] &
\Omega_{K/k} \ar[r] \ar[u] &
0
}
$$
其两行均正合。这里使用了
$H_1(L_{K/A}) = \mathfrak m/\mathfrak m^2$，以及
$H_1(L_{K/\mathbf{F}_p}) = 0$；后者成立是因为 $K/\mathbf{F}_p$ 可分，
见 Algebra, Proposition
\ref{algebra-proposition-characterize-separable-field-extensions}。
因此显然，映射 $H_1(L_{K/k}) \to \mathfrak m/\mathfrak m^2$ 与
$\Omega_{k/\mathbf{F}_p} \otimes_k K \to \Omega_{A/\mathbf{F}_p} \otimes_A K$
的核具有相同维数。

\medskip\noindent
沿用 Faltings 的方法证明 (2) $\Rightarrow$ (4)；见
\cite{Faltings-einfacher}。设 $a_1, \ldots, a_n \in k$ 为诸元，使
$\text{d}a_1, \ldots, \text{d}a_n$ 在 $\Omega_{k/\mathbf{F}_p}$ 中线性无关。
考虑域扩张 $k' = k(a_1^{1/p}, \ldots, a_n^{1/p})$。由 Algebra, Lemma
\ref{algebra-lemma-size-extension-pth-roots}，可知
$k' = k[x_1, \ldots, x_n]/(x_1^p - a_1, \ldots, x_n^p - a_n)$。
特别地，可知 $k'/k$ 的朴素余切复形同伦于具有零微分的复形
$\bigoplus_{j = 1, \ldots, n} k' \rightarrow \bigoplus_{i = 1, \ldots, n} k'$，
因为在 $\Omega_{k[x_1, \ldots, x_n]/k}$ 中
$\text{d}(x_j^p - a_j) = 0$。如上令
$A' = A \otimes_k k'$ 与 $K' = A'/\mathfrak m'$。
由 Algebra, Lemma \ref{algebra-lemma-change-base-NL}，
$\NL_{A'/A}$ 同伦等价于具有零微分的复形
$\bigoplus_{j = 1, \ldots, n} A' \rightarrow \bigoplus_{i = 1, \ldots, n} A'$；
即 $H_1(L_{A'/A})$ 与 $\Omega_{A'/A}$ 都是秩为 $n$ 的自由模。
$\mathbf{F}_p \to A \to A'$ 的 Jacobi--Zariski 序列为
$$
H_1(L_{A'/A}) \to \Omega_{A/\mathbf{F}_p} \otimes_A A'
\to \Omega_{A'/\mathbf{F}_p} \to \Omega_{A'/A} \to 0
$$
利用核为 $(x_j^p - a_j)$ 的表现 $A[x_1, \ldots, x_n] \to A'$，
展开 Algebra, Lemma \ref{algebra-lemma-exact-sequence-NL} 中的映射，
可知 $H_1(L_{A'/A})$ 的第 $j$ 个基向量映到
$\Omega_{A/\mathbf{F}_p} \otimes A'$ 中的 $\text{d}a_j \otimes 1$。
因为 $\Omega_{A'/A}$ 自由（因而平坦），与 $K'$ 张量后得到正合列
$$
K'^{\oplus n} \to \Omega_{A/\mathbf{F}_p} \otimes_A K'
\xrightarrow{\beta} \Omega_{A'/\mathbf{F}_p} \otimes_{A'} K' \to
K'^{\oplus n} \to 0
$$
由此可知诸元 $\text{d}a_j \otimes 1$ 生成 $\Ker(\beta)$；
还需证明它们线性无关，即证明 $\dim(\Ker(\beta)) = n$。
考虑下列大图
$$
\xymatrix{
0 \ar[r] &
\mathfrak m'/(\mathfrak m')^2 \ar[r] &
\Omega_{A'/\mathbf{F}_p} \otimes K' \ar[r] &
\Omega_{K'/\mathbf{F}_p} \ar[r] & 0 \\
0 \ar[r] &
\mathfrak m/\mathfrak m^2 \otimes K' \ar[r] \ar[u]^\alpha &
\Omega_{A/\mathbf{F}_p} \otimes K' \ar[r] \ar[u]^\beta &
\Omega_{K/\mathbf{F}_p} \otimes K' \ar[r] \ar[u]^\gamma & 0
}
$$
由引理 \ref{more-algebra-lemma-cartier-equality} 以及
$\mathbf{F}_p \to K \to K'$ 的 Jacobi--Zariski 序列，
$\gamma$ 的核与余核具有相同的有限维数。由假设，$A'$ 正则
（且如上所见，其维数与 $A$ 相同），故 $\alpha$ 的核与余核具有相同维数。
从而 $\beta$ 的核与余核具有相同维数，这正是所需证明的。

\medskip\noindent
推含 (1) $\Rightarrow$ (2) 是平凡的。命题的证明完成。
\end{proof}

\begin{lemma}
\label{more-algebra-lemma-geometrically-regular-over-field}
设 $k$ 为特征 $p > 0$ 的域，$(A, \mathfrak m, K)$ 为 Noether 局部
$k$-代数。假设 $A$ 在 $k$ 上几何正则。设 $K/F/k$ 为有限生成子扩张。
设 $\varphi : k[y_1, \ldots, y_m] \to A$ 为 $k$-代数映射，使得
$y_i$ 在 $K$ 中映到 $F$ 的元，并且
$\text{d}y_1, \ldots, \text{d}y_m$ 映到 $\Omega_{F/k}$ 的一组基。
令 $\mathfrak p = \varphi^{-1}(\mathfrak m)$。则
$$
k[y_1, \ldots, y_m]_\mathfrak p \to A
$$
平坦，且 $A/\mathfrak pA$ 正则。
\end{lemma}

\begin{proof}
令 $A_0 = k[y_1, \ldots, y_m]_\mathfrak p$，其极大理想为
$\mathfrak m_0$，剩余域为 $K_0$。注意 $\Omega_{A_0/k}$ 是秩为 $m$ 的自由模，
且 $\Omega_{A_0/k} \otimes K_0 \to \Omega_{K_0/k}$ 是同构。
显然 $A_0$ 在 $k$ 上几何正则。因此
$H_1(L_{K_0/k}) \to \mathfrak m_0/\mathfrak m_0^2$ 是同构；见命题
\ref{more-algebra-proposition-characterization-geometrically-regular}。现在考虑
$$
\xymatrix{
H_1(L_{K_0/k}) \otimes K \ar[d] \ar[r] &
\mathfrak m_0/\mathfrak m_0^2 \otimes K \ar[d] \\
H_1(L_{K/k}) \ar[r] & \mathfrak m/\mathfrak m^2
}
$$
由引理 \ref{more-algebra-lemma-transitivity-gamma}，左侧竖直箭头是单射；
由命题 \ref{more-algebra-proposition-characterization-geometrically-regular}，
下方水平箭头是单射，故右侧竖直箭头也是单射。
因此 $A_0$ 中的一个正则参数系映到 $A$ 中正则参数系的一部分。
由 Algebra, Lemmas \ref{algebra-lemma-flat-over-regular} 与
\ref{algebra-lemma-regular-ring-CM} 即得结论。
\end{proof}













\section{拓扑环与拓扑模}
\label{more-algebra-section-topological-ring}

\noindent
我们简要讨论拓扑 Abel 群的一些性质。若 Abel 群 $M$ 带有一个拓扑，使加法
$M \times M \to M$、$(x, y) \mapsto x + y$ 与取逆
$M \to M$、$x \mapsto -x$ 连续，则称 $M$ 为{\it 拓扑 Abel 群}。
{\it 拓扑 Abel 群的同态}就是连续的 Abel 群同态。
交换拓扑群的范畴是加性范畴并且有核与余核，但不是 Abel 范畴
（因为公理 $\Im = \Coim$ 不成立）。若 $N \subset M$ 是子群，
我们也把 $N$ 与 $M/N$ 视为拓扑群：在 $N$ 上使用诱导拓扑，
在 $M/N$ 上使用商拓扑（即使 $M \to M/N$ 为商映射的拓扑）。
注意，若 $N \subset M$ 是开子群，则 $M/N$ 上的拓扑是离散的。

\medskip\noindent
若存在一个由子群组成的 $0$ 的邻域基，则称 $M$ 上的拓扑为{\it 线性的}。
此时这些子群也是开的。下面是一个例子。设 $I$ 为有向集，且
$G_i$ 是 $I$ 上（离散）Abel 群的逆系统。则
$$
G = \lim_{i \in I} G_i
$$
赋予逆极限拓扑后是线性拓扑化的；由 $\Ker(G \to G_i)$ 给出
$0$ 的一组邻域基。反之，设 $M$ 是线性拓扑化 Abel 群。
任取开子群的一个基本系统 $U_i \subset M$、$i \in I$
（即诸 $U_i$ 构成开邻域的基本系统，且每个 $U_i$ 都是 $M$ 的子群）。
令 $i \geq i' \Leftrightarrow U_i \subset U_{i'}$，可知 $I$ 是有向集。
得到线性拓扑化 Abel 群的同态
$$
c : M \longrightarrow \lim_{i \in I} M/U_i.
$$
显然，$M$（作为拓扑空间）{\it 分离}当且仅当 $c$ 是单射。
若 $c$ 是同构，则称 $M$ {\it 完备}\footnote{我们把分离性纳入完备性的定义，
因为希望完备群中的极限唯一。任意拓扑群都有一个完备性定义；除分离性问题外，
该定义在我们的特殊情形中与这里的定义一致。}。
留给读者验证，这一条件与上面所选的开子群基本系统
$\{U_i\}_{i \in I}$ 无关。事实上，拓扑 Abel 群
$M^\wedge = \lim_{i \in I} M/U_i$ 与这一选择无关，有时称为 $M$ 的
{\it 完备化}。任何上述 $G = \lim G_i$ 都是完备的；特别地，完备化
$M^\wedge$ 总是完备的。

\begin{definition}[拓扑环]
\label{more-algebra-definition-topological-ring}
\begin{reference}
\cite[Chapter 0, Sections 7.1 and 7.2]{EGA1}
\end{reference}
设 $R$ 为环，$M$ 为 $R$-模。
\begin{enumerate}
\item 若 $R$ 带有一个拓扑，使得加法与乘法作为映射
$R \times R \to R$ 均连续，其中 $R \times R$ 带乘积拓扑，
则称 $R$ 为{\it 拓扑环}。在此情形下，若 $M$ 带有一个拓扑，使加法
$M \times M \to M$ 与标量乘法 $R \times M \to M$ 连续，
则称 $M$ 为{\it 拓扑模}。
\item {\it 拓扑模的同态}就是连续的 $R$-模映射。
{\it 拓扑环的同态}是对给定拓扑连续的环同态。
\item 若 $0$ 有一个由子模组成的邻域基，则称 $M$ {\it 线性拓扑化}。
若 $0$ 有一个由理想组成的邻域基，则称 $R$ {\it 线性拓扑化}。
\item 若 $R$ 线性拓扑化，且 $I \subset R$ 是开理想，并且 $0$ 的每个邻域
都包含 $I^n$（对某个 $n$），则称 $I$ 为{\it 定义理想}。
\item 若 $R$ 线性拓扑化，则当 $R$ 有定义理想时，称 $R$ {\it 预容许}。
\item 若 $R$ 线性拓扑化、预容许且完备，则称 $R$ {\it 容许}\footnote{
依照我们的约定，这包括分离性。}。
\item 若 $R$ 线性拓扑化，且存在定义理想 $I$，使
$\{I^n\}_{n \geq 0}$ 构成 $0$ 的一组邻域基，则称 $R$ {\it 预 adic}。
\item 若 $R$ 线性拓扑化，则当 $R$ 预 adic 且完备时，称 $R$ {\it adic}。
\end{enumerate}
注意，（预）adic 拓扑环等价于这样一个（预）容许拓扑环：它有定义理想
$I$，且对所有 $n \geq 1$，$I^n$ 都是开的。
\end{definition}

\noindent
设 $R$ 为环，$M$ 为 $R$-模，$I \subset R$ 为理想。可以在 $R$ 上考虑
以 $\{I^n\}_{n \geq 0}$ 为 $0$ 的一组邻域基的线性拓扑。
该拓扑称为 $I$-adic 拓扑；$R$ 在 $I$-adic 拓扑下是预 adic 拓扑环\footnote{
因此，$I$-adic 拓扑有时称为 $I$-预 adic 拓扑。}。此外，在 $M$ 上以
$\{I^nM\}_{n \geq 0}$ 为 $0$ 的一组开邻域基的线性拓扑，使 $M$ 成为拓扑
$R$-模。这称为 $M$ 上的{\it $I$-adic 拓扑}。可知，$M$ 是 $I$-adic 完备的
（定义见 Algebra, Definition \ref{algebra-definition-complete}），当且仅当
$M$ 在 $I$-adic 拓扑下完备\footnote{可能出现如下情形：$I$-adic 完备化
$M^\wedge$ 并非 $I$-adic 完备，尽管 $M^\wedge$ 对极限拓扑总是完备的。
若 $I$ 有限生成，则 $M^\wedge$ 上的 $I$-adic 拓扑与极限拓扑一致；
见 Algebra, Lemma \ref{algebra-lemma-hathat-finitely-generated} 及其证明。}。
特别地，$R$ 是 $I$-adic 完备的，当且仅当 $R$ 在 $I$-adic 拓扑下是
adic 拓扑环。

\medskip\noindent
作为特殊情形，注意离散拓扑是 $0$-adic 拓扑，并且任意带离散拓扑的环都是 adic 的。

\begin{lemma}
\label{more-algebra-lemma-continuous}
设 $\varphi : R \to S$ 为环映射。设 $I \subset R$ 与 $J \subset S$ 为理想，
分别在 $R$ 与 $S$ 上赋予 $I$-adic 与 $J$-adic 拓扑。
则 $\varphi$ 是拓扑环同态，当且仅当对某个 $n \geq 1$ 有
$\varphi(I^n) \subset J$。
\end{lemma}

\begin{proof}
略。
\end{proof}

\begin{lemma}[Baire 纲定理]
\label{more-algebra-lemma-baire-category-complete-module}
设 $M$ 为拓扑 Abel 群。假设 $M$ 线性拓扑化、完备，并且有可数的
$0$ 的邻域基本系统。若 $U_n \subset M$、$n \geq 1$ 是开稠密子集，
则 $\bigcap_{n \geq 1} U_n$ 稠密。
\end{lemma}

\begin{proof}
设 $U_n$ 如引理陈述。以 $U_1 \cap \ldots \cap U_n$ 替换 $U_n$ 后，
可假设 $U_1 \supset U_2 \supset \ldots$。设
$M_n$、$n \in \mathbf{N}$ 为 $0$ 的邻域基本系统，可假设
$M_{n + 1} \subset M_n$。选取 $x \in M$。我们将证明：对每个
$k \geq 1$，存在 $y \in \bigcap_{n \geq 1} U_n$ 使 $x - y \in M_k$。

\medskip\noindent
如下构造 $y$。首先选取 $y_1 \in U_1$，使 $y_1 \in x + M_k$。
这是可能的，因为 $U_1$ 稠密，而 $x + M_k$ 是开的。
然后选取 $k_1 > k$，使 $y_1 + M_{k_1} \subset U_1$。
这是可能的，因为 $U_1$ 是开的。接着选取 $y_2 \in U_2$，使
$y_2 \in y_1 + M_{k_1}$。这是可能的，因为 $U_2$ 稠密，而
$y_2 + M_{k_1}$ 是开的。然后选取 $k_2 > k_1$，使
$y_2 + M_{k_2} \subset U_2$。这是可能的，因为 $U_2$ 是开的。

\medskip\noindent
如此继续，得到 $M$ 中收敛于 $y$ 的元的序列 $y_i$。
由构造，$x - y \in M_k$。由于
$$
y - y_i = (y_{i + 1} - y_i) + (y_{i + 2} - y_{i + 1}) + \ldots
$$
属于 $M_{k_i}$，故对所有 $i$，有
$y \in y_i + M_{k_i} \subset U_i$，正合所需。
\end{proof}

\begin{lemma}
\label{more-algebra-lemma-consequence-baire-complete-module}
在引理 \ref{more-algebra-lemma-baire-category-complete-module} 的相同假设下，
若对某些闭子群 $N_n$ 有 $M = \bigcup_{n \geq 1} N_n$，
则对某个 $n$，$N_n$ 是开的。
\end{lemma}

\begin{proof}
若不然，则对所有 $n$，$U_n = M \setminus N_n$ 都稠密，
这与引理 \ref{more-algebra-lemma-baire-category-complete-module} 矛盾。
\end{proof}

\begin{lemma}[开映射引理]
\label{more-algebra-lemma-open-mapping}
设 $u : N \to M$ 为拓扑化 Abel 群的连续映射。假设 $M$ 分离，
而 $N$ 完备、线性拓扑化，并且有可数的 $0$ 的邻域基本系统。
则下列两种情形恰有一种成立：
\begin{enumerate}
\item $u$ 是开映射；或者
\item 对某个开子群 $N' \subset N$，其像 $u(N')$ 在 $M$ 中无处稠密。
\end{enumerate}
\end{lemma}

\begin{proof}
设 $N_n$、$n \in \mathbf{N}$ 为 $0$ 的邻域基本系统。
可假设 $N_n$ 是子群，且 $N_{n + 1} \subset N_n$。若 (2) 不成立，
则对 $n = 1, 2, 3, \ldots$，$u(N_n)$ 的闭包 $M_n$ 是开子群。
因为 $u$ 连续，可知 $M_n$、$n \in \mathbf{N}$ 必构成 $M$ 中
$0$ 的开邻域基本系统。又因为 $M_n$ 是 $u(N_n)$ 的闭包，所以对所有
$n \geq 1$ 有
$$
u(N_n) + M_{n + 1} = M_n
$$
选取 $x_1 \in M_1$。于是可归纳选取 $y_i \in N_i$ 与
$x_{i + 1} \in M_{i + 1}$，使
$$
u(y_i) + x_{i + 1} = x_i
$$
由于 $N$ 完备，$N$ 中的元 $y = y_1 + y_2 + y_3 + \ldots$ 存在。
又因 $M$ 分离，便有 $x_1 = u(y)$。因此 $M_1 = u(N_1)$。
读者以完全相同的方式可证明对所有 $i \geq 2$，有
$M_i = u(N_i)$，从而 $u$ 是开映射。
\end{proof}









\section{拓扑环的形式光滑映射}
\label{more-algebra-section-formally-smooth}

\noindent
形式光滑性有一个适用于拓扑环同态的版本。

\begin{definition}
\label{more-algebra-definition-formally-smooth}
设 $R \to S$ 为拓扑环同态，且 $R$ 与 $S$ 线性拓扑化。
若对拓扑环同态的每一个实线交换图
$$
\xymatrix{
S \ar[r] \ar@{-->}[rd] & A/J \\
R \ar[r] \ar[u] & A \ar[u]
}
$$
其中 $A$ 为离散环，$J \subset A$ 是平方为零的理想，都存在使图交换的虚线箭头，
则称 $S$ 在 $R$ 上{\it 形式光滑}。
\end{definition}

\noindent
我们主要在给定理想 $\mathfrak m \subset R$ 与 $\mathfrak n \subset S$ 时使用此概念：
在 $R$ 上赋予 $\mathfrak m$-adic 拓扑，在 $S$ 上赋予 $\mathfrak n$-adic 拓扑。
$\varphi : R \to S$ 连续，当且仅当对某个 $m \geq 1$ 有
$\varphi(\mathfrak m^m) \subset  \mathfrak n$；见引理 \ref{more-algebra-lemma-continuous}。
事实表明，在此情形下只有 $S$ 上的拓扑是相关的。

\begin{lemma}
\label{more-algebra-lemma-formally-smooth}
设 $\varphi : R \to S$ 为环映射。
\begin{enumerate}
\item 若 $R \to S$ 在 Algebra, Definition
\ref{algebra-definition-formally-smooth} 的意义下形式光滑，
则对 $R$ 上任一线性拓扑和 $S$ 上任一预 adic 拓扑，只要
$R \to S$ 连续，$R \to S$ 就是形式光滑的。
\item 设 $\mathfrak n \subset S$ 与 $\mathfrak m \subset R$ 为理想，
并且 $\varphi$ 对 $R$ 上的 $\mathfrak m$-adic 拓扑和 $S$ 上的
$\mathfrak n$-adic 拓扑连续。则下列条件等价：
\begin{enumerate}
\item $\varphi$ 对 $R$ 上的 $\mathfrak m$-adic 拓扑和 $S$ 上的
$\mathfrak n$-adic 拓扑形式光滑；以及
\item $\varphi$ 对 $R$ 上的离散拓扑和 $S$ 上的 $\mathfrak n$-adic 拓扑形式光滑。
\end{enumerate}
\end{enumerate}
\end{lemma}

\begin{proof}
假设 $R \to S$ 在 Algebra, Definition
\ref{algebra-definition-formally-smooth} 的意义下形式光滑。
若 $S$ 有预 adic 拓扑，则存在理想 $\mathfrak n \subset S$，使 $S$
带有 $\mathfrak n$-adic 拓扑。设给定定义 \ref{more-algebra-definition-formally-smooth}
中的实线交换图。映射 $S \to A/J$ 的连续性意味着对某个 $k \geq 1$，
$\mathfrak n^k$ 映到 $A/J$ 中的零；见引理 \ref{more-algebra-lemma-continuous}。
由所假设的 $S$ 在 $R$ 上形式光滑，得到环映射 $\psi : S \to A$。
于是 $\psi(\mathfrak n^k) \subset J$，故因 $J^2 = 0$ 有
$\psi(\mathfrak n^{2k}) = 0$。所以由引理 \ref{more-algebra-lemma-continuous}，
$\psi$ 连续。这证明了 (1)。

\medskip\noindent
(2)(b) $\Rightarrow$ (2)(a) 的证明与 (1) 的证明相同。假设 (2)(a)。
设给定定义 \ref{more-algebra-definition-formally-smooth} 中的实线交换图，
但在 $R$ 上使用离散拓扑。因为 $\varphi$ 连续，可知对某个 $n \geq 1$，
$\varphi(\mathfrak m^n) \subset \mathfrak n$。因为 $S \to A/J$ 连续，
对某个 $k \geq 1$，$\mathfrak n^k$ 映到 $A/J$ 中的零。
故 $R \to A$ 将 $\mathfrak m^{nk}$ 映入 $J$，从而将
$\mathfrak m^{2nk}$ 映到 $A$ 中的零。因此 $R \to A$ 对
$\mathfrak m$-adic 拓扑连续，于是 (2)(a) 给出所需的虚线箭头。
\end{proof}

\begin{definition}
\label{more-algebra-definition-formally-smooth-adic}
设 $R \to S$ 为环映射，$\mathfrak n \subset S$ 为理想。
若引理 \ref{more-algebra-lemma-formally-smooth} 的等价条件 (2)(a) 与 (2)(b) 成立，
则称 $R \to S$ 对{\it $\mathfrak n$-adic 拓扑形式光滑}。
\end{definition}

\noindent
该性质由完备化继承。

\begin{lemma}
\label{more-algebra-lemma-formally-smooth-completion}
设 $(R, \mathfrak m)$ 与 $(S, \mathfrak n)$ 是带有限生成理想的环。
在 $R$ 与 $S$ 上分别赋予 $\mathfrak m$-adic 与 $\mathfrak n$-adic 拓扑。
设 $R \to S$ 为拓扑环同态。下列条件等价：
\begin{enumerate}
\item $R \to S$ 对 $\mathfrak n$-adic 拓扑形式光滑；
\item $R \to S^\wedge$ 对 $\mathfrak n^\wedge$-adic 拓扑形式光滑；
\item $R^\wedge \to S^\wedge$ 对 $\mathfrak n^\wedge$-adic 拓扑形式光滑。
\end{enumerate}
这里 $R^\wedge$ 与 $S^\wedge$ 分别是 $R$ 与 $S$ 的
$\mathfrak m$-adic 与 $\mathfrak n$-adic 完备化。
\end{lemma}

\begin{proof}
$\mathfrak m$ 有限生成这一假设推出：$R^\wedge$ 是
$\mathfrak mR^\wedge$-adic 完备的，$\mathfrak mR^\wedge = \mathfrak m^\wedge$，
并且 $R^\wedge/\mathfrak m^nR^\wedge = R/\mathfrak m^n$；见 Algebra, Lemma
\ref{algebra-lemma-hathat-finitely-generated} 及其证明。对
$(S, \mathfrak n)$ 同理。因此显然，在定义
\ref{more-algebra-definition-formally-smooth} 中，情形 (1)、(2) 与 (3) 的图一一对应。
\end{proof}

\noindent
使用 adic 环的优点是可以得到更强的提升性质。

\begin{lemma}
\label{more-algebra-lemma-lift-continuous}
设 $R \to S$ 为环映射，$\mathfrak n$ 为 $S$ 的理想。
假设 $R \to S$ 对 $\mathfrak n$-adic 拓扑形式光滑。
考虑拓扑环同态的实线交换图
$$
\xymatrix{
S \ar[r]_\psi \ar@{-->}[rd] & A/J \\
R \ar[r] \ar[u] & A \ar[u]
}
$$
其中 $A$ 是 adic 的，$A/J$ 是 $A$ 对闭理想 $J \subset A$ 的商
（作为拓扑环），并且对某个 $t \geq 1$，$J^t$ 包含于 $A$ 的一个定义理想。
则拓扑环范畴中存在使该图交换的虚线箭头。
\end{lemma}

\begin{proof}
设 $I \subset A$ 为定义理想，并使得对某个 $t$ 有 $I \supset J^t$。
则 $A = \lim A/I^n$；又因为假设 $J$ 闭，
$A/J = \lim A/J + I^n$。考虑离散 $R$-代数
$A_{n, m} = A/J^n + I^m$ 的下图：
$$
\xymatrix{
A/J^3 + I^3 \ar[r] \ar[d] &
A/J^2 + I^3 \ar[r] \ar[d] &
A/J + I^3 \ar[d] \\
A/J^3 + I^2 \ar[r] \ar[d] &
A/J^2 + I^2 \ar[r] \ar[d] &
A/J + I^2 \ar[d] \\
A/J^3 + I \ar[r] &
A/J^2 + I \ar[r] &
A/J + I
}
$$
注意，每个交换方块都定义了 $R$-代数的满射
$$
A_{n + 1, m + 1} \longrightarrow A_{n + 1, m} \times_{A_{n, m}} A_{n, m + 1}
$$
其核的平方为零。我们将归纳构造 $R$-代数映射
$\varphi_{n, m} : S \to A_{n, m}$。事实上，已有映射
$\varphi_{1, m} = \psi \bmod J + I^m$。注意，由于 $\psi$ 连续，
这些映射都连续。从选定的 $\varphi_{1, 1}$ 开始，利用 $S$ 在 $R$ 上的
形式光滑性，沿上图的底行向上提升，可归纳选取映射 $\varphi_{n, 1}$。
其余映射 $\varphi_{n, m}$ 按 $n + m$ 归纳构造。具体地，沿上面显示的满射
提升一对 $(\varphi_{n + 1, m}, \varphi_{n, m + 1})$，选取
$\varphi_{n + 1, m + 1}$（再次使用 $S$ 在 $R$ 上的形式光滑性）。
如此所得所有映射 $\varphi_{n, m}$ 都与该系统的转移映射相容。
因为 $J^t \subset I$，例如
$\varphi_n = \varphi_{nt, n} \bmod I^n$ 诱导映射 $S \to A/I^n$。
取极限 $\varphi = \lim \varphi_n$，得到映射
$S \to A = \lim A/I^n$。它与到 $A/J$ 的映射复合后等于 $\psi$，
因为已看到 $A/J = \lim A/J + I^n$。
最后证明 $\varphi$ 连续。由 $\psi$ 是拓扑环的态射这一假设及引理
\ref{more-algebra-lemma-continuous}，可知对某个 $r \geq 1$，
$\psi(\mathfrak n^r) \subset J + I/J$。因此
$\varphi(\mathfrak n^r) \subset J + I$，从而如所需，
$\varphi(\mathfrak n^{rt}) \subset I$。
\end{proof}

\begin{lemma}
\label{more-algebra-lemma-increase-ideal}
设 $R \to S$ 为环映射，$\mathfrak n \subset \mathfrak n' \subset S$ 为理想。
若 $R \to S$ 对 $\mathfrak n$-adic 拓扑形式光滑，
则 $R \to S$ 对 $\mathfrak n'$-adic 拓扑形式光滑。
\end{lemma}

\begin{proof}
略。
\end{proof}

\begin{lemma}
\label{more-algebra-lemma-compose-formally-smooth}
线性拓扑化环之间形式光滑连续同态的复合仍形式光滑。
\end{lemma}

\begin{proof}
略。（提示：这是完全形式的，由考察一个适当的图得出。）
\end{proof}

\begin{lemma}
\label{more-algebra-lemma-base-change-fs}
设 $R$、$S$ 为环，$\mathfrak n \subset S$ 为理想。
设 $R \to S$ 对 $\mathfrak n$-adic 拓扑形式光滑，并设
$R \to R'$ 为任意环映射。则
$R' \to S' = S \otimes_R R'$ 对
$\mathfrak n' = \mathfrak nS'$-adic 拓扑形式光滑。
\end{lemma}

\begin{proof}
给定定义 \ref{more-algebra-definition-formally-smooth} 中的实线图
$$
\xymatrix{
S \ar[r] \ar@{-->}[rrd] & S' \ar[r] \ar@{-->}[rd] & A/J \\
R  \ar[u] \ar[r] & R' \ar[r] \ar[u] & A \ar[u]
}
$$
则复合 $S \to S' \to A/J$ 连续。由假设，较长的虚线箭头存在。
由张量积的泛性质，得到较短的虚线箭头。
\end{proof}

\noindent
我们已经在 Algebra, Lemma \ref{algebra-lemma-descent-formally-smooth}
中见过沿忠实平坦环映射的形式光滑性下降。在当前拓扑环的语境中也有类似结果。
不过，这里只证明下面这个非常简单且易证、但已相当有用的版本。

\begin{lemma}
\label{more-algebra-lemma-descent-fs}
设 $R$、$S$ 为环，$\mathfrak n \subset S$ 为理想。设 $R \to R'$ 为环映射。
令 $S' = S \otimes_R R'$ 且 $\mathfrak n' = \mathfrak nS$。若
\begin{enumerate}
\item 映射 $R \to R'$ 将 $R$ 作为 $R'$ 的 $R$-模直和项嵌入；以及
\item $R' \to S'$ 对 $\mathfrak n'$-adic 拓扑形式光滑，
\end{enumerate}
则 $R \to S$ 对 $\mathfrak n$-adic 拓扑形式光滑。
\end{lemma}

\begin{proof}
给定定义 \ref{more-algebra-definition-formally-smooth} 中的实线图
$$
\xymatrix{
S \ar[r] & A/J \\
R \ar[u] \ar[r] & A \ar[u]
}
$$
令 $A' = A \otimes_R R'$，并令
$J' = \Im(J \otimes_R R' \to A')$。上图的基变换是具有连续箭头的图
$$
\xymatrix{
S' \ar[r] \ar@{-->}[rd]^{\psi'} & A'/J' \\
R' \ar[u] \ar[r] & A' \ar[u]
}
$$
由条件 (2)，得到虚线箭头 $\psi' : S' \to A'$。
利用条件 (1)，选取 $R$-模的直和分解 $R' = R \oplus C$。
（警告：$C$ 不是 $R'$ 中的理想。）于是
$A' = A \oplus A \otimes_R C$。令
$$
J'' = \Im(J \otimes_R C \to A \otimes_R C) \subset J' \subset A'.
$$
于是作为 $A$-模有 $J' = J \oplus J''$。将 $\psi'$ 与
$S \to S'$ 复合所得的 $\psi : S \to A'$，其像包含于
$A + J' = A \oplus J''$。然而，在环 $A + J' = A \oplus J''$ 中，
$A$-子模 $J''$ 是理想！（使用 $J^2 = 0$。）因此复合
$S \to A + J' \to (A + J')/J'' = A$ 就是所求箭头。
\end{proof}







\section{局部环的形式光滑映射}
\label{more-algebra-section-formally-smooth-local}

\noindent
对局部环之间的局部同态，可以限制需要检验提升性质的图。
请与 Algebra, Lemma \ref{algebra-lemma-smooth-test-artinian} 比较。

\begin{lemma}
\label{more-algebra-lemma-fs-local}
设 $(R, \mathfrak m) \to (S, \mathfrak n)$ 为局部环之间的局部同态。
下列条件等价：
\begin{enumerate}
\item $R \to S$ 对 $\mathfrak n$-adic 拓扑形式光滑；
\item 对局部环之间局部同态的每个实线交换图
$$
\xymatrix{
S \ar[r] \ar@{-->}[rd] & A/J \\
R \ar[r] \ar[u] & A \ar[u]
}
$$
其中 $J \subset A$ 是平方为零的理想，对某个 $n > 0$ 有
$\mathfrak m_A^n = 0$，并且 $S \to A/J$ 在剩余域上诱导同构，
都存在使该图交换的虚线箭头。
\end{enumerate}
若 $S$ 是 Noether 环，这些条件还等价于：
\begin{enumerate}
\item[(3)] 与 (2) 相同，但仅要求这样的图：附加地，$A \to A/J$ 是小扩张
（Algebra, Definition \ref{algebra-definition-small-extension}）。
\end{enumerate}
\end{lemma}

\begin{proof}
推含 (1) $\Rightarrow$ (2) 由定义得出。考虑定义
\ref{more-algebra-definition-formally-smooth} 中对 $R$ 上的 $\mathfrak m$-adic 拓扑及
$S$ 上的 $\mathfrak n$-adic 拓扑的一个图
$$
\xymatrix{
S \ar[r] \ar@{-->}[rd] & A/J \\
R \ar[r] \ar[u] & A \ar[u]
}
$$
选取 $m > 0$ 使 $\mathfrak n^m(A/J) = 0$
（由图中映射的连续性，这是可能的）。设 $A'$ 为 $A$ 的子环，
即 $S$ 在 $A/J$ 中的像的逆像。把 $J' = J$ 视为 $A'$ 中的理想。
则 $J'$ 是 $A'$ 中平方为零的理想，且 $A'/J'$ 是 $S/\mathfrak n^m$ 的商。
故 $A'$ 是局部环，并且 $\mathfrak m_{A'}^{2m} = 0$。于是得到 (2) 中的图
$$
\xymatrix{
S \ar[r] \ar@{-->}[rd] & A'/J' \\
R \ar[r] \ar[u] & A' \ar[u]
}
$$
若能在此图中构造虚线箭头，再与 $A' \to A$ 复合，便得到原图中的虚线箭头。
这样便看出 (2) 推出 (1)。

\medskip\noindent
假设 $S$ Noether。推含 (1) $\Rightarrow$ (3) 是立即的。
假设 (3)，并设给定 (2) 中的图。则对某个 $n > 0$ 有
$\mathfrak m_A^n J = 0$。考虑映射
$$
A \to A/\mathfrak m_A^{n - 1}J \to \ldots \to A/\mathfrak mJ \to A/J
$$
可知，只需在 $\mathfrak m_A J = 0$ 时给出提升。
假设 $\mathfrak m_A J = 0$，并令 $A' \subset A$ 为上面构造的环。
则 $A'/J'$ 是 Artin 的，因为它是 Artin 局部环
$S/\mathfrak n^m$ 的商。因此只需证明：给定性质 (3)，
在 (2) 中满足 $A/J$ Artin 且 $\mathfrak m_A J = 0$ 的图中，
可以找到虚线箭头。设 $\kappa$ 为 $A$、$A/J$ 与 $S$ 的共同剩余域。
由 (3)，若 $J_0 \subset J$ 是满足 $\dim_\kappa(J/J_0) = 1$ 的理想，
则可以给出虚线箭头 $S \to A/J_0$。取乘积得到
$$
S \longrightarrow \prod\nolimits_{J_0 \text{ 如上}} A/J_0
$$
显然，该箭头的像包含于这样的元组成的 $R$-子代数 $A'$：这些元映入小对角
$A/J \subset \prod_{J_0} A/J$。设 $J' \subset A'$ 为映到 $A/J$ 中零元的元。
则 $J'$ 是平方为零的理想，并且作为 $\kappa$-向量空间等于
$$
J' = \prod\nolimits_{J_0 \text{ 如上}} J/J_0
$$
因此映射 $J \to J'$ 是单射。由向量空间理论，可选取分裂
$J' = J \oplus M$。于是作为 $R$-代数有
$$
A' = A \oplus M
$$
故可将映射 $S \to A'$ 与投影 $A' \to A$ 复合，得到所需的虚线箭头，
从而完成引理的证明。
\end{proof}

\noindent
下列引理将在第 \ref{more-algebra-section-regular-fs} 节中得到改进。

\begin{lemma}
\label{more-algebra-lemma-fs-implies-regular}
设 $k$ 为域，$(A, \mathfrak m, K)$ 为 Noether 局部 $k$-代数。
若 $k \to A$ 对 $\mathfrak m$-adic 拓扑形式光滑，
则 $A$ 是正则局部环。
\end{lemma}

\begin{proof}
设 $k_0 \subset k$ 为素域。则 $k_0$ 完美，故 $k / k_0$ 可分，
从而由 Algebra, Lemma \ref{algebra-lemma-formally-smooth-extensions-easy}
形式光滑。由引理 \ref{more-algebra-lemma-formally-smooth} 与
\ref{more-algebra-lemma-compose-formally-smooth}，可知
$k_0 \to A$ 对 $A$ 上的 $\mathfrak m$-adic 拓扑形式光滑。
因此可假设 $k = \mathbf{Q}$ 或 $k = \mathbf{F}_p$。

\medskip\noindent
由 Algebra, Lemmas \ref{algebra-lemma-completion-faithfully-flat} 与
\ref{algebra-lemma-flat-under-regular}，只需证明完备化 $A^\wedge$ 正则。
由引理 \ref{more-algebra-lemma-formally-smooth-completion}，可用 $A^\wedge$ 替换 $A$。
因此可假设 $A$ 是 Noether 完备局部环。由 Cohen 结构定理
（Algebra, Theorem \ref{algebra-theorem-cohen-structure-theorem}），
存在映射 $K \to A$。由于 $k$ 是素域，$K \to A$ 是 $k$-代数映射。

\medskip\noindent
设 $x_1, \ldots, x_n \in \mathfrak m$ 为诸元，其像构成
$\mathfrak m/\mathfrak m^2$ 的一组基。令 $T = K[[X_1, \ldots, X_n]]$。
注意
$$
A/\mathfrak m^2 \cong K[x_1, \ldots, x_n]/(x_ix_j)
$$
以及
$$
T/\mathfrak m_T^2 \cong K[X_1, \ldots, X_n]/(X_iX_j).
$$
令 $A/\mathfrak m^2 \to T/m_T^2$ 为局部 $K$-代数同构，
它将 $x_i$ 的类映到 $X_i$ 的类。以
$f_1 : A \to T/\mathfrak m_T^2$ 表示该同构与商映射
$A \to A/\mathfrak m^2$ 的复合。$k \to A$ 对 $\mathfrak m$-adic 拓扑
形式光滑这一假设意味着，可以把 $f_1$ 提升为映射
$f_2 : A \to T/\mathfrak{m}_T^3$，再提升为映射
$f_3 : A \to T/\mathfrak{m}_T^4$，如此继续，对所有 $n \geq 1$。
警告：映射 $f_n$ 是连续 $k$-代数映射，但未必是 $K$-代数映射。
得到局部 $k$-代数的诱导映射
$f : A \to T = \lim T/\mathfrak m_T^n$。由我们对 $f_1$ 的选择，
$f$ 诱导同构 $\mathfrak m/\mathfrak m^2 \to \mathfrak m_T/\mathfrak m_T^2$，
故每个 $f_n$ 都是满射。因为 $A$ 完备，可推出 $f$ 是满射。
这意味着 $\dim(A) \geq \dim(T) = n$。故按定义，$A$ 正则。
（也可推出 $f$ 是同构。）
\end{proof}

\begin{lemma}
\label{more-algebra-lemma-lift-residue-field}
设 $k$ 为域，$(A, \mathfrak m, \kappa)$ 为完备局部 $k$-代数。
若 $\kappa/k$ 可分，则存在 $k$-代数映射 $\kappa \to A$，使
$\kappa \to A \to \kappa$ 是 $\text{id}_\kappa$。
\end{lemma}

\begin{proof}
由 Algebra, Proposition
\ref{algebra-proposition-characterize-separable-field-extensions}，
扩张 $\kappa/k$ 形式光滑。由引理 \ref{more-algebra-lemma-formally-smooth}，
$k \to \kappa$ 在定义 \ref{more-algebra-definition-formally-smooth} 的意义下形式光滑。
再由引理 \ref{more-algebra-lemma-lift-continuous} 得到 $\kappa \to A$。
\end{proof}

\begin{lemma}
\label{more-algebra-lemma-power-series-over-residue-field}
设 $k$ 为域，$(A, \mathfrak m, \kappa)$ 为完备局部 $k$-代数。
若 $\kappa/k$ 可分且 $A$ 正则，则存在 $k$-代数同构
$A \cong \kappa[[t_1, \ldots, t_d]]$。
\end{lemma}

\begin{proof}
按引理 \ref{more-algebra-lemma-lift-residue-field} 选取 $\kappa \to A$，
并应用 Algebra, Lemma
\ref{algebra-lemma-regular-complete-containing-coefficient-field}。
\end{proof}

\noindent
下列结果将在第 \ref{more-algebra-section-regular-fs} 节中得到改进。

\begin{lemma}
\label{more-algebra-lemma-regular-implies-fs}
设 $k$ 为域，$(A, \mathfrak m, K)$ 为正则局部 $k$-代数，
且 $K/k$ 可分。则 $k \to A$ 对 $\mathfrak m$-adic 拓扑形式光滑。
\end{lemma}

\begin{proof}
只需证明 $A$ 的完备化在 $k$ 上形式光滑；见引理
\ref{more-algebra-lemma-formally-smooth-completion}。因此可假设 $A$ 是完备局部正则
$k$-代数，其剩余域 $K$ 在 $k$ 上可分。由引理
\ref{more-algebra-lemma-power-series-over-residue-field}，可知
$A = K[[x_1, \ldots, x_n]]$。

\medskip\noindent
幂级数环 $K[[x_1, \ldots, x_n]]$ 在 $k$ 上形式光滑。
事实上，$K$ 在 $k$ 上形式光滑，而作为多项式代数，
$K[x_1, \ldots, x_n]$ 在 $K$ 上形式光滑。因此由 Algebra, Lemma
\ref{algebra-lemma-compose-formally-smooth}，
$K[x_1, \ldots, x_n]$ 在 $k$ 上形式光滑。
从而由引理 \ref{more-algebra-lemma-formally-smooth}，
$k \to K[x_1, \ldots, x_n]$ 对 $(x_1, \ldots, x_n)$-adic 拓扑形式光滑。
最后，由引理 \ref{more-algebra-lemma-formally-smooth-completion}，
$k \to K[[x_1, \ldots, x_n]]$ 对 $(x_1, \ldots, x_n)$-adic 拓扑形式光滑。
\end{proof}

\begin{lemma}
\label{more-algebra-lemma-formally-smooth-finite-type}
设 $A \to B$ 为有限型环映射，且 $A$ Noether。
设 $\mathfrak q \subset B$ 是位于 $\mathfrak p \subset A$ 之上的素理想。
下列条件等价：
\begin{enumerate}
\item $A \to B$ 在 $\mathfrak q$ 处光滑；以及
\item $A_\mathfrak p \to B_\mathfrak q$ 对 $\mathfrak q$-adic 拓扑形式光滑。
\end{enumerate}
\end{lemma}

\begin{proof}
推含 (2) $\Rightarrow$ (1) 由 Algebra, Lemma
\ref{algebra-lemma-smooth-test-artinian} 得出。反之，若 $A \to B$
在 $\mathfrak q$ 处光滑，则对某个 $g \in B$、
$g \not \in \mathfrak q$，$A \to B_g$ 光滑。于是由 Algebra, Proposition
\ref{algebra-proposition-smooth-formally-smooth}，$A \to B_g$ 形式光滑。
因此 $A_\mathfrak p \to B_\mathfrak q$ 形式光滑，因为局部化保持形式光滑
（例如，由 Algebra, Proposition
\ref{algebra-proposition-characterize-formally-smooth} 的判据，
以及余切复形对局部化表现良好这一事实；见 Algebra, Lemmas
\ref{algebra-lemma-NL-localize-bottom} 与
\ref{algebra-lemma-localize-NL}）。最后，由引理
\ref{more-algebra-lemma-formally-smooth}，$A_\mathfrak p \to B_\mathfrak q$
对 $\mathfrak q$-adic 拓扑形式光滑。
\end{proof}









\section{幂级数环的一些结果}
\label{more-algebra-section-power-series}

\noindent
Noether 局部环之间形式光滑映射的问题，往往可以归结为幂级数环之间
映射的问题。本节证明若干辅助引理，以便进行这类论证。

\begin{lemma}
\label{more-algebra-lemma-power-series-ring-over-Cohen-fs}
设 $K$ 是特征为 $0$ 的域，且 $A = K[[x_1, \ldots, x_n]]$。
设 $L$ 是特征为 $p > 0$ 的域，且 $B = L[[x_1, \ldots, x_n]]$。
设 $\Lambda$ 为 Cohen 环。令 $C = \Lambda[[x_1, \ldots, x_n]]$。
\begin{enumerate}
\item $\mathbf{Q} \to A$ 对 $\mathfrak m_A$-adic 拓扑形式光滑。
\item $\mathbf{F}_p \to B$ 对 $\mathfrak m_B$-adic 拓扑形式光滑。
\item $\mathbf{Z} \to C$ 对 $\mathfrak m_C$-adic 拓扑形式光滑。
\end{enumerate}
\end{lemma}

\begin{proof}
由幂级数环的泛性质，只需证明：
\begin{enumerate}
\item $\mathbf{Q} \to K$ 形式光滑。
\item $\mathbf{F}_p \to L$ 形式光滑。
\item $\mathbf{Z} \to \Lambda$ 对 $\mathfrak m_\Lambda$-adic 拓扑形式光滑。
\end{enumerate}
前两项即 Algebra, Proposition
\ref{algebra-proposition-characterize-separable-field-extensions}。
第三项由 Algebra, Lemma
\ref{algebra-lemma-cohen-ring-formally-smooth} 得出，因为对于如定义
\ref{more-algebra-definition-formally-smooth} 中的任意检验图，$p$ 的某个幂在
$A/J$ 中为零，因而 $p$ 的某个幂在 $A$ 中也为零。
\end{proof}

\begin{lemma}
\label{more-algebra-lemma-quotient-power-series-ring-over-Cohen}
设 $K$ 为域，且 $A = K[[x_1, \ldots, x_n]]$。
设 $\Lambda$ 为 Cohen 环，并令 $B = \Lambda[[x_1, \ldots, x_n]]$。
\begin{enumerate}
\item 若 $y_1, \ldots, y_n \in A$ 是一组正则参数系，则
$K[[y_1, \ldots, y_n]] \to A$ 是同构。
\item 若 $z_1, \ldots, z_r \in A$ 构成 $A$ 的一组正则参数系的一部分，
则 $r \leq n$，且
$A/(z_1, \ldots, z_r) \cong K[[y_1, \ldots, y_{n - r}]]$。
\item 若 $p, y_1, \ldots, y_n \in B$ 是一组正则参数系，则
$\Lambda[[y_1, \ldots, y_n]] \to B$ 是同构。
\item 若 $p, z_1, \ldots, z_r \in B$ 构成 $B$ 的一组正则参数系的一部分，
则 $r \leq n$，且
$B/(z_1, \ldots, z_r) \cong \Lambda[[y_1, \ldots, y_{n - r}]]$。
\end{enumerate}
\end{lemma}

\begin{proof}
证明 (1)。令 $A' = K[[y_1, \ldots, y_n]]$。显然，对每个 $n \geq 1$，
映射 $A' \to A$ 诱导同构
$A'/\mathfrak m_{A'}^n \to A/\mathfrak m_A^n$。
由于 $A$ 与 $A'$ 均完备，故 $A' \to A$ 是同构。
证明 (2)。将 $z_1, \ldots, z_r$ 扩充成 $A$ 的一组正则参数系
$z_1, \ldots, z_r, y_1, \ldots, y_{n - r}$。考虑映射
$A' = K[[z_1, \ldots, z_r, y_1, \ldots, y_{n - r}]] \to A$。
由 (1)，这是同构。又显然
$A'/(z_1, \ldots, z_r) \cong K[[y_1, \ldots, y_{n - r}]]$，
故 (2) 成立。(3) 与 (4) 的证明分别同 (1) 与 (2) 完全相同。
\end{proof}

\begin{lemma}
\label{more-algebra-lemma-embed-map-Noetherian-complete-local-rings}
设 $A \to B$ 为 Noether 完备局部环的局部同态。则存在交换图
$$
\xymatrix{
S \ar[r] & B \\
R \ar[u] \ar[r] & A \ar[u]
}
$$
并具有以下性质：
\begin{enumerate}
\item 水平箭头均为满射；
\item 若 $A/\mathfrak m_A$ 的特征为零，则 $S$ 与 $R$
都是域上的幂级数环；
\item 若 $A/\mathfrak m_A$ 的特征为 $p > 0$，则 $S$ 与 $R$
都是 Cohen 环上的幂级数环；以及
\item $R \to S$ 将 $R$ 的一组正则参数系映为 $S$
的一组正则参数系的一部分。
\end{enumerate}
特别地，$R \to S$ 平坦（见 Algebra, Lemma
\ref{algebra-lemma-flat-over-regular}），且其纤维
$S/\mathfrak m_R S$ 正则（见 Algebra, Lemma
\ref{algebra-lemma-regular-ring-CM}）。
\end{lemma}

\begin{proof}
利用 Cohen 结构定理（Algebra, Theorem
\ref{algebra-theorem-cohen-structure-theorem}），按引理所述选取满射
$S \to B$：若剩余特征为 $p > 0$，则取 $S$ 为 Cohen 环上的幂级数环；
否则取它为域上的幂级数环。设 $J \subset S$ 为 $S \to B$ 的核。
接着选取满射 $R = \Lambda[[x_1, \ldots, x_n]] \to A$：若 $A$
的剩余特征为 $p > 0$，则取 $\Lambda$ 为 Cohen 环；否则取
$\Lambda$ 为 $A$ 的剩余域。
我们将复合映射 $\Lambda[[x_1, \ldots, x_n]] \to A \to B$
提升为映射 $\varphi : R \to S$。这是可行的，因为
$\Lambda[[x_1, \ldots, x_n]]$ 对 $\mathbf{Z}$ 在
$\mathfrak m$-adic 拓扑下形式光滑（见引理
\ref{more-algebra-lemma-power-series-ring-over-Cohen-fs}）；应用引理
\ref{more-algebra-lemma-lift-continuous} 即得。
最后，以映射
$\varphi' : R = \Lambda[[x_1, \ldots, x_n]] \to S' = S[[y_1, \ldots, y_n]]$
替换 $\varphi$，其中 $\varphi'|_\Lambda = \varphi|_\Lambda$ 且
$\varphi'(x_i) = \varphi(x_i) + y_i$。同时，以将 $y_i$ 映为零的映射
$S' \to B$ 替换 $S \to B$。替换之后显然，$R$ 的一组正则参数系
被映为 $S'$ 中某个正则序列的一部分，结论得证。
\end{proof}

\noindent
下面这个引理应当有一个初等证明。

\begin{lemma}
\label{more-algebra-lemma-dominate-two-surjections}
设 $S \to R$ 与 $S' \to R$ 是完备 Noether 局部环的满射。
则 $S \times_R S'$ 是完备 Noether 局部环。
\end{lemma}

\begin{proof}
设 $k$ 为 $R$ 的剩余域。若 $k$ 的特征为 $p > 0$，则以 $\Lambda$
表示一个剩余域为 $k$ 的 Cohen 环
（Algebra, Definition \ref{algebra-definition-cohen-ring}；
Algebra, Lemma \ref{algebra-lemma-cohen-rings-exist}）。
若 $k$ 的特征为 $0$，则令 $\Lambda = k$。
选取满射 $\Lambda[[x_1, \ldots, x_n]] \to R$
（同 Cohen 结构定理中一样，见 Algebra, Theorem
\ref{algebra-theorem-cohen-structure-theorem}），并利用引理
\ref{more-algebra-lemma-power-series-ring-over-Cohen-fs} 与
\ref{more-algebra-lemma-lift-continuous} 将其提升为映射
% P02-E0167
$\Lambda[[x_1, \ldots, x_n]] \to S$ 以及
$\varphi : \Lambda[[x_1, \ldots, x_n]] \to S$ 和
$\varphi' : \Lambda[[x_1, \ldots, x_n]] \to S'$。
接着，选取 $f_1, \ldots, f_m \in S$ 生成 $S \to R$ 的核，并选取
$f'_1, \ldots, f'_{m'} \in S'$ 生成 $S' \to R$ 的核。于是映射
% P02-E0168
$$
\Lambda[[x_1, \ldots, x_n, y_1, \ldots, y_m, z_1, \ldots, z_{m'}]]
\longrightarrow S \times_R S,
$$
将 $x_i$ 映为 $(\varphi(x_i), \varphi'(x_i))$，
将 $y_j$ 映为 $(f_j, 0)$，并将
% P02-E0169
$z_{j'}$ 映为 $(0, f'_j)$；这个映射是满射。因此
$S \times_R S'$ 是某个完备局部环的商，故为完备环。
\end{proof}







\section{几何正则性与形式光滑性}
\label{more-algebra-section-regular-fs}

\noindent
本节综合前几节的结果，证明域上几何正则局部环的如下刻画。
随后，我们重新运用其中一些论证，证明 Noether 局部环之间在
$\mathfrak m$-adic 拓扑下形式光滑的映射的一种刻画。

\begin{theorem}
\label{more-algebra-theorem-regular-fs}
设 $k$ 为域。设 $(A, \mathfrak m, K)$ 为 Noether 局部
$k$-代数。若 $k$ 的特征为零，则下列条件等价：
\begin{enumerate}
\item $A$ 是正则局部环；以及
\item $k \to A$ 对 $\mathfrak m$-adic 拓扑形式光滑。
\end{enumerate}
若 $k$ 的特征为 $p > 0$，则下列条件等价：
\begin{enumerate}
\item $A$ 在 $k$ 上几何正则；
\item $k \to A$ 对 $\mathfrak m$-adic 拓扑形式光滑；
\item 对每个有限于 $k$ 的 $k \subset k' \subset k^{1/p}$，
环 $A \otimes_k k'$ 正则；
\item $A$ 正则，且典范映射
$H_1(L_{K/k}) \to \mathfrak m/\mathfrak m^2$ 为单射；以及
\item $A$ 正则，且映射
$\Omega_{k/\mathbf{F}_p} \otimes_k K \to \Omega_{A/\mathbf{F}_p} \otimes_A K$
为单射。
\end{enumerate}
\end{theorem}

\begin{proof}
若 $k$ 的特征为零，则 (1) 与 (2) 的等价性由引理
\ref{more-algebra-lemma-fs-implies-regular} 和 \ref{more-algebra-lemma-regular-implies-fs} 得出。

\medskip\noindent
若 $k$ 的特征为 $p > 0$，则由命题
\ref{more-algebra-proposition-characterization-geometrically-regular}，
(1)、(3)、(4) 与 (5) 等价。假设 (2) 成立。由引理
\ref{more-algebra-lemma-base-change-fs} 可见，
$k' \to A' = A \otimes_k k'$ 对
$\mathfrak m' = \mathfrak mA'$-adic 拓扑形式光滑。因此，若
$k \subset k'$ 为有限纯不可分扩张，则由引理
\ref{more-algebra-lemma-fs-implies-regular}，$A'$ 是正则局部环。故 (1) 成立。

\medskip\noindent
最后证明 (5) 推出 (2)。取定义 \ref{more-algebra-definition-formally-smooth}
中的如下实线图
$$
\xymatrix{
A \ar[r]_{\bar\psi} \ar@{-->}[rd] & B/J \\
k \ar[u]^i \ar[r]^\varphi & B \ar[u]_\pi
}
$$
由于 $J^2 = 0$，$J$ 具有典范的 $B/J$-模结构，并通过
$\bar\psi$ 具有 $A$-模结构。由于 $\bar\psi$ 对
$\mathfrak m$-adic 拓扑连续，故对某个 $n$ 有
$\mathfrak m^nJ = 0$。于是可以用 $B/J$-子模给 $J$ 作过滤
$0 \subset J_1 \subset J_2 \subset \ldots \subset J_n = J$，
使每个商 $J_{t + 1}/J_t$ 均被 $\mathfrak m$ 零化。
考察环映射序列
$B \to B/J_1 \to B/J_2 \to \ldots \to B/J$，
可见只需在 $J$ 被 $\mathfrak m$ 零化时，即 $J$ 是
$K$-向量空间时，证明虚线箭头存在。

\medskip\noindent
现设给定如上的图，并且 $J$ 被 $\mathfrak m$ 零化。
由引理 \ref{more-algebra-lemma-regular-implies-fs} 可见，
$\mathbf{F}_p \to A$ 对 $\mathfrak m$-adic 拓扑形式光滑。
因此存在环映射 $\psi : A \to B$，使得
$\pi \circ \psi = \bar \psi$。于是
$\psi \circ i, \varphi : k \to B$ 是两个与 $\pi$ 复合后相等的映射。
故 $D = \psi \circ i - \varphi : k \to J$ 是导子。
由 Algebra, Lemma \ref{algebra-lemma-universal-omega}，
可对某个 $k$-线性映射
$\xi : \Omega_{k/\mathbf{F}_p} \to J$ 写成
$D = \xi \circ \text{d}$。利用 $J$ 的 $K$-向量空间结构，
将 $\xi$ 延拓为 $K$-线性映射
$\xi' : \Omega_{k/\mathbf{F}_p} \otimes_k K \to J$。
利用 (5)，可以找到一个 $K$-线性映射
% P02-E0170
$\xi'' : \Omega_{A/\mathbf{F}_p} \otimes_A K$，其在
$\Omega_{k/\mathbf{F}_p} \otimes_k K$ 上的限制为 $\xi'$。记
$$
D' : A \xrightarrow{\text{d}} \Omega_{A/\mathbf{F}_p}
\to \Omega_{A/\mathbf{F}_p} \otimes_A K \xrightarrow{\xi''} J.
$$
最后令 $\psi' = \psi - D' : A \to B$。读者可验证，$\psi'$
是满足 $\pi \circ \psi' = \bar \psi$ 且
$\psi' \circ i = \varphi$ 的环映射，正如所需。
\end{proof}

\begin{example}
\label{more-algebra-example-fs}
设 $k$ 是特征为 $p > 0$ 的域。设 $a \in k$ 不是 $p$ 次幂。
一个剩余域在 $k$ 上纯不可分的几何正则局部 $k$-代数的标准例子是
$$
A = k[x, y]_{(x, y^p - a)}/(y^p - a - x)
$$
事实上，$A$ 是 $k$ 上某个光滑代数的局部化，故 $k \to A$
形式光滑，从而 $k \to A$ 对 $\mathfrak m$-adic 拓扑形式光滑。
下面是一个密切相关的例子。令
$k = \mathbf{F}_p(s)$ 且 $K = \mathbf{F}_p(t)^{perf}$。
我们断言环映射
$$
k \longrightarrow A = K[[x]],\quad s \longmapsto t + x
$$
对 $A$ 上的 $(x)$-adic 拓扑形式光滑。事实上，
$\Omega_{k/\mathbf{F}_p}$ 是以 $\text{d}s$ 为基的 $1$ 维向量空间。
它被映到 $\Omega_{A/\mathbf{F}_p}$ 中的元素
$\text{d}x + \text{d}t = \text{d}x$。
留给读者证明，$\Omega_{A/\mathbf{F}_p}$ 作为 $A$-模以
$\text{d}x$ 为基自由。因此，定理 \ref{more-algebra-theorem-regular-fs}
的条件 (5) 成立，从而 $k \to A$ 对 $(x)$-adic 拓扑形式光滑。
\end{example}

\begin{lemma}
\label{more-algebra-lemma-formally-smooth-flat}
设 $A \to B$ 为 Noether 局部环的局部同态。
假设 $A \to B$ 对 $\mathfrak m_B$-adic 拓扑形式光滑。
则 $A \to B$ 平坦。
\end{lemma}

\begin{proof}
% P02-E0171
由引理 \ref{more-algebra-lemma-formally-smooth-completion} 以及
Algebra, Lemma
\ref{algebra-lemma-completion-Noetherian-Noetherian}，
可以假设 $A$ 与 $B$ 是 Noether 完备局部环
（这里还用到 Algebra, Lemmas
\ref{algebra-lemma-flatness-descends-more-general} 和
\ref{algebra-lemma-completion-faithfully-flat}，以说明完备化之间的映射
平坦蕴含 $A \to B$ 平坦）。
选取引理 \ref{more-algebra-lemma-embed-map-Noetherian-complete-local-rings}
中的交换图
$$
\xymatrix{
S \ar[r] & B \\
R \ar[u] \ar[r] & A \ar[u]
}
$$
其中 $R \to S$ 平坦。设 $I \subset R$ 为 $R \to A$ 的核。
由于 $B$ 在 $A$ 上形式光滑，可见 $A$-代数映射
$$
S/IS \longrightarrow B
$$
有截面；见引理 \ref{more-algebra-lemma-lift-continuous}。
因此，$B$ 是平坦 $A$-模 $S/IS$ 的直和因子
（由平坦性的基变换，见 Algebra, Lemma
\ref{algebra-lemma-flat-base-change}），故为平坦模。
\end{proof}

\begin{lemma}
\label{more-algebra-lemma-formally-smooth-JZ}
设 $A \to B$ 为 Noether 局部环的局部同态。
假设 $A \to B$ 对 $\mathfrak m_B$-adic 拓扑形式光滑。
设 $K$ 为 $B$ 的剩余域。则 $A \to B \to K$ 的
Jacobi--Zariski 序列给出正合序列
$$
0 \to H_1(\NL_{K/A}) \to \mathfrak m_B/\mathfrak m_B^2
\to \Omega_{B/A} \otimes_B K \to \Omega_{K/A} \to 0
$$
\end{lemma}

\begin{proof}
由 Algebra, Lemma \ref{algebra-lemma-NL-surjection}，注意到
$\mathfrak m_B/\mathfrak m_B^2 = H_1(\NL_{K/B})$。
根据 Algebra, Lemma \ref{algebra-lemma-exact-sequence-NL}，
只需证明
$H_1(\NL_{K/A}) \to \mathfrak m_B/\mathfrak m_B^2$ 为单射。
设 $k$ 为 $A$ 的剩余域。$A \to k \to K$ 的
Jacobi--Zariski 序列给出 $\Omega_{K/A} = \Omega_{K/k}$ 以及正合序列
$$
\mathfrak m_A/\mathfrak m_A^2 \otimes_k K \to
H_1(\NL_{K/A}) \to
H_1(\NL_{K/k}) \to 0
$$
令 $\overline{B} = B \otimes_A k$。由于 $\overline{B}$ 正则，
理想 $\mathfrak m_{\overline{B}}$ 由一个正则序列生成。
对 $\mathfrak m_A B \subset \mathfrak m_B$ 应用引理
\ref{more-algebra-lemma-conormal-sequence-H1-regular} 和
\ref{more-algebra-lemma-noetherian-finite-all-equivalent}，得到
$\mathfrak m_A B / (\mathfrak m_AB \cap \mathfrak m_B^2) =
\mathfrak m_A B / \mathfrak m_A \mathfrak m_B$。
又由引理 \ref{more-algebra-lemma-formally-smooth-flat}，$A \to B$ 平坦，
所以上式等于 $\mathfrak m_A/\mathfrak m_A^2 \otimes_k K$。
因此得到短正合序列
$$
0 \to
\mathfrak m_A/\mathfrak m_A^2 \otimes_k K \to
\mathfrak m_B/\mathfrak m_B^2 \to
\mathfrak m_{\overline{B}}/\mathfrak m_{\overline{B}}^2 \to 0
$$
Jacobi--Zariski 序列的函子性表明，我们得到交换图
$$
\xymatrix{
&
\mathfrak m_A/\mathfrak m_A^2 \otimes_k K \ar[d] \ar[r] &
H_1(\NL_{K/A}) \ar[d] \ar[r] &
H_1(\NL_{K/k}) \ar[d] \ar[r] & 0 \\
0 \ar[r] &
\mathfrak m_A/\mathfrak m_A^2 \otimes_k K \ar[r] &
\mathfrak m_B/\mathfrak m_B^2 \ar[r] &
\mathfrak m_{\overline{B}}/\mathfrak m_{\overline{B}}^2 \ar[r] & 0
}
$$
由定理 \ref{more-algebra-theorem-regular-fs}，左侧竖直箭头为单射，因为由引理
\ref{more-algebra-lemma-base-change-fs}，$k \to \overline{B}$ 对
$\mathfrak m_{\overline{B}}$-adic 拓扑形式光滑。
蛇形引理即完成证明。
\end{proof}

\begin{proposition}
\label{more-algebra-proposition-fs-flat-fibre-fs}
设 $A \to B$ 为 Noether 局部环的局部同态。
设 $k$ 为 $A$ 的剩余域，并设
$\overline{B} = B \otimes_A k$ 为特殊纤维。下列条件等价：
\begin{enumerate}
\item $A \to B$ 平坦，且 $\overline{B}$ 在 $k$ 上几何正则；
\item $A \to B$ 平坦，且 $k \to \overline{B}$ 对
$\mathfrak m_{\overline{B}}$-adic 拓扑形式光滑；以及
\item $A \to B$ 对 $\mathfrak m_B$-adic 拓扑形式光滑。
\end{enumerate}
\end{proposition}

\begin{proof}
(1) 与 (2) 的等价性由定理 \ref{more-algebra-theorem-regular-fs} 得出。

\medskip\noindent
假设 (3)。由引理 \ref{more-algebra-lemma-formally-smooth-flat}，$A \to B$ 平坦。
由引理 \ref{more-algebra-lemma-base-change-fs}，$k \to \overline{B}$ 对
$\mathfrak m_{\overline{B}}$-adic 拓扑形式光滑。故 (2) 成立。

\medskip\noindent
假设 (2)。引理 \ref{more-algebra-lemma-formally-smooth-completion}
说明形式光滑性在完备化下保持。由 Algebra, Lemma
\ref{algebra-lemma-completion-faithfully-flat}，平坦性亦然。
因此，可分别用 $A$ 与 $B$ 的完备化替换它们，并假设
$A$ 与 $B$ 是 Noether 完备局部环。在这种情况下，选取引理
\ref{more-algebra-lemma-embed-map-Noetherian-complete-local-rings} 中的图
$$
\xymatrix{
S \ar[r] & B \\
R \ar[u] \ar[r] & A \ar[u]
}
$$
下文将使用此图的全部性质而不再一一说明。
固定 $R$ 的一组正则参数系 $t_1, \ldots, t_d$；当 $k$
的特征为 $p > 0$ 时，令 $t_1 = p$。
令 $\overline{S} = S \otimes_R k$。考察短正合序列
$$
0 \to J \to S \to B \to 0
$$
由于 $\overline{B}$ 与 $\overline{S}$ 均正则，
$\overline{S} \to \overline{B}$ 的核由元素
$\overline{x}_1, \ldots, \overline{x}_r$ 生成，且这些元素构成
$\overline{S}$ 的一组正则参数系的一部分；见 Algebra, Lemma
\ref{algebra-lemma-regular-quotient-regular}。
将这些元素提升为 $x_1, \ldots, x_r \in J$。于是
$t_1, \ldots, t_d, x_1, \ldots, x_r$ 构成 $S$
的一组正则参数系的一部分。因此，若 $k$ 的特征为零，则
$S/(x_1, \ldots, x_r)$ 是域上的幂级数环；若 $k$ 的特征为
$p > 0$，则它是 Cohen 环上的幂级数环；见引理
\ref{more-algebra-lemma-quotient-power-series-ring-over-Cohen}。
此外，$R \to S/(x_1, \ldots, x_r)$ 仍将
$t_1, \ldots, t_d$ 映为 $S/(x_1, \ldots, x_r)$
的一组正则参数系的一部分。换言之，可以用
$S/(x_1, \ldots, x_r)$ 替换 $S$，并假设我们有引理
\ref{more-algebra-lemma-embed-map-Noetherian-complete-local-rings} 中的图
$$
\xymatrix{
S \ar[r] & B \\
R \ar[u] \ar[r] & A \ar[u]
}
$$
而且 $\overline{S} = \overline{B}$。此时映射
$$
S \otimes_R A \longrightarrow B
$$
是同构，因为它是满射、在特殊纤维上是同构，并且源与目标均在
$A$ 上平坦（例如可用 Algebra, Lemma
\ref{algebra-lemma-mod-injective}；也可将短正合序列
$0 \to I \to S \otimes_R A \to B \to 0$ 在 $A$ 上与 $k$
张量，得到 $I \otimes_A k = 0$，从而由 Nakayama 引理有 $I = 0$）。
于是根据引理 \ref{more-algebra-lemma-base-change-fs}，只需证明
$R \to S$ 对 $\mathfrak m_S$-adic 拓扑形式光滑。
当然，由于 $\overline{S} = \overline{B}$，可知
$\overline{S}$ 在 $k = R/\mathfrak m_R$ 上形式光滑。

\medskip\noindent
选取元素 $y_1, \ldots, y_m \in S$，使
$t_1, \ldots, t_d, y_1, \ldots, y_m$ 是 $S$ 的一组正则参数系。
若 $k$ 的特征为零，则选取系数域 $K \subset S$；若 $k$
的特征为 $p > 0$，则选取剩余域为 $K$ 的 Cohen 环
$\Lambda \subset S$。此时，映射
$K[[t_1, \ldots, t_d, y_1, \ldots, y_m]] \to S$
（特征为零的情形）或
$\Lambda[[t_2, \ldots, t_d, y_1, \ldots, y_m]] \to S$
（特征为 $p > 0$ 的情形）是同构；见引理
\ref{more-algebra-lemma-quotient-power-series-ring-over-Cohen}。
从现在起，将 $S$ 视为上述幂级数环。

\medskip\noindent
余下的证明类似于定理 \ref{more-algebra-theorem-regular-fs} 证明中的论证。
取定义 \ref{more-algebra-definition-formally-smooth} 中的如下实线图
$$
\xymatrix{
S \ar[r]_{\bar\psi} \ar@{-->}[rd] & N/J \\
R \ar[u]^i \ar[r]^\varphi & N \ar[u]_\pi
}
$$
由于 $J^2 = 0$，$J$ 具有典范的 $N/J$-模结构，并通过
$\bar\psi$ 具有 $S$-模结构。由于 $\bar\psi$ 对
$\mathfrak m_S$-adic 拓扑连续，故对某个 $n$ 有
$\mathfrak m_S^nJ = 0$。于是可以用 $N/J$-子模给 $J$ 作过滤
$0 \subset J_1 \subset J_2 \subset \ldots \subset J_n = J$，
使每个商 $J_{t + 1}/J_t$ 均被 $\mathfrak m_S$ 零化。
考察环映射序列
$N \to N/J_1 \to N/J_2 \to \ldots \to N/J$，
可见只需在 $J$ 被 $\mathfrak m_S$ 零化时，即 $J$ 是
$K$-向量空间时，证明虚线箭头存在。

\medskip\noindent
现设给定如上的图，并且 $J$ 被 $\mathfrak m_S$ 零化。
由于 $\mathbf{Q} \to S$（特征为零的情形）或
$\mathbf{Z} \to S$（特征为 $p > 0$ 的情形）对
$\mathfrak m_S$-adic 拓扑形式光滑（见引理
\ref{more-algebra-lemma-power-series-ring-over-Cohen-fs}），故存在环映射
$\psi : S \to N$，使得 $\pi \circ \psi = \bar \psi$。
由于 $S$ 是某个子环上以 $t_1, \ldots, t_d$
（特征为零的情形）或以 $t_2, \ldots, t_d$
（特征为 $p > 0$ 的情形）为变量的幂级数环，故由幂级数环的泛性质，
可改变 $\psi$ 的选择，使 $\psi(t_i)$ 等于 $\varphi(t_i)$
（在特征为 $p$ 的情形，对 $t_1 = p$ 这是自动成立的）。
于是 $\psi \circ i$ 与 $\varphi : R \to N$ 是两个与 $\pi$
复合后相等、并且在 $t_1, \ldots, t_d$ 上相同的映射。
因此，$D = \psi \circ i - \varphi : R \to J$ 是将
$t_1, \ldots, t_d$ 零化的导子。
由 Algebra, Lemma \ref{algebra-lemma-universal-omega}，
可对某个将 $\text{d}t_1, \ldots, \text{d}t_d$
（依构造）和 $\mathfrak m_R \Omega_{R/\mathbf{Z}}$
（因为 $J$ 被 $\mathfrak m_R$ 零化）零化的 $R$-线性映射
$\xi : \Omega_{R/\mathbf{Z}} \to J$，写成
$D = \xi \circ \text{d}$。故 $\xi$ 分解为复合
$$
\Omega_{R/\mathbf{Z}} \to \Omega_{k/\mathbf{Z}} \xrightarrow{\xi'} J
$$
其中 $\xi'$ 为 $k$-线性。利用 $J$ 的 $K$-向量空间结构，
将 $\xi'$ 延拓为 $K$-线性映射
$$
\xi'' : \Omega_{k/\mathbf{Z}} \otimes_k K \longrightarrow J.
$$
由于 $\overline{S}/k$ 形式光滑，由定理
\ref{more-algebra-theorem-regular-fs} 可见
$$
\Omega_{k/\mathbf{Z}} \otimes_k K \to
\Omega_{\overline{S}/\mathbf{Z}} \otimes_S K
$$
为单射（在特征为零的情形也成立，因为此时甚至
$\Omega_{k/\mathbf{Z}} \to \Omega_{K/\mathbf{Z}}$
也是单射；见 Algebra, Proposition
\ref{algebra-proposition-characterize-separable-field-extensions}）。
因此，存在 $K$-线性映射
$\xi''' : \Omega_{\overline{S}/\mathbf{Z}} \otimes_S K \to J$，
其在 $\Omega_{k/\mathbf{Z}} \otimes_k K$ 上的限制为 $\xi''$。记
$$
D' : S \xrightarrow{\text{d}} \Omega_{S/\mathbf{Z}}
\to \Omega_{\overline{S}/\mathbf{Z}} \to
\Omega_{\overline{S}/\mathbf{Z}} \otimes_S K \xrightarrow{\xi'''} J.
$$
最后令 $\psi' = \psi - D' : S \to N$。读者可验证，$\psi'$
是满足 $\pi \circ \psi' = \bar \psi$ 且
$\psi' \circ i = \varphi$ 的环映射，正如所需。
\end{proof}

\noindent
作为上述结果的一个应用，我们证明形式光滑代数的形变没有障碍。

\begin{lemma}
\label{more-algebra-lemma-lift-fs}
设 $A$ 是剩余域为 $k$ 的 Noether 完备局部环。
设 $B$ 是 Noether 完备局部 $k$-代数。假设 $k \to B$
对 $\mathfrak m_B$-adic 拓扑形式光滑。则存在 Noether 完备局部环
$C$ 以及对 $\mathfrak m_C$-adic 拓扑形式光滑的局部同态
$A \to C$，使得 $C \otimes_A k \cong B$。
\end{lemma}

\begin{proof}
选取引理 \ref{more-algebra-lemma-embed-map-Noetherian-complete-local-rings} 中的图
$$
\xymatrix{
S \ar[r] & B \\
R \ar[u] \ar[r] & A \ar[u]
}
$$
设 $t_1, \ldots, t_d$ 为 $R$ 的一组正则参数系；当 $k$
的特征为 $p > 0$ 时，令 $t_1 = p$。
由于 $B$ 与 $\overline{S} = S \otimes_R k$ 均正则，
$\Ker(\overline{S} \to B)$ 由元素
$\overline{x}_1, \ldots, \overline{x}_r$ 生成，且这些元素构成
$\overline{S}$ 的一组正则参数系的一部分；见 Algebra, Lemma
\ref{algebra-lemma-regular-quotient-regular}。
将这些元素提升为 $x_1, \ldots, x_r \in S$。于是
$t_1, \ldots, t_d, x_1, \ldots, x_r$ 构成 $S$
的一组正则参数系的一部分。因此，若 $k$ 的特征为零，则
$S/(x_1, \ldots, x_r)$ 是域上的幂级数环；若 $k$ 的特征为
$p > 0$，则它是 Cohen 环上的幂级数环；见引理
\ref{more-algebra-lemma-quotient-power-series-ring-over-Cohen}。
此外，$R \to S/(x_1, \ldots, x_r)$ 仍将
$t_1, \ldots, t_d$ 映为 $S/(x_1, \ldots, x_r)$
的一组正则参数系的一部分。换言之，可以用
$S/(x_1, \ldots, x_r)$ 替换 $S$，并假设我们有引理
\ref{more-algebra-lemma-embed-map-Noetherian-complete-local-rings} 中的图
$$
\xymatrix{
S \ar[r] & B \\
R \ar[u] \ar[r] & A \ar[u]
}
$$
而且 $\overline{S} = B$。此时，由命题
\ref{more-algebra-proposition-fs-flat-fibre-fs}，$R \to S$ 对
$\mathfrak m_S$-adic 拓扑形式光滑。因此，由引理
\ref{more-algebra-lemma-base-change-fs}，基变换 $C = S \otimes_R A$
在 $A$ 上对 $\mathfrak m_C$-adic 拓扑形式光滑。
\end{proof}

\begin{remark}
\label{more-algebra-remark-what-does-it-mean}
引理 \ref{more-algebra-lemma-lift-fs} 的断言相当强。具体而言，设有 Noether
完备局部环的局部同态图
$$
\xymatrix{
& B \\
A \ar[r] & A' \ar[u]
}
$$
其中 $A \to A'$ 诱导剩余域同构
$k = A/\mathfrak m_A = A'/\mathfrak m_{A'}$，并且
$B \otimes_{A'} k$ 在 $k$ 上形式光滑。则可将其扩充为 Noether
完备局部环的局部同态交换图
$$
\xymatrix{
C \ar[r] & B \\
A \ar[r] \ar[u] & A' \ar[u]
}
$$
其中 $A \to C$ 对 $\mathfrak m_C$-adic 拓扑形式光滑，并且
$C \otimes_A k \cong B \otimes_{A'} k$。
事实上，根据引理 \ref{more-algebra-lemma-lift-fs}，选取 $A \to C$，
使其提升 $k$ 上的 $B \otimes_{A'} k$。由形式光滑性可得到箭头
$C \to B$；见引理 \ref{more-algebra-lemma-lift-continuous}。
以 $C \otimes_A^\wedge A'$ 表示 $C \otimes_A A'$ 关于理想
$C \otimes_A \mathfrak m_{A'}$ 的完备化。注意，
$C \otimes_A^\wedge A'$ 是 Noether 完备局部环
（见 Algebra, Lemma \ref{algebra-lemma-completion-Noetherian}），
并且在 $A'$ 上平坦（见 Algebra, Lemma
\ref{algebra-lemma-flat-module-powers}）。此外还有：
\begin{enumerate}
\item $C \otimes_A^\wedge A' \to B$ 是满射；
\item 若 $A \to A'$ 是满射，则 $C \to B$ 是满射；
\item 若 $A \to A'$ 有限，则 $C \to B$ 有限；以及
\item 若 $A' \to B$ 平坦，则 $C \otimes_A^\wedge A' \cong B$。
\end{enumerate}
事实上，由关于幂零理想的 Nakayama 引理
（见 Algebra, Lemma \ref{algebra-lemma-NAK}），从
$C \otimes_A k \cong B \otimes_{A'} k$ 可得，对每个 $n$，
$C \otimes_A A'/\mathfrak m_{A'}^n \to B/\mathfrak m_{A'}^nB$
均为满射。这证明了 (1)。(2) 与 (3) 由 (1) 得出。
(4) 由 Algebra, Lemma \ref{algebra-lemma-mod-injective} 得出。
\end{remark}




\section{正则环映射}
\label{more-algebra-section-regular}

\noindent
设 $k$ 为域。回忆：一个 Noether $k$-代数 $A$ 称为在 $k$ 上
{\it 几何正则}，当且仅当对于 $k$ 的每个有限纯不可分扩张 $k'$，
$A \otimes_k k'$ 都正则；见 Algebra, Definition
\ref{algebra-definition-geometrically-regular}。
此外，若此条件成立，则对每个有限生成域扩张 $k'/k$，
$A \otimes_k k'$ 都正则；见 Algebra, Lemma
\ref{algebra-lemma-geometrically-regular}。
我们在以下定义中使用这一概念。

\begin{definition}
\label{more-algebra-definition-regular}
环映射 $R \to \Lambda$ 称为{\it 正则}，如果它平坦，并且对于每个素理想
$\mathfrak p \subset R$，纤维环
$$
\Lambda \otimes_R \kappa(\mathfrak p) =
\Lambda_\mathfrak p/\mathfrak p\Lambda_\mathfrak p
$$
是 Noether 环，且在 $\kappa(\mathfrak p)$ 上几何正则。
\end{definition}

\noindent
若 $R \to \Lambda$ 是环映射且 $\Lambda$ Noether，
则纤维环总是 Noether 环。

\begin{lemma}[正则性是局部性质]
\label{more-algebra-lemma-regular-local}
设 $R \to \Lambda$ 为环映射，且 $\Lambda$ Noether。
下列条件等价：
\begin{enumerate}
\item $R \to \Lambda$ 正则；
\item 对每个在 $\mathfrak p \subset R$ 之上的
$\mathfrak q \subset \Lambda$，$R_\mathfrak p \to \Lambda_\mathfrak q$
正则；以及
\item 对每个在 $R$ 中的 $\mathfrak m$ 之上的极大理想
$\mathfrak m' \subset \Lambda$，
$R_\mathfrak m \to \Lambda_{\mathfrak m'}$ 正则。
\end{enumerate}
\end{lemma}

\begin{proof}
这是因为一个 Noether 环正则，当且仅当它的所有局部环都是正则局部环；
见 Algebra, Definition \ref{algebra-definition-regular}；而一个环映射平坦，
当且仅当所有诱导的局部环映射平坦；见 Algebra, Lemma
\ref{algebra-lemma-flat-localization}。
\end{proof}

\begin{lemma}[正则映射与基变换]
\label{more-algebra-lemma-regular-base-change}
设 $R \to \Lambda$ 为正则环映射。对于任意有限型环映射
$R \to R'$，基变换 $R' \to \Lambda \otimes_R R'$ 也正则。
\end{lemma}

\begin{proof}
平坦性在任意基变换下保持；见 Algebra, Lemma
\ref{algebra-lemma-flat-base-change}。
考察位于 $\mathfrak p \subset R$ 之上的素理想
$\mathfrak p' \subset R'$。由于 $R'$ 在 $R$ 上有限型，
剩余域扩张 $\kappa(\mathfrak p')/\kappa(\mathfrak p)$ 有限生成。
因此，纤维环
$$
(\Lambda \otimes_R R') \otimes_{R'} \kappa(\mathfrak p') =
\Lambda \otimes_R \kappa(\mathfrak p) \otimes_{\kappa(\mathfrak p)}
\kappa(\mathfrak p')
$$
是 Noether 环，这是由 Algebra, Lemma
\ref{algebra-lemma-Noetherian-field-extension} 以及对
$R \to \Lambda$ 的纤维环的假设得出的。
纤维的几何正则性由 Algebra, Lemma
\ref{algebra-lemma-geometrically-regular} 保持。
\end{proof}

\begin{lemma}[正则映射的复合]
\label{more-algebra-lemma-regular-composition}
设 $A \to B$ 与 $B \to C$ 为正则环映射。
若 $A \to C$ 的纤维环是 Noether 环，则 $A \to C$ 正则。
\end{lemma}

\begin{proof}
设 $\mathfrak p \subset A$ 为素理想。设
$\kappa(\mathfrak p) \subset k$ 为有限纯不可分扩张。
我们必须证明 $C \otimes_A k$ 正则。由引理
\ref{more-algebra-lemma-regular-base-change}，可以假设 $A = k$，
从而归结为证明 $C$ 正则。已知 $B$ 正则，且 $B \to C$
平坦并具有正则纤维。于是由 Algebra, Lemma
\ref{algebra-lemma-flat-over-regular-with-regular-fibre}，$C$ 正则。
略去若干细节。
\end{proof}

\begin{lemma}
\label{more-algebra-lemma-colimit-smooth-regular}
设 $R$ 为环。设 $(A_i, \varphi_{ii'})$ 为光滑 $R$-代数的有向系。
令 $\Lambda = \colim A_i$。若对所有 $\mathfrak p \subset R$，
纤维环 $\Lambda \otimes_R \kappa(\mathfrak p)$ 均为 Noether 环，
则 $R \to \Lambda$ 正则。
\end{lemma}

\begin{proof}
由 Algebra, Lemmas \ref{algebra-lemma-colimit-flat} 和
\ref{algebra-lemma-smooth-syntomic}，注意到 $\Lambda$ 在 $R$ 上平坦。
设 $\kappa(\mathfrak p) \subset k$ 为有限纯不可分扩张。注意到
$$
\Lambda \otimes_R \kappa(\mathfrak p) \otimes_{\kappa(\mathfrak p)} k =
\Lambda \otimes_R k = \colim A_i \otimes_R k
$$
是光滑 $k$-代数的余极限；见 Algebra, Lemma
\ref{algebra-lemma-base-change-smooth}。由 Algebra, Lemma
\ref{algebra-lemma-characterize-smooth-over-field}，光滑 $k$-代数的
每个局部环都正则，故由 Algebra, Lemma
\ref{algebra-lemma-colimit-regular}，$\Lambda \otimes_R k$
的所有局部环均正则。这就证明了引理。
\end{proof}

\noindent
下面考察域扩张何时定义正则环映射。

\begin{lemma}
\label{more-algebra-lemma-regular-field-extension}
设 $K/k$ 为域扩张。则 $k \to K$ 是正则环映射，当且仅当
$K$ 是 $k$ 的可分域扩张。
\end{lemma}

\begin{proof}
若 $k \to K$ 正则，则 $K$ 在 $k$ 上几何既约，因而由
Algebra, Proposition
\ref{algebra-proposition-characterize-separable-field-extensions}，
$K$ 在 $k$ 上可分。反之，若 $K/k$ 可分，则由 Algebra, Lemma
\ref{algebra-lemma-colimit-syntomic}，$K$ 是光滑 $k$-代数的余极限，
从而由引理 \ref{more-algebra-lemma-colimit-smooth-regular}，它是正则的。
\end{proof}

\begin{lemma}
\label{more-algebra-lemma-regular-permanence}
设 $A \to B \to C$ 为环映射。若 $A \to C$ 正则，并且
$B \to C$ 平坦且在谱上满射，则 $A \to B$ 正则。
\end{lemma}

\begin{proof}
由 Algebra, Lemma \ref{algebra-lemma-flat-permanence}，
$A \to B$ 平坦。设 $\mathfrak p \subset A$ 为素理想。
环映射
$B \otimes_A \kappa(\mathfrak p) \to C \otimes_A \kappa(\mathfrak p)$
平坦且在谱上满射。因此，由 Algebra, Lemma
\ref{algebra-lemma-geometrically-regular-descent}，
$B \otimes_A \kappa(\mathfrak p)$ 几何正则。
\end{proof}




\section{沿正则环映射上升的性质}
\label{more-algebra-section-ascending-properties}

\noindent
本节类似于 Algebra, Section
\ref{algebra-section-ascending-properties}，但这里环映射
$R \to S$ 是正则的。

\begin{lemma}
\label{more-algebra-lemma-reduced-goes-up}
设 $\varphi : R \to S$ 为环映射。假设
\begin{enumerate}
\item $\varphi$ 正则；
\item $S$ Noether；以及
\item $R$ Noether 且既约。
\end{enumerate}
则 $S$ 既约。
\end{lemma}

\begin{proof}
对于 Noether 环，既约等价于具有性质 $(S_1)$ 和 $(R_0)$；
见 Algebra, Lemma \ref{algebra-lemma-criterion-reduced}。
因此，可以应用 Algebra, Lemmas
\ref{algebra-lemma-Sk-goes-up} 和
\ref{algebra-lemma-Rk-goes-up}。
\end{proof}

\begin{lemma}
\label{more-algebra-lemma-normal-goes-up}
设 $\varphi : R \to S$ 为环映射。假设
\begin{enumerate}
\item $\varphi$ 正则；
\item $S$ Noether；以及
\item $R$ Noether 且正规。
\end{enumerate}
则 $S$ 正规。
\end{lemma}

\begin{proof}
对于 Noether 环，正规等价于具有性质 $(S_2)$ 和 $(R_1)$；
见 Algebra, Lemma \ref{algebra-lemma-criterion-normal}。
因此，可以应用 Algebra, Lemmas
\ref{algebra-lemma-Sk-goes-up} 和
\ref{algebra-lemma-Rk-goes-up}。
\end{proof}

\begin{lemma}
\label{more-algebra-lemma-regular-goes-up}
设 $\varphi : R \to S$ 为环映射。假设
\begin{enumerate}
\item $\varphi$ 正则；
\item $S$ Noether；以及
\item $R$ Noether 且正则。
\end{enumerate}
则 $S$ 正则。
\end{lemma}

\begin{proof}
对于 Noether 环，正则等价于对所有 $k$ 都具有性质 $(R_k)$。
因此，可以应用 Algebra, Lemma \ref{algebra-lemma-Rk-goes-up}。
\end{proof}

\begin{lemma}
\label{more-algebra-lemma-CM-goes-up}
设 $\varphi : R \to S$ 为环映射。假设
\begin{enumerate}
\item $\varphi$ 正则；
\item $S$ Noether；以及
\item $R$ Noether 且为 Cohen--Macaulay 环。
\end{enumerate}
则 $S$ 为 Cohen--Macaulay 环。
\end{lemma}

\begin{proof}
对于 Noether 环，成为 Cohen--Macaulay 环等价于对所有 $k$
都具有性质 $(S_k)$。因此，可以应用 Algebra, Lemma
\ref{algebra-lemma-Sk-goes-up}。
\end{proof}




\section{性质在完备化下的保持性}
\label{more-algebra-section-permanence-completion}

\noindent
给定 Noether 局部环 $(A, \mathfrak m)$，以 $A^\wedge$ 表示
$A$ 关于 $\mathfrak m$ 的完备化。下文将不再说明地使用如下事实：
$A^\wedge$ 是极大理想为
$\mathfrak m^\wedge = \mathfrak m A^\wedge$ 的 Noether
完备局部环，且 $A \to A^\wedge$ 忠实平坦。见 Algebra, Lemmas
\ref{algebra-lemma-completion-Noetherian-Noetherian}、
\ref{algebra-lemma-completion-complete} 和
\ref{algebra-lemma-completion-faithfully-flat}。

\begin{lemma}
\label{more-algebra-lemma-completion-dimension}
设 $A$ 为 Noether 局部环。则
$\dim(A) = \dim(A^\wedge)$。
\end{lemma}

\begin{proof}
由 Algebra, Lemma \ref{algebra-lemma-completion-complete}，
映射 $A \to A^\wedge$ 对 $n \geq 1$ 诱导同构
$A/\mathfrak m^n = A^\wedge/(\mathfrak m^\wedge)^n$。
由 Algebra, Lemma \ref{algebra-lemma-pushdown-module}，这蕴含
$$
\text{length}_A(A/\mathfrak m^n) =
\text{length}_{A^\wedge}(A^\wedge/(\mathfrak m^\wedge)^n)
$$
对所有 $n \geq 1$ 成立。因此 $d(A) = d(A^\wedge)$，
再由 Algebra, Proposition \ref{algebra-proposition-dimension} 得出结论。
另一种证明是使用 Algebra, Lemma
\ref{algebra-lemma-dimension-base-fibre-equals-total}。
\end{proof}

\begin{lemma}
\label{more-algebra-lemma-completion-depth}
设 $A$ 为 Noether 局部环。则
$\text{depth}(A) = \text{depth}(A^\wedge)$。
\end{lemma}

\begin{proof}
见 Algebra, Lemma \ref{algebra-lemma-apply-grothendieck}。
\end{proof}

\begin{lemma}
\label{more-algebra-lemma-completion-CM}
设 $A$ 为 Noether 局部环。则 $A$ 是 Cohen--Macaulay 环，
当且仅当 $A^\wedge$ 是 Cohen--Macaulay 环。
\end{lemma}

\begin{proof}
局部环 $A$ 是 Cohen--Macaulay 环，当且仅当
$\dim(A) = \text{depth}(A)$。由于这两个不变量都在完备化下保持
（引理 \ref{more-algebra-lemma-completion-dimension} 和
\ref{more-algebra-lemma-completion-depth}），断言得证。
\end{proof}

\begin{lemma}
\label{more-algebra-lemma-completion-regular}
设 $A$ 为 Noether 局部环。则 $A$ 正则，当且仅当
$A^\wedge$ 正则。
\end{lemma}

\begin{proof}
若 $A^\wedge$ 正则，则由 Algebra, Lemma
\ref{algebra-lemma-flat-under-regular}，$A$ 正则。
假设 $A$ 正则。设 $\mathfrak m$ 为 $A$ 的极大理想。则
$\dim_{\kappa(\mathfrak m)} \mathfrak m/\mathfrak m^2 =
\dim(A) = \dim(A^\wedge)$
（引理 \ref{more-algebra-lemma-completion-dimension}）。
另一方面，$\mathfrak mA^\wedge$ 是 $A^\wedge$ 的极大理想，
因而 $\mathfrak m_{A^\wedge}$ 至多由 $\dim(A^\wedge)$ 个元素生成。
故 $A^\wedge$ 正则。（也可以使用 Algebra, Lemma
\ref{algebra-lemma-flat-over-regular-with-regular-fibre}。）
\end{proof}

\begin{lemma}
\label{more-algebra-lemma-completion-dvr}
设 $A$ 为 Noether 局部环。则 $A$ 是离散赋值环，当且仅当
$A^\wedge$ 是离散赋值环。
\end{lemma}

\begin{proof}
这由引理 \ref{more-algebra-lemma-completion-dimension}、
\ref{more-algebra-lemma-completion-regular} 和 Algebra, Lemma
\ref{algebra-lemma-characterize-dvr} 得出。
\end{proof}

\begin{lemma}
\label{more-algebra-lemma-completion-reduced}
设 $A$ 为 Noether 局部环。
\begin{enumerate}
\item 若 $A^\wedge$ 既约，则 $A$ 也既约。
\item 一般而言，$A$ 既约并不蕴含 $A^\wedge$ 既约。
\item 若 $A$ 是 Nagata 环，则 $A$ 既约，当且仅当
$A^\wedge$ 既约。
\end{enumerate}
\end{lemma}

\begin{proof}
由于 $A \to A^\wedge$ 忠实平坦，由 Algebra, Lemma
\ref{algebra-lemma-descent-reduced} 得到 (1)。
关于 (2)，见 Algebra, Example
\ref{algebra-example-bad-dvr-char-p}
（特征为零时也有例子，见 Algebra, Remark
\ref{algebra-remark-resolution-dim-1}）。
关于 (3)，见 Algebra, Lemmas
\ref{algebra-lemma-local-nagata-domain-analytically-unramified} 和
\ref{algebra-lemma-analytically-unramified-easy}。
\end{proof}

\begin{lemma}
\label{more-algebra-lemma-completion-normal}
设 $A$ 为 Noether 局部环。若 $A^\wedge$ 正规，则 $A$ 也正规。
\end{lemma}

\begin{proof}
由于 $A \to A^\wedge$ 忠实平坦，这由 Algebra, Lemma
\ref{algebra-lemma-descent-normal} 得出。
\end{proof}

\begin{lemma}
\label{more-algebra-lemma-flat-completion}
设 $A \to B$ 为 Noether 局部环的局部同态。
则诱导的完备化映射 $A^\wedge \to B^\wedge$ 平坦，
当且仅当 $A \to B$ 平坦。
\end{lemma}

\begin{proof}
考察交换图
$$
\xymatrix{
A^\wedge \ar[r] & B^\wedge \\
A \ar[r] \ar[u] & B \ar[u]
}
$$
竖直箭头均忠实平坦。假设 $A^\wedge \to B^\wedge$ 平坦。
则 $A \to B^\wedge$ 平坦。因此，由 Algebra, Lemma
\ref{algebra-lemma-flatness-descends-more-general}，$B$ 在 $A$ 上平坦。

\medskip\noindent
假设 $A \to B$ 平坦。则 $A \to B^\wedge$ 平坦。
因此，对所有 $n \geq 1$，
$B^\wedge/\mathfrak m_A^n B^\wedge$ 在
$A/\mathfrak m_A^n$ 上平坦。注意，
$\mathfrak m_A^n A^\wedge$ 是 $A^\wedge$ 的极大理想
$\mathfrak m_A^\wedge$ 的第 $n$ 次幂，并且
$A/\mathfrak m_A^n = A^\wedge/(\mathfrak m_A^\wedge)^n$。
于是应用 Algebra, Lemma
\ref{algebra-lemma-flat-module-powers}
（取 $R = A^\wedge$、$I = \mathfrak m_A^\wedge$、
$S = B^\wedge$、$M = S$），可知 $B^\wedge$ 在
$A^\wedge$ 上平坦。
\end{proof}

\begin{lemma}
\label{more-algebra-lemma-flat-unramified}
设 $A \to B$ 为 Noether 局部环的平坦局部同态，并且
$\mathfrak m_A B = \mathfrak m_B$ 且
$\kappa(\mathfrak m_A) = \kappa(\mathfrak m_B)$。
则 $A \to B$ 在完备化上诱导同构 $A^\wedge \to B^\wedge$。
\end{lemma}

\begin{proof}
由 Algebra, Lemma \ref{algebra-lemma-finite-after-completion}，
$B^\wedge$ 是 $B$ 的 $\mathfrak m_A$-adic 完备化，并且
$A^\wedge \to B^\wedge$ 有限。由于 $A \to B$ 平坦，有
$\text{Tor}_1^A(B, \kappa(\mathfrak m_A)) = 0$。
因此，由引理 \ref{more-algebra-lemma-flat-after-completion}，
$B^\wedge$ 在 $A^\wedge$ 上平坦。进而由 Algebra, Lemma
\ref{algebra-lemma-finite-flat-local}，$B^\wedge$ 是自由
$A^\wedge$-模。根据我们的假设（以及 Algebra, Lemma
\ref{algebra-lemma-hathat-finitely-generated}），
$A^\wedge \to B^\wedge$ 诱导同构
$$
\kappa(\mathfrak m_A) = A^\wedge/\mathfrak m_A A^\wedge \to
B^\wedge/\mathfrak m_A B^\wedge = B^\wedge/\mathfrak m_B B^\wedge =
\kappa(\mathfrak m_B)
$$
故 $B^\wedge$ 是秩为 $1$ 的自由模。因此
$A^\wedge \to B^\wedge$ 是同构。
\end{proof}





\section{性质在 \'etale 映射下的保持性}
\label{more-algebra-section-permanence-etale}

\noindent
本节考察一个 \'etale 环映射 $\varphi : A \to B$，
研究 $A$ 的哪些性质由 $B$ 继承，以及 $B$ 在 $\mathfrak q$
处的局部环的哪些性质由 $A$ 在
$\mathfrak p = \varphi^{-1}(\mathfrak q)$ 处的局部环继承。
基本上，本节只是回顾并汇集先前的结果，不增加新内容。

\medskip\noindent
下文将不再说明地使用：\'etale 环映射平坦
（Algebra, Lemma \ref{algebra-lemma-etale}），且局部环的平坦局部同态
忠实平坦（Algebra, Lemma \ref{algebra-lemma-local-flat-ff}）。

\begin{lemma}
\label{more-algebra-lemma-Noetherian-etale-extension}
若 $A \to B$ 是 \'etale 环映射，且 $\mathfrak q$ 是 $B$
中位于 $\mathfrak p \subset A$ 之上的素理想，则
$A_{\mathfrak p}$ Noether，当且仅当 $B_{\mathfrak q}$ Noether。
\end{lemma}

\begin{proof}
由于 $A_\mathfrak p \to B_\mathfrak q$ 忠实平坦，
由 Algebra, Lemma \ref{algebra-lemma-descent-Noetherian} 可见，
$B_\mathfrak q$ Noether 蕴含 $A_\mathfrak p$ Noether。
反之，若 $A_\mathfrak p$ Noether，则 $B_\mathfrak q$ Noether，
因为它是某个有限型 $A_\mathfrak p$-代数的局部化。
\end{proof}

\begin{lemma}
\label{more-algebra-lemma-dimension-etale-extension}
若 $A \to B$ 是 \'etale 环映射，且 $\mathfrak q$ 是 $B$
中位于 $\mathfrak p \subset A$ 之上的素理想，则
$\dim(A_{\mathfrak p}) = \dim(B_{\mathfrak q})$。
\end{lemma}

\begin{proof}
事实上，由于 $A_{\mathfrak p} \to B_{\mathfrak q}$ 平坦，
下降定理成立，因而有不等式
$\dim(A_{\mathfrak p}) \leq \dim(B_{\mathfrak q})$；
见 Algebra, Lemma \ref{algebra-lemma-dimension-going-up}。
另一方面，设
$\mathfrak q_0 \subset \mathfrak q_1 \subset \ldots \subset \mathfrak q_n$
是 $B_{\mathfrak q}$ 中的一条素理想链。则对应的素理想序列
$\mathfrak p_0 \subset \mathfrak p_1 \subset \ldots \subset \mathfrak p_n$
（其中 $\mathfrak p_i = \mathfrak q_i \cap A_{\mathfrak p}$）
也确为一条链，即该序列中没有相等项。这是因为 \'etale 环映射拟有限
（见 Algebra, Lemma \ref{algebra-lemma-etale-quasi-finite}），
而拟有限环映射诱导的谱映射具有离散纤维（依定义）。
这说明 $\dim(A_{\mathfrak p}) \geq \dim(B_{\mathfrak q})$，正如所需。
\end{proof}

\begin{lemma}
\label{more-algebra-lemma-regular-etale-extension}
若 $A \to B$ 是 \'etale 环映射，且 $\mathfrak q$ 是 $B$
中位于 $\mathfrak p \subset A$ 之上的素理想，则
$A_{\mathfrak p}$ 正则，当且仅当 $B_{\mathfrak q}$ 正则。
\end{lemma}

\begin{proof}
由引理 \ref{more-algebra-lemma-Noetherian-etale-extension}，为证明等价性，
可以假设 $A_\mathfrak p$ 与 $B_\mathfrak q$ 均为 Noether 环。
设 $x_1, \ldots, x_t \in \mathfrak pA_\mathfrak p$
是一组极小生成元。由于
$A_\mathfrak p \to B_\mathfrak q$ 忠实平坦，它们在
$B_\mathfrak q$ 中的像 $y_1, \ldots, y_t$ 构成
$\mathfrak pB_\mathfrak q = \mathfrak q B_\mathfrak q$
的一组极小生成元
（Algebra, Lemma \ref{algebra-lemma-etale-at-prime}）。
依定义，$A_\mathfrak p$ 正则意味着
$t = \dim(A_\mathfrak p)$；对于 $B_\mathfrak q$ 亦然。
因此，引理由引理 \ref{more-algebra-lemma-dimension-etale-extension} 中的等式
$\dim(A_\mathfrak p) = \dim(B_\mathfrak q)$ 得出。
\end{proof}

\begin{lemma}
\label{more-algebra-lemma-Dedekind-etale-extension}
若 $A \to B$ 是 \'etale 环映射，且 $A$ 是 Dedekind 整环，
则 $B$ 是 Dedekind 整环的有限直积。特别地，对于
$\mathfrak q \subset B$ 为极大理想，局部化 $B_\mathfrak q$
是离散赋值环。
\end{lemma}

\begin{proof}
关于局部环的断言由引理
\ref{more-algebra-lemma-dimension-etale-extension}、
\ref{more-algebra-lemma-regular-etale-extension} 和 Algebra, Lemma
\ref{algebra-lemma-characterize-dvr} 得出。
由此，$B$ 是维数为 $1$ 的 Noether 正规环。
根据 Algebra, Lemma
\ref{algebra-lemma-characterize-reduced-ring-normal}，
$B$ 是维数为 $1$ 的正规整环的有限直积。
由 Algebra, Lemma \ref{algebra-lemma-characterize-Dedekind}，
这些整环都是 Dedekind 整环。
\end{proof}




\section{性质在 Hensel 化下的保持性}
\label{more-algebra-section-permanence-henselization}

\noindent
给定局部环 $R$，以 $R^h$ 和 $R^{sh}$ 分别表示 $R$ 的 Hensel 化
和严格 Hensel 化；见 Algebra, Definition
\ref{algebra-definition-henselization}。正如本节将证明的，$R$
的许多性质都会反映在 $R^h$ 与 $R^{sh}$ 中。

\begin{lemma}
\label{more-algebra-lemma-dumb-properties-henselization}
设 $(R, \mathfrak m, \kappa)$ 为局部环。则有：
\begin{enumerate}
\item $R \to R^h \to R^{sh}$ 是忠实平坦环映射；
\item $\mathfrak m R^h = \mathfrak m^h$ 且
$\mathfrak m R^{sh} = \mathfrak m^h R^{sh} = \mathfrak m^{sh}$；
\item 对所有 $n$，有 $R/\mathfrak m^n = R^h/\mathfrak m^nR^h$；
\item 存在元素 $x_i \in R^{sh}$，使得
$R^{sh}/\mathfrak m^nR^{sh}$ 是以
$x_i \bmod \mathfrak m^nR^{sh}$ 为基的自由
$R/\mathfrak m^n$-模。
\end{enumerate}
\end{lemma}

\begin{proof}
依构造，$R^h$ 是 \'etale $R$-代数的余极限；见 Algebra, Lemma
\ref{algebra-lemma-henselization}。由于 \'etale 环映射平坦
（Algebra, Lemma \ref{algebra-lemma-etale}），由 Algebra, Lemma
\ref{algebra-lemma-colimit-flat}，$R^h$ 在 $R$ 上平坦。
由于平坦局部环同态忠实平坦
（Algebra, Lemma \ref{algebra-lemma-local-flat-ff}），
$R \to R^h$ 忠实平坦。
环映射 $R^h \to R^{sh}$ 是有限 \'etale 环映射的余极限；
见 Algebra, Lemma \ref{algebra-lemma-strict-henselization} 的证明。
因此，与上述相同的论证表明 $R^h \to R^{sh}$ 忠实平坦。

\medskip\noindent
(2) 由 Algebra, Lemmas \ref{algebra-lemma-henselization} 和
\ref{algebra-lemma-strict-henselization} 得出。(3) 由
Algebra, Lemma \ref{algebra-lemma-local-artinian-basis-when-flat} 得出，
因为 $R/\mathfrak m \to R^h/\mathfrak mR^h$ 是同构，并且
$R/\mathfrak m^n \to R^h/\mathfrak m^nR^h$
是平坦环映射 $R \to R^h$ 的基变换，因而平坦
（Algebra, Lemma \ref{algebra-lemma-flat-base-change}）。
设 $\kappa^{sep}$ 为 $R^{sh}$ 的剩余域
（它是 $\kappa$ 的一个可分代数闭包）。选取 $x_i \in R^{sh}$，
使其像构成 $\kappa^{sep}$ 作为 $\kappa$-向量空间的一组基。
于是，完全如上地应用 Algebra, Lemma
\ref{algebra-lemma-local-artinian-basis-when-flat} 即得 (4)。
\end{proof}

\begin{lemma}
\label{more-algebra-lemma-henselization-formally-smooth}
设 $(R, \mathfrak m, \kappa)$ 为局部环。则
\begin{enumerate}
\item $R \to R^h$、$R^h \to R^{sh}$ 与 $R \to R^{sh}$
均为形式 \'etale 映射；
\item $R \to R^h$、$R^h \to R^{sh}$ 与 $R \to R^{sh}$
分别对 $\mathfrak m^h$、$\mathfrak m^{sh}$ 与
$\mathfrak m^{sh}$-拓扑形式光滑。
\end{enumerate}
\end{lemma}

\begin{proof}
(1) 由以下事实得出：依构造，$R^h$ 与 $R^{sh}$ 是 \'etale
代数的有向余极限；\'etale 代数是形式 \'etale 的
（Algebra, Lemma \ref{algebra-lemma-formally-etale-etale}）；
形式 \'etale 代数的余极限仍形式 \'etale
（Algebra, Lemma \ref{algebra-lemma-colimit-formally-etale}）。
(2) 由形式 \'etale 环映射形式光滑这一事实以及引理
\ref{more-algebra-lemma-formally-smooth} 得出。
\end{proof}

\begin{lemma}
\label{more-algebra-lemma-henselization-noetherian}
\begin{reference}
\cite[IV, Theorem 18.6.6 and Proposition 18.8.8]{EGA}
\end{reference}
设 $R$ 为局部环。下列条件等价：
\begin{enumerate}
\item $R$ Noether；
\item $R^h$ Noether；以及
\item $R^{sh}$ Noether。
\end{enumerate}
在这种情况下：
\begin{enumerate}
\item[(a)] $(R^h)^\wedge$ 与 $(R^{sh})^\wedge$ 是 Noether 完备局部环；
\item[(b)] $R^\wedge \to (R^h)^\wedge$ 是同构；
\item[(c)] $R^h \to (R^h)^\wedge$ 与
$R^{sh} \to (R^{sh})^\wedge$ 平坦；
\item[(d)] $R^\wedge \to (R^{sh})^\wedge$ 对
$\mathfrak m_{(R^{sh})^\wedge}$-adic 拓扑形式光滑；
\item[(e)] $(R^\wedge)^{sh} = R^\wedge \otimes_{R^h} R^{sh}$；以及
\item[(f)] $((R^\wedge)^{sh})^\wedge = (R^{sh})^\wedge$。
\end{enumerate}
\end{lemma}

\begin{proof}
由于 $R \to R^h \to R^{sh}$ 忠实平坦
（引理 \ref{more-algebra-lemma-dumb-properties-henselization}），由
Algebra, Lemma \ref{algebra-lemma-descent-Noetherian} 可见，
$R^h$ 或 $R^{sh}$ Noether 蕴含 $R$ Noether。
在余下的证明中，假设 $R$ Noether。

\medskip\noindent
由于 $\mathfrak m \subset R$ 有限生成，由引理
\ref{more-algebra-lemma-dumb-properties-henselization}，可知
$\mathfrak m^h = \mathfrak m R^h$ 与
$\mathfrak m^{sh} = \mathfrak mR^{sh}$ 有限生成。
因此，由 Algebra, Lemma
\ref{algebra-lemma-complete-local-ring-Noetherian}，
$(R^h)^\wedge$ 与 $(R^{sh})^\wedge$ 是 Noether 环。这证明了 (a)。

\medskip\noindent
注意，(b) 立即由引理
\ref{more-algebra-lemma-dumb-properties-henselization} 得出。
特别地，由 Algebra, Lemma
\ref{algebra-lemma-completion-faithfully-flat} 可见，
$(R^h)^\wedge$ 在 $R$ 上平坦。

\medskip\noindent
接着证明 $R^h \to (R^h)^\wedge$ 平坦。将
$R^h = \colim_i R_i$ 写成 \'etale $R$-代数的局部化的有向余极限。
根据 Algebra, Lemma \ref{algebra-lemma-colimit-rings-flat}，
若 $(R^h)^\wedge$ 在每个 $R_i$ 上平坦，则
$R^h \to (R^h)^\wedge$ 平坦。注意，依构造有
$R^h = R_i^h$。故由 (b)，
$R_i^\wedge = (R^h)^\wedge$ 在 $R_i$ 上平坦，正如所需。
为了完成 (c) 的证明，我们证明
$R^{sh} \to (R^{sh})^\wedge$ 平坦。为此，使用与上述相同的极限论证，
只需证明 $(R^{sh})^\wedge$ 在 $R$ 上平坦。由引理
\ref{more-algebra-lemma-dumb-properties-henselization} 可知，
$(R^{sh})^\wedge$ 是一个自由 $R$-模的完备化。
由引理 \ref{more-algebra-lemma-completed-direct-sum-flat}，它在 $R$ 上平坦，
正如所需。这完成了 (c) 的证明。

\medskip\noindent
至此已知 (c) 成立，且 $(R^h)^\wedge$ 与
$(R^{sh})^\wedge$ Noether。由 Algebra, Lemma
\ref{algebra-lemma-descent-Noetherian}，$R^h$ 与 $R^{sh}$ Noether。

\medskip\noindent
(d) 由引理 \ref{more-algebra-lemma-henselization-formally-smooth} 和
\ref{more-algebra-lemma-formally-smooth-completion} 得出。

\medskip\noindent
(e) 由 Algebra, Lemma \ref{algebra-lemma-sh-from-h-map} 以及
$R^\wedge$ 是 Hensel 环这一事实得出；后者见 Algebra, Lemma
\ref{algebra-lemma-complete-henselian}。

\medskip\noindent
证明 (f)。利用 (e)，存在映射
$R^{sh} \to (R^\wedge)^{sh}$；完备化后，它诱导映射
$(R^{sh})^\wedge \to ((R^\wedge)^{sh})^\wedge$。
再利用 (e)，存在映射 $R^\wedge \to (R^{sh})^\wedge$。
由于 $(R^{sh})^\wedge$ 是严格 Hensel 环（见上文），由
Algebra, Lemma \ref{algebra-lemma-strictly-henselian-functorial}，
此映射诱导映射 $(R^\wedge)^{sh} \to (R^{sh})^\wedge$。
完备化后得到映射
$((R^\wedge)^{sh})^\wedge \to (R^{sh})^\wedge$。
略去验证这两个映射互为逆映射的过程。
\end{proof}

\begin{lemma}
\label{more-algebra-lemma-henselization-reduced}
\begin{slogan}
既约性传递到（严格）Hensel 化。
\end{slogan}
设 $R$ 为局部环。下列条件等价：$R$ 既约，$R$ 的 Hensel 化
$R^h$ 既约，以及 $R$ 的严格 Hensel 化 $R^{sh}$ 既约。
\end{lemma}

\begin{proof}
环映射 $R \to R^h \to R^{sh}$ 忠实平坦。因此，由
Algebra, Lemma \ref{algebra-lemma-descent-reduced}
得到一个方向的推含。反之，假设 $R$ 既约。由于 $R^h$ 与
$R^{sh}$ 是 \'etale、因而光滑的 $R$-代数的滤过余极限，
结论由 Algebra, Lemma \ref{algebra-lemma-reduced-goes-up} 得出。
\end{proof}

\begin{lemma}
\label{more-algebra-lemma-henselization-nil}
设 $R$ 为局部环。以 $nil(R)$ 表示 $R$ 的幂零元素构成的理想。
则 $nil(R)R^h = nil(R^h)$ 且
$nil(R)R^{sh} = nil(R^{sh})$。
\end{lemma}

\begin{proof}
注意，$nil(R)$ 是由幂零元素组成且使商 $R/nil(R)$ 既约的最大理想。
由 Algebra, Lemma \ref{algebra-lemma-locally-nilpotent}，
$nil(R)R^h$ 由幂零元素组成。另注意，由 Algebra, Lemma
\ref{algebra-lemma-quotient-henselization}，
$R^h/nil(R) R^h$ 是 $R/nil(R)$ 的 Hensel 化。
因此，由引理 \ref{more-algebra-lemma-henselization-reduced}，
$R^h/nil(R)R^h$ 既约。故 $nil(R) R^h = nil(R^h)$，正如所需。
严格 Hensel 化的情形类似，但使用 Algebra, Lemma
\ref{algebra-lemma-quotient-strict-henselization}。
\end{proof}

\begin{lemma}
\label{more-algebra-lemma-henselization-normal}
设 $R$ 为局部环。下列条件等价：$R$ 是正规整环，$R$ 的 Hensel 化
$R^h$ 是正规整环，以及 $R$ 的严格 Hensel 化 $R^{sh}$ 是正规整环。
\end{lemma}

\begin{proof}
先说明，局部环正规当且仅当它是正规整环
（见 Algebra, Definition \ref{algebra-definition-ring-normal}）。
环映射 $R \to R^h \to R^{sh}$ 忠实平坦。因此，由
Algebra, Lemma \ref{algebra-lemma-descent-normal}
得到一个方向的推含。反之，假设 $R$ 正规。由于 $R^h$ 与
$R^{sh}$ 是 \'etale、因而光滑的 $R$-代数的滤过余极限，
结论由 Algebra, Lemmas \ref{algebra-lemma-normal-goes-up} 和
\ref{algebra-lemma-colimit-normal-ring} 得出。
\end{proof}

\begin{lemma}
\label{more-algebra-lemma-henselization-dimension}
对任意局部环 $R$，有
$\dim(R) = \dim(R^h) = \dim(R^{sh})$。
\end{lemma}

\begin{proof}
由于 $R \to R^{sh}$ 忠实平坦
（引理 \ref{more-algebra-lemma-dumb-properties-henselization}），由下降定理可知
$\dim(R^{sh}) \geq \dim(R)$；见 Algebra, Lemma
\ref{algebra-lemma-dimension-going-up}。
为证反向不等式，将 $R^{sh} = \colim R_i$ 写成局部环的有向余极限，
其中每个 $R_i$ 都是某个 \'etale $R$-代数的局部化。
若
$\mathfrak q_0 \subset \mathfrak q_1 \subset \ldots \subset \mathfrak q_n$
是 $R^{sh}$ 中的一条素理想链，则当 $i$ 充分大时，序列
$$
R_i \cap \mathfrak q_0 \subset
R_i \cap \mathfrak q_1 \subset \ldots \subset
R_i \cap \mathfrak q_n
$$
是 $R_i$ 中的一条素理想链。因此
$\dim(R^{sh}) \leq \sup_i \dim(R_i)$。
但由引理 \ref{more-algebra-lemma-dimension-etale-extension}，对每个 $i$
都有 $\dim(R_i) = \dim(R)$，结论得证。
\end{proof}

\begin{lemma}
\label{more-algebra-lemma-henselization-depth}
对任意 Noether 局部环 $R$，有
$\text{depth}(R) = \text{depth}(R^h) = \text{depth}(R^{sh})$。
\end{lemma}

\begin{proof}
由引理 \ref{more-algebra-lemma-henselization-noetherian}，已知
$R^h$ 与 $R^{sh}$ Noether。因此，引理由 Algebra, Lemma
\ref{algebra-lemma-apply-grothendieck} 得出。
\end{proof}

\begin{lemma}
\label{more-algebra-lemma-henselization-CM}
设 $R$ 为 Noether 局部环。下列条件等价：$R$ 是 Cohen--Macaulay 环，
$R$ 的 Hensel 化 $R^h$ 是 Cohen--Macaulay 环，以及 $R$
的严格 Hensel 化 $R^{sh}$ 是 Cohen--Macaulay 环。
\end{lemma}

\begin{proof}
由引理 \ref{more-algebra-lemma-henselization-noetherian}，已知
$R^h$ 与 $R^{sh}$ Noether，故该引理的陈述有意义。
再由引理 \ref{more-algebra-lemma-henselization-depth} 和
\ref{more-algebra-lemma-henselization-dimension}，有
$\text{depth}(R) = \text{depth}(R^h) = \text{depth}(R^{sh})$
以及
$\dim(R) = \dim(R^h) = \dim(R^{sh})$，结论随即得出。
\end{proof}

\begin{lemma}
\label{more-algebra-lemma-henselization-regular}
设 $R$ 为 Noether 局部环。下列条件等价：$R$ 是正则局部环，
$R$ 的 Hensel 化 $R^h$ 是正则局部环，以及 $R$
的严格 Hensel 化 $R^{sh}$ 是正则局部环。
\end{lemma}

\begin{proof}
由引理 \ref{more-algebra-lemma-henselization-noetherian}，已知
$R^h$ 与 $R^{sh}$ Noether，故该引理的陈述有意义。
设 $\mathfrak m$ 为 $R$ 的极大理想。
设 $x_1, \ldots, x_t \in \mathfrak m$ 是 $\mathfrak m$
的一组极小生成元，即它们在 $\mathfrak m/\mathfrak m^2$
中的像构成该空间在 $\kappa = R/\mathfrak m$ 上的一组基。
由于 $R \to R^h$ 与 $R \to R^{sh}$ 忠实平坦，故它们在
$R^h$ 中的像 $x_1^h, \ldots, x_t^h$ 以及在 $R^{sh}$
中的像 $x_1^{sh}, \ldots, x_t^{sh}$，分别构成
$\mathfrak m^h = \mathfrak mR^h$ 与
$\mathfrak m^{sh} = \mathfrak mR^{sh}$ 的一组极小生成元。
依定义，$R$ 正则意味着 $t = \dim(R)$；对于 $R^h$ 和 $R^{sh}$
亦然。因此，引理由引理 \ref{more-algebra-lemma-henselization-dimension} 中的维数等式
$\dim(R) = \dim(R^h) = \dim(R^{sh})$ 得出
% P02-E0174
\end{proof}

\begin{lemma}
\label{more-algebra-lemma-henselization-dvr}
设 $R$ 为 Noether 局部环。则 $R$ 是离散赋值环，当且仅当
$R^h$ 是离散赋值环，当且仅当 $R^{sh}$ 是离散赋值环。
\end{lemma}

\begin{proof}
这由引理 \ref{more-algebra-lemma-henselization-dimension}、
\ref{more-algebra-lemma-henselization-regular} 和 Algebra, Lemma
\ref{algebra-lemma-characterize-dvr} 得出。
\end{proof}

\begin{lemma}
\label{more-algebra-lemma-filtered-colimit-etale-noetherian-fibres}
设 $A$ 为环。设 $B$ 为 \'etale $A$-代数的滤过余极限。
设 $\mathfrak p$ 为 $A$ 的素理想。若 $B$ Noether，则位于
$\mathfrak p$ 之上的素理想只有有限多个
$\mathfrak q_1, \ldots, \mathfrak q_r$，并且
$B \otimes_A \kappa(\mathfrak p) = \prod \kappa(\mathfrak q_i)$；
各域扩张 $\kappa(\mathfrak q_i)/\kappa(\mathfrak p)$ 都是可分代数扩张。
\end{lemma}

\begin{proof}
将 $B$ 写成滤过余极限 $B = \colim B_i$，其中
$A \to B_i$ 为 \'etale 映射。一方面，
$B \otimes_A \kappa(\mathfrak p) = \colim B_i \otimes_A \kappa(\mathfrak p)$
是 \'etale $\kappa(\mathfrak p)$-代数的滤过余极限；另一方面，
它是 Noether 环。一个 \'etale $\kappa(\mathfrak p)$-代数是有限可分
域扩张的有限直积
（Algebra, Lemma \ref{algebra-lemma-etale-over-field}）。
因此，在诸代数 $B_i \otimes_A \kappa(\mathfrak p)$ 的素理想之间
没有非平凡的特殊化关系（这些素理想全都是极大且极小素理想），
从而在 $B \otimes_A \kappa(\mathfrak p)$ 的素理想之间也没有
非平凡的特殊化关系。因此，$B \otimes_A \kappa(\mathfrak p)$
既约，且只有有限多个素理想，并且它们全都是极小素理想。
% P02-E0172
所以，它是域的有限直积（使用 Algebra, Lemma
\ref{algebra-lemma-total-ring-fractions-no-embedded-points}
或 Algebra, Proposition
\ref{algebra-proposition-dimension-zero-ring}）。
这些域中的每一个都是有限可分扩张的余极限，故引理的最后一个断言成立。
\end{proof}

\begin{lemma}
\label{more-algebra-lemma-fibres-henselization}
设 $R$ 为 Noether 局部环。设 $\mathfrak p \subset R$ 为素理想。
则
$$
R^h \otimes_R \kappa(\mathfrak p) =
\prod\nolimits_{i = 1, \ldots, t} \kappa(\mathfrak q_i)
\quad\text{分别}\quad
R^{sh} \otimes_R \kappa(\mathfrak p) =
\prod\nolimits_{i = 1, \ldots, s} \kappa(\mathfrak r_i)
$$
其中 $\mathfrak q_1, \ldots, \mathfrak q_t$ 和
$\mathfrak r_1, \ldots, \mathfrak r_s$ 分别是 $R^h$ 与 $R^{sh}$
中位于 $\mathfrak p$ 之上的素理想。
% P02-E0173
此外，域扩张
$\kappa(\mathfrak q_i)/\kappa(\mathfrak p)$ 和
$\kappa(\mathfrak r_i)/\kappa(\mathfrak p)$ 分别是可分代数扩张。
\end{lemma}

\begin{proof}
这可以由更一般的引理
\ref{more-algebra-lemma-filtered-colimit-etale-noetherian-fibres} 推出，
其中使用 Hensel 化与严格 Hensel 化是 Noether 环
（如上所见）这一事实。不过，我们也给出如下直接证明。

\medskip\noindent
下文将不再说明地使用引理
\ref{more-algebra-lemma-dumb-properties-henselization} 和
\ref{more-algebra-lemma-henselization-noetherian} 的结果。
注意，$R^h/\mathfrak pR^h$ 与 $R^{sh}/\mathfrak pR^{sh}$
分别是 $R/\mathfrak p$ 的 Hensel 化与严格 Hensel 化；
分别见 Algebra, Lemma \ref{algebra-lemma-quotient-henselization}
与 Algebra, Lemma
\ref{algebra-lemma-quotient-strict-henselization}。
因此，可以用 $R/\mathfrak p$ 替换 $R$，并假设 $R$
是 Noether 局部整环且 $\mathfrak p = (0)$。
由于 $R^h$ 与 $R^{sh}$ Noether，它们分别只有有限多个极小素理想
$\mathfrak q_1, \ldots, \mathfrak q_t$ 和
$\mathfrak r_1, \ldots, \mathfrak r_s$。
由于 $R \to R^h$ 与 $R \to R^{sh}$ 平坦，由下降定理，
这些恰是位于 $\mathfrak p = (0)$ 之上的素理想。
最后，由于 $R$ 是整环，由引理
\ref{more-algebra-lemma-henselization-reduced}，$R^h$ 与 $R^{sh}$ 既约。
因此，$R^h \otimes_R \kappa(\mathfrak p)$ 与
$R^{sh} \otimes_R \kappa(\mathfrak p)$ 是只有有限多个素理想的
既约 Noether 环，并且所有这些素理想都是极小素理想
（因而也是极大理想）。于是这些环是 Artin 环，并且是它们在极大理想处
的局部化的直积；每个这样的局部化必为域
（见 Algebra, Proposition
\ref{algebra-proposition-dimension-zero-ring} 和 Algebra, Lemma
\ref{algebra-lemma-minimal-prime-reduced-ring}）。

\medskip\noindent
最后一个断言由如下事实得出：$R \to R^h$ 与
$R \to R^{sh}$ 是 \'etale 环映射的余极限，故诱导的剩余域扩张
是有限可分扩张的余极限；见 Algebra, Lemma
\ref{algebra-lemma-etale-at-prime}。
\end{proof}









\section{再论域扩张}
\label{more-algebra-section-p-bases}

\noindent
本节研究特征为 $p > 0$ 的域扩张的一些特殊现象。

\begin{definition}
\label{more-algebra-definition-p-basis}
设 $p$ 为素数。设 $k \to K$ 为特征 $p$ 的域的扩张。
以 $kK^p$ 表示 $k$ 与 $K^p$ 在 $K$ 中的合成域。
\begin{enumerate}
\item 子集 $\{x_i\} \subset K$ 称为{\it 在 $k$ 上 p-无关}，
如果满足 $0 \leq e_i < p$ 的元素
$x^E = \prod x_i^{e_i}$ 在 $kK^p$ 上线性无关。
\item $K$ 的子集 $\{x_i\}$ 称为{\it $K$ 在 $k$ 上的 p-基}，
如果诸元素 $x^E$ 构成 $K$ 在 $kK^p$ 上的一组基。
\end{enumerate}
\end{definition}

\noindent
这与 $\mathbf{F}_p$-代数的 $p$-基概念有关；我们将在后文讨论它
（此处插入未来的引用）。
% P02-E0175

\begin{lemma}
\label{more-algebra-lemma-p-basis}
设 $K/k$ 为域扩张。假设 $k$ 的特征为 $p > 0$。
设 $\{x_i\}$ 为 $K$ 的一个子集。下列条件等价：
\begin{enumerate}
\item 诸元素 $\{x_i\}$ 在 $k$ 上 $p$-无关；以及
\item 诸元素 $\text{d}x_i$ 在 $\Omega_{K/k}$ 中 $K$-线性无关。
\end{enumerate}
任意 $p$-无关族都可以扩充成 $K$ 在 $k$ 上的一组 $p$-基。
特别地，域 $K$ 在 $k$ 上有一组 $p$-基。
此外，下列条件等价：
\begin{enumerate}
\item[(a)] $\{x_i\}$ 是 $K$ 在 $k$ 上的一组 $p$-基；以及
\item[(b)] $\text{d}x_i$ 是 $K$-向量空间 $\Omega_{K/k}$ 的一组基。
\end{enumerate}
\end{lemma}

\begin{proof}
假设 (2)，并设 $\sum a_E x^E = 0$ 是一个线性关系，其中
$a_E \in k K^p$。设 $\theta_i : K \to K$ 为满足
$\theta_i(x_j) = \delta_{ij}$（Kronecker delta）的 $k$-导子。
注意，$K$ 的任意 $k$-导子都将 $kK^p$ 零化。
对给定关系应用 $\theta_i$，得到新的关系
$$
\sum\nolimits_{E, e_i > 0}
e_i a_E x_1^{e_1}\ldots x_i^{e_i - 1} \ldots x_n^{e_n} = 0
$$
因此，如果在某个 $a_E \not = 0$ 的关系中，选取总次数
$|E| = \sum e_i$ 最小的 $\sum a_E x^E$，就会得到矛盾。
故 (1) 成立。

\medskip\noindent
若 $\{x_i\}$ 是 $K$ 在 $k$ 上的一组 $p$-基，则
$K \cong kK^p[X_i]/(X_i^p - x_i^p)$。
由此可见，$\text{d}x_i$ 构成 $\Omega_{K/k}$ 在 $K$ 上的一组基。
故 (a) 推出 (b)。

\medskip\noindent
设 $\{x_i\}$ 是 $K$ 在 $k$ 上的一个 $p$-无关子集。
应用 Zorn 引理，可将其扩充成 $K$ 在 $k$ 上的一个极大 $p$-无关子集。
我们断言，$K$ 的任意极大 $p$-无关子集 $\{x_i\}$
都是 $K$ 在 $k$ 上的一组 $p$-基。该断言将说明 (1) 推出 (2)，
并确立 $p$-基的存在性。为证明该断言，设 $L$ 为
$K$ 中由 $kK^p$ 与诸 $x_i$ 生成的子域。必须证明 $L = K$。
若 $x \in K$ 但 $x \not \in L$，则 $x^p \in L$，且
% P02-E0176
$L(x) \cong L[z]/(z^p - x)$。因此，
$\{x_i\} \cup \{x\}$ 在 $k$ 上 $p$-无关，矛盾。

\medskip\noindent
最后，还须证明 (b) 推出 (a)。由 (1) 与 (2) 的等价性，
$\{x_i\}$ 是 $K$ 在 $k$ 上的极大 $p$-无关子集。
故由上述断言，它是一组 $p$-基。
\end{proof}

\begin{lemma}
\label{more-algebra-lemma-intersection-subfields-subspace}
设 $K/k$ 为域扩张。设 $\{K_\alpha\}_{\alpha \in A}$
为 $K$ 的一族子域，并具有以下性质：
\begin{enumerate}
\item 对所有 $\alpha \in A$，有 $k \subset K_\alpha$；
\item $k = \bigcap_{\alpha \in A} K_\alpha$；
\item 对任意 $\alpha, \alpha' \in A$，存在 $\alpha'' \in A$，
使得 $K_{\alpha''} \subset K_\alpha \cap K_{\alpha'}$。
\end{enumerate}
则对于 $n \geq 1$ 以及作为 $K$-向量空间的
$V \subset K^{\oplus n}$，有
$V \cap k^{\oplus n} \not = 0$，当且仅当对所有 $\alpha \in A$
都有 $V \cap K_\alpha^{\oplus n} \not = 0$。
\end{lemma}

\begin{proof}
对 $n$ 作归纳。$n = 1$ 的情形由假设得出。
假设已对 $K^{\oplus n - 1}$ 的子空间证明该结果。
设 $V \subset K^{\oplus n}$ 对所有 $\alpha \in A$
都与 $K_\alpha^{\oplus n}$ 的交非零。若
$V \cap 0 \oplus k^{\oplus n - 1}$ 非零，则结论成立。
故可假设它为零。由归纳假设，可以找到一个 $\alpha$，使得
$V \cap 0 \oplus K_\alpha^{\oplus n - 1}$ 为零。
取一个非零元素
$v = (x_1, \ldots, x_n) \in V \cap K_\alpha^{\oplus n}$。
由 $\alpha$ 的选取，$x_1$ 非零。以 $x_1^{-1}v$ 替换 $v$，
使得 $v = (1, x_2, \ldots, x_n)$。注意，若
$v' = (x_1', \ldots, x'_n) \in V \cap K_\alpha^{\oplus n}$，
则由 $\alpha$ 的选取有 $v' - x_1'v = 0$。
因此 $V \cap K_\alpha^{\oplus n} = K_\alpha v$。
若选取某个满足 $K_{\alpha'} \subset K_\alpha$ 的 $\alpha'$，
则必有 $v \in V \cap K_{\alpha'}^{\oplus n}$
（对 $\alpha'$ 应用同样的论证）。因此
$$
x_2, \ldots, x_n \in
\bigcap\nolimits_{\alpha' \in A, K_{\alpha'} \subset K_\alpha} K_{\alpha'}
$$
而由 (2) 与 (3)，右端等于 $k$。
\end{proof}

\begin{lemma}
\label{more-algebra-lemma-intersection-subfields}
设 $K$ 为特征 $p$ 的域。设 $\{K_\alpha\}_{\alpha \in A}$
为 $K$ 的一族子域，并具有以下性质：
\begin{enumerate}
\item 对所有 $\alpha \in A$，有 $K^p \subset K_\alpha$；
\item $K^p = \bigcap_{\alpha \in A} K_\alpha$；
\item 对任意 $\alpha, \alpha' \in A$，存在 $\alpha'' \in A$，
使得 $K_{\alpha''} \subset K_\alpha \cap K_{\alpha'}$。
\end{enumerate}
则
\begin{enumerate}
\item 映射
$\Omega_{K/\mathbf{F}_p} \to \Omega_{K/K_\alpha}$
的诸核之交为零；
\item 对任意有限扩张 $L/K$，有
$L^p = \bigcap_{\alpha \in A} L^pK_\alpha$。
\end{enumerate}
\end{lemma}

\begin{proof}
证明 (1)。选取 $K$ 在 $\mathbf{F}_p$ 上的一组 $p$-基
$\{x_i\}$。设
$\eta = \sum_{i \in I'} y_i \text{d}x_i$ 在每个 $\alpha \in A$
下都映为 $\Omega_{K/K_\alpha}$ 中的零，其中指标集 $I'$ 有限。
由引理 \ref{more-algebra-lemma-p-basis}，这意味着对每个 $\alpha$ 都存在关系
$$
\sum\nolimits_E a_{E, \alpha} x^E = 0,\quad a_{E, \alpha} \in K_\alpha
$$
其中 $E$ 遍历满足 $0 \leq e_i < p$ 的多重指标
$E = (e_i)_{i \in I'}$。另一方面，引理 \ref{more-algebra-lemma-p-basis}
保证不存在系数 $a_E \in K^p$ 的这种关系
$\sum a_E x^E = 0$。由引理
\ref{more-algebra-lemma-intersection-subfields-subspace}，这给出矛盾。

\medskip\noindent
证明 (2)。设有有限域扩张塔 $L/M/K$。
令 $M_\alpha = M^p K_\alpha$ 且
$L_\alpha = L^p K_\alpha = L^p M_\alpha$。
可以先证明 $M^p = \bigcap_{\alpha \in A} M_\alpha$，
再证明 $L^p = \bigcap_{\alpha \in A} L_\alpha$。
因此，只需对没有非平凡子域的单扩张证明 (2)。
首先，假设 $L = K(\theta)$ 在 $K$ 上可分。则 $L$
在 $K$ 上由 $\theta^p$ 生成，故可假设 $\theta \in L^p$。
在这种情况下，
$$
L^p = K^p \oplus K^p\theta \oplus \ldots K^p\theta^{d - 1}
\quad\text{且}\quad
L^pK_\alpha =
K_\alpha \oplus K_\alpha \theta \oplus \ldots K_\alpha\theta^{d - 1}
$$
其中 $d = [L : K]$。在这种情况下，结论显然。
另一种情形是 $L = K(\theta)$，其中
$\theta^p = t \in K$ 且 $t \not \in K^p$。此时
$$
L^p = K^p \oplus K^pt \oplus \ldots K^pt^{p - 1}
\quad\text{且}\quad
L^pK_\alpha =
K_\alpha \oplus K_\alpha t \oplus \ldots K_\alpha t^{p - 1}
$$
结论同样显然。
\end{proof}

\begin{lemma}
\label{more-algebra-lemma-power-series-ring-subfields}
设 $k$ 为特征 $p > 0$ 的域。设 $\{x_i\}_{i \in I}$
为 $k$ 的一组 $p$-基。设 $n, m \geq 0$。
设 $K$ 为
$A = k[[x_1, \ldots, x_n]][y_1, \ldots, y_m]$ 的分式域。
设 $J$ 为 $I$ 的有限子集。考察由 $k^p$ 以及满足
$i \in I \setminus J$ 的 $x_i$ 生成的子域
% P02-E0177
$k/k_J/k^p$。下列环
$$
A_J = k_J[[x_1^p, \ldots, x_n^p]][y_1^p, \ldots, y_m^p]
$$
的分式域 $K_J$ 构成 $K$ 的一族子域，并满足引理
\ref{more-algebra-lemma-intersection-subfields} 的条件。此外，每个环扩张
$A_J \subset A$ 都有限。
\end{lemma}

\begin{proof}
由于 $k/k_J$ 有限，由 Algebra, Lemma
\ref{algebra-lemma-finite-after-completion}，环扩张
% P02-E0178
$k_J[[x_1^p, \ldots, x_d^p]] \subset k[[x_1, \ldots, x_d]]$
有限。这蕴含 $A_J \to A$ 有限。

\medskip\noindent
下面验证引理 \ref{more-algebra-lemma-intersection-subfields} 的性质 (1)、(2)、(3)。
证明 (1)。对 $a \in A$，有 $a^p \in A_J$。故
$K^p \subset K_J$。证明 (2)。设
$f/g^p \in K$，其中 $f, g \in A$ 且 $g \not = 0$；
并假设它对 $J$ 的每个选择都属于 $K_J$。
暂时固定 $J$。由于 $f/g^p \in K_J$，可将其写成
$f/g^p = a/b^p$，其中 $a \in A_J$ 且 $b \in A$ 非零。
故 $b^p f \in A_J$。对任意 $A_J$-导子 $D : A \to A$，
有 $0 = D(b^pf) = b^p D(f)$；由于 $A$ 是整环，故 $D(f) = 0$。
取 $D = \partial_{x_i}$ 与 $D = \partial_{y_j}$，得到
% P02-E0179
$f \in k[[x_1^p, \ldots, x_n^p]][y_1^p, \ldots, y_d^p]$。
应用 $k_J$-导子 $\theta : k \to k$，同样得到 $f$
的所有系数都属于 $k_J$，即 $f \in A_J$。
由于显然
% P02-E0180
$A^p = \bigcap\nolimits_J A_J$，其中 $J$ 遍历引理中所述的所有子域，
故 $f \in A^p$，正如所需。
证明 (3)。这是显然的，因为
$K_{J \cup J'} \subset K_J \cap K_{J'}$。
\end{proof}





\section{奇异轨迹}
\label{more-algebra-section-singular-locus}

\noindent
设 $R$ 为 Noether 环。$X = \Spec(R)$ 的{\it 正则轨迹}
$\text{Reg}(X)$ 是由满足 $R_\mathfrak p$ 为正则局部环的素理想
$\mathfrak p$ 组成的集合。$X = \Spec(R)$ 的{\it 奇异轨迹}
$\text{Sing}(X)$ 是补集
$X \setminus \text{Reg}(X)$，即由满足 $R_\mathfrak p$
不是正则局部环的素理想 $\mathfrak p$ 组成的集合。
由 Algebra, Definition \ref{algebra-definition-regular}
之前的讨论可见，$\text{Reg}(X)$ 在一般化下稳定。
本节研究保证 $\text{Reg}(X)$ 为开集的条件。

\begin{definition}
\label{more-algebra-definition-J}
\begin{reference}
\cite[(32.B)]{MatCA}
\end{reference}
设 $R$ 为 Noether 环。设 $X = \Spec(R)$。
\begin{enumerate}
\item 若 $\text{Reg}(X)$ 包含一个非空开集，则称 $R$ 为 {\it J-0}；
\item 若 $\text{Reg}(X)$ 为开集，则称 $R$ 为 {\it J-1}；
\item 若任意有限型 $R$-代数都是 J-1 的，则称 $R$ 为 {\it J-2}。
\end{enumerate}
\end{definition}

\noindent
环 $\mathbf{Q}[x]/(x^2)$ 不满足 J-0，但满足 J-1。
另一方面，对于 Noether 整环以及更一般的非零既约 Noether 环，
J-1 蕴含 J-0，因为这种环在极小素理想处正则。
下面给出 J-1 性质的一种刻画。

\begin{lemma}
\label{more-algebra-lemma-J-1}
设 $R$ 为 Noether 环。设 $X = \Spec(R)$。
则 $R$ 为 J-1，当且仅当对所有
$\mathfrak p \in \text{Reg}(X)$，
$V(\mathfrak p) \cap \text{Reg}(X)$ 都包含
$V(\mathfrak p)$ 的一个非空开子集。
\end{lemma}

\begin{proof}
这由 Topology, Lemma
\ref{topology-lemma-characterize-open-Noetherian} 以及
$\text{Reg}(X)$ 在一般化下稳定这一事实得出；后者见
Algebra, Lemma
\ref{algebra-lemma-localization-of-regular-local-is-regular}。
\end{proof}

\begin{lemma}
\label{more-algebra-lemma-intersection-regular-with-closed}
设 $R$ 为 Noether 环。设 $X = \Spec(R)$。
假设对所有素理想 $\mathfrak p \subset R$，环
$R/\mathfrak p$ 都是 J-0 的。则 $R$ 是 J-1 的。
\end{lemma}

\begin{proof}
我们将证明引理 \ref{more-algebra-lemma-J-1} 的判据适用。
设 $\mathfrak p \in \text{Reg}(X)$ 为高度 $r$ 的素理想。
选取 $f_1, \ldots, f_r \in \mathfrak p$，使其像生成
$\mathfrak pR_\mathfrak p$。由于
$\mathfrak p \in \text{Reg}(X)$，可知
% P02-E0181
$f_1, \ldots, f_r$ 的像在 $R_\mathfrak p$ 中构成正则序列；
见 Algebra, Lemma \ref{algebra-lemma-regular-ring-CM}。
因此，由 Algebra, Lemma
\ref{algebra-lemma-regular-sequence-in-neighbourhood} 可见，
用某个满足 $g \not \in \mathfrak p$ 的 $g \in R$ 对 $R$
作 $R_g$ 替换后，序列 $f_1, \ldots, f_r$ 在 $R$ 中正则。
再作一次替换，还可假设 $f_1, \ldots, f_r$ 生成 $\mathfrak p$。
接着，设 $\mathfrak p \subset \mathfrak q$ 为素理想，并且
$(R/\mathfrak p)_\mathfrak q$ 是正则局部环。
由引理的假设，$V(\mathfrak p)$ 中存在一个由这类素理想组成的非空开子集，
故只需证明 $R_\mathfrak q$ 正则。
注意，$f_1, \ldots, f_r$ 在 $R_\mathfrak q$ 中构成正则序列，
且 $R_\mathfrak q/(f_1, \ldots, f_r)R_\mathfrak q$ 正则。
因此，由 Algebra, Lemma
\ref{algebra-lemma-regular-mod-x}，$R_\mathfrak q$ 正则。
\end{proof}

\begin{lemma}
\label{more-algebra-lemma-J-0-goes-down}
设 $R \to S$ 为环映射。假设：
\begin{enumerate}
\item $R$ 是 Noether 整环；
\item $R \to S$ 是单射且为有限型；以及
\item $S$ 是整环且为 J-0。
\end{enumerate}
则 $R$ 是 J-0。
\end{lemma}

\begin{proof}
用某个非零 $g \in S$ 对 $S$ 作 $S_g$ 替换后，可假设
$S$ 是正则环。由一般平坦性，还可假设 $R \to S$ 忠实平坦；
见 Algebra, Lemma
\ref{algebra-lemma-generic-flatness-Noetherian}。
于是由 Algebra, Lemma
\ref{algebra-lemma-descent-regular}，$R$ 正则。
\end{proof}

\begin{lemma}
\label{more-algebra-lemma-J-0-goes-up}
设 $R \to S$ 为环映射。假设：
\begin{enumerate}
\item $R$ 是 Noether 整环且为 J-0；
\item $R \to S$ 是单射且为有限型；
\item $S$ 是整环；以及
\item 所诱导的分式域扩张可分。
\end{enumerate}
则 $S$ 是 J-0。
\end{lemma}

\begin{proof}
可用一个主局部化替换 $R$，并假设 $R$ 是正则环。
由 Algebra, Lemma \ref{algebra-lemma-smooth-at-generic-point}，
环映射 $R \to S$ 在 $(0)$ 处光滑。
因此，用一个主局部化替换 $S$ 后，可假设 $S$ 在 $R$ 上光滑。
于是 $S$ 也正则；见 Algebra, Lemma
\ref{algebra-lemma-regular-goes-up}。
\end{proof}

\begin{lemma}
\label{more-algebra-lemma-J-2}
设 $R$ 为 Noether 环。下列条件等价：
\begin{enumerate}
\item $R$ 是 J-2；
\item 每个作为整环的有限型 $R$-代数都是 J-0；
\item 每个有限 $R$-代数都是 J-1；
\item 对每个素理想 $\mathfrak p$ 和每个有限纯不可分扩张
$L/\kappa(\mathfrak p)$，都存在一个有限 $R$-代数 $R'$，
它是 J-0 整环，且其分式域为 $L$。
\end{enumerate}
\end{lemma}

\begin{proof}
显然有推含 (1) $\Rightarrow$ (2) 和 (2) $\Rightarrow$ (4)。
回忆：一个 J-1 整环是 J-0 的。因此也有推含
(1) $\Rightarrow$ (3) 和 (3) $\Rightarrow$ (4)。

\medskip\noindent
设 $R \to S$ 为有限型环映射，尝试证明 $S$ 是 J-1。
由引理 \ref{more-algebra-lemma-intersection-regular-with-closed}，只需证明
对 $S$ 的每个素理想 $\mathfrak q$，$S/\mathfrak q$ 都是 J-0。
由此可见 (2) $\Rightarrow$ (1)。

\medskip\noindent
假设 (4)。我们将证明 (2)，从而完成证明。
设 $R \to S$ 为有限型环映射，且 $S$ 为整环。
令 $\mathfrak p = \Ker(R \to S)$。设 $K$ 为 $S$ 的分式域。
存在域的交换图
$$
\xymatrix{
K \ar[r] & K' \\
\kappa(\mathfrak p) \ar[u] \ar[r] & L \ar[u]
}
$$
其中水平箭头均为有限纯不可分域扩张，并且 $K'/L$ 可分；
见 Algebra, Lemma
\ref{algebra-lemma-make-separably-generated}。
按 (4) 选取 $R' \subset L$，并令 $S'$ 为映射
$S \otimes_R R' \to K'$ 的像。则 $S'$ 是分式域为 $K'$
的整环，故由引理 \ref{more-algebra-lemma-J-0-goes-up} 以及对 $R'$
的选取，$S'$ 是 J-0。再应用引理 \ref{more-algebra-lemma-J-0-goes-down}，
可知 $S$ 是 J-0，正如所需。
\end{proof}



\section{正则性与导子}
\label{more-algebra-section-regularity-derivations}

\noindent
设 $R \to S$ 为环映射。设 $D : R \to R$ 为导子。
若存在导子 $D' : S \to S$，使得交换图
$$
\xymatrix{
S \ar[r]_{D'} & S \\
R \ar[u] \ar[r]^D & R \ar[u]
}
$$
成立，则称 $D$ {\it 延拓到} $S$。

\begin{lemma}
\label{more-algebra-lemma-derivation-extends}
设 $R$ 为环。设 $D : R \to R$ 为导子。
\begin{enumerate}
\item 对任意理想 $I \subset R$，导子 $D$ 典范地延拓为
$I$-adic 完备化上的导子
$D^\wedge : R^\wedge \to R^\wedge$。
\item 对任意乘法子集 $S \subset R$，导子 $D$ 唯一地延拓到
$R$ 的局部化 $S^{-1}R$。
\end{enumerate}
若 $R \subset R'$ 是有限型环扩张，且对某个在 $R'$ 中为非零因子的
$g \in R$ 有 $R_g \cong R'_g$，则对某个 $N \geq 0$，
$g^ND$ 延拓到 $R'$。
\end{lemma}

\begin{proof}
证明 (1)。由 Leibniz 法则，对 $n \geq 2$ 有
$D(I^n) \subset I^{n - 1}$。因此 $D$ 诱导映射
$D_n : R/I^n \to R/I^{n - 1}$。取极限得到 $D^\wedge$。
略去验证 $D^\wedge$ 是导子的过程。

\medskip\noindent
证明 (2)。为将 $D$ 延拓到 $S^{-1}R$，只需令
$D(r/s) = D(r)/s - rD(s)/s^2$ 并验证各公理。

\medskip\noindent
证明最后一个断言。设 $x_1, \ldots, x_n \in R'$ 为
$R'$ 在 $R$ 上的一组生成元。选取 $N$，使得 $g^Nx_i \in R$。
考察 $g^{N + 1}D$。由 (2)，它延拓到 $R_g$。
此外，由 Leibniz 法则和上面对延拓的构造，有
$$
g^{N + 1}D(x_i) = g^{N + 1}D(g^{-N} g^Nx_i) = -Ng^Nx_iD(g) +
gD(g^Nx_i)
$$
且右端两项均属于 $R$。这蕴含
$$
g^{N + 1}D(x_1^{e_1} \ldots x_n^{e_n}) =
\sum e_i x_1^{e_1} \ldots x_i^{e_i - 1} \ldots x_n^{e_n} g^{N + 1}D(x_i)
$$
是 $R'$ 的元素。因此，$R'$ 的每个元素（它可写成诸 $x_i$
的单项式的和，系数在 $R$ 中）都被 $g^{N + 1}D$
映到 $R'$ 的元素，结论得证。
\end{proof}

\begin{lemma}
\label{more-algebra-lemma-quotient-regular}
\begin{slogan}
正确表述的超曲面 Jacobian 判据。
\end{slogan}
设 $R$ 为正则环。设 $f \in R$。假设存在导子
$D : R \to R$，使得 $D(f)$ 是 $R/(f)$ 的单位。
则 $R/(f)$ 正则。
\end{lemma}

\begin{proof}
只需在 $R$ 为具有极大理想 $\mathfrak m$、剩余域 $\kappa$
的局部环时证明。在这种情况下，只需证明
$f \not \in \mathfrak m^2$；见 Algebra, Lemma
\ref{algebra-lemma-regular-ring-CM}。
然而，若 $f \in \mathfrak m^2$，则由 Leibniz 法则，
$D(f) \in \mathfrak m$，矛盾。
\end{proof}

\begin{lemma}
\label{more-algebra-lemma-quotient-sequence-regular}
设 $(R, \mathfrak m, \kappa)$ 为正则局部环。设 $m \geq 1$。
设 $f_1, \ldots, f_m \in \mathfrak m$。假设存在导子
$D_1, \ldots, D_m : R \to R$，使得
$\det_{1 \leq i, j \leq m}(D_i(f_j))$ 是 $R$ 的单位。
则 $R/(f_1, \ldots, f_m)$ 正则，且
$f_1, \ldots, f_m$ 是正则序列。
\end{lemma}

\begin{proof}
只需证明 $f_1, \ldots, f_m$ 在
$\mathfrak m/\mathfrak m^2$ 中 $\kappa$-线性无关；
见 Algebra, Lemma \ref{algebra-lemma-regular-ring-CM}。
然而，若存在非平凡线性关系，则对某些 $a_i \in R$，且并非所有
$a_i \in \mathfrak m$，有
$\sum a_i f_i \in \mathfrak m^2$。
注意，由导子的 Leibniz 法则，
$D_i(\mathfrak m^2) \subset \mathfrak m$ 且
$D_i(a_j f_j) \equiv a_j D_i(f_j) \bmod \mathfrak m$。
因此，这将蕴含
$$
\sum a_j D_i(f_j) \in \mathfrak m
$$
与关于行列式的假设矛盾。
\end{proof}

\begin{lemma}
\label{more-algebra-lemma-degree-p-extension-regular}
设 $R$ 为正则环。设 $f \in R$。假设存在导子
$D : R \to R$，使得 $D(f)$ 是 $R$ 的单位。
则对任意整数 $n \geq 1$，$R[z]/(z^n - f)$ 正则。
更一般地，对任意 $p \in \mathbf{Z}[z]$，
$R[z]/(p(z) - f)$ 正则。
\end{lemma}

\begin{proof}
由 Algebra, Lemma \ref{algebra-lemma-regular-goes-up} 可见，
$R[z]$ 是正则环。对 $D$ 到 $R[z]$ 的延拓应用引理
\ref{more-algebra-lemma-quotient-regular}，其中将 $z$ 映为零。
这是可行的，因为 $D$ 将任意整系数多项式零化，并将 $f$ 映为单位。
\end{proof}

\begin{lemma}
\label{more-algebra-lemma-find-D}
设 $p$ 为素数。设 $B$ 为整环，并且在 $B$ 中 $p = 0$。
设 $f \in B$ 在 $B$ 的分式域中不是 $p$ 次幂。
若 $B$ 在某个 Noether 完备局部环上有限型，则存在导子
$D : B \to B$，使得 $D(f)$ 非零。
\end{lemma}

\begin{proof}
设 $R$ 为 Noether 完备局部环，并且存在有限型环映射
$R \to B$。当然，可以用它在 $B$ 中的像替换 $R$，
故可假设 $R$ 是特征 $p > 0$ 的整环
（同时仍是 Noether 完备局部环）。
由 Algebra, Lemma
\ref{algebra-lemma-complete-local-Noetherian-domain-finite-over-regular}，
可将 $R$ 写成某个域 $k$ 和整数 $n$ 所给的
$k[[x_1, \ldots, x_n]]$ 的有限扩张。因此，可用
$k[[x_1, \ldots, x_n]]$ 替换 $R$。
接着，使用 Algebra, Lemma
\ref{algebra-lemma-Noether-normalization-over-a-domain}，
将 $R \to B$ 分解为
$$
R \subset R[y_1, \ldots, y_d] \subset B' \subset B
$$
其中 $B'$ 在 $R[y_1, \ldots, y_d]$ 上有限，并且对某个非零
$g \in R$ 有 $B'_g \cong B_g$。注意，对某个充分大的整数 $N$，
$f' = g^{pN} f \in B'$。显然，$f'$ 在 $B'$ 的分式域中不是
$p$ 次幂。若能找到满足 $D'(f') \not = 0$ 的导子
$D' : B' \to B'$，则引理 \ref{more-algebra-lemma-derivation-extends}
保证对某个 $M > 0$，$D = g^MD'$ 延拓到 $B$。此时
$D(f) = g^MD'(f) = g^MD'(g^{-pN}f') = g^{M - pN}D'(f')$
非零。因此，只需在 $B$ 是
$A = k[[x_1, \ldots, x_n]][y_1, \ldots, y_m]$
的有限扩张时证明引理。

\medskip\noindent
假设 $B$ 是
$A = k[[x_1, \ldots, x_n]][y_1, \ldots, y_m]$
的有限扩张。以 $L$ 表示 $B$ 的分式域。
注意，$\text{d}f$ 在 $\Omega_{L/\mathbf{F}_p}$ 中非零；
见 Algebra, Lemma
\ref{algebra-lemma-derivative-zero-pth-power}。
应用引理 \ref{more-algebra-lemma-power-series-ring-subfields}，可找到
有限指数的子域 $k' \subset k$，使得令
$A' = k'[[x_1^p, \ldots, x_n^p]][y_1^p, \ldots, y_m^p]$ 后，
元素 $\text{d}f$ 在 $\Omega_{L/K'}$ 中的像非零，其中 $K'$
是 $A'$ 的分式域。因此，可选取满足 $D'(f) \not = 0$ 的
$K'$-导子 $D' : L \to L$。
依构造，$A' \subset A$ 与 $A \subset B$ 均有限，
故 $A' \subset B$ 有限。选取
$b_1, \ldots, b_t \in B$，使其生成 $B$ 作为 $A'$-模。
则对某些 $f_i, g_i \in B$ 且 $g_i \not = 0$，有
$D'(b_i) = f_i/g_i$。令 $D = g_1 \ldots g_t D'$ 即得结论。
\end{proof}

\begin{lemma}
\label{more-algebra-lemma-complete-Noetherian-domain-J-0}
设 $A$ 为 Noether 完备局部整环。则 $A$ 是 J-0。
\end{lemma}

\begin{proof}
由 Algebra, Lemma
\ref{algebra-lemma-complete-local-Noetherian-domain-finite-over-regular}，
可找到正则子环 $A_0 \subset A$，使 $A$ 在 $A_0$ 上有限。
所诱导的分式域扩张 $K/K_0$ 有限。若 $K/K_0$ 可分，则由引理
\ref{more-algebra-lemma-J-0-goes-up} 即得结论。若不可分，则 $A_0$ 与 $A$
的特征为 $p > 0$。对于任意子扩张 $K/M/K_0$，存在有限子扩张
$A_0 \subset B \subset A$，其分式域为 $M$。
因此，对 $[K : K_0]$ 作归纳，可假设存在
$A_0 \subset B \subset A$，使得 $B$ 是 J-0，且 $K/M$
没有非平凡子扩张。在这种情况下，若 $K/M$ 可分，则由引理
\ref{more-algebra-lemma-J-0-goes-up} 可知 $A$ 是 J-0。
若不可分，则对某些 $b_1, b_2 \in B$ 且 $b_2 \not = 0$，
有 $K = M[z]/(z^p - b_1/b_2)$，其中 $b_1/b_2$ 在 $M$
中不是 $p$ 次幂。选取非零 $a \in A$，使得 $az \in A$。
以 $b_2 a^p z$ 替换 $z$ 后，得到
$K = M[z]/(z^p - b)$，其中 $z \in A$ 且 $b \in B$
在 $M$ 中不是 $p$ 次幂。由引理 \ref{more-algebra-lemma-find-D}，
可找到满足 $D(b) \not = 0$ 的导子 $D : B \to B$。
应用引理 \ref{more-algebra-lemma-degree-p-extension-regular} 可见，
对于 $A$ 中位于 $B$ 的某个正则素理想之上、且不包含
$D(b)$ 的任意素理想 $\mathfrak p$，$A_\mathfrak p$ 正则。
由于 $B$ 是 J-0，故 $A$ 也是 J-0。
\end{proof}

\begin{proposition}
\label{more-algebra-proposition-ubiquity-J-2}
下列各类环都是 J-2：
\begin{enumerate}
\item 域；
\item Noether 完备局部环；
\item $\mathbf{Z}$；
\item 维数为 $1$ 的 Noether 局部环；
\item 维数为 $1$ 的 Nagata 环；
\item 分式域特征为零的 Dedekind 整环；
\item 上述任一类环的有限型环扩张。
\end{enumerate}
\end{proposition}

\begin{proof}
对于情形 (1)、(3)、(5) 与 (6)，可通过验证引理
\ref{more-algebra-lemma-J-2} 的条件 (4) 来证明。我们只处理 $R$
是维数为 $1$ 的 Nagata 环的情形。设
$\mathfrak p \subset R$ 为素理想，并设
$L/\kappa(\mathfrak p)$ 为有限纯不可分扩张。
若 $\mathfrak p \subset R$ 是极大理想，则 $R \to L$ 有限，
且 $L$ 是正则环，故已验证该条件。
若 $\mathfrak p \subset R$ 是极小素理想，则 Nagata 条件保证
% P02-E0182
$R$ 在 $L$ 中的整闭包 $R' \subset L$ 在 $R$ 上有限。
于是 $R'$ 是维数为 $1$ 的正规整环
（Algebra, Lemma \ref{algebra-lemma-integral-dim-up}），
因而正则（Algebra, Lemma
\ref{algebra-lemma-criterion-normal}），故在此情形也已验证该条件。

\medskip\noindent
对于情形 (2)，使用引理 \ref{more-algebra-lemma-J-2} 的条件 (3)。
设 $R$ 为 Noether 完备局部环。注意，若 $R \to R'$ 有限，
则 $R'$ 是 Noether 完备局部环的直积；见 Algebra, Lemma
\ref{algebra-lemma-quotient-complete-local}。
因此，由引理 \ref{more-algebra-lemma-intersection-regular-with-closed}，
只需证明作为整环的 Noether 完备局部环是 J-0；
这就是引理 \ref{more-algebra-lemma-complete-Noetherian-domain-J-0}。

\medskip\noindent
对于情形 (4)，也使用引理 \ref{more-algebra-lemma-J-2} 的条件 (3)。
事实上，若 $R$ 是维数为 $1$ 的 Noether 局部环，且
$R \to R'$ 有限，则 $\Spec(R')$ 是有限集。
由于正则轨迹在一般化下稳定，可见 $R'$ 是 J-1。
\end{proof}





\section{形式光滑性与正则性}
\label{more-algebra-section-fs-regular}

\noindent
本节标题指的是 Andr\'e 定理 \ref{more-algebra-theorem-fs-regular}。

\begin{lemma}
\label{more-algebra-lemma-lift-derivation-through-fs}
设 $A \to B$ 为 Noether 局部环的局部同态。
设 $D : A \to A$ 为导子。假设 $B$ 完备，并且
$A \to B$ 对 $\mathfrak m_B$-adic 拓扑形式光滑。
则存在 $D$ 的延拓 $D' : B \to B$。
\end{lemma}

\begin{proof}
以 $B[\epsilon] = B[x]/(x^2)$ 表示 $B$ 上的对偶数环。
考察环映射 $\psi : A \to B[\epsilon]$，
$a \mapsto a + \epsilon D(a)$。考察交换图
$$
\xymatrix{
B \ar[r]_1 & B \\
A \ar[u] \ar[r]^\psi & B[\epsilon] \ar[u]
}
$$
由引理 \ref{more-algebra-lemma-lift-continuous} 以及 $B/A$ 形式光滑的假设，
可找到适合该图的映射 $\varphi : B \to B[\epsilon]$。
写成 $\varphi(b) = b + \epsilon D'(b)$。于是
$D' : B \to B$ 就是所需的延拓。
\end{proof}

\begin{proposition}
\label{more-algebra-proposition-fs-regular}
设 $A \to B$ 为 Noether 完备局部环的局部同态。
设 $k$ 为 $A$ 的剩余域，并设
$\overline{B} = B \otimes_A k$ 为特殊纤维。
下列条件等价：
\begin{enumerate}
\item $A \to B$ 正则；
\item $A \to B$ 平坦，且 $\overline{B}$ 在 $k$ 上几何正则；
\item $A \to B$ 平坦，且 $k \to \overline{B}$ 对
$\mathfrak m_{\overline{B}}$-adic 拓扑形式光滑；以及
\item $A \to B$ 对 $\mathfrak m_B$-adic 拓扑形式光滑。
\end{enumerate}
\end{proposition}

\begin{proof}
我们已在命题 \ref{more-algebra-proposition-fs-flat-fibre-fs} 中证明
(2)、(3) 与 (4) 等价。显然 (1) 推出 (2)。
因此，假设等价条件 (2)、(3) 与 (4) 成立，并证明 (1)。

\medskip\noindent
设 $\mathfrak p$ 为 $A$ 的素理想。我们将证明
$B \otimes_A \kappa(\mathfrak p)$ 在
$\kappa(\mathfrak p)$ 上几何正则。
由引理 \ref{more-algebra-lemma-base-change-fs}，可以用
$A/\mathfrak p$ 替换 $A$，并用 $B/\mathfrak pB$ 替换 $B$。
故可假设 $A$ 是整环且 $\mathfrak p = (0)$。

\medskip\noindent
按 Algebra, Lemma
\ref{algebra-lemma-complete-local-Noetherian-domain-finite-over-regular}
选取 $A_0 \subset A$。下文将不再说明地使用该引理所述的全部性质。
由于 $A_0 \to A$ 在剩余域上诱导同构，且
$B/\mathfrak m_A B$ 在 $A/\mathfrak m_A$ 上几何正则，
可找到交换图
$$
\xymatrix{
C \ar[r] & B \\
A_0 \ar[r] \ar[u] & A \ar[u]
}
$$
其中 $A_0 \to C$ 对 $\mathfrak m_C$-adic 拓扑形式光滑，
且 $B = C \otimes_{A_0} A$；见注记
\ref{more-algebra-remark-what-does-it-mean}。
（由于 $A_0 \to A$ 有限，张量积中不需要完备化；见 Algebra, Lemma
\ref{algebra-lemma-completion-tensor}。）
因此，只需证明 $C \otimes_{A_0} K_0$ 是 $A_0$ 的分式域
$K_0$ 上的几何正则代数。

\medskip\noindent
上一段的要点是，可以假设
$A = k[[x_1, \ldots, x_n]]$，其中 $k$ 为域；或者
$A = \Lambda[[x_1, \ldots, x_n]]$，其中 $\Lambda$ 为 Cohen 环。
在这种情况下，由 Algebra, Lemma
\ref{algebra-lemma-flat-over-regular-with-regular-fibre}，
$B$ 是正则环。因此，$B \otimes_A K$ 也是正则环
（其中 $K$ 为 $A$ 的分式域）；若 $K$ 的特征为零，结论即得。

\medskip\noindent
因此，只剩下 $A = k[[x_1, \ldots, x_n]]$ 且 $k$
是特征 $p > 0$ 的域这一情形。
设 $L/K$ 为有限纯不可分域扩张。我们对 $[L : K]$ 作归纳，
证明 $B \otimes_A L$ 正则。基础情形为 $L = K$，已在上文证明。
设 $K \subset M \subset L$ 为子域，使得 $L$ 是 $M$
的次数为 $p$ 的扩张，并由添加某个元素 $f \in M$ 的 $p$ 次根得到。
设 $A'$ 是 $M$ 中分式域为 $M$ 的有限 $A$-子代数。
清除分母后，可以且确实假设 $f \in A'$。
令 $A'' = A'[z]/(z^p -f)$，并注意 $A' \subset A''$ 有限，
且 $A''$ 的分式域为 $L$。
由归纳假设，环
% P02-E0183
$B \otimes_A M$ 正则。我们有
$$
B \otimes_A L = B \otimes_A M[z]/(z^p - f)
$$
由引理 \ref{more-algebra-lemma-find-D}，存在导子 $D : A' \to A'$，
使得 $D(f) \not = 0$。由引理 \ref{more-algebra-lemma-descent-fs}，
$A' \to B \otimes_A A'$ 对 $\mathfrak m$-adic 拓扑形式光滑，
故可使用引理 \ref{more-algebra-lemma-lift-derivation-through-fs}
将 $D$ 延拓为导子
$D' : B \otimes_A A' \to B \otimes_A A'$。
注意，$D'(f) = D(f)$ 在 $B \otimes_A M$ 中为单位，
因为 $D(f)$ 在 $A' \subset M$ 中非零。
因此，由引理 \ref{more-algebra-lemma-degree-p-extension-regular}，
$B \otimes_A L$ 正则，结论得证。
\end{proof}

\begin{theorem}[Andr\'e]
\label{more-algebra-theorem-fs-regular}
\begin{reference}
推含 (4) $\Rightarrow$ (1) 是 \cite{Andre-smooth} 的主要结果。
\end{reference}
设 $A \to B$ 为 Noether 局部环的局部同态。
设 $k$ 为 $A$ 的剩余域，并设
$\overline{B} = B \otimes_A k$ 为特殊纤维。
假设 $A \to A^\wedge$ 正则\footnote{例如，若 $A$ 是 G-环或
（拟）优环；见第 \ref{more-algebra-section-G-ring} 节和第
\ref{more-algebra-section-excellent} 节。}。下列条件等价：
\begin{enumerate}
\item $A \to B$ 正则；
\item $A \to B$ 平坦，且 $\overline{B}$ 在 $k$ 上几何正则；
\item $A \to B$ 平坦，且 $k \to \overline{B}$ 对
$\mathfrak m_{\overline{B}}$-adic 拓扑形式光滑；以及
\item $A \to B$ 对 $\mathfrak m_B$-adic 拓扑形式光滑。
\end{enumerate}
\end{theorem}

\begin{proof}
我们已在命题 \ref{more-algebra-proposition-fs-flat-fibre-fs} 中证明
(2)、(3) 与 (4) 等价。显然 (1) 推出 (2)。
因此，假设 (4) 成立，并证明 (1)。

\medskip\noindent
由引理 \ref{more-algebra-lemma-formally-smooth-completion}，
$A^\wedge \to B^\wedge$ 形式光滑。由命题
\ref{more-algebra-proposition-fs-regular}，$A^\wedge \to B^\wedge$ 正则。
依假设，$A \to A^\wedge$ 正则，故由引理
\ref{more-algebra-lemma-regular-composition}，$A \to B^\wedge$ 正则。
由于 $B \to B^\wedge$ 忠实平坦，由引理
\ref{more-algebra-lemma-regular-permanence} 可知 $A \to B$ 正则。
\end{proof}




\section{G-环}
\label{more-algebra-section-G-ring}

\noindent
设 $A$ 为 Noether 局部环。在第
\ref{more-algebra-section-permanence-completion} 节中，
我们已经看到，$A$ 的某些性质（但并非所有性质）会反映在
$A$ 的完备化 $A^\wedge$ 上。为了进一步研究这一现象，
我们引入一些术语。对于 $A$ 的素理想 $\mathfrak q$，纤维环
$$
A^\wedge \otimes_A \kappa(\mathfrak q) =
(A^\wedge)_\mathfrak q/\mathfrak q(A^\wedge)_\mathfrak q =
(A/\mathfrak q)^\wedge \otimes_{A/q} \kappa(\mathfrak q)
$$
称为 $A$ 的一个{\it 形式纤维}。我们把形式纤维看成
$\kappa(\mathfrak q)$-代数。因此，$A \to A^\wedge$ 是正则环同态，
当且仅当所有形式纤维都是几何正则代数。

\begin{definition}
\label{more-algebra-definition-G-ring}
环 $R$ 称为{\it G-环}，如果 $R$ 是 Noether 环，并且对于
$R$ 的每个素理想 $\mathfrak p$，环同态
$R_\mathfrak p \to (R_\mathfrak p)^\wedge$ 都是正则的。
\end{definition}

\noindent
由上面的讨论可知，$R$ 是 G-环，当且仅当每个局部环
$R_\mathfrak p$ 都有几何正则的形式纤维。注意，若
$\mathbf{Q} \subset R$，则只需检验形式纤维正则。
以下引理给出 G-环条件的另一种表述。

\begin{lemma}
\label{more-algebra-lemma-check-G-ring-easy}
设 $R$ 为 Noether 环。那么，$R$ 是 G-环，当且仅当对于每一对素理想
$\mathfrak q \subset \mathfrak p \subset R$，代数
$$
(R/\mathfrak q)_\mathfrak p^\wedge \otimes_{R/\mathfrak q} \kappa(\mathfrak q)
$$
在 $\kappa(\mathfrak q)$ 上几何正则。
\end{lemma}

\begin{proof}
这由以下事实立即得到：作为 $\kappa(\mathfrak q)$-代数，
$$
R_\mathfrak p^\wedge \otimes_R \kappa(\mathfrak q) =
(R/\mathfrak q)_\mathfrak p^\wedge \otimes_{R/\mathfrak q} \kappa(\mathfrak q)
$$
\end{proof}

\begin{lemma}
\label{more-algebra-lemma-G-ring-goes-up-quasi-finite}
设 $R \to R'$ 为 Noether 环之间的有限型同态，并设
$$
\xymatrix{
\mathfrak q' \ar[r] & \mathfrak p' \ar[r] & R' \\
\mathfrak q \ar[r] \ar@{-}[u] &
\mathfrak p \ar[r] \ar@{-}[u] & R \ar[u]
}
$$
为素理想图。假设 $R \to R'$ 在 $\mathfrak p'$ 处拟有限。
\begin{enumerate}
\item 若形式纤维 $R_\mathfrak p^\wedge \otimes_R \kappa(\mathfrak q)$
在 $\kappa(\mathfrak q)$ 上几何正则，则形式纤维
$(R'_{\mathfrak p'})^\wedge \otimes_{R'} \kappa(\mathfrak q')$
在 $\kappa(\mathfrak q')$ 上几何正则。
\item 若 $R_\mathfrak p$ 的形式纤维均几何正则，
则 $R'_{\mathfrak p'}$ 的形式纤维均几何正则。
\item 若 $R \to R'$ 拟有限且 $R$ 是 G-环，则 $R'$ 是 G-环。
\end{enumerate}
\end{lemma}

\begin{proof}
显然，(1) $\Rightarrow$ (2) $\Rightarrow$ (3)。假设
$R_\mathfrak p^\wedge \otimes_R \kappa(\mathfrak q)$
在 $\kappa(\mathfrak q)$ 上几何正则。由 Algebra，引理
\ref{algebra-lemma-completion-at-quasi-finite-prime}，可知
$$
R_\mathfrak p^\wedge \otimes_R R'
=
(R'_{\mathfrak p'})^\wedge \times B
$$
对某个 $R_\mathfrak p^\wedge$-代数 $B$ 成立。因此，
$R'_{\mathfrak p'} \to (R'_{\mathfrak p'})^\wedge$ 是同态
$R_\mathfrak p \to R_\mathfrak p^\wedge$ 的某个基变换的一个因子。
于是 $(R'_{\mathfrak p'})^\wedge \otimes_{R'} \kappa(\mathfrak q')$
是
$$
R_\mathfrak p^\wedge \otimes_R R' \otimes_{R'} \kappa(\mathfrak q') =
R_\mathfrak p^\wedge \otimes_R \kappa(\mathfrak q)
\otimes_{\kappa(\mathfrak q)} \kappa(\mathfrak q').
$$
的一个因子。由于基域扩张保持几何正则性，结论成立；见
Algebra，引理 \ref{algebra-lemma-geometrically-regular}。
\end{proof}

\begin{lemma}
\label{more-algebra-lemma-check-G-ring}
设 $R$ 为 Noether 环。那么，$R$ 是 G-环，当且仅当对于每个有限自由环同态
$R \to S$，$S$ 的形式纤维都是正则环。
\end{lemma}

\begin{proof}
假设对于任意有限自由环同态 $R \to S$，环 $S$ 都有正则形式纤维。
设 $\mathfrak q \subset \mathfrak p \subset R$ 为素理想，并设
$\kappa(\mathfrak q) \subset L$ 为有限纯不可分扩张。要证明 $R$ 是 G-环，
只需证明
$$
R_\mathfrak p^\wedge \otimes_R \kappa(\mathfrak q)
\otimes_{\kappa(\mathfrak q)} L
$$
为正则环。选取有限自由扩张 $R \to R'$，使得
$\mathfrak q' = \mathfrak qR'$ 是素理想，并且
$\kappa(\mathfrak q')$ 在 $\kappa(\mathfrak q)$ 上同构于 $L$；见
Algebra，引理 \ref{algebra-lemma-finite-free-given-residue-field-extension}。
由 Algebra，引理 \ref{algebra-lemma-completion-finite-extension}，有
$$
R_\mathfrak p^\wedge \otimes_R R' = \prod (R'_{\mathfrak p_i'})^\wedge
$$
其中 $\mathfrak p_i'$ 是 $R'$ 中位于 $\mathfrak p$ 上方的素理想。因此，
$$
R_\mathfrak p^\wedge \otimes_R \kappa(\mathfrak q)
\otimes_{\kappa(\mathfrak q)} L =
R_\mathfrak p^\wedge \otimes_R R'
\otimes_{R'} \kappa(\mathfrak q')
=
\prod (R'_{\mathfrak p_i'})^\wedge
\otimes_{R'_{\mathfrak p'_i}} \kappa(\mathfrak q')
$$
依假设，右侧各环正则，故左侧的环也正则。因此 $R$ 是 G-环。
反向推论由引理 \ref{more-algebra-lemma-G-ring-goes-up-quasi-finite} 得到。
\end{proof}

\begin{lemma}
\label{more-algebra-lemma-helper-G-ring}
设 $k$ 为特征 $p$ 的域。令
$A = k[[x_1, \ldots, x_n]][y_1, \ldots, y_m]$，并以 $K$
表示 $A$ 的分式域。设 $\mathfrak p \subset A$ 为素理想。那么
$A_\mathfrak p^\wedge \otimes_A K$ 在 $K$ 上几何正则。
\end{lemma}

\begin{proof}
设 $L/K$ 为有限纯不可分域扩张。我们将对 $[L : K]$ 归纳证明
$A_\mathfrak p^\wedge \otimes L$ 正则。基始情形为 $L = K$：
由于 $A$ 正则，$A_\mathfrak p^\wedge$ 正则
（引理 \ref{more-algebra-lemma-completion-regular}），故其局部化
$A_\mathfrak p^\wedge \otimes K$ 正则。取中间域
$K \subset M \subset L$，使得 $L$ 是通过添加某个 $f \in M$
的一个 $p$ 次根而得到的 $M$ 的 $p$ 次扩张。令 $B$ 为 $M$
的一个有限 $A$-子代数，其分式域为 $M$。清去分母后，
我们可以并且确实假设 $f \in B$。置 $C = B[z]/(z^p -f)$，
并注意 $B \subset C$ 有限，且 $C$ 的分式域为 $L$。由于
$A \subset B \subset C$ 均为有限扩张且 $L/M/K$ 均为纯不可分扩张，
可知 $B$ 或 $C$ 的每个元素都有某个幂属于 $A$。因此，
存在唯一素理想 $\mathfrak r \subset B$，分别地，也存在唯一素理想
$\mathfrak q \subset C$ 位于 $\mathfrak p$ 上方。注意
$$
A_\mathfrak p^\wedge \otimes_A M = B_\mathfrak r^\wedge \otimes_B M
$$
见 Algebra，引理 \ref{algebra-lemma-completion-finite-extension}。
由归纳假设，该环正则。同理有
$$
A_\mathfrak p^\wedge \otimes_A L =
C_\mathfrak r^\wedge \otimes_C L =
B_\mathfrak r^\wedge \otimes_B M[z]/(z^p - f)
$$
最后一个等式成立，是因为 $C = B[z]/(z^p - f)$ 的完备化等于
$B_\mathfrak r^\wedge[z]/(z^p -f)$。由引理 \ref{more-algebra-lemma-find-D}，
存在导子 $D : B \to B$，使得 $D(f) \not = 0$。换言之，
$g = D(f)$ 是 $M$ 中的单位！由引理 \ref{more-algebra-lemma-derivation-extends}，
$D$ 依次通过局部化、完备化和局部化延拓为 $B_\mathfrak r$、
$B_\mathfrak r^\wedge$ 以及 $B_\mathfrak r^\wedge \otimes_B M$ 上的导子。
由于这些是延拓，我们最终得到 $B_\mathfrak r^\wedge \otimes_B M$
上的一个导子，它将 $f$ 映到 $g$，而 $g$ 是环
$B_\mathfrak r^\wedge \otimes_B M$ 的单位。
于是由引理 \ref{more-algebra-lemma-degree-p-extension-regular}，
$A_\mathfrak p^\wedge \otimes_A L$ 正则，结论得证。
\end{proof}

\begin{proposition}
\label{more-algebra-proposition-Noetherian-complete-G-ring}
Noether 完备局部环是 G-环。
\end{proposition}

\begin{proof}
设 $A$ 为 Noether 完备局部环。由引理
\ref{more-algebra-lemma-check-G-ring-easy}，只需检验 $B = A/\mathfrak q$
在 $B$ 的极小素理想 $(0)$ 上方具有几何正则的形式纤维。
因此，可以假设 $A$ 为整环，并且只需检验 $A$ 的极小素理想
$(0)$ 上方的形式纤维满足该条件。令 $K$ 为 $A$ 的分式域。

\medskip\noindent
可以选取子环 $A_0 \subset A$，它是正则完备局部环且
$A$ 在 $A_0$ 上有限；见 Algebra，引理
\ref{algebra-lemma-complete-local-Noetherian-domain-finite-over-regular}。
此外，可以假设 $A_0$ 是域或 Cohen 环上的幂级数环。
由引理 \ref{more-algebra-lemma-G-ring-goes-up-quasi-finite}，只需对 $A_0$ 证明结论。

\medskip\noindent
假设 $A$ 是域或 Cohen 环上的幂级数环。由于 $A$ 正则，
各局部化 $A_\mathfrak p$ 均正则（见 Algebra，定义
\ref{algebra-definition-regular} 及其前面的讨论）。因此，
各完备化 $A_\mathfrak p^\wedge$ 均正则；见引理
\ref{more-algebra-lemma-completion-regular}。于是纤维
$A_{\mathfrak p}^\wedge \otimes_A K$ 作为
$A_\mathfrak p^\wedge$ 的局部化也正则。所以当 $K$ 的特征为 $0$
时结论成立。正特征情形就是 $A = k[[x_1, \ldots, x_d]]$ 的情形，
它是引理 \ref{more-algebra-lemma-helper-G-ring} 的特殊情形。
\end{proof}

\begin{lemma}
\label{more-algebra-lemma-check-G-ring-maximal-ideals}
设 $R$ 为 Noether 环。那么，$R$ 是 G-环，当且仅当对于 $R$
的每个极大理想 $\mathfrak m$，$R_\mathfrak m$ 都有几何正则的形式纤维。
\end{lemma}

\begin{proof}
假设对于 $R$ 的每个极大理想 $\mathfrak m$，同态
$R_\mathfrak m \to R_\mathfrak m^\wedge$ 都正则。设 $\mathfrak p$
为 $R$ 的素理想，并选取极大理想 $\mathfrak p \subset \mathfrak m$。
由于 $R_\mathfrak m \to R_\mathfrak m^\wedge$ 忠实平坦，
可以在 $R_\mathfrak m^\wedge$ 中选取一个位于
$\mathfrak pR_\mathfrak m$ 上方的素理想 $\mathfrak p'$。
考虑交换图
$$
\xymatrix{
R_\mathfrak m^\wedge \ar[r] &
(R_\mathfrak m^\wedge)_{\mathfrak p'} \ar[r] &
(R_\mathfrak m^\wedge)_{\mathfrak p'}^\wedge
\\
R_\mathfrak m \ar[u] \ar[r] & R_\mathfrak p \ar[u] \ar[r] &
R_\mathfrak p^\wedge \ar[u]
}
$$
依假设，环同态 $R_\mathfrak m \to R_\mathfrak m^\wedge$ 正则。
由命题 \ref{more-algebra-proposition-Noetherian-complete-G-ring}，
$(R_\mathfrak m^\wedge)_{\mathfrak p'} \to
(R_\mathfrak m^\wedge)_{\mathfrak p'}^\wedge$ 正则。
局部化 $R_\mathfrak m^\wedge \to
(R_\mathfrak m^\wedge)_{\mathfrak p'}$ 正则。因此，由引理
\ref{more-algebra-lemma-regular-composition}，$R_\mathfrak m \to
(R_\mathfrak m^\wedge)_{\mathfrak p'}^\wedge$ 正则。
由于它分解经过局部化 $R_\mathfrak p$，环同态
$R_\mathfrak p \to (R_\mathfrak m^\wedge)_{\mathfrak p'}^\wedge$
也正则。因此，可以应用引理 \ref{more-algebra-lemma-regular-permanence} 得到
$R_\mathfrak p \to R_\mathfrak p^\wedge$ 正则。
\end{proof}

\begin{lemma}
\label{more-algebra-lemma-henselization-G-ring}
设 $R$ 为 Noether 局部 G-环。那么，Hensel 化 $R^h$
和严格 Hensel 化 $R^{sh}$ 都是 G-环。
\end{lemma}

\begin{proof}
我们将使用引理 \ref{more-algebra-lemma-check-G-ring-maximal-ideals} 的判据。
设 $\mathfrak q \subset R^h$ 为素理想，并置
$\mathfrak p = R \cap \mathfrak q$。置 $\mathfrak q_1 = \mathfrak q$，
并令 $\mathfrak q_2, \ldots, \mathfrak q_t$ 为 $R^h$ 中位于
$\mathfrak p$ 上方的其余素理想，于是
$R^h \otimes_R \kappa(\mathfrak p) =
\prod\nolimits_{i = 1, \ldots, t} \kappa(\mathfrak q_i)$；见引理
\ref{more-algebra-lemma-fibres-henselization}。利用 $(R^h)^\wedge = R^\wedge$
（引理 \ref{more-algebra-lemma-henselization-noetherian}），可得
$$
\prod\nolimits_{i = 1, \ldots, t}
(R^h)^\wedge \otimes_{R^h} \kappa(\mathfrak q_i) =
(R^h)^\wedge \otimes_{R^h} (R^h \otimes_R \kappa(\mathfrak p)) =
R^\wedge \otimes_R \kappa(\mathfrak p)
$$
依假设，$(R^h)^\wedge \otimes_{R^h} \kappa(\mathfrak q_i)$
在 $\kappa(\mathfrak p)$ 上几何正则。由于
$\kappa(\mathfrak q_i)$ 是 $\kappa(\mathfrak p)$ 的可分代数扩张，
由 Algebra，引理
\ref{algebra-lemma-geometrically-regular-over-separable-algebraic}，可知
$(R^h)^\wedge \otimes_{R^h} \kappa(\mathfrak q_i)$
在 $\kappa(\mathfrak q_i)$ 上几何正则。

\medskip\noindent
设 $\mathfrak r \subset R^{sh}$ 为素理想，并置
$\mathfrak p = R \cap \mathfrak r$。置 $\mathfrak r_1 = \mathfrak r$，
并令 $\mathfrak r_2, \ldots, \mathfrak r_s$ 为 $R^{sh}$ 中位于
$\mathfrak p$ 上方的其余素理想，于是
$R^{sh} \otimes_R \kappa(\mathfrak p) =
\prod\nolimits_{i = 1, \ldots, s} \kappa(\mathfrak r_i)$；见引理
\ref{more-algebra-lemma-fibres-henselization}。于是可得
$$
\prod\nolimits_{i = 1, \ldots, s}
(R^{sh})^\wedge \otimes_{R^{sh}} \kappa(\mathfrak r_i) =
(R^{sh})^\wedge \otimes_{R^{sh}} (R^{sh} \otimes_R \kappa(\mathfrak p)) =
(R^{sh})^\wedge \otimes_R \kappa(\mathfrak p)
$$
注意，由引理 \ref{more-algebra-lemma-henselization-noetherian}，
$R^\wedge \to (R^{sh})^\wedge$ 对
$\mathfrak m_{(R^{sh})^\wedge}$-adic 拓扑形式光滑。因此，由命题
\ref{more-algebra-proposition-fs-regular}，$R^\wedge \to (R^{sh})^\wedge$ 正则。
依假设，$R^\wedge \otimes_R \kappa(\mathfrak p)$ 在
$\kappa(\mathfrak p)$ 上正则，故由引理
\ref{more-algebra-lemma-regular-composition}，可知
$(R^{sh})^\wedge \otimes_{R^{sh}} \kappa(\mathfrak r_i)$
在 $\kappa(\mathfrak p)$ 上正则。由于 $\kappa(\mathfrak r_i)$
是 $\kappa(\mathfrak p)$ 的可分代数扩张，由 Algebra，引理
\ref{algebra-lemma-geometrically-regular-over-separable-algebraic}，可知
$(R^{sh})^\wedge \otimes_{R^{sh}} \kappa(\mathfrak r_i)$
在 $\kappa(\mathfrak r_i)$ 上几何正则。
\end{proof}

\begin{lemma}
\label{more-algebra-lemma-another-helper-G-ring}
设 $p$ 为素数。设 $A$ 为 Noether 完备局部整环，其分式域 $K$
的特征为 $p$。设 $\mathfrak q \subset A[x]$ 为位于 $A$
的极大理想上方的极大理想，并设
$(0) \not = \mathfrak r \subset \mathfrak q$ 为位于
$(0) \subset A$ 上方的素理想。那么，
$A[x]_\mathfrak q^\wedge \otimes_{A[x]} \kappa(\mathfrak r)$ 在
$\kappa(\mathfrak r)$ 上几何正则。
\end{lemma}

\begin{proof}
注意，$K \subset \kappa(\mathfrak r)$ 是有限扩张。因此，给定有限纯不可分扩张
$L/\kappa(\mathfrak r)$，存在 Noether 完备局部整环的有限扩张
$A \subset B$，使得 $\kappa(\mathfrak r) \otimes_A B$ 满射到 $L$。
具体地，取 $B \subset L$ 为有限 $A$-子代数，其分式域为 $L$。
以 $\mathfrak r' \subset B[x]$ 表示映射
$B[x] = A[x] \otimes_A B \to \kappa(\mathfrak r) \otimes_A B \to L$
的核，于是 $\kappa(\mathfrak r') = L$。那么
$$
A[x]_\mathfrak q^\wedge \otimes_{A[x]} L =
A[x]_\mathfrak q^\wedge \otimes_{A[x]} B[x]
\otimes_{B[x]} \kappa(\mathfrak r') =
\prod B[x]_{\mathfrak q_i}^\wedge \otimes_{B[x]} \kappa(\mathfrak r')
$$
其中 $\mathfrak q_1, \ldots, \mathfrak q_t$ 是 $B[x]$ 中位于
$\mathfrak q$ 上方的素理想；见 Algebra，引理
\ref{algebra-lemma-completion-finite-extension}。因此，只需证明各环
$B[x]_{\mathfrak q_i}^\wedge \otimes_{B[x]} \kappa(\mathfrak r')$
均正则。这就把问题化约为在 $K = \kappa(\mathfrak r)$
这一特殊情形下证明
$A[x]_\mathfrak q^\wedge \otimes_{A[x]} \kappa(\mathfrak r)$ 正则。

\medskip\noindent
假设 $K = \kappa(\mathfrak r)$。在此情形下，
$\mathfrak r K[x]$ 由某个 $f \in K$ 的 $x - f$ 生成，并且
$$
A[x]_\mathfrak q^\wedge \otimes_{A[x]} \kappa(\mathfrak r)
=
(A[x]_\mathfrak q^\wedge \otimes_A K)/(x - f)
$$
$A[x]$ 上的导子 $D = \text{d}/\text{d}x$ 延拓到 $K[x]$，
并将 $x - f$ 映到 $K[x]$ 的一个单位。此外，由引理
\ref{more-algebra-lemma-derivation-extends}，$D$ 延拓到
$A[x]_\mathfrak q^\wedge \otimes_A K$。由于
$A \to A[x]_\mathfrak q^\wedge$ 形式光滑（见引理
\ref{more-algebra-lemma-formally-smooth} 和
\ref{more-algebra-lemma-formally-smooth-completion}），由命题
\ref{more-algebra-proposition-fs-regular}，环
$A[x]_\mathfrak q^\wedge \otimes_A K$ 正则（在这一特殊情形下，
该命题证明中的论证显著简化）。最后应用引理
\ref{more-algebra-lemma-quotient-regular} 即得结论。
\end{proof}

\begin{proposition}
\label{more-algebra-proposition-finite-type-over-G-ring}
设 $R$ 为 G-环。若 $R \to S$ 本质有限型，则 $S$ 是 G-环。
\end{proposition}

\begin{proof}
由于 G-环性质是局部环的性质，显然 G-环的局部化仍为 G-环。
反之，若在每个素理想处的局部化都是 G-环，则该环是 G-环。
因此，只需对每个有限型 $R$-代数 $S$ 及 $S$ 的每个素理想
$\mathfrak q$ 证明 $S_\mathfrak q$ 是 G-环。将 $S$ 写成
$R[x_1, \ldots, x_n]$ 的商，由引理
\ref{more-algebra-lemma-G-ring-goes-up-quasi-finite} 可知，只需证明
$R[x_1, \ldots, x_n]$ 是 G-环。对 $n$ 归纳，只需证明
$R[x]$ 是 G-环。设 $\mathfrak q \subset R[x]$ 为极大理想。
由引理 \ref{more-algebra-lemma-check-G-ring-maximal-ideals}，只需证明
$$
R[x]_\mathfrak q \longrightarrow R[x]_\mathfrak q^\wedge
$$
正则。若 $\mathfrak q$ 位于 $\mathfrak p \subset R$ 上方，
则可以用 $R_\mathfrak p$ 代替 $R$。因此，可以假设 $R$
是极大理想为 $\mathfrak m$ 的 Noether 局部 G-环，并且
$\mathfrak q \subset R[x]$ 位于 $\mathfrak m$ 上方。注意，存在唯一的素理想
$\mathfrak q' \subset R^\wedge[x]$ 位于 $\mathfrak q$ 上方。考虑交换图
$$
\xymatrix{
R[x]_\mathfrak q^\wedge \ar[r] &
(R^\wedge[x]_{\mathfrak q'})^\wedge \\
R[x]_\mathfrak q \ar[r] \ar[u] & R^\wedge[x]_{\mathfrak q'} \ar[u]
}
$$
由于 $R$ 是 G-环，下方水平箭头正则（它是正则环同态
$R \to R^\wedge$ 的某个基变换的局部化）。假设可以证明右侧竖直箭头
正则。那么复合 $R[x]_\mathfrak q \to
(R^\wedge[x]_{\mathfrak q'})^\wedge$ 正则，从而由引理
\ref{more-algebra-lemma-regular-permanence}，左侧竖直箭头正则。
因此，可以假设 $R$ 是 Noether 完备局部环，并且 $\mathfrak q$
是位于 $R$ 的极大理想上方的素理想。

\medskip\noindent
设 $R$ 为 Noether 完备局部环，并设 $\mathfrak q \subset R[x]$
为位于 $R$ 的极大理想上方的极大理想。设
$\mathfrak r \subset \mathfrak q$ 为素理想。我们要证明
$R[x]_\mathfrak q^\wedge \otimes_{R[x]} \kappa(\mathfrak r)$
是 $\kappa(\mathfrak r)$ 上的几何正则代数。置
$\mathfrak p = R \cap \mathfrak r$。由引理
\ref{more-algebra-lemma-check-G-ring-easy}，可以用 $R/\mathfrak p$ 代替 $R$，
并用 $\mathfrak q$ 和 $\mathfrak r$ 在 $R/\mathfrak p[x]$ 中的像
代替它们。因此，可以假设 $R$ 是整环且 $\mathfrak r \cap R = (0)$。

\medskip\noindent
由 Algebra，引理
\ref{algebra-lemma-complete-local-Noetherian-domain-finite-over-regular}，
可以找到正则子环 $R_0 \subset R$，使得 $R$ 在 $R_0$ 上有限。
应用引理 \ref{more-algebra-lemma-G-ring-goes-up-quasi-finite} 可知，只需在上述假设之外
再假设 $R$ 是正则完备局部环，并证明
$R[x]_\mathfrak q^\wedge \otimes_{R[x]} \kappa(\mathfrak r)$
在 $\kappa(\mathfrak r)$ 上几何正则。

\medskip\noindent
现在 $R$ 是正则完备局部环，且有
$\mathfrak r \subset \mathfrak q \subset R[x]$、
$(0) = R \cap \mathfrak r$，并且 $\mathfrak q$ 是位于 $R$
的极大理想上方的极大理想。由于 $R$ 正则，由 Algebra，引理
\ref{algebra-lemma-regular-goes-up}，环 $R[x]$ 正则。
因此，局部化 $R[x]_\mathfrak q$ 正则。于是，完备化
$R[x]_\mathfrak q^\wedge$ 正则；见引理
\ref{more-algebra-lemma-completion-regular}。因此，纤维
$R[x]_{\mathfrak q}^\wedge \otimes_{R[x]} \kappa(\mathfrak r)$
作为 $R[x]_\mathfrak q^\wedge$ 的局部化也正则。所以，当 $R$
的分式域的特征为 $0$ 时，结论成立。

\medskip\noindent
若 $R$ 的特征为正，则 $R = k[[x_1, \ldots, x_n]]$。
在这种情况下，我们将论证分为两个子情形：
\begin{enumerate}
\item 情形 $\mathfrak r = (0)$。结论直接由引理
\ref{more-algebra-lemma-helper-G-ring} 得到。
\item 情形 $\mathfrak r \not = (0)$。这就是引理
\ref{more-algebra-lemma-another-helper-G-ring}。
\end{enumerate}
\end{proof}

\begin{remark}
\label{more-algebra-remark-G-does-not-survive-completion}
设 $R$ 为 G-环，并设 $I \subset R$ 为理想。一般而言，
$I$-adic 完备化 $R^\wedge$ 不一定是 G-环。Nishimura 在
\cite{Nishimura} 中给出了一个例子。该例子的推广以及某种意义上的澄清，
可见 \cite{Dumitrescu} 的最后一节。
\end{remark}

\begin{proposition}
\label{more-algebra-proposition-ubiquity-G-ring}
下列各类环都是 G-环：
\begin{enumerate}
\item 域；
\item Noether 完备局部环；
\item $\mathbf{Z}$；
\item 分式域特征为零的 Dedekind 整环；
\item 上述任一类环的有限型环扩张。
\end{enumerate}
\end{proposition}

\begin{proof}
对于域、$\mathbf{Z}$ 和特征为零的 Dedekind 整环，
结论立即由定义以及离散赋值环的完备化仍为离散赋值环这一事实得到。
由命题 \ref{more-algebra-proposition-Noetherian-complete-G-ring}，
Noether 完备局部环是 G-环。关于有限型扩张环的陈述就是命题
\ref{more-algebra-proposition-finite-type-over-G-ring}。
\end{proof}

\begin{lemma}
\label{more-algebra-lemma-henselian-local-limit-G-rings}
设 $(A, \mathfrak m)$ 为 Hensel 局部环。那么，$A$ 是一个由
Hensel 局部 G-环构成且转移同态均为局部同态的系统的滤过余极限。
\end{lemma}

\begin{proof}
将 $A = \colim A_i$ 写成有限型 $\mathbf{Z}$-代数的滤过余极限。
设 $\mathfrak p_i$ 为 $A_i$ 中位于 $\mathfrak m$ 下方的素理想。
可以用 $A_i$ 在 $\mathfrak p_i$ 处的局部化代替 $A_i$。
此时 $A_i$ 是 Noether 局部 G-环（命题
\ref{more-algebra-proposition-ubiquity-G-ring}）。由引理
\ref{more-algebra-lemma-henselization-colimit}，可知 $A = \colim A_i^h$。
由引理 \ref{more-algebra-lemma-henselization-G-ring}，各环 $A_i^h$ 都是 G-环。
\end{proof}

\begin{lemma}
\label{more-algebra-lemma-map-G-ring-to-completion-regular}
\begin{reference}
\cite[Theorem 79]{MatCA}
\end{reference}
设 $A$ 为 G-环。设 $I \subset A$ 为理想，并设 $A^\wedge$
为 $A$ 关于 $I$ 的完备化。那么 $A \to A^\wedge$ 正则。
\end{lemma}

\begin{proof}
由 Algebra，引理 \ref{algebra-lemma-completion-flat}，环同态
$A \to A^\wedge$ 平坦。由 Algebra，引理
\ref{algebra-lemma-completion-Noetherian-Noetherian}，环 $A^\wedge$
是 Noether 环。因此，只需检验引理 \ref{more-algebra-lemma-regular-local}
的第三个条件。设 $\mathfrak m' \subset A^\wedge$ 为位于
$\mathfrak m \subset A$ 上方的极大理想。由 Algebra，引理
\ref{algebra-lemma-radical-completion}，有
$IA^\wedge \subset \mathfrak m'$。由于 $A^\wedge/IA^\wedge = A/I$，
可知 $I \subset \mathfrak m$、
$\mathfrak m/I = \mathfrak m'/IA^\wedge$，并且
$A/\mathfrak m = A^\wedge/\mathfrak m'$。由于
$A^\wedge/\mathfrak m'$ 是域，可知 $\mathfrak m$ 也是极大理想。
于是，$A_\mathfrak m \to A^\wedge_{\mathfrak m'}$ 是 Noether
局部环之间的平坦局部同态，它诱导剩余域的同构，并满足
$\mathfrak m A^\wedge_{\mathfrak m'} =
\mathfrak m'A^\wedge_{\mathfrak m'}$。因此，由引理
\ref{more-algebra-lemma-flat-unramified}，它在完备局部环上诱导同构。
设 $(A_\mathfrak m)^\wedge$ 为 $A_\mathfrak m$ 关于其极大理想的完备化。
环同态
$$
(A^\wedge)_{\mathfrak m'} \to
((A^\wedge)_{\mathfrak m'})^\wedge = (A_\mathfrak m)^\wedge
$$
忠实平坦（Algebra，引理
\ref{algebra-lemma-completion-faithfully-flat}）。因此，可以对环同态
$$
A_\mathfrak m \to (A^\wedge)_{\mathfrak m'} \to (A_\mathfrak m)^\wedge
$$
应用引理 \ref{more-algebra-lemma-regular-permanence} 得到结论，因为 $A$ 是 G-环，
所以 $A_\mathfrak m \to (A_\mathfrak m)^\wedge$ 正则。
\end{proof}

\begin{lemma}
\label{more-algebra-lemma-henselization-pair-G-ring}
\begin{slogan}
G-环性质在沿理想的 Hensel 化下稳定
\end{slogan}
\begin{reference}
\cite[Theorem 5.3 i)]{Greco}
\end{reference}
设 $A$ 为 G-环。设 $I \subset A$ 为理想。设 $(A^h, I^h)$
为偶对 $(A, I)$ 的 Hensel 化；见引理 \ref{more-algebra-lemma-henselization}。
那么 $A^h$ 是 G-环。
\end{lemma}

\begin{proof}
设 $\mathfrak m^h \subset A^h$ 为极大理想。由引理
\ref{more-algebra-lemma-check-G-ring-maximal-ideals}，我们必须证明从
$A^h_{\mathfrak m^h}$ 到其完备化的同态具有几何正则的纤维。
设 $\mathfrak m$ 为 $\mathfrak m^h$ 在 $A$ 中的逆像。注意，
$I^h \subset \mathfrak m^h$，从而 $I \subset \mathfrak m$，
因为 $(A^h, I^h)$ 是 Hensel 偶对。回顾：$A^h$ 是 Noether 环，
$I^h = IA^h$，并且 $A \to A^h$ 在 $I$-adic 完备化上诱导同构；
见引理 \ref{more-algebra-lemma-henselization-Noetherian-pair}。于是，Noether
局部环之间的局部同态
$$
A_\mathfrak m \to A^h_{\mathfrak m^h}
$$
由引理 \ref{more-algebra-lemma-flat-unramified} 在极大理想完备化上诱导同构
（略去细节）。设 $\mathfrak q^h$ 为 $A^h_{\mathfrak m^h}$
中位于 $\mathfrak q \subset A_\mathfrak m$ 上方的素理想。
置 $\mathfrak q_1 = \mathfrak q^h$，并令
$\mathfrak q_2, \ldots, \mathfrak q_t$ 为 $A^h$ 中位于
$\mathfrak q$ 上方的其余素理想，于是
$A^h \otimes_A \kappa(\mathfrak q) =
\prod\nolimits_{i = 1, \ldots, t} \kappa(\mathfrak q_i)$；见引理
\ref{more-algebra-lemma-filtered-colimit-etale-noetherian-fibres}。利用上面讨论过的
$(A^h)_{\mathfrak m^h}^\wedge = (A_\mathfrak m)^\wedge$，可得
$$
\prod\nolimits_{i = 1, \ldots, t}
(A^h_{\mathfrak m^h})^\wedge \otimes_{A^h_{\mathfrak m^h}}
\kappa(\mathfrak q_i) =
(A^h_{\mathfrak m^h})^\wedge \otimes_{A^h_{\mathfrak m^h}}
(A^h_{\mathfrak m^h} \otimes_{A_{\mathfrak m}} \kappa(\mathfrak q)) =
(A_{\mathfrak m})^\wedge \otimes_{A_{\mathfrak m}} \kappa(\mathfrak q)
$$
因此，作为其中一个因子，环
$$
(A^h_{\mathfrak m^h})^\wedge \otimes_{A^h_{\mathfrak m^h}}
\kappa(\mathfrak q^h)
$$
依关于 $A$ 的假设在 $\kappa(\mathfrak q)$ 上几何正则。
由于 $\kappa(\mathfrak q^h)$ 是 $\kappa(\mathfrak q)$
的可分代数扩张，由 Algebra，引理
\ref{algebra-lemma-geometrically-regular-over-separable-algebraic}，可知
$$
(A^h_{\mathfrak m^h})^\wedge \otimes_{A^h_{\mathfrak m^h}}
\kappa(\mathfrak q^h)
$$
如所需地在 $\kappa(\mathfrak q^h)$ 上几何正则。
\end{proof}








\section{形式纤维的性质}
\label{more-algebra-section-properties-formal-fibres}

\noindent
在本节中，我们重新处理第 \ref{more-algebra-section-G-ring} 节中的一些论证，
以便能够准确讨论 Noether 环的形式纤维所具有的性质。

\medskip\noindent
设 $P$ 是域 $k$ 到 Noether 环 $R$ 的环同态 $k \to R$
的一种性质。对于 $B$ 为 Noether 环的环同态 $A \to B$，
若对 $A$ 的每个素理想 $\mathfrak q$，同态
$\kappa(\mathfrak q) \to B \otimes_A \kappa(\mathfrak q)$ 都具有 $P$，
则称 $P$ 对该环同态的纤维成立。以下将使用这些断言：
\begin{enumerate}
\item[(A)] 对有限生成域扩张 $k'/k$，有
$P(k \to R) \Rightarrow P(k' \to R \otimes_k k')$；
\item[(B)] $P(k \to R_\mathfrak p),\ \forall \mathfrak p \in \Spec(R)
\Leftrightarrow P(k \to R)$；
\item[(C)] 给定 Noether 环的平坦同态 $A \to B \to C$，
若 $A \to B$ 的纤维具有 $P$ 且 $B \to C$ 正则，
则 $A \to C$ 的纤维具有 $P$；
\item[(D)] 给定 Noether 环的平坦同态 $A \to B \to C$，
若 $A \to C$ 的纤维具有 $P$ 且 $B \to C$ 忠实平坦，
则 $A \to B$ 的纤维具有 $P$；
\item[(E)] 给定 $k \to k' \to R$，其中 $R$ 为 Noether 环，
若 $k'/k$ 是可分代数扩张且 $P(k \to R)$ 成立，
则 $P(k' \to R)$ 成立；以及
\item[(F)] 在此添加更多内容。
\end{enumerate}
给定 Noether 局部环 $A$，若 $P$ 对 $A \to A^\wedge$ 的纤维成立，
则称“{\it $A$ 的形式纤维具有 $P$}”。若 $R$ 是 Noether 环，
且对 $R$ 的所有素理想 $\mathfrak p$，$R_\mathfrak p$ 的形式纤维
都具有 $P$，则称{\it $R$ 是 $P$-环}。

\begin{lemma}
\label{more-algebra-lemma-check-P-ring-easy}
设 $R$ 为 Noether 环，并设 $P$ 为上述性质。那么，$R$ 是 $P$-环，
当且仅当对于每一对素理想 $\mathfrak q \subset \mathfrak p \subset R$，
$\kappa(\mathfrak q)$-代数
$$
(R/\mathfrak q)_\mathfrak p^\wedge \otimes_{R/\mathfrak q} \kappa(\mathfrak q)
$$
都具有性质 $P$。
\end{lemma}

\begin{proof}
这由以下事实立即得到：作为 $\kappa(\mathfrak q)$-代数，
$$
R_\mathfrak p^\wedge \otimes_R \kappa(\mathfrak q) =
(R/\mathfrak q)_\mathfrak p^\wedge \otimes_{R/\mathfrak q} \kappa(\mathfrak q)
$$
\end{proof}

\begin{lemma}
\label{more-algebra-lemma-P-local}
设 $R \to \Lambda$ 为 Noether 环同态。假设 $P$ 具有性质 (B)。
下列条件等价：
\begin{enumerate}
\item $R \to \Lambda$ 的纤维具有 $P$；
\item 对每个位于 $\mathfrak p \subset R$ 上方的素理想
$\mathfrak q \subset \Lambda$，$R_\mathfrak p \to \Lambda_\mathfrak q$
的纤维具有 $P$；以及
\item 对每个位于 $R$ 的 $\mathfrak m$ 上方的极大理想
$\mathfrak m' \subset \Lambda$，
$R_\mathfrak m \to \Lambda_{\mathfrak m'}$ 的纤维具有 $P$。
\end{enumerate}
\end{lemma}

\begin{proof}
设 $\mathfrak p \subset R$ 为素理想。那么，$\mathfrak p$ 上方的纤维是环
$\Lambda \otimes_R \kappa(\mathfrak p)$，其谱双射地映到
$\Spec(\Lambda)$ 中由位于 $\mathfrak p$ 上方的素理想
$\mathfrak q$ 构成的子集；见 Algebra，注
\ref{algebra-remark-fundamental-diagram}。对于这样的素理想
$\mathfrak q$，选取极大理想 $\mathfrak q \subset \mathfrak m'$，
并置 $\mathfrak m = R \cap \mathfrak m'$。那么
$\mathfrak p \subset \mathfrak m$，且有
$$
(\Lambda \otimes_R \kappa(\mathfrak p))_\mathfrak q \cong
(\Lambda_{\mathfrak m'} \otimes_{R_\mathfrak m}
\kappa(\mathfrak p))_\mathfrak q
$$
作为 $\kappa(\mathfrak q)$-代数。因此，(1)、(2) 和 (3) 等价，
因为由 (B)，可以在局部环上检验性质 $P$。
\end{proof}

\begin{lemma}
\label{more-algebra-lemma-P-ring-goes-up-quasi-finite}
设 $R \to R'$ 为 Noether 环之间的有限型同态，并设
$$
\xymatrix{
\mathfrak q' \ar[r] & \mathfrak p' \ar[r] & R' \\
\mathfrak q \ar[r] \ar@{-}[u] &
\mathfrak p \ar[r] \ar@{-}[u] & R \ar[u]
}
$$
为素理想图。假设 $R \to R'$ 在 $\mathfrak p'$ 处拟有限。
假设 $P$ 满足 (A) 和 (B)。
\begin{enumerate}
\item 若 $\kappa(\mathfrak q) \to
R_\mathfrak p^\wedge \otimes_R \kappa(\mathfrak q)$
具有 $P$，则
$\kappa(\mathfrak q') \to
(R'_{\mathfrak p'})^\wedge \otimes_{R'} \kappa(\mathfrak q')$
具有 $P$。
\item 若 $R_\mathfrak p$ 的形式纤维具有 $P$，
则 $R'_{\mathfrak p'}$ 的形式纤维具有 $P$。
\item 若 $R \to R'$ 拟有限且 $R$ 是 $P$-环，则 $R'$ 是 $P$-环。
\end{enumerate}
\end{lemma}

\begin{proof}
显然，(1) $\Rightarrow$ (2) $\Rightarrow$ (3)。假设 $P$ 对
$\kappa(\mathfrak q) \to R_\mathfrak p^\wedge \otimes_R \kappa(\mathfrak q)$
成立。由 Algebra，引理
\ref{algebra-lemma-completion-at-quasi-finite-prime}，可知
$$
R_\mathfrak p^\wedge \otimes_R R'
=
(R'_{\mathfrak p'})^\wedge \times B
$$
对某个 $R_\mathfrak p^\wedge$-代数 $B$ 成立。因此，
$R'_{\mathfrak p'} \to (R'_{\mathfrak p'})^\wedge$ 是同态
$R_\mathfrak p \to R_\mathfrak p^\wedge$ 的某个基变换的一个因子。
于是 $(R'_{\mathfrak p'})^\wedge \otimes_{R'} \kappa(\mathfrak q')$
是
$$
R_\mathfrak p^\wedge \otimes_R R' \otimes_{R'} \kappa(\mathfrak q') =
R_\mathfrak p^\wedge \otimes_R \kappa(\mathfrak q)
\otimes_{\kappa(\mathfrak q)} \kappa(\mathfrak q').
$$
的一个因子。因此，结论由关于 $P$ 的假设得到。
\end{proof}

\begin{lemma}
\label{more-algebra-lemma-check-P-ring-maximal-ideals}
设 $R$ 为 Noether 环。假设 $P$ 满足 (C) 和 (D)。那么，
$R$ 是 $P$-环，当且仅当对 $R$ 的每个极大理想 $\mathfrak m$，
$R_\mathfrak m$ 的形式纤维都具有 $P$。
\end{lemma}

\begin{proof}
假设对 $R$ 的所有极大理想 $\mathfrak m$，$R_\mathfrak m$
的形式纤维都具有 $P$。设 $\mathfrak p$ 为 $R$ 的素理想，
并选取极大理想 $\mathfrak p \subset \mathfrak m$。由于
$R_\mathfrak m \to R_\mathfrak m^\wedge$ 忠实平坦，
可以选取 $R_\mathfrak m^\wedge$ 中位于
$\mathfrak pR_\mathfrak m$ 上方的素理想 $\mathfrak p'$。
考虑交换图
$$
\xymatrix{
R_\mathfrak m^\wedge \ar[r] &
(R_\mathfrak m^\wedge)_{\mathfrak p'} \ar[r] &
(R_\mathfrak m^\wedge)_{\mathfrak p'}^\wedge
\\
R_\mathfrak m \ar[u] \ar[r] & R_\mathfrak p \ar[u] \ar[r] &
R_\mathfrak p^\wedge \ar[u]
}
$$
依假设，环同态 $R_\mathfrak m \to R_\mathfrak m^\wedge$ 的纤维
具有 $P$。由命题 \ref{more-algebra-proposition-Noetherian-complete-G-ring}，
$(R_\mathfrak m^\wedge)_{\mathfrak p'} \to
(R_\mathfrak m^\wedge)_{\mathfrak p'}^\wedge$ 正则。
局部化 $R_\mathfrak m^\wedge \to
(R_\mathfrak m^\wedge)_{\mathfrak p'}$ 正则。因此，由引理
\ref{more-algebra-lemma-regular-composition}，
$R_\mathfrak m^\wedge \to
(R_\mathfrak m^\wedge)_{\mathfrak p'}^\wedge$ 正则。
于是由 (C)，$R_\mathfrak m \to
(R_\mathfrak m^\wedge)_{\mathfrak p'}^\wedge$ 的纤维具有 $P$。
由于 $R_\mathfrak m \to
(R_\mathfrak m^\wedge)_{\mathfrak p'}^\wedge$ 分解经过局部化
$R_\mathfrak p$，同态 $R_\mathfrak p \to
(R_\mathfrak m^\wedge)_{\mathfrak p'}^\wedge$ 的纤维也具有 $P$。
因此，可以应用 (D) 得到
$R_\mathfrak p \to R_\mathfrak p^\wedge$ 的纤维具有 $P$。
\end{proof}

\begin{proposition}
\label{more-algebra-proposition-finite-type-over-P-ring}
设 $R$ 为 $P$-环，其中 $P$ 满足 (A)、(B)、(C) 和 (D)。
若 $R \to S$ 本质有限型，则 $S$ 是 $P$-环。
\end{proposition}

\begin{proof}
由于 $P$-环性质是局部环的性质，显然 $P$-环的局部化仍为 $P$-环。
反之，若在每个素理想处的局部化都是 $P$-环，则该环是 $P$-环。
因此，只需对每个有限型 $R$-代数 $S$ 及 $S$ 的每个素理想
$\mathfrak q$ 证明 $S_\mathfrak q$ 是 $P$-环。将 $S$ 写成
$R[x_1, \ldots, x_n]$ 的商，由引理
\ref{more-algebra-lemma-P-ring-goes-up-quasi-finite} 可知，只需证明
$R[x_1, \ldots, x_n]$ 是 $P$-环。对 $n$ 归纳，只需证明
$R[x]$ 是 $P$-环。设 $\mathfrak q \subset R[x]$ 为极大理想。
由引理 \ref{more-algebra-lemma-check-P-ring-maximal-ideals}，只需证明
$$
R[x]_\mathfrak q \longrightarrow R[x]_\mathfrak q^\wedge
$$
的纤维具有 $P$。若 $\mathfrak q$ 位于 $\mathfrak p \subset R$ 上方，
则可以用 $R_\mathfrak p$ 代替 $R$。因此，可以假设 $R$ 是极大理想为
$\mathfrak m$ 的 Noether 局部 $P$-环，并且
$\mathfrak q \subset R[x]$ 位于 $\mathfrak m$ 上方。注意，
存在唯一的素理想 $\mathfrak q' \subset R^\wedge[x]$
位于 $\mathfrak q$ 上方。考虑交换图
$$
\xymatrix{
R[x]_\mathfrak q^\wedge \ar[r] &
(R^\wedge[x]_{\mathfrak q'})^\wedge \\
R[x]_\mathfrak q \ar[r] \ar[u] & R^\wedge[x]_{\mathfrak q'} \ar[u]
}
$$
由于 $R$ 是 $P$-环，$R[x] \to R^\wedge[x]$ 的纤维具有 $P$：
它们是 $R \to R^\wedge$ 的纤维沿有限生成域扩张的基变换，
故可应用 (A)。于是，例如由引理 \ref{more-algebra-lemma-P-local}，
下方水平箭头的纤维具有 $P$。右侧竖直箭头正则，因为
$R^\wedge$ 是 G-环（命题 \ref{more-algebra-proposition-Noetherian-complete-G-ring}
和 \ref{more-algebra-proposition-finite-type-over-G-ring}）。因此，由 (C)，复合
$R[x]_\mathfrak q \to (R^\wedge[x]_{\mathfrak q'})^\wedge$
的纤维具有 $P$。从而由 (D)，左侧竖直箭头的纤维具有 $P$，
证明完成。
\end{proof}

\begin{lemma}
\label{more-algebra-lemma-map-P-ring-to-completion-P}
设 $A$ 为 $P$-环，其中 $P$ 满足 (B) 和 (D)。设 $I \subset A$
为理想，并设 $A^\wedge$ 为 $A$ 关于 $I$ 的完备化。那么
$A \to A^\wedge$ 的纤维具有 $P$。
\end{lemma}

\begin{proof}
由 Algebra，引理 \ref{algebra-lemma-completion-flat}，环同态
$A \to A^\wedge$ 平坦。由 Algebra，引理
\ref{algebra-lemma-completion-Noetherian-Noetherian}，环 $A^\wedge$
是 Noether 环。因此，只需检验引理 \ref{more-algebra-lemma-P-local}
的第三个条件。设 $\mathfrak m' \subset A^\wedge$ 为位于
$\mathfrak m \subset A$ 上方的极大理想。由 Algebra，引理
\ref{algebra-lemma-radical-completion}，有
$IA^\wedge \subset \mathfrak m'$。由于 $A^\wedge/IA^\wedge = A/I$，
可知 $I \subset \mathfrak m$、
$\mathfrak m/I = \mathfrak m'/IA^\wedge$，并且
$A/\mathfrak m = A^\wedge/\mathfrak m'$。由于
$A^\wedge/\mathfrak m'$ 是域，可知 $\mathfrak m$ 也是极大理想。
于是，$A_\mathfrak m \to A^\wedge_{\mathfrak m'}$ 是 Noether
局部环之间的平坦局部同态，它诱导剩余域的同构，并满足
$\mathfrak m A^\wedge_{\mathfrak m'} =
\mathfrak m'A^\wedge_{\mathfrak m'}$。因此，由引理
\ref{more-algebra-lemma-flat-unramified}，它在完备局部环上诱导同构。
设 $(A_\mathfrak m)^\wedge$ 为 $A_\mathfrak m$ 关于其极大理想的完备化。
环同态
$$
(A^\wedge)_{\mathfrak m'} \to
((A^\wedge)_{\mathfrak m'})^\wedge = (A_\mathfrak m)^\wedge
$$
忠实平坦（Algebra，引理
\ref{algebra-lemma-completion-faithfully-flat}）。因此，可以对环同态
$$
A_\mathfrak m \to (A^\wedge)_{\mathfrak m'} \to (A_\mathfrak m)^\wedge
$$
应用 (D) 得到结论，因为 $A$ 是 $P$-环，所以
$A_\mathfrak m \to (A_\mathfrak m)^\wedge$ 的纤维具有 $P$。
\end{proof}

\begin{lemma}
\label{more-algebra-lemma-henselization-pair-P-ring}
\begin{slogan}
环的 Hensel 化继承形式纤维的良好性质
\end{slogan}
设 $A$ 为 $P$-环，其中 $P$ 满足 (B)、(C)、(D) 和 (E)。
设 $I \subset A$ 为理想。设 $(A^h, I^h)$ 为偶对 $(A, I)$
的 Hensel 化；见引理 \ref{more-algebra-lemma-henselization}。那么 $A^h$ 是 $P$-环。
\end{lemma}

\begin{proof}
设 $\mathfrak m^h \subset A^h$ 为极大理想。由引理
\ref{more-algebra-lemma-check-P-ring-maximal-ideals}，我们必须证明
$A^h_{\mathfrak m^h} \to (A^h_{\mathfrak m^h})^\wedge$
的纤维具有 $P$。设 $\mathfrak m$ 为 $\mathfrak m^h$ 在 $A$
中的逆像。注意，$I^h \subset \mathfrak m^h$，从而
$I \subset \mathfrak m$，因为 $(A^h, I^h)$ 是 Hensel 偶对。
回顾：$A^h$ 是 Noether 环，$I^h = IA^h$，并且
$A \to A^h$ 在 $I$-adic 完备化上诱导同构；见引理
\ref{more-algebra-lemma-henselization-Noetherian-pair}。于是，Noether 局部环之间的局部同态
$$
A_\mathfrak m \to A^h_{\mathfrak m^h}
$$
由引理 \ref{more-algebra-lemma-flat-unramified} 在极大理想完备化上诱导同构
（略去细节）。设 $\mathfrak q^h$ 为 $A^h_{\mathfrak m^h}$
中位于 $\mathfrak q \subset A_\mathfrak m$ 上方的素理想。
置 $\mathfrak q_1 = \mathfrak q^h$，并令
$\mathfrak q_2, \ldots, \mathfrak q_t$ 为 $A^h$ 中位于
$\mathfrak q$ 上方的其余素理想，于是
$A^h \otimes_A \kappa(\mathfrak q) =
\prod\nolimits_{i = 1, \ldots, t} \kappa(\mathfrak q_i)$；见引理
\ref{more-algebra-lemma-filtered-colimit-etale-noetherian-fibres}。利用上面讨论过的
$(A^h)_{\mathfrak m^h}^\wedge = (A_\mathfrak m)^\wedge$，可得
$$
\prod\nolimits_{i = 1, \ldots, t}
(A^h_{\mathfrak m^h})^\wedge \otimes_{A^h_{\mathfrak m^h}}
\kappa(\mathfrak q_i) =
(A^h_{\mathfrak m^h})^\wedge \otimes_{A^h_{\mathfrak m^h}}
(A^h_{\mathfrak m^h} \otimes_{A_{\mathfrak m}} \kappa(\mathfrak q)) =
(A_{\mathfrak m})^\wedge \otimes_{A_{\mathfrak m}} \kappa(\mathfrak q)
$$
因此，考察局部环并使用 (B)，可知
$$
\kappa(\mathfrak q) \longrightarrow
(A^h_{\mathfrak m^h})^\wedge \otimes_{A^h_{\mathfrak m^h}}
\kappa(\mathfrak q^h)
$$
具有 $P$，因为依关于 $A$ 的假设，
$\kappa(\mathfrak q) \to
(A_\mathfrak m)^\wedge \otimes_{A_\mathfrak m} \kappa(\mathfrak q)$
也具有该性质。由于 $\kappa(\mathfrak q^h)/\kappa(\mathfrak q)$
是可分代数扩张，由 (E) 可知，如所需地
$\kappa(\mathfrak q^h) \to
(A^h_{\mathfrak m^h})^\wedge \otimes_{A^h_{\mathfrak m^h}}
\kappa(\mathfrak q^h)$ 具有 $P$。
\end{proof}

\begin{lemma}
\label{more-algebra-lemma-henselization-P-ring}
设 $R$ 为 Noether 局部 $P$-环，其中 $P$ 满足 (B)、(C)、(D) 和 (E)。
那么，Hensel 化 $R^h$ 和严格 Hensel 化 $R^{sh}$ 都是 $P$-环。
\end{lemma}

\begin{proof}
对于 Hensel 化，我们已在引理 \ref{more-algebra-lemma-henselization-pair-P-ring}
中证明了结论。要对严格 Hensel 化证明结论，由引理
\ref{more-algebra-lemma-check-P-ring-maximal-ideals}，只需证明 $R^{sh}$
的形式纤维具有 $P$。设 $\mathfrak r \subset R^{sh}$ 为素理想，
并置 $\mathfrak p = R \cap \mathfrak r$。置
$\mathfrak r_1 = \mathfrak r$，并令
$\mathfrak r_2, \ldots, \mathfrak r_s$ 为 $R^{sh}$ 中位于
$\mathfrak p$ 上方的其余素理想，于是
$R^{sh} \otimes_R \kappa(\mathfrak p) =
\prod\nolimits_{i = 1, \ldots, s} \kappa(\mathfrak r_i)$；见引理
\ref{more-algebra-lemma-fibres-henselization}。于是可得
$$
\prod\nolimits_{i = 1, \ldots, t}
(R^{sh})^\wedge \otimes_{R^{sh}} \kappa(\mathfrak r_i) =
(R^{sh})^\wedge \otimes_{R^{sh}} (R^{sh} \otimes_R \kappa(\mathfrak p)) =
(R^{sh})^\wedge \otimes_R \kappa(\mathfrak p)
$$
注意，由引理 \ref{more-algebra-lemma-henselization-noetherian}，
$R^\wedge \to (R^{sh})^\wedge$ 对
$\mathfrak m_{(R^{sh})^\wedge}$-adic 拓扑形式光滑。因此，由命题
\ref{more-algebra-proposition-fs-regular}，$R^\wedge \to (R^{sh})^\wedge$ 正则。
由 (C) 以及关于 $R$ 的假设，可知性质 $P$ 对
$\kappa(\mathfrak p) \to (R^{sh})^\wedge \otimes_R \kappa(\mathfrak p)$
成立。使用性质 (B)、上面的分解并考察局部环，可知性质 $P$ 对
$\kappa(\mathfrak p) \to
(R^{sh})^\wedge \otimes_{R^{sh}} \kappa(\mathfrak r)$ 成立。
由于 $\kappa(\mathfrak r)/\kappa(\mathfrak p)$ 是可分代数扩张，
由 (E)，可知 $P$ 对
$\kappa(\mathfrak r) \to
(R^{sh})^\wedge \otimes_{R^{sh}} \kappa(\mathfrak r)$ 成立。
\end{proof}

\begin{lemma}
\label{more-algebra-lemma-formal-fibres-reduced}
性质 (A)、(B)、(C)、(D) 和 (E) 对
$P(k \to R) =$“$R$ 在 $k$ 上几何既约”成立。
\end{lemma}

\begin{proof}
(A) 由几何既约代数的定义得到（Algebra，定义
\ref{algebra-definition-geometrically-reduced}）。(B) 也成立：
一个环既约，当且仅当它的所有局部环既约。
(C) 由引理 \ref{more-algebra-lemma-reduced-goes-up} 得到。
(D) 由 Algebra，引理 \ref{algebra-lemma-descent-reduced} 得到。
(E) 由 Algebra，引理
\ref{algebra-lemma-geometrically-reduced-over-separable-algebraic} 得到。
\end{proof}

\begin{lemma}
\label{more-algebra-lemma-formal-fibres-normal}
性质 (A)、(B)、(C)、(D) 和 (E) 对
$P(k \to R) =$“$R$ 在 $k$ 上几何正规”成立。
\end{lemma}

\begin{proof}
(A) 由几何正规代数的定义得到（Algebra，定义
\ref{algebra-definition-geometrically-normal}）。(B) 也成立：
一个环正规，当且仅当它的所有局部环正规。
(C) 由引理 \ref{more-algebra-lemma-normal-goes-up} 得到。
(D) 由 Algebra，引理 \ref{algebra-lemma-descent-normal} 得到。
(E) 由 Algebra，引理
\ref{algebra-lemma-geometrically-normal-over-separable-algebraic} 得到。
\end{proof}

\begin{lemma}
\label{more-algebra-lemma-formal-fibres-Sk}
固定 $n \geq 1$。性质 (A)、(B)、(C)、(D) 和 (E) 对
$P(k \to R) =$“$R$ 满足 $(S_n)$”成立。
\end{lemma}

\begin{proof}
设 $k \to R$ 为环同态，其中 $k$ 为域，$R$ 为 Noether 环。
设 $k'/k$ 为有限生成域扩张。由 Algebra，引理
\ref{algebra-lemma-tensor-fields-CM}，环同态
$R \to R \otimes_k k'$ 的纤维是 Cohen--Macaulay 的。
因此，可以对环同态 $R \to R \otimes_k k'$ 应用 Algebra，引理
\ref{algebra-lemma-Sk-goes-up}，得到：若 $R$ 满足 $(S_n)$，
则 $R \otimes_k k'$ 也如此。这证明了 (A)。
(B) 也成立：Noether 环满足 $(S_n)$，当且仅当其所有局部环都满足
$(S_n)$。(C) 由 Algebra，引理
\ref{algebra-lemma-Sk-goes-up} 得到，因为正则同态的纤维正则，
因而尤其是 Cohen--Macaulay 的。
(D) 由 Algebra，引理 \ref{algebra-lemma-descent-Sk} 得到。
(E) 是立即的，因为该条件不涉及基域。
\end{proof}

\begin{lemma}
\label{more-algebra-lemma-formal-fibres-CM}
性质 (A)、(B)、(C)、(D) 和 (E) 对
$P(k \to R) =$“$R$ 是 Cohen--Macaulay 环”成立。
\end{lemma}

\begin{proof}
这立即由引理 \ref{more-algebra-lemma-formal-fibres-Sk} 以及以下事实得到：
Noether 环是 Cohen--Macaulay 环，当且仅当它对所有 $n$
都满足条件 $(S_n)$。
\end{proof}

\begin{lemma}
\label{more-algebra-lemma-formal-fibres-Rk}
固定 $n \geq 0$。性质 (A)、(B)、(C)、(D) 和 (E) 对
$P(k \to R) =$“对所有有限扩张 $k'/k$，
$R \otimes_k k'$ 都满足 $(R_n)$”成立。
\end{lemma}

\begin{proof}
设 $k \to R$ 为环同态，其中 $k$ 为域，$R$ 为 Noether 环。
假设 $P(k \to R)$ 成立。设 $K/k$ 为有限生成域扩张。
由 Algebra，引理 \ref{algebra-lemma-make-separable}，可以找到交换图
$$
\xymatrix{
K \ar[r] & K' \\
k \ar[u] \ar[r] & k' \ar[u]
}
$$
其中 $k'/k$、$K'/K$ 是有限纯不可分域扩张，而 $K'/k'$ 可分。
由 Algebra，引理 \ref{algebra-lemma-localization-smooth-separable}，
存在光滑 $k'$-代数 $B$，使得 $K'$ 是 $B$ 的分式域。
现在可以作如下论证：
步骤 1：$R \otimes_k k'$ 满足 $(R_n)$，因为我们假设
$P$ 对 $k \to R$ 成立。
步骤 2：$R \otimes_k k' \to R \otimes_k k' \otimes_{k'} B$
是光滑环同态（Algebra，引理 \ref{algebra-lemma-base-change-smooth}），
由 Algebra，引理 \ref{algebra-lemma-Rk-goes-up} 可知
$R \otimes_k k' \otimes_{k'} B$ 满足 $(R_n)$
（并使用 Algebra，引理
\ref{algebra-lemma-characterize-smooth-over-field} 检验其假设成立）。
步骤 3：$R \otimes_k k' \otimes_{k'} K' = R \otimes_k K'$
满足 $(R_n)$，因为它是一个满足 $(R_n)$ 的环的局部化。
步骤 4：最后，由 $(R_n)$ 沿忠实平坦环同态
$K \otimes_k R \to K' \otimes_k R$ 的下降
（Algebra，引理 \ref{algebra-lemma-descent-Rk}），
$R \otimes_k K$ 满足 $(R_n)$。这证明了 (A)。
(B) 也成立：Noether 环满足 $(R_n)$，当且仅当其所有局部环都满足
$(R_n)$。(C) 由 Algebra，引理 \ref{algebra-lemma-Rk-goes-up}
得到，因为正则同态的纤维正则（略去一个小细节）。
(D) 由 Algebra，引理 \ref{algebra-lemma-descent-Rk} 得到
（略去一个小细节）。

\medskip\noindent
对于 (E)，设 $l/k$ 为域的可分代数扩张，并设 $l \to R$ 为环同态，
其中 $R$ 为 Noether 环。假设 $k \to R$ 具有 $P$。
我们必须证明 $l \to R$ 具有 $P$。设 $l'/l$ 为有限扩张。
首先注意，存在有限子扩张 $l/m/k$ 和有限扩张 $m'/m$，使得
$l' = l \otimes_m m'$。那么 $R \otimes_l l' = R \otimes_m m'$。
因此，只需证明 $m \to R$ 具有性质 $P$；换言之，可以假设
$l/k$ 有限。若 $l/k$ 有限，则 $l'/k$ 有限，并且有
$$
l' \otimes_l R = (l' \otimes_k R) \otimes_{l \otimes_k l} l
$$
它是具有性质 $(R_n)$ 的 Noether 环 $l' \otimes_k R$ 的一个局部化
（由 Algebra，引理 \ref{algebra-lemma-separable-algebraic-diagonal}），
其中该环具有此性质是因为我们假设 $P$ 对 $k \to R$ 成立。
这证明了 $l' \otimes_l R$ 如所需地具有性质 $(R_n)$。
\end{proof}








\section{优环}
\label{more-algebra-section-excellent}

\noindent
本节讨论 Grothendieck 的优环概念。关于 G-环、J-2 环和普遍链环的定义，
我们分别参见定义 \ref{more-algebra-definition-G-ring}、定义 \ref{more-algebra-definition-J}
以及 Algebra，定义 \ref{algebra-definition-universally-catenary}。

\begin{definition}
\label{more-algebra-definition-excellent}
设 $R$ 为环。
\begin{enumerate}
\item 若 $R$ 是 Noether 环、G-环且为 J-2，则称 $R$ 是{\it 拟优环}。
\item 若 $R$ 是拟优环且为普遍链环，则称 $R$ 是{\it 优环}。
\end{enumerate}
\end{definition}

\noindent
因此，若 Noether 环具有几何正则的形式纤维，并且其上任意有限型代数的
奇异集都闭，则它是拟优环。要使这样的环成为优环，还要求（局部地）
存在良好的维数函数。稍后我们将看到（第
\ref{more-algebra-section-formally-catenary} 节），普遍链性可以表述为对 $R$
的极大理想 $\mathfrak m$，同态
$R_\mathfrak m \to R_\mathfrak m^\wedge$ 所满足的条件。

\begin{lemma}
\label{more-algebra-lemma-finite-type-over-excellent}
（拟）优环上的有限型环的任意局部化都是（拟）优环。
\end{lemma}

\begin{proof}
对于有限型代数，这由 J-2 和普遍链性的定义得到。对于 G-环，见命题
\ref{more-algebra-proposition-finite-type-over-G-ring}。我们略去局部化保持（拟）优性的证明。
\end{proof}

\begin{proposition}
\label{more-algebra-proposition-ubiquity-excellent}
下列各类环都是优环：
\begin{enumerate}
\item 域；
\item Noether 完备局部环；
\item $\mathbf{Z}$；
\item 分式域特征为零的 Dedekind 整环；
\item 上述任一类环的有限型环扩张。
\end{enumerate}
\end{proposition}

\begin{proof}
由命题 \ref{more-algebra-proposition-ubiquity-G-ring} 和
\ref{more-algebra-proposition-ubiquity-J-2} 可知，这些环是 G-环且具有 J-2。
任意 Cohen--Macaulay 环都是普遍链环；见 Algebra，引理
\ref{algebra-lemma-CM-ring-catenary}。特别地，域、Dedekind 环以及更一般的
正则环都是普遍链环。借助 Cohen 结构定理，可知完备局部环是普遍链环；
见 Algebra，注
\ref{algebra-remark-Noetherian-complete-local-ring-universally-catenary}。
\end{proof}

\noindent
上面发展的材料对 Nagata 环有一些推论。

\begin{lemma}
\label{more-algebra-lemma-Nagata-local-ring}
设 $(A, \mathfrak m)$ 为 Noether 局部环。下列条件等价：
\begin{enumerate}
\item $A$ 是 Nagata 环；以及
\item $A$ 的形式纤维几何既约。
\end{enumerate}
\end{lemma}

\begin{proof}
假设 (2)。由 Algebra，引理
\ref{algebra-lemma-local-nagata-and-analytically-unramified}，
必须证明：若 $A \to B$ 有限、$B$ 为整环且
$\mathfrak m' \subset B$ 为极大理想，则
$B_{\mathfrak m'}$ 解析无分歧。结合引理
\ref{more-algebra-lemma-formal-fibres-reduced}、
\ref{more-algebra-lemma-check-P-ring-maximal-ideals} 以及命题
\ref{more-algebra-proposition-finite-type-over-P-ring}，可知
$B_{\mathfrak m'}$ 的形式纤维几何既约。特别地，若 $L$ 为 $B$
的分式域，则 $B_{\mathfrak m'}^\wedge \otimes_B L$ 既约。
由此可知 $B_{\mathfrak m'}^\wedge$ 既约；换言之，
$B_{\mathfrak m'}$ 解析无分歧。

\medskip\noindent
假设 (1)。设 $\mathfrak q \subset A$ 为素理想，并设
$K/\kappa(\mathfrak q)$ 为有限扩张。我们必须证明
$A^\wedge \otimes_A K$ 既约。取局部子环
$A/\mathfrak q \subset B \subset K$，使得它在 $A$ 上有限且分式域为 $K$。
为了构造 $B$，选取生成 $K$ 在 $\kappa(\mathfrak q)$ 上扩张的元素
$x_1, \ldots, x_n \in K$，它们满足首一多项式
$P_i(T) = T^{d_i} + a_{i, 1} T^{d_i - 1} + \ldots + a_{i, d_i} = 0$，
其中 $a_{i, j} \in \mathfrak m$。然后令 $B$ 为 $K$ 中由
$x_1, \ldots, x_n$ 生成的 $A$-子代数。（更多细节见 Algebra，引理
\ref{algebra-lemma-local-nagata-and-analytically-unramified} 的证明。）那么
$$
A^\wedge \otimes_A K =
(A^\wedge \otimes_A B)_\mathfrak q =
B^\wedge_\mathfrak q
$$
由 Algebra，引理
\ref{algebra-lemma-local-nagata-and-analytically-unramified}，
$B^\wedge$ 既约，故证明完成。
\end{proof}

\begin{lemma}
\label{more-algebra-lemma-quasi-excellent-nagata}
拟优环是 Nagata 环。
\end{lemma}

\begin{proof}
设 $R$ 为拟优环。利用有限型 $R$-代数仍为拟优环这一事实
（引理 \ref{more-algebra-lemma-finite-type-over-excellent}），由 Algebra，引理
\ref{algebra-lemma-check-universally-japanese}，只需证明任意拟优整环都是 N-1。
应用 Algebra，引理 \ref{algebra-lemma-characterize-N-1}
（并使用拟优环为 J-2 这一事实），问题化约为证明拟优局部整环 $R$ 是 N-1。
由于 $R \to R^\wedge$ 正则，由引理 \ref{more-algebra-lemma-reduced-goes-up}，
$R^\wedge$ 既约。换言之，$R$ 解析无分歧。因此，由 Algebra，引理
\ref{algebra-lemma-analytically-unramified-easy}，$R$ 是 N-1。
\end{proof}

\begin{lemma}
\label{more-algebra-lemma-completion-normal-local-ring}
设 $(A, \mathfrak m)$ 为 Noether 局部环。若 $A$ 正规，且 $A$
的形式纤维正规（例如，若 $A$ 为优环或拟优环），则 $A^\wedge$ 正规。
\end{lemma}

\begin{proof}
这立即由 Algebra，引理
\ref{algebra-lemma-normal-goes-up-noetherian} 得到。
\end{proof}





\section{模的 Abel 范畴}
\label{more-algebra-section-abelian-categories-modules}

\noindent
设 $R$ 为环。$R$-模范畴 $\text{Mod}_R$ 是 Abel 范畴。
以下是 $\text{Mod}_R$ 的一些 Abel 子范畴的例子（我们使用 Homology，
定义 \ref{homology-definition-serre-subcategory} 以及 Homology，引理
\ref{homology-lemma-characterize-serre-subcategory} 和
\ref{homology-lemma-characterize-weak-serre-subcategory} 中引入的术语）：
\begin{enumerate}
\item 相干 $R$-模范畴是 $\text{Mod}_R$ 的弱 Serre 子范畴。
这由 Algebra，引理 \ref{algebra-lemma-coherent} 得到。
\item 设 $S \subset R$ 为乘法子集。由所有满足下述条件的 $R$-模 $M$
构成的满子范畴是 $\text{Mod}_R$ 的 Serre 子范畴：对每个 $s \in S$，
相应乘法在 $M$ 上是同构。这由 Algebra，引理
\ref{algebra-lemma-localization-and-modules} 得到。
\item 设 $I \subset R$ 为有限生成理想。$I$-幂挠模的满子范畴是
$\text{Mod}_R$ 的 Serre 子范畴。见引理 \ref{more-algebra-lemma-I-power-torsion}。
\item 在一些文献中，{\it 挠模}定义为满足下述条件的模 $M$：
对每个 $x \in M$，存在非零因子 $f \in R$，使得 $fx = 0$。
挠模的满子范畴是 $\text{Mod}_R$ 的 Serre 子范畴。
\item 若 $R$ 不是 Noether 环，则有限生成 $R$-模范畴
$\text{Mod}^{fg}_R$ {\bf 不是} Abel 范畴。事实上，若
$I \subset R$ 是非有限生成理想，则映射 $R \to R/I$
在 $\text{Mod}^{fg}_R$ 中没有核。
\item 若 $R$ 是 Noether 环，则相干 $R$-模等同于有限生成
（即有限）$R$-模；见 Algebra，引理
\ref{algebra-lemma-Noetherian-coherent}、
\ref{algebra-lemma-coherent-ring} 和
\ref{algebra-lemma-Noetherian-finite-type-is-finite-presentation}。
因此，由上面的 (1)，$\text{Mod}^{fg}_R$ 是 Abel 范畴；事实上，
在这种情况下，范畴 $\text{Mod}_R^{fg}$ 是 $\text{Mod}_R$
的（强）Serre 子范畴。
\end{enumerate}




\section{内射 Abel 群}
\label{more-algebra-section-abelian-groups}

\noindent
本节证明 Abel 群范畴有足够多的内射对象。回顾：Abel 群 $M$
是{\it 可除的}，当且仅当对每个 $x \in M$ 和每个
$n \in \mathbf{N}$，存在 $y \in M$，使得 $n y = x$。

\begin{lemma}
\label{more-algebra-lemma-injective-abelian}
Abel 群 $J$ 是 Abel 群范畴中的内射对象，当且仅当 $J$ 可除。
\end{lemma}

\begin{proof}
假设 $J$ 不可除。那么，存在 $x \in J$ 和 $n \in \mathbf{N}$，
使得不存在满足 $n y = x$ 的 $y \in J$。于是，同态
$\mathbf{Z} \to J$，$m \mapsto mx$ 不能延拓到
$\frac{1}{n}\mathbf{Z} \supset \mathbf{Z}$。因此 $J$ 不是内射对象。

\medskip\noindent
设 $A \subset B$ 为 Abel 群。假设 $J$ 为可除 Abel 群。
设 $\varphi : A \to J$ 为同态。考虑所有同态
$\varphi' : A' \to J$ 的集合，其中 $A \subset A' \subset B$
且 $\varphi'|_A = \varphi$。定义
$(A', \varphi') \geq (A'', \varphi'')$，当且仅当
$A' \supset A''$ 且 $\varphi'|_{A''} = \varphi''$。
若 $(A_i, \varphi_i)_{i \in I}$ 是这类偶对的全序族，
则得到映射 $\bigcup_{i \in I} A_i \to J$：元素 $a \in A_i$
映到 $\varphi_i(a)$。因此，可应用 Zorn 引理。为得出结论，
必须证明：若偶对 $(A', \varphi')$ 极大，则 $A' = B$。
换言之，只需证明：给定任意子群 $A \subset B$、$A \not = B$
以及任意 $\varphi : A \to J$，可以找到
$\varphi' : A' \to J$，其中 $A \subset A' \subset B$，使得
(a) 包含 $A \subset A'$ 是严格的，且 (b) 同态 $\varphi'$
延拓 $\varphi$。

\medskip\noindent
为证明这一点，选取 $x \in B$、$x \not \in A$。
若不存在 $n\in \mathbf{N}$ 使得 $nx \in A$，则
$A \oplus \mathbf{Z} \cong A + \mathbf{Z}x$。于是，可以将
$\varphi$ 延拓到 $A' = A + \mathbf{Z}x$，例如在 $A$ 上使用
$\varphi$，并将 $x$ 映到零。若确实存在 $n \in \mathbf{N}$
使得 $nx \in A$，则取最小的此类整数 $n$。取 $z \in J$，
使得 $nz = \varphi(nx)$。定义同态
$\tilde\varphi : A \oplus \mathbf{Z} \to J$ 为
$(a, m) \mapsto \varphi(a) + mz$。由 $z$ 的选取，
$\tilde \varphi$ 的核包含映射 $A \oplus \mathbf{Z} \to B$，
$(a, m) \mapsto a + mx$ 的核。因此，$\tilde \varphi$
分解经过像 $A' = A + \mathbf{Z}x$，并由此延拓同态 $\varphi$。
\end{proof}

\noindent
可以用这个引理证明每个 Abel 群都能嵌入某个内射 Abel 群。
但这是下一节结果的一个特殊情形。






\section{内射模}
\label{more-algebra-section-injectives-modules}

\noindent
这里给出关于内射模的一些引理。

\begin{definition}
\label{more-algebra-definition-projective}
设 $R$ 为环。$R$-模 $J$ 是{\it 内射的}，当且仅当函子
$\Hom_R(-, J) : \text{Mod}_R \to \text{Mod}_R$ 是正合函子。
\end{definition}

\noindent
对任意 $R$-模 $M$，函子 $\Hom_R(- , M)$ 左正合；见 Algebra，引理
\ref{algebra-lemma-hom-exact}。因此，$J$ 为内射模这一条件实际表示：
给定 $R$-模的单射 $M \to M'$，映射
$\Hom_R(M', J) \to \Hom_R(M, J)$ 满射。

\medskip\noindent
在用 ${Ext}$-模重新表述这一点之前，我们先讨论
$\Ext^1_R(M, N)$ 与 Homology，第 \ref{homology-section-extensions} 节
中的扩张之间的关系。

\begin{lemma}
\label{more-algebra-lemma-relation-ext-ext}
设 $R$ 为环。设 $\mathcal{A}$ 为 $R$-模的 Abel 范畴。
存在典范同构
$\Ext_\mathcal{A}(M, N) = \Ext^1_R(M, N)$，它与 Algebra，引理
\ref{algebra-lemma-long-exact-seq-ext} 和
\ref{algebra-lemma-reverse-long-exact-seq-ext} 的长正合序列以及
Homology，引理 \ref{homology-lemma-six-term-sequence-ext}
的 $6$ 项正合序列相容。
\end{lemma}

\begin{proof}
略。
\end{proof}

\begin{lemma}
\label{more-algebra-lemma-characterize-injective}
设 $R$ 为环。设 $J$ 为 $R$-模。下列条件等价：
\begin{enumerate}
\item $J$ 是内射模；
\item 对每个 $R$-模 $M$，都有 $\Ext^1_R(M, J) = 0$。
\end{enumerate}
\end{lemma}

\begin{proof}
设 $0 \to M'' \to M' \to M \to 0$ 为 $R$-模的短正合序列。
考虑 Algebra，引理 \ref{algebra-lemma-reverse-long-exact-seq-ext}
中的长正合序列
$$
\begin{matrix}
0
\to \Hom_R(M, J)
\to \Hom_R(M', J)
\to \Hom_R(M'', J)
\\
\phantom{0\ }
\to \Ext^1_R(M, J)
\to \Ext^1_R(M', J)
\to \Ext^1_R(M'', J)
\to \ldots
\end{matrix}
$$
由此可知 (2) 推出 (1)。反之，若 $J$ 是内射模，
则由 Homology，引理 \ref{homology-lemma-characterize-injectives}
和引理 \ref{more-algebra-lemma-relation-ext-ext}，该 $\Ext$-群为零。
\end{proof}

\begin{lemma}
\label{more-algebra-lemma-characterize-injective-bis}
设 $R$ 为环。设 $J$ 为 $R$-模。下列条件等价：
\begin{enumerate}
\item $J$ 是内射模；
\item 对每个理想 $I \subset R$，都有 $\Ext^1_R(R/I, J) = 0$；以及
\item 对理想 $I \subset R$ 和模同态 $I \to J$，存在延拓 $R \to J$。
\end{enumerate}
\end{lemma}

\begin{proof}
若 $I \subset R$ 为理想，则短正合序列
$0 \to I \to R \to R/I \to 0$ 给出正合序列
$$
\Hom_R(R, J) \to
\Hom_R(I, J) \to
\Ext^1_R(R/I, J) \to 0
$$
这是由 Algebra，引理 \ref{algebra-lemma-reverse-long-exact-seq-ext}
以及如下事实得到的：$R$ 是投射模，故
$\Ext^1_R(R, J) = 0$（Algebra，引理
\ref{algebra-lemma-characterize-projective}）。因此，(2) 与 (3) 等价。
在本证明中，我们将证明 (1) $\Leftrightarrow$ (3)，这就是著名的
Baer 判据。

\medskip\noindent
假设 (1)。给定 (3) 中的模同态 $I \to J$，由于按定义映射
$\Hom_R(R, J) \to \Hom_R(I, J)$ 满射，便得到延拓 $R \to J$。

\medskip\noindent
假设 (3)。设 $M \subset N$ 为 $R$-模的包含。设
$\varphi : M \to J$ 为同态。我们将证明 $\varphi$ 延拓到 $N$，
这就完成引理的证明。考虑所有同态 $\varphi' : M' \to J$ 的集合，
其中 $M \subset M' \subset N$ 且 $\varphi'|_M = \varphi$。
定义 $(M', \varphi') \geq (M'', \varphi'')$，当且仅当
$M' \supset M''$ 且 $\varphi'|_{M''} = \varphi''$。
若 $(M_i, \varphi_i)_{i \in I}$ 是这类偶对的全序族，
则得到映射 $\bigcup_{i \in I} M_i \to J$：元素 $a \in M_i$
映到 $\varphi_i(a)$。因此，可应用 Zorn 引理。为得出结论，
必须证明：若偶对 $(M', \varphi')$ 极大，则 $M' = N$。
换言之，只需证明：给定任意子群 $M \subset N$、$M \not = N$
以及任意 $\varphi : M \to J$，可以找到
$\varphi' : M' \to J$，其中 $M \subset M' \subset N$，使得
(a) 包含 $M \subset M'$ 是严格的，且 (b) 同态 $\varphi'$
延拓 $\varphi$。

\medskip\noindent
为证明这一点，选取 $x \in N$、$x \not \in M$。令
$I = \{f \in R \mid fx \in M\}$。这是 $R$ 的理想。
定义同态 $\psi : I \to J$ 为 $f \mapsto \varphi(fx)$。
依假设 (3)，可以将其延拓为映射 $\tilde\psi : R \to J$。
由 $I$ 的选取，同态
$M \oplus R \to J$，$(y, f) \mapsto \varphi(y) + \tilde\psi(f)$
的核包含映射 $M \oplus R \to N$，$(y, f) \mapsto y + fx$ 的核。
因此，该同态分解经过像 $M' = M + Rx$，并如所需地延拓给定同态。
\end{proof}

\noindent
本节余下部分证明环 $R$ 上有足够多的内射模。首先注意，
$\mathbf{Q}/\mathbf{Z}$ 是内射 Abel 群。这由引理
\ref{more-algebra-lemma-injective-abelian} 得到。

\begin{definition}
\label{more-algebra-definition-simple-functors}
设 $R$ 为环。
\begin{enumerate}
\item 对任意 $R$ 上的 $R$-模 $M$，记
$M^\vee = \Hom(M, \mathbf{Q}/\mathbf{Z})$，并赋予其自然的
$R$-模结构。我们把{\it $M \mapsto M^\vee$} 看成从 $R$-模范畴
到自身的反变函子。
\item 对任意 $R$-模 $M$，记
$$
F(M) = \bigoplus\nolimits_{m \in M} R[m]
$$
为以 $m \in M$ 对应的元素 $[m]$ 为基的{\it 自由模}。
令 $F(M)\to M$，$\sum f_i [m_i] \mapsto \sum f_i m_i$
为 $R$-模的自然满射。我们把 $M \mapsto (F(M) \to M)$
看成从 $R$-模范畴到 $R$-模箭头范畴的函子。
\end{enumerate}
\end{definition}

\begin{lemma}
\label{more-algebra-lemma-vee-exact}
设 $R$ 为环。函子 $M \mapsto M^\vee$ 正合。
\end{lemma}

\begin{proof}
这是因为由引理 \ref{more-algebra-lemma-injective-abelian}，
$\mathbf{Q}/\mathbf{Z}$ 是内射 Abel 群。
\end{proof}

\noindent
存在由求值给出的典范映射 $ev : M \to (M^\vee)^\vee$：
给定 $x \in M$，令
$ev(x) \in (M^\vee)^\vee = \Hom(M^\vee, \mathbf{Q}/\mathbf{Z})$
为映射 $\varphi \mapsto \varphi(x)$。

\begin{lemma}
\label{more-algebra-lemma-ev-injective}
对任意 $R$-模 $M$，求值映射 $ev : M \to (M^\vee)^\vee$ 是单射。
\end{lemma}

\begin{proof}
可以利用 $\mathbf{Q}/\mathbf{Z}$ 是内射 Abel 群这一事实来检验。
具体地，若 $x \in M$ 非零，则令 $M' \subset M$ 为它生成的循环群。
存在非零映射 $M' \to \mathbf{Q}/\mathbf{Z}$，它必不消去 $x$。
该映射延拓为 $\varphi : M \to \mathbf{Q}/\mathbf{Z}$，于是
$ev(x)(\varphi) = \varphi(x) \not = 0$。
\end{proof}

\noindent
$R$-模的典范满射 $F(M) \to M$ 经过上述构造变为 $R$-模的典范单射
$$
(M^\vee)^\vee \longrightarrow (F(M^\vee))^\vee.
$$
置 $J(M) = (F(M^\vee))^\vee$。$ev$ 与这个显示映射的复合给出
$M \to J(M)$，并且该构造关于 $M$ 具有函子性。

\begin{lemma}
\label{more-algebra-lemma-JM-injective}
设 $R$ 为环。对每个 $R$-模 $M$，$R$-模 $J(M)$ 是内射模。
\end{lemma}

\begin{proof}
注意，作为 $R$-模，$J(M) \cong \prod_{\varphi \in M^\vee} R^\vee$。
由于内射模的积仍为内射模，只需证明 $R^\vee$ 是内射模。
为此，使用等式
$$
\Hom_R(N, R^\vee) =
\Hom_R(N, \Hom_{\mathbf{Z}}(R, \mathbf{Q}/\mathbf{Z})) =
N^\vee
$$
以及由引理 \ref{more-algebra-lemma-vee-exact} 得到的事实：
$(-)^\vee$ 是正合函子。
\end{proof}

\begin{lemma}
\label{more-algebra-lemma-injectives-modules}
设 $R$ 为环。上述构造定义了从 $R$-模范畴到 $R$-模箭头范畴的协变函子
$M \mapsto (M \to J(M))$，使得对每个模 $M$，输出
$M \to J(M)$ 都是将 $M$ 嵌入内射 $R$-模 $J(M)$ 的单射。
\end{lemma}

\begin{proof}
由上文立即得到。
\end{proof}

\noindent
特别地，对任意 $R$-模映射 $M \to N$，都有相伴的同态
$J(M) \to J(N)$，使下图交换：
$$
\xymatrix{
M \ar[d] \ar[r] & N \ar[d] \\
J(M) \ar[r] & J(N) }
$$
这正是我们一般希望得到的构造。在 Homology，第
\ref{homology-section-injectives} 节中，我们引入了表述这一点的术语：
这意味着 $R$-模范畴具有函子性内射嵌入。










\section{模的导出范畴}
\label{more-algebra-section-derived-modules}

\noindent
本节给出环上模的导出范畴的一些一般性结论。

\medskip\noindent
设 $A$ 为环。以 $\text{Mod}_A$ 表示 $A$-模范畴。
我们用符号 $K(A)$ 表示 $A$-模复形的同伦范畴；换言之，
作为范畴，置 $K(A) = K(\text{Mod}_A)$；见 Derived Categories，第
\ref{derived-section-homotopy} 节。有界版本为
$K^+(A)$、$K^-(A)$ 和 $K^b(A)$。按照 Derived Categories，第
\ref{derived-section-homotopy-triangulated} 节，我们把 $K(A)$
看成三角范畴。$A$ 的{\it 导出范畴}记作 $D(A)$，它是通过在
$K(A)$ 中逆转拟同构而得到的范畴；换言之，置
$D(A) = D(\text{Mod}_A)$；见 Derived Categories，第
\ref{derived-section-derived-categories} 节\footnote{另见 Injectives，注
\ref{injectives-remark-existence-D}。}。有界版本为
$D^+(A)$、$D^-(A)$ 和 $D^b(A)$。

\medskip\noindent
设 $A$ 为环。$A$-模范畴有积，并且积正合。由引理
\ref{more-algebra-lemma-injectives-modules}，$A$-模范畴有足够多的内射对象。
因此，每个 $A$-模复形都拟同构于某个 K-内射复形
（Derived Categories，引理
\ref{derived-lemma-enough-K-injectives-Ab4-star}）。由此，$D(A)$
有可数积（Derived Categories，引理 \ref{derived-lemma-products}），
事实上还有任意积（Injectives，引理
\ref{injectives-lemma-derived-products}）。这蕴含 $D(A)$
中对象的每个逆系统都有导出极限（在同构意义下良定义）；见
Derived Categories，第 \ref{derived-section-derived-limit} 节。

\begin{lemma}
\label{more-algebra-lemma-K-injective-flat}
设 $R \to S$ 为平坦环同态。若 $I^\bullet$ 是 $S$-模的 K-内射复形，
则作为 $R$-模复形，$I^\bullet$ 也是 K-内射的。
\end{lemma}

\begin{proof}
这是因为由 Algebra，引理 \ref{algebra-lemma-adjoint-tensor-restrict}
以及张量积以 $S$ 是正合的这一事实，有
$\Hom_{K(R)}(M^\bullet, I^\bullet) =
\Hom_{K(S)}(M^\bullet \otimes_R S, I^\bullet)$。
\end{proof}

\begin{lemma}
\label{more-algebra-lemma-K-injective-epimorphism}
设 $R \to S$ 为环的满态射。设 $I^\bullet$ 为 $S$-模复形。
若作为 $R$-模复形，$I^\bullet$ 是 K-内射的，则作为 $S$-模复形，
$I^\bullet$ 也是 K-内射的。
\end{lemma}

\begin{proof}
这是因为对任意 $S$-模复形 $N^\bullet$，有
$\Hom_{K(R)}(N^\bullet, I^\bullet) =
\Hom_{K(S)}(N^\bullet, I^\bullet)$；见 Algebra，引理
\ref{algebra-lemma-epimorphism-modules}。
\end{proof}

\begin{lemma}
\label{more-algebra-lemma-hom-K-injective}
设 $A \to B$ 为环同态。若 $I^\bullet$ 是 $A$-模的 K-内射复形，
则 $\Hom_A(B, I^\bullet)$ 是 $B$-模的 K-内射复形。
\end{lemma}

\begin{proof}
这是因为由 Algebra，引理 \ref{algebra-lemma-adjoint-hom-restrict}，有
$\Hom_{K(B)}(N^\bullet, \Hom_A(B, I^\bullet)) =
\Hom_{K(A)}(N^\bullet, I^\bullet)$。
\end{proof}










\section{Tor 的计算}
\label{more-algebra-section-computing-tor}

\noindent
设 $R$ 为环。以 $D(R)$ 表示 $R$-模 Abel 范畴
$\text{Mod}_R$ 的导出范畴。注意，$\text{Mod}_R$ 有足够多的投射对象，
因为每个自由 $R$-模都是投射模。因此，可以定义从
$\text{Mod}_R$ 到任意 Abel 范畴的任何加性函子的左导出函子。

\medskip\noindent
这特别适用于函子
$ - \otimes_R M : \text{Mod}_R \to \text{Mod}_R$，
其左导出函子是 Tor 函子 $\text{Tor}_i^R(-, M)$；见 Algebra，第
\ref{algebra-section-tor} 节。还有全左导出函子
\begin{equation}
\label{more-algebra-equation-derived-tensor-module}
-\otimes_R^{\mathbf{L}} M :
D^{-}(R)
\longrightarrow
D^{-}(R)
\end{equation}
记作 $-\otimes_R^{\mathbf{L}} M$。它的卫星函子是 Tor 模；换言之，
$$
H^{-p}(N \otimes_R^{\mathbf{L}} M) = \text{Tor}_p^R(N, M).
$$

\medskip\noindent
当考虑与 $R$-代数 $A$ 的张量积时，会出现一种特殊情形。
在此情形下，我们把 $- \otimes_R A$ 看成从 $\text{Mod}_R$
到 $\text{Mod}_A$ 的函子。因此，有全右导出函子
\begin{equation}
\label{more-algebra-equation-derived-tensor-algebra}
-\otimes_R^{\mathbf{L}} A :
D^{-}(R)
\longrightarrow
D^{-}(A)
\end{equation}
记作 $-\otimes_R^{\mathbf{L}} A$。它的卫星函子是 Tor 群；换言之，
$$
H^{-p}(N \otimes_R^{\mathbf{L}} A) = \text{Tor}_p^R(N, A).
$$
特别地，这些 Tor 群自然具有 $A$-模结构。

\medskip\noindent
在接下来的几节中，我们会把本节材料推广到无界复形。






\section{复形的张量积}
\label{more-algebra-section-symmetric-monoidal}

\noindent
设 $R$ 为环。$R$-模复形范畴 $\text{Comp}(R)$ 具有对称幺半结构。
具体地，设有两个 $R$-模复形 $L^\bullet$ 和 $M^\bullet$。
利用 Homology，例 \ref{homology-example-double-complex-as-tensor-product-of}
和 Homology，定义 \ref{homology-definition-associated-simple-complex}，
得到第三个 $R$-模复形
$$
\text{Tot}(L^\bullet \otimes_R M^\bullet)
$$
显然，该构造关于 $L^\bullet$ 和 $M^\bullet$ 都具有函子性。
结合约束取 Homology，注 \ref{homology-remark-triple-complex}
中从三重复形 $K^\bullet \otimes_R L^\bullet \otimes_R M^\bullet$
构造的复形典范同构
$$
\text{Tot}(\text{Tot}(K^\bullet \otimes_R L^\bullet) \otimes_R M^\bullet)
\longrightarrow
\text{Tot}(K^\bullet \otimes_R \text{Tot}(L^\bullet \otimes_R M^\bullet))
$$
交换约束是典范同构
$$
\text{Tot}(L^\bullet \otimes_R M^\bullet) \to
\text{Tot}(M^\bullet \otimes_R L^\bullet)
$$
它在直和项 $L^p \otimes_R M^q$ 上使用符号 $(-1)^{pq}$。
为了看出这是复形映射，对 $x \in L^p$ 和 $y \in M^q$ 计算：
$$
\text{d}(x \otimes y) =
\text{d}_L(x) \otimes y + (-1)^px \otimes \text{d}_M(y)
$$
按我们的规则，右侧映到
$$
(-1)^{(p + 1)q}y \otimes \text{d}_L(x) +
(-1)^{p + p(q + 1)} \text{d}_M(y) \otimes x
$$
另一方面，
$$
\text{d}((-1)^{pq}y \otimes x) =
(-1)^{pq} \text{d}_M(y) \otimes x +
(-1)^{pq + q} y \otimes \text{d}_L(x)
$$
直接观察即可看出这两个表达式如所需地相等。

\begin{lemma}
\label{more-algebra-lemma-symmetric-monoidal-cat-complexes}
设 $R$ 为环。在赋予函子
$(L^\bullet, M^\bullet) \mapsto \text{Tot}(L^\bullet \otimes_R M^\bullet)$
以及上述结合约束和交换约束之后，$R$-模复形范畴
$\text{Comp}(R)$ 是对称幺半范畴。
\end{lemma}

\begin{proof}
略。提示：取在次数 $0$ 处为 $R$、在其他次数处为零的复形作为单位
$\mathbf{1}$，并使用显然的同构
$\text{Tot}(\mathbf{1} \otimes_R M^\bullet) = M^\bullet$ 和
$\text{Tot}(K^\bullet \otimes_R \mathbf{1}) = K^\bullet$。
为证明该引理，需要检验若干图的交换性；见 Categories，定义
\ref{categories-definition-monoidal-category} 和
\ref{categories-definition-symmetric-monoidal-category}。
每种情形中的检验都很直接。
\end{proof}

\begin{lemma}
\label{more-algebra-lemma-derived-tor-homotopy}
设 $R$ 为环。设 $P^\bullet$ 为 $R$-模复形。设
$\alpha, \beta : L^\bullet \to M^\bullet$ 为复形的同伦映射。
那么，$\alpha$ 和 $\beta$ 诱导同伦映射
$$
\text{Tot}(\alpha \otimes \text{id}_P),
\text{Tot}(\beta \otimes \text{id}_P) :
\text{Tot}(L^\bullet \otimes_R P^\bullet)
\longrightarrow
\text{Tot}(M^\bullet \otimes_R P^\bullet).
$$
特别地，构造
$L^\bullet \mapsto \text{Tot}(L^\bullet \otimes_R P^\bullet)$
定义了复形同伦范畴的自函子。
\end{lemma}

\begin{proof}
设对某个由 $h^n : L^n \to M^{n - 1}$ 定义的同伦 $h$，有
$\alpha = \beta + dh + hd$。置
$$
H^n = \bigoplus\nolimits_{a + b = n} h^a \otimes \text{id}_{P^b} :
\bigoplus\nolimits_{a + b = n} L^a \otimes_R P^b
\longrightarrow
\bigoplus\nolimits_{a + b = n} M^{a - 1} \otimes_R P^b
$$
直接计算可知，作为映射
$\text{Tot}(L^\bullet \otimes_R P^\bullet) \to
\text{Tot}(M^\bullet \otimes_R P^\bullet)$，有
$$
\text{Tot}(\alpha \otimes \text{id}_P) =
\text{Tot}(\beta \otimes \text{id}_P) + dH + Hd
$$
\end{proof}

\begin{lemma}
\label{more-algebra-lemma-symmetric-monoidal-homotopy-category}
设 $R$ 为环。在赋予函子
$(L^\bullet, M^\bullet) \mapsto \text{Tot}(L^\bullet \otimes_R M^\bullet)$
以及上述结合约束和交换约束之后，$R$-模复形的同伦范畴 $K(R)$
是对称幺半范畴。
\end{lemma}

\begin{proof}
这由引理 \ref{more-algebra-lemma-symmetric-monoidal-cat-complexes} 和
\ref{more-algebra-lemma-derived-tor-homotopy} 得到。略去细节。
\end{proof}

\begin{lemma}
\label{more-algebra-lemma-derived-tor-exact}
设 $R$ 为环。设 $P^\bullet$ 为 $R$-模复形。函子
$$
K(R) \longrightarrow K(R), \quad
L^\bullet \longmapsto \text{Tot}(P^\bullet \otimes_R L^\bullet)
$$
以及
$$
K(R) \longrightarrow K(R), \quad
L^\bullet \longmapsto \text{Tot}(L^\bullet \otimes_R P^\bullet)
$$
都是三角范畴的正合函子。
\end{lemma}

\begin{proof}
这由 Derived Categories，注
\ref{derived-remark-double-complex-as-tensor-product-of} 得到。
\end{proof}




\section{导出张量积}
\label{more-algebra-section-derived-tensor-product}

\noindent
我们可以在更一般的情形下构造导出张量积。事实上，只要选取
K-平坦解消，就不需要有界性假设。

\begin{definition}
\label{more-algebra-definition-K-flat}
设 $R$ 为环。若对每个无环复形 $M^\bullet$，全复形
$\text{Tot}(M^\bullet \otimes_R K^\bullet)$ 都无环，
则称复形 $K^\bullet$ 是{\it K-平坦的}。
\end{definition}

\begin{lemma}
\label{more-algebra-lemma-K-flat-quasi-isomorphism}
设 $R$ 为环。设 $K^\bullet$ 为 K-平坦复形。那么，函子
$$
K(R) \longrightarrow K(R), \quad
L^\bullet \longmapsto \text{Tot}(L^\bullet \otimes_R K^\bullet)
$$
将拟同构变为拟同构。
\end{lemma}

\begin{proof}
这由引理 \ref{more-algebra-lemma-derived-tor-exact} 以及以下事实得到：
$K(R)$ 中的拟同构由其锥无环刻画。
\end{proof}

\begin{lemma}
\label{more-algebra-lemma-base-change-K-flat}
设 $R \to R'$ 为环同态。若 $K^\bullet$ 是 $R$-模的 K-平坦复形，
则 $K^\bullet \otimes_R R'$ 是 $R'$-模的 K-平坦复形。
\end{lemma}

\begin{proof}
这由定义以及以下事实得到：对任意 $R'$-模复形 $L^\bullet$，有
$(K^\bullet \otimes_R R') \otimes_{R'} L^\bullet =
K^\bullet \otimes_R L^\bullet$。
\end{proof}

\begin{lemma}
\label{more-algebra-lemma-tensor-product-K-flat}
设 $R$ 为环。若 $K^\bullet$、$L^\bullet$ 是 $R$-模的 K-平坦复形，
则 $\text{Tot}(K^\bullet \otimes_R L^\bullet)$ 是 $R$-模的 K-平坦复形。
\end{lemma}

\begin{proof}
这由同构
$$
\text{Tot}(M^\bullet \otimes_R \text{Tot}(K^\bullet \otimes_R L^\bullet))
=
\text{Tot}(\text{Tot}(M^\bullet \otimes_R K^\bullet) \otimes_R L^\bullet)
$$
和定义得到。
\end{proof}

\begin{lemma}
\label{more-algebra-lemma-K-flat-two-out-of-three}
设 $R$ 为环。设 $(K_1^\bullet, K_2^\bullet, K_3^\bullet)$
为 $K(R)$ 中的可区别三角。若 $K_i^\bullet$ 中三个有两个 K-平坦，
则第三个也 K-平坦。
\end{lemma}

\begin{proof}
这由引理 \ref{more-algebra-lemma-derived-tor-exact} 以及以下事实得到：
$K(R)$ 中的可区别三角若三个对象有两个无环，则第三个也无环。
\end{proof}

\begin{lemma}
\label{more-algebra-lemma-K-flat-two-out-of-three-ses}
设 $R$ 为环。设
$0 \to K_1^\bullet \to K_2^\bullet \to K_3^\bullet \to 0$
为复形的短正合序列。若对所有 $n \in \mathbf{Z}$，$K_3^n$
平坦，并且 $K_i^\bullet$ 中三个有两个 K-平坦，则第三个也 K-平坦。
\end{lemma}

\begin{proof}
设 $L^\bullet$ 为 $R$-模复形。那么
$$
0 \to
\text{Tot}(L^\bullet \otimes_R K_1^\bullet) \to
\text{Tot}(L^\bullet \otimes_R K_2^\bullet) \to
\text{Tot}(L^\bullet \otimes_R K_3^\bullet) \to 0
$$
是复形的短正合序列。事实上，对每个 $n, m$，模序列
$0 \to L^n \otimes_R K_1^m \to
L^n \otimes_R K_2^m \to
L^n \otimes_R K_3^m \to 0$ 正合，这是由 Algebra，引理
\ref{algebra-lemma-flat-tor-zero} 得到的，而上述复形序列是这些序列的直和。
因此，引理由此以及以下事实得到：复形的短正合序列若三个有两个无环，
则第三个也无环。
\end{proof}

\begin{lemma}
\label{more-algebra-lemma-derived-tor-quasi-isomorphism}
设 $R$ 为环。设 $P^\bullet$ 为由平坦 $R$-模构成的上有界复形。
那么 $P^\bullet$ 是 K-平坦的。
\end{lemma}

\begin{proof}
设 $L^\bullet$ 为无环 $R$-模复形。设
$\xi \in H^n(\text{Tot}(L^\bullet \otimes_R P^\bullet))$。
我们必须证明 $\xi = 0$。由于
$\text{Tot}^n(L^\bullet \otimes_R P^\bullet)$ 是以
$L^a \otimes_R P^b$ 为项的直和，可知对某个 $m \in \mathbf{Z}$，
$\xi$ 来自
$H^n(\text{Tot}(\tau_{\leq m}L^\bullet \otimes_R P^\bullet))$
中的元素。由于 $\tau_{\leq m}L^\bullet$ 也无环，可以用它代替
$L^\bullet$，即用 $\tau_{\leq m}L^\bullet$ 作此替换。因此，
可以假设 $L^\bullet$ 上有界。在这种情况下，
Homology，引理 \ref{homology-lemma-first-quadrant-ss} 的谱序列满足
$$
{}'E_1^{p, q} = H^p(L^\bullet \otimes_R P^q)
$$
由于 $P^q$ 平坦且 $L^\bullet$ 无环，该项为零。因此
$H^*(\text{Tot}(L^\bullet \otimes_R P^\bullet)) = 0$。
\end{proof}

\noindent
在下一个引理中，复形系统的余极限是指逐项余极限。

\begin{lemma}
\label{more-algebra-lemma-colimit-K-flat}
设 $R$ 为环。设
$K_1^\bullet \to K_2^\bullet \to \ldots$ 为 K-平坦复形的系统。
那么 $\colim_i K_i^\bullet$ 是 K-平坦的。更一般地，
K-平坦复形的任意滤过余极限都是 K-平坦的。
\end{lemma}

\begin{proof}
由于取的是逐项余极限，由 Algebra，引理
\ref{algebra-lemma-tensor-products-commute-with-limits}，有
$$
\colim_i \text{Tot}(M^\bullet \otimes_R K_i^\bullet) =
\text{Tot}(M^\bullet \otimes_R \colim_i K_i^\bullet)
$$
因此，引理由滤过余极限正合这一事实得到；见 Algebra，引理
\ref{algebra-lemma-directed-colimit-exact}。
\end{proof}

\begin{lemma}
\label{more-algebra-lemma-universally-acyclic-K-flat}
设 $R$ 为环。设 $K^\bullet$ 为 $R$-模复形。若对所有有限表示
$R$-模 $M$，$K^\bullet \otimes_R M$ 都无环，则 $K^\bullet$ 是 K-平坦的。
\end{lemma}

\begin{proof}
我们将反复使用张量积与余极限交换这一事实（Algebra，引理
\ref{algebra-lemma-tensor-products-commute-with-limits}）。因此，
对任意 $R$-模 $M$，$K^\bullet \otimes_R M$ 都无环，因为任意
$R$-模都是有限表示 $R$-模 $M$ 的滤过余极限；见 Algebra，引理
\ref{algebra-lemma-module-colimit-fp}。设 $M^\bullet$ 为无环
$R$-模复形。我们必须证明
$\text{Tot}(M^\bullet \otimes_R K^\bullet)$ 无环。由于
$M^\bullet = \colim \tau_{\leq n} M^\bullet$（逐项余极限），有
$$
\text{Tot}(M^\bullet \otimes_R K^\bullet) =
\colim \text{Tot}(\tau_{\leq n} M^\bullet \otimes_R K^\bullet)
$$
其中截断如 Homology，第 \ref{homology-section-truncations} 节所述。
由于滤过余极限正合（Algebra，引理
\ref{algebra-lemma-directed-colimit-exact}），可以用
$\tau_{\leq n}M^\bullet$ 代替 $M^\bullet$，并假设 $M^\bullet$ 上有界。
在上有界情形下，可以写成
$M^\bullet = \colim \sigma_{\geq -n} M^\bullet$，其中复形
$\sigma_{\geq -n} M^\bullet$ 有界，但可能不再无环。
按上述方式论证，问题化约为 $M^\bullet$ 为有界复形的情形。
最后，对有界复形 $M^a \to \ldots \to M^b$，可以对复形的长度
$b - a$ 作归纳。情形 $b - a = 1$ 已在上面处理。
当 $b - a > 1$ 时，考虑复形的分裂短正合序列
$$
0 \to \sigma_{\geq a + 1}M^\bullet \to M^\bullet \to M^a[-a] \to 0
$$
并应用引理 \ref{more-algebra-lemma-derived-tor-exact} 完成归纳步骤。
略去一些细节。
\end{proof}

\begin{lemma}
\label{more-algebra-lemma-K-flat-resolution}
设 $R$ 为环。对任意复形 $M^\bullet$，都存在各项为平坦 $R$-模的
K-平坦复形 $K^\bullet$，以及逐项满射的拟同构
$K^\bullet \to M^\bullet$。
\end{lemma}

\begin{proof}
令 $\mathcal{P} \subset \Ob(\text{Mod}_R)$ 为平坦 $R$-模的类。
由 Derived Categories，引理 \ref{derived-lemma-special-direct-system}，
存在系统 $K_1^\bullet \to K_2^\bullet \to \ldots$ 和交换图
$$
\xymatrix{
K_1^\bullet \ar[d] \ar[r] &
K_2^\bullet \ar[d] \ar[r] & \ldots \\
\tau_{\leq 1}M^\bullet \ar[r] &
\tau_{\leq 2}M^\bullet \ar[r] & \ldots
}
$$
并具有该引理列出的性质 (1)、(2)、(3)。这些性质蕴含每个复形
$K_i^\bullet$ 都是由平坦模构成的上有界复形。因此，由引理
\ref{more-algebra-lemma-derived-tor-quasi-isomorphism}，$K_i^\bullet$ 是 K-平坦的。
依构造，诱导映射 $\colim_i K_i^\bullet \to M^\bullet$
是拟同构且逐项满射。由引理 \ref{more-algebra-lemma-colimit-K-flat}，复形
$\colim_i K_i^\bullet$ 是 K-平坦的。各项 $\colim K_i^n$ 平坦，
因为平坦模的滤过余极限平坦；见 Algebra，引理
\ref{algebra-lemma-colimit-flat}。
\end{proof}

\begin{remark}
\label{more-algebra-remark-P-resolution}
事实上，可以得到比引理 \ref{more-algebra-lemma-K-flat-resolution} 更强的结果。
具体地，可以找到拟同构 $P^\bullet \to M^\bullet$，其中
$P^\bullet$ 是赋有子复形滤过
$$
0 = F_{-1}P^\bullet \subset F_0P^\bullet \subset
F_1P^\bullet \subset \ldots \subset P^\bullet
$$
的 $R$-模复形，并满足
\begin{enumerate}
\item $P^\bullet = \bigcup F_pP^\bullet$；
\item 包含 $F_iP^\bullet \to F_{i + 1}P^\bullet$ 是逐项分裂单射；
\item 商 $F_{i + 1}P^\bullet/F_iP^\bullet$ 同构于平移
$R[k]$ 的直和（作为复形，因此其微分为零）。
\end{enumerate}
这将在 Differential Graded Algebra，引理 \ref{dga-lemma-resolve}
中证明（也可以像引理 \ref{more-algebra-lemma-resolution-filtered-complex}
的证明那样论证）。此外，给定这样的复形，在 $D(R)$ 中得到可区别三角
$$
\bigoplus F_iP^\bullet \to \bigoplus F_iP^\bullet \to M^\bullet
\to \bigoplus F_iP^\bullet[1]
$$
利用它，有时可以把关于一般复形的陈述化约为关于 $R[k]$
的陈述（当然，这只在该陈述于取直和下保持时有效）。更精确地，
设 $T$ 为 $D(R)$ 中对象的一种性质。假设
\begin{enumerate}
\item 若 $K_i \in D(R)$、$i \in I$ 是一族对象，且对所有
$i \in I$ 都有 $T(K_i)$，则 $T(\bigoplus K_i)$；
\item 若 $K \to L \to M \to K[1]$ 是可区别三角，并且 $T$
对其中两个对象成立，则 $T$ 对第三个对象也成立；
\item 对所有 $k$，$T(R[k])$ 成立。
\end{enumerate}
那么 $T$ 对 $D(R)$ 的所有对象成立。
\end{remark}

\begin{lemma}
\label{more-algebra-lemma-derived-tor-quasi-isomorphism-other-side}
设 $R$ 为环。设
$\alpha : P^\bullet \to Q^\bullet$ 为 $R$-模的 K-平坦复形之间的拟同构。
对每个 $R$-模复形 $L^\bullet$，诱导映射
$$
\text{Tot}(\text{id}_L \otimes \alpha) :
\text{Tot}(L^\bullet \otimes_R P^\bullet)
\longrightarrow
\text{Tot}(L^\bullet \otimes_R Q^\bullet)
$$
是拟同构。
\end{lemma}

\begin{proof}
选取拟同构 $K^\bullet \to L^\bullet$，其中 $K^\bullet$
为 K-平坦复形；见引理 \ref{more-algebra-lemma-K-flat-resolution}。考虑交换图
$$
\xymatrix{
\text{Tot}(K^\bullet \otimes_R P^\bullet) \ar[r] \ar[d] &
\text{Tot}(K^\bullet \otimes_R Q^\bullet) \ar[d] \\
\text{Tot}(L^\bullet \otimes_R P^\bullet) \ar[r] &
\text{Tot}(L^\bullet \otimes_R Q^\bullet)
}
$$
由引理 \ref{more-algebra-lemma-K-flat-quasi-isomorphism}，两条竖直箭头和上方水平箭头
都是拟同构，故结论成立。
\end{proof}

\noindent
设 $R$ 为环。设 $M^\bullet$ 为 $D(R)$ 的对象。选取 K-平坦解消
$K^\bullet \to M^\bullet$；见引理 \ref{more-algebra-lemma-K-flat-resolution}。
由引理 \ref{more-algebra-lemma-derived-tor-homotopy} 和
\ref{more-algebra-lemma-derived-tor-exact}，得到三角范畴的正合函子
$$
K(R) \longrightarrow K(R), \quad
L^\bullet \longmapsto \text{Tot}(L^\bullet \otimes_R K^\bullet)
$$
由引理 \ref{more-algebra-lemma-K-flat-quasi-isomorphism}，该函子诱导函子
$D(R) \to D(R)$，这仅仅是因为 $D(R)$ 是 $K(R)$
关于拟同构的局部化。由于 $M^\bullet$ 的 $K$-平坦解消范畴是余滤过的，
并且有引理 \ref{more-algebra-lemma-derived-tor-quasi-isomorphism-other-side}，
所得函子（在同构意义下）不依赖于 $K^\bullet$ 的选取。

\begin{definition}
\label{more-algebra-definition-derived-tor}
设 $R$ 为环。设 $M^\bullet$ 为 $D(R)$ 的对象。{\it 导出张量积}
$$
- \otimes_R^{\mathbf{L}} M^\bullet : D(R) \longrightarrow D(R)
$$
是上述三角范畴的正合函子。
\end{definition}

\noindent
该函子延拓函子（\ref{more-algebra-equation-derived-tensor-module}）。
由我们的显式构造显然可见，只要 $L^\bullet$ 和 $M^\bullet$
都属于 $D(R)$，就存在同构（涉及符号的选取，见下文）
$$
M^\bullet \otimes_R^{\mathbf{L}} L^\bullet
\cong
L^\bullet \otimes_R^{\mathbf{L}} M^\bullet
$$
因此，在写 $M^\bullet \otimes_R^{\mathbf{L}} L^\bullet$ 时，
我们通常不区分使用哪个变量来定义与之作导出张量积的函子。

\begin{lemma}
\label{more-algebra-lemma-flip-douoble-tensor-product}
设 $R$ 为环。设 $K^\bullet, L^\bullet$ 为 $R$-模复形。
存在关于两个复形都具有函子性的典范同构
$$
K^\bullet \otimes_R^\mathbf{L} L^\bullet \longrightarrow
L^\bullet \otimes_R^\mathbf{L} K^\bullet
$$
它在映射 $K^p \otimes_R L^q \to L^q \otimes_R K^p$ 上使用符号
$(-1)^{pq}$（解释见证明）。
\end{lemma}

\begin{proof}
我们可以并且确实用 K-平坦复形 $K^\bullet$ 和 $L^\bullet$
代替原复形，然后使用第 \ref{more-algebra-section-symmetric-monoidal} 节中讨论的交换约束。
\end{proof}

\begin{lemma}
\label{more-algebra-lemma-triple-tensor-product}
设 $R$ 为环。设 $K^\bullet, L^\bullet, M^\bullet$ 为 $R$-模复形。
存在关于三个复形都具有函子性的典范同构
$$
(K^\bullet \otimes_R^\mathbf{L} L^\bullet) \otimes_R^\mathbf{L} M^\bullet
=
K^\bullet \otimes_R^\mathbf{L} (L^\bullet \otimes_R^\mathbf{L} M^\bullet)
$$
\end{lemma}

\begin{proof}
用 K-平坦复形代替这些复形，并使用第
\ref{more-algebra-section-symmetric-monoidal} 节中的结合约束。
\end{proof}

\begin{lemma}
\label{more-algebra-lemma-factor-through-K-flat}
设 $R$ 为环。设 $a : K^\bullet \to L^\bullet$ 为 $R$-模复形的映射。
若 $K^\bullet$ K-平坦，则存在复形 $N^\bullet$ 以及复形映射
$b : K^\bullet \to N^\bullet$ 和 $c : N^\bullet \to L^\bullet$，使得
\begin{enumerate}
\item $N^\bullet$ K-平坦；
\item $c$ 是拟同构；
\item $a$ 与 $c \circ b$ 同伦。
\end{enumerate}
若 $K^\bullet$ 的各项平坦，则可以选取 $N^\bullet$、$b$ 和 $c$，
使得 $N^\bullet$ 也具有同样的性质。
\end{lemma}

\begin{proof}
我们将使用同伦范畴 $K(R)$ 是三角范畴这一事实；见 Derived Categories，
命题 \ref{derived-proposition-homotopy-category-triangulated}。
选取可区别三角
$K^\bullet \to L^\bullet \to C^\bullet \to K^\bullet[1]$。
选取拟同构 $M^\bullet \to C^\bullet$，其中 $M^\bullet$
K-平坦且各项平坦；见引理 \ref{more-algebra-lemma-K-flat-resolution}。
由三角范畴的公理，可以把复合
$M^\bullet \to C^\bullet \to K^\bullet[1]$ 嵌入可区别三角
$K^\bullet \to N^\bullet \to M^\bullet \to K^\bullet[1]$。
由引理 \ref{more-algebra-lemma-K-flat-two-out-of-three}，$N^\bullet$ K-平坦。
再次使用三角范畴的公理，可以选取映射
$N^\bullet \to L^\bullet$，将其嵌入下列可区别三角的态射：
$$
\xymatrix{
K^\bullet \ar[r] \ar[d] &
N^\bullet \ar[r] \ar[d] &
M^\bullet \ar[r] \ar[d] &
K^\bullet[1] \ar[d] \\
K^\bullet \ar[r] &
L^\bullet \ar[r] &
C^\bullet \ar[r] &
K^\bullet[1]
}
$$
三条箭头中有两条是拟同构，故由这些可区别三角所伴随的上同调长正合序列，
第三条箭头 $N^\bullet \to L^\bullet$ 也是拟同构（也可以考察该图在
$D(R)$ 中的像，并使用 Derived Categories，引理
\ref{derived-lemma-third-isomorphism-triangle}）。这完成了 (1)、(2)、(3)
的证明。为证明最后一个断言，可以选取 $N^\bullet$，使得
$N^n \cong M^n \oplus K^n$；见 Derived Categories，引理
\ref{derived-lemma-improve-distinguished-triangle-homotopy}。
因此，若 $K^\bullet$ 的各项平坦，就得到所需的平坦性。
\end{proof}








\section{导出换环}
\label{more-algebra-section-derived-base-change}

\noindent
设 $R \to A$ 为环同态。设 $N^\bullet$ 为 $A$-模复形。
也可以使用 K-平坦解消，把函子
$$
- \otimes_R^{\mathbf{L}} N^\bullet : D(R) \to D(A)
$$
定义为函子 $K(R) \to K(A)$，
$M^\bullet \mapsto \text{Tot}(M^\bullet \otimes_R N^\bullet)$
的左导出函子。特别地，取 $N^\bullet = A[0]$，得到导出基变换函子
$$
- \otimes_R^{\mathbf{L}} A : D(R) \to D(A)
$$
它延拓函子（\ref{more-algebra-equation-derived-tensor-algebra}）。具体地，
对每个 $R$-模复形 $M^\bullet$，可以选取 K-平坦解消
$K^\bullet \to M^\bullet$，并置
$$
M^\bullet \otimes_R^{\mathbf{L}} N^\bullet =
\text{Tot}(K^\bullet \otimes_R N^\bullet).
$$
可以使用引理 \ref{more-algebra-lemma-K-flat-resolution} 和
\ref{more-algebra-lemma-derived-tor-quasi-isomorphism-other-side} 看出该定义良好。
不过，为了把所有细节都交代周全，使用一些一般理论或许更方便。

\begin{lemma}
\label{more-algebra-lemma-derived-base-change}
上述构造不依赖于选取，并定义三角范畴的正合函子
$- \otimes_R^\mathbf{L} N^\bullet : D(R) \to D(A)$。
对 $D(R)$ 中的 $E^\bullet$，存在函子性同构
$$
E^\bullet \otimes_R^\mathbf{L} N^\bullet =
(E^\bullet \otimes_R^\mathbf{L} A) \otimes_A^\mathbf{L} N^\bullet
$$
\end{lemma}

\begin{proof}
为了证明导出函子 $- \otimes_R^\mathbf{L} N^\bullet$ 的存在性，
使用 Derived Categories，第 \ref{derived-section-derived-functors} 节
中发展的一般理论。置 $\mathcal{D} = K(R)$ 和
$\mathcal{D}' = D(A)$。记 $F : \mathcal{D} \to \mathcal{D}'$
为由规则 $F(M^\bullet) = \text{Tot}(M^\bullet \otimes_R N^\bullet)$
定义的三角范畴正合函子。要证明 $F$ 的上述性质，使用引理
\ref{more-algebra-lemma-derived-tor-homotopy} 和 \ref{more-algebra-lemma-derived-tor-exact}。
令 $S$ 为 $\mathcal{D} = K(R)$ 中所有拟同构构成的集合。
这给出 Derived Categories，情形
\ref{derived-situation-derived-functor} 中的情况，因此可应用
Derived Categories，定义
\ref{derived-definition-right-derived-functor-defined}。
我们断言 $LF$ 处处有定义。由 Derived Categories，引理
\ref{derived-lemma-find-existence-computes} 可得这一点，其中取
$\mathcal{P} \subset \Ob(\mathcal{D})$ 为 K-平坦复形的集合：
(1) 由引理 \ref{more-algebra-lemma-K-flat-resolution} 得到，(2) 由引理
\ref{more-algebra-lemma-derived-tor-quasi-isomorphism-other-side} 得到。
因此，得到导出函子
$$
LF : D(R) = S^{-1}\mathcal{D} \longrightarrow \mathcal{D}' = D(A)
$$
见 Derived Categories，等式（\ref{derived-equation-everywhere}）。
最后，Derived Categories，引理
\ref{derived-lemma-find-existence-computes} 保证：当 $K^\bullet$
K-平坦时，
$LF(K^\bullet) = F(K^\bullet) =
\text{Tot}(K^\bullet \otimes_R N^\bullet)$；换言之，$LF$
确实按上述方式计算。此外，由引理 \ref{more-algebra-lemma-base-change-K-flat}，
复形 $K^\bullet \otimes_R A$ 是 $A$-模的 K-平坦复形。因此
$$
(K^\bullet \otimes_R^\mathbf{L} A) \otimes_A^\mathbf{L} N^\bullet =
\text{Tot}((K^\bullet \otimes_R A) \otimes_A N^\bullet) =
\text{Tot}(K^\bullet \otimes_A N^\bullet) =
K^\bullet \otimes_A^\mathbf{L} N^\bullet
$$
这证明了引理的最后一个陈述。
\end{proof}

\begin{lemma}
\label{more-algebra-lemma-functoriality-derived-base-change}
设 $R \to A$ 为环同态。设 $f : L^\bullet \to N^\bullet$
为 $A$-模复形的映射。那么 $f$ 诱导函子变换
$$
1 \otimes f :
- \otimes_A^\mathbf{L} L^\bullet
\longrightarrow
- \otimes_A^\mathbf{L} N^\bullet
$$
若 $f$ 是拟同构，则 $1 \otimes f$ 是函子的同构。
\end{lemma}

\begin{proof}
由于这些函子通过在 K-平坦复形 $K^\bullet$ 上取值来计算，
可以直接使用函子性
$$
\text{Tot}(K^\bullet \otimes_R L^\bullet) \to
\text{Tot}(K^\bullet \otimes_R N^\bullet)
$$
定义该变换。最后一个陈述由引理
\ref{more-algebra-lemma-K-flat-quasi-isomorphism} 得到。
\end{proof}

\begin{lemma}
\label{more-algebra-lemma-tensor-hom-adjoint}
设 $R \to A$ 为环同态。引理 \ref{more-algebra-lemma-derived-base-change}
的函子 $D(R) \to D(A)$，$E \mapsto E \otimes_R^\mathbf{L} A$，
是限制函子 $D(A) \to D(R)$ 的左伴随。
\end{lemma}

\begin{proof}
这由 Derived Categories，引理
\ref{derived-lemma-pre-derived-adjoint-functors-general}，
以及由 Algebra，引理 \ref{algebra-lemma-adjoint-tensor-restrict}
得到的事实，即 $- \otimes_R A$ 与限制互为伴随，得到。
\end{proof}

\begin{remark}[警告]
\label{more-algebra-remark-warning-compute-base-change}
设 $R \to A$ 为环同态，并设 $N$ 和 $N'$ 为 $A$-模。
以 $N_R$ 和 $N'_R$ 表示把 $N$ 和 $N'$ 限制为 $R$-模所得的模；
见 Algebra，第 \ref{algebra-section-base-change} 节。在这种情况下，
$D(A)$ 中的对象 $N_R \otimes_R^\mathbf{L} N'$ 和
$N \otimes_R^\mathbf{L} N'_R$ 一般不同构！换言之，必须仔细分辨
两侧中的哪一侧用于给出 $A$-模结构。

\medskip\noindent
作为具体例子，置 $R = k[x, y]$、$A = R/(xy)$、$N = R/(x)$
以及 $N' = A = R/(xy)$。解消
$0 \to R \xrightarrow{xy} R \to N'_R \to 0$ 表明在 $D(A)$ 中，
$N \otimes_R^\mathbf{L} N'_R = N[1] \oplus N$。解消
$0 \to R \xrightarrow{x} R \to N_R \to 0$ 表明
$N_R \otimes_R^\mathbf{L} N'$ 由复形 $A \xrightarrow{x} A$ 表示。
要看出这两个复形不同构，可以证明第二个复形在 $D(A)$ 中不同构于
其上同调群的直和，也可以证明第一个复形不是 $D(A)$ 的完美对象，
而第二个是。略去一些细节。
\end{remark}

\begin{lemma}
\label{more-algebra-lemma-double-base-change}
设 $A \to B \to C$ 为环同态。设 $N^\bullet$ 为 $B$-模复形，
$K^\bullet$ 为 $C$-模复形。下列函子的复合
$$
D(A) \xrightarrow{- \otimes_A^\mathbf{L} N^\bullet}
D(B) \xrightarrow{- \otimes_B^\mathbf{L} K^\bullet} D(C)
$$
是函子
$- \otimes_A^\mathbf{L} (N^\bullet \otimes_B^\mathbf{L} K^\bullet) :
D(A) \to D(C)$。如果 $M$、$N$、$K$ 分别是 $A$、$B$、$C$ 上的模，
则在 $D(C)$ 中有
$$
(M \otimes_A^\mathbf{L} N) \otimes_B^\mathbf{L} K =
M \otimes_A^\mathbf{L} (N \otimes_B^\mathbf{L} K) =
(M \otimes_A^\mathbf{L} C) \otimes_C^\mathbf{L} (N \otimes_B^\mathbf{L} K)
$$
还存在典范同构
$$
(M \otimes_A^\mathbf{L} N) \otimes_B^\mathbf{L} K \longrightarrow
(M \otimes_A^\mathbf{L} K) \otimes_C^\mathbf{L} (N \otimes_B^\mathbf{L} C)
$$
其中使用符号。对复形也有类似结果。
\end{lemma}

\begin{proof}
选取 $B$-模的 K-平坦复形 $P^\bullet$ 以及拟同构
$P^\bullet \to N^\bullet$（引理 \ref{more-algebra-lemma-K-flat-resolution}）。
设 $M^\bullet$ 为表示 $D(A)$ 中任意对象的 $A$-模 K-平坦复形。
那么，由引理 \ref{more-algebra-lemma-functoriality-derived-base-change}
应用于括号内的对象可知，
$$
(M^\bullet \otimes_A^\mathbf{L} P^\bullet) \otimes_B^\mathbf{L} K^\bullet
\longrightarrow
(M^\bullet \otimes_A^\mathbf{L} N^\bullet) \otimes_B^\mathbf{L} K^\bullet
$$
是同构。由引理 \ref{more-algebra-lemma-base-change-K-flat} 和
\ref{more-algebra-lemma-tensor-product-K-flat}，复形
$$
\text{Tot}(M^\bullet \otimes_A P^\bullet) =
\text{Tot}((M^\bullet \otimes_R A) \otimes_A P^\bullet
$$
作为 $B$-模复形是 K-平坦的，并且按构造表示 $D(B)$ 中的导出张量积。
因此，
$(M^\bullet \otimes_A^\mathbf{L} P^\bullet) \otimes_B^\mathbf{L} K^\bullet$
由 $C$-模复形
$$
\text{Tot}(\text{Tot}(M^\bullet \otimes_A P^\bullet)\otimes_B K^\bullet) =
\text{Tot}(M^\bullet \otimes_A \text{Tot}(P^\bullet \otimes_B K^\bullet))
$$
表示。该等式见 Homology，评注 \ref{homology-remark-triple-complex}。
沿原路返回可见，这等于
$$
M^\bullet \otimes_A^\mathbf{L} (P^\bullet \otimes_B^\mathbf{L} K^\bullet)
\longleftarrow
M^\bullet \otimes_A^\mathbf{L} (N^\bullet \otimes_B^\mathbf{L} K^\bullet)
$$
按照函子 $-\otimes_B^\mathbf{L} K^\bullet$ 的定义，该箭头是同构。
所有这些构造关于复形 $M^\bullet$ 都具有函子性，故得到所述函子同构。

\medskip\noindent
由上文，下式中的第一个等式成立：
$$
(M \otimes_A^\mathbf{L} N) \otimes_B^\mathbf{L} K =
M \otimes_A^\mathbf{L} (N \otimes_B^\mathbf{L} K) =
(M \otimes_A^\mathbf{L} C) \otimes_C^\mathbf{L} (N \otimes_B^\mathbf{L} K)
$$
第二个等式由引理 \ref{more-algebra-lemma-derived-base-change} 的最后一个陈述得到。
同一陈述还允许我们写成
$N \otimes_B^\mathbf{L} K = (N \otimes_B^\mathbf{L} C) \otimes_C^\mathbf{L} K$，
代入即得
\begin{align*}
(M \otimes_A^\mathbf{L} N) \otimes_B^\mathbf{L} K
& =
(M \otimes_A^\mathbf{L} C) \otimes_C^\mathbf{L}
((N \otimes_B^\mathbf{L} C) \otimes_C^\mathbf{L} K) \\
& =
(M \otimes_A^\mathbf{L} C) \otimes_C^\mathbf{L}
(K \otimes_C^\mathbf{L} (N \otimes_B^\mathbf{L} C)) \\
& =
((M \otimes_A^\mathbf{L} C) \otimes_C^\mathbf{L} K)
\otimes_C^\mathbf{L} (N \otimes_B^\mathbf{L} C)) \\
& =
(M \otimes_C^\mathbf{L} K)
\otimes_C^\mathbf{L} (N \otimes_B^\mathbf{L} C)
\end{align*}
这里使用了引理 \ref{more-algebra-lemma-flip-douoble-tensor-product}、
\ref{more-algebra-lemma-triple-tensor-product} 以及前面提到的引理。
\end{proof}







\section{Tor 无关性}
\label{more-algebra-section-tor-independence}

\noindent
考虑环的交换图
$$
\xymatrix{
A \ar[r] & A' \\
R \ar[r] \ar[u] & R' \ar[u]
}
$$
给定 $D(A)$ 的对象 $K$，可以考虑它到 $D(A')$ 的对象
$K \otimes_A^\mathbf{L} A'$ 的导出基变换。也可以先把 $K$ 限制为
$D(R)$ 的对象，再考虑它到 $D(R')$ 的导出基变换，记为
$K \otimes_R^\mathbf{L} R'$。我们断言，在 $D(R')$ 中存在函子性比较映射
\begin{equation}
\label{more-algebra-equation-comparison-map}
K \otimes_R^{\mathbf{L}} R'
\longrightarrow
K \otimes_A^{\mathbf{L}} A'
\end{equation}
为构造这个比较映射，选取表示 $K$ 的 $A$-模 K-平坦复形
$K^\bullet$。接着选取拟同构 $E^\bullet \to K^\bullet$，其中
$E^\bullet$ 是 $R$-模的 K-平坦复形。上述映射就是
$$
K \otimes_R^{\mathbf{L}} R' =
E^\bullet \otimes_R R'
\longrightarrow
K^\bullet \otimes_A A' =
K \otimes_A^{\mathbf{L}} A'
$$
一般而言，该映射不可能是同构。

\medskip\noindent
不过，我们经常遇到上述图是环的“基变换”图，即
$A' = A \otimes_R R'$ 的情形。在这种情况下，对任意 $A$-模 $M$，都有
$M \otimes_A A' = M \otimes_R R'$。因此，作为函子
$\text{Mod}_A \to \text{Mod}_{A'}$，$- \otimes_R R'$ 等于
$- \otimes_A A'$。一般而言，这个等式{\bf 并不延拓}到导出张量积。
换言之，比较映射不是同构。一个简单例子是取
$R = k[x]$、$A = R' = A' = k[x]/(x) = k$ 和 $K^\bullet = A[0]$。
显然，一个必要条件是对所有 $p > 0$ 都有
$\text{Tor}_p^R(A, R') = 0$。

\begin{definition}
\label{more-algebra-definition-tor-independent}
设 $R$ 为环。设 $A$、$B$ 为 $R$-代数。如果对所有 $p > 0$ 都有
$\text{Tor}_p^R(A, B) = 0$，则称 $A$ 与 $B$ {\it 在 $R$ 上 Tor 无关}。
\end{definition}

\begin{lemma}
\label{more-algebra-lemma-base-change-comparison}
如果 $A' = A \otimes_R R'$，且 $A$ 与 $R'$ 在 $R$ 上 Tor 无关，
则比较映射（\ref{more-algebra-equation-comparison-map}）是同构。
\end{lemma}

\begin{proof}
为证明这一点，选取 $R'$ 作为 $R$-模的自由解消
$F^\bullet \to R'$。由于 $A$ 与 $R'$ 在 $R$ 上 Tor 无关，
$F^\bullet \otimes_R A$ 是 $A'$ 在 $A$ 上的自由 $A$-模解消。
由上文对导出张量积的一般构造可见
$$
K^\bullet \otimes_A A' \cong
\text{Tot}(K^\bullet \otimes_A (F^\bullet \otimes_R A)) =
\text{Tot}(K^\bullet \otimes_R F^\bullet) \cong
\text{Tot}(E^\bullet \otimes_R F^\bullet) \cong
E^\bullet \otimes_R R'
$$
这正是所需结论。
\end{proof}

\begin{lemma}
\label{more-algebra-lemma-tor-independent-flat}
考虑环的交换图
$$
\xymatrix{
A' & R' \ar[r] \ar[l] & B' \\
A \ar[u] & R \ar[l] \ar[u] \ar[r] & B \ar[u]
}
$$
假设 $R'$ 在 $R$ 上平坦，$A'$ 在 $A \otimes_R R'$ 上平坦，
且 $B'$ 在 $R' \otimes_R B$ 上平坦。那么
$$
\text{Tor}_i^R(A, B) \otimes_{(A \otimes_R B)} (A' \otimes_{R'} B') =
\text{Tor}_i^{R'}(A', B')
$$
\end{lemma}

\begin{proof}
由 Algebra，第 \ref{algebra-section-functoriality-tor} 节，存在典范映射
$$
\text{Tor}_i^R(A, B) \longrightarrow
\text{Tor}_i^{R'}(A \otimes_R R', B \otimes_R R') \longrightarrow
\text{Tor}_i^{R'}(A', B')
$$
它们诱导从引理公式左侧到右侧的映射。

\medskip\noindent
取 $A$ 作为 $R$-模的自由解消 $F_\bullet \to A$。
那么 $F_\bullet \otimes_R R'$ 是 $A \otimes_R R'$ 的解消。
因而 $\text{Tor}_i^{R'}(A \otimes_R R', B \otimes_R R')$ 由
$F_\bullet \otimes_R B \otimes_R R'$ 计算。根据 $R'$ 在 $R$ 上平坦的假设，
它计算 $\text{Tor}_i^R(A, B) \otimes_R R'$。因此
$\text{Tor}_i^{R'}(A \otimes_R R', B \otimes_R R') =
\text{Tor}_i^R(A, B) \otimes_R R'$（这里只使用了 $R'$ 在 $R$ 上的平坦性）。

\medskip\noindent
由 Lazard 定理（Algebra，定理 \ref{algebra-theorem-lazard}），可以分别把
$A'$ 和 $B'$ 写成有限自由 $A \otimes_R R'$-模和
$B \otimes_R R'$-模的滤过余极限。记
$A' = \colim M_i$ 和 $B' = \colim N_j$。上面的结果给出
$$
\text{Tor}_i^{R'}(M_i, N_j) =
\text{Tor}_i^R(A, B) \otimes_{A \otimes_R B} (M_i \otimes_{R'} N_j)
$$
把一切用基写出即可看出这一点。取余极限便得到引理的结论。
\end{proof}

\begin{lemma}
\label{more-algebra-lemma-flat-base-change-tor-independent}
设 $R \to A$ 和 $R \to B$ 为环同态。设 $R \to R'$ 为环同态，
并置 $A' = A \otimes_R R'$ 和 $B' = B \otimes_R R'$。
如果 $A$ 与 $B$ 在 $R$ 上 Tor 无关，且 $R \to R'$ 平坦，
则 $A'$ 与 $B'$ 在 $R'$ 上 Tor 无关。
\end{lemma}

\begin{proof}
这立即由引理 \ref{more-algebra-lemma-tor-independent-flat} 和定义
\ref{more-algebra-definition-tor-independent} 得到。
\end{proof}

\begin{lemma}
\label{more-algebra-lemma-lemma-tor-independent-flat-compare}
假设同引理 \ref{more-algebra-lemma-tor-independent-flat}。对 $M \in D(A)$，
存在 $A' \otimes_{R'} B'$-模的典范同构
$$
H^i((M \otimes_A^\mathbf{L} A') \otimes_{R'}^\mathbf{L} B') =
H^i(M \otimes_R^\mathbf{L} B) \otimes_{(A \otimes_R B)} (A' \otimes_{R'} B')
$$
\end{lemma}

\begin{proof}
我们来阐明等式的两侧。在左侧，有函子
$D(A) \to D(A') \to D(R') \to D(B')$ 的复合，再与函子
$H^i : D(B') \to \text{Mod}_{B'}$ 复合。由于存在从 $A'$ 到
$D(B')$ 中对象 $(M \otimes_A^\mathbf{L} A') \otimes_{R'}^\mathbf{L} B'$
的自同态环的映射，可见左侧确实是 $A' \otimes_{R'} B'$-模。
由同样的论证可见，$H^i(M \otimes_R^\mathbf{L} B)$ 具有
$A \otimes_R B$-模结构。

\medskip\noindent
先在 $B' = R' \otimes_R B$ 的情形证明该结论。
在这种情况下，选取由自由 $R$-模组成的解消 $F^\bullet \to B$。
还选取表示 $M$ 的 $A$-模 K-平坦复形 $M^\bullet$。
于是左侧由下式表示：
\begin{align*}
H^i(\text{Tot}((M^\bullet \otimes_A A') \otimes_{R'} (R' \otimes_R F^\bullet)))
& =
H^i(\text{Tot}(M^\bullet \otimes_A A' \otimes_R F^\bullet)) \\
& =
H^i(\text{Tot}(M^\bullet \otimes_R F^\bullet) \otimes_A A') \\
& =
H^i(M \otimes_R^\mathbf{L} B) \otimes_A A'
\end{align*}
最后一个等式成立，是因为 $A \to A'$ 平坦。最后一个模正是所需的模，
因为本段假设 $B' = R' \otimes_R B$，故
$A' \otimes_{R'} B' = A' \otimes_R B$。

\medskip\noindent
现在考虑一般情形。假设 $B' \to B''$ 是平坦环同态。
容易看出
$$
H^i((M \otimes_A^\mathbf{L} A') \otimes_{R'}^\mathbf{L} B'') =
H^i((M \otimes_A^\mathbf{L} A') \otimes_{R'}^\mathbf{L} B')
\otimes_{B'} B''
$$
并且
$$
H^i(M \otimes_R^\mathbf{L} B) \otimes_{(A \otimes_R B)} (A' \otimes_{R'} B'')
=
\left(
H^i(M \otimes_R^\mathbf{L} B) \otimes_{(A \otimes_R B)} (A' \otimes_{R'} B')
\right) \otimes_{B'} B''
$$
因此，对 $B'$ 的结论推出对 $B''$ 的结论。上一段已经证明了
$R' \otimes_R B$ 的情形，故这推出一般结论。
\end{proof}

\begin{lemma}
\label{more-algebra-lemma-tor-independent}
设 $R$ 为环。设 $A$、$B$ 为 $R$-代数。下列条件等价：
\begin{enumerate}
\item $A$ 与 $B$ 在 $R$ 上 Tor 无关；
\item 对位于同一素理想 $\mathfrak r \subset R$ 之上的每一对素理想
$\mathfrak p \subset A$ 和 $\mathfrak q \subset B$，环 $A_\mathfrak p$
与 $B_\mathfrak q$ 在 $R_\mathfrak r$ 上 Tor 无关；
\item 对 $A \otimes_R B$ 的每个素理想 $\mathfrak s$，模
$$
\text{Tor}_i^R(A, B)_\mathfrak s =
\text{Tor}_i^{R_\mathfrak r}(A_\mathfrak p, B_\mathfrak q)_\mathfrak s
$$
为零（其中 $\mathfrak p = A \cap \mathfrak s$、
$\mathfrak q = B \cap \mathfrak s$ 且 $\mathfrak r = R \cap \mathfrak s$）。
\end{enumerate}
\end{lemma}

\begin{proof}
设 $\mathfrak s$ 为 (3) 中 $A \otimes_R B$ 的素理想。等式
$$
\text{Tor}_i^R(A, B)_\mathfrak s =
\text{Tor}_i^{R_\mathfrak r}(A_\mathfrak p, B_\mathfrak q)_\mathfrak s
$$
其中 $\mathfrak p = A \cap \mathfrak s$、$\mathfrak q = B \cap \mathfrak s$
且 $\mathfrak r = R \cap \mathfrak s$，由引理
\ref{more-algebra-lemma-tor-independent-flat} 得到。因此 (2) 推出 (3)。
由于模的消失性可以通过在素理想处局部化来检验（Algebra，引理
\ref{algebra-lemma-characterize-zero-local}），可知 (3) 推出 (1)。
为证明 (1) $\Rightarrow$ (2)，再次使用引理
\ref{more-algebra-lemma-tor-independent-flat} 给出的等式
$$
\text{Tor}_i^{R_\mathfrak r}(A_\mathfrak p, B_\mathfrak q) =
\text{Tor}_i^R(A, B) \otimes_{(A \otimes_R B)}
(A_\mathfrak p \otimes_{R_{\mathfrak r}} B_\mathfrak q)
$$
\end{proof}








\section{Tor 的谱序列}
\label{more-algebra-section-spectral-sequence-tor}


\noindent
本节汇集考察 Tor 函子时出现的各种谱序列。

\begin{example}
\label{more-algebra-example-cohomology-complex-tensored}
设 $R$ 为环。设 $K_\bullet$ 为 $R$-模链复形，并且当 $n \ll 0$ 时
$K_n = 0$。设 $M$ 为 $R$-模。选取由自由 $R$-模组成的 $M$ 的解消
$P_\bullet \to M$。得到双链复形 $K_\bullet \otimes_R P_\bullet$。
应用 Homology，第 \ref{homology-section-double-complex} 节的材料
（特别是 Homology，引理 \ref{homology-lemma-first-quadrant-ss}），
并将其改写为链复形的语言，得到两个收敛到
$H_*(K_\bullet \otimes_R^\mathbf{L} M)$ 的谱序列。一方面，存在
$E_2$-页为
$$
(E_2)_{i, j} = \text{Tor}^R_j(H_i(K_\bullet), M)
\Rightarrow
H_{i + j}(K_\bullet \otimes^{\mathbf{L}}_R M)
$$
的谱序列，其微分 $d_2$ 由映射
$\text{Tor}^R_j(H_i(K_\bullet), M) \to
\text{Tor}^R_{j - 2}(H_{i + 1}(K_\bullet), M)$ 给出。
另一方面，存在 $E_1$-页为
$$
(E_1)_{i, j} = \text{Tor}^R_j(K_i, M)
\Rightarrow
H_{i + j}(K_\bullet \otimes^{\mathbf{L}}_R M)
$$
的谱序列，其微分 $d_1$ 由 $K_i \to K_{i - 1}$ 诱导的映射
$\text{Tor}^R_j(K_i, M) \to \text{Tor}^R_j(K_{i - 1}, M)$ 给出。
\end{example}

\begin{example}
\label{more-algebra-example-tor-change-rings}
设 $R \to S$ 为环同态。设 $M$ 为 $R$-模，$N$ 为 $S$-模。
那么存在谱序列
$$
\text{Tor}^S_n(\text{Tor}^R_m(M, S), N) \Rightarrow
\text{Tor}^R_{n + m}(M, N).
$$
为构造它，选取 $M$ 的一个 $R$-自由解消 $P^\bullet$。于是
$$
M \otimes_R^{\mathbf{L}} N = P^\bullet \otimes_R N =
(P^\bullet \otimes_R S) \otimes_S N
$$
然后应用例 \ref{more-algebra-example-cohomology-complex-tensored} 的第一个谱序列。
\end{example}

\begin{example}
\label{more-algebra-example-tor-base-change}
考虑交换图
$$
\xymatrix{
B \ar[r] & B' = B \otimes_A A' \\
A \ar[r] \ar[u] & A' \ar[u]
}
$$
以及 $B$-模 $M, N$。置 $M' = M \otimes_A A' = M \otimes_B B'$ 和
$N' = N \otimes_A A' = N \otimes_B B'$。
{\it 假设 $A \to B$ 平坦，且 $M$ 和 $N$ 都在 $A$ 上平坦。}
那么存在谱序列
$$
\text{Tor}^A_i(\text{Tor}_j^B(M, N), A')
\Rightarrow
\text{Tor}^{B'}_{i + j}(M', N')
$$
理由如下。选取一个由自由 $B$-模组成的解消 $F_\bullet \to M$。
由于 $B$ 和 $M$ 都在 $A$ 上平坦，可见
$F_\bullet \otimes_A A'$ 是 $M'$ 的自由 $B'$-解消。
因此，群 $\text{Tor}^{B'}_n(M', N')$ 由复形
$$
(F_\bullet \otimes_A A') \otimes_{B'} N' =
(F_\bullet \otimes_B N) \otimes_A A' =
(F_\bullet \otimes_B N) \otimes^{\mathbf{L}}_A A'
$$
计算。最后一个等式成立，是因为 $N$ 在 $A$ 上平坦，所以
$F_\bullet \otimes_B N$ 是由平坦 $A$-模组成的复形。
于是应用例 \ref{more-algebra-example-cohomology-complex-tensored} 的谱序列便得到所述谱序列。
\end{example}

\begin{example}
\label{more-algebra-example-tor}
设 $K^\bullet, L^\bullet$ 为 $D^{-}(R)$ 的对象。那么存在谱序列
$$
E_2^{p, q} = H^p(K^\bullet \otimes_R^{\mathbf{L}} H^q(L^\bullet))
\Rightarrow H^{p + q}(K^\bullet \otimes_R^{\mathbf{L}} L^\bullet)
$$
以及另一个谱序列
$$
E_2^{p, q} =
H^p(H^q(K^\bullet) \otimes_R^{\mathbf{L}} L^\bullet)
\Rightarrow H^{p + q}(K^\bullet \otimes_R^{\mathbf{L}} L^\bullet)
$$
两个谱序列都有
$d_2^{p, q} : E_2^{p, q} \to E_2^{p + 2, q - 1}$。
用有上界的投射模复形代替 $K^\bullet$ 和 $L^\bullet$ 后，
这些谱序列就是 Homology，第 \ref{homology-section-double-complex} 节
讨论的、用于计算 $\text{Tot}(K^\bullet \otimes L^\bullet)$
上同调的两个谱序列。
\end{example}









\section{乘积与 Tor}
\label{more-algebra-section-products-tor}

\noindent
乘积映射最简单的例子来自下列情形。假设
$K^\bullet, L^\bullet \in D(R)$。那么存在映射
\begin{equation}
\label{more-algebra-equation-simple-tor-product}
H^i(K^\bullet) \otimes_R H^j(L^\bullet)
\longrightarrow
H^{i + j}(K^\bullet \otimes_R^{\mathbf{L}} L^\bullet)
\end{equation}
为定义这些映射，可以假设 $K^\bullet, L^\bullet$ 之一是
$R$-模的 K-平坦复形（例如由自由或投射 $R$-模组成的有上界复形）。
在这种情况下，$K^\bullet \otimes_R^{\mathbf{L}} L^\bullet$
由复形 $\text{Tot}(K^\bullet \otimes_R L^\bullet)$ 表示；见第
\ref{more-algebra-section-derived-tensor-product} 节（或第
\ref{more-algebra-section-computing-tor} 节）。接着，假设
$\xi \in H^i(K^\bullet)$ 且 $\zeta \in H^j(L^\bullet)$。
选取代表 $\xi$ 和 $\zeta$ 的
$k \in \Ker(K^i \to K^{i + 1})$ 和
$l \in \Ker(L^j \to L^{j + 1})$。定义
$$
\xi \cup \zeta =
\text{下列对象的类： }k \otimes l\text{ 于 }
H^{i + j}(\text{Tot}(K^\bullet \otimes_R L^\bullet)).
$$
该定义有意义，因为全复形上的微分 $\text{d}$ 的公式（见 Homology，
定义 \ref{homology-definition-associated-simple-complex}）表明
$k \otimes l$ 是余闭元。此外，如果对某个
$k'' \in K^{i - 1}$ 有 $k' = d_K(k'')$，则由于 $l$ 是余闭元，
$k' \otimes l = \text{d}(k'' \otimes l)$。类似地，改变代表
$\zeta$ 的 $l$ 的选取不会改变 $k \otimes l$ 的类。
同样显然，$\cup$ 是双线性的，因而把
$H^i(K^\bullet) \otimes_R H^j(L^\bullet)$ 的一般元素按
$$
\sum \xi_i \otimes \zeta_i \longmapsto \sum \xi_i \cup \zeta_i
$$
映到 $H^{i + j}(\text{Tot}(K^\bullet \otimes_R L^\bullet))$ 中。

\medskip\noindent
设 $R \to A$ 为环同态。设 $K^\bullet, L^\bullet \in D(R)$。
那么在 $D(A)$ 中存在典范等同
\begin{equation}
\label{more-algebra-equation-pullback-derived-tensor-product}
(K^\bullet \otimes_R^{\mathbf{L}} A)
\otimes_A^{\mathbf{L}}
(L^\bullet \otimes_R^{\mathbf{L}} A)
=
(K^\bullet \otimes_R^{\mathbf{L}} L^\bullet) \otimes_R^{\mathbf{L}} A
\end{equation}
其构造如下。首先，在 $R$ 上选取 K-平坦解消
$P^\bullet \to K^\bullet$ 和 $Q^\bullet \to L^\bullet$。
那么左侧由复形
$\text{Tot}((P^\bullet \otimes_R A) \otimes_A (Q^\bullet \otimes_R A))$
表示，右侧由复形
$\text{Tot}(P^\bullet \otimes_R Q^\bullet) \otimes_R A$ 表示。
这些复形典范同构。因此，上述构造诱导有时很有用的乘积
$$
\text{Tor}^R_n(K^\bullet, A) \otimes_A \text{Tor}^R_m(L^\bullet, A)
\longrightarrow
\text{Tor}_{n + m}^R(K^\bullet \otimes_R^\mathbf{L} L^\bullet, A)
$$

\medskip\noindent
设 $M$、$N$ 为 $R$-模。使用上述一般构造、典范映射
$M \otimes_R^\mathbf{L} N \to M \otimes_R N$ 以及
$\text{Tor}$ 的函子性，得到典范映射
\begin{equation}
\label{more-algebra-equation-tor-product}
\text{Tor}^R_n(M, A) \otimes_A \text{Tor}^R_m(N, A)
\longrightarrow \text{Tor}_{n + m}^R(M \otimes_R N, A)
\end{equation}
下面给出使用投射解消的直接构造。首先，在 $R$ 上选取投射解消
$$
P_\bullet \to M, \quad Q_\bullet \to N, \quad T_\bullet \to M \otimes_R N
$$
由 $\otimes_R$ 的右正合性，有
$H_0(\text{Tot}(P_\bullet \otimes_R Q_\bullet)) = M \otimes_R N$。
因此，Derived Categories，引理
\ref{derived-lemma-morphisms-lift-projective} 和
\ref{derived-lemma-morphisms-equal-up-to-homotopy-projective}
保证存在唯一的复形映射
$\mu : \text{Tot}(P_\bullet \otimes_R Q_\bullet) \to T_\bullet$，使得
$H_0(\mu) = \text{id}_{M \otimes_R N}$。这诱导 $D(A)$ 中的典范映射
\begin{align*}
(M \otimes_R^{\mathbf{L}} A) \otimes_A^{\mathbf{L}}
(N \otimes_R^{\mathbf{L}} A)
& =
\text{Tot}((P_\bullet \otimes_R A) \otimes_A (Q_\bullet \otimes_R A)) \\
& =
\text{Tot}(P_\bullet \otimes_R Q_\bullet) \otimes_R A \\
& \to
T_\bullet \otimes_R A \\
& = (M \otimes_R N) \otimes_R^{\mathbf{L}} A
\end{align*}
因此，为构造上面的乘积（\ref{more-algebra-equation-tor-product}），先在 $A$ 上使用
（\ref{more-algebra-equation-simple-tor-product}）构造
$$
\text{Tor}^R_n(M, A) \otimes_A \text{Tor}^R_m(N, A) \to
H^{-n-m}((M \otimes_R^{\mathbf{L}} A) \otimes_A^{\mathbf{L}}
(N \otimes_R^{\mathbf{L}} A))
$$
再与上面的显示映射复合，最终落入
$\text{Tor}_{n + m}^R(M \otimes_R N, A)$。

\medskip\noindent
上述构造的一个有趣特例是 $M = N = B$，其中 $B$ 是 $R$-代数。
在这种情况下，得到映射
$$
\text{Tor}_n^R(B, A) \otimes_A \text{Tor}_m^R(B, A)
\longrightarrow
\text{Tor}_{n + m}^R(B \otimes_R B, A)
\longrightarrow
\text{Tor}_{n + m}^R(B, A)
$$
第二条箭头由乘法映射 $B \otimes_R B \to B$ 通过
$\text{Tor}$ 的函子性诱导。换言之，在
$\text{Tor}^R_{\star}(B, A)$ 上得到一个 $A$-代数结构。
这个代数结构具有许多有趣的性质（结合性、分次交换性、
$B$-代数结构、在某些情形下的除幂等），我们将在别处讨论这些性质
（在此插入将来的引用）。

\begin{lemma}
\label{more-algebra-lemma-functoriality-product-tor}
设 $R$ 为环。设 $A, B, C$ 为 $R$-代数，并设 $B \to C$ 为
$R$-代数同态。那么诱导映射
$$
\text{Tor}^R_{\star}(B, A)
\longrightarrow
\text{Tor}^R_{\star}(C, A)
$$
是 $A$-代数同态。
\end{lemma}

\begin{proof}
略。提示：把所有复形显式写出并逐一定义即可证明。
\end{proof}








\section{K"unneth 谱序列}
\label{more-algebra-section-kunneth-spectral-sequence}

\noindent
设 $R$ 为环。设 $K^\bullet$ 和 $L^\bullet$ 为 $R$-模的滤过复形
（见 Homology，定义 \ref{homology-definition-filtered-complex}；
注意这里的滤过递减）。那么复形
$$
T^\bullet = \text{Tot}(K^\bullet \otimes_R L^\bullet)
$$
也有由公式
$$
F^nT^\bullet =
\Im\left(
\bigoplus\nolimits_{i + j = n}
\text{Tot}(F^iK^\bullet \otimes_R F^jL^\bullet) \to
\text{Tot}(K^\bullet \otimes_R L^\bullet)
\right)
$$
定义的递减滤过。在对这些滤过复形作出一些假设后，我们将确定由
Homology，第 \ref{homology-section-filtered-complex} 节的构造所得的谱序列。

\medskip\noindent
假设每个 $K^n$、$F^iK^n$、$\text{gr}^iK^n$、$L^m$、
$F^jL^m$ 和 $\text{gr}^jL^m$ 都是平坦 $R$-模。在这种情况下，模
$F^iK^n \otimes_R L^m$、$K^n \otimes_R F^jL^m$ 和
$F^iK^n \otimes_R F^jL^m$ 都是 $K^n \otimes_R L^m$ 的子模。
类似地，模 $\text{gr}^iK^n \otimes_R \text{gr}^jL^m$ 是
$K^n/F^{i + 1}K^n \otimes_R L^m/F^{j + 1}L^m$ 的子模。考虑映射
$$
F^nT^n \longrightarrow
\bigoplus\nolimits_{i + j = n}
\text{Tot}(K^\bullet / F^{i + 1}K^\bullet
\otimes_R L^\bullet / F^{j + 1}L^\bullet)
$$
由 $F^nT^\bullet$ 的定义，最终所得确为直和而非直积；这一点留给读者证明。
对于 $a + b = n$，上述箭头在 $F^nT^\bullet$ 的子复形
$\text{Tot}(F^aK^\bullet \otimes_R F^bL^\bullet)$ 上的限制映入
$i = a$ 且 $j = b$ 的直和项。此外，由平坦性假设，该像同构于
$\text{Tot}(\text{gr}^iK^\bullet \otimes_R \text{gr}^jL^\bullet)$，
而映射
$\text{Tot}(F^iK^\bullet \otimes_R F^jL^\bullet) \to
\text{Tot}(\text{gr}^iK^\bullet \otimes_R \text{gr}^jL^\bullet)$ 的核是
$\text{Tot}(F^{i + 1}K^\bullet \otimes_R F^jL^\bullet) +
\text{Tot}(F^iK^\bullet \otimes_R F^{j + 1}L^\bullet)$。
因此，在这些假设下有复形的短正合序列
$$
0 \to F^{n + 1}T^\bullet \to F^nT^\bullet \to
\bigoplus\nolimits_{i + j = n}
\text{Tot}(\text{gr}^iK^\bullet \otimes_R \text{gr}^jL^\bullet) \to 0
$$
换言之，这说明 $\text{gr}^nT^\bullet$ 是复形
$\text{Tot}(\text{gr}^iK^\bullet \otimes_R \text{gr}^jL^\bullet)$
在 $i + j = n$ 上的直和。

\medskip\noindent
进一步假设 $R$-模复形 $K^\bullet$、$F^iK^\bullet$、
$\text{gr}^iK^\bullet$、$L^\bullet$、$F^jL^\bullet$ 和
$\text{gr}^jL^\bullet$ 都 K-平坦。在这种情况下，Homology，
第 \ref{homology-section-filtered-complex} 节中与滤过复形 $T^\bullet$
相伴的谱序列具有项
$$
E_1^{p, q} =
\bigoplus\nolimits_{i + j = p}
H^{p + q}(\text{gr}^iK^\bullet \otimes_R^\mathbf{L} \text{gr}^jL^\bullet)
$$
其微分由短正合序列
$$
0 \to \text{gr}^{n + 1}T^\bullet \to
F^nT^\bullet/F^{n + 2}T^\bullet \to
\text{gr}^nT^\bullet \to 0
$$
诱导，解释见 Homology，引理
\ref{homology-lemma-spectral-sequence-filtered-complex-d1}。
特别地，读者可以证明 $E_1^{p, q}$ 的直和项
$H^{p + q}(\text{gr}^iK^\bullet \otimes_R^\mathbf{L} \text{gr}^jL^\bullet)$
映入 $E_1^{p + 1, q}$ 中的和
$$
H^{p + q + 1}(\text{gr}^{i + 1}K^\bullet \otimes_R^\mathbf{L}
\text{gr}^jL^\bullet) \oplus
H^{p + q + 1}(\text{gr}^iK^\bullet \otimes_R^\mathbf{L}
\text{gr}^{j + 1}L^\bullet)
$$

\begin{lemma}
\label{more-algebra-lemma-kunneth-converges}
在上述假设下，如果还满足
\begin{enumerate}
\item $K^\bullet$ 上的滤过有限，且 $L^\bullet$ 上的滤过有限；或者
\item 更一般地，下列条件成立：
\begin{enumerate}
\item 当 $i \gg 0$ 时，$F^iK^\bullet$ 无环；
\item 当 $i \ll 0$ 时，$F^iK^\bullet \to K^\bullet$ 是拟同构；
\item 当 $j \gg 0$ 时，$F^jL^\bullet$ 无环；并且
\item 当 $j \ll 0$ 时，$F^jL^\bullet \to L^\bullet$ 是拟同构。
\end{enumerate}
\end{enumerate}
那么该谱序列有界，每个
$H^n(T^\bullet) = H^n(K^\bullet \otimes_R^\mathbf{L} L^\bullet)$
上的伴随滤过有限，并且有收敛
$$
E_1^{p, q} =
\bigoplus\nolimits_{i + j = p}
H^{p + q}(\text{gr}^iK^\bullet \otimes_R^\mathbf{L} \text{gr}^jL^\bullet)
\Rightarrow
H^n(K^\bullet \otimes_R^\mathbf{L} L^\bullet)
$$
\end{lemma}

\begin{proof}
在情形 (1) 中，$T^\bullet$ 上的滤过有限，引理立即由 Homology，
引理 \ref{homology-lemma-biregular-ss-converges} 得到。
在情形 (2) 中，选取 $a < b$，使得当 $i, j > b$ 时
$F^iK^\bullet$ 和 $F^jL^\bullet$ 无环，并且当 $i, j < a$ 时
$F^iK^\bullet \to K^\bullet$ 和 $F^jL^\bullet \to L^\bullet$
是拟同构。我们断言，在这种情况下，当 $n > 2b$ 时复形
$F^nT^\bullet$ 无环，而当 $n < 2a - 1$ 时
$F^nT^\bullet \to T^\bullet$ 是拟同构。该断言表明可以应用
Homology，引理 \ref{homology-lemma-ss-converges-trivial}，故本引理成立。

\medskip\noindent
证明该断言。给定 $n$ 和整数 $i_1 \leq i_2$，考虑
$$
S^{n, i_1, i_2} =
\Im\left(
\bigoplus\nolimits_{i_1 \leq i \leq i_2}
\text{Tot}(F^iK^\bullet \otimes_R F^{n - i}L^\bullet)
\to
\text{Tot}(K^\bullet \otimes_R L^\bullet)
\right)
$$
把它看作 $\text{Tot}(K^\bullet \otimes_R L^\bullet)$ 的子复形。
注意 $F^nT^\bullet = \colim S^{n, i_1, i_2}$ 是滤过余极限。
因此，为证明该断言，只需证明：当 $n > 2b$ 时
$S^{n, i_1, i_2}$ 无环；并且当 $n < 2a - 1$ 时，对一组余终的
二元组 $i_1, i_2$，映射
$S^{n, i_1, i_2} \to \text{Tot}(K^\bullet \otimes L^\bullet)$
是拟同构。

\medskip\noindent
如果 $i_1 = i_2 = i$ 且 $n > 2b$，那么
$$
S^{n, i, i} =
\text{Tot}(F^iK^\bullet \otimes_R F^{n - i}L^\bullet)
$$
无环：如果 $i > b$，则由 $F^{n - i}L^\bullet$ K-平坦的假设，
该复形无环；如果 $i \leq b$，则
$n - i \geq n - b > b$，同理可得结论。使用平坦性假设，
读者可以证明存在复形的短正合序列
$$
0 \to
S^{n + 1, i_2 + 1, i_2 + 1}
\to
S^{n, i_1, i_2} \oplus S^{n, i_2 + 1, i_2 + 1}
\to
S^{n, i_1, i_2 + 1}
\to
0
$$
因此，对 $i_2 - i_1$ 归纳可见，当 $n > 2b$ 时所有复形
$S^{n, i_1, i_2}$ 都无环。

\medskip\noindent
假设 $n < 2a - 1$。注意
$$
S^{n, a - 1, a - 1} =
\text{Tot}(F^{a - 1}K^\bullet \otimes_R F^{n - a + 1}L^\bullet)
\to \text{Tot}(K^\bullet \otimes_R L^\bullet)
$$
是拟同构，因为 $a - 1 < a$ 且 $n - a + 1 < a$，并且所有复形都
K-平坦，所以等号两侧计算 $D(R)$ 的同一个对象。注意当
$i_2 \geq a - 1$ 且 $n < 2a - 1$ 时，映射
$$
S^{n + 1, i_2 + 1, i_2 + 1} \to S^{n, i_2 + 1, i_2 + 1}
$$
是拟同构，因为左侧和右侧都拟同构于
$\text{Tot}(F^{i_2 + 1}K^\bullet \otimes L^\bullet)$。
因此，使用上一段的复形短正合序列可知，对所有
$t \geq a - 1$，映射
$S^{n, a - 1, a - 1} \to S^{n, a - 1, t}$ 是拟同构。
另一方面，类似地存在短正合序列
$$
0 \to
S^{n + 1, i_1, i_1}
\to
S^{n, i_1 - 1, i_1 - 1} \oplus
S^{n, i_1, i_2}
\to
S^{n, i_1 - 1, i_2}
\to
0
$$
同上可知，当 $i_1 \leq a - 1$ 时，映射
$$
S^{n + 1, i_1, i_1} \to S^{n, i_1 - 1, i_1 - 1}
$$
是拟同构，因为左侧和右侧都拟同构于
$\text{Tot}(K^\bullet \otimes F^{n - i_1 + 1}L^\bullet)$。
因此，对所有 $s \leq a - 1$，映射
$S^{n, s, t} \to S^{n, a - 1, t}$ 是拟同构。综合可知，当
$s \leq a - 1$ 且 $t \geq a - 1$ 时，$S^{n, s, t}$
拟同构地映到 $\text{Tot}(K^\bullet \otimes_R L^\bullet)$。
这完成了引理的证明。
\end{proof}

\begin{lemma}
\label{more-algebra-lemma-resolution-filtered-complex}
设 $R$ 为环。设 $K^\bullet$ 为滤过复形。存在滤过复形的映射
$f : P^\bullet \to K^\bullet$，使得
\begin{enumerate}
\item 每个 $P^n$、$F^iP^n$、$\text{gr}^iP^n$ 都是自由 $R$-模；
\item $R$-模复形 $P^\bullet$、$F^iP^\bullet$ 和
$\text{gr}^iP^\bullet$ 都 K-平坦；
\item $f$ 诱导拟同构
$P^\bullet \to K^\bullet$、
$F^iP^\bullet \to F^iK^\bullet$ 和
$\text{gr}^iP^\bullet \to \text{gr}^iK^\bullet$。
\end{enumerate}
\end{lemma}

\begin{proof}
如果滤过复形 $L^\bullet$ 的每个 $L^n$、$F^iL^n$、
$\text{gr}^iL^n$ 都是自由 $R$-模，并且所有微分都为零，
则称该滤过复形 {\it 基本}。

\medskip\noindent
存在基本滤过复形 $P_0^\bullet$ 以及滤过复形的映射
$f_0 : P_0^\bullet \to K^\bullet$，使得 $f_0$ 和
$F^if_0$（$i \in \mathbf{Z}$）在上同调上都是满射。为看出这一点，置
$$
P_0^n = \bigoplus\nolimits_{z \in \Ker(d_K^n)} R \xi_z
\oplus
\bigoplus\nolimits_{j \in \mathbf{Z}}
\bigoplus\nolimits_{z \in \Ker(d_{F^jK}^n)} R \xi_z
$$
取零微分，并用 $f_0(\xi_z) = z$ 定义映射 $f_0$。至于滤过，置
$$
F^iP_0^n = \bigoplus\nolimits_{j \geq i \in \mathbf{Z}}
\bigoplus\nolimits_{z \in \Ker(d_{F^jK}^n)} R \xi_z
$$
所述构造确实给出断言中的 $f_0 : P_0^\bullet \to K^\bullet$；
验证留给读者。

\medskip\noindent
我们将对 $m \geq 0$ 归纳构造滤过复形的嵌入
$$
P_0^\bullet \subset \ldots P_m^\bullet \subset P_{m + 1}^\bullet
$$
以及映射 $f_m : P_m^\bullet \to K^\bullet$，使其具有下列性质：
\begin{enumerate}
\item 滤过复形 $P_{m + 1}^\bullet / P_m^\bullet$ 是基本的；
\item $H^n(f_m)$ 的核与 $H^n(F^if_m)$ 的核在
$H^n(P_{m + 1}^\bullet)$ 和 $H^n(F^iP_{m + 1}^\bullet)$ 中都映为零；
\item 映射 $f_{m + 1} : P_{m + 1}^\bullet \to K^\bullet$
延拓映射 $f_m$。
\end{enumerate}
为此，置
$$
\Omega_{n, m} = \Ker\left(\Ker(d^n_{P_m}) \to H^n(K^\bullet)\right),
\text{ 分别 }
\Omega_{n, i, m} = \Ker\left(\Ker(d^n_{F^iP_m}) \to H^n(F^iK^\bullet)\right),
$$
注意 $\Omega_{n, m}$ 满射到 $H^n(f_m)$ 的核，而
$\Omega_{n, i, m}$ 满射到 $H^n(F^if_n)$ 的核。
对每个 $z \in \Omega_{n, m}$（相应地，$z \in \Omega_{n, i, m}$），
选取 $y_z \in K^{n - 1}$（相应地，$y_z \in F^iK^{n - 1}$），
使得 $d_K(y_z) = f_m(z)$。然后置
$$
P^n_{m + 1} = P_m^n \oplus
\bigoplus\nolimits_{z \in \Omega_{n + 1, m}} R \eta_z
\oplus
\bigoplus\nolimits_{j \in \mathbf{Z}}
\bigoplus\nolimits_{z \in \Omega_{n + 1, j, m}} R \eta_z
$$
置 $d_{P_m}(\eta_z) = z$ 并置 $f_{m + 1}(\eta_z) = y_z$。
最后，置
$$
F^iP_{m + 1}^n = F^iP_m^n \oplus
\bigoplus\nolimits_{j \geq i}
\bigoplus\nolimits_{z \in \Omega_{n + 1, j, m}} R \eta_z
$$
所述构造确实给出断言中的 $P_m^\bullet \subset P_{m + 1}^\bullet$
和 $f_{m + 1} : P_{m + 1}^\bullet \to K^\bullet$；验证留给读者。

\medskip\noindent
此时只需取
$$
P^\bullet = \bigcup P_m^\bullet
$$
作为滤过复形，并取由 $\bigcup f_m$ 给出的映射
$f : P^\bullet \to K^\bullet$。

\medskip\noindent
引理陈述的 (1) 成立，因为作为滤过模的 $P^n$ 同构于
$P_0^n$ 与所有 $m \geq 0$ 的 $P_{m + 1}^n/P_m^n$ 的直和。
略去一个小细节。

\medskip\noindent
考虑陈述的 (3)。注意
$H^n(P^\bullet) = \colim H^n(P_m^\bullet)$。映射 $f_0$
在上同调上满射；因此每个 $f_m$ 也如此。根据构造，对每个 $m$，
嵌入 $P_m^\bullet \subset P_{m + 1}^\bullet$ 杀掉
$H^n(f_m)$ 的核。综合这些事实，读者容易推出 $H^n(f)$ 是同构。
$H^n(F^if_n)$ 的情形类似。于是 $H^n(\text{gr}^if)$ 也必为同构，
因为有短正合序列 $0 \to F^{i + 1} \to F^i \to \text{gr}^i \to 0$
（例如，把它看作滤过复形范畴上的函子序列）。略去一个小细节。

\medskip\noindent
考虑陈述的 (2)。为看出 $P^\bullet$ K-平坦，由引理
\ref{more-algebra-lemma-colimit-K-flat}，只需证明 $P_m^\bullet$ K-平坦。
由引理 \ref{more-algebra-lemma-K-flat-two-out-of-three-ses} 和归纳，只需注意微分为零、
各项自由的复形 K-平坦。相同论证说明，对所有
$i \in \mathbf{Z}$，$F^iP^\bullet$ 都 K-平坦。最后，再次应用引理
\ref{more-algebra-lemma-K-flat-two-out-of-three-ses} 可知
$\text{gr}^iP^\bullet$ K-平坦。
\end{proof}

\begin{proposition}
\label{more-algebra-proposition-kunneth-filtered}
设 $R$ 为环。设 $K^\bullet$ 和 $L^\bullet$ 为 $R$-模的滤过复形。
那么存在表示 $D(R)$ 中 $K^\bullet \otimes_R^\mathbf{L} L^\bullet$
的滤过复形 $T^\bullet$，使得相伴谱序列的 $E_1$-页为
$$
E_1^{p, q} =
\bigoplus\nolimits_{i + j = p}
H^{p + q}(\text{gr}^iK^\bullet \otimes_R^\mathbf{L} \text{gr}^jL^\bullet)
$$
如果
\begin{enumerate}
\item 当 $i \gg 0$ 时，$F^iK^\bullet$ 无环；
\item 当 $i \ll 0$ 时，$F^iK^\bullet \to K^\bullet$ 是拟同构；
\item 当 $j \gg 0$ 时，$F^jL^\bullet$ 无环；并且
\item 当 $j \ll 0$ 时，$F^jL^\bullet \to L^\bullet$ 是拟同构；
\end{enumerate}
则该谱序列有界，每个
$H^n(K^\bullet \otimes_R^\mathbf{L} L^\bullet)$ 上的伴随滤过有限，
并且该谱序列收敛。
\end{proposition}

\begin{proof}
选取引理 \ref{more-algebra-lemma-resolution-filtered-complex} 中的
$P^\bullet \to K^\bullet$ 和 $Q^\bullet \to L^\bullet$。
然后使用本节正文和引理 \ref{more-algebra-lemma-kunneth-converges} 所描述的、
滤过复形
$$
T^\bullet = \text{Tot}(P^\bullet \otimes_R Q^\bullet)
$$
的谱序列。
\end{proof}

\begin{lemma}[K"unneth 谱序列]
\label{more-algebra-lemma-kunneth-spectral-sequence}
设 $R$ 为环。设 $K$ 和 $L$ 为 $D^b(R)$ 的对象。
存在双分次有界谱序列 $\{E_r\}_{r \geq 2}$，其中
$$
E_2^{p, q} = \bigoplus\nolimits_{i + j = q} \text{Tor}^R_{-p}(H^i(K), H^j(L))
$$
且 $d_r$ 的双次数为 $(r, -r + 1)$，该谱序列收敛到
$H^{p + q}(K \otimes_R^\mathbf{L} L)$。
\end{lemma}

\begin{proof}
设 $K^\bullet$ 为表示 $K$ 的 $R$-模复形，设 $L^\bullet$
为表示 $L$ 的 $R$-模复形。置
$F^iK^\bullet = \tau_{\leq -i}K^\bullet$，并对 $L$ 作类似定义；
见 Homology，第 \ref{homology-section-truncations} 节。
应用命题 \ref{more-algebra-proposition-kunneth-filtered}，并注意
$\text{gr}^iK^\bullet = H^{-i}(K)[i]$，得到有界谱序列
$\{(E')_r\}_{r \geq 1}$，其
$$
(E')^{p, q}_1 =
\bigoplus\nolimits_{i + j = p} \text{Tor}^R_{-2p - q}(H^{-i}(K), H^{-j}(L))
$$
收敛到 $K \otimes_R^\mathbf{L} L$ 的上同调。按构造，该谱序列具有微分
$d_r^{p, q} : (E')_r^{p, q} \to (E')_r^{p + r, q - r + 1}$，
其中 $r \geq 1$。为得到引理中的谱序列，对 $r \geq 2$ 置
$$
E_r^{p, q} = (E')_{r - 1}^{-q, p + 2q}
$$
所得构造确实可行；验证留给读者。
\end{proof}

\begin{example}
\label{more-algebra-example-kunneth-dedekind}
如果 $R = \mathbf{Z}$，或更一般地，如果 $R$ 是 Dedekind 整环，
那么对所有 $R$-模 $M$ 和 $N$，当 $i \not \in \{0, 1\}$ 时都有
$\text{Tor}_i^R(M, N) = 0$。因此，引理
\ref{more-algebra-lemma-kunneth-spectral-sequence} 的谱序列在 $E_2$ 页退化，
并且对所有 $n \in \mathbf{Z}$ 都得到短正合序列
$$
0 \to \bigoplus_{i + j = n} H^i(K) \otimes_R H^j(L) \to
H^n(K \otimes_R^\mathbf{L} L) \to
\bigoplus_{i + j = n + 1} \text{Tor}^R_1(H^i(K), H^j(L)) \to 0
$$
\end{example}






\section{拟相干模，I}
\label{more-algebra-section-pseudo-coherent}

\noindent
设 $R$ 为环。回顾一下，如果存在满射 $R^{\oplus a} \to M$，
则称 $R$-模 $M$ 为有限型；如果存在表示
$R^{\oplus a_1} \to R^{\oplus a_0} \to M \to 0$，
则称其为有限表示。类似地，可以考虑存在长度为 $n$ 的解消的
$R$-模：
\begin{equation}
\label{more-algebra-equation-pseudo-coherent}
R^{\oplus a_n} \to R^{\oplus a_{n - 1}} \to \ldots \to R^{\oplus a_0} \to
M \to 0
\end{equation}
其中各项是有限自由 $R$-模。如果对每个 $n$ 都能找到这样的解消，
则称模是{\it 拟相干的}。下面给出形式定义。

\begin{definition}
\label{more-algebra-definition-pseudo-coherent}
设 $R$ 为环。记其导出范畴为 $D(R)$。设 $m \in \mathbf{Z}$。
\begin{enumerate}
\item 如果存在有限自由 $R$-模组成的有界复形 $E^\bullet$ 及映射
$\alpha : E^\bullet \to K^\bullet$，使得当 $i > m$ 时
$H^i(\alpha)$ 是同构、而 $H^m(\alpha)$ 是满射，则称 $D(R)$ 的对象
$K^\bullet$ 为{\it $m$-拟相干的}。
\item 如果对象 $K^\bullet$ 是 $D(R)$ 的对象且拟同构于有限自由 $R$-模组成的上有界复形，
则称其为{\it 拟相干的}。
\item 如果 $M[0]$ 是 $D(R)$ 中的 $m$-拟相干对象，则称 $R$-模 $M$
为{\it $m$-拟相干的}。
\item 如果 $M[0]$ 是 $D(R)$ 中的拟相干对象，则称 $R$-模 $M$ 为
{\it 拟相干的}\footnote{这与 \cite{Bourbaki-CA} 中对拟相干模的含义冲突。}。
\end{enumerate}
\end{definition}

\noindent
通常，我们也把这一术语用于 $R$-模复形。由于 $D(R)$ 中的任意映射
$E^\bullet \to K^\bullet$ 都由复形的实际映射表示（见 Derived Categories，
引理 \ref{derived-lemma-morphisms-from-projective-complex}），
所以没有歧义。
事实证明，$K^\bullet$ 拟相干当且仅当对所有 $m \in \mathbf{Z}$，
$K^\bullet$ 都是 $m$-拟相干的；见引理 \ref{more-algebra-lemma-pseudo-coherent}。
此外，如果环是 Noether 环，则该条件可以理解为上同调的有限生成条件；
见引理 \ref{more-algebra-lemma-Noetherian-pseudo-coherent}。先把它与上面的非正式讨论联系起来。

\begin{lemma}
\label{more-algebra-lemma-cone-pseudo-coherent}
设 $R$ 为环且 $m \in \mathbf{Z}$。设
$(K^\bullet, L^\bullet, M^\bullet, f, g, h)$ 为 $D(R)$ 中的可区别三角。
\begin{enumerate}
\item 如果 $K^\bullet$ 是 $(m + 1)$-拟相干的且 $L^\bullet$ 是
$m$-拟相干的，则 $M^\bullet$ 是 $m$-拟相干的。
\item 如果 $K^\bullet, M^\bullet$ 都是 $m$-拟相干的，则
$L^\bullet$ 是 $m$-拟相干的。
\item 如果 $L^\bullet$ 是 $(m + 1)$-拟相干的且 $M^\bullet$ 是
$m$-拟相干的，则 $K^\bullet$ 是 $(m + 1)$-拟相干的。
\end{enumerate}
\end{lemma}

\begin{proof}
证明 (1)。选取 $\alpha : P^\bullet \to K^\bullet$，其中
$P^\bullet$ 是有限自由模组成的有界复形，并且当 $i > m + 1$ 时
$H^i(\alpha)$ 是同构、在 $i = m + 1$ 时是满射。可以用
$\sigma_{\geq m + 1}P^\bullet$ 替代 $P^\bullet$，因而不妨设
$P^i = 0$（当 $i < m + 1$）。选取 $\beta : E^\bullet \to L^\bullet$，
其中 $E^\bullet$ 是有限自由模组成的有界复形，并且当 $i > m$ 时
$H^i(\beta)$ 是同构、在 $i = m$ 时是满射。由 Derived Categories，
引理 \ref{derived-lemma-lift-map}，可以找到映射
$\gamma : P^\bullet \to E^\bullet$，使图
$$
\xymatrix{
K^\bullet \ar[r] & L^\bullet \\
P^\bullet \ar[u] \ar[r]^\gamma & E^\bullet \ar[u]_\beta
}
$$
在 $D(R)$ 中交换。锥 $C(\gamma)^\bullet$ 是有限自由 $R$-模组成的有界复形，
且图的交换性意味着存在可区别三角的态射
$$
(P^\bullet, E^\bullet, C(\gamma)^\bullet)
\longrightarrow
(K^\bullet, L^\bullet, M^\bullet).
$$
由诱导的上同调长正合序列以及 Homology，引理
\ref{homology-lemma-four-lemma} 和 \ref{homology-lemma-five-lemma}，
可知 $C(\gamma)^\bullet \to M^\bullet$ 在次数 $> m$ 的上同调上诱导同构，
在次数 $m$ 上诱导满射。因此 $M^\bullet$ 是 $m$-拟相干的。

\medskip\noindent
(2) 和 (3) 由 (1) 旋转可区别三角得到。
\end{proof}

\begin{lemma}
\label{more-algebra-lemma-finite-cohomology}
设 $R$ 为环。设 $K^\bullet$ 为 $R$-模复形。设 $m \in \mathbf{Z}$。
\begin{enumerate}
\item 如果 $K^\bullet$ 是 $m$-拟相干的且当 $i > m$ 时
$H^i(K^\bullet) = 0$，则 $H^m(K^\bullet)$ 是有限型 $R$-模。
\item 如果 $K^\bullet$ 是 $m$-拟相干的且当 $i > m + 1$ 时
$H^i(K^\bullet) = 0$，则 $H^{m + 1}(K^\bullet)$ 是有限表示的
$R$-模。
\end{enumerate}
\end{lemma}

\begin{proof}
证明 (1)。选取有限投射 $R$-模组成的有界复形 $E^\bullet$ 以及映射
$\alpha : E^\bullet \to K^\bullet$，使其在次数 $> m$ 的上同调上诱导同构、
在次数 $m$ 上诱导满射。显然只需对 $E^\bullet$ 证明结论。
设 $n$ 为使 $E^n \not = 0$ 的最大整数。如果 $n = m$，结论显然。
如果 $n > m$，则因 $H^n(E^\bullet) = 0$，映射
$E^{n - 1} \to E^n$ 是满射。由于 $E^n$ 是有限投射的，可知
$E^{n - 1} = E' \oplus E^n$。因此只需对复形 $(E')^\bullet$ 证明，
它与 $E^\bullet$ 相同，只是第 $n - 1$ 次为 $E'$、第 $n$ 次为 $0$。
对 $n$ 归纳即可。

\medskip\noindent
证明 (2)。选取有限投射 $R$-模组成的有界复形 $E^\bullet$ 以及映射
$\alpha : E^\bullet \to K^\bullet$，使其在次数 $> m$ 的上同调上诱导同构、
在次数 $m$ 上诱导满射。与 (1) 的证明一样，可以化归到
$E^i = 0$（当 $i > m + 1$）的情形。于是
$$
H^{m + 1}(K^\bullet) \cong
H^{m + 1}(E^\bullet) = \Coker(E^m \to E^{m + 1})
$$
是有限表示的。
\end{proof}

\begin{lemma}
\label{more-algebra-lemma-n-pseudo-module}
设 $R$ 为环。设 $M$ 为 $R$-模。则
\begin{enumerate}
\item $M$ 是 $0$-拟相干的，当且仅当 $M$ 是有限 $R$-模；
\item $M$ 是 $(-1)$-拟相干的，当且仅当 $M$ 是有限表示的 $R$-模；
\item $M$ 是 $(-d)$-拟相干的，当且仅当存在长度为 $d$ 的解消
$$
R^{\oplus a_d} \to R^{\oplus a_{d - 1}} \to \ldots \to R^{\oplus a_0} \to
M \to 0
$$
并且
\item $M$ 是拟相干的，当且仅当存在由有限自由 $R$-模组成的无限解消
$$
\ldots \to R^{\oplus a_1} \to R^{\oplus a_0} \to M \to 0
$$
\end{enumerate}
\end{lemma}

\begin{proof}
如果 $M$ 是有限型（分别为有限表示），则由定义
\ref{more-algebra-definition-pseudo-coherent} 之前的讨论，$M$ 是 $0$-拟相干
（分别为 $(-1)$-拟相干）的。反之，如果 $M$ 是 $0$-拟相干的，
则 $M = H^0(M[0])$ 由引理 \ref{more-algebra-lemma-finite-cohomology} 为有限型。
如果 $M$ 是 $(-1)$-拟相干的，则它也是 $0$-拟相干的，因而为有限型。
选取满射 $R^{\oplus a} \to M$，记
$K = \Ker(R^{\oplus a} \to M)$。由引理 \ref{more-algebra-lemma-cone-pseudo-coherent}，
$K$ 是 $0$-拟相干的，因而为有限型，所以 $M$ 是有限表示的。

\medskip\noindent
第三、第四个断言用归纳和类似上述的论证即可证明（细节略）。
\end{proof}

\begin{lemma}
\label{more-algebra-lemma-pseudo-coherent}
设 $R$ 为环。设 $K^\bullet$ 为 $R$-模复形。下列条件等价：
\begin{enumerate}
\item $K^\bullet$ 是拟相干的；
\item 对每个 $m \in \mathbf{Z}$，$K^\bullet$ 都是 $m$-拟相干的；以及
\item $K^\bullet$ 拟同构于有限投射 $R$-模组成的上有界复形。
\end{enumerate}
如果 (1)、(2)、(3) 成立且当 $i > b$ 时
$H^i(K^\bullet) = 0$，则可以找到拟同构 $F^\bullet \to K^\bullet$，
其中 $F^i$ 是有限自由 $R$-模且当 $i > b$ 时 $F^i = 0$。
\end{lemma}

\begin{proof}
由于有限自由模是有限投射 $R$-模，故 (1) $\Rightarrow$ (3)。
反之，设 $P^\bullet$ 是有限投射 $R$-模组成的上有界复形。
设当 $i > n_0$ 时 $P^i = 0$。选取直和分解
$F^{n_0} = P^{n_0} \oplus C^{n_0}$，其中 $F^{n_0}$ 是有限自由
$R$-模，并归纳地对 $n \leq n_0$ 取
$$
F^{n - 1} = P^{n - 1} \oplus C^n \oplus C^{n - 1}
$$
其中 $F^{n_0}$ 是有限自由 $R$-模。令复形 $F^\bullet$ 的映射
$F^{n - 1} \to F^n$ 与 $P^{n - 1} \to P^n$ 一致，在
$C^n \to C^n$ 上诱导恒等映射，并且在 $C^{n - 1}$ 上为零。
映射 $F^\bullet \to P^\bullet$ 是拟同构（事实上是同伦等价），
所以 (3) 推出 (1)。

\medskip\noindent
假设 (1)。设 $E^\bullet$ 是有限自由 $R$-模组成的上有界复形，
且 $E^\bullet \to K^\bullet$ 是拟同构。由愚蠢截断得到的
诱导的映射 $\sigma_{\geq m}E^\bullet \to K^\bullet$ 由 $E^\bullet$ 到 $K^\bullet$ 的愚蠢截断得到，
可知 $K^\bullet$ 是 $m$-拟相干的。因此 (1) 推出 (2)。

\medskip\noindent
假设 (2)。由于 $K^\bullet$ 是 $0$-拟相干的，特别可知 $K^\bullet$ 上有界。
取整数 $b$ 使得当 $i > b$ 时 $H^i(K^\bullet) = 0$。我们将对
$n \in \mathbf{Z}$ 作下降归纳，构造当 $i \geq n$ 时的有限自由
$R$-模 $F^i$、微分 $d^i : F^i \to F^{i + 1}$（当 $i \geq n$）以及与微分相容的
映射 $\alpha : F^i \to K^i$，使得 (1) $H^i(\alpha)$ 在 $i > n$
时是同构、在 $i = n$ 时是满射；(2) 当 $i > b$ 时 $F^i = 0$。
图示为
$$
\xymatrix{
& F^n \ar[r] \ar[d]^\alpha & F^{n + 1} \ar[d]^\alpha \ar[r] & \ldots \\
K^{n - 1} \ar[r] & K^n \ar[r] & K^{n + 1} \ar[r] & \ldots
}
$$
基例为 $n = b + 1$，此时对所有 $i$ 可取 $F^i = 0$。

归纳步骤。令 $C^\bullet$ 为 $\alpha$ 的锥（Derived Categories，定义
\ref{derived-definition-cone}）。上同调长正合序列
$$
0 \to H^{n - 1}(K^\bullet) \to H^{n - 1}(C^\bullet) \to
H^n(F^\bullet) \to H^n(K^\bullet) \to H^n(C^\bullet) \to \ldots
$$
表明当 $i \geq n$ 时 $H^i(C^\bullet) = 0$。由引理
\ref{more-algebra-lemma-cone-pseudo-coherent} 可知 $C^\bullet$ 是
$(n - 1)$-拟相干的。由引理 \ref{more-algebra-lemma-finite-cohomology} 可知
$H^{n - 1}(C^\bullet)$ 是有限 $R$-模。
特别地，$H^n(F^\bullet) \to H^n(K^\bullet)$ 的核是有限 $R$-模。
选取有限自由 $R$-模 $F^{n - 1}$ 以及映射
$F^{n - 1} \to \Ker(F^n \to F^{n - 1})$，使 $F^{n - 1}$ 满射到
$H^n(F^\bullet) \to H^n(K^\bullet)$ 的核。
把复形延拓为
$$
\xymatrix{
& F^{n - 1} \ar[d]^{\alpha^{n - 1}} \ar[r] &
F^n \ar[r] \ar[d]^\alpha &
F^{n + 1} \ar[d]^\alpha \ar[r] & \ldots \\
\ldots \ar[r] &
K^{n - 1} \ar[r] & K^n \ar[r] & K^{n + 1} \ar[r] & \ldots
}
$$
注意 $\alpha^{n - 1}$ 存在，是因为由构造可知
$F^{n - 1} \to F^n \to K^n$ 的像包含于 $K^{n - 1} \to K^n$ 的像中。
再次记 $C^\bullet$ 为这个延拓的复形映射的锥。此时得到正合序列
$$
H^{n - 1}(F^\bullet) \to H^{n - 1}(K^\bullet) \to H^{n - 1}(C^\bullet) \to
H^n(F^\bullet) \cong H^n(K^\bullet) \to H^n(C^\bullet) \to \ldots
$$
换言之，$H^{n - 1}(F^\bullet) \to H^{n - 1}(K^\bullet)$ 的余核是有限
$R$-模，设其由 $\xi_1, \ldots, \xi_r \in \Ker(K^{n - 1} \to K^n)$
的类生成。然后把 $F^{n - 1}$ 替换为
$F^{n - 1} \oplus R^{\oplus r}$，使自由直和中的基元素映到 $F^n$ 中的零，
并映到 $K^{n - 1}$ 中的 $\xi_i$。这完成归纳步骤的证明。
\end{proof}

\begin{lemma}
\label{more-algebra-lemma-two-out-of-three-pseudo-coherent}
设 $R$ 为环。设 $(K^\bullet, L^\bullet, M^\bullet, f, g, h)$ 为
$D(R)$ 中的可区别三角。如果 $K^\bullet, L^\bullet, M^\bullet$ 中任意
两个拟相干，则第三个也拟相干。
\end{lemma}

\begin{proof}
结合引理 \ref{more-algebra-lemma-cone-pseudo-coherent} 和
\ref{more-algebra-lemma-pseudo-coherent} 即得。
\end{proof}

\begin{lemma}
\label{more-algebra-lemma-recognize-pseudo-coherent}
设 $R$ 为环。设 $K^\bullet$ 为 $R$-模复形。设 $m \in \mathbf{Z}$。
\begin{enumerate}
\item 如果对所有 $i \geq m$ 都有 $H^i(K^\bullet) = 0$，则
$K^\bullet$ 是 $m$-拟相干的。
\item 如果当 $i > m$ 时 $H^i(K^\bullet) = 0$，且 $H^m(K^\bullet)$
是有限 $R$-模，则 $K^\bullet$ 是 $m$-拟相干的。
\item 如果当 $i > m + 1$ 时 $H^i(K^\bullet) = 0$，且模
$H^{m + 1}(K^\bullet)$ 是有限表示的、$H^m(K^\bullet)$ 是有限型的，
则 $K^\bullet$ 是 $m$-拟相干的。
\end{enumerate}
\end{lemma}

\begin{proof}
只需证明 (3)。置 $M = H^{m + 1}(K^\bullet)$。注意
$\tau_{\geq m + 1}K^\bullet$ 拟同构于 $M[- m - 1]$。由引理
\ref{more-algebra-lemma-n-pseudo-module}，$M[- m - 1]$ 是 $m$-拟相干的。
由于有可区别三角
$$
(\tau_{\leq m}K^\bullet, K^\bullet, \tau_{\geq m + 1}K^\bullet)
$$
（Derived Categories，评注
\ref{derived-remark-truncation-distinguished-triangle}），由引理
\ref{more-algebra-lemma-cone-pseudo-coherent}，只需证明
$\tau_{\leq m}K^\bullet$ 拟相干。根据假设，
$H^m(\tau_{\leq m}K^\bullet)$ 是有限型 $R$-模。
因此可以找到有限自由 $R$-模 $E$ 及映射
$E \to \Ker(d_K^m)$，使复合
$E \to \Ker(d_K^m) \to H^m(\tau_{\leq m}K^\bullet)$ 满射。
于是 $E[-m] \to \tau_{\leq m}K^\bullet$ 见证
$\tau_{\leq m}K^\bullet$ 是 $m$-拟相干的事实。
\end{proof}

\begin{lemma}
\label{more-algebra-lemma-summands-pseudo-coherent}
设 $R$ 为环。设 $m \in \mathbf{Z}$。如果 $K^\bullet \oplus L^\bullet$
是 $m$-拟相干的（分别为拟相干的），则 $K^\bullet$ 和 $L^\bullet$
也是如此。
\end{lemma}

\begin{proof}
在本证明中省略上标 ${}^\bullet$。假设 $K \oplus L$ 是
$m$-拟相干的。显然 $K, L \in D^{-}(R)$。注意有可区别三角
$$
(K \oplus L, K \oplus L, L \oplus L[1]) =
(K, K, 0) \oplus (L, L, L \oplus L[1])
$$
见 Derived Categories，引理 \ref{derived-lemma-direct-sum-triangles}。
由引理 \ref{more-algebra-lemma-cone-pseudo-coherent}，$L \oplus L[1]$ 是
$m$-拟相干的。因此 $L[1] \oplus L[2]$ 也是 $m$-拟相干的。
归纳可知 $L[n] \oplus L[n + 1]$ 是 $m$-拟相干的。
由引理 \ref{more-algebra-lemma-recognize-pseudo-coherent}，当 $n$ 足够大时
$L[n]$ 是 $m$-拟相干的。然后利用可区别三角
$$
(L[n], L[n] \oplus L[n - 1], L[n - 1])
$$
向后工作，得到 $L[n], L[n - 1], \ldots, L$ 都是 $m$-拟相干的。
拟相干的情形由此及引理 \ref{more-algebra-lemma-pseudo-coherent} 得到。
\end{proof}

\begin{lemma}
\label{more-algebra-lemma-complex-pseudo-coherent-modules}
设 $R$ 为环。设 $m \in \mathbf{Z}$。设 $K^\bullet$ 为上有界的
$R$-模复形，且对所有 $i$，$K^i$ 是 $(m - i)$-拟相干的。
则 $K^\bullet$ 是 $m$-拟相干的。特别地，如果 $K^\bullet$ 是拟相干
$R$-模组成的上有界复形，则 $K^\bullet$ 拟相干。
\end{lemma}

\begin{proof}
可以用 $\sigma_{\geq m - 1}K^\bullet$ 替代 $K^\bullet$（例如），
因而不妨设 $K^\bullet$ 有界。由复形长度归纳，使用引理
\ref{more-algebra-lemma-cone-pseudo-coherent} 和愚蠢截断，复形 $K^\bullet$ 是
$m$-拟相干的，因为每个 $K^i[-i]$ 都是 $m$-拟相干的。
对于最后一个断言，只需证明对所有 $m \in \mathbf{Z}$，
$K^\bullet$ 都是 $m$-拟相干的；见引理 \ref{more-algebra-lemma-pseudo-coherent}。
这由第一部分得到。
\end{proof}

\begin{lemma}
\label{more-algebra-lemma-cohomology-pseudo-coherent}
设 $R$ 为环。设 $m \in \mathbf{Z}$。设 $K^\bullet \in D^{-}(R)$，使得对所有 $i$，
$H^i(K^\bullet)$ 都是 $(m - i)$-拟相干（分别为拟相干）的。
则 $K^\bullet$ 是 $m$-拟相干（分别为拟相干）的。
\end{lemma}

\begin{proof}
设 $K^\bullet$ 是 $D^{-}(R)$ 的对象，并且每个 $H^i(K^\bullet)$ 都是
$(m - i)$-拟相干的。令 $n$ 为使 $H^n(K^\bullet)$ 非零的最大整数。
我们对 $n$ 归纳证明引理。
如果 $n < m$，则由引理
\ref{more-algebra-lemma-recognize-pseudo-coherent}，$K^\bullet$ 是 $m$-拟相干的。
如果 $n \geq m$，则有可区别三角
$$
(\tau_{\leq n - 1}K^\bullet, K^\bullet, H^n(K^\bullet)[-n])
$$
（Derived Categories，评注
\ref{derived-remark-truncation-distinguished-triangle}）。由于按假设
$H^n(K^\bullet)[-n]$ 是 $m$-拟相干的，应用引理
\ref{more-algebra-lemma-cone-pseudo-coherent} 可知，只需证明
$\tau_{\leq n - 1}K^\bullet$ 是 $m$-拟相干的。由对 $n$ 的归纳即可得证。
（拟相干的情形由此及引理
\ref{more-algebra-lemma-pseudo-coherent} 得到。）
\end{proof}

\begin{lemma}
\label{more-algebra-lemma-finite-push-pseudo-coherent}
设 $A \to B$ 为环同态。假设 $B$ 作为 $A$-模是拟相干的。设 $K^\bullet$ 为
$B$-模复形。下列条件等价：
\begin{enumerate}
\item $K^\bullet$ 作为 $B$-模复形是 $m$-拟相干的；以及
\item $K^\bullet$ 作为 $A$-模复形是 $m$-拟相干的。
\end{enumerate}
拟相干性的情形也有同样的等价性。
\end{lemma}

\begin{proof}
假设 (1)。选取有限自由 $B$-模组成的有界复形 $E^\bullet$ 以及映射
$\alpha : E^\bullet \to K^\bullet$，使其在次数 $> m$ 的上同调上诱导同构、
在次数 $m$ 上诱导满射。考虑可区别三角
$(E^\bullet, K^\bullet, C(\alpha)^\bullet)$。由引理
\ref{more-algebra-lemma-recognize-pseudo-coherent}，$C(\alpha)^\bullet$ 作为 $A$-模复形是
$m$-拟相干的。因此只需证明 $E^\bullet$ 作为 $A$-模复形是拟相干的，
这由引理\ref{more-algebra-lemma-complex-pseudo-coherent-modules} 得到。
(1) $\Rightarrow$ (2) 的拟相干情形由此及引理
\ref{more-algebra-lemma-pseudo-coherent} 得到。

\medskip\noindent
假设 (2)。令 $n$ 为使 $H^n(K^\bullet) \not = 0$ 的最大整数。
我们对 $n - m$ 归纳证明 $K^\bullet$ 作为 $B$-模复形是 $m$-拟相干的。
$n < m$ 的情形由引理\ref{more-algebra-lemma-recognize-pseudo-coherent} 得到。
选取有限自由 $A$-模组成的有界复形 $E^\bullet$ 以及映射
$\alpha : E^\bullet \to K^\bullet$，使其在次数 $> m$ 的上同调上诱导同构、
在次数 $m$ 上诱导满射。考虑诱导的复形映射
$$
\alpha \otimes 1 : E^\bullet \otimes_A B \to K^\bullet.
$$
注意，由于按构造 $H^n(E) \to H^n(E^\bullet \otimes_A B) \to H^n(K^\bullet)$
是满射，并且由例\ref{more-algebra-example-tor} 的谱序列，当 $i > n$ 时
$H^i(E^\bullet \otimes_A B) = 0$，故 $C(\alpha \otimes 1)^\bullet$ 在次数
$\geq n$ 上无上同调。
另一方面，由于 $K^\bullet$ 和 $E^\bullet \otimes_A B$（见引理
\ref{more-algebra-lemma-complex-pseudo-coherent-modules}）都是如此，
故由引理\ref{more-algebra-lemma-cone-pseudo-coherent}，
$C(\alpha \otimes 1)^\bullet$ 作为 $A$-模复形是 $m$-拟相干的。
因此由归纳，$C(\alpha \otimes 1)^\bullet$ 作为 $B$-模复形是 $m$-拟相干的。
最后再次应用引理\ref{more-algebra-lemma-cone-pseudo-coherent}，可知 $K^\bullet$ 作为
$B$-模复形是 $m$-拟相干的（显然 $E^\bullet \otimes_A B$ 作为 $B$-模复形是
拟相干的）。(2) $\Rightarrow$ (1) 的拟相干情形由此及引理
\ref{more-algebra-lemma-pseudo-coherent} 得到。
\end{proof}

\begin{lemma}
\label{more-algebra-lemma-pull-pseudo-coherent}
设 $A \to B$ 为环同态。设 $K^\bullet$ 为 $A$-模的 $m$-拟相干（分别为拟相干）
复形。则 $K^\bullet \otimes_A^{\mathbf{L}} B$ 是 $B$-模的
$m$-拟相干（分别为拟相干）复形。
\end{lemma}

\begin{proof}
首先注意到引理的陈述有意义：$K^\bullet$ 上有界，因而
$K^\bullet \otimes_A^{\mathbf{L}} B$ 由式
（\ref{more-algebra-equation-derived-tensor-algebra}）定义。
在此基础上，选取有限自由 $A$-模组成的有界复形 $E^\bullet$ 以及
$\alpha : E^\bullet \to K^\bullet$，使得当 $i > m$ 时 $H^i(\alpha)$ 是同构、
当 $i = m$ 时是满射。于是锥 $C(\alpha)^\bullet$ 在次数
$\geq m$ 上无上同调。由于 $-\otimes_A^{\mathbf{L}} B$ 是正合函子，得到
复形的可区别三角
$$
(E^\bullet \otimes_A^{\mathbf{L}} B, K^\bullet \otimes_A^{\mathbf{L}} B,
C(\alpha)^\bullet \otimes_A^{\mathbf{L}} B)
$$
其对象都是 $B$-模复形。由 Derived Categories，引理
\ref{derived-lemma-negative-vanishing} 的对偶可知，当 $i \geq m$ 时
$H^i(C(\alpha)^\bullet \otimes_A^{\mathbf{L}} B) = 0$。
由于 $E^\bullet$ 是有限投射 $R$-模组成的复形，可知
$E^\bullet \otimes_A^{\mathbf{L}} B = E^\bullet \otimes_A B$，因而
$$
E^\bullet \otimes_A B
\longrightarrow
K^\bullet \otimes_A^{\mathbf{L}} B
$$
是 $B$-模复形之间的映射，见证了
$K^\bullet \otimes_A^{\mathbf{L}} B$ 是 $m$-拟相干的事实。
拟相干复形的情形由 $m$-拟相干复形的情形及引理
\ref{more-algebra-lemma-pseudo-coherent} 得到。
\end{proof}

\begin{lemma}
\label{more-algebra-lemma-flat-base-change-pseudo-coherent}
设 $A \to B$ 为平坦环同态。设 $M$ 为 $A$-模，且是 $m$-拟相干的
（分别为拟相干的）。则 $M \otimes_A B$ 是 $m$-拟相干的
（分别为拟相干的）$B$-模。
\end{lemma}

\begin{proof}
这是引理\ref{more-algebra-lemma-pull-pseudo-coherent} 以及
$M \otimes_A^{\mathbf{L}} B = M \otimes_A B$（因为 $B$ 在 $A$ 上平坦）这一事实的
直接推论。
\end{proof}

\noindent
下面的引理也可由更强的引理
\ref{more-algebra-lemma-flat-descent-pseudo-coherent} 得到。

\begin{lemma}
\label{more-algebra-lemma-glue-pseudo-coherent}
设 $R$ 为环。设 $f_1, \ldots, f_r \in R$ 为生成单位理想的元素。设
$m \in \mathbf{Z}$。设 $K^\bullet$ 为 $R$-模复形。如果对每个 $i$，复形
$K^\bullet \otimes_R R_{f_i}$ 是 $m$-拟相干的（分别为拟相干的），则
$K^\bullet$ 是 $m$-拟相干的（分别为拟相干的）。
\end{lemma}

\begin{proof}
以下不再特别说明，使用 $- \otimes_R R_{f_i}$ 是正合函子、因而有
$$
H^i(K^\bullet)_{f_i} =
H^i(K^\bullet) \otimes_R R_{f_i} = H^i(K^\bullet \otimes_R R_{f_i}).
$$
假设当 $i = 1, \ldots, r$ 时 $K^\bullet \otimes_R R_{f_i}$ 是 $m$-拟相干的。
令 $n \in \mathbf{Z}$ 为使某个 $i$ 的
$H^n(K^\bullet \otimes_R R_{f_i})$ 非零的最大整数。这特别推出当 $i > n$ 时
$H^i(K^\bullet) = 0$（并且 $H^n(K^\bullet) \not = 0$），见
Algebra，引理\ref{algebra-lemma-cover}。我们对 $n - m$ 归纳证明引理。
如果 $n < m$，则由引理\ref{more-algebra-lemma-recognize-pseudo-coherent}，引理成立。
如果 $n \geq m$，则对每个 $i$，由引理
\ref{more-algebra-lemma-finite-cohomology}，$H^n(K^\bullet)_{f_i}$ 是有限的
$R_{f_i}$-模。因此由 Algebra，引理\ref{algebra-lemma-cover}，
$H^n(K^\bullet)$ 是有限 $R$-模。选取有限自由 $R$-模 $E$ 以及满射
$E \to H^n(K^\bullet)$。由于 $E$ 是投射的，可将其提升为复形映射
$\alpha : E[-n] \to K^\bullet$。于是锥 $C(\alpha)^\bullet$ 在次数
$\geq n$ 上的上同调为零。另一方面，对每个 $i$，复形
$C(\alpha)^\bullet \otimes_R R_{f_i}$ 是 $m$-拟相干的，见引理
\ref{more-algebra-lemma-cone-pseudo-coherent}。因此由归纳，$C(\alpha)^\bullet$ 作为
$R$-模复形是 $m$-拟相干的。再次应用引理
\ref{more-algebra-lemma-cone-pseudo-coherent} 即得结论。
\end{proof}

\begin{lemma}
\label{more-algebra-lemma-flat-descent-pseudo-coherent}
设 $R$ 为环。设 $m \in \mathbf{Z}$。设 $K^\bullet$ 为 $R$-模复形。设
$R \to R'$ 为忠实平坦的环同态。如果复形
$K^\bullet \otimes_R R'$ 是 $m$-拟相干的（分别为拟相干的），则
$K^\bullet$ 是 $m$-拟相干的（分别为拟相干的）。
\end{lemma}

\begin{proof}
以下不再特别说明，使用 $- \otimes_R R'$ 是正合函子，因而有
$$
H^i(K^\bullet) \otimes_R R' = H^i(K^\bullet \otimes_R R').
$$
假设 $K^\bullet \otimes_R R'$ 是 $m$-拟相干的。令
$n \in \mathbf{Z}$ 为使 $H^n(K^\bullet)$ 非零的最大整数；于是 $n$ 也是使
$H^n(K^\bullet \otimes_R R')$ 非零的最大整数。
我们对 $n - m$ 归纳证明引理。如果 $n < m$，则由引理
\ref{more-algebra-lemma-recognize-pseudo-coherent}，引理成立。
如果 $n \geq m$，则 $H^n(K^\bullet) \otimes_R R'$ 是有限
$R'$-模，见引理\ref{more-algebra-lemma-finite-cohomology}。
因此由 Algebra，引理\ref{algebra-lemma-descend-properties-modules}，
$H^n(K^\bullet)$ 是有限 $R$-模。选取有限自由 $R$-模 $E$ 以及满射
$E \to H^n(K^\bullet)$。由于 $E$ 是投射的，可将其提升为复形映射
$\alpha : E[-n] \to K^\bullet$。于是锥 $C(\alpha)^\bullet$ 在次数
$\geq n$ 上的上同调为零。另一方面，复形
$C(\alpha)^\bullet \otimes_R R'$ 是 $m$-拟相干的，见引理
\ref{more-algebra-lemma-cone-pseudo-coherent}。
因此由归纳，$C(\alpha)^\bullet$ 作为 $R$-模复形是 $m$-拟相干的。
再次应用引理\ref{more-algebra-lemma-cone-pseudo-coherent} 即得结论。
\end{proof}

\begin{lemma}
\label{more-algebra-lemma-tensor-pseudo-coherent}
设 $R$ 为环。设 $K, L$ 为 $D(R)$ 中的对象。
\begin{enumerate}
\item 如果 $K$ 是 $n$-拟相干的且当 $i > a$ 时 $H^i(K) = 0$，并且
$L$ 是 $m$-拟相干的且当 $j > b$ 时 $H^j(L) = 0$，则
$K \otimes_R^\mathbf{L} L$ 是 $t$-拟相干的，其中
$t = \max(m + a, n + b)$。
\item 如果 $K$ 和 $L$ 都是拟相干的，则
$K \otimes_R^\mathbf{L} L$ 是拟相干的。
\end{enumerate}
\end{lemma}

\begin{proof}
证明 (1)。不妨设存在有限自由 $R$-模组成的有界复形 $K^\bullet$ 和
$L^\bullet$，以及映射 $\alpha : K^\bullet \to K$ 和
$\beta : L^\bullet \to L$，使得当 $i > n$ 时 $H^i(\alpha)$ 是同构、
当 $i = n$ 时是满射，并且当 $i > m$ 时 $H^i(\beta)$ 是同构、
当 $i = m$ 时是满射。于是映射
$$
\alpha \otimes^\mathbf{L} \beta :
\text{Tot}(K^\bullet \otimes_R L^\bullet)
\to K \otimes_R^\mathbf{L} L
$$
在次数 $i$ 的上同调上，当 $i > t$ 时诱导同构，在 $i = t$ 时诱导满射。
这是由 tors 的谱序列得到的（细节略）。(2) 由 (1) 及引理
\ref{more-algebra-lemma-pseudo-coherent} 得到。
\end{proof}

\begin{lemma}
\label{more-algebra-lemma-Noetherian-pseudo-coherent}
设 $R$ 为 Noether 环。则
\begin{enumerate}
\item $R$-模复形 $K^\bullet$ 是 $m$-拟相干的，当且仅当
$K^\bullet \in D^{-}(R)$ 且当 $i \geq m$ 时 $H^i(K^\bullet)$ 是有限
$R$-模。
\item $R$-模复形 $K^\bullet$ 是拟相干的，当且仅当
$K^\bullet \in D^{-}(R)$ 且对所有 $i$，$H^i(K^\bullet)$ 是有限 $R$-模。
\item $R$-模是拟相干的，当且仅当它是有限的。
\end{enumerate}
\end{lemma}

\begin{proof}
在 Algebra，引理\ref{algebra-lemma-resolution-by-finite-free} 中，
我们已经看到每个有限 $R$-模都是拟相干的。
另一方面，拟相干模是有限的，见引理\ref{more-algebra-lemma-n-pseudo-module}。
因此 (3) 成立。假设 $K^\bullet$ 是 $m$-拟相干复形。
则存在有限自由 $R$-模组成的有界复形 $E^\bullet$，使得当 $i > m$ 时
$H^i(K^\bullet)$ 与 $H^i(E^\bullet)$ 同构，并且
$H^m(K^\bullet)$ 是 $H^m(E^\bullet)$ 的商。因此显然对每个 $i \geq m$，
$H^i(K^\bullet)$ 都是有限模。(1) 的逆向蕴含由引理
\ref{more-algebra-lemma-cohomology-pseudo-coherent} 及 (3) 得到。
(2) 由 (1) 及引理\ref{more-algebra-lemma-pseudo-coherent} 得到。
\end{proof}

\begin{lemma}
\label{more-algebra-lemma-coherent-pseudo-coherent}
设 $R$ 为相容环（Algebra，定义
\ref{algebra-definition-coherent}）。设 $K \in D^-(R)$。下列条件等价：
\begin{enumerate}
\item $K$ 是 $m$-拟相干的；
\item $H^m(K)$ 是有限 $R$-模，且当 $i > m$ 时 $H^i(K)$ 是相容模；以及
\item $H^m(K)$ 是有限 $R$-模，且当 $i > m$ 时
$H^i(K)$ 是有限表示的。
\end{enumerate}
因此，$K$ 拟相干当且仅当对所有 $i$，$H^i(K)$ 是相容模。
\end{lemma}

\begin{proof}
回顾：$R$-模 $M$ 相容，当且仅当它是有限表示的（Algebra，引理
\ref{algebra-lemma-coherent-ring}）。这说明 (2) 与 (3) 等价。
在此情形下，若选取正合序列
$0 \to N \to R^{\oplus m} \to M \to 0$，则由 Algebra，引理
\ref{algebra-lemma-coherent}，$N$ 是相容模。因此重复对 $N$ 作此过程，得到解消
$$
\ldots \to R^{\oplus n} \to R^{\oplus m} \to M \to 0
$$
其各项都是有限自由 $R$-模。换言之，$M$ 是拟相干的。
(1) 与 (2) 的等价性由此以及引理
\ref{more-algebra-lemma-cohomology-pseudo-coherent} 和
\ref{more-algebra-lemma-n-pseudo-module} 得到。
最后一个断言由 (1) 与 (2) 的等价性以及引理
\ref{more-algebra-lemma-pseudo-coherent} 得到。
\end{proof}









\section{拟相干模，II}
\label{more-algebra-section-pseudo-coherent-bis}

\noindent
我们继续第\ref{more-algebra-section-pseudo-coherent}节开始的讨论。

\begin{lemma}
\label{more-algebra-lemma-pseudo-coherence-colimit-ext}
设 $R$ 为环。设 $M = \colim M_i$ 为 $R$-模的滤过余极限。
设 $K \in D(R)$ 为 $m$-拟相干的。则当 $n < -m$ 时
$\colim \Ext^n_R(K, M_i) = \Ext^n_R(K, M)$，并且
$\colim \Ext^{-m}_R(K, M_i) \to \Ext^{-m}_R(K, M)$ 是单射。
\end{lemma}

\begin{proof}
由定义，可以找到可区别三角
$$
E \to K \to L \to E[1]
$$
在 $D(R)$ 中，其中 $E$ 由有限自由 $R$-模组成的有界复形表示，并且当 $i \geq m$ 时
$H^i(L) = 0$。于是对任意 $R$-模 $N$，当 $n \leq -m$ 时
$\Ext^n_R(L, N) = 0$，见 Derived Categories，引理
\ref{derived-lemma-negative-exts}。
由与该可区别三角关联的 $\Ext$ 长正合序列可知，
$\Ext^n_R(K, N) \to \Ext^n_R(E, N)$ 在 $n < -m$ 时是同构，
在 $n = -m$ 时是单射。因此只需证明，当 $E$ 可由有限自由 $R$-模组成的
有界复形 $E^\bullet$ 表示时，函子 $M \mapsto \Ext_R^n(E, M)$ 与滤过余极限交换。
$\Ext^n_R(E, M)$ 由复形 $\Hom_R(E^\bullet, M)$ 计算，见 Derived Categories，
引理\ref{derived-lemma-morphisms-from-projective-complex}。
由于 $E^p$ 是有限自由的，函子 $M \mapsto \Hom_R(E^p, M)$ 与滤过余极限交换。
因此作为复形有
$\Hom_R(E^\bullet, M) = \colim \Hom_R(E^\bullet, M_i)$。
由于滤过余极限是正合的（Algebra，引理
\ref{algebra-lemma-directed-colimit-exact}），结论成立。
\end{proof}

\begin{lemma}
\label{more-algebra-lemma-characterize-pseudo-coherent-colimit-ext}
设 $R$ 为环。设 $K \in D^-(R)$。设 $m \in \mathbf{Z}$。
则 $K$ 是 $m$-拟相干的，当且仅当对任意 $R$-模的滤过余极限
$M = \colim M_i$，都有
$\colim \Ext^n_R(K, M_i) = \Ext^n_R(K, M)$（当 $n < -m$），并且
$\colim \Ext^{-m}_R(K, M_i) \to \Ext^{-m}_R(K, M)$ 是单射。
\end{lemma}

\begin{proof}
一个方向已在引理\ref{more-algebra-lemma-pseudo-coherence-colimit-ext} 中证明。
现在假设对任意 $R$-模的滤过余极限 $M = \colim M_i$，都有
$$
\colim \Ext^n_R(K, M_i) = \Ext^n_R(K, M)
$$
（当 $n < -m$），并且
$$
\colim \Ext^{-m}_R(K, M_i) \to \Ext^{-m}_R(K, M)
$$
是单射。
我们证明 $K$ 是 $m$-拟相干的。

\medskip\noindent
令 $t$ 为使 $H^t(K)$ 非零的最大整数。我们对 $t$ 归纳。
如果 $t < m$，则由引理\ref{more-algebra-lemma-recognize-pseudo-coherent}，$K$ 是
$m$-拟相干的。如果 $t \geq m$，由于
$\Hom_R(H^t(K), M) = \Ext^{-t}_R(K, M)$，可知对任意滤过余极限
$M = \colim M_i$，映射
$\colim \Hom_R(H^t(K), M_i) \to \Hom_R(H^t(K), M)$ 是单射。
由 Algebra，引理\ref{algebra-lemma-characterize-finite-module-hom}，
这推出 $H^t(K)$ 是有限 $R$-模。选取有限自由 $R$-模 $F$ 以及满射
$F \to H^t(K)$。可将其提升为 $D(R)$ 中的态射 $F[-t] \to K$，并选取可区别三角
$$
F[-t] \to K \to L \to F[-t + 1]
$$
在 $D(R)$ 中。于是当 $i \geq t$ 时 $H^i(L) = 0$。此外，与该可区别三角关联的
$\Ext$ 长正合序列表明，经过一个略去的简短论证，$L$ 继承我们对 $K$ 所作的假设。
由对 $t$ 的归纳，$L$ 是 $m$-拟相干的。因此由引理
\ref{more-algebra-lemma-cone-pseudo-coherent}，$K$ 是 $m$-拟相干的。
\end{proof}

\begin{lemma}
\label{more-algebra-lemma-pseudo-coherence-and-ext}
设 $R$ 为环。设 $L$、$M$、$N$ 为 $R$-模。
\begin{enumerate}
\item 如果 $M$ 是有限表示的且 $L$ 平坦，则典范映射
$\Hom_R(M, N) \otimes_R L \to \Hom_R(M, N \otimes_R L)$
是同构。
\item 如果 $M$ 是 $(-m)$-拟相干的且 $L$ 平坦，则典范映射
$\Ext^i_R(M, N) \otimes_R L \to \Ext^i_R(M, N \otimes_R L)$
在 $i < m$ 时是同构。
\end{enumerate}
\end{lemma}

\begin{proof}
选取各项为自由 $R$-模的解消 $F_\bullet \to M$，见 Algebra，引理
\ref{algebra-lemma-resolution-by-finite-free}。复形
$\Hom_R(F_\bullet, N)$ 计算 $\Ext^i_R(M, N)$，而复形
$\Hom_R(F_\bullet, N \otimes_R L)$ 计算 $\Ext^i_R(M, N \otimes_R L)$。
总有复形映射
$$
\Hom_R(F_\bullet, N) \otimes_R L
\longrightarrow
\Hom_R(F_\bullet, N \otimes_R L)
$$
它对所有 $i \geq 0$ 诱导典范映射
$\Ext^i_R(M, N) \otimes_R L \to \Ext^i_R(M, N \otimes_R L)$
（例如，在这些映射不依赖于解消 $F_\bullet$ 的选取这一意义下是典范的）。
如果 $L$ 平坦，则复形 $\Hom_R(F_\bullet, N) \otimes_R L$ 计算
$\Ext^i_R(M, N) \otimes_R L$，因为取上同调与张量 $L$ 可交换。

\medskip\noindent
在上述事实基础上，如果 $M$ 是 $(-m)$-拟相干的，则可以选取 $F_\bullet$，使得
当 $i = 0, \ldots, m$ 时 $F_i$ 是有限自由的。于是上面显示的复形映射在次数
$\leq m$ 上是同构，因而在次数 $< m$ 上对上同调群诱导同构。这就证明了 (2)。
如果 $M$ 是有限表示的，则由引理\ref{more-algebra-lemma-n-pseudo-module}，$M$ 是
$(-1)$-拟相干的；再由 $\Hom = \Ext^0$ 得到结论。
\end{proof}

\begin{lemma}
\label{more-algebra-lemma-pseudo-coherence-and-base-change-ext}
设 $R \to R'$ 为平坦环同态。设 $M$、$N$ 为 $R$-模。
\begin{enumerate}
\item 如果 $M$ 是有限表示的 $R$-模，则
$\Hom_R(M, N) \otimes_R R' = \Hom_{R'}(M \otimes_R R', N \otimes_R R')$。
\item 如果 $M$ 是 $(-m)$-拟相干的，则
$\Ext^i_R(M, N) \otimes_R R' = \Ext^i_{R'}(M \otimes_R R', N \otimes_R R')$
当 $i < m$ 时成立。
\end{enumerate}
特别地，如果 $R$ 是 Noether 环且 $M$ 是有限模，则上述等式对所有 $i$ 都成立。
\end{lemma}

\begin{proof}
由 Algebra，引理\ref{algebra-lemma-flat-base-change-ext}，有
$\Ext^i_{R'}(M \otimes_R R', N \otimes_R R') =
\Ext^i_R(M, N \otimes_R R')$。
结合引理\ref{more-algebra-lemma-pseudo-coherence-and-ext}，即得 (1) 和 (2)。
最后一个断言由此及引理\ref{more-algebra-lemma-Noetherian-pseudo-coherent} 得到。
\end{proof}

\begin{lemma}
\label{more-algebra-lemma-pseudo-coherent-tensor}
设 $R$ 为环。设 $K \in D^-(R)$。下列条件等价：
\begin{enumerate}
\item $K$ 是拟相干的；
\item 对 $R$-模的任意族 $(Q_{\alpha})_{\alpha \in A}$，典范映射
$$
\alpha :
K \otimes_R^\mathbf{L} \left( \prod\nolimits_\alpha Q_{\alpha} \right)
\longrightarrow
\prod\nolimits_\alpha (K \otimes_R^\mathbf{L} Q_{\alpha})
$$
是 $D(R)$ 中的同构；
\item 对每个 $R$-模 $Q$ 和每个集合 $A$，典范映射
$$
\beta : K \otimes_R^\mathbf{L} Q^A \longrightarrow (K \otimes_R^\mathbf{L} Q)^A
$$
是 $D(R)$ 中的同构；以及
\item 对每个集合 $A$，典范映射
$$
\gamma : K \otimes_R^\mathbf{L} R^A \longrightarrow K^A
$$
是 $D(R)$ 中的同构。
\end{enumerate}
给定 $m \in \mathbf{Z}$，下列条件等价：
\begin{enumerate}
\item[(a)] $K$ 是 $m$-拟相干的；
\item[(b)] 对 $R$-模的任意族 $(Q_{\alpha})_{\alpha \in A}$，并令 $\alpha$ 如上，
$H^i(\alpha)$ 当 $i > m$ 时是同构、当 $i = m$ 时是满射；
\item[(c)] 对每个 $R$-模 $Q$ 和每个集合 $A$，并令 $\beta$ 如上，
$H^i(\beta)$ 当 $i > m$ 时是同构、当 $i = m$ 时是满射；
\item[(d)] 对每个集合 $A$，并令 $\gamma$ 如上，
$H^i(\gamma)$ 当 $i > m$ 时是同构、当 $i = m$ 时是满射。
\end{enumerate}
\end{lemma}

\begin{proof}
如果 $K$ 拟相干，则 $K$ 可由有限自由 $R$-模组成的上有界复形表示。
于是导出张量积由与此复形张量计算。并且，$D(R)$ 中的乘积可通过对任意
代表复形的选取取乘积来给出。因此由模的相应事实，(1) 推出 (2)、(3)、(4)，见
Algebra，命题\ref{algebra-proposition-fp-tensor}。

\medskip\noindent
同理（使用张量积是右正合的），读者可证明 (a) 推出 (b)、(c) 和 (d)。

\medskip\noindent
假设 (4) 成立。要证明 $K$ 拟相干，只需证明对所有 $m$，$K$ 都是
$m$-拟相干的（引理\ref{more-algebra-lemma-pseudo-coherent}）。因此只需证明
(d) 推出 (a)。

\medskip\noindent
假设 (d)。令 $i$ 为使 $H^i(K)$ 非零的最大整数。如果 $i < m$，则结论成立。
否则，由 (d) 以及上面对 $D(R)$ 中乘积的描述，可知
$H^i(K) \otimes_R R^A \to H^i(K)^A$ 是满射。因此由 Algebra，命题
\ref{algebra-proposition-fg-tensor}，$H^i(K)$ 是有限生成 $R$-模。
于是可以选取一个复形 $L$，其仅由位于次数 $i$ 的有限自由模组成，并选取复形映射
$L \to K$，使 $H^i(L) \to H^i(K)$ 满射。特别地，$L$ 满足 (1)、(2)、(3)、(4)。
选取可区别三角
$$
L \to K \to M \to L[1]
$$
则 $H^j(M) = 0$（当 $j \geq i$）。另一方面，经过一个略去的简短论证，$M$
仍具有性质 (d)。由对 $i$ 的归纳，$M$ 是 $m$-拟相干的。因此由引理
\ref{more-algebra-lemma-cone-pseudo-coherent}，$K$ 是 $m$-拟相干的。
\end{proof}

\begin{lemma}
\label{more-algebra-lemma-detect-cohomology-pseudo-coherent}
设 $R$ 为环。设 $K \in D(R)$ 为拟相干对象。设 $i \in \mathbf{Z}$。
存在有限表示的 $R$-模 $M$ 以及 $D(R)$ 中的映射 $K \to M[-i]$，
该映射在 $H^i(K) \to M$ 上诱导单射。
\end{lemma}

\begin{proof}
由定义\ref{more-algebra-definition-pseudo-coherent}，可以用有限自由 $R$-模组成的复形
$P^\bullet$ 表示 $K$。置 $M = \Coker(P^{i - 1} \to P^i)$。
\end{proof}

\begin{lemma}
\label{more-algebra-lemma-detect-cohomology}
设 $A$ 为 Noether 环。设 $K \in D(A)$ 为拟相干对象，即
$K \in D^-(A)$ 且上同调模有限。设 $\mathfrak m$ 为 $A$ 的极大理想。
如果 $H^i(K)/\mathfrak m H^i(K) \not = 0$，则存在被 $\mathfrak m$ 的某个幂
湮灭的有限 $A$-模 $E$，以及映射 $K \to E[-i]$，该映射在 $H^i(K)$ 上非零。
\end{lemma}

\begin{proof}
（引理陈述中拟相干性的等价表述是引理
\ref{more-algebra-lemma-Noetherian-pseudo-coherent}。）
按引理\ref{more-algebra-lemma-detect-cohomology-pseudo-coherent}，选取
$K \to M[-i]$。由 Artin-Rees（Algebra，引理
\ref{algebra-lemma-Artin-Rees}），可找到 $n$ 使得
$H^i(K) \cap \mathfrak m^n M \subset \mathfrak m H^i(K)$。
取 $E = M/\mathfrak m^n M$。
\end{proof}








\section{Tor 维数}
\label{more-algebra-section-tor}

\noindent
我们不再用投射模作解消，而是考察平坦模的解消。这引出如下概念。

\begin{definition}
\label{more-algebra-definition-tor-amplitude}
设 $R$ 为环。记 $D(R)$ 为其导出范畴。设 $a, b \in \mathbf{Z}$。
\begin{enumerate}
\item 如果对所有 $R$-模 $M$ 以及所有 $i \not \in [a, b]$，都有
$H^i(K^\bullet \otimes_R^\mathbf{L} M) = 0$，则称 $D(R)$ 中的对象
$K^\bullet$ 具有 $[a, b]$ 内的{\it Tor-振幅}。
\item 如果某些 $a, b$ 使其具有 $[a, b]$ 内的 Tor-振幅，则称 $D(R)$ 中的对象
$K^\bullet$ 具有{\it 有限 Tor 维数}。
\item 如果 $R$-模 $M$ 作为 $D(R)$ 中的对象 $M[0]$ 具有
$[-d, 0]$ 内的 Tor-振幅，则称该模具有{\it Tor 维数 $\leq d$}。
\item 如果 $R$-模 $M$ 作为 $D(R)$ 中的对象 $M[0]$ 具有有限 Tor 维数，
则称该模具有{\it 有限 Tor 维数}。
\end{enumerate}
\end{definition}

\noindent
注意，如果 $K^\bullet$ 具有有限 Tor 维数，则
$K^\bullet \in D^b(R)$。

\begin{lemma}
\label{more-algebra-lemma-last-one-flat}
设 $R$ 为环。设 $K^\bullet$ 为由平坦 $R$-模组成的上有界复形，
且具有 $[a, b]$ 内的 Tor-振幅。则
$\Coker(d_K^{a - 1})$ 是平坦 $R$-模。
\end{lemma}

\begin{proof}
由于 $K^\bullet$ 是平坦模组成的上有界复形，可知
$K^\bullet \otimes_R M = K^\bullet \otimes_R^{\mathbf{L}} M$。
因此对每个 $R$-模 $M$，序列
$$
K^{a - 2} \otimes_R M \to K^{a - 1} \otimes_R M \to K^a \otimes_R M
$$
在中间项处正合。由于
$K^{a - 2} \to K^{a - 1} \to K^a \to \Coker(d_K^{a - 1}) \to 0$
是平坦解消，这推出
$\text{Tor}_1^R(\Coker(d_K^{a - 1}), M) = 0$
对所有 $R$-模 $M$ 都成立。这意味着
$\Coker(d_K^{a - 1})$ 是平坦的，见
Algebra，引理\ref{algebra-lemma-characterize-flat}。
\end{proof}

\begin{lemma}
\label{more-algebra-lemma-tor-amplitude}
设 $R$ 为环。设 $K^\bullet$ 为 $D(R)$ 中的对象。设 $a, b \in \mathbf{Z}$。
下列条件等价：
\begin{enumerate}
\item $K^\bullet$ 具有 $[a, b]$ 内的 Tor-振幅；
\item $K^\bullet$ 拟同构于平坦 $R$-模复形 $E^\bullet$，且当
$i \not \in [a, b]$ 时 $E^i = 0$。
\end{enumerate}
\end{lemma}

\begin{proof}
如果 (2) 成立，则可以计算
$K^\bullet \otimes_R^\mathbf{L} M = E^\bullet \otimes_R M$，故 (1) 显然成立。
假设 (1)。可以用一个满足当 $i > b$ 时 $K^i = 0$ 的投射解消替代
$K^\bullet$。见 Derived Categories，引理
\ref{derived-lemma-projective-resolutions-exist}。置
$E^\bullet = \tau_{\geq a}K^\bullet$。除 $E^a$ 平坦外一切都显然；其
平坦性立即由引理\ref{more-algebra-lemma-last-one-flat} 及定义得到。
\end{proof}

\begin{lemma}
\label{more-algebra-lemma-bounded-below-tor-amplitude}
设 $R$ 为环。设 $a \in \mathbf{Z}$，且设 $K$ 为 $D(R)$ 中的对象。
下列条件等价：
\begin{enumerate}
\item $K$ 具有 $[a, \infty]$ 内的 Tor-振幅；
\item $K$ 拟同构于一个 K-平坦复形 $E^\bullet$，其各项为平坦 $R$-模，且
$E^i = 0$ 当 $i \not \in [a, \infty]$ 时成立。
\end{enumerate}
\end{lemma}

\begin{proof}
(2) $\Rightarrow$ (1) 立即成立。设 (1) 成立。
首先选取一个各项平坦且表示 $K$ 的 K-平坦复形 $K^\bullet$，
见引理\ref{more-algebra-lemma-K-flat-resolution}。
对任意 $R$-模 $M$，复形
$$
K^{n - 1} \otimes_R M \to K^n \otimes_R M \to K^{n + 1} \otimes_R M
$$
的上同调计算 $H^n(K \otimes_R^\mathbf{L} M)$。当 $n < a$ 时它总为零。
因此将引理\ref{more-algebra-lemma-last-one-flat} 应用于复形
$\ldots \to K^{a - 1} \to K^a \to K^{a + 1}$，
可得 $N = \Coker(K^{a - 1} \to K^a)$ 是平坦 $R$-模。置
$$
E^\bullet = \tau_{\geq a}K^\bullet =
(\ldots \to 0 \to N \to K^{a + 1} \to \ldots )
$$
复形 $K^\bullet \to E^\bullet$ 的核 $L^\bullet$ 为复形
$$
L^\bullet = (\ldots \to K^{a - 1} \to I \to 0 \to \ldots)
$$
其中 $I \subset K^a$ 是 $K^{a - 1} \to K^a$ 的像。
由于有短正合列 $0 \to I \to K^a \to N \to 0$，
可知 $I$ 是平坦 $R$-模。因此 $L^\bullet$ 是由平坦模组成的上有界复形，
从而由引理\ref{more-algebra-lemma-derived-tor-quasi-isomorphism} 可知它是 K-平坦的。
由引理\ref{more-algebra-lemma-K-flat-two-out-of-three-ses}，$E^\bullet$ 也是 K-平坦的。
\end{proof}

\begin{lemma}
\label{more-algebra-lemma-cone-tor-amplitude}
设 $R$ 为环。
设 $(K^\bullet, L^\bullet, M^\bullet, f, g, h)$ 为 $D(R)$ 中的可区别三角。
设 $a, b \in \mathbf{Z}$。
\begin{enumerate}
\item 若 $K^\bullet$ 具有 $[a + 1, b + 1]$ 内的 Tor-振幅，且
$L^\bullet$ 具有 $[a, b]$ 内的 Tor-振幅，则 $M^\bullet$ 具有
$[a, b]$ 内的 Tor-振幅。
\item 若 $K^\bullet, M^\bullet$ 具有 $[a, b]$ 内的 Tor-振幅，则
$L^\bullet$ 具有 $[a, b]$ 内的 Tor-振幅。
\item 若 $L^\bullet$ 具有 $[a + 1, b + 1]$ 内的 Tor-振幅，且
$M^\bullet$ 具有 $[a, b]$ 内的 Tor-振幅，则
$K^\bullet$ 具有 $[a + 1, b + 1]$ 内的 Tor-振幅。
\end{enumerate}
\end{lemma}

\begin{proof}
略去证明。结论直接来自可区别三角所关联的上同调长正合列，以及
$- \otimes_R^{\mathbf{L}} M$ 保持可区别三角这一事实。
最容易证明的是 (2)，其余结论由平移从 (2) 得到。
\end{proof}

\begin{lemma}
\label{more-algebra-lemma-tor-dimension}
设 $R$ 为环。设 $M$ 为一个 $R$-模。
设 $d \geq 0$。下列条件等价：
\begin{enumerate}
\item $M$ 具有 Tor 维数 $\leq d$；
\item 存在解消
$$
0 \to F_d \to \ldots \to F_1 \to F_0 \to M \to 0
$$
其中 $F_i$ 为平坦 $R$-模。
\end{enumerate}
特别地，$R$-模的 Tor 维数为 $0$ 当且仅当它是平坦 $R$-模。
\end{lemma}

\begin{proof}
设 (2) 成立。取 $E^\bullet$ 为满足 $E^{-i} = F_i$ 的复形，
则它与 $M$ 拟同构。因此由引理\ref{more-algebra-lemma-tor-amplitude}，
$M$ 的 Tor 维数至多为 $d$。
反之，设 (1) 成立。令 $P^\bullet \to M$ 为 $M$ 的投射解消。
由引理\ref{more-algebra-lemma-last-one-flat} 可知，$\tau_{\geq -d}P^\bullet$
是 $M$ 的长度为 $d$ 的平坦解消，即 (2) 成立。
\end{proof}

\begin{lemma}
\label{more-algebra-lemma-summands-tor-amplitude}
设 $R$ 为环。设 $a, b \in \mathbf{Z}$。
若 $K^\bullet \oplus L^\bullet$ 具有 $[a, b]$ 内的 Tor 振幅，
则 $K^\bullet$ 和 $L^\bullet$ 也具有该 Tor 振幅。
\end{lemma}

\begin{proof}
这由 Tor 函子具有可加性立即得到。
\end{proof}

\begin{lemma}
\label{more-algebra-lemma-complex-finite-tor-dimension-modules}
设 $R$ 为环。设 $K^\bullet$ 为由 $R$-模组成的有界复形，
且对每个 $i$，$K^i$ 具有 $[a - i, b - i]$ 内的 Tor 振幅。
则 $K^\bullet$ 具有 $[a, b]$ 内的 Tor 振幅。特别地，
若 $K^\bullet$ 是由有限个具有有限 Tor 维数的 $R$-模组成的复形，
则 $K^\bullet$ 具有有限 Tor 维数。
\end{lemma}

\begin{proof}
对有限复形的长度作归纳即可；使用引理\ref{more-algebra-lemma-cone-tor-amplitude}
以及朴素截断。
\end{proof}

\begin{lemma}
\label{more-algebra-lemma-cohomology-tor-amplitude}
设 $R$ 为环。设 $a, b \in \mathbf{Z}$。设 $K^\bullet \in D^b(R)$，
使得对所有 $i$，$H^i(K^\bullet)$ 具有 $[a - i, b - i]$ 内的 Tor 振幅。
则 $K^\bullet$ 具有 $[a, b]$ 内的 Tor 振幅。特别地，
若 $K^\bullet \in D^b(R)$ 且其所有上同调群都具有有限 Tor
维数，则 $K^\bullet$ 具有有限 Tor 维数。
\end{lemma}

\begin{proof}
对有限复形的长度作归纳即可；使用引理\ref{more-algebra-lemma-cone-tor-amplitude}
以及典范截断。
\end{proof}

\begin{lemma}
\label{more-algebra-lemma-push-tor-amplitude}
设 $A \to B$ 为环映射。设 $K^\bullet$ 和 $L^\bullet$ 为由
$B$-模组成的复形。设 $a, b, c, d \in \mathbf{Z}$。若
\begin{enumerate}
\item $K^\bullet$ 作为 $B$-模复形具有 $[a, b]$ 内的 Tor 振幅；
\item $L^\bullet$ 作为 $A$-模复形具有 $[c, d]$ 内的 Tor 振幅，
\end{enumerate}
则 $K^\bullet \otimes^\mathbf{L}_B L^\bullet$ 作为 $A$-模复形
具有 $[a + c, b + d]$ 内的 Tor 振幅。
\end{lemma}

\begin{proof}
不妨设 $K^\bullet$ 是由平坦 $B$-模组成且当
$i \not \in [a, b]$ 时 $K^i = 0$ 的复形，见引理\ref{more-algebra-lemma-tor-amplitude}。
令 $M$ 为 $A$-模。选取自由解消 $F^\bullet \to M$。则
$$
(K^\bullet \otimes_B^\mathbf{L} L^\bullet) \otimes_A^{\mathbf{L}} M =
\text{Tot}(\text{Tot}(K^\bullet \otimes_B L^\bullet) \otimes_A F^\bullet) =
\text{Tot}(K^\bullet \otimes_B \text{Tot}(L^\bullet \otimes_A F^\bullet))
$$
关于第二个等式，见 Homology，引理\ref{homology-remark-triple-complex}。
由假设 (2)，复形
$\text{Tot}(L^\bullet \otimes_A F^\bullet)$
仅在次数 $[c, d]$ 中有非零上同调。因此，Homology，
引理\ref{homology-lemma-ss-double-complex} 关于双复形
$K^\bullet \otimes_B \text{Tot}(L^\bullet \otimes_A F^\bullet)$
的谱序列表明，
$(K^\bullet \otimes_B^\mathbf{L} L^\bullet) \otimes_A^{\mathbf{L}} M$
仅在次数 $[a + c, b + d]$ 中有非零上同调。
\end{proof}

\begin{lemma}
\label{more-algebra-lemma-flat-push-tor-amplitude}
设 $A \to B$ 为环映射。假设 $B$ 作为 $A$-模是平坦的。设
$K^\bullet$ 为由 $B$-模组成的复形。设 $a, b \in \mathbf{Z}$。若
$K^\bullet$ 作为 $B$-模复形具有 $[a, b]$ 内的 Tor 振幅，
则 $K^\bullet$ 作为 $A$-模复形具有 $[a, b]$ 内的 Tor 振幅。
\end{lemma}

\begin{proof}
这是引理\ref{more-algebra-lemma-push-tor-amplitude} 的特例，但也可直接如下看出。我们有
$K^\bullet \otimes_A^{\mathbf{L}} M =
K^\bullet \otimes_B^{\mathbf{L}} (M \otimes_A B)$，
因为 $K^\bullet$ 作为 $B$-模复形的任意投射解消都是
$K^\bullet$ 作为 $A$-模复形的平坦解消，并可用于计算
$K^\bullet \otimes_A^{\mathbf{L}} M$。
\end{proof}

\begin{lemma}
\label{more-algebra-lemma-finite-tor-dimension-push-tor-amplitude}
设 $A \to B$ 为环映射。假设 $B$ 作为 $A$-模具有 Tor 维数 $\leq d$。
设 $K^\bullet$ 为由 $B$-模组成的复形。设 $a, b \in \mathbf{Z}$。若
$K^\bullet$ 作为 $B$-模复形具有 $[a, b]$ 内的 Tor 振幅，
则 $K^\bullet$ 作为 $A$-模复形具有 $[a - d, b]$ 内的 Tor 振幅。
\end{lemma}

\begin{proof}
这是引理\ref{more-algebra-lemma-push-tor-amplitude} 的特例，但也可直接如下看出。
令 $M$ 为 $A$-模。选取自由解消 $F^\bullet \to M$。则
$$
K^\bullet \otimes_A^{\mathbf{L}} M =
\text{Tot}(K^\bullet \otimes_A F^\bullet) =
\text{Tot}(K^\bullet \otimes_B (F^\bullet \otimes_A B)) =
K^\bullet \otimes_B^{\mathbf{L}} (M \otimes_A^{\mathbf{L}} B).
$$
由 $B$ 作为 $A$-模所满足的假设可知，
$M \otimes_A^{\mathbf{L}} B$ 的上同调只可能位于
$-d, -d + 1, \ldots, 0$ 这些次数。由于 $K^\bullet$ 具有
$[a, b]$ 内的 Tor 振幅，由 Example\ref{more-algebra-example-tor} 中的谱序列可知，
$K^\bullet \otimes_B^{\mathbf{L}} (M \otimes_A^{\mathbf{L}} B)$
的上同调只可能位于 $[-d + a, b]$，正如所需。
\end{proof}

\begin{lemma}
\label{more-algebra-lemma-pull-tor-amplitude}
设 $A \to B$ 为环映射。
设 $a, b \in \mathbf{Z}$。
设 $K^\bullet$ 为具有 $[a, b]$ 内 Tor 振幅的 $A$-模复形。
则 $K^\bullet \otimes_A^{\mathbf{L}} B$ 作为 $B$-模复形
具有 $[a, b]$ 内的 Tor 振幅。
\end{lemma}

\begin{proof}
由引理\ref{more-algebra-lemma-tor-amplitude}，可找到拟同构
$E^\bullet \to K^\bullet$，其中 $E^\bullet$ 是由平坦 $A$-模组成且当
$i \not \in [a, b]$ 时 $E^i = 0$ 的复形。于是按构造，
$E^\bullet \otimes_A B$ 计算 $K^\bullet \otimes_A ^{\mathbf{L}} B$，
并且由 Algebra，引理\ref{algebra-lemma-flat-base-change}，
每个 $E^i \otimes_A B$ 都是平坦 $B$-模。
因此由引理\ref{more-algebra-lemma-tor-amplitude} 得到结论。
\end{proof}

\begin{lemma}
\label{more-algebra-lemma-flat-base-change-finite-tor-dimension}
设 $A \to B$ 为平坦环映射。设 $d \geq 0$。
设 $M$ 为 Tor 维数 $\leq d$ 的 $A$-模。
则 $M \otimes_A B$ 是 Tor 维数 $\leq d$ 的 $B$-模。
\end{lemma}

\begin{proof}
这是引理\ref{more-algebra-lemma-pull-tor-amplitude} 的直接推论，
以及 $M \otimes_A^{\mathbf{L}} B = M \otimes_A B$ 这一事实的结果，
其中使用了 $B$ 在 $A$ 上平坦。
\end{proof}

\begin{lemma}
\label{more-algebra-lemma-tor-amplitude-localization}
设 $A  \to B$ 为环映射。设 $K^\bullet$ 为由 $B$-模组成的复形。
设 $a, b \in \mathbf{Z}$。下列条件等价：
\begin{enumerate}
\item $K^\bullet$ 作为 $A$-模复形具有 $[a, b]$ 内的 Tor 振幅；
\item 对每个素理想 $\mathfrak q \subset B$，令
$\mathfrak p = A \cap \mathfrak q$，则 $K^\bullet_\mathfrak q$
作为 $A_\mathfrak p$-模复形具有 $[a, b]$ 内的 Tor 振幅；
\item 对每个极大理想 $\mathfrak m \subset B$，令
$\mathfrak p = A \cap \mathfrak m$，则 $K^\bullet_\mathfrak m$
作为 $A_\mathfrak p$-模复形具有 $[a, b]$ 内的 Tor 振幅。
\end{enumerate}
\end{lemma}

\begin{proof}
设 (3) 成立，令 $M$ 为 $A$-模。则
$H^i = H^i(K^\bullet \otimes_A^\mathbf{L} M)$ 是 $B$-模，且
$(H^i)_\mathfrak m =
H^i(K^\bullet_\mathfrak m \otimes_{A_\mathfrak p}^\mathbf{L} M_\mathfrak p)$。
因此由 Algebra，引理\ref{algebra-lemma-characterize-zero-local}，
当 $i \not \in [a, b]$ 时 $H^i = 0$。故
(3) $\Rightarrow$ (1)。我们略去 (1) $\Rightarrow$ (2) 以及
(2) $\Rightarrow$ (3) 的证明。
\end{proof}

\begin{lemma}
\label{more-algebra-lemma-glue-tor-amplitude}
设 $R$ 为环。设 $f_1, \ldots, f_r \in R$ 为生成单位理想的一组元素。
设 $a, b \in \mathbf{Z}$。设 $K^\bullet$ 为由 $R$-模组成的复形。
若对每个 $i$，复形 $K^\bullet \otimes_R R_{f_i}$ 具有
$[a, b]$ 内的 Tor 振幅，则 $K^\bullet$ 具有 $[a, b]$ 内的 Tor 振幅。
\end{lemma}

\begin{proof}
这立即由引理\ref{more-algebra-lemma-tor-amplitude-localization} 得到，但也可直接如下看出。
注意 $- \otimes_R R_{f_i}$ 是正合函子，因此
$$
H^i(K^\bullet)_{f_i} =
H^i(K^\bullet) \otimes_R R_{f_i} = H^i(K^\bullet \otimes_R R_{f_i}).
$$
类似地，对每个 $R$-模 $M$，有
$$
H^i(K^\bullet \otimes_R^{\mathbf{L}} M)_{f_i} =
H^i(K^\bullet \otimes_R^{\mathbf{L}} M) \otimes_R R_{f_i} =
H^i(K^\bullet \otimes_R R_{f_i} \otimes_{R_{f_i}}^{\mathbf{L}} M_{f_i}).
$$
因此结果来自如下事实：$R$-模 $N$ 为零当且仅当对每个 $i$，
$N_{f_i}$ 为零，见 Algebra，引理\ref{algebra-lemma-cover}。
\end{proof}

\begin{lemma}
\label{more-algebra-lemma-flat-descent-tor-amplitude}
设 $R$ 为环。设 $a, b \in \mathbf{Z}$。设 $K^\bullet$ 为由 $R$-模组成的复形。
设 $R \to R'$ 为忠实平坦环映射。若复形 $K^\bullet \otimes_R R'$
具有 $[a, b]$ 内的 Tor 振幅，则 $K^\bullet$ 具有 $[a, b]$ 内的 Tor 振幅。
\end{lemma}

\begin{proof}
令 $M$ 为 $R$-模。由于 $R \to R'$ 平坦，有
$$
(M \otimes_R^{\mathbf{L}} K^\bullet) \otimes_R R'
=
((M \otimes_R R') \otimes_{R'}^{\mathbf{L}} (K^\bullet \otimes_R R')
$$
并且取上同调与张量 $R'$ 可交换。
因此
$\text{Tor}_i^R(M, K^\bullet) \otimes_R R' =
\text{Tor}_i^{R'}(M \otimes_R R', K^\bullet \otimes_R R')$。
由于 $R \to R'$ 忠实平坦，当
$i \not \in [a, b]$ 时
$\text{Tor}_i^{R'}(M \otimes_R R', K^\bullet \otimes_R R')$ 的消灭
蕴含 $\text{Tor}_i^R(M, K^\bullet)$ 也消灭。
\end{proof}

\begin{lemma}
\label{more-algebra-lemma-no-change-tor-amplitude}
给定环映射 $R \to A \to B$，其中 $A \to B$ 忠实平坦，且
$K \in D(A)$，则 $K$ 在 $R$ 上的 Tor 振幅与
$K \otimes_A^\mathbf{L} B$ 在 $R$ 上的 Tor 振幅相同。
\end{lemma}

\begin{proof}
这是因为对任意 $R$-模 $M$，都有
$H^i(K \otimes_R^\mathbf{L} M) \otimes_A B =
H^i((K \otimes_A^\mathbf{L} B) \otimes_R^\mathbf{L} M)$
对所有 $i$ 成立。
具体地，用一个 $A$-模复形 $K^\bullet$ 表示 $K$，并选取自由解消
$F^\bullet \to M$。则有等式
$$
\text{Tot}(K^\bullet \otimes_A B \otimes_R F^\bullet) =
\text{Tot}(K^\bullet \otimes_R F^\bullet) \otimes_A B
$$
左端的上同调群为
$H^i((K \otimes_A^\mathbf{L} B) \otimes_R^\mathbf{L} M)$，
而右端得到
$H^i(K \otimes_R^\mathbf{L} M) \otimes_A B$。
\end{proof}

\begin{lemma}
\label{more-algebra-lemma-finite-gl-dim-tor-dimension}
设 $R$ 为有限整体维数为 $d$ 的环。则
\begin{enumerate}
\item 每个模都具有 Tor 维数 $\leq d$；
\item 若由 $R$-模组成的复形 $K^\bullet$ 仅在
$i \in [a, b]$ 时满足 $H^i(K^\bullet) \not = 0$，
则它具有 $[a - d, b]$ 内的 Tor 振幅；
\item 由 $R$-模组成的复形 $K^\bullet$ 具有有限 Tor 维数，当且仅当
$K^\bullet \in D^b(R)$。
\end{enumerate}
\end{lemma}

\begin{proof}
$R$ 的假设意味着每个模都有长度至多为 $d$ 的有限投射解消，
特别地，每个模的 Tor 维数 $\leq d$。第二个断言由
引理\ref{more-algebra-lemma-cohomology-tor-amplitude} 及定义得到。第三个断言是第二个的改述。
\end{proof}

\begin{lemma}
\label{more-algebra-lemma-check-tor-dimension-modulo-nilpotent}
设 $R' \to R$ 为满射环映射，其核为幂零理想。
设 $K' \in D(R')$，并置 $K = K' \otimes_{R'}^\mathbf{L} R$。
设 $a, b \in \mathbf{Z}$。
则 $K$ 具有 $[a, b]$ 内的 Tor 振幅，当且仅当 $K'$ 具有该振幅。
\end{lemma}

\begin{proof}
一个方向由引理\ref{more-algebra-lemma-pull-tor-amplitude} 得到。
反之，设 $K$ 具有 $[a, b]$ 内的 Tor 振幅，令 $M'$ 为 $R'$-模。
需证明 $K' \otimes_{R'}^\mathbf{L} M'$ 的非零上同调只出现在
区间 $[a, b]$ 内。

\medskip\noindent
令 $I = \Ker(R' \to R)$。则某个 $n$ 使 $I^n = 0$。若 $IM' = 0$，
则可将 $M'$ 视为 $R$-模，并如下论证：
$$
K' \otimes_{R'}^\mathbf{L} M' =
K' \otimes_{R'}^\mathbf{L} (R \otimes_R^\mathbf{L} M') =
(K' \otimes_{R'}^\mathbf{L} R) \otimes_R^\mathbf{L} M' =
K \otimes_R^\mathbf{L} M'
$$
由 $K$ 的假设，其上同调只在区间 $[a, b]$ 内非零。若
$I^{t + 1}M' = 0$，则考虑短正合列
$$
0 \to IM' \to M' \to M'/IM' \to 0
$$
对 $t$ 作归纳可知，$K' \otimes_{R'}^\mathbf{L} IM'$ 与
$K' \otimes_{R'}^\mathbf{L} M'/IM'$ 的上同调都只在区间 $[a, b]$ 内非零。
于是可区别三角
$$
K' \otimes_{R'}^\mathbf{L} IM' \to
K' \otimes_{R'}^\mathbf{L} M' \to
K' \otimes_{R'}^\mathbf{L} M'/IM' \to
(K' \otimes_{R'}^\mathbf{L} IM')[1]
$$
证明 $K' \otimes_{R'}^\mathbf{L} M'$ 也具有所需性质。
\end{proof}

\section{Ext 的谱序列}
\label{more-algebra-section-spectral-sequence-ext}

\noindent
本节汇集考察 Ext 函子时出现的各种谱序列。对于环 $R$ 的导出范畴
$D(R)$ 中的任意一对对象 $L$、$K$，记
$$
\Ext^n_R(L, K) = \Hom_{D(R)}(L, K[n])
$$
这遵循我们在 Derived Categories，第
\ref{derived-section-ext} 节中的一般约定。

\medskip\noindent
若 $M$ 为 $R$-模且 $K \in D^+(R)$，则存在谱序列
\begin{equation}
\label{more-algebra-equation-first-ss-ext}
E_2^{i, j} = \Ext_R^i(M, H^j(K)) \Rightarrow \Ext_R^{i + j}(M, K)
\end{equation}
若 $K$ 由下有界的 $R$-模复形 $K^\bullet$ 表示，则还有谱序列
\begin{equation}
\label{more-algebra-equation-second-ss-ext}
E_1^{i, j} = \Ext_R^j(M, K^i) \Rightarrow \Ext_R^{i + j}(M, K)
\end{equation}
这些谱序列来自将 Derived Categories，引理
\ref{derived-lemma-two-ss-complex-functor} 应用于函子
$\Hom_R(M, -)$。

\section{投射维数}
\label{more-algebra-section-projective-dimension}

\noindent
我们在 Algebra，定义\ref{algebra-definition-finite-proj-dim} 中定义了模的投射维数。

\begin{definition}
\label{more-algebra-definition-projective-dimension}
设 $R$ 为环。设 $K$ 为 $D(R)$ 中的对象。若 $K$ 可由投射模组成的
有界复形表示，则称 $K$ 具有{\it 有限投射维数}。若 $K$ 拟同构于复形
$$
\ldots \to 0 \to P^a \to P^{a + 1} \to \ldots \to
P^{b - 1} \to P^b \to 0 \to \ldots
$$
则称 $K$ 具有{\it $[a, b]$ 内的投射振幅}，其中对所有
$i \in \mathbf{Z}$，$P^i$ 都是投射 $R$-模。
\end{definition}

\noindent
显然，$K$ 具有有限投射维数，当且仅当对某些
$a, b \in \mathbf{Z}$，$K$ 具有 $[a, b]$ 内的投射振幅。
此外，若 $K$ 具有有限投射维数，则 $K$ 是有界的。下面给出一个
用于检测 $D(R)$ 中这类对象的引理。

\begin{lemma}
\label{more-algebra-lemma-projective-amplitude}
设 $R$ 为环。设 $K$ 为 $D(R)$ 中的对象。设 $a, b \in \mathbf{Z}$。
下列条件等价：
\begin{enumerate}
\item $K$ 具有 $[a, b]$ 内的投射振幅；
\item 对所有 $R$-模 $N$ 以及所有 $i \not \in [-b, -a]$，
$\Ext^i_R(K, N) = 0$；
\item 当 $n > b$ 时 $H^n(K) = 0$，且对所有 $R$-模 $N$ 以及所有
$i > -a$，都有 $\Ext^i_R(K, N) = 0$；
\item 当 $n \not \in [a - 1, b]$ 时 $H^n(K) = 0$，且对所有
$R$-模 $N$，都有 $\Ext^{-a + 1}_R(K, N) = 0$。
\end{enumerate}
\end{lemma}

\begin{proof}
设 (1) 成立。不妨设 $K$ 是复形
$$
\ldots \to 0 \to P^a \to P^{a + 1} \to \ldots \to
P^{b - 1} \to P^b \to 0 \to \ldots
$$
其中对所有 $i \in \mathbf{Z}$，$P^i$ 都是投射 $R$-模。
此时可用复形
$$
\ldots \to 0 \to \Hom_R(P^b, N) \to \ldots \to
\Hom_R(P^a, N) \to 0 \to \ldots
$$
计算 Ext 群，从而得到 (2)。

\medskip\noindent
设 (2) 成立。选取单射 $H^n(K) \to I$，其中 $I$ 是单射 $R$-模。
由于 $\Hom_R(-, I)$ 是正合函子，可知
$\Ext^{-n}(K, I) = \Hom_R(H^n(K), I)$。
特别地，当 $n > b$ 时 $H^n(K)$ 为零。因此 (2) 推出 (3)。

\medskip\noindent
与 (2) 推出 (3) 时相同的论证表明 (3) 推出 (4)。

\medskip\noindent
设 (4) 成立。(2) 推出 (3) 的相同论证表明
$H^{a - 1}(K) = 0$，即 $H^i(K) = 0$，除非 $i \in [a, b]$。
特别地，$K$ 上有界，因此可选取一个表示 $K$ 的复形
$P^\bullet$，使得对所有 $i \in \mathbf{Z}$，$P^i$ 为投射
（例如自由）模，且当 $i > b$ 时 $P^i = 0$。
见 Derived Categories，引理
\ref{derived-lemma-subcategory-left-resolution}。
令 $Q = \Coker(P^{a - 1} \to P^a)$。则 $K$ 拟同构于复形
$$
\ldots \to 0 \to Q \to P^{a + 1} \to \ldots \to P^b \to 0 \to \ldots
$$
因为当 $i < a$ 时 $H^i(K) = 0$。
记 $K' = (P^{a + 1} \to \ldots \to P^b)$ 为 $D(R)$ 中相应的对象。
我们得到可区别三角
$$
K' \to K \to Q[-a] \to K'[1]
$$
位于 $D(R)$ 中。因此对每个 $R$-模 $N$，有正合列
$$
\Ext^{-a}(K', N) \to \Ext^1(Q, N) \to \Ext^{1 - a}(K, N)
$$
由假设，右端项消失；由蕴含 (1) $\Rightarrow$ (2)，左端项也消失。
因此由 Algebra，引理\ref{algebra-lemma-characterize-projective}，
$Q$ 是投射 $R$-模。于是 (1) 成立，证明完毕。
\end{proof}

\begin{example}
\label{more-algebra-example-ext-not-bounded}
设 $k$ 为域，令 $R$ 为 $k$ 上的对偶数环，即 $R = k[x]/(x^2)$。
记 $\epsilon \in R$ 为 $x$ 的类。令 $M = R/(\epsilon)$。
则 $M$ 拟同构于复形
$$
R \xrightarrow{\epsilon} R \xrightarrow{\epsilon} R \to \ldots
$$
但按照 Algebra，定义\ref{algebra-definition-finite-proj-dim} 中的定义，
$M$ 不具有有限投射维数。
这说明在有限投射维数的定义中，我们考虑 $D(R)$ 的对象所对应的
在两个方向上均有界的投射模复形。
\end{example}

\section{内射维数}
\label{more-algebra-section-injective-dimension}

\noindent
本节是投射维数一节的对偶版本。

\begin{definition}
\label{more-algebra-definition-injective-dimension}
设 $R$ 为环。设 $K$ 为 $D(R)$ 中的对象。
若 $K$ 可由内射 $R$-模组成的有限复形表示，则称 $K$ 具有
{\it 有限内射维数}。若 $K$ 同构于复形
$$
\ldots \to 0 \to I^a \to I^{a + 1} \to \ldots \to
I^{b - 1} \to I^b \to 0 \to \ldots
$$
则称 $K$ 具有{\it $[a, b]$ 内的内射振幅}，其中对所有
$i \in \mathbf{Z}$，$I^i$ 都是内射 $R$-模。
\end{definition}

\noindent
显然，$K$ 具有有界内射维数，当且仅当对某些
$a, b \in \mathbf{Z}$，$K$ 具有 $[a, b]$ 内的内射振幅。
此外，若 $K$ 具有有界内射维数，则 $K$ 是有界的。下面给出必需的引理。

\begin{lemma}
\label{more-algebra-lemma-injective-amplitude}
设 $R$ 为环。设 $K$ 为 $D(R)$ 中的对象。设 $a, b \in \mathbf{Z}$。
下列条件等价：
\begin{enumerate}
\item $K$ 具有 $[a, b]$ 内的内射振幅；
\item 对所有 $R$-模 $N$ 以及所有 $i \not \in [a, b]$，
$\Ext^i_R(N, K) = 0$；
\item 对所有理想 $I \subset R$ 以及所有 $i \not \in [a, b]$，
$\Ext^i(R/I, K) = 0$。
\end{enumerate}
\end{lemma}

\begin{proof}
设 (1) 成立。不妨设 $K$ 是复形
$$
\ldots \to 0 \to I^a \to I^{a + 1} \to \ldots \to
I^{b - 1} \to I^b \to 0 \to \ldots
$$
其中对所有 $i \in \mathbf{Z}$，$I^i$ 都是内射 $R$-模。
此时可用复形
$$
\ldots \to 0 \to \Hom_R(N, I^a) \to \ldots \to
\Hom_R(N, I^b) \to 0 \to \ldots
$$
计算 Ext 群，从而得到 (2)。(2) 显然推出 (3)。

\medskip\noindent
设 (3) 成立。选取非零映射 $R \to H^n(K)$。由于
$\Hom_R(R, -)$ 是正合函子，可知
$\Ext^n_R(R, K) = \Hom_R(R, H^n(K)) = H^n(K)$。
由此可知，当 $n \not \in [a, b]$ 时 $H^n(K)$ 为零。
特别地，$K$ 下有界，因此可选取拟同构
$$
K \to I^\bullet
$$
其中对所有 $i \in \mathbf{Z}$，$I^i$ 为内射模，且当 $i < a$ 时
$I^i = 0$。见 Derived Categories，引理
\ref{derived-lemma-subcategory-right-resolution}。
令 $J = \Ker(I^b \to I^{b + 1})$。则 $K$ 拟同构于复形
$$
\ldots \to 0 \to I^a \to \ldots \to I^{b - 1} \to J \to 0 \to \ldots
$$
记 $K' = (I^a \to \ldots \to I^{b - 1})$ 为 $D(R)$ 中相应的对象。
我们得到可区别三角
$$
J[-b] \to K \to K' \to J[1 - b]
$$
位于 $D(R)$ 中。因此对每个理想 $I \subset R$，有正合列
$$
\Ext^b(R/I, K') \to \Ext^1(R/I, J) \to \Ext^{1 + b}(R/I, K)
$$
由假设右端项消失；由蕴含 (1) $\Rightarrow$ (2)，左端项也消失。
因此由引理\ref{more-algebra-lemma-characterize-injective-bis}，$J$ 是内射 $R$-模。
\end{proof}

\begin{example}
\label{more-algebra-example-finite-injective-finite-global-dimension}
设 $R$ 为 Dedekind 整环。则每个非零理想 $I$ 都是有限投射模，
见引理\ref{more-algebra-lemma-dedekind-torsion-free-flat}。
因此 $R/I$ 的投射维数为 $1$。由引理\ref{more-algebra-lemma-injective-amplitude}，
每个 $R$-模 $M$ 的内射维数都满足 $\leq 1$。
于是对任意一对 $R$-模 $M, N$，当 $i \geq 2$ 时
$\Ext^i_R(M, N) = 0$。
因此 $D^b(R)$ 中任意对象 $K$ 都同构于其上同调的直和：
$K \cong \bigoplus H^i(K)[-i]$，见 Derived Categories，
引理\ref{derived-lemma-ext-2-zero}。
\end{example}

\begin{example}
\label{more-algebra-example-ext-not-bounded-reversed}
设 $k$ 为域，令 $R$ 为 $k$ 上的对偶数环，即 $R = k[x]/(x^2)$。
记 $\epsilon \in R$ 为 $x$ 的类。令 $M = R/(\epsilon)$。
则 $M$ 拟同构于复形
$$
\ldots \to R \xrightarrow{\epsilon} R \xrightarrow{\epsilon} R
$$
且 $R$ 是内射 $R$-模。然而在这种情形下，通常不认为 $M$ 具有有限内射维数。
这说明在有界内射维数的定义中，我们考虑 $D(R)$ 的对象所对应的
在两个方向上均有界的内射模复形。
\end{example}

\begin{lemma}
\label{more-algebra-lemma-finite-injective-dimension}
设 $R$ 为环。设 $K \in D(R)$。
\begin{enumerate}
\item 若 $K$ 属于 $D^b(R)$ 且对所有 $i$，$H^i(K)$ 都具有有限内射维数，
则 $K$ 具有有限内射维数；
\item 若表示 $K$ 的 $K^\bullet$ 是由 $R$-模组成的有界复形，
且对所有 $i$，$K^i$ 都具有有限内射维数，则 $K$ 具有有限内射维数。
\end{enumerate}
\end{lemma}

\begin{proof}
略去证明。提示：将 Derived Categories，引理
\ref{derived-lemma-two-ss-complex-functor} 应用于函子
$F = \Hom_R(N, -)$，以计算 $\Ext^i_A(N, K)$，并使用引理
\ref{more-algebra-lemma-injective-amplitude} 的判据。
\end{proof}

\begin{lemma}
\label{more-algebra-lemma-finite-injective-dimension-Noetherian-radical}
设 $R$ 为 Noether 环。设 $I \subset R$ 为包含于 $R$ 的 Jacobson 根中的理想。
设 $K \in D^+(R)$ 且具有有限上同调模。则下列条件等价：
\begin{enumerate}
\item $K$ 具有有限内射维数；
\item 存在某个 $b$，使得对任意理想 $J \supset I$，当 $i > b$ 时
$\Ext^i_R(R/J, K) = 0$。
\end{enumerate}
\end{lemma}

\begin{proof}
(1) $\Rightarrow$ (2) 立即成立。设 (2) 成立。
设 $H^i(K) = 0$ 当 $i < a$。则对所有 $R$-模 $M$，当 $i < a$ 时
$\Ext^i(M, K) = 0$。因此只需证明：对任意有限 $R$-模 $M$，当
$i > b$ 时 $\Ext^i(M, K) = 0$，见引理\ref{more-algebra-lemma-injective-amplitude}。
由 Algebra，引理\ref{algebra-lemma-filter-Noetherian-module}，
模 $M$ 有有限滤过，其相继商形如 $R/\mathfrak p$，其中
$\mathfrak p$ 为素理想。若
$0 \to M_1 \to M \to M_2 \to 0$ 是短正合列，且对 $j = 1, 2$，
当 $i > b$ 时 $\Ext^i(M_j, K) = 0$，则当 $i > b$ 时
$\Ext^i(M, K) = 0$。因此不妨设 $M = R/\mathfrak p$。
若 $I \subset \mathfrak p$，则消灭性由假设得到。否则选取
$f \in I$，且 $f \not \in \mathfrak p$。考虑短正合列
$$
0 \to R/\mathfrak p \xrightarrow{f} R/\mathfrak p \to R/(\mathfrak p, f) \to 0
$$
$R$-模 $R/(\mathfrak p, f)$ 有有限滤过，其相继商为
$R/\mathfrak q$，且 $(\mathfrak p, f) \subset \mathfrak q$。
因此由 Noether 归纳及上述论证，可设消灭性对 $R/(\mathfrak p, f)$ 成立。
另一方面，模
$E^i = \Ext^i(R/\mathfrak p, K)$
由关于 $K$ 的假设（下有界且具有有限上同调模）、谱序列
(\ref{more-algebra-equation-first-ss-ext}) 以及 Algebra，引理
\ref{algebra-lemma-ext-noetherian}，都是有限模。
于是当 $i > b$ 时，$E^i$ 是满足 $E^i/fE^i = 0$ 的有限 $R$-模。
由 Nakayama 引理（Algebra，引理\ref{algebra-lemma-NAK}），$E^i$ 为零。
\end{proof}

\begin{lemma}
\label{more-algebra-lemma-finite-injective-dimension-Noetherian-local}
设 $(R, \mathfrak m, \kappa)$ 为局部 Noether 环。
设 $K \in D^+(R)$ 且具有有限上同调模。
则下列条件等价：
\begin{enumerate}
\item $K$ 具有有限内射维数；
\item 当 $i \gg 0$ 时 $\Ext^i_R(\kappa, K) = 0$。
\end{enumerate}
\end{lemma}

\begin{proof}
这是引理\ref{more-algebra-lemma-finite-injective-dimension-Noetherian-radical} 的特例。
\end{proof}

\section{接近投射的模}
\label{more-algebra-section-near-projective}

\noindent
文献中似乎有许多不同的“几乎投射模”定义。本节只讨论其中一种可能性。

\begin{lemma}
\label{more-algebra-lemma-ideal-factor-through-projective}
设 $R$ 为环。设 $M$、$N$ 为 $R$-模。
\begin{enumerate}
\item 给定 $R$-模映射 $\varphi : M \to N$，下列条件等价：
(a) $\varphi$ 经由某个投射 $R$-模分解；(b) $\varphi$ 经由某个自由 $R$-模分解。
\item 满足 (1) 中等价条件的所有 $\varphi : M \to N$ 构成
$\Hom_R(M, N)$ 的一个 $R$-子模。
\item 给定映射 $\psi : M' \to M$ 和 $\xi : N \to N'$，若
$\varphi : M \to N$ 满足 (1) 的等价条件，则
$\xi \circ \varphi \circ \psi : M' \to N'$ 也满足。
\end{enumerate}
\end{lemma}

\begin{proof}
(1)(a) 与 (1)(b) 的等价性由 Algebra，引理
\ref{algebra-lemma-characterize-projective} 得到。
若 $\varphi : M \to N$ 和 $\varphi' : M \to N$ 分别经由模
$P$ 和 $P'$ 分解，则 $\varphi + \varphi'$ 经由
$P \oplus P'$ 分解，而 $\lambda \varphi$ 对所有 $\lambda \in R$ 经由
$P$ 分解。这证明了 (2)。若 $\varphi : M \to N$ 经由 $P$ 分解，且
$\psi$、$\xi$ 如 (3) 所述，则 $\xi \circ \varphi \circ \psi$ 也经由
$P$ 分解。这证明了 (3)。
\end{proof}

\begin{lemma}
\label{more-algebra-lemma-factor-through-finite-projective}
设 $R$ 为环。设 $\varphi : M \to N$ 为 $R$-模映射。
若 $\varphi$ 经由投射模分解且 $M$ 是有限 $R$-模，则 $\varphi$
经由有限投射模分解。
\end{lemma}

\begin{proof}
由引理\ref{more-algebra-lemma-ideal-factor-through-projective}，可将
$\varphi = \tau \circ \sigma$ 分解，其中 $\sigma$ 的目标为某个集合
$I$ 所指标的 $\bigoplus_{i \in I} R$。选取 $M$ 的生成元
$x_1, \ldots, x_n$。写成 $\sigma(x_j) = (a_{ji})_{i \in I}$。
对每个 $j$，只有有限多个 $a_{ij}$ 非零。
因此 $\sigma$ 的像包含于有限自由 $R$-模中，结论成立。
\end{proof}

\noindent
设 $R$ 为环。注意，一个 $R$-模是投射的，当且仅当 $R$ 上的恒等映射
经由某个投射模分解。

\begin{lemma}
\label{more-algebra-lemma-near-projective}
设 $R$ 为环。设 $I \subset R$ 为理想。设 $M$ 为 $R$-模。
下列条件等价：
\begin{enumerate}
\item 对每个 $a \in I$，映射 $a : M \to M$ 经由投射 $R$-模分解；
\item 对每个 $a \in I$，映射 $a : M \to M$ 经由自由 $R$-模分解；以及
\item 对每个 $R$-模 $N$，$\Ext^1_R(M, N)$ 被 $I$ 湮灭。
\end{enumerate}
\end{lemma}

\begin{proof}
(1) 与 (2) 的等价性由引理\ref{more-algebra-lemma-ideal-factor-through-projective} 得到。
若 (1) 成立，则 (3) 成立，因为对任意 $N$ 和任意投射模 $P$，
$\Ext^1_R(P, N)$ 为零。
反之，设 (3) 成立。选取短正合列
$0 \to N \to P \to M \to 0$，其中 $P$ 投射（甚至可取自由）。
由假设，该列在 $\Ext^1_R(M, N)$ 中的相应元素被 $I$ 湮灭。
因此对每个 $a \in I$，映射 $a : M \to M$ 可经由满射 $P \to M$
分解，从而 (1) 成立。
\end{proof}

\noindent
为了方便讨论满足引理\ref{more-algebra-lemma-near-projective} 中等价条件的模，
我们为这一性质命名。

\begin{definition}
\label{more-algebra-definition-near-projective}
设 $R$ 为环。设 $I \subset R$ 为理想。设 $M$ 为 $R$-模。
若引理\ref{more-algebra-lemma-near-projective} 的等价条件成立，则称 $M$ 为
{\it $I$-投射模}\footnote{这是非标准记号。}
\end{definition}

\noindent
被 $I$ 湮灭的模都是 $I$-投射模。

\begin{lemma}
\label{more-algebra-lemma-torsion-near-projective}
设 $R$ 为环。设 $I \subset R$ 为理想。设 $M$ 为 $R$-模。
若 $M$ 被 $I$ 湮灭，则 $M$ 是 $I$-投射模。
\end{lemma}

\begin{proof}
这直接由定义以及零模是投射模这一事实得到。
\end{proof}

\begin{lemma}
\label{more-algebra-lemma-ses-near-projective}
设 $R$ 为环。设 $I \subset R$ 为理想。设
$$
0 \to K \to P \to M \to 0
$$
为 $R$-模的短正合列。
若 $M$ 是 $I$-投射模且 $P$ 是投射模，则 $K$ 是 $I$-投射模。
\end{lemma}

\begin{proof}
元素 $\text{id}_K \in \Hom_R(K, K)$ 映射到给定扩张在
$\Ext^1_R(M, K)$ 中的类。由假设，该类被每个 $a \in I$ 湮灭，
因此 $a : K \to K$ 经由 $K \to P$ 分解，结论成立。
\end{proof}

\begin{lemma}
\label{more-algebra-lemma-dual-near-projective}
设 $R$ 为环。设 $I \subset R$ 为理想。
若 $M$ 是有限的 $I$-投射 $R$-模，则
$M^\vee = \Hom_R(M, R)$ 是 $I$-投射模。
\end{lemma}

\begin{proof}
设 $M$ 有限且为 $I$-投射模。选取短正合列
$0 \to K \to R^{\oplus r} \to M \to 0$。
这给出单射 $M^\vee \to R^{\oplus r} = (R^{\oplus r})^\vee$。
由于该短正合列在 $\Ext^1_R(M, K)$ 中对应的扩张类被 $I$ 湮灭，
可知对每个 $a \in I$，存在映射 $M \to R^{\oplus r}$，使其与给定映射
$R^{\oplus r} \to M$ 的复合等于 $a : M \to M$。
取对偶可知，$a : M^\vee \to M^\vee$ 经由上面给出的
$M^\vee \to R^{\oplus r}$ 分解，结论成立。
\end{proof}

\section{Hom 复形}
\label{more-algebra-section-hom-complexes}

\noindent
设 $R$ 为环。设 $L^\bullet$ 和 $M^\bullet$ 为由 $R$-模组成的两个复形。
我们构造复形 $\Hom^\bullet(L^\bullet, M^\bullet)$。具体地，对每个 $n$，置
$$
\Hom^n(L^\bullet, M^\bullet) =
\prod\nolimits_{n = p + q} \Hom_R(L^{-q}, M^p)
$$
可以将 $\Hom^n$ 看作从 $L^\bullet$ 到 $M^\bullet$（视为分次模）的
所有齐次次数为 $n$ 的 $R$-线性映射所成的 $R$-模。
在此术语下，用如下规则定义微分：
$$
\text{d}(f) = \text{d}_M \circ f - (-1)^n f \circ \text{d}_L
$$
其中 $f \in \Hom^n(L^\bullet, M^\bullet)$。
我们略去验证 $\text{d}^2 = 0$。符号规则见第
\ref{more-algebra-section-sign-rules} 节。
此构造是 Differential Graded Algebra，例
\ref{dga-example-category-complexes} 的特例。
由构造立即得到
\begin{equation}
\label{more-algebra-equation-cohomology-hom-complex}
H^n(\Hom^\bullet(L^\bullet, M^\bullet)) = \Hom_{K(R)}(L^\bullet, M^\bullet[n])
\end{equation}
对所有 $n \in \mathbf{Z}$ 都成立。

\begin{lemma}
\label{more-algebra-lemma-compose}
设 $R$ 为环。给定由 $R$-模组成的复形
$K^\bullet, L^\bullet, M^\bullet$，存在典范同构
$$
\Hom^\bullet(K^\bullet, \Hom^\bullet(L^\bullet, M^\bullet))
=
\Hom^\bullet(\text{Tot}(K^\bullet \otimes_R L^\bullet), M^\bullet)
$$
其为由 $R$-模组成的复形的同构。
\end{lemma}

\begin{proof}
设左端有一个次数为 $n$ 的元素 $\alpha$。于是
$$
\alpha = (\alpha^{p, q}) \in
\prod\nolimits_{p + q = n} \Hom_R(K^{-q}, \Hom^p(L^\bullet, M^\bullet))
$$
每个 $\alpha^{p, q}$ 都是元素
$$
\alpha^{p, q} = (\alpha^{r, s, q}) \in
\prod\nolimits_{r + s + q = n} \Hom_R(K^{-q}, \Hom_R(L^{-s}, M^r))
$$
若作如下恒等识别：
\begin{equation}
\label{more-algebra-equation-identification}
\Hom_R(K^{-q}, \Hom_R(L^{-s}, M^r)) = \Hom_R(K^{-q} \otimes_R L^{-s}, M^r)
\end{equation}
则按符号规则有
\begin{align*}
\text{d}(\alpha^{r, s, q})
& =
\text{d}_{\Hom^\bullet(L^\bullet, M^\bullet)} \circ \alpha^{r, s, q}
- (-1)^n \alpha^{r, s, q} \circ \text{d}_K \\
& =
\text{d}_M \circ \alpha^{r, s, q}
- (-1)^{r + s} \alpha^{r, s, q} \circ \text{d}_L
- (-1)^{r + s + q} \alpha^{r, s, q} \circ \text{d}_K
\end{align*}
另一方面，若 $\beta$ 是右端次数为 $n$ 的元素，则
$$
\beta = (\beta^{r, s, q}) \in \prod\nolimits_{r + s + q = n}
\Hom_R(K^{-q} \otimes_R L^{-s}, M^r)
$$
并且由符号规则（Homology，定义
\ref{homology-definition-associated-simple-complex}）有
\begin{align*}
\text{d}(\beta^{r, s, q})
& =
\text{d}_M \circ \beta^{r, s, q}
- (-1)^n \beta^{r, s, q} \circ
\text{d}_{\text{Tot}(K^\bullet \otimes L^\bullet)} \\
& =
\text{d}_M \circ \beta^{r, s, q}
- (-1)^{r + s + q} \left(
\beta^{r, s, q} \circ \text{d}_K + (-1)^{-q} \beta^{r, s, q} \circ \text{d}_L
\right)
\end{align*}
因此，由恒等识别 (\ref{more-algebra-equation-identification}) 所诱导的映射确实是复形态射。
\end{proof}

\begin{remark}
\label{more-algebra-remark-internal-hom}
设 $R$ 为环。由 $R$-模复形构成的范畴 $\text{Comp}(R)$
是对称幺半范畴，其张量积由 $\text{Tot}(- \otimes_R -)$ 给出，见
引理\ref{more-algebra-lemma-symmetric-monoidal-cat-complexes}。
给定 $\text{Comp}(R)$ 中的 $L^\bullet$ 和 $M^\bullet$，
元素 $f \in \Hom^0(L^\bullet, M^\bullet)$ 定义复形映射
$f : L^\bullet \to M^\bullet$，当且仅当 $\text{d}(f) = 0$。
因此引理\ref{more-algebra-lemma-compose} 还告诉我们
$$
\Mor_{\text{Comp}(R)}(K^\bullet, \Hom^\bullet(L^\bullet, M^\bullet))
=
\Mor_{\text{Comp}(R)}(\text{Tot}(K^\bullet \otimes_R L^\bullet), M^\bullet)
$$
对 $\text{Comp}(R)$ 中的 $K^\bullet, L^\bullet, M^\bullet$ 具有函子性。
这意味着 $\Hom^\bullet( - , -)$ 是对称幺半范畴
$\text{Comp}(R)$ 的内 Hom，正如 Categories，
注\ref{categories-remark-internal-hom-monoidal} 所讨论的。
\end{remark}

\begin{lemma}
\label{more-algebra-lemma-composition}
设 $R$ 为环。给定由 $R$-模组成的复形
$K^\bullet, L^\bullet, M^\bullet$，存在典范复形态射
$$
\text{Tot}\left(
\Hom^\bullet(L^\bullet, M^\bullet) \otimes_R \Hom^\bullet(K^\bullet, L^\bullet)
\right)
\longrightarrow
\Hom^\bullet(K^\bullet, M^\bullet)
$$
其为由 $R$-模组成的复形的态射。
\end{lemma}

\begin{proof}
由注\ref{more-algebra-remark-internal-hom} 中的讨论，这样的典范映射的存在性
由 Categories，注\ref{categories-remark-internal-hom-monoidal} 得到。
我们也给出直接构造。

\medskip\noindent
左端的次数为 $n$ 的元素 $\alpha$ 为
$$
\alpha = (\alpha^{p, q}) \in
\bigoplus\nolimits_{p + q = n}
\Hom^p(L^\bullet, M^\bullet) \otimes_R \Hom^q(K^\bullet, L^\bullet)
$$
元素 $\alpha^{p, q}$ 是有限和
$\alpha^{p, q} = \sum \beta^p_i \otimes \gamma^q_i$，
其中
$$
\beta^p_i = (\beta^{r, s}_i)
\in \prod\nolimits_{r + s = p} \Hom_R(L^{-s}, M^r)
$$
以及
$$
\gamma^q_i = (\gamma^{u, v}_i)
\in \prod\nolimits_{u + v = q} \Hom_R(K^{-v}, L^u)
$$
该映射将 $\alpha$ 送到 $\delta = (\delta^{r, v})$，其中
$$
\delta^{r, v} =
\sum\nolimits_{i, s}  \beta^{r, s}_i
\circ \gamma^{-s, v}_i
\in \Hom_R(K^{-v}, M^r)
$$
给定 $r + v = n$ 时，此和是有限的，因为非零的
$\alpha^{p, q}$ 只有有限多个，从而非零的 $\beta^p_i$ 和
$\gamma^q_i$ 也只有有限多个。
按符号规则有
\begin{align*}
\text{d}(\alpha^{p, q})
& =
\text{d}_{\Hom^\bullet(L^\bullet, M^\bullet)}(\alpha^{p, q})
+ (-1)^p \text{d}_{\Hom^\bullet(K^\bullet, L^\bullet)}(\alpha^{p, q}) \\
& =
\sum \Big( \text{d}_M \circ \beta^p_i \circ \gamma^q_i
- (-1)^p \beta^p_i \circ \text{d}_L \circ \gamma^q_i \Big) \\
&
\quad + (-1)^p \sum \Big( \beta^p_i \circ \text{d}_L \circ \gamma^q_i
- (-1)^q \beta^p_i \circ \gamma^q_i \circ \text{d}_K \Big) \\
& =
\sum \Big( \text{d}_M \circ \beta^p_i \circ \gamma^q_i
-(-1)^n \beta^p_i \circ \gamma^q_i \circ \text{d}_K \Big)
\end{align*}
因此规则 $\alpha \mapsto \delta$ 与微分相容，结论得证。
\end{proof}


\begin{lemma}
\label{more-algebra-lemma-diagonal-better}
设 $R$ 为一个环。给定复形 $K^\bullet, L^\bullet, M^\bullet$
（它们都是 $R$-模复形），存在一个典范态射
$$
\text{Tot}(K^\bullet \otimes_R \Hom^\bullet(M^\bullet, L^\bullet))
\longrightarrow
\Hom^\bullet(M^\bullet, \text{Tot}(K^\bullet \otimes_R L^\bullet))
$$
它是 $R$-模复形之间的态射，并且关于这三个复形具有函子性。
\end{lemma}

\begin{proof}
由注\ref{more-algebra-remark-internal-hom} 中的讨论，这样一个典范映射的存在性
由 Categories，注\ref{categories-remark-internal-hom-monoidal} 得到。
我们也给出直接构造。

\medskip\noindent
设 $\alpha$ 是右端中次数为 $n$ 的元素。于是
$$
\alpha = (\alpha^{p, q}) \in \prod\nolimits_{p + q = n}
\Hom_R(M^{-q}, \text{Tot}^p(K^\bullet \otimes_R L^\bullet))
$$
每个 $\alpha^{p, q}$ 都是元素
$$
\alpha^{p, q} = (\alpha^{r, s, q}) \in
\Hom_R(M^{-q}, \bigoplus\nolimits_{r + s + q = n} K^r \otimes_R L^s)
$$
其中我们把 $\alpha^{r, s, q}$ 看作一族映射，使得对每个
$x \in M^{-q}$，只有有限多个 $\alpha^{r, s, q}(x)$ 非零。
根据符号规则有
\begin{align*}
\text{d}(\alpha^{r, s, q})
& =
\text{d}_{\text{Tot}(K^\bullet \otimes_R L^\bullet)} \circ \alpha^{r, s, q}
- (-1)^n \alpha^{r, s, q} \circ \text{d}_M \\
& =
\text{d}_K \circ \alpha^{r, s, q} + (-1)^r \text{d}_L \circ \alpha^{r, s, q}
- (-1)^n \alpha^{r, s, q} \circ \text{d}_M
\end{align*}
另一方面，若 $\beta$ 是左端中次数为 $n$ 的元素，则
$$
\beta = (\beta^{p, q}) \in
\bigoplus\nolimits_{p + q = n} K^p \otimes_R \Hom^q(M^\bullet, L^\bullet)
$$
我们可以写成 $\beta^{p, q} = \sum \gamma_i^p \otimes \delta_i^q$，其中
$\gamma_i^p \in K^p$，且
$$
\delta_i^q = (\delta_i^{r, s}) \in
\prod\nolimits_{r + s = q} \Hom_R(M^{-s}, L^r)
$$
根据符号规则有
\begin{align*}
\text{d}(\beta^{p, q})
& =
\text{d}_K(\beta^{p, q}) +
(-1)^p \text{d}_{\Hom^\bullet(M^\bullet, L^\bullet)}(\beta^{p, q}) \\
& =
\sum \text{d}_K(\gamma_i^p) \otimes \delta_i^q +
(-1)^p \sum \gamma_i^p \otimes
(\text{d}_L \circ \delta_i^q - (-1)^q \delta_i^q \circ \text{d}_M)
\end{align*}
我们将元素 $\beta$ 送到 $\alpha$，其中
$$
\alpha^{r, s, q} = c^{r, s, q}(\sum \gamma_i^r \otimes \delta_i^{s, q})
$$
其中 $c^{r, s, q} : K^r \otimes_R \Hom_R(M^{-q}, L^s) \to
\Hom_R(M^{-q}, K^r \otimes_R L^s)$ 是典范映射。
对于给定的 $\beta$ 和 $r$，只有有限多个 $\gamma_i^r$ 非零，
因此（对给定的 $r$）只有有限多个 $\alpha^{r, s, q}$
非零。于是这族映射满足上述条件，且该映射定义良好。
比较符号可知，这与微分相容。
\end{proof}

\begin{lemma}
\label{more-algebra-lemma-diagonal}
设 $R$ 为一个环。给定 $R$-模复形 $K^\bullet, L^\bullet$，存在一个典范态射
$$
K^\bullet
\longrightarrow
\Hom^\bullet(L^\bullet, \text{Tot}(K^\bullet \otimes_R L^\bullet))
$$
它是 $R$-模复形之间的态射，并且关于这两个复形具有函子性。
\end{lemma}

\begin{proof}
由注\ref{more-algebra-remark-internal-hom} 中的讨论，这样一个典范映射的存在性
由 Categories，注\ref{categories-remark-internal-hom-monoidal} 得到。
我们也给出直接构造。

\medskip\noindent
设 $\alpha$ 是右端中次数为 $n$ 的元素。于是
$$
\alpha = (\alpha^{p, q}) \in \prod\nolimits_{p + q = n}
\Hom_R(L^{-q}, \text{Tot}^p(K^\bullet \otimes_R L^\bullet))
$$
每个 $\alpha^{p, q}$ 都是元素
$$
\alpha^{p, q} = (\alpha^{r, s, q}) \in
\Hom_R(L^{-q}, \bigoplus\nolimits_{r + s + q = n} K^r \otimes_R L^s)
$$
其中我们把 $\alpha^{r, s, q}$ 看作一族映射，使得对每个
$x \in L^{-q}$，只有有限多个 $\alpha^{r, s, q}(x)$ 非零。
根据符号规则有
\begin{align*}
\text{d}(\alpha^{r, s, q})
& =
\text{d}_{\text{Tot}(K^\bullet \otimes_R L^\bullet)} \circ \alpha^{r, s, q}
- (-1)^n \alpha^{r, s, q} \circ \text{d}_L \\
& =
\text{d}_K \circ \alpha^{r, s, q} + (-1)^r \text{d}_L \circ \alpha^{r, s, q}
- (-1)^n \alpha^{r, s, q} \circ \text{d}_L
\end{align*}
现在将元素 $\beta \in K^n$ 送到 $\alpha$，令
$\alpha^{n, -q, q} = \beta \otimes \text{id}_{L^{-q}}$，
并在 $r \not = n$ 时令 $\alpha^{r, s, q} = 0$。
这是上述意义下的元素，因为对于固定的 $q$，只有一个
$\alpha^{r, s, q}$ 非零。对微分的描述表明这与微分相容。
\end{proof}


\begin{lemma}
\label{more-algebra-lemma-evaluate-and-more}
设 $R$ 为一个环。给定 $R$-模复形 $K^\bullet, L^\bullet, M^\bullet$，存在一个典范态射
$$
\text{Tot}(\Hom^\bullet(L^\bullet, M^\bullet) \otimes_R K^\bullet)
\longrightarrow
\Hom^\bullet(\Hom^\bullet(K^\bullet, L^\bullet), M^\bullet)
$$
它是 $R$-模复形之间的态射，并且关于这三个复形具有函子性。
\end{lemma}

\begin{proof}
由注\ref{more-algebra-remark-internal-hom} 中的讨论，这样一个典范映射的存在性
由 Categories，注\ref{categories-remark-internal-hom-monoidal} 得到。
我们也给出直接构造。

\medskip\noindent
考虑右端中次数为 $n$ 的元素 $\beta$。则
$$
\beta = (\beta^{p, s}) \in
\prod\nolimits_{p + s = n} \Hom_R(\Hom^{-s}(K^\bullet, L^\bullet), M^p)
$$
根据符号规则，
\begin{align*}
\text{d}(\beta^{p, s})
& =
\text{d}_M \circ \beta^{p, s}
- (-1)^n \beta^{p, s} \circ
\text{d}_{\Hom^\bullet(K^\bullet, L^\bullet)}
\end{align*}
我们可以如下描述最后一项：
$$
(\beta^{p, s} \circ \text{d}_{\Hom^\bullet(K^\bullet, L^\bullet)})(f) =
\beta^{p, s}(\text{d}_L \circ f - (-1)^{s + 1} f \circ \text{d}_K)
$$
若 $f \in \Hom^{-s - 1}(K^\bullet, L^\bullet)$。因此，在某种未具体说明的意义下，
$\text{d}(\beta^{p, s})$ 是三个项之和，其符号如下：
\begin{equation}
\label{more-algebra-equation-beta}
\text{d}(\beta^{p, s}) =
\text{d}_M(\beta^{p, s})
-(-1)^n\text{d}_L(\beta^{p, s}) +
(-1)^{p + 1}\text{d}_K(\beta^{p, s})
\end{equation}

\noindent
接下来考虑左端中次数为 $n$ 的元素 $\alpha$。
我们可以将它写成
$$
\alpha = (\alpha^{t, r}) \in
\bigoplus\nolimits_{t + r = n} \Hom^t(L^\bullet, M^\bullet) \otimes K^r
$$
每个 $\alpha^{t, r}$ 映到元素
$$
\alpha^{t, r} \mapsto (\alpha^{p, q, r}) \in
\prod\nolimits_{p + q = t} \Hom_R(L^{-q}, M^p) \otimes_R K^r
$$
根据符号规则，
\begin{align*}
\text{d}(\alpha^{p, q, r})
& =
\text{d}_{\Hom^\bullet(L^\bullet, M^\bullet)}(\alpha^{p, q, r})
+ (-1)^{p + q} \text{d}_K(\alpha^{p, q, r})
\end{align*}
若进一步写 $\alpha^{p, q, r} = \sum g_i^{p, q} \otimes k_i^r$，则
$$
\text{d}_{\Hom^\bullet(L^\bullet, M^\bullet)}(\alpha^{p, q, r}) =
\sum (\text{d}_M \circ g_i^{p, q}) \otimes k_i^r
- (-1)^{p + q} \sum (g_i^{p, q} \circ \text{d}_L) \otimes k_i^r
$$
因此，在某种未具体说明的意义下，$\text{d}(\alpha^{p, q, r})$
是三个项之和，其符号如下：
\begin{equation}
\label{more-algebra-equation-alpha}
\text{d}(\alpha^{p, q, r}) =
\text{d}_M(\alpha^{p, q, r})
-(-1)^{p + q}\text{d}_L(\alpha^{p, q, r}) +
(-1)^{p + q}\text{d}_K(\alpha^{p, q, r})
\end{equation}

\noindent
为了定义我们的映射，将使用典范映射
$$
c_{p, q, r} : 
\Hom_R(L^{-q}, M^p) \otimes_R K^r
\longrightarrow
\Hom_R(\Hom_R(K^r, L^{-q}), M^p)
$$
它将 $\varphi \otimes k$ 送到映射 $\psi \mapsto \varphi(\psi(k))$。
该映射关于三个变量均具有函子性。令 $s = q + r$，有包含
$$
\Hom_R(\Hom_R(K^r, L^{-q}), M^p) \subset
\Hom_R(\Hom^{-s}(K^\bullet, L^\bullet), M^p)
$$
这个包含来自投影
$\Hom^{-s}(K^\bullet, L^\bullet) \to \Hom_R(K^r, L^{-q})$。
由于 $\alpha^{p, q, r}$ 只对有限多个 $r$ 非零，
可知对给定的 $s$，满足 $q + r = s$ 的 $q, r$ 中只有有限多个。
因此可将 $\alpha$ 送到元素 $\beta$，其中
$$
\beta^{p, s} =
\sum\nolimits_{q + r = s} \epsilon_{p, q, r} c_{p, q, r}(\alpha^{p, q, r})
$$
% P02-E0247
其中上述和使用前面给出的包含映射，且
$\epsilon_{p, q, r} \in \{\pm 1\}$。比较
方程 (\ref{more-algebra-equation-beta}) 和 (\ref{more-algebra-equation-alpha}) 中的符号可知，
\begin{enumerate}
\item $\epsilon_{p, q, r} = \epsilon_{p + 1, q, r}$
\item $-(-1)^n\epsilon_{p, q, r} = -(-1)^{p + q}\epsilon_{p, q - 1, r}$
或等价地 $\epsilon_{p, q, r} = (-1)^r\epsilon_{p, q - 1, r}$
\item $(-1)^{p + 1}\epsilon_{p, q, r} = (-1)^{p + q}\epsilon_{p, q, r + 1}$
或等价地
$(-1)^{q + 1}\epsilon_{p, q, r} = \epsilon_{p, q, r + 1}$。
\end{enumerate}
一个好的解是取
% P02-E0248
$$
\epsilon_{p, r, s} = (-1)^{r + qr}
$$
这一符号的选取将在证明之后的备注中解释。
\end{proof}


\begin{remark}
\label{more-algebra-remark-sign-explanation}
现在说明，评价映射引理\ref{more-algebra-lemma-evaluate-and-more} 的直接构造中所使用的符号，
为何与由注\ref{more-algebra-remark-internal-hom} 以及 Categories，
注\ref{categories-remark-internal-hom-monoidal} 中的讨论所得到的构造相一致。
记
$- \otimes - = \text{Tot}(- \otimes_R -)$，并记
$hom(-, -) = \Hom^\bullet(-, -)$。由幺半范畴语言给出的构造要求在
$\text{Comp}(R)$ 中使用箭头
$$
hom(L^\bullet, M^\bullet) \otimes K^\bullet
\longrightarrow
hom(hom(K^\bullet, L^\bullet), M^\bullet)
$$
该箭头对应于箭头
$$
hom(L^\bullet, M^\bullet) \otimes K^\bullet \otimes hom(K^\bullet, L^\bullet)
\longrightarrow
M^\bullet
$$
后者先通过交换最后两个张量积的次序，再使用评价映射
$hom(K^\bullet, L^\bullet) \otimes K^\bullet \to L^\bullet$
和
% P02-E0249
$hom(L^\bullet, K^\bullet) \otimes L^\bullet \to M^\bullet$ 得到。
只有交换时会出现符号。具体而言，在同构
$$
K^\bullet \otimes hom(K^\bullet, L^\bullet)
\to
hom(K^\bullet, L^\bullet) \otimes K^\bullet
$$
中，在 $K^r \otimes_R \Hom_R(K^{-r'}, L^q)$ 上出现符号
$(-1)^{r(q + r')}$，见符号规则一节
\ref{more-algebra-section-sign-rules} 的条目
(\ref{more-algebra-item-flip-tensor-product})。
读者可以自行验证，由于我们用来描述进入内 Hom 的映射所采用的对应关系，
该符号只有在 $r = r'$ 时才有影响；此时得到
$(-1)^{r(q + r)} = (-1)^{r + qr}$，正如直接证明中所示。
\end{remark}


\section{符号规则}
\label{more-algebra-section-sign-rules}

\noindent
本节回顾目前使用的符号规则，并讨论它们的一些推论。
在环上的模复形范畴中讨论这些问题似乎也很合适，因为大多数有趣现象
已经在这一情形中出现。
我们诚挚希望读者不必使用本节中较为深奥的部分。

\medskip\noindent
在本节其余部分，固定一个环 $R$，并记
$M^\bullet$ 为一个 $R$-模复形，其微分为
$d^n_M : M^n \to M^{n + 1}$。
\begin{enumerate}
\item 第 $k$ 个平移复形 $M^\bullet[k]$ 的项为
$(M^\bullet[k])^n = M^{n + k}$，其微分为
$d_{M[k]}^n = (-1)^kd^{n + k}_M$，见
Homology，定义\ref{homology-definition-shift-cochain}。
\item 对于复形之间的映射 $f : M^\bullet \to N^\bullet$，我们定义
$f[k] : M^\bullet[k] \to N^\bullet[k]$，不引入符号，见
Homology，定义\ref{homology-definition-shift-cochain}。
\item 我们将 $H^n(M^\bullet[k])$ 与 $H^{n + k}(M^\bullet)$
视为相同，不引入符号，见
Homology，定义\ref{homology-definition-cohomology-shift}。
\item 复形短正合序列的边界映射按蛇引理定义，不引入符号，见
Homology，引理\ref{homology-lemma-long-exact-sequence-cochain}。
\item 与逐项分裂短正合序列
$0 \to K^\bullet \to L^\bullet \to M^\bullet \to 0$
相关的复形的三角为
$$
K^\bullet \to L^\bullet \to M^\bullet \to K^\bullet[1]
$$
其中，若 $s$ 和 $\pi$ 是相容的逐项分裂，则
$M^n \to K^{n + 1}$ 是映射 $\pi^{n + 1} \circ d^n_L \circ s^n$。
换言之，不引入符号。见
Derived Categories，定义\ref{derived-definition-distinguished-triangle}
及\ref{derived-definition-split-ses}。
\item 总复形 $\text{Tot}(M^\bullet \otimes_R N^\bullet)$
具有满足 Leibniz 规则的微分 $d$
$$
d(x \otimes y) = d(x) \otimes y + (-1)^{\deg(x)}x \otimes d(y)
$$
见 Homology，例
\ref{homology-example-double-complex-as-tensor-product-of} 和
Homology，定义\ref{homology-definition-associated-simple-complex}。
\item
\label{more-algebra-item-shift-tensor}
存在典范同构
$$
\text{Tot}(M^\bullet \otimes_R N^\bullet)[a + b] \to
\text{Tot}(M^\bullet[a] \otimes_R N^\bullet[b])
$$
其在直和项 $M^p \otimes_R N^q$ 上使用符号 $(-1)^{pb}$，
见 Homology，注\ref{homology-remark-shift-double-complex}。
通常考虑相应的平移映射会更方便：
$\text{Tot}(M^\bullet \otimes_R N^\bullet) \to
\text{Tot}(M^\bullet[a] \otimes_R N^\bullet[b])[-a - b]$。
\item 存在复形的典范同构
$$
\text{Tot}(
\text{Tot}(K^\bullet \otimes_R L^\bullet) \otimes_R M^\bullet) \to
\text{Tot}(K^\bullet \otimes_R \text{Tot}(L^\bullet \otimes_R M^\bullet))
$$
其定义不引入符号。见
\ref{more-algebra-section-symmetric-monoidal} 节。
\item
\label{more-algebra-item-flip-tensor-product}
存在典范同构
$$
\text{Tot}(L^\bullet \otimes_R M^\bullet) \to
\text{Tot}(M^\bullet \otimes_R L^\bullet)
$$
其在直和项 $L^p \otimes_R M^q$ 上使用符号 $(-1)^{pq}$。
见\ref{more-algebra-section-symmetric-monoidal} 节。
\end{enumerate}

在讨论关于 Hom 复形的符号约定之前，我们先依据上述约定
构造一个复形的对偶。

\begin{lemma}
\label{more-algebra-lemma-left-dual-module}
设 $R$ 为一个环，$M$ 为一个 $R$-模。设
$N, \eta, \epsilon$ 是 $R$-模的幺半范畴中 $M$ 的一个左对偶，见
Categories，定义\ref{categories-definition-dual}。则
\begin{enumerate}
\item $M$ 和 $N$ 都是有限投射 $R$-模；
\item 映射
$e : \Hom_R(M, R) \to N$，$\lambda \mapsto (\lambda \otimes 1)(\eta)$
是同构；
\item 对 $n \in N$ 和 $m \in M$ 有
$\epsilon(n, m) = e^{-1}(n)(m)$。
\end{enumerate}
\end{lemma}

\begin{proof}
假设意味着
$$
M \xrightarrow{\eta \otimes 1} M \otimes_R N \otimes_R M
\xrightarrow{1 \otimes \epsilon} M
\quad\text{且}\quad
N \xrightarrow{1 \otimes \eta} N \otimes_R M \otimes_R N
\xrightarrow{\epsilon \otimes 1} N
$$
均为恒等映射。可以选取有限自由模 $F$、一个 $R$-模映射
$F \to M$，以及 $\eta$ 的一个提升
$\tilde \eta : R \to F \otimes_R N$。得到交换图
$$
\xymatrix{
M \ar[rr]_-{\eta \otimes 1} \ar[rrd]_-{\tilde \eta \otimes 1} & &
M \otimes_R N \otimes_R M \ar[r]_-{1 \otimes \epsilon} &
M \\
& &
F \otimes_R N \otimes_R M \ar[u] \ar[r]^-{1 \otimes \epsilon} &
F \ar[u]
}
$$
这表明 $M$ 上的恒等映射经过一个有限自由模分解，因此
$M$ 是有限投射模。由对称性可知 $N$ 也是有限投射模。
这证明了 (1)。(2) 由 Categories，引理
\ref{categories-lemma-left-dual} 及其证明得到。
(3) 由该证明中的第一个等式得到。
\end{proof}

\begin{lemma}
\label{more-algebra-lemma-left-dual-complex}
设 $R$ 为一个环，$M^\bullet$ 为一个 $R$-模复形。
设 $N^\bullet, \eta, \epsilon$ 是 $R$-模复形的幺半范畴中
$M^\bullet$ 的一个左对偶。则
\begin{enumerate}
\item $M^\bullet$ 和 $N^\bullet$ 都有界；
\item $M^n$ 和 $N^n$ 都是有限投射 $R$-模；
\item 写成 $\epsilon = \sum \epsilon_n$，其中
$\epsilon_n : N^{-n} \otimes_R M^n \to R$，并写成
$\eta = \sum \eta_n$，其中
$\eta_n : R \to  M^n \otimes_R N^{-n}$，则
$(N^{-n}, \eta_n, \epsilon_n)$ 是引理
\ref{more-algebra-lemma-left-dual-module} 意义下 $M^n$ 的左对偶；
\item 微分 $d_N^n : N^n \to N^{n + 1}$ 等于映射
$$
N^n = \Hom_R(M^{-n}, R)
\xrightarrow{d_M^{-n - 1}}
\Hom_R(M^{-n - 1}, R) = N^{n + 1}
$$
的 $-(-1)^n$ 倍，其中等号是引理
\ref{more-algebra-lemma-left-dual-module} (2) 中给出的恒等。
\end{enumerate}
反之，给定由有限投射 $R$-模组成的有界复形 $M^\bullet$，
令 $N^n = \Hom_R(M^{-n}, R)$，微分如上；令
$\epsilon = \sum \epsilon_n$，其中
$\epsilon_n : N^{-n} \otimes_R M^n \to R$ 由评价给出；并令
$\eta = \sum \eta_n$，其中 $\eta_n : R \to M^n \otimes_R N^{-n}$
将 $1$ 映到
% P02-E0250
$\text{id}_{M_n}$，则得到 $R$-模复形的幺半范畴中
$M^\bullet$ 的一个左对偶。
\end{lemma}

\begin{proof}
由 Categories，定义\ref{categories-definition-dual}，
$(1 \otimes \epsilon) \circ (\eta \otimes 1) = \text{id}_{M^\bullet}$
且
$(\epsilon \otimes 1) \circ (1 \otimes \eta) = \text{id}_{N^\bullet}$，
因而立即有
$(1 \otimes \epsilon_n) \circ (\eta_n \otimes 1) = \text{id}_{M^n}$
以及
$(\epsilon_n \otimes 1) \circ (1 \otimes \eta_n) = \text{id}_{N^{-n}}$，
这证明了 (3)。由引理\ref{more-algebra-lemma-left-dual-module} 得到 (2)。
由于和 $\eta = \sum \eta_n$ 是有限的，得到 (1)。
由于 $\eta = \sum \eta_n$ 是复形映射
$R \to \text{Tot}(M^\bullet \otimes_R N^\bullet)$，由我们对
$\text{Tot}(M^\bullet \otimes_R N^\bullet)$ 的微分所选取的符号可知
$$
(d_M^{-n - 1} \otimes 1) \circ \eta_{-n - 1}  +
(-1)^n (1 \otimes d_N^{-n}) \circ \eta_{-n} = 0
$$
展开定义即证明了 (4)。要得到引理的最后一个陈述，只需将上述论证反向读取。
\end{proof}


\noindent
我们将使用引理\ref{more-algebra-lemma-left-dual-complex} 中对复形左对偶的描述，
% P02-E0251
作为 $\Hom$ 复形符号规则的动机。我们选择符号，使得
(\ref{more-algebra-item-compatible}) 成立。
我们继续讨论上述各种符号规则。
\begin{enumerate}
\setcounter{enumi}{9}
\item 给定复形 $K^\bullet$、$M^\bullet$，令
$\Hom^\bullet(M^\bullet, K^\bullet)$ 为如下复形，其项为
$$
\Hom^n(M^\bullet, K^\bullet) = \prod\nolimits_{n = p + q} \Hom_R(M^{-q}, K^p)
$$
微分由规则
$$
d(f) = d_K \circ f - (-1)^n f \circ d_M
$$
给出。
\item
\label{more-algebra-item-compatible}
上述选择满足：若 $M^\bullet$ 具有引理
\ref{more-algebra-lemma-left-dual-complex} 中的左对偶 $N^\bullet$，则存在典范同构
$$
\text{Tot}(K^\bullet \otimes_R N^\bullet)
\longrightarrow
\Hom^\bullet(M^\bullet, K^\bullet)
$$
其定义不引入符号，并通过
$N^q = \Hom_R(M^{-q}, R)$ 以及典范映射
$K^p \otimes_R \Hom_R(M^{-q}, R) \to \Hom_R(M^{-q}, K^p)$，
将直和项 $K^p \otimes_R N^q$ 送到直和项
$\Hom_R(M^{-q}, K^p)$。
\item 存在复合映射
$$
\text{Tot}(
\Hom^\bullet(L^\bullet, K^\bullet) \otimes_R
\Hom^\bullet(M^\bullet, L^\bullet))
\longrightarrow
\Hom^\bullet(M^\bullet, K^\bullet)
$$
其定义不引入符号，见引理\ref{more-algebra-lemma-composition}。
\item 存在典范同构
$$
\Hom^\bullet(K^\bullet, \Hom^\bullet(L^\bullet, M^\bullet))
=
\Hom^\bullet(\text{Tot}(K^\bullet \otimes_R L^\bullet), M^\bullet)
$$
其定义不引入符号，见引理\ref{more-algebra-lemma-compose}。
\item 存在典范映射
$$
\text{Tot}(K^\bullet \otimes_R \Hom^\bullet(M^\bullet, L^\bullet))
\longrightarrow
\Hom^\bullet(M^\bullet, \text{Tot}(K^\bullet \otimes_R L^\bullet))
$$
其定义不引入符号，见引理\ref{more-algebra-lemma-diagonal-better}。
\item 存在典范映射
$$
K^\bullet
\longrightarrow
\Hom^\bullet(L^\bullet, \text{Tot}(K^\bullet \otimes_R L^\bullet))
$$
其定义不引入符号，见引理\ref{more-algebra-lemma-diagonal}。
\item
由引理\ref{more-algebra-lemma-evaluate-and-more} 可得一个典范映射
% P02-E0252
$$
\text{Tot}(\Hom^\bullet(L^\bullet, M^\bullet) \otimes_R K^\bullet)
\longrightarrow
\Hom^\bullet(\Hom^\bullet(K^\bullet, L^\bullet), M^\bullet)
$$
它在 R-模
$\Hom_R(L^{-q}, M^p) \otimes_R K^r$ 上使用符号
$(-1)^{r + qr}$，其原因在备注\ref{more-algebra-remark-sign-explanation} 中解释。
\item
\label{more-algebra-item-evaluation}
取 $L^\bullet = M^\bullet$，并使用
$R \to \Hom^\bullet(M^\bullet, M^\bullet)$，上一项的映射变为评价映射
$$
ev : K^\bullet \longrightarrow
\Hom^\bullet(\Hom^\bullet(K^\bullet, M^\bullet), M^\bullet)
$$
它将 $x \in K^n$ 送到如下映射：将
$f \in \Hom^m(K^\bullet, M^\bullet)$ 送到 $(-1)^{nm}f(x)$。
\item
\label{more-algebra-item-shift-hom}
存在典范识别
$$
\Hom^\bullet(M^\bullet, K^\bullet)[a - b] \to
\Hom^\bullet(M^\bullet[b], K^\bullet[a])
$$
其中使用了符号。它定义为如下映射：相应的平移映射
$$
\Hom^\bullet(M^\bullet, K^\bullet) \to
\Hom^\bullet(M^\bullet[b], K^\bullet[a])[b - a]
$$
在模 $\Hom_R(M^{-q}, K^p)$ 上（其中 $p + q = n$）
使用符号 $(-1)^{nb}$。具体地，若
$f \in \Hom^n(M^\bullet, K^\bullet)$，则源上的
$$
d(f) = d_K \circ f - (-1)^n f \circ d_M
$$
而在目标上，$f$ 属于
$$
\left(\Hom^\bullet(M^\bullet[b], K^\bullet[a])[b - a]\right)^n =
\Hom^{n + b -a}(M^\bullet[b], K^\bullet[a])
$$
因此有
\begin{align*}
d(f) & =
(-1)^{b - a}
\left(d_{K[a]} \circ f - (-1)^{n + b - a} f \circ d_{M[b]}\right) \\
& =
(-1)^{b - a}
\left((-1)^a d_K \circ f - (-1)^{n + b - a} f \circ (-1)^b d_M \right) \\
& =
(-1)^b d_K \circ f - (-1)^{n + b} f \circ d_M
\end{align*}
由此可见，在次数 $n$ 中选取的符号 $(-1)^{nb}$
会针对这些微分产生一个复形映射。
\end{enumerate}


\section{导出 Hom}
\label{more-algebra-section-RHom}

\noindent
设 $R$ 为一个环。本节定义的导出 Hom 是一个函子
$$
D(R)^{opp} \times D(R) \longrightarrow D(R),\quad
(K, L) \longmapsto R\Hom_R(K, L)
$$
它是 $R$-模导出范畴中的内 Hom，亦即由公式
\begin{equation}
\label{more-algebra-equation-internal-hom}
\Hom_{D(R)}(K, R\Hom_R(L, M))
=
\Hom_{D(R)}(K \otimes_R^\mathbf{L} L, M)
\end{equation}
刻画。这里 $K, L, M$ 是 $D(R)$ 的对象。注意，根据 Yoneda
引理，这个公式在唯一同构意义下刻画这些对象。其构造如下。
选取表示 $M$ 的 $R$-模 K-内射复形 $I^\bullet$，选取表示 $L$
的复形 $L^\bullet$，并置
$$
R\Hom_R(L, M) = \Hom^\bullet(L^\bullet, I^\bullet)
$$
记号如 Hom 复形一节\ref{more-algebra-section-hom-complexes} 中所述。
这一构造的推广见 Differential Graded Algebra，
一节\ref{dga-section-restriction}。
由(\ref{more-algebra-equation-cohomology-hom-complex}) 以及
Derived Categories，引理\ref{derived-lemma-K-injective} 可知
% P02-E0253
\begin{equation}
\label{more-algebra-equation-h0-RHom}
H^n(R\Hom_R(L, M)) = \Hom_{D(R)}(L, M[n])
\end{equation}
对所有 $n \in \mathbf{Z}$ 成立。特别地，$D(R)$ 中的对象
$R\Hom_R(L, M)$ 是良定义的，即与 K-内射复形
$I^\bullet$ 的选取无关。

\begin{lemma}
\label{more-algebra-lemma-internal-hom}
设 $R$ 为一个环。设 $K, L, M$ 是 $D(R)$ 的对象。存在典范同构
$$
R\Hom_R(K, R\Hom_R(L, M)) = R\Hom_R(K \otimes_R^\mathbf{L} L, M)
$$
它位于 $D(R)$ 中，并且关于 $K, L, M$ 具有函子性；取 $H^0$
后恢复(\ref{more-algebra-equation-internal-hom})。
\end{lemma}

\begin{proof}
选取表示 $M$ 的 K-内射复形 $I^\bullet$，以及表示 $L$ 的
$R$-模 K-平坦复形 $L^\bullet$。对于任意 $R$-模复形 $K^\bullet$，有
$$
\Hom^\bullet(K^\bullet, \Hom^\bullet(L^\bullet, I^\bullet)) =
\Hom^\bullet(\text{Tot}(K^\bullet \otimes_R L^\bullet), I^\bullet)
$$
由引理\ref{more-algebra-lemma-compose} 得到。根据 $R\Hom$ 的定义，并且由于
$\text{Tot}(K^\bullet \otimes_R L^\bullet)$ 表示导出张量积，
即得本引理。
\end{proof}


\begin{lemma}
\label{more-algebra-lemma-RHom-out-of-projective}
设 $R$ 为一个环。设 $P^\bullet$ 是由投射 $R$-模组成的上有界复形，
设 $L^\bullet$ 是 $R$-模复形。则 $R\Hom_R(P^\bullet, L^\bullet)$
由复形 $\Hom^\bullet(P^\bullet, L^\bullet)$ 表示。
\end{lemma}

\begin{proof}
由(\ref{more-algebra-equation-cohomology-hom-complex}) 以及
Derived Categories，引理\ref{derived-lemma-morphisms-from-projective-complex}，
该复形的上同调群是“正确的”。因此，若选取拟同构
$L^\bullet \to I^\bullet$，其中 $I^\bullet$ 是 $R$-模的 K-内射复形，
则诱导映射
$$
\Hom^\bullet(P^\bullet, L^\bullet)
\longrightarrow
\Hom^\bullet(P^\bullet, I^\bullet)
$$
是拟同构。由于右端正是我们对
$R\Hom_R(P^\bullet, L^\bullet)$ 的定义，结论成立。
\end{proof}

\begin{lemma}
\label{more-algebra-lemma-internal-hom-evaluate}
设 $R$ 为一个环。设 $K, L, M$ 是 $D(R)$ 的对象。
存在典范态射
$$
R\Hom_R(L, M) \otimes_R^\mathbf{L} K
\longrightarrow
R\Hom_R(R\Hom_R(K, L), M)
$$
它位于 $D(R)$ 中，并且关于 $K$、$L$、$M$ 具有函子性。
\end{lemma}

\begin{proof}
选取表示 $M$ 的 K-内射复形 $I^\bullet$，
表示 $L$ 的 K-内射复形 $J^\bullet$，以及表示 $K$ 的
K-平坦复形 $K^\bullet$。该映射使用引理
\ref{more-algebra-lemma-evaluate-and-more} 的映射定义：
$$
\text{Tot}(\Hom^\bullet(J^\bullet, I^\bullet) \otimes_R K^\bullet)
\longrightarrow
\Hom^\bullet(\Hom^\bullet(K^\bullet, J^\bullet), I^\bullet)
$$
我们略去该映射关于 $D(R)$ 中三个对象均具有函子性的证明。
\end{proof}


\begin{lemma}
\label{more-algebra-lemma-internal-hom-composition}
设 $R$ 为一个环。给定 $D(R)$ 中的 $K, L, M$，存在典范态射
$$
R\Hom_R(L, M) \otimes_R^\mathbf{L} R\Hom_R(K, L) \longrightarrow R\Hom_R(K, M)
$$
它位于 $D(R)$ 中，并且关于 $K, L, M$ 具有函子性。
\end{lemma}

\begin{proof}
选取表示 $M$ 的 K-内射复形 $I^\bullet$，表示 $L$ 的
K-内射复形 $J^\bullet$，以及任意表示 $K$ 的 $R$-模复形 $K^\bullet$。
由引理\ref{more-algebra-lemma-composition}，存在复形映射
$$
\text{Tot}\left(
\Hom^\bullet(J^\bullet, I^\bullet) \otimes_R \Hom^\bullet(K^\bullet, J^\bullet)
\right)
\longrightarrow
\Hom^\bullet(K^\bullet, I^\bullet)
$$
$R$-模复形
$\Hom^\bullet(J^\bullet, I^\bullet)$、
$\Hom^\bullet(K^\bullet, J^\bullet)$ 和
$\Hom^\bullet(K^\bullet, I^\bullet)$ 分别表示
$R\Hom_R(L, M)$、$R\Hom_R(K, L)$ 和 $R\Hom_R(K, M)$。
若选取 K-平坦复形 $H^\bullet$ 以及拟同构
$H^\bullet \to \Hom^\bullet(K^\bullet, J^\bullet)$，则存在映射
$$
\text{Tot}\left(
\Hom^\bullet(J^\bullet, I^\bullet) \otimes_R H^\bullet
\right)
\longrightarrow
\text{Tot}\left(
\Hom^\bullet(J^\bullet, I^\bullet) \otimes_R \Hom^\bullet(K^\bullet, J^\bullet)
\right)
$$
其源表示
$R\Hom_R(L, M) \otimes_R^\mathbf{L} R\Hom_R(K, L)$。
复合上述两个显示的箭头即得到所需映射。我们略去证明该构造
具有函子性。
\end{proof}


\begin{lemma}
\label{more-algebra-lemma-internal-hom-diagonal-better}
设 $R$ 为一个环。给定 $D(R)$ 中的复形 $K, L, M$，
存在典范态射
$$
K \otimes_R^\mathbf{L} R\Hom_R(M, L)
\longrightarrow
R\Hom_R(M, K \otimes_R^\mathbf{L} L)
$$
它位于 $D(R)$ 中，并且关于 $K, L, M$ 具有函子性。
\end{lemma}

\begin{proof}
选取表示 $K$ 的 K-平坦复形 $K^\bullet$、表示 $L$ 的
K-内射复形 $I^\bullet$，以及任意表示 $M$ 的复形 $M^\bullet$。
选取拟同构
$\text{Tot}(K^\bullet \otimes_R I^\bullet) \to J^\bullet$，
其中 $J^\bullet$ 是 K-内射的。然后使用映射
$$
\text{Tot}\left(
K^\bullet \otimes_R \Hom^\bullet(M^\bullet, I^\bullet)
\right)
\to
\Hom^\bullet(M^\bullet, \text{Tot}(K^\bullet \otimes_R I^\bullet))
\to
\Hom^\bullet(M^\bullet, J^\bullet)
$$
其中第一个映射是引理\ref{more-algebra-lemma-diagonal-better} 中的映射。
\end{proof}

\begin{lemma}
\label{more-algebra-lemma-internal-hom-diagonal}
设 $R$ 为一个环。给定 $D(R)$ 中的复形 $K, L$，存在典范态射
$$
K \longrightarrow R\Hom_R(L, K \otimes_R^\mathbf{L} L)
$$
它位于 $D(R)$ 中，并且关于 $K$ 和 $L$ 具有函子性。
\end{lemma}

\begin{proof}
这是引理\ref{more-algebra-lemma-internal-hom-diagonal-better} 的特殊情形，
但我们也将直接证明。
选取表示 $K$ 的 K-平坦复形 $K^\bullet$ 和任意表示 $L$ 的复形
$L^\bullet$。选取拟同构
$\text{Tot}(K^\bullet \otimes_R L^\bullet) \to J^\bullet$，
其中 $J^\bullet$ 是 K-内射的。然后使用映射
$$
K^\bullet \to
\Hom^\bullet(L^\bullet, \text{Tot}(K^\bullet \otimes_R L^\bullet))
\to \Hom^\bullet(L^\bullet, J^\bullet)
$$
其中第一个映射是引理\ref{more-algebra-lemma-diagonal} 中的映射。
\end{proof}


\section{完美复形}
\label{more-algebra-section-perfect}

\noindent
完美复形是具有有限 Tor 维数的拟相干复形。
我们不会采用此作为定义，而是直接如下定义环上的完美复形。

\begin{definition}
\label{more-algebra-definition-perfect}
设 $R$ 为一个环。记 $D(R)$ 为 $R$-模阿贝尔范畴的导出范畴。
\begin{enumerate}
\item 若 $D(R)$ 中的对象 $K$ 拟同构于有限投射 $R$-模组成的有界复形，
则称其为{\it 完美}；
\item 若 $R$-模 $M$ 的 $M[0]$ 是 $D(R)$ 中的完美对象，
则称其为{\it 完美}。
\end{enumerate}
\end{definition}

\noindent
例如，在 Noether 环上，一个有限模是完美的，当且仅当它具有有限投射维数，见
引理\ref{more-algebra-lemma-perfect-module} 以及
Algebra，定义\ref{algebra-definition-finite-proj-dim}。

\begin{lemma}
\label{more-algebra-lemma-perfect}
设 $K^\bullet$ 是 $D(R)$ 中的对象。以下条件等价：
\begin{enumerate}
\item $K^\bullet$ 是完美的；
\item $K^\bullet$ 是拟相干的并且具有有限 Tor 维数。
\end{enumerate}
若 (1) 和 (2) 成立，且 $K^\bullet$ 的 Tor 振幅位于 $[a, b]$ 中，
则 $K^\bullet$ 拟同构于由有限投射 $R$-模组成的复形
$E^\bullet$，并且当 $i \not \in [a, b]$ 时 $E^i = 0$。
\end{lemma}

\begin{proof}
(1) 推出 (2) 是显然的，见引理
\ref{more-algebra-lemma-pseudo-coherent} 和\ref{more-algebra-lemma-tor-amplitude}。
设 (2) 成立且 $K^\bullet$ 的 Tor 振幅位于 $[a, b]$ 中。
特别地，当 $i > b$ 时 $H^i(K^\bullet) = 0$。
选取由有限自由 $R$-模组成的复形 $F^\bullet$，使得当 $i > b$ 时
$F^i = 0$，并取一个拟同构 $F^\bullet \to K^\bullet$
（引理\ref{more-algebra-lemma-pseudo-coherent}）。
置 $E^\bullet = \tau_{\geq a}F^\bullet$。注意，除 $E^a$ 是有限表示
$R$-模外，$E^i$ 都是有限自由的。由引理
\ref{more-algebra-lemma-last-one-flat}，$E^a$ 是平坦的。因此由
Algebra，引理\ref{algebra-lemma-finite-projective} 可知
$E^a$ 是有限投射的。
\end{proof}


\begin{lemma}
\label{more-algebra-lemma-perfect-module}
设 $M$ 是环 $R$ 上的模。以下条件等价：
\begin{enumerate}
\item $M$ 是完美模；
\item 存在解消
$$
0 \to F_d \to \ldots \to F_1 \to F_0 \to M \to 0
$$
其中每个 $F_i$ 都是有限投射 $R$-模。
\end{enumerate}
\end{lemma}

\begin{proof}
设 (2) 成立。令复形 $E^\bullet$ 满足 $E^{-i} = F_i$，
则其与 $M[0]$ 拟同构。因此 $M$ 是完美的。
反之，设 (1) 成立。由引理
\ref{more-algebra-lemma-perfect} 和\ref{more-algebra-lemma-n-pseudo-module}，可找到解消
$E^\bullet \to M$，其中 $E^{-i}$ 是有限自由 $R$-模。由
引理\ref{more-algebra-lemma-last-one-flat} 可知，对于充分大的某个 $d$，
% P02-E0254
$F_d = \Coker(E^{d - 1} \to E^d)$ 是平坦的。由
Algebra，引理\ref{algebra-lemma-finite-projective} 可知
$F_d$ 是有限投射的。因此
$$
0 \to F_d \to E^{-d+1} \to \ldots \to E^0 \to M \to 0
$$
是所需的解消。
\end{proof}

\begin{lemma}
\label{more-algebra-lemma-two-out-of-three-perfect}
设 $R$ 为一个环。设 $(K^\bullet, L^\bullet, M^\bullet, f, g, h)$
是 $D(R)$ 中的一个三角。若 $K^\bullet, L^\bullet, M^\bullet$
中任意两个是完美的，则第三个也是完美的。
\end{lemma}

\begin{proof}
结合引理\ref{more-algebra-lemma-perfect}、
\ref{more-algebra-lemma-two-out-of-three-pseudo-coherent} 以及
\ref{more-algebra-lemma-cone-tor-amplitude} 即得。
\end{proof}

\begin{lemma}
\label{more-algebra-lemma-summands-perfect}
设 $R$ 为一个环。若 $K^\bullet \oplus L^\bullet$ 是完美的，
则 $K^\bullet$ 和 $L^\bullet$ 也都是完美的。
\end{lemma}

\begin{proof}
由引理\ref{more-algebra-lemma-perfect}、
\ref{more-algebra-lemma-summands-pseudo-coherent} 以及
\ref{more-algebra-lemma-summands-tor-amplitude} 得到。
\end{proof}


\begin{lemma}
\label{more-algebra-lemma-complex-perfect-modules}
设 $R$ 为一个环。若 $K^\bullet$ 是由完美 $R$-模组成的有界复形，
则 $K^\bullet$ 是完美复形。
\end{lemma}

\begin{proof}
对有限复形的长度作归纳即可证明：使用引理
\ref{more-algebra-lemma-two-out-of-three-perfect} 以及朴素截断。
\end{proof}

\begin{lemma}
\label{more-algebra-lemma-cohomology-perfect}
设 $R$ 为一个环。若 $K^\bullet \in D^b(R)$ 且其所有上同调模都是完美的，
则 $K^\bullet$ 是完美的。
\end{lemma}

\begin{proof}
对有限复形的长度作归纳即可证明：使用引理
\ref{more-algebra-lemma-two-out-of-three-perfect} 以及典范截断。
\end{proof}

\begin{lemma}
\label{more-algebra-lemma-perfect-push-perfect}
设 $A \to B$ 为环映射。假设 $B$ 作为 $A$-模是完美的。
设 $K^\bullet$ 为完美的 $B$-模复形。则 $K^\bullet$
作为 $A$-模复形也是完美的。
\end{lemma}

\begin{proof}
使用引理\ref{more-algebra-lemma-perfect}，这可转化为关于拟相干模及
有限 Tor 维数模的相应结果。关于这些结果，见引理
\ref{more-algebra-lemma-finite-tor-dimension-push-tor-amplitude} 以及
\ref{more-algebra-lemma-finite-push-pseudo-coherent}。
\end{proof}

\begin{lemma}
\label{more-algebra-lemma-pull-perfect}
设 $A \to B$ 为环映射。
设 $K^\bullet$ 是由 $A$-模组成的完美复形。则
$K^\bullet \otimes_A^{\mathbf{L}} B$ 是由 $B$-模组成的完美复形。
\end{lemma}

\begin{proof}
使用引理\ref{more-algebra-lemma-perfect}，这可转化为关于拟相干模及
有限 Tor 维数模的相应结果。关于这些结果，见引理
\ref{more-algebra-lemma-pull-tor-amplitude} 以及
\ref{more-algebra-lemma-pull-pseudo-coherent}。
\end{proof}


\begin{lemma}
\label{more-algebra-lemma-flat-base-change-perfect}
设 $A \to B$ 为平坦环映射。设 $M$ 是完美 $A$-模。
则 $M \otimes_A B$ 是完美 $B$-模。
\end{lemma}

\begin{proof}
由引理\ref{more-algebra-lemma-perfect-module}，假设意味着 $M$ 具有由有限投射
% P02-E0255
$R$-模组成的有限解消 $F_\bullet$。由于 $A \to B$ 平坦，
复形 $F_\bullet \otimes_A B$ 是 $M \otimes_A B$ 的有限长度解消，
其各项为 $B$-上的有限投射模。因此 $M \otimes_A B$ 是完美的。
\end{proof}


\begin{lemma}
\label{more-algebra-lemma-tensor-perfect}
设 $R$ 为一个环。若 $K$ 和 $L$ 是 $D(R)$ 中的完美对象，
则 $K \otimes_R^\mathbf{L} L$ 也是完美对象。
\end{lemma}

\begin{proof}
可以如下使用定义证明。分别用由有限投射 $R$-模组成的有界复形
$K^\bullet$ 和 $L^\bullet$ 表示 $K$ 和 $L$。于是
$K \otimes_R^\mathbf{L} L$ 由有界复形
$\text{Tot}(K^\bullet \otimes_R L^\bullet)$ 表示。
该复形的项是模
% P02-E0256
$M^a \otimes_R L^b$ 的直和。由于 $M^a$ 和 $L^b$
是有限自由 $R$-模的直和因子，$M^a \otimes_R L^b$ 也是如此。因此
我们得出复形 $\text{Tot}(K^\bullet \otimes_R L^\bullet)$ 的各项
都是有限投射的。

\medskip\noindent
也可以使用引理\ref{more-algebra-lemma-perfect} 中对完美复形的刻画，
以及关于拟相干复形（引理\ref{more-algebra-lemma-tensor-pseudo-coherent}）
和 Tor 振幅（引理\ref{more-algebra-lemma-push-tor-amplitude}，取 $A = B = R$）
的相应引理来证明。
\end{proof}


\begin{lemma}
\label{more-algebra-lemma-glue-perfect}
设 $R$ 为一个环。设 $f_1, \ldots, f_r \in R$ 生成单位理想。
设 $K^\bullet$ 为 $R$-模复形。若对每个 $i$，复形
$K^\bullet \otimes_R R_{f_i}$ 都是完美的，则 $K^\bullet$ 是完美的。
\end{lemma}

\begin{proof}
使用引理\ref{more-algebra-lemma-perfect}，这可转化为关于拟相干模及
有限 Tor 维数模的相应结果。关于这些结果，见引理
\ref{more-algebra-lemma-glue-tor-amplitude} 以及
\ref{more-algebra-lemma-glue-pseudo-coherent}。
\end{proof}

\begin{lemma}
\label{more-algebra-lemma-flat-descent-perfect}
设 $R$ 为一个环。设 $K^\bullet$ 为 $R$-模复形。
设 $R \to R'$ 为忠实平坦环映射。若复形
$K^\bullet \otimes_R R'$ 是完美的，则 $K^\bullet$ 是完美的。
\end{lemma}

\begin{proof}
使用引理\ref{more-algebra-lemma-perfect}，这可转化为关于拟相干模及
有限 Tor 维数模的相应结果。关于这些结果，见引理
\ref{more-algebra-lemma-flat-descent-tor-amplitude} 以及
\ref{more-algebra-lemma-flat-descent-pseudo-coherent}。
\end{proof}


\begin{lemma}
\label{more-algebra-lemma-regular-perfect}
设 $R$ 为正则环。则
\begin{enumerate}
\item 一个 $R$-模是完美的，当且仅当它是有限 $R$-模；
\item $R$-模复形 $K^\bullet$ 是完美的，当且仅当
$K^\bullet \in D^b(R)$ 且每个 $H^i(K^\bullet)$ 都是有限 $R$-模。
\end{enumerate}
\end{lemma}

\begin{proof}
根据定义，任何完美 $R$-模都是有限的。反之，设 $M$ 是有限
$R$-模。选取解消
$$
\ldots \to F_2 \xrightarrow{d_2} F_1 \xrightarrow{d_1} F_0 \to M \to 0
$$
其中 $F_i$ 是有限自由 $R$-模
（Algebra，引理\ref{algebra-lemma-resolution-by-finite-free}）。
置 $M_i = \Ker(d_i)$。记
$U_i \subset \Spec(R)$ 为满足 $M_{i, \mathfrak p}$ 自由的素理想
$\mathfrak p$ 所成的集合；由 Algebra，引理
\ref{algebra-lemma-finitely-presented-localization-free}，$U_i$ 是开集。
% P02-E0257
有一个正合序列 $0 \to M_{i + 1} \to F_{i + 1} \to M_i \to 0$。
若 $\mathfrak p \in U_i$，则
$0 \to M_{i + 1, \mathfrak p} \to F_{i + 1, \mathfrak p} \to
M_{i, \mathfrak p} \to 0$ 分裂。因此 $M_{i + 1, \mathfrak p}$
是有限投射模，因而是自由模
（Algebra，引理\ref{algebra-lemma-finite-projective}）。
这说明 $U_i \subset U_{i + 1}$。我们断言
$\Spec(R) = \bigcup U_i$。事实上，对每个素理想
$\mathfrak p$，正则局部环 $R_\mathfrak p$ 由
Algebra，命题\ref{algebra-proposition-regular-finite-gl-dim} 知具有有限
整体维数。因此，当 $i \gg 0$ 时，例如由 Algebra，引理
\ref{algebra-lemma-independent-resolution}，$M_{i, \mathfrak p}$ 是有限投射的
（因而是自由的）。由于 $R$ 的谱是 Noether 的
（Algebra，引理\ref{algebra-lemma-Noetherian-topology}），我们得到
对某个 $n$ 有 $U_n = \Spec(R)$。于是由
Algebra，引理\ref{algebra-lemma-finite-projective}，$M_n$ 是投射
$R$-模。因此
$$
0 \to M_n \to F_n \to \ldots \to F_1 \to M \to 0
$$
是由有限投射模组成的有界解消，从而 $M$ 是完美的。
这证明了 (1)。

\medskip\noindent
设 $K^\bullet$ 是 $R$-模复形。若 $K^\bullet$ 完美，则它属于
$D^b(R)$，并且拟同构于由有限投射 $R$-模组成的有限复形，
所以每个 $H^i(K^\bullet)$ 都是有限 $R$-模（因为 $R$ 是
Noether 的）。反之，设 $K^\bullet$ 属于 $D^b(R)$，且每个
$H^i(K^\bullet)$ 都是有限 $R$-模。由 (1)，每个
$H^i(K^\bullet)$ 都是完美 $R$-模，因此由引理
\ref{more-algebra-lemma-cohomology-perfect}，$K^\bullet$ 是完美的。
\end{proof}
\begin{lemma}
\label{more-algebra-lemma-dual-perfect-complex}
设 $A$ 为一个环，设 $K \in D(A)$ 是完美对象。则
$K^\vee = R\Hom_A(K, A)$ 是一个完美复形，并且
$K \cong (K^\vee)^\vee$。存在函子性的同构
$$
L \otimes_A^\mathbf{L} K^\vee = R\Hom_A(K, L)
\quad\text{且}\quad
H^0(L \otimes_A^\mathbf{L} K^\vee) = \Ext_A^0(K, L)
$$
其中 $L \in D(A)$。
\end{lemma}

\begin{proof}
我们可以用有限投射 $A$-模组成的有界复形 $K^\bullet$ 表示 $K$。
由引理 \ref{more-algebra-lemma-RHom-out-of-projective}，对象 $K^\vee$ 由复形
$E^\bullet = \Hom^\bullet(K^\bullet, A)$ 表示。注意
$E^n = \Hom_A(K^{-n}, A)$ 同样是有限投射的。因此 $E^\bullet$
是由有限投射模组成的有界复形，从而 $K^\vee$ 是完美的。存在典范映射
$$
K^\bullet = \text{Tot}(\Hom^\bullet(A, A) \otimes_A K^\bullet)
\longrightarrow
\Hom^\bullet(\Hom^\bullet(K^\bullet, A), A)
$$
该映射在每个次数上（差一个符号）使用求值映射，参见引理
\ref{more-algebra-lemma-evaluate-and-more}。（符号规则参见
第 \ref{more-algebra-section-sign-rules} 节。）因此，由有限投射模的双对偶仍为自身，
该映射给出典范同构 $(K^\vee)^\vee \cong K$。

\medskip\noindent
第二个等式由第一个等式、引理 \ref{more-algebra-lemma-internal-hom} 以及
Derived Categories，见引理
\ref{derived-lemma-morphisms-from-projective-complex} 推出；还要结合
Ext 群的定义，参见 Derived Categories，第
\ref{derived-section-ext} 节。设 $L^\bullet$ 是表示 $L$ 的
$A$-模复形。根据第 \ref{more-algebra-section-sign-rules} 节的
第 \ref{more-algebra-item-compatible} 项，存在复形的典范同构
$$
\text{Tot}(L^\bullet \otimes_A E^\bullet)
\longrightarrow
\Hom^\bullet(K^\bullet, L^\bullet)
$$
这里使用了引理 \ref{more-algebra-lemma-left-dual-complex}，以说明在由 $A$-模组成的复形范畴中
$E^\bullet$ 是 $K^\bullet$ 的左对偶。于是得到第一个显示等式，证明完成。
\end{proof}

\begin{remark}
\label{more-algebra-remark-dual-perfect-complex}
设 $A$ 为一个环，设 $K \in D(A)$ 是完美对象。若 $K$ 的
tor-振幅为 $[a, b]$，则其对偶 $K^\vee$ 的 tor-振幅为 $[-b, -a]$。
这由引理 \ref{more-algebra-lemma-perfect} 的最后一个断言以及引理
\ref{more-algebra-lemma-dual-perfect-complex} 证明中对 $K^\vee$ 的描述得到。
\end{remark}

\begin{lemma}
\label{more-algebra-lemma-colim-and-lim-of-duals}
\begin{slogan}
完美对象系统的平凡对偶性。
\end{slogan}
设 $A$ 为一个环，设 $(K_n)_{n \in \mathbf{N}}$ 是 $D(A)$ 中完美对象
组成的系统。设 $K = \text{hocolim} K_n$ 是导出余极限
（Derived Categories，定义 \ref{derived-definition-derived-colimit}）。
则对 $D(A)$ 中任意对象 $E$，有
$$
R\Hom_A(K, E) = R\lim E \otimes^\mathbf{L}_A K_n^\vee
$$
其中 $(K_n^\vee)$ 是完美复形对偶组成的逆系统。
\end{lemma}

\begin{proof}
由引理 \ref{more-algebra-lemma-dual-perfect-complex}，有
$R\lim E \otimes^\mathbf{L}_A K_n^\vee =
R\lim R\Hom_A(K_n, E)$，
它嵌入如下区别三角
$$
R\lim R\Hom_A(K_n, E) \to
\prod R\Hom_A(K_n, E) \to
\prod R\Hom_A(K_n, E)
$$
由于 $K$ 同样嵌入区别三角
$\bigoplus K_n \to \bigoplus K_n \to K$，只需证明
$\prod R\Hom_A(K_n, E) = R\Hom_A(\bigoplus K_n, E)$。
这由（\ref{more-algebra-equation-internal-hom}）以及导出张量积与直和交换这一事实
形式地推出。
\end{proof}

\begin{lemma}
\label{more-algebra-lemma-colimit-perfect-complexes}
设 $R = \colim_{i \in I} R_i$ 是环的滤子余极限。
\begin{enumerate}
\item 给定 $D(R)$ 中的完美对象 $K$，存在 $i \in I$ 以及
$D(R_i)$ 中的完美对象 $K_i$，使得
$K \cong K_i \otimes_{R_i}^\mathbf{L} R$（在 $D(R)$ 中）。
\item 给定 $0 \in I$ 以及 $K_0, L_0 \in D(R_0)$，其中 $K_0$ 完美，
有
$$
\Hom_{D(R)}(K_0 \otimes_{R_0}^\mathbf{L} R, L_0 \otimes_{R_0}^\mathbf{L} R) =
\colim_{i \geq 0}
\Hom_{D(R_i)}(K_0 \otimes_{R_0}^\mathbf{L} R_i,
L_0 \otimes_{R_0}^\mathbf{L} R_i)
$$
\end{enumerate}
换言之，$R$ 上的完美复形三角范畴是 $R_i$ 上完美复形三角范畴的余极限。
\end{lemma}

\begin{proof}
我们将使用 Algebra，引理
\ref{algebra-lemma-module-map-property-in-colimit} 和
\ref{algebra-lemma-colimit-category-fp-modules} 的结果，以下不再特别说明。
这些引理特别表明，有限表示 $R$-模范畴是有限表示 $R_i$-模范畴的余极限。
由于有限投射模可以刻画为有限自由模的直和项
（Algebra，引理 \ref{algebra-lemma-finite-projective}），
有限投射模范畴也具有同样性质。这就由完美对象在 $D(R)$ 中的定义证明了
（1）。

\medskip\noindent
为证明（2），可用有限投射 $R_0$-模组成的有界复形 $K_0^\bullet$
表示 $K_0$。可用一个 K-平坦复形 $L_0^\bullet$ 表示 $L_0$
（引理 \ref{more-algebra-lemma-K-flat-resolution}）。于是
$$
\Hom_{D(R)}(K_0 \otimes_{R_0}^\mathbf{L} R, L_0 \otimes_{R_0}^\mathbf{L} R) =
\Hom_{K(R)}(K_0^\bullet \otimes_{R_0} R, L_0^\bullet \otimes_{R_0} R)
$$
由 Derived Categories，引理
\ref{derived-lemma-morphisms-from-projective-complex} 得到。对于 $\Hom$，将 $R$ 换成
$R_i$ 时也同样适用。由于右端只涉及有限多个项，并且
$$
\Hom_R(K_0^p \otimes_{R_0} R, L_0^q \otimes_{R_0} R) =
\colim_{i \geq 0}
\Hom_{R_i}(K_0^p \otimes_{R_0} R_i, L_0^q \otimes_{R_0} R_i)
$$
由证明开头所引的引理，再由于滤子余极限是正合的
（Algebra，引理 \ref{algebra-lemma-directed-colimit-exact}），
可得（2）。
\end{proof}
\section{复形的提升}
\label{more-algebra-section-lifting-complexes}

\noindent
设 $R$ 为一个环，设 $I \subset R$ 为一个理想。我们将考虑的提升问题如下。
假设给定 $D(R)$ 中的对象 $K$ 以及由 $R/I$-模组成的复形 $E^\bullet$，
且 $E^\bullet$ 在 $D(R)$ 中表示 $K \otimes_R^\mathbf{L} R/I$。
问题：是否存在一个由 $R$-模组成的复形 $P^\bullet$，提升 $E^\bullet$
并在 $D(R)$ 中表示 $K$？
一般来说这个问题的答案是否定的，但在良好情形下仍可有所作为。
我们先讨论无环复形的提升。

\begin{lemma}
\label{more-algebra-lemma-lift-acyclic-complex}
设 $R$ 为一个环，设 $I \subset R$ 为一个理想。设 $\mathcal{P}$
为一类 $R$-模。假设
\begin{enumerate}
\item 每个 $P \in \mathcal{P}$ 都是投射 $R$-模；
\item 若 $P_1 \in \mathcal{P}$ 且 $P_1 \oplus P_2 \in \mathcal{P}$，则
$P_2 \in \mathcal{P}$；
\item 若 $f : P_1 \to P_2$（其中 $P_1, P_2 \in \mathcal{P}$）模 $I$
满射，则 $f$ 是满射。
\end{enumerate}
则对任意项形如 $P/IP$（其中 $P \in \mathcal{P}$）的有界上无环复形
$E^\bullet$，存在一个项属于 $\mathcal{P}$ 的有界上无环复形
$P^\bullet$，提升 $E^\bullet$。
\end{lemma}

\begin{proof}
设 $E^i = 0$（当 $i > b$）。假设给定 $n$ 以及复形态射
$$
\xymatrix{
& & P^n \ar[r] \ar[d] & P^{n + 1} \ar[r] \ar[d] & \ldots \ar[r]
& P^b \ar[r] \ar[d] & 0 \ar[r] \ar[d] & \ldots \\
\ldots \ar[r] & E^{n - 1} \ar[r] &
E^n \ar[r] & E^{n + 1} \ar[r] & \ldots \ar[r] &
E^b \ar[r] & 0 \ar[r] & \ldots
}
$$
其中 $P^i \in \mathcal{P}$，且
$P^n \to P^{n + 1} \to \ldots \to P^b$ 在次数 $\geq n + 1$ 时无环，
并且竖直映射诱导同构 $P^i/IP^i \to E^i$。
在这种情况下，可以归纳地选取同构
$P^i = Z^i \oplus Z^{i + 1}$，使得映射 $P^i \to P^{i + 1}$
由
$Z^i \oplus Z^{i + 1} \to Z^{i + 1} \to Z^{i + 1} \oplus Z^{i + 2}$
给出。由性质（2）并作归纳论证可知 $Z^i \in \mathcal{P}$。
选取 $P^{n - 1} \in \mathcal{P}$ 以及同构
$P^{n - 1}/IP^{n - 1} \to E^{n - 1}$。由于 $P^{n - 1}$ 是投射的，且
$Z^n/IZ^n = \Im(E^{n - 1} \to E^n)$，
我们可以将映射 $P^{n - 1} \to E^{n - 1} \to E^n$ 提升为映射
$P^{n - 1} \to Z^n$。由性质（3)，映射 $P^{n - 1} \to Z^n$
是满射。因此，只需将 $P^{n - 1}$ 及刚构造的映射添加到 $P^n$ 的左侧，
就得到该图表的延拓。由于当 $n > b$ 时存在所需形式的图表，
对 $n$ 作归纳即可。
\end{proof}

\begin{lemma}
\label{more-algebra-lemma-lift-complex}
设 $R$ 为一个环，设 $I \subset R$ 为一个理想。设 $\mathcal{P}$
为一类 $R$-模。设 $K \in D(R)$，设 $E^\bullet$ 是由 $R/I$-模组成的
复形，表示 $K \otimes_R^\mathbf{L} R/I$。假设
\begin{enumerate}
\item 每个 $P \in \mathcal{P}$ 都是投射 $R$-模；
\item $P_1 \in \mathcal{P}$ 且 $P_1 \oplus P_2 \in \mathcal{P}$
当且仅当 $P_1, P_2 \in \mathcal{P}$；
\item 若 $f : P_1 \to P_2$（其中 $P_1, P_2 \in \mathcal{P}$）模 $I$
满射，则 $f$ 是满射；
\item $E^\bullet$ 有界上，且 $E^i$ 的形式为 $P/IP$，
其中 $P \in \mathcal{P}$；
\item $K$ 可由项属于 $\mathcal{P}$ 的有界上复形表示。
\end{enumerate}
则存在一个项属于 $\mathcal{P}$ 的有界上复形 $P^\bullet$，使得
$P^\bullet/IP^\bullet$ 与 $E^\bullet$ 同构，并在 $D(R)$ 中表示 $K$。
\end{lemma}

\begin{proof}
由假设（5），可用项属于 $\mathcal{P}$ 的有界上复形 $K^\bullet$ 表示 $K$。
于是 $K \otimes_R^\mathbf{L} R/I$ 由 $K^\bullet/IK^\bullet$ 表示。
由于根据（4），$E^\bullet$ 是由投射 $R/I$-模组成的有界上复形，
我们可以选取拟同构
$\delta : E^\bullet \to K^\bullet/IK^\bullet$
（Derived Categories，引理
\ref{derived-lemma-morphisms-from-projective-complex}）。
设 $C^\bullet$ 为 $\delta$ 的映射锥
（Derived Categories，定义 \ref{derived-definition-cone}）。
模 $C^i$ 是直和 $K^i/IK^i \oplus E^{i + 1}$，
因此由于（2）特别说明 $\mathcal{P}$ 对取直和封闭，
其形式为某个 $P/IP$（其中 $P \in \mathcal{P}$）。
% P02-E0258
由于 $C^\bullet$ 无环，可以应用引理
\ref{more-algebra-lemma-lift-acyclic-complex}，找到 $C^\bullet$ 的无环提升
$A^\bullet$。复形 $A^\bullet$ 有界上，且其项属于 $\mathcal{P}$。在
$$
\xymatrix{
K^\bullet \ar@{..>}[r] \ar[d] & A^\bullet \ar[d] \\
K^\bullet/IK^\bullet \ar[r] & C^\bullet \ar[r] & E^\bullet[1]
}
$$
中，根据 Derived Categories，引理
\ref{derived-lemma-morphisms-lift-projective}，可以找到使图表交换的虚线箭头。
下面将证明由（1）、（2）、（3）可知，对每个 $i$，映射
$K^i \to A^i$ 是直和项的包含映射。由性质（2），
$P^i = \Coker(K^i \to A^i)$ 属于 $\mathcal{P}$。
因此可取
$P^\bullet = \Coker(K^\bullet \to A^\bullet)[-1]$，从而得到结论。

\medskip\noindent
为完成证明，需证明如下命题：设 $f : P_1 \to P_2$，
其中 $P_1, P_2 \in \mathcal{P}$，且
$P_1/IP_1 \to P_2/IP_2$ 是带分裂的单射，其余核形式为某个
$P_3/IP_3$（其中 $P_3 \in \mathcal{P}$），则 $f$ 是带分裂的单射。
记 $E_i = P_i/IP_i$，则 $E_2 = E_1 \oplus E_3$。
由于 $P_2$ 是投射的，可以选取映射 $g : P_2 \to P_3$，
提升映射 $E_2 \to E_3$。由条件（3），映射 $g$ 是满射，
% P02-E0259
因而由于 $P_3$ 是投射的，它是可分裂的。置 $P_1' = \Ker(g)$，
并选取分裂 $P_2 = P'_1 \oplus P_3$。由（2），$P'_1 \in \mathcal{P}$。
我们不知道 $g \circ f = 0$，但可以考虑映射
$$
P_1 \xrightarrow{f} P_2 \xrightarrow{projection} P'_1
$$
模 $I$ 后的复合是同构。由于 $P'_1$ 是投射的，可以分裂
$P_1 = T \oplus P'_1$。若 $T = 0$，则完成证明，因为此时
$P_2 \to P'_1$ 是 $f$ 的分裂。由（2）可知 $T \in \mathcal{P}$。
模 $I$ 计算可得 $T/IT = 0$。由于 $0 \in \mathcal{P}$
（它是 $\mathcal{P}$ 中任意 $P$ 的直和项），可知映射 $0 \to T$
是满射，从而得到 $T = 0$，如所需。
\end{proof}
\begin{lemma}
\label{more-algebra-lemma-lift-complex-projectives}
设 $R$ 为一个环，设 $I \subset R$ 为一个理想。设 $E^\bullet$
为 $R/I$-模复形，设 $K$ 为 $D(R)$ 中的对象。假设
\begin{enumerate}
\item $E^\bullet$ 是由投射 $R/I$-模组成的有界上复形；
\item $K \otimes_R^\mathbf{L} R/I$ 在 $D(R/I)$ 中由
$E^\bullet$ 表示；
\item $I$ 是幂零理想。
\end{enumerate}
则存在一个由投射 $R$-模组成的有界上复形 $P^\bullet$，它在
$D(R)$ 中表示 $K$，使得 $P^\bullet \otimes_R R/I$
与 $E^\bullet$ 同构。
\end{lemma}

\begin{proof}
对所有投射 $R$-模组成的类 $\mathcal{P}$ 应用引理
\ref{more-algebra-lemma-lift-complex}。该引理的性质（1）和（2）是显然的。
性质（3）由 Nakayama 引理
（Algebra，引理 \ref{algebra-lemma-NAK}）得到。
性质（4）由投射 $R/I$-模可提升为投射 $R$-模这一事实得到，参见
Algebra，引理 \ref{algebra-lemma-lift-projective-module}。
为说明（5）成立，只需证明 $K$ 属于 $D^{-}(R)$。
由于 $E^\bullet$ 有界上，已知
$K \otimes_R^\mathbf{L} R/I$ 属于 $D^{-}(R/I)$；
结论由引理 \ref{more-algebra-lemma-check-tor-dimension-modulo-nilpotent} 得到。
\end{proof}

\begin{lemma}
\label{more-algebra-lemma-check-pseudo-coherent-modulo-nilpotent}
设 $R' \to R$ 是满的环映射，其核为幂零理想。设 $K' \in D(R')$，
并置 $K = K' \otimes_{R'}^\mathbf{L} R$。
则 $K$ 伪凝聚当且仅当 $K'$ 伪凝聚。
\end{lemma}

\begin{proof}
一个方向由引理 \ref{more-algebra-lemma-pull-pseudo-coherent} 得到。
对另一个方向，假设 $K$ 伪凝聚。由引理
\ref{more-algebra-lemma-pseudo-coherent}，可以用有限自由 $R$-模组成的有界上复形
$E^\bullet$ 表示 $K$。由引理
\ref{more-algebra-lemma-lift-complex-projectives}，可以用投射 $R'$-模组成的有界上复形
$P^\bullet$ 表示 $K'$，使得 $P^n \otimes_{R'} R = E^n$。
由 Nakayama 引理可知 $P^n$ 是有限自由的，从而 $K'$ 也是伪凝聚的。
\end{proof}

\begin{lemma}
\label{more-algebra-lemma-lift-complex-stably-frees}
设 $R$ 为一个环，设 $I \subset R$ 为一个理想。设 $E^\bullet$
为 $R/I$-模复形，设 $K$ 为 $D(R)$ 中的对象。假设
\begin{enumerate}
\item $E^\bullet$ 是由有限稳定自由 $R/I$-模组成的有界上复形；
\item $K \otimes_R^\mathbf{L} R/I$ 在 $D(R/I)$ 中由 $E^\bullet$ 表示；
% P02-E0260
\item $K^\bullet$ 是伪凝聚的；
\item $1 + I$ 中的每个元素都可逆。
\end{enumerate}
则存在一个由有限稳定自由 $R$-模组成的有界上复形 $P^\bullet$，
它在 $D(R)$ 中表示 $K$，使得 $P^\bullet \otimes_R R/I$
与 $E^\bullet$ 同构。此外，若 $E^i$ 自由，则 $P^i$ 自由。
\end{lemma}

\begin{proof}
对所有有限稳定自由 $R$-模组成的类 $\mathcal{P}$ 应用引理
\ref{more-algebra-lemma-lift-complex}。该引理的性质（1）是显然的。
性质（2）由引理 \ref{more-algebra-lemma-exact-category-stably-free} 得到。
性质（3）由 Nakayama 引理
（Algebra，引理 \ref{algebra-lemma-NAK}）得到。
性质（4）由有限稳定自由 $R/I$-模可提升为有限稳定自由 $R$-模这一事实
得到，参见引理 \ref{more-algebra-lemma-lift-stably-free}。
性质（5）成立，因为伪凝聚复形可由有限自由 $R$-模组成的有界上复形表示。
引理的最后一个断言由引理
\ref{more-algebra-lemma-isomorphic-finite-projective-lifts} 得到。
\end{proof}

\begin{lemma}
\label{more-algebra-lemma-lift-pseudo-coherent-from-residue-field}
设 $(R, \mathfrak m, \kappa)$ 是局部环，设 $K \in D(R)$ 伪凝聚。置
$d_i = \dim_\kappa H^i(K \otimes_R^\mathbf{L} \kappa)$。
则 $d_i < \infty$，并且对某个 $b \in \mathbf{Z}$，当 $i > b$ 时
$d_i = 0$。于是存在复形
$$
\ldots \to
R^{\oplus d_{b - 2}} \to
R^{\oplus d_{b - 1}} \to
R^{\oplus d_b} \to 0 \to \ldots
$$
在 $D(R)$ 中表示 $K$。此外，该复形在同构意义下唯一（！）。
\end{lemma}

\begin{proof}
注意，作为 $D(\kappa)$ 中的对象，
$K \otimes_R^\mathbf{L} \kappa$ 是伪凝聚的，参见引理
\ref{more-algebra-lemma-pull-pseudo-coherent}。因此其上同调空间是有限维的，
并在某个截断点以上消失。$D(\kappa)$ 的每个对象都在 $D(\kappa)$ 中
同构于一个微分为零的复形 $E^\bullet$。特别地，
$E^i \cong \kappa^{\oplus d_i}$ 是有限自由的。应用引理
\ref{more-algebra-lemma-lift-complex-stably-frees} 即得存在性。

\medskip\noindent
% P02-E0261
若有两个复形 $F^\bullet$ 和 $G^\bullet$，其中 $F^i$ 与 $G^i$
都是秩为 $d_i$ 的自由模，并且它们表示 $K$。
则可以选取复形映射 $\beta : F^\bullet \to G^\bullet$，表示同构
$F^\bullet \cong K \cong G^\bullet$，参见
Derived Categories，引理
\ref{derived-lemma-morphisms-from-projective-complex}。
所诱导的复形映射
$\beta \otimes 1 : F^\bullet \otimes_R^\mathbf{L} \kappa \to
G^\bullet \otimes_R^\mathbf{L} \kappa$
必为复形同构，因为
$F^\bullet \otimes_R^\mathbf{L} \kappa$ 和
$G^\bullet \otimes_R^\mathbf{L} \kappa$ 中的微分均为零。
因此 $\beta^i : F^i \to G^i$ 是有限自由 $R$-模之间的映射，
其模 $\mathfrak m$ 的约化是同构。故 $\beta^i$ 是同构，结论成立。
\end{proof}

\begin{lemma}
\label{more-algebra-lemma-lift-perfect-from-residue-field}
设 $R$ 为一个环，设 $\mathfrak p \subset R$ 为一个素理想，
设 $K \in D(R)$ 是完美的。置
$d_i =
\dim_{\kappa(\mathfrak p)} H^i(K \otimes_R^\mathbf{L} \kappa(\mathfrak p))$。
则 $d_i < \infty$，并且只有有限多个非零。
于是存在 $f \in R$、$f \not \in \mathfrak p$ 以及复形
$$
\ldots \to 0 \to R_f^{\oplus d_a} \to R_f^{\oplus d_{a + 1}} \to
\ldots \to
R_f^{\oplus d_{b - 1}} \to
R_f^{\oplus d_b} \to 0
\to \ldots
$$
在 $D(R_f)$ 中表示 $K \otimes_R^\mathbf{L} R_f$。
\end{lemma}

\begin{proof}
注意，作为 $D(\kappa(\mathfrak p))$ 中的对象，
$K \otimes_R^\mathbf{L} \kappa(\mathfrak p)$ 是完美的，参见引理
\ref{more-algebra-lemma-pull-perfect}。因此只有有限多个 $d_i$ 非零，且它们都是有限的。
应用引理 \ref{more-algebra-lemma-lift-pseudo-coherent-from-residue-field}，
可在局部环 $R_\mathfrak p$ 上得到一个具有所需形状、表示 $K$ 的复形。
对 $f \in R$、$f \not \in \mathfrak p$，有
$R_\mathfrak p = \colim R_f$
（Algebra，引理 \ref{algebra-lemma-localization-colimit}）。
由引理 \ref{more-algebra-lemma-colimit-perfect-complexes} 得到结论。
略去一些细节。
\end{proof}

\begin{lemma}
\label{more-algebra-lemma-compare-representatives-perfect}
设 $R$ 为一个环，设 $\mathfrak p \subset R$ 为一个素理想。设
$M^\bullet$ 与 $N^\bullet$ 是由有限投射 $R$-模组成的有界复形，
它们表示 $D(R)$ 中的同一个对象。则存在
$f \in R$、$f \not \in \mathfrak p$，使得存在复形的同构（！）
$$
M^\bullet_f \oplus P^\bullet \cong N^\bullet_f \oplus Q^\bullet
$$
% P02-E0262
其中 $P^\bullet$ 与 $Q^\bullet$ 是平凡复形的有限直和，即形如
$\ldots \to 0 \to R_f \xrightarrow{1} R_f \to 0 \to \ldots$
的复形（置于任意次数）。
\end{lemma}

\begin{proof}
若在局部化 $R_\mathfrak p$ 上有上述类型的同构，则利用
$R_\mathfrak p = \colim R_f$
（Algebra，引理 \ref{algebra-lemma-localization-colimit}），
可以将该同构下降为某个 $f$ 对应的 $R_f$ 上的同构。
因此可假设 $R$ 是局部环且 $\mathfrak p$ 是极大理想。
此时结论由表示完美对象的“极小”复形的唯一性得到，参见引理
\ref{more-algebra-lemma-lift-pseudo-coherent-from-residue-field}；还要使用任意复形都是
一个平凡复形与一个极小复形的直和这一事实
（Algebra，引理 \ref{algebra-lemma-add-trivial-complex}）。
\end{proof}
\begin{lemma}
\label{more-algebra-lemma-lift-complex-finite-projectives}
设 $R$ 为一个环，设 $I \subset R$ 为一个理想。设 $E^\bullet$
为 $R/I$-模复形，设 $K$ 为 $D(R)$ 中的对象。假设
\begin{enumerate}
\item $E^\bullet$ 是由有限投射 $R/I$-模组成的有界上复形；
\item $K \otimes_R^\mathbf{L} R/I$ 在 $D(R/I)$ 中由 $E^\bullet$ 表示；
\item $K$ 是伪凝聚的；
\item $(R, I)$ 是 Hensel 对。
\end{enumerate}
则存在一个由有限投射 $R$-模组成的有界上复形 $P^\bullet$，
它在 $D(R)$ 中表示 $K$，使得 $P^\bullet \otimes_R R/I$
与 $E^\bullet$ 同构。此外，若 $E^i$ 自由，则 $P^i$ 自由。
\end{lemma}

\begin{proof}
对所有有限投射 $R$-模组成的类 $\mathcal{P}$ 应用引理
\ref{more-algebra-lemma-lift-complex}。该引理的性质（1）和（2）是显然的。
性质（3）由 Nakayama 引理
（Algebra，引理 \ref{algebra-lemma-NAK}）得到。
性质（4）由有限投射 $R/I$-模可提升为有限投射 $R$-模这一事实得到，
参见引理 \ref{more-algebra-lemma-lift-finite-projective-module}。
性质（5）成立，因为伪凝聚复形可由有限自由 $R$-模组成的有界上复形表示。
因此引理 \ref{more-algebra-lemma-lift-complex} 可用，并得到所需的 $P^\bullet$。
引理的最后一个断言由引理
\ref{more-algebra-lemma-isomorphic-finite-projective-lifts} 得到。
\end{proof}



\section{复形的分裂}
\label{more-algebra-section-splitting}

\noindent
本节讨论使环的导出范畴中一个对象成为其截断之直和的条件。
我们的方法是在适当假设下使用下述引理，分裂典范区别三角
$$
\tau_{\leq i}K \to K \to \tau_{\geq i + 1}K \to (\tau_{\leq i}K)[1]
$$
（在 $D(R)$ 中），参见 Derived Categories，注
\ref{derived-remark-truncation-distinguished-triangle}。

\begin{lemma}
\label{more-algebra-lemma-splitting-unique}
设 $R$ 为一个环，设 $K$ 和 $L$ 是 $D(R)$ 中的对象。
假设 $L$ 的投射振幅为 $[a, b]$；例如，若 $L$ 是 Tor-振幅为
$[a, b]$ 的完美对象，则满足此条件。
\begin{enumerate}
\item 若对 $i \geq a$ 有 $H^i(K) = 0$，则
$\Hom_{D(R)}(L, K) = 0$。
\item 若对 $i \geq a + 1$ 有 $H^i(K) = 0$，则对任意给定的区别三角
$K \to M \to L \to K[1]$，存在 $D(R)$ 中的同构
$M \cong K \oplus L$，且与该区别三角中的映射相容。
\item 若对 $i \geq a$ 有 $H^i(K) = 0$，则（2）中的同构存在且唯一。
\end{enumerate}
\end{lemma}

\begin{proof}
$L$ 的投射振幅为 $[a, b]$ 这一假设意味着，我们可以用投射
$R$-模组成的复形 $L^\bullet$ 表示 $L$，且当
$i \not \in [a, b]$ 时 $L^i = 0$，参见定义
\ref{more-algebra-definition-projective-dimension}。
若 $L$ 是 Tor-振幅为 $[a, b]$ 的完美对象，则可以用有限投射
$R$-模组成的复形 $L^\bullet$ 表示 $L$，且当
$i \not \in [a, b]$ 时 $L^i = 0$，参见引理 \ref{more-algebra-lemma-perfect}。
若对 $i \geq a$ 有 $H^i(K) = 0$，则 $K$ 拟同构于
$\tau_{\leq a - 1}K$。因此可以用 $R$-模复形 $K^\bullet$ 表示 $K$，
且当 $i \geq a$ 时 $K^i = 0$。于是由 Derived Categories，引理
\ref{derived-lemma-morphisms-from-projective-complex} 得到
$$
\Hom_{D(R)}(L, K) = \Hom_{K(R)}(L^\bullet, K^\bullet) = 0
$$
这证明了（1）。在（2）的假设下，由（1）有
$\Hom_{D(R)}(L, K[1]) = 0$，所以根据 Derived Categories，引理
\ref{derived-lemma-split}，该区别三角是分裂的。
（3）中的唯一性由（1）得到。
\end{proof}

\begin{lemma}
\label{more-algebra-lemma-better-cut-complex-in-two}
设 $R$ 为一个环，设 $\mathfrak p \subset R$ 为一个素理想。
设 $K^\bullet$ 是伪凝聚 $R$-模复形。假设对某个
$i \in \mathbf{Z}$，映射
$$
H^i(K^\bullet) \otimes_R \kappa(\mathfrak p)
\longrightarrow
H^i(K^\bullet \otimes_R^{\mathbf{L}} \kappa(\mathfrak p))
$$
满射。则存在 $f \in R$、$f \not \in \mathfrak p$，使得
$\tau_{\geq i + 1}(K^\bullet \otimes_R R_f)$ 是 $D(R_f)$ 中
Tor-振幅为 $[i + 1, \infty]$ 的完美对象，并且在 $D(R_f)$ 中
存在典范同构
$$
K^\bullet \otimes_R R_f \cong
\tau_{\leq i}(K^\bullet \otimes_R R_f) \oplus
\tau_{\geq i + 1}(K^\bullet \otimes_R R_f)
$$
\end{lemma}

\begin{proof}
在本证明中，所有张量积均取在 $R$ 上，并记
$\kappa = \kappa(\mathfrak p)$。可假设 $K^\bullet$
是由有限自由 $R$-模组成的有界上复形。考察次数 $i$ 中的情形：
$$
\ldots \to K^{i - 1} \xrightarrow{d^{i - 1}} K^i \xrightarrow{d^i}
K^{i + 1} \to \ldots
$$
用下列公式定义 $0 \subset V \subset W \subset K^i \otimes \kappa$：
$$
V = \Im\left(
K^{i - 1} \otimes \kappa \to K^i \otimes \kappa
\right)
\quad\text{且}\quad
W = \Ker\left(
K^i \otimes \kappa \to K^{i + 1} \otimes \kappa
\right)
$$
置 $\dim(V) = r$、$\dim(W/V) = s$ 以及
$\dim(K^i \otimes \kappa/W) = t$。
可选取 $x_1, \ldots, x_r \in K^{i - 1}$，使它们在
$d^{i - 1}$ 下映为 $V$ 的一组基。由假设，可选取
$y_1, \ldots, y_s \in \Ker(d^i)$，其像构成 $W/V$ 的一组基。
最后，选取 $z_1, \ldots, z_t \in K^i$，其像构成
$K^i \otimes \kappa/W$ 的一组基。于是元素
$d^i(z_1), \ldots, d^i(z_t) \in K^{i + 1}$
在 $K^{i + 1} \otimes_R \kappa$ 中线性无关。
由 Algebra，引理 \ref{algebra-lemma-cokernel-flat}，将 $R$
替换为某个 $f \in R$、$f \not \in \mathfrak p$ 对应的 $R_f$
之后，可以假设
\begin{enumerate}
% P02-E0263
\item $d^i(x_a), y_b, z_c$ 是 $K^i$ 的一组 $R$-基；
\item $d^i(z_1), \ldots, d^i(z_t)$ 在 $K^{i + 1}$ 中
$R$-线性无关；
\item 商 $E^{i + 1} = K^{i + 1}/\sum Rd^i(z_c)$ 是有限投射的。
\end{enumerate}
由于 $d^i$ 将 $d^{i - 1}(x_a)$ 和 $y_b$ 映为零，
由条件（2）可知
$E^{i + 1} = \Coker(d^i : K^i \to K^{i + 1})$。
于是
$$
\tau_{\geq i + 1}K^\bullet =
(\ldots \to 0 \to E^{i + 1} \to K^{i + 2} \to \ldots)
$$
是由有限投射模组成的有界复形，其非零次数位于
$[i + 1, b]$，其中 $b$ 为某个整数。因此
$\tau_{\geq i + 1}K^\bullet$ 是振幅为 $[i + 1, b]$ 的完美对象。
由于 $\tau_{\leq i}K^\bullet$ 在次数 $> i$ 中没有上同调，
可以对区别三角
$$
\tau_{\leq i}K^\bullet \to K^\bullet \to \tau_{\geq i + 1}K^\bullet \to
(\tau_{\leq i}K^\bullet)[1]
$$
应用引理 \ref{more-algebra-lemma-splitting-unique}
（Derived Categories，注
\ref{derived-remark-truncation-distinguished-triangle}），得到结论。
\end{proof}
\begin{lemma}
\label{more-algebra-lemma-isolate-a-cohomology-group}
设 $R$ 为一个环，设 $\mathfrak p \subset R$ 为一个素理想。
设 $K^\bullet$ 是伪凝聚 $R$-模复形。假设对某个
$i \in \mathbf{Z}$，映射
$$
H^i(K^\bullet) \otimes_R \kappa(\mathfrak p)
\longrightarrow
H^i(K^\bullet \otimes_R^{\mathbf{L}} \kappa(\mathfrak p))
\quad\text{且}\quad
H^{i - 1}(K^\bullet) \otimes_R \kappa(\mathfrak p)
\longrightarrow
H^{i - 1}(K^\bullet \otimes_R^{\mathbf{L}} \kappa(\mathfrak p))
$$
均为满射。则存在 $f \in R$、$f \not \in \mathfrak p$，使得
\begin{enumerate}
\item $\tau_{\geq i + 1}(K^\bullet \otimes_R R_f)$ 是 $D(R_f)$ 中
Tor-振幅为 $[i + 1, \infty]$ 的完美对象；
\item $H^i(K^\bullet)_f$ 是有限自由 $R_f$-模；
\item 在 $D(R_f)$ 中存在典范直和分解
$$
K^\bullet \otimes_R R_f \cong
\tau_{\leq i - 1}(K^\bullet \otimes_R R_f) \oplus
H^i(K^\bullet)_f[-i] \oplus
\tau_{\geq i + 1}(K^\bullet \otimes_R R_f)
$$
。
\end{enumerate}
\end{lemma}

\begin{proof}
由引理 \ref{more-algebra-lemma-better-cut-complex-in-two} 得到（1），以及
$D(R_f)$ 中的分裂
$K^\bullet \otimes_R R_f = \tau_{\leq i}K^\bullet \otimes_R R_f \oplus
\tau_{\geq i + 1}K^\bullet \otimes_R R_f$。
再对 $\tau_{\leq i}K^\bullet \otimes_R R_f$ 应用一次引理
\ref{more-algebra-lemma-better-cut-complex-in-two}，可在适当选取 $f$ 后得到分裂
% P02-E0264
$\tau_{\leq i}K^\bullet \otimes_R R_f =
\tau_{\leq i - 1}K^\bullet \otimes_R R_f \oplus H^i(K^\bullet)_f$（在
$D(R_f)$ 中）；还得到 $H^i(K)_f$ 是平坦完美模，即有限投射模。
% P02-E0265
\end{proof}

\begin{lemma}
\label{more-algebra-lemma-cut-complex-in-two}
设 $R$ 为一个环，设 $\mathfrak p \subset R$ 为一个素理想，
设 $i \in \mathbf{Z}$。设 $K^\bullet$ 是伪凝聚 $R$-模复形，满足
$H^i(K^\bullet \otimes_R^{\mathbf{L}} \kappa(\mathfrak p)) = 0$。
则存在 $f \in R$、$f \not \in \mathfrak p$，以及在 $D(R_f)$ 中的
典范直和分解
$$
K^\bullet \otimes_R R_f =
\tau_{\geq i + 1}(K^\bullet \otimes_R R_f) \oplus
\tau_{\leq i - 1}(K^\bullet \otimes_R R_f)
$$
其中 $\tau_{\geq i + 1}(K^\bullet \otimes_R R_f)$ 是
Tor-振幅为 $[i + 1, \infty]$ 的完美复形。
\end{lemma}

\begin{proof}
这是引理 \ref{more-algebra-lemma-better-cut-complex-in-two} 的一个常用特殊情形。
直接证明如下。可假设 $K^\bullet$ 是由有限自由 $R$-模组成的有界上复形。
考察次数 $i$ 中的情形：
$$
\ldots \to K^{i - 2} \to R^{\oplus l}
\to R^{\oplus m} \to R^{\oplus n} \to K^{i + 2} \to \ldots
$$
设 $A$ 为对应于 $K^{i - 1} \to K^i$ 的 $m \times l$ 矩阵，
设 $B$ 为对应于 $K^i \to K^{i + 1}$ 的 $n \times m$ 矩阵。
假设说明 $A \bmod \mathfrak p$ 的秩为 $r$，而
$B \bmod \mathfrak p$ 的秩为 $m - r$。换言之，$A$ 有某个
$r \times r$ 子式 $a$ 不属于 $\mathfrak p$，而 $B$ 有某个
$(m - r) \times (m - r)$-子式 $b$ 不属于 $\mathfrak p$。
置 $f = ab$。则在逆转 $f$ 后，可以找到直和分解
$K^{i - 1} = R^{\oplus l - r} \oplus R^{\oplus r}$、
$K^i = R^{\oplus r} \oplus R^{\oplus m - r}$、
$K^{i + 1} = R^{\oplus m - r} \oplus R^{\oplus n - m + r}$，
使得模映射 $K^{i - 1} \to K^i$ 将
% P02-E0266
$R^{\oplus l - r}$ 映为零，并诱导 $R^{\oplus r}$ 到 $K^i$
相应直和项上的同构；且模映射 $K^i \to K^{i + 1}$ 将
$R^{\oplus r}$ 映为零，并诱导 $R^{\oplus m - r}$ 到
$K^{i + 1}$ 相应直和项上的同构。因此 $K^\bullet$ 与
$$
\ldots \to K^{i - 2} \to R^{\oplus l - r}
\to 0 \to R^{\oplus n - m + r} \to K^{i + 2} \to \ldots
$$
拟同构，一切显然。
\end{proof}

\begin{lemma}
\label{more-algebra-lemma-split-using-ext-injective}
设 $R$ 为一个环，设 $K \in D^-(R)$，设 $a \in \mathbf{Z}$。
假设对任意单射 $R$-模映射 $M \to M'$，映射
$\Ext^{-a}_R(K, M) \to \Ext^{-a}_R(K, M')$ 都是单射。
则存在唯一的直和分解
$K \cong \tau_{\leq a}K \oplus \tau_{\geq a + 1}K$，并且对某个
$b$，$\tau_{\geq a + 1}K$ 的投射振幅为 $[a + 1, b]$。
\end{lemma}

\begin{proof}
考虑 $D(R)$ 中的区别三角
$$
\tau_{\leq a}K \to K \to \tau_{\geq a + 1}K \to (\tau_{\leq a}K)[1]
$$
参见 Derived Categories，注
\ref{derived-remark-truncation-distinguished-triangle}。
注意
$\Ext^{-a}_R(\tau_{\leq a}K, M) = \Hom_R(H^a(K), M)$
以及
$\Ext^{-a - 1}_R(\tau_{\leq a}K, M) = 0$，参见
Derived Categories，引理 \ref{derived-lemma-negative-exts}。
因此 $\Ext$ 的长正合序列给出正合序列
$$
0 \to
\Ext^{-a}_R(\tau_{\geq a + 1}K, M) \to
\Ext^{-a}_R(K, M) \to
\Hom_R(H^a(K), M)
$$
它关于 $R$-模 $M$ 是函子性的。现在若 $I$ 是内射 $R$-模，则例如由
Derived Categories，引理
\ref{derived-lemma-compute-ext-resolutions}，有
$\Ext^{-a}_R(\tau_{\geq a + 1}K, I) = 0$。
由于每个模都可嵌入一个内射模，得到对每个 $R$-模 $M$ 都有
$\Ext^{-a}_R(\tau_{\geq a + 1}K, M) = 0$。
由引理 \ref{more-algebra-lemma-projective-amplitude} 可知，对某个 $b$，
$\tau_{\geq a + 1}K$ 的投射振幅为 $[a + 1, b]$
（此处使用了 $K$ 有界上这一条件）。
由引理 \ref{more-algebra-lemma-splitting-unique} 得到分裂。
\end{proof}

\begin{lemma}
\label{more-algebra-lemma-split-using-ext-zero}
设 $R$ 为一个环，设 $K \in D^-(R)$，设 $a \in \mathbf{Z}$。
假设对任意 $R$-模 $M$ 都有 $\Ext^{-a}_R(K, M) = 0$。
则存在唯一的直和分解
$K \cong \tau_{\leq a - 1}K \oplus \tau_{\geq a + 1}K$，
并且对某个 $b$，$\tau_{\geq a + 1}K$ 的投射振幅为 $[a + 1, b]$。
\end{lemma}

\begin{proof}
由引理 \ref{more-algebra-lemma-split-using-ext-injective}，有直和分解
$K \cong \tau_{\leq a}K \oplus \tau_{\geq a + 1}K$，
并且对某个 $b$，$\tau_{\geq a + 1}K$ 的投射振幅为 $[a + 1, b]$。
显然必有 $H^a(K) = 0$，从而在 $D(R)$ 中
$\tau_{\leq a}K = \tau_{\leq a - 1}K$。
\end{proof}













\section{识别完美复形}
\label{more-algebra-section-recognizing-perfect}

\noindent
下面给出一些可用来证明某些复形完美的引理。

\begin{lemma}
\label{more-algebra-lemma-lift-bounded-pseudo-coherent-to-perfect}
设 $R$ 为一个环，设 $\mathfrak p \subset R$ 为一个素理想。
设 $K$ 伪凝聚且有界下。置
$d_i = \dim_{\kappa(\mathfrak p)}
H^i(K \otimes_R^\mathbf{L} \kappa(\mathfrak p))$。
若存在 $a \in \mathbf{Z}$，使得当 $i < a$ 时 $d_i = 0$，
则存在 $f \in R$、$f \not \in \mathfrak p$ 以及复形
$$
\ldots \to 0 \to R_f^{\oplus d_a} \to R_f^{\oplus d_{a + 1}} \to
\ldots \to
R_f^{\oplus d_{b - 1}} \to
R_f^{\oplus d_b} \to 0
\to \ldots
$$
在 $D(R_f)$ 中表示 $K \otimes_R^\mathbf{L} R_f$。
特别地，$K \otimes_R^\mathbf{L} R_f$ 是完美的。
\end{lemma}

\begin{proof}
% P02-E0267
减小 $a$ 后，可以进一步假设对 $i < a$ 有
$H^i(K^\bullet) = 0$。由引理 \ref{more-algebra-lemma-cut-complex-in-two}，
将 $R$ 替换为某个 $f \in R$、$f \not \in \mathfrak p$
对应的 $R_f$ 后，可以在 $D(R)$ 中写成
$K^\bullet = \tau_{\leq a - 1}K^\bullet \oplus
\tau_{\geq a}K^\bullet$，其中 $\tau_{\geq a}K^\bullet$ 完美。
由于对 $i < a$ 有 $H^i(K^\bullet) = 0$，可知在 $D(R)$ 中
$\tau_{\leq a - 1}K^\bullet = 0$。因此 $K^\bullet$ 是完美的。
再应用引理 \ref{more-algebra-lemma-lift-perfect-from-residue-field} 即得结论。
\end{proof}
\begin{lemma}
\label{more-algebra-lemma-check-perfect-pointwise}
设 $R$ 为一个环，设 $a, b \in \mathbf{Z}$。
设 $K^\bullet$ 是伪凝聚 $R$-模复形。下列条件等价：
\begin{enumerate}
\item $K^\bullet$ 是 Tor-振幅为 $[a, b]$ 的完美对象；
\item 对每个素理想 $\mathfrak p$ 以及所有
$i \not \in [a, b]$，有
$H^i(K^\bullet \otimes_R^{\mathbf{L}} \kappa(\mathfrak p)) = 0$；
\item 对每个极大理想 $\mathfrak m$ 以及所有
$i \not \in [a, b]$，有
$H^i(K^\bullet \otimes_R^{\mathbf{L}} \kappa(\mathfrak m)) = 0$。
\end{enumerate}
\end{lemma}

\begin{proof}
略去蕴含（1）$\Rightarrow$（2）$\Rightarrow$（3）的证明。
假设（3）成立。设 $i \in \mathbf{Z}$ 且 $i \not \in [a, b]$。
由引理 \ref{more-algebra-lemma-cut-complex-in-two}，该假设说明对 $R$ 的所有极大理想，
都有 $H^i(K^\bullet)_{\mathfrak m} = 0$。因此
$H^i(K^\bullet) = 0$，参见 Algebra，引理
\ref{algebra-lemma-characterize-zero-local}。
此外，引理 \ref{more-algebra-lemma-cut-complex-in-two} 还说明，对每个极大理想
$\mathfrak m$，存在 $f \in R$、$f \not \in \mathfrak m$，
使得 $K^\bullet \otimes_R R_f$ 是 Tor-振幅为 $[a, b]$ 的完美对象。
由引理 \ref{more-algebra-lemma-glue-perfect} 和
\ref{more-algebra-lemma-glue-tor-amplitude} 得到结论。
\end{proof}

\begin{lemma}
\label{more-algebra-lemma-check-perfect-stalks}
设 $R$ 为一个环，设 $K^\bullet$ 是伪凝聚 $R$-模复形。
考虑下列条件：
\begin{enumerate}
\item $K^\bullet$ 是完美的；
\item 对每个素理想 $\mathfrak p$，复形
$K^\bullet \otimes_R R_{\mathfrak p}$ 是完美的；
\item 对每个极大理想 $\mathfrak m$，复形
$K^\bullet \otimes_R R_{\mathfrak m}$ 是完美的；
\item 对每个素理想 $\mathfrak p$，当所有 $i \ll 0$ 时，有
$H^i(K^\bullet \otimes_R^{\mathbf{L}} \kappa(\mathfrak p)) = 0$；
\item 对每个极大理想 $\mathfrak m$，当所有 $i \ll 0$ 时，有
$H^i(K^\bullet \otimes_R^{\mathbf{L}} \kappa(\mathfrak m)) = 0$。
\end{enumerate}
总有下列蕴含：
% P02-E0268
$$
(1) \Rightarrow (2) \Leftrightarrow (3) \Leftrightarrow (3)
\Leftrightarrow (4) \Leftrightarrow (5)
$$
若 $K^\bullet$ 有界下，则所有条件均等价。
\end{lemma}

\begin{proof}
由引理 \ref{more-algebra-lemma-pull-perfect} 可知（1）蕴含（2）。
（2）$\Rightarrow$（3）是显然的。由于每个素理想
$\mathfrak p$ 都包含于一个极大理想 $\mathfrak m$，对映射
$R_\mathfrak m \to R_\mathfrak p$ 应用引理
\ref{more-algebra-lemma-pull-perfect}，可知（3）蕴含（2）。
对剩余域映射
$R_\mathfrak p \to \kappa(\mathfrak p)$ 和
$R_\mathfrak m \to \kappa(\mathfrak m)$ 应用引理
\ref{more-algebra-lemma-pull-perfect}，可知（2）蕴含（4），而（3）蕴含（5）。

\medskip\noindent
假设 $R$ 是极大理想为 $\mathfrak m$、剩余域为 $\kappa$ 的局部环。
我们将证明，若对某个 $a$ 以及 $i < a$ 有
$H^i(K^\bullet \otimes^\mathbf{L} \kappa) = 0$，
则 $K$ 是完美的。
% P02-E0269
这将证明（4）蕴含（2）且（5）蕴含（3），从而得到引理的第一部分。
首先对 $i = a - 1$ 应用引理 \ref{more-algebra-lemma-cut-complex-in-two}，可知
在 $D(R)$ 中
$K^\bullet = \tau_{\leq a - 1}K^\bullet \oplus
\tau_{\geq a}K^\bullet$，其中 $\tau_{\geq a}K^\bullet$ 是
Tor-振幅包含于 $[a, \infty]$ 的完美对象。
为完成证明，需要说明 $\tau_{\leq a - 1}K$ 为零，
即其上同调群均为零。否则，设 $i$ 是使
$M = H^i(\tau_{\leq a - 1}K)$ 非零的最大指标。
由于 $\tau_{\leq a - 1}K^\bullet$ 是伪凝聚的，
$M$ 是有限 $R$-模（引理 \ref{more-algebra-lemma-finite-cohomology} 和
\ref{more-algebra-lemma-summands-pseudo-coherent}）。
因此由 Nakayama 引理（Algebra，引理 \ref{algebra-lemma-NAK}），
$M \otimes_R \kappa$ 非零。这说明
$$
H^i((\tau_{\leq a - 1}K^\bullet) \otimes_R^\mathbf{L} \kappa) =
H^i(K^\bullet \otimes_R^\mathbf{L} \kappa)
$$
非零，矛盾。

\medskip\noindent
假设等价条件（2）--（5）成立，且 $K^\bullet$ 有界下。
设对 $i < a$ 有 $H^i(K^\bullet) = 0$。选取 $R$ 的一个极大理想
$\mathfrak m$。只需证明存在 $f \in R$、$f \not \in \mathfrak m$，
使得 $K^\bullet \otimes_R^\mathbf{L} R_f$ 是完美的
（引理 \ref{more-algebra-lemma-glue-perfect} 和
Algebra，引理 \ref{algebra-lemma-quasi-compact}）。
这由引理 \ref{more-algebra-lemma-lift-bounded-pseudo-coherent-to-perfect} 得到。
\end{proof}

\begin{lemma}
\label{more-algebra-lemma-projective-amplitude-pseudo-coherent}
设 $R$ 为一个环，设 $K$ 是 $D(R)$ 中的伪凝聚对象。
设 $a, b \in \mathbf{Z}$。下列条件等价：
\begin{enumerate}
\item $K$ 的投射振幅为 $[a, b]$；
\item $K$ 是 Tor-振幅为 $[a, b]$ 的完美对象；
\item 对所有有限表示 $R$-模 $N$ 以及所有
$i \not \in [-b, -a]$，有 $\Ext^i_R(K, N) = 0$；
\item 对 $n > b$ 有 $H^n(K) = 0$，并且对所有有限表示
$R$-模 $N$ 以及所有 $i > -a$，有 $\Ext^i_R(K, N) = 0$；
\item 对 $n \not \in [a - 1, b]$ 有 $H^n(K) = 0$，并且对所有
有限表示 $R$-模 $N$，有 $\Ext^{-a + 1}_R(K, N) = 0$。
\end{enumerate}
\end{lemma}

\begin{proof}
由引理 \ref{more-algebra-lemma-perfect} 的最后一个断言可知（2）蕴含（1）。
若（1）成立，则 $K$ 可由投射模 $P^i$ 组成的复形表示，其中当
$i \not \in [a, b]$ 时 $P^i = 0$。
由于投射模是平坦模（它们是自由模的直和项），由引理
\ref{more-algebra-lemma-tor-amplitude} 可知 $K$ 的 Tor-振幅为 $[a, b]$。
于是由引理 \ref{more-algebra-lemma-perfect} 可知（2）成立。

\medskip\noindent
在条件（3）、（4）、（5）中，假设的 Ext 群
$\Ext^i_R(K, M)$ 对有限表示模 $M$ 的消失，等价于其对所有
$R$-模 $M$ 的消失；这由引理
\ref{more-algebra-lemma-pseudo-coherence-colimit-ext} 和
Algebra，引理 \ref{algebra-lemma-module-colimit-fp} 得到。
因此（1）、（3）、（4）、（5）的等价性由引理
\ref{more-algebra-lemma-projective-amplitude} 得到。
\end{proof}

\noindent
% P02-E0270
下面的引理可用于在多项式环
$B = A[x_1, \ldots, x_d]$ 上寻找完美复形。

\begin{lemma}
\label{more-algebra-lemma-perfect-over-polynomial-ring}
设 $A \to B$ 为一个环映射，设 $a, b \in \mathbf{Z}$，设
$d \geq 0$。设 $K^\bullet$ 是 $B$-模复形。假设
\begin{enumerate}
\item 环映射 $A \to B$ 平坦；
\item 对每个素理想 $\mathfrak p \subset A$，环
$B \otimes_A \kappa(\mathfrak p)$ 的整体维数有限且 $\leq d$；
\item 作为 $B$-模复形，$K^\bullet$ 是伪凝聚的；
\item 作为 $A$-模复形，$K^\bullet$ 的 Tor-振幅为 $[a, b]$。
\end{enumerate}
则作为 $B$-模复形，$K^\bullet$ 是完美的，且 Tor-振幅为
$[a - d, b]$。
\end{lemma}

\begin{proof}
可假设 $K^\bullet$ 是由有限自由 $B$-模组成的有界上复形。
特别地，作为 $A$-模复形，$K^\bullet$ 是平坦的，并且对任意
$A$-模 $M$ 有
$K^\bullet \otimes_A M = K^\bullet \otimes_A^{\mathbf{L}} M$。
对 $A$ 的每个素理想 $\mathfrak p$，复形
$$
K^\bullet \otimes_A \kappa(\mathfrak p)
$$
是由有限自由 $B \otimes_A \kappa(\mathfrak p)$-模组成的有界上复形，
且除 $i \in [a, b]$ 外所有 $H^i$ 均消失。
由于 $B \otimes_A \kappa(\mathfrak p)$ 的整体维数为 $d$，
由引理 \ref{more-algebra-lemma-finite-gl-dim-tor-dimension} 可知
$K^\bullet \otimes_A \kappa(\mathfrak p)$ 的 Tor-振幅为
$[a - d, b]$。设 $\mathfrak q$ 是 $B$ 的一个位于
$\mathfrak p$ 之上的素理想。由于
$K^\bullet \otimes_A \kappa(\mathfrak p)$ 是由自由
$B \otimes_A \kappa(\mathfrak p)$-模组成的有界上复形，有
\begin{align*}
K^\bullet \otimes_B^{\mathbf{L}} \kappa(\mathfrak q)
& = K^\bullet \otimes_B \kappa(\mathfrak q) \\
& = (K^\bullet \otimes_A \kappa(\mathfrak p))
\otimes_{B \otimes_A \kappa(\mathfrak p)} \kappa(\mathfrak q) \\
& = (K^\bullet \otimes_A \kappa(\mathfrak p))
\otimes^{\mathbf{L}}_{B \otimes_A \kappa(\mathfrak p)} \kappa(\mathfrak q)
\end{align*}
所以以上论证说明，当 $i \not \in [a - d, b]$ 时，
$H^i(K^\bullet \otimes_B^{\mathbf{L}} \kappa(\mathfrak q)) = 0$。
由引理 \ref{more-algebra-lemma-check-perfect-pointwise} 得到结论。
\end{proof}
\noindent
下述引理是引理 \ref{more-algebra-lemma-perfect-over-polynomial-ring} 的局部版本。
它可用于在正则局部环上寻找完美复形。

\begin{lemma}
\label{more-algebra-lemma-perfect-over-regular-local-ring}
设 $A \to B$ 是局部环同态，设 $a, b \in \mathbf{Z}$，设
$d \geq 0$。设 $K^\bullet$ 是 $B$-模复形。假设
\begin{enumerate}
\item 环映射 $A \to B$ 平坦；
\item 环 $B/\mathfrak m_AB$ 是维数为 $d$ 的正则环；
\item 作为 $B$-模复形，$K^\bullet$ 是伪凝聚的；
\item 作为 $A$-模复形，$K^\bullet$ 的 Tor-振幅为 $[a, b]$；
事实上，只需
$H^i(K^\bullet \otimes_A^\mathbf{L} \kappa(\mathfrak m_A))$
仅对 $i \in [a, b]$ 非零即可。
\end{enumerate}
则作为 $B$-模复形，$K^\bullet$ 是完美的，且 Tor-振幅为
$[a - d, b]$。
\end{lemma}

\begin{proof}
由（3），可假设 $K^\bullet$ 是由有限自由 $B$-模组成的有界上复形。
计算得
\begin{align*}
K^\bullet \otimes_B^{\mathbf{L}} \kappa(\mathfrak m_B)
& = K^\bullet \otimes_B \kappa(\mathfrak m_B) \\
& = (K^\bullet \otimes_A \kappa(\mathfrak m_A))
\otimes_{B/\mathfrak m_A B} \kappa(\mathfrak m_B) \\
& = (K^\bullet \otimes_A \kappa(\mathfrak m_A))
\otimes^{\mathbf{L}}_{B/\mathfrak m_A B} \kappa(\mathfrak m_B)
\end{align*}
% P02-E0271
第一个等式成立，是因为 $K^\bullet$ 是由平坦 $B$-模组成的有界上复形。
第二个等式由张量积的基本性质得到。第三个等式成立，是因为
$K^\bullet \otimes_A \kappa(\mathfrak m_A) =
K^\bullet/ \mathfrak m_A K^\bullet$
是由平坦 $B/\mathfrak m_A B$-模组成的有界上复形。
由（1），$K^\bullet$ 是由平坦 $A$-模组成的有界上复形；由假设（4），
复形 $K^\bullet \otimes_A \kappa(\mathfrak m_A)$ 的上同调模 $H^i$
仅对 $i \in [a, b]$ 非零。因此，例
\ref{more-algebra-example-cohomology-complex-tensored} 的谱序列，以及
$B/\mathfrak m_AB$ 的整体维数有限且为 $d$ 这一事实
（由（2）和 Algebra，命题
\ref{algebra-proposition-regular-finite-gl-dim}），说明当
$j \not \in [a - d, b]$ 时，
$H^j(K^\bullet \otimes_B^{\mathbf{L}} \kappa(\mathfrak m_B))$ 为零。
% P02-E0272
由引理 \ref{more-algebra-lemma-check-perfect-pointwise} 即完成证明。
\end{proof}





\section{完美复形的刻画}
\label{more-algebra-section-perfect-compact}

\noindent
本节证明，完美复形恰为环的导出范畴中的紧对象。先证明如下结论。

\begin{lemma}
\label{more-algebra-lemma-perfect-ring-classical-generator}
设 $R$ 为一个环。完美对象组成的全子范畴
$D_{perf}(R) \subset D(R)$ 是包含 $R = R[0]$ 的最小严格全、饱和三角子范畴。
换言之，$D_{perf}(R) = \langle R \rangle$。
特别地，$R$ 是 $D_{perf}(R)$ 的经典生成元。
\end{lemma}

\begin{proof}
关于该断言的含义，参见 Derived Categories，定义
\ref{derived-definition-saturated} 和
\ref{derived-definition-generators}。
引理 \ref{more-algebra-lemma-two-out-of-three-perfect} 和
\ref{more-algebra-lemma-summands-perfect} 已说明
$D_{perf}(R) \subset D(R)$ 是 $D(R)$ 的严格全、饱和三角子范畴。
当然，$R \in D_{perf}(R)$。

\medskip\noindent
回忆 $\langle R \rangle = \bigcup \langle R \rangle_n$。
为完成证明，将说明：若 $M \in D_{perf}(R)$ 由
$$
\ldots \to 0 \to M^a \to M^{a + 1} \to \ldots \to M^b \to 0 \to \ldots
$$
表示，其中 $M^i$ 为有限投射模，则
$M \in \langle R \rangle_{b - a + 1}$。对 $b - a$ 作归纳。
按定义，$\langle R \rangle_1$ 包含置于任意次数的每个有限投射
$R$-模；这处理了归纳的基础情形 $b - a = 0$。
一般地，考虑区别三角
% P02-E0273
$$
M_b[-b] \to M^\bullet \to \sigma_{\leq b - 1}M^\bullet \to M_b[-b + 1]
$$
由归纳假设，截断复形 $\sigma_{\leq b - 1}M^\bullet$ 属于
$\langle R \rangle_{b - a}$，而 $M_b[-b]$ 属于
$\langle R \rangle_1$。因此按定义，
$M^\bullet \in \langle R \rangle_{b - a + 1}$。
\end{proof}

\noindent
设 $R$ 为一个环。回忆 $D(R)$ 有直和，且它们由复形的直和直接给出，
参见 Derived Categories，引理 \ref{derived-lemma-direct-sums}。
本节各引理中将不再特别说明这一点。

\begin{lemma}
\label{more-algebra-lemma-commutes-with-countable-sums}
设 $R$ 为一个环。设 $K \in D(R)$ 是满足如下条件的对象：
对每个可数对象族 $E_n \in D(R)$，典范映射
$$
\bigoplus \Hom_{D(R)}(K, E_n) \longrightarrow \Hom_{D(R)}(K, \bigoplus E_n)
$$
是双射。则对任意以 $\mathbf{N}$ 为指标的复形系统 $L_n^\bullet$，映射
$$
\colim \Hom_{D(R)}(K, L^\bullet_n) \longrightarrow \Hom_{D(R)}(K, L^\bullet)
$$
是双射，其中 $L^\bullet$ 是逐项余极限，即对所有
$m \in \mathbf{Z}$ 有 $L^m = \colim L_n^m$。
\end{lemma}

\begin{proof}
考虑复形的短正合序列
$$
0 \to \bigoplus L_n^\bullet \to \bigoplus L_n^\bullet \to L^\bullet \to 0
$$
其中第一个映射在次数 $n$ 中由 $1 - t_n$ 给出，
$t_n : L_n^\bullet \to L_{n + 1}^\bullet$ 为转移映射。
由 Derived Categories，引理
\ref{derived-lemma-derived-canonical-delta-functor}，
这是 $D(R)$ 中的区别三角。应用同调函子
$\Hom_{D(R)}(K, -)$，参见 Derived Categories，引理
\ref{derived-lemma-representable-homological}。
% P02-E0274
于是得到长正合上同调序列
$$
\xymatrix{
& \ldots \ar[r] & \Hom_{D(R)}(K, \colim L^\bullet_n[-1]) \ar[lld] \\
\Hom_{D(R)}(K, \bigoplus L^\bullet_n) \ar[r] &
\Hom_{D(R)}(K, \bigoplus L^\bullet_n) \ar[r] &
\Hom_{D(R)}(K, \colim L^\bullet_n) \ar[lld] \\
\Hom_{D(R)}(K, \bigoplus L^\bullet_n[1]) \ar[r] & \ldots
}
$$
由于已假设 $\Hom_{D(R)}(K, \bigoplus L^\bullet_n)$ 等于
$\bigoplus \Hom_{D(R)}(K, L^\bullet_n)$，图表每一行的第一个映射都是单射
（这由该映射作为 $1 - t_n$ 所诱导映射之和的显式描述得到）。
因此 $\Hom_{D(R)}(K, \colim L^\bullet_n)$ 是上述图表中间一行
第一个映射的余核，这正是所需证明的。
\end{proof}
\noindent
下述命题将完美复形刻画为导出范畴的紧对象
（Derived Categories，定义 \ref{derived-definition-compact-object}），
它在多处出现。例如参见
\cite[命题 6.3 的证明]{Rickard}（这里处理有界情形）、
\cite[定理 2.4.3]{TT}（陈述并不完全一致），以及
\cite[命题 6.4]{Bokstedt-Neeman}（注意其极其糟糕的记号约定）。

\begin{proposition}
\label{more-algebra-proposition-perfect-is-compact}
设 $R$ 为一个环。对 $D(R)$ 中的对象 $K$，下列条件等价：
\begin{enumerate}
\item $K$ 是完美的；
\item $K$ 是 $D(R)$ 中的紧对象。
\end{enumerate}
\end{proposition}

\begin{proof}
假设 $K$ 是完美的，即 $K$ 拟同构于由有限投射模组成的有界复形
$P^\bullet$，参见定义 \ref{more-algebra-definition-perfect}。
若 $E_i$ 由复形 $E_i^\bullet$ 表示，则 $\bigoplus E_i$
由次数 $n$ 的项为 $\bigoplus E_i^n$ 的复形表示。
另一方面，对所有 $n$，由于 $P^n$ 是投射的，对每个 $R$-模复形
$K^\bullet$，有
$\Hom_{D(R)}(P^\bullet, K^\bullet) = \Hom_{K(R)}(P^\bullet, K^\bullet)$，
参见 Derived Categories，引理
\ref{derived-lemma-morphisms-from-projective-complex}。
因此，$\Hom_{D(R)}(P^\bullet, E^\bullet)$ 是下述复形的上同调：
$$
\prod \Hom_R(P^n, E^{n - 1}) \to
\prod \Hom_R(P^n, E^n) \to
\prod \Hom_R(P^n, E^{n + 1}).
$$
由于 $P^\bullet$ 有界，在上述复形中可以将 $\prod$ 号换成
$\bigoplus$ 号。由于每个 $P^n$ 都是有限 $R$-模，对所有
$n, m$ 有
$\Hom_R(P^n, \bigoplus_i E_i^m) = \bigoplus_i \Hom_R(P^n, E_i^m)$。
综合这些观察可知，Derived Categories，定义
\ref{derived-definition-compact-object} 中的映射是双射。

\medskip\noindent
反之，假设 $K$ 是紧的。用复形 $K^\bullet$ 表示 $K$，并考虑映射
$$
K^\bullet
\longrightarrow
\bigoplus\nolimits_{n \geq 0} \tau_{\geq n} K^\bullet
$$
其中使用了典范截断，参见 Homology，第
\ref{homology-section-truncations} 节。该映射有意义，因为在每个次数上，
右端的直和都是有限的。由假设，该映射经过一个有限直和分解。
因此，至少对某个 $n$，映射 $K \to \tau_{\geq n} K$ 为零；
换言之，$K$ 属于 $D^{-}(R)$。

\medskip\noindent
由于 $K \in D^{-}(R)$，且每个 $R$-模都是某个自由模的商，
可用自由 $R$-模组成的有界上复形 $K^\bullet$ 表示 $K$，参见
Derived Categories，引理
\ref{derived-lemma-subcategory-left-resolution}。注意
$$
K^\bullet = \bigcup\nolimits_{n \leq 0} \sigma_{\geq n}K^\bullet
$$
其中使用了硬截断，参见 Homology，第
\ref{homology-section-truncations} 节。因此由引理
\ref{more-algebra-lemma-commutes-with-countable-sums} 可知，$D(R)$ 中的
$1 : K^\bullet \to K^\bullet$ 经过
$\sigma_{\geq n}K^\bullet \to K^\bullet$ 分解。
于是对某个有界且各项为自由 $R$-模的复形 $L^\bullet$，
$D(R)$ 中的 $1 : K^\bullet \to K^\bullet$ 分解为
$$
K^\bullet \xrightarrow{\varphi} L^\bullet \xrightarrow{\psi} K^\bullet
$$
设当 $i \not \in [a, b]$ 时 $L^i = 0$。此后固定 $a, b$。
设 $c$ 是满足下列条件的最大整数 $\leq b + 1$：存在上述
$1_{K^\bullet}$ 的分解，且当 $i < c$ 时 $L^i$ 是有限自由的。
% P02-E0275
我们将归纳证明 $c = b + 1$。具体地，写成
$L^c = \bigoplus_{\lambda \in \Lambda} R$。
由于 $L^{c - 1}$ 是有限自由的，可以找到有限子集
$\Lambda' \subset \Lambda$，使得 $L^{c - 1} \to L^c$ 经过
$\bigoplus_{\lambda \in \Lambda'} R \subset L^c$ 分解。
考虑复形映射
$$
\pi :
L^\bullet
\longrightarrow
(\bigoplus\nolimits_{\lambda \in \Lambda \setminus \Lambda'} R)[-c]
$$
它在次数 $c$ 中由投影到对应于
$\Lambda \setminus \Lambda'$ 的因子给出。
由关于 $K$ 的假设，在必要时用更大的有限子集替换 $\Lambda'$ 后，
可以假设在 $D(R)$ 中 $\pi \circ \varphi = 0$。
设 $(L')^\bullet \subset L^\bullet$ 为 $\pi$ 的核。
由于 $\pi$ 满射，得到复形的短正合序列，进而得到 $D(R)$ 中的
区别三角（参见 Derived Categories，引理
\ref{derived-lemma-derived-canonical-delta-functor}）。
由于 $\Hom_{D(R)}(K, -)$ 是同调函子（参见
Derived Categories，引理
\ref{derived-lemma-representable-homological}），并且
$\pi \circ \varphi = 0$，可以找到 $D(R)$ 中的态射
$\varphi' : K^\bullet \to (L')^\bullet$，其与
$(L')^\bullet \to L^\bullet$ 的复合等于 $\varphi$。
令 $\psi'$ 等于 $\psi$ 与 $(L')^\bullet \to L^\bullet$ 的复合，
便得到新的分解。由于 $(L')^\bullet$ 除次数 $c$ 外与
$L^\bullet$ 一致，并且
$(L')^c = \bigoplus_{\lambda \in \Lambda'} R$，归纳步骤得证。

\medskip\noindent
上一段讨论的结论是，$1_K : K \to K$ 在 $D(R)$ 中分解为
$$
K \xrightarrow{\varphi} L \xrightarrow{\psi} K
$$
其中 $L$ 可由有限自由 $R$-模复形表示。特别地，$L$ 是完美的。
注意，$e = \varphi \circ \psi \in \text{End}_{D(R)}(L)$ 是幂等元。
由 Derived Categories，引理
\ref{derived-lemma-projectors-have-images-triangulated} 可知
$L = \Ker(e) \oplus \Ker(1 - e)$。
映射 $\varphi : K \to L$ 在 $D(R)$ 中诱导到
$\Ker(1 - e)$ 的同构。因此最后由引理
\ref{more-algebra-lemma-summands-perfect} 得到 $K$ 是完美的。
\end{proof}

\begin{lemma}
\label{more-algebra-lemma-perfect-modulo-nilpotent-ideal}
设 $R$ 为一个环，设 $I \subset R$ 为一个理想，设 $K$ 为
$D(R)$ 中的对象。假设
\begin{enumerate}
\item $K \otimes_R^\mathbf{L} R/I$ 在 $D(R/I)$ 中是完美的；
\item $I$ 是幂零理想。
\end{enumerate}
则 $K$ 在 $D(R)$ 中是完美的。
\end{lemma}

\begin{proof}
选取由有限投射 $R/I$-模组成的有限复形
$\overline{P}^\bullet$，表示
$K \otimes_R^\mathbf{L} R/I$，参见定义
\ref{more-algebra-definition-perfect}。由引理
\ref{more-algebra-lemma-lift-complex-projectives}，存在由投射 $R$-模组成的复形
$P^\bullet$，表示 $K$，使得
$\overline{P}^\bullet = P^\bullet/IP^\bullet$。
由 Nakayama 引理（Algebra，引理 \ref{algebra-lemma-NAK}）可知，
$P^\bullet$ 是由有限投射 $R$-模组成的有限复形。
\end{proof}

\begin{lemma}
\label{more-algebra-lemma-perfect-modulo-two-ideals}
设 $R$ 为一个环。设 $I, J \subset R$ 为理想。
设 $K$ 为 $D(R)$ 中的对象。假设
\begin{enumerate}
\item $K \otimes_R^\mathbf{L} R/I$ 在 $D(R/I)$ 中是完美的；
\item $K \otimes_R^\mathbf{L} R/J$ 在 $D(R/J)$ 中是完美的。
\end{enumerate}
则 $K \otimes_R^\mathbf{L} R/IJ$ 在 $D(R/IJ)$ 中是完美的。
\end{lemma}

\begin{proof}
% P02-E0276
显然，可以用 $R/IJ$ 替换 $R$，并用
$K \otimes_R^\mathbf{L} R/IJ$ 替换 $K$。此时
$R \to R/(I \cap J)$ 是核的平方为零的满射。因此由引理
\ref{more-algebra-lemma-perfect-modulo-nilpotent-ideal}，只需证明
$K \otimes_R^\mathbf{L} R/(I \cap J)$ 是完美的。
故可以假设 $I \cap J = 0$。

\medskip\noindent
我们在 $I \cap J = 0$ 的情形证明该引理。首先，可以用一个
K-平坦复形 $K^\bullet$ 表示 $K$，其中每个 $K^n$ 都是平坦的；
参见引理 \ref{more-algebra-lemma-K-flat-resolution}。于是有复形的短正合序列
$$
0 \to
K^\bullet \to K^\bullet/IK^\bullet \oplus K^\bullet/JK^\bullet \to
K^\bullet/(I + J)K^\bullet \to 0
$$
按导出张量积的构造，$K^\bullet/IK^\bullet$ 表示
$K \otimes^\mathbf{L}_R R/I$。对
$K^\bullet/JK^\bullet$ 和 $K^\bullet/(I + J)K^\bullet$ 亦然。
注意，$K^\bullet/(I + J)K^\bullet$ 是
$R/(I + J)$-模的完美复形；参见引理 \ref{more-algebra-lemma-pull-perfect}。
% P02-E0277
因此，复形 $K^\bullet/IK^\bullet$、
$K^\bullet/JK^\bullet$ 和 $K^\bullet/(I + J)K^\bullet$
只有有限多个非零上同调群（因为完美复形具有有限 Tor-振幅；
参见引理 \ref{more-algebra-lemma-perfect}）。
% P02-E0278
由与上面所示复形短正合序列相伴的长正合上同调序列，
可知 $K \in D^b(R)$。特别地，可以假设 $K^\bullet$
是自由 $R$-模组成的有界上复形（参见 Derived Categories，引理
\ref{derived-lemma-subcategory-left-resolution}）。

\medskip\noindent
现在用命题 \ref{more-algebra-proposition-perfect-is-compact} 的判据证明
$K$ 是完美的。
% P02-E0279
为此，令 $E_j \in D(R)$ 为由集合 $J$ 参数化的对象族。
选取逐项平坦且表示 $E_j$ 的复形 $E_j^\bullet$；例如参见引理
\ref{more-algebra-lemma-K-flat-resolution}。显然，
$$
0 \to
E_j^\bullet \to
E_j^\bullet/IE_j^\bullet \oplus E_j^\bullet/JE_j^\bullet \to
E_j^\bullet/(I + J)E_j^\bullet \to 0
$$
是复形的短正合序列。取直和，得到类似的短正合序列
$$
0 \to
\bigoplus E_j^\bullet \to
\bigoplus E_j^\bullet/IE_j^\bullet \oplus E_j^\bullet/JE_j^\bullet \to
\bigoplus E_j^\bullet/(I + J)E_j^\bullet \to 0
$$
（注意，$- \otimes_R R/I$ 与直和交换。）
该短正合序列在 $D(R)$ 中确定一个区别三角；参见
Derived Categories，引理
\ref{derived-lemma-derived-canonical-delta-functor}。
应用同调函子 $\Hom_{D(R)}(K, -)$（参见 Derived Categories，引理
\ref{derived-lemma-representable-homological}），得到交换图
% P02-E0280
$$
\xymatrix{
\bigoplus \Hom_{D(R)}(K^\bullet, E_j^\bullet/(I + J))[-1] \ar[r] \ar[d] &
\Hom_{D(R)}(K^\bullet, \bigoplus E_j^\bullet/(I + J))[-1] \ar[d] \\
\bigoplus \Hom_{D(R)}(K^\bullet, E_j^\bullet/I \oplus E_j^\bullet/J)[-1]
\ar[r] \ar[d] &
\Hom_{D(R)}(K^\bullet, \bigoplus E_j^\bullet/I \oplus E_j^\bullet/J)[-1]
\ar[d] \\
\bigoplus \Hom_{D(R)}(K^\bullet, E_j^\bullet) \ar[r] \ar[d] &
\Hom_{D(R)}(K^\bullet, \bigoplus E_j^\bullet) \ar[d] \\
\bigoplus \Hom_{D(R)}(K^\bullet, E_j^\bullet/I \oplus E_j^\bullet/J)
\ar[r] \ar[d] &
\Hom_{D(R)}(K^\bullet, \bigoplus E_j^\bullet/I \oplus E_j^\bullet/J)
\ar[d] \\
\bigoplus \Hom_{D(R)}(K^\bullet, E_j^\bullet/(I + J)) \ar[r] &
\Hom_{D(R)}(K^\bullet, \bigoplus E_j^\bullet/(I + J))
}
$$
其各列正合。显然，对任意 $R$-模复形 $E^\bullet$，有
\begin{align*}
\Hom_{D(R)}(K^\bullet, E^\bullet/I) & =
\Hom_{K(R)}(K^\bullet, E^\bullet/I) \\
& =
\Hom_{K(R/I)}(K^\bullet/IK^\bullet, E^\bullet/I) \\
& =
\Hom_{D(R/I)}(K^\bullet/IK^\bullet, E^\bullet/I)
\end{align*}
除以 $J$ 或 $I + J$ 时也同样成立；参见 Derived Categories，引理
\ref{derived-lemma-morphisms-from-projective-complex}。
% P02-E0281
因此，除中间的水平箭头可能例外以外，其余水平箭头都是同构，
因为复形 $K^\bullet/IK^\bullet$、$K^\bullet/JK^\bullet$、
$K^\bullet/(I + J)K^\bullet$ 分别是 $R/I$、$R/J$、
$R/(I + J)$-模的完美复形；参见命题
\ref{more-algebra-proposition-perfect-is-compact}。
由 $5$-引理（Homology，引理 \ref{homology-lemma-five-lemma}），
中间的映射是同构，再由命题 \ref{more-algebra-proposition-perfect-is-compact}
即得结论。
\end{proof}






\section{强生成元与正则环}
\label{more-algebra-section-strong-generators-regular}

\noindent
设 $R$ 为一个环。
% P02-E0282
用 $D(R)_c$ 表示 $D(R)$ 的饱和满三角子范畴。已经知道
$$
\langle R \rangle = D_{perf}(R) = D(R)_c
$$
参见引理 \ref{more-algebra-lemma-perfect-ring-classical-generator} 和命题
\ref{more-algebra-proposition-perfect-is-compact}。事实证明，若 $R$ 是正则的，
则 $R$ 是强生成元（Derived Categories，定义
\ref{derived-definition-generators}）。

\begin{lemma}
\label{more-algebra-lemma-ghost-lemma}
\begin{reference}
\cite{Kelly}
\end{reference}
设 $R$ 为一个环，设 $n \geq 1$，并设
$K \in \langle R \rangle_n$；记号见 Derived Categories，第
\ref{derived-section-generators} 节。考虑 $D(R)$ 中的映射
$$
K \xrightarrow{f_1} K_1 \xrightarrow{f_2} K_2
\xrightarrow{f_3} \ldots \xrightarrow{f_n} K_n
$$
若对所有 $i, j$ 都有 $H^i(f_j) = 0$，则
$f_n \circ \ldots \circ f_1 = 0$。
\end{lemma}

\begin{proof}
若 $n = 1$，则 $K$ 是 $D(R)$ 中某个有界复形
$P^\bullet$ 的直和项，该复形的各项是有限自由 $R$-模，且微分均为零。
因此，只需证明 $D(R)$ 中满足对所有 $i$ 都有 $H^i(f) = 0$ 的任意态射
$f : P^\bullet \to K_1$ 为零。由于 $P^\bullet$ 是有限直和
$P^\bullet = \bigoplus R[m_j]$，只需证明 $D(R)$ 中满足
$H^{-m}(g) = 0$ 的任意态射 $g : R[m] \to K_1$ 为零。
这由等式 $\Hom_{D(R)}(R[-m], K) = H^m(K)$ 得到。

\medskip\noindent
当 $n > 1$ 时，对 $n$ 作归纳。我们知道，$K$ 是 $D(R)$ 中某个对象
$P$ 的直和项，而该对象位于区别三角
$$
P' \xrightarrow{i} P \xrightarrow{p} P'' \to P'[1]
$$
中，其中 $P' \in \langle R \rangle_1$ 且
$P'' \in \langle R \rangle_{n - 1}$。与上面一样，可以用 $P$ 替换
$K$，并假设在 $D(R)$ 中有
$$
P \xrightarrow{f_1} K_1 \xrightarrow{f_2} K_2
\xrightarrow{f_3} \ldots \xrightarrow{f_n} K_n
$$
且 $f_j$ 在上同调上为零。由 $n = 1$ 的情形，复合
$f_1 \circ i$ 为零。因此由 Derived Categories，引理
\ref{derived-lemma-representable-homological}，可以找到态射
$h : P'' \to K_1$，使得 $f_1 = h \circ p$。注意
$f_2 \circ h$ 在上同调上为零。于是由归纳假设，
$f_n \circ \ldots \circ f_2 \circ h = 0$，从而
$f_n \circ \ldots \circ f_1 =
f_n \circ \ldots \circ f_2 \circ h \circ p = 0$，正合所需。
\end{proof}

\begin{lemma}
\label{more-algebra-lemma-not-regular-not-strong}
设 $R$ 为 Noether 环。若 $R$ 是 $D_{perf}(R)$ 的强生成元，
则 $R$ 是有限维正则环。
\end{lemma}

\begin{proof}
假设对某个 $n \geq 1$ 有
$D_{perf}(R) = \langle R \rangle_n$。对任意有限 $R$-模 $M$，
可以选取有限自由 $R$-模复形
% P02-E0283
$$
P = (
P^{-n - 1} \xrightarrow{d^{-n - 1}}
P^{-n} \xrightarrow{d^{-n}}
P^{-n + 1} \xrightarrow{d^1}
\ldots \xrightarrow{d^{-1}} P^0)
$$
使得当 $i = -n, \ldots, - 1$ 时有 $H^i(P) = 0$，并且
$M \cong \Coker(d^{-1})$。注意 $P$ 属于 $D_{perf}(R)$。
% P02-E0284
对任意 $R$-模 $N$，可以通过 $M$ 的有限自由解消 $P$ 计算
$\Ext^n_R(M, N)$；参见 Algebra，第 \ref{algebra-section-ext} 节，
并与 Derived Categories，第 \ref{derived-section-ext} 节比较。
特别地，上述序列定义元素
$$
\xi \in \Ext^n_R(\Coker(d^{-1}), \Coker(d^{-n - 1})) =
\Ext^n_R(M, \Coker(d^{-n - 1}))
$$
并且对 $\Ext^n_R(M, N)$ 中任意元素 $\overline{\xi}$，存在
$R$-模映射 $\varphi : \Coker(d^{-n - 1}) \to N$，使得
$\varphi$ 把 $\xi$ 映到 $\overline{\xi}$。
对 $j = 1, \ldots, n - 1$，考虑复形
$$
K_j = (\Coker(d^{-n - 1}) \to P^{-n + 1} \to \ldots \to P^{-j})
$$
其中 $\Coker(d^{-n - 1})$ 位于次数 $-n$，而 $P^t$ 位于次数 $t$。
还令 $K_n = \Coker(d^{-n - 1})[n]$。于是有映射
$$
P \to K_1 \to K_2 \to \ldots \to K_n
$$
它们在上同调上诱导零映射。由于
$P \in D_{perf}(R) = \langle R \rangle_n$，引理
\ref{more-algebra-lemma-ghost-lemma} 表明这些映射的复合在 $D(R)$ 中为零。
% P02-E0285
又由 Derived Categories，引理
\ref{derived-lemma-morphisms-from-projective-complex}，有
$\Hom_{D(R)}(P, K_n) = \Hom_{K(R)}(P, K_n)$，故 $\xi = 0$。
因此对所有 $R$-模 $N$ 都有 $\Ext^n_R(M, N) = 0$；参见上面的讨论。
由 Algebra，引理 \ref{algebra-lemma-projective-dimension-ext}，
$M$ 的投射维数 $\leq n - 1$。这对所有有限 $R$-模 $M$ 都成立，
故由 Algebra，引理 \ref{algebra-lemma-finite-gl-dim}，$R$ 具有有限整体维数。
最后由 Algebra，引理
\ref{algebra-lemma-finite-gl-dim-finite-dim-regular} 得到结论。
\end{proof}

\begin{lemma}
\label{more-algebra-lemma-ext-regular}
设 $R$ 为维数 $d < \infty$ 的 Noether 正则环，设
$K, L \in D^-(R)$。
% P02-E0286
假设存在一个 $k$，使得当 $i \leq k$ 时 $H^i(K) = 0$，
且当 $i \geq k - d + 1$ 时 $H^i(L) = 0$。则
$\Hom_{D(R)}(K, L) = 0$。
\end{lemma}

\begin{proof}
设 $K^\bullet$ 为表示 $K$ 的有界上复形，例如当
$i \geq n + 1$ 时 $K^i = 0$。用
$\tau_{\geq k + 1}K^\bullet$ 替换 $K^\bullet$ 后，可以假设当
$i \leq k$ 时 $K^i = 0$。于是利用区别三角
$$
K^n[-n] \to K^\bullet \to \sigma_{\leq n - 1}K^\bullet
$$
可知，只需对 $K^n[-n]$ 和 $\sigma_{\leq n - 1}K^\bullet$
证明该引理。对 $n$ 作归纳，可知只需处理下述情形：$K$ 由复形
$M[-m]$ 表示，其中 $M$ 为某个 $R$-模，且 $m \geq k + 1$。
由 Algebra，引理
\ref{algebra-lemma-finite-gl-dim-finite-dim-regular}，$R$ 的整体维数为
$d$，故 $M$ 有投射解消
$0 \to P_d \to \ldots \to P_0 \to M \to 0$。
令复形 $P^\bullet$ 在次数 $m - i$ 的项为 $P_i$，则它是表示
$M[-m]$ 的有界投射复形。另一方面，可以选取表示 $L$ 的复形
$L^\bullet$，使得当 $i \geq k - d  + 1$ 时 $L^i = 0$。
因此任意复形映射 $P^\bullet \to L^\bullet$ 都为零。
由 Derived Categories，引理
\ref{derived-lemma-morphisms-from-projective-complex} 即得结论。
\end{proof}

\begin{lemma}
\label{more-algebra-lemma-split-complex-regular}
设 $R$ 为维数 $1 \leq d < \infty$ 的 Noether 正则环。
设 $K \in D(R)$ 是完美的，并设 $k \in \mathbf{Z}$，使得当
$i = k - d + 2, \ldots, k$ 时有 $H^i(K) = 0$
（当 $d = 1$ 时此条件为空）。则
$K = \tau_{\leq k - d + 1}K \oplus \tau_{\geq k + 1}K$。
\end{lemma}

\begin{proof}
上同调消失表明存在区别三角
$$
\tau_{\leq k - d + 1}K \to K \to \tau_{\geq k + 1}K \to
(\tau_{\leq k - d + 1}K)[1]
$$
由 Derived Categories，引理 \ref{derived-lemma-split}，
只需证明第三个箭头为零。因此，只需证明
$\Hom_{D(R)}(\tau_{\geq k + 1}K, (\tau_{\leq k - d + 1}K)[1]) = 0$；
这由引理 \ref{more-algebra-lemma-ext-regular} 得到。
\end{proof}

\begin{lemma}
\label{more-algebra-lemma-regular-strong-generator}
设 $R$ 为有限维 Noether 正则环。则 $R$ 是完美对象所成满子范畴
$D_{perf}(R) \subset D(R)$ 的强生成元。
\end{lemma}

\begin{proof}
我们将用到：$D(R)$ 的对象 $K$ 是完美的，当且仅当 $K$ 有界且
各上同调模均为有限模；参见引理 \ref{more-algebra-lemma-regular-perfect}。
三角范畴的强生成元定义于 Derived Categories，定义
\ref{derived-definition-generators}。令 $d = \dim(R)$。

\medskip\noindent
设 $K \in D_{perf}(R)$。我们将证明
$K \in \langle R \rangle_{d + 1}$。由 Algebra，引理
\ref{algebra-lemma-finite-gl-dim-finite-dim-regular}，每个有限
$R$-模的投射维数都 $\leq d$。我们将对 $0 \leq i \leq d$ 归纳证明：
若对所有 $n \in \mathbf{Z}$，$H^n(K)$ 的投射维数都
$\leq i$，则 $K$ 属于 $\langle R \rangle_{i + 1}$。

\medskip\noindent
基础情形 $i = 0$。此时 $H^n(K)$ 是投射维数为 $0$ 的有限
$R$-模。换言之，每个上同调模都是投射 $R$-模。因此，对所有
$i > 0$ 及 $m, n \in \mathbf{Z}$，都有
$\Ext^i_R(H^n(K), H^m(K)) = 0$。由 Derived Categories，引理
\ref{derived-lemma-ext-2-zero-pre}，可知 $K$ 同构于其各上同调模之平移的直和。
由于每个上同调模都是有限投射 $R$-模，它是若干个 $R$ 副本之直和的直和项。
因此根据定义，$K$ 包含在 $\langle R \rangle_1$ 中。

\medskip\noindent
归纳步骤。假设断言对 $i < d$ 成立，并设
$K \in D_{perf}(R)$ 满足：对所有 $n \in \mathbf{Z}$，
$H^n(K)$ 的投射维数都 $\leq i + 1$。选取 $a \leq b$，使得当
$n \not \in [a, b]$ 时 $H^n(K)$ 为零。对每个 $n \in [a, b]$，
选取满射 $F^n \to H^n(K)$，其中 $F^n$ 是有限自由 $R$-模。
由于 $F^n$ 是投射的，可以在 $D(R)$ 中把
$F^n \to H^n(K)$ 提升为映射 $F^n[-n] \to K$
（略去一个小细节）。于是得到态射
$\bigoplus_{a \leq n \leq b} F^n[-n] \to K$，
它在上同调模上是满射。选取 $D(R)$ 中的区别三角
$$
K' \to \bigoplus\nolimits_{a \leq n \leq b} F^n[-n] \to K \to K'[1]
$$
当然，对象 $K'$ 有界，且各上同调模均为有限模。根据所作选择，
上同调长正合序列分解成短正合序列
$$
0 \to H^n(K') \to F^n \to H^n(K) \to 0
$$
由 Algebra，引理
\ref{algebra-lemma-exact-sequence-projective-dimension}，
$H^n(K')$ 的投射维数 $\leq \max(0, i)$。因此
$K' \in \langle R \rangle_{i + 1}$。根据定义，这意味着
$K$ 属于 $\langle R \rangle_{i + 1 + 1}$，正合所需。
\end{proof}

\begin{proposition}
\label{more-algebra-proposition-regular-strong-generator}
设 $R$ 为 Noether 环。下列条件等价：
\begin{enumerate}
\item $R$ 是有限维正则环；
\item $D_{perf}(R)$ 有强生成元；
\item $R$ 是 $D_{perf}(R)$ 的强生成元。
\end{enumerate}
\end{proposition}

\begin{proof}
这是引理 \ref{more-algebra-lemma-perfect-ring-classical-generator}、
\ref{more-algebra-lemma-not-regular-not-strong}、
\ref{more-algebra-lemma-regular-strong-generator} 以及 Derived Categories，引理
\ref{derived-lemma-classical-generator-strong-generator} 的形式推论。
\end{proof}







\section{相对有限表示模}
\label{more-algebra-section-relative-finite-presentation}

\noindent
设 $R$ 为一个环，设 $A \to B$ 为有限型 $R$-代数之间的有限同态，
并设 $M$ 为有限 $B$-模。在这种情形，下面的等价关系{\bf 并不成立}：
$$
M\text{ of finite presentation over }B
\Leftrightarrow
M\text{ of finite presentation over }A
$$
% P02-E0287
一个反例是 $R = k[x_1, x_2, x_3, \ldots]$、$A = R$、
$B = R/(x_i)$ 以及 $M = B$。为了“修正”这一点，
我们引入有限表示的相对概念。

\begin{lemma}
\label{more-algebra-lemma-relatively-finitely-presented}
设 $R \to A$ 为有限型环同态，并设 $M$ 为 $A$-模。
下列条件等价：
\begin{enumerate}
\item 对某个表示 $\alpha : R[x_1, \ldots, x_n] \to A$，
模 $M$ 是有限表示 $R[x_1, \ldots, x_n]$-模；
\item 对所有表示 $\alpha : R[x_1, \ldots, x_n] \to A$，
模 $M$ 是有限表示 $R[x_1, \ldots, x_n]$-模；
\item 对任意满射 $A' \to A$，其中 $A'$ 是有限表示 $R$-代数，
模 $M$ 作为 $A'$-模是有限表示的。
\end{enumerate}
在这种情形，$M$ 是有限表示 $A$-模。
\end{lemma}

\begin{proof}
% P02-E0288
设 $\alpha : R[x_1, \ldots, x_n] \to A$ 与
$\beta : R[y_1, \ldots, y_m] \to A$ 为两个表示。
选取 $f_j \in R[x_1, \ldots, x_n]$，使得
$\alpha(f_j) = \beta(y_j)$；并选取
$g_i \in R[y_1, \ldots, y_m]$，使得
$\beta(g_i) = \alpha(x_i)$。于是得到交换图
$$
\xymatrix{
R[x_1, \ldots, x_n, y_1, \ldots, y_m]
\ar[d]^{x_i \mapsto g_i} \ar[rr]_-{y_j \mapsto f_j} & &
R[x_1, \ldots, x_n] \ar[d] \\
R[y_1, \ldots, y_m] \ar[rr] & & A
}
$$
对 (1) 与 (2) 的等价性应用 Algebra，引理
\ref{algebra-lemma-finitely-presented-over-subring} 和
\ref{algebra-lemma-finite-finitely-presented-extension} 即可。
选取表示 $A' = R[x_1, \ldots, x_n]/(f_1, \ldots, f_m)$，
并应用 Algebra，引理
\ref{algebra-lemma-finite-finitely-presented-extension}，可知
$M$ 作为 $A'$-模有限表示，当且仅当 $M$ 作为
$R[x_1, \ldots, x_n]$-模有限表示；这证明了 (2) 与 (3) 等价。
\end{proof}

\begin{definition}
\label{more-algebra-definition-relatively-finitely-presented}
设 $R \to A$ 为有限型环同态，设 $M$ 为 $A$-模。
若引理 \ref{more-algebra-lemma-relatively-finitely-presented} 中的等价条件成立，
则称 $M$ 是{it 相对于 $R$ 有限表示的 $A$-模}。
\end{definition}

\noindent
注意，若 $R \to A$ 是有限表示的，则 $M$ 是相对于 $R$ 有限表示的
$A$-模，当且仅当 $M$ 是有限表示 $A$-模。同样显然，$A$ 作为
$A$-模相对于 $R$ 有限表示，当且仅当 $A$ 在 $R$ 上有限表示。
若 $R$ 是 Noether 的，该概念没有新内容。现在可以表述我们所寻求的结果。

\begin{lemma}
\label{more-algebra-lemma-finite-extension}
设 $R$ 为一个环，设 $A \to B$ 为有限型 $R$-代数之间的有限同态，
并设 $M$ 为 $B$-模。则 $M$ 作为 $A$-模相对于 $R$ 有限表示，
当且仅当 $M$ 作为 $B$-模相对于 $R$ 有限表示。
\end{lemma}

\begin{proof}
选取满射 $R[x_1, \ldots, x_n] \to A$。选取
$y_1, \ldots, y_m \in B$，使它们在 $A$ 上生成 $B$。
由于 $A \to B$ 是有限的，每个 $y_i$ 都满足一个系数在 $A$ 中的首一方程。
因此，可以找到首一多项式
$P_j(T) \in R[x_1, \ldots, x_n][T]$，使得在 $B$ 中
$P_j(y_j) = 0$。于是得到交换图
$$
\xymatrix{
R[x_1, \ldots, x_n] \ar[d] \ar[r] &
R[x_1, \ldots, x_n, y_1, \ldots, y_m]/(P_j(y_j)) \ar[d] \\
A \ar[r] & B
}
$$
上方箭头是有限且有限表示的环同态，故由 Algebra，引理
\ref{algebra-lemma-finite-finitely-presented-extension} 和定义即得结论。
\end{proof}

\noindent
由这个结果可知，相对概念是合理的，并且对于 $R$ 上有限型环之间的有限同态
表现良好。它还在局部化和基变换下稳定，并且能够良好地“粘合”。

\begin{lemma}
\label{more-algebra-lemma-localize-relative-finite-presentation}
设 $R$ 为一个环，$f \in R$ 为元素，$R_f \to A$ 为有限型环同态，
$g \in A$，并设 $M$ 为 $A$-模。
% P02-E0289
若 $M$ 相对于 $R_f$ 有限表示，则 $M_g$ 是相对于 $R$ 有限表示的
$A_g$-模。
\end{lemma}

\begin{proof}
选取表示 $R_f[x_1, \ldots, x_n] \to A$。写成
$R_f = R[x_0]/(fx_0 - 1)$。考虑表示
$R[x_0, x_1, \ldots, x_n, x_{n + 1}] \to A_g$；
它延拓给定映射，把 $x_0$ 映到 $1/f$ 的像，并把
$x_{n + 1}$ 映到 $1/g$。选取映到 $g$ 的
$g' \in R[x_0, x_1, \ldots, x_n]$（这是可能的）。假设
$$
R_f[x_1, \ldots, x_n]^{\oplus s} \to
R_f[x_1, \ldots, x_n]^{\oplus t} \to M \to 0
$$
是由矩阵 $(h_{ij})$ 给出的 $M$ 的表示。选取映到 $h_{ij}$ 的
$h'_{ij} \in R[x_0, x_1, \ldots, x_n]$。则
$$
R[x_0, x_1, \ldots, x_n, x_{n + 1}]^{\oplus s + 2t} \to
R[x_0, x_1, \ldots, x_n, x_{n + 1}]^{\oplus t} \to M_g \to 0
$$
% P02-E0290
是 $M_f$ 的一个表示。这里，定义该映射的
$t \times (s + 2t)$ 矩阵的第一个 $t \times s$ 分块为矩阵
$h'_{ij}$；
% P02-E0291
第二个 $t \times t$ 分块为 $(x_0f - )I_t$；第三个分块为
$(x_{n + 1}g' - 1)I_t$。
\end{proof}

\begin{lemma}
\label{more-algebra-lemma-base-change-relative-finite-presentation}
设 $R \to A$ 为有限型环同态，并设 $M$ 是相对于 $R$ 有限表示的
$A$-模。对任意环同态 $R \to R'$，$A \otimes_R R'$-模
% P02-E0292
$$
M \otimes_A A' = M \otimes_R R'
$$
相对于 $R'$ 有限表示。
\end{lemma}

\begin{proof}
选取满射 $R[x_1, \ldots, x_n] \to A$，并选取表示
$$
R[x_1, \ldots, x_n]^{\oplus s} \to
R[x_1, \ldots, x_n]^{\oplus t} \to M \to 0
$$
则
$$
R'[x_1, \ldots, x_n]^{\oplus s} \to
R'[x_1, \ldots, x_n]^{\oplus t} \to M \otimes_R R' \to 0
$$
是基变换的一个表示，结论得证。
\end{proof}

\begin{lemma}
\label{more-algebra-lemma-pull-relative-finite-presentation}
设 $R \to A$ 为有限型环同态，设 $M$ 是相对于 $R$ 有限表示的
$A$-模，并设 $A \to A'$ 为有限表示环同态。则 $A'$-模
$M \otimes_A A'$ 相对于 $R$ 有限表示。
\end{lemma}

\begin{proof}
选取满射 $R[x_1, \ldots, x_n] \to A$，并选取表示
$A' = A[y_1, \ldots, y_m]/(g_1, \ldots, g_l)$。
选取映到 $g_i$ 的
$g'_i \in R[x_1, \ldots, x_n, y_1, \ldots, y_m]$。设
$$
R[x_1, \ldots, x_n]^{\oplus s} \to
R[x_1, \ldots, x_n]^{\oplus t} \to M \to 0
$$
是由矩阵 $(h_{ij})$ 给出的 $M$ 的表示。则
% P02-E0293
$$
R[x_1, \ldots, x_n, y_1, \ldots, y_m]^{\oplus s + tl} \to
R[x_0, x_1, \ldots, x_n, y_1, \ldots, y_m]^{\oplus t} \to M \otimes_A A' \to 0
$$
是 $M \otimes_A A'$ 的一个表示。这里，定义该映射的
$t \times (s + lt)$ 矩阵的第一个 $t \times s$ 分块为矩阵
$h_{ij}$，随后是 $l$ 个大小为 $t \times t$ 的分块
$g'_iI_t$。
\end{proof}

\begin{lemma}
\label{more-algebra-lemma-composition-relative-finite-presentation}
设 $R \to A \to B$ 为有限型环同态，设 $M$ 为 $B$-模。
若 $M$ 相对于 $A$ 有限表示，且 $A$ 在 $R$ 上有限表示，
则 $M$ 相对于 $R$ 有限表示。
\end{lemma}

\begin{proof}
选取满射 $A[x_1, \ldots, x_n] \to B$，并选取由矩阵
$(h_{ij})$ 给出的表示
$$
A[x_1, \ldots, x_n]^{\oplus s} \to
A[x_1, \ldots, x_n]^{\oplus t} \to M \to 0
$$
再选取表示
$$
A = R[y_1, \ldots, y_m]/(g_1, \ldots, g_u).
$$
选取映到 $h_{ij}$ 的
$h'_{ij} \in R[y_1, \ldots, y_m, x_1, \ldots, x_n]$。
于是得到表示
$$
R[y_1, \ldots, y_m, x_1, \ldots, x_n]^{\oplus s + tu} \to
R[y_1, \ldots, y_m, x_1, \ldots, x_n]^{\oplus t} \to M \to 0
$$
其中 $t \times (s + tu)$ 矩阵的第一个 $t \times s$ 分块由
$h'_{ij}$ 构成，随后是 $u$ 个大小为 $t \times t$ 的分块
$g_iI_t$，$i = 1, \ldots, u$。
\end{proof}

\begin{lemma}
\label{more-algebra-lemma-glue-relative-finite-presentation}
设 $R \to A$ 为有限型环同态，设 $M$ 为 $A$-模，并设
$f_1, \ldots, f_r \in A$ 生成单位理想。下列条件等价：
\begin{enumerate}
\item 每个 $M_{f_i}$ 都相对于 $R$ 有限表示；
\item $M$ 相对于 $R$ 有限表示。
\end{enumerate}
\end{lemma}

\begin{proof}
(2) $\Rightarrow$ (1) 是引理
\ref{more-algebra-lemma-localize-relative-finite-presentation} 的内容。
假设 (1)。在 $A$ 中写成 $1 = \sum f_ig_i$。
% P02-E0294
选取满射
$R[x_1, \ldots, x_n, y_1, \ldots, y_r, z_1, \ldots, z_r] \to A$，
使得 $y_i$ 映到 $f_i$，而 $z_i$ 映到 $g_i$。于是可见存在满射
$$
P = R[x_1, \ldots, x_n, y_1, \ldots, y_r, z_1, \ldots, z_r]/(\sum y_iz_i - 1)
\longrightarrow
A.
$$
由引理 \ref{more-algebra-lemma-relatively-finitely-presented}，$M_{f_i}$ 是有限表示
$A_{f_i}$-模；故由 Algebra，引理 \ref{algebra-lemma-cover}，
$M$ 是有限表示 $A$-模。因此，$M$ 是有限 $P$-模（$P$ 如上）。
选取满射 $P^{\oplus t} \to M$。必须证明该映射的核 $K$ 是有限
$P$-模。由于 $P_{y_i}$ 满射到 $A_{f_i}$，由引理
\ref{more-algebra-lemma-relatively-finitely-presented} 和 Algebra，引理
\ref{algebra-lemma-extension} 可知，局部化 $K_{y_i}$ 是有限生成
$P_{y_i}$-模。选取元素
$k_{i, j} \in K$，$i = 1, \ldots, r$，$j = 1, \ldots, s_i$，
使得 $k_{i, j}$ 在 $K_{y_i}$ 中的像生成该模。令
$K' \subset K$ 为这些元素 $k_{i, j}$ 生成的 $P$-模。
则对所有 $i$，模 $K/K'$ 在 $y_i$ 处的局部化都为零。由于
$(y_1, \ldots, y_r) = P$，可知 $K/K' = 0$，正合所需。
\end{proof}

\begin{lemma}
\label{more-algebra-lemma-ses-relatively-finite-presentation}
设 $R \to A$ 为有限型环同态，并设
$0 \to M' \to M \to M'' \to 0$ 为 $A$-模的短正合序列。
\begin{enumerate}
\item 若 $M', M''$ 相对于 $R$ 有限表示，则 $M$ 也如此。
\item 若 $M'$ 是有限型 $A$-模，且 $M$ 相对于 $R$ 有限表示，
则 $M''$ 相对于 $R$ 有限表示。
\end{enumerate}
\end{lemma}

\begin{proof}
% P02-E0295
这立即由 Algebra，引理 \ref{algebra-lemma-extension} 得到。
\end{proof}

\begin{lemma}
\label{more-algebra-lemma-sum-relatively-finite-presentation}
设 $R \to A$ 为有限型环同态，设 $M, M'$ 为 $A$-模。
若 $M \oplus M'$ 相对于 $R$ 有限表示，则 $M$ 与 $M'$ 也如此。
\end{lemma}

\begin{proof}
略。
\end{proof}






\section{相对伪凝聚模}
\label{more-algebra-section-relative-pseudo-coherent}

\noindent
本节给出第 \ref{more-algebra-section-relative-finite-presentation} 节关于伪凝聚性的类似理论。

\begin{lemma}
\label{more-algebra-lemma-pull-push}
设 $R$ 为一个环，设 $K^\bullet$ 为 $R$-模复形。
考虑把 $x$ 映到零的 $R$-代数同态 $R[x] \to R$。则在 $D(R)$ 中
$$
K^\bullet \otimes_{R[x]}^{\mathbf{L}} R \cong K^\bullet \oplus K^\bullet[1]
$$
\end{lemma}

\begin{proof}
在 $R$ 上选取 K-平坦解消 $P^\bullet \to K^\bullet$，使得对所有
$n$，$P^n$ 都是平坦 $R$-模；参见引理
\ref{more-algebra-lemma-K-flat-resolution}。于是 $P^\bullet \otimes_R R[x]$ 是
$R[x]$-模的 K-平坦复形，且各项都是平坦 $R[x]$-模；参见引理
\ref{more-algebra-lemma-base-change-K-flat} 和 Algebra，引理
\ref{algebra-lemma-flat-base-change}。特别地，
$x : P^n \otimes_R R[x] \to P^n \otimes_R R[x]$ 是单射，
其余核同构于 $P^n$。因此
$$
P^\bullet \otimes_R R[x] \xrightarrow{x} P^\bullet \otimes_R R[x]
$$
是 $R[x]$-模的双复形，其相伴全复形拟同构于 $P^\bullet$，
从而拟同构于 $K^\bullet$。此外，例如由引理
\ref{more-algebra-lemma-tensor-product-K-flat} 或引理
\ref{more-algebra-lemma-K-flat-two-out-of-three}，该相伴全复形是
$R[x]$-模的 K-平坦复形。因此
\begin{align*}
K^\bullet \otimes_{R[x]}^{\mathbf{L}} R
& \cong
\text{Tot}(P^\bullet \otimes_R R[x] \xrightarrow{x} P^\bullet \otimes_R R[x])
\otimes_{R[x]} R =
\text{Tot}(P^\bullet \xrightarrow{0} P^\bullet) \\
& = P^\bullet \oplus P^\bullet[1] \cong K^\bullet \oplus K^\bullet[1]
\end{align*}
正合所需。
\end{proof}

\begin{lemma}
\label{more-algebra-lemma-add-variable-pseudo-coherent}
设 $R$ 为一个环，$K^\bullet$ 为 $R$-模复形，并设
$m \in \mathbf{Z}$。考虑把 $x$ 映到零的 $R$-代数同态
$R[x] \to R$。则 $K^\bullet$ 作为 $R$-模复形是
$m$-伪凝聚的，当且仅当 $K^\bullet$ 作为 $R[x]$-模复形是
$m$-伪凝聚的。
\end{lemma}

\begin{proof}
这是引理 \ref{more-algebra-lemma-finite-push-pseudo-coherent} 的特例。
下面也给出另一种证明。

\medskip\noindent
注意 $0 \to R[x] \to R[x] \to R \to 0$ 正合。
因此 $R$ 作为 $R[x]$-模是伪凝聚的。故该引理的一个方向由引理
\ref{more-algebra-lemma-finite-push-pseudo-coherent} 得到。
为证明另一方向，假设 $K^\bullet$ 作为 $R[x]$-模复形是
$m$-伪凝聚的。由引理 \ref{more-algebra-lemma-pull-pseudo-coherent}，
$K^\bullet \otimes^{\mathbf{L}}_{R[x]} R$ 作为 $R$-模复形是
$m$-伪凝聚的。由引理 \ref{more-algebra-lemma-pull-push}，
$K^\bullet \oplus K^\bullet[1]$ 作为 $R$-模复形是 $m$-伪凝聚的。
最后，由引理 \ref{more-algebra-lemma-summands-pseudo-coherent} 可知，
$K^\bullet$ 作为 $R$-模复形是 $m$-伪凝聚的。
\end{proof}

\begin{lemma}
\label{more-algebra-lemma-relatively-pseudo-coherent}
设 $R \to A$ 为有限型环同态，设 $K^\bullet$ 为 $A$-模复形，
并设 $m \in \mathbf{Z}$。下列条件等价：
\begin{enumerate}
\item 对某个表示 $\alpha : R[x_1, \ldots, x_n] \to A$，
复形 $K^\bullet$ 是 $R[x_1, \ldots, x_n]$-模的
$m$-伪凝聚复形；
\item 对所有表示 $\alpha : R[x_1, \ldots, x_n] \to A$，
复形 $K^\bullet$ 是 $R[x_1, \ldots, x_n]$-模的
$m$-伪凝聚复形。
\end{enumerate}
特别地，对伪凝聚性也有同样的等价性。
\end{lemma}

\begin{proof}
% P02-E0296 P02-E0297
设 $\alpha : R[x_1, \ldots, x_n] \to A$ 与
$\beta : R[y_1, \ldots, y_m] \to A$ 为两个表示。
选取 $f_j \in R[x_1, \ldots, x_n]$，使得
$\alpha(f_j) = \beta(y_j)$；并选取
$g_i \in R[y_1, \ldots, y_m]$，使得
$\beta(g_i) = \alpha(x_i)$。于是得到交换图
$$
\xymatrix{
R[x_1, \ldots, x_n, y_1, \ldots, y_m]
\ar[d]^{x_i \mapsto g_i} \ar[rr]_-{y_j \mapsto f_j} & &
R[x_1, \ldots, x_n] \ar[d] \\
R[y_1, \ldots, y_m] \ar[rr] & & A
}
$$
作坐标变换后，环同态
$R[x_1, \ldots, x_n, y_1, \ldots, y_m] \to R[x_1, \ldots, x_n]$
同构于把每个 $y_i$ 映到零的环同态。图中的左侧竖直映射亦然。
因此对变量数作归纳，该引理由引理
\ref{more-algebra-lemma-add-variable-pseudo-coherent} 得到。
伪凝聚情形则由此和引理 \ref{more-algebra-lemma-pseudo-coherent} 得到。
\end{proof}

\begin{definition}
\label{more-algebra-definition-relatively-pseudo-coherent}
设 $R \to A$ 为有限型环同态，设 $K^\bullet$ 为 $A$-模复形，
设 $M$ 为 $A$-模，并设 $m \in \mathbf{Z}$。
\begin{enumerate}
\item 若引理 \ref{more-algebra-lemma-relatively-pseudo-coherent} 中的等价条件成立，
则称 $K^\bullet$ {\it 相对于 $R$ 是 $m$-伪凝聚的}。
\item 若对所有 $m \in \mathbf{Z}$，$K^\bullet$ 相对于 $R$ 都是
$m$-伪凝聚的，则称 $K^\bullet$ {\it 相对于 $R$ 是伪凝聚的}。
\item 若 $M[0]$ 相对于 $R$ 是 $m$-伪凝聚的，则称 $M$
{\it 相对于 $R$ 是 $m$-伪凝聚的}。
\item 若 $M[0]$ 相对于 $R$ 是伪凝聚的，则称 $M$
{\it 相对于 $R$ 是伪凝聚的}。
\end{enumerate}
\end{definition}

\noindent
% P02-E0298
第 (2) 项意指：对任意满射 $R[y_1, \ldots, y_m] \to A$，
$K^\bullet$ 作为 $R[x_1, \ldots, x_n]$-模复形是伪凝聚的；
参见引理 \ref{more-algebra-lemma-pseudo-coherent}。该定义具有下述良好性质。

\begin{lemma}
\label{more-algebra-lemma-finite-extension-pseudo-coherent}
设 $R$ 为一个环，设 $A \to B$ 为有限型 $R$-代数之间的有限同态，
设 $m \in \mathbf{Z}$，并设 $K^\bullet$ 为 $B$-模复形。
则 $K^\bullet$ 相对于 $R$ 是 $m$-伪凝聚的（相应地，伪凝聚的），
当且仅当把 $K^\bullet$ 视为 $A$-模复形后，它相对于 $R$ 是
$m$-伪凝聚的（伪凝聚的）。
\end{lemma}

\begin{proof}
选取满射 $R[x_1, \ldots, x_n] \to A$。
% P02-E0299
选取 $y_1, \ldots, y_m \in B$，使它们在 $A$ 上生成 $B$。
由于 $A \to B$ 是有限的，每个 $y_i$ 都满足系数在 $A$ 中的首一方程。
因此，可以找到首一多项式
$P_j(T) \in R[x_1, \ldots, x_n][T]$，使得在 $B$ 中
$P_j(y_j) = 0$。于是得到交换图
$$
\xymatrix{
& R[x_1, \ldots, x_n, y_1, \ldots, y_m] \ar[d] \\
R[x_1, \ldots, x_n] \ar[d] \ar[r] &
R[x_1, \ldots, x_n, y_1, \ldots, y_m]/(P_j(y_j)) \ar[d] \\
A \ar[r] & B
}
$$
顶部水平箭头和右上方竖直箭头满足引理
\ref{more-algebra-lemma-finite-push-pseudo-coherent} 的假设。因此，
$K^\bullet$ 作为 $R[x_1, \ldots, x_n]$-模复形是
$m$-伪凝聚的（相应地，伪凝聚的），当且仅当 $K^\bullet$ 作为
$R[x_1, \ldots, x_n, y_1, \ldots, y_m]$-模复形是
$m$-伪凝聚的（相应地，伪凝聚的）。
\end{proof}

\begin{lemma}
\label{more-algebra-lemma-cone-relatively-pseudo-coherent}
设 $R$ 为一个环，设 $R \to A$ 为有限型环同态，设
$m \in \mathbf{Z}$，并设 $(K^\bullet, L^\bullet, M^\bullet, f, g, h)$
为 $D(A)$ 中的区别三角。
\begin{enumerate}
\item 若 $K^\bullet$ 相对于 $R$ 是 $(m + 1)$-伪凝聚的，且
$L^\bullet$ 相对于 $R$ 是 $m$-伪凝聚的，则 $M^\bullet$
相对于 $R$ 是 $m$-伪凝聚的。
\item 若 $K^\bullet, M^\bullet$ 相对于 $R$ 是 $m$-伪凝聚的，
则 $L^\bullet$ 相对于 $R$ 是 $m$-伪凝聚的。
\item 若 $L^\bullet$ 相对于 $R$ 是 $(m + 1)$-伪凝聚的，且
$M^\bullet$ 相对于 $R$ 是 $m$-伪凝聚的，则 $K^\bullet$
相对于 $R$ 是 $(m + 1)$-伪凝聚的。
\end{enumerate}
此外，若 $K^\bullet, L^\bullet, M^\bullet$ 中有两个相对于 $R$
是伪凝聚的，则第三个也如此。
% P02-E0300
\end{lemma}

\begin{proof}
% P02-E0301
这立即由引理 \ref{more-algebra-lemma-cone-pseudo-coherent} 和定义得到。
\end{proof}

\begin{lemma}
\label{more-algebra-lemma-rel-n-pseudo-module}
设 $R \to A$ 为有限型环同态，设 $M$ 为 $A$-模。则
\begin{enumerate}
\item $M$ 相对于 $R$ 是 $0$-伪凝聚的，当且仅当 $M$ 是有限型
$A$-模；
% P02-E0302
\item $M$ 相对于 $R$ 是 $(-1)$-伪凝聚的，当且仅当 $M$
相对于 $R$ 有限表示；
% P02-E0303
\item $M$ 相对于 $R$ 是 $(-d)$-伪凝聚的，当且仅当对每个满射
$R[x_1, \ldots, x_n] \to A$，都存在长度为 $d$ 的解消
$$
R[x_1, \ldots, x_n]^{\oplus a_d} \to R[x_1, \ldots, x_n]^{\oplus a_{d - 1}}
\to \ldots \to R[x_1, \ldots, x_n]^{\oplus a_0} \to M \to 0
$$
；
\item $M$ 相对于 $R$ 是伪凝聚的，当且仅当对每个表示
$R[x_1, \ldots, x_n] \to A$，都存在由有限自由
$R[x_1, \ldots, x_n]$-模组成的无穷解消
$$
\ldots \to R[x_1, \ldots, x_n]^{\oplus a_1} \to
R[x_1, \ldots, x_n]^{\oplus a_0} \to M \to 0
$$
。
\end{enumerate}
\end{lemma}

\begin{proof}
% P02-E0304
这立即由引理 \ref{more-algebra-lemma-n-pseudo-module} 和定义得到。
\end{proof}

\begin{lemma}
\label{more-algebra-lemma-summands-relative-pseudo-coherent}
设 $R \to A$ 为有限型环同态，设 $m \in \mathbf{Z}$，并设
$K^\bullet, L^\bullet \in D(A)$。若 $K^\bullet \oplus L^\bullet$
相对于 $R$ 是 $m$-伪凝聚的（相应地，伪凝聚的），则
$K^\bullet$ 与 $L^\bullet$ 也如此。
\end{lemma}

\begin{proof}
% P02-E0304
这立即由引理 \ref{more-algebra-lemma-summands-pseudo-coherent} 和定义得到。
\end{proof}

\begin{lemma}
\label{more-algebra-lemma-complex-relative-pseudo-coherent-modules}
设 $R \to A$ 为有限型环同态，设 $m \in \mathbf{Z}$，并设
$K^\bullet$ 为有界上 $A$-模复形，使得对所有 $i$，$K^i$
相对于 $R$ 是 $(m - i)$-伪凝聚的。则 $K^\bullet$ 相对于
$R$ 是 $m$-伪凝聚的。
% P02-E0305
特别地，若 $K^\bullet$ 是一个有界上复形，其各项都是相对于 $R$
伪凝聚的 $A$-模，则 $K^\bullet$ 相对于 $R$ 是伪凝聚的。
\end{lemma}

\begin{proof}
% P02-E0304
这立即由引理 \ref{more-algebra-lemma-complex-pseudo-coherent-modules} 和定义得到。
\end{proof}

\begin{lemma}
\label{more-algebra-lemma-cohomology-relative-pseudo-coherent}
设 $R \to A$ 为有限型环同态，设 $m \in \mathbf{Z}$，并设
$K^\bullet \in D^{-}(A)$，使得对所有 $i$，$H^i(K^\bullet)$
相对于 $R$ 是 $(m - i)$-伪凝聚的（相应地，伪凝聚的）。
则 $K^\bullet$ 相对于 $R$ 是 $m$-伪凝聚的
（相应地，伪凝聚的）。
\end{lemma}

\begin{proof}
% P02-E0304
这立即由引理 \ref{more-algebra-lemma-cohomology-pseudo-coherent} 和定义得到。
\end{proof}

\begin{lemma}
\label{more-algebra-lemma-localize-relative-pseudo-coherent}
% P02-E0306
设 $R$ 为一个环，$f \in R$ 为元素，$R_f \to A$ 为有限型环同态，
$g \in A$，并设 $K^\bullet$ 为 $A$-模复形。若 $K^\bullet$
相对于 $R_f$ 是 $m$-伪凝聚的（相应地，伪凝聚的），则
$K^\bullet \otimes_A A_g$ 相对于 $R$ 是 $m$-伪凝聚的
（相应地，伪凝聚的）。
\end{lemma}

\begin{proof}
首先证明 $K^\bullet$ 相对于 $R$ 是 $m$-伪凝聚的。
具体地，设 $R_f[x_1, \ldots, x_n] \to A$ 为满射。写成
$R_f = R[x_0]/(fx_0 - 1)$。则 $R[x_0, x_1, \ldots, x_n] \to A$
是满射，而 $R_f[x_1, \ldots, x_n]$ 作为
$R[x_0, \ldots, x_n]$-模是伪凝聚的。因此由引理
\ref{more-algebra-lemma-finite-push-pseudo-coherent}，$K^\bullet$ 作为
$R[x_0, x_1, \ldots, x_n]$-模复形是 $m$-伪凝聚的。

\medskip\noindent
选取映到 $g \in A$ 的元素
$g' \in R[x_0, x_1, \ldots, x_n]$。由引理
\ref{more-algebra-lemma-pull-pseudo-coherent}，可知
% P02-E0307
\begin{align*}
K^\bullet \otimes_{R[x_0, x_1, \ldots, x_n]}^{\mathbf{L}}
R[x_0, x_1, \ldots, x_n, \frac{1}{g'}] & =
K^\bullet \otimes_{R[x_0, x_1, \ldots, x_n]}
R[x_0, x_1, \ldots, x_n, \frac{1}{g'}] \\
& = K^\bullet \otimes_A A_f
\end{align*}
作为 $R[x_0, x_1, \ldots, x_n, \frac{1}{g'}]$-模复形是
$m$-伪凝聚的。
% P02-E0308
写成
$$
R[x_0, x_1, \ldots, x_n, \frac{1}{g'}] =
R[x_0, \ldots, x_n, x_{n + 1}]/(x_{n + 1}g' - 1).
$$
由于 $R[x_0, x_1, \ldots, x_n, \frac{1}{g'}]$ 作为
$R[x_0, \ldots, x_n, x_{n + 1}]$-模是伪凝聚的，由引理
\ref{more-algebra-lemma-finite-push-pseudo-coherent} 可知，
$K^\bullet \otimes_A A_g$ 作为
$R[x_0, \ldots, x_n, x_{n + 1}]$-模复形是 $m$-伪凝聚的，
正合所需。
\end{proof}

\begin{lemma}
\label{more-algebra-lemma-base-change-relative-pseudo-coherent}
设 $R \to A$ 为有限型环同态，设 $m \in \mathbf{Z}$，并设
$K^\bullet$ 为相对于 $R$ 是 $m$-伪凝聚的（相应地，伪凝聚的）
$A$-模复形。设 $R \to R'$ 为环同态，使得 $A$ 与 $R'$ 在 $R$
上 Tor 独立。令 $A' = A \otimes_R R'$。则
$K^\bullet \otimes_A^{\mathbf{L}} A'$ 相对于 $R'$ 是
$m$-伪凝聚的（相应地，伪凝聚的）。
\end{lemma}

\begin{proof}
选取满射 $R[x_1, \ldots, x_n] \to A$。注意
$$
K^\bullet \otimes_A^{\mathbf{L}} A' =
K^\bullet \otimes_R^{\mathbf{L}} R' =
K^\bullet \otimes_{R[x_1, \ldots, x_n]}^{\mathbf{L}} R'[x_1, \ldots, x_n]
$$
这是两次应用引理 \ref{more-algebra-lemma-base-change-comparison} 得到的。
于是由引理 \ref{more-algebra-lemma-pull-pseudo-coherent} 即得结论。
\end{proof}

\begin{lemma}
\label{more-algebra-lemma-pull-relative-pseudo-coherent}
设 $R \to A \to B$ 为有限型环同态，设 $m \in \mathbf{Z}$，
并设 $K^\bullet$ 为 $A$-模复形。假设 $B$ 作为 $B$-模相对于
$A$ 是伪凝聚的。若 $K^\bullet$ 相对于 $R$ 是 $m$-伪凝聚的
（相应地，伪凝聚的），则 $K^\bullet \otimes_A^{\mathbf{L}} B$
相对于 $R$ 是 $m$-伪凝聚的（相应地，伪凝聚的）。
\end{lemma}

\begin{proof}
% P02-E0309
选取满射 $A[y_1, \ldots, y_m] \to B$，并选取满射
$R[x_1, \ldots, x_n] \to A$。合起来得到满射
% P02-E0310
$R[x_1, \ldots, x_n, y_1, \ldots y_m] \to B$。
% P02-E0311
选取 $B$ 的解消 $E^\bullet \to B$，其中它是有限自由
$A[y_1, \ldots, y_n]$-模复形（由关于环同态 $A \to B$ 的假设，
这是可能的）。可以假设 $K^\bullet$ 是平坦 $A$-模组成的有界上复形。
于是
\begin{align*}
K^\bullet \otimes_A^{\mathbf{L}} B & =
\text{Tot}(K^\bullet \otimes_A B[0]) \\
& = \text{Tot}(K^\bullet \otimes_A A[y_1, \ldots, y_m]
\otimes_{A[y_1, \ldots, y_m]} B[0]) \\
& \cong
\text{Tot}\left(
(K^\bullet \otimes_A A[y_1, \ldots, y_m])
\otimes_{A[y_1, \ldots, y_m]} E^\bullet
\right) \\
& =
\text{Tot}(K^\bullet \otimes_A E^\bullet)
\end{align*}
在 $D(A[y_1, \ldots, y_m])$ 中成立。同构 $\cong$ 来自引理
\ref{more-algebra-lemma-derived-tor-quasi-isomorphism} 的一次应用。
因此，必须证明 $\text{Tot}(K^\bullet \otimes_A E^\bullet)$ 作为
$R[x_1, \ldots, x_n, y_1, \ldots y_m]$-模复形是 $m$-伪凝聚的。
注意，$\text{Tot}(K^\bullet \otimes_A E^\bullet)$ 有一个由子复形组成的滤过，
其逐次商为复形 $K^\bullet \otimes_A E^i[-i]$。当 $i \ll 0$ 时，
复形 $K^\bullet \otimes_A E^i[-i]$ 在次数 $\leq m$ 的上同调为零，
因而是 $m$-伪凝聚的（在任意环上皆然）。因此，应用引理
\ref{more-algebra-lemma-cone-relatively-pseudo-coherent} 并作归纳，只需证明
% P02-E0312
$K^\bullet \otimes_A E^i[-i]$ 对所有 $i$ 都相对于 $R$ 是伪凝聚的。
注意当 $i > 0$ 时 $E^i = 0$。又因 $E^i$ 是有限自由的，
这归结为证明 $K^\bullet \otimes_A A[y_1, \ldots, y_m]$
相对于 $R$ 是 $m$-伪凝聚的；例如这由引理
\ref{more-algebra-lemma-base-change-relative-pseudo-coherent} 得到。
\end{proof}

\begin{lemma}
\label{more-algebra-lemma-pull-relative-pseudo-coherent-module}
设 $R \to A \to B$ 为有限型环同态，设 $m \in \mathbf{Z}$，
并设 $M$ 为 $A$-模。假设 $B$ 在 $A$ 上平坦，且 $B$ 作为
$B$-模相对于 $A$ 是伪凝聚的。若 $M$ 相对于 $R$ 是
$m$-伪凝聚的（相应地，伪凝聚的），则 $M \otimes_A B$
相对于 $R$ 是 $m$-伪凝聚的（相应地，伪凝聚的）。
\end{lemma}

\begin{proof}
% P02-E0313
这立即由引理 \ref{more-algebra-lemma-pull-relative-pseudo-coherent} 得到。
\end{proof}

\begin{lemma}
\label{more-algebra-lemma-composition-relative-pseudo-coherent}
设 $R$ 为一个环，设 $A \to B$ 为有限型 $R$-代数之间的同态，
设 $m \in \mathbf{Z}$，并设 $K^\bullet$ 为 $B$-模复形。
假设 $A$ 相对于 $R$ 是伪凝聚的。则下列条件等价：
\begin{enumerate}
\item $K^\bullet$ 相对于 $A$ 是 $m$-伪凝聚的
（相应地，伪凝聚的）；
\item $K^\bullet$ 相对于 $R$ 是 $m$-伪凝聚的
（相应地，伪凝聚的）。
\end{enumerate}
\end{lemma}

\begin{proof}
选取满射 $R[x_1, \ldots, x_n] \to A$。
% P02-E0314
选取满射 $A[y_1, \ldots, y_m] \to B$。于是得到满射
$$
R[x_1, \ldots, x_n, y_1, \ldots, y_m] \to A[y_1, \ldots, y_m]
$$
它是 $R[x_1, \ldots, x_n] \to A$ 的平坦基变换。
根据假设，$A$ 是伪凝聚 $R[x_1, \ldots, x_n]$-模；故由引理
\ref{more-algebra-lemma-flat-base-change-pseudo-coherent}，可知
$A[y_1, \ldots, y_m]$ 在
$R[x_1, \ldots, x_n, y_1, \ldots, y_m]$ 上是伪凝聚的。
因此，该引理由引理 \ref{more-algebra-lemma-finite-push-pseudo-coherent} 和定义得到。
\end{proof}

\begin{lemma}
\label{more-algebra-lemma-glue-relative-pseudo-coherent}
设 $R \to A$ 为有限型环同态，设 $K^\bullet$ 为 $A$-模复形，
设 $m \in \mathbf{Z}$，并设 $f_1, \ldots, f_r \in A$ 生成单位理想。
下列条件等价：
\begin{enumerate}
\item 每个 $K^\bullet \otimes_A A_{f_i}$ 都相对于 $R$ 是
$m$-伪凝聚的；
\item $K^\bullet$ 相对于 $R$ 是 $m$-伪凝聚的。
\end{enumerate}
对于相对于 $R$ 的伪凝聚性，也有同样的等价性。
\end{lemma}

\begin{proof}
(2) $\Rightarrow$ (1) 是引理
\ref{more-algebra-lemma-localize-relative-pseudo-coherent} 的内容。
假设 (1)。在 $A$ 中写成 $1 = \sum f_ig_i$。
% P02-E0315
选取满射
$R[x_1, \ldots, x_n, y_1, \ldots, y_r, z_1, \ldots, z_r] \to A$，
使得 $y_i$ 映到 $f_i$，而 $z_i$ 映到 $g_i$。于是可见存在满射
$$
P = R[x_1, \ldots, x_n, y_1, \ldots, y_r, z_1, \ldots, z_r]/(\sum y_iz_i - 1)
\longrightarrow
A.
$$
注意，$P$ 作为
$R[x_1, \ldots, x_n, y_1, \ldots, y_r, z_1, \ldots, z_r]$-模是伪凝聚的，
而 $P[1/y_i]$ 作为
$R[x_1, \ldots, x_n, y_1, \ldots, y_r, z_1, \ldots, z_r, 1/y_i]$-模
是伪凝聚的。因此，由引理 \ref{more-algebra-lemma-finite-push-pseudo-coherent}，
对每个 $i$，$K^\bullet \otimes_A A_{f_i}$ 都是
$P[1/y_i]$-模的 $m$-伪凝聚复形。于是由引理
\ref{more-algebra-lemma-glue-pseudo-coherent}，可知 $K^\bullet$ 作为
% P02-E0316
$P$-模复形是伪凝聚的；再由引理
\ref{more-algebra-lemma-finite-push-pseudo-coherent}，可知 $K^\bullet$ 作为
$R[x_1, \ldots, x_n, y_1, \ldots, y_r, z_1, \ldots, z_r]$-模复形
是伪凝聚的。
\end{proof}

\begin{lemma}
\label{more-algebra-lemma-Noetherian-relative-pseudo-coherent}
设 $R$ 为 Noether 环，并设 $R \to A$ 为有限型环同态。则
\begin{enumerate}
\item $A$-模复形 $K^\bullet$ 相对于 $R$ 是 $m$-伪凝聚的，
当且仅当 $K^\bullet \in D^{-}(A)$，并且当 $i \geq m$ 时
$H^i(K^\bullet)$ 是有限 $A$-模。
\item $A$-模复形 $K^\bullet$ 相对于 $R$ 是伪凝聚的，
当且仅当 $K^\bullet \in D^{-}(A)$，并且对所有 $i$，
$H^i(K^\bullet)$ 都是有限 $A$-模。
\item 一个 $A$-模相对于 $R$ 是伪凝聚的，当且仅当它是有限的。
\end{enumerate}
\end{lemma}

\begin{proof}
% P02-E0317
这是引理 \ref{more-algebra-lemma-Noetherian-pseudo-coherent} 和定义的直接推论。
\end{proof}










\section{伪凝聚与完美环同态}
\label{more-algebra-section-pseudo-coherent-perfect-ring-map}

\noindent
可以如下定义这些类型的环同态。

\begin{definition}
\label{more-algebra-definition-pseudo-coherent-perfect}
设 $A \to B$ 为环同态。
\begin{enumerate}
\item 若 $A \to B$ 是有限型的，且 $B$ 作为 $B$-模相对于 $A$
是伪凝聚的，则称该同态为{\it 伪凝聚环同态}。
\item 若 $A \to B$ 是伪凝聚环同态，并且 $B$ 作为 $A$-模具有有限
Tor 维数，则称该同态为{\it 完美环同态}。
\end{enumerate}
\end{definition}

\noindent
该术语可能并不标准。由引理 \ref{more-algebra-lemma-rel-n-pseudo-module}，
可知 $A \to B$ 是伪凝聚的，当且仅当
$B = A[x_1, \ldots, x_n]/I$，且 $B$ 作为
$A[x_1, \ldots, x_n]$-模具有由有限自由
$A[x_1, \ldots, x_n]$-模组成的解消。完美环同态定义的动机来自引理
\ref{more-algebra-lemma-perfect}。
% P02-E0318
下述引理给出完美环同态的一个更有用且更直观的刻画。

\begin{lemma}
\label{more-algebra-lemma-perfect-ring-map}
环同态 $A \to B$ 是完美的，当且仅当
$B = A[x_1, \ldots, x_n]/I$，且 $B$ 作为
$A[x_1, \ldots, x_n]$-模具有由有限投射
$A[x_1, \ldots, x_n]$-模组成的有限解消。
\end{lemma}

\begin{proof}
若 $A \to B$ 是完美的，则 $B = A[x_1, \ldots, x_n]/I$，
$B$ 作为 $A[x_1, \ldots, x_n]$-模是伪凝聚的，并且作为
$A$-模具有有限 Tor 维数。因此，引理
\ref{more-algebra-lemma-perfect-over-polynomial-ring} 表明 $B$ 作为
% P02-E0319
$A[x_1, \ldots, x_n]$-模是完美的，即它具有由有限投射
$A[x_1, \ldots, x_n]$-模组成的有限解消
（引理 \ref{more-algebra-lemma-perfect-module}）。
反之，若 $B = A[x_1, \ldots, x_n]/I$，且 $B$ 作为
$A[x_1, \ldots, x_n]$-模具有由有限投射
$A[x_1, \ldots, x_n]$-模组成的有限解消，则 $B$ 作为
$A[x_1, \ldots, x_n]$-模是伪凝聚的，故 $A \to B$ 是伪凝聚的。
此外，给定的 $A[x_1, \ldots, x_n]$ 上解消是由平坦 $A$-模组成的
有限解消，故 $B$ 作为 $A$-模具有有限 Tor 维数。
\end{proof}

\noindent
前几节的许多结果都可以用这一术语重新表述。
% P02-E0320
关于概形态射的相应讨论，还可参见 More on Morphisms，第
\ref{more-morphisms-section-pseudo-coherent} 节和第
\ref{more-morphisms-section-perfect} 节。

\begin{lemma}
\label{more-algebra-lemma-Noetherian-pseudo-coherent-ring-map}
Noether 环之间的有限型环同态是伪凝聚的。
\end{lemma}

\begin{proof}
参见引理 \ref{more-algebra-lemma-Noetherian-relative-pseudo-coherent}。
\end{proof}

\begin{lemma}
\label{more-algebra-lemma-flat-finite-presentation-perfect}
平坦且有限表示的环同态是完美的。
\end{lemma}

\begin{proof}
设 $A \to B$ 为平坦且有限表示的环同态。显然，$B$ 具有有限
Tor 维数。由 Algebra，引理
\ref{algebra-lemma-flat-finite-presentation-limit-flat}，
存在有限型 $\mathbf{Z}$-代数 $A_0 \subset A$ 和平坦有限型环同态
$A_0 \to B_0$，使得 $B = B_0 \otimes_{A_0} A$。由引理
\ref{more-algebra-lemma-Noetherian-relative-pseudo-coherent}，可知
$A_0 \to B_0$ 是伪凝聚的。由于 $A_0 \to B_0$ 平坦，
$B_0$ 与 $A$ 在 $A_0$ 上 Tor 独立，故可以应用引理
\ref{more-algebra-lemma-base-change-relative-pseudo-coherent}，得到
$A \to B$ 是伪凝聚的。
\end{proof}

\begin{lemma}
\label{more-algebra-lemma-regular-perfect-ring-map}
设 $A \to B$ 为有限型环同态，其中 $A$ 是有限维正则环。
则 $A \to B$ 是完美的。
\end{lemma}

\begin{proof}
由 Algebra，引理
\ref{algebra-lemma-finite-gl-dim-finite-dim-regular}，关于 $A$ 的假设
意指 $A$ 具有有限整体维数。因此，每个模都具有有限 Tor 维数；
参见引理 \ref{more-algebra-lemma-finite-gl-dim-tor-dimension}，特别地 $B$ 也如此。
由引理 \ref{more-algebra-lemma-Noetherian-pseudo-coherent-ring-map}，该同态是伪凝聚的。
\end{proof}

\begin{lemma}
\label{more-algebra-lemma-lci-perfect}
局部完全交同态是完美的。
\end{lemma}

\begin{proof}
设 $A \to B$ 为局部完全交同态。由定义
\ref{more-algebra-definition-local-complete-intersection}，这意指
$B = A[x_1, \ldots, x_n]/I$，其中 $I$ 是
$A[x_1, \ldots, x_n]$ 中的 Koszul 理想。
由引理 \ref{more-algebra-lemma-perfect-ring-map} 和 \ref{more-algebra-lemma-perfect-module}，
只需证明 $I$ 是完美 $A[x_1, \ldots, x_n]$-模。
由引理 \ref{more-algebra-lemma-glue-perfect}，这是局部问题。因此由定义
\ref{more-algebra-definition-regular-ideal}，可以假设 $I$ 由一个 Koszul-正则序列生成。
这当然意指 $I$ 具有有限自由解消，结论得证。
\end{proof}

\begin{lemma}
\label{more-algebra-lemma-relative-pseudo-coherent-is-moot}
设 $R \to A$ 为伪凝聚环同态，并设 $K \in D(A)$。下列条件等价：
\begin{enumerate}
\item $K$ 相对于 $R$ 是 $m$-伪凝聚的（相应地，伪凝聚的）；
\item $K$ 在 $D(A)$ 中是 $m$-伪凝聚的（相应地，伪凝聚的）。
\end{enumerate}
\end{lemma}

\begin{proof}
这是引理 \ref{more-algebra-lemma-composition-relative-pseudo-coherent} 一个特例的重述。
\end{proof}

\begin{lemma}
\label{more-algebra-lemma-more-relative-pseudo-coherent-is-moot}
设 $R \to B \to A$ 为环同态，其中
$\varphi : B \to A$ 满射，并且 $R \to B$ 与 $R \to A$
平坦且有限表示。
% P02-E0321
对 $K \in D(A)$，用 $\varphi_*K \in D(B)$ 表示限制。下列条件等价：
\begin{enumerate}
\item $K$ 是伪凝聚的；
\item $K$ 相对于 $R$ 是伪凝聚的；
\item $K$ 相对于 $A$ 是伪凝聚的；
\item $\varphi_*K$ 是伪凝聚的；
\item $\varphi_*K$ 相对于 $R$ 是伪凝聚的。
\end{enumerate}
% P02-E0322
对于 $m$-伪凝聚性也有类似结论。
\end{lemma}

\begin{proof}
注意，$R \to A$ 和 $R \to B$ 是完美环同态
（引理 \ref{more-algebra-lemma-flat-finite-presentation-perfect}），
因而更是伪凝聚环同态。因此，由引理
\ref{more-algebra-lemma-relative-pseudo-coherent-is-moot}，有
(1) $\Leftrightarrow$ (2) 和 (4) $\Leftrightarrow$ (5)。

\medskip\noindent
利用 $A$ 相对于 $R$ 是伪凝聚的这一事实，应用引理
\ref{more-algebra-lemma-composition-relative-pseudo-coherent} 可知
(2) $\Leftrightarrow$ (3)。
% P02-E0323 P02-E0324
另一方面，由于 $A \to B$ 是满射，从定义
\ref{more-algebra-definition-relatively-pseudo-coherent} 直接可见
(3) 与 (4) 等价。
\end{proof}




\section{相对完美模}
\label{more-algebra-section-relatively-perfect}

\noindent
本节给出第 \ref{more-algebra-section-relative-pseudo-coherent} 节关于导出范畴中
完美对象的类似理论。
% P02-E0325
由于尚不确定一般情形下的正确概念是什么，我们只在有限的一般性下定义这一概念。
相关讨论见 Derived Categories of Schemes，评注
\ref{perfect-remark-discuss-rel-perfect}。

\begin{definition}
\label{more-algebra-definition-relatively-perfect}
设 $R \to A$ 为平坦有限表示环同态。若 $D(A)$ 的对象 $K$
是伪凝聚的（定义 \ref{more-algebra-definition-pseudo-coherent}），并且在
$R$ 上具有有限 Tor 维数（定义 \ref{more-algebra-definition-tor-amplitude}），
则称 $K$ 是{\it $R$-完美的}，或称它{\it 相对于 $R$ 完美}。
\end{definition}

\noindent
由引理 \ref{more-algebra-lemma-more-relative-pseudo-coherent-is-moot}，
要求 $K$ 相对于 $R$ 伪凝聚会得到同一个概念。
下面给出一些必备引理。

\begin{lemma}
\label{more-algebra-lemma-cone-relatively-perfect}
设 $R \to A$ 为平坦有限表示环同态。$D(A)$ 中的 $R$-完美对象组成
饱和的\footnote{Derived Categories，定义
\ref{derived-definition-saturated}。}三角严格满子范畴。
\end{lemma}

\begin{proof}
这由引理 \ref{more-algebra-lemma-cone-pseudo-coherent}、
\ref{more-algebra-lemma-summands-pseudo-coherent}、
\ref{more-algebra-lemma-cone-tor-amplitude} 和
\ref{more-algebra-lemma-summands-tor-amplitude} 得到。
\end{proof}

\begin{lemma}
\label{more-algebra-lemma-perfect-relatively-perfect}
设 $R \to A$ 为平坦有限表示环同态。$D(A)$ 中的完美对象是
$R$-完美的。若 $K, M \in D(A)$，$K$ 是完美的且 $M$ 是
$R$-完美的，则 $K \otimes_A^\mathbf{L} M$ 是 $R$-完美的。
\end{lemma}

\begin{proof}
在第二个断言中取 $M = A$，即得第一个断言。
第二个断言由引理 \ref{more-algebra-lemma-perfect}、
\ref{more-algebra-lemma-push-tor-amplitude} 和 \ref{more-algebra-lemma-tensor-pseudo-coherent} 得到。
\end{proof}

\begin{lemma}
\label{more-algebra-lemma-structure-relatively-perfect}
设 $R \to A$ 为平坦有限表示环同态，并设 $K \in D(A)$。
下列条件等价：
\begin{enumerate}
\item $K$ 是 $R$-完美的；
\item $K$ 同构于一个由 $R$-平坦有限表示 $A$-模组成的有限复形。
\end{enumerate}
\end{lemma}

\begin{proof}
为证明 (2) 蕴含 (1)，由引理 \ref{more-algebra-lemma-cone-relatively-perfect}，
只需证明一个 $R$-平坦有限表示 $A$-模 $M$ 定义 $D(A)$ 中的
$R$-完美对象。由于 $M$ 在 $R$ 上具有有限 Tor 维数，
只需证明 $M$ 是伪凝聚的。由 Algebra，引理
\ref{algebra-lemma-flat-finite-presentation-limit-flat}，
存在有限型 $\mathbf{Z}$-代数 $R_0 \subset R$、平坦有限型环同态
$R_0 \to A_0$，以及在 $R_0$ 上平坦的有限 $A_0$-模 $M_0$，
使得 $A = A_0 \otimes_{R_0} R$ 且
$M = M_0 \otimes_{R_0} R$。由引理
\ref{more-algebra-lemma-Noetherian-pseudo-coherent}，可知
% P02-E0326
$M_0$ 是伪凝聚 $A_0$-模。选取由有限自由 $A_0$-模
$P_0^n$ 组成的解消 $P_0^\bullet \to M_0$。由于 $A_0$ 在
$R_0$ 上平坦，这是一个平坦解消。由于 $M_0$ 在 $R_0$ 上平坦，
可知 $P^\bullet = P_0^\bullet \otimes_{R_0} R$ 仍然解消
$M = M_0 \otimes_{R_0} R$。（可以应用引理
\ref{more-algebra-lemma-base-change-comparison} 看出这一点。）
因此 $P^\bullet$ 是 $M$ 在 $A$ 上的有限自由解消，从而
$M$ 是伪凝聚的。

\medskip\noindent
假设 (1)。可以用有限自由 $A$-模组成的有界上复形
$P^\bullet$ 表示 $K$。假设把 $K$ 视为 $D(R)$ 中的对象时，
其 Tor 振幅包含在 $[a, b]$ 中。由引理
\ref{more-algebra-lemma-last-one-flat}，可知 $\tau_{\geq a}P^\bullet$
是表示 $K$ 的复形，其各项为 $R$-平坦有限表示 $A$-模。
\end{proof}

\begin{lemma}
\label{more-algebra-lemma-base-change-relatively-perfect}
设 $R \to A$ 为平坦有限表示环同态。设 $R \to R'$ 为环同态，
并令 $A' = A \otimes_R R'$。若 $K \in D(A)$ 是 $R$-完美的，
则 $K \otimes_A^\mathbf{L} A'$ 是 $R'$-完美的。
\end{lemma}

\begin{proof}
由引理 \ref{more-algebra-lemma-pull-pseudo-coherent}，可知
$K \otimes_A^\mathbf{L} A'$ 是伪凝聚的。
由引理 \ref{more-algebra-lemma-base-change-comparison}，可知
$K \otimes_A^\mathbf{L} A'$ 在 $D(R')$ 中等于
$K \otimes_R^\mathbf{L} R'$。然后可以应用引理
\ref{more-algebra-lemma-pull-tor-amplitude}，得到 $D(R')$ 中的
$K \otimes_R^\mathbf{L} R'$ 具有有限 Tor 维数。
\end{proof}

\begin{lemma}
\label{more-algebra-lemma-compute-RHom-relatively-perfect}
设 $R \to A$ 为平坦环同态，设 $K, L \in D(A)$，其中 $K$
是伪凝聚的，而 $L$ 在 $R$ 上具有有限 Tor 维数。可以选取：
\begin{enumerate}
\item 表示 $K$ 的有限自由 $A$-模有界上复形 $P^\bullet$；
\item 表示 $L$ 的 $R$-平坦 $A$-模有界复形 $F^\bullet$。
\end{enumerate}
在作出这些选择后，有：
\begin{enumerate}
\item[(a)] $E^\bullet = \Hom^\bullet(P^\bullet, F^\bullet)$
是表示 $R\Hom_A(K, L)$ 的 $R$-平坦 $A$-模有界下复形；
\item[(b)] 对任意环同态 $R \to R'$，令 $A' = A \otimes_R R'$，
则复形 $E^\bullet \otimes_R R'$ 表示
$R\Hom_{A'}(K \otimes_A^\mathbf{L} A',
L \otimes_A^\mathbf{L} A')$。
\end{enumerate}
若还假设 $R \to A$ 是有限表示的，且 $L$ 是 $R$-完美的，
则可以选取各 $F^p$ 为有限表示 $A$-模，因而各 $E^n$ 也是有限表示
$A$-模。
\end{lemma}

\begin{proof}
$P^\bullet$ 的存在性就是伪凝聚复形的定义。首先用自由 $A$-模组成的
有界上复形 $F^\bullet$ 表示 $L$（这是可能的，因为有界 Tor 维数
尤其蕴含有界）。再设把 $L$ 视为 $D(R)$ 中的对象时，其 Tor 振幅
包含在 $[a, b]$ 中。用 $\tau_{\geq a}F^\bullet$ 替换
$F^\bullet$ 后，得到 (2) 中的复形。这由引理
\ref{more-algebra-lemma-last-one-flat} 得到。

\medskip\noindent
(a) 的证明。
% P02-E0327
由于 $F^\bullet$ 有界且 $P^\bullet$ 有界上，可知当
$n \ll 0$ 时 $E^n = 0$，并且 $E^n$ 是有限（！）直和
$$
E^n = \bigoplus\nolimits_{p + q = n} \Hom_A(P^{-q}, F^p)
$$
由于 $P^{-q}$ 是有限自由的，这的确是 $R$-平坦 $A$-模。
$E^\bullet$ 表示 $R\Hom_A(K, L)$ 这一事实由引理
\ref{more-algebra-lemma-RHom-out-of-projective} 得到。

\medskip\noindent
(b) 的证明。
设 $R \to R'$ 为环同态，并令 $A' = A \otimes_R R'$。
由引理 \ref{more-algebra-lemma-base-change-comparison}，对象
$L \otimes_A^\mathbf{L} A'$ 由视为 $A'$-模复形的
$F^\bullet \otimes_R R'$ 表示（因为各 $F^p$ 在 $R$ 上平坦）。
对 $P^\bullet \otimes_R R'$ 亦然。与上面一样，
$R\Hom_{A'}(K \otimes_A^\mathbf{L} A',
L \otimes_A^\mathbf{L} A')$ 由
$$
\Hom^\bullet(P^\bullet \otimes_R R', F^\bullet \otimes_R R') =
E^\bullet \otimes_R R'
$$
表示。逐项考察该复形，并应用
$\Hom_{A'}(P^{-q} \otimes_R R', F^p \otimes_R R') =
\Hom_A(P^{-q}, F^p) \otimes_R R'$，即得该等式。
\end{proof}

\begin{lemma}
\label{more-algebra-lemma-colimit-relatively-perfect}
设 $R = \colim_{i \in I} R_i$ 为环的滤过余极限。
取 $0 \in I$，并设 $R_0 \to A_0$ 为平坦有限表示环同态。
对 $i \geq 0$，令 $A_i = R_i \otimes_{R_0} A_0$，并令
$A = R \otimes_{R_0} A_0$。
\begin{enumerate}
\item 给定 $D(A)$ 中的 $R$-完美对象 $K$，存在 $i \in I$
以及 $D(A_i)$ 中的 $R_i$-完美对象 $K_i$，使得在 $D(A)$ 中
$K \cong K_i \otimes_{A_i}^\mathbf{L} A$。
\item 给定 $K_0, L_0 \in D(A_0)$，其中 $K_0$ 伪凝聚，
且 $L_0$ 在 $R_0$ 上具有有限 Tor 维数，则有
$$
\Hom_{D(A)}(K_0 \otimes_{A_0}^\mathbf{L} A, L_0 \otimes_{A_0}^\mathbf{L} A) =
\colim_{i \geq 0}
\Hom_{D(A_i)}(K_0 \otimes_{A_0}^\mathbf{L} A_i,
L_0 \otimes_{A_0}^\mathbf{L} A_i)
$$
\end{enumerate}
特别地，$A$ 上 $R$-完美复形的三角范畴是各
$A_i$ 上 $R_i$-完美复形三角范畴的余极限。
\end{lemma}

\begin{proof}
由 Algebra，引理 \ref{algebra-lemma-colimit-category-fp-modules}，
有限表示 $A$-模范畴是各有限表示 $A_i$-模范畴的余极限。
在此基础上，由 Algebra，引理
\ref{algebra-lemma-flat-finite-presentation-limit-flat} 可知，
% P02-E0328
$R$-平坦有限表示 $A$-模范畴是各
$R_i$-平坦有限表示 $A_i$-模范畴的余极限。
因此，引理 \ref{more-algebra-lemma-structure-relatively-perfect} 中的刻画
证明了 (1)。

\medskip\noindent
为证明 (2)，按引理 \ref{more-algebra-lemma-compute-RHom-relatively-perfect}
选取表示 $K_0$ 的 $P_0^\bullet$ 和表示 $L_0$ 的
$F_0^\bullet$。于是
$E_0^\bullet = \Hom^\bullet(P_0^\bullet, F_0^\bullet)$ 满足
$$
H^0(E_0^\bullet \otimes_{R_0} R_i) =
\Hom_{D(A_i)}(K_0 \otimes_{A_0}^\mathbf{L} A_i,
L_0 \otimes_{A_0}^\mathbf{L} A_i)
$$
以及
$$
H^0(E_0^\bullet \otimes_{R_0} R) =
\Hom_{D(A)}(K_0 \otimes_{A_0}^\mathbf{L} A, L_0 \otimes_{A_0}^\mathbf{L} A)
$$
这是该引理所给出的。% P02-E0329
因此结论成立，因为张量积与余极限交换，并且滤过余极限是正合的
（Algebra，引理 \ref{algebra-lemma-directed-colimit-exact}）。
\end{proof}

\begin{lemma}
\label{more-algebra-lemma-thickening-relatively-perfect}
设 $R' \to A'$ 为平坦有限表示环同态。设 $R' \to R$
为满射环同态，其核为幂零理想。令 $A = A' \otimes_{R'} R$。
设 $K' \in D(A')$，并令
$K = K' \otimes_{A'}^\mathbf{L} A$，视为 $D(A)$ 中的对象。
若 $K$ 是 $R$-完美的，则 $K'$ 是 $R'$-完美的。
\end{lemma}

\begin{proof}
注意，$A' \to A$ 的核是幂零的，并且由 $R' \to A'$ 的平坦性，
有 $K = K' \otimes_{R'}^\mathbf{L} R$（见第
\ref{more-algebra-section-tor-independence} 节）。因此，结合引理
\ref{more-algebra-lemma-check-pseudo-coherent-modulo-nilpotent} 与
\ref{more-algebra-lemma-check-tor-dimension-modulo-nilpotent} 即得结论。
\end{proof}

\begin{lemma}
\label{more-algebra-lemma-lift-from-fibre-relatively-perfect}
设 $R$ 为环。设 $A = R[x_1, \ldots, x_d]/I$ 在 $R$ 上平坦且有限表示。
设素理想 $\mathfrak q \subset A$ 位于素理想
$\mathfrak p \subset R$ 之上。设 $K \in D(A)$ 伪凝聚，
并设 $a, b \in \mathbf{Z}$。若
$H^i(K_\mathfrak q \otimes_{R_\mathfrak p}^\mathbf{L} \kappa(\mathfrak p))$
仅在 $i \in [a, b]$ 时非零，则 $K_\mathfrak q$ 在 $R$ 上
具有 $[a - d, b]$ 内的 Tor 振幅。
\end{lemma}

\begin{proof}
由引理 \ref{more-algebra-lemma-more-relative-pseudo-coherent-is-moot}，
$K$ 作为 $R[x_1, \ldots, x_d]$-模复形是伪凝聚的。
因此可以假设 $A = R[x_1, \ldots, x_d]$。
对 $R_\mathfrak p \to A_\mathfrak q$ 和复形 $K_\mathfrak q$
应用引理 \ref{more-algebra-lemma-perfect-over-regular-local-ring}，并利用假设，
可知 $K_\mathfrak q$ 在 $D(A_\mathfrak q)$ 中是完美的，
且具有 $[a - d, b]$ 内的 Tor 振幅。由于
$R_\mathfrak p \to A_\mathfrak q$ 平坦，由引理
\ref{more-algebra-lemma-flat-push-tor-amplitude} 即得结论。
\end{proof}

\begin{lemma}
\label{more-algebra-lemma-bounded-on-fibres-relatively-perfect}
设 $R \to A$ 为平坦有限表示环同态，并设 $K \in D(A)$ 伪凝聚。
下列条件等价：
\begin{enumerate}
\item $K$ 是 $R$-完美的；
\item $K$ 下有界，并且对每个素理想 $\mathfrak p \subset R$，
对象 $K \otimes_R^\mathbf{L} \kappa(\mathfrak p)$ 下有界。
\end{enumerate}
\end{lemma}

\begin{proof}
注意，(1) 蕴含 (2)，因为按照定义，$R$-完美复形作为
$R$-模复形具有有界 Tor 维数。下面证明反向蕴含。

\medskip\noindent
写成 $A = R[x_1, \ldots, x_d]/I$。% P02-E0330
在 $D(R[x_1, \ldots, x_d])$ 中，用 $L$ 表示 $K$ 的限制。
由引理 \ref{more-algebra-lemma-more-relative-pseudo-coherent-is-moot}，
可知 $L$ 伪凝聚。由于 $L$ 与 $K$ 在 $D(R)$ 中具有相同的像，
故 $L$ 是 $R$-完美的，当且仅当 $K$ 是 $R$-完美的。
此外，$L \otimes_R^\mathbf{L} \kappa(\mathfrak p)$ 与
$K \otimes_R^\mathbf{L} \kappa(\mathfrak p)$ 是
$D(\kappa(\mathfrak p))$ 中相同的对象。这把问题归约到
$A = R[x_1, \ldots, x_d]$ 的情形。

\medskip\noindent
设 $A = R[x_1, \ldots, x_d]$，并设 $K$ 满足 (2)。
令素理想 $\mathfrak q \subset A$ 位于素理想
$\mathfrak p \subset R$ 之上。对
$R_\mathfrak p \to A_\mathfrak q$ 和复形 $K_\mathfrak q$
应用引理 \ref{more-algebra-lemma-perfect-over-regular-local-ring}，并利用假设，
可知 $K_\mathfrak q$ 在 $D(A_\mathfrak q)$ 中是完美的。
由于 $K$ 下有界，由引理 \ref{more-algebra-lemma-check-perfect-stalks}，
可知 $K$ 在 $D(A)$ 中是完美的。再由引理
\ref{more-algebra-lemma-perfect-relatively-perfect}，这蕴含 $K$ 是 $R$-完美的，
证明完成。
\end{proof}

\section{二项复形}
\label{more-algebra-section-two-term}

\noindent
本节证明若干关于模的二项复形的结果，这些结果将帮助我们理解
朴素余切复形所满足的条件。

\begin{lemma}
\label{more-algebra-lemma-ext-1-zero}
设 $R$ 为环。设 $K \in D(R)$，且对
$i \not \in \{-1, 0\}$ 有 $H^i(K) = 0$。下列条件等价：
\begin{enumerate}
\item $H^{-1}(K) = 0$，且 $H^0(K)$ 是投射模；
\item 对每个 $R$-模 $M$，有 $\Ext^1_R(K, M) = 0$。
\end{enumerate}
若 $R$ 是 Noether 环，并且对 $i = -1, 0$，$H^i(K)$ 是有限
$R$-模，则它们还等价于：
\begin{enumerate}
\item[(3)] 对每个有限 $R$-模 $M$，有 $\Ext^1_R(K, M) = 0$。
\end{enumerate}
\end{lemma}

\begin{proof}
(1) 与 (2) 的等价性由引理 \ref{more-algebra-lemma-projective-amplitude} 得到。
若 $R$ 是 Noether 环，并且对 $i = -1, 0$，$H^i(K)$ 是有限
$R$-模，则 $K$ 是伪凝聚的；见引理
\ref{more-algebra-lemma-Noetherian-pseudo-coherent}。因此，(1) 与 (3)
的等价性由引理 \ref{more-algebra-lemma-projective-amplitude-pseudo-coherent} 得到。
\end{proof}

\begin{remark}
\label{more-algebra-remark-smoothness-ext-1-zero}
下列两个陈述由引理 \ref{more-algebra-lemma-ext-1-zero}、Algebra，定义
\ref{algebra-definition-smooth} 和 Algebra，命题
\ref{algebra-proposition-characterize-formally-smooth} 得到。
\begin{enumerate}
\item 环同态 $A \to B$ 光滑，当且仅当 $A \to B$ 有限表示，且对每个
$B$-模 $N$，有 $\Ext^1_B(\NL_{B/A}, N) = 0$。
\item 环同态 $A \to B$ 形式光滑，当且仅当对每个
$B$-模 $N$，有 $\Ext^1_B(\NL_{B/A}, N) = 0$。
\end{enumerate}
\end{remark}

\begin{lemma}
\label{more-algebra-lemma-represent-two-term-complex}
设 $R$ 为环。设 $K$ 为 $D(R)$ 的对象，且对
$i \not \in \{-1, 0\}$ 有 $H^i(K) = 0$。则：
\begin{enumerate}
\item $K$ 可以由二项复形 $K^{-1} \to K^0$ 表示，其中
$K^0$ 是自由模；
\item 若 $R$ 是 Noether 环，并且对 $i = -1, 0$，$H^i(K)$
是有限 $R$-模，则 $K$ 可以由二项复形 $K^{-1} \to K^0$ 表示，
其中 $K^0$ 是有限自由模，且 $K^{-1}$ 有限。
\end{enumerate}
\end{lemma}

\begin{proof}
证明 (1)。假设 $K$ 由模复形 $M^\bullet$ 给出。首先可以用
$\tau_{\leq 0}M^\bullet$ 替换 $M^\bullet$。因此可以假设对
$i > 0$ 有 $M^i = 0$。% P02-E0331
接着，可以选取自由解消 $P^\bullet \to M^\bullet$，使得对
$i > 0$ 有 $P^i = 0$；见 Derived Categories，引理
\ref{derived-lemma-subcategory-left-resolution}。最后，可以令
$K^\bullet = \tau_{\geq -1}P^\bullet$。

\medskip\noindent
证明 (2)。假设 $R$ 是 Noether 环，并且对 $i = -1, 0$，
$H^i(K)$ 是有限 $R$-模。由引理 \ref{more-algebra-lemma-pseudo-coherent}，
可以选取拟同构 $F^\bullet \to M^\bullet$，使得对 $i > 0$
有 $F^i = 0$，且各 $F^i$ 有限自由。然后可以令
$K^\bullet = \tau_{\geq -1}F^\bullet$。
\end{proof}

\noindent
由下面的引理容易理解导出范畴中以朴素余切复形
$\NL_{B/A}$ 或 $\NL(\alpha)$ 为源的映射（见
Algebra，第 \ref{algebra-section-netherlander} 节）。

\begin{lemma}
\label{more-algebra-lemma-map-out-of-almost-free}
设 $R$ 为环。设 $M^\bullet$ 为 $R$ 上的模复形，对
$i > 0$ 有 $M^i = 0$，且 $M^0$ 是投射 $R$-模。
设 $K^\bullet$ 为另一个复形。
\begin{enumerate}
\item 假设对 $i \leq -2$ 有 $K^i = 0$。则
$\Hom_{D(R)}(M^\bullet, K^\bullet) = \Hom_{K(R)}(M^\bullet, K^\bullet)$。
\item 假设对 $i \not \in [-1, 0]$ 有 $K^i = 0$，且
$K^0$ 是投射 $R$-模。对复形映射
$a^\bullet : M^\bullet \to K^\bullet$，下列条件等价：
\begin{enumerate}
\item 对所有 $R$-模 $N$，$a^\bullet$ 诱导零映射
$\Ext^1_R(K^\bullet, N) \to \Ext^1_R(M^\bullet, N)$；
\item 存在映射 $h^0 : M^0 \to K^{-1}$，使得
$a^{-1} + h^0 \circ d^{-1}_K = 0$。
\end{enumerate}
\item 假设对 $i \leq -3$ 有 $K^i = 0$。设
$\alpha \in \Hom_{D(R)}(M^\bullet, K^\bullet)$。若 $\alpha$
与 $K^\bullet \to K^{-2}[2]$ 的复合来自 $R$-模映射
$a : M^{-2} \to K^{-2}$，且 $a \circ d_M^{-3} = 0$，则
$\alpha$ 可以由复形映射 $a^\bullet : M^\bullet \to K^\bullet$
表示，并且 $a^{-2} = a$。
\item % P02-E0332
在 (2) 中，若另一个复形映射
$(a')^\bullet : M^\bullet \to K^\bullet$ 表示 $\alpha$，且
$a = (a')^{-2}$，则对 $i = 0, -1$，存在
$h^i : M^i \to K^{i - 1}$，使得
$$
h^{-1} \circ d_M^{-2} = 0, \quad
(a')^{-1} = a^{-1} + d_K^{-2} \circ h^{-1} + h^0 \circ d_M^{-1},\quad
(a')^0 = a^0 + d_K^{-1} \circ h^0
$$
\end{enumerate}
\end{lemma}

\begin{proof}
令 $F^0 = M^0$。选取自由 $R$-模 $F^{-1}$ 以及满射
$F^{-1} \to M^{-1}$。选取自由 $R$-模 $F^{-2}$ 以及满射
$F^{-2} \to M^{-2} \times_{M^{-1}} F^{-1}$。如此继续，得到拟同构
$p^\bullet : F^\bullet \to M^\bullet$，它逐项满射，并且对所有
$i$，$F^i$ 都是投射的。

\medskip\noindent
证明 (1)。由 Derived Categories，引理
\ref{derived-lemma-morphisms-from-projective-complex}，有
$$
\Hom_{D(R)}(M^\bullet, K^\bullet) = \Hom_{K(R)}(F^\bullet, K^\bullet)
$$
若对 $i \leq -2$ 有 $K^i = 0$，则任意复形态射
$F^\bullet \to K^\bullet$ 均经由 $p^\bullet$ 分解。类似地，
任意同伦 $\{h^i : F^i \to K^{i - 1}\}$ 均经由
$p^\bullet$ 分解。因此 (1) 成立。

\medskip\noindent
证明 (2)。若 (2)(b) 成立，则 $a^\bullet$ 同伦于一个在次数
$-1$ 上为零的复形映射
$(a')^\bullet : M^\bullet \to K^\bullet$。
另一方面，设 $N \to I^\bullet$ 为内射解消。由 Derived Categories，
引理 \ref{derived-lemma-morphisms-into-injective-complex}，有
$$
\Ext^1_R(K^\bullet, N) = \Hom_{D(R)}(K^\bullet, I^\bullet[1]) =
\Hom_{K(R)}(K^\bullet, I^\bullet[1])
$$
设 $b^\bullet : K^\bullet \to I^\bullet[1]$ 为复形映射。
由于 $K^1 = 0$，映射 $b^0 : K^0 \to I^1$ 的像落在
$I^1 \to I^2$ 的核中，而该核是 $I^0 \to I^1$ 的像。
由于 $K^0$ 是投射的，可以把 $b^0$ 提升为映射
$h : K^0 \to I^0$。因此，$b^\bullet$ 同伦于复形映射
$(b')^\bullet$，且 $(b')^0 = 0$。由于对
$i \not \in [-1, 0]$ 有 $K^i = 0$，可知作为复形映射，
$(b')^\bullet \circ (a')^\bullet = 0$。故映射
$\Ext^1_R(K^\bullet, N) \to \Ext^1_R(M^\bullet, N)$ 为零。
由此可见，(2)(b) 蕴含 (2)(a)。反之，假设 (2)(a)。
可知 $\Ext^1_R(K^\bullet, K^{-1})$ 中的典范元素映到
$\Ext^1_R(M^\bullet, K^{-1})$ 中的零。利用 (1)，立即得到
(2)(b) 中的映射 $h^0$。

\medskip\noindent
证明 (3)。选取表示 $\alpha$ 的
$b^\bullet : F^\bullet \to K^\bullet$。$\alpha$ 与
$K^\bullet \to K^{-2}[2]$ 的复合由
$b^{-2} : F^{-2} \to K^{-2}$ 表示。由于它同伦于
$a \circ p^{-2} : F^{-2} \to M^{-2} \to K^{-2}$，存在映射
$h : F^{-1} \to K^{-2}$，使得
$b^{-2} = a \circ p^{-2} + h \circ d_F^{-2}$。
把 $h$ 看作从 $F^\bullet$ 到 $K^\bullet$ 的同伦，并据此调整
$b^\bullet$，可得 $b^{-2} = a \circ p^{-2}$。因此
$b^{-2}$ 经由 $p^{-2}$ 分解。由于 $F^0 = M^0$，
$p^{-2}$ 的核满射到 $p^{-1}$ 的核（例如，因为
$p^\bullet$ 的核是非周期复形，或通过追图可见）。因此
$b^{-1}$ 也必然经由 $p^{-1}$ 分解；由这些分解以及
$a^0 = b^0$，可知 (3) 成立。

\medskip\noindent
省略 (4) 的证明。提示：$a^\bullet \circ p^\bullet$ 与
$(a')^\bullet \circ p^\bullet$ 之间存在同伦，并且与前面一样论证，
该同伦经由 $p^\bullet$ 分解。
\end{proof}

\noindent
设 $A \to B$ 为有限表示环同态。给定理想 $I \subset B$，
可以考虑如下条件：
\begin{enumerate}
\item[(*)] 对所有 $B$-模 $N$，$\Ext^1_B(\NL_{B/A}, N)$ 被 $I$ 零化。
\end{enumerate}
这一条件是概念“$A \to B$ 的奇异轨迹在概形论意义下包含于
$V(I)$ 中”的一种可能的精确数学表述。请将其与评注
\ref{more-algebra-remark-smoothness-ext-1-zero} 及下面的引理作比较。

\begin{lemma}
\label{more-algebra-lemma-ext-1-annihilated-definite}
设 $R$ 为环，$I \subset R$ 为理想。设 $K \in D(R)$。
假设对 $i \not \in \{-1, 0\}$ 有 $H^i(K) = 0$。下列条件等价：
\begin{enumerate}
\item 对所有 $R$-模 $N$，$\Ext^1_R(K, N)$ 被 $I$ 零化；
\item $K$ 可以由复形 $K^{-1} \to K^0$ 表示，其中 $K^0$ 自由，
并且对任意 $a \in I$，映射 $a : K^{-1} \to K^{-1}$ 经由
$d_K^{-1} : K^{-1} \to K^0$ 分解；
\item 每当 $K$ 由二项复形 $K^{-1} \to K^0$ 表示且
$K^0$ 投射时，对任意 $a \in I$，映射
$a : K^{-1} \to K^{-1}$ 经由
$d_K^{-1} : K^{-1} \to K^0$ 分解。
\end{enumerate}
若 $R$ 是 Noether 环，并且对 $i = -1, 0$，$H^i(K)$ 是有限
$R$-模，则它们还等价于：
\begin{enumerate}
\item[(4)] 对每个有限 $R$-模 $N$，$\Ext^1_R(K, N)$ 被 $I$ 零化；
\item[(5)] $K$ 可以由复形 $K^{-1} \to K^0$ 表示，其中
$K^0$ 有限自由且 $K^{-1}$ 有限，并且对任意 $a \in I$，映射
$a : K^{-1} \to K^{-1}$ 经由
$d_K^{-1} : K^{-1} \to K^0$ 分解。
\end{enumerate}
\end{lemma}

\begin{proof}

\medskip\noindent
假设 (1)，并设 $K^{-1} \to K^0$ 是表示 $K$ 的二项复形，
其中 $K^0$ 投射。以下将直接使用引理
\ref{more-algebra-lemma-map-out-of-almost-free} 给出的以 $K^\bullet$
为源的 $D(R)$ 中映射之描述，不再另行说明。取
$N = K^{-1}$，并考虑 $\Ext^1_R(K, N)$ 中由
$\text{id}_{K^{-1}} : K^{-1} \to K^{-1}$ 给出的元素 $\xi$。
% P02-E0333
由于它被 $a \in I$ 零化，得到下列交换图中的虚线箭头：
$$
\xymatrix{
K^{-1} \ar[d]_a \ar[r]_{d_K^{-1}} & K^0 \ar@{..>}[ld]^h \\
K^{-1}
}
$$
这证明了 (3)。由引理 \ref{more-algebra-lemma-represent-two-term-complex} 的
(1)，(3) 蕴含 (2)。假设 $K^\bullet$ 如 (2) 所述，且 $N$
为任意 $R$-模。$\Ext^1_R(K, N)$ 的任意元素 $\xi$ 都由映射
$\varphi : K^{-1} \to N$ 的类给出。于是，对 $a \in I$，
由假设可以选取上图中的映射 $h$，从而
$$
a\varphi = \varphi \circ a = \varphi \circ h \circ \text{d}_K^{-1}
$$
这证明 $a \xi$ 在 $\Ext^1_R(K, N)$ 中为零。
因此 (1)、(2) 和 (3) 等价。

\medskip\noindent
假设 $R$ 是 Noether 环，并且对 $i = -1, 0$，$H^i(K)$ 是有限
$R$-模。由引理 \ref{more-algebra-lemma-represent-two-term-complex} 的 (2)，
(3) 蕴含 (5)。显然 (5) 蕴含 (2)，而 (1) 显然蕴含 (4)。
因此，为完成证明，只需证明 (4) 蕴含其余任一条件。
由引理 \ref{more-algebra-lemma-represent-two-term-complex} 的 (2)，
设 $K^{-1} \to K^0$ 是表示 $K$ 的复形，其中 $K^0$
有限自由且 $K^{-1}$ 有限。证明 (2) $\Rightarrow$ (1) 时所用的论证表明，
若 $\Ext^1_R(K, K^{-1})$ 被 $I$ 零化，则 (1) 成立。
由此可见 (4) 蕴含 (1)，证明完成。
\end{proof}

\begin{lemma}
\label{more-algebra-lemma-two-term-base-change}
设 $R$ 为环。设 $K$ 为 $D(R)$ 的对象，且对
$i \not \in \{-1, 0\}$ 有 $H^i(K) = 0$。设
$K^{-1} \to K^0$ 是表示 $K$ 的 $R$-模二项复形，并且
$K^0$ 是平坦 $R$-模（例如投射模或自由模）。设
$R \to R'$ 为环同态。则复形 $K^\bullet \otimes_R R'$
表示 $\tau_{\geq -1}(K \otimes_R^\mathbf{L} R')$。
\end{lemma}

\begin{proof}
在 $D(R)$ 中有区分三角
$$
K^0 \to K^\bullet \to K^{-1}[1] \to K^0[1]
$$
它确定区分三角之间的映射
$$
\xymatrix{
K^0 \otimes_R^\mathbf{L} R' \ar[d] \ar[r] &
K^\bullet \otimes_R^\mathbf{L} R' \ar[r] \ar[d] &
K^{-1} \otimes_R^\mathbf{L} R'[1] \ar[r] \ar[d] &
K^0 \otimes_R^\mathbf{L} R'[1] \ar[d] \\
K^0 \otimes_R R' \ar[r] &
K^\bullet \otimes_R R' \ar[r] &
K^{-1} \otimes_R R'[1] \ar[r] &
K^0 \otimes_R R'[1]
}
$$
由于 $K^0$ 平坦，左、右两个竖直箭头都是同构。映射
$K^{-1} \otimes_R^\mathbf{L} R' \to K^{-1} \otimes_R R'$
在次数 $0$ 的上同调上是同构，故得结论。
\end{proof}

\begin{lemma}
\label{more-algebra-lemma-base-change-property-ext-1-annihilated}
设 $I$ 为环 $R$ 的理想。设 $K$ 为 $D(R)$ 的对象，且对
$i \not \in \{-1, 0\}$ 有 $H^i(K) = 0$。设
$R \to R'$ 为环同态。若 $K$ 关于 $(R, I)$ 满足引理
\ref{more-algebra-lemma-ext-1-annihilated-definite} 的等价条件 (1)、(2) 和 (3)，
则 $\tau_{\geq -1}(K \otimes_R^\mathbf{L} R')$ 关于 $(R', IR')$
满足引理 \ref{more-algebra-lemma-ext-1-annihilated-definite} 的等价条件
(1)、(2) 和 (3)。
\end{lemma}

\begin{proof}
可以假设 $K$ 由二项复形 $K^{-1} \to K^0$ 表示，其中 $K^0$
自由，并且对任意 $a \in I$，映射 $a : K^{-1} \to K^{-1}$
等于某个映射 $h_a : K^0 \to K^{-1}$ 所给出的复合
$h_a \circ d_K^{-1}$。由引理 \ref{more-algebra-lemma-two-term-base-change}，
可知 $\tau_{\geq -1}(K \otimes_R^\mathbf{L} R')$ 由
$K^\bullet \otimes_R R'$ 表示。于是对每个 $a \in I$，显然有
$a \otimes 1 : K^{-1} \otimes_R R' \to K^{-1} \otimes_R R'$
等于 $(h_a \otimes 1) \circ (d_K^{-1} \otimes 1)$。
经由 $\text{d}_K^{-1} \otimes 1$ 分解的映射
$K^{-1} \otimes_R R' \to K^{-1} \otimes_R R'$ 全体组成一个
$R'$-模，故得结论。
\end{proof}

\begin{lemma}
\label{more-algebra-lemma-two-term-surjection-map-zero}
设 $R$ 为环。设 $\alpha : K \to K'$ 为 $D(R)$ 中的态射。假设：
\begin{enumerate}
\item 对 $i \not \in \{-1, 0\}$，有 $H^i(K) = H^i(K') = 0$；
\item $H^0(\alpha)$ 是同构，且 $H^{-1}(\alpha)$ 满射。
\end{enumerate}
对任意 $f \in R$，若 $f : K \to K$ 为 $0$，则
$f : K' \to K'$ 为 $0$。
\end{lemma}

\begin{proof}
令 $M = \Ker(H^{-1}(\alpha))$。于是 $\alpha$ 嵌入区分三角
$$
M[1] \to K \to K' \to M[2]
$$
由假设，$K \to K' \xrightarrow{f} K'$ 为零，故
$f : K' \to K'$ 经由某个映射 $M[2] \to K'$ 分解。然而
$\Hom(M[2], K') = 0$；例如可由 Derived Categories，引理
\ref{derived-lemma-negative-exts} 得到。
\end{proof}

\begin{lemma}
\label{more-algebra-lemma-surjection-property-ext-1-annihilated}
设 $I$ 为环 $R$ 的理想。设 $\alpha : K \to K'$ 为
$D(R)$ 中的态射。假设：
\begin{enumerate}
\item 对 $i \not \in \{-1, 0\}$，有 $H^i(K) = H^i(K') = 0$；
\item $H^0(\alpha)$ 是同构，且 $H^{-1}(\alpha)$ 满射。
\end{enumerate}
若 $K$ 满足引理 \ref{more-algebra-lemma-ext-1-annihilated-definite} 的等价条件
(1)、(2) 和 (3)，则 $K'$ 也满足这些条件。
\end{lemma}

\begin{proof}
令 $M = \Ker(H^{-1}(\alpha))$。于是 $\alpha$ 嵌入区分三角
$$
M[1] \to K \to K' \to M[2]
$$
对任意 $R$-模 $N$，它确定正合序列
$$
\Ext^0_R(M[1], N) \to
\Ext^1_R(K', N) \to
\Ext^1_R(K, N)
$$
由于 $\Ext^0_R(M[1], N) = \Ext^{-1}_R(M, N) = 0$，可知
$\Ext^1_R(K', N)$ 是 $\Ext^1_R(K, N)$ 的子模。因此，若
$\Ext^1_R(K, N)$ 被 $I$ 零化，则 $\Ext^1_R(K', N)$ 也如此。
\end{proof}

\begin{lemma}
\label{more-algebra-lemma-ext-1-annihilated}
设 $R$ 为环，% P02-E0334
并设 $I \subset R$ 为理想。设 $K \in D(R)$，且对
$i \not \in \{-1, 0\}$ 有 $H^i(K) = 0$。下列条件等价：
\begin{enumerate}
\item 存在 $c \geq 0$，使得对 $K$ 和理想 $I^c$，
引理 \ref{more-algebra-lemma-ext-1-annihilated-definite} 的等价条件
(1)、(2)、(3) 成立；
\item 存在 $c \geq 0$，使得 (a) $I^c$ 零化 $H^{-1}(K)$，
且 (b) $H^0(K)$ 是 $I^c$-投射模（见第
\ref{more-algebra-section-near-projective} 节）。
\end{enumerate}
若 $R$ 是 Noether 环，并且对 $i = -1, 0$，$H^i(K)$ 是有限
$R$-模，则它们还等价于：
\begin{enumerate}
\item[(3)] 存在 $c \geq 0$，使得对 $K$ 和理想 $I^c$，
引理 \ref{more-algebra-lemma-ext-1-annihilated-definite} 的等价条件 (4)、(5) 成立；
\item[(4)] $H^{-1}(K)$ 是 $I$-幂挠的，并且存在
$f_1, \ldots, f_s \in R$，满足
$V(f_1, \ldots, f_s) \subset V(I)$，使得各局部化
$H^0(K)_{f_i}$ 是投射 $R_{f_i}$-模；
\item[(5)] $H^{-1}(K)$ 是 $I$-幂挠的，并且存在
$f_1, \ldots, f_s \in I$，满足
$V(f_1, \ldots, f_s) = V(I)$，使得各局部化
$H^0(K)_{f_i}$ 是投射 $R_{f_i}$-模；
\item[(6)] $H^{-1}(K)$ 是 $I$-幂挠的，并且对任意满足
$V(f_1, \ldots, f_s) = V(I)$ 的
$f_1, \ldots, f_s \in I$，各局部化 $H^0(K)_{f_i}$
都是投射 $R_{f_i}$-模。
\end{enumerate}
\end{lemma}

\begin{proof}
区分三角
$H^{-1}(K)[1] \to K \to H^0(K)[0] \to H^{-1}(K)[2]$
确定正合序列
$$
0 \to \Ext^1_R(H^0(K), N) \to \Ext^1_R(K, N) \to \Hom_R(H^{-1}(K), N) \to
\Ext^2_R(H^0(K), N)
$$
因此，(2) 蕴含对每个 $R$-模 $N$，$I^{2c}$ 零化
$\Ext^1_R(K, N)$。假设 (1)，立即可知 $H^0(K)$ 是
$I^c$-投射的。另一方面，可以选取某个内射 $R$-模 $N$
以及内射映射 $H^{-1}(K) \to N$。由于上述序列中的
$\Ext^2$ 消失，该映射是 $\Ext^1_R(K, N)$ 中某个元素的像，
故可知 $H^{-1}(K)$ 被 $I^c$ 零化。

\medskip\noindent
假设 $R$ 是 Noether 环，并且对 $i = -1, 0$，$H^i(K)$ 是有限
$R$-模。由引理 \ref{more-algebra-lemma-ext-1-annihilated-definite}，
可知 (3) 等价于 (1) 和 (2)。此外，若 (3) 成立，则对
$f \in I$，$f$ 在 $H^0(K)$ 上的乘法经由投射模分解；
% P02-E0335
这蕴含 $H^0(K)_f$ 是某个投射 $R_f$-模的直和项，因而自身也是投射
$R_f$-模。因此等价条件 (1)、(2) 和 (3) 蕴含 (6)。
当然，(6) 蕴含 (5)，而 (5) 显然蕴含 (4)。

\medskip\noindent
假设 (4)。由于 $H^{-1}(K)$ 是有限 $R$-模且为 $I$-幂挠，
可知对某个 $c_1 \geq 0$，$I^{c_1}$ 零化 $H^{-1}(K)$。
选取短正合序列
$$
0 \to M \to R^{\oplus r} \to H^0(K) \to 0
$$
它确定元素 $\xi \in \Ext^1_R(H^0(K), M)$。对任意 $f \in I$，
由引理 \ref{more-algebra-lemma-pseudo-coherence-and-base-change-ext}，有
$\Ext^1_R(H^0(K), M)_f = \Ext^1_{R_f}(H^0(K)_f, M_f)$。
因此，若 $H^0(K)_f$ 投射，则 $f$ 的某个幂零化 $\xi$。
故对某个 $c_2 \geq 0$，$(f_1, \ldots, f_s)^{c_2}$ 零化 $\xi$。
由于 $V(f_1, \ldots, f_s) \subset V(I)$，有
% P02-E0336
$\sqrt{I} \subset (f_1, \ldots, f_s)$
（Algebra，引理 \ref{algebra-lemma-Zariski-topology}）。
由于 $R$ 是 Noether 环，对某个 $c_3 \geq 0$，有
$I^{c_3} \subset (f_1, \ldots, f_s)$
（Algebra，引理 \ref{algebra-lemma-Noetherian-power}）。
% P02-E0337
因此 $I^{c2c3}$ 零化 $\xi$。这又说明 $H^0(K)$ 是
$I^{c_2c_3}$-投射的（因为对零化 $\xi$ 的 $a \in I$，
% P02-E0338
乘法经由 $R^{\oplus r}$ 分解）。因此取
$c = \max(c_1, c_2c_3)$，可知 (2) 成立。
\end{proof}

\begin{lemma}
\label{more-algebra-lemma-zero-in-derived}
设 $R$ 为环。对 $j = 1, 2, 3$，设 $K_j \in D(R)$，且对
$i \not \in \{-1, 0\}$ 有 $H^i(K_j) = 0$。设
$\varphi : K_1 \to K_2$ 和 $\psi : K_2 \to K_3$ 为
$D(R)$ 中的映射。若 $H^0(\varphi) = 0$ 且
$H^{-1}(\psi) = 0$，则
% P02-E0339
$\varphi \circ \psi = 0$。
\end{lemma}

\begin{proof}
应用 Derived Categories，引理
\ref{derived-lemma-trick-vanishing-composition}，可知
% P02-E0339
$\varphi \circ \psi$ 经由 $\tau_{\leq -2}K_2 = 0$ 分解。
\end{proof}

\begin{lemma}
\label{more-algebra-lemma-silly}
设 $R$ 为环。设 $K \in D(R)$ 由形如
$R^{\oplus n} \to R^{\oplus n}$ 的二项复形给出。
% P02-E0340
用 $A \in \text{Mat}(n \times n, R)$ 表示微分的矩阵。
% P02-E0341
则 $\det(a) : K \to K$ 在 $D(R)$ 中为零。
\end{lemma}

\begin{proof}
省略。这是一个很好的练习。
\end{proof}

\section{朴素余切复形}
\label{more-algebra-section-netherlander}

\noindent
本节继续 Algebra，第 \ref{algebra-section-netherlander} 节中开始的讨论。
先讨论基变换。第一个引理表明，把朴素余切复形与环扩张作朴素张量积，
并不像人们可能以为的那样朴素。

\begin{lemma}
\label{more-algebra-lemma-tensor-NL}
设 $R \to S$ 和 $S \to S'$ 为环同态。典范映射
$\NL_{S/R} \otimes_S^\mathbf{L} S' \to \NL_{S/R} \otimes_S S'$
在 $D(S')$ 中诱导同构
$\tau_{\geq -1}(\NL_{S/R} \otimes_S^\mathbf{L} S') \to \NL_{S/R} \otimes_S S'$。
类似地，给定 $S$ 在 $R$ 上的一个表示 $\alpha$，典范映射
$\NL(\alpha) \otimes_S^\mathbf{L} S' \to \NL(\alpha) \otimes_S S'$
在 $D(S')$ 中诱导同构
$\tau_{\geq -1}(\NL(\alpha) \otimes_S^\mathbf{L} S') \to
\NL(\alpha) \otimes_S S'$。
\end{lemma}

\begin{proof}
这是引理 \ref{more-algebra-lemma-two-term-base-change} 的特例。% P02-E0342
\end{proof}

\begin{lemma}
\label{more-algebra-lemma-base-change-NL}
设 $R \to S$ 和 $R \to R'$ 为环同态。设
$\alpha : P \to S$ 为 $S$ 在 $R$ 上的一个表示。则
$\alpha' : P \otimes_R R' \to S \otimes_R R'$ 是
$S' = S \otimes_R R'$ 在 $R'$ 上的一个表示。典范映射
% P02-E0343
$$
NL(\alpha) \otimes_S S' \to \NL(\alpha')
$$
在 $H^0$ 上是同构，在 $H^{-1}$ 上是满射。特别地，典范映射
$$
\NL_{S/R} \otimes_S S' \to \NL_{S'/R'}
$$
在 $H^0$ 上是同构，在 $H^{-1}$ 上是满射。
\end{lemma}

\begin{proof}
% P02-E0344
记 $I = \Ker(P \to S)$。记 $P' = P \otimes_R R'$，并记
$I' = \Ker(P' \to S')$。假设 $P$ 是以 $j \in J$ 为指标的
$x_j$ 上的多项式代数。引理中显示的映射变为
$$
\xymatrix{
\bigoplus_{j \in J} S' \text{d}x_j \ar[r] &
\bigoplus_{j \in J} S' \text{d}x_j \\
I/I^2 \otimes_S S' \ar[r] \ar[u] &
I'/(I')^2 \ar[u]
}
$$
其中左列为 $\NL(\alpha) \otimes_S S'$，右列为
$\NL(\alpha')$。由张量积的右正合性，可知
$I \otimes_R R' \to I'$ 满射。因此底部箭头满射。
这证明了引理的第一个陈述。关于
$\NL_{S/R} \otimes_S S' \to \NL_{S'/R'}$ 的陈述随之成立，
因为这些复形分别同伦于 $\NL(\alpha) \otimes_S S'$ 和
$\NL(\alpha')$。
\end{proof}

\begin{lemma}
\label{more-algebra-lemma-base-change-NL-flat}
考虑环的余笛卡尔图
$$
\xymatrix{
B \ar[r] & B' \\
A \ar[r] \ar[u] & A' \ar[u]
}
$$
若 $B$ 在 $A$ 上平坦，则典范映射
$\NL_{B/A} \otimes_B B' \to \NL_{B'/A'}$ 是拟同构。
若还假设 $\NL_{B/A}$ 的 Tor 振幅位于 $[-1, 0]$ 中，则
$\NL_{B/A} \otimes_B^\mathbf{L} B' \to \NL_{B'/A'}$
也是拟同构。
\end{lemma}

\begin{proof}
按 Algebra，第 \ref{algebra-section-netherlander} 节选取表示
$\alpha : P \to B$。令 $I = \Ker(\alpha)$，置
$P' = P \otimes_A A'$，并用
% P02-E0345
$\alpha' : P' \to B'$ 表示 $B'$ 在 $A'$ 上的相应表示。
由于 $B$ 在 $A$ 上平坦，可知 $I' = \Ker(\alpha')$ 等于
$I \otimes_A A'$。因此
$$
I'/(I')^2 = \Coker(I^2 \otimes_A A' \to I \otimes_A A') =
I/I^2 \otimes_A A' = I/I^2 \otimes_B B'
$$
由于两边具有相同的基，有
$\Omega_{P'/A'} = \Omega_{P/A} \otimes_A A'$。于是
$\Omega_{P'/A'} \otimes_{P'} B' = \Omega_{P/A} \otimes_P B \otimes_B B'$。
这证明 $\NL(\alpha) \otimes_B B' \to \NL(\alpha')$ 是复形同构，
从而第一个陈述成立。

\medskip\noindent
作为 $B$-模复形，有
$$
\NL(\alpha) = I/I^2 \longrightarrow \Omega_{P/A} \otimes_P B
$$
其中 $I/I^2$ 置于次数 $-1$。由于次数 $0$ 的项自由，
由引理 \ref{more-algebra-lemma-last-one-flat}，该复形的 Tor 振幅位于
$[-1, 0]$ 中，当且仅当 $I/I^2$ 是平坦 $B$-模。
若此条件成立，则
$\NL(\alpha) \otimes_B^\mathbf{L} B' =
\NL(\alpha) \otimes_B B'$，因而得到第二个陈述。
\end{proof}

\begin{lemma}
\label{more-algebra-lemma-lci-NL}
设 $A \to B$ 是定义 \ref{more-algebra-definition-local-complete-intersection}
意义下的局部完全交。则 $\NL_{B/A}$ 是 $D(B)$ 中 Tor 振幅位于
$[-1, 0]$ 的完美对象。
\end{lemma}

\begin{proof}
写成 $B = A[x_1, \ldots, x_n]/I$。则 $\NL_{B/A}$ 由
$B$-模复形
$$
I/I^2 \longrightarrow \bigoplus B \text{d}x_i
$$
表示，其中 $I/I^2$ 置于次数 $-1$。由于次数 $0$ 的项有限自由，
由引理 \ref{more-algebra-lemma-last-one-flat}，该复形的 Tor 振幅位于
$[-1, 0]$ 中，当且仅当 $I/I^2$ 是平坦 $B$-模。
按照定义，$I$ 是 Koszul 正则理想，因而是拟正则理想；见第
\ref{more-algebra-section-ideals} 节。因此 $I/I^2$ 是有限投射 $B$-模
（引理 \ref{more-algebra-lemma-quasi-regular-ideal-finite-projective}），
从而 $\NL_{B/A}$ 既是完美的，又具有 $[-1, 0]$ 内的 Tor 振幅。
\end{proof}

\begin{lemma}
\label{more-algebra-lemma-base-change-lci-bis}
考虑环的余笛卡尔图
$$
\xymatrix{
B \ar[r] & B' \\
A \ar[r] \ar[u] & A' \ar[u]
}
$$
若 $A \to B$ 与 $A' \to B'$ 都是定义
\ref{more-algebra-definition-local-complete-intersection} 意义下的局部完全交，
则映射
$H^{-1}(\NL_{B/A} \otimes_B B') \to H^{-1}(\NL_{B'/A'})$
的核是有限投射 $B'$-模。
\end{lemma}

\begin{proof}
由引理 \ref{more-algebra-lemma-lci-NL}，复形 $\NL_{B/A}$ 和
$\NL_{B'/A'}$ 都是 Tor 振幅位于 $[-1, 0]$ 的完美复形。
结合引理 \ref{more-algebra-lemma-tensor-NL}、\ref{more-algebra-lemma-pull-perfect}
和 \ref{more-algebra-lemma-pull-tor-amplitude}，有
$\NL_{B/A} \otimes_B B' = \NL_{B/A} \otimes_B^\mathbf{L} B'$，
并且该复形也是 Tor 振幅位于 $[-1, 0]$ 的完美复形。
在 $D(B')$ 中选取区分三角
$$
C \to \NL_{B/A} \otimes_B B'  \to \NL_{B'/A'} \to C[1]
$$
由引理 \ref{more-algebra-lemma-two-out-of-three-perfect} 和
\ref{more-algebra-lemma-cone-tor-amplitude}，可知 $C$ 是 Tor 振幅位于
$[-1, 1]$ 的完美复形。由引理 \ref{more-algebra-lemma-base-change-NL}，
复形 $C$ 仅有一个非零上同调模，即引理中的模，位于次数 $-1$。
该模有限表示（引理 \ref{more-algebra-lemma-n-pseudo-module}）且平坦
（引理 \ref{more-algebra-lemma-tor-dimension}）。因此，由 Algebra，引理
\ref{algebra-lemma-finite-projective}，它是有限投射的。
\end{proof}

\section{Abel 群的 Rlim}
\label{more-algebra-section-Rlim}

\noindent
我们简要讨论 Abel 群上的 $R\lim$。本节用
% P02-E0346
$\textit{Ab}(\mathbf{N})$ 表示 Abel 群逆系组成的 Abel 范畴。
这一记号与位点上 Abel 群层的记号相容，因为 Abel 群逆系等同于范畴
$\mathbf{N}$ 上的群层（当 $i \leq j$ 时，恰有一个态射
$i \to j$）；见评注 \ref{more-algebra-remark-Rlim-cohomology}。
本节的许多论证重复了构造位点上 Abel 群层的上同调机制时所用的论证。

\begin{lemma}
\label{more-algebra-lemma-compute-Rlim}
函子 $\lim : \textit{Ab}(\mathbf{N}) \to \textit{Ab}$ 具有右导出函子
\begin{equation}
\label{more-algebra-equation-Rlim}
R\lim : D(\textit{Ab}(\mathbf{N})) \longrightarrow D(\textit{Ab})
\end{equation}
按惯例，置 $R^p\lim(K) = H^p(R\lim(K))$。此外：
\begin{enumerate}
\item 对 $\textit{Ab}(\mathbf{N})$ 中任意 $(A_n)$，当
$p > 1$ 时有 $R^p\lim A_n = 0$；
\item $D(\textit{Ab})$ 的对象 $R\lim A_n$ 由复形
$$
\prod A_n \to \prod A_n,\quad (x_n) \mapsto (x_n - f_{n + 1}(x_{n + 1}))
$$
表示，该复形位于次数 $0$ 和 $1$；
\item 若 $(A_n)$ 是 ML 的，则 $R^1\lim A_n = 0$，即
$(A_n)$ 对 $\lim$ 是右非周期的；
\item 每个 $K^\bullet \in D(\textit{Ab}(\mathbf{N}))$
都拟同构于一个逐项对 $\lim$ 右非周期的复形；
\item 若每个 $K^p = (K^p_n)$ 都对 $\lim$ 右非周期，即
% P02-E0347
$R^1\lim_n K^p_n = 0$，则 $R\lim K$ 由次数 $p$ 的项为
$\lim_n K_n^p$ 的复形表示。
\end{enumerate}
\end{lemma}

\begin{proof}
设 $(A_n)$ 为任意逆系。定义逆系 $(B_n)$，其中
$$
B_n = A_n \oplus A_{n - 1} \oplus \ldots \oplus A_1
$$
转移映射由投影给出。映射 $A_n \to B_n$ 由
$(1, f_n, f_{n - 1} \circ f_n, \ldots, f_2 \circ \ldots \circ f_n)$
给出，其中 $f_i : A_i \to A_{i - 1}$ 是转移映射。
由此可见，每个逆系都是某个 ML 系统的子对象
（Homology，第 \ref{homology-section-inverse-systems} 节）。
由 Derived Categories，引理
\ref{derived-lemma-subcategory-right-acyclics} 和 Homology，引理
\ref{homology-lemma-Mittag-Leffler}，可知每个 ML 系统都对
$\lim$ 右非周期，即 (3) 成立。这已经蕴含
% P02-E0348
$RF$ 定义在 $D^+(\textit{Ab}(\mathbf{N}))$ 上；见
Derived Categories，命题 \ref{derived-proposition-enough-acyclics}。
对 $n > 1$，置
$C_n = A_{n - 1} \oplus \ldots \oplus A_1$，并置 $C_1 = 0$，
转移映射同样由投影给出。于是有逆系的短正合序列
$0 \to (A_n) \to (B_n) \to (C_n) \to 0$，其中
$B_n \to C_n$ 由
$(x_i) \mapsto (x_i - f_{i + 1}(x_{i + 1}))$ 给出。
由于 $(C_n)$ 也是 ML 的，由上面引用的命题可知 (2) 成立，
而这也蕴含 (1)。最后，由 Derived Categories，引理
\ref{derived-lemma-unbounded-right-derived}，可知 $R\lim$
事实上定义在整个 $D(\textit{Ab}(\mathbf{N}))$ 上。
事实上，Derived Categories，引理
\ref{derived-lemma-unbounded-right-derived} 的证明正是通过证明
(4) 和 (5) 来进行的。
\end{proof}

\begin{lemma}
\label{more-algebra-lemma-six-term-Rlim}
设
$$
0 \to (A_i) \to (B_i) \to (C_i) \to 0
$$
为 Abel 群逆系的短正合序列。则有相伴的 $6$ 项正合序列
$0 \to \lim A_i \to \lim B_i \to \lim C_i \to
R^1\lim A_i \to R^1\lim B_i \to R^1\lim C_i \to 0$。
\end{lemma}

\begin{proof}
这由引理 \ref{more-algebra-lemma-compute-Rlim} 中的消失性得到。% P02-E0349
\end{proof}

\noindent
下面是 Homology，引理
\ref{homology-lemma-apply-Mittag-Leffler-again} 的“正确”表述。

\begin{lemma}
\label{more-algebra-lemma-apply-Mittag-Leffler-again}
设
$$
(A^{-2}_n \to A^{-1}_n \to A^0_n \to A^1_n)
$$
为 Abel 群复形的逆系，并用
% P02-E0350
$A^{-2} \to A^{-1} \to A^0 \to A^1$ 表示其极限。
用 $(H_n^{-1})$、$(H_n^0)$ 表示上同调的逆系，并用
$H^{-1}$、$H^0$ 表示复形
$A^{-2} \to A^{-1} \to A^0 \to A^1$ 的上同调。若：
\begin{enumerate}
\item $(A^{-2}_n)$ 和 $(A^{-1}_n)$ 的 $R^1\lim$ 均消失；
\item $(H^{-1}_n)$ 的 $R^1\lim$ 消失；
\end{enumerate}
则 $H^0 = \lim H_n^0$。
\end{lemma}

\begin{proof}
设 $K \in D(\textit{Ab}(\mathbf{N}))$ 为由下述复形系表示的对象：
其第 $n$ 个分量是复形
$A^{-2}_n \to A^{-1}_n \to A^0_n \to A^1_n$。
我们将使用 Derived Categories，引理
\ref{derived-lemma-two-ss-complex-functor} 的两个谱序列来计算
$H^0(R\lim K)$\footnote{为使用这些谱序列，必须证明
$\textit{Ab}(\mathbf{N})$ 有足够多的内射对象。% P02-E0351
Abel 群逆系 $(I_n)$ 是内射的，当且仅当每个 $I_n$
都是内射 Abel 群，并且转移映射是分裂满射。每个系统都嵌入某个
这样的系统中。细节省略。}。第一个谱序列的 $E_1$-页为
$$
\begin{matrix}
0 & 0 & R^1\lim A^0_n  & R^1\lim A^1_n \\
A^{-2} & A^{-1} & A^0  & A^1
\end{matrix}
$$
其微分是横向的，且所有高阶微分均为零。% P02-E0352
第二个谱序列的 $E_2$ 页为
$$
\begin{matrix}
R^1\lim H^{-2}_n & 0 & R^1\lim H^0_n & R^1 \lim H^1_n \\
\lim H^{-2}_n &
\lim H^{-1}_n &
\lim H^0_n &
\lim H^1_n
\end{matrix}
$$
并在此页退化。结论随即得到。
\end{proof}

\begin{lemma}
\label{more-algebra-lemma-pro-isomorphism-iso-Rlim}
设 $(A_n)$ 和 $(B_n)$ 为 Abel 群逆系。pro-系的态射
$\varphi : (A_n) \to (B_n)$ 确定映射
$\lim A_n \to \lim B_n$ 和 $R^1\lim A_n \to R^1\lim B_n$。
若 $\varphi$ 是 pro-同构，则这些映射都是同构。
\end{lemma}

\begin{proof}
关于 pro-系态射的讨论，请见 Categories，例
\ref{categories-example-pro-morphism-inverse-systems}。
映射 $\varphi$ 由 $1 \leq m_1 < m_2 < m_3 < \ldots$ 以及与转移映射
相容的映射 $\varphi_n : A_{m_n} \to B_n$ 给出。置
$\varphi', \varphi'' : \prod A_n \to \prod B_n$ 分别等于
$\varphi'((a_n)) = (\varphi_n(a_{m_n}))$ 以及
$$
\varphi''((a_n)) =
(\varphi_n(a_{m_n} + f(a_{m_n + 1}) + \ldots + f(a_{m_{n + 1} - 1})))
$$
其中每次出现的 $f$ 都表示逆系 $(A_n)$ 的一个适当转移映射。
% P02-E0353
于是交换图
$$
\xymatrix{
\prod A_n \ar[r]_\delta \ar[d]^{\varphi'} &
\prod A_n \ar[d]^{\varphi''} \\
\prod B_n \ar[r]^\delta & \prod B_n
}
$$
中的横向箭头如引理 \ref{more-algebra-lemma-compute-Rlim} 所述。
由此得到所需的映射。该构造具有如下函子性：如果给定逆系
$(C_n)$、整数 $1 \leq m'_1 < m'_2 < m'_3 < \ldots$，
以及与转移映射相容的映射 $\psi_n : B_{m'_n} \to C_n$，
% P02-E0354
则
$(\psi \circ \varphi)' = \psi' \circ \varphi'$ 且
$(\psi \circ \varphi)'' = \psi'' \circ \varphi''$，
其中复合 $\psi \circ \varphi$ 指整数
$1 \leq m_{m'_1} < m_{m'_2} < \ldots$ 以及映射
$\psi_n \circ \varphi_{m'_n} : A_{m_{m'_n}} \to C_n$。
下面证明：若 $B_n = A_n$ 且 $\varphi_n$ 是转移映射，
则所得映射为恒等映射。这将同时证明该构造与
$\varphi$ 的代表元 $(m_n, \varphi_n)$ 的选择无关，并证明引理的
最后一个陈述。

\medskip\noindent
因此，令 $B_n = A_n$，并令 $\varphi_n$ 为转移映射。
设 $(a_n) \in \lim A_n$ 为一个元素。则
$\varphi'((a_n))$ 在次数 $n$ 的分量是 $a_{m_n}$ 的像，
而它等于 $a_n$。这证明了关于 $\lim A_n$ 的陈述。
设 $\xi \in R^1\lim A_n$ 是 $\prod A_n$ 中元素 $(a_n)$ 的类。
考虑元素 $(b_n)$，其中
$$
b_i = a_i + f(a_{i + 1}) + \ldots + f(a_{m_n - 1})
$$
若 $m_{n - 1} < i < m_n$，并且若对某个 $n$ 有 $i = m_n$，
则取 $0$。于是
$(a'_n) = (a_n) - \delta((b_n))$ 定义 $R^1\lim A_n$ 中相同的类，
且计算表明，除非对某个 $n$ 有 $i = m_n$，否则 $a'_i = 0$。
因此可以并且确实假设，除非对某个 $n$ 有 $i = m_n$，否则
$a_i = 0$。注意，
$$
(0, \ldots, -f(a_{m_j}), 0, \ldots, 0, a_{m_j}, 0, \ldots)
$$
在位置 $j$ 和 $m_j$ 上非零，它是
$$
c_j = (0, \ldots, 0, f(a_{m_j}), f(a_{m_j}), \ldots, a_{m_j}, 0, \ldots)
$$
在 $\delta$ 下的像，其中非零项位于位置
$j + 1, \ldots, m_j$。和 $c = \sum c_j$ 在 $\prod A_n$ 中有意义。
回忆 $\varphi''((a_n)) = (a_{m_1}, a_{m_2}, \ldots)$，可知
$$
(a_n) - \varphi''((a_n)) = \delta(c)
$$
证明完成。
\end{proof}

\begin{lemma}
\label{more-algebra-lemma-map-into-Rlim}
设 $\mathcal{D}$ 为三角范畴，$(K_n)$ 为 $\mathcal{D}$ 中对象的逆系，
并设 $K$ 为系统 $(K_n)$ 的一个导出极限。则对
$\mathcal{D}$ 中每个 $L$，有短正合序列
$$
0 \to R^1\lim \Hom_\mathcal{D}(L, K_n[-1]) \to
\Hom_\mathcal{D}(L, K) \to
\lim \Hom_\mathcal{D}(L, K_n) \to 0
$$
\end{lemma}

\begin{proof}
这由 Derived Categories，定义
\ref{derived-definition-derived-limit}、引理
\ref{derived-lemma-representable-homological}，以及上面引理
\ref{more-algebra-lemma-compute-Rlim} 中对 $\lim$ 和 $R^1\lim$ 的描述得到。
\end{proof}

\begin{lemma}
\label{more-algebra-lemma-pro-isomorphism-Rlim}
设 $\mathcal{D}$ 为三角范畴，$(K_n)$ 与 $(M_n)$ 为
$\mathcal{D}$ 中对象的逆系，其导出极限分别为 $K$ 和 $M$。
设 $a : (K_n) \to (M_n)$ 为 pro-对象的 pro-同构。
则可以用 $a$ 构造一个（非典范）同构 $K \to M$。
\end{lemma}

\begin{proof}
由 Derived Categories，评注
\ref{derived-remark-functoriality-derived-limit}，得到一个嵌入定义区分三角之态射
$$
\xymatrix{
K \ar[r] \ar[d] & \prod K_n \ar[r] \ar[d]^{a'} &
\prod K_n \ar[d]^{a''} \\
M \ar[r] & \prod M_n \ar[r] & \prod M_n
}
$$
的箭头 $K \to M$。因此，对 $\mathcal{D}$ 的每个对象 $L$，
得到短正合序列之间的映射
$$
\xymatrix{
0 \ar[r] &
R^1\lim \Hom_\mathcal{D}(L, K_n[-1]) \ar[r] \ar[d] &
\Hom_\mathcal{D}(L, K) \ar[r] \ar[d] &
\lim \Hom_\mathcal{D}(L, K_n) \ar[r] \ar[d] &
0 \\
0 \ar[r] &
R^1\lim \Hom_\mathcal{D}(L, M_n[-1]) \ar[r] &
\Hom_\mathcal{D}(L, M) \ar[r] &
\lim \Hom_\mathcal{D}(L, M_n) \ar[r] &
0
}
$$
参见引理 \ref{more-algebra-lemma-map-into-Rlim} 及其证明。
由引理 \ref{more-algebra-lemma-pro-isomorphism-iso-Rlim}，左、右两个竖直箭头
都是同构\footnote{若 pro-同构 $a$ 由映射
$a_n : K_{m_n} \to M_n$ 给出，则使用相应的映射
$\Hom_\mathcal{D}(L, K_{m_n}) \to \Hom_\mathcal{D}(L, M_n)$，
% P02-E0355
得到 $\Hom$ 上的 pro-同构。因此，引理
\ref{more-algebra-lemma-pro-isomorphism-iso-Rlim} 的证明中构造
$\lim$ 与 $R^1\lim$ 上箭头的方法，与 Derived Categories，评注
\ref{derived-remark-functoriality-derived-limit} 中构造
$a'$ 和 $a''$ 的方法一致。}。因此中间箭头是同构。
由 Yoneda 引理可知，映射 $K \to M$ 是同构。
\end{proof}

\begin{lemma}
\label{more-algebra-lemma-map-from-hocolim}
设 $\mathcal{D}$ 为三角范畴，$(K_n)$ 为 $\mathcal{D}$ 中的对象系，
并设 $K$ 为系统 $(K_n)$ 的一个导出余极限。则对
$\mathcal{D}$ 中每个 $L$，有短正合序列
$$
0 \to R^1\lim \Hom_\mathcal{D}(K_n, L[-1]) \to
\Hom_\mathcal{D}(K, L) \to
\lim \Hom_\mathcal{D}(K_n, L) \to 0
$$
\end{lemma}

\begin{proof}
这由 Derived Categories，定义
\ref{derived-definition-derived-colimit}、引理
\ref{derived-lemma-representable-homological}，以及上面引理
\ref{more-algebra-lemma-compute-Rlim} 中对 $\lim$ 和 $R^1\lim$ 的描述得到。
\end{proof}

\begin{remark}[作为上同调的 Rlim]
\label{more-algebra-remark-Rlim-cohomology}
考虑对象为自然数的范畴 $\mathbf{N}$；若 $j \geq i$，
则从 $i \to j$ 恰有一个箭头。赋予 $\mathbf{N}$ 混沌拓扑
（Sites，例 \ref{sites-example-indiscrete}），于是层
$\mathcal{F}$ 等同于 $\mathbf{N}$ 上的集合逆系
$$
\mathcal{F}_1 \leftarrow \mathcal{F}_2 \leftarrow \mathcal{F}_3
\leftarrow \ldots
$$
注意 $\Gamma(\mathbf{N}, \mathcal{F}) = \lim \mathcal{F}_n$。
对 Abel 群逆系 $\mathcal{F}_n$，有
$$
R^p\lim \mathcal{F}_n = H^p(\mathbf{N}, \mathcal{F})
$$
因为两边都是函子
$\mathcal{F} \mapsto \lim \mathcal{F}_n = H^0(\mathbf{N}, \mathcal{F})$
的高阶右导出函子。因此，$R\lim$ 的存在性也可由
Cohomology on Sites，第
\ref{sites-cohomology-section-cohomology-sheaves} 节和第
\ref{sites-cohomology-section-unbounded} 节的一般理论得到。
\end{remark}

\noindent
下一个引理中的积既可以看作复形的逐项积，也可以看作导出范畴
$D(\textit{Ab})$ 中的积；见 Derived Categories，引理
\ref{derived-lemma-products}。

\begin{lemma}
\label{more-algebra-lemma-distinguished-triangle-Rlim}
设 $K = (K_n^\bullet)$ 为 $D(\textit{Ab}(\mathbf{N}))$ 的对象。
在 $D(\textit{Ab})$ 中存在典范区分三角
$$
R\lim K \to \prod\nolimits_n K_n^\bullet \to \prod\nolimits_n K_n^\bullet
\to R\lim K[1]
$$
换言之，$R\lim K$ 是 $D(\textit{Ab})$ 中逆系
$(K_n^\bullet)$ 的一个导出极限；见 Derived Categories，定义
\ref{derived-definition-derived-limit}。
\end{lemma}

\begin{proof}
假设对每个 $p$，逆系 $(K_n^p)$ 都对 $\lim$ 右非周期。
由引理 \ref{more-algebra-lemma-compute-Rlim}，对每个 $p$ 都得到短正合序列
$$
0 \to \lim_n K^p_n \to \prod\nolimits_n K^p_n \to \prod\nolimits_n K^p_n \to 0
$$
由于由 $\lim_n K^p_n$ 组成的复形按引理
\ref{more-algebra-lemma-compute-Rlim} 计算 $R\lim K$，故该引理在此情形成立。

\medskip\noindent
接着，设 $K = (K_n^\bullet)$ 为一般对象。由引理
\ref{more-algebra-lemma-compute-Rlim}，在 $D(\textit{Ab}(\mathbf{N}))$
中存在拟同构 $K \to L$，使得对每个 $p$，$(L_n^p)$ 非周期。
由于 $\textit{Ab}$ 中积是正合的，$\prod K_n^\bullet$
拟同构于 $\prod L_n^\bullet$；因而上面已证明的关于 $L$ 的结论
蕴含关于 $K$ 的结论。
\end{proof}

\begin{lemma}
\label{more-algebra-lemma-break-long-exact-sequence}
沿用引理 \ref{more-algebra-lemma-distinguished-triangle-Rlim} 的记号，
与该区分三角相伴的上同调长正合序列分解为短正合序列
$$
0 \to R^1\lim_n H^{p - 1}(K_n^\bullet) \to
H^p(R\lim K) \to
\lim_n H^p(K_n^\bullet) \to 0
$$
\end{lemma}

\begin{proof}
该区分三角的长正合序列为
$$
\ldots \to H^p(R\lim K) \to \prod\nolimits_n H^p(K_n^\bullet)
\to \prod\nolimits_n H^p(K_n^\bullet) \to
H^{p + 1}(R\lim K) \to \ldots
$$
根据引理中给出的显式描述，中间映射的核是
$\lim_n H^p(K_n^\bullet)$。由引理
\ref{more-algebra-lemma-compute-Rlim}，该映射的余核是
$R^1\lim_n H^p(K_n^\bullet)$。
\end{proof}

\noindent
{\bf 警告。}$D(\textit{Ab}(\mathbf{N}))$ 的对象是 Abel 群逆系的复形。
也可以把它看作复形的逆系 $(K_n^\bullet)$。然而，这与
$D(\textit{Ab})$ 的对象之逆系{\bf 并不}相同；下一个引理与评注说明二者的差别。

\begin{lemma}
\label{more-algebra-lemma-lift-to-system-complexes-Ab}
设 $(K_n)$ 为 $D(\textit{Ab})$ 的对象之逆系。则存在
$D(\textit{Ab}(\mathbf{N}))$ 的对象 $M = (M_n^\bullet)$，以及
$D(\textit{Ab})$ 中的同构 $M_n^\bullet \to K_n$，使得图
$$
\xymatrix{
M_{n + 1}^\bullet \ar[d] \ar[r] &
M_n^\bullet \ar[d] \\
K_{n + 1} \ar[r] & K_n
}
$$
在 $D(\textit{Ab})$ 中交换。
\end{lemma}

\begin{proof}
具体地，设 $M_1^\bullet$ 为表示 $K_1$ 的 Abel 群复形。
假设已经构造
$M_e^\bullet \to M_{e - 1}^\bullet \to \ldots \to M_1^\bullet$
以及映射 $\psi_i : M_i^\bullet \to K_i$，使得对所有
$n < e$，引理陈述中的图都交换。然后考虑
% P02-E0356
$$
\xymatrix{
& M_n^\bullet \ar[d]^{\psi_n} \\
K_{n + 1} \ar[r] & K_n
}
$$
作为 $D(\textit{Ab})$ 中的图。根据 $D(\textit{Ab})$ 中态射的定义，
可以找到 Abel 群复形 $M_{n + 1}^\bullet$、
$D(\textit{Ab})$ 中的同构 $M_{n + 1}^\bullet \to K_{n + 1}$，
以及复形态射 $M_{n + 1}^\bullet \to M_n^\bullet$；它在
$D(\textit{Ab})$ 中表示复合
$$
K_{n + 1} \to K_n \xrightarrow{\psi_n^{-1}} M_n^\bullet
$$
。
故由归纳法，引理成立。
\end{proof}

\begin{remark}
\label{more-algebra-remark-compare-derived-limit}
设 $(K_n)$ 为 $D(\textit{Ab})$ 的对象之逆系，并设
$K = R\lim K_n$ 为该系统的一个导出极限（见 Derived Categories，
第 \ref{derived-section-derived-limit} 节）。这样的导出极限存在，
因为 $D(\textit{Ab})$ 具有可数积
（Derived Categories，引理 \ref{derived-lemma-products}）。
由引理 \ref{more-algebra-lemma-lift-to-system-complexes-Ab}，也可以把
$(K_n)$ 提升为 $D(\textit{Ab}(\mathbf{N}))$ 的对象 $M$。
于是 $K \cong R\lim M$，其中 $R\lim$ 是函子
(\ref{more-algebra-equation-Rlim})；这是因为由引理
\ref{more-algebra-lemma-distinguished-triangle-Rlim}，$R\lim M$ 也是系统
$(K_n)$ 的导出极限。因此，尽管系统 $(K_n)$ 的提升 $M$
可能有许多同构类，$R\lim M$ 的同构型却与选择无关，
因为它同构于该系统的导出极限 $K = R\lim K_n$。
因此，可以把本节证明的关于 $R\lim$ 的结果应用于导出极限。
例如，对每个 $p \in \mathbf{Z}$，都有典范短正合序列
$$
0 \to R^1\lim H^{p - 1}(K_n) \to H^p(K) \to \lim H^p(K_n) \to 0
$$
因为可以对 $M$ 应用引理 \ref{more-algebra-lemma-break-long-exact-sequence}。
% P02-E0357
不借助 $M$ 的存在性，也可以直接看出这一点：把引理
\ref{more-algebra-lemma-break-long-exact-sequence} 证明中的论证应用于（定义性的）
区分三角 $K \to \prod K_n \to \prod K_n \to K[1]$ 即可。
\end{remark}

\begin{lemma}
\label{more-algebra-lemma-Rlim-pro-equal}
设 $E \to D$ 为 $D(\textit{Ab}(\mathbf{N}))$ 的态射。
设 $(E_n)$、相应地 $(D_n)$ 为分别与 $E$、相应地 $D$ 关联的
$D(\textit{Ab})$ 对象之系统。若
$(E_n) \to (D_n)$ 是 pro-对象的同构，则
$R\lim E \to R\lim D$ 是 $D(\textit{Ab})$ 中的同构。
\end{lemma}

\begin{proof}
该假设特别蕴含 pro-对象 $H^p(E_n)$ 与 $H^p(D_n)$ 同构。
结论由引理 \ref{more-algebra-lemma-break-long-exact-sequence} 的短正合序列
以及引理 \ref{more-algebra-lemma-pro-isomorphism-iso-Rlim} 得到。
\end{proof}

\begin{lemma}[Emmanouil]
\label{more-algebra-lemma-emmanouil}
\begin{reference}
取自 \cite{Emmanouil}。
\end{reference}
设 $(A_n)$ 为 Abel 群逆系。下列条件等价：
\begin{enumerate}
\item $(A_n)$ 是 Mittag-Leffler 的；
\item $R^1\lim A_n = 0$，并且
$\bigoplus_{i \in \mathbf{N}} (A_n)$ 也满足相同结论。
\end{enumerate}
\end{lemma}

\begin{proof}
置 $B = \bigoplus_{i \in \mathbf{N}} (A_n)$，因而
$B = (B_n)$，其中 $B_n = \bigoplus_{i \in \mathbf{N}} A_n$。
若 $(A_n)$ 是 ML 的，则 $B$ 是 ML 的，故由引理
\ref{more-algebra-lemma-compute-Rlim}，有 $R^1\lim A_n = 0$ 且
$R^1\lim B_n = 0$。

\medskip\noindent
反之，假设 $(A_n)$ 不是 ML 的。于是可以选取 $m$、整数列
$m < m_1 < m_2 < \ldots$，以及元素 $x_i \in A_{m_i}$，
使得其像 $y_i \in A_m$ 不在 $A_{m_i + 1} \to A_m$ 的像中。
我们将用元素 $x_i$ 和 $y_i$ 以两种方式证明
$R^1\lim B_n \not = 0$。这将完成引理的证明。

\medskip\noindent
第一种证明。
置 $C = (C_n)$，其中 $C_n = \prod_{i \in \mathbf{N}} A_n$。
有典范单射 $B_n \to C_n$，其余核为 $Q_n$。置 $Q = (Q_n)$。
可以并且确实把 $Q_n$ 的元素 $q_n$ 看作元素序列
$q_n = (q_{n, 1}, q_{n, 2}, \ldots)$，其中 $q_{n, i} \in A_n$，
并对尾部为零的序列取模（换言之，差仅在有限个位置的序列被等同）。
有逆系的短正合序列
$$
0 \to (B_n) \to (C_n) \to (Q_n) \to 0
$$
考虑由下式给出的元素 $q_n \in Q_n$：
$$
q_{n, i} =
\left\{
\begin{matrix}
\text{image of }x_i &\text{若}& m_i \geq n \\
0 & \text{否则}
\end{matrix}
\right.
$$
显然，$q_{n + 1}$ 映到 $q_n$。因此得到
$q = (q_n) \in \lim Q_n$。另一方面，我们断言 $q$ 不在
$\lim C_n \to \lim Q_n$ 的像中。事实上，假设 $c = (c_n)$
映到 $q$。可以写成 $c_n = (c_{n, i})$；由于对
$n' \geq n$ 有 $c_{n', i} \mapsto c_{n, i}$，可知对所有
$n, i, n' \geq n$，有
$c_{n, i} \in \Im(C_{n'} \to C_n)$。特别地，
% P02-E0359
$c_{m, i}$ 在 $A_m$ 中的像属于
$\Im(A_{m_i + 1} \to A_m)$，故不可能等于 $y_i$。
因此 $c_m$ 与 $q_m = (y_1, y_2, y_3, \ldots)$ 在无穷多个位置上不同，
这是矛盾。考察上同调长正合序列
$$
0 \to \lim B_n \to \lim C_n \to \lim Q_n \to R^1\lim B_n
$$
可知最后一个群非零，如所需。

\medskip\noindent
第二种证明。对 $n' \geq n$，用
% P02-E0358
$A_{n, n'} = \Im(A_{n'} \to A_n)$ 表示其像。于是
$y_i \in A_m$，且 $y_i \not \in A_{m, m_i + 1}$。
设 $\xi = (\xi_n) \in \prod B_n$ 为如下元素：
除非 $n = m_i$，否则 $\xi_n = 0$；而
$\xi_{m_i} = (0, \ldots, 0, x_i, 0, \ldots)$，其中 $x_i$
置于第 $i$ 个直和项中。我们断言，$\xi$ 不在引理
\ref{more-algebra-lemma-compute-Rlim} 的映射 $\prod B_n \to \prod B_n$ 的像中。
这表明 $R^1\lim B_n$ 非零，并完成证明。事实上，假设 $\xi$
是 $\eta = (z_1, z_2, \ldots)$ 的像，其中
$z_n = \sum z_{n, i} \in \bigoplus_i A_n$。注意
$x_i = z_{m_i, i} \bmod A_{m_i, m_i + 1}$。
于是 $z_{m_i - 1, i}$ 是 $z_{m_i, i}$ 在
$A_{m_i} \to A_{m_i - 1}$ 下的像，依此类推，故
$z_{m, i}$ 是 $z_{m_i, i}$ 在 $A_{m_i} \to A_m$ 下的像。
因此，$z_{m, i}$ 模 $A_{m, m_i + 1}$ 同余于 $y_i$。
特别地，$z_{m, i} \not = 0$。但
$\sum z_{m, i} \in \bigoplus_i A_m$，所以只有有限多个
$z_{m, i}$ 可以非零，矛盾。
\end{proof}

\begin{lemma}
\label{more-algebra-lemma-Mittag-Leffler}
设
$$
0 \to (A_i) \to (B_i) \to (C_i) \to 0
$$
为 Abel 群逆系的短正合序列。若 $(A_i)$ 和 $(C_i)$ 是 ML 的，
则 $(B_i)$ 也是 ML 的。
\end{lemma}

\begin{proof}
这由引理 \ref{more-algebra-lemma-emmanouil}、无穷直和保持正合这一事实，
以及与 $R\lim$ 相伴的上同调长正合序列得到。
\end{proof}

\begin{lemma}
\label{more-algebra-lemma-Rlim-zero-of-direct-sums}
设 $(A_n)$ 为 Abel 群逆系。下列条件等价：
\begin{enumerate}
\item $(A_n)$ 作为 pro-对象为零；
\item $\lim A_n = 0$ 且 $R^1\lim A_n = 0$，并且
$\bigoplus_{i \in \mathbf{N}} (A_n)$ 也满足相同结论。
\end{enumerate}
\end{lemma}

\begin{proof}
由引理 \ref{more-algebra-lemma-pro-isomorphism-iso-Rlim}，(1) 蕴含 (2)。
假设 (2)。由引理 \ref{more-algebra-lemma-emmanouil}，$(A_n)$ 是 ML 的。
对 $m \geq n$，令 $A_{n, m} = \Im(A_m \to A_n)$，于是
$A_n = A_{n, n} \supset A_{n, n + 1} \supset \ldots$。
注意，$(A_n)$ 作为 pro-对象为零，当且仅当对每个 $n$，
存在 $m \geq n$ 使得 $A_{n, m} = 0$。还要注意，$(A_n)$ 是 ML 的，
当且仅当对每个 $n$，存在 $m_n \geq n$，使得
$A_{n, m} = A_{n, m + 1} = \ldots$。在 ML 情形，显然
$\lim A_n = 0$ 蕴含 $A_{n, m_n} = 0$，因为映射
% P02-E0360
$A_{n + 1, m_{n + 1}} \to A_{n, m}$ 都是满射。
证明完成。
\end{proof}

\section{模的 Rlim}
\label{more-algebra-section-Rlim-modules}

\noindent
我们简要讨论模上的 $R\lim$。本节的许多论证重复了构造环化位点上
模的上同调机制时所用的论证。

\medskip\noindent
设 $(A_n)$ 为环的逆系。我们用
% P02-E0361
$\textit{Mod}(\mathbf{N}, (A_n))$ 表示 Abel 群逆系 $(M_n)$ 的范畴，
其中每个 $M_n$ 都赋予 $A_n$-模结构，% P02-E0362
且转移映射 $M_{n + 1} \to M_n$ 是 $A_{n + 1}$-模映射。
这是一个 Abel 范畴。置 $A = \lim A_n$。给定
$\textit{Mod}(\mathbf{N}, (A_n))$ 的对象 $(M_n)$，极限
$\lim M_n$ 是 $A$-模。

\begin{lemma}
\label{more-algebra-lemma-compute-Rlim-modules}
在上述情形中，% P02-E0363
函子 $\lim : \textit{Mod}(\mathbf{N}, (A_n)) \to \text{Mod}_A$
具有右导出函子
$$
R\lim :
D(\textit{Mod}(\mathbf{N}, (A_n)))
\longrightarrow
D(A)
$$
按惯例，置 $R^p\lim(K) = H^p(R\lim(K))$。此外：
\begin{enumerate}
\item 对 $\textit{Mod}(\mathbf{N}, (A_n))$ 中任意 $(M_n)$，
当 $p > 1$ 时有 $R^p\lim M_n = 0$；
\item $D(\text{Mod}_A)$ 的对象 $R\lim M_n$ 由复形
$$
\prod M_n \to \prod M_n,\quad
(x_n) \mapsto (x_n - f_{n + 1}(x_{n + 1}))
$$
表示，该复形位于次数 $0$ 和 $1$；
\item 若 $(M_n)$ 是 ML 的，则 $R^1\lim M_n = 0$，即
$(M_n)$ 对 $\lim$ 右非周期；
\item 每个 $K^\bullet \in D(\textit{Mod}(\mathbf{N}, (A_n)))$
都拟同构于一个逐项对 $\lim$ 右非周期的复形；
\item 若每个 $K^p = (K^p_n)$ 都对 $\lim$ 右非周期，即
% P02-E0364
$R^1\lim_n K^p_n = 0$，则 $R\lim K$ 由次数 $p$ 的项为
$\lim_n K_n^p$ 的复形表示。
\end{enumerate}
\end{lemma}

\begin{proof}
其证明与引理 \ref{more-algebra-lemma-compute-Rlim} 的证明逐字相同。
\end{proof}

\begin{remark}
\label{more-algebra-remark-Rlim-cohomology-modules}
本评注续接评注 \ref{more-algebra-remark-Rlim-cohomology}。
$\mathbf{N}$ 上的环层正是环的逆系 $(A_n)$。
$(A_n)$ 上的模层正好等同于上面定义的范畴
$\textit{Mod}(\mathbf{N}, (A_n))$ 的对象。引理
\ref{more-algebra-lemma-compute-Rlim-modules} 的导出函子 $R\lim$
就是从模的导出范畴到结构层整体截面环上模的导出范畴的
$R\Gamma(\mathbf{N}, -)$。一般而言，群和模的上同调相同；见
Cohomology on Sites，引理
\ref{sites-cohomology-lemma-cohomology-modules-abelian-agree}。
\end{remark}

\noindent
下一个引理中的积既可以看作复形的逐项积，也可以看作导出范畴
$D(A)$ 中的积；见 Derived Categories，引理
\ref{derived-lemma-products}。

\begin{lemma}
\label{more-algebra-lemma-distinguished-triangle-Rlim-modules}
设 $K = (K_n^\bullet)$ 为
$D(\textit{Mod}(\mathbf{N}, (A_n)))$ 的对象。
在 $D(A)$ 中存在典范区分三角
$$
R\lim K \to \prod\nolimits_n K_n^\bullet \to \prod\nolimits_n K_n^\bullet
\to R\lim K[1]
$$
换言之，$R\lim K$ 是 $D(A)$ 中逆系 $(K_n^\bullet)$ 的一个导出极限；
见 Derived Categories，定义 \ref{derived-definition-derived-limit}。
\end{lemma}

\begin{proof}
其证明与引理 \ref{more-algebra-lemma-distinguished-triangle-Rlim} 的证明完全相同，
但使用引理 \ref{more-algebra-lemma-compute-Rlim-modules} 代替
% P02-E0365
引理 \ref{more-algebra-lemma-compute-Rlim}。
\end{proof}

\begin{lemma}
\label{more-algebra-lemma-break-long-exact-sequence-modules}
沿用引理 \ref{more-algebra-lemma-distinguished-triangle-Rlim-modules} 的记号，
与该区分三角相伴的上同调长正合序列分解为 $A$-模的短正合序列
$$
0 \to R^1\lim_n H^{p - 1}(K_n^\bullet) \to
H^p(R\lim K) \to
\lim_n H^p(K_n^\bullet) \to 0
$$
\end{lemma}

\begin{proof}
其证明与引理 \ref{more-algebra-lemma-break-long-exact-sequence} 的证明完全相同，
但使用引理 \ref{more-algebra-lemma-compute-Rlim-modules} 代替
% P02-E0365
引理 \ref{more-algebra-lemma-compute-Rlim}。
\end{proof}

\noindent
{\bf 警告。}与 Abel 群的情形一样，
$D(\textit{Mod}(\mathbf{N}, (A_n)))$ 的对象
$M = (M_n^\bullet)$ 是模复形的逆系，这与各导出范畴中对象的逆系
{\bf 并不}相同。在下一个引理中，我们证明导出范畴中对象的逆系
总能提升为 $D(\textit{Mod}(\mathbf{N}, (A_n)))$ 的对象。

\begin{lemma}
\label{more-algebra-lemma-lift-to-system-complexes}
设 $(A_n)$ 为环的逆系。假设给定：
\begin{enumerate}
\item 对每个 $n$，给定 $D(A_n)$ 的对象 $K_n$；
\item 对每个 $n$，给定 $D(A_{n + 1})$ 中的映射
$\varphi_n : K_{n + 1} \to K_n$，其中通过
$A_{n + 1} \to A_n$ 作限制，把 $K_n$ 看作
$D(A_{n + 1})$ 的对象。
\end{enumerate}
则存在对象
$M = (M_n^\bullet) \in D(\textit{Mod}(\mathbf{N}, (A_n)))$
以及 $D(A_n)$ 中的同构 $\psi_n : M_n^\bullet \to K_n$，
使得图
$$
\xymatrix{
M_{n + 1}^\bullet \ar[d]_{\psi_{n + 1}} \ar[r] & M_n^\bullet \ar[d]^{\psi_n} \\
K_{n + 1} \ar[r]^{\varphi_n} & K_n
}
$$
在 $D(A_{n + 1})$ 中交换。
\end{lemma}

\begin{proof}
我们详细写出证明。对 $A_n$-模 $T$，用
$T_{A_{n + 1}}$ 表示视为 $A_{n + 1}$-模的同一个模。
% P02-E0366
假设 $K_n^\bullet$ 是表示 $K_n$ 的 $A_n$-模复形。
则 $K_{n, A_{n + 1}}^\bullet$ 是同一个复形，但视为
$A_{n + 1}$-模复形。根据导出范畴的构造，映射
% P02-E0367
$\psi_n$ 可以写成
$$
\psi_n = \tau_n \circ \sigma_n^{-1}
$$
其中 $\sigma_n : L_{n + 1}^\bullet \to K_{n + 1}^\bullet$
是 $A_{n + 1}$-模复形的拟同构，而
$\tau_n : L_{n + 1}^\bullet \to K_{n, A_{n + 1}}^\bullet$
是 $A_{n + 1}$-模复形的映射。

\medskip\noindent
现在归纳构造复形 $M_n^\bullet$。作为基础情形，令
$M_1^\bullet = K_1^\bullet$。假设已经构造
$M_e^\bullet \to M_{e - 1}^\bullet \to \ldots \to M_1^\bullet$
以及复形映射 $\psi_i : M_i^\bullet \to K_i^\bullet$，使得对所有
$n < e$，上面的图
$$
\xymatrix{
M_{n + 1}^\bullet \ar[d]_{\psi_{n + 1}} \ar[rr] & &
M_{n, A_{n + 1}}^\bullet \ar[d]^{\psi_{n, A_{n + 1}}} \\
K_{n + 1}^\bullet & L_{n + 1}^\bullet \ar[l]_{\sigma_n} \ar[r]^{\tau_n} &
K_{n, A_{n + 1}}^\bullet
}
$$
在 $D(A_{n + 1})$ 中交换。然后考虑 $D(A_{e + 1})$ 中的图
$$
\xymatrix{
& & M_{e, A_{e + 1}}^\bullet \ar[d]^{\psi_{e, A_{e + 1}}} \\
K_{e + 1}^\bullet & L_{e + 1}^\bullet \ar[r]^{\tau_e} \ar[l]_{\sigma_e} &
K_{e, A_{e + 1}}^\bullet
}
$$
由于 $\psi_e$ 是拟同构，$\psi_{e, A_{e + 1}}$ 也是拟同构。
根据 $D(A_{e + 1})$ 中态射的定义，可以找到
$A_{e + 1}$-模复形的拟同构
$\psi_{e + 1} : M_{e + 1}^\bullet \to K_{e + 1}^\bullet$，
使得存在 $A_{e + 1}$-模的复形态射
$M_{e + 1}^\bullet \to M_{e, A_{e + 1}}^\bullet$，它在
$D(A_{e + 1})$ 中表示复合
$\psi_{e, A_{e + 1}}^{-1} \circ \tau_e \circ \sigma_e^{-1}$。
故由归纳法，引理成立。
\end{proof}

\begin{remark}
\label{more-algebra-remark-how-unique}
沿用引理 \ref{more-algebra-lemma-lift-to-system-complexes} 的假设。% P02-E0368
先验地，$D(\textit{Mod}(\mathbf{N}, (A_n)))$ 中有许多同构类的对象
$M$ 能给出引理中的系统 $(K_n, \varphi_n)$。对每个这样的 $M$，
可以考察复形 $R\lim M \in D(A)$，其中 $A = \lim A_n$。
由引理 \ref{more-algebra-lemma-distinguished-triangle-Rlim-modules}，
$R\lim M$ 是 $D(A)$ 中逆系 $(K_n)$ 的导出极限。
因此，$R\lim M$ 在 $D(A)$ 中的同构类与构造 $M$ 时所作的选择无关。
特别地，可以把本节证明的关于 $R\lim$ 的结果应用于
$D(A)$ 中逆系的导出极限。例如，对每个 $p \in \mathbf{Z}$，
都有典范短正合序列
$$
0 \to R^1\lim H^{p - 1}(K_n) \to H^p(R\lim K_n) \to \lim H^p(K_n) \to 0
$$
因为可以对 $M$ 应用引理
\ref{more-algebra-lemma-distinguished-triangle-Rlim-modules}。
% P02-E0369
不借助 $M$ 的存在性，也可以直接看出这一点：把引理
\ref{more-algebra-lemma-distinguished-triangle-Rlim-modules} 证明中的论证应用于
导出极限的（定义性）区分三角
$R\lim K_n \to \prod K_n \to \prod K_n \to (R\lim K_n)[1]$
即可。
\end{remark}

\begin{lemma}
\label{more-algebra-lemma-get-ML-system}
设 $(A_n)$ 为环的逆系。每个
$K \in D(\textit{Mod}(\mathbf{N}, (A_n)))$
都可以由复形系 $(M_n^\bullet)$ 表示，使得所有转移映射
$M_{n + 1}^\bullet \to M_n^\bullet$ 都满射。
\end{lemma}

\begin{proof}
设 $K$ 由系统 $(K_n^\bullet)$ 表示。置
$M_1^\bullet = K_1^\bullet$。假设已经构造复形的满射
$M_n^\bullet \to M_{n - 1}^\bullet \to \ldots \to M_1^\bullet$
以及同伦等价 $\psi_e : K_e^\bullet \to M_e^\bullet$，使得对所有
$e < n$，图
$$
\xymatrix{
K_{e + 1}^\bullet \ar[d] \ar[r] & K_e^\bullet \ar[d] \\
M_{e + 1}^\bullet \ar[r] & M_e^\bullet
}
$$
交换。然后考虑图
$$
\xymatrix{
K_{n + 1}^\bullet \ar[r] & K_n^\bullet \ar[d] \\
& M_n^\bullet
}
$$
由 Derived Categories，引理 \ref{derived-lemma-make-surjective}，
可以把复合 $K_{n + 1}^\bullet \to M_n^\bullet$ 分解为
$K_{n + 1}^\bullet \to M_{n + 1}^\bullet \to M_n^\bullet$，
使得第一个箭头是同伦等价，第二个箭头是逐项分裂满射。
由此并结合归纳法，引理成立。
\end{proof}

\begin{lemma}
\label{more-algebra-lemma-get-K-flat-system}
设 $(A_n)$ 为环的逆系。每个
$K \in D(\textit{Mod}(\mathbf{N}, (A_n)))$
都可以由复形系 $(K_n^\bullet)$ 表示，使得每个
$K_n^\bullet$ 都是 K-平坦的。
\end{lemma}

\begin{proof}
首先使用引理 \ref{more-algebra-lemma-get-ML-system}，以复形系
$(M_n^\bullet)$ 表示 $K$，使得所有转移映射
$M_{n + 1}^\bullet \to M_n^\bullet$ 都满射。
接着，取拟同构 $K_1^\bullet \to M_1^\bullet$，其中
$K_1^\bullet$ 是 K-平坦 $A_1$-模复形
（引理 \ref{more-algebra-lemma-K-flat-resolution}）。
假设已经构造
$K_n^\bullet \to K_{n - 1}^\bullet \to \ldots \to K_1^\bullet$
以及复形映射 $\psi_e : K_e^\bullet \to M_e^\bullet$，使得对所有
$e < n$，图
$$
\xymatrix{
K_{e + 1}^\bullet \ar[d] \ar[r] & K_e^\bullet \ar[d] \\
M_{e + 1}^\bullet \ar[r] & M_e^\bullet
}
$$
交换。然后考虑 $D(A_{n + 1})$ 中的图
$$
\xymatrix{
C^\bullet \ar@{..>}[d] \ar@{..>}[r] & K_n^\bullet \ar[d]^{\psi_n} \\
M_{n + 1}^\bullet \ar[r]^{\varphi_n} & M_n^\bullet
}
$$
由于 $M_{n + 1}^\bullet \to M_n^\bullet$ 逐项满射，
置于左上角、各项为
$$
C^p = M_{n + 1}^p \times_{M_n^p} K_n^p
$$
的复形 $C^\bullet$ 拟同构于 $M_{n + 1}^\bullet$（细节省略）。
选取拟同构 $K_{n + 1}^\bullet \to C^\bullet$，其中
$K_{n +1}^\bullet$ K-平坦。故由归纳法，引理成立。
\end{proof}

\begin{lemma}
\label{more-algebra-lemma-derived-tensor-product-systems}
设 $(A_n)$ 为环的逆系。给定
$K, L \in D(\textit{Mod}(\mathbf{N}, (A_n)))$，在
% P02-E0370
$D(\mathbf{N}, (A_n))$ 中有典范导出张量积
$K \otimes^\mathbf{L} L$，它与到 $D(A_n)$ 的映射相容。
该构造关于 $K$ 和 $L$ 对称，并且在每个变量上都是三角范畴的正合函子。
\end{lemma}

\begin{proof}
为 $K$ 选取代表元 $(K_n^\bullet)$，使得每个 $K_n^\bullet$
都是 K-平坦复形（引理 \ref{more-algebra-lemma-get-K-flat-system}）。
然后，对 $L$ 的任意代表元选择 $(L_n^\bullet)$，可以把
$K \otimes^\mathbf{L} L$ 定义为由复形系
$$
(\text{Tot}(K_n^\bullet \otimes_{A_n} L_n^\bullet))
$$
表示的对象。由引理 \ref{more-algebra-lemma-K-flat-quasi-isomorphism}
和 \ref{more-algebra-lemma-derived-tor-quasi-isomorphism-other-side}，
这在两个变量上都是良定义的。它与到 $D(A_n)$ 的映射相容是显然的。
正合性与引理 \ref{more-algebra-lemma-derived-tor-exact} 中完全相同。
\end{proof}

\begin{remark}
\label{more-algebra-remark-constructing-tensor-with-limits-functorially}
设 $A$ 为环，$(E_n)$ 为 $D(A)$ 的对象之逆系。上面已经看到，
导出极限 $R\lim E_n$ 存在。因此，对 $D(A)$ 的每个对象 $K$，
导出极限 $R\lim( K \otimes_A^\mathbf{L} E_n )$ 也存在。
事实上，可以关于 $K$ 函子性地构造这些导出极限，并得到三角范畴的正合函子
$$
R\lim(- \otimes_A^\mathbf{L} E_n) : D(A) \longrightarrow D(A)
$$
具体地，首先把 $(E_n)$ 提升为 $D(\mathbf{N}, A)$ 的对象 $E$；
见引理 \ref{more-algebra-lemma-lift-to-system-complexes}。
（该函子将依赖于这一提升的选择。）接着，注意有“对角”或“常值”函子
$$
\Delta : D(A) \longrightarrow D(\mathbf{N}, A)
$$
它把复形 $K^\bullet$ 映到值为 $K^\bullet$ 的常值复形逆系。
然后直接定义
$$
R\lim(K \otimes_A^\mathbf{L} E_n) = R\lim(\Delta(K)\otimes^\mathbf{L} E)
$$
右边使用引理 \ref{more-algebra-lemma-compute-Rlim-modules} 的函子 $R\lim$
以及引理 \ref{more-algebra-lemma-derived-tensor-product-systems} 的函子
$- \otimes^\mathbf{L} -$。
\end{remark}

\begin{lemma}
\label{more-algebra-lemma-tensor-Rlim-exact}
设 $A$ 为环。设 $E \to D \to F \to E[1]$ 为
$D(\mathbf{N}, A)$ 的区分三角。设 $(E_n)$、相应地 $(D_n)$、
相应地 $(F_n)$ 为分别与 $E$、相应地 $D$、相应地 $F$
关联的 $D(A)$ 对象之系统。则对每个 $K \in D(A)$，
都有典范区分三角
$$
R\lim (K \otimes^\mathbf{L}_A E_n) \to
R\lim (K \otimes^\mathbf{L}_A D_n) \to
R\lim (K \otimes^\mathbf{L}_A F_n) \to
R\lim (K \otimes^\mathbf{L}_A E_n)[1]
$$
位于 $D(A)$ 中；记号如评注
\ref{more-algebra-remark-constructing-tensor-with-limits-functorially}。
\end{lemma}

\begin{proof}
这由评注 \ref{more-algebra-remark-constructing-tensor-with-limits-functorially}
中的构造，以及 $\Delta : D(A) \to D(\mathbf{N}, A)$、
$- \otimes^\mathbf{L} -$ 和 $R\lim$ 都是三角范畴的正合函子这一事实
直接得到。
\end{proof}

\begin{lemma}
\label{more-algebra-lemma-tensor-Rlim-pro-equal}
设 $A$ 为环。设 $E \to D$ 为 $D(\mathbf{N}, A)$ 的态射。
设 $(E_n)$、相应地 $(D_n)$ 为分别与 $E$、相应地 $D$
关联的 $D(A)$ 对象之系统。若 $(E_n) \to (D_n)$ 是 pro-对象的同构，
则对每个 $K \in D(A)$，相应映射
$$
R\lim (K \otimes^\mathbf{L}_A E_n)
\longrightarrow
R\lim (K \otimes^\mathbf{L}_A D_n)
$$
是 $D(A)$ 中的同构（记号如评注
\ref{more-algebra-remark-constructing-tensor-with-limits-functorially}）。
\end{lemma}

\begin{proof}
这由定义和引理 \ref{more-algebra-lemma-Rlim-pro-equal} 得到。% P02-E0371
\end{proof}

\section{挠模}
\label{more-algebra-section-torsion}

\noindent
本节中的“挠模”指支撑在仿射概形 $\Spec(R)$ 的给定闭子集
$V(I)$ 上的模。这不同于、但类似于整环上的挠模概念
（定义 \ref{more-algebra-definition-torsion}）。

\begin{definition}
\label{more-algebra-definition-f-power-torsion}
设 $R$ 为环，$M$ 为 $R$-模。
\begin{enumerate}
\item 设 $I \subset R$ 为理想。若对每个 $m \in M$，存在
$n > 0$ 使得 $I^n m = 0$，则称 $M$ 为{\it $I$-幂挠模}。
\item 设 $f \in R$。若对每个 $m \in M$，存在 $n > 0$
使得 $f^n m = 0$，则称 $M$ 为{\it $f$-幂挠模}。
\end{enumerate}
\end{definition}

\noindent
因此，$f$-幂挠模等同于取 $I = (f)$ 时的 $I$-幂挠模。
对 $R$-模 $M$，将使用记号
$$
M[I^n] = \{m \in M \mid I^nm = 0\}
$$
以及
$$
M[I^\infty] = \bigcup M[I^n]
$$
因此，$M$ 是 $I$-幂挠的，当且仅当
$M = M[I^\infty]$，当且仅当 $M = \bigcup M[I^n]$。

\begin{lemma}
\label{more-algebra-lemma-I-power-torsion-presentation}
设 $R$ 为环，$I$ 为 $R$ 的理想，并设 $M$ 为 $I$-幂挠模。
则 $M$ 具有解消
$$
\ldots \to K_2 \to K_1 \to K_0 \to M \to 0
$$
其中每个 $K_i$ 都是若干 $R/I^n$ 副本的直和，而 $n$ 可变。
% P02-E0372
\end{lemma}

\begin{proof}
有典范满射
$$
\oplus_{m \in M} R/I^{n_m} \to M \to 0
$$
其中 $n_m$ 是满足 $I^{n_m} \cdot m = 0$ 的最小正整数。
上述满射的核也是 $I$-幂挠模。归纳进行，即构造出所需的
$M$ 的解消。
\end{proof}

\begin{lemma}
\label{more-algebra-lemma-torsion-free}
设 $R$ 为环，$I$ 为 $R$ 的理想。对任意 $R$-模 $M$，置
$M[I^n] = \{m \in M \mid I^nm = 0\}$。若 $I$ 有限生成，
则下列条件等价：
\begin{enumerate}
\item $M[I] = 0$；
\item 对所有 $n \geq 1$，有 $M[I^n] = 0$；
\item 若 $I = (f_1, \ldots, f_t)$，则映射
$M \to \bigoplus M_{f_i}$ 是单射。
\end{enumerate}
\end{lemma}

\begin{proof}
这由 Algebra，引理 \ref{algebra-lemma-when-injective-covering} 得到。
\end{proof}

\begin{lemma}
\label{more-algebra-lemma-divide-by-torsion}
设 $R$ 为环，$I$ 为 $R$ 的有限生成理想。
\begin{enumerate}
\item 对任意 $R$-模 $M$，有 $(M/M[I^\infty])[I] = 0$。
\item $I$-幂挠模的扩张仍是 $I$-幂挠的。
\end{enumerate}
\end{lemma}

\begin{proof}
设 $m \in M$。若 $m$ 映到 $(M/M[I^\infty])[I]$ 的元素，
则 $Im \subset M[I^\infty]$。写成 $I = (f_1, \ldots, f_t)$。
于是 $f_i m \in M[I^\infty]$，即对某个 $n_i > 0$，
有 $I^{n_i}f_i m = 0$。因此 $I^Nm = 0$，其中
$N = \sum n_i + 2$。故 $m$ 在 $(M/M[I^\infty])$ 中映到零，
这证明了引理的第一个陈述。

\medskip\noindent
对第二个陈述，设
$0 \to M' \to M \to M'' \to 0$ 为模的短正合序列，其中
$M'$ 和 $M''$ 都是 $I$-幂挠模。则
$M[I^\infty] \supset M'$，因而 $M/M[I^\infty]$ 是 $M''$
的商，所以是 $I$-幂挠的。结合第一个陈述和引理
\ref{more-algebra-lemma-torsion-free}，可知它为零。
\end{proof}

\begin{lemma}
\label{more-algebra-lemma-I-power-torsion}
设 $I$ 为环 $R$ 的有限生成理想。$I$-幂挠模组成 Abel 范畴
$\text{Mod}_R$ 的一个 Serre 子范畴；见 Homology，定义
\ref{homology-definition-serre-subcategory}。
\end{lemma}

\begin{proof}
显然，$I$-幂挠模的子模与商模都是 $I$-幂挠的。% P02-E0373
此外，由引理 \ref{more-algebra-lemma-divide-by-torsion}，两个 $I$-幂挠模的扩张
也是 $I$-幂挠的。因此，由 Homology，引理
\ref{homology-lemma-characterize-serre-subcategory}，引理的陈述成立。
% P02-E0374
\end{proof}

\begin{lemma}
\label{more-algebra-lemma-local-cohomology-closed}
设 $R$ 为环，$I \subset R$ 为有限生成理想。子范畴
$I^\infty\text{-torsion} \subset \text{Mod}_R$
仅依赖于闭子集 $Z = V(I) \subset \Spec(R)$。
事实上，$R$-模 $M$ 是 $I$-幂挠的，当且仅当其支撑包含于 $Z$。
\end{lemma}

\begin{proof}
设 $M$ 为 $R$-模，$x \in M$。若 $x \in M[I^\infty]$，
则对所有 $f \in I$，$x$ 在 $M_f$ 中映到零。因此，对所有
$\mathfrak p \not \supset I$，$x$ 在 $M_\mathfrak p$ 中映到零。
反之，若对所有 $\mathfrak p \not \supset I$，$x$ 在
$M_\mathfrak p$ 中映到零，则对所有 $f \in I$，$x$ 在
$M_f$ 中映到零。因此，若 $I = (f_1, \ldots, f_r)$，
则对某个 $n_i \geq 1$，有 $f_i^{n_i}x = 0$。
于是 $x \in M[I^{\sum n_i}]$。故 $M[I^\infty]$ 是映射
$M \to \prod_{\mathfrak p \not \in Z} M_\mathfrak p$ 的核。
引理的第二个陈述成立，并且它蕴含第一个陈述。
\end{proof}

\noindent
下面两个引理或许应放在别处。

\begin{lemma}
\label{more-algebra-lemma-bounded-above-mod-I}
设 $R$ 为环，$I \subset R$ 为理想。设 $K$ 为 $D(R)$ 的对象，
且当 $i > 0$ 时，有 $H^i(K \otimes_R^\mathbf{L} R/I) = 0$。则：
\begin{enumerate}
\item 对所有 $n \geq 1$ 和 $i > 0$，有
$H^i(K \otimes_R^\mathbf{L} R/I^n) = 0$；
\item 对任意 $I$-幂挠 $R$-模 $N$ 和 $i > 0$，有
$H^i(K \otimes_R^\mathbf{L} N) = 0$；
\item 对任意 $M \in D^b(R)$，若其上同调模 $H^i(M)$
都是 $I$-幂挠的，并且在 $i > 0$ 时为 $0$，则在 $i > 0$ 时有
$H^i(K \otimes_R^\mathbf{L} M) = 0$。
\end{enumerate}
\end{lemma}

\begin{proof}
证明 (2)。可以写成 $N = \bigcup N[I^n]$。由于张量积与余极限交换，
有
$K \otimes_R^\mathbf{L} N = \text{hocolim}_n K \otimes_R^\mathbf{L} N[I^n]$
（细节省略；提示：用 K-平坦复形表示 $K$ 并直接计算）。
因此可以假设 $N$ 被 $I^n$ 零化。考虑 $R$-代数
$R' = R/I^n \oplus N$，其中 $N$ 是平方为零的理想。
只需证明 $D(R')$ 的对象 $K' = K \otimes_R^\mathbf{L} R'$
在正次数上的上同调消失。存在 $R$-代数满射
$R' \to R/I$，其核 $J$ 为幂零的
% P02-E0375
（核中任意 $n$ 个元素的乘积为零）。有
$$
K \otimes_R^\mathbf{L} R/I =
(K \otimes_R^\mathbf{L} R') \otimes_{R'}^\mathbf{L} R/I =
K' \otimes_{R'}^\mathbf{L} R/I
$$
由假设，复形 $K \otimes_R^\mathbf{L} R/I$ 的 Tor 振幅位于
$[-\infty, 0]$ 中。因此，由引理
\ref{more-algebra-lemma-check-tor-dimension-modulo-nilpotent} 得到结论。
% P02-E0376

\medskip\noindent
(1) 从 (2) 平凡地得到。(3) 由 (2)、对 $M$ 的非零上同调模个数的归纳，
以及 Derived Categories，评注
\ref{derived-remark-truncation-distinguished-triangle} 的截断区分三角得到。
细节省略。
\end{proof}

\begin{lemma}
\label{more-algebra-lemma-derived-vanishing-mod-I}
设 $R$ 为环，$I \subset R$ 为理想。设 $K$ 为 $D(R)$ 的对象，
且 $K \otimes_R^\mathbf{L} R/I = 0$ 在 $D(R)$ 中成立。则：
\begin{enumerate}
\item 对所有 $n \geq 1$，有 $K \otimes_R^\mathbf{L} R/I^n = 0$；
\item 对任意 $I$-幂挠 $R$-模 $N$，有
$K \otimes_R^\mathbf{L} N = 0$；
\item 对任意 $M \in D^b(R)$，若其上同调模都是 $I$-幂挠的，
则 $K \otimes_R^\mathbf{L} M = 0$。
\end{enumerate}
\end{lemma}

\begin{proof}
这是引理 \ref{more-algebra-lemma-bounded-above-mod-I} 的推论（省略一些细节）。
% P02-E0377
\end{proof}

\begin{lemma}
\label{more-algebra-lemma-torsion-completion}
\begin{reference}
\cite[Lemme~1]{Beauville-Laszlo} 的轻微推广。
\end{reference}
设 $R \to R'$ 为环同态。设 $I \subset R$ 为理想，使得对
$n > 0$，$R/I^n \to R'/I^nR'$ 是同构。对任意 $I$-幂挠
$R$-模 $M$，映射 $M \to M \otimes_R R'$ 是同构。
例如，若 $I$ 有限生成，且 $R^\wedge$ 是 $R$ 关于 $I$ 的完备化，
则有 $M \cong M \otimes_R R^\wedge$。
\end{lemma}

\begin{proof}
若 $M$ 被 $I^n$ 零化，则
$$
M \otimes_R R' \cong
M \otimes_{R/I^n} R'/I^n R' \cong
M \otimes_{R/I^n} R/I^n \cong M.
$$
若 $M$ 是 $I$-幂挠的，则 $M = \bigcup M[I^n]$。
由于张量积与直极限交换
（Algebra，引理 \ref{algebra-lemma-tensor-products-commute-with-limits}），
得到所需同构。最后一个陈述由 Algebra，引理
\ref{algebra-lemma-hathat-finitely-generated} 可知是第一个陈述的特例。
\end{proof}

\section{模范畴的形式粘合}
\label{more-algebra-section-formal-glueing}

\noindent
固定一个 Noether 概形 $X$，以及一个闭子概形 $Z$，其补集为 $U$。
我们的目标是说明，如何从 $X$ 沿 $Z$ 的形式完备化上的凝聚层，
以及 $U$ 上在交叠处满足适当相容性的凝聚层，（唯一地）构造出
$X$ 上的凝聚层。我们先仅用交换代数完成此事（本节），随后在
代数空间的框架下说明它
（Pushouts of Spaces，第 \ref{spaces-pushouts-section-formal-glueing} 节）。

\medskip\noindent
以下文献讨论了本节的部分内容：
\cite[Section 2]{ArtinII}、
\cite[Appendix]{Ferrand-Raynaud}、
\cite{Beauville-Laszlo}、
\cite{MB} 以及
\cite[Section 4.6]{dJ-crystalline}。

\begin{lemma}
\label{more-algebra-lemma-characterize-flatness-on-torsion}
设 $\varphi : R \to S$ 为环映射，$I \subset R$ 为理想。
下列条件等价：
\begin{enumerate}
\item $\varphi$ 平坦，且 $R/I \to S/IS$ 忠实平坦；
\item $\varphi$ 平坦，且映射
$\Spec(S/IS) \to \Spec(R/I)$ 满射；
\item $\varphi$ 平坦，且基变换函子
$M \mapsto M \otimes_R S$ 在被 $I$ 零化的模上是忠实的；
\item $\varphi$ 平坦，且基变换函子
$M \mapsto M \otimes_R S$ 在 $I$-幂挠模上是忠实的。
\end{enumerate}
\end{lemma}

\begin{proof}
若 $R \to S$ 平坦，则对每个 $n$，$R/I^n \to S/I^nS$ 都平坦；见
Algebra，引理 \ref{algebra-lemma-flat-base-change}。
故由 Algebra，引理 \ref{algebra-lemma-ff-rings}，(1) 与 (2) 等价。
将 $I$-挠 $R$-模等同于 $R/I$-模，并利用对被 $I$ 零化的
$R$-模 $M$ 有
$$
M \otimes_R S = M \otimes_{R/I} S/IS
$$
再由 Algebra，引理 \ref{algebra-lemma-easy-ff}，可得 (1) 与 (3) 等价。
(4) $\Rightarrow$ (3) 是立即的。假设 (3)。上面已经看到
$R/I^n \to S/I^nS$ 平坦；由假设，它在谱上诱导满射，因为
$\Spec(R/I^n) = \Spec(R/I)$，对 $S$ 亦然。故基变换函子在被
$I^n$ 零化的模上是忠实的。任何 $I$-幂挠模 $M$ 都是并
$M = \bigcup M_n$，其中 $M_n$ 被 $I^n$ 零化，因此基变换函子
在所有 $I$-幂挠模组成的范畴上是忠实的
（因为张量积与余极限交换）。
% P02-E0378
\end{proof}

\begin{lemma}
\label{more-algebra-lemma-neighbourhood-isomorphism}
假设 $(\varphi : R \to S, I)$ 满足引理
\ref{more-algebra-lemma-characterize-flatness-on-torsion} 的等价条件。
下列条件等价：
\begin{enumerate}
\item 对任意 $I$-幂挠模 $M$，自然映射
$M \to M \otimes_R S$ 是同构；
\item $R/I \to S/IS$ 是同构。
\end{enumerate}
\end{lemma}

\begin{proof}
(1) $\Rightarrow$ (2) 是立即的。假设 (2)。先假设 $M$ 被 $I$ 零化。
此时 $M$ 是 $R/I$-模，故有同构
$$
M \otimes_R S = M \otimes_{R/I} S/IS = M \otimes_{R/I} R/I = M
$$
从而结论成立。接着归纳证明：对任意被 $I^n$ 零化的模 $M$，
$M \to M \otimes_R S$ 是同构。
% P02-E0379
假设归纳假设对 $n$ 成立，并设 $M$ 被 $I^{n + 1}$ 零化。于是有短正合列
$$
0 \to I^nM \to M \to M/I^nM \to 0
$$
由于 $R \to S$ 平坦，这给出短正合列
$$
0 \to I^nM \otimes_R S \to M \otimes_R S \to M/I^nM \otimes_R S \to 0
$$
对 $M' = I^nM$ 和 $M'' = M/I^nM$，规范映射由归纳假设是同构，
故对 $M$ 亦然。最后，设 $M$ 为一般的 $I$-幂挠模。则
$M = \bigcup M_n$，其中 $M_n$ 被 $I^n$ 零化；再利用张量积与
余极限交换即可得出结论。
\end{proof}

\begin{lemma}
\label{more-algebra-lemma-neighbourhood-equivalence}
设 $\varphi : R \to S$ 为平坦环映射，$I \subset R$ 为有限生成理想，
且 $R/I \to S/IS$ 是同构。则：
\begin{enumerate}
\item 对任意 $R$-模 $M$，映射 $M \to M \otimes_R S$ 在
$I$-幂挠子模上诱导同构
$M[I^\infty] \to (M \otimes_R S)[(IS)^\infty]$；
\item 若 $M$ 或 $N$ 是 $I$-幂挠的，则自然映射
$$
\Hom_R(M, N) \longrightarrow \Hom_S(M \otimes_R S, N \otimes_R S)
$$
是同构；
\item 基变换函子 $M \mapsto M \otimes_R S$ 给出 $I$-幂挠模范畴与
$IS$-幂挠模范畴之间的范畴等价。
\end{enumerate}
\end{lemma}

\begin{proof}
引理 \ref{more-algebra-lemma-characterize-flatness-on-torsion} 与引理
\ref{more-algebra-lemma-neighbourhood-isomorphism} 的等价条件都成立。
下文将不再说明而直接使用这些条件。
先证明 (1)。设 $M$ 为任意 $R$-模。置
$M' = M/M[I^\infty]$，并考虑正合列
$$
0 \to M[I^\infty] \to M \to M' \to 0
$$
由于 $M[I^\infty] = M[I^\infty] \otimes_R S$，只需证明
$(M' \otimes_R S)[(IS)^\infty] = 0$。
写成 $I = (f_1, \ldots, f_t)$。由引理
\ref{more-algebra-lemma-divide-by-torsion}，有 $M'[I^\infty] = 0$。
故对每个 $n > 0$，映射
$$
M' \longrightarrow \bigoplus\nolimits_{i = 1, \ldots t} M',
\quad
x \longmapsto (f_1^n x, \ldots, f_t^n x)
$$
是单射。由于 $S$ 在 $R$ 上平坦，相应的映射
$M' \otimes_R S \to \bigoplus_{i = 1, \ldots t} M' \otimes_R S$
也为单射。这意味着 $(M' \otimes_R S)[I^n] = 0$，如所需。

\medskip\noindent
接着证明 (2)。若 $N$ 是 $I$-幂挠的，则
$N \otimes_R S = N$，且由 Algebra，引理
\ref{algebra-lemma-adjoint-tensor-restrict}，(2) 中所显示的映射是同构。
若 $M$ 是 $I$-幂挠的，则任意映射 $M \to N$ 的像都经由
$N[I^\infty]$ 分解，而任意映射
$M \otimes_R S \to N \otimes_R S$ 的像都经由
$(N \otimes_R S)[(IS)^\infty]$ 分解。因此在此情形，(1) 保证我们
可以用 $N[I^\infty]$ 替换 $N$；结论随即归结为刚讨论过的
$N$ 是 $I$-幂挠的情形。

\medskip\noindent
最后证明 (3)。由 (2)，该函子是全忠实的。关于本质满，只需注意：
对任意 $IS$-幂挠 $S$-模 $N$，自然映射 $N \otimes_R S \to N$ 是同构。
\end{proof}

\begin{lemma}
\label{more-algebra-lemma-map-identifies-koszul-and-cech-complexes}
设 $\varphi : R \to S$ 为平坦环映射，$I \subset R$ 为有限生成理想，
且 $R/I \to S/IS$ 是同构。对任意满足
$V(f_1, \ldots, f_r) = V(I)$ 的 $f_1, \ldots, f_r \in R$，有：
\begin{enumerate}
\item Koszul 复形的映射
$K(R, f_1, \ldots, f_r) \to K(S, f_1, \ldots, f_r)$ 是拟同构；
\item 扩充交错 {\v C}ech 复形的映射
$$
\xymatrix{
R \to \prod_{i_0} R_{f_{i_0}} \to \prod_{i_0 < i_1} R_{f_{i_0}f_{i_1}}
\to \ldots \to R_{f_1\ldots f_r} \ar[d] \\
S \to \prod_{i_0} S_{f_{i_0}} \to \prod_{i_0 < i_1} S_{f_{i_0}f_{i_1}}
\to \ldots \to S_{f_1\ldots f_r}
}
$$
是拟同构。
\end{enumerate}
\end{lemma}

\begin{proof}
在两种情形下都有一个由 $R$ 模组成的复形 $K_\bullet$，并且要证明
$K_\bullet \to K_\bullet \otimes_R S$ 是拟同构。
% P02-E0380
由引理 \ref{more-algebra-lemma-neighbourhood-isomorphism} 及 $R \to S$ 的平坦性，
只要 $K$ 的所有同调群都是 $I$-幂挠的，结论即成立。
对 Koszul 复形，这由引理 \ref{more-algebra-lemma-homotopy-koszul} 成立；
对扩充交错 {\v C}ech 复形，这由引理
\ref{more-algebra-lemma-extended-alternating-torsion} 成立。
\end{proof}

\begin{lemma}
\label{more-algebra-lemma-naive-Koszul-complex}
设 $R$ 为环，$I = (f_1, \ldots, f_n)$ 为 $R$ 的有限生成理想。
设 $M$ 为由元素 $e_1, \ldots, e_n$ 生成、以关系
$f_i e_j - f_j e_i = 0$ 为约束的 $R$-模。存在短正合列
$$
0 \to K \to M \to I \to 0
$$
其中 $K$ 被 $I$ 零化。
\end{lemma}

\begin{proof}
这只是 Koszul 复形的一个截断。映射 $M \to I$ 由规则
$e_i \mapsto f_i$ 决定。若 $m = \sum a_i e_i$ 属于 $M \to I$ 的核，
即 $\sum a_i f_i = 0$，则
$f_j m = \sum f_j a_i e_i = (\sum f_i a_i) e_j = 0$。
\end{proof}

\begin{lemma}
\label{more-algebra-lemma-explicit-ext}
设 $R$ 为环，$I = (f_1, \ldots, f_n)$ 为 $R$ 的有限生成理想。
对任意 $R$-模 $N$，置
$$
H_1(N, f_\bullet) =
\frac{\{(x_1, \ldots, x_n) \in N^{\oplus n} \mid f_i x_j = f_j x_i \}}
{\{f_1x, \ldots, f_nx) \mid x \in N\}}
$$
% P02-E0381
对任意 $R$-模 $N$，存在规范短正合列
% P02-E0382
$$
0 \to \Ext_R(R/I, N) \to H_1(N, f_\bullet) \to \Hom_R(K, N)
$$
其中 $K$ 如引理 \ref{more-algebra-lemma-naive-Koszul-complex} 中所述。
\end{lemma}

\begin{proof}
上面的记号表示 $\text{Mod}_R$ 中的 $\Ext$-群，其定义见
Homology，第 \ref{homology-section-extensions} 节。
它们记作 $\Ext_R(M, N)$。利用短正合列
$0 \to I \to R \to R/I \to 0$ 所对应的
Homology，引理 \ref{homology-lemma-six-term-sequence-ext} 的长正合列，
以及 $\Ext_R(R, N) = 0$，可得
$$
\Ext_R(R/I, N) =
\Coker(N \longrightarrow \Hom(I, N))
$$
利用引理 \ref{more-algebra-lemma-naive-Koszul-complex} 的短正合列，可得到复形
$$
N \to \Hom(M, N) \to \Hom_R(K, N)
$$
其中间同调规范同构于 $\Ext_R(R/I, N)$。由于第一个映射的余核
规范同构于 $H_1(N, f_\bullet)$，引理得证。
\end{proof}

\begin{lemma}
\label{more-algebra-lemma-koszul-homology-annihilated}
设 $R$ 为环，$I = (f_1, \ldots, f_n)$ 为 $R$ 的有限生成理想。
对任意 $R$-模 $N$，引理 \ref{more-algebra-lemma-explicit-ext} 中定义的 Koszul 同调群
$H_1(N, f_\bullet)$ 被 $I$ 零化。
\end{lemma}

\begin{proof}
设 $(x_1, \ldots, x_n) \in N^{\oplus n}$，且 $f_i x_j = f_j x_i$。
则 $f_i(x_1, \ldots, x_n) = (f_i x_i, \ldots, f_i x_n)$。
% P02-E0383
换言之，$f_i$ 零化 $H_1(N, f_\bullet)$。
\end{proof}

\noindent
我们可以加强引理 \ref{more-algebra-lemma-neighbourhood-equivalence} 的全忠实性：
若 $\Ext$-群的源是 $I$-幂挠的，则传到 $S$ 同样不改变该群。
下一个引理的导出版本见 Dualizing Complexes，引理
\ref{dualizing-lemma-neighbourhood-extensions}。

\begin{lemma}
\label{more-algebra-lemma-neighbourhood-extensions}
假设 $\varphi : R \to S$ 为平坦环映射，$I \subset R$ 为有限生成理想，
且 $R/I \to S/IS$ 是同构。设 $M$、$N$ 为 $R$-模，并假设 $M$ 是
$I$-幂挠的。给定短正合列
% P02-E0384
$$
0 \to N \otimes_R S \to \tilde E \to M \otimes_R S \to 0
$$
则存在交换图
$$
\xymatrix{
0 \ar[r] &
N \ar[r] \ar[d] &
E \ar[r] \ar[d] &
M \ar[r] \ar[d] &
0 \\
0 \ar[r] &
N \otimes_R S \ar[r] &
\tilde E \ar[r] &
M \otimes_R S \ar[r] &
0
}
$$
且两行均正合。
\end{lemma}

\begin{proof}
因为 $M$ 是 $I$-幂挠的，由引理
\ref{more-algebra-lemma-neighbourhood-isomorphism} 可知 $M \otimes_R S = M$。
下文将不再说明而使用这一等同。由于 $R \to S$ 平坦，基变换函子正合，
从而得到 $\Ext$-群的函子性映射
$$
\Ext_R(M, N)
\longrightarrow
\Ext_S(M \otimes_R S, N \otimes_R S),
$$
参见 Homology，引理 \ref{homology-lemma-exact-functor-ext}。
引理的断言是：当 $M$ 是 $I$-幂挠的时，此映射满射。
事实上，我们将证明它是同构。由引理
\ref{more-algebra-lemma-I-power-torsion-presentation}，可以找到满射 $M' \to M$，
其中 $M'$ 是形如 $R/I^n$ 的模的直和。利用 Homology，引理
\ref{homology-lemma-six-term-sequence-ext} 的长正合列以及引理
\ref{more-algebra-lemma-neighbourhood-equivalence}，可知只需对 $M'$ 证明该引理。
利用 $\Ext$ 与直和的相容性（略去细节），可约化到
$M = R/I^n$ 的情形，其中 $n$ 为某个整数。

\medskip\noindent
设 $f_1, \ldots, f_t$ 为 $I^n$ 的生成元。由引理
\ref{more-algebra-lemma-explicit-ext}，有交换图
$$
\xymatrix{
0 \ar[r] &
\Ext_R(R/I^n, N) \ar[r] \ar[d] &
H_1(N, f_\bullet) \ar[r] \ar[d] &
\Hom_R(K, N) \ar[d] \\
0 \ar[r] &
\Ext_S(S/I^nS, N \otimes S) \ar[r] &
H_1(N \otimes S, f_\bullet) \ar[r] &
\Hom_S(K \otimes S, N \otimes S)
}
$$
% P02-E0385
两行正合，其中 $K$ 如引理
\ref{more-algebra-lemma-naive-Koszul-complex} 中所述。
因此只需证明右侧两个竖直箭头是同构。由于 $K$ 被 $I^n$ 零化，
由引理 \ref{more-algebra-lemma-neighbourhood-equivalence}，有
$\Hom_R(K, N) = \Hom_S(K \otimes_R S, N \otimes_R S)$。
由于 $R \to S$ 平坦，有
$H_1(N, f_\bullet) \otimes_R S = H_1(N \otimes_R S, f_\bullet)$。
由引理 \ref{more-algebra-lemma-koszul-homology-annihilated}，
$H_1(N, f_\bullet)$ 被 $I^n$ 零化，故由引理
\ref{more-algebra-lemma-neighbourhood-isomorphism}，
$H_1(N, f_\bullet) \otimes_R S = H_1(N, f_\bullet)$。
\end{proof}

\noindent
设 $R \to S$ 为环映射。设 $f_1, \ldots, f_t \in R$，并置
$I = (f_1, \ldots, f_t)$。于是对任意 $R$-模 $M$，可定义复形
\begin{equation}
\label{more-algebra-equation-glueing-complex}
0 \to M \xrightarrow{\alpha}
M \otimes_R S \times \prod M_{f_i} \xrightarrow{\beta}
\prod (M \otimes_R S)_{f_i}
\times
\prod M_{f_if_j}
\end{equation}
其中 $\alpha(m) = (m \otimes 1, m/1, \ldots, m/1)$，且
$$
\beta(m', m_1, \ldots, m_t) =
((m'/1 - m_1 \otimes 1, \ldots, m'/1 - m_t \otimes 1),
(m_1 - m_2, \ldots, m_{t - 1} - m_t).
$$
% P02-E0386
% P02-E0387
我们希望知道此复形何时正合。

\begin{lemma}
\label{more-algebra-lemma-recover-module-from-glueing-data}
假设 $\varphi : R \to S$ 为平坦环映射，且
$I = (f_1, \ldots, f_t) \subset R$ 是满足
$R/I \to S/IS$ 为同构的理想。设 $M$ 为 $R$-模。
则复形 (\ref{more-algebra-equation-glueing-complex}) 正合。
\end{lemma}

\begin{proof}
第一种证明。用
$\check{\mathcal{C}}_R \to \check{\mathcal{C}}_S$ 表示引理
\ref{more-algebra-lemma-map-identifies-koszul-and-cech-complexes} 中扩充交错
{\v C}ech 复形的拟同构。由于这些复形有界且各项平坦，
$M \otimes_R \check{\mathcal{C}}_R \to M \otimes_R \check{\mathcal{C}}_S$
也是拟同构（引理 \ref{more-algebra-lemma-derived-tor-quasi-isomorphism} 与
\ref{more-algebra-lemma-derived-tor-quasi-isomorphism-other-side}）。
复形 (\ref{more-algebra-equation-glueing-complex}) 是映射
$M \otimes_R \check{\mathcal{C}}_R \to M \otimes_R \check{\mathcal{C}}_S$
的锥的一个截断，结论即得。

\medskip\noindent
第二种计算性证明。
设 $m \in M$。若 $\alpha(m) = 0$，则由引理
\ref{more-algebra-lemma-torsion-free}，$m \in M[I^\infty]$。选取 $n$ 使
$I^n m = 0$，并考虑映射 $\varphi : R/I^n \to M$。
若 $m \otimes 1 = 0$，则 $\varphi \otimes 1_S = 0$，故
$\varphi = 0$（见引理 \ref{more-algebra-lemma-neighbourhood-equivalence}），
从而 $m = 0$。这说明 $\alpha$ 是单射。

\medskip\noindent
设 $(m', m'_1, \ldots, m'_t) \in \Ker(\beta)$。
对某个 $n > 0$ 及 $m_i \in M$，写成 $m'_i = m_i/f_i^n$。
必要时增大 $n$，可以假设
% P02-E0388
$f_i^n m' = m_i \otimes 1$ 在 $M \otimes_R S$ 中成立，且
$f_j^nm_i - f_i^nm_j = 0$ 在 $M$ 中成立。特别地，
$(m_1, \ldots, m_t)$ 定义了
$H_1(M, (f_1^n, \ldots, f_t^n))$ 的一个元素 $\xi$。
由于 $H_1(M, (f_1^n, \ldots, f_t^n))$ 被 $I^{tn + 1}$ 零化
（见引理 \ref{more-algebra-lemma-koszul-homology-annihilated}），且
$R \to S$ 平坦，由引理
\ref{more-algebra-lemma-neighbourhood-isomorphism} 可得
$$
H_1(M, (f_1^n, \ldots, f_t^n)) =
H_1(M, (f_1^n, \ldots, f_t^n)) \otimes_R S =
H_1(M \otimes_R S, (f_1^n, \ldots, f_t^n))
$$
% P02-E0389
$m'$ 的存在性蕴含 $\xi$ 在最后一个群中的像为零，即元素
$\xi$ 为零。因此存在 $m \in M$，使得 $m_i = f_i^n m$。于是对某个
$m'' \in (M \otimes_R S)[(IS)^\infty]$，有
$(m', m'_1, \ldots, m'_t) - \alpha(m)
= (m'', 0, \ldots, 0)$。由引理
\ref{more-algebra-lemma-neighbourhood-equivalence}，可知
$m'' \in M[I^\infty]$，结论即得。
\end{proof}

\begin{remark}
\label{more-algebra-remark-glueing-data}
本评注定义粘合数据的一个范畴。设 $R \to S$ 为环映射。
设 $f_1, \ldots, f_t \in R$，并置 $I = (f_1, \ldots, f_t)$。
范畴 $\text{Glue}(R \to S, f_1, \ldots, f_t)$ 定义如下：
\begin{enumerate}
\item 其对象为系统 $(M', M_i, \alpha_i, \alpha_{ij})$，其中
$M'$ 是 $S$-模，$M_i$ 是 $R_{f_i}$-模，
$\alpha_i : (M')_{f_i} \to M_i \otimes_R S$ 是同构，并且
$\alpha_{ij} : (M_i)_{f_j} \to (M_j)_{f_i}$ 是满足下列条件的同构：
\begin{enumerate}
\item 作为 $(M')_{f_if_j} \to (M_j)_{f_i}$ 的映射，
$\alpha_{ij} \circ \alpha_i = \alpha_j$；
\item 作为 $(M_i)_{f_jf_k} \to (M_k)_{f_if_j}$ 的映射，
$\alpha_{jk} \circ \alpha_{ij} = \alpha_{ik}$（上链条件）。
\end{enumerate}
\item 态射
$(M', M_i, \alpha_i, \alpha_{ij}) \to (N', N_i, \beta_i, \beta_{ij})$
由映射 $\varphi' : M' \to N'$ 与 $\varphi_i : M_i \to N_i$ 给出，
并与给定映射 $\alpha_i, \beta_i, \alpha_{ij}, \beta_{ij}$ 相容。
\end{enumerate}
有规范函子
$$
\text{Can} : \text{Mod}_R
\longrightarrow
\text{Glue}(R \to S, f_1, \ldots, f_t),
\quad
M \longmapsto (M \otimes_R S, M_{f_i}, \text{can}_i, \text{can}_{ij})
$$
其中 $\text{can}_i : (M \otimes_R S)_{f_i} \to M_{f_i} \otimes_R S$
与 $\text{can}_{ij} : (M_{f_i})_{f_j} \to (M_{f_j})_{f_i}$
是规范同构。对范畴
$\text{Glue}(R \to S, f_1, \ldots, f_t)$ 的任意对象
$\mathbf{M} = (M', M_i, \alpha_i, \alpha_{ij})$，定义
$$
H^0(\mathbf{M}) =
\{(m', m_i) \mid \alpha_i(m') = m_i \otimes 1, \alpha_{ij}(m_i) = m_j\}
$$
换言之，它由与 (\ref{more-algebra-equation-glueing-complex}) 类似的正合列
$$
0 \to H^0(\mathbf{M}) \to
M' \times \prod M_i \to
\prod M'_{f_i}
\times
\prod (M_i)_{f_j}
$$
定义。把 $H^0(\mathbf{M})$ 看作 $R$-模。这样也得到函子
$$
H^0 :
\text{Glue}(R \to S, f_1, \ldots, f_t)
\longrightarrow
\text{Mod}_R
$$
接下来的目标是说明：函子 $\text{Can}$ 与 $H^0$ 有时互为拟逆。
\end{remark}

\begin{lemma}
\label{more-algebra-lemma-H0-Can-adjoint}
在评注 \ref{more-algebra-remark-glueing-data} 中，函子
$H^0 : \text{Glue}(R \to S, f_1, \ldots, f_t) \to \text{Mod}_R$
是函子
$\text{Can} : \text{Mod}_R \to \text{Glue}(R \to S, f_1, \ldots, f_t)$
的右伴随。
\end{lemma}

\begin{proof}
设 $\mathbf{M} = (M', M_i, \alpha_i, \alpha_{ij})$ 是
$\text{Glue}(R \to S, f_1, \ldots, f_t)$ 的对象。对任意 $R$-模 $N$，
有映射
$$
\Hom_{\text{Glue}(R \to S, f_1, \ldots, f_t)}(\text{Can}(N), \mathbf{M})
\to
\Hom_R(N, H^0(\mathbf{M}))
$$
它把 $\psi$ 映到 $H^0(\psi)$ 与显然映射
$N \to H^0(\text{Can}(N))$ 的复合。由构造，所显示的映射对
$N = R$ 是同构（即使一般而言 $R \to H^0(\text{Can}(R))$
不是同构）。范畴 $\text{Glue}(R \to S, f_1, \ldots, f_t)$
有直和与余核。函子 $\text{Can}$ 与直和及余核交换。
将 $N$ 写成自由 $R$-模之间一个映射的余核，由这些观察可知
所显示的映射是双射。略去细节。
\end{proof}

\begin{lemma}
\label{more-algebra-lemma-H0-inverse}
假设 $\varphi : R \to S$ 为平坦环映射，且
$I = (f_1, \ldots, f_t) \subset R$ 是满足
$R/I \to S/IS$ 为同构的理想。则函子 $H^0$ 是评注
\ref{more-algebra-remark-glueing-data} 中函子 $\text{Can}$ 的左拟逆。
\end{lemma}

\begin{proof}
这是引理 \ref{more-algebra-lemma-recover-module-from-glueing-data} 的改述。
\end{proof}

\begin{lemma}
\label{more-algebra-lemma-exact}
假设 $\varphi : R \to S$ 为平坦环映射，并设
$I = (f_1, \ldots, f_t) \subset R$ 为理想。
则 $\text{Glue}(R \to S, f_1, \ldots, f_t)$ 是 Abel 范畴，
且函子 $\text{Can}$ 正合并与任意余极限交换。
\end{lemma}

\begin{proof}
给定范畴 $\text{Glue}(R \to S, f_1, \ldots, f_t)$ 中的态射
$$
(\varphi', \varphi_i) :
(M', M_i, \alpha_i, \alpha_{ij})
\to
(N', N_i, \beta_i, \beta_{ij})
$$
其核存在，且等于对象
$(\Ker(\varphi'), \Ker(\varphi_i), \alpha_i, \alpha_{ij})$；
其余核也存在，且等于对象
$(\Coker(\varphi'), \Coker(\varphi_i), \beta_i, \beta_{ij})$。
这是因为 $R \to S$ 平坦，故取核或余核与 $- \otimes_R S$ 交换。
略去细节。正合性来自 $R_{f_i}$ 与 $S$ 在 $R$ 上平坦；
与余极限交换则来自张量积与余极限交换。
\end{proof}

\begin{lemma}
\label{more-algebra-lemma-equivalence-I-unit}
设 $\varphi : R \to S$ 为平坦环映射，且
$(f_1, \ldots, f_t) = R$。则 $\text{Can}$ 与 $H^0$ 是互为拟逆的
范畴等价
$$
\text{Mod}_R = \text{Glue}(R \to S, f_1, \ldots, f_t)
$$
\end{lemma}

\begin{proof}
考虑 $\text{Glue}(R \to S, f_1, \ldots, f_t)$ 的对象
$\mathbf{M} = (M', M_i, \alpha_i, \alpha_{ij})$。由 Algebra，引理
\ref{algebra-lemma-glue-modules}，存在唯一的模 $M$ 及同构
$M_{f_i} \to M_i$，它们恢复粘合数据 $\alpha_{ij}$。
于是 $M'$ 与 $M \otimes_R S$ 都是 $S$-模，并且在各 $f_i$ 处局部化后
都恢复模 $M_i \otimes_R S$。因此有规范同构
$M \otimes_R S \to M'$。这说明 $\mathbf{M}$ 属于 $\text{Can}$ 的本质像。
结合引理 \ref{more-algebra-lemma-H0-inverse} 即得本引理。
\end{proof}

\begin{lemma}
\label{more-algebra-lemma-base-change-glue}
设 $\varphi : R \to S$ 为平坦环映射，并设
$I = (f_1, \ldots, f_t)$ 为理想。设 $R \to R'$ 为平坦环映射，并置
$S' = S \otimes_R R'$。
% P02-E0390
则得到如下范畴与函子的交换图：
$$
\xymatrix{
\text{Mod}_R \ar[r]_-{\text{Can}} \ar[d]_{-\otimes_R R'} &
\text{Glue}(R \to S, f_1, \ldots, f_t) \ar[r]_-{H^0} \ar[d]^{-\otimes_R R'} &
\text{Mod}_R \ar[d]^{-\otimes_R R'} \\
\text{Mod}_{R'} \ar[r]^-{\text{Can}} &
\text{Glue}(R' \to S', f_1, \ldots, f_t) \ar[r]^-{H^0} &
\text{Mod}_{R'}
}
$$
\end{lemma}

\begin{proof}
略。
\end{proof}

\begin{proposition}
\label{more-algebra-proposition-equivalence}
假设 $\varphi : R \to S$ 为平坦环映射，且
$I = (f_1, \ldots, f_t) \subset R$ 是满足
$R/I \to S/IS$ 为同构的理想。则 $\text{Can}$ 与 $H^0$
是互为拟逆的范畴等价
$$
\text{Mod}_R = \text{Glue}(R \to S, f_1, \ldots, f_t)
$$
\end{proposition}

\begin{proof}
我们已经看到 $H^0 \circ \text{Can}$ 同构于恒等函子；见引理
\ref{more-algebra-lemma-H0-inverse}。考虑
$\text{Glue}(R \to S, f_1, \ldots, f_t)$ 的对象
$\mathbf{M} = (M', M_i, \alpha_i, \alpha_{ij})$。
得到自然态射
$$
\Psi :
(H^0(\mathbf{M}) \otimes_R S, H^0(\mathbf{M})_{f_i},
\text{can}_i, \text{can}_{ij})
\longrightarrow
(M', M_i, \alpha_i, \alpha_{ij}).
$$
具体而言，按定义，$H^0(\mathbf{M})$ 带有相容的 $R$-模映射
$H^0(\mathbf{M}) \to M'$ 与 $H^0(\mathbf{M}) \to M_i$。
必须证明此映射是同构。

\medskip\noindent
选取指标 $i$ 并置 $R' = R_{f_i}$。结合引理
\ref{more-algebra-lemma-base-change-glue} 与 \ref{more-algebra-lemma-equivalence-I-unit}，
可知 $\Psi \otimes_R R'$ 是同构。因此 $\Psi$ 的核与余核分别是
形如 $(K, 0, 0, 0)$ 与 $(Q, 0, 0, 0)$ 的系统。注意
$H^0((K, 0, 0, 0)) = K$，函子 $H^0$ 左正合，并且由构造
$H^0(\Psi)$ 是双射。因此 $K = 0$，即 $\Psi$ 的核为零。

\medskip\noindent
上述结论给出短正合列
$$
0 \to H^0(\mathbf{M}) \otimes_R S \to M' \to Q \to 0
$$
并且 $M_i = H^0(\mathbf{M})_{f_i}$。可以把 $Q$ 看作 $I$-幂挠
$R$-模，因而 $Q = Q \otimes_R S$。由引理
\ref{more-algebra-lemma-neighbourhood-extensions}，存在交换图
$$
\xymatrix{
0 \ar[r] &
H^0(\mathbf{M}) \ar[r] \ar[d] &
E \ar[r] \ar[d] &
Q \ar[r] \ar[d] &
0 \\
0 \ar[r] &
H^0(\mathbf{M}) \otimes_R S \ar[r] &
M' \ar[r] &
Q \ar[r] &
0
}
$$
且两行正合。这显然在范畴
$\text{Glue}(R \to S, f_1, \ldots, f_t)$ 中确定同构
$\text{Can}(E) \to (M', M_i, \alpha_i, \alpha_{ij})$，结论即得。
（当然，事后可知 $Q = 0$。）
\end{proof}

\begin{lemma}
\label{more-algebra-lemma-application-formal-glueing}
设 $\varphi : R \to S$ 为平坦环映射，$I \subset R$ 为有限生成理想，
且 $R/I \to S/IS$ 是同构。
% P02-E0391
\begin{enumerate}
\item 给定 $R$-模 $N$、$S$-模 $M'$ 以及核和余核均为 $I$-幂挠的
$S$-模映射 $\varphi : M' \to N \otimes_R S$，则存在 $R$-模映射
$\psi : M \to N$ 及与 $\varphi$、$\psi$ 相容的同构
$M \otimes_R S = M'$。
\item 给定 $R$-模 $M$、$S$-模 $N'$ 以及核和余核均为 $I$-幂挠的
$S$-模映射 $\varphi : M \otimes_R S \to N'$，则存在 $R$-模映射
$\psi : M \to N$ 及与 $\varphi$、$\psi$ 相容的同构
$N \otimes_R S = N'$。
\end{enumerate}
在两种情形下均有 $\Ker(\varphi) \cong \Ker(\psi)$ 及
$\Coker(\varphi) \cong \Coker(\psi)$。
\end{lemma}

\begin{proof}
(1) 的证明。设 $I = (f_1, \ldots, f_t)$。显然，局部化
$\varphi_{f_i}$ 是同构。因此
$(M', N_{f_i}, \varphi_{f_i}, can_{ij})$ 是
$\text{Glue}(R \to S, f_1, \ldots, f_t)$ 的对象；见评注
\ref{more-algebra-remark-glueing-data}。
% P02-E0392
由命题 \ref{more-algebra-proposition-equivalence}，存在 $R$-模 $M$，使得
$M' = M \otimes_R S$ 及 $N_{f_i} = M_{f_i}$ 与同构
$\varphi_{f_i}$、$can_{ij}$ 相容。有
$$
(M \otimes_R S, M_{f_i}, can_i, can_{ij}) =
(M', N_{f_i}, \varphi_{f_i}, can_{ij})
\to
(N \otimes_R S, N_{f_i}, can_i, can_{ij})
$$
这一 $\text{Glue}(R \to S, f_1, \ldots, f_t)$ 中的态射，
其第一个分量使用 $\varphi$。由命题
\ref{more-algebra-proposition-equivalence} 的范畴等价，它对应于 $R$-模映射
$\psi : M \to N$。$M \to N$ 的基变换与同构
$M' \cong M \otimes_R S$ 的复合是 $\varphi$；换言之，
$M \to N$ 与 $\varphi$ 相容。

\medskip\noindent
(2) 的证明。这只是上述论证的对偶。具体而言，局部化
$\varphi_{f_i}$ 是同构。因此
$(N', M_{f_i}, \varphi_{f_i}^{-1}, can_{ij})$ 是
$\text{Glue}(R \to S, f_1, \ldots, f_t)$ 的对象；见评注
\ref{more-algebra-remark-glueing-data}。由命题
\ref{more-algebra-proposition-equivalence}，存在 $R$-模 $N$，使得
$N' = N \otimes_R S$ 及 $N_{f_i} = M_{f_i}$ 与同构
$\varphi_{f_i}^{-1}$、$can_{ij}$ 相容。有态射
$$
(M \otimes_R S, M_{f_i}, can_i, can_{ij}) \to
(N', M_{f_i}, \varphi_{f_i}, can_{ij}) =
(N \otimes_R S, N_{f_i}, can_i, can_{ij})
$$
% P02-E0393
这一 $\text{Glue}(R \to S, f_1, \ldots, f_t)$ 中的态射，
其第一个分量使用 $\varphi$。由命题
\ref{more-algebra-proposition-equivalence} 的范畴等价，它对应于 $R$-模映射
$\psi : M \to N$。$M \to N$ 的基变换与同构
$N' \cong N \otimes_R S$ 的复合是 $\varphi$；换言之，
$M \to N$ 与 $\varphi$ 相容。

\medskip\noindent
最后一个陈述例如可由引理
\ref{more-algebra-lemma-neighbourhood-equivalence} 得到。
\end{proof}

\noindent
接下来将命题 \ref{more-algebra-proposition-equivalence} 特化，以得到更便于使用的形式。
具体而言，若 $I = (f)$ 是主理想，则
$\text{Glue}(R \to S, f)$ 的对象不过是三元组
$(M', M_1, \alpha_1)$，而且{\it 无需}检验上链条件！

\begin{theorem}
\label{more-algebra-theorem-formal-glueing}
设 $R$ 为环，$f \in R$。设 $\varphi : R \to S$ 为平坦环映射，
并诱导同构 $R/fR \to S/fS$。则函子
$$
\text{Mod}_R
\longrightarrow
\text{Mod}_S \times_{\text{Mod}_{S_f}} \text{Mod}_{R_f},
\quad
M
\longmapsto
(M \otimes_R S, M_f, \text{can})
$$
是等价。
\end{theorem}

\begin{proof}
箭头右侧的范畴是三元组 $(M', M_1, \alpha_1)$ 组成的范畴，其中
$M'$ 是 $S$-模，$M_1$ 是 $R_f$-模，并且
$\alpha_1 : M'_f \to M_1 \otimes_R S$ 是 $S_f$-同构；见
Categories，例 \ref{categories-example-2-fibre-product-categories}。
% P02-E0394
故该定理是命题 \ref{more-algebra-proposition-equivalence} 的特例。
\end{proof}

\noindent
定理 \ref{more-algebra-theorem-formal-glueing} 的一个有用特例是：$R$ 为 Noether 环，
而 $S$ 是 $R$ 关于元素 $f$ 的完备化。完备化映射 $R \to S$ 平坦；
当 $M$ 有限生成时，函子 $M \mapsto M \otimes_R S$ 可等同于
$f$-进完备化函子。为更精确地陈述这一点，用
$\text{Mod}^{fg}_R$ 表示有限生成 $R$-模的范畴。

\begin{proposition}
\label{more-algebra-proposition-formal-glueing}
设 $R$ 为 Noether 环，$f \in R$ 为元素，并设 $R^\wedge$ 为
$R$ 的 $f$-进完备化。则函子
$M \mapsto (M^\wedge, M_f, \text{can})$ 定义等价
$$
\text{Mod}^{fg}_R
\longrightarrow
\text{Mod}^{fg}_{R^\wedge}
\times_{\text{Mod}^{fg}_{(R^\wedge)_f}}
\text{Mod}^{fg}_{R_f}
$$
\end{proposition}

\begin{proof}
由 Algebra，引理 \ref{algebra-lemma-completion-flat}，环映射
$R \to R^\wedge$ 平坦。显然 $R/fR = R^\wedge/fR^\wedge$。
由 Algebra，引理 \ref{algebra-lemma-completion-tensor}，有限
$R$-模 $M$ 的完备化等于 $M \otimes_R R^\wedge$。
故命题所显示的函子等于定理 \ref{more-algebra-theorem-formal-glueing} 中的函子，
特别地它全忠实。设 $(M_1, M_2, \psi)$ 为右侧的对象。
由定理 \ref{more-algebra-theorem-formal-glueing}，存在 $R$-模 $M$，使得
$M_1 = M \otimes_R R^\wedge$ 及 $M_2 = M_f$。
由于 $R \to R^\wedge \times R_f$ 忠实平坦，由 Algebra，引理
\ref{algebra-lemma-cover} 可知 $M$ 有限生成，即
$M \in \text{Mod}^{fg}_R$。命题得证。
\end{proof}

\begin{remark}
\label{more-algebra-remark-formal-glueing-algebras}
命题 \ref{more-algebra-proposition-equivalence}、定理
\ref{more-algebra-theorem-formal-glueing} 及命题
\ref{more-algebra-proposition-formal-glueing} 的等价保持模的性质。例如，若
$M$ 对应于 $\mathbf{M} = (M', M_i, \alpha_i, \alpha_{ij})$，
则 $M$ 在 $R$ 上有限、有限表现、平坦或投射，当且仅当
$M'$ 与 $M_i$ 分别在 $S$ 与 $R_{f_i}$ 上具有相应性质。
这是因为 $R \to S \times \prod R_{f_i}$ 忠实平坦，并且这些性质
可沿忠实平坦映射下降与上升；见 Algebra，引理
\ref{algebra-lemma-descend-properties-modules} 及定理
\ref{algebra-theorem-ffdescent-projectivity}。
% P02-E0395
这些函子还保持两侧的 $\otimes$-结构。因此，它定义由
% P02-E0396
偶 $(\text{Mod}_R, \otimes)$ 构造的各类范畴之间的等价，
例如代数范畴之间的等价。
\end{remark}

\begin{remark}
\label{more-algebra-remark-topological-analogue}
给定微分流形 $X$ 及紧闭子流形 $Z$，其补集为 $U$，指定 $X$ 上的层
等价于指定 $U$ 上的层、$Z$ 在 $X$ 中某个未指定的管状邻域 $T$ 上的层，
以及这两个层在 $T \cap U$ 上所得层之间的同构。
% P02-E0397
代数几何中并不存在这种意义下的管状邻域，但命题
\ref{more-algebra-proposition-equivalence}、定理
\ref{more-algebra-theorem-formal-glueing} 及命题
\ref{more-algebra-proposition-formal-glueing} 等结果允许我们改用形式邻域。
\end{remark}

\section{Beauville--Laszlo 定理}
\label{more-algebra-section-beauville-laszlo}

\noindent
设 $R$ 为环，$f$ 为 $R$ 的元素。用
$R^\wedge = \lim R/f^n R$ 表示 $R$ 的 $f$-进完备化。
% P02-E0398
本节讨论并稍作推广 Beauville 与 Laszlo 的一个定理；见
\cite{Beauville-Laszlo}。该定理断言：在适当条件下，可以把
$R^\wedge$ 上的模与 $R_f$ 上的模，沿着它们基扩张到
$(R^\wedge)_f$ 后所得模之间的一个同构“粘合起来”，从而构造
$R$ 上的模。

\medskip\noindent
\cite{Beauville-Laszlo} 假设 $f$ 在 $R$ 与 $M$ 上都是非零因子。
事实上，只需假设
$$
R[f^\infty] \longrightarrow R^\wedge[f^\infty]
$$
是双射，并且
$$
M[f^\infty] \longrightarrow M \otimes_R R^\wedge
$$
是单射。这一优化部分受到 \cite[\S 1.3]{Kedlaya-Liu-I}
为非 Archimedes 解析空间理论而引入的另一种粘合方法的启发。

\medskip\noindent
事实上，我们将在更一般的环映射
$$
R \longrightarrow R'
$$
的框架下建立 Beauville--Laszlo 定理；它对每个 $n > 0$ 都诱导同构
$R/f^nR \to R'/f^nR'$，并诱导同构
$R[f^\infty] \to R'[f^\infty]$。这个框架更适合整体化，而且不能从
$R'$ 是 $R$ 的完备化的情形形式地推出。例如，条件
$R[f^\infty] \to R'[f^\infty]$ 为双射并不蕴含
$R[f^\infty] \to R^\wedge[f^\infty]$ 为双射。

\medskip\noindent
本节所证明的 Beauville 与 Laszlo 定理可看作定理
\ref{more-algebra-theorem-formal-glueing} 的非平坦版本；当
$R' = R^\wedge$ 时，也可看作命题
\ref{more-algebra-proposition-formal-glueing} 的非 Noether 版本。
它与平坦下降的比较见评注 \ref{more-algebra-remark-not-descent}。

\medskip\noindent
还可以建立更强的结果（例如不对 $M$ 施加限制），但为此必须在
导出范畴层面工作。详见 \cite[\S 5]{Bhatt-Algebraize}。

\begin{lemma}
\label{more-algebra-lemma-same-quotients}
设 $R$ 为环，$f \in R$。对每个正整数 $n$，映射
$R/f^nR \to R^\wedge/f^n R^\wedge$ 是同构。
\end{lemma}

\begin{proof}
这是 Algebra，引理
\ref{algebra-lemma-hathat-finitely-generated} 的特例。
\end{proof}

\noindent
我们将使用第 \ref{more-algebra-section-torsion} 节引入的记号。因此，对
$R$-模 $M$，用 $M[f^n]$ 表示 $M$ 中被 $f^n$ 零化的子模，并置
% P02-E0399
$$
M[f^\infty] = \bigcup\nolimits_{n = 1}^\infty M[f^n] = \Ker(M \to M_f).
$$
若 $M = M[f^\infty]$，则称 $M$ 为 $f$-幂挠模。

\begin{lemma}
\label{more-algebra-lemma-BL-faithful}
设 $R$ 为环，$f \in R$，且 $R \to R'$ 为对 $n > 0$ 诱导同构
$R/f^nR \to R'/f^nR'$ 的环映射。$R$-模 $R' \oplus R_f$ 是忠实的：
对每个非零 $R$-模 $M$，模 $M \otimes_R (R' \oplus R_f)$ 也非零。
例如，若 $M$ 非零，则 $M \otimes_R (R^\wedge \oplus R_f)$ 非零。
\end{lemma}

\noindent
然而，映射 $M \to M \otimes_R (R' \oplus R_f)$ 未必是单射；
见例 \ref{more-algebra-example-not-glueing-pair}。

\begin{proof}
若 $M \neq 0$ 但 $M \otimes_R R_f = 0$，则 $M$ 是 $f$-幂挠的。
由引理 \ref{more-algebra-lemma-torsion-completion}，
$M \otimes_R R' \cong M \neq 0$。由引理
\ref{more-algebra-lemma-same-quotients}，最后一个陈述是第一个陈述的特例。
\end{proof}

\begin{lemma}
\label{more-algebra-lemma-cover-spec}
设 $R$ 为环，$f \in R$，且 $R \to R'$ 为诱导同构
$R/fR \to R'/fR'$ 的环映射。映射
$\Spec(R') \amalg \Spec(R_f) \to \Spec(R)$ 是满射。
例如，映射
$\Spec(R^\wedge) \amalg \Spec(R_f) \to \Spec(R)$ 是满射。
\end{lemma}

\begin{proof}
回忆 $\Spec(R) = V(f) \amalg D(f)$，其中
$V(f) = \Spec(R/fR)$ 且 $D(f) = \Spec(R_f)$；见
Algebra，第 \ref{algebra-section-spectrum-ring} 节，尤其是引理
\ref{algebra-lemma-spec-closed} 与
\ref{algebra-lemma-standard-open}。由于映射 $R \to R/fR$
经由 $R'$ 分解，引理成立。由引理
\ref{more-algebra-lemma-same-quotients}，最后一个陈述是第一个陈述的特例。
\end{proof}

\begin{lemma}
\label{more-algebra-lemma-faithful-descent}
\begin{reference}
\cite[Lemme~2(a)]{Beauville-Laszlo} 的轻微推广。
\end{reference}
设 $R$ 为环，$f \in R$，且 $R \to R'$ 为对 $n > 0$ 诱导同构
$R/f^nR \to R'/f^nR'$ 的环映射。$R$-模 $M$ 有限生成，当且仅当
（$R' \oplus R_f$）-模 $M \otimes_R (R' \oplus R_f)$ 有限生成。
例如，若 $M \otimes_R (R^\wedge \oplus R_f)$ 作为
$R^\wedge \oplus R_f$ 上的模有限生成，则 $M$ 是有限生成 $R$-模。
\end{lemma}

\begin{proof}
“仅当”方向是清楚的，故假设
$M \otimes_R (R' \oplus R_f)$ 有限生成。将每个生成元写成纯张量之和，
可知 $M \otimes_R (R' \oplus R_f)$ 有一个由 $M$ 中元素组成的
有限生成集。换言之，存在从有限自由 $R$-模到 $M$ 的态射，其余核
在与 $R' \oplus R_f$ 作张量后被零化；对该余核应用引理
\ref{more-algebra-lemma-BL-faithful}，可推出 $M$ 有限生成。
% P02-E0400
由引理 \ref{more-algebra-lemma-same-quotients}，最后一个陈述是第一个陈述的特例。
\end{proof}

\begin{remark}
\label{more-algebra-remark-not-descent}
虽然 $R \to R_f$ 总是平坦的，但除非 $R$ 是 Noether 环，
$R \to R^\wedge$ 通常并不平坦（见 Algebra，引理
\ref{algebra-lemma-completion-flat} 以及 Examples，第
\ref{examples-section-nonflat} 节中的讨论）。因此，一般不能把
Descent，第 \ref{descent-section-descent-modules} 节所讨论的忠实平坦下降
应用于态射 $R \to R^\wedge \oplus R_f$。此外，即使在 Noether 情形，
该态射的通常下降数据定义也涉及环
$R^\wedge \otimes_R R^\wedge$，而本节将避免考虑此环。
\end{remark}

\noindent
{\bf 粘合对。}设 $R \to R'$ 为对 $n > 0$ 诱导同构
$R/f^nR \to R'/f^nR'$ 的环映射。考虑序列
\begin{equation}
\label{more-algebra-equation-BL-cech-re}
0 \to R \to R' \oplus R_f \to R'_f \to 0,
\end{equation}
其中右侧映射是两个规范同态之差。若此序列正合，则称
$(R \to R', f)$ 为一个\emph{粘合对}。若
$(R \to R^\wedge, f)$ 是粘合对，则称 $(R, f)$ 为
\emph{粘合对}；由引理 \ref{more-algebra-lemma-same-quotients}，这一定义有意义。
因此，$(R, f)$ 是粘合对，当且仅当序列
\begin{equation}
\label{more-algebra-equation-BL-cech}
0 \to R \to R^\wedge \oplus R_f \to (R^\wedge)_f \to 0,
\end{equation}
正合。

\begin{lemma}
\label{more-algebra-lemma-same-f-torsion}
设 $R$ 为环，$f \in R$，且 $R \to R'$ 为对 $n > 0$ 诱导同构
$R/f^nR \to R'/f^nR'$ 的环映射。序列
(\ref{more-algebra-equation-BL-cech-re})：
\begin{enumerate}
\item 在右端正合；
\item 在左端正合，当且仅当
$R[f^\infty] \to R'[f^\infty]$ 是单射；
\item 在中间正合，当且仅当
$R[f^\infty] \to R'[f^\infty]$ 是满射。
\end{enumerate}
特别地，$(R \to R', f)$ 是粘合对，当且仅当
$R[f^\infty] \to R'[f^\infty]$ 是双射。例如，$(R, f)$ 是粘合对，
当且仅当 $R[f^\infty] \to R^\wedge[f^\infty]$ 是双射。
\end{lemma}

\begin{proof}
设 $x \in R'_f$。写成 $x = x'/f^n$，其中 $x' \in R'$。
再写成 $x' = x'' + f^n y$，其中 $x'' \in R$ 且 $y \in R'$。
于是 $(y, -x''/f^n)$ 映到 $x$，故 (1) 成立。

\medskip\noindent
(2) 来自等式 $\Ker(R \to R_f) = R[f^\infty]$。

\medskip\noindent
若序列在中间正合，则形如 $(x, 0)$、其中
$x \in R'[f^\infty]$ 的元素都在第一个箭头的像中。
这蕴含 $R[f^\infty] \to R'[f^\infty]$ 是满射。
反之，假设 $R[f^\infty] \to R'[f^\infty]$ 满射。设
$(x, y)$ 是映到右端零的中间元素。对某个 $y' \in R$，写成
$y = y'/f^n$。于是 $f^n x - y'$ 在 $R'$ 中被 $f$ 的某个幂零化。
由假设，可对某个 $z \in R[f^\infty]$ 写成
$f^nx - y' = z$。于是 $y = y''/f^n$，其中
$y'' = y' + z$ 属于 $R \to R/f^nR$ 的核。因此对某个
$y''' \in R$，$y$ 可表示为 $y'''/1$。此时
$x - y'''$ 属于 $R'[f^\infty]$，故
$x - y''' = z' \in R[f^\infty]$。于是
$(x, y'''/1) = (y''' + z', (y''' + z')/1)$，如所需。

\medskip\noindent
由引理 \ref{more-algebra-lemma-same-quotients}，引理的最后一个陈述是倒数第二个
陈述的特例。
\end{proof}

\begin{remark}
\label{more-algebra-remark-BL-special-case}
假设 $f$ 是非零因子。则 Algebra，引理
\ref{algebra-lemma-completion-differ-by-torsion} 说明
$f$ 在 $R^\wedge$ 中是非零因子。故 $(R, f)$ 是粘合对。
\end{remark}

\begin{remark}
\label{more-algebra-remark-noetherian-case}
若 $R \to R^\wedge$ 平坦，则对每个正整数 $n$，将序列
$0 \to R[f^n] \to R \to R$ 与 $R^\wedge$ 作张量，得到序列
$0 \to R[f^n] \otimes_R R^\wedge \to R^\wedge \to R^\wedge$。
结合引理 \ref{more-algebra-lemma-torsion-completion}，可得
$R[f^n] \to R^\wedge[f^n]$ 是同构。故 $(R, f)$ 是粘合对。
特别地，当 $R$ 为 Noether 环时这成立；见 Algebra，引理
\ref{algebra-lemma-completion-flat}。
\end{remark}

\begin{example}
\label{more-algebra-example-not-glueing-pair}
设 $k$ 为域，并置
$$
R = k[f, T_1, T_2, \ldots]/(fT_1, fT_2 - T_1, fT_3 - T_2, \ldots).
$$
则 $(R, f)$ 不是粘合对，因为映射
$R[f^\infty] \to R^\wedge[f^\infty]$ 不是单射：$T_1$ 的像在
$R^\wedge$ 中可被 $f$ 整除。对于
$$
R = k[f, T_1, T_2, \ldots]/(fT_1, f^2T_2, \ldots),
$$
映射 $R[f^\infty] \to R^\wedge[f^\infty]$ 不是满射，因为元素
$T_1 + fT_2 + f^2 T_3 + \ldots$ 不在像中。特别地，由评注
\ref{more-algebra-remark-noetherian-case}，这二者都是
$R \to R^\wedge$ 不平坦的例子。
\end{example}

\noindent
{\bf 可粘合模。}
设 $R \to R'$ 为对 $n > 0$ 诱导同构
$R/f^nR \to R'/f^nR'$ 的环映射。对任意 $R$-模 $M$，可将
(\ref{more-algebra-equation-BL-cech-re}) 与 $M$ 作张量，得到序列
\begin{equation}
\label{more-algebra-equation-BL-cech-mod-re}
0 \to M \to (M \otimes_R R') \oplus (M \otimes_R R_f) \to
M \otimes_R R'_f \to 0
\end{equation}
注意 $M \otimes_R R_f = M_f$，并且
$M \otimes_R R'_f = (M \otimes_R R')_f$。
若此序列正合，则称 $M$ \emph{对 $(R \to R', f)$ 可粘合}。
若 $R$ 为环且 $f \in R$，并且 $M$ 对
$(R \to R^\wedge, f)$ 可粘合，则称 $R$-模\emph{可粘合}。
% P02-E0401
因此，$M$ 可粘合，当且仅当序列
\begin{equation}
\label{more-algebra-equation-BL-cech-mod}
0 \to M \to (M \otimes_R R^\wedge) \oplus (M \otimes_R R_f) \to
M \otimes_R (R^\wedge)_f \to 0
\end{equation}
正合。

\begin{lemma}
\label{more-algebra-lemma-same-f-torsion-module}
设 $R$ 为环，$f \in R$，且 $R \to R'$ 为对 $n > 0$ 诱导同构
$R/f^nR \to R'/f^nR'$ 的环映射。序列
(\ref{more-algebra-equation-BL-cech-mod-re})：
\begin{enumerate}
\item 在右端正合；
\item 在左端正合，当且仅当
$M[f^\infty] \to (M \otimes_R R')[f^\infty]$ 是单射；
\item 在中间正合，当且仅当
$M[f^\infty] \to (M \otimes_R R')[f^\infty]$ 是满射。
\end{enumerate}
因此，$M$ 对 $(R \to R', f)$ 可粘合，当且仅当
$M[f^\infty] \to (M \otimes_R R')[f^\infty]$ 是双射。
若 $(R \to R', f)$ 是粘合对，则 $M$ 对 $(R \to R', f)$ 可粘合，
当且仅当 $M[f^\infty] \to (M \otimes_R R')[f^\infty]$ 是单射。
例如，若 $(R, f)$ 是粘合对，则 $M$ 可粘合，当且仅当
$M[f^\infty] \to (M \otimes_R R^\wedge)[f^\infty]$ 是单射。
\end{lemma}

\begin{proof}
下文将不再说明而使用引理 \ref{more-algebra-lemma-same-f-torsion} 的结果。
函子 $M \otimes_R -$ 右正合
（Algebra，引理 \ref{algebra-lemma-tensor-product-exact}），故得到 (1)。

\medskip\noindent
$M \to M \otimes_R R_f = M_f$ 的核为 $M[f^\infty]$，故 (2) 成立。

\medskip\noindent
若序列在中间正合，则形如 $(x, 0)$、其中
$x \in (M \otimes_R R')[f^\infty]$ 的元素都在第一个箭头的像中。
这蕴含
$M[f^\infty] \to (M \otimes_R R')[f^\infty]$ 是满射。
反之，假设
$M[f^\infty] \to (M \otimes_R R')[f^\infty]$ 是满射。
设 $(x, y)$ 为映到右端零的中间元素。
对某个 $y' \in M$，写成 $y = y'/f^n$。于是
$f^n x - y'$ 在 $M \otimes_R R'$ 中被 $f$ 的某个幂零化。
由假设，可对某个 $z \in M[f^\infty]$ 写成
$f^nx - y' = z$。于是 $y = y''/f^n$，其中
$y'' = y' + z$ 属于 $M \to M/f^nM$ 的核。因此对某个
$y''' \in M$，$y$ 可表示为 $y'''/1$。此时
$x - y'''$ 属于 $(M \otimes_R R')[f^\infty]$，故
$x - y''' = z' \in M[f^\infty]$。于是
$(x, y'''/1) = (y''' + z', (y''' + z')/1)$，如所需。

\medskip\noindent
若 $(R \to R', f)$ 是粘合对，则对任意 $M$，序列
(\ref{more-algebra-equation-BL-cech-mod-re}) 由 Algebra，引理
\ref{algebra-lemma-tensor-product-exact} 在中间正合。
这给出引理的倒数第二个陈述。最后一个陈述由此以及如下事实得到：
$(R, f)$ 是粘合对，当且仅当
$(R \to R^\wedge, f)$ 是粘合对。
\end{proof}

\begin{remark}
\label{more-algebra-remark-glueable}
设 $(R \to R', f)$ 为粘合对，$M$ 为 $R$-模。以下观察可用来判定
$M$ 是否对 $(R \to R', f)$ 可粘合。
\begin{enumerate}
\item 由引理 \ref{more-algebra-lemma-same-f-torsion-module}，$M$ 对
$(R \to R^\wedge, f)$ 可粘合，当且仅当
$M[f^\infty] \to M \otimes_R R^\wedge$ 是单射。
若 $M[f] \to M^\wedge$ 是单射，即
$M[f] \cap \bigcap_{n = 1}^\infty f^n M = 0$，则该条件成立。
\item 若 $\text{Tor}_1^R(M, R'_f) = 0$，则 $M$ 对
$(R \to R', f)$ 可粘合
（使用 Algebra，引理 \ref{algebra-lemma-long-exact-sequence-tor}）。
这等价于 $\text{Tor}_1^R(M, R')$ 是 $f$-幂挠的。
特别地，任意平坦 $R$-模都对 $(R \to R', f)$ 可粘合。
\item 若 $R \to R'$ 平坦，则对每个 $R$-模都有
$\text{Tor}_1^R(M, R') = 0$，故每个 $R$-模都对
$(R \to R', f)$ 可粘合。特别地，当 $R$ 为 Noether 环且
$R' = R^\wedge$ 时，该条件成立；见 Algebra，引理
\ref{algebra-lemma-completion-flat}
% P02-E0402
\end{enumerate}
\end{remark}

\begin{example}[不可粘合模]
\label{more-algebra-example-not-glueable-module}
% P02-E0403
\begin{reference}
\cite[\S 4, Remarques]{Beauville-Laszlo}
\end{reference}
设 $R$ 为 $\mathbf{R}$ 上 $0$ 处的 $C^\infty$ 函数芽之环。
设 $f \in R$ 为函数 $f(x) = x$。则 $f$ 在 $R$ 中是非零因子，
故 $(R, f)$ 是粘合对，且 $R^\wedge \cong \mathbf{R}[[x]]$。
设 $\varphi \in R$ 为函数
$\varphi(x) = \text{exp}(-1/x^2)$。
% P02-E0404
则 $\varphi$ 的 Taylor 级数为零，故
$\varphi \in \Ker(R \to R^\wedge)$。由于当 $x \neq 0$ 时
$\varphi(x) \neq 0$，可知 $\varphi$ 在 $R$ 中是非零因子。
函数 $\varphi/f$ 的 Taylor 级数也为零，故它在
$M = R/\varphi R$ 中的像是 $M[f]$ 的非零元素，而该元素映到
$$
M \otimes_R R^\wedge = R^\wedge/\varphi R^\wedge = R^\wedge
$$
中的零。因此 $M$ 不可粘合。
\end{example}

\noindent
接下来计算一些 Tor 群。

\begin{lemma}
\label{more-algebra-lemma-first-tor}
设 $(R \to R', f)$ 为粘合对。则对每个 $n > 0$，
$\text{Tor}^R_1(R', f^n R) = 0$。
\end{lemma}

\begin{proof}
由正合列 $0 \to R[f^n] \to R \to f^n R \to 0$，只需检验
$R[f^n] \otimes_R R' \to R'$ 是单射。由引理
\ref{more-algebra-lemma-torsion-completion}，有
$R[f^n] \otimes_R R' = R[f^n]$；又由引理
\ref{more-algebra-lemma-same-f-torsion} 及 $(R \to R', f)$ 是粘合对，
$R[f^n] \to R'$ 是单射。
\end{proof}

\begin{lemma}
\label{more-algebra-lemma-first-tor-total}
设 $(R \to R',f)$ 为粘合对。
则 $\text{Tor}^R_1(R', R/R[f^\infty]) = 0$。
\end{lemma}

\begin{proof}
有 $R/R[f^\infty] = \colim R/R[f^n] = \colim f^nR$。
由于 Tor 群的形成与滤过余极限交换
（Algebra，引理 \ref{algebra-lemma-tor-commutes-filtered-colimits}），
可应用引理 \ref{more-algebra-lemma-first-tor}。
\end{proof}

\begin{lemma}
\label{more-algebra-lemma-BL3}
\begin{reference}
\cite[Lemme 3(a)]{Beauville-Laszlo} 的轻微推广。
\end{reference}
设 $(R \to R', f)$ 为粘合对。对每个 $R$-模 $M$，有
$\text{Tor}^R_1(R', \Coker(M \to M_f)) = 0$。
\end{lemma}

\begin{proof}
置 $\overline{M} = M/M[f^\infty]$。则
$\Coker(M \to M_f) \cong \Coker(\overline{M} \to \overline{M}_f)$，
故不妨假设 $f$ 在 $M$ 上是非零因子。此时 $M \subset M_f$，且
$M_f/M = \colim M/f^nM$，其中转移映射由乘以 $f$ 给出。
由于 Tor 群的形成与余极限交换
（Algebra，引理 \ref{algebra-lemma-tor-commutes-filtered-colimits}），
只需证明 $\text{Tor}^R_1(R', M/f^n M) = 0$。

\medskip\noindent
先处理 $M = R/R[f^\infty]$ 的情形。由引理
\ref{more-algebra-lemma-same-f-torsion}，有
$M \otimes_R R' = R'/R'[f^\infty]$。从短正合列
$0 \to M \to M \to M/f^nM \to 0$，由 Algebra，引理
\ref{algebra-lemma-long-exact-sequence-tor} 得到正合列
$$
\xymatrix{
\text{Tor}_1^R(R', R/R[f^\infty]) \ar[r] &
\text{Tor}_1^R(R', M/f^n M) \ar[r] &
R'/R'[f^\infty] \ar[dll]_{f^n} \\
R'/R'[f^\infty] \ar[r] &
(R'/R'[f^\infty])/(f^n
(R'/R'[f^\infty])) \ar[r] & 0
}
$$
这里对角箭头是单射。由引理 \ref{more-algebra-lemma-first-tor-total}，第一个群
$\text{Tor}_1^R(R', R/R[f^\infty])$ 为零，故推出
$\text{Tor}_1^R(R', M/f^nM) = 0$，如所需。

\medskip\noindent
为处理一般情形，选取满射 $F \to M$，其中 $F$ 是自由
$R/R[f^\infty]$-模，并构造正合列
$$
0 \to N \to F/f^n F \to M/f^n M \to 0.
$$
由引理 \ref{more-algebra-lemma-torsion-completion}，与 $R'$ 作张量后此序列不变，
因而仍正合。由上一段，
$\text{Tor}^R_1(R', F/f^n F) = 0$，故推出
$\text{Tor}^R_1(R', M/f^n M) = 0$，如所需。
\end{proof}

\noindent
设 $(R \to R', f)$ 为粘合对。这意味着对 $n > 0$，
$R/f^nR \to R'/f^nR'$ 是同构，且序列
$$
0 \to R \to R' \oplus R_f \to R_f' \to 0
$$
正合。考虑评注 \ref{more-algebra-remark-glueing-data} 中引入的范畴
$\text{Glue}(R \to R', f)$。我们称
$\text{Glue}(R \to R', f)$ 的对象 $(M', M_1, \alpha_1)$ 为一个
\emph{粘合数据}。它由一个 $R'$-模 $M'$、一个 $R_f$-模 $M_1$
及同构 $\alpha_1 : (M')_f \to M_1 \otimes_R R'$ 组成。
有显然的函子
$$
\text{Can} : \text{Mod}_R \longrightarrow \text{Glue}(R \to R', f),\quad
M \longmapsto (M \otimes_R R', M_f, \text{can}),
$$
并且反方向有函子
$$
H^0 : \text{Glue}(R \to R', f) \longrightarrow \text{Mod}_R,\quad
(M', M_1, \alpha_1) \longmapsto \Ker(M' \oplus M_1 \to (M')_f)
$$
其精确定义见评注 \ref{more-algebra-remark-glueing-data}。

\begin{theorem}
\label{more-algebra-theorem-BL-glueing}
\begin{reference}
\cite{Beauville-Laszlo} 主定理的轻微推广。
\end{reference}
设 $(R \to R',f)$ 为粘合对。函子
$\text{Can} : \text{Mod}_R \longrightarrow \text{Glue}(R \to R', f)$
确定如下两个范畴之间的等价：对 $(R \to R', f)$ 可粘合的
$R$-模范畴，与粘合数据的范畴
$\text{Glue}(R \to R', f)$。
% P02-E0407
\end{theorem}

\begin{proof}
设 $(M', M_1, \alpha_1)$ 为粘合数据。我们将证明
$M = H^0((M', M_1, \alpha_1))$ 是对 $(R \to R', f)$ 可粘合的模，
且 $(M', M_1, \alpha_1) \cong \text{Can}(M)$。

\medskip\noindent
先检验函子 $H^0$ 定义中所用的映射
$\text{d} : M' \oplus M_1 \to (M')_f$ 是满射。注意，
$(x, y) \in M' \oplus M_1$ 映到 $(M')_f$ 中的
$\text{d}(x, y) = x/1 - \alpha_1^{-1}(y \otimes 1)$。
若 $z \in (M')_f$，则可写成
$\alpha_1(z) = \sum y_i \otimes g_i$，其中 $g_i \in R'$ 且
$y_i \in M_1$。对某个 $y'_i \in M'$ 与 $n \geq 0$，写成
$\alpha_1^{-1}(y_i \otimes 1) = y_i'/f^n$
（可以对所有 $i$ 选取同一个 $n$）。再写成
$g_i = a_i + f^n b_i$，其中 $a_i \in R$ 且 $b_i \in R'$。
于是取 $y = \sum a_i y_i \in M_1$ 及
$x = \sum b_i y'_i \in M'$，即有
$\text{d}(x, -y) = z$，如所需。

\medskip\noindent
由于 $M = H^0((M', M_1, \alpha_1)) = \Ker(\text{d})$，
得到 $R$-模的正合列
\begin{equation}
\label{more-algebra-equation-define-M}
0 \to M \to M' \oplus M_1 \to (M')_f \to 0.
\end{equation}
我们将证明映射 $M \to M'$ 与 $M \to M_1$ 诱导同构
$M \otimes_R R' \to M'$ 与 $M \otimes_R R_f \to M_1$。
这将蕴含 $M$ 对 $(R \to R', f)$ 可粘合，且
$\text{Can}(M) \cong (M', M_1, \alpha_1)$，如所需。

\medskip\noindent
由于 $f$ 在 $M_1$ 上是非零因子，有
$M[f^\infty] \cong M'[f^\infty]$。这给出正合列
\begin{equation}
\label{more-algebra-equation-exact-mod-torsion}
0 \to M/M[f^\infty] \to M_1 \to (M')_f/M' \to 0.
\end{equation}
由于 $R \to R_f$ 平坦，可将该正合列与 $R_f$ 作张量，推出
$M \otimes_R R_f = (M/M[f^\infty]) \otimes_R R_f \to M_1$
是同构。

\medskip\noindent
由引理 \ref{more-algebra-lemma-BL3}，有
$\text{Tor}_1^R(R', \Coker(M' \to (M')_f)) = 0$。
因此，序列 (\ref{more-algebra-equation-exact-mod-torsion}) 与 $R'$ 在 $R$ 上
作张量后仍正合。利用 $\alpha_1$ 与引理
\ref{more-algebra-lemma-torsion-completion}，所得正合列可写成
\begin{equation}
\label{more-algebra-equation-mod-torsion-sequence}
0 \to (M/M[f^\infty]) \otimes_R R' \to
(M')_f \to (M')_f/M' \to 0
\end{equation}
这给出同构
$(M/M[f^\infty]) \otimes_R R' \cong M'/M'[f^\infty]$。
由此，在图
$$
\xymatrix{
& M[f^\infty] \otimes_R R' \ar[r] \ar[d] &
M \otimes_R R'  \ar[r] \ar[d] &
(M/M[f^\infty]) \otimes_R R' \ar[r] \ar[d] & 0 \\
0 \ar[r] &
M'[f^\infty] \ar[r] &
M' \ar[r] &
M'/M'[f^\infty] \ar[r] & 0,
}
$$
中，第三个竖直箭头是同构。两行均正合；由引理
\ref{more-algebra-lemma-torsion-completion} 及
$M[f^\infty] = M'[f^\infty]$，第一个竖直箭头是同构。
五引理因而蕴含 $M \otimes_R R' \to M'$ 是同构。

\medskip\noindent
以上说明 $\text{Can}$ 本质满，且函子 $H^0$ 的像落在可粘合模范畴中。
由于对可粘合模，(\ref{more-algebra-equation-BL-cech-mod-re}) 正合，所以在可粘合
模范畴上有 $H^0 \circ \text{Can} = \text{id}$。结合引理
\ref{more-algebra-lemma-H0-Can-adjoint} 与 Categories，引理
\ref{categories-lemma-adjoint-fully-faithful}，这蕴含
$\text{Can}$ 全忠实。证明完成。
\end{proof}

\begin{remark}
\label{more-algebra-remark-what-you-get-for-general-modules}
设 $(R \to R', f)$ 为粘合对。设 $M$ 为不一定对
$(R \to R', f)$ 可粘合的 $R$-模。置
$M' = M \otimes_R R'$ 及 $M_1 = M_f$，得到粘合数据
$\text{Can}(M) = (M', M_1, \text{can})$。则
$\tilde M = H^0(M', M_1, \text{can})$ 是对
$(R \to R', f)$ 可粘合的 $R$-模，并且规范映射
$M \to \tilde M$ 给出同构
$M \otimes_R R' \to \tilde M \otimes_R R'$ 及
$M_f \to \tilde M_f$；见定理 \ref{more-algebra-theorem-BL-glueing}。
由序列
$$
M \to (M \otimes_R R' )\oplus M_f \to M \otimes_R (R')_f  \to 0
$$
及
$$
0 \to \tilde M \to (\tilde M \otimes_R R') \oplus \tilde M_f \to
\tilde M \otimes_R  (R')_f \to 0
$$
的正合性，可知映射 $M \to \tilde M$ 是满射。
\end{remark}

\noindent
回忆：粘合对 $(R \to R', f)$ 上的平坦 $R$-模是可粘合的
（评注 \ref{more-algebra-remark-glueable}）。因此，下一个引理说明定理
\ref{more-algebra-theorem-BL-glueing} 确定如下两个范畴之间的等价：
平坦 $R$-模的范畴，与粘合数据 $(M', M_1, \alpha_1)$ 的范畴，
其中 $M'$、$M_1$ 分别在 $R'$、$R_f$ 上平坦。

\begin{lemma}
\label{more-algebra-lemma-BL-flat}
设 $(R \to R', f)$ 为粘合对。设 $M$ 为不一定对
$(R \to R', f)$ 可粘合的 $R$-模。则 $M$ 在 $R$ 上平坦，
当且仅当 $M \otimes_R R'$ 在 $R'$ 上平坦，且
$M_f$ 在 $R_f$ 上平坦。
\end{lemma}

\begin{proof}
引理的一个方向来自 Algebra，引理
\ref{algebra-lemma-flat-base-change}。为证明另一个方向，假设
$M \otimes_R R'$ 在 $R'$ 上平坦，且 $M_f$ 在 $R_f$ 上平坦。
设 $\tilde M$ 如评注 \ref{more-algebra-remark-what-you-get-for-general-modules} 中所述。
若 $\tilde M$ 在 $R$ 上平坦，则对短正合列
$0 \to \Ker(M \to \tilde M) \to M \to \tilde M \to 0$
应用 Algebra，引理 \ref{algebra-lemma-flat-tor-zero}，可知
$\Ker(M \to \tilde M) \otimes_R (R' \oplus R_f)$ 为零。
故由引理 \ref{more-algebra-lemma-BL-faithful}，$M = \tilde M$，从而得到结论。
换言之，可以用 $\tilde M$ 替换 $M$，并假设 $M$ 对
$(R \to R', f)$ 可粘合。设 $N$ 为另一个 $R$-模。
由 Algebra，引理 \ref{algebra-lemma-characterize-flat}，只需证明
$\text{Tor}_1^R(M, N) = 0$。

\medskip\noindent
短正合列 $0 \to R \to R' \oplus R_f \to (R')_f \to 0$
与 $N$ 所对应的 Tor 长正合列给出正合列
$$
0 \to \text{Tor}_1^R(R', N) \to \text{Tor}_1^R((R')_f, N)
$$
以及对 $i \geq 2$ 的同构
$\text{Tor}_i^R(R', N) = \text{Tor}_i^R((R')_f, N)$。
由于
$\text{Tor}_i^R((R')_f, N) = \text{Tor}_i^R(R', N)_f$，
可知 $f$ 在 $\text{Tor}_1^R(R', N)$ 上是非零因子，并且在
$i \geq 2$ 时于 $\text{Tor}_i^R(R', N)$ 上可逆。
由于 $M \otimes_R R'$ 在 $R'$ 上平坦，由例
\ref{more-algebra-example-tor-change-rings} 的谱序列，有
$$
\text{Tor}_i^R(M \otimes_R R', N) =
(M \otimes_R R') \otimes_{R'} \text{Tor}_i^R(R', N)
$$
将 $M \otimes_R R'$ 写成有限自由 $R'$-模的滤过余极限
（Algebra，定理 \ref{algebra-theorem-lazard}），可知 $f$ 在
$\text{Tor}_1^R(M \otimes_R R', N)$ 上是非零因子，并在
$\text{Tor}_i^R(M \otimes_R R', N)$ 上可逆。
% P02-E0405
接着考虑由 $M$ 可粘合而来的正合列
$0 \to M \to M \otimes_R R' \oplus M_f \to M \otimes_R (R')_f \to 0$
及其所对应的 $\text{Tor}$ 长正合列。相关部分为
$$
\xymatrix{
\text{Tor}_1^R(M, N) \ar[r] &
\text{Tor}_1^R(M \otimes_R R', N) \ar[r] &
\text{Tor}_1^R(M \otimes_R (R')_f, N) \\
& \text{Tor}_2^R(M \otimes_R R', N) \ar[r] &
\text{Tor}_2^R(M \otimes_R (R')_f, N) \ar[llu]
}
$$
由上面对 $f$ 在
$\text{Tor}_i^R(M \otimes_R R', N)$ 上作用的讨论，可得
$\text{Tor}_1^R(M, N) = 0$。
% P02-E0406
\end{proof}

\noindent
注意，对于“有限生成”这一性质，我们已经在引理
\ref{more-algebra-lemma-faithful-descent} 中看到下一个引理的结论。

\begin{lemma}
\label{more-algebra-lemma-BL-properties}
设 $(R \to R', f)$ 为粘合对。设 $M$ 为不一定对
$(R \to R', f)$ 可粘合的 $R$-模。则 $M$ 是有限投射 $R$-模，
当且仅当 $M \otimes_R R'$ 在 $R'$ 上有限投射，且
$M_f$ 在 $R_f$ 上有限投射。
\end{lemma}

\begin{proof}
假设 $M \otimes_R R'$ 是有限投射 $R'$-模，且 $M_f$ 是有限投射
$R_f$-模。我们的任务是证明 $M$ 在 $R$ 上有限投射。下文将不再说明
而使用 Algebra，引理 \ref{algebra-lemma-finite-projective}。
由引理 \ref{more-algebra-lemma-BL-flat}，$M$ 平坦。由引理
\ref{more-algebra-lemma-faithful-descent}，$M$ 有限。选取短正合列
$0 \to K \to R^{\oplus n} \to M \to 0$。由于有限投射模是有限表现的，
且该序列与 $R'$ 作张量后（由 Algebra，引理
\ref{algebra-lemma-flat-tor-zero}）以及与 $R_f$ 作张量后仍正合，
可知 $K \otimes_R R'$ 与 $K_f$ 是有限模。利用上面的引理，
可知 $K$ 有限生成。因此 $M$ 有限表现，从而有限投射。
\end{proof}

\begin{remark}
\label{more-algebra-remark-compare-BL}
\cite{Beauville-Laszlo} 假设 $f$ 在 $R$ 中是非零因子，并且
$R' = R^\wedge$；由引理 \ref{more-algebra-lemma-same-f-torsion}，这给出粘合对。
即使在该框架下，定理 \ref{more-algebra-theorem-BL-glueing} 仍给出新内容：
\cite{Beauville-Laszlo} 的结果只适用于 $f$ 为非零因子的模
（因而按我们的定义可粘合；见引理
\ref{more-algebra-lemma-same-f-torsion-module}）。引理
\ref{more-algebra-lemma-BL-properties} 也轻微推广了
\cite{Beauville-Laszlo} 的结果：不仅可允许 $M$ 具有非零的
$f$-幂挠，甚至不要求它可粘合。
\end{remark}

\section{导出完备化}
\label{more-algebra-section-derived-completion}

\noindent
本节内容的一些参考文献为
\cite{Dwyer-Greenlees}、\cite{Greenlees-May}、\cite{PSY}、\cite{dag12}
（尤其是第 4 章）。我们的论述遵循 \cite{BS}。
本节对于挠模的类似物（或“对偶”）见 Dualizing Complexes，第
\ref{dualizing-section-local-cohomology} 节。具有挠同调的复形之导出范畴
与导出完备复形之间的关系见 Dualizing Complexes，第
\ref{dualizing-section-torsion-and-complete} 节。

\medskip\noindent
设 $K \in D(A)$，$f \in A$。用 $T(K, f)$ 表示系
% P02-E0408
$$
\ldots \to K \xrightarrow{f} K \xrightarrow{f} K
$$
在 $D(A)$ 中的一个导出极限。

\begin{lemma}
\label{more-algebra-lemma-hom-from-Af}
设 $A$ 为环，$f \in A$，$K \in D(A)$。下列条件等价：
\begin{enumerate}
\item 对所有 $n$，$\Ext^n_A(A_f, K) = 0$；
\item 对 $D(A_f)$ 中所有 $E$，$\Hom_{D(A)}(E, K) = 0$；
\item $T(K, f) = 0$；
\item 对每个 $p \in \mathbf{Z}$，$T(H^p(K), f) = 0$；
\item 对每个 $p \in \mathbf{Z}$，都有
$\Hom_A(A_f, H^p(K)) = 0$ 及 $\Ext^1_A(A_f, H^p(K)) = 0$；
\item $R\Hom_A(A_f, K) = 0$；
\item 映射 $\prod_{n \geq 0} K \to \prod_{n \geq 0} K$，
$(x_0, x_1, \ldots) \mapsto (x_0 - fx_1, x_1 - fx_2, \ldots)$
是 $D(A)$ 中的同构；
\item 此处添加更多条件。
% P02-E0409
\end{enumerate}
\end{lemma}

\begin{proof}
显然 (2) 蕴含 (1)，且 (1) 等价于 (6)。假设 (1)。
设 $I^\bullet$ 为表示 $K$ 的 K-内射 $A$-模复形。
条件 (1) 表示 $\Hom_A(A_f, I^\bullet)$ 非周期。
设 $M^\bullet$ 为表示 $E$ 的 $A_f$-模复形。则由 Algebra，引理
\ref{algebra-lemma-adjoint-hom-restrict}，
$$
\Hom_{D(A)}(E, K) =
\Hom_{K(A)}(M^\bullet, I^\bullet) =
\Hom_{K(A_f)}(M^\bullet, \Hom_A(A_f, I^\bullet))
$$
由引理 \ref{more-algebra-lemma-hom-K-injective}，
$\Hom_A(A_f, I^\bullet)$ 是 K-内射 $A_f$-模复形；
% P02-E0410
它非周期这一事实蕴含它同伦等价于零
（Derived Categories，引理 \ref{derived-lemma-K-injective}）。
故得到 (2)。

\medskip\noindent
$A$-模 $A_f$ 的一个自由消解为
$$
0 \to \bigoplus\nolimits_{n \in \mathbf{N}} A \to
\bigoplus\nolimits_{n \in \mathbf{N}} A
\to A_f \to 0
$$
其中第一个映射把 $(a_0, a_1, a_2, \ldots)$ 映到
% P02-E0411
$(a_0, a_1 - fa_0, a_2 - fa_1, \ldots)$，第二个映射把
$(a_0, a_1, a_2, \ldots)$ 映到
$a_0 + a_1/f + a_2/f^2 + \ldots$。应用
$\Hom_A(-, I^\bullet)$，得到
$$
0 \to \Hom_A(A_f, I^\bullet) \to \prod I^\bullet \to \prod I^\bullet \to 0
$$
由于 $\prod I^\bullet$ 表示 $\prod_{n \geq 0} K$，这证明
(1) 与 (7) 等价。另一方面，由 Derived Categories，第
\ref{derived-section-derived-limit} 节中导出极限的构造，
所显示的正合列说明对象 $T(K, f)$ 是 $D(A)$ 中
$R\Hom_A(A_f, K)$ 的一个代表。因此 (1) 与 (3) 等价。
% P02-E0412

\medskip\noindent
有谱序列
$$
E_2^{p, q} = \Ext^p_A(A_f, H^q(K)) \Rightarrow
\Ext^{p + q}_A(A_f, K)
$$
见等式 (\ref{more-algebra-equation-first-ss-ext})\footnote{若 $K$ 不有界下，
则不能使用所给参考。此时，选择一个表示 $K$ 的复形 $K^\bullet$
以及一个两项自由消解
$0 \to F^{-1} \to F^0 \to A_f \to 0$（见上文）。
% P02-E0413
于是得到各项为 $\Hom_A(F^{-q}, K^p)$ 的双复形，其总化计算
$R\Hom(A_f, K)$。与此双复形相联系的第二个谱序列
（Homology，第 \ref{homology-section-double-complex} 节）
给出所需结论。略去细节。}。
因为 $A_f$ 具有长度为 $1$ 的投射 $A$-模消解（见上文），
$E_2$ 页只有两个非零列，故此谱序列在 $E_2$ 处退化。
因此得到短正合列
$$
0 \to \Ext^1_A(A_f, H^{p - 1}(K)) \to
\Ext^p_A(A_f, K) \to
\Hom_A(A_f, H^p(K)) \to 0
$$
这证明 (4) 与 (5) 均等价于 (1)。
% P02-E0414
\end{proof}

\begin{lemma}
\label{more-algebra-lemma-ideal-of-elements-complete-wrt}
设 $A$ 为环，$K \in D(A)$。使 $T(K, f) = 0$ 的所有
$f \in A$ 组成的集合 $I$ 是 $A$ 的根理想。
\end{lemma}

\begin{proof}
下文将不再说明而使用引理 \ref{more-algebra-lemma-hom-from-Af} 的结果。
若 $f \in I$ 且 $g \in A$，则 $A_{gf}$ 是 $A_f$-模，故对所有
$n$ 有 $\Ext^n_A(A_{gf}, K) = 0$，从而 $gf \in I$。
设 $f, g \in I$。则有短正合列
$$
0 \to A_{f + g} \to A_{f(f + g)} \oplus A_{g(f + g)} \to A_{gf(f + g)} \to 0
$$
因为 $f, g$ 在 $A_{f + g}$ 中生成单位理想。这来自 Algebra，引理
\ref{algebra-lemma-standard-covering} 以及末箭头为满射这一简单事实。
由 $\Ext$ 的长正合列及对所有 $n$ 均有
$\Ext^n_A(A_{f(f + g)}, K)$、
$\Ext^n_A(A_{g(f + g)}, K)$ 与
$\Ext^n_A(A_{gf(f + g)}, K)$ 为零，可推出对所有 $n$，
$\Ext^n_A(A_{f + g}, K)$ 为零。最后，若对某个 $n > 0$ 有
$f^n \in I$，则 $f \in I$，因为
$T(K, f) = T(K, f^n)$，也可因为 $A_f \cong A_{f^n}$。
\end{proof}

\begin{lemma}
\label{more-algebra-lemma-complete-derived-complete}
设 $A$ 为环，$I \subset A$ 为理想，$M$ 为 $A$-模。
\begin{enumerate}
\item 若 $M$ 是 $I$-进完备的，则对所有 $f \in I$，
$T(M, f) = 0$。
\item 反之，若对所有 $f \in I$ 有 $T(M, f) = 0$，且 $I$ 有限生成，
则 $M \to \lim M/I^nM$ 是满射。
\end{enumerate}
\end{lemma}

\begin{proof}
(1) 的证明。假设 $M$ 是 $I$-进完备的。由引理
\ref{more-algebra-lemma-hom-from-Af}，只需证明
$\Ext^1_A(A_f, M) = 0$ 及 $\Hom_A(A_f, M) = 0$。
由于 $M = \lim M/I^nM$，并且
$\Hom_A(A_f, M/I^nM) = 0$，可得
$\Hom_A(A_f, M) = 0$。设有扩张
$$
0 \to M \to E \to A_f \to 0
$$
对 $n \geq 0$，选取映到 $1/f^n$ 的 $e_n \in E$。
对 $n \geq 0$，置 $\delta_n = fe_{n + 1} - e_n \in M$。
将 $e_n$ 替换为
$$
e'_n = e_n + \delta_n + f\delta_{n + 1} + f^2 \delta_{n + 2} + \ldots
$$
由于 $M$ 关于 $I$ 完备且 $f \in I$，此无穷和存在。
简单计算表明 $fe'_{n + 1} = e'_n$。因此，将 $1/f^n$ 映到
$e'_n$，即得到该扩张的分裂。

\medskip\noindent
(2) 的证明。假设 $I = (f_1, \ldots, f_r)$，并且对
$i = 1, \ldots, r$ 有 $T(M, f_i) = 0$。由 Algebra，引理
\ref{algebra-lemma-when-surjective-to-completion}，可以假设
$I = (f)$ 且 $T(M, f) = 0$。设对 $n \geq 0$ 有 $x_n \in M$。
考虑扩张
$$
0 \to M \to E \to A_f \to 0
$$
其中
$$
E = M \oplus \bigoplus Ae_n\Big/\langle x_n - fe_{n + 1} + e_n\rangle
$$
并把 $e_n$ 映到 $A_f$ 中的 $1/f^n$（见上文）。
由假设及引理 \ref{more-algebra-lemma-hom-from-Af}，此扩张分裂，故得到元素
$x + e_0$，它在 $E$ 中生成 $A_f$ 的一个副本。于是
$$
x + e_0 = x - x_0 + fe_1 = x - x_0 - f x_1 + f^2 e_2 = \ldots
$$
由蛇引理，$M/f^nM = E/f^nE$，故模 $f^nM$ 有
$x = x_0 + fx_1 + \ldots + f^{n - 1}x_{n - 1}$。
换言之，映射 $M \to \lim M/f^nM$ 满射，如所需。
\end{proof}

\noindent
受上述结果启发，作如下定义。

\begin{definition}
\label{more-algebra-definition-derived-complete}
设 $A$ 为环，$K \in D(A)$，$I \subset A$ 为理想。若对每个
$f \in I$ 都有 $T(K, f) = 0$，则称 $K$
{\it 关于 $I$ 导出完备}。若 $M$ 是 $A$-模，且
$M[0] \in D(A)$ 关于 $I$ 导出完备，则称 $M$
{\it 关于 $I$ 导出完备}。
\end{definition}

\noindent
由导出完备对象组成的全子范畴
$D_{comp}(A) = D_{comp}(A, I) \subset D(A)$ 是严格全、饱和的三角
子范畴；见 Derived Categories，定义
\ref{derived-definition-triangulated-subcategory} 与
\ref{derived-definition-saturated}。由引理
\ref{more-algebra-lemma-ideal-of-elements-complete-wrt}，子范畴
$D_{comp}(A, I)$ 只依赖于 $I$ 的根理想 $\sqrt{I}$；换言之，
它只依赖于 $\Spec(A)$ 的闭子集 $Z = V(I)$。
子范畴 $D_{comp}(A, I)$ 在 $D(A)$ 中的乘积与同伦极限下保持，
但一般不在可数直和下保持。若理想 $I$ 的选择由语境清楚，
我们常简称 $M$ 为导出完备模。

\begin{proposition}
\label{more-algebra-proposition-derived-complete-modules}
设 $I \subset A$ 为环 $A$ 的有限生成理想，$M$ 为 $A$-模。
下列条件等价：
\begin{enumerate}
\item $M$ 是 $I$-进完备的；
\item $M$ 关于 $I$ 导出完备，且 $\bigcap I^nM = 0$。
\end{enumerate}
\end{proposition}

\begin{proof}
这由引理 \ref{more-algebra-lemma-complete-derived-complete} 的结果立即得到。
\end{proof}

\noindent
下一个引理说明导出完备模的范畴 $\mathcal{C}$ 是 Abel 范畴。
$\mathcal{C}$ 并非 Grothendieck Abel 范畴；见 Examples，第
\ref{examples-section-derived-complete-modules} 节。

\begin{lemma}
\label{more-algebra-lemma-serre-subcategory}
设 $I$ 为环 $A$ 的理想。
\begin{enumerate}
\item 导出完备 $A$-模组成 $\text{Mod}_A$ 的弱 Serre 子范畴
$\mathcal{C}$。
\item $D_\mathcal{C}(A) \subset D(A)$ 是导出完备对象的全子范畴。
\end{enumerate}
\end{lemma}

\begin{proof}
(2) 由引理 \ref{more-algebra-lemma-hom-from-Af} 及定义立即得到。对 (1)，设
$M \to N$ 为导出完备模之间的映射。用 $K = (M \to N)$ 表示
% P02-E0415
$D(A)$ 中相应的对象。选取 $f \in I$。对所有 $n$，
$\Ext_A^n(A_f, K)$ 为零，因为对所有 $n$，
$\Ext_A^n(A_f, M)$ 与 $\Ext_A^n(A_f, N)$ 都为零。
故 $K$ 导出完备。由 (2)，$\Ker(M \to N)$ 与
$\Coker(M \to N)$ 是 $\mathcal{C}$ 的对象。
最后，设 $0 \to M_1 \to M_2 \to M_3 \to 0$ 是 $A$-模的短正合列，
且 $M_1$、$M_3$ 导出完备。则由 $\Ext$ 的长正合列，
$M_2$ 导出完备。因此，由 Homology，引理
\ref{homology-lemma-characterize-weak-serre-subcategory}，
$\mathcal{C}$ 是弱 Serre 子范畴。
\end{proof}

\noindent
我们将在引理 \ref{more-algebra-lemma-derived-complete-zero-bis} 中推广下一个引理。

\begin{lemma}
\label{more-algebra-lemma-derived-complete-zero}
设 $I$ 为环 $A$ 的有限生成理想，$M$ 为导出完备 $A$-模。
若 $M/IM = 0$，则 $M = 0$。
\end{lemma}

\begin{proof}
假设 $M/IM$ 为零。设 $I = (f_1, \ldots, f_r)$。
令 $i < r$ 为使
$N = M/(f_1, \ldots, f_i)M$ 非零的最大整数。
若这样的 $i$ 不存在，则 $M = 0$，这正是要证明的。
由这一选择及引理 \ref{more-algebra-lemma-serre-subcategory}，$N$ 作为
导出完备模之间一个映射的余核是导出完备的。又由 $i$ 的选择，
$f_{i + 1} : N \to N$ 是满射。故
$$
\lim (\ldots \to N \xrightarrow{f_{i + 1}} N \xrightarrow{f_{i + 1}} N)
$$
非零，这与 $N$ 的导出完备性矛盾。
\end{proof}

\noindent
若环是 $I$-进完备的，则可得到丰富的导出完备复形。

\begin{lemma}
\label{more-algebra-lemma-pseudo-coherent-is-derived-complete}
设 $A$ 为环，$I \subset A$ 为理想。若 $A$ 导出完备
（例如，$I$-进完备），则 $D(A)$ 的任意伪凝聚对象都导出完备。
\end{lemma}

\begin{proof}
（引理 \ref{more-algebra-lemma-complete-derived-complete} 解释了引理的括号内陈述。）
设 $K$ 为 $D(A)$ 的伪凝聚对象。按定义，这意味着 $K$ 可由有限自由
$A$-模组成的有界上复形 $K^\bullet$ 表示。由于 $A$ 导出完备，
由引理 \ref{more-algebra-lemma-serre-subcategory} 的 (1)，对所有 $n$，
$H^n(K)$ 导出完备。再由同一引理的 (2)，这蕴含 $K$ 导出完备。
\end{proof}

\begin{lemma}
\label{more-algebra-lemma-double-localize}
设 $A$ 为环，$f, g \in A$。则对 $K \in D(A)$，有
$R\Hom_A(A_f, R\Hom_A(A_g, K)) = R\Hom_A(A_{fg}, K)$。
\end{lemma}

\begin{proof}
这由引理 \ref{more-algebra-lemma-internal-hom} 得出。
\end{proof}

\begin{lemma}
\label{more-algebra-lemma-derived-completion}
\begin{slogan}
沿有限生成理想的导出完备化存在，并且可用 {\v C}ech 方法计算。
\end{slogan}
设 $I$ 为环 $A$ 的有限生成理想。
包含函子 $D_{comp}(A, I) \to D(A)$ 有左伴随；也就是说，给定
$D(A)$ 的任意对象 $K$，存在从 $K$ 到 $D(A)$ 的一个导出完备对象
% P02-E0416
$K \to K^\wedge$ 的映射，使得只要 $E$ 是 $D(A)$ 的导出完备对象，映射
$$
\Hom_{D(A)}(K^\wedge, E) \longrightarrow \Hom_{D(A)}(K, E)
$$
就是双射。事实上，若 $I$ 由 $f_1, \ldots, f_r \in A$ 生成，则有
$$
K^\wedge = R\Hom\left((A \to \prod\nolimits_{i_0} A_{f_{i_0}} \to
\prod\nolimits_{i_0 < i_1} A_{f_{i_0}f_{i_1}}
\to \ldots \to A_{f_1\ldots f_r}), K\right)
$$
且该等式关于 $K$ 函子性成立。
\end{lemma}

\begin{proof}
用引理中最后一个展示公式定义 $K^\wedge$。存在复形映射
$$
(A \to \prod\nolimits_{i_0} A_{f_{i_0}} \to
\prod\nolimits_{i_0 < i_1} A_{f_{i_0}f_{i_1}} \to
\ldots \to A_{f_1\ldots f_r}) \longrightarrow A
$$
它诱导映射 $K \to K^\wedge$。只需证明 $K^\wedge$ 导出完备，并且
若 $K$ 导出完备，则 $K \to K^\wedge$ 是同构\footnote{具体而言，
若 $E$ 导出完备且 $a : K \to E$ 是映射，则交换图
$$
\xymatrix{
K \ar[r]_a \ar[d] & E \ar[d]^{\cong} \\
K^\wedge \ar[r]^{a^\wedge} & E^\wedge
}
$$
表明到 $E$ 的任意映射都经 $K^\wedge$ 分解。选取一个杰出三角形
$K \to K^\wedge \to C$ 并施以 ${}^\wedge$。于是得到
$C^\wedge = 0$。由上文，这意味着对导出完备的 $E$，
$\Hom(C, E) = 0$，从而按所需证明了
$\Hom(K, E) = \Hom(K^\wedge, E)$。}。

\medskip\noindent
设 $f \in A$。由引理 \ref{more-algebra-lemma-double-localize}，对象
$R\Hom_A(A_f, K^\wedge)$ 等于
$$
R\Hom\left((A_f \to \prod\nolimits_{i_0} A_{ff_{i_0}} \to
\prod\nolimits_{i_0 < i_1} A_{ff_{i_0}f_{i_1}} \to
\ldots \to A_{ff_1\ldots f_r}), K\right)
$$
若 $f \in I$，则 $f_1, \ldots, f_r$ 在 $A_f$ 中生成单位理想，
故增广交错 {\v C}ech 复形
$$
A_f \to \prod\nolimits_{i_0} A_{ff_{i_0}} \to
\prod\nolimits_{i_0 < i_1} A_{ff_{i_0}f_{i_1}} \to
\ldots \to A_{ff_1\ldots f_r}
$$
由引理 \ref{more-algebra-lemma-extended-alternating-torsion} 在 $D(A)$ 中为零。
（事实上，若对某个 $i$ 有 $f = f_i$，则由引理
\ref{more-algebra-lemma-extended-alternating-homotopy-zero}，此复形零伦；
这是我们所需的唯一情形。）
因此 $R\Hom_A(A_f, K^\wedge) = 0$，再由引理
\ref{more-algebra-lemma-hom-from-Af} 得出 $K^\wedge$ 导出完备。

\medskip\noindent
反之，若 $K$ 导出完备，则对所有
$f = f_{i_0} \ldots f_{i_p}$、$p \geq 0$，
$R\Hom_A(A_f, K)$ 都为零。因此 $K \to K^\wedge$ 在
$D(A)$ 中是同构。
\end{proof}

\begin{remark}
\label{more-algebra-remark-derived-completion}
设 $A$ 为环，$I \subset A$ 为有限生成理想。由引理
\ref{more-algebra-lemma-derived-completion} 存在的包含函子
$D_{comp}(A, I) \to D(A)$ 的左伴随称为{\it 导出完备化}。
为表明这一点，我们会说“令 $K^\wedge$ 为 $K$ 的导出完备化”。
请记住，伴随的单位是函子性映射 $K \to K^\wedge$。
\end{remark}

\begin{lemma}
\label{more-algebra-lemma-derived-completion-vanishes}
设 $A$ 为环，$I \subset A$ 为有限生成理想。设 $K^\bullet$ 为
$A$-模复形，且对某个 $f \in I$，映射
$f : K^\bullet \to K^\bullet$ 是同构，也就是说，$K^\bullet$
是 $A_f$-模复形。则 $K^\bullet$ 的导出完备化为零。
\end{lemma}

\begin{proof}
事实上，在这种情形下，对任意导出完备复形 $L$，
% P02-E0417
$R\Hom_A(K, L)$ 都为零，见引理 \ref{more-algebra-lemma-hom-from-Af}。
因此，由引理 \ref{more-algebra-lemma-derived-completion} 中的泛性质，
$K^\wedge$ 为零。
\end{proof}

\begin{lemma}
\label{more-algebra-lemma-completion-RHom}
设 $A$ 为环，$I \subset A$ 为有限生成理想。设
$K, L \in D(A)$。则
$$
R\Hom_A(K, L)^\wedge = R\Hom_A(K, L^\wedge) = R\Hom_A(K^\wedge, L^\wedge)
$$
\end{lemma}

\begin{proof}
由引理 \ref{more-algebra-lemma-derived-completion}，导出完备化由某个
$C \in D(A)$ 的 $R\Hom_A(C, -)$ 给出。于是
\begin{align*}
R\Hom_A(C, R\Hom_A(K, L))
& =
R\Hom_A(C \otimes_A^\mathbf{L} K, L) \\
& =
R\Hom_A(K, R\Hom_A(C, L))
\end{align*}
这由引理 \ref{more-algebra-lemma-internal-hom} 得出。它证明第一个等式。
映射 $K \to K^\wedge$ 诱导映射
$$
R\Hom_A(K^\wedge, L^\wedge) \to R\Hom_A(K, L^\wedge)
$$
按照导出完备化作为包含函子的左伴随这一良定义，它在 $D(A)$
中是同构。
\end{proof}

\begin{lemma}
\label{more-algebra-lemma-naive-derived-completion}
设 $A$ 为环，$I \subset A$ 为理想。设 $(K_n)$ 为 $D(A)$ 中对象的
逆系，使得对所有 $f \in I$ 与 $n$，都存在
$e = e(n, f)$，使 $f^e$ 在 $K_n$ 上为零。则对 $K \in D(A)$，
对象 $K' = R\lim (K \otimes_A^\mathbf{L} K_n)$ 关于 $I$ 导出完备。
\end{lemma}

\begin{proof}
由于导出完备对象的范畴在 $R\lim$ 下保持，只需证明每个
$K \otimes_A^\mathbf{L} K_n$ 都导出完备。由假设，对所有
$f \in I$，都存在 $e$，使 $f^e$ 在
$K \otimes_A^\mathbf{L} K_n$ 上为零。这当然蕴含
$T(K \otimes_A^\mathbf{L} K_n, f) = 0$，结论即得。
\end{proof}

\begin{situation}
\label{more-algebra-situation-koszul}
设 $A$ 为环，$I = (f_1, \ldots, f_r) \subset A$。令
$K_n^\bullet = K_\bullet(A, f_1^n, \ldots, f_r^n)$ 为关于
$f_1^n, \ldots, f_r^n$ 的 Koszul 复形，并把它看作位于次数
$-r, -r + 1, \ldots, 0$ 的上链复形。利用引理
\ref{more-algebra-lemma-functorial} 的函子性，得到逆系
$$
\ldots \to K_3^\bullet \to K_2^\bullet \to K_1^\bullet
$$
它与逆系
$H^0(K_n^\bullet) = A/(f_1^n, \ldots, f_r^n)$
以及映射 $A \to K_n^\bullet$ 相容。
\end{situation}

\noindent
% P02-E0418
下文讨论的一个关键特征在于：若 $m > n$，则映射
$$
K_m^{-p} = \wedge^p(A^{\oplus r}) \to \wedge^p(A^{\oplus r}) = K_n^{-p}
$$
在基元素 $e_{i_1} \wedge \ldots \wedge e_{i_p}$ 上由乘以
$f_{i_1}^{m - n} \ldots f_{i_p}^{m - n}$ 给出。

\begin{lemma}
\label{more-algebra-lemma-koszul-derived-completion-complete}
在情形 \ref{more-algebra-situation-koszul} 中。对 $K \in D(A)$，对象
$K' = R\lim (K \otimes_A^\mathbf{L} K_n^\bullet)$ 关于 $I$ 导出完备。
\end{lemma}

\begin{proof}
这是引理 \ref{more-algebra-lemma-naive-derived-completion} 的特殊情形，因为
$f_i^n$ 在 $K_n^\bullet$ 上的作用是一个自同态，而由引理
\ref{more-algebra-lemma-homotopy-koszul}，该自同态零伦。
\end{proof}

\begin{lemma}
\label{more-algebra-lemma-characterize-derived-complete-Koszul}
在情形 \ref{more-algebra-situation-koszul} 中。设 $K \in D(A)$。
以下条件等价：
\begin{enumerate}
\item $K$ 关于 $I$ 导出完备；以及
\item 典范映射 $K \to R\lim (K \otimes_A^\mathbf{L} K_n^\bullet)$
是 $D(A)$ 中的同构。
\end{enumerate}
\end{lemma}

\begin{proof}
若 (2) 成立，则由引理
\ref{more-algebra-lemma-koszul-derived-completion-complete}，$K$ 关于 $I$
导出完备。反之，假设 $K$ 关于 $I$ 导出完备。考虑由愚蠢截断
（Homology，第 \ref{homology-section-truncations} 节）给出的滤过
$$
K_n^\bullet \supset
\sigma_{\geq -r + 1}K_n^\bullet \supset
\sigma_{\geq -r + 2}K_n^\bullet \supset \ldots \supset
\sigma_{\geq -1}K_n^\bullet \supset
\sigma_{\geq 0}K_n^\bullet = A
$$
。因为构造 $R\lim(K \otimes E)$ 关于第二个变元是正合的
（引理 \ref{more-algebra-lemma-tensor-Rlim-exact}），所以只需证明
$$
R\lim \left(
K \otimes_A^\mathbf{L}
(\sigma_{\geq p}K_n^\bullet/ \sigma_{\geq p + 1}K_n^\bullet)
\right) = 0
$$
对 $p < 0$ 成立。上面对 Koszul 复形的显式描述表明
$$
R\lim \left(
K \otimes_A^\mathbf{L}
(\sigma_{\geq p}K_n^\bullet/ \sigma_{\geq p + 1}K_n^\bullet)
\right) =
\bigoplus\nolimits_{i_1, \ldots, i_{-p}}
T(K, f_{i_1}\ldots f_{i_{-p}})
$$
，而由关于 $K$ 的假设，这在 $p < 0$ 时为零。
\end{proof}

\begin{lemma}
\label{more-algebra-lemma-derived-completion-koszul}
在情形 \ref{more-algebra-situation-koszul} 中。
把 $K \in D(A)$ 送到导出极限
$K' = R\lim( K \otimes_A^\mathbf{L} K_n^\bullet )$ 的函子，是引理
\ref{more-algebra-lemma-derived-completion} 中所构造的包含函子
$D_{comp}(A) \to D(A)$ 的左伴随。
\end{lemma}

\begin{proof}[第一种证明]
赋值 $K \leadsto K'$ 是函子，而且由引理
\ref{more-algebra-lemma-koszul-derived-completion-complete}，$K'$ 关于 $I$
导出完备。由一个形式论证（略去），只需证明当 $K$ 关于 $I$
导出完备时，$K \to K'$ 是同构。这正是引理
\ref{more-algebra-lemma-characterize-derived-complete-Koszul}。
\end{proof}

\begin{proof}[第二种证明]
% P02-E0419
用 $K \mapsto K^\wedge$ 表示引理
\ref{more-algebra-lemma-derived-completion} 中构造的伴随。由该引理，
$$
K^\wedge = R\Hom\left((A \to \prod\nolimits_{i_0} A_{f_{i_0}} \to
\prod\nolimits_{i_0 < i_1} A_{f_{i_0}f_{i_1}}
\to \ldots \to A_{f_1\ldots f_r}), K\right)
$$
在引理 \ref{more-algebra-lemma-extended-alternating-Cech-is-colimit-koszul} 中，
我们已经看到增广交错 {\v C}ech 复形
$$
A \to \prod\nolimits_{i_0} A_{f_{i_0}} \to
\prod\nolimits_{i_0 < i_1} A_{f_{i_0}f_{i_1}}
\to \ldots \to A_{f_1\ldots f_r}
$$
是位于次数 $0, \ldots, r$ 的 Koszul 复形
$K^n = K(A, f_1^n, \ldots, f_r^n)$ 的余极限。注意，$K^n$
是由有限自由 $A$-模组成的有限链复形，其对偶（如引理
\ref{more-algebra-lemma-dual-perfect-complex} 中所述）为
$R\Hom_A(K^n, A) = K_n$，其中 $K_n$ 是位于次数
$-r, \ldots, 0$ 的 Koszul 上链复形（同往常一样）。
因此只需证明
$$
R\Hom_A(\text{hocolim} K^n, K) =  R\lim (K \otimes_A^\mathbf{L} K_n)
$$
这由引理 \ref{more-algebra-lemma-colim-and-lim-of-duals} 得出。
\end{proof}

\begin{lemma}
\label{more-algebra-lemma-derived-complete-bounded}
设 $I = (f_1, \ldots, f_r)$ 为环 $A$ 的有限生成理想。
设 $K$ 为 $D(A)$ 的导出完备对象。以下条件等价：
\begin{enumerate}
\item 对 $i > 0$，有 $H^i(K) = 0$；
\item 对 $i > 0$，有 $H^i(K \otimes_A^\mathbf{L} A/I) = 0$；
\item 对 $i > 0$，有 $H^i(K \otimes_A^\mathbf{L} K_1^\bullet) = 0$，
其中 $K_1^\bullet$ 如情形 \ref{more-algebra-situation-koszul} 中所述。
\end{enumerate}
\end{lemma}

\begin{proof}
蕴含关系 (1) $\Rightarrow$ (2) 总成立。蕴含关系
(2) $\Rightarrow$ (3) 由引理 \ref{more-algebra-lemma-bounded-above-mod-I} 得出。
假设 (3)。对 $s = 0, \ldots, r$，考虑复形
$$
K(s) =  K \otimes_A^\mathbf{L} K_\bullet(A, f_1, \ldots, f_s)
$$
，其中 $K_\bullet(A, f_1, \ldots, f_s)$ 是置于上同调次数
$-s, \ldots, 0$ 的 Koszul 复形。由引理
\ref{more-algebra-lemma-cone-koszul}，我们有一个
% P02-E0420
杰出三角形
$$
K(s - 1) \xrightarrow{f_s} K(s - 1) \to K(s) \to K(s - 1)[1]
$$
。由于 $K(0) = K$ 导出完备，向上归纳得出每个 $K(s)$ 都是
$D(A)$ 的导出完备对象。我们将向下归纳证明 $K(s)$ 在次数
$> 0$ 没有非零上同调。事实上，由假设 (3)，这对 $s = r$
成立。若它对 $K(s)$ 成立且 $s > 0$，则对 $i > 0$，
$f_s : H^i(K(s - 1)) \to H^i(K(s - 1))$ 是满射；再由引理
\ref{more-algebra-lemma-serre-subcategory} 与
\ref{more-algebra-lemma-derived-complete-zero}，得出 $H^i(K(s - 1))$ 为零。
\end{proof}

\begin{lemma}
\label{more-algebra-lemma-derived-complete-zero-bis}
\begin{slogan}
导出 Nakayama
\end{slogan}
\begin{reference}
一个相关结果是 \cite[Proposition 6.5]{Dwyer-Greenlees}。
例如，在 Bhatt 于 2018 年秋季在哥伦比亚大学所作关于棱柱上同调的
第三讲中，可在第 2 节性质 (2) 找到导出 Nakayama 引理。Leonid
Positselski 在 \url{https://mathoverflow.net/a/331501} 中提出了一份证明。
不过，我们采用评论中 Anonymous 建议的证明。
\end{reference}
设 $I$ 为环 $A$ 的有限生成理想。设 $K$ 为 $D(A)$ 的导出完备对象。
若 $K \otimes_A^\mathbf{L} A/I = 0$，则 $K = 0$。
\end{lemma}

\begin{proof}
这是引理 \ref{more-algebra-lemma-derived-complete-bounded} 的一个特殊情形，
但我们也给出一个直接证明。选取 $I$ 的生成元
$f_1, \ldots, f_r$。
% P02-E0421
用 $K_n$ 表示 $A$ 上关于 $f_1^n, \ldots, f_r^n$ 的 Koszul 复形。
回忆 $K_n$ 有界，并且 $K_n$ 的上同调模被
$f_1^n, \ldots, f_r^n$ 消去，因而被 $I^{nr}$ 消去。
由引理 \ref{more-algebra-lemma-derived-vanishing-mod-I}，得到
$K \otimes_A^\mathbf{L} K_n = 0$。由于 $K$ 导出完备，由引理
\ref{more-algebra-lemma-derived-completion-koszul} 有
$K = R\lim K \otimes_A^\mathbf{L} K_n = 0$，如所需。
\end{proof}

\noindent
作为与 Koszul 复形之关系的一个应用，我们得到导出完备化具有
有限上同调维数。

\begin{lemma}
\label{more-algebra-lemma-derived-completion-finite-cohomological-dimension}
设 $A$ 为环，$I \subset A$ 为可由 $r$ 个元素生成的理想。
则导出完备化具有有限上同调维数：
\begin{enumerate}
\item 设 $K \to L$ 为 $D(A)$ 中的态射，使得
$H^i(K) \to H^i(L)$ 在 $i \geq 1$ 时是同构，在 $i = 0$ 时是满射。
则 $H^i(K^\wedge) \to H^i(L^\wedge)$ 在 $i \geq 1$ 时是同构，
在 $i = 0$ 时是满射。
\item 设 $K \to L$ 为 $D(A)$ 的态射，使得
$H^i(K) \to H^i(L)$ 在 $i \leq -1$ 时是同构，在 $i = 0$ 时是单射。
则 $H^i(K^\wedge) \to H^i(L^\wedge)$ 在 $i \leq -r - 1$ 时是同构，
在 $i = -r$ 时是单射。
\end{enumerate}
\end{lemma}

\begin{proof}
设 $I$ 由 $f_1, \ldots, f_r$ 生成。对任意 $K \in D(A)$，
由引理 \ref{more-algebra-lemma-derived-completion-koszul} 有
$K^\wedge = R\lim K \otimes_A^\mathbf{L} K_n$，其中 $K_n$ 是关于
$f_1^n, \ldots, f_r^n$ 的 Koszul 复形。因此，由引理
\ref{more-algebra-lemma-break-long-exact-sequence-modules} 得到短正合列
$$
0 \to R^1\lim H^{i - 1}(K \otimes_A^\mathbf{L} K_n)
\to H^i(K^\wedge) \to \lim H^i(K \otimes_A^\mathbf{L} K_n) \to 0
$$
。

\medskip\noindent
(1) 的证明。选取杰出三角形 $K \to L \to C \to K[1]$。
则对 $i \geq 0$，有 $H^i(C) = 0$。由于 $K_n$ 位于次数
$\leq 0$，可见对 $i \geq 0$，有
$H^i(C \otimes_A^\mathbf{L} K_n) = 0$，并且
$H^{-1}(C \otimes_A^\mathbf{L} K_n) =
H^{-1}(C) \otimes_A A/(f_1^n, \ldots, f_r^n)$ 是转移映射均为满射的
一个系。上面的展示等式表明对 $i \geq 0$，
$H^i(C^\wedge) = 0$。应用杰出三角形
$K^\wedge \to L^\wedge \to C^\wedge \to K^\wedge[1]$，即得 (1)。

\medskip\noindent
(2) 的证明。选取杰出三角形 $K \to L \to C \to K[1]$。
则对 $i < 0$，有 $H^i(C) = 0$。由于 $K_n$ 位于次数
$-r, \ldots, 0$，可见对 $i < -r$，有
$H^i(C \otimes_A^\mathbf{L} K_n) = 0$。上面的展示等式表明
% P02-E0422
对 $i < r$，有 $H^i(C^\wedge) = 0$。应用杰出三角形
$K^\wedge \to L^\wedge \to C^\wedge \to K^\wedge[1]$，即得 (2)。
\end{proof}

\begin{lemma}
\label{more-algebra-lemma-derived-completion-spectral-sequence}
设 $A$ 为环，$I \subset A$ 为有限生成理想。设 $K^\bullet$ 为
$A$-模的滤过复形。存在双分次导出完备 $A$-模的典范谱序列
$(E_r, \text{d}_r)_{r \geq 1}$，其中
% P02-E0423
$d_r$ 的双次数为 $(r, -r + 1)$，并且
$$
E_1^{p, q} = H^{p + q}((\text{gr}^pK^\bullet)^\wedge)
$$
若每个 $K^n$ 上的滤过有限，则该谱序列有界并收敛于
$H^*((K^\bullet)^\wedge)$。
\end{lemma}

\begin{proof}
由引理 \ref{more-algebra-lemma-derived-completion}，导出完备化由某个
$C \in D^b(A)$ 的 $R\Hom_A(C, -)$ 给出。由引理
\ref{more-algebra-lemma-derived-completion-finite-cohomological-dimension} 和
\ref{more-algebra-lemma-projective-amplitude}，可见 $C$ 具有有限投射维数。
因此可选取一个表示 $C$ 的有界投射模复形 $P^\bullet$。于是
$$
M^\bullet = \Hom^\bullet(P^\bullet, K^\bullet)
$$
是一个表示 $(K^\bullet)^\wedge$ 的 $A$-模复形。它带有由
$F^pM^\bullet = \Hom^\bullet(P^\bullet, F^pK^\bullet)$ 给出的滤过。
可见 $F^pM^\bullet$ 表示 $(F^pK^\bullet)^\wedge$，因而
% P02-E0424
$\text{gr}^pM^\bullet$ 表示 $(\text{gr}K^\bullet)^\wedge$。
于是取滤过复形 $M^\bullet$ 的谱序列，便得到我们的谱序列；
见 Homology，第 \ref{homology-section-filtered-complex} 节。
若每个 $K^n$ 上的滤过有限，则因为 $P^\bullet$ 是有界复形，
每个 $M^n$ 上的滤过也有限。因此，最后一个陈述由 Homology，
引理 \ref{homology-lemma-biregular-ss-converges} 得出。
\end{proof}

\begin{example}
\label{more-algebra-example-derived-completion-spectral-sequence}
设 $A$ 为环，$I \subset A$ 为有限生成理想。设 $K^\bullet$ 为
$A$-模复形。可以在引理
\ref{more-algebra-lemma-derived-completion-spectral-sequence} 中取
$F^pK^\bullet = \tau_{\leq -p}K^\bullet$。于是得到有界谱序列
$$
E_1^{p, q} = H^{p + q}(H^{-p}(K^\bullet)^\wedge[p]) =
H^{2p + q}(H^{-p}(K^\bullet)^\wedge)
$$
，它收敛于 $H^{p + q}((K^\bullet)^\wedge)$。重新编号
$p = -j$ 与 $q = i + 2j$ 后，得到：对任意 $K \in D(A)$，
存在双分次导出完备模的有界谱序列
$(E'_r, d'_r)_{r \geq 2}$，其中 $d'_r$ 的双次数为
$(r, -r + 1)$，满足
$$
(E'_2)^{i, j} = H^i(H^j(K)^\wedge)
$$
，并收敛于 $H^{i + j}(K^\wedge)$。
\end{example}

\begin{lemma}
\label{more-algebra-lemma-restriction-derived-complete}
设 $A \to B$ 为环映射，$I \subset A$ 为理想。$D_{comp}(A, I)$
在限制函子 $D(B) \to D(A)$ 下的原像是 $D_{comp}(B, IB)$。
\end{lemma}

\begin{proof}
利用引理 \ref{more-algebra-lemma-ideal-of-elements-complete-wrt} 可见，
$L \in D(B)$ 属于 $D_{comp}(B, IB)$，当且仅当对每个
% P02-E0425
局部截面 $f \in I$，$T(L, f)$ 为零。注意，$T(L, f)$ 的上同调
在阿贝尔群范畴中计算，所以无论把 $f$ 看作 $A$ 的元素，还是取
$f$ 在 $B$ 中的像，都没有区别。由此及导出完备对象的定义，
立即得到该引理。
\end{proof}

\begin{lemma}
\label{more-algebra-lemma-restriction-derived-complete-equivalence}
设 $A \to B$ 为环映射，$I \subset A$ 为有限生成理想。
若 $A \to B$ 平坦且 $A/I \cong B/IB$，则限制函子
$D(B) \to D(A)$ 诱导等价
$D_{comp}(B, IB) \to D_{comp}(A, I)$。
\end{lemma}

\begin{proof}
选取 $I$ 的生成元 $f_1, \ldots, f_r$。
% P02-E0426
用 $\check{\mathcal{C}}^\bullet_A \to \check{\mathcal{C}}^\bullet_B$
表示引理 \ref{more-algebra-lemma-map-identifies-koszul-and-cech-complexes} 中
增广交错 {\v C}ech 复形的拟同构。设
$K \in D_{comp}(A, I)$。设 $I^\bullet$ 为表示 $K$ 的
K-内射 $A$-模复形。由于对所有 $f \in I$ 与
$n \in \mathbf{Z}$，$\Ext^n_A(A_f, K)$ 和
$\Ext^n_A(B_f, K)$ 均为零（引理 \ref{more-algebra-lemma-hom-from-Af}），
我们得出 $\check{\mathcal{C}}^\bullet_A \to A$ 与
$\check{\mathcal{C}}^\bullet_B \to B$ 诱导拟同构
$$
I^\bullet = \Hom_A(A, I^\bullet) \longrightarrow
\text{Tot}(\Hom_A(\check{\mathcal{C}}^\bullet_A, I^\bullet))
$$
以及
$$
\Hom_A(B, I^\bullet) \longrightarrow
\text{Tot}(\Hom_A(\check{\mathcal{C}}^\bullet_B, I^\bullet))
$$
。略去一些细节。
由于 $\check{\mathcal{C}}^\bullet_A \to \check{\mathcal{C}}^\bullet_B$
是拟同构且 $I^\bullet$ 是 K-内射的，得出
$\Hom_A(B, I^\bullet) \to I^\bullet$ 是拟同构。
因为复形 $\Hom_A(B, I^\bullet)$ 是 $B$-模复形，得出 $K$
属于限制映射的像，也就是说，
% P02-E0427
该函子本质满。

\medskip\noindent
事实上，该论证表明
$F : D_{comp}(A, I) \to D_{comp}(B, IB)$，
$K \mapsto \Hom_A(B, I^\bullet)$ 是限制函子的左逆。最后，假设
$L \in D_{comp}(B, IB)$。用 $B$-模的 K-内射复形 $J^\bullet$
表示 $L$。则 $J^\bullet$ 作为 $A$-模复形也是 K-内射的
（引理 \ref{more-algebra-lemma-K-injective-flat}），因而
$F(\text{下列对象的限制： }L) = \Hom_A(B, J^\bullet)$。
存在 $B$-模复形的映射
$J^\bullet \to \Hom_A(B, J^\bullet)$，它与
$\Hom_A(B, J^\bullet) \to J^\bullet$ 的复合是恒等映射。
我们得出 $F$ 也是限制函子的右逆，证明完成。
\end{proof}

\section{导出完备模的范畴}
\label{more-algebra-section-derived-complete-modules}

\noindent
设 $A$ 为环，$I$ 为理想。
% P02-E0428
用 $\mathcal{C}$ 表示导出完备模的范畴，见定义
\ref{more-algebra-definition-derived-complete}。本节讨论此范畴的一些性质。
在 Examples，第 \ref{examples-section-derived-complete-modules} 节中，
我们说明 $\mathcal{C}$ 一般不是 Grothendieck 阿贝尔范畴。

\medskip\noindent
由引理 \ref{more-algebra-lemma-serre-subcategory}，范畴 $\mathcal{C}$ 是阿贝尔的，
且包含函子 $\mathcal{C} \to \text{Mod}_A$ 正合。

\medskip\noindent
由于 $D_{comp}(A) \subset D(A)$ 在乘积下封闭
（见定义 \ref{more-algebra-definition-derived-complete} 后的讨论），并且
$D(A)$ 中的乘积在复形层面计算，可见 $\mathcal{C}$ 有乘积，
且它们与 $\text{Mod}_A$ 中的乘积一致。因此，$\mathcal{C}$
事实上有任意极限，包含函子
$\mathcal{C} \to \text{Mod}_A$ 与它们交换；见 Categories，
引理 \ref{categories-lemma-limits-products-equalizers}。

\medskip\noindent
假设 $I$ 有限生成。用 ${}^\wedge : D(A) \to D(A)$ 表示引理
\ref{more-algebra-lemma-derived-completion} 的导出完备化函子。我们来证明函子
$$
\text{Mod}_A \longrightarrow \mathcal{C},\quad
M \longmapsto H^0(M^\wedge)
$$
是包含函子 $\mathcal{C} \to \text{Mod}_A$ 的左伴随。
例如，由引理
\ref{more-algebra-lemma-derived-completion-finite-cohomological-dimension}，
对 $i > 0$，有 $H^i(M^\wedge) = 0$。因此，若 $N$ 是导出完备
$A$-模，则有
\begin{align*}
\Hom_\mathcal{C}(H^0(M^\wedge), N)
& =
\Hom_{D_{comp}(A)}(M^\wedge, N)\\
& =
\Hom_{D(A)}(M, N) \\
& =
\Hom_A(M, N)
\end{align*}
，如所需。

\medskip\noindent
设 $T$ 为预序集，$t \mapsto M_t$ 为导出完备 $A$-模的一个系，
也就是 $\mathcal{C}$ 中以 $T$ 为指标的一个系；见 Categories，
第 \ref{categories-section-posets-limits} 节。
% P02-E0429
用 $\colim_{t \in T} M_t$ 表示该系在 $\text{Mod}_A$ 中的余极限。
由上文形式地得到
$$
H^0((\colim_{t \in T} M_t)^\wedge)
$$
是该系在 $\mathcal{C}$ 中的余极限。这样便看出
$\mathcal{C}$ 有所有余极限。一般而言，包含函子
$\mathcal{C} \to \text{Mod}_A$ 不与余极限交换；见 Examples，
第 \ref{examples-section-derived-complete-modules} 节。

\begin{lemma}
\label{more-algebra-lemma-derived-complete-modules}
设 $A$ 为环，$I \subset A$ 为理想。导出完备模的范畴
$\mathcal{C}$ 是阿贝尔的并有任意极限，且包含函子
$F : \mathcal{C} \to \text{Mod}_A$ 正合并与极限交换。
若 $I$ 有限生成，则 $\mathcal{C}$ 有任意余极限，且
% P02-E0430
$F$ 有左伴随
\end{lemma}

\begin{proof}
这是对上文讨论的总结。
\end{proof}

\section{主理想的导出完备化}
\label{more-algebra-section-derived-completion-principal}

\noindent
本节讨论理想由单个元素生成时导出完备化的情形。

\begin{lemma}
\label{more-algebra-lemma-lift-universally}
设 $A$ 为环，$f \in A$。若存在整数 $c \geq 1$，使
$A[f^c] = A[f^{c + 1}] = A[f^{c + 2}] = \ldots$
（例如，当 $A$ 为 Noether 环时），则对所有 $n \geq 1$，
在 $D(A)$ 中存在映射
$$
(A \xrightarrow{f^n} A) \longrightarrow A/(f^n),
\quad\text{且}\quad
A/(f^{n + c}) \longrightarrow (A \xrightarrow{f^n} A)
$$
，它们诱导 $D(A)$ 中 pro-对象 $\{A/(f^n)\}$ 与
$\{(f^n : A \to A)\}$ 的同构。
\end{lemma}

\begin{proof}
第一个展示箭头是显然的。可以用图
$$
\xymatrix{
A/A[f^c] \ar[r]_-{f^{n + c}} \ar[d]_{f^c} & A \ar[d]^1 \\
A \ar[r]^{f^n} & A
}
$$
定义引理中的第二个箭头。由于上方水平箭头是单射，上行中的复形
拟同构于 $A/f^{n + c}A$。我们略去为证明关于 pro 对象的陈述而需
进行的复合计算。
\end{proof}

\begin{lemma}
\label{more-algebra-lemma-when-does-it-work}
设 $A$ 为环且 $f \in A$。置 $I = (f)$。在此情形下，有朴素导出完备化
$K \mapsto K' = R\lim (K \otimes_A^\mathbf{L} A/f^nA)$，以及引理
\ref{more-algebra-lemma-derived-completion-koszul} 的导出完备化
$$
K \mapsto K^\wedge = R\lim (K \otimes_A^\mathbf{L} (A \xrightarrow{f^n} A))
$$
。函子的自然变换 $K^\wedge \to K'$ 是同构，当且仅当 $A$ 的
$f$-幂挠有界。
\end{lemma}

\begin{proof}
若 $f$-幂挠有界，则由引理 \ref{more-algebra-lemma-lift-universally}，pro-对象
$\{(f^n : A \to A)\}$ 与 $\{A/f^nA\}$ 同构。因此，由引理
\ref{more-algebra-lemma-Rlim-pro-equal}，这些函子同构。反之，由引理
\ref{more-algebra-lemma-tensor-Rlim-exact} 可见，该条件恰好是
$$
R\lim (K \otimes_A^\mathbf{L} A[f^n])
$$
对所有 $K \in D(A)$ 都为零。这里，系 $(A[f^n])$ 的映射由乘以
$f$ 给出。取 $K = A$ 以及 $K = \bigoplus_{i \in \mathbf{N}} A$，
由引理 \ref{more-algebra-lemma-Rlim-zero-of-direct-sums} 可见，这蕴含
$(A[f^n])$ 作为 pro-对象为零，也就是说，对某个 $n$，
$f^{n - 1}A[f^n] = 0$，即 $A[f^{n - 1}] = A[f^n]$，也即
$f$-幂挠有界。
\end{proof}

\begin{example}
\label{more-algebra-example-derived-complete-modules}
设 $A$ 为环，$f \in A$ 为非零因子。应当牢记的一个例子是
$A = \mathbf{Z}_p$ 与 $f = p$。设 $M$ 为 $A$-模。断言：
$M$ 关于 $f$ 导出完备，当且仅当存在短正合列
$$
0 \to K \to L \to M \to 0
$$
，其中 $K, L$ 是 $f$-进完备模，且其 $f$-挠为零。具体而言，
% P02-E0431
若存在这样的短正合列，则
$$
M \otimes_A^\mathbf{L} (A \xrightarrow{f^n} A) = (K/f^nK \to L/f^nL)
$$
，因为 $f$ 在 $K$ 和 $L$ 上都是非零因子；于是得出
$R\lim (M \otimes_A^\mathbf{L} (A \xrightarrow{f^n} A))$
拟同构于 $K \to L$，也就是 $M$。由引理
\ref{more-algebra-lemma-characterize-derived-complete-Koszul}，这表明 $M$
导出完备。

反之，假设 $M$ 导出完备。选取满射 $F \to M$，其中 $F$ 是自由
$A$-模。由于 $f$ 在 $F$ 上是非零因子，$F$ 的导出完备化为
$L = \lim F/f^nF$。注意 $L$ 无 $f$-挠：若 $x_n \in F$ 的
$(x_n)$ 表示 $L$ 的一个元素 $\xi$ 且 $f\xi = 0$，则对某些
$z_n, y_n \in F$，有 $x_n = x_{n + 1} + f^nz_n$ 与
$fx_n = f^ny_n$。于是
$f^n y_n = fx_n =
fx_{n + 1} + f^{n + 1}z_n = f^{n + 1}y_{n + 1} + f^{n + 1}z_n$；
由于 $f$ 在 $F$ 上是非零因子，可见 $y_n \in fF$，这蕴含
$x_n \in f^nF$，即 $\xi = 0$。由于 $L$ 是导出完备化，
泛性质给出分解 $F \to M$ 的映射 $L \to M$。令
$K = \Ker(L \to M)$ 为其核。仍然有 $K$ 无 $f$-挠，故 $K$
的导出完备化为 $\lim K/f^nK$。另一方面，$M$ 与 $L$
都导出完备，所以由引理 \ref{more-algebra-lemma-serre-subcategory}，$K$
也导出完备。因此 $K = \lim K/f^nK$，断言得证。
\end{example}

\begin{example}
\label{more-algebra-example-derived-complete-not-complete}
设 $p$ 为素数。考虑将 $x$ 送到 $py$ 的多项式代数映射
$\mathbf{Z}_p[x] \to \mathbf{Z}_p[y]$。考虑诱导的（通常）
$p$-进完备化映射的余核
$M = \Coker(\mathbf{Z}_p[x]^\wedge \to \mathbf{Z}_p[y]^\wedge)$。
由命题 \ref{more-algebra-proposition-derived-complete-modules} 和引理
\ref{more-algebra-lemma-serre-subcategory}，$M$ 是导出完备
$\mathbf{Z}_p$-模；另见例 \ref{more-algebra-example-derived-complete-modules}
中的讨论。然而，$M$ 不是 $p$-进完备的，因为
$1 + py + p^2 y^2 + \ldots$ 映到 $M$ 的一个非零元素，而该元素
属于 $\bigcap p^nM$。
\end{example}

\begin{example}
\label{more-algebra-example-spectral-sequence-principal}
设 $A$ 为环且 $f \in A$。
% P02-E0432
用 $K \mapsto K^\wedge$ 表示关于 $(f)$ 的导出完备化。设 $M$
为 $A$-模。利用
$$
M^\wedge = R\lim (M \xrightarrow{f^n} M)
$$
（引理 \ref{more-algebra-lemma-derived-completion-koszul}）以及引理
\ref{more-algebra-lemma-break-long-exact-sequence-modules}，得到
$$
H^{-1}(M^\wedge) = \lim M[f^n] = T_f(M)
$$
，即{\it $M$ 的 $f$-进 Tate 模}。这里，映射
$M[f^n] \to M[f^{n - 1}]$ 由乘以 $f$ 给出。于是有描述
$H^0(M^\wedge)$ 的短正合列
$$
0 \to R^1\lim M[f^n] \to H^0(M^\wedge) \to \lim M/f^n M \to 0
$$
。因为转移映射是满射（引理 \ref{more-algebra-lemma-compute-Rlim-modules}），
有 $H^1(M^\wedge) = R^1\lim M/f^nM = 0$。由于平凡的理由，
$M^\wedge$ 的所有其他上同调也为零。最后，对
$K \in D(A)$ 与 $p \in \mathbf{Z}$，有短正合列
$$
0 \to H^0(H^p(K)^\wedge) \to H^p(K^\wedge) \to T_f(H^{p + 1}(K)) \to 0
$$
。这由例 \ref{more-algebra-example-derived-completion-spectral-sequence} 的谱序列
得出，因为它在 $E_2$ 退化（只有 $i = -1, 0$ 给出非零项）；
下一个引理给出更多信息。
\end{example}

\begin{lemma}
\label{more-algebra-lemma-compare-completion-derived-completion}
设 $A$ 为环，$f \in A$，$K$ 为 $D(A)$ 的对象。
% P02-E0433
置 $K_n = K \otimes_A^\mathbf{L} (A \xrightarrow{f^n} A)$。
对所有 $p \in \mathbf{Z}$，存在交换图
$$
\xymatrix{
& 0 & 0 \\
0 \ar[r] &
\widehat{H^p(K)} \ar[r] \ar[u] &
\lim H^p(K_n) \ar[r] \ar[u] &
T_f(H^{p + 1}(K)) \ar[r] &
0 \\
0 \ar[r] &
H^0(H^p(K)^\wedge) \ar[r] \ar[u] &
H^p(K^\wedge) \ar[r] \ar[u] &
T_f(H^{p + 1}(K)) \ar[r] \ar@{=}[u] &
0 \\
&
R^1\lim H^p(K)[f^n] \ar[u] \ar[r]^\cong &
R^1\lim H^{p - 1}(K_n) \ar[u] \\
& 0 \ar[u] & 0 \ar[u]
}
$$
，其行与列均正合，其中
$\widehat{H^p(K)} = \lim H^p(K)/f^nH^p(K)$ 是通常的 $f$-进完备化。
左侧竖直短正合列与中间水平短正合列取自例
% P02-E0434
\ref{more-algebra-example-spectral-sequence-principal}。中间竖直短正合列来自引理
\ref{more-algebra-lemma-break-long-exact-sequence-modules}。
\end{lemma}

\begin{proof}
为构造上方水平短正合列，注意我们有如下逆系的短正合列
$$
0 \to H^p(K)/f^nH^p(K) \to H^p(K_n) \to H^{p + 1}(K)[f^n] \to 0
$$
，它来自把 $K_n$ 构造为 $f^n : K \to K$ 的锥的移位。
取这些逆极限，便得到上方水平短正合列；见 Homology，引理
\ref{homology-lemma-Mittag-Leffler}。

\medskip\noindent
我们来证明有引理中的交换图。考虑映射
$L = \tau_{\leq p}K \to K$。置
$L_n = L \otimes_A^\mathbf{L} (A \xrightarrow{f^n} A)$，
得到逆系的映射 $(L_n) \to (K_n)$，它诱导短正合列的映射
$$
\xymatrix{
0 & 0 \\
\lim H^p(L_n) \ar[r] \ar[u] &
\lim H^p(K_n) \ar[u] \\
H^p(L^\wedge) \ar[r] \ar[u] &
H^p(K^\wedge) \ar[u] \\
R^1\lim H^{p - 1}(L_n) \ar[r] \ar[u] &
R^1\lim H^{p - 1}(K_n) \ar[u] \\
0 \ar[u] & 0 \ar[u]
}
$$
。由于对 $i > p$ 有 $H^i(L) = 0$，且
$H^p(L) = H^p(K)$，利用引理陈述中的参考结果进行计算可得
$H^p(L^\wedge) = H^0(H^p(K)^\wedge)$，以及
$H^p(L_n) = H^p(K)/f^nH^p(K)$。另一方面，有
$H^{p - 1}(L_n) = H^{p - 1}(K_n)$，所以得到引理陈述中所示的
同构，因为已经知道
$H^0(H^p(K)^\wedge) \to \widehat{H^p(K)}$ 的核等于
$R^1\lim H^p(K)[f^n]$。若用证明第一段的构造定义顶行，
我们略去验证图中最右方正方形交换。
\end{proof}

\begin{remark}
\label{more-algebra-remark-ML-condition}
沿用引理 \ref{more-algebra-lemma-compare-completion-derived-completion} 的记号，
我们还看到，逆系 $H^p(K_n)$ 具有 ML 性质，当且仅当逆系
$H^{p + 1}(K)[f^n]$ 具有 ML 性质。这由逆系的短正合列
$0 \to H^p(K)/f^nH^p(K) \to H^p(K_n) \to H^{p + 1}(K)[f^n] \to 0$
（见该引理的证明），结合 Homology，引理
\ref{homology-lemma-Mittag-Leffler} 与引理
\ref{more-algebra-lemma-Mittag-Leffler} 得出。
\end{remark}

\begin{lemma}[Bhatt]
\label{more-algebra-lemma-torsion-and-derived-complete}
设 $I$ 为环 $A$ 中的有限生成理想。设 $M$ 为导出完备 $A$-模。
若 $M$ 是 $I$-幂挠模，则对某个 $n$ 有 $I^nM = 0$。
\end{lemma}

\begin{proof}
设 $I = (f_1, \ldots, f_r)$。只需证明对每个 $i$，都存在
$n_i$，使 $f_i^{n_i}M = 0$。因此可假设 $I = (f)$ 为主理想。
设 $B = \mathbf{Z}[x] \to A$ 为将 $x$ 送到 $f$ 的环映射。
由引理 \ref{more-algebra-lemma-restriction-derived-complete}，$M$ 作为
$B$-模关于理想 $(x)$ 导出完备。用 $B$ 代替 $A$ 后，可假设
$f$ 是 $A$ 中的非零因子。

\medskip\noindent
假设 $I = (f)$，其中 $f \in A$ 为非零因子。根据例
\ref{more-algebra-example-derived-complete-modules}，存在短正合列
$$
0 \to K \xrightarrow{u} L \to M \to 0
$$
，其中 $K$ 与 $L$ 是 $I$-进完备 $A$-模，且其 $f$-挠为零
\footnote{对该证明，只需说明存在序列
$K \xrightarrow{u} L \to M \to 0$，其中 $K$ 与 $L$ 是
$I$-进完备 $A$-模。为此可选取一个表示
$F_1 \to F_0 \to M \to 0$，其中 $F_i$ 自由，然后分别令
$K$ 与 $L$ 等于 $F_1$ 与 $F_0$ 的 $f$-进完备化。事实上，
由于 $f$ 是非零因子，这些完备化就是导出完备化，且该序列
仍然正合。}。把 $K$ 与 $L$ 看作带 $I$-进拓扑的拓扑模。
则 $u$ 连续。令
$$
L_n = \{x \in L \mid f^n x \in u(K)\}
$$
。由于 $M$ 是 $f$-幂挠的，可见 $L = \bigcup L_n$。
设 $N_n$ 为 $L_n$ 在 $L$ 中的闭包。由引理
\ref{more-algebra-lemma-consequence-baire-complete-module}，对某个 $n$，
$N_n$ 在 $L$ 中开。固定这样的 $n$。由于
$f^{n + m} : L \to L$ 是连续开映射，且
$f^{n + m} L_n \subset u(f^m K)$，得出对所有 $m \geq 1$，
$u(f^mK)$ 的闭包都是开的。因此，由引理
\ref{more-algebra-lemma-open-mapping}，得出 $u$ 是开映射。故对某个 $t$，
$f^tL \subset \Im(u)$，从而 $f^t$ 消去 $M$，如所需。
\end{proof}

\begin{lemma}
\label{more-algebra-lemma-kernel-to-completion-square-zero}
设 $f \in A$ 为环中的元素，置 $J = \bigcap f^nA$。设 $M$
为关于 $f$ 导出完备的 $A$-模。则 $JM' = 0$，其中
$M' = \Ker(M \to \lim M/f^nM)$。特别地，若 $A$ 导出完备，
则 $J$ 是平方零理想。
\end{lemma}

\begin{proof}
取 $x \in M'$ 与 $g \in J$。对每个 $n \geq 1$，可写
$x = f^n x_n$。由于 $g$ 属于 $f^nA$，可见 $M'$ 中的元素
$y_n = gx_n$ 与 $x_n$ 的选择无关。特别地，可取
$x_n = fx_{n + 1}$，于是得到 $y_n = fy_{n + 1}$。因此得到映射
$A_f \to M$，它把 $1/f^n$ 送到 $y_n$。由于 $M$ 导出完备，
此映射必须为零（引理 \ref{more-algebra-lemma-hom-from-Af}）
% P02-E0435
，因而对所有 $n$ 都有 $y_n = 0$。由于
$gx = gfx_1 = fy_1$，证明完成。
\end{proof}

\begin{lemma}
\label{more-algebra-lemma-derived-complete-henselian}
设 $A$ 为关于理想 $I$ 导出完备的环。则 $(A, I)$ 是 Hensel 对。
\end{lemma}

\begin{proof}
设 $f \in I$。由引理 \ref{more-algebra-lemma-largest-ideal-henselian}，
只需证明 $(A, fA)$ 是 Hensel 对。注意 $A$ 关于 $fA$ 导出完备
（这立即由定义 \ref{more-algebra-definition-derived-complete} 得出）。
由引理 \ref{more-algebra-lemma-complete-derived-complete}，从 $A$ 到
$A$ 的 $f$-进完备化 $A'$ 的映射是满射。由引理
\ref{more-algebra-lemma-complete-henselian}，对 $(A', fA')$ 是 Hensel 对。
因此，只需证明 $(A, \bigcap f^nA)$ 是 Hensel 对；见引理
\ref{more-algebra-lemma-henselian-henselian-pair}。这由引理
\ref{more-algebra-lemma-kernel-to-completion-square-zero} 与
\ref{more-algebra-lemma-locally-nilpotent-henselian} 得出。
\end{proof}

\begin{lemma}
\label{more-algebra-lemma-kernel-to-completion-nilpotent}
设 $A$ 为关于理想 $I$ 导出完备的环。置 $J = \bigcap I^n$。
若 $I$ 可由 $r$ 个元素生成，则 $J^N = 0$，其中 $N = 2^r$。
\end{lemma}

\begin{proof}
当 $r = 1$ 时，这是引理
\ref{more-algebra-lemma-kernel-to-completion-square-zero}。设
$I = (f_1, \ldots, f_r)$，其中 $r > 1$。由引理
\ref{more-algebra-lemma-serre-subcategory}，环 $A_t = A/f_r^tA$ 关于 $I$
导出完备，因而更关于
$I_t = (f_1, \ldots, f_{r - 1})A_t$ 导出完备。注意
$A \to A_t$ 把 $J$ 送入 $J_t = \bigcap I_t^n$。由归纳，
$J_t^{N/2} = 0$，其中 $N = 2^r$。理想
$\bigcap \Ker(A \to A_t) = \bigcap f_r^t A$ 由 $r = 1$
的情形是平方零的。这就完成证明。
\end{proof}

\begin{lemma}
\label{more-algebra-lemma-reduced-derived-complete-complete}
设 $A$ 为关于有限生成理想 $I$ 导出完备的既约环。
则 $A$ 是 $I$-进完备的。
\end{lemma}

\begin{proof}
% P02-E0436
这由引理 \ref{more-algebra-lemma-kernel-to-completion-nilpotent} 与命题
\ref{more-algebra-proposition-derived-complete-modules} 得出。
\end{proof}

\section{Noether 环的导出完备化}
\label{more-algebra-section-derived-completion-noetherian}

\noindent
设 $A$ 为环，$I \subset A$ 为理想。对任意 $K \in D(A)$，
可考虑导出极限
$$
K' = R\lim (K \otimes_A^\mathbf{L} A/I^n)
$$
。这是关于 $K$ 的函子，见评注
\ref{more-algebra-remark-constructing-tensor-with-limits-functorially}。映射系
$A \to A/I^n$ 诱导映射 $K \to K'$，且 $K'$ 关于 $I$
导出完备（引理 \ref{more-algebra-lemma-naive-derived-completion}）。
一般而言，这一“朴素”导出完备化构造并不等于引理
\ref{more-algebra-lemma-derived-completion} 的伴随。例如，若
$A = \mathbf{Z}_p \oplus \mathbf{Q}_p/\mathbf{Z}_p$，且第二个直和项
为平方零理想，$K = A[0]$，$I = (p)$，则朴素导出完备化给出
$\mathbf{Z}_p[0]$，而引理 \ref{more-algebra-lemma-derived-completion} 的构造给出
$K^\wedge \cong \mathbf{Z}_p[1] \oplus \mathbf{Z}_p[0]$
（略去计算）。当 $I$ 由单个元素生成时，引理
\ref{more-algebra-lemma-when-does-it-work} 刻画了这两个函子何时一致。

\medskip\noindent
% P02-E0437
本节的主要目标是证明：若 $A$ 为 Noether 环，则朴素导出完备化
等于导出完备化。

\begin{lemma}
\label{more-algebra-lemma-sequence-Koszul-complexes}
在情形 \ref{more-algebra-situation-koszul} 中。若 $A$ 为 Noether 环，则
$D(A)$ 的 pro-对象 $\{K_n^\bullet\}$ 与
$\{A/(f_1^n, \ldots, f_r^n)\}$ 同构
\footnote{特别地，对每个 $n$，都存在 $m \geq n$，使
$K_m^\bullet \to K_n^\bullet$ 经映射
$K_m^\bullet \to A/(f_1^m, \ldots, f_r^m)$ 分解。}。
\end{lemma}

\begin{proof}
% P02-E0438
我们有杰出三角形的一个逆系
$$
\tau_{\leq -1}K_n^\bullet \to K_n^\bullet \to
A/(f_1^m, \ldots, f_r^m) \to
(\tau_{\leq -1}K_n^\bullet)[1]
$$
，见 Derived Categories，评注
\ref{derived-remark-truncation-distinguished-triangle}。由 Derived
Categories，引理 \ref{derived-lemma-pro-isomorphism}，只需证明逆系
$\tau_{\leq -1}K_n^\bullet$ 是 pro-零的。回忆 $K_n^\bullet$
仅在满足 $-r \leq i \leq 0$ 的次数 $i$ 上有非零项。因此，由
Derived Categories，引理
\ref{derived-lemma-essentially-constant-cohomology}，只需证明对
$p \leq -1$，$H^p(K_n^\bullet)$ 是 pro-零的。换言之，对每个
$n \in \mathbf{N}$，必须证明存在 $m \geq n$，使
$H^p(K_m^\bullet) \to H^p(K_n^\bullet)$ 为零。由于 $A$
为 Noether 环，可见
$$
H^p(K_n^\bullet) =
\frac{\Ker(K_n^p \to K_n^{p + 1})}{\Im(K_n^{p - 1} \to K_n^p)}
$$
是有限 $A$-模。此外，当 $p < 0$ 时，映射
$K_m^p \to K_n^p$ 由一个对角矩阵给出，其各项属于理想
$(f_1^{m - n}, \ldots, f_r^{m - n})$。注意
$H^p(K_n^\bullet)$ 被
$J = (f_1^n, \ldots, f_r^n)$ 消去，见引理
\ref{more-algebra-lemma-homotopy-koszul}。现在，当 $m - n \geq tn$ 时，
有 $(f_1^{m - n}, \ldots, f_r^{m - n}) \subset J^t$。因此，把
Algebra，引理 \ref{algebra-lemma-Artin-Rees}（Artin--Rees）
应用于理想 $J$ 和模 $M = K_n^p$ 及其子模
% P02-E0439
$N = \Ker(K_n^p \to K_n^{p + 1})$，则当 $m$ 充分大时，
$K_m^p \to K_n^p$ 的像与
$\Ker(K_n^p \to K_n^{p + 1})$ 的交包含于
$J \Ker(K_n^p \to K_n^{p + 1})$。对这样的 $m$，得到零映射。
\end{proof}

\begin{proposition}
\label{more-algebra-proposition-noetherian-naive-completion-is-completion}
设 $A$ 为 Noether 环，$I \subset A$ 为理想。把
$K \in D(A)$ 送到导出极限
$K' = R\lim( K \otimes_A^\mathbf{L} A/I^n )$ 的函子，是引理
\ref{more-algebra-lemma-derived-completion} 中所构造的包含函子
$D_{comp}(A) \to D(A)$ 的左伴随。
\end{proposition}

\begin{proof}
设 $(f_1, \ldots, f_r) = I$，并令 $K_n^\bullet$ 为关于
$f_1^n, \ldots, f_r^n$ 的 Koszul 复形。由引理
\ref{more-algebra-lemma-derived-completion-koszul}，只需证明
$$
R\lim (K \otimes_A^\mathbf{L} K_n^\bullet) =
R\lim (K \otimes_A^\mathbf{L} A/(f_1^n, \ldots, f_r^n) ) =
R\lim (K \otimes_A^\mathbf{L} A/I^n ).
$$
由引理 \ref{more-algebra-lemma-sequence-Koszul-complexes}，$D(A)$ 的 pro-对象
$\{K_n^\bullet\}$ 与 $\{A/(f_1^n, \ldots, f_r^n)\}$ 同构。
显然，pro-对象 $\{A/(f_1^n, \ldots, f_r^n)\}$ 与
$\{A/I^n\}$ 同构。因此，由引理
\ref{more-algebra-lemma-tensor-Rlim-pro-equal}，从左到右的映射是同构。
\end{proof}

\begin{lemma}
\label{more-algebra-lemma-noetherian-calculate}
设 $I$ 为 Noether 环 $A$ 的理想，$M$ 为 $A$-模，其导出完备化为
$M^\wedge$。则存在短正合列
$$
0 \to R^1\lim \text{Tor}_{i + 1}^A(M, A/I^n) \to
H^{-i}(M^\wedge) \to \lim \text{Tor}_i^A(M, A/I^n) \to 0
$$
。对 $M \in D^-(A)$ 也有类似结果。
\end{lemma}

\begin{proof}
% P02-E0440
这是命题 \ref{more-algebra-proposition-noetherian-naive-completion-is-completion}
与引理 \ref{more-algebra-lemma-break-long-exact-sequence-modules} 的直接推论。
\end{proof}

\noindent
作为上述命题的一个应用，我们在 Noether 情形下为伪凝聚复形
识别导出完备化。

\begin{lemma}
\label{more-algebra-lemma-derived-completion-pseudo-coherent}
设 $A$ 为 Noether 环，$I \subset A$ 为理想。设 $K$ 为 $D(A)$
的对象，使得对所有
% P02-E0441
$n \in \mathbf{Z}$，$H^n(K)$ 都是有限 $A$-模。则导出完备化的
上同调模 $H^n(K^\wedge)$ 是上同调模 $H^n(K)$
的 $I$-进完备化。
\end{lemma}

\begin{proof}
由引理 \ref{more-algebra-lemma-Noetherian-pseudo-coherent}，对所有 $m$，
复形 $\tau_{\leq m}K$ 都是伪凝聚的。因此
$\tau_{\leq m}K$ 可由有限自由 $A$-模的有界上复形
$P^\bullet$ 表示。于是
$\tau_{\leq m}K \otimes_A^\mathbf{L} A/I^n = P^\bullet/I^nP^\bullet$。
故
$(\tau_{\leq m}K)^\wedge = R\lim P^\bullet/I^nP^\bullet$
（命题 \ref{more-algebra-proposition-noetherian-naive-completion-is-completion}）；
又因为 $R\lim$ 由逐项的 $\lim$ 给出
（引理 \ref{more-algebra-lemma-compute-Rlim-modules}），且 $I$-进完备化在有限
$A$-模上是正合函子
（Algebra，引理 \ref{algebra-lemma-completion-flat}），得出该结果
对 $\tau_{\leq m}K$ 成立。由于导出完备化具有有限上同调维数，
该结果因而对 $K$ 成立；见引理
\ref{more-algebra-lemma-derived-completion-finite-cohomological-dimension}。
\end{proof}

\begin{lemma}
\label{more-algebra-lemma-derived-complete-finite}
设 $I$ 为 Noether 环 $A$ 的理想。设 $M$ 为导出完备 $A$-模。
若 $M/IM$ 是有限 $A/I$-模，则
$M = \lim M/I^nM$，且 $M$ 是有限 $A^\wedge$-模。
\end{lemma}

\begin{proof}
假设 $M/IM$ 有限。选取 $x_1, \ldots, x_t \in M$，使其像生成
$M/IM$。得到映射 $A^{\oplus t} \to M$，它把第 $i$ 个基向量
送到 $x_i$。由命题
\ref{more-algebra-proposition-noetherian-naive-completion-is-completion}，$A$
的导出完备化为 $A^\wedge = \lim A/I^n$。由于 $M$ 导出完备，
可见我们的映射经一个映射
$q : (A^\wedge)^{\oplus t} \to M$ 分解。由引理
\ref{more-algebra-lemma-derived-complete-zero}，模 $\Coker(q)$ 为零。
因此 $M$ 是有限 $A^\wedge$-模。由于 $A^\wedge$ 是 Noether 环，
且关于 $IA^\wedge$ 完备，得出 $M$ 是 $I$-进完备的
（使用 Algebra，引理
\ref{algebra-lemma-completion-Noetherian}、
\ref{algebra-lemma-when-finite-module-complete-over-complete-ring} 以及
\ref{algebra-lemma-Artin-Rees}）。
\end{proof}

\begin{lemma}
\label{more-algebra-lemma-when-derived-completion-is-completion}
设 $I$ 为 Noether 环 $A$ 中的理想。
\begin{enumerate}
\item 若 $M$ 是有限 $A$-模且 $N$ 是平坦 $A$-模，则
$M \otimes_A N$ 的导出 $I$-进完备化就是
$M \otimes_A N$ 的通常 $I$-进完备化。
\item 若 $M$ 是有限 $A$-模且 $f \in A$，则 $M_f$ 的导出
$I$-进完备化就是 $M_f$ 的通常 $I$-进完备化。
\end{enumerate}
\end{lemma}

\begin{proof}
对 $A$-模 $M$，
% P02-E0442
用 $M^\wedge$ 表示导出完备化，用 $\lim M/I^nM$ 表示通常完备化。
假设 $M$ 有限。对 $i > 0$，系
$\text{Tor}^A_i(M, A/I^n)$ 是 pro-零的，见引理
\ref{more-algebra-lemma-tor-strictly-pro-zero}。因为 $N$ 平坦，
$\text{Tor}_i^A(M \otimes_A N, A/I^n) =
\text{Tor}_i^A(M, A/I^n) \otimes_A N$，所以系
$\text{Tor}^A_i(M \otimes_A N, A/I^n)$ 也如此。
由引理 \ref{more-algebra-lemma-noetherian-calculate}，得出
$R\lim (M \otimes_A N) \otimes_A^\mathbf{L} A/I^n$
仅在次数 $0$ 有上同调，且该上同调由通常完备化
$\lim M \otimes_A N/ I^n(M \otimes_A N)$ 给出。这证明 (1)。
(2) 由 (1) 以及 $M_f = M \otimes_A A_f$ 得出。
\end{proof}

\begin{lemma}
\label{more-algebra-lemma-derived-completion-tensor-finite}
设 $I$ 为 Noether 环 $A$ 中的理想。令
${}^\wedge$ 表示关于 $I$ 的导出完备化。设 $K \in D^-(A)$。
\begin{enumerate}
\item 若 $M$ 是有限 $A$-模，则
$(K \otimes_A^\mathbf{L} M)^\wedge = K^\wedge \otimes_A^\mathbf{L} M$。
\item 若 $L \in D(A)$ 是伪凝聚的，则
$(K \otimes_A^\mathbf{L} L)^\wedge = K^\wedge \otimes_A^\mathbf{L} L$。
\end{enumerate}
\end{lemma}

\begin{proof}
设 $L$ 如 (2) 中所述。可以用自由 $A$-模的有界上复形
$P^\bullet$ 表示 $K$，并用有限自由 $A$-模的有界上复形
$F^\bullet$ 表示 $L$。由于
$\text{Tot}(P^\bullet \otimes_A F^\bullet)$ 表示
$K \otimes_A^\mathbf{L} L$，可见
$(K \otimes_A^\mathbf{L} L)^\wedge$ 由
$$
\text{Tot}((P^\bullet)^\wedge \otimes_A F^\bullet)
$$
表示，其中 $(P^\bullet)^\wedge$ 的各项是
% P02-E0443
通常的 $=$ 导出完备化 $(P^n)^\wedge$；例如见命题
\ref{more-algebra-proposition-noetherian-naive-completion-is-completion} 与引理
\ref{more-algebra-lemma-when-derived-completion-is-completion}。这证明 (2)。
(1) 是 (2) 的特殊情形。
\end{proof}

\section{Berthelot 与 Ogus 引入的算子}
\label{more-algebra-section-eta}

\noindent
本节讨论 \cite[Section 8]{Berthelot-Ogus} 中引入并在
\cite[Section 6]{BMS} 中推广的一个构造。我们强烈建议读者查看
讨论这一概念的原始论文。

\medskip\noindent
设 $A$ 为环，$f \in A$ 为非零因子。若 $M$ 为 $A$-模，则由引理
% P02-E0444
\ref{more-algebra-lemma-torsion-free}，以下条件等价：
\begin{enumerate}
\item $f$ 在 $M$ 上是非零因子；
\item $M[f] = 0$；
\item 对所有 $n \geq 1$，有 $M[f^n] = 0$；以及
\item 映射 $M \to M_f$ 是单射。
\end{enumerate}
若这些等价条件成立，则（在本节中）我们称{\it $M$ 无 $f$-挠}。
在这种情况下，
% P02-E0445
用 $f^iM \subset M_f$ 表示由形如 $f^ix$（其中 $x \in M$）
的元素组成的子模。当然，$f^iM$ 作为 $A$-模同构于 $M$。
设 $M^\bullet$ 为由无 $f$-挠 $A$-模组成的复形，其微分为
$d^i : M^i \to M^{i + 1}$。在这种情形下，定义
$\eta_fM^\bullet$ 为具有各项
$$
(\eta_fM)^i = \{x \in f^iM^i \mid d^i(x) \in f^{i + 1}M^{i + 1}\}
$$
以及由 $d^i$ 诱导的微分的复形。注意 $\eta_fM^\bullet$
是另一个由无 $f$-挠 $A$-模组成的复形。若
$a^\bullet : M^\bullet \to N^\bullet$ 是无 $f$-挠 $A$-模复形的
映射，则由映射 $f^iM^i \to f^iN^i$ 诱导复形映射
$$
\eta_fa^\bullet : \eta_fM^\bullet \longrightarrow \eta_fN^\bullet
$$
。读者可检验，我们在无 $f$-挠 $A$-模复形的范畴上得到一个自函子。
若 $a^\bullet, b^\bullet : M^\bullet \to N^\bullet$ 是无
$f$-挠 $A$-模复形的两个映射，而
$h = \{h^i : M^i \to N^{i - 1}\}$ 是 $a^\bullet$ 与 $b^\bullet$
之间的伦，则定义 $\eta_fh$ 为映射族
$(\eta_fh)^i : (\eta_fM)^i \to (\eta_fN)^{i - 1}$，它把 $x$
送到 $h^i(x)$；这是有意义的，因为 $x \in f^iM^i$ 蕴含
$h^i(x) \in f^iN^{i - 1}$，而后者当然包含于
$(\eta_fN)^{i - 1}$。读者可检验，$\eta_fh$ 是
$\eta_fa^\bullet$ 与 $\eta_fb^\bullet$ 之间的伦。总之，可见在无
$f$-挠 $A$-模这一加性范畴的伦范畴
（Derived Categories，第 \ref{derived-section-homotopy} 节）上，
我们得到函子
$$
\eta_f :
K(f\text{-torsion free }A\text{-模})
\longrightarrow
K(f\text{-torsion free }A\text{-模})
$$
。$\eta_f$ 在任何意义下都不是三角范畴的正合函子；见例
\ref{more-algebra-example-eta-not-distinguished}。

\begin{example}
\label{more-algebra-example-eta-not-distinguished}
设 $A$ 为环，$f \in A$ 为非零因子。考虑函子
$\eta_f : K(f\text{-torsion free }A\text{-模})
\to K(f\text{-torsion free }A\text{-模})$。设
$M^\bullet$ 为无 $f$-挠 $A$-模的复形。乘以 $f$ 定义同构
$\eta_f(M^\bullet[1]) \to (\eta_fM^\bullet)[1]$，所以在这一意义下
$\eta_f$ 与移位相容。然而，考虑图
$$
\xymatrix{
A \ar[r]_f &
A \ar[r]_1 &
A \ar[r] &
0 \\
0 \ar[r] \ar[u] &
0 \ar[r] \ar[u] &
A \ar[r]^{-1} \ar[u]^f &
A \ar[u]
}
$$
。把每一列看作无 $f$-挠 $A$-模的复形，其中上方的模在次数 $1$，
其下的模在次数 $0$。于是，此图在
$K(f\text{-torsion free }A\text{-模})$ 中给出一个杰出三角形，
其三角结构如 Derived Categories，第
\ref{derived-section-homotopy-triangulated} 节中所述。具体而言，
第三个复形是前两个复形之间映射的锥。然而，对每一列施以
$\eta_f$，得到
$$
\xymatrix{
fA \ar[r]_f &
fA \ar[r]_1 &
fA \ar[r] &
0 \\
0 \ar[r] \ar[u] &
0 \ar[r] \ar[u] &
A \ar[r]^{-1} \ar[u]^f &
A \ar[u]
}
$$
。然而，第三个复形是非周期的，甚至零伦。因此，若这是杰出三角形，
则第一个箭头必须是伦范畴中的同构；除非 $f$ 是单位，否则这不成立。
\end{example}

\begin{lemma}
\label{more-algebra-lemma-eta-first-property}
设 $A$ 为环，$f \in A$ 为非零因子。设 $M^\bullet$ 为无
$f$-挠 $A$-模的复形。存在由乘以 $f^i$ 给出的典范同构
$$
f^i : H^i(M^\bullet)/H^i(M^\bullet)[f] \longrightarrow H^i(\eta_fM^\bullet)
$$
。
\end{lemma}

\begin{proof}
注意
$\Ker(d^i : (\eta_fM)^i \to (\eta_fM)^{i + 1})$ 等于
$\Ker(d^i : f^iM^i \to f^iM^{i + 1}) =
f^i\Ker(d^i : M^i \to M^{i + 1})$。因此，
% P02-E0446
把 $z \in \Ker(d^i : M^i \to M^{i + 1})$ 的类送到 $f^iz$ 的类，
得到满射
$f^i : H^i(M^\bullet) \to H^i(\eta_fM^\bullet)$。若由此在
$H^i(\eta_fM^\bullet)$ 中得到零类，则对某个
$y \in M^{i - 1}$，有
$f^i z = d^{i - 1}(f^{i - 1}y)$。因为 $f$ 在所涉及的所有模上
都是非零因子，这意味着 $f z = d^{i - 1}(y)$，而这恰好意味着
$z$ 的类是 $f$-挠的，如所需。
\end{proof}

\begin{lemma}
\label{more-algebra-lemma-eta-qis}
设 $A$ 为环，$f \in A$ 为非零因子。若
$M^\bullet \to N^\bullet$ 是无 $f$-挠 $A$-模复形的拟同构，
则诱导映射 $\eta_fM^\bullet \to \eta_fN^\bullet$ 也是拟同构。
\end{lemma}

\begin{proof}
这是因为引理 \ref{more-algebra-lemma-eta-first-property} 的同构与复形映射相容。
\end{proof}

\begin{lemma}
\label{more-algebra-lemma-Leta}
设 $A$ 为环，$f \in A$ 为非零因子。存在加性函子
\footnote{注意，此函子不正合，也就是说，它不把杰出三角形变换为
杰出三角形。见例 \ref{more-algebra-example-eta-not-distinguished}。}
$L\eta_f : D(A) \to D(A)$，使得若 $M \in D(A)$ 由无
$f$-挠 $A$-模复形 $M^\bullet$ 表示，则
$L\eta_fM = \eta_fM^\bullet$；对态射也类似。
\end{lemma}

\begin{proof}
% P02-E0447
用 $\mathcal{T} \subset \text{Mod}_A$ 表示无 $f$-挠 $A$-模的
全子范畴。相应地，有 $K(\mathcal{T})$ 作为 $K(A)$ 的全三角
子范畴的包含
$$
K(\mathcal{T}) \quad\subset\quad K(\text{Mod}_A) = K(A)
$$
。设 $S \subset \text{Arrows}(K(\mathcal{T}))$ 为拟同构。
我们将应用 Derived Categories，引理
\ref{derived-lemma-localization-subcategory} 证明映射
$$
S^{-1}K(\mathcal{T}) \longrightarrow D(A)
$$
是三角范畴的等价。该引理表明，只需证明：给定 $A$-模复形
$M^\bullet$，存在拟同构 $K^\bullet \to M^\bullet$，其中
$K^\bullet$ 是无 $f$-挠模的复形。由引理
\ref{more-algebra-lemma-K-flat-resolution}，可找到拟同构
$K^\bullet \to M^\bullet$，使复形 $K^\bullet$ 是 K-平坦的
（我们不会用到这一点），且由平坦 $A$-模 $K^i$ 组成。
特别地，对所有 $i$，$f$ 在 $K^i$ 上都是非零因子，如所需。

\medskip\noindent
准备工作完成后，可以定义 $L\eta_f$。事实上，由本节开头的讨论，
我们已有良定义函子
$$
K(\mathcal{T}) \xrightarrow{\eta_f} K(\mathcal{T}) \to K(A) \to D(A)
$$
，而根据引理 \ref{more-algebra-lemma-eta-qis}，它把拟同构送到拟同构。因此，
% P02-E0448
由 Categories，引理
\ref{categories-lemma-properties-left-localization}，此函子经
$S^{-1}K(\mathcal{T}) = D(A)$ 分解。
\end{proof}

\begin{remark}
\label{more-algebra-remark-eta-BZ}
设 $A$ 为环，$f \in A$ 为非零因子。设 $M^\bullet$ 为无
$f$-挠 $A$-模的复形。对每个 $i$，置
$\overline{M}^i = M^i/fM^i$。
% P02-E0449
用 $B^i \subset Z^i \subset \overline{M}^i$ 表示复形
$\overline{M}^\bullet = M^\bullet \otimes_A A/fA$ 的微分的
上边界与上循环。我们断言存在行正合的交换图
$$
\xymatrix{
0 \ar[r] &
B^{i + 1} \ar[r] \ar@{=}[d] &
B^{i + 1} \oplus B^i \ar[r] \ar[d]^{s, s'} &
B^i \ar[r] \ar[d] & 0 \\
0 \ar[r] &
B^{i + 1} \ar[r]^-s &
(\eta_fM)^i /f(\eta_fM)^i \ar[r]^-t &
Z^i \ar[r] & 0
}
$$
。这些映射构造如下：
\begin{enumerate}
\item 若 $x \in (\eta_fM)^i$，则 $x = f^ix'$，且
$d^i(x') = 0$ 在 $\overline{M}^{i + 1}$ 中成立。因此，可把 $x$
送到 $x'$ 的类来定义映射 $t$。
\item 若 $y \in M^{i + 1}$ 在
$B^{i + 1} \subset \overline{M}^{i + 1}$ 中的类为
$\overline{y}$，则可对某些 $y' \in M^{i + 1}$ 与 $x \in M^i$
写成 $y = fy' + d^i(x)$。因此，可把 $\overline{y}$ 送到
$(\eta_fM)^i /f(\eta_fM)^i$ 中 $f^{i + 1}x$ 的类来定义映射
$s$；我们略去该映射良定义的验证。
\item 若 $x \in M^i$ 在 $B^i \subset \overline{M}^i$ 中的类为
$\overline{x}$，则可对某些 $x' \in M^i$ 与
$z \in M^{i - 1}$ 写成
$x = fx' + d^{i - 1}(z)$。定义映射 $s'$，把 $\overline{x}$
送到 $(\eta_fM)^i/f(\eta_fM)^i$ 中 $f^i d^{i - 1}(z)$ 的类。
这是良定义的，因为若 $fx' + d^{i - 1}(z) = 0$，则
$f^ix'$ 属于 $(\eta_fM)^i$，因而
$f^id^{i - 1}(z)$ 属于 $f(\eta_fM)^i$。
\end{enumerate}
我们略去验证展示图的下行是模的短正合列。由这些构造立即可见，
我们有交换图
$$
\xymatrix{
B^{i + 1} \oplus B^i \ar[d]^{s, s'} \ar[r] &
B^{i + 2} \oplus B^{i + 1} \ar[d]^{s, s'} \\
(\eta_fM)^i /f(\eta_fM)^i \ar[r] &
(\eta_fM)^{i + 1} /f(\eta_fM)^{i + 1}
}
$$
，其中上方水平箭头由源与靶中直和项 $B^{i + 1}$ 的同一给出。
换言之，我们找到了
$\eta_fM^\bullet / f(\eta_fM^\bullet) =
\eta_fM^\bullet \otimes_A A/fA$ 的一个非周期子复形，而以此子复形
取商所得复形的各项 $Z^i/B^i$，是复形
$\overline{M}^\bullet = M^\bullet \otimes_A A/fA$ 的上同调模。
\end{remark}

\noindent
为以更典范的方式解释评注 \ref{more-algebra-remark-eta-BZ} 中观察到的现象，
我们将构造 Bockstein 算子。设 $A$ 为环，$f \in A$ 为非零因子。
设 $M^\bullet$ 为无 $f$-挠 $A$-模的复形。对每个
$i \in \mathbf{Z}$，都有交换图（张量积均在 $A$ 上取）
$$
\xymatrix{
0 \ar[r] &
M^\bullet \otimes f^{i + 1}A \ar[r] \ar[d] &
M^\bullet \otimes f^iA \ar[r] \ar[d] &
M^\bullet \otimes f^iA/f^{i + 1}A \ar[r] \ar@{=}[d] &
0 \\
0 \ar[r] &
M^\bullet \otimes f^{i + 1}A/f^{i + 2}A \ar[r] &
M^\bullet \otimes f^iA/f^{i + 2}A \ar[r] &
M^\bullet \otimes f^iA/f^{i + 1}A \ar[r] &
0
}
$$
，其两行都是复形的短正合列。当然，通过适当地乘以 $f$ 的幂，
不同 $i$ 的这些短正合列都彼此同构。特别地，下方序列的上同调
长正合列对所有 $i \in \mathbf{Z}$ 确定{\it Bockstein 算子}
$$
\beta = \beta^i : H^i(M^\bullet \otimes f^iA/f^{i + 1}A) \to
H^{i + 1}(M^\bullet \otimes f^{i + 1}A/f^{i + 2}A)
$$
。为供后文使用，我们在此记录：由上面的交换图，Bockstein 算子
具有分解
\begin{equation}
\label{more-algebra-equation-factorization-bockstein}
\vcenter{
\xymatrix{
H^i(M^\bullet \otimes f^iA/f^{i + 1}A)
\ar[r]_\delta \ar[rd]_\beta &
H^{i + 1}(M^\bullet \otimes f^{i + 1}A) \ar[d] \\
&
H^{i + 1}(M^\bullet \otimes f^{i + 1}A/f^{i + 2}A)
}
}
\end{equation}
，其中 $\delta$ 是来自上述交换图顶行的边界算子。
我们来证明得到一个复形
\begin{equation}
\label{more-algebra-equation-complex-bocksteins}
H^\bullet(M^\bullet/f) =
\left[
\begin{matrix}
\ldots \\
\downarrow \\
H^{i - 1}(M^\bullet \otimes f^{i - 1}A/f^iA) \\
\downarrow \beta \\
H^i(M^\bullet \otimes f^iA/f^{i + 1}A) \\
\downarrow \beta \\
H^{i + 1}(M^\bullet \otimes f^{i + 1}A/f^{i + 2}A) \\
\downarrow \\
\ldots
\end{matrix}
\right]
\end{equation}
，也就是说，$\beta \circ \beta = 0$
\footnote{另一种方法是论证 $\beta$ 作为复形
$(M^\bullet)_f$ 关于子复形 $f^iM^\bullet$ 的滤过之谱序列的
微分出现。还有一个证明更强结论的论证：先考虑
$M^\bullet = A$ 的情形。此时短正合列
$0 \to f^{i + 1}A/f^{i + 2}A \to f^iA/f^{i + 2}A \to
f^iA/f^{i + 1}A \to 0$ 在 $D(A)$ 中定义映射
$\beta^i : f^iA/f^{i + 1}A \to f^{i + 1}A/f^{i + 2}A[1]$。
然后可计算（论证与正文类似）出复合
$f^iA/f^{i + 1}A \to f^{i + 1}A/f^{i + 2}A[1] \to
f^{i + 2}A/f^{i + 3}A[2]$ 在 $D(A)$ 中为零。由于关于
$M^\bullet$ 的各项无 $f$-挠这一假设，有
$M^\bullet \otimes f^iA/f^{i + 1}A =
M^\bullet \otimes^\mathbf{L} f^iA/f^{i + 1}A$，所以得出
% P02-E0450
复合
$$
(M^\bullet \otimes f^iA/f^{i + 1}A)
\to
(M^\bullet \otimes f^{i + 1}A/f^{i + 2}A)[1]
\to
(M^\bullet \otimes f^{i + 2}A/f^{i + 3}A)[2]
$$
在 $D(A)$ 中也为零。}。事实上，利用分解
(\ref{more-algebra-equation-factorization-bockstein})，只需证明
$$
H^{i + 1}(M^\bullet \otimes f^{i + 1}A)
\to
H^{i + 1}(M^\bullet \otimes f^{i + 1}A/f^{i + 2}A)
\xrightarrow{\beta^{i + 1}}
H^{i + 2}(M^\bullet \otimes f^{i + 2}A/f^{i + 3}A)
$$
为零。这成立是因为 $\beta^{i + 1}$ 的核由可提升到
$H^{i + 1}(M^\bullet \otimes f^{i + 1}A/f^{i + 3}A)$ 的
上同调类组成，而第一个映射像中的类当然可以提升！

\begin{lemma}
\label{more-algebra-lemma-eta-second-property}
设 $A$ 为环，$f \in A$ 为非零因子。设 $M^\bullet$ 为无
$f$-挠 $A$-模的复形。存在复形的典范映射
$$
\eta_fM^\bullet \otimes_A A/fA
\longrightarrow
H^\bullet(M^\bullet/f)
$$
，它是拟同构，其中右端是复形
(\ref{more-algebra-equation-complex-bocksteins})。
\end{lemma}

\begin{proof}
设 $x \in (\eta_fM)^i$。则
% P02-E0451
$x = f^ix' \in f^iM$，且
$d^i(x) = f^{i + 1}y \in f^{i + 1}M^{i + 1}$。于是 $d^i$
把 $x' \otimes f^i$ 映到
$M^{i + 1} \otimes f^iA/f^{i + 1}A$ 中的零。本证明中的所有
张量积都在 $A$ 上取。因此，可把 $x$ 映到
$H^i(M^\bullet \otimes f^iA/f^{i + 1}A)$ 中
$x' \otimes f^i$ 的类。显然，此规则定义 $A/fA$-模的映射
$$
(\eta_fM)^i \otimes A/fA
\longrightarrow
H^i(M^\bullet \otimes f^iA/f^{i + 1}A)
$$
。注意，在上述情形中，可把 $x' \otimes f^i$ 看作
$M^i \otimes f^iA/f^{i + 2}A$ 的元素，其微分为
$d^i(x' \otimes f^i) = y \otimes f^{i + 1}$。由上面对
$\beta$ 的构造，得到
$\beta(x' \otimes f^i) = y \otimes f^{i + 1}$，因而我们的映射
与微分相容，也就是说，得到复形的映射。

\medskip\noindent
为完成证明，注意上一段给出的构造与评注 \ref{more-algebra-remark-eta-BZ}
中讨论的映射
$(\eta_fM)^i \otimes A/fA \to Z^i/B^i$ 一致。由于已经看到这些
映射的核是 $\eta_fM^\bullet \otimes A/fA$ 的非周期子复形，
引理得证。
\end{proof}

\begin{lemma}
\label{more-algebra-lemma-vanishing-beta}
设 $A$ 为环，$f \in A$ 为非零因子。设 $M^\bullet$ 为无
$f$-挠 $A$-模的复形。对 $i \in \mathbf{Z}$，以下条件等价：
\begin{enumerate}
\item $\Ker(d^i \bmod f^2)$ 满射到 $\Ker(d^i \bmod f)$；
\item 映射
$\beta : H^i(M^\bullet \otimes_A f^iA/f^{i + 1}A) \to
H^{i + 1}(M^\bullet \otimes_A f^{i + 1}A/f^{i + 2}A)$ 为零。
\end{enumerate}
条件 $H^{i + 1}(M^\bullet)[f] = 0$ 蕴含这些等价条件。
\end{lemma}

\begin{proof}
(1) 与 (2) 的等价性，由 $\beta$ 定义为一个复形短正合列的
上同调边界映射，以及该短正合列同构于复形短正合列
$0 \to fM^\bullet/f^2M^\bullet \to M^\bullet/f^2M^\bullet \to
M^\bullet/fM^\bullet \to 0$ 得出。若 $\beta \not = 0$，则由分解
(\ref{more-algebra-equation-factorization-bockstein})，
$H^{i + 1}(M^\bullet)[f] \not = 0$。
\end{proof}

\begin{lemma}
\label{more-algebra-lemma-eta-vanishing-beta}
设 $A$ 为环，$f \in A$ 为非零因子。设 $M^\bullet$ 为无
$f$-挠 $A$-模的复形。若 $\Ker(d^i \bmod f^2)$ 满射到
$\Ker(d^i \bmod f)$，则典范映射
$$
(1, d^i) :
(\eta_fM)^i / f(\eta_fM)^i \longrightarrow
f^iM^i/f^{i + 1}M^i \oplus f^{i + 1}M^{i + 1}/f^{i + 2}M^{i + 1}
$$
把左端同一为右端的若干子模的直和。
\end{lemma}

\begin{proof}
沿用评注 \ref{more-algebra-remark-eta-BZ} 的记号，定义映射
$t^{-1} : Z^i \to (\eta_fM)^i / f(\eta_fM)^i$。具体而言，
对满足 $d^i(x) = f^2y$ 的 $x \in M^i$，把 $x$ 在 $Z^i$ 中的类
送到 $(\eta_fM)^i / f(\eta_fM)^i$ 中 $f^ix$ 的类。我们略去验证
这是良定义的；引理的假设恰好意味着此运算的定义域是整个
$Z^i$。于是 $t \circ t^{-1} = \text{id}_{Z^i}$。故 $t^{-1}$
定义了评注 \ref{more-algebra-remark-eta-BZ} 中短正合列的一个分裂，所得直和分解
$$
(\eta_fM)^i / f(\eta_fM)^i = Z^i \oplus B^{i + 1}
$$
与引理中展示的映射相容。
\end{proof}

\begin{lemma}
\label{more-algebra-lemma-eta-third-property}
设 $A$ 为环，$f, g \in A$ 为非零因子。设 $M^\bullet$ 为
$A$-模复形，使 $fg$ 在所有 $M^i$ 上都是非零因子。则
$\eta_f\eta_gM^\bullet = \eta_{fg}M^\bullet$。
\end{lemma}

\begin{proof}
该陈述意味着，在次数 $i$，我们得到局部化
$M^i_{fg} = (M^i_g)_f$ 的同一子模。略去细节。
\end{proof}

\begin{lemma}
\label{more-algebra-lemma-eta-flat-base-change}
设 $A$ 为环，$f \in A$ 为非零因子。设 $A \to B$ 为平坦环映射，
并设
% P02-E0452
$g \in B$ 为 $f$ 的像。设 $M^\bullet$ 为无 $f$-挠 $A$-模的复形。
则 $g$ 是非零因子，$M^\bullet \otimes_A B$ 是无 $g$-挠模的复形，
并且
$\eta_fM^\bullet \otimes_A B = \eta_g(M^\bullet \otimes_A B)$。
\end{lemma}

\begin{proof}
略。
\end{proof}

\section{完美复形与 eta 算子}
\label{more-algebra-section-perfect-eta}

\noindent
本节将作若干代数准备，以用于我们版本的 Macpherson 图构造；见
《关于平坦性的进一步结果》第
\ref{flat-section-blowup-complexes-III} 节。我们将使用第
\ref{more-algebra-section-eta} 节引入的 $\eta_f$ 算子。

\medskip\noindent
设 $A$ 为环，$f \in A$ 为非零因子。设 $M^\bullet$ 为有限自由
$A$-模的有界复形。对每个 $i$，令 $r_i$ 为 $M^i$ 的秩，并置
$$
I_i(M^\bullet, f) = \text{ideal generated by the }
r_i \times r_i\text{-minors of }
(f, d^i) : M^i \to M^i \oplus M^{i + 1}
$$
。注意 $f^{r_i} \in I_i(M^\bullet, f)$。

\begin{lemma}
\label{more-algebra-lemma-ideal-well-defined}
设 $A$ 为环，$f \in A$ 为非零因子。设 $M^\bullet$ 与
$N^\bullet$ 为有限自由 $A$-模的两个有界复形，它们代表
$D(A)$ 中同一个对象。则
$$
f^m I_i(M^\bullet, f) = f^n I_i(N^\bullet, f)
$$
是 $A$ 的理想，其中整数 $n, m \geq 0$ 满足
$$
m + \sum\nolimits_{j \geq i} (-1)^{j - i}rk(M^j) =
n + \sum\nolimits_{j \geq i} (-1)^{j - i}rk(N^j)
$$
。
\end{lemma}

\begin{proof}
只需在 $A$ 的每个素理想处局部化后证明该等式。因此，由引理
\ref{more-algebra-lemma-compare-representatives-perfect} 以及我们略去的归纳论证，
% P02-E0453
可以假设 $N^\bullet = M^\bullet \oplus Q^\bullet$，其中
$Q^\bullet$ 为某个平凡复形，即
$$
Q^\bullet = \ldots \to 0 \to A \xrightarrow{1} A \to 0 \to \ldots
$$
，两个 $A$ 分别置于次数 $j$ 与 $j + 1$。若
$j \not = i - 1, i, i + 1$，则显然
$I_i(M^\bullet, f) = I_i(N^\bullet, f)$ 且 $m = n$，从而得到所需
等式。若 $j = i + 1$，则映射
$$
(f, d^i) : M^i \to M^i \oplus M^{i + 1}
\quad\text{且}\quad
(f, d^i, 0) : M^i \to M^i \oplus M^{i + 1} \oplus A
$$
具有相同的非零子式，故此时同样有
$I_i(M^\bullet, f) = I_i(N^\bullet, f)$ 且 $m = n$。若 $j = i$，
则 $I_i(M^\bullet, f)$ 是由映射
$$
(f, d^i) : M^i \to M^i \oplus M^{i + 1}
$$
的 $r_i \times r_i$ 子式生成的理想，而 $I_i(N^\bullet, f)$
是由映射
$$
(f \oplus f, d^i \oplus 1) :
(M^i \oplus A) \to (M^i \oplus A) \oplus (M^{i + 1} \oplus A)
$$
的 $(r_i + 1) \times (r_i + 1)$ 子式生成的理想。适当选取坐标后，
可见第二个映射的矩阵具有分块形式
% P02-E0454
$$
T =
\left(
\begin{matrix}
T_1 & 0 \\
0 & T_2
\end{matrix}
\right),
\quad
T_1 = \text{matrix of first map},
\quad
T_2 =
\left(
\begin{matrix}
f \\
1
\end{matrix}
\right)
$$
。沿用引理 \ref{more-algebra-lemma-ideals-generated-by-minors} 的记号，
$I_0(T_2) = A$，$I_1(T_2) = A$，且对 $p \geq 2$ 有
$I_p(T_2) = 0$，因而
$I_{r_i + 1}(T) = I_{r_i + 1}(T_1) + I_{r_i}(T_1) =
I_{r_i}(T_1)$；这意味着
$I_i(M^\bullet, f) = I_i(N^\bullet, f)$。又有 $m = n$，故
$j = i$ 的情形得证。

最后，设 $j = i - 1$。此时可见 $m = n + 1$，故须证明
$fI_i(M^\bullet, f) = I_i(N^\bullet, f)$。在此情形下，
$I_i(M^\bullet, f)$ 是由映射
$$
(f, d^i) : M^i \to M^i \oplus M^{i + 1}
$$
的 $r_i \times r_i$ 子式生成的理想，而 $I_i(N^\bullet, f)$
是由映射
$$
(f \oplus f, d^i) : (M^i \oplus A) \to (M^i \oplus A) \oplus M^{i + 1}
$$
的 $(r_i + 1) \times (r_i + 1)$ 子式生成的理想。适当选取坐标后，
可见第二个映射的矩阵具有分块形式
% P02-E0454
$$
T =
\left(
\begin{matrix}
T_1 & 0 \\
0 & T_2
\end{matrix}
\right),
\quad
T_1 = \text{matrix of first map},
\quad
T_2 =
\left(
\begin{matrix}
f
\end{matrix}
\right)
$$
。与上面同样论证，便得到
$fI_i(M^\bullet, f) = I_i(N^\bullet, f)$。
\end{proof}

\begin{lemma}
\label{more-algebra-lemma-eta-change-unit}
设 $f \in A$ 为环 $A$ 的非零因子，$u \in A$ 为单位。设
$M^\bullet$ 为有限自由 $A$-模的有界复形。则
$I_i(M^\bullet, f) = I_i(M^\bullet, uf)$。
\end{lemma}

\begin{proof}
略。
\end{proof}

\begin{lemma}
\label{more-algebra-lemma-eta-base-change-pre}
设 $A \to B$ 为环映射，$f \in A$ 为非零因子。设 $M^\bullet$
为有限自由 $A$-模的有界复形。假设 $f$ 在 $B$ 中的像 $g$
是非零因子。则
$I_i(M^\bullet, f)B = I_i(M^\bullet \otimes_A B, g)$。
\end{lemma}

\begin{proof}
映射 $(f, d^i) : M^i \to M^i \oplus M^{i + 1}$ 的各子式映到
相应映射
$(g, d^i) : M^i \otimes_A B \to M^i \otimes_A B \oplus M^{i + 1} \otimes_A B$
的对应子式。
\end{proof}

\begin{lemma}
\label{more-algebra-lemma-eta-cohomology-free}
设 $A$ 为环，$\mathfrak p \subset A$ 为素理想，且 $f \in A$
为非零因子。设 $M^\bullet$ 为有限自由 $A$-模的有界复形。若对所有
$i$，$H^i(M^\bullet)_\mathfrak p$ 都自由，则对所有 $i$，
$I_i(M^\bullet, f)_\mathfrak p$ 都是主理想；事实上，它由 $f$
的一个幂生成。
\end{lemma}

\begin{proof}
由引理 \ref{more-algebra-lemma-eta-base-change-pre}，可以假设 $A$ 为极大理想
$\mathfrak p$ 下的局部环。由引理 \ref{more-algebra-lemma-ideal-well-defined}，
还可用一个拟同构复形替换 $M^\bullet$。根据上同调模自由这一假设，
可见 $M^\bullet$ 拟同构于这样的复形：其次数 $i$ 项是
$H^i(M^\bullet)$，且微分全为零；例如见《导出范畴》引理
\ref{derived-lemma-ext-2-zero-pre}。换言之，可以假设复形
$M^\bullet$ 的所有微分都为零。此时显然
$I_i(M^\bullet, f) = (f^{r_i})$ 是主理想。
\end{proof}

\begin{lemma}
\label{more-algebra-lemma-eta-locally-free}
设 $A$ 为环，$f \in A$ 为非零因子。设 $M^\bullet$ 为有限自由
$A$-模的有界复形。假设 $I_i(M^\bullet, f)$ 是主理想。则
$(\eta_fM)^i$ 是秩为 $r_i$ 的局部自由模，且映射
$(1, d^i) : (\eta_fM)^i \to f^iM^i \oplus f^{i + 1}M^{i + 1}$
是一个直和项的嵌入。
\end{lemma}

\begin{proof}
选取 $I_i(M^\bullet, f)$ 的一个生成元 $g$。由于
$f^{r_i} \in I_i(M^\bullet, f)$，可见 $g$ 整除 $f$ 的某个幂。
特别地，$g$ 是 $A$ 中的非零因子。映射
$(f, d^i) : M^i \to M^i \oplus M^{i + 1}$ 的
$r_i \times r_i$ 子式生成理想 $I_i(M^\bullet, f)$，而
$(f, d^i)$ 的 $(r_i + 1) \times (r_i + 1)$ 子式为零：这可在
$f$ 处局部化后检验，因为局部化后该映射的秩等于 $r_i$。考虑满射
$$
M^i \oplus M^{i + 1}
\longrightarrow
Q = \Coker(f, d^i)/g\text{-torsion}
$$
。由引理 \ref{more-algebra-lemma-fitting-ideals-and-pd1}，模 $Q$ 是秩为
$r_{i + 1}$ 的有限局部自由模。因此 $Q$ 无 $f$-挠，从而
$(f, d^i)$ 的余核模去 $f$-幂挠后同样为 $Q$。

\medskip\noindent
考虑有限自由 $A$-模复形
$$
0 \to f^{i + 1}M^i \xrightarrow{1, d^i} f^iM^i \oplus f^{i + 1}M^{i + 1}
\xrightarrow{d^i, -1} f^iM^{i + 1} \to 0
$$
；它在 $f$ 处局部化后成为分裂正合复形。映射
$(1, d^i) : f^{i + 1}M^i \to f^iM^i \oplus f^{i + 1}M^{i + 1}$
同构于上面研究过的映射
$(f, d^i) : M^i \to M^i \oplus M^{i + 1}$。因此像
$$
Q' = \Im(f^iM^i \oplus f^{i + 1}M^{i + 1} \xrightarrow{d^i, -1} f^iM^{i + 1})
$$
同构于 $Q$，特别地是投射模。
% P02-E0455
另一方面，由第 \ref{more-algebra-section-eta} 节中 $\eta_f$ 的构造，
单射
$(1, d^i) : (\eta_fM)^i \to f^iM^i \oplus f^{i + 1}M^{i + 1}$
的像是 $(d^i, -1)$ 的核。于是得到同构
$(\eta_fM)^i \oplus Q' = f^iM^i \oplus f^{i + 1}M^{i + 1}$，
并且可见
% P02-E0456
$\eta_fM^i$ 确为秩 $r_i$ 的有限局部自由模，而 $(1, d^i)$
是一个直和项的嵌入。
\end{proof}

\begin{lemma}
\label{more-algebra-lemma-eta-base-change}
设 $A \to B$ 为环映射，$f \in A$ 为非零因子。设 $M^\bullet$
为有限自由 $A$-模的有界复形。假设 $f$ 在 $B$ 中的像 $g$
是非零因子，并且对所有 $i \in \mathbf{Z}$，
$I_i(M^\bullet, f)$ 都是主理想。则有典范同构
$\eta_fM^\bullet \otimes_A B = \eta_g(M^\bullet \otimes_A B)$。
\end{lemma}

\begin{proof}
置 $N^i = M^i \otimes_A B$。注意，作为 $(N^i)_g$ 的子模，
$f^iM^i \otimes_A B = g^iN^i$。映射
% P02-E0457
$$
(\eta_fM)^i \otimes_A B \to g^iN^i \otimes g^{i + 1}N^{i + 1}
\quad\text{且}\quad
(\eta_gN)^i \to g^iN^i \otimes g^{i + 1}N^{i + 1}
$$
由引理 \ref{more-algebra-lemma-eta-locally-free} 都是直和项的嵌入。它们的像在
$g$ 处局部化后相同，结论随即得出。
\end{proof}

\begin{lemma}
\label{more-algebra-lemma-ideal-direct-summand}
设 $A$ 为环，$M$、$N_1$、$N_2$ 为有限投射 $A$-模。设
$s : M \to N_1 \oplus N_2$ 为分裂单射。存在具有以下性质的有限生成
理想 $J \subset A$：环映射 $A \to B$ 经过 $A/J$ 分解，当且仅当
$s \otimes \text{id}_B$ 把 $M \otimes_A B$ 同一为
$N_1 \otimes_A B \oplus N_2 \otimes_A B$ 的若干子模的直和。
\end{lemma}

\begin{proof}
选取 $s$ 的一个分裂
$\pi : N_1 \oplus N_2 \to M$。
% P02-E0458
以 $q_i : N_1 \oplus N_2 \to N_1 \oplus N_2$ 表示到 $N_i$
上的投影算子。置 $p_i = \pi \circ q_i \circ s$。注意
$p_1 + p_2 = \text{id}_M$。我们断言，$M$ 是
$N_1 \oplus N_2$ 的若干子模的直和，当且仅当 $p_1$ 与 $p_2$
是正交投影算子。因此，$J$ 是 $A$ 的满足下列条件的最小理想：
$p_1 \circ p_1 - p_1$、$p_2 \circ p_2 - p_2$、
$p_1 \circ p_2$ 与
$p_2 \circ p_1$ 都包含在
$J \otimes_A \text{End}_A(M)$ 中。略去一些细节。
\end{proof}

\noindent
设 $A$ 为环，$f \in A$ 为非零因子。设 $M^\bullet$ 为有限自由
$A$-模的有界复形。假设对所有 $i \in \mathbf{Z}$，理想
$I_i(M^\bullet, f)$ 都是主理想。于是，由引理
\ref{more-algebra-lemma-eta-locally-free}，映射
$$
(1, d^i) : (\eta_fM)^i / f(\eta_fM)^i \longrightarrow
f^iM^i/f^{i + 1}M^i \oplus f^{i + 1}M^{i + 1}/f^{i + 2}M^{i + 1}
$$
是分裂单射。
% P02-E0459
以 $J_i(M^\bullet, f) \subset A/fA$ 表示引理
\ref{more-algebra-lemma-ideal-direct-summand} 中对应于上面所示分裂单射
$(1, d^i)$ 的有限生成理想。

\begin{lemma}
\label{more-algebra-lemma-eta-vanishing-beta-plus-pre}
设 $A$ 为环，$f \in A$ 为非零因子。设 $M^\bullet$ 与
$N^\bullet$ 为有限自由 $A$-模的两个有界复形，它们代表
$D(A)$ 中同一个对象。假设对所有 $i \in \mathbf{Z}$，
$I_i(M^\bullet, f)$ 都是主理想。则作为 $A/fA$ 的理想，
$J_i(M^\bullet, f) = J_i(N^\bullet, f)$。
\end{lemma}

\begin{proof}
注意，
% P02-E0460
$I_i(M^\bullet, f)$ 是主理想这一事实，由引理
\ref{more-algebra-lemma-ideal-well-defined} 蕴含 $I_i(M^\bullet, f)$ 是主理想，
故该陈述有意义。如同引理 \ref{more-algebra-lemma-ideal-well-defined} 的证明，
可以假设 $N^\bullet = M^\bullet \oplus Q^\bullet$，其中
$Q^\bullet$ 为某个平凡复形，即
$$
Q^\bullet = \ldots \to 0 \to A \xrightarrow{1} A \to 0 \to \ldots
$$
，两个 $A$ 分别置于次数 $j$ 与 $j + 1$。由于 $\eta_f$ 与直和相容，
可见映射
$$
(1, d^i) : (\eta_fN)^i / f(\eta_fN)^i \longrightarrow
f^iN^i/f^{i + 1}N^i \oplus f^{i + 1}N^{i + 1}/f^{i + 2}N^{i + 1}
$$
是 $M^\bullet$ 的相应映射与 $Q^\bullet$ 的相应映射的直和。由定义
所论理想的泛性质，得到
$J_i(N^\bullet, f) = J_i(M^\bullet, f) + J_i(Q^\bullet, f)$。
因此，只需证明对所有 $i$ 都有 $J_i(Q^\bullet, f) = 0$。这是一个
我们略去的计算。
\end{proof}

\begin{lemma}
\label{more-algebra-lemma-eta-vanishing-beta-plus}
设 $A$ 为环，$f \in A$ 为非零因子。设 $M^\bullet$ 为有限自由
$A$-模的有界复形。假设对所有 $i \in \mathbf{Z}$，
$I_i(M^\bullet, f)$ 都是主理想。考虑 $A/fA$ 的理想
$J(M^\bullet, f) = \sum_i J_i(M^\bullet, f)$。考虑素理想集合
\begin{align*}
E
& =
\{f \in \mathfrak p \subset A \mid
\Ker(d^i \bmod f^2)_\mathfrak p \text{ surjects onto }
\Ker(d^i \bmod f)_\mathfrak p \text{ 对所有 }i \in \mathbf{Z}\} \\
& =
\{f \in \mathfrak p \subset A \mid
\text{the localizations }\beta_\mathfrak p
\text{ of the Bockstein operators are zero}\}
\end{align*}
。则有：
\begin{enumerate}
\item $J(M^\bullet, f)$ 是有限生成的；
\item $A/fA \to C = (A/fA)/J(M^\bullet, f)$ 是有限表示满射；
\item 对 $\mathfrak p \in E$，有
$J(M^\bullet, f)_\mathfrak p = 0$；
\item 若 $f \in \mathfrak p$，且对所有 $i \in \mathbf{Z}$，
$H^i(M^\bullet)_\mathfrak p$ 都自由，则 $\mathfrak p \in E$；以及
\item $\eta_f M^\bullet \otimes_A C$ 的各上同调模都是有限局部自由
$C$-模。
\end{enumerate}
\end{lemma}

\begin{proof}
$E$ 定义中两个集合的相等性由引理 \ref{more-algebra-lemma-vanishing-beta} 得出；
此外，该引理的最后一个陈述蕴含 (4)。

\medskip\noindent
(1) 成立，因为理想 $J_i(M^\bullet, f)$ 都有限生成，而且
$M^\bullet$ 有界，故几乎所有 $i$ 对应的 $J_i(M^\bullet, f)$ 都为零。
(2) 只是 (1) 的改述。

\medskip\noindent
证明 (3)。由引理 \ref{more-algebra-lemma-eta-locally-free}，对所有 $i$，
$(\eta_fM)^i$ 都是秩为 $r_i$ 的有限局部自由模。考虑映射
$$
(1, d^i) :
(\eta_fM)^i / f(\eta_fM)^i
\longrightarrow
f^iM^i/f^{i + 1}M^i \oplus f^{i + 1}M^{i + 1}/f^{i + 2}M^{i + 1}
$$
。取 $\mathfrak p \in E$。由引理
\ref{more-algebra-lemma-eta-vanishing-beta} 以及各模 $(\eta_fM)^i$ 的局部自由性，
可以写成
$$
\left((\eta_fM)^i / f(\eta_fM)^i\right)_\mathfrak p =
(A/fA)_\mathfrak p^{\oplus m_i} \oplus (A/fA)_\mathfrak p^{\oplus n_i}
$$
，并与上面的箭头 $(1, d^i)$ 相容。由理想
$J_i(M^\bullet, f)$ 的泛性质，得到
$J_i(M^\bullet, f)_\mathfrak p = 0$。因此，
% P02-E0461
$I_\mathfrak p = fA_\mathfrak p$ 对 $\mathfrak p \in E$ 成立。

\medskip\noindent


\medskip\noindent
证明 (5)。注意，$\eta_fM^\bullet$ 上的微分嵌入交换图
% P02-E0462
$$
\xymatrix{
(\eta_fM)^i \ar[d] \ar[r] &
f^iM^i \oplus f^{i + 1}M^{i + 1}
\ar[d]^{\left(
\begin{matrix}
0 & 1 \\
0 & 0
\end{matrix}
\right)} \\
(\eta_fM)^{i + 1} \ar[r] &
f^{i + 1}M^i \oplus f^{i + 2}M^{i + 2}
}
$$
。由构造，在与 $C$ 张量后，左边各模是右边各直和项的若干直和项
之直和。考察图
% P02-E0462
$$
\xymatrix{
(\eta_fM)^i \otimes_A C \ar[d] \ar@{=}[r] &
K^i \oplus L^i \ar[r] \ar[d] &
f^iM^i \otimes_A C \oplus f^{i + 1}M^{i + 1} \otimes_A C
\ar[d]^{\left(
\begin{matrix}
0 & 1 \\
0 & 0
\end{matrix}
\right)} \\
(\eta_fM)^{i + 1} \otimes_A C \ar@{=}[r] &
K^{i + 1} \oplus L^{i + 1} \ar[r] &
f^{i + 1}M^i \otimes_A C \oplus f^{i + 2}M^{i + 2} \otimes_A C
}
$$
，其中各水平箭头既与直和分解相容，又是直和项的嵌入。由此，微分把
$L^i$ 同一为 $K^{i + 1}$ 的一个直和项，因而
$\eta_fM^\bullet \otimes_A C$ 在次数 $i$ 的上同调为模
% P02-E0463
$K^{i + 1}/L^i$；它按要求是有限投射模。
\end{proof}

\section{复形的极限}
\label{more-algebra-section-limits}

\noindent
本节讨论一个复形具有“形式变形”并对其取极限时会发生什么。我们将
考虑两种情形：
\begin{enumerate}
\item 有环的一个逆系统的极限 $A = \lim A_n$，其过渡映射为满射，
核局部幂零；并有对象 $K_n \in D(A_n)$，它们按
$K_n = K_{n + 1} \otimes_{A_{n + 1}}^\mathbf{L} A_n$ 的意义相容；或者
\item 有环 $A$、理想 $I$ 以及对象 $K_n \in D(A/I^n)$，它们按
$K_n = K_{n + 1} \otimes_{A/I^{n + 1}}^\mathbf{L} A/I^n$ 的意义相容。
\end{enumerate}
在附加假设下，可以证明 $K = R\lim K_n$ 重现该系统，即
$K_n = K \otimes_A^\mathbf{L} A_n$ 或
$K_n = K \otimes_A^\mathbf{L} A/I^n$。

\begin{lemma}
\label{more-algebra-lemma-Rlim-pseudo-coherent-gives-pseudo-coherent}
设 $A = \lim A_n$ 为环的一个逆系统 $(A_n)$ 的极限。给定
$K_n \in D(A_n)$ 以及 $D(A_{n + 1})$ 中的映射
$K_{n + 1} \to K_n$。假设：
\begin{enumerate}
\item 过渡映射 $A_{n + 1} \to A_n$ 是满射且核局部幂零；
\item 或者所有 $K_n$ 都伪凝聚，或者存在整数 $n_0$，使得
$K_{n_0}$ 伪凝聚，且当 $n \geq n_0$ 时，
$A_{n + 1} \to A_n$ 的核是幂零理想；
\item 这些映射诱导同构
$K_{n + 1} \otimes_{A_{n + 1}}^\mathbf{L} A_n \to K_n$。
\end{enumerate}
则 $K = R\lim K_n$ 是 $D(A)$ 的伪凝聚对象，并且对所有 $n$，
$K \otimes_A^\mathbf{L} A_n \to K_n$ 都是同构。
\end{lemma}

\begin{proof}
由假设，可找到一个代表 $K_1$ 的有限自由 $A_1$-模上有界复形
$P_1^\bullet$；见定义 \ref{more-algebra-definition-pseudo-coherent}。由引理
\ref{more-algebra-lemma-check-pseudo-coherent-modulo-nilpotent}，得到对所有 $n$，
$K_n$ 都伪凝聚。于是，由引理 \ref{more-algebra-lemma-lift-complex-stably-frees}，
对 $n > 1$ 归纳可找到代表 $K_n$ 的有限自由 $A_n$-模复形
$P_n^\bullet$，以及代表映射 $K_n \to K_{n - 1}$ 的映射
$P_n^\bullet \to P_{n - 1}^\bullet$；它们诱导复形同构（！）
$P_n^\bullet \otimes_{A_n} A_{n - 1} \to P_{n - 1}^\bullet$。
因此，由引理 \ref{more-algebra-lemma-compute-Rlim-modules} 与评注
\ref{more-algebra-remark-how-unique}，
$K = R\lim K_n$ 由 $P^\bullet = \lim P_n^\bullet$ 代表。
对每个 $n$，$P_n^i$ 是有限自由 $A_n$-模，又因
$A = \lim A_n$，可见对每个 $i$，$P^i$ 都是与 $P_1^i$ 同秩的
有限自由模。这意味着 $K$ 伪凝聚。又有
$K \otimes_A^\mathbf{L} A_n$ 由
$P^\bullet \otimes_A A_n = P_n^\bullet$ 代表，从而证明最后一个断言。
\end{proof}

\begin{lemma}
\label{more-algebra-lemma-Rlim-pseudo-coherent-gives-complete-pseudo-coherent}
设 $A$ 为环，$I \subset A$ 为理想。给定 $K_n \in D(A/I^n)$
以及 $D(A/I^{n + 1})$ 中的映射 $K_{n + 1} \to K_n$。假设：
\begin{enumerate}
\item $A$ 是 $I$-进完备的；
\item $K_1$ 伪凝聚；且
\item 这些映射诱导同构
$K_{n + 1} \otimes_{A/I^{n + 1}}^\mathbf{L} A/I^n \to K_n$。
\end{enumerate}
则 $K = R\lim K_n$ 是 $D(A)$ 的伪凝聚导出完备对象，并且对所有
$n$，$K \otimes_A^\mathbf{L} A/I^n \to K_n$ 都是同构。
\end{lemma}

\begin{proof}
由引理 \ref{more-algebra-lemma-Rlim-pseudo-coherent-gives-pseudo-coherent}，已经知道
$K$ 伪凝聚，并且对所有 $n$，
$K \otimes_A^\mathbf{L} A/I^n \to K_n$ 都是同构。最后，由引理
\ref{more-algebra-lemma-naive-derived-completion}，$K$ 导出完备。
\end{proof}

\begin{lemma}
\label{more-algebra-lemma-Rlim-perfect-gives-perfect}
\begin{reference}
\cite[Lemma 4.2]{Bhatt-Algebraize}
\end{reference}
设 $A = \lim A_n$ 为环的一个逆系统 $(A_n)$ 的极限。给定
$K_n \in D(A_n)$ 以及 $D(A_{n + 1})$ 中的映射
$K_{n + 1} \to K_n$。假设：
\begin{enumerate}
\item 过渡映射 $A_{n + 1} \to A_n$ 是满射且核局部幂零；
\item 或者所有 $K_n$ 都完美，或者存在整数 $n_0$，使得
$K_{n_0}$ 完美，且当 $n \geq n_0$ 时，
$A_{n + 1} \to A_n$ 的核为幂零理想；且
\item 这些映射诱导同构
$K_{n + 1} \otimes_{A_{n + 1}}^\mathbf{L} A_n \to K_n$。
\end{enumerate}
则 $K = R\lim K_n$ 是 $D(A)$ 的完美对象，并且对所有 $n$，
$K \otimes_A^\mathbf{L} A_n \to K_n$ 都是同构。
\end{lemma}

\begin{proof}
由引理 \ref{more-algebra-lemma-Rlim-pseudo-coherent-gives-pseudo-coherent}，已经知道
$K$ 伪凝聚，并且对所有 $n$，
$K \otimes_A^\mathbf{L} A_n \to K_n$ 都是同构。考虑一个满射
$A \to \kappa$，其核为极大理想 $\mathfrak m$。由《代数》引理
\ref{algebra-lemma-locally-nilpotent-unit}，$A$ 中任何在 $A_1$ 中
映为单位的元素在 $A$ 中都是单位；因而，由《代数》引理
\ref{algebra-lemma-contained-in-radical}，
$\Ker(A \to A_1)$ 包含在 $A$ 的 Jacobson 根中。故
$A \to \kappa$ 分解为 $A \to A_1 \to \kappa$。于是
$$
K \otimes_A^\mathbf{L} \kappa =
K \otimes_A^\mathbf{L} A_1 \otimes_{A_1}^\mathbf{L} \kappa =
K_1 \otimes_{A_1}^\mathbf{L} \kappa
$$
。
% P02-E0464
注意，由假设 (2)，$K_1$ 完美。因此，由引理
\ref{more-algebra-lemma-perfect}，$K_1$ 具有有限 Tor 维数。故存在
$a, b \in \mathbf{Z}$，使得对所有 $i \not \in [a, b]$，
$H^i(K \otimes_A^\mathbf{L} \kappa) = 0$。由引理
\ref{more-algebra-lemma-check-perfect-pointwise}，得到 $K$ 完美。
\end{proof}

\begin{lemma}
\label{more-algebra-lemma-Rlim-perfect-gives-complete}
设 $A$ 为环，$I \subset A$ 为理想。给定 $K_n \in D(A/I^n)$
以及 $D(A/I^{n + 1})$ 中的映射 $K_{n + 1} \to K_n$。假设：
\begin{enumerate}
\item $A$ 是 $I$-进完备的；
\item $K_1$ 是完美对象；且
\item 这些映射诱导同构
$K_{n + 1} \otimes_{A/I^{n + 1}}^\mathbf{L} A/I^n \to K_n$。
\end{enumerate}
则 $K = R\lim K_n$ 是 $D(A)$ 的完美导出完备对象，并且对所有
$n$，$K \otimes_A^\mathbf{L} A/I^n \to K_n$ 都是同构。
\end{lemma}

\begin{proof}
合并引理 \ref{more-algebra-lemma-Rlim-perfect-gives-perfect} 与
\ref{more-algebra-lemma-Rlim-pseudo-coherent-gives-complete-pseudo-coherent}
（后者给出导出完备性）。
\end{proof}

\noindent
我们不知道下一个引理对无界复形是否成立。

\begin{lemma}
\label{more-algebra-lemma-Rlim-gives-complete}
设 $A$ 为环，$I \subset A$ 为理想。给定 $K_n \in D(A/I^n)$
以及 $D(A/I^{n + 1})$ 中的映射 $K_{n + 1} \to K_n$。若
\begin{enumerate}
\item $A$ 是 Noether 环；
\item $K_1$ 上有界；且
\item 这些映射诱导同构
$K_{n + 1} \otimes_{A/I^{n + 1}}^\mathbf{L} A/I^n \to K_n$，
\end{enumerate}
则 $K = R\lim K_n$ 是 $D^-(A)$ 的导出完备对象，并且对所有 $n$，
$K \otimes_A^\mathbf{L} A/I^n \to K_n$ 都是同构。
\end{lemma}

\begin{proof}
由引理 \ref{more-algebra-lemma-naive-derived-completion}，$D(A)$ 的对象 $K$
导出完备。

\medskip\noindent
假设当 $i > b$ 时 $H^i(K_1) = 0$。则可找到代表 $K_1$ 的自由
$A/I$-模复形 $P_1^\bullet$，且当 $i > b$ 时 $P_1^i = 0$。由引理
\ref{more-algebra-lemma-lift-complex-projectives}，对 $n > 1$ 归纳可找到代表
$K_n$ 的自由 $A/I^n$-模复形 $P_n^\bullet$，以及代表映射
$K_n \to K_{n - 1}$ 的映射
$P_n^\bullet \to P_{n - 1}^\bullet$；它们诱导复形同构（！）
$P_n^\bullet/I^{n - 1}P_n^\bullet \to P_{n - 1}^\bullet$。

\medskip\noindent
于是得到如下情形：$R\lim K_n$ 由
$P^\bullet = \lim P_n^\bullet$ 代表；见引理
\ref{more-algebra-lemma-compute-Rlim-modules} 与评注 \ref{more-algebra-remark-how-unique}。
各复形 $P_n^\bullet$ 是平坦 $A/I^n$-模的一致上有界复形，且过渡映射
逐项满射。于是，由引理 \ref{more-algebra-lemma-limit-flat}，$P^\bullet$
是平坦 $A$-模的上有界复形。因此
$K \otimes_A^\mathbf{L} A/I^t$ 由
$P^\bullet \otimes_A A/I^t$ 代表。由引理 \ref{more-algebra-lemma-limit-flat}，
逐项有
$P^\bullet \otimes_A A/I^t = \lim P_n^\bullet \otimes_A A/I^t$。
根据对 $P_n^\bullet$ 的选取，当 $n \geq t$ 时，过渡映射
$P_{n + 1}^\bullet \otimes_A A/I^t \to P_n^\bullet \otimes_A A/I^t$
都是同构，因而
$\lim P_n^\bullet \otimes_A A/I^t =
P_t^\bullet \otimes_A A/I^t = P_t^\bullet$。由于 $P_t^\bullet$
代表 $K_t$，可见
$K \otimes_A^\mathbf{L} A/I^t \to K_t$ 是同构。
\end{proof}

\noindent
下面给出另一类结果。

\begin{lemma}[Koll\'ar-Kov\'acs]
\label{more-algebra-lemma-kollar-kovacs}
\begin{reference}
% P02-E0465
Kovacs 于 23/02/2018 发来的电子邮件。
\end{reference}
设 $I$ 为 Noether 环 $A$ 的理想，$K \in D(A)$。置
$K_n = K \otimes_A^\mathbf{L} A/I^n$。假设对所有
$i \in \mathbf{Z}$ 都有：
\begin{enumerate}
\item $H^i(K)$ 是有限 $A$-模；且
\item 系统 $H^i(K_n)$ 满足 Mittag-Leffler 性质。
\end{enumerate}
则对所有 $i \in \mathbf{Z}$，
$\lim H^i(K)/I^nH^i(K)$ 等于 $\lim H^i(K_n)$。
\end{lemma}

\begin{proof}
回忆，$K^\wedge = R\lim K_n$ 是 $K$ 的导出完备化；见命题
\ref{more-algebra-proposition-noetherian-naive-completion-is-completion}。由引理
\ref{more-algebra-lemma-derived-completion-pseudo-coherent}，有
$H^i(K^\wedge) = \lim H^i(K)/I^nH^i(K)$。由引理
\ref{more-algebra-lemma-break-long-exact-sequence-modules}，得到短正合列
% P02-E0464
$$
0 \to R^1\lim H^{i - 1}(K_n) \to H^i(K^\wedge) \to \lim H^i(K_n) \to 0
$$
。Mittag-Leffler 性质保证左端项为零（引理
\ref{more-algebra-lemma-compute-Rlim-modules}），故引理得证。
\end{proof}

\section{若干求值映射}
\label{more-algebra-section-evaluation}

\noindent
本节证明，对适当类型的复形，某些典范的 $R\Hom$ 映射是同构。

\begin{lemma}
\label{more-algebra-lemma-internal-hom-evaluate-isomorphism}
设 $R$ 为环，$K, L, M$ 为 $D(R)$ 的对象。
% P02-E0466
引理 \ref{more-algebra-lemma-internal-hom-evaluate} 的映射
$$
R\Hom_R(L, M) \otimes_R^\mathbf{L} K \longrightarrow R\Hom_R(R\Hom_R(K, L), M)
$$
在以下两种情形中是同构：
\begin{enumerate}
\item $K$ 完美；或者
\item $K$ 伪凝聚，$L \in D^+(R)$，并且
% P02-E0467
$M$ 具有有限内射维数。
\end{enumerate}
\end{lemma}

\begin{proof}
选取代表 $M$ 的 K-内射复形 $I^\bullet$、代表 $L$ 的 K-内射复形
$J^\bullet$，以及代表 $K$ 的有限投射模上有界复形 $K^\bullet$。
考虑引理 \ref{more-algebra-lemma-evaluate-and-more} 的复形映射
$$
\text{Tot}(\Hom^\bullet(J^\bullet, I^\bullet) \otimes_R K^\bullet)
\longrightarrow
\Hom^\bullet(\Hom^\bullet(K^\bullet, J^\bullet), I^\bullet)
$$
。注意
$$
\left(\prod\nolimits_{p + r = t} \Hom_R(J^{-r}, I^p)\right) \otimes_R K^s =
\prod\nolimits_{p + r = t} \Hom_R(J^{-r}, I^p) \otimes_R K^s
$$
，因为 $K^s$ 有限投射。该映射由映射
$$
c_{p, r, s} :
\Hom_R(J^{-r}, I^p) \otimes_R K^s
\longrightarrow
\Hom_R(\Hom_R(K^s, J^{-r}), I^p)
$$
给出；由于 $K^s$ 有限投射，它们都是同构。对于左端次数 $n$
的任意元素 $\alpha = (\alpha^{p, r, s})$，只有有限多个 $s$ 的值
使 $\alpha^{p, r, s}$ 非零（这里对某些满足
$n = p + r + s$ 的 $p, r$ 而言）。因此，若右端元素
$\beta = (\beta^{p, r, s})$ 也被迫满足同一消失条件，则该映射是同构。
若 $K^\bullet$ 是有限投射模的有界复形，这一点显然。另一方面，若可取
$I^\bullet$ 有界而 $J^\bullet$ 下有界，则当 $p$ 落在某个固定范围之外、
$s \gg 0$ 或 $r \gg 0$ 时，$\beta^{p, r, s}$ 为零。因此，在满足
$n = p + r + s$ 且 $\beta^{p, r, s}$ 非零的解中，只出现有限多个
$s$ 的值。
\end{proof}

\begin{lemma}
\label{more-algebra-lemma-internal-hom-evaluate-isomorphism-technical}
设 $R$ 为环，$K, L, M$ 为 $D(R)$ 的对象。
% P02-E0466
若以下三个条件成立，则引理 \ref{more-algebra-lemma-internal-hom-evaluate} 的映射
$$
R\Hom_R(L, M) \otimes_R^\mathbf{L} K \longrightarrow R\Hom_R(R\Hom_R(K, L), M)
$$
是同构：
\begin{enumerate}
\item $L, M$ 具有有限内射维数；
\item $R\Hom_R(L, M)$ 具有有限 Tor 维数；
% P02-E0468
\item 对每个 $n \in \mathbf{Z}$，截断 $\tau_{\leq n}K$ 都伪凝聚。
\end{enumerate}
\end{lemma}

\begin{proof}
选取整数 $n$，并考虑区分三角
% P02-E0469
$$
\tau_{\leq n}K \to K \to \tau_{\geq n + 1}K \to \tau_{\leq n}K[1]
$$
；见《导出范畴》评注
\ref{derived-remark-truncation-distinguished-triangle}。由假设 (3)
与引理 \ref{more-algebra-lemma-internal-hom-evaluate-isomorphism}，该映射对
$\tau_{\leq n}K$ 是同构。因此，只需证明
$$
R\Hom_R(L, M) \otimes_R^\mathbf{L} \tau_{\geq n + 1}K
\quad\text{且}\quad
R\Hom_R(R\Hom_R(\tau_{\geq n + 1}K, L), M)
$$
的上同调在次数 $\leq n - c$ 中消失，其中 $c$ 为某个整数。这立即由
假设 (2) 与 (1) 得出。
\end{proof}

\begin{lemma}
\label{more-algebra-lemma-internal-hom-evaluate-tensor-isomorphism}
设 $R$ 为环，$K, L, M$ 为 $D(R)$ 的对象。引理
\ref{more-algebra-lemma-internal-hom-diagonal-better} 的映射
$$
K \otimes_R^\mathbf{L} R\Hom_R(M, L) \longrightarrow
R\Hom_R(M, K \otimes_R^\mathbf{L} L)
$$
在以下情形中是同构：
\begin{enumerate}
\item $M$ 完美；或者
\item $K$ 完美；或者
\item $M$ 伪凝聚，$L \in D^+(R)$，并且 $K$ 的 Tor 振幅位于
% P02-E0470
$[a, \infty]$ 中。
\end{enumerate}
\end{lemma}

\begin{proof}
先证明 $M$ 完美的情形。注意，该箭头的两端都把 $M$ 中的区分三角
变为区分三角，并与直和交换。因此，只需检验 $M = R[n]$ 时该断言
成立；见《导出范畴》评注
\ref{derived-remark-check-on-generator} 与引理
\ref{more-algebra-lemma-perfect-ring-classical-generator}。在此情形下，结果显然。

\medskip\noindent
再证明 $K$ 完美的情形。论证与前一种情形相同。

\medskip\noindent
最后证明 (3)。可以用 $R$-模下有界复形 $K^\bullet$ 与
$L^\bullet$ 分别代表 $K$ 与 $L$。可以假设 $K^\bullet$ 是由平坦
$R$-模组成的 K-平坦复形；见引理
\ref{more-algebra-lemma-bounded-below-tor-amplitude}。可以用有限自由 $R$-模的
上有界复形 $M^\bullet$ 代表 $M$；见定义
\ref{more-algebra-definition-pseudo-coherent}。于是左端对象由
% P02-E0469
$$
\text{Tot}(K^\bullet \otimes_R \Hom^\bullet(M^\bullet, L^\bullet))
$$
代表，而右端对象由
$$
\Hom^\bullet(M^\bullet, \text{Tot}(K^\bullet \otimes_R L^\bullet))
$$
代表。这使用引理 \ref{more-algebra-lemma-RHom-out-of-projective}。两个复形在次数
$n$ 的模都是
$$
\bigoplus\nolimits_{p + q + r = n} K^p \otimes \Hom_R(M^{-r}, L^q) =
\bigoplus\nolimits_{p + q + r = n} \Hom_R(M^{-r}, K^p \otimes_R L^q)
$$
，因为 $M^{-r}$ 有限自由
% P02-E0471
（这些也都是有限直和）。引理
\ref{more-algebra-lemma-internal-hom-diagonal-better} 中定义的映射来自引理
\ref{more-algebra-lemma-diagonal-better} 中定义的复形映射；后者使用这些模之间的
典范同构。
\end{proof}

\begin{lemma}
\label{more-algebra-lemma-hom-complex-K-flat}
设 $R$ 为环，$P^\bullet$ 为投射 $R$-模的上有界复形，
$K^\bullet$ 为 $R$-模的 K-平坦复形。若 $P^\bullet$ 是 $D(R)$
的完美对象，则 $\Hom^\bullet(P^\bullet, K^\bullet)$ 是 K-平坦的，
并代表 $R\Hom_R(P^\bullet, K^\bullet)$。
\end{lemma}

\begin{proof}
最后一个陈述即引理 \ref{more-algebra-lemma-RHom-out-of-projective}。由于
$P^\bullet$ 代表完美对象，存在有限投射 $R$-模的有限复形
$F^\bullet$，使得 $P^\bullet$ 与 $F^\bullet$ 在 $D(R)$ 中同构；
见定义 \ref{more-algebra-definition-perfect}。于是，由《导出范畴》引理
\ref{derived-lemma-morphisms-from-projective-complex}，
$P^\bullet$ 与 $F^\bullet$ 伦等价。因此
$\Hom^\bullet(P^\bullet, K^\bullet)$ 与
$\Hom^\bullet(F^\bullet, K^\bullet)$ 伦等价。故前者 K-平坦当且仅当
后者 K-平坦（这由定义 \ref{more-algebra-definition-K-flat} 与引理
\ref{more-algebra-lemma-derived-tor-homotopy} 得出）。显然，
% P02-E0469
$$
\Hom^\bullet(F^\bullet, K^\bullet) =
\text{Tot}(E^\bullet \otimes_R K^\bullet)
$$
，其中 $E^\bullet$ 是 $F^\bullet$ 的对偶复形，其项为
$E^n = \Hom_R(F^{-n}, R)$；见引理
\ref{more-algebra-lemma-dual-perfect-complex} 及其证明。由于 $E^\bullet$
是投射模的有界复形，由引理
\ref{more-algebra-lemma-derived-tor-quasi-isomorphism}，它是 K-平坦的。最后由引理
\ref{more-algebra-lemma-tensor-product-K-flat} 得出结论。
\end{proof}

\section{导出 Hom 的基变换}
\label{more-algebra-section-base-change-RHom}

\noindent
关于这一点，我们已经在引理
\ref{more-algebra-lemma-pseudo-coherence-and-base-change-ext} 与《代数》第
\ref{algebra-section-functoriality-ext} 节中看到一些内容。

\begin{lemma}
\label{more-algebra-lemma-upgrade-adjoint-tensor-RHom}
设 $R \to R'$ 为环映射。对 $K \in D(R)$ 与 $M \in D(R')$，
存在典范同构
$$
R\Hom_R(K, M) = R\Hom_{R'}(K \otimes_R^\mathbf{L} R', M)
$$
。
\end{lemma}

\begin{proof}
选取代表 $M$ 的 $R'$-模 K-内射复形 $J^\bullet$。选取拟同构
$J^\bullet \to I^\bullet$，其中 $I^\bullet$ 是 $R$-模的 K-内射
复形。选取代表 $K$ 的 $R$-模 K-平坦复形 $K^\bullet$。考虑映射
% P02-E0472
$$
\Hom^\bullet(K^\bullet \otimes_R R', J^\bullet)
\longrightarrow
\Hom^\bullet(K^\bullet, I^\bullet)
$$
。次数 $n$ 项上的映射由映射
$$
\prod\nolimits_{n = p + q} \Hom_{R'}(K^{-q} \otimes_R R', J^p)
\longrightarrow
\prod\nolimits_{n = p + q} \Hom_R(K^{-q}, I^p)
$$
给出；后者来自与 $K^{-q} \to K^{-q} \otimes_R R'$ 预复合并与
$J^p \to I^p$ 后复合。为完成证明，只需证明在上同调群上得到同构：
% P02-E0472
$$
\Hom_{D(R)}(K, M) = \Hom_{D(R')}(K \otimes_R^\mathbf{L} R', M)
$$
。这成立是因为，由引理 \ref{more-algebra-lemma-tensor-hom-adjoint}，基变换
$- \otimes_R^\mathbf{L} R' : D(R) \to D(R')$ 左伴随于限制函子
$D(R') \to D(R)$。
\end{proof}

\noindent
设 $R \to R'$ 为环映射。有基变换映射
\begin{equation}
\label{more-algebra-equation-base-change-RHom}
R\Hom_R(K, M) \otimes_R^\mathbf{L} R'
\longrightarrow
R\Hom_{R'}(K \otimes_R^\mathbf{L} R', M \otimes_R^\mathbf{L} R')
\end{equation}
，它位于 $D(R')$ 中，并对 $K, M \in D(R)$ 函子性。事实上，由
$- \otimes_R^\mathbf{L} R' : D(R) \to D(R')$ 与限制函子
$D(R') \to D(R)$ 之间的伴随性，这等价于映射
$$
R\Hom_R(K, M)
\longrightarrow
R\Hom_{R'}(K \otimes_R^\mathbf{L} R', M \otimes_R^\mathbf{L} R') =
R\Hom_R(K, M \otimes_R^\mathbf{L} R')
$$
（等号由引理 \ref{more-algebra-lemma-upgrade-adjoint-tensor-RHom} 给出）；为此可使用
典范映射 $M \to M \otimes_R^\mathbf{L} R'$（伴随的单位）。

\begin{lemma}
\label{more-algebra-lemma-base-change-RHom}
设 $R \to R'$ 为环映射，$K, M \in D(R)$。映射
(\ref{more-algebra-equation-base-change-RHom})
$$
R\Hom_R(K, M) \otimes_R^\mathbf{L} R'
\longrightarrow
R\Hom_{R'}(K \otimes_R^\mathbf{L} R', M \otimes_R^\mathbf{L} R')
$$
在以下情形中是 $D(R')$ 中的同构：
\begin{enumerate}
\item $K$ 完美；
\item $R'$ 作为 $R$-模完美；
\item $R \to R'$ 平坦，$K$ 伪凝聚，且 $M \in D^{+}(R)$；或者
% P02-E0473
\item $R'$ 作为 $R$-模具有有限 Tor 维数，$K$ 伪凝聚，且
$M \in D^{+}(R)$。
\end{enumerate}
\end{lemma}

\begin{proof}
可以在应用限制函子 $D(R') \to D(R)$ 后检验该映射是同构。应用此
函子后，我们的映射成为
% P02-E0474
$$
R\Hom_R(K, L) \otimes_R^\mathbf{L} R'
\longrightarrow
R\Hom_R(K, L \otimes_R^\mathbf{L} R')
$$
，即引理 \ref{more-algebra-lemma-internal-hom-diagonal-better} 的映射。为匹配左右
两端，见该引理上方的讨论；特别地，这里使用引理
\ref{more-algebra-lemma-upgrade-adjoint-tensor-RHom}。于是由引理
\ref{more-algebra-lemma-internal-hom-evaluate-tensor-isomorphism} 得出结论。
\end{proof}

\section{模的系统}
\label{more-algebra-section-systems}

\noindent
设 $I$ 为 Noether 环 $A$ 的理想。本节将补充我们关于有限 $A$-模与
有限 $A/I^n$-模系统之间关系的认识。

\begin{lemma}
\label{more-algebra-lemma-consequence-Artin-Rees}
设 $I$ 为 Noether 环 $A$ 的理想。设
$
K \xrightarrow{\alpha} L \xrightarrow{\beta} M
$
为有限 $A$-模的复形。置 $H = \Ker(\beta)/\Im(\alpha)$。对
$n \geq 0$，设
$$
K/I^nK \xrightarrow{\alpha_n} L/I^nL \xrightarrow{\beta_n} M/I^nM
$$
为诱导复形。置 $H_n = \Ker(\beta_n)/\Im(\alpha_n)$。则存在典范
$A$-模映射，给出交换图
$$
\xymatrix{
& & & H \ar[lld] \ar[ld] \ar[d] \\
\ldots \ar[r] & H_3 \ar[r] & H_2 \ar[r] & H_1
}
$$
。此外，存在 $c > 0$，并且当 $n \geq c$ 时存在典范 $A$-模映射
$H_n \to H/I^{n - c}H$，使复合
$$
H/I^n H \to H_n \to  H/I^{n - c}H
\quad\text{且}\quad
H_n \to H/I^{n - c}H \to H_{n - c}
$$
都是典范映射。再者：
\begin{enumerate}
\item $(H_n)$ 与 $(H/I^nH)$ 作为 $\text{Mod}_A$ 的 pro-对象同构；
\item $\lim H_n = \lim H/I^n H$；
\item 逆系统 $(H_n)$ 满足 Mittag-Leffler 性质；
\item $H_{n + c} \to H_n$ 的像等于 $H \to H_n$ 的像；
\item 复合
$I^cH_n \to H_n \to H/I^{n - c}H \to H_n/I^{n - c}H_n$
是嵌入 $I^cH_n \to H_n$ 后接商映射
$H_n \to H_n/I^{n - c}H_n$；且
% P02-E0478
\item $H/I^nH \to H_n$ 的核与余核被 $I^c$ 零化。
\end{enumerate}
\end{lemma}

\begin{proof}
注意，
% P02-E0475
$H_n = \beta^{-1}(I^nM)/\Im(\alpha) + I^nL$。对 $n \geq 2$，有
$\beta^{-1}(I^nM) \subset \beta^{-1}(I^{n - 1}M)$ 以及
$\Im(\alpha) + I^nL \subset \Im(\alpha) + I^{n - 1}L$。因此得到典范
映射 $H_n \to H_{n - 1}$。类似地，
$\Ker(\beta) \subset \beta^{-1}(I^nM)$ 且
$\Im(\alpha) \subset \Im(\alpha) + I^nL$，从而产生典范映射
$H \to H_n$。我们略去该图交换的验证。

\medskip\noindent
由 Artin-Rees，可以选取 $c_1, c_2 \geq 0$，使得当
$n \geq c_1$ 时
$\beta^{-1}(I^nM) \subset \Ker(\beta) + I^{n - c_1}L$，且当
$n \geq c_2$ 时
$\Ker(\beta) \cap I^nL \subset I^{n - c_2}\Ker(\beta)$；见《代数》
引理 \ref{algebra-lemma-map-AR} 与
\ref{algebra-lemma-Artin-Rees}。置 $c = c_1 + c_2$。

\medskip\noindent
设 $n \geq c$。如下定义
$\psi_n : H_n \to H/I^{n - c}H$。取 $x \in H_n$，并选取代表 $x$ 的
$y \in \beta^{-1}(I^nM)$。写成 $y = z + w$，其中
$z \in \Ker(\beta)$ 且 $w \in I^{n - c_1}L$（由 $c_1$ 的选取可行）。
定义 $\psi_n(x)$ 为 $z$ 在 $H/I^{n - c}H$ 中的类。为证明这是良定义
的，假设对 $x$ 有另一组选取 $y', z', w'$，记号含义同上。此时
$y' - y \in \Im(\alpha) + I^nL$；写
$y' - y = \alpha(v) + u$，其中 $v \in K$ 且 $u \in I^nL$。于是
$$
y' = z' + w' = \alpha(v) + u + z + w
\Rightarrow
z' = z + \alpha(v) + u + w - w'
$$
。因为 $\beta(z' - z - \alpha(v)) = 0$，得到
$u + w - w' \in \Ker(\beta) \cap I^{n - c_1}L$；由 $c_2$ 的选取，
后者包含于
$I^{n - c_1 - c_2}\Ker(\beta) = I^{n - c}\Ker(\beta)$。故按要求，
$z'$ 与 $z$ 在 $H/I^{n - c}H$ 中有相同的像。

\medskip\noindent
复合 $H/I^n H \to H_n \to  H/I^{n - c}H$ 是典范映射。事实上，若
$z \in \Ker(\beta)$ 代表
% P02-E0475
$H/I^nH = \Ker(\beta)/\Im(\alpha) + I^n\Ker(\beta)$ 中的元素 $x$，
则由上面的构造，$x$ 在所得映射下映到 $z$ 在
$H/I^{n - c}H$ 中的类。

\medskip\noindent
考虑复合 $H_n \to H/I^{n - c}H \to H_{n - c}$。给定上面构造
$\psi_n$ 时的 $x, y, z, w$，可见 $x$ 映到 $z$ 在 $H_{n - c}$ 中的
% P02-E0476
类。另一方面，证明第一段的典范映射 $H_n \to H_{n - c}$ 把 $x$
送到 $y$ 的类。因此须证明
$y - z \in \Im(\alpha) + I^{n - c}L$；这成立，因为
$y - z = w \in I^{n - c_1}L \subset I^{n - c}L$。

\medskip\noindent
(1)--(4) 是刚才所证内容的形式推论。具体地，(1) 由这些映射的存在性
以及《范畴》评注 \ref{categories-remark-pro-category} 中 pro-对象
态射的定义得出。(2) 成立，因为同构的 pro-对象具有同构的极限。
(3) 立即由 (4) 得出。(4) 由典范映射 $H_{n + c} \to H_n$ 的分解
$H_{n + c} \to H/I^nH \to H_n$ 得出。

\medskip\noindent
证明 (5)。设 $x \in I^cH_n$。写成 $x = \sum f_i x_i$，其中
$x_i \in H_n$ 且 $f_i \in I^c$。对 $x_i$，按 $\psi_n$ 的构造选取
$y_i, z_i, w_i$。则计算 $x$ 的 $\psi_n$ 时，可取
$y = \sum f_iy_i$、$z = \sum f_i z_i$ 与 $w = \sum f_i w_i$，从而
可见 $\psi_n(x)$ 由 $z$ 的类给出。它在
$H_n/I^{n - c}H_n$ 中的像等于 $y$ 的类，因为
$w = \sum f_i w_i$ 位于 $I^nL$ 中。这证明 (5)。

\medskip\noindent
证明 (6)。设 $y \in \Ker(\beta)$，其在 $H$ 中的类为 $x$。若 $x$
在 $H_n$ 中映为零，则 $y \in I^nL + \Im(\alpha)$。故对某个
$v \in K$，有
$y - \alpha(v) \in \Ker(\beta) \cap I^nL$。于是
$y - \alpha(v) \in I^{n - c_2}\Ker(\beta)$，从而 $y$ 在 $H/I^nH$
中的类被 $I^{c_2}$ 零化。最后，设 $x \in H_n$ 是
$y \in \beta^{-1}(I^nM)$ 的类。仍如上写成 $y = z + w$，其中
$z \in \Ker(\beta)$ 且 $w \in I^{n - c_1}L$。显然，若
$f \in I^{c_1}$，则
% P02-E0477
$fx$ 是 $fy + fw \equiv fy$ 模 $\Im(\alpha) + I^nL$ 的类，因而
按要求，$fx$ 是 $H$ 中 $fy$ 的类之像。
\end{proof}

\begin{lemma}
\label{more-algebra-lemma-kollar-kovacs-pseudo-coherent}
\begin{reference}
% P02-E0479
Kovacs 于 23/02/2018 发来的电子邮件。
\end{reference}
设 $I$ 为 Noether 环 $A$ 的理想。设 $K \in D(A)$ 伪凝聚。置
$K_n = K \otimes_A^\mathbf{L} A/I^n$。则对所有
$i \in \mathbf{Z}$，系统 $H^i(K_n)$ 满足 Mittag-Leffler 性质，并且
$\lim H^i(K)/I^nH^i(K)$ 等于 $\lim H^i(K_n)$。
\end{lemma}

\begin{proof}
可以用有限自由 $A$-模的上有界复形 $P^\bullet$ 代表 $K$。于是
$K_n$ 由 $P^\bullet/I^nP^\bullet$ 代表。因此，
% P02-E0480
由引理 \ref{more-algebra-lemma-consequence-Artin-Rees} 得到 Mittag-Leffler 性质。
最后一个陈述随即由引理 \ref{more-algebra-lemma-kollar-kovacs} 得出。
\end{proof}

\begin{lemma}
\label{more-algebra-lemma-derived-completion-plain-completion}
设 $A$ 为 Noether 环，$I \subset A$ 为理想。设 $M^\bullet$ 为有限
$A$-模的有界复形。映射的逆系统
% P02-E0481
$$
M^\bullet \otimes_A^\mathbf{L} A/I^n \longrightarrow M^\bullet/I^nM^\bullet
$$
定义 $D(A)$ 的 pro-对象的一个同构。
\end{lemma}

\begin{proof}
设 $I = (f_1, \ldots, f_r)$。令 $K_n \in D(A)$ 为由
$f_1^n, \ldots, f_r^n$ 上的 Koszul 复形所代表的对象。回忆，存在
映射 $K_n \to A/I^n$，它们诱导逆系统的 pro-同构；见引理
\ref{more-algebra-lemma-sequence-Koszul-complexes}。因此，只需证明
$$
M^\bullet \otimes_A^\mathbf{L} K_n \longrightarrow M^\bullet/I^nM^\bullet
$$
定义 $D(A)$ 的 pro-对象的一个同构。由于 $K_n$ 由位于次数
$-r, \ldots, 0$ 的有限自由 $A$-模复形代表，存在
$a, b \in \mathbf{Z}$，使得对所有 $n$，所示箭头的源与靶的上同调
都在 $[a, b]$ 之外消失。因此，可应用《导出范畴》引理
\ref{derived-lemma-pro-isomorphism-bis}，得到只需证明：对任意
$i \in \mathbf{Z}$，映射
$$
H^i(M^\bullet \otimes_A^\mathbf{L} A/I^n) \to H^i(M^\bullet/I^nM^\bullet)
$$
定义 $A$-模的 pro-系统的同构。为此，选取拟同构
$P^\bullet \to M^\bullet$，其中 $P^\bullet$ 是有限自由 $A$-模的
上有界复形。上述箭头由映射
% P02-E0481
$$
H^i(P^\bullet/I^nP^\bullet) \to H^i(M^\bullet/I^nM^\bullet)
$$
给出。由引理 \ref{more-algebra-lemma-consequence-Artin-Rees}，它们定义
pro-系统的同构。具体而言，该引理表明两者都同构于 pro-系统
$H^i/I^nH^i$，其中
$H^i = H^i(M^\bullet) = H^i(P^\bullet)$。
\end{proof}

\begin{lemma}
\label{more-algebra-lemma-hom-systems-ML}
设 $A$ 为 Noether 环，$I \subset A$ 为理想。设 $M$、$N$ 为有限
$A$-模。置 $M_n = M/I^nM$ 与 $N_n = N/I^nN$。则：
\begin{enumerate}
\item 系统 $(\Hom_A(M_n, N_n))$ 与
$(\text{Isom}_A(M_n, N_n))$ 满足 Mittag-Leffler 性质；
\item 存在 $c \geq 0$，使得对所有 $n$，映射
$$
\Hom_A(M, N)/I^n\Hom_A(M, N) \to \Hom_A(M_n, N_n)
$$
的核与余核被 $I^c$ 零化；
\item 有
$\lim \Hom_A(M_n, N_n) =\Hom_A(M, N)^\wedge =
\Hom_{A^\wedge}(M^\wedge, N^\wedge)$；且
% P02-E0482
\item $\lim \text{Isom}_A(M_n, N_n) =
\text{Isom}_{A^\wedge}(M^\wedge, N^\wedge)$。
\end{enumerate}
这里 ${}^\wedge$ 表示通常的 $I$-进完备化。
\end{lemma}

\begin{proof}
注意 $\Hom_A(M_n, N_n) = \Hom_A(M, N_n)$。选取表示
$$
A^{\oplus t} \to A^{\oplus s} \to M \to 0
$$
。应用
% P02-E0483
右正合函子 $\Hom_A(-, N)$，得到复形
$$
0 \xrightarrow{\alpha} N^{\oplus s} \xrightarrow{\beta} N^{\oplus t}
$$
，其在中间项的上同调为 $\Hom_A(M, N)$；并且对 $n \geq 0$，
复形
$$
0 \xrightarrow{\alpha_n} N_n^{\oplus s} \xrightarrow{\beta_n} N_n^{\oplus t}
$$
的上同调为 $\Hom_A(M_n, N_n)$。对这个 $A$、$I$、$\alpha$ 与
$\beta$，令 $c \geq 0$ 如引理
\ref{more-algebra-lemma-consequence-Artin-Rees} 所给。由该引理 (3)，得到系统
$(\Hom_A(M_n, N_n))$ 的 Mittag-Leffler 性质。映射
$\Hom_A(M, N)/I^n\Hom_A(M, N) \to \Hom_A(M_n, N_n)$ 的核与余核
% P02-E0484
由该引理 (6) 被 $I^c$ 零化。由该引理 (2)，得到
$\lim \Hom_A(M_n, N_n) = \Hom_A(M, N)^\wedge$。等式
% P02-E0481
$$
\Hom_{A^\wedge}(M^\wedge, N^\wedge) = \lim \Hom_A(M_n, N_n)
$$
形式地由 $M^\wedge = \lim M_n$、$M_n = M^\wedge/I^nM^\wedge$
以及 $N$ 的相应事实得出；见《代数》引理
\ref{algebra-lemma-completion-complete}。

\medskip\noindent
同构的结果来自对 $(\Hom(M_n, N_n))$ 与 $(\Hom(N_n, M_n))$ 都应用
同态的情形，以及如下事实：若 $n > m > 0$，且有映射
$\alpha : M_n \to N_n$ 与 $\beta : N_n \to M_n$，它们诱导
% P02-E0485
同构 $M_m \to N_m$ 与 $N_m \to M_m$，则 $\alpha$ 与 $\beta$
是同构。事实上，由 Nakayama 引理（《代数》引理
\ref{algebra-lemma-NAK}），$\alpha \circ \beta$ 是满射，因而由
《代数》引理 \ref{algebra-lemma-fun}，$\alpha \circ \beta$ 是同构。
\end{proof}

\begin{lemma}
\label{more-algebra-lemma-isomorphic-completions}
设 $A$ 为 Noether 环，$I \subset A$ 为理想。设 $M$、$N$ 为有限
$A$-模。置 $M_n = M/I^nM$ 与 $N_n = N/I^nN$。若对所有 $n$，
$M_n \cong N_n$，则作为 $A^\wedge$-模，
$M^\wedge \cong N^\wedge$。
\end{lemma}

\begin{proof}
由引理 \ref{more-algebra-lemma-hom-systems-ML}，系统
$(\text{Isom}_A(M_n, N_n))$ 满足 Mittag-Leffler 性质。由假设，
每个集合 $\text{Isom}_A(M_n, N_n)$ 都非空，故
$\lim \text{Isom}_A(M_n, N_n)$ 非空。由于
$\lim \text{Isom}_A(M_n, N_n) =
\text{Isom}_{A^\wedge}(M^\wedge, N^\wedge)$，得到一个同构。
\end{proof}

\begin{remark}
\label{more-algebra-remark-weird-systems}
设 $I$ 为 Noether 环 $A$ 的理想。对 $n \geq 1$，置
$A_n = A/I^n$。考虑如下范畴：
\begin{enumerate}
\item 对象是序列 $\{E_n\}_{n \geq 1}$，其中 $E_n$ 是有限
$A_n$-模。
\item 态射 $\{E_n\} \to \{E'_n\}$ 由映射
$$
\varphi_n : I^cE_n \longrightarrow E'_n/E'_n[I^c]
\quad\text{对 }n \geq c
$$
给出，其中 $E'_n[I^c]$ 是挠子模（第 \ref{more-algebra-section-torsion} 节）；
并对以下等价关系取商：称 $(c, \varphi_n)$ 与
$(c + 1, \overline{\varphi}_n)$ 相同，其中
$\overline{\varphi}_n : I^{c + 1}E_n \longrightarrow E'_n/E'_n[I^{c + 1}]$
是诱导映射。
\end{enumerate}
态射 $(c, \varphi_n) : \{E_n\} \to \{E'_n\}$ 与
$(c', \varphi'_n) : \{E'_n\} \to \{E''_n\}$ 的复合由显然的复合
% P02-E0486
$$
I^{c + c'}E_n \to I^{c'}E'_n/E'_n[I^{c}] \to E''_n/E''_n[I^{c + c'}]
$$
定义，其中 $n \geq c + c'$。我们略去它确为一个范畴的验证。
\end{remark}

\begin{lemma}
\label{more-algebra-lemma-iso}
评注 \ref{more-algebra-remark-weird-systems} 的范畴中的态射
$(c, \varphi_n)$ 是同构，当且仅当存在 $c' \geq 0$，使得对所有
$n \gg 0$，$\Ker(\varphi_n)$ 与 $\Coker(\varphi_n)$ 都是
$I^{c'}$-挠的。
\end{lemma}

\begin{proof}
可以并且确实假设 $c' \geq c$，且对所有 $n$，
$\Ker(\varphi_n)$ 与 $\Coker(\varphi_n)$ 都是 $I^{c'}$-挠的。对
$n \geq c'$ 及 $x \in I^{c'}E'_n$，由于
$\Coker(\varphi_n)$ 被 $I^{c'}$ 零化，可以选取 $y \in I^cE_n$，
使得 $x = \varphi_n(y) \bmod E'_n[I^c]$。定义 $\psi_n(x)$ 为 $y$
在 $E_n/E_n[I^{c'}]$ 中的类。若另有选取 $y' \in I^cE_n$ 满足
$x = \varphi_n(y') \bmod E'_n[I^c]$，则差 $y - y'$ 在
$E'_n/E'_n[I^c]$ 中映为零，因而它在 $I^cE_n$ 中被 $I^{c'}$
零化。因此，映射
$\psi_n : I^{c'}E'_n \to E_n/E_n[I^{c'}]$ 良定义。我们略去验证
$(c', \psi_n)$ 是该范畴中 $(c, \varphi_n)$ 的逆。
\end{proof}

\begin{lemma}
\label{more-algebra-lemma-dejong-kollar-kovacs}
\begin{reference}
% P02-E0479
Janos Kollar、Sandor Kovacs 与 Johan de Jong 于 23/02/2018
之间的电子邮件通信。
\end{reference}
设 $I$ 为 Noether 环 $A$ 的理想。设 $M$ 与 $N$ 为有限 $A$-模。
写 $A_n = A/I^n$、$M_n = M/I^nM$ 以及 $N_n = N/I^nN$。对每个
$i \geq 0$，对象
$$
\{\Ext^i_A(M, N)/I^n\Ext^i_A(M, N)\}_{n \geq 1}
\quad\text{且}\quad
\{\Ext^i_{A_n}(M_n, N_n)\}_{n \geq 1}
$$
在评注 \ref{more-algebra-remark-weird-systems} 的范畴 $\mathcal{C}$ 中同构。
\end{lemma}

\begin{proof}
选取短正合列
$$
0 \to K \to A^{\oplus r} \to M \to 0
$$
，并置 $K_n = K/I^nK$。对 $n \geq 1$，定义
$K(n) = \Ker(A_n^{\oplus r} \to M_n)$，于是有正合列
$$
0 \to K(n) \to A_n^{\oplus r} \to M_n \to 0
$$
以及满射 $K_n \to K(n)$。事实上，由引理
\ref{more-algebra-lemma-consequence-Artin-Rees}，存在 $c \geq 0$ 与“几乎互逆”
的映射 $K(n) \to K_n/I^{n - c}K_n$。由于
$I^{n - c}K_n \subset K_n[I^c]$，
% P02-E0487
这些映射见证了系统 $\{K(n)\}_{n \geq 1}$ 与
$\{K_n\}_{n \geq 1}$ 在 $\mathcal{C}$ 中同构。

\medskip\noindent
我们断言系统
$$
\{\Ext^i_{A_n}(K(n), N_n)\}_{n \geq 1}
\quad\text{且}\quad
\{\Ext^i_{A_n}(K_n, N_n)\}_{n \geq 1}
$$
在范畴 $\mathcal{C}$ 中同构。事实上，满射 $K_n \to K(n)$ 的核
被 $I^c$ 零化，因此确定映射
$$
\Ext^i_{A_n}(K(n), N_n) \to \Ext^i_{A_n}(K_n, N_n)
$$
，其核与余核被 $I^c$ 零化。于是该断言由引理
\ref{more-algebra-lemma-iso} 得出。

\medskip\noindent
当 $i \geq 2$ 时，有同构
$$
\Ext^{i - 1}_A(K, N) = \Ext^i_A(M, N)
\quad\text{且}\quad
\Ext^{i - 1}_{A_n}(K(n), N_n) = \Ext^i_{A_n}(M_n, N_n)
$$
。由此可见，只需对 $i = 0, 1$ 证明引理。

\medskip\noindent
对 $i = 0, 1$，考虑交换图
{\footnotesize
$$
\xymatrix{
0 \ar[r] &
\Hom(M, N) \ar[r] \ar[dd] &
N^{\oplus r} \ar[r]_-\varphi \ar[dd] &
\Hom(K, N) \ar[r] \ar[d] &
\Ext^1(M, N) \ar[r] &
0 \\
& & &
\Hom(K_n, N_n)
\\
0 \ar[r] &
\Hom(M_n, N_n) \ar[r] &
N_n^{\oplus r} \ar[r] &
\Hom(K(n), N_n) \ar[r] \ar[u] &
\Ext^1(M_n, N_n) \ar[r] &
0
}
$$
}
。由引理 \ref{more-algebra-lemma-hom-systems-ML}，可见映射
$$
\Hom(M, N)/I^n \Hom(M, N) \to \Hom(M_n, N_n)
$$
与
% P02-E0488
$$
\Hom(K, N)/I^n \Hom(K, N) \to \Hom(K_n, N_n)
$$
的核与余核，被某个
$c \geq 0$ 所对应且与 $n$ 无关的 $I^c$ 零化。上面已经看到，在必要时增大 $c$ 后，
单射 $\Hom(K(n), N_n) \to \Hom(K_n, N_n)$ 的余核被 $I^c$ 零化。
对这样的 $c$，得到嵌入图中的映射
$$
\delta_n : I^c\Hom(K, N)/I^n\Hom(K, N) \to \Hom(K(n), N_n)
$$
（略去精确定式）。必要时增大 $c$ 后，$\delta_n$ 的核与余核被
$I^c$ 零化，因为已知
$$
\Hom(K, N)/I^n \Hom(K, N) \to \Hom(K_n, N_n)
$$
及
$$
\Hom(K(n), N_n) \to \Hom(K_n, N_n)
$$
具有同一性质。然后利用实线图
的交换性
{\scriptsize
$$
\xymatrix{
\varphi^{-1}(I^c\Hom(K, N)) \ar[r]_-\varphi \ar[d] &
I^c\Hom(K, N)/I^n\Hom(K, N) \ar[r] \ar[d]^{\delta_n} &
I^c\Ext^1(M, N)/I^n\Ext^1(M, N) \ar[r] \ar@{..>}[d] & 0 \\
N_n^{\oplus r} \ar[r] &
\Hom(K(n), N_n) \ar[r] &
\Ext^1(M_n, N_n) \ar[r] & 0
}
$$
}
来定义虚线箭头。直接追图（略去）表明，在最后一次必要地增大 $c$
之后，虚线箭头的核与余核
% P02-E0489
被 $I^c$ 零化。
\end{proof}

\begin{remark}
\label{more-algebra-remark-awkward}
引理 \ref{more-algebra-lemma-dejong-kollar-kovacs} 陈述中的别扭，部分源于对变化的
$n$，模 $\Ext^i_{A_n}(M_n, N_n)$ 之间没有明显映射这一事实。由该
引理可以推出：存在 $c \geq 0$，使得当 $m \gg n \gg 0$ 时，有
（典范）映射
% P02-E0490
$$
I^c\Ext^i_{A_n}(M_m, N_m)/I^n\Ext^i_{A_n}(M_m, N_m) \to
\Ext^i_{A_n}(M_n, N_n)/\Ext^i_{A_n}(M_n, N_n)[I^c]
$$
，其核与余核被 $I^c$ 零化。这正是我们得到模的系统的（弱）意义。
\end{remark}

\begin{example}
\label{more-algebra-example-has-to-be-awkward}
设 $k$ 为域，$A = k[[x, y]]/(xy)$。稍微滥用记号，
% P02-E0491
以 $x$ 与 $y$ 表示 $x$ 与 $y$ 在 $A$ 中的像。设 $I = (x)$，
$M = A/(y)$。有自由分解
% P02-E0481
$$
\ldots \to
A \xrightarrow{y}
A \xrightarrow{x}
A \xrightarrow{y} A \to M \to 0
$$
。由此得到
$$
\Ext^2_A(M, N) = N[y]/xN
$$
，其中 $N[y] = \Ker(y : N \to N)$。以
$A_n = A/I^n$、$M_n = M/I^nM$ 以及 $N_n = N/I^nN$ 表示相应对象。
对每个 $n$，有自由分解
$$
\ldots \to
A_n^{\oplus 2} \xrightarrow{y, x^{n - 1}}
A_n \xrightarrow{x}
A_n \xrightarrow{y} A_n \to M_n \to 0
$$
。由此得到
$$
\Ext^2_{A_n}(M_n, N_n) = (N_n[y] \cap N_n[x^{n - 1}])/xN_n
$$
，其中 $N_n[y] = \Ker(y : N_n \to N_n)$，且
% P02-E0492
$N[x^{n - 1}] = \Ker(x^{n - 1} : N_n \to N_n)$。取
$N = A/(y)$。则
$$
\Ext^2_A(M, N) = N[y]/xN = N/xN \cong k
$$
，但
$$
\Ext^2_{A_n}(M_n, N_n) = (N_n[y] \cap N_n[x^{n - 1}])/xN_n =
N_n[x^{n - 1}]/xN_n = 0
$$
对所有
% P02-E0493
$r$ 都成立，因为 $N_n = k[x]/(x^n)$，且序列
$$
N_n \xrightarrow{x} N_n \xrightarrow{x^{n - 1}} N_n
$$
正合。因此，为得到引理 \ref{more-algebra-lemma-dejong-kollar-kovacs} 那样的结果，
必须忽略某种 $I$-幂挠。
\end{example}

\begin{lemma}
\label{more-algebra-lemma-not-awkward}
\begin{reference}
% P02-E0479
Janos Kollar、Sandor Kovacs 与 Johan de Jong 于 23/02/2018
之间的电子邮件通信。
\end{reference}
设 $A \to B$ 为 Noether 环的平坦同态，$I \subset A$ 为理想。
% P02-E0494
设 $M, N$ 为 $A$-模。置 $B_n = B/I^nB$、
$M_n = M/I^nM$、$N_n = N/I^nN$。若 $M$ 在 $A$ 上平坦，则对所有
$i \in \mathbf{Z}$ 有
$$
\lim \Ext^i_B(M, N)/I^n \Ext^i_B(M, N) =
\lim \Ext^i_{B_n}(M_n, N_n)
$$
。
\end{lemma}

\begin{proof}
选取有限自由 $B$-模 $P_i$ 构成的分解
% P02-E0495
$$
\ldots \to P_2 \to P_1 \to P_0 \to M \to 0
$$
。置 $P_{i, n} = P_i/I^nP_i$。由于 $M$ 与 $B$ 在 $A$ 上平坦，
序列
$$
\ldots \to P_{2, n} \to P_{1, n} \to P_{0, n} \to M_n \to 0
$$
正合。一方面，复形
$$
\Hom_B(P_0, N) \to \Hom_B(P_1, N) \to \Hom_B(P_2, N) \to \ldots
$$
计算各模 $\Ext^i_B(M, N)$；另一方面，复形
$$
\Hom_{B_n}(P_{0, n}, N_n) \to \Hom_{B_n}(P_{1, n}, N_n) \to
\Hom_{B_n}(P_{2, n}, N_n) \to \ldots
$$
计算各模 $\Ext^i_{B_n}(M_n, N_n)$。由于
$$
\Hom_{B_n}(P_{i, n}, N_n) = \Hom_B(P_i, N)/I^n \Hom_B(P_i, N)
$$
，由引理 \ref{more-algebra-lemma-consequence-Artin-Rees} 的 (2) 得到结果。
\end{proof}

\section{模的系统（续）}
\label{more-algebra-section-systems-bis}

\noindent
设 $I$ 为 Noether 环 $A$ 的理想。在第 \ref{more-algebra-section-systems} 节中，
我们考察了有限 $A$-模 $M$ 的 $M/I^nM$ 型系统。本节改为考察系统
$I^nM$。

\begin{lemma}
\label{more-algebra-lemma-consequence-Artin-Rees-bis}
设 $I$ 为 Noether 环 $A$ 的理想。设
$
K \xrightarrow{\alpha} L \xrightarrow{\beta} M
$
为有限 $A$-模的复形。置 $H = \Ker(\beta)/\Im(\alpha)$。对
$n \geq 0$，设
$$
I^nK \xrightarrow{\alpha_n} I^nL \xrightarrow{\beta_n} I^nM
$$
为诱导复形。置 $H_n = \Ker(\beta_n)/\Im(\alpha_n)$。则存在典范
$A$-模映射
$$
\ldots \to H_3 \to H_2 \to H_1 \to H
$$
。存在 $c > 0$，使得当 $n \geq c$ 时，$H_n \to H$ 的像包含在
$I^{n - c}H$ 中；并存在典范 $A$-模映射
$I^nH \to H_{n - c}$，使复合
$$
I^n H \to H_{n - c} \to  I^{n - 2c}H
\quad\text{且}\quad
H_n \to I^{n - c}H \to H_{n - 2c}
$$
都是典范映射。特别地，逆系统 $(H_n)$ 与 $(I^nH)$ 作为
$\text{Mod}_A$ 的 pro-对象同构。
\end{lemma}

\begin{proof}
有
% P02-E0496
$H_n = \Ker(\beta) \cap I^nL/\alpha(I^nK)$。由于
$\Ker(\beta) \cap I^nL \subset \Ker(\beta) \cap I^{n - 1}L$ 且
$\alpha(I^nK) \subset \alpha(I^{n - 1}K)$，得到映射
$H_n \to H_{n - 1}$。映射 $H_1 \to H$ 的情形类似。

\medskip\noindent
由 Artin-Rees，可以选取 $c_1, c_2 \geq 0$，使得当
$n \geq c_1$ 时
$\Im(\alpha) \cap I^nL \subset \alpha(I^{n - c_1}K)$，且当
$n \geq c_2$ 时
$\Ker(\beta) \cap I^nL \subset I^{n - c_2}\Ker(\beta)$；见《代数》
引理 \ref{algebra-lemma-map-AR} 与
\ref{algebra-lemma-Artin-Rees}。置 $c = c_1 + c_2$。

\medskip\noindent
由 $c \geq c_2$ 的选取立即得到：当 $n \geq c$ 时，
$H_n \to H$ 的像包含在 $I^{n - c}H$ 中。

\medskip\noindent
设 $n \geq c$。如下定义 $\psi_n : I^nH \to H_{n - c}$。取
$x \in I^nH$，并选取代表 $x$ 的 $y \in I^n\Ker(\beta)$。定义
$\psi_n(x)$ 为 $y$ 在 $H_{n - c}$ 中的类。为证明其良定义，假设
$y'$ 是 $x$ 的另一个上述代表。则
$y' - y \in \Im(\alpha)$。由 $c \geq c_1$ 的选取，得到
$y' - y \in \alpha(I^{n - c}K)$，这意味着 $y$ 与 $y'$ 代表
$H_{n - c}$ 的同一个元素。因此 $\psi_n$ 良定义。

\medskip\noindent
关于复合 $I^n H \to H_{n - c} \to  I^{n - 2c}H$ 与
$H_n \to I^{n - c}H \to H_{n - 2c}$ 的陈述立即由定义得出。
\end{proof}

\begin{lemma}
\label{more-algebra-lemma-ext-factors}
设 $A$ 为 Noether 环，$I \subset A$ 为理想。设 $M$、$N$ 为
$A$-模，其中 $M$ 有限。对每个 $p > 0$，存在 $c \geq 0$，使得当
$n \geq c$ 时，映射
$\Ext_A^p(M, N) \to \Ext_A^p(I^nM, N)$ 经过映射
$\Ext^p_A(I^nM, I^{n - c}N) \to \Ext_A^p(I^nM, N)$ 分解。
\end{lemma}

\begin{proof}
对 $p = 0$，若 $\varphi : M \to N$ 是 $A$-线性映射，则对
% P02-E0497
$f_i \in A$ 与 $m_i \in M$，有
$\varphi(\sum f_i m_i) = \sum f_i \varphi(m_i)$。因此，对所有 $n$，
$\varphi$ 诱导映射 $I^nM \to I^nN$，结论对 $c = 0$ 成立。

\medskip\noindent
选取短正合列 $0 \to K \to A^{\oplus t} \to M \to 0$。对每个 $n$，
选取短正合列 $0 \to L_n \to A^{\oplus s_n} \to I^nM \to 0$。
显然可以构造短正合列之间的映射
$$
\xymatrix{
0 \ar[r] &
L_n \ar[r] \ar[d] &
A^{\oplus s_n} \ar[r] \ar[d] &
I^nM \ar[r] \ar[d] & 0 \\
0 \ar[r] &
K \ar[r] &
A^{\oplus t} \ar[r] &
M \ar[r] & 0
}
$$
，使得 $A^{\oplus s_n} \to A^{\oplus t}$ 的像位于
$(I^n)^{\oplus t}$ 中。由 Artin-Rees（《代数》引理
\ref{algebra-lemma-Artin-Rees}），存在 $c \geq 0$，使得当
$n \geq c$ 时，$L_n \to K$ 经过 $I^{n - c}K$ 分解。

\medskip\noindent
对 $p = 1$，上述选取诱导实线交换图
{\footnotesize
$$
\xymatrix{
\Hom_A(A^{\oplus s_n}, N) \ar[r] &
\Hom_A(L_n, N) \ar[r] &
\Ext_A^1(I^nM, N) \ar[r] & 0 \\
\Hom_A((I^n)^{\oplus t}, I^{n - c}N) \ar[r] \ar[u] &
\Hom_A(K \cap (I^n)^{\oplus t}, I^{n - c}N) \ar[r] \ar[u] &
\Ext_A^1(I^nM, I^{n - c}N) \ar[u] \\
\Hom_A(A^{\oplus t}, N) \ar[r] \ar[u] &
\Hom_A(K, N) \ar[r] \ar[u] &
\Ext_A^1(M, N) \ar@{..>}[u] \ar[r] & 0
}
$$
}
，其水平箭头正合。下方中间的竖直箭头来自
$K \cap (I^n)^{\oplus t} \subset I^{n - c}K$；因此，按第一段的论证，
% P02-E0498
任意 $A$-线性映射 $K \to N$ 都诱导 $A$-线性映射
$(I^n)^{\oplus t} \to I^{n - c}N$。于是按要求得到虚线箭头。

\medskip\noindent
对 $p > 1$，得到交换图
$$
\xymatrix{
\Ext^{p - 1}_A(I^{n - c}K, N) \ar[r] &
\Ext^{p - 1}_A(L_n, N) \ar[r] &
\Ext_A^p(I^nM, N) \\
\Ext^{p - 1}_A(K, N) \ar[rr] \ar[u] & &
\Ext_A^p(M, N) \ar[u]
}
$$
，其下方水平箭头是同构。对 $p$ 归纳，存在某个 $c' \geq 0$，
使得对所有 $n \geq c + c'$，左侧竖直映射经过
$\Ext^{p - 1}_A(I^{n - c}K, I^{n - c - c'}N)$ 分解。使用复合
$\Ext^{p - 1}_A(I^{n - c}K, I^{n - c - c'}N) \to
\Ext^{p - 1}_A(L_n, I^{n - c - c'}N) \to
\Ext^p_A(I^nM, I^{n - c - c'}N)$，便得到所需分解（对
$n \geq c + c'$，并以 $c + c'$ 替换 $c$）。
\end{proof}

\begin{lemma}
\label{more-algebra-lemma-ext-annihilated}
设 $A$ 为 Noether 环，$I \subset A$ 为理想。设 $M$、$N$ 为
$A$-模，其中 $M$ 有限且 $N$ 被 $I$ 的某个幂零化。对每个
$p > 0$，存在 $n$，使映射
$\Ext_A^p(M, N) \to \Ext_A^p(I^nM, N)$ 为零。
\end{lemma}

\begin{proof}
这立即由引理 \ref{more-algebra-lemma-ext-factors} 以及存在 $m > 0$ 使
$I^mN = 0$ 的事实得出。
\end{proof}

\begin{lemma}
\label{more-algebra-lemma-ext-induced-toplogy}
设 $A$ 为 Noether 环，$I \subset A$ 为理想。设 $K \in D(A)$
伪凝聚，$M$ 为有限 $A$-模。对每个 $p \in \mathbf{Z}$，存在
% P02-E0499
$c$，使得当 $n \geq c$ 时，
$\Ext_A^p(K, I^nM) \to \Ext_A^p(K, M)$ 的像包含在
$I^{n - c}\Ext_A^p(K, M)$ 中。
\end{lemma}

\begin{proof}
选取代表 $K$ 的有限自由 $A$-模上有界复形 $P^\bullet$。则
$\Ext_A^p(K, M)$ 是
% P02-E0500
$$
\Hom_A(F^{-p + 1}, M) \xrightarrow{a}
\Hom_A(F^{-p}, M) \xrightarrow{b}
\Hom_A(F^{-p - 1}, M)
$$
的上同调，而 $\Ext_A^p(K, I^nM)$ 由把这些有限 $A$-模替换为它们的
$I^n$ 倍所得复形计算。因此，
% P02-E0501
结果由引理 \ref{more-algebra-lemma-consequence-Artin-Rees-bis} 得出
（事实上还有更强的结论）。
\end{proof}

\noindent
在情形 \ref{more-algebra-situation-koszul} 中，我们定义复形 $I_n^\bullet$，使得
在 $A$-模复形模伦的三角范畴 $K(A)$ 中有区分三角
% P02-E0502
$$
I_n^\bullet \to A \to K_n^\bullet \to I_n^\bullet[1]
$$
。具体地，置 $I_n^\bullet = \sigma_{\leq -1}K_n^\bullet[-1]$。
存在逐项分裂的复形短正合列
$$
0 \to A \to K_n^\bullet \to I_n^\bullet[1] \to 0
$$
，它们由 $K(A)$ 上三角结构的定义给出区分三角。旋转这些三角便得到
上面所需的区分三角。注意，$I_n^0 = A^{\oplus r} \to A$ 在第
$i$ 个因子上由乘以 $f_i^n$ 给出。因此，$I_n^\bullet \to A$
分解为
$$
I_n^\bullet \to (f_1^n, \ldots, f_r^n) \to A
$$
。事实上，有短正合列
$$
0 \to H^{-1}(K_n^\bullet) \to H^0(I_n^\bullet) \to (f_1^n, \ldots, f_r^n) \to 0
$$
，并且对每个 $i < 0$，有
$H^i(I_n^\bullet) = H^{i - 1}(K_n^\bullet)$。映射
$K_{n + 1}^\bullet \to K_n^\bullet$ 诱导映射
$I_{n + 1}^\bullet \to I_n^\bullet$，并得到交换图
$$
\xymatrix{
\ldots \ar[r] &
I_3^\bullet \ar[d] \ar[r] &
I_2^\bullet \ar[d] \ar[r] &
I_1^\bullet \ar[d] \\
\ldots \ar[r] &
(f_1^3, \ldots, f_r^3) \ar[r] &
(f_1^2, \ldots, f_r^2) \ar[r] &
(f_1, \ldots, f_r)
}
$$
，它位于 $K(A)$ 中。

\begin{lemma}
\label{more-algebra-lemma-sequence-powers-pro-bounded}
在情形 \ref{more-algebra-situation-koszul} 中，假设 $A$ 为 Noether 环。沿用上述
记号，逆系统 $(I^n)$ 在 $D(A)$ 中 pro-同构于逆系统
$(I_n^\bullet)$。
\end{lemma}

\begin{proof}
容易证明，逆系统 $I^n$ 在 $A$-模范畴中 pro-同构于逆系统
$(f_1^n, \ldots, f_r^n)$。考虑区分三角的逆系统
$$
I_n^\bullet \to (f_1^n, \ldots, f_r^n) \to C_n^\bullet \to I_n^\bullet[1]
$$
，其中 $C_n^\bullet$ 是第一个箭头的锥。由《导出范畴》引理
\ref{derived-lemma-pro-isomorphism}，只需证明逆系统
$C_n^\bullet$ 是 pro-零的。复形 $I_n^\bullet$ 仅在满足
$-r + 1 \leq i \leq 0$ 的次数 $i$ 上有非零项，故
$C_n^\bullet$ 也类似地有界。因此，由《导出范畴》引理
\ref{derived-lemma-essentially-constant-cohomology}，只需证明
$H^p(C_n^\bullet)$ 是 pro-零的。由上面的讨论，当 $p \leq -1$ 时，
$H^p(C_n^\bullet) = H^p(K_n^\bullet)$；当 $p \geq 0$ 时，
$H^p(C_n^\bullet) = 0$。在引理
\ref{more-algebra-lemma-sequence-Koszul-complexes} 的证明中，已经证明各逆系统
$H^p(K_n^\bullet)$ 是 pro-零的（也正是在这里使用了 $A$ 为
Noether 环这一假设）。
\end{proof}

\begin{lemma}
\label{more-algebra-lemma-tensoring-Deligne-system}
设 $A$ 为 Noether 环，$I \subset A$ 为理想。设 $M^\bullet$ 为有限
$A$-模的有界复形。映射的逆系统
$$
I^n \otimes_A^\mathbf{L} M^\bullet \longrightarrow I^nM^\bullet
$$
定义 $D(A)$ 的 pro-对象的一个同构。
\end{lemma}

\begin{proof}
选取 $I$ 的生成元 $f_1, \ldots, f_r \in I$。逆系统 $I^n$ 在
$A$-模范畴中 pro-同构于逆系统
$(f_1^n, \ldots, f_r^n)$。沿用引理
\ref{more-algebra-lemma-sequence-powers-pro-bounded} 的记号，可见只需证明映射的
逆系统
$$
I_n^\bullet \otimes_A^\mathbf{L} M^\bullet
\longrightarrow
(f_1^n, \ldots, f_r^n)M^\bullet
$$
定义 $D(A)$ 的 pro-对象的一个同构。设 $a \leq b$ 满足：若
$i \not \in [a, b]$，则 $M^i = 0$。于是上述箭头的源与靶仅在次数
$[-r + a, b]$ 中有上同调。因此，只需证明：对任意
$p \in \mathbf{Z}$，映射的逆系统
$$
H^p(I_n^\bullet \otimes_A^\mathbf{L} M^\bullet)
\longrightarrow
H^p((f_1^n, \ldots, f_r^n)M^\bullet)
$$
定义 $A$-模的 pro-对象的一个同构；见《导出范畴》引理
\ref{derived-lemma-pro-isomorphism-bis}。利用
$I_n^\bullet \otimes_A^\mathbf{L} M^\bullet$ 与
$I^n \otimes_A^\mathbf{L} M^\bullet$ 之间的 pro-同构，以及
$(f_1^n, \ldots, f_r^n)M^\bullet$ 与 $I^nM^\bullet$ 之间的
pro-同构，这等价于证明映射的逆系统
$$
H^p(I^n \otimes_A^\mathbf{L} M^\bullet)
\longrightarrow
H^p(I^nM^\bullet)
$$
定义 $A$-模的 pro-对象的一个同构。
% P02-E0502
选取有限自由 $A$-模的上有界复形 $P^\bullet$ 及拟同构
$P^\bullet \to M^\bullet$。于是，只需证明映射的逆系统
$$
H^p(I^nP^\bullet)
\longrightarrow
H^p(I^nM^\bullet)
$$
是 pro-同构。这由引理
\ref{more-algebra-lemma-consequence-Artin-Rees-bis} 以及
$H^p(P^\bullet) = H^p(M^\bullet)$ 得出。
\end{proof}

\begin{lemma}
\label{more-algebra-lemma-factor-through-derived-tensor-product}
设 $A$ 为 Noether 环，$I \subset A$ 为理想，$M$ 为有限 $A$-模。
存在整数 $n > 0$，使得 $I^nM \to M$ 在 $D(A)$ 中经过映射
$I \otimes_A^\mathbf{L} M \to M$ 分解。
\end{lemma}

\begin{proof}
这由引理 \ref{more-algebra-lemma-tensoring-Deligne-system} 得出。
% P02-E0503
也可如下直接看出。考虑区分三角
% P02-E0502
$$
I \otimes_A^\mathbf{L} M \to M \to A/I \otimes_A^\mathbf{L} M \to
I \otimes_A^\mathbf{L} M[1]
$$
。由三角范畴的公理，只需证明对某个 $n$，
$I^nM \to A/I \otimes_A^\mathbf{L} M$ 在 $D(A)$ 中为零。选取 $I$
的生成元 $f_1, \ldots, f_r$，令
$K = K_\bullet(A, f_1, \ldots, f_r)$ 为 Koszul 复形，并考虑商映射
的分解 $A \to K \to A/I$。于是只需证明：对某个 $n > 0$，
$I^nM \to K \otimes_A M$ 在 $D(A)$ 中为零。假设已找到
$n > 0$，使 $I^nM \to K \otimes_A M$ 在 $D(A)$ 中经过
$\tau_{\geq t}(K \otimes_A M)$ 分解。将其进一步经过
$\tau_{\geq t + 1}(K \otimes_A M)$ 分解的障碍，是
$\Ext^t(I^nM, H_t(K \otimes_A M))$ 的一个元素。有限 $A$-模
$H_t(K \otimes_A M)$ 被 $I$ 零化。由引理
\ref{more-algebra-lemma-ext-annihilated}，增大 $n$ 后可假设该障碍元素为零。
有限次重复此过程，便找到 $n$，使得 $I^nM \to K \otimes_A M$
在 $D(A)$ 中经过
$0 = \tau_{\geq r + 1}(K \otimes_A M)$ 分解，证明完成。
\end{proof}

\section{杂项}
\label{more-algebra-section-misc}

\noindent
下面给出一些不适合归入其他章节的结果。

\begin{lemma}
\label{more-algebra-lemma-ext-annihilated-into}
设 $A$ 为 Noether 环，$I \subset A$ 为理想，$K \in D(A)$ 伪凝聚，
$a \in \mathbf{Z}$。假设对每个有限 $A$-模 $M$，当 $i \geq a$ 时，
模 $\Ext^i_A(K, M)$ 都是 $I$-幂挠的。则对 $i \geq a$ 及有限模
$M$，系统 $\Ext^i_A(K, M/I^nM)$ 本质常值，其值为
% P02-E0504
$$
\Ext^i_A(K, M) = \lim \Ext^i_A(K, M/I^nM)
$$
。
\end{lemma}

\begin{proof}
设 $M$ 为有限 $A$-模。由于 $K$ 伪凝聚，$\Ext^i_A(K, M)$ 是有限
$A$-模。因此，当 $i \geq a$ 时，它被某个 $I^t$（$t \geq 0$）
零化。由引理 \ref{more-algebra-lemma-ext-induced-toplogy}，对某个 $n > 0$，
$\Ext^i_A(K, I^nM) \to \Ext^i_A(K, M)$ 的像为零。短正合列
$0 \to I^nM \to M \to M/I^n M \to 0$ 给出长正合列
$$
\Ext^i_A(K, I^nM) \to \Ext^i_A(K, M) \to
\Ext^i_A(K, M/I^nM) \to \Ext^{i + 1}_A(K, I^nM)
$$
。由刚才的论证（将其应用于有限 $A$-模 $I^mM$），系统
$\Ext^i_A(K, I^nM)$ 与 $\Ext^{i + 1}_A(K, I^nM)$ 都本质常值，
其值为 $0$。追图可知，$\Ext^i_A(K, M/I^nM)$ 本质常值，其值为
$\Ext^i_A(K, M)$。
\end{proof}

\begin{lemma}
\label{more-algebra-lemma-tor-annihilated}
设 $A$ 为 Noether 环，$I \subset A$ 为理想，$M$ 为有限 $A$-模，
$N$ 为被 $I$ 零化的 $A$-模。存在整数 $n > 0$，使得对所有
$p \geq 0$，映射
$\text{Tor}^A_p(I^nM, N) \to \text{Tor}^A_p(M, N)$ 为零。
\end{lemma}

\begin{proof}
由引理 \ref{more-algebra-lemma-factor-through-derived-tensor-product}，映射
$I^nM \to M$ 可分解为
$I^nM \to M \otimes_A^\mathbf{L} I \to M$。我们断言复合
$$
I^nM \otimes_A^\mathbf{L} N \to
(M \otimes_A^\mathbf{L} I) \otimes_A^\mathbf{L} N
\to M \otimes_A^\mathbf{L} N
$$
为零。事实上，图
$$
\xymatrix{
(M \otimes_A^\mathbf{L} I) \otimes_A^\mathbf{L} N \ar[rr] \ar[rd] & &
M \otimes_A^\mathbf{L} (I \otimes_A^\mathbf{L} N) \ar[ld] \\
& M \otimes_A^\mathbf{L} N
}
$$
交换（略去细节），而映射 $I \otimes_A^\mathbf{L} N \to N$
为零，因为 $N$ 被 $I$ 零化。
\end{proof}

\begin{lemma}
\label{more-algebra-lemma-pseudo-coherent-tensor-limit}
设 $R$ 为环，$K \in D(R)$ 伪凝聚，$(M_n)$ 为 $R$-模的逆系统。则
$R\lim K \otimes_R^\mathbf{L} M_n = K \otimes_R^\mathbf{L} R\lim M_n$。
\end{lemma}

\begin{proof}
考虑定义性的区分三角
$$
R\lim M_n \to \prod M_n \to \prod M_n \to R\lim M_n[1]
$$
，并应用引理 \ref{more-algebra-lemma-pseudo-coherent-tensor}。
\end{proof}

\begin{lemma}
\label{more-algebra-lemma-additivity-of-pd}
设 $R$ 为 Noether 局部环，$I \subset R$ 为理想，$E$ 为非零有限
$R/I$-模。若 $R/I$ 具有有限投射维数，且 $E$ 作为 $R/I$-模具有
有限投射维数，则 $E$ 作为 $R$-模具有有限投射维数，并且
$$
\text{pd}_R(E) = \text{pd}_R(R/I) + \text{pd}_{R/I}(E)
$$
。
\end{lemma}

\begin{proof}
我们将使用如下事实：对有限模而言，在 $R$（相应地，在 $R/I$）上
具有有限投射维数，等价于它是完美模；见定义
\ref{more-algebra-definition-perfect} 后的讨论。由引理
\ref{more-algebra-lemma-cohomology-perfect}，$E$ 在 $R$ 上具有有限投射维数。
因此可以应用 Auslander-Buchsbaum（《代数》命题
\ref{algebra-proposition-Auslander-Buchsbaum}），得到
$$
\text{pd}_R(E) + \text{depth}(E) = \text{depth}(R),\quad
\text{pd}_{R/I}(E) + \text{depth}(E) = \text{depth}(R/I),
$$
以及
$$
\text{pd}_R(R/I) + \text{depth}(R/I) = \text{depth}(R)
$$
。注意，第一个等式中的 $E$ 的深度按 $R$-模计算，第二个等式中
则按 $R/I$-模计算。不过这两个深度相同（这很显然，也可由《代数》
引理 \ref{algebra-lemma-depth-goes-down-finite} 得出）。证明完毕。
\end{proof}

\begin{lemma}
\label{more-algebra-lemma-enlarge}
设 $A \to B$ 为环映射。存在基数
$\kappa = \kappa(A \to B)$，具有如下性质。设 $M^\bullet$ 为
$A$-模复形，$N^\bullet$ 为 $B$-模复形。设
$a : M^\bullet \to N^\bullet$ 为 $A$-模复形的映射，并且它在
$D(B)$ 中诱导同构
$M^\bullet \otimes_A^\mathbf{L} B \to N^\bullet$。设
$M_1^\bullet \subset M^\bullet$ 为 $A$-模子复形，
$N_1^\bullet \subset N^\bullet$ 为 $B$-模子复形，并满足
$a(M_1^\bullet) \subset N_1^\bullet$。则存在子复形
$$
M_1^\bullet \subset M_2^\bullet \subset M^\bullet
\quad\text{且}\quad
N_1^\bullet \subset N_2^\bullet \subset N^\bullet
$$
，使得 $a(M_2^\bullet) \subset N_2^\bullet$，且具有以下性质：
\begin{enumerate}
\item $\Ker(H^i(M_1^\bullet \otimes_A^\mathbf{L} B) \to H^i(N_1^\bullet))$
在 $H^i(M_2^\bullet \otimes_A^\mathbf{L} B)$ 中的像为零；
\item
% P02-E0505
$\Im(H^i(N_1^\bullet) \to H^i(N_2^\bullet))$ 包含在
$\Im(H^i(M_2^\bullet \otimes_A^\mathbf{L} B) \to H^2(N_2^\bullet))$
中；
\item $|\bigcup M_2^i \cup \bigcup N_2^i| \leq
\max(\kappa, |\bigcup M_1^i \cup \bigcup N_1^i|)$。
\end{enumerate}
\end{lemma}

\begin{proof}
置 $\kappa = \max(|A|, |B|, \aleph_0)$。记
$|M^\bullet| = |\bigcup M^i|$，其他复形也采用类似记号。于是，对
引理陈述中使用的量，有
% P02-E0506
$$
\max(\kappa, |\bigcup M_1^i \cup \bigcup N_1^i|) =
\max(\kappa, |M_1^\bullet|, |M_2^\bullet|)
$$
。下文将不再特别说明地使用此等由基数算术得到的观察。

\medskip\noindent
首先说明存在足够多的“小”子复形。对有限子集族
$E = \{E^i\}$ 与 $F = \{F^i\}$ 的每一对，其中
$E^i \subset M^i$、$i \in \mathbf{Z}$，且
$F^i \subset N^i$、$i \in \mathbf{Z}$，令
$$
M_1^\bullet \subset M_1(E, F)^\bullet \subset M^\bullet
\quad\text{且}\quad
N_1^\bullet \subset N_1(E, F)^\bullet \subset N^\bullet
$$
为分别由 $A$-模和 $B$-模构成的最小子复形，使得
$a(M_1(E, F)^\bullet) \subset N_1(E, F)^\bullet$，并且
% P02-E0507
$E^i \subset M_1(E, F)^i$ 与 $F^i \subset M_2(E, F)^i$。容易看出
% P02-E0508
$$
|M_1(E, F)^\bullet| \leq \max(\kappa, |M_1^\bullet|)
\quad\text{且}\quad
|M_2(E, F)^\bullet| \leq \max(\kappa, |M_2^\bullet|)
$$
。略去细节。显然有
$$
M^\bullet = \colim_{(E, F)} M_1(E, F)^\bullet
\quad\text{且}\quad
N^\bullet = \colim_{(E, F)} N_1(E, F)^\bullet
$$
，并且这些余极限是（逐项的）滤过余极限。

\medskip\noindent
存在由自由 $A$-模 $F_i$ 构成的解消
% P02-E0509
$\ldots \to F^{-1} \to F^0 \to B$，满足 $|F_i| \leq \kappa$
（略去细节）。$M_1^\bullet \otimes_A^\mathbf{L} B$ 的上同调模由
$\text{Tot}(M_1^\bullet \otimes_A F^\bullet)$ 计算。因此
$|H^i(M_1^\bullet \otimes_A^\mathbf{L} B)| \leq
\max(\kappa, |M_1^\bullet|)$。

\medskip\noindent
设 $i \in \mathbf{Z}$，并设
$\xi \in H^i(M_1^\bullet \otimes_A^\mathbf{L} B)$ 是一个在
$H^i(N_1^\bullet)$ 中映为零的元素。则 $\xi$ 在
$H^i(N^\bullet)$ 中映为零，从而 $\xi$ 在
$H^i(M^\bullet \otimes_A^\mathbf{L} B)$ 中也映为零。由于导出张量积
与滤过余极限交换，可以找到上述有限子集族 $E_\xi$ 与 $F_\xi$，
使得 $\xi$ 在
$H^i(M_1(E_\xi, F_\xi)^\bullet \otimes_A^\mathbf{L} B)$ 中映为零。

\medskip\noindent
设 $i \in \mathbf{Z}$，并设 $\eta \in H^i(N_1^\bullet)$。则
$\eta$ 在 $H^i(N^\bullet)$ 中的像属于
$H^i(M^\bullet \otimes_A^\mathbf{L} B) \to H^i(N^\bullet)$ 的像。
因此，与前面一样，可以找到上述有限子集族 $E_\eta$ 与 $F_\eta$，
使得 $\eta$ 映到
% P02-E0510
$H^i(N_1(E_\eta, F_\eta)$ 的一个元素，而该元素属于映射
$H^i(M_1(E_\eta, F_\eta)^\bullet \otimes_A^\mathbf{L} B) \to
H^i(N_1(E_\eta, F_\eta)$ 的像。

\medskip\noindent
现在只需定义
$$
M_2^\bullet =
\sum\nolimits_\xi M_1(E_\xi, F_\xi)^\bullet +
\sum\nolimits_\eta M_1(E_\eta, F_\eta)^\bullet
$$
，其中求和遍历前两段中的 $\xi$ 与 $\eta$，并在 $M^\bullet$ 内
取和。类似地，置
$$
N_2^\bullet =
\sum\nolimits_\xi N_1(E_\xi, F_\xi)^\bullet +
\sum\nolimits_\eta N_1(E_\eta, F_\eta)^\bullet
$$
，其中在 $N^\bullet$ 内取和。由构造，这些选取满足性质 (1) 与
(2)。性质 (3) 的界也随之成立，因为 $\xi$ 与 $\eta$ 所成集合的
基数至多为 $\max(\kappa, |M_1^\bullet|, |N_1^\bullet|)$。
\end{proof}

\begin{lemma}
\label{more-algebra-lemma-f-bounded}
设 $R$ 为环，$f \in R$，$M \in D(R)$，并设 $C$ 为
$f : M \to M$ 的锥。若当 $i < 0$ 时 $H^i(M)_f = 0$，且当
$i < -1$ 时 $H^i(C) = 0$，则当 $i < 0$ 时 $H^i(M) = 0$。
\end{lemma}

\begin{proof}
记 $M_f = M \otimes_R^\mathbf{L} R_f$，并选取区分三角
$F \to M \to M_f$。由假设，当 $i < 0$ 时，
$H^i(M_f) = H^i(M)_f = 0$。因此，只需证明当 $i < 0$ 时
$H^i(F) = 0$。注意，对所有 $i \in \mathbf{Z}$，$H^i(F)$ 都是
$f$-幂挠的。另一方面，由于 $f : M_f \to M_f$ 是同构，可见
$C$ 同构于 $f : F \to F$ 的锥（使用《导出范畴》命题
\ref{derived-proposition-9}）。现在，若 $H^i(F) \not = 0$，则
$f : H^i(F) \to H^i(F)$ 的核非零，这意味着
$H^{i - 1}(C)$ 非零。由假设，当 $i - 1 < -1$ 时不可能发生这种
情形，证明完毕。
\end{proof}

\begin{lemma}
\label{more-algebra-lemma-f-flat}
设 $R$ 为环，$f \in R$，$M \in D(R)$。假设
\begin{enumerate}
\item 当 $i > 0$ 时，$H^i(M) = 0$；
\item $M \otimes_R^\mathbf{L} R_f$ 同构于置于次数 $0$ 的平坦
$R_f$-模；
\item $M \otimes_R^\mathbf{L} R/fR$ 同构于置于次数 $0$ 的平坦
$R/fR$-模。
\end{enumerate}
则 $M$ 同构于置于次数 $0$ 的平坦 $R$-模。
\end{lemma}

\begin{proof}
设 $K$ 为 $R$-模。只需证明
$N = M \otimes_R^\mathbf{L} K$ 同构于置于次数 $0$ 的
$R$-模；见第 \ref{more-algebra-section-tor} 节。由假设 (1)，$N$ 仅在
$\leq 0$ 的次数上可能有非零上同调。因此，由引理
\ref{more-algebra-lemma-f-bounded}，只需证明当 $i < 0$ 时
$H^i(N)_f = 0$，并证明当 $i < -1$ 时 $H^i(C) = 0$；这里
$C$ 是 $f : N \to N$ 的锥。对于第一点，注意
% P02-E0504
$$
H^i(N)_f = H^i(N \otimes_R^\mathbf{L} R_f) =
H^i(M \otimes_R^\mathbf{L} K \otimes_R^\mathbf{L} R_f ) =
H^i((M \otimes_R^\mathbf{L} R_f) \otimes_{R_f}^\mathbf{L}
(K \otimes_R^\mathbf{L} R_f))
$$
。由假设 (2)，除非 $i = 0$，否则它为零；此时得到
$H^0(M)_f \otimes_R K$。对于第二点，设 $C'$ 为
$f : K \to K$ 的锥。则
$$
C = M \otimes_R^\mathbf{L} C'
$$
。复形 $C'$ 仅在次数 $0$ 与 $-1$ 上有非零上同调，分别等于
$K/fK$ 与 $K[f]$；换言之，存在区分三角
$$
(K[f])[1] \to C' \to K/fK
$$
。另一方面，对任意被 $f$ 零化的 $R$-模 $K'$，有
$$
M \otimes_R^\mathbf{L} K' =
M \otimes_R^\mathbf{L} R/fR \otimes_{R/fR}^\mathbf{L} K'
$$
。由假设 (3)，它等于置于次数 $0$ 的模
$H^0(M \otimes_R^\mathbf{L} R/fR) \otimes_{R/fR} K'$。
结合以上结论，$C$ 仅在次数 $0$ 与 $-1$ 上有非零上同调，证明完成。
\end{proof}

\section{双复形的若干技巧}
\label{more-algebra-section-tricks}

\noindent
本节继续《同调》第
\ref{homology-section-double-complexes-abelian-groups} 节的讨论。

\begin{lemma}
\label{more-algebra-lemma-vanishing-coh-prod-totalization}
设 $A_0^\bullet \to A_1^\bullet \to A_2^\bullet \to \ldots$
为阿贝尔群复形所成的复形。假设对所有 $p \geq 0$，有
$H^{-p}(A_p^\bullet) = 0$。置 $A^{p, q} = A_p^q$，并把
$A^{\bullet, \bullet}$ 视为双复形。则
$H^0(\text{Tot}_\pi(A^{\bullet, \bullet})) = 0$。
\end{lemma}

\begin{proof}
记 $f_p : A_p^\bullet \to A_{p + 1}^\bullet$ 为给定的复形映射。
回忆 $\text{Tot}_\pi(A^{\bullet, \bullet})$ 上的微分在次数 $n$
的元素上由
$$
\prod\nolimits_{p + q = n} (f^q_p + (-1)^p\text{d}^q_{A_p^\bullet})
$$
给出。设 $\xi \in H^0(\text{Tot}_\pi(A^{\bullet, \bullet}))$
为上同调类。我们将证明 $\xi$ 为零。用上循环
% P02-E0511
$x = (x_p) \in \prod A^{p, -p}$ 的类表示 $\xi$。由
$\text{d}(x) = 0$，得到
$\text{d}_{A_0^\bullet}(x_0) = 0$。由于
$H^0(A_0^\bullet) = 0$，存在 $y_{-1} \in A^{0, -1}$，满足
$\text{d}_{A_0^\bullet}(y_{-1}) = x_0$。于是
$\text{d}_{A_1^\bullet}(x_1 + f_0(y_{-1})) = 0$。由于
$H^{-1}(A_1^\bullet) = 0$，可以找到 $y_{-2} \in A^{1, -2}$，
使得 $-\text{d}_{A_1^\bullet}(y_{-2}) = x_1 + f_0(y_{-1})$。
归纳可找到 $y_{-p - 1} \in A^{p, -p - 1}$，使得
% P02-E0513
$$
(-1)^p\text{d}_{A_p^\bullet}(y_{-p - 1}) = x_p + f_{p - 1}(y_{-p})
$$
。这意味着 $\text{d}(y) = x$，其中 $y = (y_{-p - 1})$。
\end{proof}

\begin{lemma}
\label{more-algebra-lemma-prod-qis-gives-qis}
设
$$
(A_0^\bullet \to A_1^\bullet \to A_2^\bullet \to \ldots)
\longrightarrow
(B_0^\bullet \to B_1^\bullet \to B_2^\bullet \to \ldots)
$$
为两个阿贝尔群复形所成复形之间的映射。置
$A^{p, q} = A_p^q$、$B^{p, q} = B_p^q$，得到双复形。设
$\text{Tot}_\pi(A^{\bullet, \bullet})$ 与
$\text{Tot}_\pi(B^{\bullet, \bullet})$ 为这些双复形相伴的乘积
全复形。若每个 $A_p^\bullet \to B_p^\bullet$ 都是拟同构，则
$\text{Tot}_\pi(A^{\bullet, \bullet}) \to
\text{Tot}_\pi(B^{\bullet, \bullet})$ 是拟同构。
\end{lemma}

\begin{proof}
回忆 $\text{Tot}_\pi(A^{\bullet, \bullet})$ 在次数 $n$ 上由
% P02-E0512
$\prod_{p + q = n} A^{p, q} = \prod_{p + 1 = n} A^q_p$ 给出。
设 $C_p^\bullet$ 为映射 $A_p^\bullet \to B_p^\bullet$ 的锥；
见《导出范畴》第 \ref{derived-section-cones} 节。由锥构造的函子性，
得到复形所成的复形
$$
C_0^\bullet \to C_1^\bullet \to C_2^\bullet \to \ldots
$$
。可见 $\text{Tot}_\pi(C^{\bullet, \bullet})$ 在次数 $n$ 上由
$$
\prod_{p + q = n} C^{p, q} = \prod_{p + q = n} C^q_p =
\prod_{p + q = n} (B^q_p \oplus A^{q + 1}_p) =
\prod_{p + q = n} B^q_p \oplus \prod_{p + q = n} A^{q + 1}_p
$$
给出。因此，$\text{Tot}_\pi(C^{\bullet, \bullet})$ 是映射
$\text{Tot}_\pi(A^{\bullet, \bullet}) \to
\text{Tot}_\pi(B^{\bullet, \bullet})$ 的锥（
% P02-E0514
略去验证微分一致的过程）。所以只需证明：若每个
$A_p^\bullet$ 都无环，则 $\text{Tot}_\pi(A^{\bullet, \bullet})$
无环。这由引理 \ref{more-algebra-lemma-vanishing-coh-prod-totalization} 得出。
\end{proof}

\section{弱 \'{e}tale 环映射}
\label{more-algebra-section-weakly-etale}

\noindent
本节的大多数结果来自 Olivier 的论文 \cite{Olivier-AF}。另见相关论文
\cite{Ferrand-epi}。

\begin{definition}
\label{more-algebra-definition-weakly-etale}
若每个 $A$-模都在 $A$ 上平坦，则称环 $A$ {\it 绝对平坦}。若
$A \to B$ 与 $B \otimes_A B \to B$ 都平坦，则称环映射
$A \to B$ {\it 弱 \'{e}tale}或{\it 绝对平坦}。
\end{definition}

\noindent
绝对平坦环有时称为 von Neumann 正则环（这一名称常用于非交换环的
语境）。局部化是弱 \'{e}tale 环映射。\'{e}tale 环映射是弱
\'{e}tale 的。下面是一个简单而关键的性质。

\begin{lemma}
\label{more-algebra-lemma-key}
设 $A \to B$ 为环映射，并且 $B \otimes_A B \to B$ 平坦。设 $N$
为 $B$-模。若 $N$ 作为 $A$-模平坦，则 $N$ 作为 $B$-模平坦。
\end{lemma}

\begin{proof}
假设 $N$ 作为 $A$-模平坦。
% P02-E0515
则函子
$$
\text{Mod}_B \longrightarrow \text{Mod}_{B \otimes_A B},\quad
N' \mapsto N \otimes_A N'
$$
正合。由于 $B \otimes_A B \to B$ 平坦，可知函子
$$
\text{Mod}_B \longrightarrow \text{Mod}_B,\quad
N' \mapsto (N \otimes_A N') \otimes_{B \otimes_A B} B = N \otimes_B N'
$$
正合，因此 $N$ 在 $B$ 上平坦。
\end{proof}

\begin{definition}
\label{more-algebra-definition-weak-dimension}
设 $A$ 为环，$d \geq 0$ 为整数。若每个 $A$-模的 Tor 维数
$\leq d$，则称 $A$ 的{\it 弱维数 $\leq d$}。
\end{definition}

\begin{lemma}
\label{more-algebra-lemma-weak-dimension-goes-up}
设 $A \to B$ 为弱 \'{e}tale 环映射。若 $A$ 的弱维数至多为 $d$，
则 $B$ 也如此。
\end{lemma}

\begin{proof}
设 $N$ 为 $B$-模。若 $d = 0$，则 $N$ 作为 $A$-模平坦，从而由
引理 \ref{more-algebra-lemma-key}，它作为 $B$-模平坦。假设 $d > 0$。选取
自由 $B$-模解消 $F_\bullet \to N$。我们的假设蕴含
$K = \Im(F_d \to F_{d - 1})$ 是 $A$-平坦的；见引理
\ref{more-algebra-lemma-last-one-flat}。故由引理 \ref{more-algebra-lemma-key}，它是
$B$-平坦的。因此
$0 \to K \to F_{d - 1} \to \ldots \to F_0 \to N \to 0$
是长度为 $d$ 的平坦解消，可见 $N$ 的 Tor 维数至多为 $d$。
\end{proof}

\begin{lemma}
\label{more-algebra-lemma-absolutely-flat}
设 $A$ 为环。下列条件等价：
\begin{enumerate}
\item $A$ 的弱维数 $\leq 0$；
\item $A$ 绝对平坦；
\item $A$ 既约，且每个素理想都是极大理想；
\item $A$ 的每个局部环都是域。
\end{enumerate}
\end{lemma}

\begin{proof}
(1) 与 (2) 的等价性由定义立即得到。

\medskip\noindent
证明 (1) $\Rightarrow$ (3)。
假设 $A$ 绝对平坦。则 $A$ 的每个理想都是纯理想；见《代数》定义
\ref{algebra-definition-pure-ideal}。因此，由《代数》引理
\ref{algebra-lemma-finitely-generated-pure-ideal}，每个有限生成理想
都由一个幂等元生成。若 $f \in A$，则对某个幂等元 $e \in A$，
有 $(f) = (e)$，且 $D(f) = D(e)$ 既开又闭（《代数》引理
\ref{algebra-lemma-idempotent-spec}）。由《代数》引理
\ref{algebra-lemma-ring-with-only-minimal-primes}，这蕴含 $A$ 的每个
素理想都是极大理想。此外，若 $f$ 幂零，则 $e = 0$，故 $f = 0$。
因此 $A$ 既约。

\medskip\noindent
证明 (3) $\Rightarrow$ (4)。若 $A$ 既约且每个素理想都是极大理想，
则由《代数》引理 \ref{algebra-lemma-minimal-prime-reduced-ring}，
每个局部环都是域。

\medskip\noindent
证明 (4) $\Rightarrow$ (2)。若 $A$ 的每个局部环都是域，则由
《代数》引理 \ref{algebra-lemma-flat-localization}，每个 $A$-模
都平坦。
\end{proof}

\begin{lemma}
\label{more-algebra-lemma-product-fields-absolutely-flat}
域的乘积是绝对平坦环。
\end{lemma}

\begin{proof}
设 $K_i$ 为一族域。若 $f = (f_i) \in \prod K_i$，则 $f$ 生成的理想
与幂等元 $e = (e_i)$ 生成的理想相同，其中根据 $f_i$ 是否为 $0$，
$e_i = 0, 1$。因此 $D(f) = D(e)$ 既开又闭；结论由引理
\ref{more-algebra-lemma-absolutely-flat} 与《代数》引理
\ref{algebra-lemma-ring-with-only-minimal-primes} 得出。
\end{proof}

\begin{lemma}
\label{more-algebra-lemma-base-change-weakly-etale}
设 $A \to B$ 与 $A \to A'$ 为环映射，设
$B' = B \otimes_A A'$ 为 $B$ 的基变换。
\begin{enumerate}
\item 若 $B \otimes_A B \to B$ 平坦，则
$B' \otimes_{A'} B' \to B'$ 平坦。
\item 若 $A \to B$ 弱 \'{e}tale，则 $A' \to B'$ 弱 \'{e}tale。
\end{enumerate}
\end{lemma}

\begin{proof}
假设 $B \otimes_A B \to B$ 平坦。环映射
$B' \otimes_{A'} B' \to B'$ 是 $B \otimes_A B \to B$ 沿
$A \to A'$ 的基变换。因此，由《代数》引理
\ref{algebra-lemma-flat-base-change}，它是平坦的。这证明了 (1)。
(2) 由 (1) 以及刚才使用的平坦环映射的基变换仍平坦这一事实得出。
\end{proof}

\begin{lemma}
\label{more-algebra-lemma-absolutely-flat-over-absolutely-flat}
设 $A \to B$ 为环映射，并且 $B \otimes_A B \to B$ 平坦。
\begin{enumerate}
\item 若 $A$ 是绝对平坦环，则 $B$ 也如此。
\item 若 $A$ 既约且 $A \to B$ 弱 \'{e}tale，则 $B$ 既约。
\end{enumerate}
\end{lemma}

\begin{proof}
(1) 立即由引理 \ref{more-algebra-lemma-key} 与定义得出。若 $A$ 既约，则存在
单射
$$
A \to A' = \prod_{\mathfrak p \subset A\text{ minimal}} A_\mathfrak p
$$
，把 $A$ 嵌入一个绝对平坦环中（《代数》引理
\ref{algebra-lemma-reduced-ring-sub-product-fields} 与引理
\ref{more-algebra-lemma-product-fields-absolutely-flat}）。若 $A \to B$ 平坦，
则诱导映射 $B \to B' = B \otimes_A A'$ 也是单射。由引理
\ref{more-algebra-lemma-base-change-weakly-etale}，环映射 $A' \to B'$ 弱
\'{e}tale。由 (1)，$B'$ 绝对平坦。再由引理
\ref{more-algebra-lemma-absolutely-flat}，环 $B'$ 既约。因此 $B$ 既约。
\end{proof}

\begin{lemma}
\label{more-algebra-lemma-composition-weakly-etale}
设 $A \to B$ 与 $B \to C$ 为环映射。
\begin{enumerate}
\item 若 $B \otimes_A B \to B$ 与 $C \otimes_B C \to C$ 平坦，则
$C \otimes_A C \to C$ 平坦。
\item 若 $A \to B$ 与 $B \to C$ 弱 \'{e}tale，则 $A \to C$
弱 \'{e}tale。
\end{enumerate}
\end{lemma}

\begin{proof}
(1) 由乘法映射的分解
$$
C \otimes_A C \longrightarrow C \otimes_B C \longrightarrow C
$$
、等式
$$
C \otimes_B C = (C \otimes_A C) \otimes_{B \otimes_A B} B,
$$
、平坦映射的基变换仍平坦以及平坦环映射的复合仍平坦这些事实得出。
见《代数》引理 \ref{algebra-lemma-flat-base-change} 与
\ref{algebra-lemma-composition-flat}。(2) 由 (1) 以及刚才使用的
平坦环映射的复合仍平坦这一事实得出。
\end{proof}

\begin{lemma}
\label{more-algebra-lemma-go-down}
设 $A \to B \to C$ 为环映射。
\begin{enumerate}
\item 若 $B \to C$ 忠实平坦，且 $C \otimes_A C \to C$ 平坦，则
$B \otimes_A B \to B$ 平坦。
\item 若 $B \to C$ 忠实平坦，且 $A \to C$ 弱 \'{e}tale，则
$A \to B$ 弱 \'{e}tale。
\end{enumerate}
\end{lemma}

\begin{proof}
假设 $B \to C$ 忠实平坦且 $C \otimes_A C \to C$ 平坦。考虑交换图
% P02-E0522
$$
\xymatrix{
C \otimes_A C \ar[r] & C \\
B \otimes_A B \ar[r] \ar[u] & B \ar[u]
}
$$
。竖直箭头平坦，上方水平箭头平坦。因此 $C$ 作为
$B \otimes_A B$-模平坦。映射 $B \to C$ 忠实平坦，且
$C = B \otimes_B C$。故由《代数》引理
\ref{algebra-lemma-flatness-descends-more-general}，$B$ 作为
$B \otimes_A B$-模平坦。这证明了 (1)。(2) 由 (1) 以及如下事实
得出：若 $A \to C$ 平坦且 $B \to C$ 忠实平坦，则 $A \to B$
平坦（《代数》引理
\ref{algebra-lemma-flatness-descends-more-general}）。
\end{proof}

\begin{lemma}
\label{more-algebra-lemma-weakly-etale-permanence}
设 $A$ 为环。设 $B \to C$ 为弱 \'{e}tale $A$-代数之间的
$A$-代数映射。则 $B \to C$ 弱 \'{e}tale。
\end{lemma}

\begin{proof}
由引理 \ref{more-algebra-lemma-key}，环映射 $B \to C$ 平坦。环映射
$C \otimes_A C \to C \otimes_B C$ 满射，因而是上满射。于是，
% P02-E0516
由引理 \ref{more-algebra-lemma-key}，由于 $C$ 在 $C \otimes_A C$ 上平坦，
$C$ 在 $C \otimes_B C$ 上也平坦。
\end{proof}

\begin{lemma}
\label{more-algebra-lemma-formally-unramified}
设 $A \to B$ 为环映射，并且 $B \otimes_A B \to B$ 平坦。则
$\Omega_{B/A} = 0$，即 $B$ 在 $A$ 上形式非分歧。
\end{lemma}

\begin{proof}
设 $I \subset B \otimes_A B$ 为平坦满射
$B \otimes_A B \to B$ 的核。则 $I$ 是纯理想（《代数》定义
\ref{algebra-definition-pure-ideal}），故 $I^2 = I$（《代数》引理
\ref{algebra-lemma-pure}）。由于 $\Omega_{B/A} = I/I^2$
（《代数》引理 \ref{algebra-lemma-differentials-diagonal}），得到
所述消失。由《代数》引理
\ref{algebra-lemma-characterize-formally-unramified}，这意味着 $B$
在 $A$ 上形式非分歧。
\end{proof}

\begin{lemma}
\label{more-algebra-lemma-weakly-etale-finite-type}
设 $A \to B$ 为环映射，并且 $B \otimes_A B \to B$ 平坦。
\begin{enumerate}
\item 若 $A \to B$ 为有限型，则 $A \to B$ 非分歧。
\item 若 $A \to B$ 有限表现且平坦，则 $A \to B$ 是 \'{e}tale 的。
\end{enumerate}
特别地，有限表现的弱 \'{e}tale 环映射是 \'{e}tale 的。
\end{lemma}

\begin{proof}
(1) 由引理 \ref{more-algebra-lemma-formally-unramified} 与《代数》定义
\ref{algebra-definition-unramified} 得出。(2) 由 (1) 与《代数》引理
\ref{algebra-lemma-etale-flat-unramified-finite-presentation} 得出。
\end{proof}

\begin{lemma}
\label{more-algebra-lemma-when-weakly-etale}
设 $A \to B$ 为环映射。在以下每种情形中，$A \to B$ 都弱
\'{e}tale：
\begin{enumerate}
\item $B = S^{-1}A$ 是 $A$ 的局部化；
\item $A \to B$ 是 \'{e}tale 的；
\item $B$ 是弱 \'{e}tale $A$-代数的滤过余极限。
\end{enumerate}
\end{lemma}

\begin{proof}
\'{e}tale 环映射平坦，而且作为 \'{e}tale $A$-代数之间的映射，
$B \otimes_A B \to B$ 也是 \'{e}tale 的（《代数》引理
\ref{algebra-lemma-map-between-etale}）。这证明了 (2)。

\medskip\noindent
设 $B_i$ 为弱 \'{e}tale $A$-代数的有向系统。由《代数》引理
\ref{algebra-lemma-colimit-flat}，$B = \colim B_i$ 在 $A$ 上平坦。
注意，由引理 \ref{more-algebra-lemma-weakly-etale-permanence}，过渡映射
$B_i \to B_{i'}$ 平坦。因此，对每个 $i$，$B$ 在 $B_i$ 上平坦；
由《代数》引理 \ref{algebra-lemma-composition-flat}，可见 $B$
在 $B_i \otimes_A B_i$ 上平坦。故由《代数》引理
\ref{algebra-lemma-colimit-rings-flat}，$B$ 在
$B \otimes_A B = \colim B_i \otimes_A B_i$ 上平坦。

\medskip\noindent
(1) 可以直接证明，也可由 (2) 与 (3) 合并得到。
\end{proof}

\begin{lemma}
\label{more-algebra-lemma-absolutely-flat-fields}
设 $L/K$ 为域扩张。若 $L \otimes_K L \to L$ 平坦，则 $L$ 是
$K$ 的代数可分扩张。
\end{lemma}

\begin{proof}
由引理 \ref{more-algebra-lemma-go-down} 可见，
% P02-E0517
对任意中间域 $K \subset L' \subset L$，映射
$L' \otimes_K L' \to L'$ 都平坦。因此可假设 $L$ 是 $K$ 的有限生成
域扩张。在这种情形下，$L/K$ 形式非分歧（引理
\ref{more-algebra-lemma-formally-unramified}）这一事实蕴含 $L/K$ 有限可分；
见《代数》引理
\ref{algebra-lemma-characterize-separable-algebraic-field-extensions}。
\end{proof}

\begin{lemma}
\label{more-algebra-lemma-absolutely-flat-over-field}
设 $B$ 为域 $K$ 上的代数。下列条件等价：
\begin{enumerate}
\item $B \otimes_K B \to B$ 平坦；
\item $K \to B$ 弱 \'{e}tale；
\item $B$ 是 \'{e}tale $K$-代数的滤过余极限。
\end{enumerate}
此外，$B$ 的每个有限生成 $K$-子代数都在 $K$ 上 \'{e}tale。
\end{lemma}

\begin{proof}
(1) 与 (2) 等价，因为每个 $K$-代数都在 $K$ 上平坦。由引理
\ref{more-algebra-lemma-when-weakly-etale}，(3) 蕴含 (1) 与 (2)。
% P02-E0518

\medskip\noindent
假设 (1) 与 (2) 成立。我们将证明 (3) 以及引理中的有限性陈述。
% P02-E0519
域是绝对平坦环，故由引理
\ref{more-algebra-lemma-absolutely-flat-over-absolutely-flat}，$B$ 是绝对平坦环。
因此 $B$ 既约，且每个局部环都是域；见引理
\ref{more-algebra-lemma-absolutely-flat}。

\medskip\noindent
设 $\mathfrak q \subset B$ 为素理想。环映射
$B \to B_\mathfrak q$ 弱 \'{e}tale，故
$B_\mathfrak q$ 在 $K$ 上弱 \'{e}tale（引理
\ref{more-algebra-lemma-composition-weakly-etale}）。因此，由引理
\ref{more-algebra-lemma-absolutely-flat-fields}，$B_\mathfrak q$ 是 $K$ 的
可分代数扩张。

\medskip\noindent
设 $K \subset A \subset B$ 为有限生成 $K$-子代数。
% P02-E0520
我们将证明 $A$ 在 $K$ 上 \'{e}tale，这将完成引理的证明。
于是，每个极小素理想 $\mathfrak p \subset A$ 都是 $B$ 的
某个素理想 $\mathfrak q$ 的像；见《代数》引理
\ref{algebra-lemma-injective-minimal-primes-in-image}。因此，作为
$B_\mathfrak q = \kappa(\mathfrak q)$ 的子域，
$\kappa(\mathfrak p)$ 是 $K$ 的可分代数扩张。故
$\Spec(A)$ 的每个泛点都是闭点（《代数》引理
\ref{algebra-lemma-finite-residue-extension-closed}），从而
$\dim(A) = 0$。于是 $A$ 是其各局部环的乘积，例如可由《代数》命题
\ref{algebra-proposition-dimension-zero-ring} 得到。此外，由于 $A$
既约，所有局部环都等于各自的剩余域，而这些剩余域都是 $K$ 的有限
可分扩张。由《代数》引理 \ref{algebra-lemma-etale-over-field}，
这意味着 $A$ 在 $K$ 上 \'{e}tale，并完成证明。
\end{proof}

\begin{lemma}
\label{more-algebra-lemma-weakly-etale-residue-field-extensions}
设 $A \to B$ 为环映射。若 $A \to B$ 弱 \'{e}tale，则
$A \to B$ 诱导可分代数的剩余域扩张。
\end{lemma}

\begin{proof}
设 $\mathfrak p$ 为 $A$ 的素理想。由引理
\ref{more-algebra-lemma-base-change-weakly-etale}，
$\kappa(\mathfrak p) \to B \otimes_A \kappa(\mathfrak p)$ 弱
\'{e}tale。因此，由引理 \ref{more-algebra-lemma-absolutely-flat-over-field}，
$B \otimes_A \kappa(\mathfrak p)$ 是 \'{e}tale
$\kappa(\mathfrak p)$-代数的滤过余极限。故对位于
$\mathfrak p$ 上方的 $\mathfrak q \subset B$，由《代数》引理
\ref{algebra-lemma-etale-over-field}，扩张
$\kappa(\mathfrak q)/\kappa(\mathfrak p)$ 是有限可分扩张的滤过余极限。
\end{proof}

\begin{lemma}
\label{more-algebra-lemma-weak-dimension-at-most-1}
设 $A$ 为环。下列条件等价：
\begin{enumerate}
\item $A$ 的弱维数 $\leq 1$；
\item $A$ 的每个理想都平坦；
\item $A$ 的每个有限生成理想都平坦；
\item 平坦 $A$-模的每个子模都平坦；
\item $A$ 的每个局部环都是赋值环。
\end{enumerate}
\end{lemma}

\begin{proof}
若 $A$ 的弱维数 $\leq 1$，则解消
$0 \to I \to A \to A/I \to 0$ 与引理 \ref{more-algebra-lemma-last-one-flat}
表明每个理想 $I$ 都平坦。因此 (1) $\Rightarrow$ (2)。

\medskip\noindent
假设 (4)。设 $M$ 为 $A$-模。选取满射 $F \to M$，其中 $F$ 是自由
$A$-模。由假设，$\Ker(F \to M)$ 平坦；由引理
\ref{more-algebra-lemma-tor-dimension}，可见 $M$ 的 Tor 维数 $\leq 1$。
因此 (4) $\Rightarrow$ (1)。

\medskip\noindent
每个理想都是其中各有限生成理想的并。因此，由《代数》引理
\ref{algebra-lemma-colimit-flat}，(3) 蕴含 (2)。从而
(3) $\Leftrightarrow$ (2)。

\medskip\noindent
假设 (2)。设 $N \subset M$，其中 $M$ 为平坦 $A$-模。我们将证明
$N$ 平坦。由《代数》定理 \ref{algebra-theorem-lazard}，可写
$M = \colim M_i$，其中每个 $M_i$ 都是有限自由模。令
$N_i \subset M_i$ 为 $N$ 的逆像，可见 $N = \colim N_i$。由
% P02-E0521
《代数》引理 \ref{algebra-lemma-colimit-flat}，只需证明 $N_i$
平坦，于是归约到 $M = R^{\oplus n}$ 的情形。
% P02-E0523
在这种情形下，模 $N$ 具有由子模 $R^{\oplus j} \cap N$ 构成的
有限滤过，其逐次商都是理想。由 (2)，这些理想平坦；故由《代数》
引理 \ref{algebra-lemma-flat-ses}，$N$ 平坦。因此
(2) $\Rightarrow$ (4)。

\medskip\noindent
假设 $A$ 满足 (1)，并设 $\mathfrak p \subset A$ 为素理想。由引理
\ref{more-algebra-lemma-when-weakly-etale} 与
\ref{more-algebra-lemma-weak-dimension-goes-up}，可见 $A_\mathfrak p$ 满足 (1)。
下面证明：若 $A$ 是满足 (3) 的局部环，则 $A$ 是赋值环。设
$f \in \mathfrak m$ 为非零元素。则 $(f)$ 是由一个元素生成的
非零平坦模。故由《代数》引理
\ref{algebra-lemma-finite-flat-local}，它是自由 $A$-模。于是 $f$
是非零因子，且 $A$ 是整环。若 $I \subset A$ 是有限生成理想，
则同理可见 $I$ 是有限自由 $A$-模，因而（考察秩）是秩 $1$ 的自由模，
且 $I$ 是主理想。由《代数》引理
\ref{algebra-lemma-characterize-valuation-ring}，$A$ 是赋值环。
因此 (1) $\Rightarrow$ (5)。

\medskip\noindent
假设 (5)。设 $I \subset A$ 为有限生成理想。则
$I_\mathfrak p \subset A_\mathfrak p$ 是赋值环中的有限生成理想，
故为主理想（《代数》引理
\ref{algebra-lemma-characterize-valuation-ring}），从而平坦。因此，
由《代数》引理 \ref{algebra-lemma-flat-localization}，$I$ 平坦。
所以 (5) $\Rightarrow$ (3)。引理得证。
\end{proof}

\begin{lemma}
\label{more-algebra-lemma-product-weak-dimension-at-most-1}
设 $J$ 为集合。对每个 $j \in J$，设 $A_j$ 为分式域为 $K_j$ 的
赋值环。置 $A = \prod A_j$ 与 $K = \prod K_j$。则 $A$ 的弱维数
至多为 $1$，并且 $A \to K$ 是局部化。
\end{lemma}

\begin{proof}
设 $I \subset A$ 为有限生成理想。由引理
\ref{more-algebra-lemma-weak-dimension-at-most-1}，只需证明 $I$ 是平坦
$A$-模。设 $I_j \subset A_j$ 为 $I$ 的像。注意
$I_j = I \otimes_A A_j$，故由《代数》命题
\ref{algebra-proposition-fg-tensor}，$I \to \prod I_j$ 满射。因此
$I = \prod I_j$。由于 $A_j$ 是赋值环，理想 $I_j$ 由单个元素生成
（《代数》引理 \ref{algebra-lemma-characterize-valuation-ring}）。
设 $I_j = (f_j)$。则 $I$ 由元素 $f = (f_j)$ 生成。设 $e \in A$
为如下幂等元：根据 $f_j$ 是否为 $0$，它在 $A_j$ 中的分量为
$0$ 或 $1$。于是对某个非零因子 $g \in A$，有 $f = g e$：
取 $g = (g_j)$，其中若 $f_j = 0$，则 $g_j = 1$；否则
% P02-E0524
$g_j = f_j$。因此，作为模有 $I \cong (e)$。由于 $(e)$ 是 $A$
的直和项，故 $I$ 平坦。最后一个陈述成立，因为
$K = S^{-1}A$，其中 $S = \prod (A_j \setminus \{0\})$。
\end{proof}

\begin{lemma}
\label{more-algebra-lemma-product-found-valuation-rings}
设 $A$ 为分式域为 $K$ 的正规整环。存在环的笛卡尔图
$$
\xymatrix{
A \ar[d] \ar[r] & K \ar[d] \\
V \ar[r] & L
}
$$
，其中 $V$ 的弱维数至多为 $1$，且 $V \to L$ 是平坦、单射的
环满态射。
\end{lemma}

\begin{proof}
对每个 $x \in K$、$x \not \in A$，选取《代数》引理
% P02-E0525
\ref{algebra-lemma-find-valuation-rings} 中的 $V_x \subset K$。
置 $V = \prod_{x \in K \setminus A} V_x$ 与
$L = \prod_{x \in K \setminus A} K$。由引理
\ref{more-algebra-lemma-product-weak-dimension-at-most-1}，环 $V$ 的弱维数至多为
$1$；同一引理也表明 $V \to L$ 是局部化。局部化平坦且为满态射；
见《代数》引理 \ref{algebra-lemma-flat-localization} 与
\ref{algebra-lemma-epimorphism-local}。
\end{proof}

\begin{lemma}
\label{more-algebra-lemma-weak-dimension-at-most-1-integrally-closed}
设 $A$ 为弱维数至多为 $1$ 的环。若 $A \to B$ 是平坦、单射的
环满态射，则 $A$ 在 $B$ 中整闭。
\end{lemma}

\begin{proof}
设 $x \in B$ 在 $A$ 上整。令 $A' = A[x] \subset B$。由《代数》引理
\ref{algebra-lemma-characterize-finite-in-terms-of-integral}，
$A'$ 是 $A$ 的有限环扩张。由《代数》引理
\ref{algebra-lemma-finite-epimorphism-surjective}，为证明
$A = A'$，只需证明 $A \to A'$ 是满态射。由关于 $A$ 的假设、
$B$ 在 $A$ 上平坦这一事实以及引理
\ref{more-algebra-lemma-weak-dimension-at-most-1}，$A'$ 在 $A$ 上平坦。因此复合
$$
A' \otimes_A A' \to B \otimes_A A' \to B \otimes_A B \to B
$$
是单射，即 $A' \otimes_A A' \cong A'$，引理得证。
\end{proof}

\begin{lemma}
\label{more-algebra-lemma-normality-goes-up}
设 $A$ 为分式域为 $K$ 的正规整环。设 $A \to B$ 弱
\'{e}tale。则 $B$ 在 $B \otimes_A K$ 中整闭。
\end{lemma}

\begin{proof}
选取引理 \ref{more-algebra-lemma-product-found-valuation-rings} 中的图。由于
$A \to B$ 平坦，基变换给出环的笛卡尔图
$$
\xymatrix{
B \ar[d] \ar[r] & B \otimes_A K \ar[d] \\
B \otimes_A V \ar[r] & B \otimes_A L
}
$$
。注意，$V \to B \otimes_A V$ 弱 \'{e}tale（引理
\ref{more-algebra-lemma-base-change-weakly-etale}），故由引理
\ref{more-algebra-lemma-weak-dimension-goes-up}，$B \otimes_A V$ 的弱维数至多
为 $1$。又注意，$B \otimes_A V \to B \otimes_A L$ 是某个
平坦、单射的环满态射的平坦基变换，因而仍是平坦、单射的环满态射
（《代数》引理 \ref{algebra-lemma-flat-base-change} 与
\ref{algebra-lemma-base-change-epimorphism}）。由引理
\ref{more-algebra-lemma-weak-dimension-at-most-1-integrally-closed}，
$B \otimes_A V$ 在 $B \otimes_A L$ 中整闭。由该图的笛卡尔性质，
可知 $B$ 在 $B \otimes_A K$ 中整闭。
\end{proof}

\begin{lemma}
\label{more-algebra-lemma-integral-over-henselian}
设 $A \to B$ 为环同态。假设
\begin{enumerate}
\item $A$ 是 Hensel 局部环；
\item $A \to B$ 是整的；
\item $B$ 是整环。
\end{enumerate}
则 $B$ 是 Hensel 局部环，且 $A \to B$ 是局部同态。若 $A$
严格 Hensel，则 $B$ 是严格 Hensel 局部环，而且剩余域扩张
$\kappa(\mathfrak m_B)/\kappa(\mathfrak m_A)$ 纯不可分。
\end{lemma}

\begin{proof}
把 $B$ 写成有限 $A$-子代数的滤过余极限
$B = \colim B_i$。
% P02-E0520
若对每个 $B_i$ 都证明这些结论，则 $B$ 的结论随之成立；见《代数》
引理 \ref{algebra-lemma-colimit-henselian}。若 $A \to B$ 有限，则
由《代数》引理 \ref{algebra-lemma-finite-over-henselian}，$B$
是若干 Hensel 局部环的乘积。由于 $B$ 是整环，可见 $B$ 是局部环。
由 $A \to B$ 的上升定理（《代数》引理
\ref{algebra-lemma-integral-going-up}），$B$ 的极大理想位于 $A$
的极大理想上方。若 $A$ 严格 Hensel，则代数域扩张
$\kappa(\mathfrak m_B)/\kappa(\mathfrak m_A)$ 必为纯不可分。
于是当然 $\kappa(\mathfrak m_B)$ 可分代数闭，且 $B$ 严格 Hensel。
\end{proof}

\begin{theorem}[Olivier]
\label{more-algebra-theorem-olivier}
设 $A \to B$ 为局部环之间的局部同态。若 $A$ 严格 Hensel，且
$A \to B$ 弱 \'{e}tale，则 $A = B$。
\end{theorem}

\begin{proof}
我们将证明：对所有 $\mathfrak p \subset A$，存在唯一的
$\mathfrak q \subset B$ 位于 $\mathfrak p$ 上方，并且
$\kappa(\mathfrak p) = \kappa(\mathfrak q)$。这蕴含
$B \otimes_A B \to B$ 在谱上双射，而且该映射满射且平坦。因此，
例如由《代数》引理
\ref{algebra-lemma-pure-open-closed-specializations} 中对纯理想的
描述，它是同构。故 $A \to B$ 是忠实平坦的环满态射。由《代数》
引理 \ref{algebra-lemma-faithfully-flat-epimorphism}，得到
$A = B$。

\medskip\noindent
注意，由引理 \ref{more-algebra-lemma-base-change-weakly-etale} 与
\ref{more-algebra-lemma-absolutely-flat-over-field}，纤维环
$B \otimes_A \kappa(\mathfrak p)$ 是
$\kappa(\mathfrak p)$ 的 \'{e}tale 扩张的余极限。因此，若有不止
一个素理想位于 $\mathfrak p$ 上方，或对某个 $\mathfrak q$ 有
$\kappa(\mathfrak p) \not = \kappa(\mathfrak q)$，则对某个（可分）
代数域扩张 $L/\kappa(\mathfrak p)$，$B \otimes_A L$ 具有非平凡幂等元。

\medskip\noindent
设 $L/\kappa(\mathfrak p)$ 为代数域扩张。令 $A' \subset L$ 为
$A/\mathfrak p$ 在 $L$ 中的整闭包。由引理
\ref{more-algebra-lemma-integral-over-henselian}，$A'$ 是严格 Hensel 局部环，
其剩余域是 $A$ 的剩余域的纯不可分扩张。故由《代数》引理
\ref{algebra-lemma-local-tensor-with-integral}，
$B \otimes_A A'$ 是局部环。另一方面，由引理
\ref{more-algebra-lemma-normality-goes-up}，$B \otimes_A A'$ 在
$B \otimes_A L$ 中整闭。由于 $B \otimes_A A'$ 是局部环，可知
$B \otimes_A L$ 没有非平凡幂等元，这正是所要证明的。
\end{proof}

\section{域上的弱 \'{e}tale 代数}
\label{more-algebra-section-weakly-etale-over-field}

\noindent
若 $K$ 为域，则代数 $B$ 在 $K$ 上弱 \'{e}tale，当且仅当它是
\'{e}tale $K$-代数的滤过余极限。这就是引理
\ref{more-algebra-lemma-absolutely-flat-over-field}。

\begin{lemma}
\label{more-algebra-lemma-class-weakly-etale-over-field}
设 $K$ 为域。若 $B$ 在 $K$ 上弱 \'{e}tale，则
\begin{enumerate}
\item $B$ 既约；
\item $B$ 在 $K$ 上整；
\item $B$ 的任意有限生成 $K$-子代数都是 $K$ 的有限可分扩张的
有限乘积；
\item $B$ 是域，当且仅当 $B$ 没有非平凡幂等元；在这种情形下，
它是 $K$ 的可分代数扩张；
\item
% P02-E0526
$B$ 的任意 $K$-子代数或其商代数都在 $K$ 上弱 \'{e}tale；
\item 若 $B'$ 在 $K$ 上弱 \'{e}tale，则
$B \otimes_K B'$ 在 $K$ 上弱 \'{e}tale。
\end{enumerate}
\end{lemma}

\begin{proof}
(1) 由引理 \ref{more-algebra-lemma-absolutely-flat-over-absolutely-flat} 得出，当然
它也由 (3) 得出。(3) 由引理
\ref{more-algebra-lemma-absolutely-flat-over-field} 以及如下事实得出：
\'{e}tale $K$-代数是 $K$ 的有限可分扩张的有限乘积；见《代数》引理
\ref{algebra-lemma-etale-over-field}。(3) 蕴含 (2)。(4) 由 (3)
得出，因为域的乘积是域，当且仅当它没有非平凡幂等元。

\medskip\noindent
若 $S \subset B$ 是子代数，则它是其各有限生成子代数的滤过
余极限；由上文，这些子代数都在 $K$ 上 \'{e}tale，故由引理
\ref{more-algebra-lemma-absolutely-flat-over-field}，$S$ 在 $K$ 上弱
\'{e}tale。若 $B \to Q$ 是商代数，则 $Q$ 是有限可分扩张
$L_i/K$ 的有限乘积 $\prod_{i \in I} L_i$ 的 $K$-代数商的滤过
余极限。这种商具有形式 $\prod_{i \in J} L_i$，其中
$J \subset I$；故以同一论证，商代数的结论也成立。

\medskip\noindent
关于张量积的陈述可以类似证明，也可合并引理
\ref{more-algebra-lemma-base-change-weakly-etale} 与
\ref{more-algebra-lemma-composition-weakly-etale} 得到。
\end{proof}

\begin{lemma}
\label{more-algebra-lemma-max-weakly-etale-subalgebra}
设 $K$ 为域，$A$ 为 $K$-代数。存在极大的弱 \'{e}tale
$K$-子代数 $B_{max} \subset A$。
\end{lemma}

\begin{proof}
设 $B_1, B_2 \subset A$ 为弱 \'{e}tale $K$-子代数。则
$B_1 \otimes_K B_2$ 在 $K$ 上弱 \'{e}tale，而且
$B_1 \otimes_K B_2 \to A$ 的像也如此（引理
\ref{more-algebra-lemma-class-weakly-etale-over-field}）。因此，弱 \'{e}tale
$K$-子代数 $B \subset A$ 所成集合 $\mathcal{B}$ 是有向的；由引理
\ref{more-algebra-lemma-when-weakly-etale}，余极限
$B_{max} = \colim_{B \in \mathcal{B}} B$ 是弱 \'{e}tale $K$-代数。
因此，$B_{max} \to A$ 的像在 $K$ 上弱 \'{e}tale（引用前一个引理）。
于是该像属于 $\mathcal{B}$，故 $\mathcal{B}$ 有极大元（而该像
与 $B_{max}$ 相同）。
\end{proof}

\begin{lemma}
\label{more-algebra-lemma-properties-of-max-weakly-etale-subalgebra}
设 $K$ 为域。对 $K$-代数 $A$，以 $B_{max}(A)$ 表示引理
% P02-E0527
\ref{more-algebra-lemma-max-weakly-etale-subalgebra} 中 $A$ 的极大弱
\'{e}tale $K$-子代数。则
\begin{enumerate}
\item 任意 $K$-代数映射 $A' \to A$ 都诱导 $K$-代数映射
$B_{max}(A') \to B_{max}(A)$；
\item 若 $A' \subset A$，则
$B_{max}(A') = B_{max}(A) \cap A'$；
\item 若 $A = \colim A_i$ 为滤过余极限，则
$B_{max}(A) = \colim B_{max}(A_i)$；
\item 映射 $B_{max}(A) \to B_{max}(A_{red})$ 是同构；
\item $B_{max}(A_1 \times \ldots \times A_n) =
B_{max}(A_1) \times \ldots \times B_{max}(A_n)$；
\item 若 $A$ 没有非平凡幂等元，则 $B_{max}(A)$ 是域，并且是
$K$ 的可分代数扩张；
\item
% P02-E0528
在此添加更多内容。
\end{enumerate}
\end{lemma}

\begin{proof}
证明 (1)。这是因为，由引理
\ref{more-algebra-lemma-class-weakly-etale-over-field}，
$B_{max}(A') \to A$ 的像在 $K$ 上弱 \'{e}tale。

\medskip\noindent
证明 (2)。由 (1)，有 $B_{max}(A') \subset B_{max}(A)$。反之，由引理
\ref{more-algebra-lemma-class-weakly-etale-over-field}，
$B_{max}(A) \cap A'$ 是弱 \'{e}tale $K$-代数，因而包含在
$B_{max}(A')$ 中。

\medskip\noindent
证明 (3)。由 (1)，存在映射 $\colim B_{max}(A_i) \to A$；由于系统
滤过且 $B_{max}(A_i) \subset A_i$，该映射单射。由引理
\ref{more-algebra-lemma-when-weakly-etale}，余极限
$\colim B_{max}(A_i)$ 在 $K$ 上弱 \'{e}tale。因此得到单射
$\colim B_{max}(A_i) \to B_{max}(A)$。设
$a \in B_{max}(A)$。则 $a$ 生成有限表现 $K$-子代数
$B \subset B_{max}(A)$。由《代数》引理
\ref{algebra-lemma-characterize-finite-presentation}，存在 $i$ 以及
$K$-代数映射 $f : B \to A_i$，它提升给定映射 $B \to A$。由引理
\ref{more-algebra-lemma-class-weakly-etale-over-field}，$B$ 弱 \'{e}tale，
故 $f(B) \subset B_{max}(A_i)$；于是 $a$ 属于
$\colim B_{max}(A_i) \to B_{max}(A)$ 的像。

\medskip\noindent
证明 (4)。把 $B_{max}(A_{red}) = \colim B_i$ 写成
\'{e}tale $K$-代数的滤过余极限（引理
\ref{more-algebra-lemma-absolutely-flat-over-field}）。由《代数》引理
\ref{algebra-lemma-smooth-strong-lift}，对每个 $i$，存在
$K$-代数映射 $f_i : B_i \to A$，提升给定映射
$B_i \to A_{red}$。因此，
% P02-E0529
典范映射 $B_{max}(A_{red}) \to B_{max}(A)$ 满射。其核由幂零元组成，
因而为零，因为
% P02-E0530
$B_{max}(A_{red})$ 既约（引理
\ref{more-algebra-lemma-class-weakly-etale-over-field}）。

\medskip\noindent
证明 (5)。略。

\medskip\noindent
证明 (6)。这由引理 \ref{more-algebra-lemma-class-weakly-etale-over-field}
的 (4) 得出。
\end{proof}

\begin{lemma}
\label{more-algebra-lemma-change-fields-max-weakly-etale-subalgebra}
设 $L/K$ 为域扩张，$A$ 为 $K$-代数。设 $B \subset A$ 为引理
\ref{more-algebra-lemma-max-weakly-etale-subalgebra} 中 $A$ 的极大弱
\'{e}tale $K$-子代数。则 $B \otimes_K L$ 是
$A \otimes_K L$ 的极大弱 \'{e}tale $L$-子代数。
\end{lemma}

\begin{proof}
对 $K$ 上的代数 $A$，以 $B_{max}(A/K)$ 表示 $A$ 的极大弱
\'{e}tale $K$-子代数。类似地，若 $A'$ 是 $L$-代数，则以
$B_{max}(A'/L)$ 表示 $A'$ 的极大弱 \'{e}tale $L$-子代数。由于
$B_{max}(A/K) \otimes_K L$ 在 $L$ 上弱 \'{e}tale（引理
\ref{more-algebra-lemma-base-change-weakly-etale}），并且
$B_{max}(A/K) \otimes_K L \subset A \otimes_K L$，得到典范单射
% P02-E0533
$$
B_{max}(A/K) \otimes_K L \to B_{max}((A \otimes_K L)/L)
$$
。引理断言该映射是同构。

\medskip\noindent
为对 $L$ 与当前 $K$-代数 $A$ 证明引理，只需对 $L$ 的任意域扩张
$L'$ 证明引理。事实上，有分解
$$
B_{max}(A/K) \otimes_K L' \to
B_{max}((A \otimes_K L)/L) \otimes_L L' \to
B_{max}((A \otimes_K L')/L')
$$
，故若
$B_{max}(A/K) \otimes_K L \to B_{max}((A \otimes_K L)/L)$
不满射，则该复合不可能满射。因此可假设 $L$ 代数闭。

\medskip\noindent
归约到有限型 $K$-代数。可以把 $A$ 写成其各有限型
$K$-子代数的滤过余极限。
% P02-E0531
利用引理 \ref{more-algebra-lemma-properties-of-max-weakly-etale-subalgebra}，
可见只需对有限型 $K$-代数证明引理。

\medskip\noindent
假设 $A$ 是有限型 $K$-代数。由于
$A \to A_{red}$ 的核幂零，
$A \otimes_K L \to A_{red} \otimes_K L$ 的核也幂零。于是
$$
B_{max}((A \otimes_K L)/L) \to
B_{max}((A_{red} \otimes_K L)/L)
$$
单射，因为其核幂零，而弱 \'{e}tale $L$-代数
$B_{max}((A \otimes_K L)/L)$ 既约（引理
\ref{more-algebra-lemma-class-weakly-etale-over-field}）。由引理
\ref{more-algebra-lemma-properties-of-max-weakly-etale-subalgebra}，
$B_{max}(A/K) = B_{max}(A_{red}/K)$，故只需对 $A_{red}$ 证明引理。

\medskip\noindent
假设 $A$ 是既约有限型 $K$-代数。令 $Q = Q(A)$ 为 $A$ 的全商环。
则 $A \subset Q$，且
% P02-E0532
$A \otimes_K L \subset Q \otimes_A L$，从而
$$
B_{max}(A/K) = A \cap B_{max}(Q/K)
$$
以及
$$
B_{max}((A \otimes_K L)/L) =
(A \otimes_K L) \cap B_{max}((Q \otimes_K L)/L)
$$
；这是引理 \ref{more-algebra-lemma-properties-of-max-weakly-etale-subalgebra}
的结论。由于 $-\otimes_K L$ 是正合函子，若对 $Q$ 证明了结论，
则 $A$ 的结论随之成立。由于 $Q$ 是域的有限乘积（《代数》引理
\ref{algebra-lemma-total-ring-fractions-no-embedded-points}、
\ref{algebra-lemma-minimal-prime-reduced-ring}、
\ref{algebra-lemma-Noetherian-irreducible-components} 与
\ref{algebra-lemma-Noetherian-permanence}），而 $B_{max}$ 与乘积交换
（引理 \ref{more-algebra-lemma-properties-of-max-weakly-etale-subalgebra}），
只需在 $A$ 为域时证明引理。

\medskip\noindent
假设 $A$ 为域。利用证明第三段中的论证，可归约到 $A$ 在 $K$ 上
有限生成的情形。（事实上，我们归约到域的方式产生了 $K$ 的
有限生成域扩张。）

\medskip\noindent
假设 $A$ 是 $K$ 的有限生成域扩张。由引理
\ref{more-algebra-lemma-properties-of-max-weakly-etale-subalgebra} 的 (6)，
$K' = B_{max}(A/K)$ 是 $K$ 的可分代数域扩张。因此，$K'$ 是 $K$
的有限可分域扩张，并且由《代数》引理
\ref{algebra-lemma-make-geometrically-irreducible}，$A$ 在 $K'$
上几何不可约。由于 $L$ 代数闭且 $K'/K$ 有限可分，可见
$$
K' \otimes_K L \to \prod\nolimits_{\sigma \in \Hom_K(K', L)} L,\quad
\alpha \otimes \beta \mapsto (\sigma(\alpha)\beta)_\sigma
$$
是同构（《域》引理
\ref{fields-lemma-finite-separable-tensor-alg-closed}）。因此
$$
A \otimes_K L = A \otimes_{K'} (K' \otimes_K L) =
\prod\nolimits_{\sigma \in \Hom_K(K', L)} A \otimes_{K', \sigma} L
$$
。由于 $A$ 在 $K'$ 上几何不可约，可见
$A \otimes_{K', \sigma} L$ 有唯一的极小素理想。由于 $L$ 代数闭，
由引理 \ref{more-algebra-lemma-properties-of-max-weakly-etale-subalgebra} 的 (6)，
$B_{max}((A \otimes_{K', \sigma} L)/L) = L$，因为这个 $L$-代数是
$L$ 的代数域扩张。于是，$A \otimes_K L$ 的极大弱
\'{e}tale $K' \otimes_K L$-子代数是 $K' \otimes_K L$，因为可以
像上面一样把这些子代数分解为乘积。故包含
$K' \otimes_K L \subset B_{max}((A \otimes_K L)/L)$ 是等式：
由引理 \ref{more-algebra-lemma-weakly-etale-permanence}，环映射
$K' \otimes_K L \to B_{max}((A \otimes_K L)/L)$ 弱 \'{e}tale。
\end{proof}
\section{局部不可约性}
\label{more-algebra-section-unibranch}

\noindent
下述定义看来是普遍接受的定义。为理解它，注意若 $A \subset B$
是局部整环的整扩张，则由上升定理，$A \to B$ 是局部环同态
（《代数》引理 \ref{algebra-lemma-integral-going-up}）。

\begin{definition}
\label{more-algebra-definition-unibranch}
\begin{reference}
\cite[Chapter 0 (23.2.1)]{EGA4}
\end{reference}
设 $A$ 为局部环。若约化环 $A_{red}$ 是整环，并且 $A_{red}$
在其分式域中的整闭包 $A'$ 是局部环，则称 $A$ 为{\it 单支的}。
若 $A$ 单支，并且 $A'$ 的剩余域是 $A$ 的剩余域的纯不可分扩张，
则称 $A$ 为{\it 几何单支的}。
\end{definition}

\noindent
设 $A$ 为局部环。以下是等价的表述：
\begin{enumerate}
\item $A$ 单支，当且仅当 $A$ 有唯一的极小素理想 $\mathfrak p$，
并且 $A/\mathfrak p$ 在其分式域中的整闭包是局部环；
\item $A$ 几何单支，当且仅当 $A$ 有唯一的极小素理想
$\mathfrak p$，并且 $A/\mathfrak p$ 在其分式域中的整闭包是局部环，
而且该局部环的剩余域是 $A$ 的剩余域的纯不可分扩张。
\end{enumerate}
正规局部环是几何单支的（这由定义 \ref{more-algebra-definition-unibranch}
与《代数》定义 \ref{algebra-definition-ring-normal} 得出）。
引理 \ref{more-algebra-lemma-unibranch} 与 \ref{more-algebra-lemma-geometrically-unibranch}
表明，研究（几何）单支这一性质是合理的。

\begin{lemma}
\label{more-algebra-lemma-branches}
设 $A$ 为局部环，并假设 $A$ 只有有限多个极小素理想。令 $A'$
为 $A$ 在 $A_{red}$ 的全商环中的整闭包，令 $A^h$ 为 $A$
的 Hensel 化。考虑映射
$$
\Spec(A') \leftarrow \Spec((A')^h) \rightarrow \Spec(A^h)
$$
，其中 $(A')^h = A' \otimes_A A^h$。则：
\begin{enumerate}
\item 左侧箭头在极大理想上双射；
\item 右侧箭头在极小素理想上双射；
\item $(A')^h$ 的每个极小素理想都包含于唯一的极大理想中，
且每个极大理想恰好包含一个极小素理想。
\end{enumerate}
\end{lemma}

\begin{proof}
令 $I \subset A$ 为幂零元构成的理想。由《代数》引理
\ref{algebra-lemma-quotient-henselization}，有
$(A/I)^h = A^h/IA^h$。环 $A$、$A^h$、$A'$ 与 $(A')^h$ 的谱分别与
$A/I$、$A^h/IA^h$、$A'$ 及
$(A')^h = A' \otimes_{A/I} A^h/IA^h$ 的谱相同。因此可以用
$A_{red} = A/I$ 替换 $A$，并假设 $A$ 既约。于是 $A \subset A'$；
以下将不再特别说明这一点。

\medskip\noindent
证明 (1)。由于 $A'$ 在 $A$ 上整，$(A')^h$ 在 $A^h$ 上整。
由上升定理（《代数》引理 \ref{algebra-lemma-integral-going-up}），
$A'$ 与 $(A')^h$ 的每个极大理想分别位于 $A$ 的极大理想
$\mathfrak m$ 与 $A^h$ 的极大理想 $\mathfrak m^h$ 之上。
因此 (1) 由同构
$$
(A')^h \otimes_{A^h} \kappa^h =
A' \otimes_A A^h \otimes_{A^h} \kappa^h = A' \otimes_A \kappa
$$
得出，因为 $A \to A^h$ 所诱导的剩余域扩张
$\kappa^h/\kappa$ 是平凡的。下文将用到：所显示的环在一个域上整，
故其谱为 profinite 空间；见《代数》引理
\ref{algebra-lemma-integral-over-field} 与
\ref{algebra-lemma-ring-with-only-minimal-primes}。
% P02-E0534

\medskip\noindent
证明 (3)。环 $A'$ 是正规环；事实上，它是正规整环的有限乘积，
见《代数》引理
\ref{algebra-lemma-characterize-reduced-ring-normal}。由于 $A^h$
是 \'etale $A$-代数的滤过余极限，$(A')^h$ 是 \'etale
$A'$-代数的滤过余极限；故由《代数》引理
\ref{algebra-lemma-normal-goes-up} 与
\ref{algebra-lemma-colimit-normal-ring}，$(A')^h$ 是正规环。
% P02-E0535
因此 $(A')^h$ 的每个局部环都是正规整环，从而每个极大理想都包含
唯一的极小素理想。对 $A^h \to (A')^h$ 应用引理
\ref{more-algebra-lemma-integral-over-henselian-pair}，可见
$((A')^h, \mathfrak m(A')^h)$ 是 Hensel 对。若
$\mathfrak q \subset (A')^h$ 是极小素理想（或任意素理想），则由
引理 \ref{more-algebra-lemma-irreducible-henselian-pair-connected}，
$V(\mathfrak q)$ 与 $V(\mathfrak m (A')^h)$ 的交连通。又因
$V(\mathfrak m (A')^h) = \Spec((A')^h \otimes \kappa^h)$ 是
profinite 空间，可见存在唯一的极大理想包含 $\mathfrak q$。
% P02-E0536

\medskip\noindent
证明 (2)。$A'$ 的极小素理想恰为位于 $A$ 的极小素理想之上的素理想
（由构造可知）。由于 $A' \to (A')^h$ 平坦，由下降定理
（《代数》引理 \ref{algebra-lemma-flat-going-down}），$(A')^h$
的每个极小素理想都位于 $A'$ 的某个极小素理想之上。反之，
$(A')^h$ 中位于 $A'$ 的极小素理想之上的任意素理想都是极小的，
因为 $(A')^h$ 是 $A'$ 上 \'etale、因而拟有限的代数的滤过余极限
（略去一个小细节）。由此，$(A')^h$ 的极小素理想恰为位于 $A$
的极小素理想之上的那些素理想。类似地，$A^h$ 的极小素理想恰为
位于 $A$ 的极小素理想之上的素理想。由构造，
$A' \otimes_A Q(A) = Q(A)$，其中 $Q(A)$ 是既约局部环 $A$
的全商环。当然，$Q(A)$ 是 $A$ 的各极小素理想的剩余域的有限乘积。
因此
% P02-E0537
$$
(A')^h \otimes_A Q(A) = A^h \otimes_A A' \otimes_A Q(A) = A^h \otimes_A Q(A)
$$
。上述讨论表明，左侧环的谱是 $(A')^h$ 的极小素理想集合，
右侧环的谱是 $A^h$ 的极小素理想集合。这就完成了证明。
% P02-E0538
\end{proof}

\begin{lemma}
\label{more-algebra-lemma-unibranch}
\begin{reference}
\cite[Chapter IV Proposition 18.6.12]{EGA4}
\end{reference}
设 $A$ 为局部环，$A^h$ 为 $A$ 的 Hensel 化。以下条件等价：
\begin{enumerate}
\item $A$ 单支；
\item $A^h$ 有唯一的极小素理想。
\end{enumerate}
\end{lemma}

\begin{proof}
这由引理 \ref{more-algebra-lemma-branches} 得出，不过我们也给出直接证明。
以 $\mathfrak m$ 表示环 $A$ 的极大理想。回忆剩余域
$\kappa = A/\mathfrak m$ 与 $A^h$ 的剩余域相同。
% P02-E0539

\medskip\noindent
假设 (2) 成立。令 $\mathfrak p^h$ 为 $A^h$ 的唯一极小素理想。
$A \to A^h$ 的平坦性蕴含
$\mathfrak p = A \cap \mathfrak p^h$ 是 $A$ 的唯一极小素理想
（由下降定理，见《代数》引理
\ref{algebra-lemma-flat-going-down}）。又因
$A^h/\mathfrak pA^h = (A/\mathfrak p)^h$（见《代数》引理
\ref{algebra-lemma-quotient-henselization}）由引理
\ref{more-algebra-lemma-henselization-reduced} 是既约的，所以
$\mathfrak p^h = \mathfrak pA^h$。令 $A'$ 为
$A/\mathfrak p$ 在其分式域中的整闭包。我们须证明 $A'$ 是局部环。
由于 $A \to A'$ 是整映射，$A'$ 的每个极大理想都位于
$\mathfrak m$ 之上（由整环映射的上升定理，见《代数》引理
\ref{algebra-lemma-integral-going-up}）。若 $A'$ 不是局部环，则可取
两个不同的极大理想 $\mathfrak m_1$、$\mathfrak m_2$。选取元素
$f_1, f_2 \in A'$，使得 $f_i \in \mathfrak m_i$ 且
$f_i \not \in \mathfrak m_{3 - i}$。我们得到有限子代数
% P02-E0540
$B = A/\mathfrak p[f_1, f_2] \subset A'$，其极大理想
$B \cap \mathfrak m_i$（$i = 1, 2$）彼此不同。注意包含关系
$$
A/\mathfrak p \subset B \subset \kappa(\mathfrak p)
$$
与平坦环映射 $A \to A^h$ 张量后给出包含关系
$$
A^h/\mathfrak p^h \subset
B \otimes_A A^h \subset
\kappa(\mathfrak p) \otimes_A A^h \subset
\kappa(\mathfrak p^h)
$$
；最后一个包含成立，是因为
$\kappa(\mathfrak p) \otimes_A A^h =
\kappa(\mathfrak p) \otimes_{A/\mathfrak p} A^h/\mathfrak p^h$
是整环 $A^h/\mathfrak p^h$ 的局部化。注意
$B \otimes_A \kappa$ 至少有两个极大理想，因为
$B/\mathfrak mB$ 有两个极大理想。因此，由于 $A^h$ 是 Hensel 环，
$B \otimes_A A^h$ 是 $\geq 2$ 个局部环的乘积；见《代数》引理
\ref{algebra-lemma-mop-up}。但我们刚刚看到
$B \otimes_A A^h$ 是一个整环的子环，由此得到矛盾。

\medskip\noindent
假设 (1) 成立。令 $\mathfrak p \subset A$ 为唯一的极小素理想，
并令 $A'$ 为 $A/\mathfrak p$ 在其分式域中的整闭包。设
$A \to B$ 为局部环之间的局部映射，它诱导剩余域的同构，并且是
某个 \'etale $A$-代数的局部化。特别地，$\mathfrak m_B$ 是包含
$\mathfrak m B$ 的唯一素理想。于是 $B' = A' \otimes_A B$
在 $B$ 上整，而 $A \to A'$ 是局部映射这一假设蕴含 $B'$ 是局部环
（《代数》引理 \ref{algebra-lemma-local-tensor-with-integral}）。
另一方面，$A' \to B'$ 是 \'etale 环映射的局部化，故 $B'$ 正规；
见《代数》引理 \ref{algebra-lemma-normal-goes-up}。因此 $B'$
是（局部）正规整环。最后，有
% P02-E0537
$$
B/\mathfrak pB \subset B \otimes_A \kappa(\mathfrak p)
= B' \otimes_{A'} (\text{fraction field of }A') \subset
\text{fraction field of }B'
$$
。故 $B/\mathfrak pB$ 是整环，从而 $B$ 有唯一的极小素理想
（因为由 $A \to B$ 的平坦性，这些极小素理想全都必须位于
$\mathfrak p$ 之上）。由于 $A^h$ 是这些局部环 $B$ 的滤过余极限，
$A^h$ 有唯一的极小素理想。具体而言，若某些非幂零元
$f, g$ 在 $A^h$ 中满足 $fg = 0$，则可找到一个如上的 $B$
同时包含 $f$ 与 $g$，这会导致矛盾。
\end{proof}

\begin{lemma}
\label{more-algebra-lemma-geometric-branches}
设 $(A, \mathfrak m, \kappa)$ 为局部环，并假设 $A$ 只有有限多个
极小素理想。令 $A'$ 为 $A$ 在 $A_{red}$ 的全商环中的整闭包。
选取 $\kappa$ 的代数闭包 $\overline{\kappa}$，并以
$\kappa^{sep} \subset \overline{\kappa}$ 表示 $\kappa$
的可分代数闭包。令 $A^{sh}$ 为 $A$ 关于 $\kappa^{sep}$
的严格 Hensel 化。考虑映射
$$
\Spec(A') \xleftarrow{c} \Spec((A')^{sh}) \xrightarrow{e} \Spec(A^{sh})
$$
，其中 $(A')^{sh} = A' \otimes_A A^{sh}$。则：
\begin{enumerate}
\item 对极大理想 $\mathfrak m' \subset A'$，剩余域
$\kappa'$ 是 $\kappa$ 的代数扩张，且 $c$ 在 $\mathfrak m'$
上的纤维可典范地与
$\Hom_\kappa(\kappa', \overline{\kappa})$ 等同；
\item 右侧箭头在极小素理想上双射；
\item $(A')^{sh}$ 的每个极小素理想都包含于唯一的极大理想中，
且每个极大理想包含唯一的极小素理想。
\end{enumerate}
\end{lemma}

\begin{proof}
本证明几乎与引理 \ref{more-algebra-lemma-branches} 的证明完全相同。令
$I \subset A$ 为幂零元构成的理想。由《代数》引理
\ref{algebra-lemma-quotient-henselization}，有
$(A/I)^{sh} = A^{sh}/IA^{sh}$。环
% P02-E0541
$A$、$A^{sh}$、$A'$ 与 $(A')^h$ 的谱分别与
$A/I$、$A^{sh}/IA^{sh}$、$A'$ 及
$(A')^{sh} = A' \otimes_{A/I} A^{sh}/IA^{sh}$ 的谱相同。
因此可以用 $A_{red} = A/I$ 替换 $A$，并假设 $A$ 既约。于是
$A \subset A'$；以下将不再特别说明这一点。

\medskip\noindent
证明 (1)。域扩张 $\kappa'/\kappa$ 是代数的，因为 $A'$
在 $A$ 上整。由于 $A'$ 在 $A$ 上整，$(A')^{sh}$ 在
$A^{sh}$ 上整。由上升定理（《代数》引理
\ref{algebra-lemma-integral-going-up}），$A'$ 与 $(A')^{sh}$
的每个极大理想分别位于 $A$ 的极大理想 $\mathfrak m$ 与
% P02-E0542
$A^h$ 的极大理想 $\mathfrak m^{sh}$ 之上。我们有
$$
(A')^{sh} \otimes_{A^{sh}} \kappa^{sep} =
A' \otimes_A A^h \otimes_{A^h} \kappa^{sep} =
(A' \otimes_A \kappa) \otimes_{\kappa} \kappa^{sep}
$$
，因为 $A^{sh}$ 的剩余域是 $\kappa^{sep}$。因此 $c$
在 $\mathfrak m'$ 上的纤维是
$\kappa' \otimes_\kappa \kappa^{sep}$ 的谱。由双射
% P02-E0537
$$
\Hom_\kappa(\kappa', \overline{\kappa}) \to
\Spec(\kappa' \otimes_\kappa \kappa^{sep}),\quad
\sigma \mapsto \Ker(
\sigma \otimes 1 : \kappa' \otimes_\kappa \kappa^{sep} \to \overline{\kappa}
)
$$
，可知 (1) 成立。下文将用到：所显示的环在一个域上整，故其谱为
profinite 空间；见《代数》引理
\ref{algebra-lemma-integral-over-field} 与
\ref{algebra-lemma-ring-with-only-minimal-primes}。

\medskip\noindent
证明 (3)。环 $A'$ 是正规环；事实上，它是正规整环的有限乘积，
见《代数》引理
\ref{algebra-lemma-characterize-reduced-ring-normal}。由于
$A^{sh}$ 是 \'etale $A$-代数的滤过余极限，$(A')^{sh}$ 是
\'etale $A'$-代数的滤过余极限；故由《代数》引理
\ref{algebra-lemma-normal-goes-up} 与
\ref{algebra-lemma-colimit-normal-ring}，$(A')^{sh}$ 是正规环。
% P02-E0535
因此 $(A')^{sh}$ 的每个局部环都是正规整环，从而每个极大理想都包含
唯一的极小素理想。
对 $A^{sh} \to (A')^{sh}$ 应用引理
\ref{more-algebra-lemma-integral-over-henselian-pair}，可见
$((A')^{sh}, \mathfrak m(A')^{sh})$ 是 Hensel 对。若
$\mathfrak q \subset (A')^{sh}$ 是极小素理想（或任意素理想），
则由引理 \ref{more-algebra-lemma-irreducible-henselian-pair-connected}，
$V(\mathfrak q)$ 与 $V(\mathfrak m (A')^{sh})$ 的交连通。又因
% P02-E0543
$V(\mathfrak m (A')^{sh}) = \Spec((A')^{sh} \otimes \kappa^{sh})$
是 profinite 空间，可见存在唯一的极大理想包含 $\mathfrak q$。

\medskip\noindent
证明 (2)。$A'$ 的极小素理想恰为位于 $A$ 的极小素理想之上的素理想
（由构造可知）。由于 $A' \to (A')^{sh}$ 平坦，由下降定理
（《代数》引理 \ref{algebra-lemma-flat-going-down}），
$(A')^{sh}$ 的每个极小素理想都位于 $A'$ 的某个极小素理想之上。
反之，$(A')^{sh}$ 中位于 $A'$ 的极小素理想之上的任意素理想都是
极小的，因为 $(A')^{sh}$ 是 $A'$ 上 \'etale、因而拟有限的代数的
滤过余极限（略去一个小细节）。由此，$(A')^{sh}$ 的极小素理想
恰为位于 $A$ 的极小素理想之上的那些素理想。类似地，$A^{sh}$
的极小素理想恰为位于 $A$ 的极小素理想之上的素理想。由构造，
$A' \otimes_A Q(A) = Q(A)$，其中 $Q(A)$ 是既约局部环 $A$
的全商环。当然，$Q(A)$ 是 $A$ 的各极小素理想的剩余域的有限乘积。
因此
% P02-E0537
$$
(A')^{sh} \otimes_A Q(A) =
A^{sh} \otimes_A A' \otimes_A Q(A) = A^{sh} \otimes_A Q(A)
$$
。上述讨论表明，左侧环的谱是 $(A')^{sh}$ 的极小素理想集合，
右侧环的谱是 $A^{sh}$ 的极小素理想集合。这就完成了证明。
% P02-E0538
\end{proof}

\begin{lemma}
\label{more-algebra-lemma-geometrically-unibranch}
\begin{reference}
\cite[Lemma 2.2]{Etale-coverings} and
\cite[Chapter IV Proposition 18.8.15]{EGA4}
\end{reference}
设 $A$ 为局部环，$A^{sh}$ 为 $A$ 的一个严格 Hensel 化。
以下条件等价：
\begin{enumerate}
\item $A$ 几何单支；
\item $A^{sh}$ 有唯一的极小素理想。
\end{enumerate}
\end{lemma}

\begin{proof}
这由引理 \ref{more-algebra-lemma-geometric-branches} 得出，不过我们也给出一个
直接证明；它与引理 \ref{more-algebra-lemma-unibranch} 的直接证明几乎完全相同。
以 $\mathfrak m$ 表示环 $A$ 的极大理想，以
$\kappa$、$\kappa^{sh}$ 分别表示 $A$、$A^{sh}$ 的剩余域。
% P02-E0539

\medskip\noindent
假设 (2) 成立。令 $\mathfrak p^{sh}$ 为 $A^{sh}$ 的唯一极小素理想。
$A \to A^{sh}$ 的平坦性蕴含
$\mathfrak p = A \cap \mathfrak p^{sh}$ 是 $A$ 的唯一极小素理想
（由下降定理，见《代数》引理
\ref{algebra-lemma-flat-going-down}）。又因
$A^{sh}/\mathfrak pA^{sh} = (A/\mathfrak p)^{sh}$
（见《代数》引理
\ref{algebra-lemma-quotient-strict-henselization}）由引理
\ref{more-algebra-lemma-henselization-reduced} 是既约的，所以
$\mathfrak p^{sh} = \mathfrak pA^{sh}$。令 $A'$ 为
$A/\mathfrak p$ 在其分式域中的整闭包。我们须证明 $A'$ 是局部环，
且其剩余域是 $\kappa$ 的纯不可分扩张。由于 $A \to A'$ 是整映射，
$A'$ 的每个极大理想都位于 $\mathfrak m$ 之上（由整环映射的
上升定理，见《代数》引理
\ref{algebra-lemma-integral-going-up}）。若 $A'$ 不是局部环，则可取
两个不同的极大理想 $\mathfrak m_1$、$\mathfrak m_2$。选取元素
$f_1, f_2 \in A'$，使得 $f_i \in \mathfrak m_i, f_i \not \in \mathfrak m_{3 - i}$，
便得到有限子代数
$B = A[f_1, f_2] \subset A'$，其极大理想
$B \cap \mathfrak m_i$（$i = 1, 2$）彼此不同。
% P02-E0544
若 $A'$ 是以 $\mathfrak m'$ 为极大理想的局部环，但
$A/\mathfrak m \subset A'/\mathfrak m'$ 不是纯不可分扩张，
则可取 $f \in A'$，使它在 $A'/\mathfrak m'$ 中的像生成
$A/\mathfrak m$ 的一个有限但非纯不可分扩张；由此得到有限局部
子代数 $B = A[f] \subset A'$，其剩余域不是
$A/\mathfrak m$ 的纯不可分扩张。注意包含关系
$$
A/\mathfrak p \subset B \subset \kappa(\mathfrak p)
$$
与平坦环映射 $A \to A^{sh}$ 张量后给出包含关系
$$
A^{sh}/\mathfrak p^{sh} \subset
B \otimes_A A^{sh} \subset
\kappa(\mathfrak p) \otimes_A A^{sh} \subset
\kappa(\mathfrak p^{sh})
$$
；最后一个包含成立，是因为
$\kappa(\mathfrak p) \otimes_A A^{sh} =
\kappa(\mathfrak p) \otimes_{A/\mathfrak p} A^{sh}/\mathfrak p^{sh}$
是整环 $A^{sh}/\mathfrak p^{sh}$ 的局部化。注意
$B \otimes_A \kappa^{sh}$ 至少有两个极大理想，因为
$B/\mathfrak mB$ 或者有两个极大理想，或者只有一个极大理想而其
剩余域不是 $\kappa$ 的纯不可分扩张，并且 $\kappa^{sh}$
是可分代数闭的。因此，由于 $A^{sh}$ 是严格 Hensel 环，
$B \otimes_A A^{sh}$ 是 $\geq 2$ 个局部环的乘积；见《代数》引理
\ref{algebra-lemma-mop-up-strictly-henselian}。但我们刚刚看到
$B \otimes_A A^{sh}$ 是一个整环的子环，由此得到矛盾。

\medskip\noindent
假设 (1) 成立。令 $\mathfrak p \subset A$ 为唯一的极小素理想，
并令 $A'$ 为 $A/\mathfrak p$ 在其分式域中的整闭包。设
$A \to B$ 为局部环之间的局部映射，它是某个
\'etale $A$-代数的局部化。特别地，$\mathfrak m_B$ 是包含
$\mathfrak m_AB$ 的唯一素理想。于是 $B' = A' \otimes_A B$
在 $B$ 上整，而 $A \to A'$ 是局部映射且诱导纯不可分剩余域扩张
这一假设蕴含 $B'$ 是局部环（《代数》引理
\ref{algebra-lemma-local-tensor-with-integral}）。另一方面，
$A' \to B'$ 是 \'etale 环映射的局部化，故 $B'$ 正规；
见《代数》引理 \ref{algebra-lemma-normal-goes-up}。因此 $B'$
是（局部）正规整环。最后，有
% P02-E0537
$$
B/\mathfrak pB \subset B \otimes_A \kappa(\mathfrak p)
= B' \otimes_{A'} (\text{fraction field of }A') \subset
\text{fraction field of }B'
$$
。故 $B/\mathfrak pB$ 是整环，从而 $B$ 有唯一的极小素理想
（因为由 $A \to B$ 的平坦性，这些极小素理想全都必须位于
$\mathfrak p$ 之上）。由于 $A^{sh}$ 是这些局部环 $B$ 的滤过余极限，
$A^{sh}$ 有唯一的极小素理想。具体而言，若某些非幂零元
$f, g$ 在 $A^{sh}$ 中满足 $fg = 0$，则可找到一个如上的 $B$
同时包含 $f$ 与 $g$，这会导致矛盾。
\end{proof}

\begin{definition}
\label{more-algebra-definition-number-of-branches}
设 $A$ 为局部环，其 Hensel 化与严格 Hensel 化分别为
$A^h$、$A^{sh}$。若 $A^h$ 的极小素理想数有限，则定义
{\it $A$ 的分支数}为这一数目；否则定义为 $\infty$。
若 $A^{sh}$ 的极小素理想数有限，则定义
{\it $A$ 的几何分支数}为这一数目；否则定义为 $\infty$。
\end{definition}

\noindent
下面具体说明它与定义 \ref{more-algebra-definition-unibranch} 的关系。

\begin{lemma}
\label{more-algebra-lemma-number-of-branches-1}
设 $(A, \mathfrak m, \kappa)$ 为局部环。
\begin{enumerate}
\item 若 $A$ 有无限多个极小素理想，则 $A$ 的（几何）分支数为
$\infty$。
\item $A$ 的分支数为 $1$，当且仅当 $A$ 单支。
\item $A$ 的几何分支数为 $1$，当且仅当 $A$ 几何单支。
\end{enumerate}
假设 $A$ 只有有限多个极小素理想，并令 $A'$ 为 $A$ 在
$A_{red}$ 的全商环中的整闭包。则：
\begin{enumerate}
\item[(4)] $A$ 的分支数等于 $A'$ 的极大理想 $\mathfrak m'$ 的数目；
\item[(5)] 为得到 $A$ 的几何分支数，须按
$\kappa(\mathfrak m')/\kappa$ 的可分次数给出的重数，计数 $A'$
的每个极大理想 $\mathfrak m'$。
\end{enumerate}
\end{lemma}

\begin{proof}
本引理由各定义、引理 \ref{more-algebra-lemma-branches}、引理
\ref{more-algebra-lemma-geometric-branches} 以及《域》引理
\ref{fields-lemma-separable-degree} 立即得出。
\end{proof}

\begin{lemma}
\label{more-algebra-lemma-invariance-number-branches-smooth}
设 $A \to B$ 为局部环之间的局部同态，并且是某个光滑环映射的
局部化。
\begin{enumerate}
\item $A$ 的几何分支数等于 $B$ 的几何分支数。
\item 若 $A \to B$ 诱导剩余域的纯不可分扩张，则 $A$ 的分支数
等于 $B$ 的分支数。
\end{enumerate}
\end{lemma}

\begin{proof}
我们将用到：光滑环映射平坦（《代数》引理
\ref{algebra-lemma-smooth-syntomic}）；局部化平坦（《代数》引理
\ref{algebra-lemma-flat-localization}）；平坦环映射的复合平坦
（《代数》引理 \ref{algebra-lemma-composition-flat}）；平坦环映射
经基变换后仍平坦（《代数》引理
\ref{algebra-lemma-flat-base-change}）；平坦局部同态忠实平坦
（《代数》引理 \ref{algebra-lemma-local-flat-ff}）；（严格）
Hensel 化平坦（引理 \ref{more-algebra-lemma-dumb-properties-henselization}）；
以及平坦环映射的下降定理（《代数》引理
\ref{algebra-lemma-flat-going-down}）。
% P02-E0545

\medskip\noindent
证明 (2)。令 $A^h$、$B^h$ 分别为 $A$、$B$ 的 Hensel 化。
则 $B^h$ 是 $A^h \otimes_A B$ 在位于 $\mathfrak m_B$
之上的唯一极大理想处的 Hensel 化；见《代数》引理
\ref{algebra-lemma-henselian-functorial-improve}。因此可以并且确实
假设 $A$ 是 Hensel 环。由于 $A \to B \to B^h$ 平坦，
$B^h$ 的每个极小素理想都位于 $A$ 的某个极小素理想之上；
又由于 $A \to B^h$ 忠实平坦，$A$ 的每个极小素理想之上确有
$B^h$ 的某个极小素理想；两处都使用平坦环映射的下降定理。
因此只须证明：给定极小素理想 $\mathfrak p \subset A$，
$B^h$ 中至多有一个极小素理想位于 $\mathfrak p$ 之上。
用 $A/\mathfrak p$ 替换 $A$、用 $B/\mathfrak p B$ 替换 $B$ 后，
可假设 $A$ 为整环；此时由《代数》引理
\ref{algebra-lemma-quotient-henselization}，$A$ 仍是 Hensel 环。
% P02-E0546
由引理 \ref{more-algebra-lemma-unibranch}，$A$ 在其分式域中的整闭包
$A'$ 是局部整环。当然，$A'$ 是正规整环。由《代数》引理
\ref{algebra-lemma-normal-goes-up}，
$A' \otimes_A B^h$ 是正规环（该引理直接给出
$A' \otimes_A B$ 的结论；为推到 $A' \otimes_A B^h$，
使用 $B^h$ 是 \'etale $B$-代数的余极限以及《代数》引理
\ref{algebra-lemma-colimit-normal-ring}）。由《代数》引理
\ref{algebra-lemma-local-tensor-with-integral}，
$A' \otimes_A B^h$ 是局部环（这里使用了关于 $A$ 与 $B$
的剩余域的假设）。因此 $A' \otimes_A B^h$ 是局部正规环，
从而是局部整环。由 $A \to B^h$ 的平坦性，
$B^h \subset A' \otimes_A B^h$，故 $B^h$ 是整环，如所求。

\medskip\noindent
证明 (1)。令 $A^{sh}$、$B^{sh}$ 分别为 $A$、$B$ 的严格
Hensel 化。则 $B^{sh}$ 是
% P02-E0547
$A^h \otimes_A B$ 在某个同时位于 $\mathfrak m_B$
与 $\mathfrak m_{A^h}$ 之上的极大理想处的严格 Hensel 化；
见《代数》引理
\ref{algebra-lemma-strictly-henselian-functorial-improve}。
因此可以并且确实假设 $A$ 是严格 Hensel 环。由于
$A \to B \to B^{sh}$ 平坦，$B^{sh}$ 的每个极小素理想都位于
$A$ 的某个极小素理想之上；又由于 $A \to B^{sh}$ 忠实平坦，
$A$ 的每个极小素理想之上确有 $B^{sh}$ 的某个极小素理想；
两处都使用平坦环映射的下降定理。因此只须证明：给定极小素理想
$\mathfrak p \subset A$，$B^{sh}$ 中至多有一个极小素理想位于
$\mathfrak p$ 之上。用 $A/\mathfrak p$ 替换 $A$、用
$B/\mathfrak p B$ 替换 $B$ 后，可假设 $A$ 为整环；此时由
《代数》引理 \ref{algebra-lemma-quotient-strict-henselization}，
$A$ 仍是严格 Hensel 环。由引理
\ref{more-algebra-lemma-geometrically-unibranch}，$A$ 在其分式域中的整闭包
$A'$ 是局部整环，而且其剩余域是 $A$ 的剩余域的纯不可分扩张。
当然，$A'$ 是正规整环。由《代数》引理
\ref{algebra-lemma-normal-goes-up}，
$A' \otimes_A B^{sh}$ 是正规环（该引理直接给出
$A' \otimes_A B$ 的结论；为推到 $A' \otimes_A B^{sh}$，
使用 $B^{sh}$ 是 \'etale $B$-代数的余极限以及《代数》引理
\ref{algebra-lemma-colimit-normal-ring}）。由《代数》引理
\ref{algebra-lemma-local-tensor-with-integral}，
$A' \otimes_A B^{sh}$ 是局部环（因为 $A \subset A'$
诱导纯不可分剩余域扩张）。因此 $A' \otimes_A B^{sh}$
是局部正规环，从而是局部整环。由 $A \to B^{sh}$ 的平坦性，
$B^{sh} \subset A' \otimes_A B^{sh}$，故 $B^{sh}$
是整环，如所求。
\end{proof}
\section{分支杂论}
\label{more-algebra-section-branches}

\noindent
本节给出一些与第 \ref{more-algebra-section-unibranch} 节所定义的局部环分支有关的
结果。

\begin{lemma}
\label{more-algebra-lemma-ideal-inverse-nonzero}
设 $A$ 与 $B$ 为整环，$A \to B$ 为环映射。进一步假设
$A \to B$ 至少具有以下性质之一：
% P02-E0548
\begin{enumerate}
\item 它是某个 \'etale 环映射的局部化；
\item 它平坦，并且是某个非分歧环映射的局部化；
\item 它平坦，并且是某个拟有限环映射的局部化；
\item 它平坦，并且是某个整环映射的局部化；
\item 它平坦，并且 $\Spec(B) \to \Spec(A)$ 的各纤维中的点之间
不存在非平凡特殊化；
\item $\Spec(B) \to \Spec(A)$ 把泛点映到泛点，并且各纤维中的点
之间不存在非平凡特殊化；或者
\item $\Spec(B)$ 中恰有一个点映到 $\Spec(A)$ 的泛点。
\end{enumerate}
则对 $B$ 的每个非零理想 $J$，$A \cap J$ 都非零。
\end{lemma}

\begin{proof}
先证明情形 (7)。分别令 $K$ 与 $L$ 为 $A$ 与 $B$ 的分式域。
由《代数》引理
\ref{algebra-lemma-minimal-prime-image-minimal-prime}，
$\Spec(B)$ 中映到泛点
$(0) \in \Spec(A)$ 的唯一点是 $(0) \in \Spec(B)$。由此，
$B \otimes_A K$ 是仅有一个素理想且该素理想的剩余域为 $L$
的环（事实上它等于 $L$，但我们不需要这一点）。选取非零的
$b \in J$。则 $b$ 在 $L$ 中的像是单位。因此 $b$ 在
$B \otimes_A K$ 中的像是单位（《代数》引理
\ref{algebra-lemma-surjective-on-spec-units}）。由于
$B \otimes_A K = \colim_{f \in A \setminus \{0\}} B_f$，
存在非零的 $f \in A$，使 $b$ 在 $B_f$ 中的像为单位。
这意味着对某个 $b' \in B$ 与 $n \geq 1$，有
$b b' = f^n$。故 $f^n \in A \cap J$，如所求。

\medskip\noindent
在证明的其余部分，我们说明另外每项假设都蕴含 (7)。
% P02-E0549
在假设 (1)--(5) 下，环映射 $A \to B$ 平坦，因而
$A \to B$ 单射（因为由《代数》引理
\ref{algebra-lemma-local-flat-ff}，平坦局部同态忠实平坦）。
所以 $\Spec(B)$ 的泛点映到 $\Spec(A)$ 的泛点。现在，若
$\Spec(B) \to \Spec(A)$ 的各纤维中的点之间不存在非平凡特殊化，
则 $\Spec(B)$ 的这个泛点当然必须是映到 $\Spec(A)$ 泛点的唯一点。
故 (6) 蕴含 (7)。最后，为完成证明，我们说明在情形 (1)--(5) 中，
$\Spec(B) \to \Spec(A)$ 的各纤维中的点之间不存在非平凡特殊化。
具体而言，整映射情形见《代数》引理
\ref{algebra-lemma-integral-no-inclusion}，拟有限情形见《代数》定义
\ref{algebra-definition-quasi-finite}；并使用非分歧环映射与
\'etale 环映射都是拟有限的这一事实（《代数》引理
\ref{algebra-lemma-unramified-quasi-finite} 与
\ref{algebra-lemma-etale-quasi-finite}）。
\end{proof}

\begin{lemma}
\label{more-algebra-lemma-local-unramified-extension-unibranch-domain-is-etale}
设 $A \to B$ 为环映射，$\mathfrak q \subset B$ 为位于素理想
$\mathfrak p \subset A$ 之上的素理想。假设：
\begin{enumerate}
\item $A$ 是整环；
\item $A_\mathfrak p$ 几何单支；
\item $A \to B$ 在 $\mathfrak q$ 处非分歧；
\item $A_\mathfrak p \to B_\mathfrak q$ 单射。
\end{enumerate}
则存在 $g \in B$、$g \not \in \mathfrak q$，使得
$B_g$ 在 $A$ 上 \'etale。
\end{lemma}

\begin{proof}
由《代数》命题
\ref{algebra-proposition-unramified-locally-standard}，用一个
主局部化替换 $B$ 后，可以找到标准 \'etale 环映射
$A \to B'$ 与满射 $B' \to B$。以 $\mathfrak q' \subset B'$
表示 $\mathfrak q$ 的逆像。
% P02-E0550
我们将证明，必要时再用一个主局部化替换 $B'$ 后，
$B' \to B$ 是单射。

\medskip\noindent
本段将问题归约到 $B'$ 是整环的情形。因为 $A$ 是整环，
由《代数》引理 \ref{more-algebra-lemma-reduced-goes-up}，环 $B'$ 既约。
令 $K$ 为 $A$ 的分式域。则 $B' \otimes_A K$ 在一个域上
\'etale，因而是域的有限乘积；见《代数》引理
\ref{algebra-lemma-etale-over-field}。由于 $A \to B'$ 是
\'etale 的（因而平坦），$B'$ 的极小素理想位于
$(0) \subset A$ 之上（由平坦环映射的下降定理）。
% P02-E0551
由此 $B'$ 只有有限多个极小素理想，记作
$\mathfrak r_1, \ldots, \mathfrak r_r \subset B'$。由于
$A_\mathfrak p$ 几何单支且 $A \to B'$ 是 \'etale 的，
由引理 \ref{more-algebra-lemma-invariance-number-branches-smooth} 与
\ref{more-algebra-lemma-number-of-branches-1}，环 $B'_{\mathfrak q'}$
是整环。因此恰对一个 $i = i_0$，有
$\mathfrak q' \supset \mathfrak r_i$。由《代数》引理
\ref{algebra-lemma-silly}，可以选取 $g' \in B'$，使
$g' \not \in \mathfrak r_{i_0}$，但对 $i \not = i_0$ 有
$g' \in \mathfrak r_i$。分别用 $B'_{g'}$ 与 $B_{g'}$
替换 $B'$ 与 $B$ 后，$B'$ 成为整环。

\medskip\noindent
假设 $B'$ 是整环；特别地，$B' \subset B'_{\mathfrak q'}$。
若 $B' \to B$ 不单射，则
$J = \Ker(B'_{\mathfrak q'} \to B_\mathfrak q)$ 非零。
对 $A_\mathfrak p \to B'_{\mathfrak q'}$ 应用引理
\ref{more-algebra-lemma-ideal-inverse-nonzero}，得到非零元素
$a \in A_\mathfrak p$，它在 $B_\mathfrak q$ 中的像为零；
这与假设 (4) 矛盾。证明完成。
\end{proof}

\begin{lemma}
\label{more-algebra-lemma-unramified-extension-unibranch-domain-is-etale}
\begin{reference}
Generalization of \cite[Expose I, Theorem 9.5 part (ii)]{SGA1}
\end{reference}
设 $(A, \mathfrak m)$ 为几何单支局部整环。设
$A \to B$ 为局部环之间的单射局部同态，并且本质有限型。
若 $\mathfrak m B$ 是 $B$ 的极大理想，且所诱导的剩余域扩张可分，
则 $A \to B$ 是某个 \'etale 环映射的局部化。
\end{lemma}

\begin{proof}
可以写成 $B = C_\mathfrak q$，其中 $A \to C$ 是有限型环映射，
$\mathfrak q \subset C$ 是位于 $\mathfrak m$ 之上的素理想。
由《代数》引理 \ref{algebra-lemma-characterize-unramified}，
环映射 $A \to C$ 在 $\mathfrak q$ 处非分歧。由《代数》命题
\ref{algebra-proposition-unramified-locally-standard}，用一个
主局部化替换 $C$ 后，可以找到标准 \'etale 环映射
$A \to C'$ 与满射 $C' \to C$。以 $\mathfrak q' \subset C'$
表示 $\mathfrak q$ 的逆像，并置 $B' = C'_{\mathfrak q'}$。
则 $B' \to B$ 满射。只须证明 $B' \to B$ 也单射。

\medskip\noindent
由于 $A$ 是整环，由《代数》引理
\ref{more-algebra-lemma-reduced-goes-up}，环 $C'$ 与 $B'$ 既约。由于
$A$ 几何单支，由引理
\ref{more-algebra-lemma-invariance-number-branches-smooth} 与
\ref{more-algebra-lemma-number-of-branches-1}，环 $B'$ 是整环。
% P02-E0552
若 $B' \to B$ 不单射，则由引理
\ref{more-algebra-lemma-ideal-inverse-nonzero}，$A \cap \Ker(B' \to B)$ 非零，
这与 $A \to B$ 单射的假设矛盾。
\end{proof}

\begin{lemma}
\label{more-algebra-lemma-minimal-primes-tensor-strictly-henselian}
设 $k$ 为代数闭域。设 $A$、$B$ 为剩余域等于 $k$ 的严格
Hensel 局部 $k$-代数。令 $C$ 为 $A \otimes_k B$ 在极大理想
$\mathfrak m_A \otimes_k B + A \otimes_k \mathfrak m_B$
处的严格 Hensel 化。则 $C$ 的极小素理想与 $A$、$B$
的极小素理想对按 $1$ 对 $1$ 的方式对应。
% P02-E0553
\end{lemma}

\begin{proof}
首先注意，$C$ 的极小素理想 $\mathfrak r$ 在 $A$ 中的像是
极小素理想 $\mathfrak p$，在 $B$ 中的像是极小素理想
$\mathfrak q$，因为环映射 $A \to C$ 与 $B \to C$ 平坦
（由平坦环映射的下降定理，《代数》引理
\ref{algebra-lemma-flat-going-down}）。
% P02-E0554
因此只须证明
% P02-E0555
$(A/\mathfrak p \otimes_k B/\mathfrak q)_{
\mathfrak m_A \otimes_k B + A \otimes_k \mathfrak m_B}$
的严格 Hensel 化有唯一的极小素理想。由《代数》引理
\ref{algebra-lemma-quotient-strict-henselization}，环
$A/\mathfrak p$、$B/\mathfrak q$ 严格 Hensel。因此可假设
$A$ 与 $B$ 是严格 Hensel 局部整环；我们的目标是证明 $C$
有唯一的极小素理想。由引理
\ref{more-algebra-lemma-geometrically-unibranch}，$A$ 在其分式域中的整闭包
$A'$ 是剩余域为 $k$ 的正规局部整环。$B$ 在其分式域中的整闭包
$B'$ 也一样。
% P02-E0556
由《代数》引理
\ref{algebra-lemma-geometrically-normal-tensor-normal}，
$A' \otimes_k B'$ 是正规环。因此其局部化
% P02-E0558
$$
R = (A' \otimes_k B')_{
\mathfrak m_{A'} \otimes_k B' + A' \otimes_k \mathfrak m_{B'}}
$$
是正规局部整环。注意 $A \otimes_k B \to A' \otimes_k B'$
是整映射（因而上升定理成立——《代数》引理
\ref{algebra-lemma-integral-going-up}），并且
% P02-E0557
$\mathfrak m_{A'} \otimes_k B' + A' \otimes_k \mathfrak m_{B'}$
是 $A' \otimes_k B'$ 中位于
$\mathfrak m_A \otimes_k B + A \otimes_k \mathfrak m_B$
之上的唯一极大理想。因此，由《代数》引理
\ref{algebra-lemma-unique-prime-over-localize-below}，
% P02-E0558
$$
R = (A' \otimes_k B')_{
\mathfrak m_A \otimes_k B + A \otimes_k \mathfrak m_B}
$$
。由此，映射
% P02-E0558
$$
(A \otimes_k B)_{
\mathfrak m_A \otimes_k B + A \otimes_k \mathfrak m_B}
\longrightarrow
R
$$
是整映射。因此，$R$ 是
$(A \otimes_k B)_{
\mathfrak m_A \otimes_k B + A \otimes_k \mathfrak m_B}$
在其分式域中的整闭包；再次应用引理
\ref{more-algebra-lemma-geometrically-unibranch}，可知 $C$ 有唯一的素理想。
% P02-E0559
\end{proof}
\section{完备化的分支}
\label{more-algebra-section-branches-completion}

\noindent
设 $(A, \mathfrak m)$ 为 Noether 局部环。考虑映射
$A \to A^h \to A^\wedge$。一般而言，映射
$A^h \to A^\wedge$ 未必诱导极小素理想之间的双射；见《例》一书
第 \ref{examples-section-bad} 节。换言之，$A$ 的分支数
（定义见 \ref{more-algebra-definition-number-of-branches}）可能不同于
$A^\wedge$ 的分支数。不过，在某些条件下两者相同；例如，
$A$ 的维数为 $1$ 时即如此。

\begin{lemma}
\label{more-algebra-lemma-nr-branches-completion}
设 $(A, \mathfrak m)$ 为 Noether 局部环。
\begin{enumerate}
\item 映射 $A^h \to A^\wedge$ 定义了从 $A^\wedge$ 的极小素理想
到 $A^h$ 的极小素理想的满射。
\item $A$ 的分支数不大于 $A^\wedge$ 的分支数。
\item $A$ 的几何分支数不大于 $A^\wedge$ 的几何分支数。
\end{enumerate}
\end{lemma}

\begin{proof}
由引理 \ref{more-algebra-lemma-henselization-noetherian}，映射
$A^h \to A^\wedge$ 平坦且单射。结合下降定理（《代数》引理
\ref{algebra-lemma-flat-going-down}）与《代数》引理
\ref{algebra-lemma-injective-minimal-primes-in-image}，可得 (1)。
(2) 由此、定义 \ref{more-algebra-definition-number-of-branches} 以及
$A^\wedge$ 是 Hensel 环这一事实得出（《代数》引理
\ref{algebra-lemma-complete-henselian}）。由引理
\ref{more-algebra-lemma-henselization-noetherian}，有
$(A^\wedge)^{sh} = A^{sh} \otimes_{A^h} A^\wedge$。因此，使用平坦
单射映射 $A^{sh} \to (A^\wedge)^{sh}$ 重复上述论证，即得 (3)。
\end{proof}

\begin{lemma}
\label{more-algebra-lemma-equal-nr-branches-completion}
设 $(A, \mathfrak m)$ 为 Noether 局部环。$A$ 的分支数等于
$A^\wedge$ 的分支数，当且仅当对 Hensel 化的每个极小素理想
$\mathfrak q \subset A^h$，$\sqrt{\mathfrak qA^\wedge}$ 都是素理想。
\end{lemma}

\begin{proof}
这由引理 \ref{more-algebra-lemma-nr-branches-completion} 以及如下事实得出：
由《代数》引理
\ref{algebra-lemma-Noetherian-irreducible-components} 及
$A^h$ 与 $A^\wedge$ 都是 Noether 环这一事实（引理
\ref{more-algebra-lemma-henselization-noetherian}），$A$ 与 $A^\wedge$
都只有有限多个分支。
% P02-E0560
\end{proof}

\noindent
下面是一个简单的黏合引理。

\begin{lemma}
\label{more-algebra-lemma-glueing-sum-components-open}
设 $A$ 为环，$I$ 为有限生成理想。设 $A \to C$ 为环映射，
使得对每个 $f \in I$，环映射 $A_f \to C_f$ 都是关于某个幂等元
的局部化。则存在满射 $A \to C'$，使得对每个 $f \in I$，
$A_f \to (C \times C')_f$ 都是同构。
\end{lemma}

\begin{proof}
选取 $I$ 的一组生成元 $f_1, \ldots, f_r$。写成
$$
C_{f_i} = (A_{f_i})_{e_i}
$$
，其中 $e_i \in A_{f_i}$ 为某个幂等元。对某个
$a_i \in A$ 与 $n \geq 0$，写 $e_i = a_i/f_i^n$；对所有
$i = 1, \ldots, r$ 可以使用同一个 $n$。取适当的
$m \gg 0$，用 $f_i^ma_i$ 替换 $a_i$、用 $n + m$ 替换 $n$
之后，可以假设对所有 $i$ 都有 $a_i^2 = f_i^n a_i$。由于
$e_i$ 在 $C_{f_if_j} = (A_{f_if_j})_{e_j} = A_{f_if_ja_j}$
中的像为 $1$，可见对某个 $N$ 有
$$
(f_if_ja_j)^N(f_j^n a_i  - f_i^na_j) = 0
$$
（对所有序对 $i, j$ 可以选取同一个 $N$）。利用
$a_j^2 = f_j^na_j$，得到
% P02-E0561
$$
f_i^{N + n} f_j^{N + nN} a_j = f_i^N f_j^{N + n} a_ia_j^N
$$
。把 $n$ 增大为 $n + N + nN$ 并用 $f_i^{N + nN}a_i$
替换 $a_i$ 后，可见对所有序对 $i, j$，$f_i^n a_j$
都属于由 $a_i$ 生成的理想。令
$C' = A/(a_1, \ldots, a_r)$。则
$$
C'_{f_i} = A_{f_i}/(a_i) = A_{f_i}/(e_i)
$$
，因为反演 $f_i$ 后，$a_j$ 属于由 $a_i$ 生成的理想。对环 $B$
的幂等元 $e$，有 $B = B_e \times B/(e)$，故当 $f$ 等于
$f_1, \ldots, f_r$ 之一时，引理的结论成立。利用函数的黏合，
具体使用《代数》引理 \ref{algebra-lemma-cover}，可知结论对所有
$f \in I$ 都成立。事实上，对 $f \in I$，元素
$f_1, \ldots, f_r$ 在 $A_f$ 中生成单位理想，故
$A_f \to (C \times C')_f$ 是同构，当且仅当在分别关于
$f_1, \ldots, f_r$ 局部化后它是同构。
\end{proof}

\noindent
可以利用引理 \ref{more-algebra-lemma-quotient-by-idempotent}，从形式完备化的
给定有限型扩张构造有限型扩张。我们将在《形式空间的代数化》引理
\ref{restricted-lemma-approximate-by-etale-over-complement}
中推广这一引理。

\begin{lemma}
\label{more-algebra-lemma-quotient-by-idempotent}
设 $A$ 为 Noether 环，$I$ 为理想，$B$ 为有限型 $A$-代数。
设 $B^\wedge \to C$ 为核是 $J$ 的满射，其中 $B^\wedge$
是 $I$-进完备化。若对某个 $c \geq 0$，$J/J^2$ 被 $I^c$ 零化，
则 $C$ 同构于某个有限型 $A$-代数的完备化。
\end{lemma}

\begin{proof}
令 $f \in I$。由于 $B^\wedge$ 是 Noether 环（《代数》引理
\ref{algebra-lemma-completion-Noetherian-Noetherian}），
$J$ 是有限生成理想。因此，由《代数》引理
\ref{algebra-lemma-ideal-is-squared-union-connected} 可知
$$
C_f = ((B^\wedge)_f)_e
$$
，其中 $e \in (B^\wedge)_f$ 是某个幂等元。由引理
\ref{more-algebra-lemma-glueing-sum-components-open}，可以找到满射
$B^\wedge \to C'$，使得反演任意 $f \in I$ 后，
$B^\wedge \to C \times C'$ 都成为同构。注意
$C \times C'$ 是有限 $B^\wedge$-代数。

\medskip\noindent
选取生成元 $f_1, \ldots, f_r \in I$。以
$\alpha_i : (C \times C')_{f_i} \to B_{f_i} \otimes_B B^\wedge$
表示上面得到的 $(B^\wedge)_{f_i}$-代数同构的逆，以
$\alpha_{ij} : (B_{f_i})_{f_j} \to (B_{f_j})_{f_i}$
表示显然的 $B$-代数同构。
% P02-E0562
考虑注记 \ref{more-algebra-remark-glueing-data} 中所引入的范畴
$\text{Glue}(B \to B^\wedge, f_1, \ldots, f_r)$ 的对象
$$
(C \times C', B_{f_i}, \alpha_i, \alpha_{ij})
$$
。我们略去条件 (1)(a) 与 (1)(b) 的验证。因为
$B \to B^\wedge$ 是平坦映射（《代数》引理
\ref{algebra-lemma-completion-flat}），并诱导同构
$B/IB \to B^\wedge/IB^\wedge$，所以可以应用命题
\ref{more-algebra-proposition-equivalence} 与注记
\ref{more-algebra-remark-formal-glueing-algebras}。由此，对某个有限
$B$-代数 $D$，$C \times C'$ 同构于 $D \otimes_B B^\wedge$。
于是 $D/ID \cong C/IC \times C'/IC'$。令
$\overline{e} \in D/ID$ 为对应于因子 $C/IC$ 的幂等元。由引理
\ref{more-algebra-lemma-lift-idempotent-upstairs}，存在 \'etale 环映射
$B \to B'$，它诱导同构 $B/IB \to B'/IB'$，且
$D' = D \otimes_B B'$ 包含提升 $\overline{e}$ 的幂等元 $e$。
由于 $C \times C'$ 是 $I$-进完备的，对
$(C \times C', IC \times IC')$ 是 Hensel 对（引理
\ref{more-algebra-lemma-complete-henselian}）。因此，可以把映射
$B \to C \times C'$ 分解经过 $B'$（见引理
\ref{more-algebra-lemma-characterize-henselian-pair}）。这样做以后，可以用
$B'$ 替换 $B$、用 $D'$ 替换 $D$（因为这不改变 $I$-进完备化）。
% P02-E0563
于是
$D = D_e \times D_{1 - e} = D/(1 - e) \times D/(e)$
是有限型 $A$-代数的乘积；第一个因子的完备化是 $C$，第二个因子的
完备化是 $C'$。
\end{proof}

\begin{lemma}
\label{more-algebra-lemma-one-dimensional-formal-branch}
设 $(A, \mathfrak m)$ 为 Noether 局部环，其 Hensel 化为 $A^h$。
设 $\mathfrak q \subset A^\wedge$ 为极小素理想，且
$\dim(A^\wedge/\mathfrak q) = 1$。则存在 $A^h$ 的极小素理想
$\mathfrak q^h$，使得
$\mathfrak q = \sqrt{\mathfrak q^hA^\wedge}$。
\end{lemma}

\begin{proof}
由于 $A$ 与 $A^h$ 的完备化相同，可以假设 $A$ 是 Hensel 环
（引理 \ref{more-algebra-lemma-henselization-noetherian}）。我们将对映射
$A^\wedge \to A^\wedge/J$ 应用引理
\ref{more-algebra-lemma-quotient-by-idempotent}，其中
$J = \Ker(A^\wedge \to (A^\wedge)_{\mathfrak q})$。由于
$\dim((A^\wedge)_\mathfrak q) = 0$，对某个 $n$ 有
$\mathfrak q^n \subset J$。所以 $J/J^2$ 被 $\mathfrak q^n$
零化。另一方面，因 $J_\mathfrak q = 0$，有
$(J/J^2)_\mathfrak q = 0$。故 $\mathfrak m$ 是 $J/J^2$
唯一的伴随素理想，从而 $\mathfrak m$ 的某个幂零化 $J/J^2$。
于是该引理适用，得到
$A^\wedge/J = C^\wedge$，其中 $C$ 是某个有限型 $A$-代数。

\medskip\noindent
由于 $A^\wedge/J$ 具有相应性质，有
$C/\mathfrak m C = A/\mathfrak m$。因此
$\mathfrak m_C = \mathfrak m C$ 是极大理想，且
$A \to C$ 在 $\mathfrak m_C$ 处非分歧（《代数》引理
\ref{algebra-lemma-characterize-unramified}）。用一个主局部化替换
$C$ 后，可以假设 $C$ 是某个 \'etale $A$-代数 $B$ 的商；
见《代数》命题
\ref{algebra-proposition-unramified-locally-standard}。不过，
由于 $A \to C_{\mathfrak m_C}$ 的剩余域扩张平凡且 $A$
是 Hensel 环，再作一次局部化后可知 $B = A$。因此对某个理想
$I \subset A$，有 $C = A/I$，从而
$J = IA^\wedge$（因为在当前情形下完备化正合，见《代数》引理
\ref{algebra-lemma-completion-flat}），并且
$I = J \cap A$（由 $A \to A^\wedge$ 的平坦性）。由于
$\mathfrak q^n \subset J \subset \mathfrak q$，可知
$\mathfrak p = \mathfrak q \cap A$ 满足
$\mathfrak p^n \subset I \subset \mathfrak p$。于是
$\sqrt{\mathfrak p A^\wedge} = \mathfrak q$，证明完成。
\end{proof}

\begin{lemma}
\label{more-algebra-lemma-completion-disconnected}
设 $(A, \mathfrak m)$ 为 Noether 局部环。$A^\wedge$ 的去心谱不连通，
当且仅当 $A^h$ 的去心谱不连通。
\end{lemma}

\begin{proof}
由于 $A$ 与 $A^h$ 的完备化相同，可以假设 $A$ 是 Hensel 环
（引理 \ref{more-algebra-lemma-henselization-noetherian}）。

\medskip\noindent
由于 $A \to A^\wedge$ 忠实平坦（见刚才给出的参考文献），
从 $A^\wedge$ 的去心谱到 $A$ 的去心谱的映射是满射
（见《代数》引理 \ref{algebra-lemma-ff-rings}）。因此，若
$A$ 的去心谱不连通，则 $A^\wedge$ 的去心谱也不连通。

\medskip\noindent
假设 $A^\wedge$ 的去心谱不连通。这意味着
% P02-E0564
$$
\Spec(A^\wedge) \setminus \{\mathfrak m^\wedge\} = Z \amalg Z'
$$
，其中 $Z$ 与 $Z'$ 闭。令
% P02-E0565
$\overline{Z}, \overline{Z}' \subset \Spec(A^\wedge)$ 为其闭包。
设对某些理想 $J, J' \subset A^\wedge$ 有
$\overline{Z} = V(J)$、$\overline{Z}' = V(J')$。则
$V(J + J') = \{\mathfrak m^\wedge\}$ 且
$V(JJ') = \Spec(A^\wedge)$。第一个等式意味着
$\mathfrak m^\wedge = \sqrt{J + J'}$，从而对某个
$e \geq 1$ 有
$(\mathfrak m^\wedge)^e \subset J + J'$。第二个等式蕴含
$JJ'$ 的每个元素都幂零，故对某个 $n \geq 1$ 有
$(JJ')^n = 0$。合在一起，这意味着 $J^n/J^{2n}$ 被
$J^n$ 与 $(J')^n$ 零化，因而被
$(\mathfrak m^\wedge)^{2en}$ 零化。
% P02-E0566
于是可以应用引理 \ref{more-algebra-lemma-quotient-by-idempotent}，得到有限型
$A$-代数 $C$ 与同构 $A^\wedge/J^n = C^\wedge$。

\medskip\noindent
证明的其余部分与引理
\ref{more-algebra-lemma-one-dimensional-formal-branch} 证明的第二部分完全相同；
当然，该引理正是本引理的一个特殊情形！由于 $A^\wedge/J^n$
具有相应性质，有 $C/\mathfrak m C = A/\mathfrak m$。因此
$\mathfrak m_C = \mathfrak m C$ 是极大理想，且
$A \to C$ 在 $\mathfrak m_C$ 处非分歧（《代数》引理
\ref{algebra-lemma-characterize-unramified}）。用一个主局部化替换
$C$ 后，可以假设 $C$ 是某个 \'etale $A$-代数 $B$ 的商；
见《代数》命题
\ref{algebra-proposition-unramified-locally-standard}。不过，
由于 $A \to C_{\mathfrak m_C}$ 的剩余域扩张平凡且 $A$
是 Hensel 环，再作一次局部化后可知 $B = A$。因此对某个理想
$I \subset A$，有 $C = A/I$，从而
$J^n = IA^\wedge$（因为在当前情形下完备化正合，见《代数》引理
\ref{algebra-lemma-completion-flat}），并且
$I = J^n \cap A$（由 $A \to A^\wedge$ 的平坦性）。
由对称性，$I' = (J')^n \cap A$ 满足 $(J')^n = I'A^\wedge$。
于是 $\mathfrak m^e \subset I + I'$ 且 $II' = 0$，从而
$V(I)$ 与 $V(I')$ 是闭子概形，并给出 $A$ 的去心谱所需的不交并分解。
% P02-E0567
\end{proof}

\begin{lemma}
\label{more-algebra-lemma-one-dimensional-number-of-branches}
设 $(A, \mathfrak m)$ 为维数 $1$ 的 Noether 局部环。则
$A$ 与 $A^\wedge$ 的（几何）分支数相同。
\end{lemma}

\begin{proof}
为说明分支数的结论，结合引理
\ref{more-algebra-lemma-nr-branches-completion}、
\ref{more-algebra-lemma-equal-nr-branches-completion} 与
\ref{more-algebra-lemma-one-dimensional-formal-branch}，并使用
$A^\wedge$ 的维数为一这一事实；见引理
\ref{more-algebra-lemma-completion-dimension}。为说明几何分支数的结论，
使用刚证明的分支数结论、严格 Hensel 化不改变维数这一事实
（引理 \ref{more-algebra-lemma-henselization-dimension}），以及由引理
\ref{more-algebra-lemma-henselization-noetherian} 得到的
$(A^{sh})^\wedge = ((A^\wedge)^{sh})^\wedge$。
\end{proof}
\begin{lemma}
\label{more-algebra-lemma-geometrically-normal-formal-fibres-number-of-branches}
\begin{reference}
\cite[Theorem 2.3]{Beddani}
\end{reference}
设 $(A, \mathfrak m)$ 为 Noether 局部环。若 $A$ 的形式纤维
几何正规（例如 $A$ 优秀或拟优秀），则 $A$ 是 Nagata 环，
并且 $A$ 与 $A^\wedge$ 的（几何）分支数相同。
\end{lemma}

\begin{proof}
由于正规环既约，由引理 \ref{more-algebra-lemma-Nagata-local-ring}，
$A$ 是 Nagata 环。在证明的其余部分，我们将使用引理
\ref{more-algebra-lemma-formal-fibres-normal}、命题
\ref{more-algebra-proposition-finite-type-over-P-ring} 与引理
\ref{more-algebra-lemma-check-P-ring-maximal-ideals}。这些结果告诉我们，
$A$ 是 P-环，其中 $P(k \to R) = $``$R$ is geometrically
normal over $k$''，而且任意（本质）有限型 $A$-代数也具有同样性质。

\medskip\noindent
令 $\mathfrak q \subset A$ 为极小素理想。则
$A^\wedge/\mathfrak q A^\wedge = (A/\mathfrak q)^\wedge$
且 $A^h/\mathfrak qA^h = (A/\mathfrak q)^h$
（《代数》引理 \ref{algebra-lemma-quotient-henselization}）。
因此，$A$ 的分支数是各环 $A/\mathfrak q$ 的分支数之和；
$A^\wedge$ 也同样。这样便归约到 $A$ 是整环的情形。

\medskip\noindent
假设 $A$ 是整环。令 $A'$ 为 $A$ 在 $A$ 的分式域 $K$ 中的整闭包。
由于 $A$ 是 Nagata 环，$A \to A'$ 有限。回忆 $A$ 的分支数等于
$A'$ 的极大理想 $\mathfrak m'$ 的数目（引理
\ref{more-algebra-lemma-branches}）。还回忆
% P02-E0568
$$
(A')^\wedge = A' \otimes_A A^\wedge =
\prod\nolimits_{\mathfrak m' \subset A'} (A'_{\mathfrak m'})^\wedge
$$
，见《代数》引理
\ref{algebra-lemma-completion-finite-extension}。因为
$A'_{\mathfrak m'}$ 是形式纤维几何正规的局部环，由引理
\ref{more-algebra-lemma-completion-normal-local-ring}，
$(A'_{\mathfrak m'})^\wedge$ 正规。因此
$A' \otimes_A A^\wedge$ 的极小素理想与上述分解中的因子按
$1$ 对 $1$ 的方式对应。由 $A \to A^\wedge$ 的平坦性，有
% P02-E0568
$$
A^\wedge \subset A' \otimes_A A^\wedge \subset K \otimes_A A^\wedge
$$
。由于左侧环与右侧环具有同一个极小素理想集合，中间的环也如此
（略去一个小细节），这就完成了证明。
% P02-E0569

\medskip\noindent
为说明几何分支数的结论，使用刚证明的分支数结论、
$A^{sh}$ 的形式纤维几何正规这一事实（引理
\ref{more-algebra-lemma-formal-fibres-normal} 与
\ref{more-algebra-lemma-henselization-P-ring}），以及由引理
\ref{more-algebra-lemma-henselization-noetherian} 得到的
$(A^{sh})^\wedge = ((A^\wedge)^{sh})^\wedge$。
\end{proof}
\section{形式链环}
\label{more-algebra-section-formally-catenary}

\noindent
本节证明 Ratliff 的定理 \cite{Ratliff}：Noether 局部环是普遍链环，
当且仅当它是形式链环。

\begin{definition}
\label{more-algebra-definition-formally-catenary}
若对每个极小素理想 $\mathfrak p \subset A$，
$A^\wedge/\mathfrak p A^\wedge$ 的谱都等维，则称 Noether
局部环 $A$ 为{\it 形式链环}。
\end{definition}

\noindent
设 $A$ 为形式链 Noether 局部环。由 Ratliff 的结果（命题
\ref{more-algebra-proposition-ratliff}），$A$ 的任意商也形式链
（因为普遍链环这一类对取商稳定）。于是，对 $A$ 的每个素理想
$\mathfrak p$，$A^\wedge/\mathfrak p A^\wedge$ 的谱都等维。

\begin{lemma}
\label{more-algebra-lemma-not-formally-catenary}
设 $(A, \mathfrak m)$ 为非形式链的 Noether 局部环。则 $A$
不是普遍链环。
\end{lemma}

\begin{proof}
由假设，存在极小素理想 $\mathfrak p \subset A$，使
$A^\wedge /\mathfrak p A^\wedge$ 的谱不等维。用
$A/\mathfrak p$ 替换 $A$ 后，可以假设 $A$ 是整环，且
$A^\wedge$ 的谱不等维。令 $\mathfrak q$ 为 $A^\wedge$
的一个极小素理想，使 $d = \dim(A^\wedge/\mathfrak q)$ 最小；
于是 $0 < d < \dim(A)$。我们对 $d$ 作归纳以证明本引理。

\medskip\noindent
先考虑 $d = 1$ 的情形。此时
$\dim(A^\wedge_\mathfrak q) = 0$。故
$A^\wedge_\mathfrak q$ 是 Artin 局部环，从而对某个 $n > 0$，
理想 $J = \mathfrak q^n$ 在 $A^\wedge_\mathfrak q$ 中的像为零。
由此，$\mathfrak m$ 是 $J/J^2$ 唯一的伴随素理想，所以对某个
$m > 0$，$\mathfrak m^m$ 零化 $J/J^2$。因此，可以利用引理
\ref{more-algebra-lemma-quotient-by-idempotent} 找到有限型的 $A \to B$，
使得 $B^\wedge \cong A^\wedge/J$。
% P02-E0570
于是 $\mathfrak m_B = \sqrt{\mathfrak mB}$ 是极大理想，
其剩余域与 $\mathfrak m$ 的剩余域相同，且 $B^\wedge$
是 $\mathfrak m_B$-进完备化（《代数》引理
\ref{algebra-lemma-finite-after-completion}）。于是
$$
\dim(B_{\mathfrak m_B}) = \dim(B^\wedge) = 1 = d.
$$
由于有分解 $A \to B \to A^\wedge/J$，$\mathfrak q/J$
的逆像是素理想 $\mathfrak q' \subset \mathfrak m_B$，它位于
$A$ 的 $(0)$ 之上。因此，若 $A$ 是普遍链环，则维数公式
（《代数》引理 \ref{algebra-lemma-dimension-formula}）将给出
\begin{align*}
\dim(B_{\mathfrak m_B})
& \geq
\dim((B/\mathfrak q')_{\mathfrak m_B}) \\
& =
\dim(A) + \text{trdeg}_A(B/\mathfrak q') -
\text{trdeg}_{\kappa(\mathfrak m)}(\kappa(\mathfrak m_B)) \\
& =
\dim(A) + \text{trdeg}_A(B/\mathfrak q')
\end{align*}
。这个矛盾完成了 $d = 1$ 情形的论证。
% P02-E0571

\medskip\noindent
假设 $d > 1$。令 $Z \subset \Spec(A^\wedge)$ 为除
$V(\mathfrak q)$ 外其余不可约分支之并。令
$\mathfrak r_1, \ldots, \mathfrak r_m \subset A^\wedge$
为对应于 $V(\mathfrak q) \cap Z$ 中维数 $> 0$ 的不可约分支的
素理想。利用素理想回避（《代数》引理
\ref{algebra-lemma-silly}），选取
$f \in \mathfrak m$、$f \not \in A \cap \mathfrak r_j$。
则 $\dim(A/fA) = \dim(A) - 1$，且 $V(\mathfrak q, f)$
有某个维数为 $d - 1$ 的不可约分支。因此 $A/fA$ 不是形式链环，
且不变量 $d$ 已减小。由归纳假设，$A/fA$ 不是普遍链环，
故 $A$ 也不是普遍链环。
\end{proof}

\begin{lemma}
\label{more-algebra-lemma-flat-under-catenary-equidimensional}
设 $A \to B$ 为局部 Noether 环之间的平坦局部环映射。假设
$B$ 是链环，且 $\Spec(B)$ 等维。则：
% P02-E0572
\begin{enumerate}
\item 对所有 $\mathfrak p \subset A$，
$\Spec(B/\mathfrak p B)$ 都等维；
\item $A$ 是链环且 $\Spec(A)$ 等维。
\end{enumerate}
\end{lemma}

\begin{proof}
令 $\mathfrak p \subset A$ 为素理想，令 $\mathfrak q \subset B$
为 $\mathfrak pB$ 上的极小素理想。由 $A \to B$ 的下降定理
（《代数》引理 \ref{algebra-lemma-flat-going-down}），
$\mathfrak q \cap A = \mathfrak p$。故
$A_\mathfrak p \to B_\mathfrak q$ 是特殊纤维维数为 $0$
的平坦局部环映射，因此
$$
\dim(A_\mathfrak p) = \dim(B_\mathfrak q) = \dim(B) - \dim(B/\mathfrak q)
$$
（《代数》引理
\ref{algebra-lemma-dimension-base-fibre-equals-total}）。
第二个等式成立，是因为 $\Spec(B)$ 等维且 $B$ 是链环。
% P02-E0573
所以 $\dim(B/\mathfrak q)$ 与 $\mathfrak q$ 的选择无关，从而
$\Spec(B/\mathfrak p B)$ 等维，其维数为
$\dim(B) - \dim(A_\mathfrak p)$。另一方面，由平坦性
（见刚引用的引理），有
$\dim(B/\mathfrak p B) = \dim(A/\mathfrak p) + \dim(B/\mathfrak m_A B)$
以及
$\dim(B) = \dim(A) + \dim(B/\mathfrak m_A B)$，故得到
% P02-E0574
$$
\dim(A_\mathfrak p) = \dim(A) - \dim(A/\mathfrak p)
$$
，这对 $A$ 中所有 $\mathfrak p$ 都成立。把它应用于 $A$
的所有极小素理想，可见 $A$ 等维。若
$\mathfrak p \subset \mathfrak p'$ 是中间没有素理想的严格包含，
则可把上述结论应用于 $A/\mathfrak p$ 中的素理想
$\mathfrak p'/\mathfrak p$，因为
$A/\mathfrak p \to B/\mathfrak p B$ 平坦且
$\Spec(B/\mathfrak p B)$ 等维，从而得到
% P02-E0575
$$
1 = \dim((A/\mathfrak p)_{\mathfrak p'}) =
\dim(A/\mathfrak p) - \dim(A/\mathfrak p')
$$
。所以 $\mathfrak p \mapsto \dim(A/\mathfrak p)$ 是维数函数，
由此 $A$ 是链环。
\end{proof}

\begin{lemma}
\label{more-algebra-lemma-formally-catenary}
设 $A$ 为形式链 Noether 局部环。则 $A$ 是普遍链环。
\end{lemma}

\begin{proof}
可以用 $A/\mathfrak p$ 替换 $A$，其中 $\mathfrak p$ 是 $A$
的极小素理想；见《代数》引理
\ref{algebra-lemma-catenary-check-irreducible}。因此，可以假设
$A^\wedge$ 的谱等维。只须证明 $A$ 上本质有限型的每个局部环都是
链环（例如见《代数》引理
\ref{algebra-lemma-catenary-check-local}）。所以只须证明
$A[x_1, \ldots, x_n]_\mathfrak m$ 是链环，其中
$\mathfrak m \subset A[x_1, \ldots, x_n]$ 是位于
$\mathfrak m_A$ 之上的极大理想；见《代数》引理
\ref{algebra-lemma-localization-at-closed-point-special-fibre}
（以及《代数》引理 \ref{algebra-lemma-quotient-catenary} 与
\ref{algebra-lemma-localization-catenary}）。令
$\mathfrak m' \subset A^\wedge[x_1, \ldots, x_n]$ 为位于
$\mathfrak m$ 之上的唯一极大理想。则
% P02-E0576
$$
A[x_1, \ldots, x_n]_\mathfrak m \to A^\wedge[x_1, \ldots, x_n]_{\mathfrak m'}
$$
是局部平坦映射（《代数》引理
\ref{algebra-lemma-completion-flat}）。因此，由引理
\ref{more-algebra-lemma-flat-under-catenary-equidimensional}，只须证明右侧的环
是谱等维的链环。它是链环，因为完备局部环都是普遍链环
（《代数》注记
\ref{algebra-remark-Noetherian-complete-local-ring-universally-catenary}）。
任取 $A^\wedge[x_1, \ldots, x_n]_{\mathfrak m'}$ 的极小素理想
$\mathfrak q$。则对 $A^\wedge$ 的某个极小素理想 $\mathfrak p$，
有
$\mathfrak q = \mathfrak p A^\wedge[x_1, \ldots, x_n]_{\mathfrak m'}$
（略去一个小细节）。所以
% P02-E0577
$$
\dim(A^\wedge[x_1, \ldots, x_n]_{\mathfrak m'}/\mathfrak q) =
\dim(A^\wedge/\mathfrak p) + n = \dim(A^\wedge) + n
$$
；第一个等式由《代数》引理
\ref{algebra-lemma-dimension-base-fibre-equals-total} 得到，
第二个等式成立是因为 $A^\wedge$ 的谱等维。这就完成了证明。
\end{proof}

\begin{proposition}[Ratliff]
\label{more-algebra-proposition-ratliff}
\begin{reference}
\cite{Ratliff}
\end{reference}
Noether 局部环是普遍链环，当且仅当它是形式链环。
\end{proposition}

\begin{proof}
结合引理 \ref{more-algebra-lemma-not-formally-catenary} 与
\ref{more-algebra-lemma-formally-catenary}。
\end{proof}

\begin{lemma}
\label{more-algebra-lemma-geometrically-normal-fibres-universally-catenary}
\begin{reference}
\cite[Corollary 2.3]{Heinzer-Rotthaus-Wiegand}
\end{reference}
设 $(A, \mathfrak m)$ 为形式纤维几何正规的 Noether 局部环。则：
\begin{enumerate}
\item $A^h$ 是普遍链环；
\item 若 $A$ 单支（例如为正规环时），则 $A$ 是普遍链环。
\end{enumerate}
\end{lemma}

\begin{proof}
由引理
\ref{more-algebra-lemma-geometrically-normal-formal-fibres-number-of-branches}，
$A$ 与 $A^\wedge$ 的分支数相同，故引理
\ref{more-algebra-lemma-equal-nr-branches-completion} 适用。于是，对任意极小素理想
$\mathfrak q \subset A^h$，$A^\wedge/\mathfrak q A^\wedge$
有唯一的极小素理想。因此，$A^h$ 依定义是形式链环，进而由命题
\ref{more-algebra-proposition-ratliff} 是普遍链环。若 $A$ 单支，则 $A^h$
有唯一的极小素理想，故 $A^\wedge$ 有唯一的极小素理想，
于是 $A$ 是形式链环；同理可得结论。
\end{proof}
\section{群作用与整闭包}
\label{more-algebra-section-group-actions-integral}

\noindent
在某种意义下，本节是《代数》第
\ref{algebra-section-going-down-integral-over-normal} 节的延续。
在《基本群》第 \ref{pione-section-group-actions-integral} 节中还可找到
更多类似材料。
% P02-E0578

\begin{lemma}
\label{more-algebra-lemma-pol-lifting}
设 $\varphi : A \to B$ 为环满射，$G$ 为作用在
$\varphi : A \to B$ 上的 $n$ 阶有限群。若 $b \in B^G$，
则存在首一多项式 $P \in A^G[T]$，它在 $B^G[T]$ 中的像为
$(T - b)^n$。
\end{lemma}

\begin{proof}
选取提升 $b$ 的 $a \in A$，并置
$P = \prod_{\sigma \in G} (T - \sigma(a))$。
\end{proof}

\begin{lemma}
\label{more-algebra-lemma-invariants-modulo}
设 $R$ 为环，$G$ 为作用在 $R$ 上的有限群。设 $I \subset R$
为理想，使得对所有 $\sigma \in G$ 都有 $\sigma(I) \subset I$。
则 $R^G/I^G \subset (R/I)^G$ 是整环扩张；它在谱上诱导同胚，
并诱导剩余域的纯不可分扩张。
\end{lemma}

\begin{proof}
由于 $I^G = R^G \cap I$，该映射显然单射。引理
\ref{more-algebra-lemma-pol-lifting} 表明《代数》引理
\ref{algebra-lemma-universally-bijective} 适用。
\end{proof}

\begin{lemma}
\label{more-algebra-lemma-pol-in-ideal}
设 $G$ 为作用在环 $R$ 上的 $n$ 阶有限群，$J \subset R^G$
为理想。对 $x \in JR$，有
$\prod_{\sigma \in G} (T - \sigma(x)) =
T^n + a_1 T^{n - 1} + \ldots + a_n$，其中 $a_i \in J$。
\end{lemma}

\begin{proof}
注意该多项式确实首一，且系数属于 $R^G$。可以写
$x = f_1 b_1 + \ldots + f_m b_m$，其中 $f_j \in J$、
$b_j \in R$。因此，对 $m$ 作归纳，可以假设
$x = y - fb$，其中 $f \in J$、$b \in R$ 且 $y \in JR$，
并且结论对 $y$ 成立。于是
$$
\prod\nolimits_{\sigma \in G} (T - \sigma(x)) =
\prod\nolimits_{\sigma \in G} (T - \sigma(y) + f\sigma(b)) =
\prod\nolimits_{\sigma \in G} (T - \sigma(y)) +
\sum_{i = 1, \ldots, n} f^i a_i
$$
，其中
% P02-E0579
$$
a_i =
\sum\nolimits_{S \subset G,\ |S| = i}
\prod\nolimits_{\sigma \in S} \sigma(b)
\prod\nolimits_{\sigma \not \in S} (T - \sigma(y))
$$
。我们略去的一个计算表明 $a_i \in R^G$（提示：所给表达式是
对称的）。因此，引理陈述中关于 $x$ 的多项式模 $J$
同余于关于 $y$ 的多项式，这就完成了归纳步骤。
\end{proof}

\begin{lemma}
\label{more-algebra-lemma-invariants-modulo-bis}
设 $R$ 为环，$G$ 为作用在 $R$ 上的 $n$ 阶有限群，
$J \subset R^G$ 为理想。则 $R^G/J \to (R/JR)^G$ 是环映射，
并且：
% P02-E0580
\begin{enumerate}
\item 对 $b \in (R/JR)^G$，存在首一多项式
$P \in R^G/J[T]$，它在 $(R/JR)^G[T]$ 中的像为 $(T - b)^n$；
\item 对 $a \in \Ker(R^G/J \to (R/JR)^G)$，在
$R^G/J[T]$ 中有 $(T - a)^n = T^n$。
\end{enumerate}
特别地，$R^G/J \to (R/JR)^G$ 是整环映射；它在谱上诱导同胚，
并诱导剩余域的纯不可分扩张。
\end{lemma}

\begin{proof}
(1) 由取 $I = JR$ 的引理 \ref{more-algebra-lemma-pol-lifting} 得出。
% P02-E0581
若 $a$ 如 (2) 中所述，则 $a$ 是某个
$x \in R^G \cap JR$ 的像。因此
$(T - x)^n = \prod_{\sigma \in G} (T - \sigma(x))$
由引理 \ref{more-algebra-lemma-pol-in-ideal} 模 $J$ 同余于 $T^n$。
这就证明了 (2)。为得到最后的陈述，可以应用《代数》引理
\ref{algebra-lemma-universally-bijective}。
\end{proof}

\begin{remark}
\label{more-algebra-remark-when-iso}
在引理 \ref{more-algebra-lemma-invariants-modulo-bis} 中可见，若 $n$ 在 $R$
中可逆，则映射 $R^G/J \to (R/JR)^G$ 是同构。
\end{remark}

\begin{lemma}
\label{more-algebra-lemma-functor-invariants-tensor}
设 $R$ 为环，$G$ 为作用在 $R$ 上的 $n$ 阶有限群，
$A$ 为 $R^G$-代数。
\begin{enumerate}
\item 对 $b \in (A \otimes_{R^G} R)^G$，存在首一多项式
$P \in A[T]$，它在 $(A \otimes_{R^G} R)^G[T]$ 中的像为
$(T - b)^n$；
\item 对 $a \in \Ker(A \to (A \otimes_{R^G} R)^G)$，在
$A[T]$ 中有 $(T - a)^n = T^n$。
\end{enumerate}
\end{lemma}

\begin{proof}
选取满射 $E \to A$，其中 $E$ 是 $R^G$ 上的多项式代数。由于
$E$ 是自由 $R^G$-模，有
$(E \otimes_{R^G} R)^G = E$。以
$J = \Ker(E \to A)$ 表示其核。
% P02-E0582
由于张量积右正合，$A \otimes_{R^G} R$ 是
$E \otimes_{R^G} R$ 关于 $J$ 生成的理想的商。这样可见，
本引理是引理 \ref{more-algebra-lemma-invariants-modulo-bis} 的特殊情形。
\end{proof}

\begin{lemma}
\label{more-algebra-lemma-base-change-invariants}
设 $R$ 为环，$G$ 为作用在 $R$ 上的有限群，$R^G \to A$
为环映射。映射
$$
A \to (A \otimes_{R^G} R)^G
$$
在 $R^G \to A$ 平坦时是同构。一般而言，该映射是整的，在谱上
诱导同胚，并诱导纯不可分的剩余域扩张。
\end{lemma}

\begin{proof}
为证明第一项陈述，考虑正合列
$0 \to R^G \to R \to \bigoplus_{\sigma \in G} R$，其中第二个映射
把 $x$ 送到 $(\sigma(x) - x)_{\sigma \in G}$。若
$R^G \to A$ 平坦，与 $A$ 张量后该序列仍正合。因此，$A$
是 $(A \otimes_{R^G} R)^G$ 中的 $G$-不变量。
% P02-E0583

\medskip\noindent
第二项陈述由引理 \ref{more-algebra-lemma-functor-invariants-tensor} 与《代数》引理
\ref{algebra-lemma-universally-bijective} 得出。
\end{proof}

\begin{lemma}
\label{more-algebra-lemma-one-orbit}
设 $G$ 为作用在环 $R$ 上的有限群。对位于 $R^G$ 中同一个素理想
之上的任意两个素理想 $\mathfrak q, \mathfrak q' \subset R$，
存在 $\sigma \in G$，使得
$\sigma(\mathfrak q) = \mathfrak q'$。
\end{lemma}

\begin{proof}
扩张 $R^G \subset R$ 是整的，因为每个 $x \in R$ 都是
$R^G[T]$ 中首一多项式
$\prod_{\sigma \in G}(T - \sigma(x))$ 的根。因此，位于给定素理想
$\mathfrak p$ 之上的素理想之间不存在包含关系（《代数》引理
\ref{algebra-lemma-integral-no-inclusion}）。若本引理不成立，
则可选取 $x \in \mathfrak q'$，使得对所有 $\sigma \in G$
都有 $x \not \in \sigma(\mathfrak q)$；见《代数》引理
\ref{algebra-lemma-silly}。于是
$y = \prod_{\sigma \in G} \sigma(x)$ 属于 $R^G$，并属于
$\mathfrak p = R^G \cap \mathfrak q'$。另一方面，对所有
$\sigma$ 都有 $x \not \in \sigma(\mathfrak q)$，意味着对所有
$\sigma$ 都有 $\sigma(x) \not \in \mathfrak q$。由于
$\mathfrak q$ 是素理想，$y \not \in \mathfrak q$。
这不可能，因为 $y \in \mathfrak p \subset \mathfrak q$。
\end{proof}

\begin{lemma}
\label{more-algebra-lemma-one-orbit-geometric}
设 $G$ 为作用在环 $R$ 上的有限群。设素理想
$\mathfrak q \subset R$ 位于 $\mathfrak p \subset R^G$ 之上。
则 $\kappa(\mathfrak q)/\kappa(\mathfrak p)$ 是代数正规扩张，
并且映射
$$
D = \{\sigma \in G \mid \sigma(\mathfrak q) = \mathfrak q\}
\longrightarrow
\text{Aut}(\kappa(\mathfrak q)/\kappa(\mathfrak p))
$$
满射\footnote{回忆我们只在 Galois 扩张的情形使用记号
$\text{Gal}$。}。
\end{lemma}

\begin{proof}
取 $A = (R^G)_\mathfrak p$ 与 $B = A \otimes_{R^G} R$。由于局部化
平坦，由引理 \ref{more-algebra-lemma-base-change-invariants}，
$A = B^G$。注意 $\mathfrak pA$ 与 $\mathfrak qB$ 是素理想，
$D$ 是 $\mathfrak qB$ 的稳定子，并且
$\kappa(\mathfrak p) = \kappa(\mathfrak pA)$、
$\kappa(\mathfrak q) = \kappa(\mathfrak qB)$。因此可以用 $B$
替换 $R$，并假设 $\mathfrak p$ 是极大理想。由于
$R^G \subset R$ 是整环扩张，$R$ 的极大理想恰为位于
$\mathfrak p$ 之上的素理想（这由《代数》引理
\ref{algebra-lemma-integral-no-inclusion} 与
\ref{algebra-lemma-integral-going-up} 得出）。由引理
\ref{more-algebra-lemma-one-orbit}，它们只有有限多个，记作
$\mathfrak q = \mathfrak q_1, \mathfrak q_2, \ldots, \mathfrak q_m$，
且构成 $G$ 的一个轨道。由中国剩余定理（《代数》引理
\ref{algebra-lemma-chinese-remainder}），映射
$R \to \prod_{j = 1, \ldots, m} R/\mathfrak q_j$ 满射。

\medskip\noindent
首先证明该扩张正规。取元素
$\alpha \in \kappa(\mathfrak q)$。我们须证明 $\alpha$
在 $\kappa(\mathfrak p)$ 上的极小多项式 $P$ 完全分裂。
由上述结论，可选取
$a \in \mathfrak q_2 \cap \ldots \cap \mathfrak q_m$，使它在
$\kappa(\mathfrak q)$ 中的像为 $\alpha$。考虑 $R^G[T]$
中的多项式 $Q = \prod_{\sigma \in G} (T - \sigma(a))$。
$Q$ 在 $R[T]$ 中的像完全分裂为一次因子，故其在
$\kappa(\mathfrak q)[T]$ 中的像也如此。由于 $P$ 整除 $Q$
在 $\kappa(\mathfrak p)[T]$ 中的像，可知 $P$ 在
$\kappa(\mathfrak q)$ 上完全分裂为一次因子，如所求。

\medskip\noindent
由于 $\kappa(\mathfrak q)/\kappa(\mathfrak p)$ 正规，可以假设
$\kappa(\mathfrak q) = \kappa_1 \otimes_{\kappa(\mathfrak p)} \kappa_2$，
其中 $\kappa_1/\kappa(\mathfrak p)$ 纯不可分，而
$\kappa_2/\kappa(\mathfrak p)$ 为 Galois 扩张；见《域》引理
\ref{fields-lemma-normal-case}。选取 $\alpha \in \kappa_2$：
若该扩张有限，使其在 $\kappa(\mathfrak p)$ 上生成
$\kappa_2$；若该扩张无限，使它生成次数 $> |G|$ 的子域
（以得到矛盾）。
% P02-E0584
由《域》引理 \ref{fields-lemma-primitive-element}，这是可能的。
如前段一样选取 $a$、$P$ 与 $Q$。若
$\alpha' \in \kappa_2$ 是 $\alpha$ 在 $\kappa(\mathfrak p)$ 上的
Galois 共轭元，则
% P02-E0585
$P$ 整除 $P$ 在 $\kappa(\mathfrak p)[T]$ 中的像这一事实表明，
存在 $\sigma \in G$，使得 $\sigma(a)$ 映到 $\alpha'$。
由 $a$ 的选取方式（它在其他极大理想处消失），这蕴含
$\sigma \in D$，且 $\sigma$ 在
$\text{Aut}(\kappa(\mathfrak q)/\kappa(\mathfrak p))$
中的像把 $\alpha$ 映到 $\alpha'$。因此得到所需的满射；
若 $\alpha$ 在 $\kappa(\mathfrak p)$ 上的次数 $> |G|$，
则得到所需的荒谬结论。
% P02-E0586
\end{proof}

\begin{lemma}
\label{more-algebra-lemma-one-orbit-geometric-galois}
设 $A$ 为分式域是 $K$ 的正规整环，$L/K$ 为（可能无限的）
Galois 扩张。令 $G = \text{Gal}(L/K)$，并令 $B$ 为 $A$
在 $L$ 中的整闭包。
\begin{enumerate}
\item 对位于 $A$ 中同一个素理想之上的任意两个素理想
$\mathfrak q, \mathfrak q' \subset B$，存在
$\sigma \in G$，使得 $\sigma(\mathfrak q) = \mathfrak q'$。
\item 设素理想 $\mathfrak q \subset B$ 位于
$\mathfrak p \subset A$ 之上。则
$\kappa(\mathfrak q)/\kappa(\mathfrak p)$ 是代数正规扩张，
并且映射
$$
D = \{\sigma \in G \mid \sigma(\mathfrak q) = \mathfrak q\}
\longrightarrow
\text{Aut}(\kappa(\mathfrak q)/\kappa(\mathfrak p))
$$
满射。
\end{enumerate}
\end{lemma}

\begin{proof}
证明 (1)。考虑序对 $(M, \sigma)$，其中
$K \subset M \subset L$ 是子域，$M/K$ 为 Galois 扩张，
$\sigma \in \text{Gal}(M/K)$，并且
$\sigma(\mathfrak q \cap M) = \mathfrak q' \cap M$。
定义 $(M', \sigma') \geq (M, \sigma)$，当且仅当
$M \subset M'$ 且 $\sigma'|_M = \sigma$。注意
$(K, \text{id}_K)$ 是这样的序对，因为 $A$ 是正规整环，所以
$A = K \cap B$。这些序对构成的集合满足 Zorn 引理的假设，
故存在极大序对 $(M, \sigma)$。若 $M \not = L$，则可以找到
$M \subset M' \subset L$，使得 $M'/M$ 非平凡且有限，并且
$M'/K$ 为 Galois 扩张（《域》引理
\ref{fields-lemma-normal-closure-inside-normal}）。选取
$\sigma' \in \text{Gal}(M'/K)$，使它在 $M$ 上的限制为
$\sigma$（《域》引理 \ref{fields-lemma-galois-infinite}）。
于是素理想 $\sigma'(\mathfrak q \cap M')$ 与
$\mathfrak q' \cap M'$ 在 $B \cap M$ 中的收缩相同。由于
$B \cap M = (B \cap M')^{\text{Gal}(M'/M)}$，可以应用引理
\ref{more-algebra-lemma-one-orbit}，找到 $\tau \in \text{Gal}(M'/M)$，使得
$\tau(\sigma'(\mathfrak q \cap M')) = \mathfrak q' \cap M'$。
于是 $(M', \tau \circ \sigma') > (M, \sigma)$，这与
$(M, \sigma)$ 的极大性矛盾。

\medskip\noindent
(2) 的证明方式与 (1) 完全相同。下面写出细节。选取
$\overline{\sigma} \in
\text{Aut}(\kappa(\mathfrak q)/\kappa(\mathfrak p))$。
考虑序对 $(M, \sigma)$，其中 $K \subset M \subset L$ 是子域，
$M/K$ 为 Galois 扩张，$\sigma \in \text{Gal}(M/K)$，
$\sigma(\mathfrak q \cap M) = \mathfrak q \cap M$，并且图
$$
\xymatrix{
\kappa(\mathfrak q \cap M) \ar[r] \ar[d]_\sigma &
\kappa(\mathfrak q) \ar[d]_{\overline{\sigma}} \\
\kappa(\mathfrak q \cap M) \ar[r] & \kappa(\mathfrak q)
}
$$
交换。定义 $(M', \sigma') \geq (M, \sigma)$，当且仅当
$M \subset M'$ 且 $\sigma'|_M = \sigma$。与上面一样，
$(K, \text{id}_K)$ 是这样的序对。这些序对构成的集合满足
Zorn 引理的假设，故存在极大序对 $(M, \sigma)$。若
$M \not = L$，则可以找到 $M \subset M' \subset L$，使得
$M'/M$ 有限，并且 $M'/K$ 为 Galois 扩张（《域》引理
\ref{fields-lemma-normal-closure-inside-normal}）。选取
$\sigma' \in \text{Gal}(M'/K)$，使它在 $M$ 上的限制为
$\sigma$（《域》引理 \ref{fields-lemma-galois-infinite}）。
于是素理想 $\sigma'(\mathfrak q \cap M')$ 与
$\mathfrak q \cap M'$ 在 $B \cap M$ 中的收缩相同。像第一段那样
调整 $\sigma'$ 的选择后，可以假设
$\sigma'(\mathfrak q \cap M') = \mathfrak q \cap M'$。
于是 $\sigma'$ 与 $\overline{\sigma}$ 定义映射
$\kappa(\mathfrak q \cap M') \to \kappa(\mathfrak q)$，
且它们在 $\kappa(\mathfrak q \cap M)$ 上相同。由于
$B \cap M = (B \cap M')^{\text{Gal}(M'/M)}$，可以应用引理
\ref{more-algebra-lemma-one-orbit-geometric}，找到
$\tau \in \text{Gal}(M'/M)$，使得
$\tau(\mathfrak q \cap M') = \mathfrak q \cap M'$，并使
% P02-E0587
$\tau \circ \sigma$ 与 $\overline{\sigma}$ 在
$\kappa(\mathfrak q \cap M')$ 上诱导同一个映射。这里有一个小细节：
该引理首先保证
$\kappa(\mathfrak q \cap M')/\kappa(\mathfrak q \cap M)$ 正规，
由此可知两个映射之差是该扩张的自同构（《域》引理
\ref{fields-lemma-normal-embeddings-differ-by-aut}），再对它应用该引理
即可得到 $\tau$。因此
$(M', \tau \circ \sigma') > (M, \sigma)$，这与
$(M, \sigma)$ 的极大性矛盾。
\end{proof}

\begin{lemma}
\label{more-algebra-lemma-one-orbit-geometric-galois-compare}
设 $A$ 为分式域是 $K$ 的正规整环，
$M/L/K$ 为 $K$ 的（可能无限的）Galois 扩张塔。令
$H = \text{Gal}(M/K)$、$G = \text{Gal}(L/K)$，并分别令
$C$ 与 $B$ 为 $A$ 在 $M$ 与 $L$ 中的整闭包。
% P02-E0588
设 $\mathfrak r \subset C$，并令
$\mathfrak q = B \cap \mathfrak r$。置
$D_\mathfrak r = \{\tau \in H \mid \tau(\mathfrak r) = \mathfrak r\}$
以及
$I_\mathfrak r = \{\tau \in D_\mathfrak r \mid
\tau \bmod \mathfrak r = \text{id}_{\kappa(\mathfrak r)}\}$，
并类似地定义 $D_\mathfrak q$ 与 $I_\mathfrak q$。
在映射 $H \to G$ 下，所诱导的映射
$D_\mathfrak r \to D_\mathfrak q$ 与
$I_\mathfrak r \to I_\mathfrak q$ 都满射。
\end{lemma}

\begin{proof}
令 $\sigma \in D_\mathfrak q$。选取映到 $\sigma$ 的
$\tau \in H$；由《域》引理
\ref{fields-lemma-galois-infinite}，这是可能的。于是
$\tau(\mathfrak r)$ 与 $\mathfrak r$ 都位于 $\mathfrak q$ 之上。
故由引理 \ref{more-algebra-lemma-one-orbit-geometric-galois}，存在
$\sigma' \in \text{Gal}(M/L)$，使得
$\sigma'(\tau(\mathfrak r)) = \mathfrak r$。于是
$\sigma'\tau \in D_\mathfrak r$ 映到 $\sigma$。
惰性群的情形完全相同，只须利用到自同构群的满射性。
\end{proof}
\section{离散赋值环的扩张}
\label{more-algebra-section-discrete-valuation-rings}

\noindent
在本节及随后几节中，我们使用以下定义。

\begin{definition}
\label{more-algebra-definition-extension-discrete-valuation-rings}
若 $A$ 与 $B$ 都是离散赋值环，并且 $A \to B$ 单射且为局部映射，
则称 $A \to B$ 或 $A \subset B$ 为{\it 离散赋值环的扩张}。
特别地，若 $\pi_A$ 与 $\pi_B$ 分别为 $A$ 与 $B$ 的一致化元，
则对某个 $e \geq 1$ 以及 $B$ 的某个单位 $u$，有
$\pi_A = u \pi_B^e$。整数 $e$ 与一致化元的选择无关，因为它也是
使下式成立的唯一整数 $\geq 1$：
% P02-E0589
$$
\mathfrak m_A B = \mathfrak m_B^e
$$
。整数 $e$ 称为 $B$ 在 $A$ 上的{\it 分歧指数}。若 $e = 1$，
则称 $B$ 在 $A$ 上{\it 弱非分歧}。若剩余域扩张
$\kappa_A = A/\mathfrak m_A \subset \kappa_B = B/\mathfrak m_B$
有限，则置 $f = [\kappa_B : \kappa_A]$，并称它为扩张
$A \subset B$ 的{\it 剩余次数}。
\end{definition}

\noindent
注意，我们不要求分式域扩张有限。

\begin{lemma}
\label{more-algebra-lemma-inequality}
设 $A \subset B$ 为离散赋值环的扩张，其分式域为
$K \subset L$。若扩张 $L/K$ 有限，则剩余域扩张有限，并且
$ef \leq [L : K]$。
\end{lemma}

\begin{proof}
% P02-E0595
剩余域扩张的有限性见《代数》引理
\ref{algebra-lemma-finite-extension-residue-fields-dimension-1}。
不等式由《代数》引理 \ref{algebra-lemma-finite-length} 与
\ref{algebra-lemma-pushdown-module} 得出。
\end{proof}

\begin{lemma}
\label{more-algebra-lemma-multiplicative-e-f}
设 $A \subset B \subset C$ 为离散赋值环的扩张。则 $B/A$
与 $C/B$ 的分歧指数之积等于 $C/A$ 的分歧指数，即
% P02-E0596
$e_{C/A} = e_{B/A} e_{C/B}$。若各剩余次数有限，它们也满足同样的
乘法性。
\end{lemma}

\begin{proof}
这由定义及《域》引理
\ref{fields-lemma-multiplicativity-degrees} 立即得出。
\end{proof}

\begin{lemma}
\label{more-algebra-lemma-ramification-index-a-power-of-p}
设 $A \subset B$ 为离散赋值环的扩张，所诱导的域扩张为
$K \subset L$。若 $K$ 的特征为 $p > 0$，且 $L$ 在 $K$
上纯不可分，则分歧指数 $e$ 是 $p$ 的幂。
\end{lemma}

\begin{proof}
对某个 $u \in B^*$，写 $\pi_A = u \pi_B^e$。另一方面，
对某个 $p$ 的幂 $q$，有 $\pi_B^q \in K$。对某个
$v \in A^*$ 与 $k \in \mathbf{Z}$，写
$\pi_B^q = v \pi_A^k$。于是
$\pi_A^q = u^q \pi_B^{qe} = u^q v^e \pi_A^{ke}$。
在 $B$ 中取赋值，得到 $ke = q$。
\end{proof}

\noindent
在下一个引理中，我们讨论离散赋值环的扩张 $A \subset B$
为“非分歧”的含义，即分歧指数为 $1$，且剩余域扩张可分
（可能非代数）。不过，我们不能直接使用“非分歧”一词，因为已经有
非分歧环映射这一概念；见《代数》第
\ref{algebra-section-unramified} 节。

\begin{lemma}
\label{more-algebra-lemma-extension-dvrs-formally-smooth}
设 $A \subset B$ 为离散赋值环的扩张。以下条件等价：
\begin{enumerate}
\item $A \to B$ 关于 $\mathfrak m_B$-进拓扑形式光滑；
\item $A \to B$ 弱非分歧，且 $\kappa_B/\kappa_A$
是可分域扩张。
\end{enumerate}
\end{lemma}

\begin{proof}
这由命题 \ref{more-algebra-proposition-fs-flat-fibre-fs} 与《代数》命题
\ref{algebra-proposition-characterize-separable-field-extensions}
得出。
\end{proof}

\begin{remark}
\label{more-algebra-remark-finite-separable-extension}
设 $A$ 为分式域是 $K$ 的离散赋值环，$L/K$ 为有限可分域扩张，
$B \subset L$ 为 $A$ 在 $L$ 中的整闭包。图为：
% P02-E0590
$$
\xymatrix{
B \ar[r] & L \\
A \ar[u] \ar[r] & K \ar[u]
}
$$
。由《代数》引理
\ref{algebra-lemma-Noetherian-normal-domain-finite-separable-extension}，
环扩张 $A \subset B$ 有限，故 $B$ 是 Noether 环。由《代数》引理
\ref{algebra-lemma-integral-sub-dim-equal}，$B$ 的维数为 $1$，
故 $B$ 是 Dedekind 整环；见《代数》引理
\ref{algebra-lemma-characterize-Dedekind}。令
$\mathfrak m_1, \ldots, \mathfrak m_n$ 为 $B$ 的极大理想
（即位于 $\mathfrak m_A$ 之上的素理想）。我们得到离散赋值环的扩张
% P02-E0590
$$
A \subset B_{\mathfrak m_i}
$$
，从而得到分歧指数 $e_i$ 与剩余次数 $f_i$。由《代数》引理
\ref{algebra-lemma-finite-extension-dim-1} 对 $A$ 中的一致化元的应用，
有
% P02-E0590
$$
[L : K] = \sum\nolimits_{i = 1, \ldots, n} e_i f_i
$$
。注意，若 $A$ 是 Hensel 环，则 $n = 1$（由《代数》引理
\ref{algebra-lemma-finite-over-henselian}）；例如 $A$ 完备时即如此。
\end{remark}

\begin{definition}
\label{more-algebra-definition-types-of-extensions}
设 $A$ 为分式域是 $K$ 的离散赋值环，$L/K$ 为有限可分扩张。
沿用注记 \ref{more-algebra-remark-finite-separable-extension} 中的 $B$ 与
$\mathfrak m_i$（$i = 1, \ldots, n$），称扩张 $L/K$：
\begin{enumerate}
\item {\it 关于 $A$ 非分歧}，若对所有 $i$ 都有 $e_i = 1$，
且扩张 $\kappa(\mathfrak m_i)/\kappa_A$ 可分；
\item {\it 关于 $A$ 驯分歧}，若 $\kappa_A$ 的特征为 $0$，或者
$\kappa_A$ 的特征为 $p > 0$、各域扩张
$\kappa(\mathfrak m_i)/\kappa_A$ 可分且各分歧指数 $e_i$
与 $p$ 互素；
\item {\it 关于 $A$ 完全分歧}，若 $n = 1$，且剩余域扩张
$\kappa(\mathfrak m_1)/\kappa_A$ 平凡。
\end{enumerate}
若从上下文可明确离散赋值环 $A$，则有时简称 $L/K$ 非分歧、
完全分歧或驯分歧。
\end{definition}

\noindent
对非分歧扩张，有以下基本引理。

\begin{lemma}
\label{more-algebra-lemma-permanence-unramified}
设 $A$ 为分式域是 $K$ 的离散赋值环。
\begin{enumerate}
\item 若 $M/L/K$ 是有限可分扩张，且 $M$ 关于 $A$ 非分歧，
则 $L$ 关于 $A$ 非分歧。
\item 若 $L/K$ 是关于 $A$ 非分歧的有限可分扩张，则存在包含
$L$ 且关于 $A$ 非分歧的 Galois 扩张 $M/K$。
\item 若 $L_1/K$、$L_2/K$ 是关于 $A$ 非分歧的有限可分扩张，
则存在包含 $L_1$ 与 $L_2$ 且关于 $A$ 非分歧的有限可分扩张
$L/K$。
% P02-E0591
\end{enumerate}
\end{lemma}

\begin{proof}
以下不再特别说明地使用注记
\ref{more-algebra-remark-finite-separable-extension} 中讨论的结果。

\medskip\noindent
证明 (1)。令 $C/B/A$ 分别为 $A$ 在 $M/L/K$ 中的整闭包。
% P02-E0592
由于 $C$ 是 $B$ 的有限环扩张，
$\Spec(C) \to \Spec(B)$ 满射。因此，对每个极大理想
$\mathfrak m \subset B$，都有一个位于 $\mathfrak m$ 之上的极大理想
$\mathfrak m' \subset C$。由分歧指数的乘法性（引理
\ref{more-algebra-lemma-multiplicative-e-f}）与假设，$B_\mathfrak m$
在 $A$ 上的分歧指数为 $1$。由于
$\kappa(\mathfrak m')/\kappa_A$ 有限可分，
$\kappa(\mathfrak m)/\kappa_A$ 也如此。

\medskip\noindent
证明 (2)。令 $M$ 为 $L$ 在 $K$ 上的正规闭包；见《域》定义
\ref{fields-definition-normal-closure}。由《域》引理
\ref{fields-lemma-normal-closure-galois}，$M/K$ 是 Galois 扩张。
另一方面，存在 $K$-代数的满射
% P02-E0593
$$
L \otimes_K \ldots \otimes_K L \longrightarrow M
$$
；见《域》引理
\ref{fields-lemma-normal-closure-tensor-product}。令 $B$ 为
$A$ 在 $L$ 中的整闭包，如注记
\ref{more-algebra-remark-finite-separable-extension} 所述。$L$ 关于 $A$
非分歧这一条件恰好意味着 $A \to B$ 是 \'etale 环映射；
见《代数》引理 \ref{algebra-lemma-characterize-etale}。
由 \'etale 环映射的恒久性质，
$$
B \otimes_A \ldots \otimes_A B
$$
在 $A$ 上 \'etale；见《代数》引理
\ref{algebra-lemma-etale}。因此，所显示的环是 Dedekind 整环的
乘积；见引理 \ref{more-algebra-lemma-Dedekind-etale-extension}。由此，
$M$ 是某个在 $A$ 上有限 \'etale 的 Dedekind 整环的分式域。
这意味着 $M$ 如所求地关于 $A$ 非分歧。

\medskip\noindent
证明 (3)。令 $B_i \subset L_i$ 分别为 $A$ 在相应域中的整闭包。
% P02-E0594
采用与上面相同的论证，可知
$B_1 \otimes_A B_2$ 在 $A$ 上有限 \'etale。略去细节。
\end{proof}

\begin{lemma}
\label{more-algebra-lemma-composition-unramified}
设 $A$ 为分式域是 $K$ 的离散赋值环，$M/L/K$ 为有限可分扩张。
令 $B$ 为 $A$ 在 $L$ 中的整闭包。若 $L/K$ 关于 $A$ 非分歧，
并且对 $B$ 的每个极大理想 $\mathfrak m$，$M/L$
都关于 $B_\mathfrak m$ 非分歧，则 $M/K$ 关于 $A$ 非分歧。
\end{lemma}

\begin{proof}
令 $C$ 为 $A$ 在 $M$ 中的整闭包。$C$ 的每个极大理想
$\mathfrak m'$ 都位于 $B$ 的某个极大理想 $\mathfrak m$ 之上。
于是，本引理由分歧指数的乘法性（引理
\ref{more-algebra-lemma-multiplicative-e-f}）以及有限域扩张塔
$\kappa(\mathfrak m')/\kappa(\mathfrak m)/\kappa_A$ 得出。
\end{proof}
\section{Galois 扩张与分歧}
\label{more-algebra-section-ramification}

\noindent
在 Galois 扩张的情形，可以进一步展开第
\ref{more-algebra-section-discrete-valuation-rings} 节的讨论。

\begin{lemma}
\label{more-algebra-lemma-galois}
设 $A$ 为分式域是 $K$ 的离散赋值环，$L/K$ 为 Galois 群是
$G$ 的有限 Galois 扩张。则 $G$ 作用在注记
\ref{more-algebra-remark-finite-separable-extension} 的环 $B$ 上，并传递地作用在
$B$ 的极大理想集合上。
\end{lemma}

\begin{proof}
注意 $A = B^G$，因为 $A$ 在 $K$ 中整闭且 $K = L^G$。
故本引理是引理 \ref{more-algebra-lemma-one-orbit} 的特殊情形。
\end{proof}

\begin{lemma}
\label{more-algebra-lemma-galois-conclusion}
设 $A$ 为分式域是 $K$ 的离散赋值环，$L/K$ 为有限 Galois 扩张。
则存在 $e \geq 1$ 与 $f \geq 1$，使得对所有 $i$ 都有
$e_i = e$ 与 $f_i = f$（记号同注记
\ref{more-algebra-remark-finite-separable-extension}）。特别地，
$[L : K] = n e f$。
\end{lemma}

\begin{proof}
这是引理 \ref{more-algebra-lemma-galois} 与各定义的直接推论。
\end{proof}

\begin{definition}
\label{more-algebra-definition-decomposition-inertia}
设 $A$ 为分式域是 $K$ 的离散赋值环，$L/K$ 为 Galois 群是
$G$ 的有限 Galois 扩张。令 $B$ 为 $A$ 在 $L$ 中的整闭包，
$\mathfrak m \subset B$ 为极大理想。
\begin{enumerate}
\item $\mathfrak m$ 的{\it 分解群}是子群
$D = \{\sigma \in G \mid \sigma(\mathfrak m) = \mathfrak m\}$。
\item $\mathfrak m$ 的{\it 惰性群}是映射
$D \to \text{Aut}(\kappa(\mathfrak m)/\kappa_A)$ 的核 $I$。
\end{enumerate}
\end{definition}

\noindent
注意，域 $\kappa(\mathfrak m)$ 在 $\kappa_A$ 上可能不可分。
特别地，域扩张 $\kappa(\mathfrak m)/\kappa_A$ 未必是 Galois 扩张。
若 $\kappa_A$ 完美，则它是 Galois 扩张。

\begin{lemma}
\label{more-algebra-lemma-galois-galois}
设 $A$ 为分式域是 $K$、剩余域是 $\kappa$ 的离散赋值环，
$L/K$ 为 Galois 群是 $G$ 的有限 Galois 扩张。令 $B$ 为
$A$ 在 $L$ 中的整闭包，$\mathfrak m$ 为 $B$ 的极大理想。则：
\begin{enumerate}
\item 域扩张 $\kappa(\mathfrak m)/\kappa$ 正规；
\item $D \to \text{Aut}(\kappa(\mathfrak m)/\kappa)$ 满射。
\end{enumerate}
若对某个（等价地，每个）极大理想 $\mathfrak m \subset B$，
域扩张 $\kappa(\mathfrak m)/\kappa$ 可分，则：
\begin{enumerate}
\item[(3)] $\kappa(\mathfrak m)/\kappa$ 是 Galois 扩张；
\item[(4)] $D \to \text{Gal}(\kappa(\mathfrak m)/\kappa)$ 满射。
\end{enumerate}
这里 $D \subset G$ 是 $\mathfrak m$ 的分解群。
\end{lemma}

\begin{proof}
注意 $A = B^G$，因为 $A$ 在 $K$ 中整闭且 $K = L^G$。故
(1) 与 (2) 由引理 \ref{more-algebra-lemma-one-orbit-geometric} 得出。
引理中的“等价地，每个”部分由引理 \ref{more-algebra-lemma-galois} 得出。
假设 $\kappa(\mathfrak m)/\kappa$ 可分，则 (3) 与 (4)
立即由 (1) 与 (2) 得出。
\end{proof}

\begin{lemma}
\label{more-algebra-lemma-galois-inertia}
设 $A$ 为分式域是 $K$ 的离散赋值环，$L/K$ 为 Galois 群是
$G$ 的有限 Galois 扩张。令 $B$ 为 $A$ 在 $L$ 中的整闭包，
$\mathfrak m \subset B$ 为极大理想。$\mathfrak m$ 的惰性群
$I$ 位于典范正合列
$$
1 \to P \to I \to I_t \to 1
$$
中，并满足：
\begin{enumerate}
\item 若 $D$ 是分解群，则
$$
P = \{\sigma \in I \mid \sigma|_{\mathfrak m/\mathfrak m^2} =
\text{id}_{\mathfrak m/\mathfrak m^2}\}
=
\{\sigma \in D \mid \sigma \text{ acts trivially on }\text{Gr}_\mathfrak m(B)\}
$$
；
\item $P$ 是 $D$ 的正规子群；
\item 若 $\kappa_A$ 的特征为 $p > 0$，则 $P$ 是 $p$-群；
若 $\kappa_A$ 的特征为零，则 $P = \{1\}$；
\item $I_t$ 是循环群，其阶为整数 $e$ 的与 $p$ 互素部分；
% P02-E0597
\item 存在典范同构
$\theta : I_t \to \mu_e(\kappa(\mathfrak m))$；
\item 若 $\kappa(m)$ 在 $A$ 的剩余域上可分，则
% P02-E0598
$P =
\{\sigma \in D \mid \sigma|_{B/\mathfrak m^2} = \text{id}_{B/\mathfrak m^2}\}$。
\end{enumerate}
这里 $e$ 是引理 \ref{more-algebra-lemma-galois-conclusion} 中的整数。
\end{lemma}

\begin{proof}
回忆 $|G| = [L : K] = nef$；见引理
\ref{more-algebra-lemma-galois-conclusion}。由于 $G$ 传递地作用在 $B$
的极大理想集合
$\{\mathfrak m_1, \ldots, \mathfrak m_n\}$ 上
（引理 \ref{more-algebra-lemma-galois}），且 $D$ 是一个元素的稳定子，
可见 $|D| = ef$。由引理 \ref{more-algebra-lemma-galois-galois}，有
$$
ef = |D| = |I| \cdot |\text{Aut}(\kappa(\mathfrak m)/\kappa)|
$$
，其中 $\kappa$ 是 $A$ 的剩余域。由于
$\kappa(\mathfrak m)$ 在 $\kappa$ 上正规，
$\text{Aut}(\kappa(\mathfrak m)/\kappa)$ 的阶与 $f$
相差 $p$ 的一个幂（见《域》引理
\ref{fields-lemma-normal-and-automorphisms} 以及《域》定义
\ref{fields-definition-insep-degree} 后的讨论）。因此，$|I|$
的与 $p$ 互素部分等于 $e$ 的与 $p$ 互素部分。

\medskip\noindent
置 $C = B_\mathfrak m$。则 $I$ 在 $A$ 上作用于 $C$，并且
平凡作用于 $C$ 的剩余域。令 $\pi_A \in A$ 与 $\pi_C \in C$
为一致化元。对 $C$ 中某个单位 $u$，写
$\pi_A = u \pi_C^e$。对 $\sigma \in I$，对 $C$ 中某个单位
$\theta_\sigma$，写
$\sigma(\pi_C) = \theta_\sigma \pi_C$。于是
% P02-E0599
$$
\pi_A = \sigma(\pi_A) = \sigma(u) (\theta_\sigma \pi_C)^e
= \sigma(u) \theta_\sigma^e \pi_C^e = \frac{\sigma(u)}{u} \theta_\sigma^e \pi_A
$$
。由于 $\sigma \in I$，有
$\sigma(u) \equiv u \bmod \mathfrak m_C$，故
$\theta_\sigma$ 在 $\kappa_C = \kappa(\mathfrak m)$ 中的像
$\overline{\theta}_\sigma$ 是 $e$ 次单位根。我们得到映射
\begin{equation}
\label{more-algebra-equation-inertia-character}
\theta : I \longrightarrow \mu_e(\kappa(\mathfrak m)),\quad
\sigma \mapsto \overline{\theta}_\sigma
\end{equation}
。我们断言 $\theta$ 是群同态，且与一致化元 $\pi_C$ 的选择无关。
事实上，若 $\tau$ 是 $I$ 的另一个元素，则
$\tau(\sigma(\pi_C)) = \tau(\theta_\sigma \pi_C) =
\tau(\theta_\sigma) \theta_\tau \pi_C$，故
$\theta_{\tau \sigma} = \tau(\theta_\sigma) \theta_\tau$。由于
$\tau \in I$，可知
$\overline{\theta}_{\tau \sigma} =
\overline{\theta}_\sigma \overline{\theta}_\tau$。若 $\pi'_C$
是另一个一致化元，则对 $C$ 中某个单位 $w$，有
$\pi'_C = w \pi_C$，并且
$\sigma(\pi'_C) = w^{-1}\sigma(w)\theta_\sigma \pi'_C$。
因此 $\theta'_\sigma = w^{-1}\sigma(w)\theta_\sigma$，
从而 $\theta'_\sigma$ 与 $\theta_\sigma$ 如前一样映到剩余域中的
同一个元素。

\medskip\noindent
由于 $\kappa(\mathfrak m)$ 的特征为 $p$，群
$\mu_e(\kappa(\mathfrak m))$ 是循环群，其阶至多为 $e$
的与 $p$ 互素部分（见《域》第
\ref{fields-section-roots-of-1} 节）。
% P02-E0597

\medskip\noindent
令 $P = \Ker(\theta)$。由构造，$P$ 的元素恰为 $I$
中在 $\mathfrak m/\mathfrak m^2 =
\text{Gr}_\mathfrak m^1(B)$ 上平凡作用的元素，即
$P = \{\sigma \in I \mid \sigma|_{\mathfrak m/\mathfrak m^2} =
\text{id}_{\mathfrak m/\mathfrak m^2}\}$。此外，$I$ 由 $D$
中在 $\kappa(\mathfrak m) = B/\mathfrak m$ 上平凡作用的元素构成。
由于分次环 $\text{Gr}_\mathfrak m(B)$ 由
$\text{Gr}^0_\mathfrak m(B) = B/\mathfrak m$ 上的
$\text{Gr}^1_\mathfrak m(B) = \mathfrak m/\mathfrak m^2$ 生成，
可知
$P = \{\sigma \in D \mid \sigma \text{ acts trivially on }
\text{Gr}_\mathfrak m(B)\}$。所以 (1) 成立。由此也得到 (2)，
因为 $P$ 是同态
$D \to \text{Aut}(\text{Gr}_\mathfrak m(B))$ 的核。
若能证明 (3)，则 (4) 与 (5) 也随之成立：$I_t$
将同构于 $\mu_e(\kappa(\mathfrak m))$，因为上述论证表明
$|I_t| \geq |\mu_e(\kappa(\mathfrak m))|$。

\medskip\noindent
因此，只须证明 $\theta$ 的核 $P$ 是 $p$-群。令 $\sigma$
为该核的非平凡元素。则对所有 $i$，$\sigma - \text{id}$
把 $\mathfrak m_C^i$ 映入 $\mathfrak m_C^{i + 1}$。
令 $m$ 为 $\sigma$ 的阶。选取 $c \in C$，使得
$\sigma(c) \not = c$。则对某个 $i$，有
$\sigma(c) - c \in \mathfrak m_C^i$ 但
$\sigma(c) - c \not \in \mathfrak m_C^{i + 1}$，且
\begin{align*}
0
& =
\sigma^m(c) - c \\
& =
\sigma^m(c) - \sigma^{m - 1}(c) + \ldots + \sigma(c) - c \\
& =
\sum\nolimits_{j = 0, \ldots, m - 1} \sigma^j(\sigma(c) - c) \\
& \equiv
m(\sigma(c) - c) \bmod \mathfrak m_C^{i + 1}
\end{align*}
。由此 $p | m$（若 $p = 1$，则 $m = 0$）。因此，
$\theta$ 的核中每个元素的阶都被 $p$ 整除，即
$\Ker(\theta)$ 是 $p$-群。
% P02-E0600

\medskip\noindent
证明 (6)。假设 $\kappa(\mathfrak m)/\kappa$ 可分。此时，由引理
\ref{more-algebra-lemma-galois-galois}，$f = |D/I|$ 且 $I$ 的阶为 $e$。
若 $e = 1$，则 $P = I = \{\text{id}\}$，结论成立。若
$e > 1$，则 $B/\mathfrak m^2 = C/\pi_C^2C$ 是
$\kappa$-代数，也是 $\kappa(\mathfrak m)$ 的一个小扩张。
由于 $\kappa(\mathfrak m)$ 在 $\kappa$ 上形式
\'etale（《代数》引理
\ref{algebra-lemma-characterize-separable-algebraic-field-extensions}），
存在唯一的 $\kappa$-代数映射
$\kappa(\mathfrak m) \to B/\mathfrak m^2$，它是
$B/\mathfrak m^2 \to \kappa(\mathfrak m)$ 的右逆。换言之，
存在 $\kappa$-代数的 $D$-等变同构
$B/\mathfrak m^2 = \mathfrak m/\mathfrak m^2 \oplus \kappa(\mathfrak m)$。
这立即表明，在此情形下
$P =
\{\sigma \in D \mid \sigma|_{B/\mathfrak m^2} = \text{id}_{B/\mathfrak m^2}\}$。
\end{proof}

\begin{example}
\label{more-algebra-example-counter-description-P}
若没有关于 $\kappa(\mathfrak m)$ 的假设，引理
\ref{more-algebra-lemma-galois-inertia} 的 (6) 中的等式不成立。一个反例是由
$B = A[x]/(x^2 - t)$ 给出的 $A = Z_2[t]_{(2)}$ 的扩张。
事实上，此时 $\sigma(x) = -x = x - 2x$，而
$2x$ 不属于 $\mathfrak m^2$。
\end{example}

\begin{definition}
\label{more-algebra-definition-wild-inertia}
沿用引理 \ref{more-algebra-lemma-galois-inertia} 的假设与记号。
% P02-E0601
\begin{enumerate}
\item $\mathfrak m$ 的{\it 狂惰性群}是子群 $P$。
\item $\mathfrak m$ 的{\it 驯惰性群}是商 $I \to I_t$。
\end{enumerate}
以 $\theta : I \to \mu_e(\kappa(\mathfrak m))$ 表示满射
(\ref{more-algebra-equation-inertia-character})；其核为 $P$，并诱导同构
$I_t \to \mu_e(\kappa(\mathfrak m))$。
% P02-E0602
\end{definition}

\begin{lemma}
\label{more-algebra-lemma-inertia-character}
沿用引理 \ref{more-algebra-lemma-galois-inertia} 的假设与记号。惰性特征标
$\theta : I \to \mu_e(\kappa(\mathfrak m))$ 满足：
% P02-E0603
$$
\theta(\tau \sigma \tau^{-1}) = \tau(\theta(\sigma))
$$
，其中 $\tau \in D$ 且 $\sigma \in I$。
\end{lemma}

\begin{proof}
该公式有意义，因为 $I$ 是 $D$ 的正规子群，并且 $\tau$
通过例如引理 \ref{more-algebra-lemma-galois-galois} 中讨论的映射
$D \to \text{Aut}(\kappa(\mathfrak m))$ 作用在
$\kappa(\mathfrak m)$ 上。回忆 $\theta$ 的构造。选取
$B_\mathfrak m$ 的一致化元 $\pi$，并对 $\sigma \in I$ 写
$\sigma(\pi) = \theta_\sigma \pi$。则 $\theta(\sigma)$ 是
$\theta_\sigma$ 在剩余域中的像
$\overline{\theta}_\sigma$。对任意 $\tau \in D$，可以对某个单位
$\theta_\tau$ 写成 $\tau(\pi) = \theta_\tau \pi$。于是
$\theta_{\tau^{-1}} = \tau^{-1}(\theta_\tau^{-1})$。计算得到
\begin{align*}
\theta_{\tau \sigma \tau^{-1}}
& =
\tau(\sigma(\tau^{-1}(\pi)))/\pi \\
& =
\tau(\sigma(\tau^{-1}(\theta_\tau^{-1}) \pi))/\pi \\
& =
\tau(\sigma(\tau^{-1}(\theta_\tau^{-1})) \theta_\sigma \pi)/\pi \\
& =
\tau(\sigma(\tau^{-1}(\theta_\tau^{-1}))) \tau(\theta_\sigma) \theta_\tau
\end{align*}
。不过，由于 $\sigma$ 模 $\pi$ 平凡作用，乘积
$\tau(\sigma(\tau^{-1}(\theta_\tau^{-1}))) \theta_\tau$
在剩余域中的像为 $1$。这就证明了本引理。
% P02-E0604
\end{proof}

\noindent
我们将在《基本群》引理
\ref{pione-lemma-inertial-invariants-unramified} 中推广下一个引理。

\begin{lemma}
\label{more-algebra-lemma-inertial-invariants-unramified}
设 $A$ 为分式域是 $K$ 的离散赋值环，$L/K$ 为有限 Galois 扩张。
设 $\mathfrak m \subset B$ 为 $A$ 在 $L$ 中的整闭包的一个极大理想，
$I \subset G$ 为 $\mathfrak m$ 的惰性群。则 $B^I$ 是 $A$
在 $L^I$ 中的整闭包，且
$A \to (B^I)_{B^I \cap \mathfrak m}$ 是 \'etale 映射。
\end{lemma}

\begin{proof}
写 $B' = B^I$。由定义可知，$B' = B^I$ 是 $A$ 在 $L^I$
中的整闭包。写
$\mathfrak m' = B^I \cap \mathfrak m = B' \cap \mathfrak m \subset B'$。
由引理 \ref{more-algebra-lemma-one-orbit}，$\mathfrak m$ 是 $B$
中位于 $\mathfrak m'$ 之上的唯一素理想。由于 $I$
平凡作用在 $\kappa(\mathfrak m)$ 上，由引理
\ref{more-algebra-lemma-invariants-modulo} 可见扩张
$\kappa(\mathfrak m)/\kappa(\mathfrak m')$ 纯不可分
（一个也许更简单的替代办法是应用引理
\ref{more-algebra-lemma-one-orbit-geometric} 的结果）。由于 $D/I$
忠实作用在 $\kappa(\mathfrak m')$ 上，可知 $D/I$
忠实作用在 $\kappa(\mathfrak m)$ 上。当然，$A$ 的剩余域
$\kappa$ 的元素在此作用下不变。由 Galois 理论，
$[\kappa(\mathfrak m') : \kappa] \geq |D/I|$；
见《域》引理 \ref{fields-lemma-galois-over-fixed-field}。

\medskip\noindent
令 $\pi$ 为 $A$ 的一致化元。由于
$\text{Norm}_{L/K}(\pi) = \pi^{[L : K]}$，由《代数》引理
\ref{algebra-lemma-finite-extension-dim-1} 可知
% P02-E0605
$$
|G| = [L : K] = [L : K]\ \text{ord}_A(\pi) =
|G/D|\ [\kappa(\mathfrak m) : \kappa]\ \text{ord}_{B_\mathfrak m}(\pi)
$$
；这是因为 $B$ 有 $n = |G/D|$ 个极大理想，且它们在 $G$
下全都共轭；见注记 \ref{more-algebra-remark-finite-separable-extension}
与引理 \ref{more-algebra-lemma-galois}。对次数为 $|I|$ 的有限扩张
$L/L^I$ 应用同样的论证，得到
% P02-E0606
$$
|I|\ \text{ord}_{B'_{\mathfrak m'}}(\pi) =
[\kappa(\mathfrak m) : \kappa(\mathfrak m')]\ \text{ord}_{B_\mathfrak m}(\pi)
$$
。由此
% P02-E0605
$$
\text{ord}_{B'_{\mathfrak m'}}(\pi) =
\frac{|D/I|}{[\kappa(\mathfrak m') : \kappa]}
$$
。由于左侧是正整数，而由上述结论右侧 $\leq 1$，可知必有等号，
即 $\text{ord}_{B'_{\mathfrak m'}}(\pi) = 1$，且
$\kappa(\mathfrak m')/\kappa$ 的次数为 $|D/I|$。因此
% P02-E0607
$\pi B'_{\mathfrak m'} = \mathfrak m' B_\mathfrak m'$，且
$\kappa(\mathfrak m')$ 在 $\kappa$ 上是 Galois 扩张，其
Galois 群为 $D/I$；特别地，它可分，见《域》引理
\ref{fields-lemma-finite-Galois}。由《代数》引理
\ref{algebra-lemma-characterize-etale}，
$A \to B'_{\mathfrak m'}$ 是 \'etale 映射，如所求。
\end{proof}

\begin{remark}
\label{more-algebra-remark-tower-of-rings}
设 $A$ 为分式域是 $K$ 的离散赋值环，$L/K$ 为有限 Galois 扩张。
设 $\mathfrak m \subset B$ 为 $A$ 在 $L$ 中的整闭包的一个极大理想。
令
$$
P \subset I \subset D \subset G
$$
依次为 $\mathfrak m$ 的狂惰性群、惰性群与分解群。
% P02-E0608
考虑图
$$
\xymatrix{
\mathfrak m \ar@{-}[d] \ar@{-}[r] &
\mathfrak m^P \ar@{-}[d] \ar@{-}[r] &
\mathfrak m^I \ar@{-}[d] \ar@{-}[r] &
\mathfrak m^D \ar@{-}[d] \ar@{-}[r] &
A \cap \mathfrak m \ar@{-}[d] \\
B & B^P \ar[l] & B^I \ar[l] & B^D \ar[l] & A \ar[l]
}
$$
。注意，$B^P, B^I, B^D$ 分别是 $A$ 在域
$L^P$、$L^I$、$L^D$ 中的整闭包。因此还可看出，$B^P$
是 $B^I$ 在 $L^P$ 中的整闭包，依此类推。注意
$\mathfrak m^P = \mathfrak m \cap B^P$、
$\mathfrak m^I = \mathfrak m \cap B^I$ 且
$\mathfrak m^D = \mathfrak m \cap B^D$。所以，图的顶行对应于
$\mathfrak m \in \Spec(B)$ 在所诱导的谱映射下的各个像。
综上，有：
\begin{enumerate}
\item 扩张 $L^I/L^D$ 是 Galois 扩张，其群为 $D/I$；
\item 扩张 $L^P/L^I$ 是 Galois 扩张，其群为 $I_t = I/P$；
\item 扩张 $L^P/L^D$ 是 Galois 扩张，其群为 $D/P$；
\item $\mathfrak m^I$ 是 $B^I$ 中位于 $\mathfrak m^D$
之上的唯一素理想；
\item $\mathfrak m^P$ 是 $B^P$ 中位于 $\mathfrak m^I$
之上的唯一素理想；
\item $\mathfrak m$ 是 $B$ 中位于 $\mathfrak m^P$ 之上的唯一素理想；
\item $\mathfrak m^P$ 是 $B^P$ 中位于 $\mathfrak m^D$
之上的唯一素理想；
\item $\mathfrak m$ 是 $B$ 中位于 $\mathfrak m^I$ 之上的唯一素理想；
\item $\mathfrak m$ 是 $B$ 中位于 $\mathfrak m^D$ 之上的唯一素理想；
\item $A \to B^D_{\mathfrak m^D}$ 是 \'etale 映射，并诱导平凡的
剩余域扩张；
\item $B^D_{\mathfrak m^D} \to B^I_{\mathfrak m^I}$ 是
\'etale 映射，并诱导 Galois 群为 $D/I$ 的 Galois 剩余域扩张；
\item $A \to B^I_{\mathfrak m^I}$ 是 \'etale 映射；
\item $B^I_{\mathfrak m^I} \to B^P_{\mathfrak m^P}$
的分歧指数为与 $p$ 互素的 $|I/P|$，并诱导平凡的剩余域扩张；
\item $B^D_{\mathfrak m^D} \to B^P_{\mathfrak m^P}$
的分歧指数为与 $p$ 互素的 $|I/P|$，并诱导可分剩余域扩张；
\item $A \to B^P_{\mathfrak m^P}$ 的分歧指数为与 $p$
互素的 $|I/P|$，并诱导可分剩余域扩张。
% P02-E0609
\end{enumerate}
(1)、(2) 与 (3) 由 Galois 理论（《域》第
\ref{fields-section-galois-theory} 节）及引理
\ref{more-algebra-lemma-galois-inertia} 立即得出。(4)--(9) 由引理
\ref{more-algebra-lemma-galois} 显然成立。(12) 是引理
\ref{more-algebra-lemma-inertial-invariants-unramified}。由于有分解
$A \to B^D_{\mathfrak m^D} \to B^I_{\mathfrak m^I}$，
可随之得到 (10) 与 (11) 中的 \'etale 性。(10) 中的剩余域扩张
必须平凡，因为它可分，并且如引理 \ref{more-algebra-lemma-galois-galois}
所示，$D/I$ 映满
$\text{Aut}(\kappa(\mathfrak m)/\kappa_A)$。同一论证给出 (11)
中关于剩余域的陈述。为证明 (13)、(14) 与 (15)，只须证明 (13)。
由上述结论，扩张 $L^P/L^I$ 是 Galois 扩张，其 Galois 群循环且
阶与 $p$ 互素；$\mathfrak m^P$ 是位于 $\mathfrak m^I$
之上的唯一素理想，且 $I/P$ 在剩余域上的作用平凡。因此，可以对
该扩张及离散赋值环 $B^I_{\mathfrak m^I}$ 应用引理
\ref{more-algebra-lemma-galois-inertia}，得到 (13)。
\end{remark}

\begin{lemma}
\label{more-algebra-lemma-compare-inertia}
设 $A$ 为分式域是 $K$ 的离散赋值环，$M/L/K$ 为扩张塔，其中
$M/K$ 与 $L/K$ 都是有限 Galois 扩张。令 $C$、$B$ 分别为
$A$ 在 $M$、$L$ 中的整闭包。
% P02-E0610
设 $\mathfrak m' \subset C$ 为极大理想，并置
$\mathfrak m = \mathfrak m' \cap B$。令
$$
P \subset I \subset D \subset \text{Gal}(L/K)
\quad\text{且}\quad
P' \subset I' \subset D' \subset \text{Gal}(M/K)
$$
依次为 $\mathfrak m$ 与 $\mathfrak m'$ 的狂惰性群、惰性群及
分解群。则典范满射
$\text{Gal}(M/K) \to \text{Gal}(L/K)$ 诱导满射
$P' \to P$、$I' \to I$ 与 $D' \to D$。此外，它们组成交换图
$$
\vcenter{
\xymatrix{
D' \ar[r] \ar[d] &
\text{Aut}(\kappa(\mathfrak m')/\kappa_A) \ar[d] \\
D \ar[r] &
\text{Aut}(\kappa(\mathfrak m)/\kappa_A)
}
}
\quad\text{且}\quad
\vcenter{
\xymatrix{
I' \ar[r]_-{\theta'} \ar[d] &
\mu_{e'}(\kappa(\mathfrak m')) \ar[d]^{(-)^{e'/e}} \\
I \ar[r]^-\theta &
\mu_e(\kappa(\mathfrak m))
}
}
$$
，其中 $e'$ 与 $e$ 分别是
$A \to C_{\mathfrak m'}$ 与 $A \to B_\mathfrak m$ 的分歧指数。
\end{lemma}

\begin{proof}
在映射 $\text{Gal}(M/K) \to \text{Gal}(L/K)$ 下，各群
$P', I', D'$ 分别映入 $P, I, D$，这由这些群的定义立即得出。
第一个图的交换性是清楚的（注意，由于
$\kappa(\mathfrak m)/\kappa_A$ 正规，
$\kappa(\mathfrak m')$ 在 $\kappa_A$ 上的每个自同构确实诱导
$\kappa(\mathfrak m)$ 在 $\kappa_A$ 上的自同构，故得到第一个图
中的右侧竖直箭头；见引理 \ref{more-algebra-lemma-galois-galois} 与《域》引理
\ref{fields-lemma-lift-maps}）。

\medskip\noindent
由引理 \ref{more-algebra-lemma-one-orbit-geometric-galois-compare}，映射
$I' \to I$ 与 $D' \to D$ 满射。$P' \to P$ 的满射性由
$P'$ 与 $P$ 分别是 $I'$ 与 $I$ 的 Sylow p-子群得出。

\medskip\noindent
为证明第二个图交换，选取 $C_{\mathfrak m'}$ 的一致化元
$\pi'$ 与 $B_\mathfrak m$ 的一致化元 $\pi$。则对
$C_{\mathfrak m'}$ 中某个单位 $c'$，有
$\pi = c' (\pi')^{e'/e}$。对 $\sigma' \in I'$，其像
$\sigma \in I$ 就是 $\sigma'$ 在 $L$ 上的限制。对某个单位
$c \in C_{\mathfrak m'}$ 写
$\sigma'(\pi') = c \pi'$，并对 $B_\mathfrak m$ 中某个单位 $b$
写 $\sigma(\pi) = b \pi$。于是
$\sigma'(\pi) = b \pi$，并得到
% P02-E0611
$$
b \pi = \sigma'(\pi) = \sigma'(c' (\pi')^{e'/e}) =
\sigma'(c') c^{e'/e} (\pi')^{e'/e} =
\frac{\sigma'(c')}{c'} c^{e'/e} \pi
$$
。由于 $\sigma' \in I'$，$b$ 与 $c^{e'/e}$ 在剩余域中的像相同，
这就证明了所需结论。
\end{proof}

\begin{remark}
\label{more-algebra-remark-canonical-inertia-character}
为把惰性特征标
$\theta : I \to \mu_e(\kappa(\mathfrak m))$ 用于无限 Galois 扩张，
对它作缩放是方便的。沿用引理 \ref{more-algebra-lemma-galois-inertia}
与定义 \ref{more-algebra-definition-wild-inertia} 中的
$A, K, L, B, \mathfrak m, G, P, I, D, e, \theta$。
则 $e = q |I_t|$，其中 $q$ 是 $\kappa(\mathfrak m)$ 的特征
$p$ 的一个幂（若该特征为正），或在特征为零时为 $1$。
% P02-E0612
注意，因 $\kappa(\mathfrak m)$ 的特征为 $p$，有
$\mu_e(\kappa(\mathfrak m)) = \mu_{|I_t|}(\kappa(\mathfrak m))$。
考虑映射
% P02-E0613
$$
\theta_{can} = q\theta : I \longrightarrow \mu_{|I_t|}(\kappa(\mathfrak m))
$$
。该映射诱导同构
$\theta_{can} : I_t \to \mu_{|I_t|}(\kappa(\mathfrak m))$。
由引理 \ref{more-algebra-lemma-inertia-character}，对 $\tau \in D$ 与
$\sigma \in I$，有
$\theta_{can}(\tau \sigma \tau^{-1}) = \tau(\theta_{can}(\sigma))$。
最后，若 $M/L$ 是使 $M/K$ 成为 Galois 扩张的扩张，且
$\mathfrak m'$ 是 $A$ 在 $M$ 中的整闭包中位于
$\mathfrak m$ 之上的素理想，则得到交换图
% P02-E0613
$$
\xymatrix{
I' \ar[r]_-{\theta'_{can}} \ar[d] &
\mu_{|I'_t|}(\kappa(\mathfrak m')) \ar[d]^{(-)^{|I'_t|/|I_t|}} \\
I \ar[r]^-{\theta_{can}} &
\mu_{|I_t|}(\kappa(\mathfrak m))
}
$$
；这是引理 \ref{more-algebra-lemma-compare-inertia} 的结论。
\end{remark}

\section{Krasner 引理}
\label{more-algebra-section-krasner}

\noindent
下面给出离散赋值域情形下的 Krasner 引理。

\begin{lemma}[Krasner 引理]
\label{more-algebra-lemma-krasner}
设 $A$ 是维数为 $1$ 的完备局部整环。设 $P(t) \in A[t]$
是系数在 $A$ 中的多项式。设 $\alpha \in A$ 是 $P$ 的根，
但不是导数 $P' = \text{d}P/\text{d}t$ 的根。对每个
$c \geq 0$，存在整数 $n$，使得对任意系数均属于
$\mathfrak m_A^n$ 的 $Q \in A[t]$，多项式 $P + Q$ 都有根
$\beta \in A$，并且 $\beta - \alpha \in \mathfrak m_A^c$。
\end{lemma}

\begin{proof}
选取非零元 $\pi \in \mathfrak m$。
% P02-E0614
由于 $A$ 的维数为 $1$，有 $\mathfrak m = \sqrt{(\pi)}$。
由假设，可对某个 $m \geq 0$ 与 $a \in A$ 写成
$P'(\alpha)^{-1} = \pi^{-m} a$。不妨设
$c \geq m + 1$。选取 $n$，使得
$\mathfrak m_A^n \subset (\pi^{c + m})$。任取命题中的
$Q$。为后面使用，注意可写成
% P02-E0618
$$
P(x + y) = P(x) + P'(x)y + R(x, y)y^2
$$
，其中某个 $R(x, y) \in A[x, y]$。我们将用归纳法证明，
可以找到序列 $\alpha_m, \alpha_{m + 1}, \alpha_{m + 2}, \ldots$，
使得
% P02-E0615
\begin{enumerate}
\item $\alpha_k \equiv \alpha \bmod \pi^c$，
\item $\alpha_{k + 1} - \alpha_k \in (\pi^k)$，并且
\item $(P + Q)(\alpha_k) \in (\pi^{m + k})$。
\end{enumerate}
置 $\beta = \lim \alpha_k$ 即可完成证明。

\medskip\noindent
初始情形。由于 $Q$ 的系数均属于 $(\pi^{c + m})$，有
$(P + Q)(\alpha) \in (\pi^{c + m})$。因此
$\alpha_m = \alpha$ 可行。这个选择保证对所有
$k \geq m$ 都有 $\alpha_k \equiv \alpha \bmod \pi^c$。
% P02-E0615

\medskip\noindent
归纳步骤。给定 $\alpha_k$，对某个 $\delta \in (\pi^k)$ 写
$\alpha_{k + 1} = \alpha_k + \delta$。于是
% P02-E0618
$$
(P + Q)(\alpha_{k + 1}) =
P(\alpha_k + \delta) + Q(\alpha_k + \delta)
$$
。由于 $Q$ 的系数均属于 $(\pi^{c + m})$，可知
$Q(\alpha_k + \delta) \equiv Q(\alpha_k) \bmod \pi^{c + m + k}$。
另一方面，
% P02-E0618
$$
P(\alpha_k + \delta) =
P(\alpha_k) + P'(\alpha_k)\delta + R(\alpha_k, \delta)\delta^2
$$
。注意，由于 $\alpha_k \equiv \alpha \bmod \pi^{m + 1}$，
有 $P'(\alpha_k) \equiv P'(\alpha) \bmod (\pi^{m + 1})$。
因此
% P02-E0618
$$
P(\alpha_k + \delta) \equiv P(\alpha_k) + P'(\alpha) \delta
\bmod \pi^{k + m + 1}
$$
。重新合并两项，得到
% P02-E0618
$$
(P + Q)(\alpha_{k + 1}) \equiv (P + Q)(\alpha_k) + P'(\alpha) \delta
\bmod \pi^{k + m + 1}
$$
。因此可取
$\delta = -P'(\alpha)^{-1} (P + Q)(\alpha_k) =  - \pi^{-m} a (P + Q)(\alpha_k)$；
由归纳假设，它属于 $(\pi^k)$。
\end{proof}

\begin{lemma}
\label{more-algebra-lemma-approximate-separable-extension}
设 $A$ 是分式域为 $K$ 的离散赋值环。设 $A^\wedge$ 是
$A$ 的完备化，其分式域为 $K^\wedge$。若 $M/K^\wedge$
是有限可分扩张，则存在有限可分扩张 $L/K$，使得
$M = K^\wedge \otimes_K L$。
\end{lemma}

\begin{proof}
注意，$A^\wedge$ 也是离散赋值环（由引理
\ref{more-algebra-lemma-completion-regular} 与
\ref{more-algebra-lemma-completion-dimension}）。特别地，$A^\wedge$
是整环。以下证明更一般地适用于满足下述条件的 Noether 局部环
$A$：其 $A^\wedge$ 是维数为 $1$ 的局部整环。

\medskip\noindent
取在 $K^\wedge$ 上生成 $M$ 的元素 $\theta \in M$。
（本原元定理。）设 $P(t) \in K^\wedge[t]$ 是
$\theta$ 在 $K^\wedge$ 上的极小多项式。取非零元
$\pi \in \mathfrak m_A$。把 $\theta$ 替换为
$\pi^n\theta$ 后，可设 $P(t)$ 的系数均属于
$A^\wedge$。令
$B = A^\wedge[\theta] = A^\wedge[t]/(P(t))$。注意，
$B$ 是维数为 $1$ 的完备局部整环，因为它在 $A$ 上有限且包含于
$M$。
% P02-E0616
由于 $M$ 在 $K$ 上可分，元素 $\theta$ 不是 $P$ 的导数的根。
% P02-E0617
对任意整数 $n$，可找到首一多项式 $P_1 \in A[t]$，使
$P - P_1$ 的系数均属于 $\pi^nA^\wedge[t]$。由 Krasner 引理
（引理 \ref{more-algebra-lemma-krasner}），当 $n$ 足够大时，$P_1$
在 $B$ 中有根 $\beta$。此外，可设（若 $n$ 取得足够大）
$\theta - \beta \in \pi B$。考虑把 $t$ 映到 $\beta$ 的
$A^\wedge$-代数映射
$\Phi : A^\wedge[t]/(P_1) \to B$。由于
$B = \pi B + \sum_{i < \deg(P)} A^\wedge \theta^i$，
由 Nakayama 引理，映射 $\Phi$ 是满射。因
$\deg(P_1) = \deg(P)$，可知 $\Phi$ 是同构。于是环扩张
$L = K[t]/(P_1(t))$ 满足 $K^\wedge \otimes_K L \cong M$。
这蕴含 $L$ 是域，证明完成。
\end{proof}

\begin{definition}
\label{more-algebra-definition-mixed}
设 $A$ 是离散赋值环。若 $A$ 的剩余域的特征为
$p > 0$，而 $A$ 的分式域的特征为 $0$，则称
$A$ 具有{\it 混合特征}。在这种情形下，得到离散赋值环扩张
$\mathbf{Z}_{(p)} \subset A$；$A$ 的{\it 绝对分歧指数}
是这个扩张的分歧指数。
\end{definition}
\section{Abhyankar 引理与驯分歧}
\label{more-algebra-section-abhyankar-tame}

\noindent
本节证明我们认为对离散赋值环而言最一般形式的 Abhyankar 引理。
随后应用它证明离散赋值环分式域的驯分歧扩张的一些结果。

\begin{remark}
\label{more-algebra-remark-construction}
设 $A \to B$ 是分式域分别为 $K \subset L$ 的离散赋值环扩张。
设 $K_1/K$ 是有限域扩张。设 $A_1 \subset K_1$ 是
$A$ 在 $K_1$ 中的整闭包。另一方面，令
$L_1 = (L \otimes_K K_1)_{red}$。则 $L_1$ 是由 $L$
的有限域扩张组成的非空有限乘积。设 $B_1$ 是 $B$ 在
$L_1$ 中的整闭包。得到相容的交换图
% P02-E0619
$$
\vcenter{
\xymatrix{
L \ar[r] & L_1 \\
K \ar[u] \ar[r] & K_1 \ar[u]
}
}
\quad\text{且}\quad
\vcenter{
\xymatrix{
B \ar[r] & B_1 \\
A \ar[u] \ar[r] & A_1 \ar[u]
}
}
$$
。在这种情形下，有如下结论：
\begin{enumerate}
\item 由 Algebra，引理
\ref{algebra-lemma-integral-closure-Dedekind}，环 $A_1$
是 Dedekind 整环，而 $B_1$ 是 Dedekind 整环的有限乘积。
\item 注意，$L \otimes_K K_1 = (B \otimes_A A_1)_\pi$，
其中 $\pi \in A$ 是一致化元，并且 $\pi$ 是
$B \otimes_A A_1$ 上的非零因子。因此环映射
$B \otimes_A A_1 \to B_1$ 是整的，其核由幂零元组成。
故 $\Spec(B_1) \to \Spec(B \otimes_A A_1)$ 在谱上是满射
（Algebra，引理
\ref{algebra-lemma-integral-overring-surjective}）。
映射 $\Spec(B \otimes_A A_1) \to \Spec(A_1)$ 是满射，因为
$A_1/\mathfrak m_A A_1 \to
B/\mathfrak m_AB \otimes_{\kappa_A} A_1/\mathfrak m_A A_1$
是单射环映射，且 $A_1/\mathfrak m_A A_1$ 是 Artin 环。
于是 $\Spec(B_1) \to \Spec(A_1)$ 是满射。
\item 设 $\mathfrak m_i$（$i = 1, \ldots n$，其中 $n \geq 1$）
是 $A_1$ 的极大理想。
% P02-E0626
对每个 $i = 1, \ldots, n$，设
$\mathfrak m_{ij}$（$j = 1, \ldots, m_i$，其中 $m_i \geq 1$）
是 $B_1$ 中位于 $\mathfrak m_i$ 之上的极大理想。得到离散赋值环扩张图
% P02-E0619
$$
\xymatrix{
B \ar[r] & (B_1)_{\mathfrak m_{ij}} \\
A \ar[u] \ar[r] & (A_1)_{\mathfrak m_i} \ar[u]
}
$$
。
\item 若 $A$ 是 Hensel 环（例如完备），则 $A_1$ 是离散赋值环，
即 $n = 1$。事实上，$A_1$ 是 $A$ 的有限扩张的并，且这些扩张均为整环，
故由 Algebra，引理
\ref{algebra-lemma-finite-over-henselian}，它们都是局部环。
\item 若 $B$ 是 Hensel 环（例如完备），则 $B_1$
是离散赋值环的乘积，即对 $i = 1, \ldots, n$ 有 $m_i = 1$。
\item 若 $K \subset K_1$ 是纯不可分扩张，则 $A_1$ 与 $B_1$
都是离散赋值环，即 $n = 1$ 且 $m_1 = 1$。这是因为对每个
$b \in B_1$，$b$ 的某个 $p$-幂属于 $B$，故 $B_1$
只能有一个极大理想。
% P02-E0620
\item 若 $K \subset K_1$ 是有限可分扩张，则
$L_1 = L \otimes_K K_1$，而且它也是有限可分扩张的有限乘积。
因此由 Algebra，引理
\ref{algebra-lemma-Noetherian-normal-domain-finite-separable-extension}，
$A \subset A_1$ 与 $B \subset B_1$ 都是有限的。
\item 若 $A$ 是 Nagata 环，则 $A \subset A_1$ 是有限的。
\item 若 $B$ 是 Nagata 环，则 $B \subset B_1$ 是有限的。
\end{enumerate}
\end{remark}

\begin{lemma}
\label{more-algebra-lemma-pull-root-uniformizer}
设 $A$ 是一致化元为 $\pi$ 的离散赋值环。设 $n \geq 2$。
令 $K_1 = K[\pi^{1/n}]$。则
\begin{enumerate}
\item $[K_1 : K] = n$；
\item $A$ 在 $K_1$ 中的整闭包 $A_1$ 是环 $A[\pi^{1/n}]$；
\item $A_1$ 是离散赋值环；
\item $A_1$ 在 $A$ 上的分歧指数为 $n$；
\item $K_1$ 关于 $A$ 完全分歧；并且
\item 若 $n$ 与 $A$ 的剩余特征互素，则 $K_1/K$ 驯分歧，
且 $K_1/K$ 的任意子扩张都由 $\pi^{1/d}$ 生成，其中 $d$
是 $n$ 的某个因子。
\end{enumerate}
\end{lemma}

\begin{proof}
考虑环 $A' = A[x]/(x^n - \pi)$，并用 $\pi'$ 表示 $x$ 在
$A_1$ 中的像。
% P02-E0621
% P02-E0622
由于 $A'$ 作为 $A$-模是秩为 $n$ 的有限自由模，元素
$\pi$ 在 $A'$ 中是非零因子，因而 $\pi'$ 也是非零因子。
此外，$A'/\pi'A'$ 同构于域 $A/\pi A$。故 $A'$ 是以
$\pi'$ 为一致化元的离散赋值环。因此，$A'$ 的分式域 $K'$
是次数为 $n$ 的扩张。显然，$K'$ 由在 $K$ 上添加 $\pi$
的一个 $n$ 次根 $\pi' = \pi^{1/n}$ 得到，即
$K' \cong K_1$。
% P02-E0623
由于 $A'$ 正规，可知 $A'$ 是 $A$ 在 $K'$ 中的整闭包，
即 $A' = A_1$。这证明了 (1)--(5)。

\medskip\noindent
由定义 \ref{more-algebra-definition-types-of-extensions} 及 (3)、(4)，
若 $n$ 在剩余域中可逆，则 $K_1/K$ 驯分歧的断言是直接的。
为证明最后一个陈述，只须证明每个 $n$ 次单位根
$\zeta \in K'$ 都属于 $K$；见 Fields，引理
\ref{fields-lemma-subfields-kummer}。显然
$\zeta \in (A')^*$。可写 $\zeta = a + b$，其中
$a \in A$ 且
$b \in A \pi' \oplus \ldots \oplus A (\pi')^{n - 1}$。
若 $b$ 非零，则可写 $b = u(\pi')^t$，其中
$t \geq 1$，且 $b$ 是 $A'$ 的单位。
% P02-E0624
于是
% P02-E0619
$$
1 = \zeta^n \equiv a^n + n a^{n - 1} u (\pi')^t \bmod (\pi')^{t + 1}
$$
。由于 $n$ 与 $a$ 都是 $A$ 的单位，这将推出
$b = u(\pi')^t$ 模 $(\pi')^{t + 1}$ 同余于
$(1 - a^n)/(n a^{n - 1}) \in A$。直和分解
$A' = A \oplus \ldots \oplus A(\pi')^{n - 1}$ 表明这是不可能的，
因为乘以 $(\pi')^{t + 1}$ 的像是相应的直和，而它与
$A(\pi')^t$ 的交为 $(\pi A)(\pi')^t$。故 $b = 0$，
并如愿得到 $\zeta \in A$。
\end{proof}

\begin{lemma}
\label{more-algebra-lemma-formally-smooth-goes-up}
设 $A \to B$ 是分式域为 $K \subset L$ 的离散赋值环扩张。
假设 $A \to B$ 在 $\mathfrak m_B$-进拓扑下形式光滑。
则对任意有限扩张 $K_1/K$，有
$L_1 = L \otimes_K K_1$、$B_1 = B \otimes_A A_1$，而且每个扩张
$(A_1)_{\mathfrak m_i} \subset (B_1)_{\mathfrak m_{ij}}$
（见评注 \ref{more-algebra-remark-construction}）在
$\mathfrak m_{ij}$-进拓扑下形式光滑。
\end{lemma}

\begin{proof}
我们将使用引理 \ref{more-algebra-lemma-extension-dvrs-formally-smooth} 的等价条件，
不再另行说明。设 $\pi \in A$ 与
$\pi_i \in (A_1)_{\mathfrak m_i}$ 分别为一致化元。
由于 $\kappa_A \subset \kappa_B$ 可分，环
% P02-E0619
$$
(B \otimes_A (A_1)_{\mathfrak m_i})/\pi_i (B \otimes_A (A_1)_{\mathfrak m_i}) =
B/\pi B \otimes_{A/\pi A} (A_1)_{\mathfrak m_i}/\pi_i (A_1)_{\mathfrak m_i}
$$
是若干域的乘积，其中每个域在 $\kappa_{\mathfrak m_i}$ 上可分。
因此 $B \otimes_A (A_1)_{\mathfrak m_i}$ 中的元素 $\pi_i$
是非零因子，且对它取商得到域的乘积。由此
$B \otimes_A A_1$ 是 Dedekind 整环，特别地是既约的。
% P02-E0625
因此 $B \otimes_A A_1 \subset B_1$ 是等式。
\end{proof}


\noindent
下面的引理是我们针对离散赋值环所用的 Abhyankar 引理版本。
注意，并不假设 $\kappa_B/\kappa_A$ 是代数域扩张。

\begin{lemma}[Abhyankar 引理]
\label{more-algebra-lemma-abhyankar}
设 $A \subset B$ 是离散赋值环扩张。假设 $A$ 的剩余特征为
$0$；或者其剩余特征为 $p$，分歧指数 $e$ 与 $p$ 互素，且
$\kappa_B/\kappa_A$ 是可分域扩张。设 $K_1/K$ 是有限扩张。
使用评注 \ref{more-algebra-remark-construction} 的记号，并假设对某个 $i$，
$e$ 整除 $A \subset (A_1)_{\mathfrak m_i}$ 的分歧指数。
则对所有 $j = 1, \ldots, m_i$，扩张
$(A_1)_{\mathfrak m_i} \subset (B_1)_{\mathfrak m_{ij}}$
在 $\mathfrak m_{ij}$-进拓扑下形式光滑。
\end{lemma}

\begin{proof}
设 $\pi \in A$ 是一致化元。设 $\pi_1$ 是
$(A_1)_{\mathfrak m_i}$ 的一致化元。写成
$\pi = u \pi_1^{e_1}$，其中 $u$ 是
$(A_1)_{\mathfrak m_i}$ 的单位，而 $e_1$ 是
$A \subset (A_1)_{\mathfrak m_i}$ 的分歧指数。

\medskip\noindent
断言：可设 $u$ 是 $K_1$ 中的 $e$ 次幂。事实上，设
$K_2$ 是通过添加 $x^e = u$ 的一个根而得到的 $K_1$ 的扩张；
于是 $K_2$ 是 $K_1[x]/(x^e - u)$ 的一个因子。由关于 $e$
的假设，$K_2/K_1$ 是有限可分扩张，因而 $A_1 \subset A_2$
是有限的。由于
$(A_1)_{\mathfrak m_i} \to (A_1)_{\mathfrak m_i}[x]/(x^e - u)$
是有限 \'etale 的（因为 $e$ 与剩余特征互素且 $u$ 是单位），
可知 $(A_2)_{\mathfrak m_i}$ 是
$(A_1)_{\mathfrak m_i}$ 的某个有限 \'etale 扩张的因子，
故其本身在 $(A_1)_{\mathfrak m_i}$ 上有限 \'etale。
% P02-E0627
同样的推理表明，$B_1 \subset B_2$ 诱导有限 \'etale 扩张
$(B_1)_{\mathfrak m_{ij}} \subset (B_2)_{\mathfrak m_{ij}}$。
% P02-E0627
选取极大理想 $\mathfrak m'_{ij} \subset B_2$，使其位于
$\mathfrak m_{ij} \subset B_1$ 之上（当然可能不止一个），并考虑
% P02-E0630
$$
\xymatrix{
(B_1)_{\mathfrak m_{ij}} \ar[r] & (B_2)_{\mathfrak m'_{ij}} \\
(A_1)_{\mathfrak m_i} \ar[u] \ar[r] &
(A_2)_{\mathfrak m'_i} \ar[u]
}
$$
，其中 $\mathfrak m'_i \subset A_2$ 是其像。
% P02-E0628
现在，水平箭头的分歧指数为 $1$，并诱导有限可分剩余域扩张。
因此，使用引理 \ref{more-algebra-lemma-extension-dvrs-formally-smooth} 的等价条件，
只须证明右侧竖直箭头在 $\mathfrak m'_{ij}$-进拓扑下形式光滑。
由于 $u$ 在 $K_2$ 中有一个 $e$ 次根，断言得证。
% P02-E0629

\medskip\noindent
假设 $u$ 在 $K_1$ 中有一个 $e$ 次根。由于 $e | e_1$，且
$u$ 在 $K_1$ 中有一个 $e$ 次根，可知对某个
$\theta \in K_1$ 有 $\pi = \theta^e$。
% P02-E0629
设 $K'_1 = K[\theta] \subset K_1$ 是由 $\theta$ 生成的子域。
由引理 \ref{more-algebra-lemma-pull-root-uniformizer}，$A$ 在
$K[\theta]$ 中的整闭包 $A'_1$ 是离散赋值环
$A'_1 = A[\theta]$，它在 $A$ 上的分歧指数为 $e$。
若能对扩张 $K'_1/K$ 证明本引理，则对图
% P02-E0630
$$
\xymatrix{
(B'_1)_{B'_1 \cap \mathfrak m_{ij}} \ar[r] & (B_1)_{\mathfrak m_{ij}} \\
A'_1 \ar[u] \ar[r] &
(A_1)_{\mathfrak m_i} \ar[u]
}
$$
应用引理 \ref{more-algebra-lemma-formally-smooth-goes-up}，即可对所有
$j = 1, \ldots, m_i$ 得到结论。这把问题约化为下一段讨论的情形。

\medskip\noindent
假设 $K_1 = K[\pi^{1/e}]$，并置 $\theta = \pi^{1/e}$。
设 $\pi_B$ 是 $B$ 的一致化元，并对 $B$ 的某个单位 $w$
写成 $\pi = w \pi_B^e$。于是可见
$L_1 = L \otimes_K K_1$ 通过添加 $\theta / \pi_B$ 得到，
后者是单位 $w$ 的一个 $e$ 次根。因此 $B \subset B_1$
是有限 \'etale 的。于是对任意极大理想
$\mathfrak m \subset B_1$，考虑交换图
% P02-E0630
$$
\xymatrix{
B \ar[r]_1 & (B_1)_{\mathfrak m} \\
A \ar[u]^e \ar[r]^e & A_1 \ar[u]_{e_\mathfrak m}
}
$$
。箭头旁的数是分歧指数。由分歧指数的乘法性
（引理 \ref{more-algebra-lemma-multiplicative-e-f}），得到
$e_\mathfrak m = 1$。考察剩余域扩张，可知
$\kappa(\mathfrak m)$ 是 $\kappa_B$ 的有限可分扩张，
而后者在 $\kappa_A$ 上可分。因此
$\kappa(\mathfrak m)$ 在 $\kappa_A$ 上可分；后者
等于 $A_1$ 的剩余域。于是由引理
\ref{more-algebra-lemma-extension-dvrs-formally-smooth} 得到所需结论。
\end{proof}

\begin{lemma}
\label{more-algebra-lemma-composition-tame}
设 $A$ 是分式域为 $K$ 的离散赋值环。设 $M/L/K$
是有限可分扩张。设 $B$ 是 $A$ 在 $L$ 中的整闭包。
若 $L/K$ 关于 $A$ 驯分歧，并且对 $B$ 的每个极大理想
$\mathfrak m$，$M/L$ 关于 $B_\mathfrak m$ 驯分歧，
则 $M/K$ 关于 $A$ 驯分歧。
\end{lemma}

\begin{proof}
设 $C$ 是 $A$ 在 $M$ 中的整闭包。$C$ 的每个极大理想
$\mathfrak m'$ 都位于 $B$ 的某个极大理想 $\mathfrak m$ 之上。
于是本引理由分歧指数的乘法性（引理
\ref{more-algebra-lemma-multiplicative-e-f}）以及有限域扩张塔
$\kappa(\mathfrak m')/\kappa(\mathfrak m)/\kappa_A$ 得出。
\end{proof}

\begin{lemma}
\label{more-algebra-lemma-subextension-tame}
设 $A$ 是分式域为 $K$ 的离散赋值环。若 $M/L/K$
是有限可分扩张，且 $M$ 关于 $A$ 驯分歧，则 $L$
关于 $A$ 驯分歧。
\end{lemma}

\begin{proof}
我们将使用评注 \ref{more-algebra-remark-finite-separable-extension} 中讨论的结果，
不再另行说明。设 $C/B/A$ 是 $A$ 在 $M/L/K$ 中的整闭包。
% P02-E0631
由于 $C$ 是 $B$ 的有限环扩张，可知
$\Spec(C) \to \Spec(B)$ 是满射。因此，对每个极大理想
$\mathfrak m \subset B$，都存在极大理想
$\mathfrak m' \subset C$ 位于 $\mathfrak m$ 之上。
% P02-E0632
由分歧指数的乘法性（引理
\ref{more-algebra-lemma-multiplicative-e-f}）及假设，得到
$B_\mathfrak m$ 在 $A$ 上的分歧指数与剩余特征互素。
由于 $\kappa(\mathfrak m')/\kappa_A$ 有限可分，
$\kappa(\mathfrak m)/\kappa_A$ 也具有同一性质。
\end{proof}


\begin{lemma}
\label{more-algebra-lemma-characterize-tame}
设 $A$ 是分式域为 $K$ 的离散赋值环。设
$\pi \in A$ 是一致化元。设 $L/K$ 是有限可分扩张。
下列条件等价：
\begin{enumerate}
\item $L$ 关于 $A$ 驯分歧；
\item 存在一个在 $\kappa_A$ 中可逆的 $e \geq 1$，以及一个关于
$A' = A[\pi^{1/e}]$ 非分歧的扩张
$L'/K' = K[\pi^{1/e}]$，使得 $L$ 包含于 $L'$；并且
\item 存在一个在 $\kappa_A$ 中可逆的 $e_0 \geq 1$，使得对每个
在 $\kappa_A$ 中可逆的 $d \geq 1$，(2) 对
$e = de_0$ 成立。
\end{enumerate}
\end{lemma}

\begin{proof}
注意，$A'$ 是分式域为 $K'$ 的离散赋值环；见引理
\ref{more-algebra-lemma-pull-root-uniformizer}。当然，$A'$ 在 $A$ 上的分歧指数为
$e$。因此，若 (2) 成立，则由引理 \ref{more-algebra-lemma-composition-tame}，
$L'$ 关于 $A$ 驯分歧。再由引理
\ref{more-algebra-lemma-subextension-tame}，$L$ 关于 $A$ 驯分歧。

\medskip\noindent
蕴含 (3) $\Rightarrow$ (2) 是直接的。

\medskip\noindent
假设 (1) 成立。设 $B$ 是 $A$ 在 $L$ 中的整闭包，并设
$\mathfrak m_1, \ldots, \mathfrak m_n$ 是它的极大理想。
用 $e_i$ 表示 $A \to B_{\mathfrak m_i}$ 的分歧指数。
% P02-E0633
设 $e_0$ 是 $e_1, \ldots, e_r$ 的最小公倍数。
% P02-E0634
由假设 (1)，它在 $\kappa_A$ 中可逆。按 (3) 令
$e = de_0$。置 $A' = A[\pi^{1/e}]$。则
$A \to A'$ 是分式域为 $K' = K[\pi^{1/e}]$ 的离散赋值环扩张；
见引理 \ref{more-algebra-lemma-pull-root-uniformizer}。选取乘积分解
% P02-E0636
$$
L \otimes_K K' = \prod L'_j
$$
，其中 $L'_j$ 是域。设 $B'_j$ 是 $A$ 在 $L'_j$ 中的整闭包。
设 $\mathfrak m_{ijk}$ 是 $B'_j$ 中位于 $\mathfrak m_i$
之上的极大理想。注意，$(B'_j)_{\mathfrak m_i}$ 是
$B_{\mathfrak m_i}$ 在 $L'_j$ 中的整闭包。
% P02-E0635
把 Abhyankar 引理（引理 \ref{more-algebra-lemma-abhyankar}）应用于
$A \subset B_{\mathfrak m_i}$ 与扩张 $K'/K$，可知
$A' \to (B'_j)_{\mathfrak m_{ijk}}$ 在
$\mathfrak m_{ijk}$-进拓扑下形式光滑。这蕴含分歧指数为
$1$ 且剩余域扩张可分（引理
\ref{more-algebra-lemma-extension-dvrs-formally-smooth}）。
由此可见，$L'_j$ 关于 $A'$ 非分歧。证明完成：任取某个
$j$，令 $L' = L'_j$。
\end{proof}


\begin{lemma}
\label{more-algebra-lemma-permanence-tame}
设 $A$ 是分式域为 $K$ 的离散赋值环。
\begin{enumerate}
\item 若 $L/K$ 是关于 $A$ 驯分歧的有限可分扩张，则存在
包含 $L$ 且关于 $A$ 驯分歧的 Galois 扩张 $M/K$。
\item 若 $L_1/K$、$L_2/K$ 是关于 $A$ 驯分歧的有限可分扩张，
则存在包含 $L_1$ 与 $L_2$ 且关于 $A$ 驯分歧的有限可分扩张
$L/K$。
\end{enumerate}
\end{lemma}

\begin{proof}
证明 (2)。选取一致化元 $\pi \in A$。可以选取在
$\kappa_A$ 中可逆的整数 $e$，以及关于
$A' = A[\pi^{1/e}]$ 非分歧的扩张
$L_i'/K' = K[\pi^{1/e}]$，使得作为 $K$ 的扩张有
$L'_i/L_i$；见引理 \ref{more-algebra-lemma-characterize-tame}。
% P02-E0637
由引理 \ref{more-algebra-lemma-permanence-unramified}，可以找到关于 $A'$
非分歧的扩张 $L'/K'$，使得对 $i = 1, 2$，
$L'_i/K$ 同构于 $L'/K'$ 的一个子扩张。
% P02-E0638
由于 $L'/K$ 驯分歧，这完成了 (3) 的证明
（使用上面同一个引理）。
% P02-E0639

\medskip\noindent
证明 (1)。首先可以把 $L$ 替换为更大的扩张，并假设
$L$ 是 $K' = K[\pi^{1/e}]$ 的扩张，关于
$A' = A[\pi^{1/e}]$ 非分歧，其中 $e$ 在 $\kappa_A$
中可逆；见引理 \ref{more-algebra-lemma-characterize-tame}。
设 $M$ 是 $L$ 在 $K$ 上的正规闭包；见 Fields，定义
\ref{fields-definition-normal-closure}。由 Fields，引理
\ref{fields-lemma-normal-closure-galois}，$M/K$ 是 Galois 扩张。
另一方面，有 $K$-代数满射
% P02-E0640
% P02-E0642
$$
L \otimes_K \ldots \otimes_K L \longrightarrow M
$$
；见 Fields，引理
\ref{fields-lemma-normal-closure-tensor-product}。
如评注 \ref{more-algebra-remark-finite-separable-extension}，设 $B$ 是
$A$ 在 $L$ 中的整闭包。$L$ 关于
$A' = A[\pi^{1/e}]$ 非分歧这一条件，恰好意味着
$A' \to B$ 是 \'etale 环映射；见 Algebra，引理
\ref{algebra-lemma-characterize-etale}。
断言：
% P02-E0640
% P02-E0642
$$
K' \otimes_K \ldots \otimes_K K' = \prod K'_i
$$
是关于 $A$ 驯分歧的域扩张 $K'_i/K$ 的乘积。于是，若
$A'_i$ 是 $A$ 在 $K'_i$ 中的整闭包，则
% P02-E0640
% P02-E0642
$$
\prod A'_i \otimes_{(A' \otimes_A \ldots \otimes_A A')}
(B \otimes_A \ldots \otimes_A B)
$$
在 $\prod A'_i$ 上有限 \'etale，因而是 Dedekind 整环的乘积
（引理 \ref{more-algebra-lemma-Dedekind-etale-extension}）。
可知 $M$ 是这些 Dedekind 整环之一的分式域，而该整环在
$A'_i$ 上有限 \'etale，其中 $i$ 为某个指标。因此 $M/K'_i$ 关于
$A'_i$ 的每个极大理想都非分歧，从而由引理
\ref{more-algebra-lemma-composition-tame}，$M/K$ 驯分歧。

\medskip\noindent
还需证明该断言。为此写成
$A' = A[x]/(x^e - \pi)$，并且可见
% P02-E0641
% P02-E0640
% P02-E0642
$$
A' \otimes_A \ldots \otimes_A A' =
A'[x_1, \ldots, x_r]/(x_1^e - \pi, \ldots, x_r^e - \pi)
$$
该环的正规化当然包含元素 $y_i = x_i/x_1$（其中
$i = 2, \ldots, r$），它们满足关系 $y_i^e - 1 = 0$，并得到
% P02-E0642
$$
A[x_1, y_2, \ldots, y_r]/(x_1^e - \pi, y_2^e - 1, \ldots, y_r^e - 1) =
A'[y_2, \ldots, y_r]/(y_2^e - 1, \ldots, y_r^e - 1)
$$
。由于 $e$ 在 $A'$ 中可逆，该环在 $A'$ 上有限 \'etale。
因此它是 Dedekind 整环的乘积，每个因子在 $A'$ 上非分歧，
如愿得证（若有疑惑，参见上面给出的参考文献）。
\end{proof}


\begin{lemma}
\label{more-algebra-lemma-tame-goes-up}
设 $A \subset B$ 是离散赋值环扩张。用 $L/K$ 表示相应的分式域扩张。
设 $K'/K$ 是有限可分扩张。则
% P02-E0644
$$
K' \otimes_K L = \prod L'_i
$$
是域的有限乘积，并且下列结论成立：
\begin{enumerate}
\item 若 $K'$ 关于 $A$ 非分歧，则每个 $L'_i$ 关于 $B$ 非分歧。
\item 若 $K'$ 关于 $A$ 驯分歧，则每个 $L'_i$ 关于 $B$ 驯分歧。
\end{enumerate}
\end{lemma}

\begin{proof}
代数 $K' \otimes_K L$ 是 $L$ 上的有限 \'etale 代数，
故为域的有限乘积。设 $A'$ 是 $A$ 在 $K'$ 中的整闭包。

\medskip\noindent
在情形 (1) 中，环映射 $A \to A'$ 有限 \'etale。因此
$B' = B \otimes_A A'$ 在 $B$ 上有限 \'etale，而且是
Dedekind 整环的有限乘积（引理
\ref{more-algebra-lemma-Dedekind-etale-extension}）。故 $B'$ 是
$B$ 在 $K' \otimes_K L$ 中的整闭包。立即可知每个
$L'_i$ 关于 $B$ 非分歧。

\medskip\noindent
选取一致化元 $\pi \in A$。为证明 (2)，可把 $K'$ 替换为
关于 $A$ 驯分歧的更大扩张（略去细节；提示：使用引理
\ref{more-algebra-lemma-subextension-tame}）。
% P02-E0643
因此，由引理 \ref{more-algebra-lemma-characterize-tame}，可设存在某个在
$\kappa_A$ 中可逆的 $e \geq 1$，使得 $K'$ 包含
$K[\pi^{1/e}]$，并且 $K'$ 关于 $A[\pi^{1/e}]$ 非分歧。
选取乘积分解
% P02-E0644
$$
K[\pi^{1/e}] \otimes_K L = \prod L_{e, j}
$$
。对每个 $i$，存在 $j_i$，使得 $L'_i/L_{e, j_i}$
是有限可分扩张。设 $B_{e, j}$ 是 $B$ 在 $L_{e, j}$
中的整闭包。把 (1) 应用于 $K'/K[\pi^{1/e}]$ 与
$A[\pi^{1/e}] \subset (B_{e, j_i})_\mathfrak m$，可知
对每个极大理想 $\mathfrak m \subset B_{e, j_i}$，
$L'_i$ 关于 $(B_{e, j_i})_\mathfrak m$ 非分歧。
因此，由引理 \ref{more-algebra-lemma-composition-tame}，只须证明
$L_{e, j}$ 关于 $B$ 驯分歧。

\medskip\noindent
在 $B$ 中选取一致化元 $\theta$。写成
$\pi = u \theta^t$，其中 $u$ 是 $B$ 的单位且 $t \geq 1$。
于是
% P02-E0644
$$
A[\pi^{1/e}] \otimes_A B = B[x]/(x^e - u \theta^t)
\subset B[y, z]/(y^{e'} - \theta, z^e - u)
$$
，其中 $e' = e/\gcd(e, t)$。该映射把 $x$ 映到
$z y^{t/\gcd(e, t)}$。由于右侧是 Dedekind 整环的乘积，
其中每个因子在 $B$ 上驯分歧，证明完成（略去细节）。
\end{proof}


\section{消除分歧}
\label{more-algebra-section-eliminating-ramification}

\noindent
本节讨论 Helmut Epp 的一个结果，见 \cite{Epp}。我们强烈建议读者阅读原文。
我们的方法略有不同，因为我们试图用同一种方法处理混合特征与等特征情形。
相关结果还可参见
\cite{Ponomarev}、\cite{Ponomarev-Abhyankar}、\cite{Kuhlmann}
与 \cite{ZK}。

\medskip\noindent
设 $A \subset B$ 是分式域为 $K \subset L$ 的离散赋值环扩张。
本节的目标是找到有限扩张 $K_1/K$，使得在
% P02-E0645
$$
\vcenter{
\xymatrix{
L \ar[r] & L_1 \\
K \ar[u] \ar[r] & K_1 \ar[u]
}
}
\quad\text{且}\quad
\vcenter{
\xymatrix{
B \ar[r] & B_1 \ar[r] & (B_1)_{\mathfrak m_{ij}} \\
A \ar[u] \ar[r] & A_1 \ar[r] \ar[u] & (A_1)_{\mathfrak m_i} \ar[u]
}
}
$$
如评注 \ref{more-algebra-remark-construction} 所示时，扩张
$(A_1)_{\mathfrak m_i} \subset (B_1)_{\mathfrak m_{ij}}$
全都弱非分歧，乃至在相应进拓扑下形式光滑。
最简单（但非平凡）的例子是 Abhyankar 引理，见引理
\ref{more-algebra-lemma-abhyankar}。

\begin{definition}
\label{more-algebra-definition-solution}
设 $A \to B$ 是分式域为 $K \subset L$ 的离散赋值环扩张。
\begin{enumerate}
\item 若评注 \ref{more-algebra-remark-construction} 的所有扩张
$(A_1)_{\mathfrak m_i} \subset (B_1)_{\mathfrak m_{ij}}$
都弱非分歧，则称有限域扩张 $K_1/K$ 是
$A \subset B$ 的一个{\it 弱解}。
\item 若评注 \ref{more-algebra-remark-construction} 的每个扩张
$(A_1)_{\mathfrak m_i} \subset (B_1)_{\mathfrak m_{ij}}$
都在 $\mathfrak m_{ij}$-进拓扑下形式光滑，则称有限域扩张
$K_1/K$ 是 $A \subset B$ 的一个{\it 解}。
\end{enumerate}
若解 $K_1/K$ 满足 $K_1/K$ 可分，则称它是一个{\it 可分的解}。
\end{definition}

\noindent
一般而言，弱解或解未必存在；\cite{Epp} 中有一个例子。
在对剩余域扩张施加温和假设后，我们将依照 \cite{Epp} 在定理
\ref{more-algebra-theorem-epp} 中证明弱解的存在性。下一节将从中推出几何上有意义的
情形中解的存在性，有时还可推出可分的解的存在性；见命题
\ref{more-algebra-proposition-epp-essentially-finite-type} 与引理
\ref{more-algebra-lemma-epp-essentially-finite-type-separable}。
不过，下面的例子表明，一般而言，即使为了得到弱解也需要不可分扩张。

\begin{example}
\label{more-algebra-example-inseparable-necessary}
设 $k$ 是特征为 $p > 0$ 的完美域。令 $A = k[[x]]$ 且
$K = k((x))$。令 $B = A[x^{1/p}]$。$A \to B$ 的任意弱解
$K_1/K$ 都不可分（而 $K$ 的任意有限不可分扩张都是解）。
略去证明。
\end{example}

\noindent
解在进一步扩张下稳定，见引理 \ref{more-algebra-lemma-solution-goes-up}。
弱解未必具有这一性质。弱解在某种意义下对完全分歧扩张稳定，见引理
\ref{more-algebra-lemma-weakly-unramified-goes-up-along-totally-ramified}。

\begin{lemma}
\label{more-algebra-lemma-weakly-unramified-goes-up-along-totally-ramified}
设 $A \to B$ 是分式域为 $K \subset L$ 的离散赋值环扩张。
假设 $A \to B$ 弱非分歧。则对任意关于 $A$ 完全分歧的有限可分扩张
$K_1/K$，$L_1 = L \otimes_K K_1$ 是域，$A_1$ 与
$B_1 = B \otimes_A A_1$ 是离散赋值环，而且扩张
$A_1 \subset B_1$（见评注 \ref{more-algebra-remark-construction}）弱非分歧。
\end{lemma}

\begin{proof}
设 $\pi \in A$ 与 $\pi_1 \in A_1$ 是一致化元。由于
$K_1/K$ 关于 $A$ 完全分歧，对 $A_1$ 中的某个单位 $u_1$ 有
$\pi_1^e = u_1 \pi$。因此 $A_1$ 在 $A$ 上由 $\pi_1$ 生成，
而 $\pi_1$ 在 $K$ 上的极小多项式 $P(t)$ 具有形式
% P02-E0645
$$
P(t) = t^e + a_{e - 1} t^{e - 1} + \ldots + a_0
$$
，其中 $a_i \in (\pi)$，且对 $A$ 的某个单位 $u$ 有
$a_0 = u\pi$。还要注意，$e = [K_1 : K]$。由于
$A \to B$ 弱非分歧，$\pi$ 是 $B$ 的一致化元，故
$B_1 = B[t]/(P(t))$ 是离散赋值环，其一致化元是 $t$ 的类。
引理于是显然。
\end{proof}

\begin{lemma}
\label{more-algebra-lemma-solutions-go-down}
设 $A \to B \to C$ 是分式域为
$K \subset L \subset M$ 的离散赋值环扩张。设 $K_1/K$ 是有限扩张。
\begin{enumerate}
\item 若 $K_1$ 是 $A \to C$ 的弱解或解，则 $K_1$ 是
$A \to B$ 的弱解或解。
\item 若 $K_1$ 是 $A \to B$ 的弱解或解，且
$L_1 = (L \otimes_K K_1)_{red}$ 是若干域的乘积，其中每个域都是
$B \to C$ 的弱解或解，则 $K_1$ 是 $A \to C$ 的弱解或解。
\end{enumerate}
\end{lemma}

\begin{proof}
令 $L_1 = (L \otimes_K K_1)_{red}$ 与
$M_1 = (M \otimes_K K_1)_{red}$，并设
$B_1 \subset L_1$ 与 $C_1 \subset M_1$ 是 $B$ 与 $C$ 的整闭包。
% P02-E0646
注意，$M_1 = (M \otimes_L L_1)_{red}$，而 $L_1$ 是由
$L$ 的有限扩张组成的非空有限乘积。因此环映射
$B_1 \to C_1$ 是评注 \ref{more-algebra-remark-construction} 中所讨论形式的环映射的
有限乘积。特别地，映射 $\Spec(C_1) \to \Spec(B_1)$ 是满射。
选取极大理想 $\mathfrak m \subset C_1$，并考虑离散赋值环扩张
% P02-E0645
$$
(A_1)_{A_1 \cap \mathfrak m} \to
(B_1)_{B_1 \cap \mathfrak m} \to
(C_1)_\mathfrak m
$$
。若复合映射弱非分歧，则映射
$(A_1)_{A_1 \cap \mathfrak m} \to
(B_1)_{B_1 \cap \mathfrak m}$ 也弱非分歧。
若剩余域扩张
$\kappa_{A_1 \cap \mathfrak m} \to \kappa_\mathfrak m$
可分，则子扩张
$\kappa_{A_1 \cap \mathfrak m} \to
\kappa_{B_1 \cap \mathfrak m}$ 也可分。结合引理
\ref{more-algebra-lemma-extension-dvrs-formally-smooth}，这证明了 (1)。
类似论证适用于 (2)。
\end{proof}

\begin{lemma}
\label{more-algebra-lemma-solution-after-strict-henselization}
设 $A \to B$ 是离散赋值环扩张。存在离散赋值环扩张的交换图
% P02-E0645
$$
\xymatrix{
B \ar[r] & B' \\
A \ar[r] \ar[u] & A' \ar[u]
}
$$
，使得
\begin{enumerate}
\item 分式域扩张 $K'/K$ 与 $L'/L$ 可分代数；
\item $A'$ 与 $B'$ 的剩余域分别是 $A$ 与 $B$ 的剩余域的
可分代数闭包；并且
\item 若 $A' \to B'$ 存在解、弱解或可分的解，则
$A \to B$ 存在解、弱解或可分的解。
\end{enumerate}
\end{lemma}

\begin{proof}
由 Algebra，引理
\ref{algebra-lemma-colimit-finite-etale-given-residue-field}，
存在扩张 $A \subset A'$，它是有限 \'etale 扩张的滤过余极限，
且 $A'$ 的剩余域是 $A$ 的剩余域的可分代数闭包。
于是 $A \subset A'$ 是离散赋值环扩张，所诱导的分式域扩张
$K'/K$ 可分代数。

\medskip\noindent
设 $B \subset B'$ 是 $B$ 的严格 Hensel 化。则
$B \subset B'$ 是离散赋值环扩张，其分式域扩张可分代数。
由 Algebra，引理
\ref{algebra-lemma-strictly-henselian-functorial-prepare}，
存在引理陈述中的交换图。引理的 (1)、(2) 是清楚的。

\medskip\noindent
设 $K'_1/K'$ 是 $A' \to B'$ 的弱解或解。由于 $A'$ 是余极限，
可找到有限 \'etale 扩张 $A \subset A_1'$，以及 $A_1'$
的分式域 $F$ 的有限扩张 $K_1$，使得
$K'_1 = K' \otimes_F K_1$。由于 $A \subset A_1'$ 有限
\'etale 且 $B'$ 严格 Hensel，$B' \otimes_A A_1'$ 是与
$B'$ 同构的环的有限乘积。因此
% P02-E0645
$$
L' \otimes_K K_1 = L' \otimes_K F \otimes_F K_1
$$
是与 $L' \otimes_{K'} K'_1$ 同构的环的有限乘积。
由此可见，$K_1/K$ 是 $A \to B'$ 的弱解或解。
故由引理 \ref{more-algebra-lemma-solutions-go-down}，它也是
$A \to B$ 的弱解或解。
% P02-E0647
\end{proof}

\begin{lemma}
\label{more-algebra-lemma-galois-relative}
设 $A \to B$ 是分式域为 $K \subset L$ 的离散赋值环扩张。设
$K_1/K$ 是正规扩张，并令 $G = \text{Aut}(K_1/K)$。则 $G$
作用在评注 \ref{more-algebra-remark-construction} 的环 $K_1$、$L_1$、$A_1$
与 $B_1$ 上，并且在 $B_1$ 的极大理想集合上传递作用。
% P02-E0648
\end{lemma}

\begin{proof}
除最后一个断言外，一切都是清楚的。若该作用有两个或更多轨道，则可找到
$b \in B_1$，它在一个轨道的所有极大理想处为零，而在另一个轨道的所有
极大理想处的剩余类都是 $1$。于是
$b' = \prod_{\sigma \in G} \sigma(b)$ 是
$B_1 \subset L_1 = (L \otimes_K K_1)_{red}$ 的一个 $G$-不变元素；
它属于 $B_1$ 的某些极大理想，却不属于 $B_1$ 的所有极大理想。将它提升为
$L \otimes_K K_1$ 的一个元素并取充分高次幂，可得到
$L \otimes_K K_1$ 的一个 $G$-不变元素 $b''$，它对某个 $N > 0$
映到 $(b')^N$；事实上，只有在特征为 $p > 0$ 时才需要这样做，而在这种情形，
取充分大的 $p$ 次幂 $q$ 给出典范映射
$(L \otimes_K K_1)_{red} \to L \otimes_K K_1$。
由于 $K = (K_1)^G$，可得 $b'' \in L$。又因为 $b''$ 映到 $B_1$
的一个元素，故 $b'' \in B$（因为 $B$ 正规）。一方面，由于 $b'$ 属于
$B_1$ 的某个极大理想，必有 $b'' \in \mathfrak m_B$；另一方面，由于
$b'$ 不属于 $B_1$ 的所有极大理想，又必有
$b'' \not \in \mathfrak m_B$。这一矛盾完成了引理的证明。
\end{proof}

\begin{lemma}
\label{more-algebra-lemma-make-degree-q-extension}
设 $A$ 是以 $\pi$ 为一致化元的离散赋值环。若 $A$ 的剩余特征为
$p > 0$，则对每个 $n > 1$ 和 $p$ 的幂 $q$，都存在关于 $A$
完全分歧的次数为 $q$ 的可分扩张 $L/K$，使得 $A$ 在 $L$ 中的整闭包
$B$ 的分歧指数为 $q$，并且有一致化元 $\pi_B$，满足
$\pi_B^q = \pi + \pi^n b$ 和
$\pi_B^q = \pi + (\pi_B)^{nq}b'$，其中某些 $b, b' \in B$。
\end{lemma}

\begin{proof}
若 $K$ 的特征为零，则可取由 $\pi_B^q = \pi$ 给出的扩张，见引理
\ref{more-algebra-lemma-pull-root-uniformizer}。若 $K$ 的特征为 $p > 0$，则可取由
$z^q - \pi^n z = \pi^{1 - q}$ 给出的 $K$ 的扩张。事实上，令
$y = \pi z$，便有 $y^q - \pi^{n + q - 1} y = \pi$。取
$\pi_B = y$ 即得所求结论。
\end{proof}

\begin{lemma}
\label{more-algebra-lemma-pre-purely-inseparable-case}
% P02-E0649
设 $A$ 是离散赋值环。假设剩余域 $\kappa_A$ 的特征为 $p > 0$，并且
$a \in A$ 的剩余类在 $\kappa_A$ 中不是 $p$ 次幂。则 $a$ 在 $K$
中也不是 $p$ 次幂，而且 $A$ 在 $K[a^{1/p}]$ 中的整闭包是环
$A[a^{1/p}]$；这是 $A$ 上弱非分歧的离散赋值环。
\end{lemma}

\begin{proof}
本引理不证自明。
\end{proof}

\begin{lemma}
\label{more-algebra-lemma-purely-inseparable-case}
% P02-E0650
设 $A \subset B \subset C$ 是分式域为 $K \subset L \subset M$
的离散赋值环扩张，并设 $\pi \in A$ 是一致化元。假设
\begin{enumerate}
\item $B$ 是 Nagata 环；
\item $A \subset B$ 弱非分歧；
\item $M$ 是 $L$ 的次数为 $p$ 的纯不可分扩张。
\end{enumerate}
则下列情形之一成立：
\begin{enumerate}
\item $A \to C$ 弱非分歧；
\item $C = B[\pi^{1/p}]$；
\item 存在关于 $A$ 完全分歧的次数为 $p$ 的可分扩张 $K_1/K$，使得
$L_1 = L \otimes_K K_1$ 与 $M_1 = M \otimes_K K_1$ 都是域，
并且整闭包之间的映射 $A_1 \to B_1 \to C_1$ 都是离散赋值环的
弱非分歧扩张。
\end{enumerate}
\end{lemma}

\begin{proof}
设 $e$ 是 $C$ 在 $B$ 上的分歧指数。若 $e = 1$，则结论成立。
否则，由引理 \ref{more-algebra-lemma-inequality} 与
\ref{more-algebra-lemma-ramification-index-a-power-of-p} 可知 $e = p$。
这又蕴含 $B$ 与 $C$ 的剩余域相同。选取 $C$ 的一致化元 $\pi_C$，并写成
$\pi_C^p = u \pi$，其中 $u$ 是 $C$ 的某个单位。由于
$\pi_C^p \in L$，可知 $u \in B^*$。此外，$M = L[\pi_C]$。

\medskip\noindent
假设存在整数 $m \geq 0$，使得
$$
u = \sum\nolimits_{0 \leq i < m} b_i^p \pi^i + b \pi^m
$$
，其中 $b_i \in B$，且 $b \in B$ 在 $\kappa_B$ 中的像不是
$p$ 次幂。按引理 \ref{more-algebra-lemma-make-degree-q-extension} 选取扩张
$K_1/K$，其中取 $n = m + 2$；并把 $A$ 在 $K_1$ 中的整闭包
$A_1$ 的一致化元记作 $\pi'$，使得对某个 $a \in A_1$ 有
$\pi = (\pi')^p + (\pi')^{np} a$。设 $B_1$ 是 $B$ 在
$L \otimes_K K_1$ 中的整闭包。由引理
\ref{more-algebra-lemma-weakly-unramified-goes-up-along-totally-ramified}，
$A_1 \to B_1$ 弱非分歧。在 $B_1$ 中有
$$
u \pi =
\left(\sum\nolimits_{0 \leq i < m} b_i (\pi')^{i + 1}\right)^p +
b (\pi')^{(m + 1)p} + (\pi')^{np} b_1
$$
，其中某个 $b_1 \in B_1$（计算从略）。由此可知，$M_1$ 是由在
$L_1$ 上添加下式的一个 $p$ 次根而得到的：
$$
b + (\pi')^{n - m - 1} b_1
$$
。由于 $B_1$ 的剩余域等于 $B$ 的剩余域，由引理
\ref{more-algebra-lemma-pre-purely-inseparable-case} 可知 $M_1/L_1$ 的次数为
$p$，并且 $B_1$ 的整闭包 $C_1$ 在 $B_1$ 上弱非分歧。因此在此情形
得出结论。

\medskip\noindent
% P02-E0651
若不存在前一段中的整数 $m$，则 $u$ 在 $B_1$ 的 $\pi$-进完备化中
是 $p$ 次幂。由于 $B$ 是 Nagata 环，由 Algebra，引理
\ref{algebra-lemma-nagata-pth-roots}，这意味着 $u$ 在 $B_1$ 中是
$p$ 次幂。因此，引理陈述中的第二种情形成立。
\end{proof}

\begin{lemma}
\label{more-algebra-lemma-cohen}
设 $A$ 是由素数 $p$ 零化的局部环，且其极大理想是幂零的。存在环映射
$\sigma : \kappa_A \to A$，它是剩余映射 $A \to \kappa_A$ 的截面。
若 $A \to A'$ 是局部环之间的局部同态，则只要扩张
$\kappa_{A'}/\kappa_A$ 可分，就可以选取与 $\sigma$ 相容的类似环映射
$\sigma' : \kappa_{A'} \to A'$。
\end{lemma}

\begin{proof}
由 Algebra，命题
\ref{algebra-proposition-characterize-separable-field-extensions}，
可分扩张形式光滑。因此，由 $\mathbf{F}_p \to \kappa_A$ 可分可得
$\sigma$ 的存在性。与 $\sigma$ 相容的 $\sigma'$ 的存在性同理。
\end{proof}

\begin{lemma}
\label{more-algebra-lemma-pre-characteristic-p-case}
设 $A$ 是分式域 $K$ 的特征为 $p > 0$ 的离散赋值环。设
$\xi \in K$，并设 $L$ 是在 $K$ 上添加 $z^p - z = \xi$ 的一个根
而得到的扩张。则 $L/K$ 是 Galois 扩张，并且下列情形之一成立：
\begin{enumerate}
\item $L = K$；
\item $L/K$ 关于 $A$ 非分歧，次数为 $p$；
\item $L/K$ 关于 $A$ 完全分歧，分歧指数为 $p$；
\item $A$ 在 $L$ 中的整闭包 $B$ 是离散赋值环，$A \subset B$
弱非分歧，而且 $A \to B$ 诱导次数为 $p$ 的纯不可分剩余域扩张。
\end{enumerate}
设 $\pi$ 是 $A$ 的一致化元。则有以下蕴含：
\begin{enumerate}
\item[(A)] 若 $\xi \in A$，则处于情形 (1) 或 (2)。
\item[(B)] 若 $\xi = \pi^{-n}a$，其中 $n > 0$ 不被 $p$ 整除，且
$a$ 是 $A$ 中的单位，则处于情形 (3)。
\item[(C)] 若 $\xi = \pi^{-n} a$，其中 $n > 0$ 被 $p$ 整除，且
$a$ 在 $\kappa_A$ 中的像不是 $p$ 次幂，则处于情形 (4)。
\end{enumerate}
\end{lemma}

\begin{proof}
由 Fields，章节 \ref{fields-section-Artin-Schreier} 中的讨论，
该扩张是阶数整除 $p$ 的 Galois 扩张。由章节
\ref{more-algebra-section-ramification} 中的讨论立即可知，我们处于引理所列的
(1)--(4) 中的某种情形。

\medskip\noindent
情形 (A)。此时 $A \to A[x]/(x^p - x - \xi)$ 是有限
\'etale 环扩张，故我们处于情形 (1) 或 (2)。

\medskip\noindent
情形 (B)。写成 $\xi = \pi^{-n}a$，其中 $p$ 不整除 $n$。设
$B \subset L$ 是 $A$ 在 $L$ 中的整闭包。若对某个极大理想
$\mathfrak m$ 有 $C = B_\mathfrak m$，则显然
$p \text{ord}_C(z) = -n \text{ord}_C(\pi)$。特别地，
$A \subset C$ 的分歧指数被 $p$ 整除。于是该指数等于 $p$，且
$B = C$。

\medskip\noindent
情形 (C)。令 $k = n/p$。于是可将方程改写为
$$
(\pi^kz)^p - \pi^{n - k} (\pi^kz) = a
$$
。由于 $A[y]/(y^p - \pi^{n - k}y - a)$ 是 $A$ 上弱非分歧的
离散赋值环，引理得证。
\end{proof}

\begin{lemma}
\label{more-algebra-lemma-characteristic-p-case}
设 $A \subset B \subset C$ 是分式域为 $K \subset L \subset M$
的离散赋值环扩张。假设
\begin{enumerate}
\item $A \subset B$ 弱非分歧；
\item $K$ 的特征为 $p$；
\item $M$ 是 $L$ 的次数为 $p$ 的 Galois 扩张；并且
\item $\kappa_A = \bigcap_{n \geq 1} \kappa_B^{p^n}$。
\end{enumerate}
则存在关于 $A$ 完全分歧的有限 Galois 扩张 $K_1/K$，它是
$A \to C$ 的弱解。
\end{lemma}

\begin{proof}
由于 $L$ 的特征为 $p$，我们知道 $M$ 是 $L$ 的 Artin--Schreier
扩张（Fields，引理 \ref{fields-lemma-Artin-Schreier}）。因此可选取
$z \in M$、$z \not \in L$，使得 $\xi = z^p - z \in L$。选取
$n \geq 0$，使得 $\pi^n\xi \in B$。在所有这些选取中，选取使
$n$ 最小的 $z$。若 $n = 0$，则 $M/L$ 关于 $B$ 非分歧
（引理 \ref{more-algebra-lemma-pre-characteristic-p-case}），结论成立。因此
$n > 0$。

\medskip\noindent
假设 (4) 蕴含 $\kappa_A$ 是完美域。因而可按引理
\ref{more-algebra-lemma-cohen} 选取相容的环映射
% P02-E0652
$\overline{\sigma} : \kappa_A \to A/\pi^n A$ 与
$\overline{\sigma} : \kappa_B \to B/\pi^n B$。把其中第二个映射提升为
集合映射 $\sigma : \kappa_B \to B$\footnote{若 $B$ 完备，则可选取
$\sigma$ 为环映射。若 $A$ 也完备且 $\sigma$ 是环映射，则
$\sigma$ 把 $\kappa_A$ 映入 $A$。}。于是可写成
$$
\xi = \sum\nolimits_{i = n, \ldots, 1} \sigma(\lambda_i) \pi^{-i} + b
$$
，其中某些 $\lambda_i \in \kappa_B$ 与 $b \in B$。令
$$
I = \{i \in \{n, \ldots, 1\} \mid \lambda_i \in \kappa_A\}
$$
以及
% P02-E0653
$$
J = \{j \in \{n, \ldots, 1\} \mid \lambda_i \not \in \kappa_A\}
$$
。我们将对有限集 $J$ 的基数作归纳。

\medskip\noindent
情形 $J = \emptyset$。此时，对每个 $i \in \{n, \ldots, 1\}$，
由 $\sigma$ 的选取，可写
$\sigma(\lambda_i) = a_i + \pi^n b_i$，其中某些 $a_i \in A$ 与
$b_i \in B$。因此，对某些 $a \in A$ 与 $b \in B$ 有
$\xi = \pi^{-n} a + b$。若 $p | n$，则可写
$a = a_0^p + \pi a_1$，其中某些 $a_0, a_1 \in A$
（因为 $A$ 的剩余域完美）。计算得
$$
(z - \pi^{-n/p}a_0)^p - (z - \pi^{-n/p}a_0) =
\pi^{-(n - 1)}(a_1 + \pi^{n - 1 - n/p}a_0) + b'
$$
，其中某个 $b' \in B$。这与 $n$ 的最小性矛盾。因此 $p$ 不整除
$n$。考虑由 $w^p - w = \pi^{-n}a$ 给出的 $K$ 的次数为 $p$
的扩张 $K_1$。由引理 \ref{more-algebra-lemma-pre-characteristic-p-case}，该扩张是
Galois 扩张，且关于 $A$ 完全分歧。因此
$L_1 = L \otimes_K K_1$ 是域，而且 $A_1 \subset B_1$ 弱非分歧
（引理 \ref{more-algebra-lemma-weakly-unramified-goes-up-along-totally-ramified}）。
由引理 \ref{more-algebra-lemma-pre-characteristic-p-case}，环
$M_1 = M \otimes_K K_1$ 或者是 $L_1$ 的 $p$ 个副本的乘积
（此时结论成立），或者是 $L_1$ 的次数为 $p$ 的域扩张。此外，在第二种
情形中，或者 $C_1$ 在 $B_1$ 上弱非分歧（此时结论成立），或者
$M_1/L_1$ 的次数为 $p$、是 Galois 扩张，并且关于 $B_1$ 完全分歧。
在最后这种情形中，扩张 $M_1/L_1$ 由元素 $z - w$ 生成，并且
$$
(z - w)^p - (z - w) = z^p - z - (w^p - w) = b
$$
，其中 $b \in B$（见上文）。再用一次引理
\ref{more-algebra-lemma-pre-characteristic-p-case}，扩张 $M_1/L_1$ 关于 $B_1$
非分歧，故可断定 $K_1$ 是 $A \to C$ 的弱解。从现在起假设
$J \not = \emptyset$。

\medskip\noindent
假设存在 $j', j \in J$ 与某个 $r > 0$，使得 $j' = p^r j$。
把所选的 $z$ 改为
$$
z' = z -
(\sigma(\lambda_j) \pi^{-j} + \sigma(\lambda_j^p) \pi^{-pj} + \ldots +
\sigma(\lambda_j^{p^{r - 1}}) \pi^{-p^{r - 1}j})
$$
。于是 $\xi$ 按下式变为 $\xi' = (z')^p - (z')$：
$$
\xi' =
\xi - \sigma(\lambda_j) \pi^{-j} + \sigma(\lambda_j^{p^r}) \pi^{-j'}
+ \text{something in }B
$$
。像先前一样写成
$\xi' = \sum\nolimits_{i = n, \ldots, 1} \sigma(\lambda'_i) \pi^{-i} + b'$，
便得到：当 $i \not = j, j'$ 时有 $\lambda'_i = \lambda_i$，并且
$\lambda'_j = 0$。因此集合 $J$ 变小了。由对 $J$ 的基数所作归纳，
可以假设不存在这样的指标对 $j, j'$。（请注意，在此过程中可能回到
上面已经处理过的 $J = \emptyset$ 情形。）

\medskip\noindent
对 $j \in J$，写
$\lambda_j = \mu_j^{p^{r_j}}$，其中某个 $r_j \geq 0$，而且
$\mu_j \in \kappa_B$ 不是 $p$ 次幂。由假设 (4)，这是可能的。
令 $j \in J$ 是使 $j p^{-r_j}$ 最大的唯一指标（由前一段的结论，
该指标唯一）。选取
% P02-E0654
$r > \max(r_j + 1)$，并使得对 $j \in J$ 有
$j p^{r - r_j} > n$。选取关于 $A$ 完全分歧、次数为 $p^r$ 的可分扩张
$K_1/K$，使得相应离散赋值环 $A_1 \subset K_1$ 有一致化元 $\pi'$，
且对某个 $a \in A_1$ 有
$(\pi')^{p^r} = \pi + \pi^{n + 1}a$
（引理 \ref{more-algebra-lemma-make-degree-q-extension}）。注意，
$L_1 = L \otimes_K K_1$ 是域，而且 $L_1/L$ 关于 $B$ 完全分歧
（引理 \ref{more-algebra-lemma-weakly-unramified-goes-up-along-totally-ramified}）。
在整闭包 $B_1$ 中计算得
$$
\xi = \sum\nolimits_{i \in I} \sigma(\lambda_i) (\pi')^{-i p^r} +
\sum\nolimits_{j \in J} \sigma(\mu_j)^{p^{r_j}} (\pi')^{-j p^r} + b_1
$$
，其中某个 $b_1 \in B_1$。注意，对 $i \in I$，
% P02-E0655
$\sigma(\lambda_i)$ 模 $\pi^n$ 是 $q$ 次幂，换言之模
$(\pi')^{n p^r}$ 也有上述性质。因此可将上式改写为
$$
\xi = \sum\nolimits_{i \in I} x_i^{p^r} (\pi')^{-i p^r} +
\sum\nolimits_{j \in J} \sigma(\mu_j)^{p^{r_j}} (\pi')^{- j p^r}
+ b_1
$$
。如前一段那样，把所选的 $z$ 改为
% P02-E0656
\begin{align*}
z' & = z \\
& -
\sum\nolimits_{i \in I}
\left(x_i (\pi')^{-i} + \ldots + x_i^{p^{r - 1}} (\pi')^{-i p^{r - 1}}\right)
\\
& -
\sum\nolimits_{j \in J}
\left(
\sigma(\mu_j) (\pi')^{- j p^{r - r_j}}
+ \ldots +
\sigma(\mu_j)^{p^{r_j - 1}} (\pi')^{- j p^{r - 1}}
\right)
\end{align*}
，便得到
$$
(z')^p - z' =
\sum\nolimits_{i \in I} x_i (\pi')^{-i} +
\sum\nolimits_{j \in J} \sigma(\mu_j) (\pi')^{- j p^{r - r_j}} + b_1'
$$
，其中某个 $b'_1 \in B_1$。由于恰有一个 $j$ 使得
$j p^{r - r_j}$ 最大，而且 $j p^{r - r_j}$ 大于每个 $i \in I$
并被 $p$ 整除，可知 $M_1 / L_1$ 属于引理
\ref{more-algebra-lemma-pre-characteristic-p-case} 的情形 (C)。证明完成。
\end{proof}

\begin{lemma}
\label{more-algebra-lemma-prepare}
设 $A$ 是包含本原 $p$ 次单位根 $\zeta$ 的环。令 $w = 1 - \zeta$。则
$$
P(z) = \frac{(1 + wz)^p - 1}{w^p} =
z^p - z + \sum\nolimits_{0 < i < p} a_i z^i
$$
是 $A[z]$ 的元素，并且事实上 $a_i \in (w)$。此外，在多项式环
$A[z_1, z_2]$ 中有
$$
P(z_1 + z_2 + w z_1 z_2) = P(z_1) + P(z_2) + w^p P(z_1) P(z_2)
$$
。
\end{lemma}

\begin{proof}
只需在
$$
A = \mathbf{Z}[\zeta] = \mathbf{Z}[x]/(x^{p - 1} + \ldots + x + 1)
$$
是圆分域的整数环时证明。多项式恒等式
$t^p - 1 = (t - 1)(t - \zeta) \ldots (t - \zeta^{p - 1})$
（比较两边的根即可证明）表明
$t^{p - 1} + \ldots + t + 1 = (t - \zeta) \ldots (t - \zeta^{p - 1})$。
代入 $t = 1$，得到
$p = (1 - \zeta)(1 - \zeta^2) \ldots (1 - \zeta^{p - 1})$。
极大理想 $(p, w) = (w)$ 是 $A$ 中位于 $p$ 之上的唯一素理想
（因为特征为 $p$ 的域没有非平凡的 $1$ 的 $p$ 次根）。因此对某个单位
$u$ 有 $p = u w^{p - 1}$。这蕴含
$$
a_i = \frac{1}{p} {p \choose i} u w^{i - 1}
$$
对 $p > i > 1$ 成立，而且 $- 1 + a_1 = pw/w^p = u$。由于
$P(-1) = 0$，可知模 $(w)$ 有 $0 = (-1)^p - u$。因此
$a_1 \in (w)$，第一部分的证明完成。第二部分由从略的直接计算得到。
% P02-E0657
\end{proof}

\begin{lemma}
% P02-E0658
\label{more-algebra-lemma-extension-defined-by-nice-polynial}
设 $A$ 是包含 $1$ 的一个本原 $p$ 次根的混合特征 $(0, p)$ 离散赋值环。
设 $P(t) \in A[t]$ 是引理 \ref{more-algebra-lemma-prepare} 的多项式。设
$\xi \in K$，并设 $L$ 是在 $K$ 上添加 $P(z) = \xi$ 的一个根
而得到的扩张。则 $L/K$ 是 Galois 扩张，并且下列情形之一成立：
\begin{enumerate}
\item $L = K$；
\item $L/K$ 关于 $A$ 非分歧，次数为 $p$；
\item $L/K$ 关于 $A$ 完全分歧，分歧指数为 $p$；
\item $A$ 在 $L$ 中的整闭包 $B$ 是离散赋值环，$A \subset B$
弱非分歧，而且 $A \to B$ 诱导次数为 $p$ 的纯不可分剩余域扩张。
\end{enumerate}
设 $\pi$ 是 $A$ 的一致化元。则有以下蕴含：
\begin{enumerate}
\item[(A)] 若 $\xi \in A$，则处于情形 (1) 或 (2)。
\item[(B)] 若 $\xi = \pi^{-n}a$，其中 $n > 0$ 不被 $p$ 整除，且
$a$ 是 $A$ 中的单位，则处于情形 (3)。
\item[(C)] 若 $\xi = \pi^{-n} a$，其中 $n > 0$ 被 $p$ 整除，且
$a$ 在 $\kappa_A$ 中的像不是 $p$ 次幂，则处于情形 (4)。
\end{enumerate}
\end{lemma}

\begin{proof}
% P02-E0659
添加 $P(z) = \xi$ 的一个根等价于添加
$y^p = w^p(1 + \xi)$ 的一个根。由于 $K$ 包含 $1$ 的一个本原
$p$ 次根，由 Fields，章节 \ref{fields-section-Kummer} 中的讨论，
该扩张是阶数整除 $p$ 的 Galois 扩张。由章节
\ref{more-algebra-section-ramification} 中的讨论立即可知，我们处于引理所列的
(1)--(4) 中的某种情形。

\medskip\noindent
情形 (A)。此时 $A \to A[x]/(P(x) - \xi)$ 是有限
\'etale 环扩张，故我们处于情形 (1) 或 (2)。

\medskip\noindent
情形 (B)。写成 $\xi = \pi^{-n}a$，其中 $p$ 不整除 $n$。设
$B \subset L$ 是 $A$ 在 $L$ 中的整闭包。若对某个极大理想
$\mathfrak m$ 有 $C = B_\mathfrak m$，则显然
$p \text{ord}_C(z) = -n \text{ord}_C(\pi)$。特别地，
$A \subset C$ 的分歧指数被 $p$ 整除。于是该指数等于 $p$，且
$B = C$。

\medskip\noindent
情形 (C)。令 $k = n/p$。于是可将方程改写为
$$
(\pi^kz)^p - \pi^{n - k} (\pi^kz) + \sum a_i \pi^{n - ik} (\pi^kz)^i = a
$$
。由于
% P02-E0660
$A[y]/(y^p - \pi^{n - k}y - \sum  a_i \pi^{n - ik} y^i - a)$
是 $A$ 上弱非分歧的离散赋值环，引理得证。
\end{proof}

\noindent
设 $A$ 是包含 $1$ 的一个本原 $p$ 次根的混合特征 $(0, p)$
离散赋值环。设 $w \in A$ 与 $P(t) \in A[t]$ 如引理
\ref{more-algebra-lemma-prepare} 中所述，并设 $L$ 是 $K$ 的有限扩张。若 $L$
是在 $K$ 上添加方程 $P(z) = \xi$ 的一个根而得到的次数为 $p$
的扩张，其中 $\xi \in K$ 且 $w^p \xi \in \mathfrak m_A$，
则称 $L/K$ 是一个\textit{有限层级的次数为 $p$ 的扩张}。

\medskip\noindent
由于下面这个直接的引理，这一定义与本节的讨论有关。

\begin{lemma}
\label{more-algebra-lemma-make-finite-level}
设 $A \subset B \subset C$ 是分式域为 $K \subset L \subset M$
的离散赋值环扩张。假设
\begin{enumerate}
\item $A$ 具有混合特征 $(0, p)$；
\item $A \subset B$ 弱非分歧；
\item $B$ 包含 $1$ 的一个本原 $p$ 次根；并且
\item $M/L$ 是次数为 $p$ 的 Galois 扩张。
\end{enumerate}
则存在关于 $A$ 完全分歧的有限 Galois 扩张 $K_1/K$，它或者是
$A \to C$ 的弱解，或者使得 $M_1/L_1$ 是有限层级的次数为 $p$
的扩张。
\end{lemma}

\begin{proof}
设 $\pi \in A$ 是一致化元。由 Kummer 理论（Fields，引理
\ref{fields-lemma-Kummer}），$M$ 是从 $L$ 添加某个 $b \in L$
的方程 $y^p = b$ 的根而得到的。

\medskip\noindent
若 $\text{ord}_B(b)$ 与 $p$ 互素，则选取关于 $A$ 完全分歧的
次数为 $p$ 的可分扩张 $K_1/K$（例如使用引理
\ref{more-algebra-lemma-make-degree-q-extension}）。设 $A_1$ 是 $A$ 在 $K_1$
中的整闭包。由引理
\ref{more-algebra-lemma-weakly-unramified-goes-up-along-totally-ramified}，
$B$ 在 $L_1 = L \otimes_K K_1$ 中的整闭包 $B_1$ 是 $A_1$ 上
弱非分歧的离散赋值环。若 $K_1/K$ 不是 $A \to C$ 的弱解，则
$C$ 在 $M_1 = M \otimes_K K_1$ 中的整闭包 $C_1$ 是离散赋值环，
而 $B_1 \to C_1$ 的分歧指数为 $p$。在这种情形中，域 $M_1$
由在 $L_1$ 上添加 $b$ 的 $p$ 次根而得到，而
$\text{ord}_{B_1}(b)$ 被 $p$ 整除。用 $A_1$ 替换 $A$ 等对象，
便可假设 $b = \pi^n u$，其中 $u \in B$ 是单位，且 $n$ 被 $p$
整除。当然，在此情形，扩张 $M$ 是由在 $L$ 上添加某个单位的
$p$ 次根而得到的。

\medskip\noindent
假设 $M$ 是从 $L$ 添加方程 $y^p = u$ 的根而得到的，其中 $u$
是 $B$ 的某个单位。若 $u$ 在 $\kappa_B$ 中的剩余类不是 $p$ 次幂，
则 $B \subset C$ 弱非分歧（引理
\ref{more-algebra-lemma-pre-purely-inseparable-case}），结论成立。否则，可用
$y/v$ 替换所选的 $y$，其中 $v^p$ 与 $u$ 在 $\kappa_B$ 中的像相同。
作此替换后有
$$
y^p = 1 + \pi b
$$
，其中某个 $b \in B$。于是由 $z = (y - 1)/w$ 可知
$P(z) = \pi b/ w^p$。因此可见，该扩张是有限层级的次数为 $p$
的扩张，其中 $\xi = \pi b / w^p$。
\end{proof}

\noindent
设 $A$ 是包含 $1$ 的一个本原 $p$ 次根的混合特征 $(0, p)$
离散赋值环。设 $w \in A$ 与 $P(t) \in A[t]$ 如引理
\ref{more-algebra-lemma-prepare} 中所述。设 $L$ 是 $K$ 的有限层级的次数为 $p$
的扩张。选取在 $K$ 上生成 $L$ 的 $z \in L$，使得
$\xi = P(z) \in K$。选取 $A$ 的一致化元 $\pi$，并写成
$w = u \pi^{e_1}$，其中整数 $e_1 = \text{ord}_A(w)$，而 $u \in A$
是单位。最后，选取 $n \geq 0$，使得
$$
\pi^n \xi \in A
$$
。当 $z$ 取遍所有满足 $\xi = P(z) \in K$ 且生成扩张 $L/K$
的元素时，量 $n/e_1$ 所取的最小值称为 $L/K$ 的\textit{层级}。
% P02-E0661

\medskip\noindent
作几点说明。由于该扩张具有有限层级，我们知道可选取 $z$ 使得
$n < pe_1$。因此层级是包含在 $[0, p)$ 中的有理数。若层级为零，
则由引理 \ref{more-algebra-lemma-extension-defined-by-nice-polynial}，$L/K$ 关于
$A$ 非分歧。我们的下一个目标是降低层级。

\begin{lemma}
\label{more-algebra-lemma-lowering-the-level}
设 $A \subset B \subset C$ 是分式域为 $K \subset L \subset M$
的离散赋值环扩张。假设
\begin{enumerate}
\item $A$ 具有混合特征 $(0, p)$；
\item $A \subset B$ 弱非分歧；
\item $B$ 包含 $1$ 的一个本原 $p$ 次根；
\item $M/L$ 是有限层级 $l > 0$ 的次数为 $p$ 的扩张；
\item $\kappa_A = \bigcap_{n \geq 1} \kappa_B^{p^n}$。
\end{enumerate}
则存在 $K$ 的关于 $A$ 完全分歧的有限可分扩张 $K_1$，使得
$K_1$ 或者是 $A \to C$ 的弱解，或者扩张 $M_1/L_1$ 是层级
$\leq \max(0, l - 1, 2l - p)$ 的次数为 $p$ 的有限层级扩张。
\end{lemma}

\begin{proof}
设 $\pi \in A$ 是一致化元。设 $w \in B$ 与 $P \in B[t]$ 如引理
\ref{more-algebra-lemma-prepare} 中所述（对 $B$ 而言）。令
$e_1 = \text{ord}_B(w)$，于是 $w$ 与 $\pi^{e_1}$ 在 $B$ 中相伴。
选取在 $L$ 上生成 $M$ 的 $z \in M$，使得 $\xi = P(z) \in K$，
并按照 $M$ 在 $L$ 上的层级定义选取 $n$，使得
$\pi^n\xi \in B$；即 $l = n/e_1$。

\medskip\noindent
本引理的证明与引理 \ref{more-algebra-lemma-characteristic-p-case} 的证明完全类似。
为说明其中的情形，注意到，对任意满足
% P02-E0663
$\pi^{-n} P(z) \in B$ 的 $z \in L$（利用 $z$ 的赋值最坏为
$-n/p$ 以及多项式 $P$ 的形状），都有
\begin{equation}
\label{more-algebra-equation-first-congruence}
P(z) \equiv z^p - z \bmod \pi^{-n + e_1}B
\end{equation}
。此外，对 $\xi_1, \xi_2 \in \pi^{-n}B$ 有
\begin{equation}
\label{more-algebra-equation-second-congruence}
\xi_1 + \xi_2 + w^p \xi_1 \xi_2 \equiv \xi_1 + \xi_2 \bmod \pi^{-2n + pe_1}B 
\end{equation}
。最后，注意
% P02-E0664
$n - e_1 = (l - 1)/e_1$ 与
$-2n + pe_1 = -(2l - p)e_1$。写
$m = n - e_1 \max(0, l - 1, 2l - p)$。上述事实表明，在
$\pi^{-n}B / \pi^{-n + m}B$ 中计算时，多项式 $P$ 的表现与多项式
$z^p -  z$ 完全相同。这解释了引理为何成立，不过下面仍给出细节。

\medskip\noindent
% P02-E0662
假设 (4) 蕴含 $\kappa_A$ 是完美域。注意 $m \leq e_1$，所以
$A/\pi^m$ 被 $w$ 零化，从而被 $p$ 零化。因此可按引理
\ref{more-algebra-lemma-cohen} 选取相容的环映射
$\overline{\sigma} : \kappa_A \to A/\pi^mA$ 与
$\overline{\sigma} : \kappa_B \to B/\pi^mB$。把其中第二个映射提升为
集合映射 $\sigma : \kappa_B \to B$。于是可写成
% P02-E0665
$$
\xi =
\sum\nolimits_{i = n, \ldots, n - m + 1} \sigma(\lambda_i) \pi^{-i} +
\pi^{-n + m)} b
$$
，其中某些 $\lambda_i \in \kappa_B$ 与 $b \in B$。令
$$
I = \{i \in \{n, \ldots, n - m + 1\} \mid \lambda_i \in \kappa_A\}
$$
以及
% P02-E0666
$$
J = \{j \in \{n, \ldots, n - m + 1\} \mid \lambda_i \not \in \kappa_A\}
$$
。我们将对有限集 $J$ 的基数作归纳。

\medskip\noindent
情形 $J = \emptyset$。此时，对每个
$i \in \{n, \ldots, n - m + 1\}$，由 $\overline{\sigma}$ 的选取，
% P02-E0667
可写 $\sigma(\lambda_i) = a_i + \pi^{n - m}b_i$，其中某些
$a_i \in A$ 与 $b_i \in B$。因此，对某些 $a \in A$ 与 $b \in B$
有 $\xi = \pi^{-n} a + \pi^{-n + m} b$。若 $p | n$，则可写
$a = a_0^p + \pi a_1$，其中某些 $a_0, a_1 \in A$
（因为 $A$ 的剩余域完美）。令 $z_1 = - \pi^{-n/p} a_0$。注意
$P(z_1) \in \pi^{-n}B$，而且 $z + z_1 + w z z_1$ 是在 $L$ 上
生成 $M$ 的元素（注意，由于 $n < pe_1$，有
$wz_1 \not = -1$）。此外，由引理 \ref{more-algebra-lemma-prepare} 有
$$
P(z + z_1 + w z z_1) = P(z) + P(z_1) + w^p P(z) P(z_1) \in K
$$
，而由方程 (\ref{more-algebra-equation-first-congruence}) 与
(\ref{more-algebra-equation-second-congruence}) 有
$$
P(z) + P(z_1) + w^p P(z) P(z_1)
\equiv
\xi + z_1^p - z_1 \bmod \pi^{-n + m}B
$$
。这里某个 $b' \in B$。
% P02-E0668
这与 $n$ 的最小性矛盾！因此 $p$ 不整除 $n$。考虑由
$P(y) = -\pi^{-n}a$ 给出的 $K$ 的次数为 $p$ 的扩张 $K_1$。
由引理 \ref{more-algebra-lemma-extension-defined-by-nice-polynial}，该扩张可分，
且关于 $A$ 完全分歧。因此 $L_1 = L \otimes_K K_1$ 是域，且
$A_1 \subset B_1$ 弱非分歧
（引理 \ref{more-algebra-lemma-weakly-unramified-goes-up-along-totally-ramified}）。
由引理 \ref{more-algebra-lemma-extension-defined-by-nice-polynial}，环
$M_1 = M \otimes_K K_1$ 或者是 $L_1$ 的 $p$ 个副本的乘积
（此时结论成立），或者是 $L_1$ 的次数为 $p$ 的域扩张。此外，在第二种
情形中，或者 $C_1$ 在 $B_1$ 上弱非分歧（此时结论成立），或者
$M_1/L_1$ 的次数为 $p$、是 Galois 扩张并且关于 $B_1$ 完全分歧。
在最后一种情形，扩张 $M_1/L_1$ 由元素 $z + y + wzy$ 生成，
而且 $P(z + y + wzy) \in L_1$，并且
\begin{align*}
P(z + y + wzy)
& = P(z) + P(y) + w^p P(z) P(y) \\
& \equiv
\xi - \pi^{-n}a \bmod \pi^{-n + m}B_1 \\
& \equiv
0 \bmod \pi^{-n + m}B_1
\end{align*}
，方式与上文完全相同。由 $m$ 的选取，这恰好意味着 $M_1/L_1$
的层级至多为 $\max(0, l - 1, 2l - p)$。从现在起假设
$J \not = \emptyset$。

\medskip\noindent
假设存在 $j', j \in J$ 与某个 $r > 0$，使得 $j' = p^r j$。
令
$$
z_1 = - \sigma(\lambda_j) \pi^{-j} - \sigma(\lambda_j^p) \pi^{-pj} -
\ldots - \sigma(\lambda_j^{p^{r - 1}}) \pi^{-p^{r - 1}j}
$$
，并把 $z$ 改为 $z' = z + z_1 + wzz_1$。注意，$z' \in M$
在 $L$ 上生成 $M$，并且有
% P02-E0669
$\xi' = P(z') = P(z) + P(z_1) + wP(z)P(z_1) \in L$，且
$$
\xi' \equiv
\xi - \sigma(\lambda_j) \pi^{-j} + \sigma(\lambda_j^{p^r}) \pi^{-j'}
\bmod \pi^{-n + m}B
$$
；这是像上面一样使用方程 (\ref{more-algebra-equation-first-congruence}) 与
(\ref{more-algebra-equation-second-congruence}) 得到的。像先前一样写成
$$
\xi' = \sum\nolimits_{i = n, \ldots, n - m + 1} \sigma(\lambda'_i) \pi^{-i}
+ \pi^{-n + m}b'
$$
，便得到：当 $i \not = j, j'$ 时有 $\lambda'_i = \lambda_i$，并且
$\lambda'_j = 0$。因此集合 $J$ 变小了。由对 $J$ 的基数所作归纳，
可以假设 $J$ 中不存在使 $j'/j$ 为 $p$ 的幂的指标对 $j, j'$。
（请注意，在此过程中可能回到上面已经处理过的
$J = \emptyset$ 情形。）

\medskip\noindent
对 $j \in J$，写
$\lambda_j = \mu_j^{p^{r_j}}$，其中某个 $r_j \geq 0$，而且
$\mu_j \in \kappa_B$ 不是 $p$ 次幂。由假设 (4)，这是可能的。
令 $j \in J$ 是使 $j p^{-r_j}$ 最大的唯一指标（由前一段的结论，
该指标唯一）。选取
% P02-E0670
$r > \max(r_j + 1)$，并使得对 $j \in J$ 有
$j p^{r - r_j} > n$。令 $K_1/K$ 是由
$(\pi')^{p^r} = \pi$ 定义的、关于 $A$ 完全分歧的次数为 $p^r$
的扩张。注意，$\pi'$ 是相应离散赋值环 $A_1 \subset K_1$
的一致化元。还要注意，$L_1 = L \otimes_K K_1$ 是域，而且
$L_1/L$ 关于 $B$ 完全分歧
（引理 \ref{more-algebra-lemma-weakly-unramified-goes-up-along-totally-ramified}）。
在整闭包 $B_1$ 中计算得
$$
\xi = \sum\nolimits_{i \in I} \sigma(\lambda_i) (\pi')^{-i p^r} +
\sum\nolimits_{j \in J} \sigma(\mu_j)^{p^{r_j}} (\pi')^{-j p^r} +
\pi^{-n + m} b_1
$$
，其中某个 $b_1 \in B_1$。注意，对 $i \in I$，
% P02-E0671
$\sigma(\lambda_i)$ 模 $\pi^m$ 是 $q$ 次幂，换言之模
$(\pi')^{m p^r}$ 也有上述性质。因此可将上式改写为
$$
\xi = \sum\nolimits_{i \in I} x_i^{p^r} (\pi')^{-i p^r} +
\sum\nolimits_{j \in J} \sigma(\mu_j)^{p^{r_j}} (\pi')^{- j p^r}
+ \pi^{-n + m}b_1
$$
。与前一段中的选取类似，令
% P02-E0672
% P02-E0673
\begin{align*}
z_1 & - \sum\nolimits_{i \in I}
\left(x_i (\pi')^{-i} + \ldots + x_i^{p^{r - 1}} (\pi')^{-i p^{r - 1}}\right)
\\
& - \sum\nolimits_{j \in J}
\left(
\sigma(\mu_j) (\pi')^{- j p^{r - r_j}}
+ \ldots +
\sigma(\mu_j)^{p^{r_j - 1}} (\pi')^{- j p^{r - 1}}
\right)
\end{align*}
，并把所选的 $z$ 改为 $z' = z + z_1 + wzz_1$。于是 $z'$ 在
$L_1$ 上生成 $M_1$，并且
$\xi' = P(z') = P(z) + P(z_1) + w^p P(z) P(z_1) \in L_1$，
而计算表明
$$
\xi' \equiv
\sum\nolimits_{i \in I} x_i (\pi')^{-i} +
\sum\nolimits_{j \in J} \sigma(\mu_j) (\pi')^{- j p^{r - r_j}} +
(\pi')^{(-n + m)p^r}b'_1
$$
，其中某个 $b'_1 \in B_1$。恰有一个 $j$ 使得
$j p^{r - r_j}$ 最大，而且 $j p^{r - r_j}$ 大于每个 $i \in I$。
若 $j p^{r - r_j} \leq (n - m)p^r$，则扩张 $M_1/L_1$ 的层级小于
% P02-E0674
$\max(0, l - 1, 2l - p)$。否则，由于 $p$ 整除
$j p^{r - r_j}$，可知 $M_1 / L_1$ 属于引理
\ref{more-algebra-lemma-extension-defined-by-nice-polynial} 的情形 (C)。证明完成。
\end{proof}

\begin{lemma}
\label{more-algebra-lemma-special-case}
设 $A \subset B \subset C$ 是分式域为 $K \subset L \subset M$
的离散赋值环扩张。假设
\begin{enumerate}
\item $A$ 的剩余域 $k$ 是特征为 $p > 0$ 的代数闭域；
\item $A$ 与 $B$ 完备；
\item $A \to B$ 弱非分歧；
\item $M$ 是 $L$ 的有限扩张；
\item $k = \bigcap\nolimits_{n \geq 1} \kappa_B^{p^n}$。
\end{enumerate}
则存在 $A \to C$ 的弱解，即有限扩张 $K_1/K$。
\end{lemma}

\begin{proof}
设 $M'$ 是 $L$ 的任意有限扩张，并考虑 $B$ 在 $M'$ 中的整闭包
$C'$。由 Algebra，引理
\ref{algebra-lemma-Noetherian-complete-local-Nagata}，$B$ 是 Nagata
环，故 $C'$ 在 $B$ 上有限。此外，$C'$ 是离散赋值环，见评注
\ref{more-algebra-remark-construction} 中的讨论。又因为 $C'$ 作为 $B$-模完备，
所以它作为离散赋值环完备，见 Algebra，章节
\ref{algebra-section-completion}。特别地，这表明 $C$ 是 $B$ 在 $M$
中的整闭包（由赋值环是支配关系下的极大元这一定义）。

\medskip\noindent
设 $M \subset M'$ 是有限扩张，并设 $C' \subset M'$ 是如上所述的
$B$ 的整闭包。由引理 \ref{more-algebra-lemma-solutions-go-down}，只需对
$A \to B \to C'$ 证明结论。因此，由 Fields，引理
\ref{fields-lemma-normal-closure}，可以假设 $M/L$ 正规。

\medskip\noindent
若 $M / L$ 正规，则可找到有限扩张链
$$
L = L^0 \subset L^1 \subset L^2 \subset \ldots \subset L^r = M
$$
，使得每个扩张 $L^{j + 1}/L^j$ 都属于以下类型之一：
\begin{enumerate}
\item[(a)] 次数为 $p$ 的纯不可分扩张；
\item[(b)] 关于 $B^j$ 完全分歧、次数为 $p$ 的 Galois 扩张；
\item[(c)] 关于 $B^j$ 完全分歧、阶数与 $p$ 互素的循环 Galois 扩张；
\item[(d)] 关于 $B^j$ 非分歧的 Galois 扩张。
\end{enumerate}
这里 $B^j$ 是 $B$ 在 $L^j$ 中的整闭包。事实上，由于 $M/L$
正规，可将它写成一个 Galois 扩张与一个纯不可分扩张的合成
（Fields，引理 \ref{fields-lemma-normal-case}）。对纯不可分扩张，
滤过的存在性是清楚的。在 Galois 情形，注意 $G$ 是“该”分解群，
并令 $I \subset G$ 为惯性群。一方面，由引理
\ref{more-algebra-lemma-galois-inertia}，$I$ 可解；另一方面，由引理
\ref{more-algebra-lemma-inertial-invariants-unramified}，扩张 $M^I/L$ 关于 $B$
非分歧。这就证明了所述滤过的存在性。

\medskip\noindent
我们将对整数 $r$ 作归纳。假设可找到有限扩张 $K_1/K$，它是
$A \to B^1$ 的弱解，其中 $B^1$ 是 $B$ 在 $L^1$ 中的整闭包。
设 $K'_1$ 是 $K_1/K$ 的正规闭包
（Fields，引理 \ref{fields-lemma-normal-closure}）。由于 $A$ 完备且
$A$ 的剩余域代数闭，可知 $K'_1/K_1$ 关于 $A_1$ 可分并完全分歧
（略去一些细节）。因此，由引理
\ref{more-algebra-lemma-weakly-unramified-goes-up-along-totally-ramified}，
$K'_1/K$ 也是 $A \to B^1$ 的弱解。换言之，我们可以并且确实假设
$K_1$ 是 $K$ 的正规扩张。在此之后，考虑序列
$$
L^0_1 = (L^0 \otimes_K K_1)_{red} \subset
L^1_1 = (L^1 \otimes_K K_1)_{red} \subset \ldots \subset
L^r_1 = (L^r \otimes_K K_1)_{red}
$$
以及相应的整闭包 $B^i_1$。注意，$C_1 = B^r_1$ 是若干离散赋值环
的乘积；由引理 \ref{more-algebra-lemma-galois-relative}，
$G = \text{Aut}(K_1/K)$ 传递地置换这些离散赋值环。特别地，所有离散
赋值环扩张 $A_1 \to (C_1)_\mathfrak m$ 都彼此同构，而其中一个的
弱解就是所有这些扩张的弱解。可将归纳假设用于序列
$$
A_1 \to (B^1_1)_{B^1_1 \cap \mathfrak m} \to
(B^2_1)_{B^2_1 \cap \mathfrak m} \to
\ldots \to
(B^r_1)_{B^r_1 \cap \mathfrak m} =
(C_1)_\mathfrak m
$$
，得到 $A_1 \to (C_1)_\mathfrak m$ 的弱解 $K_2/K_1$。于是，由前述
讨论，扩张 $K_2/K$ 是 $A \to C$ 的弱解。注意，归纳假设确实适用：
由 $K_1$ 的选取，环映射
$A_1 \to (B^1_1)_{B^1_1 \cap \mathfrak m}$ 弱非分歧；而分式域扩张
的序列中的每一项仍具有上述 (a)、(b)、(c) 或 (d) 中的一种性质。
此外，注意对任意有限扩张 $\kappa_B \subset \kappa$，仍有
$k = \bigcap \kappa^{p^n}$。

\medskip\noindent
因此，一切都归结为：当域扩张 $M/L$ 具有 (a)、(b)、(c) 或 (d)
中的一种性质时，为 $A \subset C$ 找到弱解。

\medskip\noindent
情形 (d)。此情形平凡，因为此时 $B \to C$ 已经非分歧。

\medskip\noindent
情形 (c)。设 $M/L$ 是阶数为 $n$ 且其阶数与 $p$ 互素的循环扩张。
因为 $M/L$ 关于 $B$ 完全分歧，所以 $B \subset C$ 的分歧指数为
$n$，从而 $A \subset C$ 的分歧指数也是 $n$。选取一致化元
$\pi \in A$，并令 $K_1 = K[\pi^{1/n}]$。由 Abhyankar 引理
（引理 \ref{more-algebra-lemma-abhyankar}），$K_1/K$ 是 $A \subset C$ 的解。

\medskip\noindent
情形 (b)。把此情形分成混合特征情形与等特征情形。在等特征情形，
结论是引理 \ref{more-algebra-lemma-characteristic-p-case}。在混合特征情形，
先用引理 \ref{more-algebra-lemma-make-finite-level}，以一个有限扩张替换 $K$，
从而达到 $M/L$ 是有限层级的次数为 $p$ 的扩张这一情形。此时层级是
有理数 $l \in [0, p)$，见引理 \ref{more-algebra-lemma-lowering-the-level} 之前的
讨论。若层级为 $0$，则 $B \to C$ 弱非分歧，结论成立。否则，由引理
\ref{more-algebra-lemma-lowering-the-level}，可以用一个有限扩张替换域 $K$，得到
% P02-E0675
层级为 $l' \leq \max(0, l - 1, 2l - p)$ 的新情形。若
$l = p - \epsilon$ 且 $\epsilon < 1$，则
$l' \leq p - 2\epsilon$。因此，经过有限次替换，就得到层级
$\leq p - 1$ 的情形。再作至多 $p - 1$ 次这样的替换，便达到层级为零
的情形。

\medskip\noindent
情形 (a) 是引理 \ref{more-algebra-lemma-purely-inseparable-case}。这是唯一可能需要
$K$ 的纯不可分扩张的情形；也就是说，在该引理陈述的情形 (2) 中，
通过添加元素 $\pi$ 的一个 $p$ 次幂即可取胜。引理证明完成。
% P02-E0676
\end{proof}

\noindent
至此，我们已经收集了证明本节主要结果所需的所有引理。

\begin{theorem}[Epp]
\label{more-algebra-theorem-epp}
设 $A \subset B$ 是分式域为 $K \subset L$ 的离散赋值环扩张。
若 $\kappa_A$ 的特征为 $p > 0$，则假设
$$
\bigcap\nolimits_{n \geq 1} \kappa_B^{p^n}
$$
中的每个元素在 $\kappa_A$ 上都可分且代数。于是存在有限扩张
$K_1/K$，它是定义 \ref{more-algebra-definition-solution} 所定义的 $A \to B$
的弱解。
\end{theorem}

\begin{proof}
若 $\kappa_A$ 的特征为零，或者剩余特征为 $p$、分歧指数与 $p$
互素且剩余域扩张可分，则结论由 Abhyankar 引理
（引理 \ref{more-algebra-lemma-abhyankar}）得到。事实上，设分歧指数为 $e$。
选取一致化元 $\pi \in A$。令 $K_1/K$ 为添加 $\pi$ 的一个 $e$
次根而得到的扩张。由引理 \ref{more-algebra-lemma-pull-root-uniformizer}，
$A$ 在 $K_1$ 中的整闭包 $A_1$ 是在 $A$ 上具有分歧指数的离散赋值环。
% P02-E0677
因此，由引理 \ref{more-algebra-lemma-abhyankar}，对 $B_1$ 的每个极大理想
$\mathfrak m$，$A_1 \to (B_1)_\mathfrak m$ 在
$\mathfrak m$-进拓扑中形式光滑；从而更有，这些都是离散赋值环的
弱非分歧扩张。

\medskip\noindent
从现在起，设 $p$ 是素数，并假设 $\kappa_A$ 的特征为 $p$。首先使用
引理 \ref{more-algebra-lemma-solution-after-strict-henselization}，约化到 $A$ 与
$B$ 都有可分闭剩余域的情形。由于此过程分别用其可分代数闭包替换
$\kappa_A$ 与 $\kappa_B$，从定理的条件可知，我们得到
$$
\kappa_A \supset \bigcap\nolimits_{n \geq 1} \kappa_B^{p^n}
$$
。

\medskip\noindent
设 $\pi \in A$ 是一致化元，并设 $A^\wedge$ 与 $B^\wedge$ 分别为
$A$ 与 $B$ 的完备化。我们有离散赋值环扩张的交换图
$$
\xymatrix{
B \ar[r] & B^\wedge \\
A \ar[u] \ar[r] & A^\wedge \ar[u]
}
$$
。设 $K^\wedge$ 是 $A^\wedge$ 的分式域。假设可找到有限扩张
$M/K^\wedge$，使它 (a) 是 $A^\wedge \to B^\wedge$ 的弱解，并且
(b) 是一个可分扩张与一个通过添加 $\pi$ 的某个 $p$-幂次根而得到的
扩张的合成。于是由引理 \ref{more-algebra-lemma-approximate-separable-extension}，
可找到有限扩张 $K_1/K$，使得
$K^\wedge \otimes_K K_1 = M$。令 $A_1$（相应地，$A_1^\wedge$）
为 $A$（相应地，$A^\wedge$）在 $K_1$（相应地，$M$）中的整闭包。
由于 $A \to A^\wedge$ 在 $\mathfrak m^\wedge$-进拓扑中形式光滑
（引理 \ref{more-algebra-lemma-extension-dvrs-formally-smooth}），可知
$A_1 \to A_1^\wedge$ 在 $\mathfrak m_1^\wedge$-进拓扑中形式光滑
（引理 \ref{more-algebra-lemma-formally-smooth-goes-up}；由评注
\ref{more-algebra-remark-construction} 中的讨论，$A_1$ 与 $A_1^\wedge$ 是离散赋值环）。
由引理 \ref{more-algebra-lemma-solutions-go-down} 的 (2) 可得，$K_1/K$ 是
$A \to B^\wedge$ 的弱解。再用引理 \ref{more-algebra-lemma-solutions-go-down}
的 (1)，可知 $K_1/K$ 是 $A \to B$ 的弱解。

\medskip\noindent
因此，可以假设 $A$ 与 $B$ 是完备离散赋值环，具有特征为 $p$ 的
可分闭剩余域，并且
$\kappa_A \supset \bigcap\nolimits_{n \geq 1} \kappa_B^{p^n}$。
此外，给定一致化元 $\pi \in A$，我们必须找到 $A \to B$ 的弱解，
它是一个可分扩张与一个通过取 $\pi$ 的 $p$-幂次根而得到的域的合成。
注意，若 $A$ 具有混合特征，则第二个条件自动成立。

\medskip\noindent
令 $k = \bigcap\nolimits_{n \geq 1} \kappa_B^{p^n}$。注意，$k$
是特征为 $p$ 的代数闭域。若 $A$ 具有混合特征，则令 $\Lambda$ 为
$k$ 的 Cohen 环；在等特征情形，令 $\Lambda = k[[t]]$。可以选取环映射
$\Lambda \to A$，它在等特征情形把 $t$ 映到 $\pi$。在等特征情形，
这由 Cohen 结构定理（Algebra，定理
\ref{algebra-theorem-cohen-structure-theorem}）得到；在混合特征情形，
这由 $\mathbf{Z}_p \to \Lambda$ 在进拓扑中形式光滑得到
（引理 \ref{more-algebra-lemma-extension-dvrs-formally-smooth} 与
\ref{more-algebra-lemma-lift-continuous}）。应用引理 \ref{more-algebra-lemma-solutions-go-down}，
可知只需证明 $\Lambda \to B$ 存在弱解，并要求在等特征 $p$ 情形，
该弱解是一个可分扩张与一个通过取 $t$ 的 $p$-幂次根而得到的域的合成。
然而，由于在等特征情形 $\Lambda = k[[t]]$，且 $k((t))$ 的任意扩张
都是这样的合成，现在可以去掉这一要求！

\medskip\noindent
这样便达到如下情形：$A$ 与 $B$ 完备；$A$ 的剩余域 $k$ 是特征为
$p > 0$ 的代数闭域；有 $k = \bigcap \kappa_B^{p^n}$；并且在混合特征
情形，$p$ 是 $A$ 的一致化元（即 $A$ 是 $k$ 的 Cohen 环）。若 $A$
具有混合特征，则选取 $\kappa_B$ 的 Cohen 环 $\Lambda$；在等特征情形，
令 $\Lambda = \kappa_B[[t]]$。像上面一样论证，可以选取环映射
$A \to \Lambda$，它提升 $k \to \kappa_B$ 并把一致化元映到一致化元。
由于 $k \subset \kappa_B$ 可分，环映射 $A \to \Lambda$ 在进拓扑中
形式光滑（引理 \ref{more-algebra-lemma-extension-dvrs-formally-smooth}）。因此可找到
环映射 $\Lambda \to B$，使得复合映射
$A \to \Lambda \to B$ 是给定的环映射 $A \to B$
（见引理 \ref{more-algebra-lemma-lift-continuous}）。由于 $\Lambda$ 与 $B$ 是具有
相同剩余域的完备离散赋值环，由 Algebra，引理
\ref{algebra-lemma-finite-over-complete-ring}，$B$ 在 $\Lambda$ 上有限。
这就约化到引理 \ref{more-algebra-lemma-special-case} 中讨论的特殊情形。
\end{proof}

\section{消除分歧，II}
\label{more-algebra-section-eliminating-ramification-bis}

\noindent
本节使用章节 \ref{more-algebra-section-eliminating-ramification} 的结果，在若干情形
得到（可分的）解。

\begin{lemma}
\label{more-algebra-lemma-solution-goes-up}
设 $A \to B$ 是分式域为 $K \subset L$ 的离散赋值环扩张。若
$K_1/K$ 是 $A \subset B$ 的解，则对任意有限扩张 $K_2/K_1$，
扩张 $K_2/K$ 都是 $A \subset B$ 的解。
\end{lemma}

\begin{proof}
这由引理 \ref{more-algebra-lemma-formally-smooth-goes-up} 得到。细节从略。
\end{proof}

\begin{lemma}
\label{more-algebra-lemma-nagata-goes-down}
设 $A \subset B$ 是离散赋值环扩张。若 $B$ 是 Nagata 环，且分式域
扩张 $L/K$ 可分，则 $A$ 是 Nagata 环。
\end{lemma}

\begin{proof}
离散赋值环是 Nagata 环当且仅当它是 N-2 的。设 $K_1/K$ 是有限纯不
可分域扩张。我们必须证明 $A$ 在 $K_1$ 中的整闭包 $A_1$ 在 $A$
上有限，见 Algebra，引理 \ref{algebra-lemma-domain-char-p-N-1-2}。
由于 $L/K$ 可分而 $K_1/K$ 纯不可分，代数 $L \otimes_K K_1$ 是域
（由 Algebra，引理
\ref{algebra-lemma-separable-extension-preserves-reducedness} 与
\ref{algebra-lemma-radicial-integral-bijective}）。设 $B_1$ 是 $B$ 在
$L \otimes_K K_1$ 中的整闭包。由于 $B$ 是 Nagata 环，$B_1$ 在
$B$ 上有限。由 $B \otimes_A A_1 \subset B_1$ 以及 $B$ 是
Noether 环，可知 $B \otimes_A A_1$ 在 $B$ 上有限。由于
$A \to B$ 忠实平坦，这蕴含 $A_1$ 在 $A$ 上有限，见 Algebra，
引理 \ref{algebra-lemma-descend-properties-modules}。
\end{proof}

\begin{lemma}
\label{more-algebra-lemma-construct-extension}
设 $A' \subset A$ 是环扩张，并设 $f \in A'$。假设 (a) $A$ 在
$A'$ 上有限，(b) $f$ 是 $A$ 上的非零因子，并且 (c)
$A'_f = A_f$。则存在整数 $n_0 > 0$，使得对所有 $n \geq n_0$，
以下结论成立：给定环 $B'$、非零因子 $g \in B'$ 以及满足
$\varphi'(f) \equiv g$ 的同构
$\varphi' : A'/f^n A' \to B'/g^n B'$，存在有限扩张
$B' \subset B$ 以及与 $\varphi'$ 相容的同构
$\varphi : A/fA \to B/gB$。
\end{lemma}

\begin{proof}
% P02-E0678
由于 $A$ 在 $A'$ 上有限且 $A'_f = A_f$，可选取 $t > 0$，使得
$f^t A \subset A'$。令 $n_0 = 2t$。给定引理陈述中的
$n, B', g, \varphi'$，以 $N \subset B'$ 表示由满足
$b \bmod g^nB' \in \varphi'(f^tA)$ 的 $b \in B'$ 组成的集合。
令 $B = g^{-t}N$。由于 $f^tA' \subset f^tA$，且 $\varphi'$ 把
$f$ 映到 $g$，有 $g^tB' \subset N$，故 $B' \subset B$。由于
$f^tA \cdot f^tA \subset f^t \cdot f^tA$，且 $\varphi'$ 把 $f$
映到 $g$，可知 $N \cdot N \subset g^t N$。因此在 $B$ 上得到一个
延拓 $B'$ 之乘法的乘法。我们有 $A'/f^nA'$-模的同构
% P02-E0679
$$
A/f^tA' \xrightarrow{f^t} f^tA/f^nA' \xrightarrow{\varphi'}
g^tB/g^nB' \xrightarrow{g^{-t}} B/g^tB'
$$
，其中右边的模结构用 $\varphi'$ 定义。由于 $A/f^tA'$ 是有限
$A'$-模，可得 $B/g^tB'$ 是有限 $B'$-模，从而 $B' \to B$ 有限。
最后，留给读者验证：所显示的模同构把 $fA$ 映入 $gB$，并诱导与
$\varphi'$ 相容的环同构 $\varphi : A/fA \to B/gB$（它甚至诱导同构
$A/f^tA \to B/g^tB$，但我们不需要这一点）。
\end{proof}

\begin{remark}
\label{more-algebra-remark-functoriality-construct-extension}
引理 \ref{more-algebra-lemma-construct-extension} 中的构造满足以下“函子性”。假设有
交换图
$$
\xymatrix{
A'_2 \ar[r] & A_2 \\
A'_1 \ar[r] \ar[u] & A_1 \ar[u]
}
$$
，其中水平箭头单射。假设给定元素 $f \in A'_1$，使得
$(A'_1 \subset A_1, f)$ 与 $(A'_2 \subset A_2, f)$ 满足引理
\ref{more-algebra-lemma-construct-extension} 的性质 (a)、(b)、(c)。设
$n_{0, 1}$ 与 $n_{0, 2}$ 是引理对这两个情形所给出的整数。最后，
设 $B'_1 \to B'_2$ 是环映射，设
% P02-E0680
$g \in B'_1$ 是 $B_1$ 与 $B_2$ 上的非零因子，设
$n \geq \max(n_{0, 1}, n_{0, 2})$，并给定交换图
% P02-E0681
$$
\xymatrix{
A'_2/f^nA'_2 \ar[r]_{\varphi'_2} & B'_2/g^nB'_2 \\
A'_1/f^nA'_1 \ar[r]^{\varphi'_1} \ar[u] & B'_2/g^nB'_2 \ar[u]
}
$$
，其中水平箭头是同构，且 $\varphi'_1(f) \equiv g$。于是得到交换图
% P02-E0682
$$
\vcenter{
\xymatrix{
B'_2 \ar[r] & B_2 \\
B'_1 \ar[r] \ar[u] & B_1 \ar[u]
}
}
\quad\text{且}\quad
\vcenter{
\xymatrix{
A_2/fA_2 \ar[r]_{\varphi_2} & B_2/gB_2 \\
A_1/fA_1 \ar[r]^{\varphi_1} \ar[u] & B_2/gB_2 \ar[u]
}
}
$$
，其中 $(B'_1 \subset B_1, \varphi_1)$ 与
$(B'_2 \subset B_2, \varphi_2)$ 按引理
\ref{more-algebra-lemma-construct-extension} 的证明构造。详细验证从略。
\end{remark}

\begin{lemma}
\label{more-algebra-lemma-approximate-solution}
设 $p$ 是素数。设 $A \subset B$ 是分式域扩张为 $L/K$ 的离散赋值环
扩张，并设 $K_2/K_1/K$ 是有限域扩张塔。假设
\begin{enumerate}
\item $K$ 的特征为 $p$；
\item $L/K$ 可分；
\item $B$ 是 Nagata 环；
\item $K_2$ 是 $A \subset B$ 的解；
\item $K_2/K_1$ 是次数为 $p$ 的纯不可分扩张。
\end{enumerate}
则存在可分扩张 $K_3/K_1$，它是 $A \subset B$ 的解。
\end{lemma}

\begin{proof}
使用评注 \ref{more-algebra-remark-construction} 中的记号；我们将使用那里所作的
全部观察。由于 $L/K$ 可分，代数 $L_1 = L \otimes_K K_1$ 约化
（Algebra，引理
\ref{algebra-lemma-separable-extension-preserves-reducedness}）。
由于 $B$ 是 Nagata 环，环扩张 $B \subset B_1$ 有限，其中 $B_1$
是 $B$ 在 $L_1$ 中的整闭包，并且 $B_1$ 是 Nagata 环。类似地，由
引理 \ref{more-algebra-lemma-nagata-goes-down}，环 $A$ 是 Nagata 环，因而
$A \subset A_1$ 有限，且 $A_1$ 也是 Nagata 环。此外，对 $K_2$
也有同样断言；即 $L_2 = L \otimes_K K_2$ 约化，环扩张
$A_1 \subset A_2$ 与 $B_1 \subset B_2$ 有限，其中
$A_2$（相应地，$B_2$）是 $A$（相应地，$B$）在 $K_2$
（相应地，$L_2$）中的整闭包。

\medskip\noindent
设 $\pi \in A$ 是一致化元。注意，$\pi$ 是 $K_1$、$K_2$、$A_1$、
$A_2$、$L_1$、$L_2$、$B_1$ 与 $B_2$ 上的非零因子，并且
$K_1 = (A_1)_\pi$、$K_2 = (A_2)_\pi$、$L_1 = (B_1)_\pi$ 且
$L_2 = (B_2)_\pi$。由 Fields，引理
\ref{fields-lemma-finite-purely-inseparable}，可写
$K_2 = K_1(\alpha)$，其中 $\alpha^p = a_1 \in K_1$。把 $\alpha$
乘以 $\pi$ 的某个幂后，可以并且确实假设 $a_1 \in A_1$。在证明的
其余部分，为方便起见写成
$K_2 = K_1[x]/(x^p - a_1)$ 与 $L_2 = L_1[x]/(x^p - a_1)$。
考虑环扩张
$$
A'_2 = A_1[x]/(x^p - a_1) \subset A_2
\quad\text{且}\quad
B'_2 = B_1[x]/(x^p - a_1) \subset B_2
$$
。可将引理 \ref{more-algebra-lemma-construct-extension} 分别用于
$A'_2 \subset A_2$、$f = \pi^2$ 以及
$B'_2 \subset B_2$、$f = \pi^2$。选取对这两者都适用的充分大整数
$n$。

\medskip\noindent
考虑代数
% P02-E0685
$$
K_3 = K_1[x]/(x^p - \pi^{2n} x - a_1)
\quad\text{且}\quad
L_3 = L_1[x]/(x^p - \pi^{2n} x - a_1)
$$
。注意，$K_3/K_1$ 与 $L_3/L_1$ 是次数为 $p$ 的有限
\'etale 代数扩张。考虑子环
% P02-E0683
$$
A'_3 = A_1[x]/(x^p - \pi^n x - a_1)
\quad\text{且}\quad
B'_3 = B_1[x]/(x^p - \pi^n x - a_1)
$$
；它们分别是
% P02-E0684
$K_3 = (A'_2)_\pi$ 与 $L_3 = (B'_3)_\pi$ 的子环。我们将构造交换图
$$
\xymatrix{
B'_2/\pi^{2n} B'_2 \ar[r]_{\psi'} & B'_3/\pi^{2n} B'_3 \\
A'_2/\pi^{2n} A'_2 \ar[r]^{\varphi'} \ar[u] & A'_3/\pi^{2n} A'_3 \ar[u]
}
$$
。事实上，$\varphi'$ 是把 $x$ 的类映到 $x$ 的类的唯一 $A_1$-代数
同构。类似地，$\psi'$ 是把 $x$ 的类映到 $x$ 的类的唯一 $B_1$-代数
同构。由 $n$ 的选取，通过引理 \ref{more-algebra-lemma-construct-extension} 与评注
\ref{more-algebra-remark-functoriality-construct-extension}，得到有限环扩张
$A'_3 \subset A_3$ 与 $B'_3 \subset B_3$，使得
$A'_3 \to B'_3$ 延拓为环映射 $A_3 \to B_3$，并得到交换图
$$
\xymatrix{
B_2/\pi^2 B_2 \ar[r]_\psi & B_3/\pi^2B_3 \\
A_2/\pi^2 A_2 \ar[r]^\varphi \ar[u] & A_3/\pi^2A_3 \ar[u]
}
$$
，它具有所引用结果中断言的所有性质（特别地，$\varphi$ 与 $\psi$
是同构）。

\medskip\noindent
有了所有这些数据，就可以完成证明。首先注意，$A_3$ 与 $B_3$ 是
若干 Dedekind 整环的有限乘积，并且 $\pi$ 属于它们的所有极大理想。
事实上，若 $\mathfrak p \subset A_3$ 是极大理想，则由于
$A \to A_3$ 有限，有 $\pi \in \mathfrak p$。于是
$\mathfrak p/\pi^2 A_3$ 通过 $\varphi$ 对应于
$A_2 / \pi^2A_2$ 中的极大理想；后者是主理想，因为 $A_2$ 是若干
Dedekind 整环的有限乘积。可得 $\mathfrak p/\pi^2 A_3$ 是主理想，
因而由 Nakayama 引理，$\mathfrak p (A_3)_\mathfrak p$ 是主理想。
同样的论证适用于 $B_3$。由此可得，$A_3$ 是 $A$ 在 $K_3$ 中的
整闭包，而 $B_3$ 是 $B$ 在 $L_3$ 中的整闭包。设
$\mathfrak q \subset B_3$ 是位于 $\mathfrak p \subset A_3$ 之上的
极大理想。为完成证明，必须证明
$(A_3)_\mathfrak p \to (B_3)_\mathfrak q$ 在
$\mathfrak q$-进拓扑中形式光滑。由引理
\ref{more-algebra-lemma-extension-dvrs-formally-smooth} 的判据，只需证明
$\mathfrak p (B_3)_\mathfrak q = \mathfrak q (B_3)_\mathfrak q$，并且
域扩张 $\kappa(\mathfrak q)/\kappa(\mathfrak p)$ 可分。这是真的，
因为两个断言都可以通过观察环映射
$A_3/\pi^2 A_3 \to B_3/\pi^2 B_3$ 来检验，而该环映射同构于环映射
$A_2/\pi^2 A_2 \to B_2/\pi^2 B_2$；由于假设 $K_2$ 是
$A \subset B$ 的解，相应断言对后一个映射成立。略去一些细节。
\end{proof}

\begin{lemma}
\label{more-algebra-lemma-separable-solution-separable-solution}
设 $A \subset B$ 是离散赋值环扩张。假设
\begin{enumerate}
\item 分式域扩张 $L/K$ 可分；
\item $B$ 是 Nagata 环；并且
\item $A \subset B$ 存在解。
\end{enumerate}
则 $A \subset B$ 存在可分的解。
\end{lemma}

\begin{proof}
若 $K$ 的特征为零，引理平凡；因此可以并且确实假设 $K$ 的特征为
$p > 0$。

\medskip\noindent
设 $K_2/K$ 是 $A \to B$ 的解。我们将对 $K_2/K$ 的不可分次数
$[K_2 : K]_i$（Fields，定义 \ref{fields-definition-insep-degree}）
作归纳。若 $[K_2 : K]_i = 1$，则 $K_2$ 在 $K$ 上可分，结论成立。
否则，存在子域塔 $K_2/K_1/K$，使得 $K_2/K_1$ 是次数为 $p$ 的
纯不可分扩张（Fields，引理 \ref{fields-lemma-separable-first} 与
\ref{fields-lemma-finite-purely-inseparable}）。由引理
\ref{more-algebra-lemma-approximate-solution}，存在可分扩张 $K_3/K_1$，它是
$A \subset B$ 的解。于是
$[K_3 : K]_i = [K_1 : K]_i = [K_2 : K]_i/p$
（Fields，引理 \ref{fields-lemma-multiplicativity-all-degrees}）更小，
由归纳得出结论。
\end{proof}

\begin{lemma}
\label{more-algebra-lemma-big-extension-is-ok}
设 $A \to B$ 是分式域为 $K \subset L$ 的离散赋值环扩张。假设 $B$
在 $A$ 上本质有限型。设 $K'/K$ 是代数域扩张，使得 $A$ 在 $K'$
中的整闭包 $A'$ 是 Noether 环。则 $B$ 在
$L' = (L \otimes_K K')_{red}$ 中的整闭包 $B'$ 也是 Noether 环。
此外，映射 $\Spec(B') \to \Spec(A')$ 满射，并且相应的剩余域扩张
是有限生成域扩张。
\end{lemma}

\begin{proof}
设 $A \to C$ 是有限型环映射，使得 $B$ 是 $C$ 在素理想
$\mathfrak p$ 处的局部化。则 $C' = C \otimes_A A'$ 是有限型
$A'$-代数，特别地是 Noether 环。由于 $A \to A'$ 整，
$C \to C'$ 也整。因此 $B = C_\mathfrak p \subset C'_\mathfrak p$
也整。由此，$C'_\mathfrak p$ 的维数为 $1$
（Algebra，引理 \ref{algebra-lemma-integral-sub-dim-equal}）。当然，
$C'_\mathfrak p$ 是 Noether 环。设
$\mathfrak q_1, \ldots, \mathfrak q_n$ 是 $C'_\mathfrak p$ 的极小
素理想。设 $B'_i$ 是 $B = C_\mathfrak p$ 的整闭包；或者由上文等价地，
% P02-E0686
是 $C'_\mathfrak p$ 在 $C'_{\mathfrak p'}/\mathfrak q_i$ 的分式域
中的整闭包。由 Krull--Akizuki 定理（将 Algebra，引理
\ref{algebra-lemma-krull-akizuki} 用于 $C'_\mathfrak p$ 在其极大理想
处的有限多个局部化），每个 $B'_i$ 都是 Noether 环。此外，由
Algebra，引理
\ref{algebra-lemma-finite-extension-residue-fields-dimension-1}，
$C'_\mathfrak p \to B'_i$ 中的剩余域扩张有限。最后注意，
$B' = \prod B'_i$ 是 $B$ 在
$L' = (L \otimes_K K')_{red}$ 中的整闭包。
\end{proof}

\begin{proposition}
\label{more-algebra-proposition-epp-essentially-finite-type}
\begin{reference}
特殊情形见 \cite[Lemma 2.13]{alterations}。
\end{reference}
设 $A \to B$ 是分式域为 $K \subset L$ 的离散赋值环扩张。若 $B$
在 $A$ 上本质有限型，则存在有限扩张 $K_1/K$，它是定义
\ref{more-algebra-definition-solution} 所定义的 $A \to B$ 的解。
\end{proposition}

\begin{proof}
注意，若 $A$ 的剩余域完美，则弱解就是解，见引理
\ref{more-algebra-lemma-extension-dvrs-formally-smooth}。因此，若 $A$ 的剩余特征为
$0$，该命题立即由定理 \ref{more-algebra-theorem-epp} 得到（事实上，无需假设
$A \to B$ 本质有限型）。若 $A$ 的剩余特征为 $p > 0$，也将从 Epp
定理推出该命题。

\medskip\noindent
设 $x_i \in A$（$i \in I$）是一族元素，它们映到 $A$ 的剩余域
$\kappa$ 的一个 $p$-基。令
% P02-E0688
$$
A' = \bigcup\nolimits_{n \geq 1} A[t_{i, n}]/(t_{i, n}^{p^n} - x_i)
$$
，其中转移映射把 $t_{i, n + 1}$ 映到 $t_{i, n}^p$。注意，$A'$
是 $A$ 上离散赋值环的弱非分歧有限扩张的滤过余极限。因此 $A'$ 是
离散赋值环，且 $A \to A'$ 弱非分歧。由构造，剩余域
$\kappa' = A'/\mathfrak m_A A'$ 是 $\kappa$ 的完美化。

\medskip\noindent
设 $K'$ 是 $A'$ 的分式域。可将引理 \ref{more-algebra-lemma-big-extension-is-ok}
用于扩张 $K'/K$。因此 $B'$ 是若干 Dedekind 整环的有限乘积。设
$\mathfrak m_1, \ldots, \mathfrak m_n$ 是 $B'$ 的极大理想。使用 Epp
定理（定理 \ref{more-algebra-theorem-epp}），对每个扩张
$A' \subset B'_{\mathfrak m_i}$ 找到弱解 $K'_i/K'$。由于 $A'$
的剩余域完美，这些实际上都是解。设 $K'_1/K'$ 是包含每个 $K'_i$
的有限扩张。由引理 \ref{more-algebra-lemma-solution-goes-up}，$K'_1/K'$ 仍是每个
$A' \subset B'_{\mathfrak m_i}$ 的解。

\medskip\noindent
设 $A'_1$ 是 $A$ 在 $K'_1$ 中的整闭包。注意，把评注
\ref{more-algebra-remark-construction} 中的讨论用于 $K' \subset K'_1$，可知
$A'_1$ 是 Dedekind 整环。因此，引理 \ref{more-algebra-lemma-big-extension-is-ok}
适用于 $K'_1/K$。从而 $B$ 在
$L'_1 = (L \otimes_K K'_1)_{red}$ 中的整闭包 $B'_1$ 是 Dedekind
整环；并且由于 $K'_1/K'$ 是每个 $A' \subset B'_{\mathfrak m_i}$
的解，可知对每个极大理想 $\mathfrak m \subset B'_1$，
% P02-E0687
$(A'_1)_{A'_1 \cap \mathfrak m} \to (B'_1)_{\mathfrak m}$ 在
$\mathfrak m$-进拓扑中形式光滑。

\medskip\noindent
由构造，域 $K'_1$ 是 $K$ 的有限扩张的滤过余极限。写
$K'_1 = \colim_{i \in I} K_i$。对每个 $i$，令 $A_i$（相应地，
$B_i$）为 $A$（相应地，$B$）在 $K_i$（相应地，
$L_i = (L \otimes_K K_i)_{red}$）中的整闭包。于是显然
$$
A'_1 =  \colim A_i\quad\text{且}\quad B'_1 = \colim B_i
$$
。由于环映射 $A_i \to A'_1$ 与 $B_i \to B'_1$ 是单射整环映射，
而 $A'_1$ 与 $B'_1$ 的谱有限，可知对所有充分大的 $i$，环映射
$A_i \to A'_1$ 与 $B_i \to B'_1$ 在谱上双射。一旦如此，对所有充分
大的 $i$，映射 $A_i \to A'_1$ 与 $B_i \to B'_1$ 将弱非分歧
（只要一致化元已在像中）。由分歧指数的乘法性可知，对这样的 $i$，
$A_i \to B_i$ 在 $B_i$ 的所有极大理想局部化上诱导弱非分歧映射。
再稍微增大 $i$，可知
$$
B_i \otimes_{A_i} A'_1 \longrightarrow B'_1
$$
在剩余域上诱导满射（因为由引理 \ref{more-algebra-lemma-big-extension-is-ok}，
$B'_1$ 的剩余域在 $A'_1$ 的相应剩余域上有限生成）。在 $B'_1$
的一个选定极大理想之下的极大理想处，剩余域的图为
$$
\xymatrix{
\kappa_{B_i} \ar[r] &
\kappa_{B_{i'}} \ar[r] &
 & \ldots &
\kappa_{B'_1} \\
\kappa_{A_i} \ar[r] \ar[u] &
\kappa_{A_{i'}} \ar[r] \ar[u] &
 & \ldots &
\kappa_{A'_1} \ar[u]
}
$$
。因此，$\kappa_{B_i}$ 是 $\kappa_{A_i}$ 的有限生成扩张，使得
$\kappa_{B_i}$ 与 $\kappa_{A'_1}$ 在 $\kappa_{B'_1}$ 中的合成在
$\kappa_{A'_1}$ 上可分。此事已经在某个有限阶段发生：例如，设
$\kappa_{B'_1}$ 在 $\kappa_{A'_1}(x_1, \ldots, x_n)$ 上有限可分；
只需增大 $i$，使得 $x_1, \ldots, x_n$ 属于 $\kappa_{B_i}$，并使
所有生成元都满足系数在 $\kappa_{A_i}(x_1, \ldots, x_n)$ 中的可分
多项式方程。这意味着，对 $B_i$ 的所有极大理想 $\mathfrak m$，
$A_i \to (B_i)_\mathfrak m$ 在 $\mathfrak m$-进拓扑中形式光滑，
证明完成。
\end{proof}

\begin{lemma}
\label{more-algebra-lemma-epp-essentially-finite-type-separable}
设 $A \to B$ 是分式域为 $K \subset L$ 的离散赋值环扩张。假设
\begin{enumerate}
\item $B$ 在 $A$ 上本质有限型；
\item $A$ 或 $B$ 是 Nagata 环；并且
\item $L/K$ 可分。
\end{enumerate}
则 $A \to B$ 存在可分的解（定义 \ref{more-algebra-definition-solution}）。
\end{lemma}

\begin{proof}
注意，若 $A$ 是 Nagata 环，则 $B$ 也是 Nagata 环
（Algebra，引理 \ref{algebra-lemma-nagata-localize} 与命题
\ref{algebra-proposition-nagata-universally-japanese}）。因此，结合命题
\ref{more-algebra-proposition-epp-essentially-finite-type} 与引理
\ref{more-algebra-lemma-separable-solution-separable-solution} 即得本引理。
\end{proof}

\section{环的 Picard 群}
\label{more-algebra-section-picard}

\noindent
首先如下定义可逆模。

\begin{definition}
\label{more-algebra-definition-invertible}
设 $R$ 是环。若函子
$$
\text{Mod}_R \longrightarrow \text{Mod}_R,\quad
N \longmapsto M \otimes_R N
$$
是范畴等价，则称 $R$-模 $M$ \textit{可逆}。若可逆 $R$-模作为
$R$-模同构于 $R$，则称它\textit{平凡}。
\end{definition}

\begin{lemma}
\label{more-algebra-lemma-invertible}
设 $R$ 是环，并设 $M$ 是 $R$-模。下列条件等价：
\begin{enumerate}
\item $M$ 是秩为 $1$ 的有限局部自由模；
\item $M$ 可逆；
\item 存在 $R$-模 $N$，使得 $M \otimes_R N \cong R$。
\end{enumerate}
此外，在这种情形，(3) 中的模 $N$ 同构于 $\Hom_R(M, R)$。
\end{lemma}

\begin{proof}
假设 (1)。考虑模 $N = \Hom_R(M, R)$ 与求值映射
$M \otimes_R N = M \otimes_R \Hom_R(M, R) \to R$。若
$f \in R$ 使得 $M_f \cong R_f$，则求值映射在 $f$ 处局部化后成为
同构（细节从略）。因此，由 Algebra，引理 \ref{algebra-lemma-cover}，
求值映射是同构。故 (1) $\Rightarrow$ (3)。

\medskip\noindent
假设 (3)。则函子 $K \mapsto K \otimes_R N$ 是函子
$K \mapsto K \otimes_R M$ 的拟逆。因此 (3) $\Rightarrow$ (2)。
反之，若 (2) 成立，则 $K \mapsto K \otimes_R M$ 本质满，从而 (3)
成立。

\medskip\noindent
假设等价条件 (2) 与 (3) 成立。以
$\psi : M \otimes_R N \to R$ 表示 (3) 中的同构。选取元素
$\xi = \sum_{i = 1, \ldots, n} x_i \otimes y_i$，使得
$\psi(\xi) = 1$。考虑同构
$$
M \to M \otimes_R M \otimes_R N \to M
$$
，其中第一个箭头把 $x$ 映到 $\sum x_i \otimes x \otimes y_i$，
第二个箭头把 $x \otimes x' \otimes y$ 映到
$\psi(x' \otimes y)x$。由此可得，
$x \mapsto \sum \psi(x \otimes y_i)x_i$ 是 $M$ 的自同构。该自同构
分解为
$$
M \to R^{\oplus n} \to M
$$
，其中第一个箭头由
$x \mapsto (\psi(x \otimes y_1), \ldots, \psi(x \otimes y_n))$
给出，第二个箭头由
$(a_1, \ldots, a_n) \mapsto \sum a_i x_i$ 给出。这样便可断定 $M$
是有限自由 $R$-模的直和项。这意味着 $M$ 有限局部自由
（Algebra，引理 \ref{algebra-lemma-finite-projective}）。由对称性，
$N$ 也如此；又因为 $M \otimes_R N \cong R$，可知 $M$ 与 $N$
的秩都必须为 $1$。
\end{proof}

\noindent
这些模的同构类集合通常称为 $R$ 的\textit{类群}或\textit{Picard 群}。
群结构由如下规则确定：把可逆模 $L$ 与 $L'$ 的同构类映到
$L \otimes_R L'$ 的同构类。可逆模 $L$ 的逆元是模
$$
L^{\otimes -1} = \Hom_R(L, R),
$$
因为如引理 \ref{more-algebra-lemma-invertible} 的证明所示，求值映射
$L \otimes_R L^{\otimes -1} \to R$ 是同构。以 $\Pic(R)$ 表示 $R$
的 Picard 群。

\begin{lemma}
\label{more-algebra-lemma-UFD-Pic-trivial}
设 $R$ 是 UFD。则 $\Pic(R)$ 平凡。
\end{lemma}

\begin{proof}
设 $L$ 是可逆 $R$-模。由引理 \ref{more-algebra-lemma-invertible}，$L$ 是有限局部
自由 $R$-模。特别地，$L$ 无挠并且在 $R$ 上有限。从对偶可逆模中
选取非零元素 $\varphi \in \Hom_R(L, R)$。则
$I = \varphi(L) \subset R$ 是作为模可逆的理想。选取非零
$f \in I$，并设
$$
f = u p_1^{e_1} \ldots p_r^{e_r}
$$
是它分解成两两不同的素元 $p_i$ 的分解。由于 $L$ 有限局部自由，
存在 $a_i \in R$、$a_i \not \in (p_i)$，使得对某个
$g_i \in R_{a_i}$ 有 $I_{a_i} = (g_i)$。于是 $p_i$ 在 UFD
$R_{a_i}$ 中仍是素元，并且可对某个不被 $p_i$ 整除的
$g'_i \in R_{a_i}$ 写成 $g_i = p_i^{c_i} g'_i$。由于
$f \in I_{a_i}$，可知 $e_i \geq c_i$。我们断言 $I$ 由
$h = p_1^{c_1} \ldots p_r^{c_r}$ 生成；这将完成证明。

\medskip\noindent
为证明该断言，只需对任意使 $I_a$ 为主理想的 $a \in R$ 证明
$I_a$ 由 $h$ 生成（Algebra，引理 \ref{algebra-lemma-cover}）。设
$I_a = (g)$。令 $J \subset \{1, \ldots, r\}$ 为使 $p_i$ 在 $R_a$
中是非单位（从而是素元）的 $i$ 的集合。由于 $f \in I_a = (g)$，
得到素因子分解
% P02-E0689
$g = v \prod_{i \in J} p_j^{b_j}$，其中 $v$ 是单位，且
$b_j \leq e_j$。对每个 $j \in J$，有
$I_{aa_j} = g R_{aa_j} = g_j R_{aa_j}$；换言之，$g$ 与 $g_j$
在 $R_{aa_j}$ 中的像相伴。由分解的唯一性，这蕴含 $b_j = c_j$，
证明完成。
\end{proof}

\section{行列式}
\label{more-algebra-section-determinants}

\noindent
设 $R$ 是环，并设 $M$ 是有限投射 $R$-模。存在乘积分解
$R = R_0 \times \ldots \times R_t$，使得在 $M$ 的相应分解
$M = M_0 \times \ldots \times M_t$ 中，$M_i$ 是 $R_i$ 上秩为 $i$
的有限局部自由模。这由 Algebra，引理
\ref{algebra-lemma-finite-projective}（用以看出秩局部常值），以及
Algebra，引理 \ref{algebra-lemma-disjoint-decomposition} 与
\ref{algebra-lemma-disjoint-implies-product}（用以把 $R$ 分解为乘积）
得到。在这种情形，定义 $R$-模
$$
\det(M) = \wedge^0_{R_0}(M_0) \times \ldots \times \wedge^t_{R_t}(M_t)
$$
。这是秩为 $1$ 的有限局部自由模，因为每一项都是秩为 $1$ 的有限局部
自由模。若 $\varphi : M \to N$ 是有限投射 $R$-模的同构，则得到秩为
$1$ 的局部自由模的典范同构
$$
\det(\varphi) : \det(M) \longrightarrow \det(N)
$$
。更一般地，若对 $R$ 的所有素理想 $\mathfrak p$，自由模
$M_\mathfrak p$ 与 $N_\mathfrak p$ 的秩相同，则任意 $R$-模同态
$\varphi : M \to N$ 都诱导 $R$-模映射
$\det(\varphi) : \det(M) \to \det(N)$。最后，若 $M = N$，则
$\det(\varphi) : \det(M) \to \det(M)$ 是可逆 $R$-模的自同态。
由于对可逆 $R$-模有 $R = \Hom_R(L, L)$，可以并且确实把
$\det(\varphi)$ 看作 $R$ 的元素。这样便得到\textit{行列式}
$$
\det : \Hom_R(M, M) \longrightarrow R
$$
；这是乘法映射。

\begin{remark}
\label{more-algebra-remark-determinant-as-socle}
设 $R$ 是环，并设 $M$ 是有限投射 $R$-模。可以考虑 $R$ 上由 $M$
生成的分次交换 $R$-代数，即外代数 $\wedge^*_R(M)$。关于
$\det(M)$ 的一个公式是：$\det(M) \subset \wedge^*_R(M)$ 是
$M \subset \wedge^*_R(M)$ 的零化子。这个公式有时很有用，因为它不
涉及把 $R$ 分解成乘积。当然，要证明它具有所需性质，或者必须把 $R$
分解成乘积（如上），或者必须观察 $R$ 的各素理想处的局部化。
\end{remark}

\noindent
接下来考虑给定有限投射模的短正合列时，行列式会怎样。
% P02-E0690

\begin{lemma}
\label{more-algebra-lemma-det-ses}
设 $R$ 是环。设
$$
0 \to M' \to M \to M'' \to 0
$$
是有限投射 $R$-模的短正合列。则存在典范同构
$$
\gamma : \det(M') \otimes \det(M'') \longrightarrow \det(M)
$$
\end{lemma}

\begin{proof}[第一种证明]
第一种证明。把 $R$ 分解成环 $R_{ij}$ 的乘积，使得
$M' = \prod M'_{ij}$ 与 $M'' = \prod M''_{ij}$，其中 $M'_{ij}$
的秩为 $i$，而 $M''_{ij}$ 的秩为 $j$。当然，此时
$M = \prod M_{ij}$，并且 $M_{ij}$ 的秩为 $i + j$。这把我们约化到
$M'$ 与 $M''$ 具有常秩的情形，设其秩分别为 $i$ 与 $j$。在这种情形，
必须构造典范映射
$$
\wedge^i(M') \otimes \wedge^j(M'') \longrightarrow \wedge^{i + j}(M)
$$
。为此，选取 $M'$ 中的 $m'_1, \ldots, m'_i$ 与 $M''$ 中的
$m''_1, \ldots, m''_j$。以 $m_1, \ldots, m_i \in M$ 表示
$m'_1, \ldots, m'_i$ 的像，并以
$m_{i + 1}, \ldots , m_{i + j} \in M$ 表示映到 $M''$ 中
$m''_1, \ldots, m''_j$ 的元素。规则为
$$
m'_1 \wedge \ldots \wedge m'_i \otimes
m''_1 \wedge \ldots \wedge m''_j
\longmapsto
m_1 \wedge \ldots \wedge m_{i + j}
$$
。关于它良定义且是同构的详细证明从略。
\end{proof}

\begin{proof}[第二种证明]
我们将使用评注 \ref{more-algebra-remark-determinant-as-socle} 中给出的
$\det(M)$、$\det(M')$ 与 $\det(M'')$ 的描述。考虑 $R$-代数映射
$\wedge^*_R(M') \to \wedge^*_R(M)$ 与
$\wedge^*_R(M) \to \wedge^*_R(M'')$。第一个单射，第二个满射。取元素
$x' \in \det(M') \subset \wedge^*_R(M')$ 与元素
$x'' \in \det(M'') \subset \wedge^*_R(M'')$。选取映到 $x''$
的元素 $y'' \in \wedge^*(M)$，并令
$$
\gamma(x' \otimes x'') = x' \wedge y'' \in \det(M) \subset \wedge^*_R(M)
$$
。通过观察各素理想处的局部化，读者容易验证它良定义且是同构。
% P02-E0691
此外，这一构造给出与第一种证明中构造相同的映射。
\end{proof}

\begin{lemma}
\label{more-algebra-lemma-det-ses-functorial}
设 $R$ 是环。设
$$
\xymatrix{
0 \ar[r] &
M' \ar[r] \ar[d]^u &
M \ar[r] \ar[d]^v &
M'' \ar[r] \ar[d]^w &
0 \\
0 \ar[r] &
K' \ar[r] &
K \ar[r] &
K'' \ar[r] &
0
}
$$
是有限投射 $R$-模的交换图，其中竖直箭头都是同构。则得到同构的交换图
$$
\xymatrix{
\det(M') \otimes \det(M'') \ar[r]_-\gamma \ar[d]_{\det(u) \otimes \det(w)} &
\det(M) \ar[d]^{\det(v)} \\
\det(K') \otimes \det(K'') \ar[r]^-\gamma & \det(K)
}
$$
，其中水平箭头是引理 \ref{more-algebra-lemma-det-ses} 中构造的箭头。
\end{lemma}

\begin{proof}
从略。提示：使用引理 \ref{more-algebra-lemma-det-ses} 中映射 $\gamma$ 的第二种构造。
\end{proof}

\begin{lemma}
\label{more-algebra-lemma-det-filtration}
设 $R$ 是环。设
$$
K \subset L \subset M
$$
是 $R$-模，使得 $K$、$L/K$ 与 $M/L$ 是有限投射 $R$-模。则图
$$
\xymatrix{
\det(K) \otimes \det(L/K) \otimes \det(M/L) \ar[r] \ar[d] &
\det(L) \otimes \det(M/L) \ar[d] \\
\det(K) \otimes \det(M/K) \ar[r] &
\det(M)
}
$$
交换，其中映射是引理 \ref{more-algebra-lemma-det-ses} 中的映射。
\end{lemma}

\begin{proof}
从略。提示：在 $R$ 的某个素理想处局部化后，可以假设
$K \subset L \subset M$ 同构于
$R^{\oplus a} \subset R^{\oplus a + b}  \subset R^{\oplus a + b + c}$；
在这种情形，结论是显然的计算。
\end{proof}

\begin{lemma}
\label{more-algebra-lemma-det-direct-sum}
设 $R$ 是环，并设 $M'$ 与 $M''$ 是两个有限投射 $R$-模。则图
$$
\xymatrix{
\det(M') \otimes \det(M'') \ar[r]
\ar[d]_{\epsilon \cdot (\text{switch tensors})} &
\det(M' \oplus M'') \ar[d]^{\det(\text{switch summands})} \\
\det(M'') \otimes \det(M') \ar[r] &
\det(M'' \oplus M')
}
$$
交换，其中 $\epsilon = \det( -\text{id}_{M' \otimes M''}) \in R^*$，
水平箭头是引理 \ref{more-algebra-lemma-det-ses} 中的箭头。
\end{lemma}

\begin{proof}
从略。
\end{proof}

\begin{lemma}
\label{more-algebra-lemma-det-switch}
设 $R$ 是环，并设 $M$、$N$ 是有限投射 $R$-模。设
$a : M \to N$ 与 $b : N \to M$ 是 $R$-线性映射。则作为 $R$ 的元素，
有
$$
\det(\text{id} + a \circ b) = \det(\text{id} + b \circ a)
$$
。
\end{lemma}

\begin{proof}
只需在以素理想局部化的环替换 $R$ 后证明该断言。因此可假设 $R$
局部，而 $M$ 与 $N$ 有限自由。在这种情形，必须证明通常的矩阵行列式
之间的等式
$$
\det(I_n + AB) = \det(I_m + BA)
$$
，其中 $A$ 的大小为 $n \times m$，而 $B$ 的大小为 $m \times n$。
这约化到环
$R = \mathbf{Z}[a_{ij}, b_{ji}; 1 \leq i \leq n, 1 \leq j \leq m]$
的情形，其中
% P02-E0693
$a_{ij}$ 与 $b_{ij}$ 是变量以及矩阵 $A$ 与 $B$ 的元素。取分式域后，
这约化到特征为零的域的情形。在特征为零时，存在一个普遍多项式，
它以所述矩阵各次幂的迹表示大小为 $\leq N$ 的矩阵的行列式。因此只需对
所有 $k \geq 1$ 证明
% P02-E0692
$$
\text{Trace}((I_n + AB)^k) = \text{Trace}((I_m + BA)^k)
$$
。展开后可知，只需对所有 $k \geq 0$ 证明
$\text{Trace}((AB)^k) = \text{Trace}((BA)^k)$。当 $k = 1$ 时，
这是熟知的事实 $\text{Trace}(AB) = \text{Trace}(BA)$。当 $k > 1$
时，把 $(AB)^k = A(BA)^{k - 1}B$ 与
$(BA)^k = (BA)^{k - 1} A B$ 写出，便由该事实得到结论。
\end{proof}

\noindent
回顾我们在 Algebra，章节 \ref{algebra-section-K-groups} 中定义了群
$K_0(R)$：它是以有限投射 $R$-模的同构类为生成元的自由群，模去关系
$[M'] + [M''] = [M' \oplus M'']$。

\begin{lemma}
\label{more-algebra-lemma-det}
设 $R$ 是环。存在映射
$$
\det : K_0(R) \longrightarrow \Pic(R)
$$
；若 $M$ 是秩为 $n$ 的有限局部自由模，则它把 $[M]$ 映到可逆模
$\wedge^n(M)$ 的类。
\end{lemma}

\begin{proof}
这立即由上述构造得到；特别地，用引理 \ref{more-algebra-lemma-det-ses} 可看出这些
关系被映到 $0$。
\end{proof}

\section{完美复形与 K-群}
\label{more-algebra-section-perfect-K-group}

\noindent
我们迅速证明：环 $R$ 的完美复形导出范畴的第零 K-群，与 Algebra，
章节 \ref{algebra-section-K-groups} 中定义的 $K_0(R)$ 相同。

\begin{lemma}
\label{more-algebra-lemma-perfect-to-K-group}
设 $R$ 是环。存在映射
$$
c : \text{perfect complexes over }R \longrightarrow K_0(R)
$$
，它具有以下性质：
\begin{enumerate}
\item 对完美复形 $K$ 有 $c(K[n]) = (-1)^nc(K)$；
\item 若 $K \to L \to M \to K[1]$ 是完美复形的可区别三角，则
$c(L) = c(K) + c(M)$；
\item 若 $K$ 由有限投射模组成的有限复形 $M^\bullet$ 表示，则
% P02-E0694
$c(K) = \sum (-1)^i[M_i]$。
\end{enumerate}
\end{lemma}

\begin{proof}
设 $K$ 是 $D(R)$ 的完美对象。由定义，可用有限投射 $R$-模的有限复形
$M^\bullet$ 表示 $K$。通过令
$$
c(K) = \sum (-1)^n[M^n]
$$
在 $K_0(R)$ 中定义 $c$。当然，必须证明它良定义；但一旦良定义，(1)
与 (3) 就立即成立。暂时把映射 $c$ 看作定义在有限投射 $R$-模的复形上。

\medskip\noindent
假设 $L^\bullet \to M^\bullet$ 是有限投射 $R$-模的有限复形之间的
满射。设 $K^\bullet$ 为其核。于是得到 $R$-模的短正合列
$$
0 \to K^n \to L^n \to M^n \to 0
$$
；这些列分裂，因为 $M^n$ 投射。因此 $K^\bullet$ 也是有限投射
$R$-模的有限复形，并且在 $K_0(R)$ 中有
$c(L^\bullet) = c(K^\bullet) + c(M^\bullet)$。

\medskip\noindent
假设给定有限投射 $R$-模的非循环有限复形 $M^\bullet$。设当
$n \not \in [a, b]$ 时 $M^n = 0$。则可把 $M^\bullet$ 分解成短正合列
% P02-E0695
$$
\begin{matrix}
0 \to M^a \to M^{a + 1} \to N^{a + 1} \to 0, \\
0 \to N^{a + 1} \to M^{a + 2} \to N^{a + 3} \to 0, \\
\ldots \\
0 \to N^{b - 3} \to M^{b - 2} \to N^{b - 2} \to 0, \\
0 \to N^{b - 2} \to M^{b - 1} \to M^b \to 0
\end{matrix}
$$
。作降序归纳，可知 $N^{b - 2}, \ldots, N^{a + 1}$ 是有限投射
$R$-模，这些列分裂正合，并且
% P02-E0696
$$
c(M^\bullet) = \sum (-1)[M^n] = \sum (-1)^n([N^{n - 1}] + [N^n]) = 0
$$
。因此，我们的构造在非循环复形上给出零。

\medskip\noindent
由前两段的结果形式地可知，$c$ 良定义并满足 (2)。事实上，假设有限投射
$R$-模的有限复形 $M^\bullet$ 与 $L^\bullet$ 表示 $D(R)$ 的同一对象。
由 Derived Categories，引理
\ref{derived-lemma-morphisms-from-projective-complex}，可以用复形映射
$f : M^\bullet \to L^\bullet$ 表示该同构。由 Derived Categories，
定义 \ref{derived-definition-cone}，得到复形的短正合列
% P02-E0697
$$
0 \to L^\bullet \to C(f)^\bullet \to K^\bullet[1] \to 0
$$
。由于 $f$ 是拟同构，锥 $C(f)^\bullet$ 非循环（例如，这由 Derived
Categories，章节 \ref{derived-section-canonical-delta-functor} 中的讨论
得到）。因此
$$
0 = c(C(f)^\bullet) = c(L^\bullet) + c(K^\bullet[1]) =
c(L^\bullet) - c(K^\bullet)
$$
，如所需。类似的 (2) 之证明从略。
\end{proof}

\noindent
下面的引理表明 $K_0(R)$ 等于 $K_0(D_{perf}(R))$。

\begin{lemma}
\label{more-algebra-lemma-perfect-to-K-group-universal}
设 $R$ 是环。设 $D_{perf}(R)$ 是完美对象的导出范畴，见引理
\ref{more-algebra-lemma-perfect-ring-classical-generator}。引理
\ref{more-algebra-lemma-perfect-to-K-group} 的映射 $c$ 给出同构
$K_0(D_{perf}(R)) = K_0(R)$。
\end{lemma}

\begin{proof}
由 $K_0(D_{perf}(R))$ 的定义（Derived Categories，定义
\ref{derived-definition-K-zero}），$c$ 诱导同态
$K_0(D_{perf}(R)) \to K_0(R)$。

\medskip\noindent
给定 $R$ 上的有限投射模 $M$，以 $M[0]$ 表示如下 $R$ 上的完美复形：
$M$ 位于次数 $0$，其他次数均为零。给定有限投射模的短正合列
$0 \to M \to M' \to M'' \to 0$，得到可区别三角
$M[0] \to M'[0] \to M''[0] \to M[1]$，见 Derived Categories，
章节 \ref{derived-section-canonical-delta-functor}。这表明，把 $[M]$
映到 $[M[0]]$ 得到映射 $K_0(R) \to K_0(D_{perf}(R))$；请原谅这糟糕的
记号。

\medskip\noindent
显然，$K_0(R) \to K_0(D_{perf}(R)) \to K_0(R)$ 是恒等映射。另一方面，
若 $M^\bullet$ 是有限投射 $R$-模的有界复形，则“粗截断”的可区别三角
的存在性（见 Homology，章节 \ref{homology-section-truncations}）
% P02-E0698
$$
\sigma_{\geq n}M^\bullet \to \sigma_{\geq n - 1}M^\bullet \to
M^{n - 1}[-n + 1] \to (\sigma_{\geq n}M^\bullet)[1]
$$
与归纳表明
$$
[M^\bullet] = \sum (-1)^i[M^i[0]]
$$
在 $K_0(D_{perf}(R))$ 中成立（再次为该记号致歉）。因此映射
$K_0(R) \to K_0(D_{perf}(R))$ 满射，证明完成。
\end{proof}



\section{有限长度模的自同态的行列式}
\label{more-algebra-section-determinants-finite-length}

\noindent
设 $(R, \mathfrak m, \kappa)$ 是局部环。考虑如下二元组 $(M, \varphi)$
所成的范畴：前一分量是有限长度 $R$-模，而
$\varphi : M \to M$ 是一个自同态。该范畴是 Abel 范畴，并且每个对象
既是 Artin 的也是 Noether 的。相关定义见 Homology，章节
\ref{homology-section-jordan-holder}。

\medskip\noindent
若 $(M, \varphi)$ 是该范畴的单对象，则 $M$ 被 $\mathfrak m$ 零化；
否则
% P02-E0699
$(\mathfrak m M, \varphi|_{\mathfrak m M})$
会是一个非平凡子对象。此外，$\dim_\kappa(M) = \text{length}_R(M)$
是有限的。因此可以用线性代数把行列式和迹
$$
\det\nolimits_\kappa(\varphi),\quad \text{Trace}_\kappa(\varphi)
$$
定义为 $\kappa$ 的元素。在此情形下，$\varphi$ 的特征多项式也可类似定义。

\medskip\noindent
由 Homology，引理 \ref{homology-lemma-finite-length}，对上述范畴的任意对象
$(M, \varphi)$，存在由 $\varphi$-稳定子模组成的有限滤过
$$
0 \subset M_1 \subset \ldots \subset M_n = M
$$
，使得 $(M_i/M_{i - 1}, \varphi_i)$ 是该范畴的单对象；这里
$\varphi_i : M_i/M_{i - 1} \to M_i/M_{i - 1}$
是诱导映射。定义 $(M, \varphi)$ 在 $\kappa$ 上的{\it 行列式}为
$$
\det\nolimits_\kappa(\varphi) = \prod \det\nolimits_\kappa(\varphi_i)
$$
，其中 $\det_\kappa(\varphi_i)$ 如上一段所定义。定义
$(M, \varphi)$ 在 $\kappa$ 上的{\it 迹}为
$$
\text{Trace}_\kappa(\varphi) = \sum \text{Trace}_\kappa(\varphi_i)
$$
，其中 $\text{Trace}_\kappa(\varphi_i)$ 如上一段所定义。类似地，
可将 $\varphi$ 在 $\kappa$ 上的特征多项式定义为上一段所定义的
$\varphi_i$ 的特征多项式之积。由 Jordan--H\"older 定理
（Homology，引理 \ref{homology-lemma-jordan-holder}），这些定义良好。

\begin{lemma}
\label{more-algebra-lemma-ses}
设 $(R, \mathfrak m, \kappa)$ 是局部环。设
$0 \to (M, \varphi) \to (M', \varphi') \to (M'', \varphi'') \to 0$
是上述范畴中的短正合列。则
$$
\det\nolimits_\kappa(\varphi') =
\det\nolimits_\kappa(\varphi)\det\nolimits_\kappa(\varphi''),\quad
\text{Trace}_\kappa(\varphi') = \text{Trace}_\kappa(\varphi) + 
\text{Trace}_\kappa(\varphi'')
$$
。此外，$\varphi'$ 在 $\kappa$ 上的特征多项式等于
$\varphi$ 与 $\varphi''$ 的特征多项式之积。
\end{lemma}

\begin{proof}
留作练习。
\end{proof}

\begin{lemma}
\label{more-algebra-lemma-multiplication}
设 $(R, \mathfrak m, \kappa) \to (R', \mathfrak m', \kappa')$
是局部环的局部同态，并假设 $\kappa'/\kappa$ 是有限扩张。设
$u \in R'$。则对任意有限长度 $R'$-模 $M'$，有
$$
\det\nolimits_\kappa(u : M' \to M') =
\text{Norm}_{\kappa'/\kappa}(u \bmod \mathfrak m')^m
$$
，其中 $m = \text{length}_{R'}(M')$。
\end{lemma}

\begin{proof}
注意到
$\text{length}_R(M') = \text{length}_{R'}(M') [\kappa' : \kappa]$，
所以命题有意义。若 $M' = \kappa'$，则等式由范数的定义成立，因为范数
就是由乘以 $u$ 给出的线性算子的行列式。一般情形下，选取适当滤过并应用
引理 \ref{more-algebra-lemma-ses} 的乘法性，即可归约到这一情形。略去若干细节。
\end{proof}

\begin{lemma}
\label{more-algebra-lemma-flat-base-change-det}
设 $(R, \mathfrak m, \kappa) \to (R', \mathfrak m', \kappa')$
是局部环的平坦局部同态，并且
$m = \text{length}_{R'}(R'/\mathfrak mR') < \infty$。
对任意上述 $(M, \varphi)$，元素 $\det_\kappa(\varphi)^m$
在 $\kappa'$ 中的像等于
$\det_{\kappa'}(\varphi \otimes 1 : M \otimes_R R' \to M \otimes_R R')$。
\end{lemma}

\begin{proof}
$R \to R'$ 的平坦性保证引理 \ref{more-algebra-lemma-ses} 中的短正合列经基变换后
仍为 $R'$ 上的短正合列。因此，由引理 \ref{more-algebra-lemma-ses} 的乘法性，
可以假设 $(M, \varphi)$ 是本范畴的单对象（见本节引言）。
在单对象情形下，$M$ 被 $\mathfrak m$ 零化。选取滤过
$$
0 \subset I_1 \subset I_2 \subset \ldots \subset I_{m - 1} \subset
R'/\mathfrak mR'
$$
，使其各相继商作为 $R'$-模都同构于 $\kappa'$。于是得到滤过
$$
0 \subset
M \otimes_\kappa I_1 \subset
M \otimes_\kappa I_2 \subset
\ldots \subset
M \otimes_\kappa I_{m - 1} \subset
M \otimes_\kappa R'/\mathfrak mR' = M \otimes_R R'
$$
，其各相继商都同构于 $M \otimes_\kappa \kappa'$。并且这些子模在
$\varphi \otimes 1$ 下不变。由引理 \ref{more-algebra-lemma-ses} 得
$$
\det\nolimits_{\kappa'}(\varphi \otimes 1 : M \otimes_R R' \to M \otimes_R R')
=
\det\nolimits_{\kappa'}(\varphi \otimes 1 :
M \otimes_\kappa \kappa' \to M \otimes_\kappa \kappa')^m =
\det\nolimits_\kappa(\varphi)^m
$$
。最后一个等式来自线性映射的行列式与域扩张的相容性。这就证明了引理。
\end{proof}







\section{正则局部环是唯一分解整环}
\label{more-algebra-section-regular-local-UFD}

\noindent
我们证明节标题中所述的结果。

\begin{lemma}
\label{more-algebra-lemma-regular-local-Pic-zero}
设 $R$ 是正则局部环，且 $f \in R$。则 $\Pic(R_f) = 0$。
\end{lemma}

\begin{proof}
设 $L$ 是可逆 $R_f$-模。特别地，$L$ 是有限 $R_f$-模。存在有限
$R$-模 $M$ 使得 $M_f \cong L$，见 Algebra，引理
\ref{algebra-lemma-construct-fp-module}。由 Algebra，命题
\ref{algebra-proposition-regular-finite-gl-dim}，$M$ 在 $R$ 上有有限自由
消解 $F_\bullet$。由此 $L$ 拟同构于自由 $R_f$-模的有限复形。因此，
由引理 \ref{more-algebra-lemma-perfect-to-K-group}，对某个 $n \in \mathbf{Z}$，
% P02-E0700
$K_0(R_f)$ 中有 $[L_f] = n[R_f]$。应用引理 \ref{more-algebra-lemma-det} 的映射，
可知 $L$ 平凡。
\end{proof}

\begin{lemma}
\label{more-algebra-lemma-regular-local-UFD}
正则局部环是唯一分解整环。
\end{lemma}

\begin{proof}
回忆正则局部环是整环，见 Algebra，引理
\ref{algebra-lemma-regular-domain}。我们对正则局部环 $R$ 的维数归纳，
证明唯一分解性质。若 $\dim(R) = 0$，则 $R$ 是域，因而尤其是唯一分解
整环。假设 $\dim(R) > 0$。取
$x \in \mathfrak m$，$x \not \in \mathfrak m^2$。由 Algebra，引理
\ref{algebra-lemma-regular-ring-CM}，$R/(x)$ 正则；再由 Algebra，
引理 \ref{algebra-lemma-regular-domain}，它是整环，故 $x$ 是素元。
设 $\mathfrak p \subset R$ 是高为 $1$ 的素理想。我们必须证明
$\mathfrak p$ 是主理想，见 Algebra，引理
\ref{algebra-lemma-characterize-UFD-height-1}。可以假设
$x \not \in \mathfrak p$，因为若 $x \in \mathfrak p$，则
$\mathfrak p = (x)$，结论已得。
对每个非极大素理想 $\mathfrak q \subset R$，局部环
$R_\mathfrak q$ 都是正则局部环，见 Algebra，引理
\ref{algebra-lemma-localization-of-regular-local-is-regular}。
由归纳假设，$\mathfrak pR_\mathfrak q$ 是主理想。特别地，
$R_x$-模 $\mathfrak p_x = \mathfrak pR_x \subset R_x$ 是有限呈示 $R_x$-模，
并且在任意素理想处的局部化都是秩 $1$ 自由模。由 Algebra，引理
\ref{algebra-lemma-finite-projective}，$\mathfrak p_x$ 是可逆
$R_x$-模。由引理 \ref{more-algebra-lemma-regular-local-Pic-zero}，对某个
$y \in R_x$ 有 $\mathfrak p_x = (y)$。对某个
$f \in \mathfrak p$ 与 $e \in \mathbf{Z}$，可以写成
$y = x^e f$。把 $f = a_1 \ldots a_r$ 分解为 $R$ 的不可约元之积
（Algebra，引理 \ref{algebra-lemma-factorization-exists}）。
由于 $\mathfrak p$ 是素理想，对某个 $i$ 有
$a_i \in \mathfrak p$。由于 $\mathfrak p_x = (y)$ 是素理想，并且在
$R_x$ 中 $a_i | y$，故 $\mathfrak p_x$ 在 $R_x$ 中由 $a_i$ 生成，
即 $a_i$ 在 $R_x$ 中的像是素元。由于 $x$ 是素元，由 Algebra，引理
\ref{algebra-lemma-invert-prime-elements} 可知 $a_i$ 在 $R$ 中是素元。
又因 $(a_i) \subset \mathfrak p$ 且 $\mathfrak p$ 的高为 $1$，
所以 $(a_i) = \mathfrak p$，如所需。
\end{proof}

\begin{lemma}
\label{more-algebra-lemma-picard-group-generic-fibre-regular}
设 $R$ 是分式域为 $K$、剩余域为 $\kappa$ 的赋值环。设
$R \to A$ 是环同态，使得
\begin{enumerate}
\item $A$ 是局部环，且 $R \to A$ 是局部同态；
\item $A$ 在 $R$ 上平坦且本质有限型；
% P02-E0701
\item $A \otimes_R \kappa$ 正则。
\end{enumerate}
则 $\Pic(A \otimes_R K) = 0$。
\end{lemma}

\begin{proof}
设 $L$ 是可逆 $A \otimes_R K$-模。特别地，$L$ 是有限模。
存在有限 $A$-模 $M$ 使得 $M \otimes_R K \cong L$，见 Algebra，引理
\ref{algebra-lemma-construct-fp-module}。可以假设 $M$ 作为 $R$-模无挠。
因此 $M$ 作为 $R$-模平坦（引理
\ref{more-algebra-lemma-valuation-ring-torsion-free-flat}）。由引理
\ref{more-algebra-lemma-flat-finite-type-valuation-ring-finite-presentation}，
$M$ 作为 $A$-模是有限呈示的，并且 $A$ 作为 $R$-代数本质有限呈示。
由引理 \ref{more-algebra-lemma-structure-relatively-perfect}，$M$ 相对于 $R$ 是完美的；
特别地，$M$ 作为 $A$-模是伪凝聚的。由引理
\ref{more-algebra-lemma-perfect-over-regular-local-ring}，$M$ 是完美的，故
$M$ 在 $A$ 上有有限自由消解 $F_\bullet$。由此 $L$ 拟同构于自由
$A \otimes_R K$-模的有限复形。因此，由引理
\ref{more-algebra-lemma-perfect-to-K-group}，对某个 $n \in \mathbf{Z}$，
$K_0(A \otimes_R K)$ 中有
$[L] = n[A \otimes_R K]$。应用引理 \ref{more-algebra-lemma-det} 的映射，可知
$L$ 平凡。
\end{proof}
\section{复形的行列式}
\label{more-algebra-section-determinants-complexes}

\noindent
在章节 \ref{more-algebra-section-perfect-K-group} 中，我们已经看到，对环 $R$ 上的完美
复形 $K$，可以关联一个可逆 $R$-模的同构类，即 $\Pic(R)$ 的一个元素。
事实上，与章节 \ref{more-algebra-section-determinants} 类似，存在函子
$$
\det :
\left\{
\begin{matrix}
\text{category of perfect complexes} \\
\text{morphisms are isomorphisms}
\end{matrix}
\right\}
\longrightarrow
\left\{
\begin{matrix}
\text{category of invertible modules} \\
\text{morphisms are isomorphisms}
\end{matrix}
\right\}
$$
。此外，若 $(L, F)$ 是 $R$ 的滤过导出范畴 $DF(R)$ 中的对象，其滤过
有限且各分次部分都是完美复形，则存在典范同构
$\det(\text{gr}L) \to \det(L)$。原始论述见 \cite{determinant}。
% P02-E0702
我们将在以后补入这些材料（插入未来的引用）。

\medskip\noindent
眼下，我们对 $D(R)$ 中 tor-幅度位于 $[-1, 0]$ 的完美对象 $L$
给出一种特设构造。这样的对象可由复形
$$
L^\bullet = \ldots \to 0 \to L^{-1} \to L^0 \to 0 \to \ldots
$$
表示，其中 $L^{-1}$ 与 $L^0$ 是有限投射 $R$-模，见引理
\ref{more-algebra-lemma-perfect}。在此情形下，置
$$
\det(L^\bullet) = \det(L^0) \otimes_R \det(L^{-1})^{\otimes -1} =
\Hom_R(\det(L^{-1}), \det(L^0))
$$
。若这种形式的复形对 $R$ 的所有素理想都有
$L^{-1}_\mathfrak p$ 与 $L^0_\mathfrak p$ 秩相同，则称其{\it 秩为 $0$}。
若 $L^\bullet$ 的秩为 $0$，则章节 \ref{more-algebra-section-determinants} 已表明，
存在典范元素
$$
\delta(L^\bullet) \in \det(L^\bullet)
$$
；它就是 $d : L^{-1} \to L^0$ 的
% P02-E0703
行列式。注意，$\delta(L^\bullet)$ 是 $\det(L^\bullet)$ 的平凡化，
当且仅当 $L^\bullet$ 非循环。

\medskip\noindent
考虑复形映射 $a^\bullet : K^\bullet \to L^\bullet$，并假设
\begin{enumerate}
\item $a^\bullet$ 是拟同构；
\item 对所有 $n$，$a^n : K^n \to L^n$ 都满射；
\item $K^n$、$L^n$ 都是有限投射 $R$-模，并且仅当
$n \in \{-1, 0\}$ 时可能非零。
\end{enumerate}
在此情形下，我们将构造同构
$$
\det(a^\bullet) : \det(K^\bullet) \longrightarrow \det(L^\bullet)
$$
。利用正合列 $0 \to \Ker(a^i) \to K^i \to L^i \to 0$，
由引理 \ref{more-algebra-lemma-det-ses}，对 $i = -1, 0$ 得到同构
$$
\gamma^i : \det(\Ker(a^i)) \otimes \det(L^i) \to \det(K^i)
$$
。由于 $a^\bullet$ 是拟同构，复形 $\Ker(a^\bullet)$ 非循环且秩为 $0$。
因此，上述典范元素 $\delta(\Ker(a^\bullet))$ 是可逆 $R$-模
$\det(\Ker(a^\bullet))$ 的一个平凡化。把
$\det(a^\bullet) : \det(K^\bullet) \to \det(L^\bullet)$ 定义为使图表
$$
\xymatrix{
\det(K^\bullet)
\ar[rr]_{\det(a^\bullet)}
& &
\det(L^\bullet)
\ar[ld]^{\delta(\Ker(a^\bullet))}
\\
& \det(L^\bullet) \otimes \det(\Ker(a^\bullet))
\ar[lu]^{\gamma^0 \otimes (\gamma^{-1})^{\otimes -1}}
}
$$
交换的唯一同构。

\begin{lemma}
\label{more-algebra-lemma-canonical-element-well-defined}
设 $R$ 是环。设 $a^\bullet : K^\bullet \to L^\bullet$ 是满足上述
(1)、(2)、(3) 的 $R$-模复形映射。若 $L^\bullet$ 秩为 $0$，则
$\det(a^\bullet)$ 把典范元素 $\delta(K^\bullet)$ 映到
$\delta(L^\bullet)$。
\end{lemma}

\begin{proof}
记 $M^i = \Ker(a^i)$。于是有短正合列的映射
$$
\xymatrix{
0 \ar[r] &
M^{-1} \ar[r] \ar[d]_{d_M} &
K^{-1} \ar[r] \ar[d]_{d_K} &
L^{-1} \ar[r] \ar[d]_{d_L} &
0 \\
0 \ar[r] &
M^0 \ar[r] &
K^0 \ar[r] &
L^0 \ar[r] &
0
}
$$
。由引理 \ref{more-algebra-lemma-det-ses-functorial}，作为映射的
$\det(d_K)$ 对应于 $\det(d_M) \otimes \det(d_L)$。展开各定义即得所需
等式。
\end{proof}

\begin{lemma}
\label{more-algebra-lemma-homotopic-surjections}
设 $R$ 是环。设 $a^\bullet : K^\bullet \to L^\bullet$ 是满足上述
(1)、(2)、(3) 的 $R$-模复形映射。设 $h : K^0 \to L^{-1}$ 是映射，
使得 $b^0 = a^0 + d \circ h$ 与
$b^{-1} = a^{-1} + h \circ d$ 都满射。则作为
$\det(K^\bullet) \to \det(L^\bullet)$ 的映射，有
$\det(a^\bullet) = \det(b^\bullet)$。
\end{lemma}

\begin{proof}
假设存在映射 $\tilde h : K^0 \to K^{-1}$，使得
$h = a^{-1} \circ \tilde h$，并且
$k^0 = \text{id} + d \circ \tilde h : K^0 \to K^0$ 与
% P02-E0704
$k^1 = \text{id} + \tilde h \circ d : K^{-1} \to K^{-1}$ 都是同构。
于是得到复形的交换图
$$
\xymatrix{
0 \ar[r] &
\Ker(b^\bullet) \ar[r] \ar[d]_{c^\bullet} &
K^\bullet \ar[r]_{b^\bullet}
\ar[d]_{k^\bullet} &
L^\bullet \ar[r] \ar[d]^{\text{id}} &
0 \\
0 \ar[r] &
\Ker(a^\bullet) \ar[r] &
K^\bullet \ar[r]^{a^\bullet} &
L^\bullet \ar[r] &
0
}
$$
，其中 $c^\bullet$ 是诱导的核之间的同构。利用引理
\ref{more-algebra-lemma-det-ses-functorial} 可知
$$
\xymatrix{
\det(\Ker(b^i)) \otimes \det(L^i) \ar[r] \ar[d]_{\det(c^i) \otimes 1} &
\det(K^i) \ar[d]^{\det(k^i)} \\
\det(\Ker(a^i)) \otimes \det(L^i) \ar[r] &
\det(K^i)
}
$$
交换。由于
% P02-E0705
$\det(c^\bullet)$ 把 $\det(\Ker(a^\bullet))$ 的典范平凡化映到
$\Ker(b^\bullet)$ 的典范平凡化
（引理 \ref{more-algebra-lemma-canonical-element-well-defined}），可知当且仅当
% P02-E0706
$$
\det(k^0) = \det(k^{-1})
$$
在 $R$ 中成立时便可得到结论；而该等式由引理
\ref{more-algebra-lemma-det-switch} 得到。

\medskip\noindent
假设存在直和因子 $U \subset K^{-1}$，使得
$a^{-1}|_U : U \to L^{-1}$ 与 $b^{-1}|_U : U \to L^{-1}$
都是同构。把 $\tilde h$ 定义为 $h$ 与 $a^{-1}|_U$ 的逆之复合。
我们断言 $\tilde h$ 是证明第一段中所需的映射。事实上，由构造有
$h = a^{-1} \circ \tilde h$。
为证明 $k^{-1} : K^{-1} \to K^{-1}$ 是同构，只需证明它满射
（Algebra，引理 \ref{algebra-lemma-fun}）。设 $u \in U$。可以选取
$u' \in U$ 使得 $b^{-1}(u') = a^{-1}(u)$。于是
$u = k^{-1}(u')$。事实上，$u$ 与 $k^{-1}(u')$ 都在 $U$ 中，并且由计算
\footnote{$a^{-1}(k^{-1}(u')) =
a^{-1}(u') + a^{-1}(\tilde h(d(u'))) =
a^{-1}(u') + h(d(u')) = b^{-1}(u') = a^{-1}(u)$}
有 $a^{-1}(u) = a^{-1}(k^{-1}(u'))$。由于 $a^{-1}|_U$ 是同构，
得到该等式。因此 $U \subset \Im(k^{-1})$。另一方面，若
$x \in \Ker(a^{-1})$，则 $x = k^{-1}(x) \bmod U$。由于
$K^{-1} = \Ker(a^{-1}) + U$，可知 $k^{-1}$ 满射。
最后证明 $k^0 : K^0 \to K^0$ 满射。首先，由
$a^0 \circ k^0 = b^0$ 可知 $a^0 \circ k^0$ 满射。若
$x \in \Ker(a^0)$，则对某个 $y \in \Ker(a^{-1})$ 有 $x = d(y)$。
由上可对某个 $z \in K^{-1}$ 写成 $y = k^{-1}(z)$。于是
$x = k^0(d(z))$，结论得证。

\medskip\noindent
证明的最后一步。只需找到上一段中的 $U$，但这未必总能直接做到。不过，
为证明两个 $R$-模映射相等，只需在 $R$ 的各素理想处局部化后证明。
因此可以假设 $R$ 是局部环。于是问题变为：假设
$$
\alpha, \beta : R^{\oplus n} \longrightarrow R^{\oplus m}
$$
是两个满 $R$-线性映射，要找一个直和因子
$U \subset R^{\oplus n}$，使得 $\alpha|_U$ 与 $\beta|_U$
都是同构。若 $R$ 是域，这可由线性代数做到。一般地，先在剩余域上取一个
解，再把它提升为局部环 $R$ 上的解。略去若干细节。
\end{proof}
\begin{lemma}
\label{more-algebra-lemma-compose-surjections}
设 $R$ 是环。设 $a^\bullet : K^\bullet \to L^\bullet$ 与
$b^\bullet : L^\bullet \to M^\bullet$ 是满足上述 (1)、(2)、(3) 的
$R$-模复形映射。则
$\det(b^\bullet) \circ \det(a^\bullet) = 
\det(b^\bullet \circ a^\bullet)$；这里两边作为
% P02-E0707
$\det(M^\bullet) \to \det(K^\bullet)$ 的映射相等。
\end{lemma}

\begin{proof}
从略。提示：直接应用引理
\ref{more-algebra-lemma-det-ses}、\ref{more-algebra-lemma-det-ses-functorial}
与 \ref{more-algebra-lemma-det-filtration}。
\end{proof}

\begin{lemma}
\label{more-algebra-lemma-determinant-two-term-complexes}
设 $R$ 是环。上述构造确定函子
$$
\det :
\left\{
\begin{matrix}
\text{category of perfect complexes} \\
\text{with tor amplitude in }[-1, 0] \\
\text{morphisms are isomorphisms}
\end{matrix}
\right\}
\longrightarrow
\left\{
\begin{matrix}
\text{category of invertible modules} \\
\text{morphisms are isomorphisms}
\end{matrix}
\right\}
$$
。此外，给定 $D(R)$ 中秩为 $0$ 且 tor-幅度位于 $[-1, 0]$ 的完美对象
$L$，存在典范元素 $\delta(L) \in \det(L)$，使得对 $D(R)$ 中任意同构
$a : L \to K$ 都有 $\det(a)(\delta(L)) = \delta(K)$。
\end{lemma}

\begin{proof}
由引理 \ref{more-algebra-lemma-perfect}，源范畴的每个对象都可由复形
$$
L^\bullet = \ldots \to 0 \to L^{-1} \to L^0 \to 0 \to \ldots
$$
表示，其中 $L^{-1}$ 与 $L^0$ 是有限投射 $R$-模。暂称这种复形为良好
复形。由 Derived Categories，引理
\ref{derived-lemma-morphisms-from-projective-complex}，导出范畴中良好
复形之间的态射就是复形映射的链同伦类。因此，可以使用良好复形来工作，
并可使用上文研究的行列式
$\det(L^\bullet) = \det(L^0) \otimes \det(L^{-1})^{\otimes -1}$。

\medskip\noindent
设 $a^\bullet : L^\bullet \to K^\bullet$ 是良好复形之间的态射，并且
它在 $D(R)$ 中是同构，即为拟同构。若图表
$$
\xymatrix{
L^\bullet \ar[rr]_{a^\bullet} & & K^\bullet \\
& M^\bullet \ar[lu]^{b^\bullet} \ar[ru]_{c^\bullet}
}
$$
在链同伦意义下交换，且 $b^\bullet$ 与 $c^\bullet$ 满足上述
(1)、(2)、(3)，则称其为良好图表。一旦有这样的图表，就可以有意义地定义
$$
\det(a^\bullet) = \det(c^\bullet) \circ \det(b^\bullet)^{-1}
$$
，其中 $\det(c^\bullet)$ 与 $\det(b^\bullet)$ 是上文构造的同构。
我们将证明良好图表总是存在，并且所得映射 $\det(a^\bullet)$ 与良好图表
的选择无关。

\medskip\noindent
证明良好复形的拟同构 $a^\bullet : L^\bullet \to K^\bullet$
存在良好图表。选取满射 $p : R^{\oplus n} \to K^{-1}$。然后考虑新的
良好复形
$$
M^\bullet = \ldots \to 0 \to
L^{-1} \oplus R^{\oplus n} \xrightarrow{d \oplus 1}
L^0 \oplus R^{\oplus n} \to 0 \to \ldots
$$
，并取投影映射 $b^\bullet : M^\bullet \to L^\bullet$；映射
$c^\bullet : M^\bullet \to K^\bullet$ 在次数 $-1$ 使用
$a^{-1} \oplus p$，在次数 $0$ 使用 $a^0 \oplus d \circ p$。
映射 $b^\bullet : M^\bullet \to L^\bullet$ 与
$c^\bullet : M^\bullet \to K^\bullet$ 满足上述 (1)、(2)、(3)，
因而得到良好图表。

\medskip\noindent
假设有良好图表
$$
\xymatrix{
L^\bullet \ar[rr]_{\text{id}^\bullet} & & L^\bullet \\
& M^\bullet \ar[lu]^{b^\bullet} \ar[ru]_{c^\bullet}
}
$$
。由引理 \ref{more-algebra-lemma-homotopic-surjections}，
$\det(c^\bullet) = \det(b^\bullet)$。因此，
$\det(\text{id}^\bullet) = \text{id}$ 与良好图表的选择无关。

\medskip\noindent
在证明一般的无关性之前，先考虑复合。设有良好复形的拟同构
$L_1^\bullet \to L_2^\bullet$ 与 $L_2^\bullet \to L_3^\bullet$，
以及良好图表
$$
\vcenter{
\xymatrix{
L_1^\bullet \ar[rr] & &
L_2^\bullet \\
& M_{12}^\bullet \ar[lu] \ar[ru]
}
}
\quad\text{且}\quad
\vcenter{
\xymatrix{
L_2^\bullet \ar[rr] & &
L_3^\bullet \\
& M_{23}^\bullet \ar[lu] \ar[ru]
}
}
$$
。可以把它扩充为图表
$$
\xymatrix{
L_1^\bullet \ar[rr] & &
L_2^\bullet \ar[rr] & &
L_3^\bullet \\
& M_{12}^\bullet \ar[lu] \ar[ru]
& & M_{23}^\bullet \ar[lu] \ar[ru] \\
& &
M_{123}^\bullet \ar[lu] \ar[ru]
}
$$
，其中 $M_{123}^\bullet \to M_{12}^\bullet$ 与
$M_{123}^\bullet \to M_{23}^\bullet$ 具有性质 (1)、(2)、(3)，
且图中的方形交换：只需取
$M_{123}^n = M_{12}^n \times_{L_2^n} M_{23}^n$。于是引理
\ref{more-algebra-lemma-compose-surjections} 表明
$$
\xymatrix{
\det(L_2^\bullet) &
\det(M_{23}^\bullet) \ar[l] \\
\det(M_{12}^\bullet) \ar[u] &
\det(M_{123}^\bullet) \ar[l] \ar[u]
}
$$
交换。追图表明，由使用 $M_{12}^\bullet$ 与 $M_{23}^\bullet$ 的两个
良好图表所关联的映射之复合
$\det(L_1^\bullet) \to \det(L_2^\bullet) \to \det(L_3^\bullet)$，
等于良好图表
$$
\xymatrix{
L_1^\bullet \ar[rr] & & L_3^\bullet \\
& M_{123}^\bullet \ar[lu] \ar[ru]
}
$$
所关联的映射。因此，如果能够证明这些映射与选择无关，那么复合律也成立，
从而得到所需函子。

\medskip\noindent
无关性。给定良好复形的拟同构
$a^\bullet : L^\bullet \to K^\bullet$。选取逆拟同构
$b^\bullet : K^\bullet \to L^\bullet$。令
% P02-E0708
$L_1^\bullet = L$、$L_2^\bullet = K^\bullet$ 且
$L_3^\bullet = L^\bullet$，可以固定对 $b^\bullet$ 所选的良好图表，
并改变 $a^\bullet$ 的良好图表。于是前几段的结果表明，无论作何选择，
其复合总是 $\det(L^\bullet)$ 上的恒等映射。这显然证明了与这些选择无关。

\medskip\noindent
关于典范元素的陈述立即由引理
\ref{more-algebra-lemma-canonical-element-well-defined} 及我们的构造得到。
\end{proof}
\section{赋值环的扩张}
\label{more-algebra-section-valuation-rings}

\noindent
本节是章节 \ref{more-algebra-section-discrete-valuation-rings} 对一般赋值环的对应版本。

\begin{definition}
\label{more-algebra-definition-extension-valuation-rings}
若 $A$ 与 $B$ 都是赋值环，并且 $A \to B$ 是单射局部同态，
则称 $A \to B$ 或 $A \subset B$ 是{\it 赋值环的扩张}。
这样的扩张诱导交换图
$$
\xymatrix{
A \setminus \{0\} \ar[r] \ar[d]_v & B \setminus \{0\} \ar[d]^v \\
\Gamma_A \ar[r] & \Gamma_B
}
$$
，其中 $\Gamma_A$ 与 $\Gamma_B$ 是值群。若下方的水平箭头是双射，
则称 $B$ 在 $A$ 上{\it 弱非分歧}。若剩余域扩张
$\kappa_A = A/\mathfrak m_A \subset \kappa_B = B/\mathfrak m_B$
有限，则置 $f = [\kappa_B : \kappa_A]$，并称其为扩张
$A \subset B$ 的{\it 剩余次数}。
\end{definition}

\noindent
注意，$\Gamma_A \to \Gamma_B$ 是单射，因为在映射 $A \to B$ 下，
% P02-E0709
$A$ 的单位元是 $B$ 的单位元的逆。还要注意，我们并不要求分式域扩张有限。

\begin{lemma}
\label{more-algebra-lemma-inequality-general}
设 $A \subset B$ 是赋值环的扩张，分式域为 $K \subset L$。
若扩张 $L/K$ 有限，则剩余域扩张有限，$\Gamma_A$ 在 $\Gamma_B$
中的指数有限，并且
$$
[\Gamma_B : \Gamma_A] [\kappa_B : \kappa_A] \leq [L : K].
$$
\end{lemma}

\begin{proof}
设 $b_1, \ldots, b_n \in B$ 是单位元，并且它们在 $\kappa_B$ 中的像
在 $\kappa_A$ 上线性无关。设 $c_1, \ldots, c_m \in B$ 是非零元素，
并且它们在 $\Gamma_B/\Gamma_A$ 中的像两两不同。我们断言
$b_i c_j$ 在 $L$ 中于 $K$ 上线性无关。也就是说，断言一个
$a_{ij} \in K$ 不全为零的和
$$
\sum a_{ij} b_i c_j
$$
不可能为零。选取 $(i_0, j_0)$ 使
$v(a_{i_0j_0}b_{i_0}c_{j_0})$ 最小。把 $a_{ij}$ 替换为
$a_{ij}/a_{i_0j_0}$，从而 $a_{i_0 j_0} = 1$。令
$$
P = \{(i, j) \mid
v(a_{ij}b_ic_j) = v(a_{i_0j_0}b_{i_0}c_{j_0}) \}
$$
。由对 $c_1, \ldots, c_m$ 的选择，$(i, j) \in P$ 蕴含
$j = j_0$。因此若 $(i, j) \in P$，则
$v(a_{ij}) = v(a_{i_0j_0}) = 0$，即 $a_{ij}$ 是单位元。
由对 $b_1, \ldots, b_n$ 的选择，
$$
\sum\nolimits_{(i, j) \in P} a_{ij}b_i
$$
是 $B$ 中的单位元。因此
$\sum\nolimits_{(i, j) \in P} a_{ij}b_ic_j$ 的赋值为
$v(c_{j_0}) = v(a_{i_0j_0}b_{i_0}c_{j_0})$。第一个陈列和中
$(i, j) \not \in P$ 的各项具有严格更大的赋值，故该和不可能为零，
从而证明了引理。
\end{proof}

\begin{lemma}
\label{more-algebra-lemma-valuation-ring-purely-inseparable}
设 $A$ 是赋值环，其分式域 $K$ 的特征为 $p > 0$。设 $L/K$
是纯不可分扩张。则 $A$ 在 $L$ 中的整闭包 $B$ 是分式域为 $L$ 的赋值环，
并且 $A \subset B$ 是赋值环的扩张。
\end{lemma}

\begin{proof}
从略。提示：例如使用 Algebra，引理
\ref{algebra-lemma-x-or-x-inverse-valuation-ring} 与
\ref{algebra-lemma-integral-overring-surjective}。
\end{proof}

\begin{lemma}
\label{more-algebra-lemma-extension-normal-domains-and-roots}
设 $A \to B$ 是 Noether 局部正规整环的平坦局部同态。设
$f \in A$、$h \in B$，并且对某个 $n > 1$ 以及 $B$ 的某个单位元
$w$ 有 $f = w h^n$。假设对每个高为 $1$ 的素理想
$\mathfrak p \subset A$，都存在位于 $\mathfrak p$ 上方的高为 $1$
的素理想 $\mathfrak q \subset B$，使扩张
$A_\mathfrak p \subset B_\mathfrak q$ 弱非分歧。则对某个
$g \in A$ 及 $A$ 的某个单位元 $u$，有 $f = u g^n$。
\end{lemma}

\begin{proof}
$A$ 与 $B$ 在高为 $1$ 的素理想处的局部环都是离散赋值环
（Algebra，引理 \ref{algebra-lemma-characterize-dvr}）。因此该假设
是有意义的（通过定义
\ref{more-algebra-definition-extension-discrete-valuation-rings}）。
设 $\mathfrak p_1, \ldots, \mathfrak p_r$ 是 $A$ 中在 $f$ 上极小的
素理想。由 Algebra，引理 \ref{algebra-lemma-minimal-over-1}，
它们的高都是 $1$。对每个 $i$，设
$\mathfrak q_{i, j} \subset B$（$j = 1, \ldots, r_i$）是 $B$ 中
位于 $\mathfrak p_i$ 上方的高为 $1$ 的素理想。给它们编号，使得
$A_{\mathfrak p_i} \to B_{\mathfrak q_{i, 1}}$ 弱非分歧。
由于 $f$ 在 $B_{\mathfrak q_{i, 1}}$ 中映为一个 $n$ 次幂乘以单位元，
$f$ 在 $A_{\mathfrak p_i}$ 中的赋值 $v_i$ 被 $n$ 整除。写成
$v_i = n w_i$，其中 $w_i \geq 0$。考虑定义理想 $I$ 的正合列
$$
0 \to I \to A \to
\prod\nolimits_{i = 1, \ldots, r}
A_{\mathfrak p_i}/\mathfrak p_i^{w_i}A_{\mathfrak p_i}
$$
。应用正合函子 $- \otimes_A B$，得到正合列
$$
0 \to I \otimes_A B \to B \to
\prod\nolimits_{i = 1, \ldots, r}
(A_{\mathfrak p_i}/\mathfrak p_i^{w_i}A_{\mathfrak p_i}) \otimes_A B
$$
。固定 $i$。我们断言典范映射
$$
(A_{\mathfrak p_i}/\mathfrak p_i^{w_i}A_{\mathfrak p_i}) \otimes_A B
\to
\prod\nolimits_{j = 1, \ldots, r_i}
B_{\mathfrak q_{i, j}}/\mathfrak q_{i, j}^{e_{i, j}w_i}B_{\mathfrak q_{i, j}}
$$
是单射。这里 $e_{i, j}$ 是
$A_{\mathfrak p_i} \to B_{\mathfrak q_{i, j}}$ 的分歧指数。
该断言说，$\mathfrak p_i^{w_i}B_{\mathfrak p_i}$ 等于
$B_{\mathfrak p_i}$ 中在 $\mathfrak q_{i, j}$ 处的赋值
$\geq e_{i, j}w_i$ 的元素 $b$ 所成的集合。选取主理想
$\mathfrak p_i^{w_i}$ 的生成元
$a \in A_{\mathfrak p_i}$。则 $a$ 在 $\mathfrak q_{i, j}$ 处的赋值
等于 $e_{i, j}w_i$。由于 $B_{\mathfrak p_i}$ 是正规整环，其高为一的
素理想恰为 $\mathfrak q_{i, j}$（$j = 1, \ldots, r_i$），
对上述 $b$，由 Algebra，引理
\ref{algebra-lemma-normal-domain-intersection-localizations-height-1}
可知 $b/a \in B_{\mathfrak p_i}$。断言得证。

\medskip\noindent
该断言与上面的第二个正合列合在一起，确定正合列
$$
0 \to I \otimes_A B \to B \to
\prod\nolimits_{i = 1, \ldots, r}
\prod\nolimits_{j = 1, \ldots, r_i}
B_{\mathfrak q_{i, j}}/\mathfrak q_{i, j}^{e_{i, j}w_i}B_{\mathfrak q_{i, j}}
$$
。由此，$I \otimes_A B$ 就是 $B$ 中在
$\mathfrak q_{i, j}$ 处赋值 $\geq e_{i, j}w_i$ 的元素 $h'$
所成的集合。由于在 $B$ 中 $f = wh^n$，可知 $h$ 在
$\mathfrak q_{i, j}$ 处的赋值为 $e_{i, j}w_i$。于是由 Algebra，引理
\ref{algebra-lemma-normal-domain-intersection-localizations-height-1}，
$h'/h \in B$。因此 $I \otimes_A B$ 是秩为 $1$ 的自由 $B$-模
（由 $h$ 生成）。所以 $I$ 是秩为 $1$ 的自由 $A$-模，见 Algebra，
引理 \ref{algebra-lemma-finite-projective-descends}。设
$g \in I$ 是生成元。则 $g$ 与 $h$ 在 $B$ 中相差一个单位元。
倒推可知，$g$ 在 $A_{\mathfrak p_i}$ 中的赋值为
$w_i = v_i/n$。所以 $g^n$ 与 $f$ 在 $A$ 中相差一个单位元
（由 Algebra，引理
\ref{algebra-lemma-normal-domain-intersection-localizations-height-1}），
如所需。
\end{proof}

\begin{lemma}
\label{more-algebra-lemma-etale-extension-valuation-ring}
设 $A$ 是赋值环。设 $A \to B$ 是 \'etale 环映射，并设
$\mathfrak m \subset B$ 是位于 $A$ 的极大理想上方的素理想。
则 $A \subset B_\mathfrak m$ 是弱非分歧的赋值环扩张。
\end{lemma}

\begin{proof}
由引理 \ref{more-algebra-lemma-weak-dimension-at-most-1}，环 $A$ 的弱维数
$\leq 1$。再由引理 \ref{more-algebra-lemma-weak-dimension-goes-up} 与
\ref{more-algebra-lemma-when-weakly-etale}，$B$ 的弱维数 $\leq 1$。
% P02-E0710
故由引理 \ref{more-algebra-lemma-weak-dimension-at-most-1}，局部环
$B_\mathfrak m$ 是赋值环。由于扩张 $A \subset B_\mathfrak m$
诱导有限的分式域扩张，可知 $\Gamma_A$ 在 $B_{\mathfrak m}$
的值群中的指数有限。因此，对每个 $h \in B_\mathfrak m$，存在
$n > 0$、元素 $f \in A$ 以及单位元 $w \in B_\mathfrak m$，
使得在 $B_\mathfrak m$ 中 $f = w h^n$。我们将证明，这蕴含对某个
$g \in A$ 及单位元 $u \in A$ 有 $f = ug^n$；这将表明
$A$ 与 $B_\mathfrak m$ 的值群相同，从而得到引理的断言。

\medskip\noindent
把 $A = \colim A_i$ 写成它的各个在 $\mathbf{Z}$ 上本质有限型的局部
子环之余极限。由于 $A$ 是正规整环（Algebra，引理
\ref{algebra-lemma-valuation-ring-normal}），可以假设每个
$A_i$ 都正规（这里使用如下事实：取正规化后，局部环仍在
$\mathbf{Z}$ 上本质有限型；见 Algebra，命题
\ref{algebra-proposition-ubiquity-nagata}）。对某个 $i$，
可以找到 \'etale 扩张 $A_i \to B_i$，使得
$B = A \otimes_{A_i} B_i$，见 Algebra，引理
\ref{algebra-lemma-etale}。令 $\mathfrak m_i$ 为
$B_i$ 与 $\mathfrak m$ 的交。于是可以对环映射
$A_i \to (B_i)_{\mathfrak m_i}$ 应用引理
\ref{more-algebra-lemma-extension-normal-domains-and-roots} 得到结论。
该引理的假设成立，因为：
\begin{enumerate}
\item $A_i$ 与 $(B_i)_{\mathfrak m_i}$ 都是 Noether 的，因为它们在
$\mathbf{Z}$ 上本质有限型；
\item $A_i \to (B_i)_{\mathfrak m_i}$ 平坦，因为 $A_i \to B_i$
是 \'etale 的；
\item $B_i$ 正规，因为 $A_i \to B_i$ 是 \'etale 的，见 Algebra，
引理 \ref{algebra-lemma-normal-goes-up}；
\item 对 $A_i$ 的每个高为 $1$ 的素理想，都存在
$(B_i)_{\mathfrak m_i}$ 中位于其上方的高为 $1$ 的素理想；这由
Algebra，引理 \ref{algebra-lemma-finite-in-codim-1} 以及
$\Spec((B_i)_{\mathfrak m_i}) \to \Spec(A_i)$ 满射得到；
\item 对每个位于素理想 $\mathfrak p$ 上方的素理想 $\mathfrak q$，
诱导扩张 $(A_i)_\mathfrak p \to (B_i)_\mathfrak q$ 非分歧，因为
$A_i \to B_i$ 是 \'etale 的。
\end{enumerate}
这就完成了引理的证明。
\end{proof}

\begin{lemma}
\label{more-algebra-lemma-henselization-valuation-ring}
设 $A$ 是赋值环。设 $A^h$、$A^{sh}$ 分别是它的 Hensel 化与严格
Hensel 化。则
$$
A \subset A^h \subset A^{sh}
$$
是诱导值群双射的赋值环扩张，即弱非分歧扩张。
\end{lemma}

\begin{proof}
把 $A^h = \colim (B_i)_{\mathfrak q_i}$ 写成如下形式：$A \to B_i$
是 \'etale 的，而 $\mathfrak q_i \subset B_i$ 是位于
$\mathfrak m_A$ 上方的素理想；见 Algebra，引理
\ref{algebra-lemma-henselization-different}。引理
\ref{more-algebra-lemma-etale-extension-valuation-ring} 告诉我们，
$(B_i)_{\mathfrak q_i}$ 是赋值环，并且诱导映射
$$
(A \setminus \{0\})/A^* \longrightarrow
((B_i)_{\mathfrak q_i} \setminus \{0\}) / (B_i)_{\mathfrak q_i}^*
$$
是双射。由 Algebra，引理
\ref{algebra-lemma-colimit-valuation-rings}，$A^h$ 是赋值环。
而且，$(A \setminus \{0\})/A^* \to (A^h \setminus \{0\})/(A^h)^*$
也是双射。这证明了包含 $A \subset A^h$ 的情形。
为证明 $A \subset A^{sh}$，可以使用完全相同的论证，只需把
Algebra，引理 \ref{algebra-lemma-henselization-different} 换成
Algebra，引理 \ref{algebra-lemma-strict-henselization-different}。
由于 $A^{sh} = (A^h)^{sh}$，这也证明了包含
$A^h \subset A^{sh}$ 的所有断言。
\end{proof}
\section{主理想整环上模的结构}
\label{more-algebra-section-structure-modules}

\noindent
我们稍作推广（遵循 Warfield 的论文 \cite{Warfield-Purity} 与
\cite{Warfield-Decomposition}），使证明也适用于赋值环。

\begin{lemma}
\label{more-algebra-lemma-characterize-PD-modules}
\begin{reference}
\cite[Corollary 1]{Warfield-Purity}
\end{reference}
设 $P$ 是环 $R$ 上的模。下列条件等价：
\begin{enumerate}
\item $P$ 是形如 $R/fR$（其中 $f \in R$ 变化）的模之直和的直和因子。
\item 对每个满足如下条件的 $R$-模短正合列
$0 \to A \to B \to C \to 0$：对所有 $f \in R$ 都有
$fA = A \cap fB$，映射 $\Hom_R(P, B) \to \Hom_R(P, C)$ 满射。
\end{enumerate}
\end{lemma}

\begin{proof}
设 $0 \to A \to B \to C \to 0$ 是 (2) 中的正合列。为证明 (1) 蕴含
(2)，只需对每个 $f \in R$ 证明
$\Hom_R(R/fR, B) \to \Hom_R(R/fR, C)$ 满射。设
$\psi : R/fR \to C$ 是映射，并设 $\psi(1)$ 是 $b \in B$ 的像。
则 $fb \in A$。因此存在 $a \in A$ 使 $fa = fb$。于是
$f(b - a) = 0$，故得到把 $1$ 映到 $b - a$ 的态射
$\varphi : R/fR \to B$，它提升 $\psi$。

\medskip\noindent
反过来，假设 (2) 成立。设 $I$ 是所有二元组 $(f, \varphi)$ 所成的集合，
其中 $f \in R$ 且 $\varphi : R/fR \to P$。对 $i \in I$，以
$(f_i, \varphi_i)$ 表示对应的二元组。考虑映射
$$
B = \bigoplus\nolimits_{i \in I} R/f_iR \longrightarrow P
$$
，它把分量 $R/f_iR$ 中的元素 $r$ 映到 $P$ 中的
$\varphi_i(r)$。令 $A = \Ker(B \to P)$。那么，只要序列
$$
0 \to A \to B \to P \to 0
$$
是 (2) 中的正合列，(1) 就成立。为此，设 $f \in R$，且
$a \in A$ 在 $B$ 中等于 $f b$。写
$b = (r_i)_{i \in I}$，其中几乎所有 $r_i = 0$。于是 $P$ 中有
$$
f\sum \varphi_i(r_i) = 0
$$
。故存在 $i_0 \in I$，使得 $f_{i_0} = f$ 且
$\varphi_{i_0}(1) = \sum \varphi_i(r_i)$。设
$x_{i_0} \in R/f_{i_0}R$ 是 $1$ 的类。则
$$
a' = (r_i)_{i \in I} - (0, \ldots, 0, x_{i_0}, 0, \ldots )
$$
属于 $A$，并且 $fa' = a$，如所需。
\end{proof}

\begin{lemma}[广义赋值环]
\label{more-algebra-lemma-generalized-valuation-ring}
\begin{reference}
\cite{Warfield-Decomposition}
\end{reference}
设 $R$ 是非零环。下列条件等价：
\begin{enumerate}
\item 对 $a, b \in R$，或者 $a$ 整除 $b$，或者 $b$ 整除 $a$。
\item 每个有限生成理想都是主理想，且 $R$ 是局部环。
\item $R$ 的理想集在包含关系下全序。
\end{enumerate}
特别地，当 $R$ 是赋值环时这些条件成立。
\end{lemma}

\begin{proof}
假设 (2)，并取 $a, b \in R$。则 $(a, b) = (c)$。若 $c = 0$，
则 $a = b = 0$，故 $a$ 整除 $b$。假设 $c \not = 0$。写成
$c = ua + vb$、$a = wc$ 及 $b = zc$。于是
$c(1 - uw - vz) = 0$。由于 $R$ 是局部环，这蕴含
$1 - uw - vz \in \mathfrak m$。所以 $w$ 或 $z$ 中有一个是单位元，
从而 $a$ 整除 $b$ 或 $b$ 整除 $a$。故 (2) 蕴含 (1)。

\medskip\noindent
假设 (1)。若 $R$ 有两个极大理想 $\mathfrak m_i$，则可以选取
$a \in \mathfrak m_1$ 且 $a \not \in \mathfrak m_2$，以及
$b \in \mathfrak m_2$ 且 $b \not \in \mathfrak m_1$。此时
$a$ 不整除 $b$，而 $b$ 也不整除 $a$。因此 $R$ 有唯一极大理想，
从而是局部环。由条件 (1) 与归纳容易推出每个有限生成理想都是主理想。
所以 (1) 蕴含 (2)。

\medskip\noindent
(1) 与 (3) 等价的证明直接。最后一个陈述是 Algebra，引理
\ref{algebra-lemma-valuation-ring-x-or-x-inverse}。
\end{proof}

\begin{lemma}
\label{more-algebra-lemma-generalized-valuation-ring-modules}
\begin{reference}
\cite[Theorem 1]{Warfield-Decomposition}
\end{reference}
设环 $R$ 满足引理 \ref{more-algebra-lemma-generalized-valuation-ring} 中的等价条件。
则每个有限呈示 $R$-模都同构于有限个形如 $R/fR$ 的模之直和。
\end{lemma}

\begin{proof}
设 $M$ 是有限呈示 $R$-模。下文将不加说明地使用引理
\ref{more-algebra-lemma-generalized-valuation-ring} 中 $R$ 的全部等价性质。
记 $\mathfrak m \subset R$ 为极大理想，$\kappa = R/\mathfrak m$
为剩余域。设 $I \subset R$ 是 $M$ 的零化子。选取有限维
$\kappa$-向量空间 $M/\mathfrak m M$ 的一组基
$y_1, \ldots, y_n$。我们对 $n$ 归纳。

\medskip\noindent
由 Nakayama 引理，任意一组提升 $M/\mathfrak m M$ 中元素
$y_1, \ldots, y_n$ 的 $x_1, \ldots, x_n \in M$ 都生成 $M$，见
Algebra，引理 \ref{algebra-lemma-NAK}。这立即证明归纳的基础情形
$n = 0$。

\medskip\noindent
我们断言，存在指标 $i$，使得对任意映到 $y_i$ 的
$x_i \in M$，$x_i$ 的零化子都是 $I$。否则，可以选取
$x_1, \ldots, x_n$，使对所有 $i$ 都有
$I_i = \text{Ann}(x_i) \not = I$。但对所有 $i$ 都有
$I \subset I_i$；理想全序蕴含对 $i = 1, \ldots, n$，
$I_i$ 严格大于 $I$，再次使用全序便得到
$\text{Ann}(M) = I_1 \cap \ldots \cap I_n$ 大于 $I$，矛盾。
重新编号后，可以假设 $y_1$ 具有该性质：对任意提升 $y_1$ 的
$x_1 \in M$，$x_1$ 的零化子都是 $I$。

\medskip\noindent
置 $A = Rx_1 \subset M$。考虑正合列
$0 \to A \to M \to M/A \to 0$。由于 $A$ 有限，
$M/A$ 是生成元更少的有限呈示 $R$-模
（Algebra，引理 \ref{algebra-lemma-extension}）。因此由归纳，
$M/A \cong \bigoplus_{j = 1, \ldots, m} R/f_jR$。另一方面，我们断言
$A \to M$ 满足如下性质：若 $f \in R$，则
$fA = A \cap fM$。包含 $fA \subset A \cap fM$ 显然。反之，
若 $x \in A \cap fM$，则对某些 $g \in R$ 与 $y \in M$ 有
$x = gx_1 = f y$。若 $f$ 整除 $g$，则如所需有 $x \in fA$。
否则，可对某个 $h \in \mathfrak m$ 写成 $f = hg$。由上一段，
元素 $x'_1 = x_1 - hy$ 的零化子是 $I$。因此 $g \in I$，从而
$x = 0$，如所需。该断言与引理
\ref{more-algebra-lemma-characterize-PD-modules} 蕴含序列
$0 \to A \to M \to M/A \to 0$ 分裂，故
$M \cong A \oplus \bigoplus_{j = 1, \ldots, m} R/f_jR$。
于是 $A = R/I$ 是有限呈示的（因为它是 $M$ 的直和因子），所以
$I$ 有限生成，进而为主理想。证明完成。
\end{proof}

\begin{lemma}
\label{more-algebra-lemma-warfield}
\begin{reference}
\cite[Theorem 3]{Warfield-Decomposition}
\end{reference}
设 $R$ 是环，并且 $R$ 在每个极大理想处的局部环都满足引理
\ref{more-algebra-lemma-generalized-valuation-ring} 中的等价条件。则每个有限呈示
$R$-模都是某个有限直和的直和因子；该直和的各项形如
$R/fR$，其中 $f$ 在 $R$ 中变化。
\end{lemma}

\begin{proof}
设 $M$ 是有限呈示 $R$-模。先证明 $M$ 是形如 $R/fR$ 的模之直和的
直和因子，最后再证明该直和可以取为有限直和。设
$$
0 \to A \to B \to C \to 0
$$
是 $R$-模短正合列，并且对所有 $f \in R$ 都有
$fA = A \cap fB$。由引理 \ref{more-algebra-lemma-characterize-PD-modules}，
必须证明 $\Hom_R(M, B) \to \Hom_R(M, C)$ 满射。只需在各极大理想
$\mathfrak m$ 处局部化后证明这一点，见 Algebra，引理
\ref{algebra-lemma-characterize-zero-local}。注意，局部化序列
$0 \to A_\mathfrak m \to B_\mathfrak m \to C_\mathfrak m \to 0$
满足如下条件：对所有 $f \in R_\mathfrak m$，都有
$fA_\mathfrak m = A_\mathfrak m \cap fB_\mathfrak m$（因为可以写
$f = uf'$，其中 $u \in R_\mathfrak m$ 是单位元而 $f' \in R$，
并且局部化正合）。由于 $M$ 有限呈示，由 Algebra，引理
\ref{algebra-lemma-hom-from-finitely-presented}，
$$
\Hom_R(M, B)_\mathfrak m = \Hom_{R_\mathfrak m}(M_\mathfrak m, B_\mathfrak m)
\quad\text{且}\quad
\Hom_R(M, C)_\mathfrak m = \Hom_{R_\mathfrak m}(M_\mathfrak m, C_\mathfrak m)
$$
。模 $M_\mathfrak m$ 是有限呈示 $R_\mathfrak m$-模。由引理
\ref{more-algebra-lemma-generalized-valuation-ring-modules}，$M_\mathfrak m$
是形如 $R_\mathfrak m/fR_\mathfrak m$ 的模之直和。因此，由引理
\ref{more-algebra-lemma-characterize-PD-modules}，局部化后的映射满射。

\medskip\noindent
此时已知 $M$ 是 $\bigoplus_{i \in I} R/f_i R$ 的直和因子。
考虑映射 $M \to \bigoplus_{i \in I} R/f_i R$。由于 $M$ 是有限
$R$-模，对某个有限子集 $I' \subset I$，其像包含于
$\bigoplus_{i \in I'} R/f_i R$。证明完成。
\end{proof}

\begin{definition}
\label{more-algebra-definition-bezout}
设 $R$ 是整环。
\begin{enumerate}
\item 若 $R$ 的每个有限生成理想都是主理想，则称 $R$ 为
{\it B\'ezout 整环}。
\item 若对所有 $n , m \geq 1$ 以及每个 $n \times m$ 矩阵 $A$，
都存在大小分别为 $n \times n, m \times m$ 的可逆矩阵 $U, V$，使得
$$
U A V =
\left(
\begin{matrix}
f_1 & 0 & 0 & \ldots \\
0 & f_2 & 0 & \ldots \\
0 & 0 & f_3 & \ldots \\
\ldots & \ldots & \ldots & \ldots
\end{matrix}
\right)
$$
，其中 $f_1, \ldots, f_{\min(n, m)} \in R$ 且
$f_1 | f_2 | \ldots$，则称 $R$ 为{\it 初等因子整环}。
\end{enumerate}
\end{definition}

\noindent
每个 B\'ezout 整环 $R$ 是否都是初等因子整环，看来仍是一个未决问题。
这等价于问每个有限呈示 $R$-模是否都是循环模的直和。反向蕴含成立。

\begin{lemma}
\label{more-algebra-lemma-elementary-divisor-is-bezout}
初等因子整环是 B\'ezout 整环。
\end{lemma}

\begin{proof}
设 $a, b \in R$ 非零。考虑 $1 \times 2$ 矩阵 $A = (a\ b)$。
则 $u(a\ b)V = (f\ 0)$，其中 $u \in R$ 可逆，而
$V = (g_{ij})$ 是可逆 $2 \times 2$ 矩阵。于是
$f = u a g_{11} + u b g_{2 1}$ 且 $(g_{11}, g_{2 1}) = R$。
所以 $(a, b) = (f)$。归纳论证（从略）进而表明 $R$ 的每个有限生成
理想都由一个元素生成。
\end{proof}

\begin{lemma}
\label{more-algebra-lemma-localize-bezout}
B\'ezout 整环的局部化是 B\'ezout 整环。B\'ezout 整环的每个局部环
都是赋值环。局部整环是 B\'ezout 整环，当且仅当它是赋值环。
\end{lemma}

\begin{proof}
关于局部化的陈述之证明从略。最后一个陈述是 Algebra，引理
\ref{algebra-lemma-characterize-valuation-ring}。第二个陈述由另外两个推出。
\end{proof}

\begin{lemma}
\label{more-algebra-lemma-split-off-free-part}
设 $R$ 是 B\'ezout 整环。
\begin{enumerate}
\item 自由模的每个有限子模都是有限自由模。
\item 每个有限呈示 $R$-模 $M$ 都是一个有限自由模与一个挠模的直和，
其中该挠模 $M_{tors}$ 是形如
$\bigoplus_{i = 1, \ldots, n} R/f_iR$ 的模之直和因子，且
$f_1, \ldots, f_n \in R$ 非零。
\end{enumerate}
\end{lemma}

\begin{proof}
证明 (1)。设 $M \subset F$ 是自由模 $F$ 的有限子模。由于 $M$ 有限，
可以假设 $F$ 是有限自由模（细节从略）。设
$F = R^{\oplus n}$。对 $n$ 归纳。若 $n = 1$，则 $M$ 是有限生成理想，
因 $R$ 是 B\'ezout 整环而为主理想。若 $n > 1$，考虑 $M$ 在投影
$R^{\oplus n} \to R$（投到最后一个分量）下的像 $I$。
若 $I = (0)$，则 $M \subset R^{\oplus n - 1}$，由归纳得证。
若 $I \not = 0$，则 $I = (f) \cong R$。因此
$M \cong R \oplus \Ker(M \to I)$，再次由归纳得证。

\medskip\noindent
设 $M$ 是有限呈示 $R$-模。由于
% P02-E0711
$R$ 的局部化是极大理想处的赋值环（引理
\ref{more-algebra-lemma-localize-bezout}），可以应用引理 \ref{more-algebra-lemma-warfield}。
因此 $M$ 是形如
$R^{\oplus r} \oplus \bigoplus_{i = 1, \ldots, n} R/f_iR$
的模之直和因子，其中 $f_i \not = 0$。由于取挠子模是函子，可知
$M_{tors}$ 是 $\bigoplus_{i = 1, \ldots, n} R/f_iR$ 的直和因子，
而 $M/M_{tors}$ 是 $R^{\oplus r}$ 的直和因子。由证明的第一部分，
$M/M_{tors}$ 是有限自由模。因此
$M \cong M_{tors} \oplus M/M_{tors}$，如所需。
\end{proof}

\begin{lemma}
\label{more-algebra-lemma-modules-PID}
设 $R$ 是主理想整环。每个有限 $R$-模 $M$
% P02-E0712
都属于同构于形如
$$
R^{\oplus r} \oplus \bigoplus\nolimits_{i = 1, \ldots, n} R/f_iR
$$
的模，其中 $r, n \geq 0$，且 $f_1, \ldots, f_n \in R$ 非零。
\end{lemma}

\begin{proof}
主理想整环是 Noether B\'ezout 环。由引理
\ref{more-algebra-lemma-split-off-free-part}，只需在 $M$ 为挠模时证明结论。
由于 $M$ 有限，这意味着 $M$ 的零化子非零。设对某个非零
$f \in R$ 有 $fM = 0$。于是可把 $M$ 视为 $R/fR$-模。
由于 $R/fR$ 是维数为 $0$ 的 Noether 环（略去一个小细节），
% P02-E0713
$R/fR = \prod R_j$ 是 Artin 局部环 $R_i$ 的有限乘积
（Algebra，命题 \ref{algebra-proposition-dimension-zero-ring}）。
每个 $R_i$ 都是局部环及主理想整环的商，因而是引理
\ref{more-algebra-lemma-generalized-valuation-ring} 意义下的广义赋值环
（略去一个小细节）。写 $M = \prod M_j$，其中
$M_j = e_j M$，而 $e_j \in R/fR$ 是对应于因子 $R_j$ 的幂等元。
由引理 \ref{more-algebra-lemma-generalized-valuation-ring-modules}，
$M_j = \bigoplus_{i = 1, \ldots, n_j} R_j/\overline{f}_{ji}R_j$，
其中 $\overline{f}_{ji} \in R_j$。选取提升 $f_{ji} \in R$，并选取
% P02-E0714
$g_{ji} \in R$ 使 $(g_{ji}) = (f_j, f_{ji})$。于是作为 $R$-模，
$$
M \cong \bigoplus R/g_{ji}R
$$
，证明完成。
\end{proof}

\noindent
% P02-E0715
还可证明主理想整环是一个初等因子整环（在此插入未来的引用）；
一种方法是证明类似于下列引理的结果。

\begin{lemma}
\label{more-algebra-lemma-unimodular-vector}
设 $R$ 是 B\'ezout 整环。设 $n \geq 1$，并设
$f_1, \ldots, f_n \in R$ 生成单位理想。则存在 $R$ 中的可逆
$n \times n$ 矩阵，其第一行为 $f_1 \ldots f_n$。
\end{lemma}

\begin{proof}
这由引理 \ref{more-algebra-lemma-split-off-free-part} 得到；也可如下直接证明。
对 $n$ 归纳。$n = 1$ 时结论成立。假设 $n > 1$。重新编号后可以假设
$f_1 \not = 0$。选取 $f \in R$，使
$(f) = (f_1, \ldots, f_{n - 1})$。设 $A$ 是
$(n - 1) \times (n - 1)$ 矩阵，其第一行为
$f_1/f, \ldots, f_{n - 1}/f$。选取 $a, b \in R$，使
$af - bf_n = 1$；这是可能的，因为
$1 \in (f_1, \ldots, f_n) = (f, f_n)$。则一个解为矩阵
$$
\left(
\begin{matrix}
f & 0 & \ldots & 0 & f_n \\
0 & 1 & \ldots & 0 & 0 \\
  &   & \ldots \\
0 & 0 & \ldots & 1 & 0 \\
b & 0 & \ldots & 0 & a
\end{matrix}
\right)
\left(
\begin{matrix}
  & & & 0 \\
  & A \\
  & & & 0 \\
0 & \ldots & 0 & 1
\end{matrix}
\right)
$$
。注意，左侧矩阵的行列式为 $1$，故可逆。
\end{proof}
\section{主根理想}
\label{more-algebra-section-principal-radical-ideals}

\noindent
% P02-E0716
本节证明，一个链环 Noether 局部正规整环存在非平凡的主根理想。
该结果可见于 \cite{Artin-Lipman}。

\begin{lemma}
\label{more-algebra-lemma-polypoly}
设 $(R,\mathfrak m)$ 是一维 Noether 局部环，并设
$x\in\mathfrak m$ 不属于 $R$ 的任何极小素理想。则
\begin{enumerate}
\item 函数 $P : n \mapsto \text{length}_R(R/x^n R)$ 对
$n \geq 0$ 满足 $P(n) \leq n P(1)$；
\item 若 $x$ 是非零因子，则对 $n \geq 0$ 有 $P(n) = nP(1)$。
\end{enumerate}
\end{lemma}

\begin{proof}
由于 $\dim(R) = 1$，有 $\dim(R/x^n R) = 0$，故对每个 $n$，
$\text{length}_R(R/x^n R)$ 都有限
（Algebra，引理 \ref{algebra-lemma-support-point}）。
我们对 $n$ 归纳以证明引理。由于 $x^0 R = R$，有
$P(0) = \text{length}_R(R/x^0R) = \text{length}_R 0 = 0$。
$n = 1$ 时陈述也成立。现设 $n \geq 2$，并假设陈述对 $n - 1$
成立。下列序列正合：
$$
R/x^{n-1}R \xrightarrow{x} R/x^nR \to R/xR \to 0
$$
，其中 $x$ 表示乘以 $x$ 的映射。由于长度可加
（Algebra，引理 \ref{algebra-lemma-length-additive}），有
$P(n) \leq P(n - 1) + P(1)$。由归纳假设，
$P(n - 1) \leq (n - 1)P(1)$，所以 $P(n) \leq nP(1)$。
这证明了归纳步骤。

\medskip\noindent
若 $x$ 是非零因子，则上面的陈列正合列在左端也正合。因此对所有
$n \geq 1$，有 $P(n) = P(n - 1) + P(1)$。
\end{proof}

\begin{lemma}
\label{more-algebra-lemma-minprimespoly}
设 $(R, \mathfrak m)$ 是维数为 $1$ 的 Noether 局部环。设
$x \in \mathfrak m$ 不属于 $R$ 的任何极小素理想。设 $t$ 是 $R$
的极小素理想数目。则 $t \leq \text{length}_R(R/xR)$。
\end{lemma}

\begin{proof}
设 $\mathfrak p_1, \ldots, \mathfrak p_t$ 是 $R$ 的极小素理想。置
$R' = R/\sqrt{0} = R/(\bigcap_{i = 1}^t \mathfrak p_i)$。
我们断言，只需对 $R'$ 证明引理。事实上，显然 $R'$ 也有 $t$ 个极小
素理想，并且
$\text{length}_{R'}(R'/xR') = \text{length}_R(R'/xR')$
小于 $\text{length}_R(R/xR)$，因为存在满射
$R/xR \to R'/xR'$。因此可以假设 $R$ 既约。

\medskip\noindent
假设 $R$ 既约，其极小素理想为
$\mathfrak p_1, \ldots, \mathfrak p_t$。这意味着存在正合列
$$
0 \to R \to
\prod\nolimits_{i = 1}^t R/\mathfrak p_i \to Q \to 0
$$
。这里 $Q$ 是第一个映射的余核。写
$M = \prod_{i = 1}^t R/\mathfrak p_i$。在 $\mathfrak p_j$ 处局部化，
可见
$$
R_{\mathfrak p_j} \to M_{\mathfrak p_j} =
\left(\prod\nolimits_{i=1}^t R/\mathfrak p_i\right)_{\mathfrak p_j} =
(R/\mathfrak p_j)_{\mathfrak p_j}
$$
满射。因此对所有 $j$ 都有 $Q_{\mathfrak p_j} = 0$。由于
$\mathfrak m$ 是 $R$ 中唯一不同于各 $\mathfrak p_i$ 的素理想，
得到 $\text{Supp}(Q) = \{\mathfrak m\}$。所以 $Q$ 长度有限
（Algebra，引理 \ref{algebra-lemma-support-point}）。由于
$\text{Supp}(Q) = \{\mathfrak m\}$，可以选取 $n \gg 0$，使
$x^n$ 在 $Q$ 上作用为 $0$
（Algebra，引理
\ref{algebra-lemma-Noetherian-power-ideal-kills-module}）。
现考虑图表
$$
\xymatrix{
0 \ar[r] & R \ar[r] \ar[d]^-{x^n} & M
\ar[r] \ar[d]^-{x^n} & Q \ar[r] \ar[d]^-{x^n} & 0 \\
0 \ar[r] & R \ar[r] & M \ar[r] & Q \ar[r] & 0
}
$$
，其中竖直映射均为乘以 $x^n$。因为 $x$ 不属于任何
$\mathfrak p_i$，该映射在 $R$ 与 $M$ 上都是单射。由蛇引理
（Algebra，引理 \ref{algebra-lemma-snake}），下列序列正合：
$$
0 \to Q \to R/x^nR \to M/x^nM \to Q \to 0
$$
。所以当 $n$ 充分大时，
$\text{length}_R(R/x^nR) = \text{length}_R(M/x^nM)$。
写 $R_i = R/\mathfrak p_i$，则
$\text{length}(M/x^nM) =
\sum_{i = 1}^t \text{length}_R(R_i/x^nR_i)$。
应用引理 \ref{more-algebra-lemma-polypoly}，并利用 $x$ 在 $R$ 及 $R_i$ 上都是
非零因子，得到
$$
n \text{length}_R(R/xR) =
\sum\nolimits_{i = 1}^t n \text{length}_{R_i}(R_i/x R_i)
$$
。由于 $\text{length}_{R_i}(R_i/x R_i) \geq 1$，引理得证。
\end{proof}

\begin{lemma}
\label{more-algebra-lemma-sopexists}
设 $(R,\mathfrak m)$ 是维数 $d > 1$ 的 Noether 局部环，
设 $f \in \mathfrak m$ 不属于 $R$ 的任何极小素理想，并设
$k\in\mathbf{N}$。则存在
$g_1, \ldots, g_{d - 1} \in \mathfrak m^k$，使得
$f, g_1, \ldots, g_{d - 1}$ 是一组参数系。
\end{lemma}

\begin{proof}
由 Algebra，引理 \ref{algebra-lemma-one-equation}，
$\dim(R/fR) = d - 1$。在 $R/fR$ 中选取参数系
$\overline{g}_1, \ldots, \overline{g}_{d - 1}$
（Algebra，命题 \ref{algebra-proposition-dimension}），并在 $R$ 中
取提升 $g_1, \ldots, g_{d - 1}$。容易看出
$f, g_1, \ldots, g_{d - 1}$ 是 $R$ 的参数系。
于是 $f, g_1^k, \ldots, g_{d - 1}^k$ 也是参数系，证明完成。
\end{proof}

\begin{lemma}
\label{more-algebra-lemma-syspar}
设 $(R,\mathfrak m)$ 是二维 Noether 局部环，并设
$f \in \mathfrak m$ 不属于 $R$ 的任何极小素理想。则存在
$g \in \mathfrak m$ 与 $N \in \mathbf{N}$，使得
\begin{enumerate}
\item[(a)] $f,g$ 构成 $R$ 的一组参数系。
\item[(b)] 若 $h \in \mathfrak m^N$，则 $f + h, g$ 是参数系，并且
$\text{length}_R (R/(f, g)) = \text{length}_R(R/(f + h, g))$。
\end{enumerate}
\end{lemma}

\begin{proof}
由引理 \ref{more-algebra-lemma-sopexists}，存在 $g \in \mathfrak m$，使
$f, g$ 是 $R$ 的参数系。于是
$\mathfrak m = \sqrt{(f, g)}$。因此存在 $n$，使
$\mathfrak m^n \subset (f, g)$，见 Algebra，引理
\ref{algebra-lemma-Noetherian-power}。我们断言 $N = n + 1$ 可用。
事实上，设 $h \in \mathfrak m^N$。由 $N$ 的选择，可以写
$h = af + bg$，其中 $a, b \in \mathfrak m$。于是
$$
(f + h, g) = (f + af + bg, g) = ((1 + a)f, g) = (f, g)
$$
，因为 $1 + a$ 是 $R$ 中的单位元。这证明了长度相等，并证明
$f + h, g$ 是参数系。
\end{proof}

\begin{lemma}
\label{more-algebra-lemma-radical-element}
设 $R$ 是维数为 $2$ 的 Noether 局部正规整环。设
$\mathfrak p_1, \ldots, \mathfrak p_r$ 是两两不同的高为 $1$
的素理想。则存在非零元素
$f \in \mathfrak p_1 \cap \ldots \cap \mathfrak p_r$，使得
$R/fR$ 既约。
\end{lemma}

\begin{proof}
取非零元素
$f \in \mathfrak p_1 \cap \ldots \cap \mathfrak p_r$。
我们将稍微修改 $f$，得到生成根理想的元素。$R$ 在每个高为一
的素理想 $\mathfrak p$ 处的局部化 $R_\mathfrak p$ 都是离散赋值环，
见 Algebra，引理 \ref{algebra-lemma-characterize-dvr} 或
Algebra，引理 \ref{algebra-lemma-criterion-normal}。记
$\text{ord}_\mathfrak p(f)$ 为 $f$ 在 $R_{\mathfrak p}$ 中的相应
赋值。设 $\mathfrak q_1, \ldots, \mathfrak q_s$ 是包含 $f$
的所有不同的高为一的素理想。对每个 $j$，写
$\text{ord}_{\mathfrak q_j}(f) = m_j \geq 1$。定义
$\text{div}(f) = \sum_{j = 1}^s m_j\mathfrak q_j$，把它看作高为一
的素理想以整数为系数的形式线性组合。留意以后会用到：每个
$\mathfrak p_i$ 都出现在某个 $\mathfrak q_j$ 中。
$R/fR$ 既约，当且仅当对 $j = 1, \ldots, s$ 都有 $m_j = 1$。
事实上，若 $m_j$ 为 $1$，则
% P02-E0717
$(R/fR)\mathfrak q_j$ 既约，并且
$R/fR \subset \prod (R/fR)_{\mathfrak q_j}$，因为
% P02-E0718
$\mathfrak q_1, \ldots, \mathfrak q_j$ 是 $R/fR$ 的伴随素理想；
见 Algebra，引理 \ref{algebra-lemma-zero-at-ass-zero} 与
\ref{algebra-lemma-normal-domain-intersection-localizations-height-1}。

\medskip\noindent
选定引理 \ref{more-algebra-lemma-syspar} 中的 $g$ 与 $N$。对非零
$y \in R$，以 $t(y)$ 表示 $y$ 上的极小素理想数目。由于 $R$
是正规整环，这些素理想的高为一，并且与 $R/yR$ 的极小素理想形成
$1$ 对 $1$ 的对应
（Algebra，引理 \ref{algebra-lemma-minimal-over-1} 与
\ref{algebra-lemma-normal-domain-intersection-localizations-height-1}）。
例如，$t(f) = s$ 就是出现在 $\text{div}(f)$ 中的素理想
$\mathfrak q_j$ 的数目。设 $h \in \mathfrak m^N$。由引理
\ref{more-algebra-lemma-minprimespoly}，
\begin{align*}
t(f + h) & \leq \text{length}_{R/(f + h)}(R/(f + h, g)) \\
& = \text{length}_R(R/(f + h, g)) \\
& = \text{length}_R(R/(f, g))
\end{align*}
；第一个等式见 Algebra，引理
\ref{algebra-lemma-length-independent}。因此 $t(f + h)$ 有一个与
$h \in \mathfrak m^N$ 无关的上界。

\medskip\noindent
由上述有界性，可以选取
$h \in \mathfrak m^N \cap \mathfrak p_1 \cap \ldots \cap \mathfrak p_r$，
使 $t(f + h)$ 在这样的 $h$ 中达到最大值。置 $f' = f + h$。
给定
$h' \in \mathfrak m^N \cap \mathfrak p_1 \cap \ldots \cap \mathfrak p_r$，
有 $t(f' + h') \leq t(f + h)$。因此，把 $f$ 换成 $f'$ 后，可以假设
对每个
$h \in \mathfrak m^N \cap \mathfrak p_1 \cap \ldots \cap \mathfrak p_r$
都有 $t(f + h) \leq s$。

\medskip\noindent
接着，假设能找到 $h \in \mathfrak m^N$，使对每个 $j$ 都有
$\text{ord}_{\mathfrak q_j}(h) \geq 1$，并且
$\text{ord}_{\mathfrak q_j}(h) = 1 \Leftrightarrow m_j > 1$。
注意，
$h \in \mathfrak m^N \cap \mathfrak p_1 \cap \ldots \cap \mathfrak p_r$。
由赋值的初等性质，对每个 $j$ 都有
$\text{ord}_{\mathfrak q_j}(f + h) = 1$。因此
$$
\text{div}(f + h) = \sum\nolimits_{j = 1}^s \mathfrak q_j +
\sum\nolimits_{k = 1}^v e_k \mathfrak r_k
$$
，其中 $\mathfrak r_1, \ldots, \mathfrak r_v$ 是两两不同的高为一
的素理想，且 $e_k \geq 1$。但是
$s = t(f) \geq t(f + h)$，所以 $v = 0$，我们找到了所需元素。

\medskip\noindent
现在选取满足上述条件的 $h$。由素理想回避
（Algebra，引理 \ref{algebra-lemma-silly}），对每个
$1 \leq j \leq s$，可以找到 $a_j \in \mathfrak q_j$，使
$a_j \not \in \mathfrak q_{j'}$（$j' \not = j$），且
$a_j \not \in \mathfrak q_j^{(2)}$。这里
$\mathfrak q_j^{(2)} = \{x \in R \mid \text{ord}_{\mathfrak q_j}(x) \geq 2\}$
是 $\mathfrak q_j$ 的第二符号幂。然后取
$$
h = \prod\nolimits_{m_j = 1} a_j^2 \times
\prod\nolimits_{m_j > 1} a_j
$$
。显然 $h$ 满足上面对赋值施加的条件。若
$h \not \in \mathfrak m^N$，则把它乘以 $\mathfrak m^N$ 中一个对所有
$j$ 都不属于 $\mathfrak q_j$ 的元素。
\end{proof}

\begin{lemma}
\label{more-algebra-lemma-divides-radical}
设 $(A, \mathfrak m, \kappa)$ 是维数为 $2$ 的 Noether 局部正规整环。
若 $a \in \mathfrak m$ 非零，则存在 $c \in A$，使得 $A/cA$
既约，并且对某个 $n$，$a$ 整除 $c^n$。
\end{lemma}

\begin{proof}
沿用引理 \ref{more-algebra-lemma-radical-element} 证明中的记号，写
$\text{div}(a) = \sum_{i = 1, \ldots, r} n_i \mathfrak p_i$。
由引理 \ref{more-algebra-lemma-radical-element}，选取
$c \in \mathfrak p_1 \cap \ldots \cap \mathfrak p_r$，使 $A/cA$
既约。对 $n \geq \max(n_i)$，
$-\text{div}(a) + \text{div}(c^n)$ 是有效除子（所有系数非负）。
因此，由 Algebra，引理
\ref{algebra-lemma-normal-domain-intersection-localizations-height-1}，
$c^n/a \in A$。
\end{proof}

\noindent
本节余下部分证明维数 $> 2$ 时的结论。

\begin{lemma}
\label{more-algebra-lemma-multiplicity}
设 $(R, \mathfrak m)$ 是维数为 $d$ 的 Noether 局部环，
$g_1, \ldots, g_d$ 是参数系，并令 $I = (g_1, \ldots, g_d)$。
若数值多项式
$n \mapsto \text{length}_R(R/I^{n+1})$ 的首项系数为 $e_I/d!$，则
$e_I \leq \text{length}_R(R/I)$。
\end{lemma}

\begin{proof}
由 Algebra，命题
\ref{algebra-proposition-hilbert-function-polynomial}，该函数是数值多项式。
由 Algebra，命题 \ref{algebra-proposition-dimension}，它的次数为 $d$。
若 $d = 0$，结论显然。若 $d = 1$，结论就是引理
\ref{more-algebra-lemma-polypoly}。为证明一般情形，注意存在满射
$$
\bigoplus\nolimits_{i_1, \ldots, i_d \geq 0,\ \sum i_j = n} R/I
\longrightarrow
I^n/I^{n + 1}
$$
，它把对应于 $i_1, \ldots, i_d$ 的基元素送到
$I^n/I^{n + 1}$ 中 $g_1^{i_1} \ldots g_d^{i_d}$ 的类。因此
$$
\text{length}_R(R/I^{n + 1}) - \text{length}_R(R/I^n)
\leq  \text{length}_R(R/I) {n + d - 1 \choose d - 1}
$$
。由于 $d \geq 2$，左边的数值多项式次数为 $d - 1$，首项系数为
$e_I / (d - 1)!$。右边的多项式次数为 $d - 1$，首项系数为
$\text{length}_R(R/I)/ (d - 1)!$。这就证明了引理。
\end{proof}

\begin{lemma}
\label{more-algebra-lemma-minprimespolyhigher}
设 $(R, \mathfrak m)$ 是维数为 $d$ 的 Noether 局部环，$t$ 是 $R$
中维数为 $d$ 的极小素理想数目，并设 $(g_1,\ldots,g_d)$ 是参数系。则
% P02-E0719
$t \leq \text{length}_R(R/(g_1,\ldots,g_n))$。
\end{lemma}

\begin{proof}
若 $d = 0$，引理显然。若 $d = 1$，引理就是引理
\ref{more-algebra-lemma-minprimespoly}。因此可以假设 $d > 1$。设
$\mathfrak p_1, \ldots, \mathfrak p_s$ 是 $R$ 的极小素理想，其中前
$t$ 个的维数为 $d$，并记
% P02-E0719
$I = (g_1, \ldots, g_n)$。完全仿照引理
\ref{more-algebra-lemma-minprimespoly} 的证明，可以假设 $R$ 既约。

\medskip\noindent
% P02-E0720
假设 $R$ 既约，其极小素理想为
$\mathfrak p_1, \ldots, \mathfrak p_t$。这意味着存在正合列
$$
0 \to R \to
\prod\nolimits_{i = 1}^t R/\mathfrak p_i \to Q \to 0
$$
。这里 $Q$ 是第一个映射的余核。记
$M = \prod_{i = 1}^t R/\mathfrak p_i$。在 $\mathfrak p_j$ 处局部化，得到
$$
R_{\mathfrak p_j} \to M_{\mathfrak p_j} =
\left(\prod\nolimits_{i=1}^t R/\mathfrak p_i\right)_{\mathfrak p_j} =
(R/\mathfrak p_j)_{\mathfrak p_j}
$$
满射。因此对所有 $j$ 都有 $Q_{\mathfrak p_j} = 0$。所以 $R$ 的高为
$0$ 的素理想都不在 $Q$ 的支集内。由此，数值多项式
$n \mapsto \text{length}_R(Q/I^nQ)$ 的次数等于
$\dim(\text{Supp}(Q)) < d$；见 Algebra，引理
\ref{algebra-lemma-support-dimension-d}。由 Algebra，引理
\ref{algebra-lemma-hilbert-ses-chi}（它适用是因为 $R$ 不是有限长度），
多项式
$$
n \longmapsto
\text{length}_R(M/I^nM) - \text{length}_R(R/I^n) - \text{length}_R(Q/I^nQ)
$$
的次数 $< d$。由于 $M = \prod R/\mathfrak p_i$，并且
% P02-E0721
$n \to \text{length}_R(R/\mathfrak p_i + I^n)$ 对
$i = 1, \ldots, t$ 是次数恰好为(!) $d$ 的数值多项式（由 Algebra，引理
\ref{algebra-lemma-support-dimension-d}），可见
$n \mapsto \text{length}_R(M/I^nM)$ 的首项系数至少为 $t/d!$。
于是由引理 \ref{more-algebra-lemma-multiplicity} 得到结论。
\end{proof}

\begin{lemma}
\label{more-algebra-lemma-sysparhigher}
设 $(R, \mathfrak m)$ 是维数为 $d$ 的 Noether 局部环，且
$f \in \mathfrak m$ 不属于 $R$ 的任何极小素理想。则存在元素
$g_1, \ldots, g_{d - 1} \in \mathfrak m$ 与 $N \in \mathbf{N}$，使得
\begin{enumerate}
\item $f, g_1, \ldots, g_{d - 1}$ 构成 $R$ 的参数系；
\item 若 $h \in \mathfrak m^N$，则
$f + h, g_1, \ldots, g_{d - 1}$ 是参数系，并且
$\text{length}_R R/(f, g_1, \ldots, g_{d-1}) =
\text{length}_R R/(f + h, g_1, \ldots, g_{d-1})$。
\end{enumerate}
\end{lemma}

\begin{proof}
由引理 \ref{more-algebra-lemma-sopexists}，存在
$g_1, \ldots, g_{d - 1} \in \mathfrak m$，使
$f, g_1, \ldots, g_{d - 1}$ 是 $R$ 的参数系。于是
$\mathfrak m = \sqrt{(f, g_1, \ldots, g_{d - 1})}$。因此存在 $n$，使
% P02-E0722
$\mathfrak m^n \subset (f, g)$；见 Algebra，引理
\ref{algebra-lemma-Noetherian-power}。我们断言 $N = n + 1$ 可用。
事实上，设 $h \in \mathfrak m^N$。由 $N$ 的选择，可以写
$h = af + \sum b_ig_i$，其中 $a, b_i \in \mathfrak m$。于是
\begin{align*}
(f + h, g_1, \ldots, g_{d - 1})
& =
(f + af + \sum b_ig_i, g_1, \ldots, g_{d - 1}) \\
& =
((1 + a)f, g_1, \ldots, g_{d - 1}) \\
& =
(f, g_1, \ldots, g_{d - 1})
\end{align*}
，因为 $1 + a$ 是 $R$ 中的单位元。这证明了长度相等，并证明
$f + h, g_1, \ldots, g_{d - 1}$ 是参数系。
\end{proof}

\begin{proposition}
\label{more-algebra-proposition-propdimd}
\begin{reference}
\cite[Lemma 3.14]{Artin-Lipman} 在不假设环是链环的情形下给出了这一结果
\end{reference}
设 $R$ 是链环 Noether 局部正规整环。设 $J \subset R$ 是根理想。
则存在非零元素 $f \in J$，使得 $R/fR$ 既约。
\end{proposition}

\begin{proof}
证明与引理 \ref{more-algebra-lemma-radical-element} 的证明相同，只需以引理
\ref{more-algebra-lemma-minprimespolyhigher} 代替引理 \ref{more-algebra-lemma-minprimespoly}，并以
引理 \ref{more-algebra-lemma-sysparhigher} 代替引理 \ref{more-algebra-lemma-syspar}。
可以使用引理 \ref{more-algebra-lemma-minprimespolyhigher}，因为 $R$ 是链整环，所以
$R$ 中每个高为一的素理想的维数都是 $d - 1$，从而
$R/(f + h)$ 的谱是等维的。为方便读者，下面写出细节。

\medskip\noindent
取非零元素 $f \in J$。我们将稍微修改 $f$，得到生成根理想的元素。
$R$ 在每个高为一的素理想 $\mathfrak p$ 处的局部化
$R_\mathfrak p$ 都是离散赋值环；见 Algebra，引理
\ref{algebra-lemma-characterize-dvr} 或 Algebra，引理
\ref{algebra-lemma-criterion-normal}。记
$\text{ord}_\mathfrak p(f)$ 为 $f$ 在 $R_{\mathfrak p}$ 中的相应赋值。
设 $\mathfrak q_1, \ldots, \mathfrak q_s$ 是包含 $f$ 的所有不同的
高为一的素理想。对每个 $j$，写
$\text{ord}_{\mathfrak q_j}(f) = m_j \geq 1$。定义
$\text{div}(f) = \sum_{j = 1}^s m_j\mathfrak q_j$，把它看作高为一
的素理想以整数为系数的形式线性组合。$R/fR$ 既约，当且仅当
对 $j = 1, \ldots, s$ 都有 $m_j = 1$。事实上，若 $m_j$ 为 $1$，则
% P02-E0717
$(R/fR)\mathfrak q_j$ 既约，并且
$R/fR \subset \prod (R/fR)_{\mathfrak q_j}$，因为
% P02-E0718
$\mathfrak q_1, \ldots, \mathfrak q_j$ 是 $R/fR$ 的伴随素理想；见
Algebra，引理 \ref{algebra-lemma-zero-at-ass-zero} 与
\ref{algebra-lemma-normal-domain-intersection-localizations-height-1}。

\medskip\noindent
% P02-E0723
选定引理 \ref{more-algebra-lemma-sysparhigher} 中的
$g_2, \ldots, g_{d - 1}$ 与 $N$。对非零 $y \in R$，以 $t(y)$ 表示
$y$ 上的极小素理想数目。由于 $R$ 是正规整环，这些素理想的高为一，
并且与 $R/yR$ 的极小素理想形成 $1$ 对 $1$ 的对应（Algebra，引理
\ref{algebra-lemma-minimal-over-1} 与
\ref{algebra-lemma-normal-domain-intersection-localizations-height-1}）。
例如，$t(f) = s$ 就是出现在 $\text{div}(f)$ 中的素理想
$\mathfrak q_j$ 的数目。设 $h \in \mathfrak m^N$。因为 $R$ 是链环，
对 $R$ 的每个高为一的素理想 $\mathfrak p$，都有
% P02-E0724
$\dim(R/\mathfrak p) = d$。因此由引理
\ref{more-algebra-lemma-minprimespolyhigher}，
\begin{align*}
t(f + h) & \leq \text{length}_{R/(f + h)}(R/(f + h, g_1, \ldots, g_{d - 1})) \\
& = \text{length}_R(R/(f + h, g_1, \ldots, g_{d - 1})) \\
& = \text{length}_R(R/(f, g_1, \ldots, g_{d - 1}))
\end{align*}
；第一个等式见 Algebra，引理
\ref{algebra-lemma-length-independent}。因此 $t(f + h)$ 有一个与
$h \in \mathfrak m^N$ 无关的上界。

\medskip\noindent
由上述有界性，可以选取 $h \in \mathfrak m^N \cap J$，使
$t(f + h)$ 在这样的 $h$ 中达到最大值。置 $f' = f + h$。
给定 $h' \in \mathfrak m^N \cap J$，有
$t(f' + h') \leq t(f + h)$。因此，把 $f$ 换成 $f'$ 后，可以假设
对每个 $h \in \mathfrak m^N \cap J$ 都有 $t(f + h) \leq s$。

\medskip\noindent
接着，假设能找到 $h \in \mathfrak m^N \cap J$，使对每个 $j$ 都有
$\text{ord}_{\mathfrak q_j}(h) \geq 1$，并且
$\text{ord}_{\mathfrak q_j}(h) = 1 \Leftrightarrow m_j > 1$。
由赋值的初等性质，对每个 $j$ 都有
$\text{ord}_{\mathfrak q_j}(f + h) = 1$。因此
$$
\text{div}(f + h) = \sum\nolimits_{j = 1}^s \mathfrak q_j +
\sum\nolimits_{k = 1}^v e_k \mathfrak r_k
$$
，其中 $\mathfrak r_1, \ldots, \mathfrak r_v$ 是两两不同的高为一
的素理想，且 $e_k \geq 1$。但是
$s = t(f) \geq t(f + h)$，所以 $v = 0$，我们找到了所需元素。

\medskip\noindent
现在选取满足上述条件的 $h$。由素理想回避（Algebra，引理
\ref{algebra-lemma-silly}），对每个 $1 \leq j \leq s$，可以找到
$a_j \in \mathfrak q_j \cap J$，使
$a_j \not \in \mathfrak q_{j'}$（$j' \not = j$）。接着可以选取
% P02-E0725
$b_j \in J \cap \mathfrak q_1 \cap \ldots \cap q_s$，使
$b_j \not \in \mathfrak q_j^{(2)}$。这里
$\mathfrak q_j^{(2)} = \{x \in R \mid \text{ord}_{\mathfrak q_j}(x) \geq 2\}$
是 $\mathfrak q_j$ 的第二符号幂。素理想回避可以应用，因为理想
% P02-E0725
$J' = J \cap \mathfrak q_1 \cap \ldots \cap q_s$ 是根理想，故
$R/J'$ 既约，故 $(R/J')_{\mathfrak q_j}$ 既约，故 $J'$ 包含一个元素
$x$，满足 $\text{ord}_{\mathfrak q_j}(x) = 1$；因此
$J' \not \subset \mathfrak q_j^{(2)}$。于是元素
$$
c = \sum\nolimits_{j = 1, \ldots, s}
b_j \times \prod\nolimits_{j' \not = j} a_{j'}
$$
属于 $J$，且由赋值的初等性质，对所有 $j = 1, \ldots, s$ 都有
$\text{ord}_{\mathfrak q_j}(c) = 1$。最后令
$$
h = c \times \prod\nolimits_{m_j = 1} a_j \times y
$$
，其中 $y \in \mathfrak m^N$ 是对所有 $j$ 都不属于 $\mathfrak q_j$
的元素。
\end{proof}

\section{导出范畴中的可逆对象}
\label{more-algebra-section-invertible-D-or-R}

\noindent
我们刻画环的导出范畴中的可逆对象。

\begin{lemma}
\label{more-algebra-lemma-symmetric-monoidal-derived}
设 $R$ 是环。$R$ 的导出范畴 $D(R)$ 是对称幺半范畴；其张量积由导出
张量积给出，结合约束与交换约束如第 \ref{more-algebra-section-sign-rules} 节所述。
\end{lemma}

\begin{proof}
略。提示：结合约束是引理 \ref{more-algebra-lemma-triple-tensor-product} 的同构，交换
约束是引理 \ref{more-algebra-lemma-flip-douoble-tensor-product} 的同构。说明这些以后，
各图表的交换性由 $R$-模复形范畴的相应结果得到；见第
\ref{more-algebra-section-symmetric-monoidal} 节。
\end{proof}

\noindent
因此，我们知道 $D(R)$ 的对象有（左）对偶或可逆是什么意思。在具体确定
这等价于什么之前，需要一个简单引理。

\begin{lemma}
\label{more-algebra-lemma-complex-bounded-above-free-colim-bounded-finite-free}
设 $R$ 是环，$F^\bullet$ 是自由 $R$-模组成的上有界复形。给定成对数据
$(n_i, f_i)$，$i = 1, \ldots, N$，其中 $n_i \in \mathbf{Z}$ 且
$f_i \in F^{n_i}$，则存在包含所有 $f_i$ 的子复形
$G^\bullet \subset F^\bullet$，它有界且由有限自由 $R$-模组成。
\end{lemma}

\begin{proof}
对 $a = \min(n_i; i = 1, \ldots, N)$ 作降归纳。若对 $n \geq a$ 都有
$F^n = 0$，则取 $G^\bullet$ 为零复形即可。一般而言，重新编号后，可以
假设存在
% P02-E0726
$1 \leq r \leq N$，使 $n_1 = \ldots = n_r = a$，并且对 $i > r$
有 $n_i > a$。为 $F^a$ 选择一组基 $b_j, j \in J$。可以选取有限子集
$J' \subset J$，使对 $i = 1, \ldots, r$ 都有
$f_i \in \bigoplus_{j \in J'} Rb_j$。为 $F^{a + 1}$ 选择一组基
$c_k, k \in K$。可以选取有限子集 $K' \subset K$，使对 $j \in J'$
都有 $\text{d}_F^a(b_j) \in \bigoplus_{k \in K'} Rc_k$。于是可以应用
归纳假设，找到子复形 $H^\bullet \subset F^\bullet$，它包含
$c_k \in F^{a + 1}$（$k \in K'$）以及 $f_i \in F^{n_i}$（$i > r$）。
令 $G^\bullet$ 在次数 $> a$ 时等于 $H^\bullet$，在次数 $a$ 时等于
$\bigoplus_{j \in J'} Rb_j$。
\end{proof}

\begin{lemma}
\label{more-algebra-lemma-have-dual-derived}
设 $R$ 是环，$M$ 是 $D(R)$ 的对象。以下条件等价：
\begin{enumerate}
\item $M$ 在 $D(R)$ 中有 Categories，定义
\ref{categories-definition-dual} 意义下的左对偶；
\item $M$ 是 $D(R)$ 的完美对象。
\end{enumerate}
此外，在这种情形下，$M$ 的左对偶是引理
\ref{more-algebra-lemma-dual-perfect-complex} 的对象 $M^\vee$。
\end{lemma}

\begin{proof}
若 $M$ 完美，则可以用有限投射 $R$-模组成的有界复形 $M^\bullet$ 表示
$M$。此时，由引理 \ref{more-algebra-lemma-left-dual-complex}，$M^\bullet$ 在复形
范畴中有左对偶，因而它当然也是 $D(R)$ 中的左对偶。

\medskip\noindent
假设 (1)。设 $N$，
$\eta : R \to M \otimes_R^\mathbf{L} N$，以及
% P02-E0727
$\epsilon : M \otimes_R^\mathbf{L} N \to R$ 构成 Categories，定义
\ref{categories-definition-dual} 意义下的左对偶。选择表示 $M$ 的复形
$M^\bullet$。选择一个各项平坦、表示 $N$ 的
% P02-E0728
K-平坦复形 $N^\bullet$；见引理 \ref{more-algebra-lemma-K-flat-resolution}。于是
$\eta$ 由复形映射
$$
\eta : R \longrightarrow \text{Tot}(M^\bullet \otimes_R N^\bullet)
$$
给出。可以把 $1$ 的像写成有限和
$$
\eta(1) = \sum\nolimits_n \sum\nolimits_i m_{n, i} \otimes n_{-n, i}
$$
，其中 $m_{n, i} \in M^n$ 且 $n_{-n, i} \in N^{-n}$。令
$K^\bullet \subset M^\bullet$ 为所有元素 $m_{n, i}$ 与
$\text{d}(m_{n, i})$ 生成的子复形。由 $N^\bullet$ 的选择，有
$\text{Tot}(K^\bullet \otimes_R N^\bullet) \subset
\text{Tot}(M^\bullet \otimes_R N^\bullet)$，而根据上面的选择，
$\eta(1)$ 位于这个子复形中。以 $K$ 表示 $K^\bullet$ 所表示的
$D(R)$ 的对象。于是 $\eta$ 通过映射
$\tilde \eta : R \longrightarrow K \otimes_R^\mathbf{L} N$ 分解。
由于
$(1 \otimes \epsilon) \circ (\eta \otimes 1) = \text{id}_M$，
由交换图表
$$
\xymatrix{
M \ar[rr]_-{\eta \otimes 1} \ar[rrd]_{\tilde \eta \otimes 1} & &
M \otimes_R^\mathbf{L} N \otimes_R^\mathbf{L} M
\ar[r]_-{1 \otimes \epsilon} &
M \\
& &
K \otimes_R^\mathbf{L} N \otimes_R^\mathbf{L} M
\ar[u] \ar[r]^-{1 \otimes \epsilon} &
K \ar[u]
}
$$
可知 $M$ 上的恒等映射通过 $K$ 分解。由于 $K$ 上有界，得到
$M \in D^-(R)$。因此，可以用自由 $R$-模组成的上有界复形
$M^\bullet$ 表示 $M$；例如见 Derived Categories，引理
\ref{derived-lemma-subcategory-left-resolution}。像以前一样写
$\eta(1) = \sum\nolimits_n \sum\nolimits_i m_{n, i} \otimes n_{-n, i}$。
由引理 \ref{more-algebra-lemma-complex-bounded-above-free-colim-bounded-finite-free}，
可以找到包含所有元素 $m_{n, i}$ 的子复形
$K^\bullet \subset M^\bullet$，它有界且由有限自由 $R$-模组成。
与上面一样，$M$ 上的恒等映射通过 $K$ 分解。由于 $K$ 完美，得到
$M$ 也完美；见引理 \ref{more-algebra-lemma-summands-perfect}。
\end{proof}

\begin{lemma}
\label{more-algebra-lemma-invertible-derived}
设 $R$ 是环，$M$ 是 $D(R)$ 的对象。以下条件等价：
\begin{enumerate}
\item $M$ 在 $D(R)$ 中可逆；见 Categories，定义
\ref{categories-definition-invertible}；
\item 对每个素理想 $\mathfrak p \subset R$，存在
$f \in R$，$f \not \in \mathfrak p$，使对某个 $n \in \mathbf{Z}$
有 $M_f \cong R_f[-n]$。
\end{enumerate}
此外，在这种情形下：
\begin{enumerate}
\item[(a)] $M$ 是 $D(R)$ 的完美对象；
\item[(b)] 在 $D(R)$ 中，$M = \bigoplus H^n(M)[-n]$；
\item[(c)] 每个 $H^n(M)$ 都是有限投射 $R$-模；
\item[(d)] 可以写成 $R = \prod_{a \leq n \leq b} R_n$，使
$H^n(M)$ 对应于可逆 $R_n$-模。
\end{enumerate}
\end{lemma}

\begin{proof}
假设 (2)。考虑对象 $R\Hom_R(M, R)$ 与复合映射
$$
R\Hom(M, R) \otimes_R^\mathbf{L} M \to R
$$
。在局部检查可知这是同构；细节从略。由于 $D(R)$ 是对称幺半范畴，
$M$ 可逆。

\medskip\noindent
假设 (1)。注意，幺半范畴中的可逆对象有左对偶，即它的逆。因此由引理
\ref{more-algebra-lemma-have-dual-derived}，$M$ 完美。考虑素理想
$\mathfrak p \subset R$，其剩余域为 $\kappa$。于是
$M \otimes_R^\mathbf{L} \kappa$ 是 $D(\kappa)$ 的可逆对象。显然，这
意味着 $\dim H^i(M \otimes_R^\mathbf{L} \kappa)$ 恰好对一个 $i$
非零，并且此时等于 $1$。由引理
\ref{more-algebra-lemma-lift-perfect-from-residue-field}，这给出 (2)。

\medskip\noindent
在上面的证明中已经看到 (a) 成立。令 $U_n \subset \Spec(R)$ 为所有形如
$D(f)$ 且满足 $M_f \cong R_f[-n]$ 的开集之并。显然，若
$n \not = n'$，则 $U_n \cap U_{n'} = \emptyset$。若 $M$ 的 tor-幅度在
$[a, b]$ 内，则当 $n \not \in [a, b]$ 时 $U_n = \emptyset$。因此得到
(d) 中的乘积分解 $R = \prod_{a \leq n \leq b} R_n$，使 $U_n$ 对应于
$\Spec(R_n)$；见 Algebra，引理
\ref{algebra-lemma-disjoint-implies-product}。由于
$D(R) = \prod_{a \leq n \leq b} D(R_n)$，模范畴也同样，
% P02-E0729
(b)、(c) 与 (d) 立即成立。
\end{proof}

\section{分裂出一个自由模}
\label{more-algebra-section-serre-splitting}

\noindent
本节中的论证归功于 Serre；见 \cite{Serre-projective}。

\begin{situation}
\label{more-algebra-situation-splitting}
这里 $R$ 是环，$M$ 是有限表现 $R$-模。以 $\Omega \subset \Spec(R)$
表示所有闭点组成的集合，并赋予诱导拓扑。对 $x \in \Omega$，以
$M(x) = M/xM$ 表示 $M$ 在 $x$ 处的纤维。这是 $x$ 处剩余域
$\kappa(x)$ 上的有限维向量空间。给定 $s \in M$，以 $s(x)$ 表示
$s$ 在 $M(x)$ 中的像。
\end{situation}

\begin{lemma}
\label{more-algebra-lemma-which-elements-split}
在情形 \ref{more-algebra-situation-splitting} 中，设 $x \in \Omega$。存在
$\kappa(x)$-向量空间的典范短正合列
$$
0 \to B(x) \to M(x) \to V(x) \to 0
$$
，它具有以下性质：
% P02-E0730
对 $s_1, \ldots, s_r \in M$，以下条件等价：
\begin{enumerate}
\item 存在 $f \in R$，$f \not \in x$，使映射
$s_1, \ldots, s_r : R^{\oplus r} \to M$ 在逆转 $f$ 后成为一个
直和项的嵌入；
\item $s_1(x), \ldots, s_r(x)$ 在 $V(x)$ 中的像线性无关。
\end{enumerate}
\end{lemma}

\begin{proof}
把 $B(x) \subset M(x)$ 定义为映射
$$
\Hom_R(M, R) \to \Hom_{\kappa(x)}(M(x), \kappa(x))
$$
之像的正交补，并令 $V(x) = M(x)/B(x)$。于是每个 $R$-线性映射
$\varphi : M \to R$ 都诱导映射
$\overline{\varphi} : V(x) \to \kappa(x)$；反之，每个
$\kappa(x)$-线性映射 $\lambda : V(x) \to \kappa(x)$ 都等于某个
$\varphi$ 的 $\overline{\varphi}$。设 $s_1, \ldots, s_r \in M$。

\medskip\noindent
假设 $s_1, \ldots, s_r$ 在 $V(x)$ 中的像线性无关。于是可以找到
$\varphi_1, \ldots, \varphi_r \in \Hom_R(M, R)$，使
$\varphi_i(s_j)$ 在 $\kappa(x)$ 中的像为
$\delta_{ij}$\footnote{Kronecker 符号。}。因此复合映射
$$
R^{\oplus r}
\xrightarrow{s_1, \ldots, s_r}
M
\xrightarrow{\varphi_1, \ldots, \varphi_r}
R^{\oplus r}
$$
的矩阵有一个行列式 $f \in R$，它在 $\kappa(x)$ 中的像为 $1$
% P02-E0731
。显然，这说明 $s_1, \ldots, s_r : R^{\oplus r} \to M$ 在逆转
$f$ 后是一个直和项的嵌入。

\medskip\noindent
反之，假设存在 $f \in R$，$f \not \in x$，使
$s_1, \ldots, s_r : R^{\oplus r} \to M$ 在逆转 $f$ 后是一个
直和项的嵌入。于是可以找到 $R_f$-线性映射
$\varphi_i : M_f \to R_f$，使
$\varphi_i(s_j) = \delta_{ij} \in R_f$。由 Algebra，引理
\ref{algebra-lemma-hom-from-finitely-presented}，有
$\Hom_R(M, R)_f = \Hom_{R_f}(M_f, R_f)$，故可以找到 $n \geq 0$ 与
$\varphi'_i \in \Hom_R(M, R)$，使
$\varphi'_i(s_j) = f^n\delta_{ij} \in R$。由此，
$s_1, \ldots, s_r$ 在 $V(x)$ 中的像线性无关，因为
% P02-E0732
$\overline{\varphi}'_i(s_j) = f^n\delta_{ij}$。
\end{proof}

\noindent
在情形 \ref{more-algebra-situation-splitting} 中，给定
$s_1, \ldots, s_r \in M$，以
$Z(s_1, \ldots, s_r) \subset \Omega$ 表示所有满足如下条件的
$x \in \Omega$ 所成的集合：$s_1(x), \ldots, s_r(x)$ 在 $V(x)$ 中
的像线性相关。由引理，这是 $\Omega$ 的闭子集。

\begin{lemma}
\label{more-algebra-lemma-choose-values}
在情形 \ref{more-algebra-situation-splitting} 中，设
$x_1, \ldots, x_n \in \Omega$ 两两不同，并设
$v_i \in V(x_i)$。则存在 $s \in M$，使对
$i = 1, \ldots, n$，$s(x_i)$ 的像为 $v_i$。
\end{lemma}

\begin{proof}
由于 $x_i$ 是 $R$ 的极大理想，可以应用 Algebra，引理
\ref{algebra-lemma-chinese-remainder}，得到
$M(x_1) \oplus \ldots \oplus M(x_n)$ 是 $M$ 的商。
\end{proof}

\begin{proposition}
\label{more-algebra-proposition-splitting}
\begin{reference}
\cite[Theorem 2]{Serre-projective}
\end{reference}
在情形 \ref{more-algebra-situation-splitting} 中，假设 $\Omega$ 是 Noether
拓扑空间。设 $s_1, \ldots, s_h \in M$。设
$Z(s_1, \ldots, s_h) \subset F \subset \Omega$ 是闭集。设
$x_1, \ldots, x_n \in F$ 两两不同，并设 $v_i  \in V(x_i)$。
设 $k \geq 0$ 是整数，满足
$$
(*)\quad h + k \leq \dim_{\kappa(x)} V(x)\text{ 对所有 }x \in \Omega
$$
。则存在 $s \in M$ 与闭集 $F' \subset \Omega$，使得：
\begin{enumerate}
\item[(a)] $s(x_i)$ 的像为 $v_i$；
\item[(b)] $Z(s_1, \ldots, s_h, s) \subset F \cup F'$；
\item[(c)] $F'$ 的每个不可约分支在 $\Omega$ 中的余维都
$\geq k$。
\end{enumerate}
\end{proposition}

\begin{proof}
我们指出，余维定义于 Topology，第
\ref{topology-section-catenary-spaces} 节；还将使用 Topology，第
\ref{topology-section-noetherian} 节所收录的若干关于 Noether
拓扑空间的结果。

\medskip\noindent
对 $k$ 作归纳。若 $k = 0$，则按引理
\ref{more-algebra-lemma-choose-values} 选取 $s \in M$，并取 $F' = \Omega$。

\medskip\noindent
假设 $k > 0$。由归纳假设，可以选取 $u \in M$ 与闭集
$G \subset \Omega$，它们对
$s_1, \ldots, s_h$、$F$、$x_1, \ldots, x_n$、
$v_1, \ldots, v_n$ 以及 $k - 1$ 满足 (a)、(b)、(c)。

\medskip\noindent
设 $G = G_1 \cup \ldots \cup G_m$ 是 $G$ 的不可约分支分解。若
$G_j \subset F$，则可从列表中去掉它。因此可以假设，对
$j = 1, \ldots, m$，$G_j$ 不包含于 $F$。对
$j = 1, \ldots, m$，选择 $y_j \in G_j$，使
$y_j \not \in F$，且当 $j' \not = j$ 时
$y_j \not \in G_{j'}$。这是可行的，因为 $G$ 的不可约分支之间没有
包含关系。选取 $w_j \in V(y_j)$，使它不属于
$s_1(y_j), \ldots, s_h(y_j)$ 诸像张成的空间；这是可行的，因为
$h + k \leq \dim V(y_j)$ 且 $k > 0$。

\medskip\noindent
把归纳假设应用于 $h + 1$ 个截面
$s_1, \ldots, s_h, u$，闭集 $F \cup G$，诸点
$x_1, \ldots, x_n, y_1, \ldots, y_m \in F \cup G$，诸元素
$0 \in V(x_i)$ 与 $w_j \in V(y_j)$，以及整数 $k - 1$。注意，我们
把 $h$ 增加了 $1$，把 $k$ 减少了 $1$，所以命题的假设 $(*)$
仍然成立。这给出 $t \in M$ 与闭集 $H \subset \Omega$；它们对
$s_1, \ldots, s_h, u$、$F \cup G$、
$x_1, \ldots, x_n, y_1, \ldots, y_m$、
$0, \ldots, 0, w_1, \ldots, w_m$ 以及 $k - 1$ 满足 (a)、(b)、(c)。

\medskip\noindent
设 $H_1, \ldots, H_p \subset H$ 是 $H$ 中不包含于
$F \cup G$ 的不可约分支。与前面一样，选取 $z_l \in H_l$，使
$z_l \not \in F \cup G$，且当 $l' \not = l$ 时
$z_l \not \in H_{l'}$。利用 Algebra，引理
\ref{algebra-lemma-chinese-remainder}，可以选取 $f \in R$，使
$f(y_j) = 1$（$j = 1, \ldots, m$），且
$f(z_l) = 0$（$l = 1, \ldots, p$）。断言：元素
$s = u + f t$ 可用。

\medskip\noindent
首先，因为 $t(x_i) = 0$，$s(x_i)$ 的值与 $u(x_i)$ 相同，所以
$s(x_i)$ 的像为 $v_i$。这证明了 (a)。为完成证明，只需说明
$Z(s_1, \ldots, s_h, s)$ 中每个不包含于 $F$ 的不可约分支 $Z$
在 $\Omega$ 中的余维都 $\geq k$。事实上，这时可以令 $F'$ 等于
这些分支之并，从而得到 (b) 与 (c)。下面说明
$Z(s_1, \ldots, s_h, s)$ 不存在余维 $\leq k - 1$ 的不可约分支
$Z$：
\begin{enumerate}
\item 注意，
% P02-E0733
$Z(s_1, \ldots, s_h, s) \subset Z(s_1, \ldots, s_h, u, t) = F \cup H$
，因为 $s = u + ft$。故 $Z \subset H$。
\item $H$ 的不可约分支的余维都 $\geq k - 1$。因为 $Z$ 的余维
$\leq k - 1$，所以 $Z$ 等于 $H$ 的某个不可约分支。因此，对某个
$l \in \{1, \ldots, p\}$ 有 $Z = H_l$，或者对某个
% P02-E0734
$j \in \{1, \ldots m\}$ 有 $Z = G_j$。
\item 但 $Z = G_j$ 不可能，因为
$s_1(y_j), \ldots, s_h(y_j)$ 在 $V(y_j)$ 中的像线性无关，并且
$s(y_j) = u(y_j) + f(y_j) t(y_j) = u(y_j) + t(y_j)$ 的像形如
% P02-E0735
$$
\text{linear combination images of }s_i(y_j) + w_j
$$
；由 $w_j$ 的选择，它与 $s_1(y_j), \ldots, s_h(y_j)$ 在
$V(y_j)$ 中的像线性无关。
\item 同样，$Z = H_l$ 也不可能。事实上，
$s_1(z_l), \ldots, s_h(z_l)$ 在 $V(z_l)$ 中的像线性无关，而
$s(z_l) = u(z_l) + f(z_l) t(z_l) = u(z_l)$ 的像在 $V(z_l)$ 中
与它们线性无关，因为 $z_l \not \in F \cup G$。
\end{enumerate}
证明完毕。
\end{proof}

\begin{theorem}
\label{more-algebra-theorem-splitting}
\begin{reference}
\cite[Theorem 1]{Serre-projective}
\end{reference}
设环 $R$ 的极大谱 $\Omega \subset \Spec(R)$ 是维数
$d < \infty$ 的 Noether 拓扑空间。设 $M$ 是有限表现 $R$-模，并且
对所有 $\mathfrak m \in \Omega$，$R_\mathfrak m$-模
$M_\mathfrak m$ 都有秩 $> d$ 的自由直和项。则
$M \cong R \oplus M'$。
\end{theorem}

\begin{proof}
设对 $\mathfrak m \in \Omega$，$R_\mathfrak m^{\oplus r}$ 是
$M_\mathfrak m$ 的直和项。由 Algebra，引理
\ref{algebra-lemma-localization-colimit} 与
\ref{algebra-lemma-colimit-category-fp-modules}，对某个
$f \in R$，$f \not \in \mathfrak m$，$R_f^{\oplus r}$ 是
$M_f$ 的直和项。因此，该假设意味着对所有 $x \in \Omega$，
$\dim V(x) > d$，其中 $V(x)$ 如引理
\ref{more-algebra-lemma-which-elements-split} 所定义。把命题
\ref{more-algebra-proposition-splitting} 应用于 $F = \emptyset$、$h = 0$ 且没有
$s_i$、$n = 0$ 且没有 $x_i, v_i$，以及 $k = d + 1$，得到
$s \in M$ 与 $F' \subset \Omega$，使 $F'$ 的每个不可约分支的余维
都 $\geq d + 1$，且 $Z(s) \subset F'$。由于
$d = \dim(\Omega)$，这迫使 $F' = \emptyset$。因此
$s : R \to M$ 在每个极大理想处都是一个直和项的嵌入。由此，$s$
是普遍单射；见 Algebra，引理
\ref{algebra-lemma-universally-injective-check-stalks}。于是，由
Algebra，引理 \ref{algebra-lemma-universally-exact-split}，$s$ 是
分裂单射。
\end{proof}

\section{大投射模是自由模}
\label{more-algebra-section-big-projective-free}

\noindent
本节讨论 \cite{Bass} 的一个结果；建议读者参阅原论文。我们的论证将使用
\cite{Akasaki} 与 \cite{Hinohara} 中给出的略为简化的证明。

\begin{lemma}[Eilenberg 引理]
\label{more-algebra-lemma-eilenberg-swindle}
\begin{slogan}
Eilenberg 骗术
\end{slogan}
\begin{history}
在 \cite{Bass} 中可以读到：“……这是一个优雅的小骗术，Eilenberg
多年前便已注意到它，而它很可能是从 Barry Mazur 的额头上迸发出来的。”
\end{history}
\begin{reference}
\cite[Eilenberg's lemma]{Bass}
\end{reference}
若 $P \oplus Q \cong F$，其中 $F$ 是非有限生成的自由模，则
$P \oplus F \cong F$。
\end{lemma}

\begin{proof}
$$
F \cong F \oplus F \oplus \ldots \cong
P \oplus Q \oplus P \oplus Q \oplus \ldots \cong
P \oplus F \oplus F \oplus \ldots \cong P \oplus F
$$
\end{proof}

\begin{lemma}
\label{more-algebra-lemma-projective-plus-free-is-free}
设 $R$ 是环，$P$ 是投射模。存在自由模 $F$，使 $P \oplus F$ 自由。
\end{lemma}

\begin{proof}
由于 $P$ 投射，对某个模 $Q$，$F_0 = P \oplus Q$ 是自由模。令
$F = \bigoplus_{n \geq 1} F_0$。由引理
\ref{more-algebra-lemma-eilenberg-swindle}，$P \oplus F \cong F$。
\end{proof}

\begin{lemma}
\label{more-algebra-lemma-element-projective}
设 $R$ 是环，$P$ 是投射模。设 $s \in P$。则存在有限自由模 $F$ 与
有限自由直和项 $K \subset F \oplus P$，使
$(0, s) \in K$。
\end{lemma}

\begin{proof}
由引理 \ref{more-algebra-lemma-projective-plus-free-is-free}，可以找到一个（可能无限的）
自由模 $F$，使 $F \oplus P$ 自由。于是 $(0, s)$ 当然包含在某个
有限自由直和项 $K \subset F \oplus P$ 中。进而，$K$ 包含在
$F' \oplus P$ 中，其中 $F' \subset F$ 是有限自由直和项。
\end{proof}

\begin{lemma}
\label{more-algebra-lemma-trick-to-find-good-element}
设 $R$ 是环，其 Jacobson 根为 $J$，且 $R/J$ 是 Noether 环。设 $P$
是投射 $R$-模，并且对 $R$ 的每个极大理想 $\mathfrak m$，
$P_\mathfrak m$ 都有无限秩。设 $s \in P$ 且 $M \subset P$，满足
$Rs + M = P$。则可以找到 $m \in M$，使 $R(s + m)$ 是 $P$ 的自由
直和项。
\end{lemma}

\begin{proof}
该陈述有意义，因为由 Algebra，定理
\ref{algebra-theorem-projective-free-over-local-ring}，
$P_\mathfrak m$ 自由。

\medskip\noindent
以 $M' \subset P/JP$ 表示 $M$ 的像，以 $s' \in P/JP$ 表示 $s$ 的像。
注意，
% P02-E0736
$R/J s' + M' = P/JP$。假设可以找到 $m' \in M'$，使
$R/J(s' + m')$ 是 $M'$ 的自由直和项。
% P02-E0737
选择给出分裂的 $\varphi' : P/JP \to R/J$，即在 $R/J$ 中有
$\varphi'(s' + m') = 1$。由于 $P$ 是投射 $R$-模，可以找到
$\varphi'$ 的提升 $\varphi : P \to R$。选取映到 $m'$ 的
$m \in M$。于是 $\varphi(s + m) \in R$ 模 $J$ 同余于 $1$，因而是
$R$ 中的单位元（Algebra，引理
\ref{algebra-lemma-contained-in-radical}）。所以 $R(s + m)$ 是 $P$
的自由直和项。这把问题约化为下一段讨论的情形。

\medskip\noindent
假设 $R$ 是 Noether 环。设 $m \in M$，并设
$\varphi_1, \ldots, \varphi_n : P \to R$ 是 $R$-线性映射。以
$$
Z(s + m, \varphi_1, \ldots, \varphi_n) \subset \Spec(R)
$$
表示 $\varphi_1(s + m), \ldots, \varphi_n(s + m) \in R$ 的零点集。

\medskip\noindent
假设 $\mathfrak m$ 是 $R$ 的极大理想，且
$\mathfrak m \in Z(s + m, \varphi_1, \ldots, \varphi_n)$。令
$K = M \cap \bigcap \Ker(\varphi_i)$。我们断言，映射
$$
K/\mathfrak mK \to P/\mathfrak m P
$$
的像有无限维。事实上，商 $P/K$ 是有限 $R$-模，因为它同构于
$P/M \oplus R^{\oplus n}$ 的一个子模。因此，该映射的核是
$\text{Tor}_1^R(P/K, \kappa(\mathfrak m))$ 的一个商；由 Algebra，
引理 \ref{algebra-lemma-tor-noetherian}，后者有限。结合
$P/\mathfrak mP$ 有无限维这一事实，便得到断言。因此可以找到
$t \in K$，它映到向量空间 $P/\mathfrak mP$ 的元素
$\overline{t}$，而该元素与 $s + m$ 的像
$\overline{s + m}$ 线性无关。由线性代数，可以找到 $R$-线性映射
% P02-E0738
$\overline{\varphi} : P \to \kappa(\mathfrak m)$，使
$\overline{\varphi}(\overline{t}) = 1$ 且
$\overline{\varphi}(\overline{s + m}) = 0$。由于 $P$ 投射，可以找到
提升 $\overline{\varphi}$ 的 $R$-线性映射 $\varphi : P \to R$。
于是零点集
$Z(s + m + t, \varphi_1, \ldots, \varphi_n, \varphi)$ 包含于
$Z(s + m, \varphi_1, \ldots, \varphi_n)$，但不包含
$\mathfrak m$；换言之，它严格小于
$Z(s + m, \varphi_1, \ldots, \varphi_n)$。

\medskip\noindent
由于 $\Spec(R)$ 是 Noether 拓扑空间，由上述论证可知，可以找到
$m \in M$ 与 $\varphi_1, \ldots, \varphi_n : P \to R$，使闭子集
$Z(s + m, \varphi_1, \ldots, \varphi_n)$ 不包含 $\Spec(R)$ 的任何
闭点。因此 $Z(s + m, \varphi_1, \ldots, \varphi_n) = \emptyset$。
所以可以找到 $r_1, \ldots, r_n \in R$，使
$\sum r_i\varphi_i(s + m) = 1$。于是
$$
R \xrightarrow{s + m} P \xrightarrow{\sum r_i \varphi_i} R
$$
就是所需的分裂。
\end{proof}

\begin{lemma}
\label{more-algebra-lemma-element-in-free-summand}
设 $R$ 是环，其 Jacobson 根为 $J$，且 $R/J$ 是 Noether 环。设 $P$
是投射 $R$-模，并且对 $R$ 的每个极大理想 $\mathfrak m$，
$P_\mathfrak m$ 都有无限秩。设 $s \in P$。则可以找到有限稳定自由
直和项 $M \subset P$，使 $s \in M$。
\end{lemma}

\begin{proof}
由引理 \ref{more-algebra-lemma-element-projective}，可以找到有限自由模 $F$ 与
有限自由直和项 $K \subset F \oplus P$，使
$(0, s) \in K$。对 $F$ 的秩作归纳，把问题约化为下一段讨论的情形。

\medskip\noindent
假设存在有限稳定自由直和项 $K \subset R \oplus P$，使
$(0, s) \in K$。选择 $K$ 的补 $K'$，即
$R \oplus P = K \oplus K'$。投影
% P02-E0739
$\pi : R \oplus P \to K'$ 是满射；因此由引理
\ref{more-algebra-lemma-trick-to-find-good-element}，可以找到 $p \in P$，使
$\pi(1, p) \in K'$ 生成自由直和项。相应地，写
$K' = R\pi(1, p) \oplus K''$。于是
$$
R \oplus P = K \oplus K' = K \oplus R\pi(1, p) \oplus K''
$$
。投影 $\pi' : P \to K''$ 是满射\footnote{事实上，若
$k'' \in K''$，则把 $k''$ 看作 $K'$ 的元素，可以写成
$k'' = \lambda \pi(1, 0) + \pi(0, q)$，其中 $\lambda \in R$ 且
$q \in P$。这意味着
$k'' = \lambda \pi(1, p) + \pi(0, q - \lambda p)$。进而，这意味着
$q - \lambda p$ 在复合映射
$P \to R \oplus P \xrightarrow{\pi} K' \to K''$ 下映到 $k''$，因为
$K' \to K''$ 零化 $\pi(1, p)$。}，因而它分裂
% P02-E0740
（因为 $K''$ 投射）。因此 $\Ker(\pi') \subset P$ 是包含 $s$ 的
直和项。最后，由构造有同构
$$
R \oplus \Ker(\pi') \cong K \oplus R\pi(1, p)
$$
；所以，由于 $K$ 有限且稳定自由，$\Ker(\pi')$ 也如此。
\end{proof}

\begin{theorem}
\label{more-algebra-theorem-countable-free}
\begin{reference}
\cite[Theorem 3.1]{Bass} 的交换情形
\end{reference}
设 $R$ 是环，其 Jacobson 根为 $J$，且 $R/J$ 是 Noether 环。设 $P$
是可数生成投射 $R$-模，并且对 $R$ 的每个极大理想
$\mathfrak m$，$P_\mathfrak m$ 都有无限秩。则 $P$ 自由。
\end{theorem}

\begin{proof}
先证明 $P$ 是有限稳定自由模的可数直和。设
$x_1, x_2, \ldots$ 是 $P$ 的一组可数生成元。归纳构造 $P$ 的有限
稳定自由直和项 $F_1, F_2, \ldots$，使对所有 $n$，
$F_1 \oplus \ldots \oplus F_n$ 都是包含
$x_1, \ldots, x_n$ 的 $P$ 的直和项。事实上，给定具有所需性质的
$F_1, \ldots, F_n$，写
$$
P = F_1 \oplus \ldots \oplus F_n \oplus P'
$$
，并令 $s \in P'$ 是 $x_{n + 1}$ 的像。由引理
\ref{more-algebra-lemma-element-in-free-summand}，可以找到包含 $s$ 的有限
稳定自由直和项 $F_{n + 1} \subset P'$。于是
$P = \bigoplus_{i = 1}^{\infty} F_i$。

\medskip\noindent
假设 $P$ 是非零有限稳定自由模的无限直和
$P = \bigoplus_{i = 1}^{\infty} F_i$。下面按如下方式使用各 $F_i$
的稳定自由性：每个 $F_i$ 的秩都是常数（且为正）。因此，对 $R$
的每个极大理想 $\mathfrak m$，$P_\mathfrak m$ 都是可数无限秩的
自由模。把引理 \ref{more-algebra-lemma-trick-to-find-good-element} 应用于
$s = 0$ 与 $M = P$，可以找到 $t_1 \in P$，使 $Rt_1$ 是 $P$ 的
自由直和项。于是对某个 $n_1 > n_0 = 0$，
$t_1$ 包含在 $F_1 \oplus \ldots \oplus F_{n_1}$ 中。把同一论证
应用于 $\bigoplus_{n > n_1} F_n$，得到 $n_1 < n_2$ 与
$t_2 \in F_{n_1 + 1} \oplus \ldots \oplus F_{n_2}$，后者生成一个
自由直和项。继续如此，得到无限秩自由直和项
$$
\bigoplus\nolimits_{i \geq 1} t_i :
\bigoplus\nolimits_{i \geq 1} R
\longrightarrow
\bigoplus\nolimits_{i \geq 1}
\bigoplus\nolimits_{n_i \geq n > n_{i - 1}} F_n = P
$$
。因此，对某个可数秩自由 $R$-模 $F$，有
$P \cong Q \oplus F$。由于 $Q$ 可数生成，对某个模 $Q'$，有
$Q \oplus Q' \cong F$。于是 Eilenberg 骗术（引理
\ref{more-algebra-lemma-eilenberg-swindle}）说明 $Q \oplus F \cong F$，从而
$P$ 自由。
\end{proof}
